% Options for packages loaded elsewhere
% Options for packages loaded elsewhere
\PassOptionsToPackage{unicode}{hyperref}
\PassOptionsToPackage{hyphens}{url}
\PassOptionsToPackage{dvipsnames,svgnames,x11names}{xcolor}
\PassOptionsToPackage{space}{xeCJK}
%
\documentclass[
  a4paper,
]{book}
\usepackage{xcolor}
\usepackage{amsmath,amssymb}
\setcounter{secnumdepth}{-\maxdimen} % remove section numbering
\usepackage{iftex}
\ifPDFTeX
  \usepackage[T1]{fontenc}
  \usepackage[utf8]{inputenc}
  \usepackage{textcomp} % provide euro and other symbols
\else % if luatex or xetex
  \usepackage{unicode-math} % this also loads fontspec
  \defaultfontfeatures{Scale=MatchLowercase}
  \defaultfontfeatures[\rmfamily]{Ligatures=TeX,Scale=1}
\fi
\usepackage{lmodern}
\ifPDFTeX\else
  % xetex/luatex font selection
  \setmainfont[]{Noto Serif CJK SC}
  \setsansfont[]{Noto Sans CJK SC}
  \setmonofont[]{Noto Sans Mono CJK SC}
  \ifXeTeX
    \usepackage{xeCJK}
    \setCJKmainfont[]{Noto Serif CJK SC}
    \setCJKsansfont[]{Noto Sans CJK SC}
    \setCJKmonofont[]{Noto Sans Mono CJK SC}
  \fi
  \ifLuaTeX
    \usepackage[]{luatexja-fontspec}
    \setmainjfont[]{Noto Serif CJK SC}
  \fi
\fi
% Use upquote if available, for straight quotes in verbatim environments
\IfFileExists{upquote.sty}{\usepackage{upquote}}{}
\IfFileExists{microtype.sty}{% use microtype if available
  \usepackage[]{microtype}
  \UseMicrotypeSet[protrusion]{basicmath} % disable protrusion for tt fonts
}{}
\makeatletter
\@ifundefined{KOMAClassName}{% if non-KOMA class
  \IfFileExists{parskip.sty}{%
    \usepackage{parskip}
  }{% else
    \setlength{\parindent}{0pt}
    \setlength{\parskip}{6pt plus 2pt minus 1pt}}
}{% if KOMA class
  \KOMAoptions{parskip=half}}
\makeatother
% Make \paragraph and \subparagraph free-standing
\makeatletter
\ifx\paragraph\undefined\else
  \let\oldparagraph\paragraph
  \renewcommand{\paragraph}{
    \@ifstar
      \xxxParagraphStar
      \xxxParagraphNoStar
  }
  \newcommand{\xxxParagraphStar}[1]{\oldparagraph*{#1}\mbox{}}
  \newcommand{\xxxParagraphNoStar}[1]{\oldparagraph{#1}\mbox{}}
\fi
\ifx\subparagraph\undefined\else
  \let\oldsubparagraph\subparagraph
  \renewcommand{\subparagraph}{
    \@ifstar
      \xxxSubParagraphStar
      \xxxSubParagraphNoStar
  }
  \newcommand{\xxxSubParagraphStar}[1]{\oldsubparagraph*{#1}\mbox{}}
  \newcommand{\xxxSubParagraphNoStar}[1]{\oldsubparagraph{#1}\mbox{}}
\fi
\makeatother

\usepackage{color}
\usepackage{fancyvrb}
\newcommand{\VerbBar}{|}
\newcommand{\VERB}{\Verb[commandchars=\\\{\}]}
\DefineVerbatimEnvironment{Highlighting}{Verbatim}{commandchars=\\\{\}}
% Add ',fontsize=\small' for more characters per line
\usepackage{framed}
\definecolor{shadecolor}{RGB}{241,243,245}
\newenvironment{Shaded}{\begin{snugshade}}{\end{snugshade}}
\newcommand{\AlertTok}[1]{\textcolor[rgb]{0.68,0.00,0.00}{#1}}
\newcommand{\AnnotationTok}[1]{\textcolor[rgb]{0.37,0.37,0.37}{#1}}
\newcommand{\AttributeTok}[1]{\textcolor[rgb]{0.40,0.45,0.13}{#1}}
\newcommand{\BaseNTok}[1]{\textcolor[rgb]{0.68,0.00,0.00}{#1}}
\newcommand{\BuiltInTok}[1]{\textcolor[rgb]{0.00,0.23,0.31}{#1}}
\newcommand{\CharTok}[1]{\textcolor[rgb]{0.13,0.47,0.30}{#1}}
\newcommand{\CommentTok}[1]{\textcolor[rgb]{0.37,0.37,0.37}{#1}}
\newcommand{\CommentVarTok}[1]{\textcolor[rgb]{0.37,0.37,0.37}{\textit{#1}}}
\newcommand{\ConstantTok}[1]{\textcolor[rgb]{0.56,0.35,0.01}{#1}}
\newcommand{\ControlFlowTok}[1]{\textcolor[rgb]{0.00,0.23,0.31}{\textbf{#1}}}
\newcommand{\DataTypeTok}[1]{\textcolor[rgb]{0.68,0.00,0.00}{#1}}
\newcommand{\DecValTok}[1]{\textcolor[rgb]{0.68,0.00,0.00}{#1}}
\newcommand{\DocumentationTok}[1]{\textcolor[rgb]{0.37,0.37,0.37}{\textit{#1}}}
\newcommand{\ErrorTok}[1]{\textcolor[rgb]{0.68,0.00,0.00}{#1}}
\newcommand{\ExtensionTok}[1]{\textcolor[rgb]{0.00,0.23,0.31}{#1}}
\newcommand{\FloatTok}[1]{\textcolor[rgb]{0.68,0.00,0.00}{#1}}
\newcommand{\FunctionTok}[1]{\textcolor[rgb]{0.28,0.35,0.67}{#1}}
\newcommand{\ImportTok}[1]{\textcolor[rgb]{0.00,0.46,0.62}{#1}}
\newcommand{\InformationTok}[1]{\textcolor[rgb]{0.37,0.37,0.37}{#1}}
\newcommand{\KeywordTok}[1]{\textcolor[rgb]{0.00,0.23,0.31}{\textbf{#1}}}
\newcommand{\NormalTok}[1]{\textcolor[rgb]{0.00,0.23,0.31}{#1}}
\newcommand{\OperatorTok}[1]{\textcolor[rgb]{0.37,0.37,0.37}{#1}}
\newcommand{\OtherTok}[1]{\textcolor[rgb]{0.00,0.23,0.31}{#1}}
\newcommand{\PreprocessorTok}[1]{\textcolor[rgb]{0.68,0.00,0.00}{#1}}
\newcommand{\RegionMarkerTok}[1]{\textcolor[rgb]{0.00,0.23,0.31}{#1}}
\newcommand{\SpecialCharTok}[1]{\textcolor[rgb]{0.37,0.37,0.37}{#1}}
\newcommand{\SpecialStringTok}[1]{\textcolor[rgb]{0.13,0.47,0.30}{#1}}
\newcommand{\StringTok}[1]{\textcolor[rgb]{0.13,0.47,0.30}{#1}}
\newcommand{\VariableTok}[1]{\textcolor[rgb]{0.07,0.07,0.07}{#1}}
\newcommand{\VerbatimStringTok}[1]{\textcolor[rgb]{0.13,0.47,0.30}{#1}}
\newcommand{\WarningTok}[1]{\textcolor[rgb]{0.37,0.37,0.37}{\textit{#1}}}

\usepackage{longtable,booktabs,array}
\newcounter{none} % for unnumbered tables
\usepackage{calc} % for calculating minipage widths
% Correct order of tables after \paragraph or \subparagraph
\usepackage{etoolbox}
\makeatletter
\patchcmd\longtable{\par}{\if@noskipsec\mbox{}\fi\par}{}{}
\makeatother
% Allow footnotes in longtable head/foot
\IfFileExists{footnotehyper.sty}{\usepackage{footnotehyper}}{\usepackage{footnote}}
\makesavenoteenv{longtable}
\usepackage{graphicx}
\makeatletter
\newsavebox\pandoc@box
\newcommand*\pandocbounded[1]{% scales image to fit in text height/width
  \sbox\pandoc@box{#1}%
  \Gscale@div\@tempa{\textheight}{\dimexpr\ht\pandoc@box+\dp\pandoc@box\relax}%
  \Gscale@div\@tempb{\linewidth}{\wd\pandoc@box}%
  \ifdim\@tempb\p@<\@tempa\p@\let\@tempa\@tempb\fi% select the smaller of both
  \ifdim\@tempa\p@<\p@\scalebox{\@tempa}{\usebox\pandoc@box}%
  \else\usebox{\pandoc@box}%
  \fi%
}
% Set default figure placement to htbp
\def\fps@figure{htbp}
\makeatother





\setlength{\emergencystretch}{3em} % prevent overfull lines

\providecommand{\tightlist}{%
  \setlength{\itemsep}{0pt}\setlength{\parskip}{0pt}}



 


\makeatletter
\@ifpackageloaded{bookmark}{}{\usepackage{bookmark}}
\makeatother
\makeatletter
\@ifpackageloaded{caption}{}{\usepackage{caption}}
\AtBeginDocument{%
\ifdefined\contentsname
  \renewcommand*\contentsname{Table of contents}
\else
  \newcommand\contentsname{Table of contents}
\fi
\ifdefined\listfigurename
  \renewcommand*\listfigurename{List of Figures}
\else
  \newcommand\listfigurename{List of Figures}
\fi
\ifdefined\listtablename
  \renewcommand*\listtablename{List of Tables}
\else
  \newcommand\listtablename{List of Tables}
\fi
\ifdefined\figurename
  \renewcommand*\figurename{Figure}
\else
  \newcommand\figurename{Figure}
\fi
\ifdefined\tablename
  \renewcommand*\tablename{Table}
\else
  \newcommand\tablename{Table}
\fi
}
\@ifpackageloaded{float}{}{\usepackage{float}}
\floatstyle{ruled}
\@ifundefined{c@chapter}{\newfloat{codelisting}{h}{lop}}{\newfloat{codelisting}{h}{lop}[chapter]}
\floatname{codelisting}{Listing}
\newcommand*\listoflistings{\listof{codelisting}{List of Listings}}
\makeatother
\makeatletter
\makeatother
\makeatletter
\@ifpackageloaded{caption}{}{\usepackage{caption}}
\@ifpackageloaded{subcaption}{}{\usepackage{subcaption}}
\makeatother
\usepackage{bookmark}
\IfFileExists{xurl.sty}{\usepackage{xurl}}{} % add URL line breaks if available
\urlstyle{same}
\hypersetup{
  pdftitle={Rebellion and Factionalism in a Chinese Province: Zhejiang{,} 1966-1976},
  colorlinks=true,
  linkcolor={Maroon},
  filecolor={Maroon},
  citecolor={Blue},
  urlcolor={Blue},
  pdfcreator={LaTeX via pandoc}}


\title{Rebellion and Factionalism in a Chinese Province: Zhejiang,
1966-1976}
\usepackage{etoolbox}
\makeatletter
\providecommand{\subtitle}[1]{% add subtitle to \maketitle
  \apptocmd{\@title}{\par {\large #1 \par}}{}{}
}
\makeatother
\subtitle{一个中国省份的造反与派性：浙江，1966-1976}
\author{}
\date{}
\begin{document}
\frontmatter
\maketitle

\renewcommand*\contentsname{Table of contents}
{
\setcounter{tocdepth}{1}
\tableofcontents
}

\mainmatter
\bookmarksetup{startatroot}

\chapter*{Rebellion and Factionalism in a Chinese Province Zhejiang,
1966-1976 /
一个中国省份的造反与派性：浙江，1966-1976}\label{rebellion-and-factionalism-in-a-chinese-province-zhejiang-1966-1976-ux4e00ux4e2aux4e2dux56fdux7701ux4efdux7684ux9020ux53cdux4e0eux6d3eux6027ux6d59ux6c5f1966-1976}
\addcontentsline{toc}{chapter}{Rebellion and Factionalism in a Chinese
Province Zhejiang, 1966-1976 /
一个中国省份的造反与派性：浙江，1966-1976}

\markboth{Rebellion and Factionalism in a Chinese Province Zhejiang,
1966-1976 / 一个中国省份的造反与派性：浙江，1966-1976}{Rebellion and
Factionalism in a Chinese Province Zhejiang, 1966-1976 /
一个中国省份的造反与派性：浙江，1966-1976}

\bookmarksetup{startatroot}

\chapter*{ACKNOWLEDGMENTS / 致谢}\label{acknowledgments-ux81f4ux8c22}
\addcontentsline{toc}{chapter}{ACKNOWLEDGMENTS / 致谢}

\markboth{ACKNOWLEDGMENTS / 致谢}{ACKNOWLEDGMENTS / 致谢}

This monograph is based on a doctoral thesis submitted to the University
of Adelaide in 1985. In the four years since that time further research
has been carried out at the Menzies Library of the Australian National
University and the Australian National Library in Canberra, at the
Institute of International Relations in Taiwan, at the libraries of
Harvard University, the Universities of Michigan and California
(Berkeley) in the United States, at the library of Hangzhou University
in the People's Republic of China and at the Institute of International
Relations in Taibei.

本专著以一篇 1985
年提交给阿德莱德大学的博士论文为基础。此后四年间，我又在堪培拉澳大利亚国立大学孟席斯图书馆和澳大利亚国家图书馆、台湾的国际关系研究所、美国哈佛大学以及密歇根大学和加利福尼亚大学（伯克利）的图书馆、中华人民共和国杭州大学图书馆以及台北的国际关系研究所开展了进一步研究。

I gratefully acknowledge financial assistance from the University of
Melbourne and the Research School of Pacific Studies, the Australian
National University as well as research grants from the Pacific Cultural
Foundation, Taiwan, the Republic of China, and The Academies'
Australia-China Exchange in the Humanities and Social Sciences provided
by The Academy of the Social Sciences and the Australian Academy of the
Humanities in Australia, which enabled me to undertake further research
for this work.

我谨感谢墨尔本大学和澳大利亚国立大学太平洋研究院提供的经费支持，也感谢中华民国台湾太平洋文化基金会的研究资助，以及由澳大利亚社会科学院和澳大利亚人文科学院提供的``学院澳中人文与社会科学交流''项目资助；这些支持使我得以为本书开展进一步研究。

Over the past three years I have attempted to clarify my thinking, hone
the methodological approach, and further tighten the argument. For their
assistance in this task I acknowledge with thanks the detailed and most
helpful referees' reports on my thesis by Bill Brugger and Fred Teiwes
and for their encouragement to publish the work. Fred Teiwes also read
an additional chapter and put forward further useful suggestions. The
Head of the Contemporary China Centre, Jon Unger, took time out from a
heavy schedule to read and pass helpful comments on a draft of the
manuscript and to take the trouble to think of the title.

过去三年中，我一直努力厘清思路、打磨方法论路径，并进一步收紧论证。在这一过程中，Bill
Brugger 和 Fred Teiwes
对我的论文所作详尽且极有帮助的评审报告，以及他们鼓励我出版此作，我在此致以谢意。Fred
Teiwes 还阅读了另一个章节，并提出了更多有益建议。当代中国研究中心主任
Jon Unger
在繁忙日程中抽出时间阅读书稿草稿并提出有益意见，还费心思考书名。

The staff of the Asian Studies section of the Menzies Library of the
Australian National University, in particular Y. S. Chan, Susan Prentice
and Renata Osborne, have helped me locate sources from the superb
collection in the library. I thank them for their uncomplaining
assistance and consideration. Andrew Gosling and his staff of the East
Asian Section, Asian Collection at the National Library provided
cheerful and efficient service. Kay Dancey of the Cartography Unit,
Research School of Pacific Studies drew the maps with great willingness
and expertise.

澳大利亚国立大学孟席斯图书馆亚洲研究部的工作人员，尤其是 Y. S.
Chan、Susan Prentice 和 Renata
Osborne，帮助我从馆内卓越藏书中查找资料。我感谢他们毫无怨言的协助和周到照顾。国家图书馆亚洲馆藏东亚部的
Andrew Gosling 及其同事提供了热情而高效的服务。太平洋研究院制图室的 Kay
Dancey 以极大的热忱和专业能力绘制了地图。

For assistance in using computers and in particular the Macintosh
Chinese language program I express my gratitude to my colleague David
Kelly of the Contemporary China Centre and to Doug Whaite of the
Department of Economic History of the Research School of Pacific
Studies. The Centre's secretary Dianne Stacey has been of inestimable
assistance to me since my arrival in Canberra and checked the final
proofs of the manuscript. It should be noted that the manuscript has
been prepared by the author using Apple Macintosh equipment together
with Macintosh Chinese Word and Microsoft Word software. The main
problem in using the two programs together can be seen in the short
lines where Chinese characters are used. This is because MS allows two
spaces for each Chinese character. The shortcoming cannot be overcome
without a major cut and paste operation.

在使用计算机，尤其是 Macintosh 中文程序方面，我感谢当代中国研究中心同事
David Kelly，以及太平洋研究院经济史系的 Doug Whaite
给予的帮助。自我抵达堪培拉以来，中心秘书 Dianne Stacey
给了我无法估量的帮助，并校阅了书稿的最后清样。需要说明的是，本书书稿由作者使用
Apple Macintosh 设备，并配合 Macintosh Chinese Word 与 Microsoft Word
软件完成。两个程序合用时的主要问题，可见于使用汉字处出现的短行；这是因为
MS 为每个汉字预留两个空格。若不进行大规模剪切和粘贴，这一缺陷无法克服。

To my many friends in Hangzhou who helped me find my way through the
labyrinthine paths of local politics and helped me with research
material, especially Jane, I express my thanks and appreciation.

我还要向杭州的许多朋友表示感谢，尤其是
Jane；他们帮助我穿行于地方政治迷宫般的路径，并在研究资料方面给予协助。

My major regret in publishing this work is that my parents have not
lived to witness its publication. The last book my father read before
his death in 1976, on the eve of our first trip to Hangzhou, contained a
chapter on that city. Fortunately, my mother was later able to make the
journey to China and see Hangzhou for herself.

出版本书时，我最大的遗憾是父母未能活着见证它问世。1976
年，在我们首次前往杭州前夕，父亲去世；他生前读的最后一本书中有一章写的正是这座城市。幸运的是，母亲后来得以前往中国，并亲眼见到杭州。

In the completion of both the thesis and the book, the chief and
overwhelming thanks are due to my wife Margaret. She has proved a loyal
partner over twenty years. She has coped with my demanding nature and
given love, affection, companionship and care in return. Margaret read
the final draft of the manuscript and made useful suggestion to improve
the expression and style. This book is dedicated to her with gratitude
and love.

无论是完成论文还是完成本书，最主要、最深切的谢意都应归于我的妻子
Margaret。二十年来，她一直是忠诚的伴侣。她包容了我苛求的一面，并以爱、深情、陪伴和关怀相待。Margaret
阅读了书稿终稿，并提出了有助于改进表达和文风的建议。本书谨献给她，满怀感激与爱意。

Canberra

堪培拉

February 1990.

1990 年 2 月。

\bookmarksetup{startatroot}

\chapter*{ABBREVIATIONS /
缩略语}\label{abbreviations-ux7f29ux7565ux8bed}
\addcontentsline{toc}{chapter}{ABBREVIATIONS / 缩略语}

\markboth{ABBREVIATIONS / 缩略语}{ABBREVIATIONS / 缩略语}

{\def\LTcaptype{none} % do not increment counter
\begin{longtable}[]{@{}
  >{\raggedright\arraybackslash}p{(\linewidth - 2\tabcolsep) * \real{0.5000}}
  >{\raggedright\arraybackslash}p{(\linewidth - 2\tabcolsep) * \real{0.5000}}@{}}
\toprule\noalign{}
\begin{minipage}[b]{\linewidth}\raggedright
Abbreviation
\end{minipage} & \begin{minipage}[b]{\linewidth}\raggedright
Full form
\end{minipage} \\
\midrule\noalign{}
\endhead
\bottomrule\noalign{}
\endlastfoot
AS & ASIAN SURVEY \\
BR & BEIJING REVIEW \\
CC & CENTRAL COMMITTEE \\
CNA & CHINA NEWS ANALYSIS \\
CNS & CHINA NEWS SUMMARY \\
CQ & CHINA QUARTERLY \\
CCP & CHINESE COMMUNIST PARTY \\
CYL & COMMUNIST YOUTH LEAGUE \\
CS & CURRENT SCENE \\
F\&F & FACTS \& FEATURES \\
FEER & FAR EASTERN ECONOMIC REVIEW \\
FBIS/CHI & FOREIGN BROADCAST INFORMATION SERVICE. DAILY REPORT. CHINA \\
FBIS/CR/PSMA & FOREIGN BROADCAST INFORMATION SERVICE. CHINA REPORT.
POLITICAL, SOCIOLOGICAL AND MILITARY AFFAIRS \\
HMC & HANGZHOU MUNICIPAL COMMITTEE \\
HMMCH & HANGZHOU MUNICIPAL MILITIA COMMAND HEADQUARTERS \\
HMRC & HANGZHOU MUNICIPAL REVOLUTIONARY COMMITTEE \\
HMWC & HANGZHOU MUNICIPAL WORKERS' CONGRESS \\
HZRB & HANGZHOU RIBAO \\
I\&S & ISSUES \& STUDIES \\
MAC & MILITARY AFFAIRS COMMISSION \\
NPC & NATIONAL PEOPLE'S CONGRESS \\
PR & PEKING REVIEW \\
PLA & PEOPLE'S LIBERATION ARMY \\
RMRB & RENMIN RIBAO \\
SCMM & UNITED STATES CONSULATE-GENERAL (HONG KONG). SELECTIONS FROM
CHINA MAINLAND MAGAZINES \\
SCMP & UNITED STATES CONSULATE-GENERAL (HONG KONG). SURVEY OF CHINA
MAINLAND PRESS \\
SWB/FE & BBC. SUMMARY OF WORLD BROADCASTS. PART THREE. THE FAR EAST. \\
URS & UNION RESEARCH INSTITUTE (HONG KONG).UNION RESEARCH SERVICE \\
ZPC & ZHEJIANG PROVINCIAL COMMITTEE \\
ZPMD & ZHEJIANG PROVINCIAL MILITARY DISTRICT \\
ZPRC & ZHEJIANG PROVINCIAL REVOLUTIONARY COMMITTEE \\
ZJRB & ZHEJIANG, RIBAO \\
ZPS & ZHEJIANG PROVINCIAL SERVICE \\
ZPTUC & ZHEJIANG PROVINCIAL TRADE UNION COUNCIL \\
\end{longtable}
}

{\def\LTcaptype{none} % do not increment counter
\begin{longtable}[]{@{}ll@{}}
\toprule\noalign{}
缩略语 & 全称 \\
\midrule\noalign{}
\endhead
\bottomrule\noalign{}
\endlastfoot
AS & 《亚洲调查》 \\
BR & 《北京周报》 \\
CC & 中央委员会 \\
CNA & 《中国新闻分析》 \\
CNS & 《中国新闻摘要》 \\
CQ & 《中国季刊》 \\
CCP & 中国共产党 \\
CYL & 共产主义青年团 \\
CS & \emph{Current Scene} \\
F\&F & 《事实与特写》 \\
FEER & 《远东经济评论》 \\
FBIS/CHI & 外国广播信息服务：《每日报告：中国》 \\
FBIS/CR/PSMA & 外国广播信息服务：《中国报告：政治、社会和军事事务》 \\
HMC & 中共杭州市委 \\
HMMCH & 杭州市民兵指挥部 \\
HMRC & 杭州市革命委员会 \\
HMWC & 杭州市工人代表大会 \\
HZRB & 《杭州日报》 \\
I\&S & 《问题与研究》 \\
MAC & 军事委员会 \\
NPC & 全国人民代表大会 \\
PR & 《北京周报》 \\
PLA & 中国人民解放军 \\
RMRB & 《人民日报》 \\
SCMM & 美国驻香港总领事馆：《中国大陆杂志选辑》 \\
SCMP & 美国驻香港总领事馆：《中国大陆报刊概览》 \\
SWB/FE & BBC：《世界广播摘要》第三部分：远东 \\
URS & 友联研究所（香港）：《友联研究服务》 \\
ZPC & 中共浙江省委 \\
ZPMD & 浙江省军区 \\
ZPRC & 浙江省革命委员会 \\
ZJRB & 《浙江日报》 \\
ZPS & 浙江省广播服务 \\
ZPTUC & 浙江省总工会 \\
\end{longtable}
}

\bookmarksetup{startatroot}

\chapter{INTRODUCTION / 导论}\label{introduction-ux5bfcux8bba}

Up until now there has been no detailed study of the politics of the
decade long Cultural Revolution at the sub-national level. Such a hiatus
has meant that our understanding of these years has largely been
conditioned by generalizations based on studies of national politics, or
research taking Guangdong province as its subject. Guangdong, with its
proximity to Hong Kong, distance from Beijing and history of antipathy
toward central rule, may represent the exception rather than the rule.
Essays describing the early years of the Cultural Revolution in a number
of provinces were published almost two decades ago.\footnote{The
  Cultural Revolution in The Provinces, Harvard East Asian Monographs,
  42 (Cambridge, Mass.: Harvard University Press, 1971).

  The Cultural Revolution in The Provinces, Harvard East Asian
  Monographs, 42 (Cambridge, Mass.: Harvard University Press, 1971)。}
These appeared in a period when China's doors were shut tight to the
outside world and access to provincial newspapers and other primary
source material was extremely limited. As for the period of the early to
mid 1970s, there is an even greater lacuna in the literature.

迄今为止，还没有关于长达十年的文化大革命在次国家层面政治的详细研究。这一空缺意味着，我们对这些岁月的理解，很大程度上受制于以全国政治研究为基础的概括，或以广东省为对象的研究。广东靠近香港、远离北京，并有反感中央统治的历史，它也许代表的是例外而非常态。大约二十年前，已有若干论文描述文化大革命早期在多个省份的情况。\footnote{The
  Cultural Revolution in The Provinces, Harvard East Asian Monographs,
  42 (Cambridge, Mass.: Harvard University Press, 1971).

  The Cultural Revolution in The Provinces, Harvard East Asian
  Monographs, 42 (Cambridge, Mass.: Harvard University Press, 1971)。}
这些论文出现时，中国的大门对外界紧闭，研究者接触省级报纸和其他第一手资料极为有限。至于
20 世纪 70 年代初至中期，文献中的空白更大。

This study, by a detailed analysis of events in the important eastern
province of Zhejiang --- a province located in the heartland of Han
Chinese political, economic and social power and culture -- attempts to
fill the gap. In essence, it presents a revised view of the influence of
the influence wielded by local rebels turned radical politicians in
provincial affairs. While not making any dogmatic assertions about the
applicability of these findings to other areas of China, the study
suggests that this interpretation may well have validity beyond the
borders of Zhejiang. The book describes and analyzes the nature and
composition of the mass organizations which were established at the end
of 1966 and traces the evolution of their structure, tactics and
activities over the ten-year period. It assesses, through a careful
sifting of the available evidence, the nature and extent of their power
and influence in the province, particularly in the early to mid 1970s.
The study also investigates and draws out the links between prominent
personnel of the mass organizations and central and provincial civilian
and military party leaders.

本研究通过对东部重要省份浙江事件的详细分析，试图填补这一空白。浙江位于汉人政治、经济、社会权力和文化的核心地带。实质上，本书对地方造反者转变为激进政治家后在省级事务中所拥有的影响力，提出一种修正性看法。尽管本研究并不武断断言这些发现适用于中国其他地区，但它提示，这一解释很可能具有超出浙江边界的有效性。本书描述并分析
1966
年底成立的群众组织的性质和构成，并追踪其结构、策略和活动在十年期间的演变。通过仔细筛选现有证据，本书评估这些组织在该省的权力和影响的性质与范围，特别是在
20 世纪 70
年代初至中期。本研究还考察并揭示群众组织中突出人物与中央及省级文职、军事党领导之间的联系。

It documents in some detail the methods which the rebels used in the
political campaigns of 1973-75 to undermine the authority of the
provincial leadership. It points to the organizational power of the
local radicals as manifest in their control over the appointments'
mechanism and their ability to use officially sanctioned organizations
as vehicles to challenge and undermine the provincial party leadership.
The study highlights the innovative and effective strategy of operating
both within the formal political system, as holders of party and
government posts, and outside it as rebels intent on weakening the very
same system. It is suggested that the change in content and style of
rebellion and factionalism from the mid 1960s to the mid 1970s goes a
long way to help explain the success, however limited and temporary, of
the rebel politicians. Thus, the twin themes of rebellion and
factionalism underpin the chronological narrative which forms the
framework of the book.

本书较为详细地记录了造反派在 1973-75
年政治运动中用来削弱省级领导权威的方法。它指出，地方激进派的组织力量，体现在他们对任命机制的控制，以及他们利用官方认可的组织作为工具来挑战和削弱省党委领导的能力。本研究突出了一种创新且有效的策略：一方面作为党政职位持有者在正式政治体系内部运作，另一方面又作为意在削弱同一体系的造反者在体系外行动。本文认为，从
20 世纪 60 年代中期到 70
年代中期，造反和派性的内容与风格发生变化，这在很大程度上有助于解释造反政治家的成功，尽管这种成功有限且短暂。因此，造反与派性这两个主题，支撑着构成本书框架的编年叙事。

Before turning to a more detailed discussion of these themes, however,
it is necessary to review the source material which has formed the
empirical basis of the book. The principal references drawn on in the
following pages are the official newspapers of the Chinese Communist
Party's (CCP) Zhejiang Provincial Committee (ZPC) and the Hangzhou
Municipal Committee (HMC), the Zhejiang Daily (浙江日报) and the
Hangzhou Daily (杭州日报) respectively. The latter newspaper is, except
for the odd issue, unavailable outside China. In 1977, when my wife and
I arrived to teach in Hangzhou, all local newspapers were categorized as
internal (内部) publications and declared off-limits to foreigners. It
was the permission I received early in 1979 to read the local press that
provided the topic for the doctoral thesis on which this book is
based.\footnote{Keith Forster, ``The Hangzhou Incident of 1975: The
  Impact of Factionalism on a Chinese Provincial Administration'', (Ph.
  D. thesis: University of Adelaide, 1985).

  Keith Forster, ``The Hangzhou Incident of 1975: The Impact of
  Factionalism on a Chinese Provincial Administration'', (Ph. D. thesis:
  University of Adelaide, 1985)。}

然而，在更详细讨论这些主题之前，有必要回顾构成本书经验基础的资料来源。以下各页主要引用的，是中国共产党（CCP）浙江省委（ZPC）和杭州市委（HMC）的机关报，分别为《浙江日报》和《杭州日报》。后一份报纸除个别期号外，在中国境外无法取得。1977
年，我和妻子到杭州任教时，所有地方报纸都被列为内部出版物，并宣布外国人不得阅读。1979
年初，我获准阅读地方报刊，正是这一许可提供了本书所依据博士论文的题目。\footnote{Keith
  Forster, ``The Hangzhou Incident of 1975: The Impact of Factionalism
  on a Chinese Provincial Administration'', (Ph. D. thesis: University
  of Adelaide, 1985).

  Keith Forster, ``The Hangzhou Incident of 1975: The Impact of
  Factionalism on a Chinese Provincial Administration'', (Ph. D. thesis:
  University of Adelaide, 1985)。}

Understandably, the reliability and usefulness of such publications,
which convey policy statements and party propaganda, is often
questioned. The author has repeatedly been reminded of the dangers which
stem from reliance on such partial and in many ways unreliable sources
of information as party newspapers. In the period 1976-79, when some of
the references used in this study were published, the party and its
propaganda arm were undertaking an extensive and thorough political
campaign against the local followers of the Gang of Four. Distortion,
exaggeration, vindictiveness and personal abuse all featured prominently
in the campaign. Yet for all their partiality and crude dogma, Chinese
newspapers provide an indispensable window into the thinking and
behavior of officialdom and the workings of the political system.
Articles published during the campaign against the Gang of Four
contained detailed, descriptive passages, tendentious in tone but
offering valuable insight into the workings of a secretive political
system and allowing the critical and discriminating reader to extract,
with the necessary qualifications, information of great value.

可以理解，这类传达政策声明和党内宣传的出版物，其可靠性和有用性常常受到质疑。作者曾一再被提醒，依赖党报这样有偏向且在许多方面并不可靠的信息来源，会带来危险。在
1976-79
年间，即本研究所用部分资料发表的时期，党及其宣传机构正在针对``四人帮''地方追随者开展一场广泛而彻底的政治运动。扭曲、夸张、报复性和人身攻击都是这场运动中的突出特征。然而，尽管中国报纸带有偏向且教条粗糙，它们仍提供了一扇不可或缺的窗口，使人得以观察官方的思想和行为，以及政治体系的运作。反``四人帮''运动期间发表的文章包含详尽的描述性段落，语气虽有倾向性，却为理解一个隐秘政治体系的运作提供了宝贵洞见，使具有批判力和辨别力的读者能够在作出必要限定后，提取极有价值的信息。

Additionally, two sets of official documents, published in 1977 and 1978
respectively, setting out in great detail the regime's case against two
of the principal rebel leaders, have been extensively drawn upon. The
documents consist of photocopies of signed statements by detained rebels
and witnesses to their alleged crimes. In effect, the file sets out the
case for the prosecution. However, due to careless editing, names and
dates blacked out from the typed version have been left untouched in the
original handwriting, testifying to the authenticity of the documents.
The invaluable publications have in many instances fleshed out or
confirmed incidents or events referred to in the press. However,
wherever I have not been able to check or verify various claims and
allegations, such qualifiers as ``reportedly'' and ``allegedly'' and so
on, inelegant as they are, have been used.

此外，本书还大量使用了分别于 1977 年和 1978
年发表的两套官方文件；这些文件极为详细地陈述了政权针对两名主要造反派领袖的案情。文件由被拘押造反派以及其所谓罪行见证人的签名陈述复印件组成。实际上，这些档案列出了控方案情。然而，由于编辑粗心，打字稿中被涂黑的姓名和日期，在原始手写材料中仍未遮盖，这证明了文件的真实性。这些极有价值的出版物在许多情况下充实或证实了报刊中提到的事件。不过，凡是我无法核查或证实的各种说法和指控，我都使用了``据报道''\,``据称''等限定语，尽管它们并不优雅。

Chinese-language sources have been supplemented by a thorough reading of
provincial broadcasts from the Zhejiang Provincial Service (ZPS) as
monitored by the BBC Summary of World Broadcasts (SWB), US Foreign
Broadcast Information Service (FBIS), Hong Kong Union Research Service
(URS) and China News Summary (CNS). Guided tours and stays varying from
a half-day to a fortnight at state-run factories and rural people's
communes provided further material for the study. Nevertheless, given my
status as a ``foreign expert'' and the political climate of the years
1977-79, so different in many respects to the more open atmosphere
prevailing until recently, I felt constrained as to how far I could
probe into politically-sensitive issues.

中文资料之外，我还通读了由
BBC《世界广播摘要》（SWB）、美国外国广播信息服务（FBIS）、香港友联研究所《友联研究服务》（URS）和《中国新闻摘要》（CNS）监测的浙江省广播服务（ZPS）的省级广播。对国营工厂和农村人民公社的参观和停留，从半天到两周不等，也为本研究提供了更多材料。不过，鉴于我``外国专家''的身份，以及
1977-79
年间的政治气候在许多方面与不久前仍存在的较开放氛围大不相同，我感到自己在探查政治敏感问题的深度上受到限制。

Return trips in 1988-89 proved much more fruitful and I was able to hold
extensive interviews and discussions on local politics in a much more
uninhibited environment. My lengthy acquaintance with Hangzhou, now
extending over twelve years, had by now enabled me to build up a group
of friends and associates who had participated in, were witness to, or
were familiar with many of the events described in the following pages.
However, one unavoidable gap in the sources relied on for this study is
the absence of Red Guard materials, available in abundance for places
like Beijing and Guangzhou. All my efforts to obtain access to such
documents have been frustrated. The lack of this material makes a
detailed analysis of the social background of members of the major mass
organizations in Zhejiang almost impossible. Informed guesswork,
speculation, and comparative reference to studies which have drawn on
such materials is thus unavoidable, particularly in chapter one.

1988-89
年的回访则收获大得多，我能够在更少拘束的环境中，就地方政治进行广泛访谈和讨论。我与杭州相识已超过十二年，这使我此时得以建立一个由朋友和熟人组成的群体，他们曾参与、见证，或熟悉下文所述许多事件。然而，本研究所依赖的资料中存在一个不可避免的缺口，即缺少红卫兵材料；这类材料在北京、广州等地则相当丰富。我为取得这类文件所作的一切努力都受挫。缺少这一材料，使得对浙江主要群众组织成员社会背景作详细分析几乎不可能。因此，基于信息的猜测、推断，以及同使用过此类材料的研究进行比较参照，便不可避免，尤其是在第一章中。

\section{Factionalism}\label{factionalism}

\section{派性}\label{ux6d3eux6027}

The word factionalism is used throughout this study and an explanation
of what it means is warranted from the outset. The CCP has always
defined factionalism (派性) pejoratively and contrasted it with party
spirit (党性), the fundamental attribute expected of all party members.
Those who practise factionalism are accused of placing loyalty to an
illegitimate group (as defined by Leninist organizational norms and
party statutes) above that to the party as a whole. Thus, by definition,
factions engage in conspiratorial activities which threaten the unity
and discipline of the larger group and may even attempt to take over its
leadership.

本研究通篇使用``派性''一词，因此有必要从一开始就说明其含义。中国共产党一直以贬义方式界定派性，并将其与党性相对照；党性被视为所有党员应具备的根本属性。搞派性者被指责为把对非法集团（按列宁主义组织规范和党章界定）的忠诚置于对整个党的忠诚之上。因此，按定义而言，派别从事威胁更大群体的团结和纪律的阴谋活动，甚至可能试图夺取其领导权。

Western political science has grappled with the concept of faction in
its many manifestations in various political environments. Even so,
according to the editors of a major collection of writings on the
subject, ``The very idea of faction comes to us with an inheritance
whose affective content is a highly negative evaluation seemingly
indelibly associated with the term''.\footnote{Frank P. Belloni and
  Dennis C. Beller, ``Party and Faction: Modes of Political
  Competition'', in Belloni and Beller (eds), Faction Politics:
  Political Parties and Factionalism in Comparative Perspective (Santa
  Barbara, California: ABC-Clio Inc., 1978), p.~445.

  Frank P. Belloni and Dennis C. Beller, ``Party and Faction: Modes of
  Political Competition'', in Belloni and Beller (eds), Faction
  Politics: Political Parties and Factionalism in Comparative
  Perspective (Santa Barbara, California: ABC-Clio Inc., 1978), p.~445。}
More important than making derogatory remarks about factions, argue
Belloni and Beller, is the need to define their structure, function and
causes. But first a definition is required. The most encompassing
definition, and the one most relevant here, has been provided by Ann
Fenwick from her reading of both empirical studies and theoretical
literature. She defines a faction as ``a noncorporate political conflict
group recruited by a leader on the basis of diverse
principles''.\footnote{Ann E. Fenwick, ``The Gang of Four and the
  Politics of Opposition: China, 1971-1976'', (Ph. D. thesis: Stanford
  University, 1984), p.~20.

  Ann E. Fenwick, ``The Gang of Four and the Politics of Opposition:
  China, 1971-1976'', (Ph. D. thesis: Stanford University, 1984),
  p.~20。}

西方政治学一直努力处理``派别''这一概念在各种政治环境中的诸多表现。即便如此，据一部关于该主题的重要文集的编者所言，``派别这一观念来到我们面前时，继承了一种情感内容，即一个似乎已与该术语不可磨灭地联系在一起的高度负面评价''。\footnote{Frank
  P. Belloni and Dennis C. Beller, ``Party and Faction: Modes of
  Political Competition'', in Belloni and Beller (eds), Faction
  Politics: Political Parties and Factionalism in Comparative
  Perspective (Santa Barbara, California: ABC-Clio Inc., 1978), p.~445.

  Frank P. Belloni and Dennis C. Beller, ``Party and Faction: Modes of
  Political Competition'', in Belloni and Beller (eds), Faction
  Politics: Political Parties and Factionalism in Comparative
  Perspective (Santa Barbara, California: ABC-Clio Inc., 1978), p.~445。}
Belloni 和 Beller
认为，比对派别作贬斥性评论更重要的，是界定其结构、功能和成因。但首先需要一个定义。Ann
Fenwick
在阅读经验研究和理论文献后，提出了最具包容性、也最适用于此处的定义。她把派别定义为``由一名领袖基于多种原则招募而成的非社团化政治冲突群体''。\footnote{Ann
  E. Fenwick, ``The Gang of Four and the Politics of Opposition: China,
  1971-1976'', (Ph. D. thesis: Stanford University, 1984), p.~20.

  Ann E. Fenwick, ``The Gang of Four and the Politics of Opposition:
  China, 1971-1976'', (Ph. D. thesis: Stanford University, 1984),
  p.~20。}

This definition is sufficiently broad to allow for different degrees of
organizational coherence and does not restrict the \emph{raison d'être}
of factions to the pursuit of power, although clearly this will be one
of their primary objectives. The attention that factions devote to
questions of power in no way detracts from their propensity and ability
to formulate policy positions and express them in certain ideological
forms. One analyst of the Chinese political scene\footnote{Lucien Pye,
  The Dynamics of Chinese Politics (Cambridge, Mass.: Oelgeschlager,
  Gunn \& Hain, 1981).

  Lucien Pye, The Dynamics of Chinese Politics (Cambridge, Mass.:
  Oelgeschlager, Gunn \& Hain, 1981)。} has attempted, unsuccessfully in
my opinion, to isolate policy and ideology from factional disputes and
to portray factionalism in the Chinese context as a cultural expression
of dependency and security on the part of political superiors and
subordinates.

这个定义足够宽泛，允许不同程度的组织凝聚力，也不把派别的存在理由局限于追求权力，尽管这显然会是其主要目标之一。派别对权力问题的关注，丝毫不削弱它们形成政策立场并以某些意识形态形式表达这些立场的倾向和能力。一位中国政治场景分析者\footnote{Lucien
  Pye, The Dynamics of Chinese Politics (Cambridge, Mass.:
  Oelgeschlager, Gunn \& Hain, 1981).

  Lucien Pye, The Dynamics of Chinese Politics (Cambridge, Mass.:
  Oelgeschlager, Gunn \& Hain, 1981)。}
曾试图把政策和意识形态同派性争端分离开来，并把中国语境中的派性描绘成政治上下级之间依赖和安全感的文化表达；在我看来，这一尝试并不成功。

As Beller and Belloni correctly point out, the causes of factionalism
are manifold and include societal, political and structural
factors.\footnote{Beller and Belloni, Faction Politics, pp.~430-7. See
  also, Norman K. Nicholson, ``The Factional Model and The Study of
  Politics'', Comparative Political Studies, 5: 3 (1972), pp.~303-05.

  Beller and Belloni, Faction Politics, pp.~430-7。另见 Norman K.
  Nicholson, ``The Factional Model and The Study of Politics'',
  Comparative Political Studies, 5: 3 (1972), pp.~303-05。} They
differentiate three types of factional organizations in terms of their
structure, recruitment pattern and members' self-awareness.\footnote{Beller
  and Belloni, Faction Politics, pp.~419-30.

  Beller and Belloni, Faction Politics, pp.~419-30。} The first type is
classified as the factional clique or tendency. It is basically an
unstructured, informal group whose membership has a low level of
cognition of its common interests or bonds. In communist societies,
where the expression of political interests is normally tightly
controlled and directed, such groups are not tolerated.

正如 Beller 和 Belloni
正确指出的，派性的成因是多重的，包括社会、政治和结构性因素。\footnote{Beller
  and Belloni, Faction Politics, pp.~430-7. See also, Norman K.
  Nicholson, ``The Factional Model and The Study of Politics'',
  Comparative Political Studies, 5: 3 (1972), pp.~303-05.

  Beller and Belloni, Faction Politics, pp.~430-7。另见 Norman K.
  Nicholson, ``The Factional Model and The Study of Politics'',
  Comparative Political Studies, 5: 3 (1972), pp.~303-05。}
他们依据结构、招募模式和成员自我意识，将派性组织区分为三种类型。\footnote{Beller
  and Belloni, Faction Politics, pp.~419-30.

  Beller and Belloni, Faction Politics, pp.~419-30。}
第一类被归为派性小集团或倾向。它基本上是一个无结构的非正式群体，其成员对共同利益或纽带的认知水平较低。在共产党社会中，政治利益的表达通常受到严密控制和引导，此类群体并不被容忍。

Beller and Belloni define the second type of faction as the personal,
client group faction whose organization is based on vertical links
between leaders and followers and whose membership is recruited
personally by these leaders. It is this type of faction, of intermediate
duration and whose existence does not outlive the political life of its
leader, that Andrew Nathan described in his important article on the
subject.\footnote{Andrew J. Nathan, ``A Factionalism Model for CCP
  Politics'', China Quarterly (CQ), 53 (1973), pp.~34-66.

  Andrew J. Nathan, ``A Factionalism Model for CCP Politics'', China
  Quarterly (CQ), 53 (1973), pp.~34-66。} Norman Nicholson seems to
refer to this type of faction when he points out that a distinction
should be drawn between the ``core members'' of a faction, whose
commitment to the leader and the faction is far stronger than that of
its ordinary members, who respond in different measure to particular
issues and whose level of factional participation is thus of a more
intermittent nature.\footnote{Nicholson, ``The Factional Model'',
  p.~298.

  Nicholson, ``The Factional Model'', p.~298。} The factions of the
mid-1970s in Zhejiang, whose structure and activities are recounted at
length in chapter seven of this study, belong to this type.

Beller 和 Belloni
将第二类派别定义为个人性、庇护关系式派别，其组织建立在领导者与追随者之间的纵向联系之上，成员由这些领导者亲自招募。Andrew
Nathan
在其关于该主题的重要文章中描述的，正是这种持续时间中等、其存在不会长于领袖政治生命的派别。\footnote{Andrew
  J. Nathan, ``A Factionalism Model for CCP Politics'', China Quarterly
  (CQ), 53 (1973), pp.~34-66.

  Andrew J. Nathan, ``A Factionalism Model for CCP Politics'', China
  Quarterly (CQ), 53 (1973), pp.~34-66。} Norman Nicholson
似乎也指涉这类派别；他指出，应区分派别的``核心成员''与普通成员，前者对领袖和派别的承诺远强于后者，后者对具体问题的回应程度不同，因此其派性参与程度更具间歇性。\footnote{Nicholson,
  ``The Factional Model'', p.~298.

  Nicholson, ``The Factional Model'', p.~298。}
本研究第七章详细叙述其结构和活动的 20 世纪 70
年代中期浙江派别，就属于这一类型。

The third type of faction, according to Beller and Belloni, is the
institutionalized or organized faction, whose membership is more open,
whose group cognitive awareness is high and which establishes and
maintains horizontal linkages. The mass organizations of the 1966-68
period of the Cultural Revolution conform to this type and were also
characterized by hierarchy and functional specialization. Red Guard
headquarters and umbrella rebel organizations established their own
publications, armed detachments and outside liaison stations to shore up
group coherence and facilitate operations. The pattern continued into
the mid-1970s with more sophistication and greater effectiveness, in
Zhejiang at least. By this time, factional groupings in Zhejiang
exhibited features attributed to both the second and third types.

据 Beller 和 Belloni
所言，第三类派别是制度化或组织化派别，其成员资格更加开放，群体认知意识较高，并建立和维持横向联系。文化大革命
1966-68
年期间的群众组织符合这一类型，并且也具有等级制和功能专业化特征。红卫兵总部和伞状造反组织建立自己的出版物、武装分队和外地联络站，以巩固群体凝聚力并便利行动。至少在浙江，这一模式以更成熟、更有效的形式延续到
20 世纪 70
年代中期。到此时，浙江的派性集团表现出第二类和第三类的双重特征。

In synthesizing the patterns of factional behavior Beller and Belloni
distinguish between closed, elitist factionalism and open factionalism.
The former is common to the first two types of factional structures
mentioned above. Open factionalism is indicative of institutionalized
factions. This study will demonstrate that it was the combination of the
open, organizational factionalism characteristic of the period 1966-69,
with the closed, elitist, clientelist factionalism widespread in the
upper echelons of the CCP, which gave the factionalism of the period
1973-76 in Zhejiang its unique and dangerous features.

Beller 和 Belloni
在综合派性行为模式时，区分了封闭的精英主义派性与公开派性。前者常见于上述前两类派性结构。公开派性则是制度化派别的标志。本研究将表明，正是
1966-69
年期间那种公开的、组织化的派性，与中国共产党高层普遍存在的封闭的、精英主义的、庇护关系式派性相结合，赋予了
1973-76 年间浙江派性以独特而危险的特征。

\section{Rebels and Conservatives in the Cultural
Revolution}\label{rebels-and-conservatives-in-the-cultural-revolution}

\section{文化大革命中的造反派与保守派}\label{ux6587ux5316ux5927ux9769ux547dux4e2dux7684ux9020ux53cdux6d3eux4e0eux4fddux5b88ux6d3e}

One issue which continues to bedevil analysts of the Cultural Revolution
relates to the nature of the mass organizations which participated in
it. In virtually every province they divided into two opposed groupings
but the basis on which these divisions occurred remains a matter of
conjecture and debate. William Hinton has observed that the divisions

一个持续困扰文化大革命研究者的问题，涉及参与其中的群众组织的性质。几乎在每一个省份，它们都分裂为两个相互对立的集团，但这种分裂发生的基础仍是推测和争论的问题。William
Hinton 观察到，这种分裂：

\begin{quote}
occurred with such regularity and persistence that it had to be
recognized as some sort of law of the political sphere as universal as
Boyle's law in chemistry or Newton's law in physics.\footnote{William
  Hinton, \emph{Shenfan} (London: Picador Books, 1983), p.~611.

  William Hinton, \emph{Shenfan} (London: Picador Books, 1983), p.~611。}
\end{quote}

\begin{quote}
以如此规律且持久的方式发生，以至于不得不承认它是政治领域的某种规律，如同化学中的波义耳定律或物理学中的牛顿定律一样具有普遍性。\footnote{William
  Hinton, \emph{Shenfan} (London: Picador Books, 1983), p.~611.

  William Hinton, \emph{Shenfan} (London: Picador Books, 1983), p.~611。}
\end{quote}

In 1984, when the CCP carried out an ideological and political critique
of the Cultural Revolution, an authoritative article in the party
journal Red Flag declared that judged from their behavior, guiding
ideology and their professed loyalty to Mao Zedong, both the so-called
``rebel'' and ``conservative'' (or ``royalist'') mass organizations were
indistinguishable.\footnote{Zhang Yun, ``增强党性克服派性''(Boost party
  spirit, overcome factionalism), Hongqi, No.~9 (1984), pp.~6-8. See
  also Keith Forster, ``The Repudiation of the Cultural Revolution'',
  Journal of Contemporary Asia, 17: 1 (1987), pp.~71-72.

  张云，``增强党性克服派性''(Boost party spirit, overcome
  factionalism)，《红旗》，1984 年第 9 期，pp.~6-8。另见 Keith Forster,
  ``The Repudiation of the Cultural Revolution'', Journal of
  Contemporary Asia, 17: 1 (1987), pp.~71-72。} According to this
analysis, which has become the prevailing CCP line on the subject, it
was the scramble for power which divided the rebels and split them into
two organizations standing opposed to or supportive of leading party
cadres and officials. There is much in the present party position that
is valid, as the discussion below will indicate. Yet it ignores the
important economic and social factors which propelled individuals into
the political arena and were instrumental in determining their
subsequent actions and behavior.

1984
年，当中国共产党对文化大革命进行意识形态和政治批判时，党的刊物《红旗》上一篇权威文章宣称，从行为、指导思想以及其自称对毛泽东的忠诚来看，所谓``造反派''和``保守派''（或``保皇派''）群众组织没有区别。\footnote{Zhang
  Yun, ``增强党性克服派性''(Boost party spirit, overcome factionalism),
  Hongqi, No.~9 (1984), pp.~6-8. See also Keith Forster, ``The
  Repudiation of the Cultural Revolution'', Journal of Contemporary
  Asia, 17: 1 (1987), pp.~71-72.

  张云，``增强党性克服派性''(Boost party spirit, overcome
  factionalism)，《红旗》，1984 年第 9 期，pp.~6-8。另见 Keith Forster,
  ``The Repudiation of the Cultural Revolution'', Journal of
  Contemporary Asia, 17: 1 (1987), pp.~71-72。}
按照这一已成为中共关于该问题主流路线的分析，正是权力争夺使造反派分裂，并分化为反对或支持党内领导干部和官员的两个组织。如下文讨论将显示，当前党的立场中有许多有效之处。然而，它忽视了将个人推入政治舞台，并在决定其随后行动和行为方面起重要作用的经济和社会因素。

Studies published in the West on this subject have examined the
composition of mass organizations in the southern city of
Guangzhou.\footnote{For example, see Hong Yung Lee, The Politics of the
  Chinese Cultural Revolution: A Case Study (Berkeley and Los Angeles:
  University of California Press, 1978); Anita Chan, Stanley Rosen and
  Jonathan Unger, ``Students and Class Warfare: The Social Roots of the
  Red Guard Conflict in Guangzhou (Canton)'', CQ, 83, (September 1980),
  pp.~397-446.

  例如，见 Hong Yung Lee, The Politics of the Chinese Cultural
  Revolution: A Case Study (Berkeley and Los Angeles: University of
  California Press, 1978)；Anita Chan, Stanley Rosen and Jonathan Unger,
  ``Students and Class Warfare: The Social Roots of the Red Guard
  Conflict in Guangzhou (Canton)'', CQ, 83, (September 1980),
  pp.~397-446。} Hong Yung Lee has found that radical mass organizations
recruited their members from social groups that held grievances against
the status quo. Among students, radicals included those with a ``bad''
family background, those sent to rural areas and those who had revolted
against the work teams sent onto campuses in May 1966. Workers who were
employed in smaller, poorer factories, contract or temporary workers,
apprentices and unskilled workers in larger factories and individual
laborers tended to join radical organizations. Among the ranks of
cadres, low or middle ranking officials, those with bad family
background (capitalist, landlord, etc.) and those who had been
disciplined or punished before the Cultural Revolution, were more likely
to side with radical organizations. Conservative mass organizations, on
the other hand, contained members with good class background (worker,
peasant, soldier) and were mobilized by political organizations such as
the CCP, Communist Youth League (CYL) and trade unions.

西方关于这一主题的研究，考察了南方城市广州群众组织的构成。\footnote{For
  example, see Hong Yung Lee, The Politics of the Chinese Cultural
  Revolution: A Case Study (Berkeley and Los Angeles: University of
  California Press, 1978); Anita Chan, Stanley Rosen and Jonathan Unger,
  ``Students and Class Warfare: The Social Roots of the Red Guard
  Conflict in Guangzhou (Canton)'', CQ, 83, (September 1980),
  pp.~397-446.

  例如，见 Hong Yung Lee, The Politics of the Chinese Cultural
  Revolution: A Case Study (Berkeley and Los Angeles: University of
  California Press, 1978)；Anita Chan, Stanley Rosen and Jonathan Unger,
  ``Students and Class Warfare: The Social Roots of the Red Guard
  Conflict in Guangzhou (Canton)'', CQ, 83, (September 1980),
  pp.~397-446。} Hong Yung Lee
发现，激进群众组织从对现状怀有不满的社会群体中招募成员。在学生中，激进派包括家庭出身``不好''者、被送往农村者，以及在
1966 年 5
月反抗派驻校园工作组者。在较小、较贫困工厂工作的工人、合同工或临时工、大厂中的学徒和非熟练工人，以及个体劳动者，往往加入激进组织。在干部队伍中，低级或中级官员、家庭出身不好（资本家、地主等）者，以及文化大革命前曾受纪律处分或惩罚者，更可能站在激进组织一边。另一方面，保守派群众组织则包含阶级出身好（工人、农民、士兵）的成员，并由中国共产党、共产主义青年团（CYL）和工会等政治组织动员。

While radicals defined class pragmatically, by examining an individual's
behavior and performance, their opponents tended to look to the existing
``caste'' structure,\footnote{For a useful discussion of class as caste,
  see R.C. Kraus, Class Conflict in Chinese Socialism (New York:
  Columbia University Press, 1981), pp.~117-39.

  关于把阶级视为种姓的有益讨论，见 R.C. Kraus, Class Conflict in Chinese
  Socialism (New York: Columbia University Press, 1981), pp.~117-39。}
based on pre-revolution categories. Class, however, was a less important
consideration for university students than it was for their secondary
school juniors.\footnote{Jonathan Unger, Education Under Mao: Class and
  Competition in Canton Schools, 1960-1980 (New York: Columbia
  University Press, 1982), pp.~130-31.

  Jonathan Unger, Education Under Mao: Class and Competition in Canton
  Schools, 1960-1980 (New York: Columbia University Press, 1982),
  pp.~130-31。} While radicals directed their attacks against party
cadres, conservatives concentrated their fire on the old bourgeois
class, thus deflecting the movement away from the real power-holders.
Radicals broadened, while conservatives restricted, the scope of attack.
The conservatives generally maintained good relations with the PLA,
while the radicals established links with the Central Cultural
Revolution Group under Chen Boda and Jiang Qing. Both groups, however,
accepted the unchallenged supremacy of Mao Zedong and his
Thought.\footnote{Lee, The Politics of the Chinese Cultural Revolution,
  pp.~340-43.

  Lee, The Politics of the Chinese Cultural Revolution, pp.~340-43。}
Rebel leaders tended to be children of intellectuals while conservative
or loyalist organizations were generally led by cadres' children. Rebels
attacked the number one leader of a unit while the loyalists attacked
its deputy.\footnote{Chan et al., ``Students and Class Warfare'',
  pp.~433-42.

  Chan et al., ``Students and Class Warfare'', pp.~433-42。}

激进派通过考察个人行为和表现，以务实方式界定阶级；其对手则倾向于依据革命前类别，诉诸既有的``种姓''结构。\footnote{For
  a useful discussion of class as caste, see R.C. Kraus, Class Conflict
  in Chinese Socialism (New York: Columbia University Press, 1981),
  pp.~117-39.

  关于把阶级视为种姓的有益讨论，见 R.C. Kraus, Class Conflict in Chinese
  Socialism (New York: Columbia University Press, 1981), pp.~117-39。}
然而，对大学生而言，阶级考虑并不像对中学生那样重要。\footnote{Jonathan
  Unger, Education Under Mao: Class and Competition in Canton Schools,
  1960-1980 (New York: Columbia University Press, 1982), pp.~130-31.

  Jonathan Unger, Education Under Mao: Class and Competition in Canton
  Schools, 1960-1980 (New York: Columbia University Press, 1982),
  pp.~130-31。}
激进派把攻击指向党内干部，保守派则集中火力攻击旧资产阶级，从而使运动偏离真正的掌权者。激进派扩大攻击范围，保守派则限制攻击范围。保守派通常与人民解放军保持良好关系，而激进派则同陈伯达、江青领导下的中央文化革命小组建立联系。不过，两派都接受毛泽东及其思想无可挑战的最高地位。\footnote{Lee,
  The Politics of the Chinese Cultural Revolution, pp.~340-43.

  Lee, The Politics of the Chinese Cultural Revolution, pp.~340-43。}
造反派领袖往往是知识分子子女，而保守派或忠诚派组织通常由干部子女领导。造反派攻击一个单位的一把手，而忠诚派则攻击其副手。\footnote{Chan
  et al., ``Students and Class Warfare'', pp.~433-42.

  Chan et al., ``Students and Class Warfare'', pp.~433-42。}

A recently-published tract from a Chinese participant in the Cultural
Revolution has divided the mass organizations along the same lines. It
differentiates between ``rebels'' and ``conservatives'' on the grounds
of the class composition of leaders and rank and file members, their
political experience before the Cultural Revolution, and whether they
belonged to the CCP or its mass organizations.\footnote{Anita Chan
  (ed.), ``A Brief Analysis of the Cultural Revolution'' by Liu Guokai,
  in Chinese Sociology and Anthropology, 19: 2 (1986-7), esp.~pp.~73-82,
  85-92. Originally published as ``文化大革命''简析 in a Chinese
  mainland underground publication and then in Taiwan in
  大陆地下刊物汇编(A Collection of Mainland underground publications),
  Vol. 17 (Taibei: Institute for the Study of Chinese Communist
  Problems, 1983). pp.~91-244.

  Anita Chan (ed.), ``A Brief Analysis of the Cultural Revolution'' by
  Liu Guokai, in Chinese Sociology and Anthropology, 19: 2
  (1986-7)，尤其 pp.~73-82,
  85-92。原以``文化大革命''简析为题发表于中国大陆地下出版物，后收入台湾出版的《大陆地下刊物汇编》(A
  Collection of Mainland underground publications), Vol. 17 (Taibei:
  Institute for the Study of Chinese Communist Problems, 1983),
  pp.~91-244。} Another participant in the Cultural Revolution, Gao
Yuan, suggests in his recollections that it was personal friendships,
grudges and other particularistic relationships that brought together or
divided individuals.\footnote{Gao Yuan, Born Red (Stanford: Stanford
  University Press, 1987).

  Gao Yuan, Born Red (Stanford: Stanford University Press, 1987)。}

一位文化大革命中国参与者近期发表的小册子，也沿同样脉络划分群众组织。它依据领导人和普通成员的阶级构成、他们在文化大革命前的政治经历，以及他们是否属于中国共产党或其群众组织，来区分``造反派''和``保守派''。\footnote{Anita
  Chan (ed.), ``A Brief Analysis of the Cultural Revolution'' by Liu
  Guokai, in Chinese Sociology and Anthropology, 19: 2 (1986-7),
  esp.~pp.~73-82, 85-92. Originally published as ``文化大革命''简析 in a
  Chinese mainland underground publication and then in Taiwan in
  大陆地下刊物汇编(A Collection of Mainland underground publications),
  Vol. 17 (Taibei: Institute for the Study of Chinese Communist
  Problems, 1983). pp.~91-244.

  Anita Chan (ed.), ``A Brief Analysis of the Cultural Revolution'' by
  Liu Guokai, in Chinese Sociology and Anthropology, 19: 2
  (1986-7)，尤其 pp.~73-82,
  85-92。原以``文化大革命''简析为题发表于中国大陆地下出版物，后收入台湾出版的《大陆地下刊物汇编》(A
  Collection of Mainland underground publications), Vol. 17 (Taibei:
  Institute for the Study of Chinese Communist Problems, 1983),
  pp.~91-244。} 另一位文化大革命参与者 Gao Yuan
在回忆中则指出，是个人友谊、怨恨和其他特殊关系，使个人走到一起或彼此分裂。\footnote{Gao
  Yuan, Born Red (Stanford: Stanford University Press, 1987).

  Gao Yuan, Born Red (Stanford: Stanford University Press, 1987)。}

This study does not pretend to provide all the answers to these
important questions but it addresses them and, from limited data,
attempts to come up with some tentative hypotheses. Evidence to prove
the validity of the above generalizations for Zhejiang is decidedly thin
and ambiguous. There is evidence that both mass organizations put up
people of unimpeachable class background as leaders ahead of others of
greater ability but susceptible to criticism on class grounds. There is
also evidence that revenge for prior injustice, suffered in such
campaigns as the Socialist Education Movement (1963-65), motivated
certain individuals in particular to rebel against their leaders. But
other factors, such as personality differences and rivalries between
different mass organization leaders as well as their ability to attract
and retain the loyalty of followers, played an equal if not more
important role in the attachment of various individuals and social
groups to mass organizations.

本研究并不自称能够回答这些重要问题的全部，但它会处理这些问题，并在有限资料基础上提出若干试探性假设。证明上述概括适用于浙江的证据显然稀薄且含混。有证据表明，两个群众组织都把阶级背景无可指摘的人推上领导岗位，而不是让能力更强但在阶级问题上易受批评者担任领导。也有证据表明，在社会主义教育运动（1963-65）等运动中遭受既往不公后的复仇心理，尤其促使某些个人反抗其领导。但其他因素，例如不同群众组织领导人之间的性格差异和竞争，以及他们吸引并保持追随者忠诚的能力，在各种个人和社会群体依附于群众组织的过程中，发挥了同等甚至更重要的作用。

By definition, a revolution pits a ``progressive'' against a
``reactionary'' force. Mao Zedong sometimes saw the Cultural Revolution
as a continuation of the civil war between the CCP and the KMT, a view
out of touch with reality but conveying in simple, comprehensible
imagery the confrontation between good and evil. In Zhejiang it seems
that if the Central Cultural Revolution Group had not actually coined
the terms ``proletarian revolutionary'' and ``conservative'' to
differentiate between the two mass organizations, the logic and rhetoric
of the times would have required such a distinction. The ``royalist'' or
``conservative'' group in Zhejiang was in its own way, both by word and
deed, every bit as radical as its protagonist, a point noted by Zhou
Enlai in 1968 when trying to mediate a settlement between the two.

按定义，革命使一股``进步''力量与一股``反动''力量相对立。毛泽东有时把文化大革命视为中国共产党与国民党之间内战的延续，这种看法脱离现实，却以简单、易懂的意象传达善恶对抗。在浙江，似乎即使中央文化革命小组并未实际创造``无产阶级革命派''和``保守派''这两个术语来区分两个群众组织，当时的逻辑和修辞也会要求作出这种区分。浙江的``保皇派''或``保守派''集团，无论在言辞还是行动上，都以自己的方式与其对手一样激进；周恩来在
1968 年试图调停双方和解时就注意到这一点。

What, in essence, divided the two organizations, as the following
chapters of this study will show, was their assessment of the
revolutionary credentials of leading cadres, an assessment called for by
the ``proletarian headquarters'' in Beijing. This issue divided rebel
and Red Guard groups across the country, particularly in relation to the
number one provincial official. The determined efforts made in Zhejiang
in 1967 to induce one of the leaders of the so-called ``conservative''
organization to change sides because of the strength of his following,
his popularity and his considerable ability, suggests that class
background may not have been as important in determining allegiances as
more pragmatic considerations. In the view of some readers, the claim
that callow youth and hard-boiled workers beat, tortured and killed each
other over an issue as seemingly remote from their daily terms of
reference as the loyalty of political leaders to an ill-defined
``revolutionary line'' of Chairman Mao may detach political action from
its social basis. Nevertheless it seemed to be true of the behavior of
the rebels of both sides in Zhejiang.

正如本研究后续章节将显示，两个组织之间的实质分歧，在于它们对领导干部革命资格的评估，而这种评估是北京的``无产阶级司令部''所要求的。这个问题使全国各地的造反派和红卫兵群体发生分裂，尤其是在涉及省级一把手时。1967
年，浙江曾坚决努力劝诱所谓``保守派''组织的一名领导人改换阵营，原因在于他追随者众、声望高且能力相当强；这表明，在决定归属时，阶级背景可能不如更务实的考虑重要。在一些读者看来，声称稚嫩青年和老练工人会为了政治领导人是否忠于毛主席界定不明的``革命路线''这样一个似乎远离其日常参照的问题而相互殴打、折磨和杀害，或许会使政治行动脱离其社会基础。尽管如此，这似乎确是浙江双方造反者行为的真实情况。

The remainder of this introduction will sketch the history of the
Cultural Revolution in Zhejiang as it relates to the question of mass
organization factionalism discussed above. This exercise also serves the
purpose of providing a synopsis of the more detailed description and
analysis contained in later chapters.

本导论的其余部分，将概述浙江文化大革命的历史，重点是它与上述群众组织派性问题的关系。这一工作也旨在为后续章节中更详细的描述和分析提供一个概要。

\section{The Cultural Revolution in
Zhejiang}\label{the-cultural-revolution-in-zhejiang}

\section{浙江的文化大革命}\label{ux6d59ux6c5fux7684ux6587ux5316ux5927ux9769ux547d}

The origin of factional groupings in Zhejiang dates to the January 1967
power-seizure phase of the Cultural Revolution. With the leaders of the
CCP ZPC united against the attacks of the rebels, and the latter divided
as to which leaders of the committee merited support and which deserved
repudiation as capitalist-roaders, a political stalemate ensued. When
the military leaders of the Zhejiang Provincial Military District (ZPMD)
were ordered by the central authorities to support one rebel
organization in its attempt to overthrow the local administration, they
demurred and escorted their civilian colleagues to the safety of the
barracks.

浙江派性集团的起源，可追溯到文化大革命 1967 年 1
月夺权阶段。中共浙江省委领导团结起来抵御造反派攻击，而造反派内部则在省委哪些领导值得支持、哪些应作为走资派受到批判的问题上发生分裂，于是政治僵局形成。当浙江省军区（ZPMD）军事领导人奉中央命令，支持一个造反组织试图推翻地方行政当局时，他们表示异议，并将文职同僚护送到军营避难。

In February 1967, the mass organization officially recognized by the
Cultural Revolution Group in Beijing, the Zhejiang Provincial
Revolutionary Rebel United Headquarters (浙江省革命造反联合总指挥部),
decided to hold a mass rally to announce the dismissal of Jiang Hua
(江华), First Secretary of the CCP ZPC. Its rival group, the Zhejiang
Provincial Red Storm Provisional Headquarters
(浙江省红色暴动派临时指挥部), reportedly backed by members of the
provincial political elite, broke up the meeting after discovering that
Mao Zedong had decided to protect Jiang. On Zhou Enlai's instructions,
Jiang was flown to the comparative safety of Beijing. The incident
provoked the split between United Headquarters and Red Storm, a rift
which escalated during the following two years. Most of Jiang Hua's
colleagues were bundled out of office and subjected to varying forms and
degrees of punishment. Very few survived the 1967 upheaval and the
prolonged absence from office of experienced administrators contributed
to the instability which afflicted Zhejiang for the greater part of the
following decade.

1967 年 2
月，获得北京文化革命小组正式承认的群众组织``浙江省革命造反联合总指挥部''，决定召开群众大会，宣布罢免中共浙江省委第一书记江华。其竞争对手``浙江省红色暴动派临时指挥部''，据称得到省级政治精英成员支持，在发现毛泽东已决定保护江华后，冲散了会议。按周恩来指示，江华被飞机送往相对安全的北京。这一事件激发了``联合总部''与``红暴''之间的分裂，并在随后两年中不断升级。江华的大多数同僚被赶下台，并受到形式和程度不一的惩罚。很少有人挺过
1967
年的动荡；有经验行政人员长期离职，也助长了此后十年大部分时间困扰浙江的不稳定。

The ``revolutionary rebels'', together with the PLA and ``revolutionary
leading cadres'', were to join together as the Maoist-inspired trinity
of forces called revolutionary committees, established as governing
bodies at all administrative levels. United Headquarters, with the
blessing and backing of the Cultural Revolution Group in Beijing,
claimed recognition as the ``core'' of the proposed revolutionary
alliance of rebels. Red Storm refused to grant United Headquarters this
status. The ambiguous stance adopted toward Jiang Hua by the center gave
credibility to Red Storm's pretensions for acceptance as a revolutionary
organization. However, from February 1967 onward, United Headquarters
became known as the provincial ``rebel'' group (造反派) while Red Storm
labored under the opprobrious epithet of ``conservative'' faction
(保守派), labels bestowed by the Cultural Revolution Group. To avoid
confusion, this study retains the terms while recognizing that the
distinction, as it was then drawn, is in many ways misleading.

``革命造反派''将同人民解放军和``革命领导干部''结合在一起，组成毛主义启发下的三结合力量，即革命委员会，并作为各级行政层面的治理机构而建立。在北京文化革命小组的赞同和支持下，``联合总部''要求被承认为拟议中的造反派革命大联合的``核心''。``红暴''拒绝给予``联合总部''这一地位。中央对江华采取的含混立场，使``红暴''自称可被接受为革命组织的姿态具有了可信性。然而，自
1967 年 2
月起，``联合总部''被称为省级``造反派''，而``红暴''则背负``保守派''这一污名化称谓；这些标签均由文化革命小组赋予。为避免混淆，本研究保留这些术语，同时承认当时作出的这种区分在许多方面具有误导性。

The failure of the local military forces garrisoned in Hangzhou to
support the rebels led to Beijing's decision in late April 1967 to
despatch main force Unit 6409 (the 20th Army) from Jiangsu province to
take military control of Zhejiang. Together with personnel from the ZPMD
and Air Force and Naval units stationed in the province, the PLA
performed ``support the left'' activities as directed by Mao and his
deputy Lin Biao. However, the freedom given to the military to decide
which mass organization was ``revolutionary'' and therefore deserving of
its support, often worked to the detriment of United Headquarters. Only
the Air Force and 20th Army units proved reliable allies to this
organization. Naval and local military forces were sympathetic to Red
Storm. Outside the provincial capital of Hangzhou, sub-district and
People's Armed Forces Department (militia) forces contingents tended to
side with local party officials, especially against outside challenges
to their rule.

驻杭州地方军事力量未能支持造反派，促使北京在 1967 年 4
月下旬决定从江苏省调派主力 6409
部队（第二十军）对浙江实行军事管制。人民解放军同浙江省军区以及驻省空军、海军部队人员一道，按毛及其副手林彪指示开展``支左''活动。然而，军队被赋予决定哪个群众组织是``革命的''、因而值得支持的自由，这常常不利于``联合总部''。只有空军和第二十军部队证明是该组织可靠盟友。海军和地方军事力量则同情``红暴''。在省会杭州之外，分区和人民武装部（民兵）部队倾向于站在地方党内官员一边，特别是在面对外来挑战其统治时。

Mass factionalism thus spilled over into the armed forces. Additionally,
each mass organization attempted to gain support from the veteran cadres
of the defunct CCP ZPC by condemning some and praising and defending
others. Only those allowed to ``pass the test'' by United Headquarters
reappeared on the political stage by the time Mao convened the CCP 9th
Congress in April 1969. Before that, continuing controversy over the
status of Jiang Hua and several of his subordinates kept the mass
organizations in Zhejiang at loggerheads, implicating the PLA. In August
1967 leaders of the military district who had prevaricated over support
for United Headquarters were dismissed and replaced by the commanders of
the 20th Army, who were more enthusiastic in their display of loyalty
and obedience to Beijing.

群众派性由此外溢到武装力量中。此外，每个群众组织都试图通过谴责某些人、赞扬和保护另一些人，从已不存在的中共浙江省委老干部中争取支持。到毛在
1969 年 4
月召开中共九大时，只有那些被``联合总部''允许``过关''的人重新出现在政治舞台上。在此之前，围绕江华及其若干下属地位的持续争议，使浙江的群众组织长期对立，并牵连人民解放军。1967
年 8
月，在支持``联合总部''问题上含糊其辞的军区领导被免职，由第二十军指挥员取代；后者在向北京表达忠诚和服从方面更为热心。

Major incidents in the factories of Hangzhou and such provincial towns
as Wenzhou (温州), Jinhua (金华), Xiaoshan (萧山) and Zhuji (诸暨) in
the summer of 1967 inflamed factional antagonisms even further. Mao
Zedong visited Hangzhou for one day in September 1967 to observe the
situation at first hand. During his tour he called for unity among the
two major organizations but the call went unheeded. In order to break
the deadlock, in December 1967 Mao issued instructions which partially
rehabilitated Red Storm as an ``old rebel organization which has
committed mistakes''. The compromise undoubtedly pleased neither party.
United Headquarters was forced to negotiate with the old rival it had
fought unrelentingly against over the previous year. Red Storm, while
seemingly justified in its refusal to kneel to United Headquarters,
remained in an unequal, tainted position.

1967
年夏，杭州工厂以及温州、金华、萧山和诸暨等省内城镇发生的重大事件，进一步激化了派性对立。毛泽东于
1967 年 9
月到杭州停留一天，亲自观察局势。视察期间，他呼吁两大组织团结，但这一呼吁未被采纳。为打破僵局，毛在
1967 年 12
月发出指示，部分恢复``红暴''的地位，称其为``犯了错误的老造反组织''。这种妥协无疑令双方都不满意。``联合总部''被迫与过去一年中一直毫不妥协斗争的老对手谈判。``红暴''虽然似乎因拒绝向``联合总部''屈膝而获得正当性，却仍处于不平等且受玷污的位置。

A further two months of hard bargaining was required, under the
direction of an impatient center, before the two groups finally signed
an agreement in February 1968. The alliance permitted the formation of
the Zhejiang Provincial Revolutionary Committee (ZPRC) to take place in
the following month. However, because the clauses allotting a specified
number of places on the standing committee to representatives of Red
Storm were not honored, the agreement soon broke down. The newly
established political institution struggled to gain legitimacy and
respect throughout 1968. Ridiculed by remnant influential officials of
the old power structure and unrecognized by Red Storm, it was powerless
to discipline the opposition forces which openly defied its edicts. The
summer of 1968 witnessed further armed struggles, especially in the
isolated southern port-city of Wenzhou.

在不耐烦的中央指导下，又经过两个月艰难讨价还价，两个集团终于在 1968 年 2
月签署协议。这一联盟使浙江省革命委员会（ZPRC）得以在次月成立。然而，由于将常委会若干指定席位分配给``红暴''代表的条款没有得到履行，协议很快破裂。这个新成立的政治机构在整个
1968
年都努力争取合法性和尊重。它受到旧权力结构中残存的有影响力官员嘲笑，又不被``红暴''承认，因此无力惩戒公开违抗其命令的反对力量。1968
年夏见证了进一步武斗，尤其是在孤立的南部港口城市温州。

Throughout the two years of chaos and strife up until the end of 1968,
supporters of Jiang Hua both in Zhejiang and Beijing strove for his
inclusion in the new political leadership. Their efforts contributed
greatly to the destabilization in Zhejiang. They eventually failed but
not before they had provoked further rifts between United Headquarters
and Red Storm. It was not until November 1968, when Jiang Hua was
publicly criticized by name in a series of polemical articles carried in
the local press, that his political fate was sealed, temporarily at
least.

在截至 1968
年底的两年混乱和冲突中，浙江和北京的江华支持者都努力使其进入新的政治领导层。他们的努力极大助长了浙江的不稳定。他们最终失败了，但在失败之前已进一步激化``联合总部''和``红暴''之间的裂痕。直到
1968 年 11
月，地方报刊刊发一系列论战文章公开点名批判江华，他的政治命运才被确定下来，至少暂时如此。

After the 9th National Congress of ``unity and victory'' a further
attempt was made to unite Zhejiang's two mass organizations and thus
effectively demobilize them. An agreement to this effect was announced
in May 1969 under the supervision of the provincial authorities. But a
substantial number of Red Storm activists refused to accept the terms of
the agreement and traveled to Beijing to lodge their complaints. Not
surprisingly, they received a cold reception and were ordered back to
Hangzhou to undergo virtual internment in ``study classes''.

在``团结的大会，胜利的大会''即九大之后，又一次尝试被用来联合浙江两个群众组织，并由此有效解除其动员。1969
年 5
月，在省级当局监督下，一项具有这一效果的协议被宣布。但相当数量的``红暴''活动分子拒绝接受协议条款，并前往北京申诉。不出所料，他们受到冷淡接待，并被命令返回杭州，在``学习班''中接受事实上的拘禁。

Despite being signatories to the agreement which officially disbanded
their organizations, both United Headquarters and Red Storm retained
their separate identities and an organizational structure which
henceforth operated underground. Thus, the mass factionalism of the
previous years was supplemented and exacerbated by the formation of
secret informal groups which linked the former leaders of United
Headquarters with sympathetic military and civilian cadres who dominated
the power structure in Zhejiang. The 5th CCP Zhejiang Provincial
Congress, which met in January 1971, formalized the military's hegemony
over the Zhejiang administration. The running of the province was now
largely in the hands of cadres who were both novices to civilian
administration and who were outsiders and hence unfamiliar with local
conditions.

尽管双方都签署了正式解散其组织的协议，``联合总部''和``红暴''仍保留各自身份以及此后转入地下运作的组织结构。因此，前几年群众派性又因秘密非正式集团的形成而得到补充和加剧；这些集团把``联合总部''前领导人与支配浙江权力结构、并同情他们的军队和文职干部联系起来。1971
年 1
月召开的中共浙江省第五次代表大会，使军队对浙江行政的霸权正式化。该省的运行此时主要掌握在一批干部手中；他们既是文职行政的新手，又是外来者，因而不熟悉地方条件。

Details of political developments between the years 1969 and 1972 are
scarce. What is clear is that the provincial leadership owed its
position primarily to the central military command led by Lin Biao. On
this basis it carried out a series of campaigns to militarize and
radicalize the civilian economy and administration and to display
loyalty to its patron. Regular contacts between Lin Biao's son, a rising
star in the airforce, and the airforce command in Zhejiang strengthened
these ties. Ousted civilian cadres returned to their posts but the
process was confined mainly to the county level and below. Mass
representatives sitting on revolutionary committees and reconstituted
party committees struggled to retain any influence and adopted a
defensive posture, hoping thereby to avoid persecution in the series of
vindictive campaigns which the military leadership directed against
Cultural Revolution activists.

1969 至 1972
年间政治发展的细节很少。清楚的是，省级领导的地位主要归功于林彪领导的中央军事指挥系统。在此基础上，它开展了一系列运动，把民用经济和行政军事化、激进化，并展示对其庇护者的忠诚。林彪之子这位空军新星与浙江空军指挥系统之间的定期接触，加强了这些联系。被赶下台的文职干部回到岗位，但这一过程主要限于县级及以下。坐在革命委员会和重建党委中的群众代表努力保留任何影响力，并采取防御姿态，希望借此避免在军事领导针对文化大革命活动分子发动的一系列报复性运动中受到迫害。

Following the demise of Lin Biao and the downfall of his provincial
supporters, the central authorities in Beijing were forced once again to
take temporary control of Zhejiang through their emissary Xu Shiyou
(许世友), Commander of the Nanjing Military Region. Xu conducted the
purge of the provincial leadership and recommended that two experienced
cadres, Tan Qilong (谭启龙) and Tie Ying (铁瑛), head the administration
in Zhejiang. Tan Qilong had accumulated previous experience in Zhejiang,
having served there in the early post-Liberation years. In 1972 and 1973
the new leadership rehabilitated leading pre-Cultural Revolution
officials and placed them in key positions in the bureaucracy. They
released members of Red Storm languishing in jail. The provincial
leaders also initiated an investigation into the deeds of the leaders of
United Headquarters during the stormy days of the late 1960s, as well as
their relationship with the disgraced military cadres. Other military
officials were transferred to Zhejiang from the Nanjing Military Region
to root out Lin Biao's political influence and factional network.

林彪死亡及其省内支持者倒台后，北京中央当局被迫再次通过其特使、南京军区司令员许世友暂时控制浙江。许世友清洗了省级领导层，并建议由两名经验丰富的干部谭启龙和铁瑛主持浙江行政。谭启龙此前曾在浙江积累经验，解放初期曾在那里任职。1972
和 1973
年，新领导层为文化大革命前的主要官员平反，并将他们安置在官僚体系关键职位上。他们释放了在狱中受苦的``红暴''成员。省领导还开始调查``联合总部''领导人在
20 世纪 60
年代末风暴年代中的行为，以及他们同失势军事干部的关系。其他军事官员从南京军区调往浙江，以根除林彪的政治影响和派性网络。

In August 1973 the CCP convened its 10th National Congress in an
atmosphere heavy with the tension that heralded the resurgence of
leftist mobilizational politics. Chairman Mao's ``continuous
revolution'' was gearing up for another campaign. On the eve of the
Congress, pre-Cultural Revolution mass organizations held congresses
across the country to signal the resumption of their activities. In
Zhejiang, leaders of United Headquarters and Red Storm gained solid
representation on the Zhejiang Provincial Trade Union Council (ZPTUC)
and CYL executives. On the former body they were balanced by the
election of veteran trade union leaders and model workers of the 1950s
and 1960s. Shortly after the closure of the CCP Congress in Beijing, new
urban militia organizations were established under the formal leadership
of municipal party committees and modelled after the first such group
set up in Shanghai. Hangzhou established such a body in February 1974.

1973 年 8
月，中国共产党在一种充满紧张的氛围中召开第十次全国代表大会；这种紧张预示着左翼动员政治的复兴。毛主席的``继续革命''正准备发动另一场运动。大会前夕，文化大革命前的群众组织在全国各地召开代表大会，以标志其活动恢复。在浙江，``联合总部''和``红暴''领导人在浙江省总工会（ZPTUC）和共青团执委会中获得了稳固代表性。在前一机构中，他们受到
20 世纪 50 和 60
年代资深工会领导人及劳动模范当选的制衡。北京党代会闭幕后不久，新的城市民兵组织在市委正式领导下建立，其模式仿照上海最早成立的此类组织。杭州于
1974 年 2 月建立了这样的机构。

Three former prominent leaders of United Headquarters had gained
official posts of some substance at this time. They were Zhang Yongsheng
(张永生), a former student of the nationally-famous Zhejiang Academy of
Fine Arts, Weng Senhe (翁森鹤), a worker at the large Hangzhou Silk
Dyeing and Printing Complex, and He Xianchun (贺贤春), a former worker
at the Hangzhou Heavy Machinery Factory. All three had achieved a great
deal of notoriety in the Cultural Revolution. When the anti-Lin Biao
anti-Confucius campaign commenced in early 1974 Zhang, Weng and He were
ready to swing into action. Encouraged by Wang Hongwen in Beijing and
tolerated or condoned in their activities by influential cadres from
both provincial and municipal party committees, the radical leaders
found themselves back in the type of environment in which they thrived.
In a series of mass meetings, rallies and demonstrations they were
successful in recreating something of the atmosphere of disorder which
had characterized the early days of the Cultural Revolution. Their solid
body of supporters in the factories of Hangzhou was mobilized into
contingents of the urban militia. He Xianchun and an associate took
charge of the Hangzhou Municipal Militia Command Headquarters
(杭州市民兵指挥部 -- HMMCH). The substantial number of industrial
workers drafted into the militia in Hangzhou acted as a highly effective
and mobile task force for use against factional opponents in the city's
factories. In addition, the militia substituted for the PLA and public
security forces in guard duty and supervision of social order. The
intimidation and harassment of selected opponents, both during and after
work hours, had disastrous effects on harmony in the work-place and on
industrial production.

此时，``联合总部''三名前著名领导人已获得若干实质性官方职位。他们是张永生，全国知名的浙江美术学院前学生；翁森鹤，杭州丝绸印染联合厂的一名工人；以及贺贤春，杭州重型机械厂前工人。三人都在文化大革命中声名狼藉。1974
年初批林批孔运动开始时，张、翁、贺已准备行动。在北京王洪文的鼓励下，并在省、市两级党委中有影响力干部的容忍或纵容下，这些激进领导人发现自己又回到了他们如鱼得水的环境中。通过一系列群众大会、集会和示威，他们成功再造了某种文化大革命初期特有的混乱气氛。他们在杭州工厂中稳固的支持者群体被动员为城市民兵分队。贺贤春和一名同伙掌管了杭州市民兵指挥部（HMMCH）。杭州大量被编入民兵的产业工人，成为一支高度有效且机动的特遣力量，用于对付城内工厂中的派性对手。此外，民兵还替代人民解放军和公安力量执行警卫和社会秩序监督任务。无论在工作时间内外，对特定对手的恐吓和骚扰，都对工作场所和谐及工业生产造成了灾难性影响。

Zhang, Weng and He, together with their allies in the bureaucracy, also
gained control of the leading group which had been established to direct
the anti-Lin Biao anti-Confucius campaign. For most of 1974 it seemed to
replace the provincial party committee as the ultimate source of
authority in Zhejiang. At a four-month conference of the provincial
party committee, revolutionary committee and the party committee of the
military district, the rebels coordinated a campaign of vilification and
humiliation of the provincial leaders. In the face of these tactics, the
administration seemed powerless and looked on helplessly as industrial
production stalled, social order deteriorated and government became
paralyzed. In July 1974 the Central Committee acted to bring the
anti-Lin Biao/Confucius campaign to an end before the economy suffered
even greater damage. The Zhejiang authorities took considerable time to
rein in the destructive activities of the rebels, but by September they
had achieved initial success. Yet the situation remained highly volatile
and, with the militia continuing to create disturbances, at the end of
the year Beijing called both its local representatives and the factional
leaders of Zhejiang to a conference in the capital.

张、翁、贺及其在官僚体系中的盟友，还控制了为领导批林批孔运动而设立的领导小组。1974
年大部分时间里，它似乎取代了省委，成为浙江最终权威来源。在省党委、革命委员会和军区党委为期四个月的会议上，造反派协调了一场诋毁和羞辱省领导的运动。面对这些策略，行政当局似乎无能为力，只能眼看着工业生产停滞、社会秩序恶化、政府陷入瘫痪。1974
年 7
月，中央委员会采取行动结束批林批孔运动，以免经济遭受更大损害。浙江当局花费相当长时间才约束住造反派的破坏活动，但到
9
月已经取得初步成功。然而，局势仍高度不稳定；随着民兵继续制造骚乱，年底北京把其地方代表和浙江派性领导人都召到首都开会。

Over the following six months, the political struggle focused on
Beijing's attempt to reassert central control over Zhejiang and to prop
up the provincial authorities against those who sought to undermine and
overthrow it. It appears that early in 1975 the decision was made to
disband the urban militia command and so remove a major source of
instability from the Zhejiang political scene. A procession of central
leaders visited Hangzhou to investigate the situation for themselves.
Deng Xiaoping was possibly one such visitor and, in February 1975, Mao
Zedong stopped there in transit on his way back to Beijing. The peasant
leader and Politburo member, Chen Yonggui, arrived in April to address a
provincial meeting concerning agriculture. In his speech Chen uttered
some stern words about the dangers of continued factionalism. However,
the problem had clearly not been settled when, at the end of June 1975,
the central leadership despatched two senior members of the Politburo,
Vice-chairman Wang Hongwen and Vice-premier Ji Dengkui, to Hangzhou.
They brought with them a high-powered delegation from the CC
Organization Department and two central ministries.

在随后六个月里，政治斗争集中于北京试图重新确立对浙江的中央控制，并支撑省级当局抵御那些企图削弱和推翻它的人。看来，到
1975
年初，解散城市民兵指挥机构的决定已经作出，以便从浙江政治场景中移除一个主要不稳定来源。一批中央领导人相继到杭州亲自调查情况。邓小平可能就是其中一位；1975
年 2 月，毛泽东在返回北京途中也在杭州停留。农民领袖、政治局委员陈永贵于
4
月抵达，在一次有关农业的省级会议上讲话。陈在讲话中对持续派性的危险说了一些严厉的话。然而，到
1975 年 6
月底，问题显然仍未解决；中央领导派遣两名政治局高级成员，即副主席王洪文和副总理纪登奎前往杭州。他们带来一个由中央组织部和两个中央部委组成的高级代表团。

After a three-week investigation the two central leaders decided upon
recommendations which were submitted to the Politburo for approval. This
body made three major decisions which were then relayed back to Zhejiang
for implementation. First, a major reshuffle of the provincial and
municipal party leaderships, as well as that of the military district,
was announced. Cadres who had been compromised by their involvement in
factional disturbances were demoted or transferred and replacements
brought in from other provinces. The rebel leaders and their supporters
were banished, sent to the countryside, or enrolled in study classes to
review their mistakes and receive ideological and political
re-education.

经过三周调查，两位中央领导决定提出建议，并提交政治局批准。政治局作出三项重大决定，随后转回浙江执行。第一，宣布对省、市两级党领导以及军区领导进行重大改组。因卷入派性骚乱而受到牵连的干部被降职或调离，替代者从外省调入。造反派领导人及其支持者被驱逐、送往农村，或被编入学习班，以检查错误并接受思想政治再教育。

Second, the central leadership decided to send at least 10,000 soldiers
from the three services into the factories of Hangzhou to pacify workers
and supervise the resumption of a normal working environment. Outside
troops, uninvolved in the lengthy and bitter local squabbles, were moved
into the city. Third, the investigation of industrial problems in
Hangzhou resulted in the propagation of an intensive emulation campaign,
known as the eight factories experience. The campaign oversaw the
restoration of order in the factories and the criticism and discipline
of unruly workers, and issued appeals for the completion of output
targets in preparation for the launching of the 6th 5-Year Plan in 1976.

第二，中央领导决定派遣三军至少一万名士兵进入杭州工厂，以安抚工人并监督正常工作环境的恢复。未卷入长期而激烈地方争吵的外来部队被调入城内。第三，对杭州工业问题的调查，导致一场被称为``八厂经验''的密集学比运动得到推广。这场运动监督工厂秩序恢复，批评并处分不守纪律的工人，同时发出号召，要求完成产量指标，为
1976 年启动第六个五年计划作准备。

These extraordinary measures -- unprecedented even in the tumultuous
years of the Cultural Revolution -- provide testimony to the havoc that
rebellion and factionalism had caused in Zhejiang over the preceding ten
years. However, the incident and the circumstances surrounding it to
this day remain shrouded in a fog of misunderstanding and
incomprehension. In the forty-year history of the People's Republic the
communist regime has been loth to involve troops in civil disturbances.
Before the massive mobilization of soldiers against the people of
Beijing in May and June of 1989, the previous occasion on which the army
had been called out to intervene in civilian unrest involving Han
Chinese was in the 1975 military occupation of the factories of
Hangzhou.

这些非同寻常的措施，即便在文化大革命动荡岁月中也前所未有，证明了造反和派性在此前十年给浙江造成的破坏。然而，这一事件及其周围情境，至今仍笼罩在误解和不可理解的迷雾中。在人民共和国四十年历史中，共产党政权一直不愿让军队卷入民事骚乱。在
1989 年 5 月和 6
月大规模动员士兵对付北京人民之前，上一次军队被调出介入涉及汉人的平民骚乱，是
1975 年对杭州工厂的军事占领。

In spite of the compromises which had been made by the leadership in
Beijing to rescue the authority of the CCP ZPC, the decisions of July
1975 contained the seeds of their own destruction. Military occupation
of factories was only a short-term solution to the unrest, and the new
leadership installed in Zhejiang was at best an interim group, comprised
as it was of representatives of different factional groupings. When the
Gang of Four launched its counter-attack on Deng Xiaoping at the end of
1975, his supporters in Zhejiang came under renewed pressure. The
correctness of the decisions made under Deng's aegis were therefore
questioned and openly challenged by the time of Mao's death in September
1976. The Cultural Revolution, with its special brand of rebellion and
factionalism, ended with the arrest of the Gang of Four in the following
month. Its followers in Zhejiang were dismissed and the rebel leaders
arrested and jailed.

尽管北京领导为挽救中共浙江省委权威作出了妥协，1975 年 7
月的决定却包含着自我毁灭的种子。对工厂实行军事占领只是针对骚乱的短期解决办法，而安置在浙江的新领导层充其量只是一个临时班子，因为它由不同派性集团的代表组成。1975
年底``四人帮''发动对邓小平的反击时，他在浙江的支持者再次承压。因此，到
1976 年 9
月毛去世时，在邓支持下作出的决定是否正确，受到质疑并被公开挑战。文化大革命及其特殊形式的造反和派性，随着次月``四人帮''被捕而告终。其在浙江的追随者被免职，造反派领袖被逮捕入狱。

\section{Notes}\label{notes}

\section{注释}\label{ux6ce8ux91ca}

\bookmarksetup{startatroot}

\chapter{Chapter One - THE OUTBREAK OF THE CULTURAL REVOLUTION, MAY 1966
- FEBRUARY 1967 / 第一章 -
文化大革命的爆发，1966年5月-1967年2月}\label{chapter-one---the-outbreak-of-the-cultural-revolution-may-1966---february-1967-ux7b2cux4e00ux7ae0---ux6587ux5316ux5927ux9769ux547dux7684ux7206ux53d11966ux5e745ux6708-1967ux5e742ux6708}

When the Cultural Revolution commenced in 1966, the provincial
leadership in Zhejiang, like its counterparts elsewhere in China, faced
an unprecedented set of circumstances. Polemics on the literary front,
initiated by Yao Wenyuan's broadside against Wu Han's play, Hai Rui
Dismissed from Office, spread to the political sphere, engulfing
Beijing's municipal leadership in the process. Student unrest on
campuses in the national capital expanded in scope and militancy. Calls
were made for the dismissal of school administrators for unspecified but
sweeping crimes including the attempt to suppress the protest. When Mao
Zedong added his voice to that of the rebels by demanding the overthrow
of the ``bourgeois headquarters'' within the CCP, it was clear that the
Cultural Revolution was to go beyond the bounds of previous mass
movements and political campaigns, both in its scope and method of
conduct.

1966年文化大革命开始时，浙江省领导层同中国其他地方的同级领导一样，面对的是一系列前所未有的局面。由姚文元猛烈批判吴晗剧作《海瑞罢官》而发端的文艺战线论战，扩展到政治领域，并在此过程中卷入了北京市领导层。首都各校园的学生动荡在规模和激烈程度上不断扩大。人们要求撤换学校管理者，理由是他们犯下了包括试图压制抗议在内的笼统而严重的罪行。当毛泽东以要求打倒中共内部``资产阶级司令部''的方式加入造反者的声浪时，显然，文化大革命无论在范围还是进行方式上，都将超出以往群众运动和政治运动的界限。

With Mao's political victory at the CCP CC 11th Plenum in August 1966, a
new social force known as the Red Guards burst upon the political scene
to spearhead the campaign. Initially, these young zealots targeted
academics in schools and colleges as well as seeking out and destroying
examples (whether architectural, written or symbolic) of the feudalistic
nature of China's culture. Next they turned their attention to party
officials who had cynically manipulated the often genuine resentment
that the young felt toward privilege and autocratic leadership and who
had diverted it from the real power-holders -- that is, the cadres
themselves.

随着毛在1966年8月中共中央十一中全会上取得政治胜利，一股被称为红卫兵的新社会力量突然登上政治舞台，成为这场运动的先锋。最初，这些年轻的狂热者把矛头指向学校和大学里的学者，同时寻找并摧毁中国文化中带有封建性质的事物，无论它们是建筑、文字还是象征物。随后，他们把注意力转向党内官员；这些官员曾老于世故地操纵年轻人对特权和专断领导的真实不满，并把这种不满从真正掌权者身上转移开去，也就是从干部自身身上转移开去。

\section{The Zhejiang Political
Scene}\label{the-zhejiang-political-scene}

\section{浙江政治图景}\label{ux6d59ux6c5fux653fux6cbbux56feux666f}

Since 1949 Zhejiang had been ruled by cadres drawn largely from the 3rd
Field Army, which had marched into the province at the beginning of that
year.\footnote{See William Whitson, The Chinese High Command: A History
  of Communist Military Politics, 1927-71 (New York: Praeger Publishers,
  1973), pp.~245-57.

  参见 William Whitson, The Chinese High Command: A History of Communist
  Military Politics, 1927-71 (New York: Praeger Publishers,
  1973)，第245-257页。} Many of the civilian cadres who arrived in the
wake of the PLA were natives of Shandong or northern Jiangsu provinces.
Local communist activists who had participated in guerrilla operations
in southern Zhejiang (浙南) in the early 1930s; in eastern Zhejiang
(浙东) first in the 1920s and later in the anti-Japanese war; or who had
worked in underground party organizations in Shanghai, were, by the end
of the 1950s, relegated to a secondary role in the new power structure.
Senior cadres of the 3rd Field Army such as Tan Zhenlin (谭震林) and Tan
Qilong served in the province during the early post-Liberation years of
reconstruction. After Tan Qilong's transfer in 1954, his deputy, Jiang
Hua,\footnote{For a biographical sketch of Jiang Hua see ``Jiang Hua -
  President of the Supreme People's Court'', I\&S, 16: 7 (1980),
  pp.~85-88; CNS, 288 (18 September 1969), pp.~B12-17. The whereabouts
  of Jiang's birthplace seems to have caused a great deal of confusion.
  The above article claims that he was born in Hubei province. Klein and
  Clark state that Jiang's birthplace was in Shandong province. See D.W.
  Klein and A.B. Clark (eds). Biographic Dictionary of Chinese
  Communism, 1921-1965 (Cambridge, Mass.: Harvard University Press,
  1971),Vol. I, p.~173. Jiang was in fact born in Jianghua county, Hunan
  province, from which place he took his name, a fact well known to the
  citizens of Zhejiang. For written proof, see the interview with Jiang
  Hua in which Jiang reminisced about Mao Zedong's frequent visits to
  Zhejiang and the interviewer noted Jiang's Hunan accent. ZJRB,
  December 26,1983, p.~3. Another fact about Jiang not previously
  published is that ethnically he is a Yao national. Jianghua county is
  now a Yao autonomous county in the south of Hunan.

  关于江华的传略，见``Jiang Hua - President of the Supreme People's
  Court'', I\&S, 16:7 (1980), 第85-88页；CNS, 288 (1969年9月18日),
  第B12-17页。江出生地所在似乎造成了很多混乱。上述文章称他出生于湖北省。Klein和Clark则称江的出生地在山东省。见
  D.W. Klein and A.B. Clark (eds). Biographic Dictionary of Chinese
  Communism, 1921-1965 (Cambridge, Mass.: Harvard University Press,
  1971),
  第I卷，第173页。事实上，江出生于湖南省江华县，并由此地取名；这一事实为浙江民众所熟知。书面证据见对江华的采访，采访中江回忆毛泽东频繁访问浙江，采访者还注意到江的湖南口音。ZJRB，1983年12月26日，第3页。另一个此前未发表的事实是，江在民族身份上属瑶族。江华县如今是湖南南部的瑶族自治县。}
became the paramount leader of Zhejiang.

自1949年以来，浙江一直由主要出身于第三野战军的干部治理；该野战军于当年年初进入浙江。\footnote{See
  William Whitson, The Chinese High Command: A History of Communist
  Military Politics, 1927-71 (New York: Praeger Publishers, 1973),
  pp.~245-57.

  参见 William Whitson, The Chinese High Command: A History of Communist
  Military Politics, 1927-71 (New York: Praeger Publishers,
  1973)，第245-257页。}
许多随解放军之后抵达的文职干部来自山东或苏北。在20世纪30年代初参与浙南游击活动、20世纪20年代及后来的抗日战争时期参与浙东活动，或曾在上海地下党组织中工作的本地共产主义活动分子，到20世纪50年代末已在新的权力结构中退居次要地位。第三野战军的高级干部，如谭震林、谭启龙，曾在解放初期的重建年代于浙江任职。1954年谭启龙调离后，他的副手江华\footnote{For
  a biographical sketch of Jiang Hua see ``Jiang Hua - President of the
  Supreme People's Court'', I\&S, 16: 7 (1980), pp.~85-88; CNS, 288 (18
  September 1969), pp.~B12-17. The whereabouts of Jiang's birthplace
  seems to have caused a great deal of confusion. The above article
  claims that he was born in Hubei province. Klein and Clark state that
  Jiang's birthplace was in Shandong province. See D.W. Klein and A.B.
  Clark (eds). Biographic Dictionary of Chinese Communism, 1921-1965
  (Cambridge, Mass.: Harvard University Press, 1971),Vol. I, p.~173.
  Jiang was in fact born in Jianghua county, Hunan province, from which
  place he took his name, a fact well known to the citizens of Zhejiang.
  For written proof, see the interview with Jiang Hua in which Jiang
  reminisced about Mao Zedong's frequent visits to Zhejiang and the
  interviewer noted Jiang's Hunan accent. ZJRB, December 26,1983, p.~3.
  Another fact about Jiang not previously published is that ethnically
  he is a Yao national. Jianghua county is now a Yao autonomous county
  in the south of Hunan.

  关于江华的传略，见``Jiang Hua - President of the Supreme People's
  Court'', I\&S, 16:7 (1980), 第85-88页；CNS, 288 (1969年9月18日),
  第B12-17页。江出生地所在似乎造成了很多混乱。上述文章称他出生于湖北省。Klein和Clark则称江的出生地在山东省。见
  D.W. Klein and A.B. Clark (eds). Biographic Dictionary of Chinese
  Communism, 1921-1965 (Cambridge, Mass.: Harvard University Press,
  1971),
  第I卷，第173页。事实上，江出生于湖南省江华县，并由此地取名；这一事实为浙江民众所熟知。书面证据见对江华的采访，采访中江回忆毛泽东频繁访问浙江，采访者还注意到江的湖南口音。ZJRB，1983年12月26日，第3页。另一个此前未发表的事实是，江在民族身份上属瑶族。江华县如今是湖南南部的瑶族自治县。}
成为浙江最高领导人。

Through his lengthy tenure in Zhejiang, Jiang Hua had, by the mid-1960s,
built up a powerful network of loyal subordinates within the provincial
party leadership. His closest and most trusted colleagues were placed in
key bureaucratic posts in the political/legal network, the propaganda
department and central office registry. Jiang's wife, Wu Zhonglian
(吴仲廉)\footnote{For a brief account of Wu's life and political
  activities, see the articles published in ZJRB, November 6,1978,
  pp.~1,3; and Kang Keqing, Zeng Zhi and Peng Ru, ``忠心耿耿,
  有胆有识--深切怀念吴仲廉同志'' (Loyal and devoted, courageous and
  knowledgeable - deeply cherish the memory of Comrade Wu Zhonglian),
  ZJRB, January 28,1987, p.~2.

  关于吴的生平和政治活动的简要叙述，见
  ZJRB，1978年11月6日，第1、3页；以及康克清、曾志、彭儒，《忠心耿耿,
  有胆有识--深切怀念吴仲廉同志》（Loyal and devoted, courageous and
  knowledgeable - deeply cherish the memory of Comrade Wu
  Zhonglian），ZJRB，1987年1月28日，第2页。} had held the important
positions of head of the ZPC Political and Legal Group and President of
the provincial Supreme Court since the mid-1950s. Wang Fang
(王芳),\footnote{For a sketch of his career, see ``Wang Fang - Newly
  Appointed Secretary of the CCP Zhejiang Provincial Committee'', I\&S
  19: 8 (1983), pp.~67-71.

  关于其仕途概述，见``Wang Fang - Newly Appointed Secretary of the CCP
  Zhejiang Provincial Committee'', I\&S 19:8 (1983), 第67-71页。} who
was Deputy-governor and entrusted with overall responsibility for public
security work in Zhejiang, was assisted by the Minister for Public
Security, Lü Jianguang (吕剑光).\footnote{For a recent account of the
  continued strength of this network see Liang Ruinian,
  ``彭真帮'公检法'势力抬头'' (Peng Zhen's clique in the public security
  network is gaining ground), 潮流 (The Tide), 3 (May 15, 1987), p.~12.

  关于这一网络持续强大的近年叙述，见梁瑞年，《彭真帮'公检法'势力抬头》（Peng
  Zhen's clique in the public security network is gaining
  ground），《潮流》(The Tide)，第3期（1987年5月15日），第12页。}

江华在浙江任职时间很长，到20世纪60年代中期，他已在省委领导层内部建立起一个强大的忠诚下属网络。他最亲近、最受信任的同事被安置在政法系统、宣传部和省委办公厅机要文书系统等关键官僚岗位。江的妻子吴仲廉\footnote{For
  a brief account of Wu's life and political activities, see the
  articles published in ZJRB, November 6,1978, pp.~1,3; and Kang Keqing,
  Zeng Zhi and Peng Ru, ``忠心耿耿, 有胆有识--深切怀念吴仲廉同志''
  (Loyal and devoted, courageous and knowledgeable - deeply cherish the
  memory of Comrade Wu Zhonglian), ZJRB, January 28,1987, p.~2.

  关于吴的生平和政治活动的简要叙述，见
  ZJRB，1978年11月6日，第1、3页；以及康克清、曾志、彭儒，《忠心耿耿,
  有胆有识--深切怀念吴仲廉同志》（Loyal and devoted, courageous and
  knowledgeable - deeply cherish the memory of Comrade Wu
  Zhonglian），ZJRB，1987年1月28日，第2页。}
自20世纪50年代中期起担任浙江省委政法小组组长和省高级人民法院院长等重要职务。副省长王芳\footnote{For
  a sketch of his career, see ``Wang Fang - Newly Appointed Secretary of
  the CCP Zhejiang Provincial Committee'', I\&S 19: 8 (1983), pp.~67-71.

  关于其仕途概述，见``Wang Fang - Newly Appointed Secretary of the CCP
  Zhejiang Provincial Committee'', I\&S 19:8 (1983), 第67-71页。}
负责浙江公安工作的全面事务，协助他的是公安厅长吕剑光。\footnote{For a
  recent account of the continued strength of this network see Liang
  Ruinian, ``彭真帮'公检法'势力抬头'' (Peng Zhen's clique in the public
  security network is gaining ground), 潮流 (The Tide), 3 (May 15,
  1987), p.~12.

  关于这一网络持续强大的近年叙述，见梁瑞年，《彭真帮'公检法'势力抬头》（Peng
  Zhen's clique in the public security network is gaining
  ground），《潮流》(The Tide)，第3期（1987年5月15日），第12页。}

Another close associate of Jiang's was Xue Ju (薛驹), who held the post
of Director of the ZPC General Office from where he could oversee the
flow of documents and correspondence. Xue had commenced work in Zhejiang
in 1947.\footnote{See Xue's preface to 浙江省经济研究中心 (Zhejiang
  provincial economic research center), 浙江省情概要 (An Outline of the
  Affairs of Zhejiang) (Hangzhou: Zhejiang Renmin Chubanshe, 1986),
  p.~1. For details of Xue's political career, see ``Xue Ju - Newly
  Elected Governor of Zhejiang'', I\&S, 19:10 (October 1983), pp.~82-7.

  见薛为浙江省经济研究中心《浙江省情概要》(An Outline of the Affairs of
  Zhejiang) (杭州：浙江人民出版社，1986)
  所作序言，第1页。关于薛的政治生涯详情，见``Xue Ju - Newly Elected
  Governor of Zhejiang'', I\&S, 19:10 (1983年10月), 第82-87页。} Chen
Bing (陈冰), known as the master of the pen (笔杆子), headed the
Propaganda Department. Chen Weida (陈伟达), who had been a member of the
provincial party secretariat since 1959, was responsible for the
implementation of agricultural policy. Zhejiang's regular grain
surpluses provided Jiang with bargaining power vis-à-vis the central
authorities to obtain funds for the development of the
province.\footnote{For the importance of grain surpluses to the
  maintenance of personal stability and to the political leverage that
  it bestowed on provincial leaders prior to the Cultural Revolution,
  see Frederick C. Teiwes, ``Provincial Politics in China: Themes and
  Variations'', in J. Lindbeck (ed), China: Management of a
  Revolutionary Society (Seattle: University of Washington, 1971),
  pp.174-6.

  关于粮食盈余对维持个人稳定及其在文革前赋予省级领导人的政治杠杆的重要性，见
  Frederick C. Teiwes, ``Provincial Politics in China: Themes and
  Variations'', 收于 J. Lindbeck 编, China: Management of a
  Revolutionary Society (Seattle: University of Washington, 1971),
  第174-176页。}

江的另一位亲信是薛驹，他担任省委办公厅主任，能够掌控文件与函件的流转。薛自1947年开始在浙江工作。\footnote{See
  Xue's preface to 浙江省经济研究中心 (Zhejiang provincial economic
  research center), 浙江省情概要 (An Outline of the Affairs of Zhejiang)
  (Hangzhou: Zhejiang Renmin Chubanshe, 1986), p.~1. For details of
  Xue's political career, see ``Xue Ju - Newly Elected Governor of
  Zhejiang'', I\&S, 19:10 (October 1983), pp.~82-7.

  见薛为浙江省经济研究中心《浙江省情概要》(An Outline of the Affairs of
  Zhejiang) (杭州：浙江人民出版社，1986)
  所作序言，第1页。关于薛的政治生涯详情，见``Xue Ju - Newly Elected
  Governor of Zhejiang'', I\&S, 19:10 (1983年10月), 第82-87页。}
被称为``笔杆子''的陈冰领导宣传部。陈伟达自1959年起任省委书记处成员，负责农业政策的执行。浙江经常性的粮食盈余，使江在同中央打交道、争取本省发展资金时拥有谈判资本。\footnote{For
  the importance of grain surpluses to the maintenance of personal
  stability and to the political leverage that it bestowed on provincial
  leaders prior to the Cultural Revolution, see Frederick C. Teiwes,
  ``Provincial Politics in China: Themes and Variations'', in J.
  Lindbeck (ed), China: Management of a Revolutionary Society (Seattle:
  University of Washington, 1971), pp.174-6.

  关于粮食盈余对维持个人稳定及其在文革前赋予省级领导人的政治杠杆的重要性，见
  Frederick C. Teiwes, ``Provincial Politics in China: Themes and
  Variations'', 收于 J. Lindbeck 编, China: Management of a
  Revolutionary Society (Seattle: University of Washington, 1971),
  第174-176页。}

Two of the remaining members of the ZPC secretariat, Li Fengping
(李丰平) and Wu Xian (吴宪), had served under Jiang since 1954, Li
having spent a three-year interregnum in the neighboring province of
Anhui between 1963 and 1966.\footnote{安徽省人民政府办公厅编 (Office of
  the Anhui provincial people's government), 安徽省情 (The Affairs of
  Anhui Province), (Hefei: Anhui Renmin Chubanshe, 1985), p.~182. At the
  2nd CCP Anhui Provincial Congress held in July 1963, Li was elected
  Secretary of the provincial secretariat and deputy to 1st Secretary Li
  Baohua.

  安徽省人民政府办公厅编，《安徽省情》(The Affairs of Anhui
  Province)（合肥：安徽人民出版社，1985），第182页。1963年7月召开的中共安徽省第二次代表大会上，李当选省委书记处书记，并任第一书记李葆华的副手。}
Lai Keke (赖可可) was the most recent addition to this elite, having
been promoted in 1964 from his post of Secretary-general of the CCP ZPC.
Lai's career had suffered a serious setback in 1953 as a result of his
involvement in conspiratorial plans drawn up by Gao Gang and Rao Shushi
to reshuffle the central leadership directly below Mao. After his
demotion in 1953-54, Lai's transfer to Zhejiang in 1961 marked a new
phase in his career.\footnote{Shen Weicai,
  ``从高饶反党联盟的得力干将到'四人帮'的忠实代理人'' (From a capable
  member of the Gao-Rao anti-party alliance to a loyal follower of the
  Gang of Four), ZJRB, August 16,1977, pp.~1,3.

  沈维才，《从高饶反党联盟的得力干将到'四人帮'的忠实代理人》（From a
  capable member of the Gao-Rao anti-party alliance to a loyal follower
  of the Gang of Four），ZJRB，1977年8月16日，第1、3页。} Significantly,
he was the only member of the ZPC secretariat to have suffered the
ignominy of demotion before the Cultural Revolution.

省委书记处其余成员中的两人，李丰平和吴宪，自1954年以来就在江的领导下工作；其中李在1963年至1966年间曾有三年调往邻省安徽。\footnote{安徽省人民政府办公厅编
  (Office of the Anhui provincial people's government), 安徽省情 (The
  Affairs of Anhui Province), (Hefei: Anhui Renmin Chubanshe, 1985),
  p.~182. At the 2nd CCP Anhui Provincial Congress held in July 1963, Li
  was elected Secretary of the provincial secretariat and deputy to 1st
  Secretary Li Baohua.

  安徽省人民政府办公厅编，《安徽省情》(The Affairs of Anhui
  Province)（合肥：安徽人民出版社，1985），第182页。1963年7月召开的中共安徽省第二次代表大会上，李当选省委书记处书记，并任第一书记李葆华的副手。}
赖可可是这个精英圈中最新加入者，1964年由中共浙江省委秘书长升任。赖的仕途曾在1953年遭受严重挫折，原因是他卷入了高岗、饶漱石策划的阴谋，意在改组毛以下的中央领导层。1953-54年被降职后，赖于1961年调往浙江，标志着其政治生涯的新阶段。\footnote{Shen
  Weicai, ``从高饶反党联盟的得力干将到'四人帮'的忠实代理人'' (From a
  capable member of the Gao-Rao anti-party alliance to a loyal follower
  of the Gang of Four), ZJRB, August 16,1977, pp.~1,3.

  沈维才，《从高饶反党联盟的得力干将到'四人帮'的忠实代理人》（From a
  capable member of the Gao-Rao anti-party alliance to a loyal follower
  of the Gang of Four），ZJRB，1977年8月16日，第1、3页。}
值得注意的是，他是浙江省委书记处成员中唯一在文革前曾遭降职羞辱的人。

Another more recent arrival in Zhejiang was Cao Xiangren (曹祥仁) who
had been appointed to the ZPC secretariat in 1959. He had previously
served as China's ambassador to Bulgaria (1950-54) and Deputy-minister
of the First Ministry of Machine-building (1955-59). His transfer to
Zhejiang perhaps signified a central decision to devote more attention
and expertise to the development of industry in the province. It is
alleged that Cao found some difficulty working with the rough, blunt,
domineering Jiang. For whatever reason, it is significant that Cao and
Lai alone broke ranks with Jiang's leadership team when the ZPC came
under direct attack from the rebels at the end of 1966. And in January
1967 Cao displayed personal malice toward Jiang by publicly demanding Wu
Zhonglian's posthumous expulsion from the CCP.\footnote{See n.~103
  below.

  见下文注释103。}

另一位较晚来到浙江的是曹祥仁，他于1959年被任命为省委书记处成员。此前，他曾任中国驻保加利亚大使（1950-54年）和第一机械工业部副部长（1955-59年）。他调任浙江，或许表明中央决定在浙江工业发展上投入更多关注和专业力量。据称，曹发现与粗犷、直率而专横的江共事颇为困难。不论原因如何，重要的是，当1966年底省委遭到造反派直接攻击时，只有曹和赖脱离了江的领导班子。1967年1月，曹公开要求在吴仲廉死后开除其党籍，显示出他对江的个人敌意。\footnote{See
  n.~103 below.

  见下文注释103。}

\textbf{TABLE ONE}

\textbf{表一}

\textbf{CCP ZPC STANDING COMMITTEE AT 6 JUNE 1966}

\textbf{1966年6月6日中共浙江省委常委会}

{\def\LTcaptype{none} % do not increment counter
\begin{longtable}[]{@{}
  >{\raggedright\arraybackslash}p{(\linewidth - 4\tabcolsep) * \real{0.3143}}
  >{\centering\arraybackslash}p{(\linewidth - 4\tabcolsep) * \real{0.2571}}
  >{\raggedright\arraybackslash}p{(\linewidth - 4\tabcolsep) * \real{0.4286}}@{}}
\toprule\noalign{}
\begin{minipage}[b]{\linewidth}\raggedright
\end{minipage} & \begin{minipage}[b]{\linewidth}\centering
\end{minipage} & \begin{minipage}[b]{\linewidth}\raggedright
Date of appointment
\end{minipage} \\
\midrule\noalign{}
\endhead
\bottomrule\noalign{}
\endlastfoot
First Secretary: & Jiang Hua & August 1950 \\
Secretaries: & Li Fengping & 1952 (1963-66 in Anhui) \\
& Cao Xiangren & June 1959 \\
& Chen Weida & April 1955 \\
& Wu Xian & April 1955 \\
& Lai Keke & April 1963 \\
Members: & Long Qian & 1966 \\
& Lu Zhixian & April 1963 \\
& Li Weixin & April 1963 \\
& Wang Qi & April 1963 \\
& Chen Bing & April 1965 \\
& Shen Ce & April 1965 \\
\end{longtable}
}

\begin{quote}
Sources: ZJRB, 28 June 1966, p.~1. Zhejiang shengqing gaiyao, pp.~24-6.
Kao Ch'ung-yen, Zhonggong renshi biandong, 1959-69, pp.~673-5.
\end{quote}

{\def\LTcaptype{none} % do not increment counter
\begin{longtable}[]{@{}
  >{\raggedright\arraybackslash}p{(\linewidth - 4\tabcolsep) * \real{0.3143}}
  >{\centering\arraybackslash}p{(\linewidth - 4\tabcolsep) * \real{0.2571}}
  >{\raggedright\arraybackslash}p{(\linewidth - 4\tabcolsep) * \real{0.4286}}@{}}
\toprule\noalign{}
\begin{minipage}[b]{\linewidth}\raggedright
\end{minipage} & \begin{minipage}[b]{\linewidth}\centering
\end{minipage} & \begin{minipage}[b]{\linewidth}\raggedright
任命日期
\end{minipage} \\
\midrule\noalign{}
\endhead
\bottomrule\noalign{}
\endlastfoot
第一书记： & 江华 & 1950年8月 \\
书记： & 李丰平 & 1952年（1963-66年在安徽） \\
& 曹祥仁 & 1959年6月 \\
& 陈伟达 & 1955年4月 \\
& 吴宪 & 1955年4月 \\
& 赖可可 & 1963年4月 \\
委员： & 龙潜 & 1966年 \\
& 卢之先 & 1963年4月 \\
& 李维新 & 1963年4月 \\
& 王起 & 1963年4月 \\
& 陈冰 & 1965年4月 \\
& 沈策 & 1965年4月 \\
\end{longtable}
}

\begin{quote}
资料来源： ZJRB，1966年6月28日，第1页。 Zhejiang shengqing
gaiyao，第24-26页。 高崇焱，Zhonggong renshi biandong,
1959-69，第673-675页。
\end{quote}

Thus, at the beginning of 1966, Jiang Hua led a team of experienced
politicians most of whom had worked together for the previous decade.
Jiang had several times displayed his commitment to upholding central
policy: in his purge of localists in 1957-58,\footnote{See the documents
  compiled in Current Background, 487 (January 10,1958), pp.~1-57.

  见 Current Background, 487（1958年1月10日）所辑文件，第1-57页。} his
enthusiastic espousal of both the Great Leap Forward program\footnote{See
  Jiang Hua, ``认清形势, 加快建设'' (Clearly understand the situation,
  speed up construction), ZJRB, November 6,1958, pp.~1,3; ``Maintain
  Uninterrupted Revolution, Maintain High Speed'', Hongqi, No.~21,
  (November 1959), translated in Joint Publications Research Service,
  1098-D (January 4, 1960), pp.~29-32.

  见江华，《认清形势, 加快建设》（Clearly understand the situation,
  speed up construction），ZJRB，1958年11月6日，第1、3页；``Maintain
  Uninterrupted Revolution, Maintain High Speed'', Hongqi,
  No.~21（1959年11月），译载 Joint Publications Research Service,
  1098-D（1960年1月4日），第29-32页。} and the policies of readjustment
which followed.\footnote{See the charges against Jiang Hua published in
  the Cultural Revolution. ZPS, 7 July 1968, SWB/FE/2823/B/12-14; ZPS, 7
  July 1968, SWB/FE/2824/B/9-10; ``打倒江华'' (Down with Jiang Hua),
  ZJRB, November 15, 1968, p.~1; ``江华'独立王国'的破产'' (The
  bankruptcy of Jiang Hua's ``independent kingdom'') ZJRB, January 13,
  1969, p.~1.

  见文化大革命期间发表的对江华的指控。ZPS，1968年7月7日，SWB/FE/2823/B/12-14；ZPS，1968年7月7日，SWB/FE/2824/B/9-10；《打倒江华》（Down
  with Jiang
  Hua），ZJRB，1968年11月15日，第1页；《江华'独立王国'的破产》（The
  bankruptcy of Jiang Hua's ``independent
  kingdom''），ZJRB，1969年1月13日，第1页。} Jiang was the perfect
example of the provincial politician of the time, the middleman par
excellence, who, in the words of the leading analyst of the phenomenon

因此，在1966年初，江华领导着一支经验丰富的政治家队伍，其中大多数人在此前十年里一直共事。江多次表现出维护中央政策的决心：1957-58年清洗地方主义者，\footnote{See
  the documents compiled in Current Background, 487 (January 10,1958),
  pp.~1-57.

  见 Current Background, 487（1958年1月10日）所辑文件，第1-57页。}
热情拥护大跃进纲领，\footnote{See Jiang Hua, ``认清形势, 加快建设''
  (Clearly understand the situation, speed up construction), ZJRB,
  November 6,1958, pp.~1,3; ``Maintain Uninterrupted Revolution,
  Maintain High Speed'', Hongqi, No.~21, (November 1959), translated in
  Joint Publications Research Service, 1098-D (January 4, 1960),
  pp.~29-32.

  见江华，《认清形势, 加快建设》（Clearly understand the situation,
  speed up construction），ZJRB，1958年11月6日，第1、3页；``Maintain
  Uninterrupted Revolution, Maintain High Speed'', Hongqi,
  No.~21（1959年11月），译载 Joint Publications Research Service,
  1098-D（1960年1月4日），第29-32页。}
以及随后实行的调整政策。\footnote{See the charges against Jiang Hua
  published in the Cultural Revolution. ZPS, 7 July 1968,
  SWB/FE/2823/B/12-14; ZPS, 7 July 1968, SWB/FE/2824/B/9-10;
  ``打倒江华'' (Down with Jiang Hua), ZJRB, November 15, 1968, p.~1;
  ``江华'独立王国'的破产'' (The bankruptcy of Jiang Hua's ``independent
  kingdom'') ZJRB, January 13, 1969, p.~1.

  见文化大革命期间发表的对江华的指控。ZPS，1968年7月7日，SWB/FE/2823/B/12-14；ZPS，1968年7月7日，SWB/FE/2824/B/9-10；《打倒江华》（Down
  with Jiang
  Hua），ZJRB，1968年11月15日，第1页；《江华'独立王国'的破产》（The
  bankruptcy of Jiang Hua's ``independent
  kingdom''），ZJRB，1969年1月13日，第1页。}
江是当时省级政治人物的典型，是最出色的中间人；用研究这一现象的主要分析者的话说，

\begin{quote}
is a man eminently sensitive to shifting winds from Beijing. Although
jealous of provincial interests, he is careful to tailor his
articulation of those interests to the prevailing policy line. He
emphasizes the primacy of national interest and shapes provincial policy
according to central directives. When the provincial politician errs, it
is more likely by going too far in implementing ``infallible'' central
policies than by opposing them.\footnote{Teiwes ``Provincial Politics'',
  p.~146. See also Parris Chang, Power and Policy in China, 2nd enlarged
  ed.~(University Park, Pennsylvania: Pennsylvania State University
  Press, 1978), p.~44, who basically agrees with Teiwes but adds that
  ``through their control of the flow of information, provincial
  officials often enjoyed a great deal of autonomy and were in a good
  position to evade central supervision'', (p.~45). In his study of
  Fujian province for the period 1949-66 Victor Falkenheim comes to the
  same conclusions as Teiwes. See Victor C. Falkenheim, ``Provincial
  Leadership in Fujian: 1949-66'', in R. Scalapino (ed), Elites in the
  People's Republic of China (Washington: University of Washington
  Press, 1972), pp.~199-244, esp.~p.~204. See also David S. G. Goodman,
  ``The Provincial First Party Secretary in the People's Republic of
  China, 1949-78: A Profile'', British Journal of Political Science,
  10:1 (1980), p.~43.

  Teiwes, ``Provincial Politics''，第146页。又见 Parris Chang, Power and
  Policy in China, 第2增订版 (University Park, Pennsylvania:
  Pennsylvania State University Press, 1978),
  第44页；他基本同意Teiwes，但补充说，``由于控制信息流动，省级官员常常享有很大自主权，并处于能够规避中央监督的有利位置''（第45页）。Victor
  Falkenheim关于1949-66年福建省的研究也得出与Teiwes相同的结论。见 Victor
  C. Falkenheim, ``Provincial Leadership in Fujian: 1949-66'', 收于 R.
  Scalapino 编, Elites in the People's Republic of China (Washington:
  University of Washington Press, 1972), 第199-244页，尤见第204页。另见
  David S. G. Goodman, ``The Provincial First Party Secretary in the
  People's Republic of China, 1949-78: A Profile'', British Journal of
  Political Science, 10:1 (1980), 第43页。}
\end{quote}

\begin{quote}
他对北京风向的变化极其敏感。尽管十分重视本省利益，他却小心地使自己对这些利益的表达符合当时的政策路线。他强调国家利益至上，并按照中央指示塑造省级政策。当省级政治人物犯错时，更可能是因为在执行``绝对正确''的中央政策时走得太远，而不是因为反对这些政策。\footnote{Teiwes
  ``Provincial Politics'', p.~146. See also Parris Chang, Power and
  Policy in China, 2nd enlarged ed.~(University Park, Pennsylvania:
  Pennsylvania State University Press, 1978), p.~44, who basically
  agrees with Teiwes but adds that ``through their control of the flow
  of information, provincial officials often enjoyed a great deal of
  autonomy and were in a good position to evade central supervision'',
  (p.~45). In his study of Fujian province for the period 1949-66 Victor
  Falkenheim comes to the same conclusions as Teiwes. See Victor C.
  Falkenheim, ``Provincial Leadership in Fujian: 1949-66'', in R.
  Scalapino (ed), Elites in the People's Republic of China (Washington:
  University of Washington Press, 1972), pp.~199-244, esp.~p.~204. See
  also David S. G. Goodman, ``The Provincial First Party Secretary in
  the People's Republic of China, 1949-78: A Profile'', British Journal
  of Political Science, 10:1 (1980), p.~43.

  Teiwes, ``Provincial Politics''，第146页。又见 Parris Chang, Power and
  Policy in China, 第2增订版 (University Park, Pennsylvania:
  Pennsylvania State University Press, 1978),
  第44页；他基本同意Teiwes，但补充说，``由于控制信息流动，省级官员常常享有很大自主权，并处于能够规避中央监督的有利位置''（第45页）。Victor
  Falkenheim关于1949-66年福建省的研究也得出与Teiwes相同的结论。见 Victor
  C. Falkenheim, ``Provincial Leadership in Fujian: 1949-66'', 收于 R.
  Scalapino 编, Elites in the People's Republic of China (Washington:
  University of Washington Press, 1972), 第199-244页，尤见第204页。另见
  David S. G. Goodman, ``The Provincial First Party Secretary in the
  People's Republic of China, 1949-78: A Profile'', British Journal of
  Political Science, 10:1 (1980), 第43页。}
\end{quote}

\section{The Beginnings of the Cultural
Revolution}\label{the-beginnings-of-the-cultural-revolution}

\section{文化大革命的开端}\label{ux6587ux5316ux5927ux9769ux547dux7684ux5f00ux7aef}

Two interpretations of the response of provincial officials to the
unfolding of events in 1966 have been put forward. One view, enunciated
by Parris Chang\footnote{Chang, ``Provincial Party Leaders' Strategies
  for Survival during the Cultural Revolution'', in Scalapino,
  pp.~501-39.

  Chang, ``Provincial Party Leaders' Strategies for Survival during the
  Cultural Revolution''，收于 Scalapino，第501-539页。} and applied to
Sichuan province by T. J. Mathews,\footnote{Mathews, ``The Cultural
  Revolution in Szechwan'' in The Cultural Revolution in The Provinces,
  pp.~97-109.

  Mathews, ``The Cultural Revolution in Szechwan''，收于 The Cultural
  Revolution in The Provinces，第97-109页。} has described the response
as, in Chang's phrase, ``strategies for survival''. Chang writes on the
assumption that provincial officials were intent on resisting the
Cultural Revolution. In the initial phase, according to Chang,
provincial leaders adopted such strategies as evasion and diversion,
deception and containment. After the August 1966 plenum these strategies
no longer remained viable and sheer survival, physical and political,
became of paramount importance. Summoning all the resources at their
disposal, provincial party officials sought support and protection from
their comrades in provincial military garrisons, lobbied patrons in
Beijing, bought off under-privileged groups among workers and peasants
by offering them bonuses and subsidies and even resorted to suppression
of their rebel challengers. Chang concludes, admiringly,

关于1966年事态展开时省级官员的反应，已有两种解释。一种观点由张柏瑞（Parris
Chang）提出，\footnote{Chang, ``Provincial Party Leaders' Strategies for
  Survival during the Cultural Revolution'', in Scalapino, pp.~501-39.

  Chang, ``Provincial Party Leaders' Strategies for Survival during the
  Cultural Revolution''，收于 Scalapino，第501-539页。} 并由T. J.
Mathews用于四川省研究，\footnote{Mathews, ``The Cultural Revolution in
  Szechwan'' in The Cultural Revolution in The Provinces, pp.~97-109.

  Mathews, ``The Cultural Revolution in Szechwan''，收于 The Cultural
  Revolution in The Provinces，第97-109页。}
将这种反应概括为张所谓的``生存策略''。张的论述基于这样一个假设：省级官员有意抵制文化大革命。按照张的说法，在最初阶段，省级领导人采取了逃避与转移、欺骗与遏制等策略。1966年8月全会之后，这些策略已不再可行，单纯的生存，无论是身体上的还是政治上的，成为最重要的问题。省级党政官员调动一切可用资源，向驻省军区的同志寻求支持和保护，游说北京的庇护者，通过发放奖金和补贴收买工农中的弱势群体，甚至诉诸镇压挑战他们的造反派。张带着赞赏意味总结道，

\begin{quote}
The fact that it was the use of military forces that finally brought
down the recalcitrant well-entrenched provincial Party leaders attests
to the remarkable effectiveness of their survival strategies.\footnote{Chang,
  ``Provincial Party Leaders' Strategies'', p.~537. ``The cohesion of
  the ruling group also contributed immeasurably to the success of this
  resistance. See Teiwes,''Provincial Politics'', p.~157.

  Chang, ``Provincial Party Leaders'
  Strategies''，第537页。``统治集团的凝聚力也对这种抵抗的成功作出了不可估量的贡献。''见
  Teiwes, ``Provincial Politics''，第157页。}
\end{quote}

\begin{quote}
最终正是动用军事力量才打倒了这些顽固而根基深厚的省级党委领导人，这一事实证明了他们生存策略的显著有效性。\footnote{Chang,
  ``Provincial Party Leaders' Strategies'', p.~537. ``The cohesion of
  the ruling group also contributed immeasurably to the success of this
  resistance. See Teiwes,''Provincial Politics'', p.~157.

  Chang, ``Provincial Party Leaders'
  Strategies''，第537页。``统治集团的凝聚力也对这种抵抗的成功作出了不可估量的贡献。''见
  Teiwes, ``Provincial Politics''，第157页。}
\end{quote}

Another, more subtle interpretation of the response of provincial
officials has been advanced by Andrew Walder.\footnote{Walder, Chang
  Ch'un-ch'iao and Shanghai's January Revolution, Michigan Papers in
  Chinese Studies, No.~32 (Ann Arbor: University of Michigan, 1978).

  Walder, Chang Ch'un-ch'iao and Shanghai's January Revolution, Michigan
  Papers in Chinese Studies, No.~32 (Ann Arbor: University of Michigan,
  1978)。} Based on a study of Shanghai for the period 1966-67, he
argues persuasively that local officials, at least in the early stages
of the Cultural Revolution, responded obediently to Mao's initiatives
even when they did not fully comprehend the raison d'être behind the
campaign. The provincial politicians did their best to move with the
constantly expanding parameters of Beijing's dictates.\footnote{Such
  seemed to be Zhao Ziyang's reaction in Guangdong province. See Zhao
  Wei, 赵紫阳传 (A Biography of Zhao Ziyang), (Beijing: Zhongguo Xinwen
  Chubanshe, 1989), pp.~129-30, 145.

  赵紫阳在广东省的反应似乎就是如此。见赵蔚，《赵紫阳传》(A Biography of
  Zhao Ziyang)（北京：中国新闻出版社，1989），第129-130、145页。}
Teiwes' description of provincial leaders, cited above, would strongly
suggest such a course of action. But, argues Walder, the momentum of the
Cultural Revolution left provincial politicians floundering in its wake
and eventually they lost all control over local events. In Walder's
view, various social groups needed no incitement from party officials to
put forward particularistic demands in an environment where social order
and industrial discipline had collapsed. Nor was it surprising that
undirected student groups should splinter and engage in increasingly
violent actions and internecine confrontation. In short, responsibility
for the social and political chaos in the localities lay not with
obstructionist acts committed by insubordinate officials, as suggested
by Chang, but rather with the increased scope allowed for
centrally-inspired and authorized attacks on local political authority.

安德鲁·沃尔德（Andrew Walder）提出了另一种更为细致的解释。\footnote{Walder,
  Chang Ch'un-ch'iao and Shanghai's January Revolution, Michigan Papers
  in Chinese Studies, No.~32 (Ann Arbor: University of Michigan, 1978).

  Walder, Chang Ch'un-ch'iao and Shanghai's January Revolution, Michigan
  Papers in Chinese Studies, No.~32 (Ann Arbor: University of Michigan,
  1978)。}
基于对1966-67年上海的研究，他有说服力地指出，地方官员至少在文化大革命初期是顺从毛的倡议的，即使他们并未完全理解这场运动背后的存在理由。省级政治人物尽力随着北京指令不断扩大的范围而行动。\footnote{Such
  seemed to be Zhao Ziyang's reaction in Guangdong province. See Zhao
  Wei, 赵紫阳传 (A Biography of Zhao Ziyang), (Beijing: Zhongguo Xinwen
  Chubanshe, 1989), pp.~129-30, 145.

  赵紫阳在广东省的反应似乎就是如此。见赵蔚，《赵紫阳传》(A Biography of
  Zhao Ziyang)（北京：中国新闻出版社，1989），第129-130、145页。}
上文引用的Teiwes对省级领导人的描述，也强烈暗示他们会采取这种行动路线。但沃尔德认为，文化大革命的势头使省级政治人物措手不及，最终完全失去了对地方事件的控制。在沃尔德看来，在社会秩序和工业纪律崩溃的环境中，各种社会群体无须党政官员煽动，也会提出特殊利益诉求。没有统一指挥的学生组织发生分裂，并采取日益暴力的行动和内部对抗，也并不令人意外。简言之，地方社会和政治混乱的责任，并不在于张所说的不服从官员所采取的阻挠行为，而在于中央鼓动并授权攻击地方政治权威的范围扩大。

The passage of time has tended to confirm the accuracy of Walder's
explanation. Certainly, it conforms more closely to the unfolding
situation in Zhejiang.\footnote{In my doctoral thesis I uncritically
  propounded Chang's interpretation. See Forster, ``The Hangzhou
  Incident of 1975'', pp.~44-50.

  在我的博士论文中，我曾未加批判地提出张的解释。见 Forster, ``The
  Hangzhou Incident of 1975''，第44-50页。} Nevertheless, it bestows on
provincial officials too large a measure of political good intention. As
a group they were grappling with a crisis of enormous dimensions. Their
most senior member, Peng Zhen, had already lost his posts in Beijing.
Precedent, established party norms,\footnote{See Frederick C. Teiwes,
  Politics and Purges in China: Rectification and the Decline of Party
  Norms, 1950-65 (White Plains, New York: M. E. Sharpe Inc., 1979).

  见 Frederick C. Teiwes, Politics and Purges in China: Rectification
  and the Decline of Party Norms, 1950-65 (White Plains, New York: M. E.
  Sharpe Inc., 1979)。} experience and common sense all seemed of little
assistance to them and it would have been surprising if desperate
circumstances had not produced novel counter-measures. What is certain
is that Jiang Hua's spirited and multi-faceted fight for survival,
whatever its motivations, had important consequences for the future of
the Cultural Revolution in Zhejiang. Moreover, Jiang, like his
counterparts in other provincial administrations, linked his fate with
that of the CCP. Was he not the party's principal representative in
Zhejiang and was not any attack on his leadership in effect an assault
on the party itself, and hence counter-revolutionary?

时间的推移趋向于证明沃尔德解释的准确性。当然，它更符合浙江事态的发展。\footnote{In
  my doctoral thesis I uncritically propounded Chang's interpretation.
  See Forster, ``The Hangzhou Incident of 1975'', pp.~44-50.

  在我的博士论文中，我曾未加批判地提出张的解释。见 Forster, ``The
  Hangzhou Incident of 1975''，第44-50页。}
然而，这种解释赋予省级官员过多的政治善意。作为一个群体，他们正在应对一场规模巨大的危机。他们中级别最高的人物彭真，已经在北京失去职务。先例、既定的党内规范、\footnote{See
  Frederick C. Teiwes, Politics and Purges in China: Rectification and
  the Decline of Party Norms, 1950-65 (White Plains, New York: M. E.
  Sharpe Inc., 1979).

  见 Frederick C. Teiwes, Politics and Purges in China: Rectification
  and the Decline of Party Norms, 1950-65 (White Plains, New York: M. E.
  Sharpe Inc., 1979)。}
经验和常识似乎都帮不上什么忙；如果绝境没有催生新的反制措施，反倒会令人惊讶。可以确定的是，无论动机为何，江华充满活力且多方面的求生斗争，对浙江文化大革命的未来产生了重要后果。此外，江同其他省级行政机关中的同行一样，把自己的命运同中共的命运联系在一起。难道他不是党在浙江的主要代表吗？任何对他领导的攻击，难道不实际上就是对党本身的攻击，因而就是反革命的吗？

Presumably, Mao's presence in Hangzhou in 1966 for such crucial meetings
as the May enlarged Politburo meeting, which drew up the May 16 Circular
directing the Cultural Revolution into the political realm, provided
Jiang with the opportunity to gauge the Chairman's mood and intentions.
This access to Mao and his entourage was not enjoyed by Jiang's
colleagues in other provinces. The question arises whether Jiang used
the Chairman's proximity,\footnote{Mao visited Hangzhou at least once
  every year between 1953 and 1966, and sporadically thereafter until
  his last visit in 1975. He was present in December 1957 when Jiang Hua
  launched, at the Chairman's behest, his savage attack on the Governor
  of Zhejiang and other leading officials. In 1963 Mao wrote the
  important paragraphs published separately as ``Where do correct ideas
  come from?'' in Hangzhou and the May 16 circular of 1966 was drawn up
  at an enlarged meeting of the Politburo Standing Committee in the
  city. According to a Taiwan report, Jiang Hua attended this meeting.
  Facts \& Features, 1:13 (April 17,1968), p.~9. Mao called the
  beautiful lake resort his second home. See the reminiscences of Jiang
  Hua, ZJRB, December 26,1983, p.~3 and Wang Fang, ``重视学习,
  善于学习'' (Paid close attention to study and was adept at study),
  ZJRB, December 23,1983, p.~1. Wang was in charge of security for Mao
  on every trip which the Chairman made to the province between 1956 and
  1966.

  1953年至1966年间，毛每年至少访问杭州一次，此后也零星来访，直到1975年最后一次来访。1957年12月，江华按主席旨意猛烈攻击浙江省长和其他主要官员时，毛就在场。1963年，毛在杭州写下后来以《人的正确思想是从哪里来的？》单独发表的重要段落；1966年的``五一六通知''也是在该市一次政治局常委扩大会议上起草的。据台湾报道，江华出席了这次会议。Facts
  \& Features,
  1:13（1968年4月17日），第9页。毛称这个美丽的湖滨胜地为他的第二故乡。见江华回忆，ZJRB，1983年12月26日，第3页；以及王芳，《重视学习,
  善于学习》（Paid close attention to study and was adept at
  study），ZJRB，1983年12月23日，第1页。1956年至1966年主席每次赴浙江，王都负责毛的安全。}
his contacts in Shanghai -- where the East China Bureau of the CCP (of
which he was a member) was located and where the Zhejiang provincial
government maintained an office\footnote{The office had been opened in
  1958 and was closed in the early stages of the Cultural Revolution.
  See, Zhejiang shengqing gaiyao, p.~680.

  该办事处于1958年开设，并在文化大革命初期关闭。见 Zhejiang shengqing
  gaiyao，第680页。} -- to discover Mao's intentions in order to promote
or sabotage the unfolding Cultural Revolution.

可以推测，毛在1966年身处杭州，并参加了诸如制定``五一六通知''、将文化大革命引入政治领域的5月政治局扩大会议等关键会议，这给江提供了揣摩主席情绪和意图的机会。其他省份的江的同级干部并不享有这种接触毛及其随行人员的机会。问题在于，江是否利用主席近在咫尺这一条件，\footnote{Mao
  visited Hangzhou at least once every year between 1953 and 1966, and
  sporadically thereafter until his last visit in 1975. He was present
  in December 1957 when Jiang Hua launched, at the Chairman's behest,
  his savage attack on the Governor of Zhejiang and other leading
  officials. In 1963 Mao wrote the important paragraphs published
  separately as ``Where do correct ideas come from?'' in Hangzhou and
  the May 16 circular of 1966 was drawn up at an enlarged meeting of the
  Politburo Standing Committee in the city. According to a Taiwan
  report, Jiang Hua attended this meeting. Facts \& Features, 1:13
  (April 17,1968), p.~9. Mao called the beautiful lake resort his second
  home. See the reminiscences of Jiang Hua, ZJRB, December 26,1983, p.~3
  and Wang Fang, ``重视学习, 善于学习'' (Paid close attention to study
  and was adept at study), ZJRB, December 23,1983, p.~1. Wang was in
  charge of security for Mao on every trip which the Chairman made to
  the province between 1956 and 1966.

  1953年至1966年间，毛每年至少访问杭州一次，此后也零星来访，直到1975年最后一次来访。1957年12月，江华按主席旨意猛烈攻击浙江省长和其他主要官员时，毛就在场。1963年，毛在杭州写下后来以《人的正确思想是从哪里来的？》单独发表的重要段落；1966年的``五一六通知''也是在该市一次政治局常委扩大会议上起草的。据台湾报道，江华出席了这次会议。Facts
  \& Features,
  1:13（1968年4月17日），第9页。毛称这个美丽的湖滨胜地为他的第二故乡。见江华回忆，ZJRB，1983年12月26日，第3页；以及王芳，《重视学习,
  善于学习》（Paid close attention to study and was adept at
  study），ZJRB，1983年12月23日，第1页。1956年至1966年主席每次赴浙江，王都负责毛的安全。}
以及他在上海的联系------中共中央华东局（他是其成员）所在地，且浙江省政府在那里设有办事处\footnote{The
  office had been opened in 1958 and was closed in the early stages of
  the Cultural Revolution. See, Zhejiang shengqing gaiyao, p.~680.

  该办事处于1958年开设，并在文化大革命初期关闭。见 Zhejiang shengqing
  gaiyao，第680页。}------来探知毛的意图，以推动或破坏正在展开的文化大革命。

Jiang had powerful patrons in Beijing to assist him assess and adapt to
the new situation. The most important of these was Tan Zhenlin. Like
Jiang, Tan was an old associate of Mao's dating back to the days on the
Jinggang mountains. At the beginning of the Cultural Revolution he was a
member of the Politburo, Vice-premier of the State Council and Director
of its office of agriculture and forestry.\footnote{For details of Tan's
  career see ``Tan Zhenlin - A Maoist `Liberated' Cadre'', I\&S, 10: 6
  (1974), pp.~101-05; the memorial speech by Hu Yaobang, ZJRB, October
  6,1983, p.~2; and the reminiscences by Tan Qilong, ``怀念谭震林同志''
  (In memory of Comrade Tan Zhenlin), 人物 (Personages), No.~6 (1984),
  pp.~43-8.

  关于谭的生涯详情，见``Tan Zhenlin - A Maoist `Liberated' Cadre'',
  I\&S, 10:6 (1974),
  第101-105页；胡耀邦悼词，ZJRB，1983年10月6日，第2页；以及谭启龙回忆《怀念谭震林同志》(In
  memory of Comrade Tan
  Zhenlin)，《人物》(Personages)，1984年第6期，第43-48页。} Tan had been
the first post-1949 provincial leader of Zhejiang until his transfer to
Jiangsu province in 1952. At that time Jiang served under Tan as a
member of the CCP ZPC standing committee and as Second
Deputy-secretary.\footnote{Zhejiang shengqing gaiyao, p.~24.

  Zhejiang shengqing gaiyao，第24页。} After his transfer to Beijing,
Tan retained strong personal and institutional links with his old
province and the whole of East China. However, the maintenance of close
relations with powerful central politicians was a two-edged sword. While
central patrons remained in good standing with Mao, the benefits were
passed on to their followers; but once they fell into disfavor all their
close associates were implicated. Zhao Ziyang basked in the glory of Tao
Zhu's meteoric rise in Beijing in 1966 and shared in full the
consequences of his disgrace at the end of that year.\footnote{Zhao Wei,
  Zhao Ziyang zhuan, pp.~139-40, 145-48.

  赵蔚，《赵紫阳传》，第139-140、145-148页。} Similarly, the
reverberations accompanying Tan Zhenlin's fall in February 1967 were
felt in East China.

江在北京有强有力的庇护者，帮助他评估并适应新形势。其中最重要的是谭震林。同江一样，谭也是毛在井冈山时期的老相识。文化大革命开始时，他是政治局委员、国务院副总理兼农林办公室主任。\footnote{For
  details of Tan's career see ``Tan Zhenlin - A Maoist `Liberated'
  Cadre'', I\&S, 10: 6 (1974), pp.~101-05; the memorial speech by Hu
  Yaobang, ZJRB, October 6,1983, p.~2; and the reminiscences by Tan
  Qilong, ``怀念谭震林同志'' (In memory of Comrade Tan Zhenlin), 人物
  (Personages), No.~6 (1984), pp.~43-8.

  关于谭的生涯详情，见``Tan Zhenlin - A Maoist `Liberated' Cadre'',
  I\&S, 10:6 (1974),
  第101-105页；胡耀邦悼词，ZJRB，1983年10月6日，第2页；以及谭启龙回忆《怀念谭震林同志》(In
  memory of Comrade Tan
  Zhenlin)，《人物》(Personages)，1984年第6期，第43-48页。}
谭曾是1949年后浙江第一任省级领导人，直到1952年调往江苏省。当时江在谭手下任中共浙江省委常委和第二副书记。\footnote{Zhejiang
  shengqing gaiyao, p.~24.

  Zhejiang shengqing gaiyao，第24页。}
调往北京后，谭仍与其旧省以及整个华东保持强大的个人和制度联系。然而，与有权势的中央政治人物保持密切关系是一把双刃剑。当中央庇护者仍受毛青睐时，其追随者能够分享好处；但一旦他们失宠，所有亲近者都会受到牵连。1966年陶铸在北京急速上升时，赵紫阳沐浴在其荣耀之中；而当陶在年底蒙羞时，赵也完全承受了后果。\footnote{Zhao
  Wei, Zhao Ziyang zhuan, pp.~139-40, 145-48.

  赵蔚，《赵紫阳传》，第139-140、145-148页。}
同样，1967年2月谭震林倒台所产生的震荡也波及华东。

The acuteness of Jiang Hua's political sensitivities may well have been
put to the test in November 1965. In that month Yao Wenyuan's article
criticizing Wu Han's play was published in Shanghai. Mao ordered other
localities to republish the piece\footnote{Gao Gao and Yan Jiaqi,
  ``文化大革命''十年史 (Ten Year History of the ``Great Cultural
  Revolution''), (Tianjin: Tianjin Renmin Chubanshe, 1986), p.~7.

  高皋、严家祺，《``文化大革命''十年史》(Ten Year History of the ``Great
  Cultural Revolution'')（天津：天津人民出版社，1986），第7页。} and
Jiang responded by ordering Zhejiang Daily to print it on 24 November, a
fortnight later. Thus, Zhejiang had beaten other localities such as
Beijing and Guangzhou in its response to Mao's instructions.\footnote{Renmin
  Ribao republished the article on November 30, the PLA daily Jiefangjun
  Bao on November 21 and the Guangzhou papers on December 1. See J.
  Domes, ``Generals and Red Guards - The Role of Huang Yung-sheng and
  the Center Military Area Command in the Kuangtung Cultural
  Revolution'', Asia Quarterly 1 (1971), p.~14.

  《人民日报》于11月30日转载该文，解放军报《解放军报》于11月21日转载，广州各报于12月1日转载。见
  J. Domes, ``Generals and Red Guards - The Role of Huang Yung-sheng and
  the Center Military Area Command in the Kuangtung Cultural
  Revolution'', Asia Quarterly 1 (1971), 第14页。} A Red Guard
publication subsequently charged Jiang with political opportunism by
using his contacts in Shanghai to discover the strategic implications
behind Yao's article.\footnote{``Down with Jiang Hua!'', Wenge tongxun
  (Cultural Revolution Despatches), No.~16 (July 1968) in SCMP, 4230
  (August 1,1968), p.~3.

  ``Down with Jiang Hua!'', Wenge tongxun (Cultural Revolution
  Despatches), 第16期（1968年7月），载 SCMP,
  4230（1968年8月1日），第3页。} Equally credible is the possibility
that by his promptness in republishing Yao's article, Jiang was ensuring
that his loyalty to Mao remained unopen to doubt.

江华政治敏感性的敏锐程度，很可能在1965年11月受到了考验。当月，姚文元批判吴晗剧作的文章在上海发表。毛命令其他地方转载此文，\footnote{Gao
  Gao and Yan Jiaqi, ``文化大革命''十年史 (Ten Year History of the
  ``Great Cultural Revolution''), (Tianjin: Tianjin Renmin Chubanshe,
  1986), p.~7.

  高皋、严家祺，《``文化大革命''十年史》(Ten Year History of the ``Great
  Cultural Revolution'')（天津：天津人民出版社，1986），第7页。}
江随即命令《浙江日报》于两周后的11月24日刊登。因此，浙江在响应毛的指示方面早于北京、广州等地。\footnote{Renmin
  Ribao republished the article on November 30, the PLA daily Jiefangjun
  Bao on November 21 and the Guangzhou papers on December 1. See J.
  Domes, ``Generals and Red Guards - The Role of Huang Yung-sheng and
  the Center Military Area Command in the Kuangtung Cultural
  Revolution'', Asia Quarterly 1 (1971), p.~14.

  《人民日报》于11月30日转载该文，解放军报《解放军报》于11月21日转载，广州各报于12月1日转载。见
  J. Domes, ``Generals and Red Guards - The Role of Huang Yung-sheng and
  the Center Military Area Command in the Kuangtung Cultural
  Revolution'', Asia Quarterly 1 (1971), 第14页。}
后来一份红卫兵出版物指控江搞政治投机，称他利用自己在上海的关系探知姚文背后的战略含义。\footnote{``Down
  with Jiang Hua!'', Wenge tongxun (Cultural Revolution Despatches),
  No.~16 (July 1968) in SCMP, 4230 (August 1,1968), p.~3.

  ``Down with Jiang Hua!'', Wenge tongxun (Cultural Revolution
  Despatches), 第16期（1968年7月），载 SCMP,
  4230（1968年8月1日），第3页。}
同样可信的另一种可能是，江迅速转载姚文，是为了确保自己对毛的忠诚不容怀疑。

The same Cultural Revolution source also claimed that Jiang obtained an
advance draft of the May 16 circular which enabled him, in collaboration
with the noted theorist Hu Qiaomu, to instruct Chen Bing to publish
three articles under the pen-name of Xin Wenbing (辛文兵). These
articles attacked the Beijing ``Three Family Village'' writing team
operating under Peng Zhen's aegis.\footnote{``Down with Jiang Hua!'',
  p.~4. See ZJRB, May 18,19 and 20, 1966 for these three articles which
  were all published prominently on p.~1.

  ``Down with Jiang Hua!''，第4页。关于这三篇文章，见
  ZJRB，1966年5月18、19、20日；三文均显著刊登于第1页。} However, the
most sensational accusation levelled at Jiang was that early in 1966 he
ordered Lü Jianguang, together with public security officials from
Beijing, to spy on Mao. They allegedly emptied Mao's wastepaper baskets
for scraps of information, eavesdropped on his conversations, bugged his
residence and tapped his telephone calls. This story, apocryphal though
it may be, was to have far-reaching consequences for those allegedly
involved.\footnote{``Down with Jiang Hua!'', pp.~2-4. To cover his
  tracks Jiang, according to this document, expelled 180 personnel from
  the Public Security Bureau. This incredible charge is repeated without
  the hair-raising detail in ZJRB, November 15, 1968. In March 1968, at
  a meeting between central leaders and a delegation from Zhejiang,
  Jiang Qing made similar allegations against Jiang. At the same meeting
  Zhou Enlai commented that his correspondence had been opened and
  referred to evidence about such practices as having been provided by
  the new leader of Zhejiang, Nan Ping, and by former governor Zhou
  Jianren. The Premier went on to demand the destruction of the
  security/judicial apparatus in the province. (SCMP, 4182 May 21,1968),
  pp.~4-5. These denunciations apparently triggered a witch-hunt in the
  public security department resulting in the death of at least one
  detective in June 1968. ZJRB, August 11,1978. He probably died under
  interrogation designed to force a confession which would implicate
  Jiang Hua. Jiang's successor as 1st Secretary of the CCP ZPC, Nan
  Ping, repeated the charges against Jiang at the CCP 5th Zhejiang
  Provincial Congress held in January 1971. One month later, at a
  national conference of public security bureaux, Zhou Enlai, in a
  remarkable volte-face, rebuked all such insinuations against the
  Zhejiang Public Security Bureau and by implication, Jiang Hua, stating
  that ``we can't say that spies exercised dictatorship over us''. To
  support his position, Zhou quoted Mao as puzzling: ``some people say
  that the Zhejiang public security organs have gone rotten. I've stayed
  in Hangzhou many times and nothing has happened. How can this be
  explained?'' (有人说浙江公安机关烂了,我去杭州住过多次也没有出什么事.
  这怎么解释). Zhou's remarks were quoted by Tie Ying in his report to
  the CCP 6th Zhejiang Provincial Congress. HZRB, June 7,1978. During
  the campaign to weed out followers of Lin Biao and the Gang of Four in
  Zhejiang, Shen Ce, a member of the pre-Cultural Revolution CCP ZPC
  Standing Committee who had been ``liberated'' as a ``revolutionary
  leading cadre'' in 1967, was charged with trying to frame Jiang Hua
  for this crime. HZRB, July 26,1978.

  ``Down with Jiang
  Hua!''，第2-4页。据该文件称，为掩盖踪迹，江把公安局180名人员开除。这一令人难以置信的指控在
  ZJRB
  1968年11月15日被重复，但没有那些令人毛骨悚然的细节。1968年3月，在中央领导人与浙江代表团的一次会议上，江青对江提出了类似指控。同次会议上，周恩来说他的信件曾被拆阅，并提到新任浙江领导人南萍和前省长周建人提供了这类做法的证据。总理随后要求摧毁该省的公安司法机构。(SCMP,
  4182，1968年5月21日)，第4-5页。这些谴责显然在公安部门引发了一场追查，导致至少一名侦查员于1968年6月死亡。ZJRB，1978年8月11日。他很可能死于逼供审讯，审讯目的在于迫使其作出牵连江华的供认。江的中共浙江省委第一书记继任者南萍，在1971年1月召开的中共浙江省第五次代表大会上重复了对江的指控。一个月后，在全国公安局会议上，周恩来以显著的转向斥责了所有这类针对浙江公安局、也暗含针对江华的影射，称``不能说特务对我们实行专政''。为支持其立场，周引用毛的疑问：``有人说浙江公安机关烂了,我去杭州住过多次也没有出什么事.
  这怎么解释。''周的话被铁瑛在中共浙江省第六次代表大会报告中引用。HZRB，1978年6月7日。在浙江清查林彪和``四人帮''追随者的运动中，文革前中共浙江省委常委、1967年作为``革命领导干部''被``解放''的沈策，被指控试图以这一罪名陷害江华。HZRB，1978年7月26日。}

同一文革资料还声称，江获得了``五一六通知''的预先草稿，使他能够同著名理论家胡乔木合作，指示陈冰以``辛文兵''为笔名发表三篇文章。这些文章攻击了在彭真庇护下运作的北京``三家村''写作班子。\footnote{``Down
  with Jiang Hua!'', p.~4. See ZJRB, May 18,19 and 20, 1966 for these
  three articles which were all published prominently on p.~1.

  ``Down with Jiang Hua!''，第4页。关于这三篇文章，见
  ZJRB，1966年5月18、19、20日；三文均显著刊登于第1页。}
然而，对江最耸人听闻的指控是，1966年初他命令吕剑光会同北京公安人员监视毛。据称，他们清空毛的废纸篓以搜寻信息碎片，偷听他的谈话，在其住处安装窃听器，并监听他的电话。这个故事或许并不可信，但它将对那些据称涉入其中的人产生深远后果。\footnote{``Down
  with Jiang Hua!'', pp.~2-4. To cover his tracks Jiang, according to
  this document, expelled 180 personnel from the Public Security Bureau.
  This incredible charge is repeated without the hair-raising detail in
  ZJRB, November 15, 1968. In March 1968, at a meeting between central
  leaders and a delegation from Zhejiang, Jiang Qing made similar
  allegations against Jiang. At the same meeting Zhou Enlai commented
  that his correspondence had been opened and referred to evidence about
  such practices as having been provided by the new leader of Zhejiang,
  Nan Ping, and by former governor Zhou Jianren. The Premier went on to
  demand the destruction of the security/judicial apparatus in the
  province. (SCMP, 4182 May 21,1968), pp.~4-5. These denunciations
  apparently triggered a witch-hunt in the public security department
  resulting in the death of at least one detective in June 1968. ZJRB,
  August 11,1978. He probably died under interrogation designed to force
  a confession which would implicate Jiang Hua. Jiang's successor as 1st
  Secretary of the CCP ZPC, Nan Ping, repeated the charges against Jiang
  at the CCP 5th Zhejiang Provincial Congress held in January 1971. One
  month later, at a national conference of public security bureaux, Zhou
  Enlai, in a remarkable volte-face, rebuked all such insinuations
  against the Zhejiang Public Security Bureau and by implication, Jiang
  Hua, stating that ``we can't say that spies exercised dictatorship
  over us''. To support his position, Zhou quoted Mao as puzzling:
  ``some people say that the Zhejiang public security organs have gone
  rotten. I've stayed in Hangzhou many times and nothing has happened.
  How can this be explained?''
  (有人说浙江公安机关烂了,我去杭州住过多次也没有出什么事. 这怎么解释).
  Zhou's remarks were quoted by Tie Ying in his report to the CCP 6th
  Zhejiang Provincial Congress. HZRB, June 7,1978. During the campaign
  to weed out followers of Lin Biao and the Gang of Four in Zhejiang,
  Shen Ce, a member of the pre-Cultural Revolution CCP ZPC Standing
  Committee who had been ``liberated'' as a ``revolutionary leading
  cadre'' in 1967, was charged with trying to frame Jiang Hua for this
  crime. HZRB, July 26,1978.

  ``Down with Jiang
  Hua!''，第2-4页。据该文件称，为掩盖踪迹，江把公安局180名人员开除。这一令人难以置信的指控在
  ZJRB
  1968年11月15日被重复，但没有那些令人毛骨悚然的细节。1968年3月，在中央领导人与浙江代表团的一次会议上，江青对江提出了类似指控。同次会议上，周恩来说他的信件曾被拆阅，并提到新任浙江领导人南萍和前省长周建人提供了这类做法的证据。总理随后要求摧毁该省的公安司法机构。(SCMP,
  4182，1968年5月21日)，第4-5页。这些谴责显然在公安部门引发了一场追查，导致至少一名侦查员于1968年6月死亡。ZJRB，1978年8月11日。他很可能死于逼供审讯，审讯目的在于迫使其作出牵连江华的供认。江的中共浙江省委第一书记继任者南萍，在1971年1月召开的中共浙江省第五次代表大会上重复了对江的指控。一个月后，在全国公安局会议上，周恩来以显著的转向斥责了所有这类针对浙江公安局、也暗含针对江华的影射，称``不能说特务对我们实行专政''。为支持其立场，周引用毛的疑问：``有人说浙江公安机关烂了,我去杭州住过多次也没有出什么事.
  这怎么解释。''周的话被铁瑛在中共浙江省第六次代表大会报告中引用。HZRB，1978年6月7日。在浙江清查林彪和``四人帮''追随者的运动中，文革前中共浙江省委常委、1967年作为``革命领导干部''被``解放''的沈策，被指控试图以这一罪名陷害江华。HZRB，1978年7月26日。}

\section{Jiang Hua and the Cultural Revolution,
1966-67}\label{jiang-hua-and-the-cultural-revolution-1966-67}

\section{江华与文化大革命，1966-67年}\label{ux6c5fux534eux4e0eux6587ux5316ux5927ux9769ux547d1966-67ux5e74}

Between June and July 1966 Jiang faithfully implemented central
directives concerning the conduct of the Cultural Revolution issued by
Liu Shaoqi and the Central Committee. Work teams were sent out in
mid-June and recalled, on Mao's instructions, at the end of
July.\footnote{See Liu Guokai, ``A Brief Analysis'', p.~20; and
  《当代中国的江苏》编委会, 江苏省档案局编 (Compilation Committee of
  ``Contemporary China's Jiangsu'' and Jiangsu Provincial Archives
  Bureau eds.), 江苏大事记 1949-1985 (Chronicle of Major Events in
  Jiangsu Province, 1949-1985), (Nanjing?: Jiangsu Renmin Chubanshe,
  1988), pp.~277, 279; Zhao Wei, Zhao Ziyang zhuan, pp.~136-37.

  见刘国凯，``A Brief
  Analysis''，第20页；以及《当代中国的江苏》编委会、江苏省档案局编，《江苏大事记
  1949-1985》(Chronicle of Major Events in Jiangsu Province,
  1949-1985)（南京？：江苏人民出版社，1988），第277、279页；赵蔚，《赵紫阳传》，第136-137页。}
In Zhejiang the work teams were under the overall leadership of Chen
Bing who had been appointed Director of the provincial Cultural
Revolution Group.\footnote{Jiangsu's Cultural Revolution group was
  established in late July 1966. Jiangsusheng dashiji, p.~278. In
  Guangdong province it was established in early May 1966 and dissolved
  on September 4, 1966. See Zhao Wei, Zhao Ziyang zhuan, pp.~128,141.

  江苏的文化革命小组成立于1966年7月下旬。Jiangsusheng
  dashiji，第278页。广东省则于1966年5月初成立，并于1966年9月4日解散。见赵蔚，《赵紫阳传》，第128、141页。}
Liu Dan, one of the leading party officials and administrators at the
prestigious Zhejiang University, and the head of its Cultural Revolution
Group, was dismissed at a public rally on 22 June. The meeting was held
on campus and attended by Cao Xiangren, representing the CCP
ZPC.\footnote{ZJRB, June 24,1966, p.~1.

  ZJRB，1966年6月24日，第1页。} Liu was charged with investigating the
class backgrounds of students who had pasted up wall posters and
excluding those lacking the right family background from further
participation in the campaign. A Taiwan report alleged that Jiang Hua
himself was behind this investigation of the students' class
background.\footnote{Facts \& Features, 1:13, p.~10. For a valuable
  discussion of this important issue, see Gordon White, The Politics of
  Class and Class Origin: The Case of the Cultural Revolution
  (Contemporary China Papers, 9 (Canberra: Australian National
  University, 1976).

  Facts \& Features, 1:13，第10页。关于这一重要问题的有价值讨论，见
  Gordon White, The Politics of Class and Class Origin: The Case of the
  Cultural Revolution (Contemporary China Papers, 9 (Canberra:
  Australian National University, 1976)。} If these accusations were
true they prove that Liu Dan and Jiang were implementing the party's
class line. That this line changed and Liu's diligence brought about his
downfall illustrated the volatility and unpredictability of the
political situation at the time. In the light of later events, the ZPC's
decision to dismiss Liu on these grounds appears to have been one that
appealed to the increasingly militant Red Guard organizations. They were
demanding that membership not be restricted to children of the ``five
red'' categories of worker, peasant, cadre, soldier and revolutionary
martyr. Selected sacrificial victims in the propaganda and cultural
systems were also dismissed and vilified at this time.\footnote{For the
  names of some of the leading officials in these departments who
  eventually died as a result of interrogation and torture see ZPS,
  November 23,1977, SWB/FE/5677/BI/7, HZRB, June 11,1978; ZPS, July
  16,1978, SWB/FE/5873/BII/18; ZPS, August 13,1978, SWB/FE/5900/BII/12;
  ZPS, August 20, 1978, SWB/FE/5900/BII/12-13; ZPS, September 7, 1978,
  SWB/FE/5917/BII/7-8. Leading officials of the same departments were
  targeted in other provinces. See Jiangsusheng dashiji, pp.~277, 279.

  关于这些部门中部分主要官员姓名，他们最终因审讯和酷刑而死亡，见
  ZPS，1977年11月23日，SWB/FE/5677/BI/7；HZRB，1978年6月11日；ZPS，1978年7月16日，SWB/FE/5873/BII/18；ZPS，1978年8月13日，SWB/FE/5900/BII/12；ZPS，1978年8月20日，SWB/FE/5900/BII/12-13；ZPS，1978年9月7日，SWB/FE/5917/BII/7-8。其他省份同类部门的主要官员也成为目标。见
  Jiangsusheng dashiji，第277、279页。} To retain the initiative, on 26
June Jiang delivered a major address before a gathering of 10,000 CCP
members and a full complement of party and military leaders. In his
oration on the topic of Mao Zedong Thought\footnote{Jiang Hua,
  高举毛泽东思想伟大红旗, 把无产阶级文化大革命进行到底 (Hold aloft the
  great red banner of Mao Zedong's thought and carry on the Great
  Cultural Revolution through to the end), ZJRB, June 28,1966, pp.~1-2.

  江华，《高举毛泽东思想伟大红旗, 把无产阶级文化大革命进行到底》(Hold
  aloft the great red banner of Mao Zedong's thought and carry on the
  Great Cultural Revolution through to the
  end)，ZJRB，1966年6月28日，第1-2页。} Jiang explained the significance
of the Cultural Revolution as an ideological and political class
struggle against both representatives of the bourgeoisie in the CCP and
bourgeois intellectuals. He cautioned those who practised political
deception by waving the red banner only to oppose it -- a warning no
doubt delivered in all sincerity but one that before long would rebound
on its speaker.

1966年6月至7月，江忠实执行刘少奇和中央委员会发布的有关开展文化大革命的中央指示。6月中旬派出工作组，并按毛的指示于7月底撤回。\footnote{See
  Liu Guokai, ``A Brief Analysis'', p.~20; and 《当代中国的江苏》编委会,
  江苏省档案局编 (Compilation Committee of ``Contemporary China's
  Jiangsu'' and Jiangsu Provincial Archives Bureau eds.), 江苏大事记
  1949-1985 (Chronicle of Major Events in Jiangsu Province, 1949-1985),
  (Nanjing?: Jiangsu Renmin Chubanshe, 1988), pp.~277, 279; Zhao Wei,
  Zhao Ziyang zhuan, pp.~136-37.

  见刘国凯，``A Brief
  Analysis''，第20页；以及《当代中国的江苏》编委会、江苏省档案局编，《江苏大事记
  1949-1985》(Chronicle of Major Events in Jiangsu Province,
  1949-1985)（南京？：江苏人民出版社，1988），第277、279页；赵蔚，《赵紫阳传》，第136-137页。}
在浙江，工作组由已被任命为省文化革命小组组长的陈冰统一领导。\footnote{Jiangsu's
  Cultural Revolution group was established in late July 1966.
  Jiangsusheng dashiji, p.~278. In Guangdong province it was established
  in early May 1966 and dissolved on September 4, 1966. See Zhao Wei,
  Zhao Ziyang zhuan, pp.~128,141.

  江苏的文化革命小组成立于1966年7月下旬。Jiangsusheng
  dashiji，第278页。广东省则于1966年5月初成立，并于1966年9月4日解散。见赵蔚，《赵紫阳传》，第128、141页。}
刘丹是著名的浙江大学主要党政干部之一，也是该校文化革命小组负责人；6月22日，他在一次公开大会上被撤职。会议在校园举行，曹祥仁代表中共浙江省委出席。\footnote{ZJRB,
  June 24,1966, p.~1.

  ZJRB，1966年6月24日，第1页。}
刘被指控调查张贴大字报学生的阶级成分，并将家庭出身不合格者排除在运动的进一步参与之外。台湾方面一份报道声称，江华本人就在这次学生阶级成分调查背后。\footnote{Facts
  \& Features, 1:13, p.~10. For a valuable discussion of this important
  issue, see Gordon White, The Politics of Class and Class Origin: The
  Case of the Cultural Revolution (Contemporary China Papers, 9
  (Canberra: Australian National University, 1976).

  Facts \& Features, 1:13，第10页。关于这一重要问题的有价值讨论，见
  Gordon White, The Politics of Class and Class Origin: The Case of the
  Cultural Revolution (Contemporary China Papers, 9 (Canberra:
  Australian National University, 1976)。}
如果这些指控属实，它们证明刘丹和江是在执行党的阶级路线。这条路线发生变化，而刘的勤勉反倒导致其垮台，说明了当时政治局势的易变和不可预测。鉴于后来的事件，省委以这些理由撤换刘的决定，似乎迎合了日益激进的红卫兵组织。红卫兵要求成员资格不得限于工人、贫农下中农、干部、军人和革命烈士这``五红''类别的子女。宣传和文化系统中被选中的牺牲品，也在此时遭到撤职和诋毁。\footnote{For
  the names of some of the leading officials in these departments who
  eventually died as a result of interrogation and torture see ZPS,
  November 23,1977, SWB/FE/5677/BI/7, HZRB, June 11,1978; ZPS, July
  16,1978, SWB/FE/5873/BII/18; ZPS, August 13,1978, SWB/FE/5900/BII/12;
  ZPS, August 20, 1978, SWB/FE/5900/BII/12-13; ZPS, September 7, 1978,
  SWB/FE/5917/BII/7-8. Leading officials of the same departments were
  targeted in other provinces. See Jiangsusheng dashiji, pp.~277, 279.

  关于这些部门中部分主要官员姓名，他们最终因审讯和酷刑而死亡，见
  ZPS，1977年11月23日，SWB/FE/5677/BI/7；HZRB，1978年6月11日；ZPS，1978年7月16日，SWB/FE/5873/BII/18；ZPS，1978年8月13日，SWB/FE/5900/BII/12；ZPS，1978年8月20日，SWB/FE/5900/BII/12-13；ZPS，1978年9月7日，SWB/FE/5917/BII/7-8。其他省份同类部门的主要官员也成为目标。见
  Jiangsusheng dashiji，第277、279页。}
为了掌握主动，6月26日，江在一万名中共党员以及党政军领导人齐集的大会上发表重要讲话。在关于毛泽东思想的演说中，\footnote{Jiang
  Hua, 高举毛泽东思想伟大红旗, 把无产阶级文化大革命进行到底 (Hold aloft
  the great red banner of Mao Zedong's thought and carry on the Great
  Cultural Revolution through to the end), ZJRB, June 28,1966, pp.~1-2.

  江华，《高举毛泽东思想伟大红旗, 把无产阶级文化大革命进行到底》(Hold
  aloft the great red banner of Mao Zedong's thought and carry on the
  Great Cultural Revolution through to the
  end)，ZJRB，1966年6月28日，第1-2页。}
江解释了文化大革命的意义，称其是针对中共内部资产阶级代表人物和资产阶级知识分子的思想政治阶级斗争。他警告那些打着红旗反红旗、进行政治欺骗的人；这个警告无疑出自真诚，但不久后就会反噬到发言者本人。

The first major challenge that Jiang Hua and his colleagues faced arose
from the formation of local Red Guards and the arrival of Red Guard
emissaries from Beijing to check and supervise the progress of the
Cultural Revolution in the province.\footnote{For the impact of Red
  Guard emissaries in Guangdong see, Zhao Wei, Zhao Ziyang zhuan,
  p.~142; and in Shanghai, see Ye Yonglie, 王洪文兴衰录 (The Rise and
  Fall of Wang Hongwen), (Changchun: Shidai Wenyi Chubanshe, 1989),
  chs.~4-5.

  关于红卫兵使者在广东的影响，见赵蔚，《赵紫阳传》，第142页；关于上海，见叶永烈，《王洪文兴衰录》(The
  Rise and Fall of Wang
  Hongwen)（长春：时代文艺出版社，1989），第4-5章。} For the CCP ZPC,
the principal dilemma at this time was how to keep in step with the
center by implementing a set of ill-defined, almost incomprehensible
directives, while at the same time carrying on with everyday
administration, which was rendered increasingly difficult by the
implementation of these very same directives. The presence of Red Guards
to check on the party committee's performance made such time-honored
tactics as prevarication, distortion or feigned compliance both
difficult and politically risky.

江华及其同事面临的第一个重大挑战，来自本地红卫兵的形成，以及北京红卫兵使者抵达浙江检查并监督本省文化大革命进展。\footnote{For
  the impact of Red Guard emissaries in Guangdong see, Zhao Wei, Zhao
  Ziyang zhuan, p.~142; and in Shanghai, see Ye Yonglie, 王洪文兴衰录
  (The Rise and Fall of Wang Hongwen), (Changchun: Shidai Wenyi
  Chubanshe, 1989), chs.~4-5.

  关于红卫兵使者在广东的影响，见赵蔚，《赵紫阳传》，第142页；关于上海，见叶永烈，《王洪文兴衰录》(The
  Rise and Fall of Wang
  Hongwen)（长春：时代文艺出版社，1989），第4-5章。}
对中共浙江省委而言，当时的主要困境是：一方面要执行一套界定模糊、几乎难以理解的指示，以跟上中央步调；另一方面又要维持日常行政，而正是这些指示的执行使日常行政越来越困难。红卫兵到场检查党委表现，使拖延、歪曲或假装服从等由来已久的策略既难以实行，又在政治上充满风险。

A Yugoslav correspondent stationed in Beijing reported that at the 11th
Plenum the East China delegates voted as a block for Mao's Cultural
Revolution, thereby helping the Chairman carry the day in a tight
situation.\footnote{Ming Bao, (H.K.), January 1-2, 1967 in SCMP, 3855,
  p.~3.

  《明报》（香港），1967年1月1-2日，载 SCMP, 3855，第3页。} This gesture
may have impressed Mao but soon after the plenum Jiang faced the
disruptive activities of Red Guards. Their first public appearance was
on 23 August when they were joined in street marches by workers and
peasants from suburban communes.\footnote{See ZJRB, August 24, 1966;
  ZPS, August 24,1966, SWB/FE/2253/B/6-7; 浙江省经济研究中心编 (Zhejiang
  provincial economic research center), 浙江省情 (The Affairs of
  Zhejiang), (Hangzhou: Zhejiang Renmin Chubanshe, 1986), p.~1052. The
  date given here of June 23 is clearly a mistake.

  见
  ZJRB，1966年8月24日；ZPS，1966年8月24日，SWB/FE/2253/B/6-7；浙江省经济研究中心编，《浙江省情》(The
  Affairs of
  Zhejiang)（杭州：浙江人民出版社，1986），第1052页。此处给出的6月23日日期显然有误。}
In their first attack on the ``four olds'' the youthful iconoclasts
marched around Hangzhou changing the names of streets, restaurants and
public buildings and destroying precious cultural relics in the historic
city.\footnote{Gao Gao and Yan Jiaqi, ``Wenhua dageming'', p.~61, lists
  some of the buildings involved. However, their reference to
  destruction of the tomb of the beautiful and talented courtesan Su
  Xiaoxiao (苏小小) near Xiling bridge is incorrect. ``Credit'' for this
  act of vandalism has previously been bestowed upon Jiang Hua and
  Hangzhou Municipal Party Secretary Wang Pingyi in a decision that they
  had made in December 1964. See HZRB, April 11, 1978, p.~2. According
  to another source the real ``villain'' behind the removal of the tombs
  of Su Xiaoxiao, the legendary strongman of the classic novel ``Water
  Margin'' Wu Song, and martyrs of the 1911 Revolution was none other
  than Mao Zedong himself. In the spring of 1956 while resting in
  Hangzhou, the Chairman remarked to a ``responsible person'' of the
  Zhejiang provincial committee (probably Jiang Hua) that ``there are
  too many tombs around the West Lake; they can be dismantled and moved
  out to the suburbs and buried there. Can't the dead also live a
  collective life?''! Mao's casual words were acted immediately, causing
  a furor among prominent citizens of the city. See Chen Xiuliang,
  ``西子湖畔风波迭起'' (Turbulent times on the banks of the West Lake),
  联合时报 (United Times), July 8,1988.

  高皋、严家祺，《Wenhua
  dageming》，第61页，列举了涉及的一些建筑。然而，他们关于西泠桥附近才貌双全名妓苏小小墓被毁的说法并不正确。此前，这一破坏行为的``功劳''曾被归于江华和杭州市委书记王平夷1964年12月作出的决定。见
  HZRB，1978年4月11日，第2页。另据一项资料，拆除苏小小墓、古典小说《水浒传》中传奇壮士武松墓以及1911年革命烈士墓的真正``罪魁''不是别人，正是毛泽东本人。1956年春，主席在杭州休息时对浙江省委一名``负责人''（可能是江华）说：``西湖周围坟墓太多，可以拆掉迁到郊外去埋。死人也可以过集体生活嘛？''毛的随口之言立即得到执行，引发杭州名流哗然。见陈修良，《西子湖畔风波迭起》(Turbulent
  times on the banks of the West Lake)，《联合时报》(United
  Times)，1988年7月8日。} By September Red Guards from Beijing and
elsewhere had arrived in Hangzhou to push along their more ``backward''
fellow-students in the provinces. At two meetings with the provincial
leadership they submitted suggestions concerning its conduct of the
campaign and gave initial approval of its performance. Red Guards from
outside Zhejiang attended the first meeting, while those from within the
province were present at the second.\footnote{ZJRB, September 11,1966,
  p.~1; ZJRB, September 14,1966, p.~1.

  ZJRB，1966年9月11日，第1页；ZJRB，1966年9月14日，第1页。} Jiang Hua's
presence at these meetings were his last recorded public appearances in
Zhejiang during the Cultural Revolution.

一名驻北京的南斯拉夫记者报道说，在十一中全会上，华东代表团作为一个整体投票支持毛的文化大革命，从而帮助主席在紧张局势中取胜。\footnote{Ming
  Bao, (H.K.), January 1-2, 1967 in SCMP, 3855, p.~3.

  《明报》（香港），1967年1月1-2日，载 SCMP, 3855，第3页。}
这一姿态或许给毛留下了印象，但全会后不久，江就面对红卫兵的破坏性活动。红卫兵首次公开亮相是在8月23日，当时郊区公社的工人和农民也加入了他们的街头游行。\footnote{See
  ZJRB, August 24, 1966; ZPS, August 24,1966, SWB/FE/2253/B/6-7;
  浙江省经济研究中心编 (Zhejiang provincial economic research center),
  浙江省情 (The Affairs of Zhejiang), (Hangzhou: Zhejiang Renmin
  Chubanshe, 1986), p.~1052. The date given here of June 23 is clearly a
  mistake.

  见
  ZJRB，1966年8月24日；ZPS，1966年8月24日，SWB/FE/2253/B/6-7；浙江省经济研究中心编，《浙江省情》(The
  Affairs of
  Zhejiang)（杭州：浙江人民出版社，1986），第1052页。此处给出的6月23日日期显然有误。}
在第一次``破四旧''行动中，这些年轻的反传统者在杭州城内游行，更改街道、饭店和公共建筑的名称，并在这座历史名城中毁坏珍贵文物。\footnote{Gao
  Gao and Yan Jiaqi, ``Wenhua dageming'', p.~61, lists some of the
  buildings involved. However, their reference to destruction of the
  tomb of the beautiful and talented courtesan Su Xiaoxiao (苏小小) near
  Xiling bridge is incorrect. ``Credit'' for this act of vandalism has
  previously been bestowed upon Jiang Hua and Hangzhou Municipal Party
  Secretary Wang Pingyi in a decision that they had made in December
  1964. See HZRB, April 11, 1978, p.~2. According to another source the
  real ``villain'' behind the removal of the tombs of Su Xiaoxiao, the
  legendary strongman of the classic novel ``Water Margin'' Wu Song, and
  martyrs of the 1911 Revolution was none other than Mao Zedong himself.
  In the spring of 1956 while resting in Hangzhou, the Chairman remarked
  to a ``responsible person'' of the Zhejiang provincial committee
  (probably Jiang Hua) that ``there are too many tombs around the West
  Lake; they can be dismantled and moved out to the suburbs and buried
  there. Can't the dead also live a collective life?''! Mao's casual
  words were acted immediately, causing a furor among prominent citizens
  of the city. See Chen Xiuliang, ``西子湖畔风波迭起'' (Turbulent times
  on the banks of the West Lake), 联合时报 (United Times), July 8,1988.

  高皋、严家祺，《Wenhua
  dageming》，第61页，列举了涉及的一些建筑。然而，他们关于西泠桥附近才貌双全名妓苏小小墓被毁的说法并不正确。此前，这一破坏行为的``功劳''曾被归于江华和杭州市委书记王平夷1964年12月作出的决定。见
  HZRB，1978年4月11日，第2页。另据一项资料，拆除苏小小墓、古典小说《水浒传》中传奇壮士武松墓以及1911年革命烈士墓的真正``罪魁''不是别人，正是毛泽东本人。1956年春，主席在杭州休息时对浙江省委一名``负责人''（可能是江华）说：``西湖周围坟墓太多，可以拆掉迁到郊外去埋。死人也可以过集体生活嘛？''毛的随口之言立即得到执行，引发杭州名流哗然。见陈修良，《西子湖畔风波迭起》(Turbulent
  times on the banks of the West Lake)，《联合时报》(United
  Times)，1988年7月8日。}
到9月，来自北京及其他地方的红卫兵抵达杭州，推动地方上较为``落后''的同学。在同省领导层的两次会议上，他们就领导层开展运动的方式提出建议，并初步认可其表现。第一次会议有浙江以外的红卫兵参加，第二次会议则有省内红卫兵参加。\footnote{ZJRB,
  September 11,1966, p.~1; ZJRB, September 14,1966, p.~1.

  ZJRB，1966年9月11日，第1页；ZJRB，1966年9月14日，第1页。}
这些会议是江华在文化大革命期间于浙江最后几次有记录的公开露面。

By the autumn of 1966 the Hangzhou Tertiary and Secondary Schools' Red
Guards Headquarters (杭州大中学校红卫兵司令部) had split into three
major groups.\footnote{See n.~40.

  见注释40。} The first of the groups was Red Storm (红暴), which was
based at Zhejiang University and led by Liu Ying (刘英) and Teng Zhu
(滕铸). Jinggang Mountain (井冈山), the least influential of the three
organizations, occupied the middle ground but tended to side with Red
Storm in the clashes which this group became involved in with its
principal rival, Red Third Headquarters (红三司). Red Third HQ found its
strongest support at Hangzhou University and the Zhejiang Academy of
Fine Arts. Its leaders were Li Xiantong (李显通) -- a student from
Hangzhou University's Department of Foreign Languages -- and Du Yingxin
(杜英信) and Zhang Yongsheng -- a young teacher and graduate student,
respectively, from the Academy. In a letter of 5 July 1966 Zhang had
described his college as a black gang's den where he had waged struggles
for six years. He also complained of the incessant persecution he had
suffered at the hands of the demons (魔鬼) who controlled the
college.\footnote{中共浙江省委清查王张江姚材料组 (Materials Group of the
  CCP Zhejiang Provincial Committee investigating Wang {[}Hongwen{]},
  Zhang {[}Chunqiao{]}, Jiang {[}Qing{]} and Yao {[}Wenyuan{]}),
  ``'四人帮'在浙江的亲信, 现行反革命分子张永生的罪证材料'' (Criminal
  Evidence concerning the ``Gang of Four's trusted follower in Zhejiang,
  the counter-revolutionary Zhang Yongsheng - hereafter Zhang
  Yongshengde zuizheng cailiao), 宣传通讯 (Propaganda Bulletin), No.~11,
  September 1,1978, p.~28.

  中共浙江省委清查王张江姚材料组（Materials Group of the CCP Zhejiang
  Provincial Committee investigating Wang {[}Hongwen{]}, Zhang
  {[}Chunqiao{]}, Jiang {[}Qing{]} and Yao
  {[}Wenyuan{]}），《'四人帮'在浙江的亲信,
  现行反革命分子张永生的罪证材料》（Criminal Evidence concerning the
  ``Gang of Four's trusted follower in Zhejiang, the
  counter-revolutionary Zhang Yongsheng；下文简称 Zhang Yongshengde
  zuizheng cailiao），《宣传通讯》(Propaganda
  Bulletin)，第11期，1978年9月1日，第28页。} It is alleged that the Red
Guard leaders at Zhejiang University held little respect for Zhang
Yongsheng and personal antagonisms had played an important part in the
disintegration of the original Red Guard HQ.\footnote{Oral source.

  口述资料。}

到1966年秋，杭州大中学校红卫兵司令部已分裂为三个主要集团。\footnote{See
  n.~40.

  见注释40。}
第一个集团是红暴，以浙江大学为基地，由刘英和滕铸领导。井冈山是三个组织中影响最小的一个，居于中间地带，但在红暴同其主要对手红三司发生冲突时，往往倾向于站在红暴一边。红三司在杭州大学和浙江美术学院获得最强支持。其领导人包括杭州大学外语系学生李显通，以及分别来自浙江美术学院的青年教师杜英信和研究生张永生。张在1966年7月5日的一封信中把自己的学院描述为一个黑帮窝点，说他已在那里斗争了六年。他还抱怨自己不断受到控制学院的``魔鬼''的迫害。\footnote{中共浙江省委清查王张江姚材料组
  (Materials Group of the CCP Zhejiang Provincial Committee
  investigating Wang {[}Hongwen{]}, Zhang {[}Chunqiao{]}, Jiang
  {[}Qing{]} and Yao {[}Wenyuan{]}), ``'四人帮'在浙江的亲信,
  现行反革命分子张永生的罪证材料'' (Criminal Evidence concerning the
  ``Gang of Four's trusted follower in Zhejiang, the
  counter-revolutionary Zhang Yongsheng - hereafter Zhang Yongshengde
  zuizheng cailiao), 宣传通讯 (Propaganda Bulletin), No.~11, September
  1,1978, p.~28.

  中共浙江省委清查王张江姚材料组（Materials Group of the CCP Zhejiang
  Provincial Committee investigating Wang {[}Hongwen{]}, Zhang
  {[}Chunqiao{]}, Jiang {[}Qing{]} and Yao
  {[}Wenyuan{]}），《'四人帮'在浙江的亲信,
  现行反革命分子张永生的罪证材料》（Criminal Evidence concerning the
  ``Gang of Four's trusted follower in Zhejiang, the
  counter-revolutionary Zhang Yongsheng；下文简称 Zhang Yongshengde
  zuizheng cailiao），《宣传通讯》(Propaganda
  Bulletin)，第11期，1978年9月1日，第28页。}
据称，浙江大学的红卫兵领导人并不怎么尊重张永生，个人敌意在原红卫兵司令部分崩离析中发挥了重要作用。\footnote{Oral
  source.

  口述资料。}

It appears that shortly after its formation Red Storm established
contact with senior provincial leaders. For example, it was alleged in a
publication from a rival group at Zhejiang University that Jiang Hua's
deputy, Li Fengping, supported the organization and a Taiwan source
included Jiang Hua and Chen Weida as backers of Red Storm.\footnote{Kao
  Ch'ung-yen, 中共人事变动 1959-69 (Changes in Personnel in Communist
  China, 1959-69) (Hong Kong: Union Research Institute, 1970), p.~674.
  Facts \& Features, 1:13, p.~10.

  高崇焱，《中共人事变动 1959-69》(Changes in Personnel in Communist
  China, 1959-69)（香港：友联研究所，1970），第674页。Facts \& Features,
  1:13，第10页。} When a contingent from Red Third Headquarters marched
upon the famous Buddhist monastery and temple of Lingyin (灵隐寺),
students from Red Storm at neighboring Zhejiang University, acting on
Zhou Enlai's instructions, blocked the road to save the building from
destruction. Presumably, Zhou had directed the provincial leaders to
mobilize Red Storm knowing of their contacts and influence with the
organization. Zhejiang University student links with the provincial
party leadership may have dated back to June 1966, when the ZPC had
supported students there in their rebellion against the university
leaders.\footnote{This hypothesis is derived from Unger, Education Under
  Mao, p.~131.

  这一假设源自 Unger, Education Under Mao，第131页。}

红暴似乎在成立不久后就同省级高级领导人建立了联系。例如，浙江大学一个敌对团体的出版物声称，江华的副手李丰平支持该组织；台湾方面一项资料则把江华和陈伟达列为红暴的支持者。\footnote{Kao
  Ch'ung-yen, 中共人事变动 1959-69 (Changes in Personnel in Communist
  China, 1959-69) (Hong Kong: Union Research Institute, 1970), p.~674.
  Facts \& Features, 1:13, p.~10.

  高崇焱，《中共人事变动 1959-69》(Changes in Personnel in Communist
  China, 1959-69)（香港：友联研究所，1970），第674页。Facts \& Features,
  1:13，第10页。}
当红三司一支队伍向著名佛教寺院灵隐寺进发时，邻近浙江大学的红暴学生按照周恩来的指示拦住道路，使建筑免遭破坏。可以推测，周知道省领导人与红暴有联系并对其有影响，因而指示省领导动员红暴。浙江大学学生同省委领导层的联系，可能要追溯到1966年6月，当时省委曾支持该校学生反抗大学领导。\footnote{This
  hypothesis is derived from Unger, Education Under Mao, p.~131.

  这一假设源自 Unger, Education Under Mao，第131页。}

Red Third Headquarters, by its choice of name, clearly identified itself
with the most radical of the Beijing Red Guard organizations, Third
Headquarters. In Julianna Heaslet's categorization of students groups
Third Headquarters belonged to the category described retrospectively as
``Mao and Lin's political hatchetmen''.\footnote{Juliana Pennington
  Heaslet, ``The Red Guards: Instruments of Destruction in the Cultural
  Revolution'', AS, 12:12, December 1972, pp.~1033-39. See also Liu
  Guokai's categorization of secondary and tertiary students, in ``A
  Brief Analysis'', pp.~29-34.

  Juliana Pennington Heaslet, ``The Red Guards: Instruments of
  Destruction in the Cultural Revolution'', AS, 12:12,
  1972年12月，第1033-1039页。另见刘国凯在``A Brief
  Analysis''中对中学生和大学生的分类，第29-34页。} Its counterpart in
Zhejiang likewise belonged to this camp. Red Storm belonged to the
category that defended provincial leaders when the latter came under
attack late in 1966.\footnote{Hao Mengbi and Duan Haoran (eds.),
  中国共产党六十年 (Sixty Years of the Chinese Communist Party) Vol. 2,
  (Beijing: Jiefangjun Chubanshe, 1984), p.~587.

  郝梦笔、段浩然编，《中国共产党六十年》(Sixty Years of the Chinese
  Communist Party) 第2卷（北京：解放军出版社，1984），第587页。}
However, such a simplistic differentiation is misleading and inaccurate.
Both Red Third HQ and Red Storm, like their counterparts across China,
claimed to be loyal followers of Chairman Mao and went to extreme
lengths to prove their loyalty. To Red Storm and other organizations
later labeled ``conservative'', their provincial leaders were the local
spokesmen and representatives of the Chairman and the CCP and to
question their revolutionary credentials was tantamount to
counter-revolutionary behavior.

红三司通过其名称选择，清楚地表明自己认同北京红卫兵组织中最激进的第三司令部。在Julianna
Heaslet对学生团体的分类中，第三司令部属于后来被描述为``毛和林的政治打手''的类别。\footnote{Juliana
  Pennington Heaslet, ``The Red Guards: Instruments of Destruction in
  the Cultural Revolution'', AS, 12:12, December 1972, pp.~1033-39. See
  also Liu Guokai's categorization of secondary and tertiary students,
  in ``A Brief Analysis'', pp.~29-34.

  Juliana Pennington Heaslet, ``The Red Guards: Instruments of
  Destruction in the Cultural Revolution'', AS, 12:12,
  1972年12月，第1033-1039页。另见刘国凯在``A Brief
  Analysis''中对中学生和大学生的分类，第29-34页。}
浙江的对应组织同样属于这一阵营。红暴则属于在1966年底省领导人遭到攻击时为其辩护的类别。\footnote{Hao
  Mengbi and Duan Haoran (eds.), 中国共产党六十年 (Sixty Years of the
  Chinese Communist Party) Vol. 2, (Beijing: Jiefangjun Chubanshe,
  1984), p.~587.

  郝梦笔、段浩然编，《中国共产党六十年》(Sixty Years of the Chinese
  Communist Party) 第2卷（北京：解放军出版社，1984），第587页。}
然而，这种简单区分具有误导性且并不准确。红三司和红暴同全国各地的同类组织一样，都声称自己是毛主席的忠诚追随者，并不惜走向极端来证明忠诚。对红暴以及后来被贴上``保守派''标签的其他组织而言，省领导人是主席和中共在地方的发言人与代表，质疑他们的革命资格就等同于反革命行为。

\textbf{TABLE TWO: RED GUARD LEADERS}

\textbf{表二：红卫兵领导人}

{\def\LTcaptype{none} % do not increment counter
\begin{longtable}[]{@{}
  >{\raggedright\arraybackslash}p{(\linewidth - 4\tabcolsep) * \real{0.5000}}
  >{\raggedright\arraybackslash}p{(\linewidth - 4\tabcolsep) * \real{0.2857}}
  >{\raggedright\arraybackslash}p{(\linewidth - 4\tabcolsep) * \real{0.2143}}@{}}
\toprule\noalign{}
\begin{minipage}[b]{\linewidth}\raggedright
Organization
\end{minipage} & \begin{minipage}[b]{\linewidth}\raggedright
Leader
\end{minipage} & \begin{minipage}[b]{\linewidth}\raggedright
Unit
\end{minipage} \\
\midrule\noalign{}
\endhead
\bottomrule\noalign{}
\endlastfoot
Red Storm (红暴) & Liu Ying & Zhejiang University \\
Red Storm (红暴) & Teng Zhu & Zhejiang University \\
Red Third Headquarters (红三司) & Du Yingxin & Zhejiang College of Fine
Arts \\
Red Third Headquarters (红三司) & Zhang Yongsheng & Zhejiang College of
Fine Arts \\
Red Third Headquarters (红三司) & Li Xiantong & Hangzhou University \\
Jinggangshan (井冈山) & & \\
\end{longtable}
}

{\def\LTcaptype{none} % do not increment counter
\begin{longtable}[]{@{}lll@{}}
\toprule\noalign{}
组织 & 领导人 & 单位 \\
\midrule\noalign{}
\endhead
\bottomrule\noalign{}
\endlastfoot
红暴 & 刘英 & 浙江大学 \\
红暴 & 滕铸 & 浙江大学 \\
红三司 & 杜英信 & 浙江美术学院 \\
红三司 & 张永生 & 浙江美术学院 \\
红三司 & 李显通 & 杭州大学 \\
井冈山 & & \\
\end{longtable}
}

At about the same time as the split among Red Guard groups, Jiang Hua
had withdrawn from public duties, leaving his principal subordinate Li
Fengping in charge of the provincial party committee. By October 1966,
provincial party committees were virtually inoperative, with the number
of leading cadres making public appearances decreasing all the time. Li
Fengping was present at the 1966 National Day rally and at a struggle
rally against a leading provincial cadre on 5 January 1967.\footnote{ZJRB,
  October 2,1966, p.~3; ZJRB, January 13,1967, p.~3.

  ZJRB，1966年10月2日，第3页；ZJRB，1967年1月13日，第3页。} Thereafter,
he also disappeared from public view. Cao Xiangren and Shen Ce (沈策),
both members of the ZPC Standing Committee were present at a report
meeting on the study of Mao's works held in late October.\footnote{ZJRB,
  October 26,1966, p.~1.

  ZJRB，1966年10月26日，第1页。} But a Red Guard from Fujian, who was in
Hangzhou at this time, reported that ``All West Lake was now the world
of the Red Guards and the military''.\footnote{Ken Ling, The Revenge of
  Heaven (New York: Ballantine Books, 1972), p.~119.

  Ken Ling, The Revenge of Heaven (New York: Ballantine Books,
  1972)，第119页。} And it was true that the military was forced
increasingly to take on responsibility for civilian administration.

大约在红卫兵团体分裂的同时，江华已停止公开履职，由其主要下属李丰平负责省委工作。到1966年10月，省级党委事实上已无法运作，公开露面的领导干部人数不断减少。李丰平出席了1966年国庆集会，并在1967年1月5日一次针对省级重要干部的斗争大会上露面。\footnote{ZJRB,
  October 2,1966, p.~3; ZJRB, January 13,1967, p.~3.

  ZJRB，1966年10月2日，第3页；ZJRB，1967年1月13日，第3页。}
此后，他也从公众视野中消失。省委常委曹祥仁和沈策出席了10月下旬一次学习毛著的报告会。\footnote{ZJRB,
  October 26,1966, p.~1.

  ZJRB，1966年10月26日，第1页。}
但一名当时在杭州的福建红卫兵报告说，``整个西湖如今都是红卫兵和军队的天下''。\footnote{Ken
  Ling, The Revenge of Heaven (New York: Ballantine Books, 1972),
  p.~119.

  Ken Ling, The Revenge of Heaven (New York: Ballantine Books,
  1972)，第119页。}
实际情况也确实是，军队被迫越来越多地承担民政管理责任。

When the ZPMD convened a provincial congress of PLA activists in the
study of Mao Zedong Thought between 7 and 16 October\footnote{ZJRB,
  October 7,1966, p.~1 and October 17,1966, p.~1.

  ZJRB，1966年10月7日，第1页；1966年10月17日，第1页。} the principal
speakers were ZPC Secretary Wu Xian, ZPMD Commander Zhang Xiulong
(张秀龙) and Political Commissar Long Qian (龙潜).\footnote{Zhang's
  speech was published in ZJRB, October 8,1966, p.~1,2; and Long's on
  October 15,1966, pp.~1, 2.

  张的讲话刊于
  ZJRB，1966年10月8日，第1、2页；龙的讲话刊于1966年10月15日，第1、2页。}
It is highly likely that the above military officials were acting as
Jiang Hua's surrogates, as there was a close relationship between
provincial civil and military authorities at this time.\footnote{Harvey
  W. Nelsen, ``Military Forces in the Cultural Revolution'', CQ, 51
  (1972), pp.~448-54.

  Harvey W. Nelsen, ``Military Forces in the Cultural Revolution'', CQ,
  51 (1972), 第448-454页。} The close institutional and personnel
relationships between the civilian and military wings of the CCP would
prove a major obstacle to the rebels when they attempted to seize power
in 1967. The rebels themselves were acutely aware of the close relations
between local military and civilian party leaders. The issue was
discussed in a meeting of the Red Guards in December 1966. Zhang
Yongsheng from Red Third Headquarters pointed out that many of the
members of work teams sent onto campuses in June and July 1966 had come
from the military district. After discussions with Red Guards at
Zhejiang University, Han Xiangdong (韩向东), a student from Qinghua
University and a representative of Beijing's Third Headquarters, argued
that the ZPMD had close ties with the ZPC and that both had close
relations with the ``conservative faction'' (Red Storm). He concluded
that the key to ``solving the problem of the ZPC'' lay with the military
district. As Zhang Yongsheng pointed out in January 1967, his faction
would find its military protectors elsewhere in the shape of the 20th
Army and 5th Air Force.\footnote{Zhang Yongshengde zuizheng cailiao,
  pp.~30, 35.

  Zhang Yongshengde zuizheng cailiao，第30、35页。} As First Political
Commissar of the ZPMD, Jiang had strong institutional and personal ties
with his military colleagues, notwithstanding his derogatory remarks
concerning their capabilities.\footnote{Jiang was reported to have
  exclaimed in irritation, concerning the number of high-ranking
  officers who were garrisoned in the ZPMD, ``what lieutenant-generals
  and major-generals - they're just broad-bean sauce''
  (什么中将少将都是豆瓣酱) - a pun on the homophone jiang, meaning both
  general and sauce.

  据报道，江曾对驻扎在浙江省军区的高级军官人数恼怒地喊道：``什么中将少将都是豆瓣酱''------这是利用``将''和``酱''同音的双关。}
Political Commissar Long Qian also served as a member of the ZPC
Standing Committee and both he and Commander Zhang Xiulong had
accumulated a substantial period of service in Zhejiang. It was not only
at the top of the provincial hierarchy that the two bureaucracies
overlapped. At the district, county and commune levels, civilian and
military leadership coexisted on intimate terms, an administrative
arrangement of momentous importance when the rebels left the relative
safety of provincial capitals to unseat ``local emperors'' in towns and
villages.

10月7日至16日，浙江省军区召开全省解放军学习毛泽东思想积极分子代表大会，\footnote{ZJRB,
  October 7,1966, p.~1 and October 17,1966, p.~1.

  ZJRB，1966年10月7日，第1页；1966年10月17日，第1页。}
主要发言者是省委书记吴宪、省军区司令员张秀龙和政治委员龙潜。\footnote{Zhang's
  speech was published in ZJRB, October 8,1966, p.~1,2; and Long's on
  October 15,1966, pp.~1, 2.

  张的讲话刊于
  ZJRB，1966年10月8日，第1、2页；龙的讲话刊于1966年10月15日，第1、2页。}
这些军方官员很可能是在充当江华的代理人，因为当时省内党政与军方之间关系密切。\footnote{Harvey
  W. Nelsen, ``Military Forces in the Cultural Revolution'', CQ, 51
  (1972), pp.~448-54.

  Harvey W. Nelsen, ``Military Forces in the Cultural Revolution'', CQ,
  51 (1972), 第448-454页。}
中共文职和军事两翼之间紧密的制度和人事关系，将在1967年造反派试图夺权时成为重大障碍。造反派自己也敏锐意识到地方军方与党政领导之间的密切关系。1966年12月一次红卫兵会议讨论了这一问题。红三司的张永生指出，1966年6月和7月派往校园的工作组中，许多成员来自军区。清华大学学生、北京第三司令部代表韩向东在同浙江大学红卫兵讨论后认为，浙江省军区同省委关系密切，二者又都同``保守派''（红暴）关系密切。他总结说，``解决省委问题''的关键在于军区。正如张永生在1967年1月所指出的，他的派别将在别处找到军事保护者，即第二十军和第五空军。\footnote{Zhang
  Yongshengde zuizheng cailiao, pp.~30, 35.

  Zhang Yongshengde zuizheng cailiao，第30、35页。}
作为浙江省军区第一政治委员，江同军方同事有牢固的制度和个人联系，尽管他曾贬低他们的能力。\footnote{Jiang
  was reported to have exclaimed in irritation, concerning the number of
  high-ranking officers who were garrisoned in the ZPMD, ``what
  lieutenant-generals and major-generals - they're just broad-bean
  sauce'' (什么中将少将都是豆瓣酱) - a pun on the homophone jiang,
  meaning both general and sauce.

  据报道，江曾对驻扎在浙江省军区的高级军官人数恼怒地喊道：``什么中将少将都是豆瓣酱''------这是利用``将''和``酱''同音的双关。}
政委龙潜同时也是省委常委，他和司令员张秀龙都在浙江任职相当长时间。两套官僚体系的重叠并不仅限于省级层级顶端。在专区、县和公社层面，文职和军事领导也以密切方式共存；当造反派离开省会相对安全的环境，前往城镇和乡村推翻``土皇帝''时，这种行政安排具有重大意义。

Nevertheless, Jiang's political position and authority were being
rapidly undermined. The October 1966 Central Work Conference in Beijing
decided to press ahead with the Cultural Revolution and to expand its
scope from the educational sphere into offices, factories and communes.
Jiang Hua continued to fight for his political life. A ``Zhejiang
delegation of veteran conservatives'' arrived in the national capital in
October to report that the situation in the province was similar to that
prevailing on the eve of the anti-rightist movement in 1957. Tan
Zhenlin's wife, Ge Huimin, reportedly relayed her husband's instructions
to the delegation.\footnote{ZPS, June 28,1968, SWB/FE/2814/B/17-18.

  ZPS，1968年6月28日，SWB/FE/2814/B/17-18。} Tan also interceded with
the central authorities to argue for Jiang Hua's continuation in
power.\footnote{ZPS, June 3,1968, SWB/FE/2790/B/18-19; ZPS, May 30,1968,
  SWB/FE/2786/B/10-11; ZPS, November 10,1968, SWB/FE/2924/B/24-5; ZPS,
  November 19, 1968, SWB/FE/2933/B/3-4; speech of Jiang Qing, March
  21,1968, SCMP, 4166 (April 29,1968), p.~2.

  ZPS，1968年6月3日，SWB/FE/2790/B/18-19；ZPS，1968年5月30日，SWB/FE/2786/B/10-11；ZPS，1968年11月10日，SWB/FE/2924/B/24-5；ZPS，1968年11月19日，SWB/FE/2933/B/3-4；江青1968年3月21日讲话，SCMP,
  4166（1968年4月29日），第2页。}

尽管如此，江的政治地位和权威正被迅速削弱。1966年10月在北京召开的中央工作会议决定推进文化大革命，并把其范围从教育领域扩大到机关、工厂和公社。江华继续为自己的政治生命而斗争。一个``浙江老保守派代表团''于10月抵达首都，报告说浙江局势类似于1957年反右运动前夕的局面。据称，谭震林的妻子葛慧敏向代表团转达了丈夫的指示。\footnote{ZPS,
  June 28,1968, SWB/FE/2814/B/17-18.

  ZPS，1968年6月28日，SWB/FE/2814/B/17-18。}
谭也向中央当局斡旋，主张江华继续掌权。\footnote{ZPS, June 3,1968,
  SWB/FE/2790/B/18-19; ZPS, May 30,1968, SWB/FE/2786/B/10-11; ZPS,
  November 10,1968, SWB/FE/2924/B/24-5; ZPS, November 19, 1968,
  SWB/FE/2933/B/3-4; speech of Jiang Qing, March 21,1968, SCMP, 4166
  (April 29,1968), p.~2.

  ZPS，1968年6月3日，SWB/FE/2790/B/18-19；ZPS，1968年5月30日，SWB/FE/2786/B/10-11；ZPS，1968年11月10日，SWB/FE/2924/B/24-5；ZPS，1968年11月19日，SWB/FE/2933/B/3-4；江青1968年3月21日讲话，SCMP,
  4166（1968年4月29日），第2页。}

\section{Power Seizures}\label{power-seizures}

\section{夺权}\label{ux593aux6743}

Rebel delegations of workers also traveled to Beijing to appeal for
support in their efforts to overthrow local power-holders. On 17
November 1966 the country's first worker-delegation was received by two
members of the Cultural Revolution Group, Yao Wenyuan and Wang
Li.\footnote{审查翁森鹤专案小组 (Special Case Group Investigating Weng
  Senhe), 现行反革命分子翁森鹤部分罪证材料 (Part of the Criminal
  Evidence concerning the Counter-revolutionary Weng Senhe - hereafter
  Weng Senhe bufen zuizheng cailiao), 宣传通讯 (Propaganda Bulletin),
  No.~15, October 1977, p.~3. 无产阶级文化大革命部分资料汇编 (A
  Compilation of Material on the Great Proletarian Cultural Revolution),
  Vol. 2, (November-December 1966), (Hangzhou: Zhejiang Teachers'
  College Cultural Revolution Preparatory Committee and Liaison Station
  in Taizhou, n.~d.), pp.~70-1. On October 29,1966, Red Guards from
  Jinhua District in central Zhejiang appealed to Zhang Chunqiao to
  endorse their attempt to bring down the district party secretary. See
  Wuchanjieji wenhua dageming bufen ziliao huibian, Vol. 1 (July-October
  1966), pp.~95-8.

  审查翁森鹤专案小组（Special Case Group Investigating Weng
  Senhe），《现行反革命分子翁森鹤部分罪证材料》(Part of the Criminal
  Evidence concerning the Counter-revolutionary Weng Senhe；下文简称
  Weng Senhe bufen zuizheng cailiao)，《宣传通讯》(Propaganda
  Bulletin)，第15期，1977年10月，第3页。《无产阶级文化大革命部分资料汇编》(A
  Compilation of Material on the Great Proletarian Cultural
  Revolution)，第2卷（1966年11-12月）（杭州：浙江师范学院文化革命筹备委员会及驻台州联络站，无日期），第70-71页。1966年10月29日，浙江中部金华地区红卫兵向张春桥求助，请其支持他们打倒地委书记的尝试。见
  Wuchanjieji wenhua dageming bufen ziliao
  huibian，第1卷（1966年7-10月），第95-98页。} It comprised employees
from the city's largest silk mill, the Hangzhou Silk Printing and Dyeing
Complex (杭州丝绸印染联合厂).\footnote{Construction of this mill began
  in 1956 and it commenced production in 1960. In 1977 it employed 4,700
  workers in three workshops, operating 536 reeling machines and 623
  looms. Visit, January 12, 1977.

  该厂于1956年开始建设，1960年投产。1977年，它在三个车间雇有4700名工人，运行536台缫丝机和623台织机。访问，1977年1月12日。}
The leader of the delegation was Weng Senhe, who was to become the most
notorious rebel from Hangzhou in the Cultural Revolution and who later
stood second in fame nationally only to Wang Hongwen.\footnote{See Keith
  Forster, ``Weng Senhe - The Cultural Revolution Rebel from Zhejiang''
  (unpublished manuscript).

  见 Keith Forster, ``Weng Senhe - The Cultural Revolution Rebel from
  Zhejiang''（未发表手稿）。} On 16 January Weng's rebel group the Red
Rebel Corps (红色造反兵团), had seized power at the silk complex. As
early as June 1966 Weng, a 30 year-old screen printer, had put up a big
character poster rebelling against the factory leadership.\footnote{Weng's
  act of rebellion occurred three days before Wang Hongwen's, which
  would date it at June 9, 1966. See, Parris H. Chang, ``Political
  Profiles: Wang Hung-wen and Li Teh-sheng'', China Quarterly (CQ), 57
  (1974), p.~125. For the similar experiences of a famous Shanxi rebel,
  the journeyman worker at the Taiyuan Steel Plant Yang Chengxiao, see
  Hinton, Shenfan, p.~592.

  翁的造反行动发生在王洪文之前三天，也就是1966年6月9日。见 Parris H.
  Chang, ``Political Profiles: Wang Hung-wen and Li Teh-sheng'', China
  Quarterly (CQ), 57 (1974),
  第125页。关于山西著名造反派、太原钢铁厂熟练工杨成孝的类似经历，见
  Hinton, Shenfan，第592页。} Weng, who had received three years'
secondary schooling, had originally been a clerk (文书) at a township
office in Zhejiang.\footnote{Oral source.

  口述资料。} For unknown reasons he was the beneficiary of the much
sought-after transfer to a large city. In 1962 Weng was assigned to work
at the Hangzhou Chemical Fibre Plant\footnote{Visit to Hangzhou Chemical
  Fibre Plant, January 27,1978.

  访问杭州化纤厂，1978年1月27日。} and then, in 1963, he was transferred
to the silk complex. In 1965, during the Socialist Education Movement,
he was reportedly reprimanded for speculative activities.\footnote{Peking
  Review (PR), (September 30,1977), p.~24; and oral source.

  Peking Review (PR)，1977年9月30日，第24页；以及口述资料。} This
experience gave rise to feelings of resentment against the factory
leadership and the political and social system in general. Weng
reportedly regarded a worker's existence as meaningless\footnote{Shen
  Chuyun, 清除``四害''杭丝联又大步前进了 (After clearing out the ``four
  pests'' the Hangzhou Silk Complex is again advancing with great
  strides), RMRB, May 17,1977, p.~2.

  沈初云，《清除``四害''杭丝联又大步前进了》(After clearing out the
  ``four pests'' the Hangzhou Silk Complex is again advancing with great
  strides)，RMRB，1977年5月17日，第2页。} and spent his leisure hours
studying the careers of such twentieth century notables as Hitler and
Mussolini.\footnote{Visit to Hangzhou Silk Complex, January 12,1977.
  Weng was apparently not alone among young Chinese in his fascination
  with such subjects. See Simon Leys, The Burning Forest (New York:
  Holt, Rinehart \& Winston, 1985), p.117, fn.

  访问杭州丝绸联合厂，1977年1月12日。翁显然并非对这些主题着迷的中国青年中的孤例。见
  Simon Leys, The Burning Forest (New York: Holt, Rinehart \& Winston,
  1985)，第117页注。} When the Cultural Revolution commenced Weng
reportedly proclaimed that ``the reds of the past seventeen years must
stand aside; those who have been oppressed, unite and
rebel''.\footnote{HZRB, December 14,1977.

  HZRB，1977年12月14日。} Weng's June rebellion took considerable time
to gather momentum. At the beginning of December 1966 only 45 workers in
the complex had committed themselves to his cause. Workers of good class
background and those on higher pay and with better skills tended to
support the factory leadership. Weng could only attract those of lower
status in the socialist class pecking order. The party leadership in
turn maligned Weng's group as rightists and, to protect its position,
suddenly admitted more workers into the CCP and trade union.\footnote{China
  Reconstructs (July 1977), p.~40; Shen Chuyun, RMRB, May 17,1977.

  China Reconstructs（1977年7月），第40页；沈初云，RMRB，1977年5月17日。}
However, by the middle of December 1966 the rebels had attracted enough
workers to hold their first meeting and, by the time they convened their
second meeting, the party leaders felt sufficiently threatened to
disrupt it.

工人造反派代表团也前往北京，请求支持他们推翻地方掌权者。1966年11月17日，全国第一个工人代表团受到文化革命小组两名成员姚文元和王力接见。\footnote{审查翁森鹤专案小组
  (Special Case Group Investigating Weng Senhe),
  现行反革命分子翁森鹤部分罪证材料 (Part of the Criminal Evidence
  concerning the Counter-revolutionary Weng Senhe - hereafter Weng Senhe
  bufen zuizheng cailiao), 宣传通讯 (Propaganda Bulletin), No.~15,
  October 1977, p.~3. 无产阶级文化大革命部分资料汇编 (A Compilation of
  Material on the Great Proletarian Cultural Revolution), Vol. 2,
  (November-December 1966), (Hangzhou: Zhejiang Teachers' College
  Cultural Revolution Preparatory Committee and Liaison Station in
  Taizhou, n.~d.), pp.~70-1. On October 29,1966, Red Guards from Jinhua
  District in central Zhejiang appealed to Zhang Chunqiao to endorse
  their attempt to bring down the district party secretary. See
  Wuchanjieji wenhua dageming bufen ziliao huibian, Vol. 1 (July-October
  1966), pp.~95-8.

  审查翁森鹤专案小组（Special Case Group Investigating Weng
  Senhe），《现行反革命分子翁森鹤部分罪证材料》(Part of the Criminal
  Evidence concerning the Counter-revolutionary Weng Senhe；下文简称
  Weng Senhe bufen zuizheng cailiao)，《宣传通讯》(Propaganda
  Bulletin)，第15期，1977年10月，第3页。《无产阶级文化大革命部分资料汇编》(A
  Compilation of Material on the Great Proletarian Cultural
  Revolution)，第2卷（1966年11-12月）（杭州：浙江师范学院文化革命筹备委员会及驻台州联络站，无日期），第70-71页。1966年10月29日，浙江中部金华地区红卫兵向张春桥求助，请其支持他们打倒地委书记的尝试。见
  Wuchanjieji wenhua dageming bufen ziliao
  huibian，第1卷（1966年7-10月），第95-98页。}
该代表团由全市最大丝绸厂杭州丝绸印染联合厂的职工组成。\footnote{Construction
  of this mill began in 1956 and it commenced production in 1960. In
  1977 it employed 4,700 workers in three workshops, operating 536
  reeling machines and 623 looms. Visit, January 12, 1977.

  该厂于1956年开始建设，1960年投产。1977年，它在三个车间雇有4700名工人，运行536台缫丝机和623台织机。访问，1977年1月12日。}
代表团领袖是翁森鹤，他后来成为杭州文革中最声名狼藉的造反派，全国知名度仅次于王洪文。\footnote{See
  Keith Forster, ``Weng Senhe - The Cultural Revolution Rebel from
  Zhejiang'' (unpublished manuscript).

  见 Keith Forster, ``Weng Senhe - The Cultural Revolution Rebel from
  Zhejiang''（未发表手稿）。}
1月16日，翁的造反组织红色造反兵团在该丝绸联合厂夺权。早在1966年6月，30岁的印花工翁就贴出一张大字报，反抗工厂领导。\footnote{Weng's
  act of rebellion occurred three days before Wang Hongwen's, which
  would date it at June 9, 1966. See, Parris H. Chang, ``Political
  Profiles: Wang Hung-wen and Li Teh-sheng'', China Quarterly (CQ), 57
  (1974), p.~125. For the similar experiences of a famous Shanxi rebel,
  the journeyman worker at the Taiyuan Steel Plant Yang Chengxiao, see
  Hinton, Shenfan, p.~592.

  翁的造反行动发生在王洪文之前三天，也就是1966年6月9日。见 Parris H.
  Chang, ``Political Profiles: Wang Hung-wen and Li Teh-sheng'', China
  Quarterly (CQ), 57 (1974),
  第125页。关于山西著名造反派、太原钢铁厂熟练工杨成孝的类似经历，见
  Hinton, Shenfan，第592页。}
翁受过三年中等教育，原本是浙江某乡镇办公室文书。\footnote{Oral source.

  口述资料。}
不知何故，他得到了人人渴望的调往大城市的机会。1962年，翁被分配到杭州化纤厂工作，\footnote{Visit
  to Hangzhou Chemical Fibre Plant, January 27,1978.

  访问杭州化纤厂，1978年1月27日。}
1963年又调到丝绸联合厂。1965年社会主义教育运动期间，据称他因投机活动受到批评。\footnote{Peking
  Review (PR), (September 30,1977), p.~24; and oral source.

  Peking Review (PR)，1977年9月30日，第24页；以及口述资料。}
这一经历使他对工厂领导以及整个政治社会制度产生怨恨。据称，翁认为工人的存在毫无意义，\footnote{Shen
  Chuyun, 清除``四害''杭丝联又大步前进了 (After clearing out the ``four
  pests'' the Hangzhou Silk Complex is again advancing with great
  strides), RMRB, May 17,1977, p.~2.

  沈初云，《清除``四害''杭丝联又大步前进了》(After clearing out the
  ``four pests'' the Hangzhou Silk Complex is again advancing with great
  strides)，RMRB，1977年5月17日，第2页。}
并在闲暇时间研究希特勒和墨索里尼等20世纪著名人物的经历。\footnote{Visit
  to Hangzhou Silk Complex, January 12,1977. Weng was apparently not
  alone among young Chinese in his fascination with such subjects. See
  Simon Leys, The Burning Forest (New York: Holt, Rinehart \& Winston,
  1985), p.117, fn.

  访问杭州丝绸联合厂，1977年1月12日。翁显然并非对这些主题着迷的中国青年中的孤例。见
  Simon Leys, The Burning Forest (New York: Holt, Rinehart \& Winston,
  1985)，第117页注。}
文化大革命开始时，据称翁宣称：``十七年来的红人要靠边站；受压迫的人，联合起来造反。''\footnote{HZRB,
  December 14,1977.

  HZRB，1977年12月14日。}
翁在6月的造反花了相当长时间才形成声势。到1966年12月初，联合厂中只有45名工人投身他的事业。阶级出身好、工资较高且技术较好的工人倾向于支持工厂领导。翁只能吸引社会主义阶级排序中地位较低者。党内领导反过来把翁的组织诬称为右派，并为保护自身地位，突然吸收更多工人加入中共和工会。\footnote{China
  Reconstructs (July 1977), p.~40; Shen Chuyun, RMRB, May 17,1977.

  China Reconstructs（1977年7月），第40页；沈初云，RMRB，1977年5月17日。}
然而，到1966年12月中旬，造反派已经吸引了足够多工人召开第一次会议；等到他们召开第二次会议时，党内领导感到威胁足够大，便出面破坏会议。

Eventually, the authorities were overwhelmed by the rebels who,
benefiting from Weng's undoubted leadership skills, finally seized
power. They were also assisted by the unprecedented nature of the
political campaign\footnote{A point commented on by Zhao Ziyang in a
  speech of June 26,1966. See Zhao Wei, Zhao Ziyang zhuan, pp.~131-32.

  赵紫阳在1966年6月26日的讲话中评论过这一点。见赵蔚，《赵紫阳传》，第131-132页。}
and the kudos that daredevils such as Weng accrued for the persecution
they had suffered several months before. Previously labeled an
``anti-party'' trouble-maker, Weng was vindicated and praised for his
actions. Ironically however, the rebels' first task on taking over the
seals of office was to try and resume normal production -- a goal that
the erstwhile management had been striving unsuccessfully to accomplish.
The rebels then issued the call popularized by Zhou Enlai to ``grasp
revolution and promote production''.\footnote{ZPS, January 26, 1967,
  SWB/FE/2381/B/33-4.

  ZPS，1967年1月26日，SWB/FE/2381/B/33-4。}

最终，当局被造反派压倒；造反派凭借翁无可置疑的领导才能，最终夺取了权力。他们还受益于这场政治运动的空前性质，\footnote{A
  point commented on by Zhao Ziyang in a speech of June 26,1966. See
  Zhao Wei, Zhao Ziyang zhuan, pp.~131-32.

  赵紫阳在1966年6月26日的讲话中评论过这一点。见赵蔚，《赵紫阳传》，第131-132页。}
以及翁这样胆大妄为者因几个月前遭受迫害而获得的声望。此前被贴上``反党''捣乱分子标签的翁，其行动得到了平反和赞扬。然而具有讽刺意味的是，造反派接管公章后的第一项任务，是试图恢复正常生产，而这正是原管理层一直努力却未能实现的目标。随后，造反派发出了由周恩来推广的``抓革命、促生产''号召。\footnote{ZPS,
  January 26, 1967, SWB/FE/2381/B/33-4.

  ZPS，1967年1月26日，SWB/FE/2381/B/33-4。}

On 9 January 1967 one of the most celebrated (or infamous) events of the
Cultural Revolution occurred with the overthrow of the party and
government leadership in Shanghai. The ``January Storm'' (一月风暴), as
this month of power seizures became known, rapidly spread to other parts
of China.\footnote{See Liu Guokai, ``A Brief Analysis'', pp.~50-7.

  见刘国凯，``A Brief Analysis''，第50-57页。} Situated in close
proximity to the giant metropolis, Zhejiang immediately felt the waves
of this extraordinary upheaval. Over the previous two months, workers at
industrial plants and employees in government offices and various trades
had established their own unit-based groups. The state of social
disorder had provided them with the unprecedented opportunity to visit
other units for the purpose of seeking out and making contacts with
like-minded organizations. Such linkages were in turn utilized to
bolster each organization's own position and were drawn on to fight off
challenges from opposing groups. Cross-unit ties led to the formation of
umbrella organizations spanning schools, colleges, factories and
offices.

1967年1月9日，文化大革命中最著名（或最臭名昭著）的事件之一发生了：上海党政领导被推翻。这个夺权之月后来被称为``一月风暴''，并迅速扩散到中国其他地区。\footnote{See
  Liu Guokai, ``A Brief Analysis'', pp.~50-7.

  见刘国凯，``A Brief Analysis''，第50-57页。}
浙江毗邻这个巨大都市，立即感受到这场非凡动荡的冲击波。在此前两个月里，工业企业工人、政府机关和各行业职员已建立起各自以单位为基础的团体。社会失序状态给他们提供了前所未有的机会，使他们能够访问其他单位，寻找并接触志同道合的组织。这些联系又被用来巩固各组织自身地位，并用于抵御敌对团体的挑战。跨单位联系导致了涵盖学校、大学、工厂和机关的伞状组织形成。

On 9 November the Shanghai Workers' Revolutionary Rebel Headquarters had
been formed under the leadership of Wang Hongwen, a security officer at
the city's No.~17 cotton mill. A similar process occurred in Zhejiang.
In fact it is claimed that Weng Senhe was the first worker rebel leader
in the Cultural Revolution to organize a General Headquarters (总司令部)
of worker rebels.\footnote{Hao Mengbi and Duan Haoran, Liushinian,
  p.~589; Ye Yonglie, Wang Hongwen xingshuailu, p.~77.

  郝梦笔、段浩然，Liushinian，第589页；叶永烈，Wang Hongwen
  xingshuailu，第77页。} By the end of 1966 Red Storm had expanded from
its base at Zhejiang University to encompass a large number of factories
in the light industrial textile district of Gongshu (拱墅) in the north
of Hangzhou. This organization took the name of Zhejiang Provincial Red
Storm Provisional Headquarters and it is alleged that the leaders of the
CCP ZPC were responsible for its establishment.\footnote{See Kao,
  Zhonggong renshi biandong, p.~674.

  见高，Zhonggong renshi biandong，第674页。} A rival organization,
whose student component was supplied by Red Third Headquarters, included
workers from heavy industrial plants such as the Hangzhou Iron and Steel
Mill (杭州钢铁厂) and the Hangzhou Glass Factory situated further to the
north of Hangzhou in the industrial district of Banshan (半山). This
organization, which was formally established on 30 December
1966\footnote{ZJRB, August 24,1969, p.~1.

  ZJRB，1969年8月24日，第1页。} as an alliance of twenty-two rebel
organizations, named itself the Zhejiang Provincial Revolutionary Rebel
United Headquarters. Each factory and college sent a delegate to the
parent organization. To proclaim its formation, Zhang Yongsheng led a
contingent to storm the downtown building housing Zhejiang Daily and
closed down the newspaper, which did not reappear until 10 January
1967.\footnote{ZJRB, January 10,1967, pp.~1, 2; ZJRB, January 11,1967,
  p.~1; Zhejiang shengqing, p.~1052. On January 14,1967 rebels in Jiande
  county in south-west Zhejiang occupied the local broadcasting station.
  See 建德县志编纂委员会 (Jiande County Gazetteer Compilation
  Committee), 建德县志 (Jiande County Gazetteer), (Hangzhou: Zhejiang
  Renmin Chubanshe, 1986), p.~22. Rebels in Jiangsu took over the
  provincial newspaper, New China Daily (新华日报) and the provincial
  and Nanjing municipal television stations on January 9, 1967. See
  Jiangsusheng dashiji, p.~285.

  ZJRB，1967年1月10日，第1、2页；ZJRB，1967年1月11日，第1页；Zhejiang
  shengqing，第1052页。1967年1月14日，浙西南建德县造反派占领地方广播站。见建德县志编纂委员会，《建德县志》(Jiande
  County
  Gazetteer)（杭州：浙江人民出版社，1986），第22页。江苏造反派于1967年1月9日接管省报《新华日报》以及省和南京市电视台。见
  Jiangsusheng dashiji，第285页。}

11月9日，在上海第十七棉纺厂保卫干部王洪文领导下，上海工人革命造反总司令部成立。浙江也发生了类似过程。事实上，有人称翁森鹤是文化大革命中第一个组织工人造反派总司令部的工人造反领袖。\footnote{Hao
  Mengbi and Duan Haoran, Liushinian, p.~589; Ye Yonglie, Wang Hongwen
  xingshuailu, p.~77.

  郝梦笔、段浩然，Liushinian，第589页；叶永烈，Wang Hongwen
  xingshuailu，第77页。}
到1966年底，红暴已从浙江大学基地扩展到杭州北部拱墅轻工业纺织区的大量工厂。该组织取名为浙江省红暴临时总指挥部，据称中共浙江省委领导人对其成立负有责任。\footnote{See
  Kao, Zhonggong renshi biandong, p.~674.

  见高，Zhonggong renshi biandong，第674页。}
一个敌对组织的学生部分由红三司提供，成员包括来自杭州更北部半山工业区的杭州钢铁厂、杭州玻璃厂等重工业企业的工人。该组织由二十二个造反组织结盟，于1966年12月30日正式成立，\footnote{ZJRB,
  August 24,1969, p.~1.

  ZJRB，1969年8月24日，第1页。}
自称浙江省革命造反联合总指挥部。每个工厂和学院都向上级组织派出一名代表。为了宣告其成立，张永生率领一支队伍冲击市中心《浙江日报》所在大楼，并使该报停刊，直到1967年1月10日才复刊。\footnote{ZJRB,
  January 10,1967, pp.~1, 2; ZJRB, January 11,1967, p.~1; Zhejiang
  shengqing, p.~1052. On January 14,1967 rebels in Jiande county in
  south-west Zhejiang occupied the local broadcasting station. See
  建德县志编纂委员会 (Jiande County Gazetteer Compilation Committee),
  建德县志 (Jiande County Gazetteer), (Hangzhou: Zhejiang Renmin
  Chubanshe, 1986), p.~22. Rebels in Jiangsu took over the provincial
  newspaper, New China Daily (新华日报) and the provincial and Nanjing
  municipal television stations on January 9, 1967. See Jiangsusheng
  dashiji, p.~285.

  ZJRB，1967年1月10日，第1、2页；ZJRB，1967年1月11日，第1页；Zhejiang
  shengqing，第1052页。1967年1月14日，浙西南建德县造反派占领地方广播站。见建德县志编纂委员会，《建德县志》(Jiande
  County
  Gazetteer)（杭州：浙江人民出版社，1986），第22页。江苏造反派于1967年1月9日接管省报《新华日报》以及省和南京市电视台。见
  Jiangsusheng dashiji，第285页。}

The leaders of United Headquarters were the students Zhang Yongsheng and
Li Xiantong, already referred to as joint supremos of Red Third
Headquarters. Because of a socially less acceptable class background and
because of his unorthodox social behavior Du Yingxin, although he was
reportedly more capable than Zhang, did not achieve the prominence his
abilities would have otherwise warranted.\footnote{Oral source.

  口述资料。} Worker rebels who were prominent in the organization were
He Xianchun from the Hangzhou Heavy Machinery Plant, Xia Genfa (夏根法)
from the Hangzhou Oxygen Generator Plant (杭州制氧机厂), Guo Zhisong
(郭志松) from the Zhejiang Construction Company and Ye Rende (叶仁德)
from the iron and steel mill.\footnote{The rebels, presumably United
  Headquarters, seized power in this mill on January 16,1967. See ZJRB,
  February 15,1967, p.~2.

  造反派，推测为联合总部，于1967年1月16日在该厂夺权。见
  ZJRB，1967年2月15日，第2页。} Red Storm's leadership included the two
student leaders Liu Ying and Teng Zhu from Zhejiang University and the
workers Weng Senhe and Huang Yintang (黄阴堂) from the Hangzhou silk
complex\footnote{Rebels from Red Storm seized power here also on January
  16. ZPS, January 26,1967, SWB/FE/2381/B/33-4.

  红暴造反派也于1月16日在此夺权。ZPS，1967年1月26日，SWB/FE/2381/B/33-4。}
and Fang Jianwen (方剑文) from the Hangzhou No.~1 Cotton Textile and
Dyeing Mill (杭州第一棉纺织印染厂). What is striking about Red Storm's
full name is the word ``provisional'' as well as the absence of the
virtually obligatory ``rebel'' qualifier. In its composition and outlook
Red Storm somewhat resembled Liu Guokai's description of organizations
comprising members of ``high political quality''. By contrast, United
Headquarters seems to have contained more members of ``low political
quality'' who were more extreme in their actions and designated
themselves ``revolutionary rebels''.\footnote{Liu Guokai, ``A Brief
  Analysis'', pp.~43-5.

  刘国凯，``A Brief Analysis''，第43-45页。} Yet the distinction was
never so neat. Whether the fanaticism of both groups reflected the
composition of their memberships and socio-economic circumstances
specific to Hangzhou, or both factors, is unclear. It is clear, however,
that neither United Headquarters nor Red Storm placed workers' economic
grievances high on their list of priorities although they may have
supported such demands when it suited them. Action in support of such
demands was condemned as a manifestation of economism, or what Walder
has aptly described as ``interest-group radicalism''.\footnote{Andrew
  Walder, ``Cultural Revolution Radicalism: An Interpretation and
  Typology'' (paper delivered at the conference ``New Perspectives on
  the Cultural Revolution'', Harvard University, May 15-17,1987). For a
  description of the workers' actions and their suppression, see Liu
  Guokai, ``A Brief Analysis'', pp.~45-49.

  Andrew Walder, ``Cultural Revolution Radicalism: An Interpretation and
  Typology''（提交于``New Perspectives on the Cultural
  Revolution''会议，哈佛大学，1987年5月15-17日）。关于工人行动及其遭镇压的描述，见刘国凯，``A
  Brief Analysis''，第45-49页。} The significance of the geographic and
sectional division between the two rebel organizations is also difficult
to assess. It may not have been as clear-cut as the above description
suggests. At first glance Red Storm seemed to have built up a more
substantial base in a city where light and textile industries were
relatively developed. But women comprised a large proportion of the
employees in this sector and in the Cultural Revolution the bulk of the
fighting was undertaken by men. The sheer weight of numbers on the
streets often counted a great deal. Additionally, there may have been
more temporary and contract workers in the heavy industrial sector,
another factor which worked to United Headquarters' advantage. Such
speculation may go some of the way to understanding the relative
strength of the two organizations.

联合总部的领导人是学生张永生和李显通，前文已称他们为红三司的共同首脑。由于阶级出身在社会上不那么容易被接受，又因社会行为不合常规，杜英信尽管据称比张更有能力，却没有获得其能力本应带来的显赫地位。\footnote{Oral
  source.

  口述资料。}
该组织中突出的工人造反派包括杭州重型机器厂的贺贤春、杭州制氧机厂的夏根法、浙江建筑公司的郭志松和钢铁厂的叶仁德。\footnote{The
  rebels, presumably United Headquarters, seized power in this mill on
  January 16,1967. See ZJRB, February 15,1967, p.~2.

  造反派，推测为联合总部，于1967年1月16日在该厂夺权。见
  ZJRB，1967年2月15日，第2页。}
红暴领导层包括来自浙江大学的两名学生领袖刘英、滕铸，以及杭州丝绸联合厂的工人翁森鹤、黄阴堂\footnote{Rebels
  from Red Storm seized power here also on January 16. ZPS, January
  26,1967, SWB/FE/2381/B/33-4.

  红暴造反派也于1月16日在此夺权。ZPS，1967年1月26日，SWB/FE/2381/B/33-4。}
和杭州第一棉纺织印染厂的方剑文。红暴全名中引人注意的是``临时''一词，以及几乎必不可少的``造反''限定语的缺席。就构成和取向而言，红暴有些类似刘国凯所描述的由``政治素质高''成员组成的组织。相比之下，联合总部似乎包含更多``政治素质低''的成员，他们行动更极端，并自称``革命造反派''。\footnote{Liu
  Guokai, ``A Brief Analysis'', pp.~43-5.

  刘国凯，``A Brief Analysis''，第43-45页。}
然而这种区分从来没有如此整齐。两个组织的狂热究竟反映了其成员构成、杭州特有的社会经济环境，还是两者共同作用，并不清楚。不过可以明确的是，联合总部和红暴都没有把工人的经济不满置于优先事项的高位，尽管在对自己有利时它们可能支持此类要求。支持这些要求的行动被谴责为经济主义表现，或沃尔德贴切称之为``利益集团激进主义''。\footnote{Andrew
  Walder, ``Cultural Revolution Radicalism: An Interpretation and
  Typology'' (paper delivered at the conference ``New Perspectives on
  the Cultural Revolution'', Harvard University, May 15-17,1987). For a
  description of the workers' actions and their suppression, see Liu
  Guokai, ``A Brief Analysis'', pp.~45-49.

  Andrew Walder, ``Cultural Revolution Radicalism: An Interpretation and
  Typology''（提交于``New Perspectives on the Cultural
  Revolution''会议，哈佛大学，1987年5月15-17日）。关于工人行动及其遭镇压的描述，见刘国凯，``A
  Brief Analysis''，第45-49页。}
两个造反组织之间地理和部门划分的意义也难以评估。它可能并不像上文描述所暗示的那样清楚。乍看之下，红暴似乎在轻工和纺织业相对发达的城市中建立了更稳固的基础。但女性在这一部门职工中占很大比例，而文化大革命中的大部分斗争由男性承担。街头人数的绝对规模往往十分重要。此外，重工业部门可能有更多临时工和合同工，这是另一个有利于联合总部的因素。这类推测或许有助于理解两个组织的相对力量。

As 1966 drew to a close the provincial authorities found themselves in a
no-win situation. Demands for pay increases, improved working conditions
and a change in the terms of employment from temporary or contract to
permanent status, with all the accompanying benefits, security and
rights, proliferated. In order to place further pressure on local
leaders, central radicals such as Jiang Qing momentarily and
opportunistically expressed support for such demands. If provincial
leaders acceded to such claims they were accused of inciting economism
or encouraging strikes to buy off the workers. On the other hand, if
local leaders refused to give in to particularistic demands they were
charged with condoning exploitation and inequity. The new
military-dominated administrations which were established in the wake of
failed power seizures early in 1967 faced the same dilemma, but they
were more ruthless in suppressing interest-group demands.\footnote{For
  an interpretation of such events in Shanghai, see Walder, Chang
  Ch'un-ch'iao and Shanghai's January Revolution, esp.~chs.~5-6, 8. For
  a description of the situation in the Zhejiang port city of Ningbo,
  see Ningbo regional service, January 13,1967, SWB/FE/2368/B/12-14. See
  also Michel Oksenberg, ``Occupational Groups in Chinese Society and
  the Cultural Revolution'', The Cultural Revolution in 1967 in Review,
  Michigan Papers in Chinese Studies, No.~2 (Ann Arbor: University of
  Michigan, 1968), pp.~1-44.

  关于上海此类事件的解释，见 Walder, Chang Ch'un-ch'iao and Shanghai's
  January Revolution，尤见第5-6、8章。关于浙江港口城市宁波的情况，见
  Ningbo regional service，1967年1月13日，SWB/FE/2368/B/12-14。另见
  Michel Oksenberg, ``Occupational Groups in Chinese Society and the
  Cultural Revolution'', The Cultural Revolution in 1967 in Review,
  Michigan Papers in Chinese Studies, No.~2 (Ann Arbor: University of
  Michigan, 1968), 第1-44页。} Nevertheless, accusations that officials
deliberately provoked and then submitted to economistic demands, so as
to bankrupt the state treasury and sabotage the Cultural Revolution,
continued to be made.\footnote{Xinhua, January 10,1967, CNS, 155,
  (January 26,1967), p.~A 8; Xinhua, January 20,1967, SWB/FE/2372/B/29.

  新华社，1967年1月10日，CNS,
  155（1967年1月26日），第A8页；新华社，1967年1月20日，SWB/FE/2372/B/29。}
Even rebel groups, vociferous in their condemnation of economism,
succumbed to the lure of emoluments in cash and kind.\footnote{ZJRB,
  January 17,1967, p.~1.

  ZJRB，1967年1月17日，第1页。} On 20 January and 1 February 1967 United
Headquarters, together with other rebel groups, issued urgent notices
calling for a cessation of economism in rural areas.\footnote{``Urgent
  Notice on striking down counter-revolutionary economism in the
  countryside'', January 20,1967, ZJRB, January 21,1967; urgent notice
  of February 1,1967, issued by eight rebel groups, ZJRB, February
  1,1967, p.~1.

  《关于打击农村反革命经济主义的紧急通知》，1967年1月20日，ZJRB，1967年1月21日；八个造反团体1967年2月1日发布的紧急通知，ZJRB，1967年2月1日，第1页。}
The notice of 20 January gave some indication of the intense struggle
then in progress in the province. It accused the power-holders of
provoking peasants to reject state requisition policies and fight among
themselves. Commune collective funds were being distributed to
individual peasants, who had also withdrawn large sums from banks and
credit cooperatives. The circular directed groups such as rusticated
educated youth, who had taken advantage of the breakdown in social order
to return to the cities, to go back to the countryside. Requests to
local military and militia forces to protect mass organizations, and not
beat and detain their members, revealed the attitude of the powerful
rural armed forces toward the rebels. At this critical stage of the
Cultural Revolution both Red Storm and United Headquarters were striving
to gain official recognition from the central Cultural Revolution Group.
With Beijing calling for the unseating of capitalist-roaders across the
country, each individual seizure of power would assist the rebel group
responsible obtain such approval.

随着1966年接近尾声，省级当局发现自己陷入无论如何都无法取胜的处境。加薪、改善工作条件、把临时或合同用工改为永久身份并享有相应福利、保障和权利的要求大量出现。为了进一步向地方领导施压，江青等中央激进派一度机会主义地表示支持这些要求。如果省级领导人同意这些要求，他们就会被指控煽动经济主义或鼓励罢工以收买工人。另一方面，如果地方领导拒绝让步于特殊利益诉求，他们又会被指责纵容剥削和不平等。1967年初夺权失败后建立的以军队为主导的新行政机构面临同样困境，但它们在镇压利益集团诉求方面更为冷酷。\footnote{For
  an interpretation of such events in Shanghai, see Walder, Chang
  Ch'un-ch'iao and Shanghai's January Revolution, esp.~chs.~5-6, 8. For
  a description of the situation in the Zhejiang port city of Ningbo,
  see Ningbo regional service, January 13,1967, SWB/FE/2368/B/12-14. See
  also Michel Oksenberg, ``Occupational Groups in Chinese Society and
  the Cultural Revolution'', The Cultural Revolution in 1967 in Review,
  Michigan Papers in Chinese Studies, No.~2 (Ann Arbor: University of
  Michigan, 1968), pp.~1-44.

  关于上海此类事件的解释，见 Walder, Chang Ch'un-ch'iao and Shanghai's
  January Revolution，尤见第5-6、8章。关于浙江港口城市宁波的情况，见
  Ningbo regional service，1967年1月13日，SWB/FE/2368/B/12-14。另见
  Michel Oksenberg, ``Occupational Groups in Chinese Society and the
  Cultural Revolution'', The Cultural Revolution in 1967 in Review,
  Michigan Papers in Chinese Studies, No.~2 (Ann Arbor: University of
  Michigan, 1968), 第1-44页。}
尽管如此，关于官员蓄意挑动并随后屈从于经济主义要求，以耗空国库、破坏文化大革命的指控仍不断出现。\footnote{Xinhua,
  January 10,1967, CNS, 155, (January 26,1967), p.~A 8; Xinhua, January
  20,1967, SWB/FE/2372/B/29.

  新华社，1967年1月10日，CNS,
  155（1967年1月26日），第A8页；新华社，1967年1月20日，SWB/FE/2372/B/29。}
即使是高声谴责经济主义的造反派组织，也屈服于现金和实物报酬的诱惑。\footnote{ZJRB,
  January 17,1967, p.~1.

  ZJRB，1967年1月17日，第1页。}
1967年1月20日和2月1日，联合总部同其他造反团体一起发布紧急通知，要求农村地区停止经济主义。\footnote{``Urgent
  Notice on striking down counter-revolutionary economism in the
  countryside'', January 20,1967, ZJRB, January 21,1967; urgent notice
  of February 1,1967, issued by eight rebel groups, ZJRB, February
  1,1967, p.~1.

  《关于打击农村反革命经济主义的紧急通知》，1967年1月20日，ZJRB，1967年1月21日；八个造反团体1967年2月1日发布的紧急通知，ZJRB，1967年2月1日，第1页。}
1月20日的通知显示了当时省内正在进行的激烈斗争。它指控掌权者煽动农民拒绝国家征购政策并相互争斗。公社集体资金被分发给农民个人，农民还从银行和信用社提取了大量存款。通知要求插队知识青年等趁社会秩序崩溃返回城市的群体回到农村。通知请求地方军队和民兵保护群众组织，不要殴打和拘留其成员，这显示出强大的农村武装力量对造反派的态度。在文化大革命这一关键阶段，红暴和联合总部都在努力争取中央文化革命小组的正式承认。既然北京号召在全国范围内打倒走资本主义道路的当权派，那么每一次个别夺权都会帮助负责的造反组织获得这种认可。

\textbf{TABLE THREE: Leading Personnel, as far as it is known, of the
Cultural Revolution Mass Organizations, 1966-68}

\textbf{表三：据已知情况，文化大革命群众组织主要人员，1966-68年}

{\def\LTcaptype{none} % do not increment counter
\begin{longtable}[]{@{}
  >{\raggedright\arraybackslash}p{(\linewidth - 4\tabcolsep) * \real{0.4118}}
  >{\raggedright\arraybackslash}p{(\linewidth - 4\tabcolsep) * \real{0.1765}}
  >{\raggedright\arraybackslash}p{(\linewidth - 4\tabcolsep) * \real{0.4118}}@{}}
\toprule\noalign{}
\begin{minipage}[b]{\linewidth}\raggedright
Organization
\end{minipage} & \begin{minipage}[b]{\linewidth}\raggedright
Name
\end{minipage} & \begin{minipage}[b]{\linewidth}\raggedright
Unit/Support
\end{minipage} \\
\midrule\noalign{}
\endhead
\bottomrule\noalign{}
\endlastfoot
United Headquarters & Zhang Yongsheng & Zhejiang College of Fine Arts \\
United Headquarters & Du Yingxin & Zhejiang College of Fine Arts \\
United Headquarters & Li Xiantong & Hangzhou University \\
United Headquarters & He Xianchun & Hangzhou Heavy Machinery Factory \\
United Headquarters & Weng Senhe (after June 1967) & Hangzhou Silk
Printing \& Dyeing Complex \\
United Headquarters & Hu Yulan & Zhejiang Daily \\
United Headquarters & Lin Jianxue & Hangzhou University \\
United Headquarters & Huang Yintang (most probably defected with Weng
Senhe) & Hangzhou Silk Printing \& Dyeing Complex \\
United Headquarters & Xia Genfa & Hangzhou Oxygen Generator Factory \\
United Headquarters & Ye Rende & Hangzhou Iron \& Steel Mill \\
United Headquarters & Supporters & 20th Army (Unit 6409); 5th Air Force
(Unit 7350) \\
Red Storm & Weng Senhe (until June 1967) & \\
Red Storm & Fang Jianwen & Hangzhou No.~1 Cotton Textile \& Dyeing
Mill \\
Red Storm & Liu Ying & Zhejiang University \\
Red Storm & Teng Zhu & Zhejiang University \\
Red Storm & Qiu Honggen & Xuejun Middle School \\
Red Storm & Supporters & Zhejiang Provincial Military District; Nanjing
Military Region; East China Fleet Units; 22nd Army (stationed on
Zhoushan island chain) \\
\end{longtable}
}

{\def\LTcaptype{none} % do not increment counter
\begin{longtable}[]{@{}
  >{\raggedright\arraybackslash}p{(\linewidth - 4\tabcolsep) * \real{0.4118}}
  >{\raggedright\arraybackslash}p{(\linewidth - 4\tabcolsep) * \real{0.1765}}
  >{\raggedright\arraybackslash}p{(\linewidth - 4\tabcolsep) * \real{0.4118}}@{}}
\toprule\noalign{}
\begin{minipage}[b]{\linewidth}\raggedright
组织
\end{minipage} & \begin{minipage}[b]{\linewidth}\raggedright
姓名
\end{minipage} & \begin{minipage}[b]{\linewidth}\raggedright
单位/支持力量
\end{minipage} \\
\midrule\noalign{}
\endhead
\bottomrule\noalign{}
\endlastfoot
联合总部 & 张永生 & 浙江美术学院 \\
联合总部 & 杜英信 & 浙江美术学院 \\
联合总部 & 李显通 & 杭州大学 \\
联合总部 & 贺贤春 & 杭州重型机器厂 \\
联合总部 & 翁森鹤（1967年6月以后） & 杭州丝绸印染联合厂 \\
联合总部 & 胡玉兰 & 浙江日报 \\
联合总部 & 林建学 & 杭州大学 \\
联合总部 & 黄阴堂（很可能随翁森鹤倒戈） & 杭州丝绸印染联合厂 \\
联合总部 & 夏根法 & 杭州制氧机厂 \\
联合总部 & 叶仁德 & 杭州钢铁厂 \\
联合总部 & 支持者 & 第二十军（6409部队）；第五空军（7350部队） \\
红暴 & 翁森鹤（至1967年6月） & \\
红暴 & 方剑文 & 杭州第一棉纺织印染厂 \\
红暴 & 刘英 & 浙江大学 \\
红暴 & 滕铸 & 浙江大学 \\
红暴 & 裘红根 & 学军中学 \\
红暴 & 支持者 &
浙江省军区；南京军区；东海舰队单位；第二十二军（驻舟山群岛） \\
\end{longtable}
}

\section{The Downfall of Jiang Hua}\label{the-downfall-of-jiang-hua}

\section{江华的倒台}\label{ux6c5fux534eux7684ux5012ux53f0}

Rebellion in factories and offices was a mere side-show compared to the
challenge of removing the CCP ZPC from power. With the successful
overthrow of the Shanghai municipal authorities in mid-January 1967, the
central authorities issued a call to ``proletarian revolutionaries'' on
20 January to ``form a great alliance to seize power from those in
authority who are taking the capitalist road''. The problem of deciding
which provincial leaders were following the capitalist road remained one
for local rebels to thrash out, and it proved as difficult to judge in
Zhejiang as elsewhere in China.

同把中共浙江省委赶下台这一挑战相比，工厂和机关中的造反不过是支线插曲。1967年1月中旬上海市当局被成功推翻后，中央当局于1月20日号召``无产阶级革命派''\,``组成大联合，向党内一小撮走资本主义道路当权派夺权''。哪些省级领导人在走资本主义道路，这一问题仍需地方造反派自己争辩解决；事实证明，在浙江作出判断同在中国其他地方一样困难。

The first struggle rally organized against a senior member of the ZPC
elite took place on 5 January 1967. The principal target was Chen Bing,
one of Jiang Hua's closest subordinates and appointed by Jiang to head
the Zhejiang Cultural Revolution Group. Chen and one of his deputies in
the propaganda department, who was also editor-in-chief of the Zhejiang
Daily, were paraded in dunces' hats and subjected to a torrent of verbal
abuse.\footnote{ZJRB, January 13,1967, p.~3; January 14,1967, p.~2.

  ZJRB，1967年1月13日，第3页；1967年1月14日，第2页。} Both the
provincial and municipal leaders came under political siege at another
large meeting on 17 January\footnote{ZJRB, January 19, 1967, p.~3.

  ZJRB，1967年1月19日，第3页。} and the seals of power were seized at
the provincial propaganda department, the office of the ZPC and at the
provincial finance department on 16, 18 and 20 January
respectively.\footnote{ZJRB, January 20,1967, p.~3; February 16,1967,
  p.~1.

  ZJRB，1967年1月20日，第3页；1967年2月16日，第1页。} United
Headquarters convened a further mass rally in Hangzhou on 25 January at
which leaders of the CCP ZPC and HMC such as Wu Xian, Wang Zida (王子达)
and Li Weixin (李维新), member of the CCP ZPC standing committee,
Deputy-governor of Zhejiang and recently-deposed Director of the
provincial Finance and Trade Department, were paraded and denounced for
approving the issuing of bonuses and funds to ``royalist''
organizations.\footnote{ZJRB, January 26,1967, p.~3.

  ZJRB，1967年1月26日，第3页。} ZPC Secretary Chen Weida, apparently
fearing that his turn would come next, fled to the safety of the ZPMD
headquarters situated at Qingbo Gate (清波门) in the south-west of
downtown Hangzhou. Central directives dated 5 October 1966, 16 November
1966, and 14 January 1967 had forbidden provincial authorities, under
the pretext of transferring confidential documents and files for
safe-keeping, from storing ``black material'' on rebels in military
headquarters. An article in Zhejiang Daily, without naming Chen, alleged
that he had defied this edict with the knowledge and complicity of the
ZPMD leadership.\footnote{ZJRB, January 18,1967, p.~3.

  ZJRB，1967年1月18日，第3页。} Once United Headquarters discovered
Chen's presence in the PLA building its members surrounded it, thus
provoking a confrontation with the military. On the following day, 26
January, Zhang Xiulong led soldiers of the ZPMD to relieve the siege.
The PLA had already received orders signed by Mao that it end its false
neutrality in the Cultural Revolution and intervene to support the left.
Like its counterparts garrisoned in other provincial capitals, the
Zhejiang military leadership was reluctant to support the anarchic
behavior of misguided and immature youth, particularly when the rebels'
wrath was directed against its colleagues of the provincial party
committee. Additionally, children of military cadres were most probably
prominent members of the rival mass organization, Red Storm. The account
of the incident in the provincial newspaper stated that twenty rebels
had sustained injuries, with two being in danger of losing their lives.
A joint editorial of Zhejiang Daily and Hangzhou Daily, now under the
control of United Headquarters, held Jiang Hua, Long Qian and Zhang
Xiulong principally responsible for the melee.\footnote{ZJRB, January
  27,1967, p.~2; 中共年报,1968年 (Communist China Yearbook 1968) (Hong
  Kong: Union Research Institute, 1968), p.~433.

  ZJRB，1967年1月27日，第2页；《中共年报,1968年》(Communist China
  Yearbook 1968)（香港：友联研究所，1968），第433页。} On 29 January, in
response to the bloody confrontation three days earlier, United
Headquarters called a rally at which the leaders of the ZPC and the ZPMD
came under fierce attack. Further violence ensued.\footnote{Zhonggong
  nianbao, 1968, p.~433; Facts \& Features, 1:13, p.~10.

  Zhonggong nianbao, 1968，第433页；Facts \& Features, 1:13，第10页。}
Not to be outdone, members of Red Storm attacked the provincial
government offices with the objective of seizing party leaders
considered unsympathetic to their organization such as Cao Xiangren and
Shen Ce, head of the ZPC Organization Department.

针对省委精英高级成员组织的第一次斗争大会发生在1967年1月5日。主要目标是陈冰，他是江华最亲近的下属之一，并由江任命为浙江文化革命小组负责人。陈和他在宣传部的一名副手，即《浙江日报》总编辑，被戴上高帽游街，并遭受滔滔不绝的辱骂。\footnote{ZJRB,
  January 13,1967, p.~3; January 14,1967, p.~2.

  ZJRB，1967年1月13日，第3页；1967年1月14日，第2页。}
1月17日的另一次大型会议上，省市领导均遭到政治围攻，\footnote{ZJRB,
  January 19, 1967, p.~3.

  ZJRB，1967年1月19日，第3页。}
1月16日、18日和20日，省宣传部、省委办公机构和省财政厅的权力印章分别被夺取。\footnote{ZJRB,
  January 20,1967, p.~3; February 16,1967, p.~1.

  ZJRB，1967年1月20日，第3页；1967年2月16日，第1页。}
1月25日，联合总部在杭州又召开一次群众大会，会上，中共浙江省委和杭州市委领导如吴宪、王子达，以及中共浙江省委常委、浙江副省长、刚被罢免的省财贸部主任李维新等，被游斗并遭谴责，原因是他们批准向``保皇派''组织发放奖金和经费。\footnote{ZJRB,
  January 26,1967, p.~3.

  ZJRB，1967年1月26日，第3页。}
省委书记陈伟达显然担心下一个就轮到自己，逃往位于杭州市区西南清波门的浙江省军区司令部寻求安全。1966年10月5日、1966年11月16日和1967年1月14日的中央指示，禁止省级当局以转移机密文件和档案以便保管为借口，在军事机关中存放有关造反派的``黑材料''。《浙江日报》一篇文章虽未点名陈，却指称他在省军区领导知情并合谋的情况下违抗了这项命令。\footnote{ZJRB,
  January 18,1967, p.~3.

  ZJRB，1967年1月18日，第3页。}
联合总部一发现陈在解放军大楼内，成员便包围大楼，从而引发同军方的对峙。次日即1月26日，张秀龙率领省军区士兵解围。解放军此前已经接到毛签署的命令，要求其结束在文化大革命中的虚假中立，介入并支持左派。像驻扎在其他省会的同级军方一样，浙江军事领导层不愿支持被误导且不成熟的青年所表现出的无政府行为，特别是当造反派的怒火指向其省党委同事时。此外，军队干部子女很可能是敌对群众组织红暴的重要成员。省报对事件的报道说，二十名造反派受伤，其中两人有生命危险。此时受联合总部控制的《浙江日报》和《杭州日报》发表联合社论，认为江华、龙潜和张秀龙应对这场混战负主要责任。\footnote{ZJRB,
  January 27,1967, p.~2; 中共年报,1968年 (Communist China Yearbook 1968)
  (Hong Kong: Union Research Institute, 1968), p.~433.

  ZJRB，1967年1月27日，第2页；《中共年报,1968年》(Communist China
  Yearbook 1968)（香港：友联研究所，1968），第433页。}
1月29日，为回应三天前的流血冲突，联合总部召开集会，省委和省军区领导在会上受到猛烈攻击。随后又发生更多暴力。\footnote{Zhonggong
  nianbao, 1968, p.~433; Facts \& Features, 1:13, p.~10.

  Zhonggong nianbao, 1968，第433页；Facts \& Features, 1:13，第10页。}
红暴成员也不甘示弱，攻击省政府机关，目标是抓捕被认为不同情其组织的党内领导人，如曹祥仁和省委组织部部长沈策。

At this critical juncture the central authorities issued a decision
placing responsibility for the 26 January incident firmly on the
shoulders of the ZPC and unnamed leaders of the ZPMD.\footnote{``中共中央国务院中央军委关于浙江最近发生事件的决定'',
  (Decision of the CCP CC, State Council and Military Affairs Commission
  concerning the recent incident in Zhejiang), January 30,1967, CCP
  Documents of the Great Proletarian Cultural Revolution, 1966-67 (Hong
  Kong: Union Research Institute, 1968), pp.~219-23.

  《中共中央国务院中央军委关于浙江最近发生事件的决定》(Decision of the
  CCP CC, State Council and Military Affairs Commission concerning the
  recent incident in Zhejiang)，1967年1月30日，CCP Documents of the
  Great Proletarian Cultural Revolution,
  1966-67（香港：友联研究所，1968），第219-223页。} The central document
demanded that the children of high-ranking cadres who had sought refuge
in garrison headquarters with Chen Weida be handed over to United
Headquarters for punishment. It ordered leaders of the ZPC and ZPMD to
make thorough self-criticisms (深刻检讨) before the masses. The ZPMD was
again commanded, in the spirit of the 23 January order to the PLA to

在这一关键时刻，中央当局发布决定，明确把1月26日事件的责任归于浙江省委和省军区某些未点名领导人。\footnote{``中共中央国务院中央军委关于浙江最近发生事件的决定'',
  (Decision of the CCP CC, State Council and Military Affairs Commission
  concerning the recent incident in Zhejiang), January 30,1967, CCP
  Documents of the Great Proletarian Cultural Revolution, 1966-67 (Hong
  Kong: Union Research Institute, 1968), pp.~219-23.

  《中共中央国务院中央军委关于浙江最近发生事件的决定》(Decision of the
  CCP CC, State Council and Military Affairs Commission concerning the
  recent incident in Zhejiang)，1967年1月30日，CCP Documents of the
  Great Proletarian Cultural Revolution,
  1966-67（香港：友联研究所，1968），第219-223页。}
中央文件要求，将同陈伟达一起在驻军司令部寻求庇护的高级干部子女移交联合总部处理。文件命令省委和省军区领导人在群众面前作深刻检讨。省军区再次被命令，按照1月23日给解放军命令中关于

\begin{quote}
``support the left'', to firmly stand on the side of the revolutionary
left, firmly support the great alliance of the proletarian
revolutionaries, and firmly support the revolutionary rebels' struggle
to seize power so as to facilitate deepening of the class struggle and
the all-round unfolding of the Great Proletarian Cultural Revolution in
Zhejiang.
\end{quote}

\begin{quote}
``支左''的精神，坚定站在革命左派一边，坚决支持无产阶级革命派的大联合，坚决支持革命造反派夺权斗争，以利于浙江阶级斗争的深入和无产阶级文化大革命的全面展开。
\end{quote}

The ``Decision'' requested United Headquarters to leave the ZPMD
building and to unite with other ``proletarian revolutionaries''. The
central authorities also appointed Du Fangping (杜方平), Political
Commissar of the Nanjing Military Region and later Vice-chairman of the
Jiangsu Provincial Military Control Commission,\footnote{Jiangsusheng
  dashiji, pp.~286-87.

  Jiangsusheng dashiji，第286-287页。} Ruan Xianbang (阮贤榜),
Deputy-commander of the ZPMD and Ding Jun (丁钧) of the CC's central
office to inquire into the affair. It may well have been that Du, who
was later dismissed for ultra-leftism, and Ding combined to pin the
blame on the local military, because Ruan was later accused of
``opposing the three Reds'' and being a ``representative of the those
who stubbornly adhere to the reactionary bourgeois line''.\footnote{HZRB,
  September 12, 1979, p.~1.

  HZRB，1979年9月12日，第1页。} The incident sealed Chen Weida's fate
and undoubtedly destroyed the radical central leaders' remaining
confidence in the civilian and military leadership of Zhejiang.

《决定》要求联合总部离开省军区大楼，并同其他``无产阶级革命派''联合起来。中央当局还任命南京军区政治委员、后来任江苏省军事管制委员会副主任的杜方平，\footnote{Jiangsusheng
  dashiji, pp.~286-87.

  Jiangsusheng dashiji，第286-287页。}
浙江省军区副司令员阮贤榜，以及中央办公厅的丁钧调查此事。很可能后来因``极左''被撤职的杜方平同丁钧一起把责任推到地方军方身上，因为阮后来被指控``反对三红''，并被称为``顽固坚持资产阶级反动路线者的代表''。\footnote{HZRB,
  September 12, 1979, p.~1.

  HZRB，1979年9月12日，第1页。}
这一事件决定了陈伟达的命运，并无疑摧毁了中央激进派领导人对浙江文职和军事领导层残存的信心。

The ZPMD's defiance of central directives and support for the ZPC
further infuriated United Headquarters. On 27 January Zhejiang Daily and
Hangzhou Daily published a joint editorial clamoring for the purging of
``capitalist-roaders'' from the ranks of the military. Entitled ``What
are a small handful in the Zhejiang Military District up to?''
(浙江省军区的一小撮人想干什么?!) the editorial was accompanied by an
item excoriating the local military leaders as ``bastards'' (混蛋们) and
``tools of the bourgeois dictatorship''.\footnote{ZJRB, January 27,1967,
  p.~2; Criticism group of the ZPMD Logistics Department, ``Those who
  try to destroy the Great Wall will themselves meet destruction - an
  exposure and criticism of the sworn followers, agents, and lackeys of
  Lin Biao and the 'Gang of Four' in Zhejiang in opposing and
  destabilizing the Army'', HZRB, September 14, 1978; Hangzhou Daily
  criticism group,''We must become the party's tool for public opinion,
  not a trumpeter for factionalism'', HZRB, September 28, 1977.

  ZJRB，1967年1月27日，第2页；浙江省军区后勤部批判组，``Those who try to
  destroy the Great Wall will themselves meet destruction - an exposure
  and criticism of the sworn followers, agents, and lackeys of Lin Biao
  and the 'Gang of Four' in Zhejiang in opposing and destabilizing the
  Army''，HZRB，1978年9月14日；《杭州日报》批判组，``We must become the
  party's tool for public opinion, not a trumpeter for
  factionalism''，HZRB，1977年9月28日。} The ferocity of this language,
accompanied by continued rebel attacks on arsenals,\footnote{Wang Xianyu
  and Guan Guodong (ZPMD Political Department), ``Attacking the Party
  and destabilizing the Army will not succeed'', HZRB, April 4,1979.

  王显玉、关国栋（浙江省军区政治部），``Attacking the Party and
  destabilizing the Army will not succeed''，HZRB，1979年4月4日。}
undoubtedly spurred the Hangzhou garrison command into passing a
resolution on the same day which backed the rebels.\footnote{Reported by
  Jurgen Domes, ``The Role of the Military in the formation of
  Revolutionary Committees, 1967-68'', CQ, 44 (1970), p.~130.

  Jurgen Domes 报道，``The Role of the Military in the formation of
  Revolutionary Committees, 1967-68''，CQ, 44 (1970)，第130页。}
Meetings and demonstrations were also held on January 27 to mobilize
support for the rebels and opposition to the PLA.\footnote{ZJRB, January
  28,1967, p.~2.

  ZJRB，1967年1月28日，第2页。} Nevertheless, United Headquarters'
principal target remained Jiang Hua. Its campaign to unseat him climaxed
in early February 1967. On the 12th of that month, United Headquarters
convened a rally at the provincial sports stadium to announce Jiang's
dismissal.\footnote{ZJRB (专刊), August 21,1969.

  ZJRB（专刊），1969年8月21日。} Unbeknown to United Headquarters,
however, the previous evening Jiang had departed Hangzhou in a special
aircraft despatched by Zhou Enlai. The decision to rescue Jiang had been
taken by Mao.\footnote{ZJRB, June 8,1968, p.~3; Gao Wenxian,
  ``艰难而光辉的最后岁月--记'文化大革命'期间的周恩来'' (Difficult but
  glorious last years -- recalling Zhou Enlai in the Cultural
  Revolution) RMRB, 5 January 1986, p.~5; translated in
  SWB/FE/8162/BII/3-13. Mao, on three separate occasions, had instructed
  that provincial and municipal party secretaries be brought to Beijing
  for their protection. See Ji Xizhen, ``一场捍卫党的原则的伟大斗争'' (A
  Great Struggle to defend the Party's principles), RMRB, February 26,
  1979, p.~2, transl. in SWB/FE/6056/BII/5. With the Chairman's approval
  Zhou Enlai drew up a list of about twenty regional and
  provincial-level party secretaries to be protected. The list included
  Jiang Hua. See Jin Chunming, ``文化大革命''论析 (A Discussion and
  Analysis of the ``Cultural Revolution'') (Shanghai: Shanghai Renmin
  Chubanshe, 1985), p.~122.

  ZJRB，1968年6月8日，第3页；高文谦，《艰难而光辉的最后岁月--记'文化大革命'期间的周恩来》(Difficult
  but glorious last years -- recalling Zhou Enlai in the Cultural
  Revolution)，RMRB，1986年1月5日，第5页；译载
  SWB/FE/8162/BII/3-13。毛曾三次指示将省市党委书记带到北京加以保护。见季锡珍，《一场捍卫党的原则的伟大斗争》(A
  Great Struggle to defend the Party's
  principles)，RMRB，1979年2月26日，第2页，译载
  SWB/FE/6056/BII/5。经主席批准，周恩来拟定了一份约二十名地区和省级党委书记的保护名单，其中包括江华。见金春明，《``文化大革命''论析》(A
  Discussion and Analysis of the ``Cultural
  Revolution'')（上海：上海人民出版社，1985），第122页。} The Nanjing
Military Region command informed Red Storm of the decision on the same
night, after Jiang's departure.\footnote{Tragically for Jiang, he left
  Hangzhou a widower. Wu Zhonglian could not endure the intense
  political and physical pressure exerted on her by the rebels and
  committed suicide on January 19,1967. ZJRB, November 6,1978, p.~1. On
  January 27 a meeting of rebels in Hangzhou debated Wu's ``crimes''. Lü
  Jianguang, Director of the provincial Public Security Bureau, attended
  as a target of criticism and his superior in the public security
  system, Wang Fang, was described as a ``counter-revolutionary
  revisionist''. Cao Xiangren, a member of the CCP ZPC secretariat,
  delivered a speech in which he supported the call for Wu's posthumous
  expulsion from the CCP. ZJRB, January 28,1967, p.~3. A more recent
  article dates Wu's death at January 17,1967. Kang Keqing et al., ZJRB,
  January 28,1987.

  对江而言，悲剧在于他以鳏夫身份离开杭州。吴仲廉无法承受造反派施加的强烈政治和身体压力，于1967年1月19日自杀。ZJRB，1978年11月6日，第1页。1月27日，杭州造反派会议讨论吴的``罪行''。省公安厅厅长吕剑光作为批判对象出席，其在公安系统中的上级王芳被称为``反革命修正主义分子''。中共浙江省委书记处成员曹祥仁发表讲话，支持在吴死后开除其党籍的呼吁。ZJRB，1967年1月28日，第3页。一篇较近的文章把吴的死亡日期定为1967年1月17日。康克清等，ZJRB，1987年1月28日。}
United Headquarters did not discover what had transpired until Jiang's
absence from the rally alerted them to the fact that something was
amiss. With the encouragement of the Nanjing Military Region, Red Storm
broke up the meeting.\footnote{ZJRB, August 21,1969; Qiu Honggen
  (裘红根) (a prominent member of Red Storm who at the time was a
  student at Hangzhou's most prestigious secondary school, called the
  ``Learn from the PLA Middle School'' (学军中学), ``From my two
  detentions see the wolfish ambition of Lin Bao and the Gang of Four'',
  HZRB, November 26, 1978.

  ZJRB，1969年8月21日；裘红根（红暴重要成员，当时为杭州最著名中学``学军中学''的学生），``From
  my two detentions see the wolfish ambition of Lin Bao and the Gang of
  Four''，HZRB，1978年11月26日。} And so a prolonged period of bitter
antagonism between the two major mass organizations was inaugurated.

省军区违抗中央指示并支持省委，进一步激怒了联合总部。1月27日，《浙江日报》和《杭州日报》发表联合社论，叫嚣要从军队队伍中清除``走资派''。社论题为《浙江省军区的一小撮人想干什么?!》，旁边还配有一篇文章，痛斥地方军方领导人为``混蛋们''和``资产阶级专政的工具''。\footnote{ZJRB,
  January 27,1967, p.~2; Criticism group of the ZPMD Logistics
  Department, ``Those who try to destroy the Great Wall will themselves
  meet destruction - an exposure and criticism of the sworn followers,
  agents, and lackeys of Lin Biao and the 'Gang of Four' in Zhejiang in
  opposing and destabilizing the Army'', HZRB, September 14, 1978;
  Hangzhou Daily criticism group,''We must become the party's tool for
  public opinion, not a trumpeter for factionalism'', HZRB, September
  28, 1977.

  ZJRB，1967年1月27日，第2页；浙江省军区后勤部批判组，``Those who try to
  destroy the Great Wall will themselves meet destruction - an exposure
  and criticism of the sworn followers, agents, and lackeys of Lin Biao
  and the 'Gang of Four' in Zhejiang in opposing and destabilizing the
  Army''，HZRB，1978年9月14日；《杭州日报》批判组，``We must become the
  party's tool for public opinion, not a trumpeter for
  factionalism''，HZRB，1977年9月28日。}
如此激烈的语言，加上造反派持续攻击军火库，\footnote{Wang Xianyu and Guan
  Guodong (ZPMD Political Department), ``Attacking the Party and
  destabilizing the Army will not succeed'', HZRB, April 4,1979.

  王显玉、关国栋（浙江省军区政治部），``Attacking the Party and
  destabilizing the Army will not succeed''，HZRB，1979年4月4日。}
无疑促使杭州警备区司令部在同日通过一项支持造反派的决议。\footnote{Reported
  by Jurgen Domes, ``The Role of the Military in the formation of
  Revolutionary Committees, 1967-68'', CQ, 44 (1970), p.~130.

  Jurgen Domes 报道，``The Role of the Military in the formation of
  Revolutionary Committees, 1967-68''，CQ, 44 (1970)，第130页。}
1月27日还举行了会议和示威，以动员支持造反派、反对解放军。\footnote{ZJRB,
  January 28,1967, p.~2.

  ZJRB，1967年1月28日，第2页。}
尽管如此，联合总部的主要目标仍然是江华。其打倒江的运动在1967年2月初达到高潮。2月12日，联合总部在省体育场召开大会，宣布罢免江。\footnote{ZJRB
  (专刊), August 21,1969.

  ZJRB（专刊），1969年8月21日。}
然而，联合总部并不知道，前一天晚上江已乘周恩来派遣的专机离开杭州。营救江的决定是毛作出的。\footnote{ZJRB,
  June 8,1968, p.~3; Gao Wenxian,
  ``艰难而光辉的最后岁月--记'文化大革命'期间的周恩来'' (Difficult but
  glorious last years -- recalling Zhou Enlai in the Cultural
  Revolution) RMRB, 5 January 1986, p.~5; translated in
  SWB/FE/8162/BII/3-13. Mao, on three separate occasions, had instructed
  that provincial and municipal party secretaries be brought to Beijing
  for their protection. See Ji Xizhen, ``一场捍卫党的原则的伟大斗争'' (A
  Great Struggle to defend the Party's principles), RMRB, February 26,
  1979, p.~2, transl. in SWB/FE/6056/BII/5. With the Chairman's approval
  Zhou Enlai drew up a list of about twenty regional and
  provincial-level party secretaries to be protected. The list included
  Jiang Hua. See Jin Chunming, ``文化大革命''论析 (A Discussion and
  Analysis of the ``Cultural Revolution'') (Shanghai: Shanghai Renmin
  Chubanshe, 1985), p.~122.

  ZJRB，1968年6月8日，第3页；高文谦，《艰难而光辉的最后岁月--记'文化大革命'期间的周恩来》(Difficult
  but glorious last years -- recalling Zhou Enlai in the Cultural
  Revolution)，RMRB，1986年1月5日，第5页；译载
  SWB/FE/8162/BII/3-13。毛曾三次指示将省市党委书记带到北京加以保护。见季锡珍，《一场捍卫党的原则的伟大斗争》(A
  Great Struggle to defend the Party's
  principles)，RMRB，1979年2月26日，第2页，译载
  SWB/FE/6056/BII/5。经主席批准，周恩来拟定了一份约二十名地区和省级党委书记的保护名单，其中包括江华。见金春明，《``文化大革命''论析》(A
  Discussion and Analysis of the ``Cultural
  Revolution'')（上海：上海人民出版社，1985），第122页。}
江离开后，当晚南京军区司令部将这一决定通知了红暴。\footnote{Tragically
  for Jiang, he left Hangzhou a widower. Wu Zhonglian could not endure
  the intense political and physical pressure exerted on her by the
  rebels and committed suicide on January 19,1967. ZJRB, November
  6,1978, p.~1. On January 27 a meeting of rebels in Hangzhou debated
  Wu's ``crimes''. Lü Jianguang, Director of the provincial Public
  Security Bureau, attended as a target of criticism and his superior in
  the public security system, Wang Fang, was described as a
  ``counter-revolutionary revisionist''. Cao Xiangren, a member of the
  CCP ZPC secretariat, delivered a speech in which he supported the call
  for Wu's posthumous expulsion from the CCP. ZJRB, January 28,1967,
  p.~3. A more recent article dates Wu's death at January 17,1967. Kang
  Keqing et al., ZJRB, January 28,1987.

  对江而言，悲剧在于他以鳏夫身份离开杭州。吴仲廉无法承受造反派施加的强烈政治和身体压力，于1967年1月19日自杀。ZJRB，1978年11月6日，第1页。1月27日，杭州造反派会议讨论吴的``罪行''。省公安厅厅长吕剑光作为批判对象出席，其在公安系统中的上级王芳被称为``反革命修正主义分子''。中共浙江省委书记处成员曹祥仁发表讲话，支持在吴死后开除其党籍的呼吁。ZJRB，1967年1月28日，第3页。一篇较近的文章把吴的死亡日期定为1967年1月17日。康克清等，ZJRB，1987年1月28日。}
直到江缺席大会使他们察觉有异，联合总部才发现发生了什么。在南京军区鼓励下，红暴冲散了会议。\footnote{ZJRB,
  August 21,1969; Qiu Honggen (裘红根) (a prominent member of Red Storm
  who at the time was a student at Hangzhou's most prestigious secondary
  school, called the ``Learn from the PLA Middle School'' (学军中学),
  ``From my two detentions see the wolfish ambition of Lin Bao and the
  Gang of Four'', HZRB, November 26, 1978.

  ZJRB，1969年8月21日；裘红根（红暴重要成员，当时为杭州最著名中学``学军中学''的学生），``From
  my two detentions see the wolfish ambition of Lin Bao and the Gang of
  Four''，HZRB，1978年11月26日。}
于是，两个主要群众组织之间长期而尖锐的敌对关系由此开启。

Although Jiang was not the only provincial-level cadre to benefit from
Mao's intervention, the question of why Mao agreed to protect Jiang
requires an answer. Several factors seem to have been involved. First,
as a national minority cadre, Jiang qualified for lenient and privileged
treatment. Second, his personal connections with Mao dated back to
Jinggangshan -- cradle of the Chinese revolution -- the Long March and
Yanan, where Jiang had served among Mao's security staff. Jiang's wife
had once worked as Mao's secretary on Jinggangshan.\footnote{Kang Keqing
  et al., ZJRB, January 28, 1987.

  康克清等，ZJRB，1987年1月28日。} Third, Jiang had maintained regular
contact with the Chairman after 1949 due to Mao's preference for
Hangzhou as a place to rest and escape the pressures of the national
capital. As a southerner Mao presumably preferred the climate, landscape
and food of Hangzhou. As the hosting provincial leader, Jiang was partly
responsible for planning inspection tours and making certain that Mao's
various requirements and comforts were well attended to. Zhejiang's
progress in agricultural production, and its regular grain and financial
contributions to the center, may well have impressed Mao.

尽管江并非唯一因毛的干预而受益的省级干部，但为何毛同意保护江这一问题需要回答。似乎涉及几个因素。首先，作为少数民族干部，江符合从宽和优待条件。其次，他同毛的个人联系可追溯到中国革命摇篮井冈山、长征和延安时期，当时江曾在毛的警卫人员中任职。江的妻子曾在井冈山担任毛的秘书。\footnote{Kang
  Keqing et al., ZJRB, January 28, 1987.

  康克清等，ZJRB，1987年1月28日。}
第三，1949年后，由于毛偏爱杭州作为休息和摆脱首都压力的地方，江与主席保持了经常接触。作为南方人，毛大概偏爱杭州的气候、风景和食物。作为接待方的省级领导，江部分负责安排视察行程，并确保毛的各种要求和舒适需求得到妥善照顾。浙江在农业生产方面的进展，以及其对中央稳定的粮食和财政贡献，很可能给毛留下了印象。

After Jiang's arrival in Beijing Tan Zhenlin showed great concern for
his fate. On 11 February he met Jiang and told him not to submit a
self-examination (检讨).\footnote{ZJRB, June 8,1968, p.~3.

  ZJRB，1968年6月8日，第3页。} In a letter of 17 February 1967 to Lin
Biao, written the day after a heated confrontation between members of
the Cultural Revolution Group and veteran central officials in Beijing,
Tan protested bitterly against the harsh treatment meted out to such
local party leaders as Tan Qilong\footnote{See Tan Qilong's grateful
  recollection of the elder Tan's protection and concern. Tan Qilong,
  Renwu, p.~48.

  见谭启龙对谭震林保护与关怀的感激回忆。谭启龙，Renwu，第48页。}
{[}Shandong{]} and Jiang Hua. Both men had served under Tan in Zhejiang
after 1949. Tan Zhenlin wrote that they had been ``struggled against,
had worn high hats and been placed in the aeroplane position with their
health broken, family scattered (妻离子散) and fortune lost
(倾家荡产)''.\footnote{Jin Chunming, A Discussion and Analysis,
  pp.~127-8 (quotation from p.~127). While Jiang undoubtedly went
  through a traumatic experience it is doubtful that his agony matched
  that endured for five years by the noted Hangzhou Beijing opera star
  Gai Jiaotian. Gai died in 1971 in a leaky, gloomy hut in Hangzhou
  after enduring years of humiliation and torture. See Shang Yiren
  ``江南活武松之死'' (The death of south China's living Wu Song), in
  十年奇冤录 (A record of ten years of strange injustices), (Beijing:
  Qunzhong Chubanshe, 1986), pp.~153-62. See also the account by Chen
  Xiuliang (陈修良), widow of former governor of Zhejiang Sha Wenhan and
  together with him branded a rightist in 1957, of her suffering in the
  years 1966-72, in 【文革】中经受的磨难 (Hardships experienced during
  the ``Cultural Revolution''), Lianhe shibao, August 19,1988, p.~3.

  金春明，A Discussion and
  Analysis，第127-128页（引文见第127页）。虽然江无疑经历了创伤性遭遇，但他的痛苦是否能比得上杭州著名京剧演员盖叫天五年间所承受的痛苦，令人怀疑。盖在多年羞辱和酷刑后，于1971年死在杭州一间漏雨阴暗的小屋中。见尚一人，《江南活武松之死》(The
  death of south China's living Wu Song)，载《十年奇冤录》(A record of
  ten years of strange
  injustices)（北京：群众出版社，1986），第153-162页。另见陈修良（浙江前省长沙文汉遗孀，1957年同沙一起被打成右派）关于其1966-72年遭受苦难的叙述，《【文革】中经受的磨难》(Hardships
  experienced during the ``Cultural
  Revolution'')，《联合时报》，1988年8月19日，第3页。} Then, in a 22
February meeting with a Zhejiang delegation which included
representatives of United Headquarters, Tan stated that although Jiang
had offended people and everyone wanted to bring him down, this was of
no great event. A little explanation would clear things up
(``过去很会骂人，得罪了一些人，人家要打倒他，这不要紧么！将来解释解释好！'').\footnote{ZJRB,
  June 8,1968, p.~3.

  ZJRB，1968年6月8日，第3页。} In other words, according to Tan Jiang's
mistakes were minor and he should be allowed to survive his ordeal,
suitably chastened by the experience.

江抵达北京后，谭震林非常关心他的命运。2月11日，他会见江，并告诉他不要作检讨。\footnote{ZJRB,
  June 8,1968, p.~3.

  ZJRB，1968年6月8日，第3页。}
1967年2月17日，即北京文化革命小组成员同中央老干部发生激烈对抗后的第二天，谭在写给林彪的信中，对谭启龙\footnote{See
  Tan Qilong's grateful recollection of the elder Tan's protection and
  concern. Tan Qilong, Renwu, p.~48.

  见谭启龙对谭震林保护与关怀的感激回忆。谭启龙，Renwu，第48页。}（山东）和江华等地方党委领导遭受的严酷待遇表示强烈抗议。两人1949年后都曾在浙江的谭手下任职。谭震林写道，他们``被斗争、戴高帽、喷气式，身体搞坏，妻离子散，倾家荡产''。\footnote{Jin
  Chunming, A Discussion and Analysis, pp.~127-8 (quotation from
  p.~127). While Jiang undoubtedly went through a traumatic experience
  it is doubtful that his agony matched that endured for five years by
  the noted Hangzhou Beijing opera star Gai Jiaotian. Gai died in 1971
  in a leaky, gloomy hut in Hangzhou after enduring years of humiliation
  and torture. See Shang Yiren ``江南活武松之死'' (The death of south
  China's living Wu Song), in 十年奇冤录 (A record of ten years of
  strange injustices), (Beijing: Qunzhong Chubanshe, 1986), pp.~153-62.
  See also the account by Chen Xiuliang (陈修良), widow of former
  governor of Zhejiang Sha Wenhan and together with him branded a
  rightist in 1957, of her suffering in the years 1966-72, in
  【文革】中经受的磨难 (Hardships experienced during the ``Cultural
  Revolution''), Lianhe shibao, August 19,1988, p.~3.

  金春明，A Discussion and
  Analysis，第127-128页（引文见第127页）。虽然江无疑经历了创伤性遭遇，但他的痛苦是否能比得上杭州著名京剧演员盖叫天五年间所承受的痛苦，令人怀疑。盖在多年羞辱和酷刑后，于1971年死在杭州一间漏雨阴暗的小屋中。见尚一人，《江南活武松之死》(The
  death of south China's living Wu Song)，载《十年奇冤录》(A record of
  ten years of strange
  injustices)（北京：群众出版社，1986），第153-162页。另见陈修良（浙江前省长沙文汉遗孀，1957年同沙一起被打成右派）关于其1966-72年遭受苦难的叙述，《【文革】中经受的磨难》(Hardships
  experienced during the ``Cultural
  Revolution'')，《联合时报》，1988年8月19日，第3页。}
随后，在2月22日同包括联合总部代表在内的浙江代表团会面时，谭表示，尽管江得罪了人，人人都想打倒他，这并不是什么大不了的事。稍作解释就会澄清（``过去很会骂人，得罪了一些人，人家要打倒他，这不要紧么！将来解释解释好！''）。\footnote{ZJRB,
  June 8,1968, p.~3.

  ZJRB，1968年6月8日，第3页。}
换言之，在谭看来，江的错误是轻微的，应让他熬过这场磨难，并从中受到足够教训。

Tan's intervention on behalf of his former subordinate may have reminded
the Chairman, who read the letter, of his past associations with Jiang
and softened his view of the gravity of Jiang's errors. For whatever
reason, the center's ambiguous attitude toward criticism of
Jiang\footnote{The public naming of leaders in the official press
  required approval from the center. Approval to name Jiang Hua was not
  forthcoming until November 1968. In May 1968, at a meeting with United
  Headquarters' leader Zhang Yongsheng, Jiang Qing asked him, ``Since
  other top capitalist-roaders of the province have been named, why is
  Jiang Hua not named?'' She went on to assert confidently: ``What you
  have to do is just give us a ring''. See 火炬通讯 (Torch Bulletin),
  No.~1, July 1968, translated in SCMM, 622 (August 6,1968), p.~10. But
  clearly it was not that simple. Chen Boda had not hesitated to name
  Jiang, Li Fengping, Chen Weida and Chen Bing as the leading
  capitalist-roaders in Zhejiang at a closed-door meeting in Beijing
  held in March 1968. See SCMP, 4182, p.~12. But Chen, by mentioning
  Jiang by name, was not conforming with central policy. See ibid, p.~2.
  Finally, after Jiang was named in the provincial press, his supporters
  argued that ``for Zhejiang Daily to print a person's name means
  nothing''. ZPS, December 9,1968, SWB/FE/2953/B1/13-15. Jiang was
  never, so far as I am aware, named in Central Committee documents.
  Zhao Ziyang's biographer claims that it was significant if a
  provincial leader, such as Zhao, was not named by the central press
  when the establishment of the provincial revolutionary committee was
  announced and that it signified that this particular person's status
  had not been determined (尚未定性). See Zhao Wei, Zhao Ziyang zhuan,
  p.~170. By this reasoning Jiang Hua's case was in the same category as
  Zhao's.

  官方报刊公开点名领导人需要中央批准。直到1968年11月，点名江华的批准才下达。1968年5月，在同联合总部领导人张永生会面时，江青问他：``既然省里其他头号走资派都点了名，为什么江华不点名？''她接着自信地断言：``你们要做的只是给我们打个电话。''见《火炬通讯》(Torch
  Bulletin)，第1期，1968年7月，译载 SCMM,
  622（1968年8月6日），第10页。但显然事情并没有这么简单。陈伯达曾在1968年3月北京一次闭门会议上毫不犹豫地点名江、李丰平、陈伟达和陈冰为浙江主要走资派。见
  SCMP,
  4182，第12页。但陈点江的名，并不符合中央政策。见同上，第2页。最终，当江在省级报刊上被点名后，他的支持者辩称``《浙江日报》登一个人的名字说明不了什么''。ZPS，1968年12月9日，SWB/FE/2953/B1/13-15。据我所知，中央委员会文件从未点名江。赵紫阳传记作者称，如果宣布成立省革命委员会时，像赵这样的省级领导人没有被中央报刊点名，这是有意义的，表明此人的性质尚未确定（尚未定性）。见赵蔚，《赵紫阳传》，第170页。按此推理，江华的情况与赵同属一类。}
contributed greatly to the subsequent upheavals in Zhejiang. United
Headquarters was thwarted in its efforts to destroy Jiang politically.
Red Storm perceived Zhou Enlai's action as central endorsement for its
backing of Jiang as a member of Mao's ``revolutionary headquarters''.
The Cultural Revolution Group clearly favored United Headquarters and
endorsed its attempts to overthrow Jiang. Yet its backing was
insufficient to carry the day. Jiang Hua's position remained
undetermined while the two mass organizations began a lengthy, violent
struggle for ascendancy. Long after Jiang's departure from the province
into temporary political oblivion, United Headquarters could cry out in
exasperation, ``The ghost of Jiang Hua still haunts Zhejiang Province''
(江华的幽灵还在浙江的土地上游荡着).\footnote{ZJRB, November 19,1968,
  p.~1.

  ZJRB，1968年11月19日，第1页。}

谭代表其旧部进行干预，或许使读过这封信的主席想起自己同江的旧日关系，并减轻了他对江错误严重性的看法。不论原因如何，中央对批判江的态度暧昧，\footnote{The
  public naming of leaders in the official press required approval from
  the center. Approval to name Jiang Hua was not forthcoming until
  November 1968. In May 1968, at a meeting with United Headquarters'
  leader Zhang Yongsheng, Jiang Qing asked him, ``Since other top
  capitalist-roaders of the province have been named, why is Jiang Hua
  not named?'' She went on to assert confidently: ``What you have to do
  is just give us a ring''. See 火炬通讯 (Torch Bulletin), No.~1, July
  1968, translated in SCMM, 622 (August 6,1968), p.~10. But clearly it
  was not that simple. Chen Boda had not hesitated to name Jiang, Li
  Fengping, Chen Weida and Chen Bing as the leading capitalist-roaders
  in Zhejiang at a closed-door meeting in Beijing held in March 1968.
  See SCMP, 4182, p.~12. But Chen, by mentioning Jiang by name, was not
  conforming with central policy. See ibid, p.~2. Finally, after Jiang
  was named in the provincial press, his supporters argued that ``for
  Zhejiang Daily to print a person's name means nothing''. ZPS, December
  9,1968, SWB/FE/2953/B1/13-15. Jiang was never, so far as I am aware,
  named in Central Committee documents. Zhao Ziyang's biographer claims
  that it was significant if a provincial leader, such as Zhao, was not
  named by the central press when the establishment of the provincial
  revolutionary committee was announced and that it signified that this
  particular person's status had not been determined (尚未定性). See
  Zhao Wei, Zhao Ziyang zhuan, p.~170. By this reasoning Jiang Hua's
  case was in the same category as Zhao's.

  官方报刊公开点名领导人需要中央批准。直到1968年11月，点名江华的批准才下达。1968年5月，在同联合总部领导人张永生会面时，江青问他：``既然省里其他头号走资派都点了名，为什么江华不点名？''她接着自信地断言：``你们要做的只是给我们打个电话。''见《火炬通讯》(Torch
  Bulletin)，第1期，1968年7月，译载 SCMM,
  622（1968年8月6日），第10页。但显然事情并没有这么简单。陈伯达曾在1968年3月北京一次闭门会议上毫不犹豫地点名江、李丰平、陈伟达和陈冰为浙江主要走资派。见
  SCMP,
  4182，第12页。但陈点江的名，并不符合中央政策。见同上，第2页。最终，当江在省级报刊上被点名后，他的支持者辩称``《浙江日报》登一个人的名字说明不了什么''。ZPS，1968年12月9日，SWB/FE/2953/B1/13-15。据我所知，中央委员会文件从未点名江。赵紫阳传记作者称，如果宣布成立省革命委员会时，像赵这样的省级领导人没有被中央报刊点名，这是有意义的，表明此人的性质尚未确定（尚未定性）。见赵蔚，《赵紫阳传》，第170页。按此推理，江华的情况与赵同属一类。}
极大地促成了浙江随后发生的动荡。联合总部在政治上摧毁江的努力受挫。红暴则把周恩来的行动理解为中央认可其支持江，并承认江是毛的``革命司令部''成员。文化革命小组显然偏向联合总部，并支持其推翻江的尝试。然而，这种支持不足以取得胜利。江华的地位仍未确定，而两个群众组织开始了一场漫长而暴力的争夺优势之战。江离开浙江、暂时陷入政治沉寂很久以后，联合总部仍可恼怒地喊出：``江华的幽灵还在浙江的土地上游荡着''。\footnote{ZJRB,
  November 19,1968, p.~1.

  ZJRB，1968年11月19日，第1页。}

The February 1967 confrontation between United Headquarters and Red
Storm was the major contributing factor to the bitter enmity which
characterized relations between the two organizations from that time
onward. Although United Headquarters assumed the mantle of ``rebel''
organization while Red Storm henceforth bore the opprobrious title
``conservative'' or ``royalist'', as Walder has pointed out these names
derived from a practical or situational orientation and depended on the
specific political situation in each locality.\footnote{Walder,
  ``Cultural Revolution Radicalism''.

  Walder, ``Cultural Revolution Radicalism''。} One of the major
difficulties in assigning labels to the two mass organizations in
Zhejiang is the fact that they were both extreme, a point later
commented on by Zhou Enlai. The principal source of their differences
seemed to hinge on their respective attitudes toward Jiang Hua. This may
not have been the origin of mass factionalism in Zhejiang, but it was
certainly the major factor behind the tension that continued over the
next three years and made the task of administering the province an
onerous one.

1967年2月联合总部与红暴之间的对抗，是此后两组织关系以深刻敌意为特征的主要原因。尽管联合总部承担了``造反派''组织的外衣，而红暴此后背负``保守派''或``保皇派''的恶名，但正如沃尔德所指出的，这些名称源于实际或情境取向，并取决于各地具体政治形势。\footnote{Walder,
  ``Cultural Revolution Radicalism''.

  Walder, ``Cultural Revolution Radicalism''。}
给浙江这两个群众组织贴标签的一大困难在于，它们二者都很极端，这一点后来周恩来也曾评论。二者差异的主要来源，似乎取决于它们各自对江华的态度。这未必是浙江群众派性产生的源头，但它无疑是此后三年持续紧张背后的主要因素，并使治理全省成为一项艰巨任务。

Outside interference did not help matters. Two Qinghua University
student emissaries from the Cultural Revolution Group in Beijing, Lin
Gang (林岗) and Han Xiangdong, encouraged the rebels of United
Headquarters to persist in their attacks on local leaders despite the
protection afforded Jiang Hua at the center. It is probably true to say
that if the Cultural Revolution radicals had not coined the tags
``proletarian revolutionary'' and ``conservative'' the logic of the
situation would have required them to do so. All revolutions require two
opposing sides -- the revolutionary and the reactionary. In Zhejiang's
case, ideological correctness and opportunistic political interests
outweighed fidelity to local reality.

外部干预无助于缓和局势。北京文化革命小组派来的两名清华大学学生使者林岗和韩向东，鼓励联合总部造反派继续攻击地方领导，尽管江华在中央受到了保护。可以说，即使文化大革命激进派没有创造出``无产阶级革命派''和``保守派''这些标签，形势逻辑也会要求他们这样做。一切革命都需要两个对立面：革命者与反动派。在浙江的情形中，意识形态正确性和机会主义政治利益压倒了对地方现实的忠实。

\section{Notes}\label{notes-1}

\section{注释}\label{ux6ce8ux91ca-1}

\bookmarksetup{startatroot}

\chapter{Chapter Two - DISUNITY AND VIOLENCE, FEBRUARY-AUGUST 1967 /
第二章 -
分裂与暴力，1967年2月至8月}\label{chapter-two---disunity-and-violence-february-august-1967-ux7b2cux4e8cux7ae0---ux5206ux88c2ux4e0eux66b4ux529b1967ux5e742ux6708ux81f38ux6708}

\section{Disunity, February-May 1967}\label{disunity-february-may-1967}

\section{分裂，1967年2月至5月}\label{ux5206ux88c21967ux5e742ux6708ux81f35ux6708}

The rebels in Zhejiang had failed to seize power; disunity among their
ranks had contributed greatly to this failure. In addition to the split
between United Headquarters and Red Storm a further threatening
development had materialized at the end of January -- the organization
of a group calling itself the Red Defence Army of Zhejiang Provincial
Worker-Peasant-Demobilized Soldiers' Revolutionary Rebel General
Headquarters (浙江省工农复转军人革命造反总部红卫军). On 20 January this
organization seized the seals from the Xiaoshan County Public Security
Bureau, located twenty kilometers or so from Hangzhou. On the following
morning, the group repeated the operation in Hangzhou. Two urgent
notices of 23 and 25 January reflected the gravity of the situation and
the alarm felt by the rebels of United Headquarters.\footnote{ZJRB,
  January 23,1967, p.~1, January 25,1967, p.~1 and January 29,1967,
  p.~3.

  ZJRB，1967年1月23日，第1页；1967年1月25日，第1页；1967年1月29日，第3页。}

浙江的造反派未能夺取权力；其内部的不团结极大促成了这一失败。除联合总部与红暴之间的分裂外，1月底又出现了一个更具威胁的新动向，即一个自称``浙江省工农复转军人革命造反总部红卫军''的组织成立。1月20日，该组织夺取了距杭州约二十公里的萧山县公安局印章。次日上午，这一集团又在杭州重复了同样行动。1月23日和25日的两份紧急通告反映出局势的严重性，也反映出联合总部造反派所感到的惊恐。\footnote{ZJRB,
  January 23,1967, p.~1, January 25,1967, p.~1 and January 29,1967,
  p.~3.

  ZJRB，1967年1月23日，第1页；1967年1月25日，第1页；1967年1月29日，第3页。}

The Zhejiang Daily published editorials on 5 and 10 February calling for
unity and warning against sabotage from the right and the
ultra-left.\footnote{ZPS, February 5,1967, SWB/FE/2389/B/15-17; ZJRB,
  February 10,1967, p.~2.

  ZPS，1967年2月5日，SWB/FE/2389/B/15-17；ZJRB，1967年2月10日，第2页。}
In a swipe at Red Storm, the 5 February editorial rejected any ``motley
alliance'' which was based on ``conciliation, eclecticism and
opportunism''. The editorial of 10 February referred to certain people
who

《浙江日报》在2月5日和10日发表社论，呼吁团结，并警告来自右的和极``左''的破坏。\footnote{ZPS,
  February 5,1967, SWB/FE/2389/B/15-17; ZJRB, February 10,1967, p.~2.

  ZPS，1967年2月5日，SWB/FE/2389/B/15-17；ZJRB，1967年2月10日，第2页。}
2月5日的社论暗批红暴，拒绝任何建立在``调和、折衷主义和机会主义''基础上的``大杂烩联盟''。2月10日的社论提到某些人：

\begin{quote}
shot arrows from behind at the revolutionary rebels at a time when the
rebels are rallying to charge the enemy, thus disrupting the order of
battle and creating confusion in our own camp.
\end{quote}

\begin{quote}
当革命造反派正在集合起来向敌人冲锋的时候，从背后向革命造反派放冷箭，从而扰乱战斗秩序，在我们自己的营垒中制造混乱。
\end{quote}

These words gave vent to the anger that United Headquarters felt toward
its factional enemies.

这些话宣泄了联合总部对其派性敌人的愤怒。

Despite the animosity, on 14 February United Headquarters and 30 other
rebel groups in Hangzhou issued a proposal for unity.\footnote{ZJRB,
  February 14,1967, p.~1.

  ZJRB，1967年2月14日，第1页。} In the aftermath of the 12 February
showdown, however, the prospects seemed remote. The proposal advocated
the establishment of a preparatory group which would make suggestions on
how to convene a congress of the workers of Hangzhou. This congress
would unite the city's workers as the first step in achieving the unity
of workers, peasants, cadres and students across the province. The
proposal also recognized that without the support of the PLA all talk of
power seizures was idle and tacit admission of the military's failure to
support the unseating of Jiang Hua.

尽管敌意深重，2月14日，联合总部和杭州另外30个造反组织仍发表了一项团结倡议。\footnote{ZJRB,
  February 14,1967, p.~1.

  ZJRB，1967年2月14日，第1页。}
然而，在2月12日对峙之后，前景显得渺茫。该倡议主张成立一个筹备小组，就如何召开杭州市工人代表大会提出建议。这个大会将团结全市工人，作为实现全省工人、农民、干部和学生团结的第一步。倡议还承认，没有人民解放军支持，所有夺权言论都是空谈；这也等于默认军方未能支持推翻江华。

Continued squabbles among United Headquarters and Red Storm, the
solidarity of the provincial leaders, and the reluctance of the local
military to support United Headquarters, finally snapped the patience of
the central authorities. On March 19 the CC Military Affairs Commission
(MAC) notified the country of its decision to order the PLA to ``support
the left, industry and agriculture and implement military control and
military training'' (the ``three supports and two
militaries'').\footnote{Zhao Wei, Zhao Ziyang zhuan, p.~160.

  赵蔚，Zhao Ziyang zhuan，第160页。} Subsequently, it ordered the
troops of the 20th Army (Unit 6409) stationed in Jiangsu province to
move into Zhejiang to assume military control of the province in
conjunction with the ZPMD.\footnote{Unit 6409 certainly arrived in
  Zhejiang before May 1967, the date given by Harvey W. Nelsen, The
  Chinese Military System: An Organizational Study of the Chinese PLA
  (Boulder Co.: Westview Press, 1977), p.~130, and well before late 1967
  as claimed by the editor of Current Scene. ``Revolutionary Committee
  Leadership - China's Current Provincial Authorities'', CS, 6:18
  (October 18,1968), p.~3. The presence in Hangzhou of the Political
  Commissar of Unit 6409, Nan Ping, was mentioned by the provincial
  daily newspaper as early as March 1967. See ZJRB, March 27, 1967,
  p.~1.

  6409部队肯定在1967年5月以前已经抵达浙江，早于 Harvey W. Nelsen 在 The
  Chinese Military System: An Organizational Study of the Chinese
  PLA（Boulder Co.: Westview Press，1977年）第130页所列日期，也远早于
  Current Scene 编辑所称的1967年末，见 ``Revolutionary Committee
  Leadership - China's Current Provincial
  Authorities''，CS，6:18（1968年10月18日），第3页。省级日报早在1967年3月就提到6409部队政委南萍在杭州。见
  ZJRB，1967年3月27日，第1页。} The decision to send in the troops was
probably made at one of a series of meetings which took place in
February 1967 in Beijing.

联合总部与红暴之间持续的争执、省级领导人的团结一致以及地方军方不愿支持联合总部，最终耗尽了中央当局的耐心。3月19日，中共中央军事委员会通知全国，决定命令人民解放军``支左、支工、支农、军管、军训''（即``三支两军''）。\footnote{Zhao
  Wei, Zhao Ziyang zhuan, p.~160.

  赵蔚，Zhao Ziyang zhuan，第160页。}
随后，中央命令驻江苏的第二十军（6409部队）开入浙江，与浙江省军区共同对该省实行军事管制。\footnote{Unit
  6409 certainly arrived in Zhejiang before May 1967, the date given by
  Harvey W. Nelsen, The Chinese Military System: An Organizational Study
  of the Chinese PLA (Boulder Co.: Westview Press, 1977), p.~130, and
  well before late 1967 as claimed by the editor of Current Scene.
  ``Revolutionary Committee Leadership - China's Current Provincial
  Authorities'', CS, 6:18 (October 18,1968), p.~3. The presence in
  Hangzhou of the Political Commissar of Unit 6409, Nan Ping, was
  mentioned by the provincial daily newspaper as early as March 1967.
  See ZJRB, March 27, 1967, p.~1.

  6409部队肯定在1967年5月以前已经抵达浙江，早于 Harvey W. Nelsen 在 The
  Chinese Military System: An Organizational Study of the Chinese
  PLA（Boulder Co.: Westview Press，1977年）第130页所列日期，也远早于
  Current Scene 编辑所称的1967年末，见 ``Revolutionary Committee
  Leadership - China's Current Provincial
  Authorities''，CS，6:18（1968年10月18日），第3页。省级日报早在1967年3月就提到6409部队政委南萍在杭州。见
  ZJRB，1967年3月27日，第1页。}
派兵入浙的决定，很可能是在1967年2月北京一系列会议中的某次会议上作出的。

The most important of the sessions, those of 14 and 16 February,
witnessed heated clashes between senior military leaders of the CC MAC
and administrators of the State Council on the one side, and members of
the Cultural Revolution Group on the other. In the middle sat an
impassive Zhou Enlai and an indecisive Mao,\footnote{Only one reference
  mentions the presence of Mao, at the first meeting of February 13. See
  Suo Guoxin, 1967年的78天--``二月逆流''记实 (78 days in 1967 -- The
  True Story of the ``February Countercurrent'') (Changsha: Hunan Wenyi
  Chubanshe, 1986), transl. in parts in Keith Forster (ed.), Chinese Law
  and Government, 22: 1 (Spring 1989).

  只有一项资料提到毛出席了2月13日的第一次会议。见 Suo
  Guoxin，1967年的78天--``二月逆流''记实（78 days in 1967 -- The True
  Story of the ``February Countercurrent''）（长沙：Hunan Wenyi
  Chubanshe，1986年），部分译文收入 Keith Forster 编 Chinese Law and
  Government，22:1（1989年春）。} responding to arguments and pressure
from both sides. The crux of the debate was whether to press on with the
Cultural Revolution --- with all the risks, including the possibility of
civil war --- or to rein in or even abort the experiment. Pressure from
those who argued for the latter option was later referred to as the
February Countercurrent (二月逆流) and those involved severely
criticized, publicly denounced, and, in the case of Tan Zhenlin, cruelly
beaten.\footnote{The most detailed and authoritative descriptions of
  these meetings have been written by Ji Xizhen, RMRB, February 26,1979,
  pp.~2, 4; and the same author's ``二月逆流始末记'' (The beginning and
  the end of the ``February Countercurrent''), in Zhou Ming (ed.),
  历史在这里沉思: 1966-1976年纪实 (Reflections on History: True Stories
  from the Years 1966-1976), Vol. 2 (Beijing: Huaxia Chubanshe, 1986)
  pp.~52-87, transl. in part in CR/PSMA 105 (August 15,1980), pp.~17-62;
  Suo Guoxin, 1967niande 78 tian; 聂荣臻回忆录 (Reminiscences of Nie
  Rongzhen), Vol. 3 (Beijing: Jiefangjun Chubanshe, 1984), pp.~852-61;
  Xu Xiangqian, 历史的回顾 (Looking back on History), (Beijing:
  Jiefangjun Chubanshe, 1987), pp.~818-50; Tong De, ``力挽狂澜正气凛然''
  (Fighting to turn the tide with awe-inspiring righteousness), 星火燎原
  (A Single Spark Can Start a Prairie Fire), No.~6 (1984), pp.~23-31.

  关于这些会议最详尽且权威的记述，见 Ji
  Xizhen，RMRB，1979年2月26日，第2、4页；同一作者《二月逆流始末记》（The
  beginning and the end of the ``February Countercurrent''），载 Zhou
  Ming 编《历史在这里沉思：1966-1976年纪实》（Reflections on History:
  True Stories from the Years 1966-1976）第2卷（北京：Huaxia
  Chubanshe，1986年），第52-87页，部分译文见 CR/PSMA
  105（1980年8月15日），第17-62页；Suo Guoxin，1967niande 78
  tian；《聂荣臻回忆录》（Reminiscences of Nie
  Rongzhen）第3卷（北京：Jiefangjun Chubanshe，1984年），第852-861页；Xu
  Xiangqian，《历史的回顾》（Looking back on History）（北京：Jiefangjun
  Chubanshe，1987年），第818-850页；Tong
  De，《力挽狂澜正气凛然》（Fighting to turn the tide with awe-inspiring
  righteousness），《星火燎原》（A Single Spark Can Start a Prairie
  Fire），1984年第6期，第23-31页。} The intervention of outside military
forces in Zhejiang may have helped to stabilize a chaotic situation, but
it also exacerbated tensions between the two mass organizations,
ensnared both local military forces and outside troops in the dispute,
and thereby brought about tense relations between local and outside PLA
units. The Zhejiang Provincial Military Control Commission
(浙江省军事管制委员会), which was set up either in March 1967\footnote{See
  HZRB, September 12,1979, p.~1. The Jiangsu Provincial Military
  Commission was established on March 10. See Jiangsusheng dashiji,
  p.~286. The Guangdong commission was established on March 15. See Zhao
  Wei, Zhao Ziyang zhuan, pp.~161-2.

  见 HZRB，1979年9月12日，第1页。江苏省军事委员会于3月10日成立。见
  Jiangsusheng dashiji，第286页。广东委员会于3月15日成立。见 Zhao
  Wei，Zhao Ziyang zhuan，第161-162页。} or more probably, given the
authoritative nature of the source, on 24 April,\footnote{See Zhejiang
  shengqing, p.~1052. The Control Commission set up a Production
  Committee to take charge of the provincial economy. David Goodman has
  incorrectly dated the establishment of the Zhejiang Military Control
  Commission from January 1967 when the CCP CC issued its document on
  the incident at the headquarters of the ZPMD, cited in fn. 95, ch.~1.
  David S. G. Goodman, ``The Provincial Revolutionary Committee in the
  People's Republic of China, 1967-1979: An Obituary'', CQ, 85 (1981),
  pp.~56, 57.

  见 Zhejiang
  shengqing，第1052页。军管会设立生产委员会，负责省内经济。David Goodman
  将浙江省军管会的成立错误地定为1967年1月，即中共中央发布关于浙江省军区司令部事件的文件之时，该文件见第1章脚注95。David
  S. G. Goodman，``The Provincial Revolutionary Committee in the
  People's Republic of China, 1967-1979: An
  Obituary''，CQ，85（1981年），第56、57页。} became the focus of
contention for control of the province between the ZPMD leaders and
outside military forces. According to its leaders, Unit 7350 (the 5th
Air Force), which was stationed in northern Zhejiang, had responded to
Mao's order to support the left in January 1967, and had most probably
played a major political role in the province from that time
onward.\footnote{See Chen Liyun (陈励耘) and Bai Zongshan (白宗善)
  (Political Commissar and Commander respectively of this unit),
  ``以斗私批修为纲, 坚决支持无产阶级革命派'' (Take ``struggle against
  self and criticize revisionism'' as the crux, and firmly support the
  proletarian revolutionaries), RMRB, October 25,1967, p.~2. Chen Liyun
  had been present at an August 1966 rally addressed by Jiang Hua. See
  ZJRB, August 20, 1966, pp.~1, 2.

  见 Chen Liyun（陈励耘）和 Bai
  Zongshan（白宗善）（分别为该部队政委和司令员），《以斗私批修为纲，坚决支持无产阶级革命派》（Take
  ``struggle against self and criticize revisionism'' as the crux, and
  firmly support the proletarian
  revolutionaries），RMRB，1967年10月25日，第2页。陈励耘曾出席1966年8月由江华讲话的一次集会。见
  ZJRB，1966年8月20日，第1、2页。} Outwardly, the immediate response
from the rebels to the decision committing more PLA forces to the
Cultural Revolution was a hearty welcome. A joint editorial published by
Zhejiang Daily and Hangzhou Daily on 17 February\footnote{ZJRB, February
  17,1967, p.~3.

  ZJRB，1967年2月17日，第3页。} expressed these sentiments. It seems
that United Headquarters had by this time received recognition from
Beijing as the authorized representative of proletarian revolutionaries
in Zhejiang. However, the editorial gave the misleading impression that
United Headquarters possessed the unanimous and total endorsement of the
ZPMD and all three branches of the centrally-directed PLA forces. This
was certainly not the case. The ZPMD continued to throw its support
behind Red Storm and the position of the naval units of the East China
fleet was at best ambiguous. Behind the ZPMD stood the might of the
Nanjing Military Region and its redoubtable commander Xu Shiyou, as well
as the 22nd Army stationed on Zhoushan Island. Nevertheless, United
Headquarters could count upon the backing of the powerful 20th Army and
the 5th Air Force, whose commanders were eager to display fealty to
their commander, Lin Biao.

其中最重要的是2月14日和16日的会议。会上，一方是中央军委高级军事领导人和国务院行政官员，另一方是中央文革小组成员，双方发生激烈冲突。周恩来态度冷静，毛泽东犹疑不决，\footnote{Only
  one reference mentions the presence of Mao, at the first meeting of
  February 13. See Suo Guoxin, 1967年的78天--``二月逆流''记实 (78 days
  in 1967 -- The True Story of the ``February Countercurrent'')
  (Changsha: Hunan Wenyi Chubanshe, 1986), transl. in parts in Keith
  Forster (ed.), Chinese Law and Government, 22: 1 (Spring 1989).

  只有一项资料提到毛出席了2月13日的第一次会议。见 Suo
  Guoxin，1967年的78天--``二月逆流''记实（78 days in 1967 -- The True
  Story of the ``February Countercurrent''）（长沙：Hunan Wenyi
  Chubanshe，1986年），部分译文收入 Keith Forster 编 Chinese Law and
  Government，22:1（1989年春）。}
坐在中间回应两边的论辩和压力。争论的核心在于，是继续推进文化大革命，承担包括内战可能性在内的一切风险，还是收束乃至中止这一试验。主张后一种选择者施加的压力，后来被称为``二月逆流''，参与者受到严厉批判和公开谴责；谭震林甚至遭到残酷殴打。\footnote{The
  most detailed and authoritative descriptions of these meetings have
  been written by Ji Xizhen, RMRB, February 26,1979, pp.~2, 4; and the
  same author's ``二月逆流始末记'' (The beginning and the end of the
  ``February Countercurrent''), in Zhou Ming (ed.), 历史在这里沉思:
  1966-1976年纪实 (Reflections on History: True Stories from the Years
  1966-1976), Vol. 2 (Beijing: Huaxia Chubanshe, 1986) pp.~52-87,
  transl. in part in CR/PSMA 105 (August 15,1980), pp.~17-62; Suo
  Guoxin, 1967niande 78 tian; 聂荣臻回忆录 (Reminiscences of Nie
  Rongzhen), Vol. 3 (Beijing: Jiefangjun Chubanshe, 1984), pp.~852-61;
  Xu Xiangqian, 历史的回顾 (Looking back on History), (Beijing:
  Jiefangjun Chubanshe, 1987), pp.~818-50; Tong De, ``力挽狂澜正气凛然''
  (Fighting to turn the tide with awe-inspiring righteousness), 星火燎原
  (A Single Spark Can Start a Prairie Fire), No.~6 (1984), pp.~23-31.

  关于这些会议最详尽且权威的记述，见 Ji
  Xizhen，RMRB，1979年2月26日，第2、4页；同一作者《二月逆流始末记》（The
  beginning and the end of the ``February Countercurrent''），载 Zhou
  Ming 编《历史在这里沉思：1966-1976年纪实》（Reflections on History:
  True Stories from the Years 1966-1976）第2卷（北京：Huaxia
  Chubanshe，1986年），第52-87页，部分译文见 CR/PSMA
  105（1980年8月15日），第17-62页；Suo Guoxin，1967niande 78
  tian；《聂荣臻回忆录》（Reminiscences of Nie
  Rongzhen）第3卷（北京：Jiefangjun Chubanshe，1984年），第852-861页；Xu
  Xiangqian，《历史的回顾》（Looking back on History）（北京：Jiefangjun
  Chubanshe，1987年），第818-850页；Tong
  De，《力挽狂澜正气凛然》（Fighting to turn the tide with awe-inspiring
  righteousness），《星火燎原》（A Single Spark Can Start a Prairie
  Fire），1984年第6期，第23-31页。}
外来军事力量介入浙江，或许有助于稳定混乱局面，但也加剧了两个群众组织之间的紧张，把地方军队和外来部队都卷入争端，并由此造成地方与外来解放军部队之间的紧张关系。浙江省军事管制委员会可能成立于1967年3月，\footnote{See
  HZRB, September 12,1979, p.~1. The Jiangsu Provincial Military
  Commission was established on March 10. See Jiangsusheng dashiji,
  p.~286. The Guangdong commission was established on March 15. See Zhao
  Wei, Zhao Ziyang zhuan, pp.~161-2.

  见 HZRB，1979年9月12日，第1页。江苏省军事委员会于3月10日成立。见
  Jiangsusheng dashiji，第286页。广东委员会于3月15日成立。见 Zhao
  Wei，Zhao Ziyang zhuan，第161-162页。}
或者更可能鉴于资料权威性，成立于4月24日；\footnote{See Zhejiang
  shengqing, p.~1052. The Control Commission set up a Production
  Committee to take charge of the provincial economy. David Goodman has
  incorrectly dated the establishment of the Zhejiang Military Control
  Commission from January 1967 when the CCP CC issued its document on
  the incident at the headquarters of the ZPMD, cited in fn. 95, ch.~1.
  David S. G. Goodman, ``The Provincial Revolutionary Committee in the
  People's Republic of China, 1967-1979: An Obituary'', CQ, 85 (1981),
  pp.~56, 57.

  见 Zhejiang
  shengqing，第1052页。军管会设立生产委员会，负责省内经济。David Goodman
  将浙江省军管会的成立错误地定为1967年1月，即中共中央发布关于浙江省军区司令部事件的文件之时，该文件见第1章脚注95。David
  S. G. Goodman，``The Provincial Revolutionary Committee in the
  People's Republic of China, 1967-1979: An
  Obituary''，CQ，85（1981年），第56、57页。}
它成为浙江省军区领导人和外来军事力量争夺全省控制权的焦点。按照该部队领导人的说法，驻扎在浙北的7350部队（第五空军）在1967年1月响应了毛的支左命令，并很可能自那时起就在全省政治中发挥了重要作用。\footnote{See
  Chen Liyun (陈励耘) and Bai Zongshan (白宗善) (Political Commissar and
  Commander respectively of this unit), ``以斗私批修为纲,
  坚决支持无产阶级革命派'' (Take ``struggle against self and criticize
  revisionism'' as the crux, and firmly support the proletarian
  revolutionaries), RMRB, October 25,1967, p.~2. Chen Liyun had been
  present at an August 1966 rally addressed by Jiang Hua. See ZJRB,
  August 20, 1966, pp.~1, 2.

  见 Chen Liyun（陈励耘）和 Bai
  Zongshan（白宗善）（分别为该部队政委和司令员），《以斗私批修为纲，坚决支持无产阶级革命派》（Take
  ``struggle against self and criticize revisionism'' as the crux, and
  firmly support the proletarian
  revolutionaries），RMRB，1967年10月25日，第2页。陈励耘曾出席1966年8月由江华讲话的一次集会。见
  ZJRB，1966年8月20日，第1、2页。}
表面上，造反派对投入更多解放军力量参加文化大革命的决定热烈欢迎。2月17日《浙江日报》和《杭州日报》联合发表的社论\footnote{ZJRB,
  February 17,1967, p.~3.

  ZJRB，1967年2月17日，第3页。}
表达了这种情绪。看起来，联合总部此时已经得到北京承认，被视为浙江无产阶级革命派的授权代表。然而，该社论给人一种误导性印象，似乎联合总部得到了浙江省军区以及中央直接指挥的解放军三大军种的一致和完全支持。事实当然并非如此。浙江省军区继续支持红暴，东海舰队海军部队的立场至多也是暧昧的。站在浙江省军区背后的是南京军区的力量及其令人敬畏的司令员许世友，还有驻舟山岛的第二十二军。不过，联合总部可以依靠强大的第二十军和第五空军支持，这些部队的指挥员急于向自己的统帅林彪表示忠诚。

Official recognition of United Headquarters certainly represented a
set-back for Red Storm. But the knowledge that the ZPMD and its
superiors in Nanjing had not deserted it was a great consolation. The
failure of United Headquarters to seize power encouraged Red Storm to
fight on, confident that its cause had a future. As discussed above, it
is difficult to discover substantial differences between the memberships
of the two mass organizations in terms of class background, occupational
category or other meaningful social criteria. The struggle for power and
survival dictated the logic of the polarization.

联合总部得到正式承认，无疑是红暴的一次挫折。但知道浙江省军区及其南京上级并未抛弃自己，对红暴来说是一大安慰。联合总部夺权失败，鼓励红暴继续斗争，并相信自己的事业仍有前途。如前所述，从阶级出身、职业类别或其他有意义的社会标准来看，很难发现两个群众组织成员之间存在实质差异。权力与生存之争决定了两极分化的逻辑。

Both organizations built up a network of subordinate branches in the
counties and towns of Zhejiang.\footnote{For example United Headquarters
  established a branch in Xiangshan county on April 19, 1967, and Red
  Storm a preparatory committee on July 14. See 象山县志编纂委员会
  (Xiangshan County Gazetteer Compilation Committee), 象山县志
  (Xiangshan County Gazetteer), (Hangzhou: Zhejiang Renmin Chubanshe,
  1988), p.~34. On June 12, United Headquarters set up a branch in Lanxi
  county. See 兰溪市志编纂委员会 (Lanxi City Gazetteer Compilation
  Committee), 兰溪市志 (Lanxi City Gazetteer), (Hangzhou: Zhejiang
  Renmin Chubanshe, 1988), p.~19. In Kaihua county the two factions
  established headquarters in the spring and summer of 1967. See
  开化县志编纂委员会 (Kaihua County Gazetteer Compilation Committee),
  开化县志 (Kaihua County Gazetteer), (Hangzhou: Zhejiang Renmin
  Chubanshe, 1988), p.~23.

  例如，联合总部于1967年4月19日在象山县成立分部，红暴于7月14日成立筹备委员会。见《象山县志》编纂委员会（Xiangshan
  County Gazetteer Compilation Committee），《象山县志》（Xiangshan
  County Gazetteer）（杭州：Zhejiang Renmin
  Chubanshe，1988年），第34页。6月12日，联合总部在兰溪县设立分部。见《兰溪市志》编纂委员会（Lanxi
  City Gazetteer Compilation Committee），《兰溪市志》（Lanxi City
  Gazetteer）（杭州：Zhejiang Renmin
  Chubanshe，1988年），第19页。在开化县，两派于1967年春夏建立总部。见《开化县志》编纂委员会（Kaihua
  County Gazetteer Compilation Committee），《开化县志》（Kaihua County
  Gazetteer）（杭州：Zhejiang Renmin Chubanshe，1988年），第23页。} They
also established contacts and liaison stations in Shanghai (in
particular) and Beijing. Central rebel organizations and politicians in
turn cultivated links in the provinces to bolster their positions, but
lacked the means to impose the kind of organizational discipline that
the CCP had always wielded over its subordinate branches. Thus, the
growth of undirected anomic factionalism proceeded unchecked. Factional
struggles were clearly upsetting the normal functioning of the economy,
the agricultural sector in particular.\footnote{See the urgent notice
  issued by United Headquarters on February 14,1967, in RMRB, February
  16,1967, in SCMP, 3887 (February 27,1967), pp.~10-13.

  见联合总部1967年2月14日发布的紧急通知，载 RMRB，1967年2月16日，转载于
  SCMP，3887（1967年2月27日），第10-13页。} With the pre-Cultural
Revolution power structure in complete disarray it was left to the ZPMD
and United Headquarters to convene a conference in agricultural
production at the end of February.\footnote{ZPS, February 28, 1967,
  SWB/FE/2406/B/24-5; ZPS, March 2, 1967, SWB/FE/2410/B/2-3; Domes,
  ``The Role of the Military'', pp.130-1. Continued and urgent appeals
  from the military illustrated the difficulty in securing the planting
  of spring rice. See ZPS, March 3, 1967, SWB/FE/2413/B/13-15 and ZPS,
  March 16, 1967, SWB/FE/2422/B/17. Between March 22 and 25, the
  military district ``production leading group'' (acting for the defunct
  provincial People's Council) convened a conference on industry,
  finance and trade and put forward an economic plan for 1967. Zhejiang
  shengqing, p.~1052. In the counties, People's Armed Forces Departments
  established similar groups to organize spring sowing. See Kaihua
  xianzhi, p.~23; 临海市志编纂委员会编 (Linhai Municipality Gazetteer
  Compilation Committee), 临海县志 (Linhai County Gazetteer), (Hangzhou:
  Zhejiang Renmin Chubanshe, 1989), p.~28.

  ZPS，1967年2月28日，SWB/FE/2406/B/24-5；ZPS，1967年3月2日，SWB/FE/2410/B/2-3；Domes，``The
  Role of the
  Military''，第130-131页。军方持续且紧急的呼吁说明，确保春稻种植存在困难。见
  ZPS，1967年3月3日，SWB/FE/2413/B/13-15；ZPS，1967年3月16日，SWB/FE/2422/B/17。3月22日至25日，军区``生产领导小组''（代行已停止运作的省人民委员会职能）召开工业、财贸会议，并提出1967年经济计划。Zhejiang
  shengqing，第1052页。在各县，人民武装部成立类似小组组织春播。见 Kaihua
  xianzhi，第23页；临海市志编纂委员会编（Linhai Municipality Gazetteer
  Compilation Committee），《临海县志》（Linhai County
  Gazetteer）（杭州：Zhejiang Renmin Chubanshe，1989年），第28页。} In
the meanwhile moves were in train to bring about an alliance of
Hangzhou's worker rebels, except for those belonging to Red Storm. An
initial breakthrough was achieved on 26 March when a meeting to forge
such an alliance opened in the city.\footnote{ZPS, March 20, 1967,
  SWB/FE/2430/B/28-30; ZJRB, March 27,1967, pp.~1, 4.

  ZPS，1967年3月20日，SWB/FE/2430/B/28-30；ZJRB，1967年3月27日，第1、4页。}
The congress lasted four days and was attended and addressed by leading
military officials including the newly-arrived Political Commissar of
Unit 6409, Nan Ping, and ZPMD Political Commissar Long Qian.\footnote{ZJRB,
  March 30,1967, p.~1. For Long's speech see ZJRB, March 27,1967, p.~3.

  ZJRB，1967年3月30日，第1页。龙的讲话见 ZJRB，1967年3月27日，第3页。}
The organization which emerged from these deliberations was called the
Hangzhou Revolutionary Committee of Workers and Staff Members
(杭州市革命职工委员会). He Xianchun from United Headquarters emerged as
one of the leaders (probably Chairman) of the committee and the head of
its fighting division (战处长) was Xia Genfa.\footnote{See HZRB, April
  17, 1979.

  见 HZRB，1979年4月17日。} The main stumbling blocks to a provincial
alliance remained United Headquarters' insistence that it form the
``core'' of the revolutionary alliance, and Red Storm's equally strong
determination that it enter on its own terms as an equal partner.
However, an editorial in Zhejiang Daily of 21 April firmly rejected any
concessions to or reconciliation with ``conservative'' organizations. It
insisted that the ``staunch revolutionary left'' {[}United
Headquarters{]} form the nucleus of a militant alliance.\footnote{ZPS,
  April 20, 1967, SWB/FE/2449/B/12-14; see also the editorial in ZJRB,
  April 24, 1967, p.~2.

  ZPS，1967年4月20日，SWB/FE/2449/B/12-14；另见 ZJRB
  社论，1967年4月24日，第2页。} Although not referred to by name, Red
Storm was accused of being manipulated by the ``bourgeois reactionary
line''. Proponents of this position, argued the editorial, had advanced
the ``reactionary theory of family lineage'' to ``deceive and poison the
masses''. This reference is further indication of the ``better'' class
composition of Red Storm and the vulnerability felt by United
Headquarters on the issue. Unnamed ``mediators'', continued the
editorial, had characterized the contradictions between mass
organizations as centering on the issue of power. The editorial rejected
such allegations, but the lengths to which it went in order to disprove
the charge did not make its argument any more persuasive. It appealed to
the hoodwinked members of Red Storm to abandon their leaders and the
organization and rejoin the revolutionary ranks. Such appeals were to be
repeated many times during the coming months, but with little apparent
success.

两个组织都在浙江各县镇建立了下属分支网络。\footnote{For example United
  Headquarters established a branch in Xiangshan county on April 19,
  1967, and Red Storm a preparatory committee on July 14. See
  象山县志编纂委员会 (Xiangshan County Gazetteer Compilation Committee),
  象山县志 (Xiangshan County Gazetteer), (Hangzhou: Zhejiang Renmin
  Chubanshe, 1988), p.~34. On June 12, United Headquarters set up a
  branch in Lanxi county. See 兰溪市志编纂委员会 (Lanxi City Gazetteer
  Compilation Committee), 兰溪市志 (Lanxi City Gazetteer), (Hangzhou:
  Zhejiang Renmin Chubanshe, 1988), p.~19. In Kaihua county the two
  factions established headquarters in the spring and summer of 1967.
  See 开化县志编纂委员会 (Kaihua County Gazetteer Compilation
  Committee), 开化县志 (Kaihua County Gazetteer), (Hangzhou: Zhejiang
  Renmin Chubanshe, 1988), p.~23.

  例如，联合总部于1967年4月19日在象山县成立分部，红暴于7月14日成立筹备委员会。见《象山县志》编纂委员会（Xiangshan
  County Gazetteer Compilation Committee），《象山县志》（Xiangshan
  County Gazetteer）（杭州：Zhejiang Renmin
  Chubanshe，1988年），第34页。6月12日，联合总部在兰溪县设立分部。见《兰溪市志》编纂委员会（Lanxi
  City Gazetteer Compilation Committee），《兰溪市志》（Lanxi City
  Gazetteer）（杭州：Zhejiang Renmin
  Chubanshe，1988年），第19页。在开化县，两派于1967年春夏建立总部。见《开化县志》编纂委员会（Kaihua
  County Gazetteer Compilation Committee），《开化县志》（Kaihua County
  Gazetteer）（杭州：Zhejiang Renmin Chubanshe，1988年），第23页。}
它们还在上海（尤其是上海）和北京建立了联络点和联络站。中央的造反组织和政治人物也反过来培植同各省的联系，以巩固自己的地位，但他们缺乏中共过去对下属分支一贯拥有的那种组织纪律控制手段。因此，无方向、失范的派性主义不受遏制地增长。派性斗争显然扰乱了经济的正常运行，尤其是农业部门。\footnote{See
  the urgent notice issued by United Headquarters on February 14,1967,
  in RMRB, February 16,1967, in SCMP, 3887 (February 27,1967),
  pp.~10-13.

  见联合总部1967年2月14日发布的紧急通知，载 RMRB，1967年2月16日，转载于
  SCMP，3887（1967年2月27日），第10-13页。}
文革前的权力结构已完全混乱，2月底只能由浙江省军区和联合总部召开农业生产会议。\footnote{ZPS,
  February 28, 1967, SWB/FE/2406/B/24-5; ZPS, March 2, 1967,
  SWB/FE/2410/B/2-3; Domes, ``The Role of the Military'', pp.130-1.
  Continued and urgent appeals from the military illustrated the
  difficulty in securing the planting of spring rice. See ZPS, March 3,
  1967, SWB/FE/2413/B/13-15 and ZPS, March 16, 1967, SWB/FE/2422/B/17.
  Between March 22 and 25, the military district ``production leading
  group'' (acting for the defunct provincial People's Council) convened
  a conference on industry, finance and trade and put forward an
  economic plan for 1967. Zhejiang shengqing, p.~1052. In the counties,
  People's Armed Forces Departments established similar groups to
  organize spring sowing. See Kaihua xianzhi, p.~23;
  临海市志编纂委员会编 (Linhai Municipality Gazetteer Compilation
  Committee), 临海县志 (Linhai County Gazetteer), (Hangzhou: Zhejiang
  Renmin Chubanshe, 1989), p.~28.

  ZPS，1967年2月28日，SWB/FE/2406/B/24-5；ZPS，1967年3月2日，SWB/FE/2410/B/2-3；Domes，``The
  Role of the
  Military''，第130-131页。军方持续且紧急的呼吁说明，确保春稻种植存在困难。见
  ZPS，1967年3月3日，SWB/FE/2413/B/13-15；ZPS，1967年3月16日，SWB/FE/2422/B/17。3月22日至25日，军区``生产领导小组''（代行已停止运作的省人民委员会职能）召开工业、财贸会议，并提出1967年经济计划。Zhejiang
  shengqing，第1052页。在各县，人民武装部成立类似小组组织春播。见 Kaihua
  xianzhi，第23页；临海市志编纂委员会编（Linhai Municipality Gazetteer
  Compilation Committee），《临海县志》（Linhai County
  Gazetteer）（杭州：Zhejiang Renmin Chubanshe，1989年），第28页。}
与此同时，杭州工人造反派的联盟也在酝酿之中，但红暴所属者除外。3月26日，旨在建立这种联盟的会议在杭州开幕，初步取得突破。\footnote{ZPS,
  March 20, 1967, SWB/FE/2430/B/28-30; ZJRB, March 27,1967, pp.~1, 4.

  ZPS，1967年3月20日，SWB/FE/2430/B/28-30；ZJRB，1967年3月27日，第1、4页。}
大会持续四天，新到任的6409部队政委南萍和浙江省军区政委龙潜等高级军事官员出席并讲话。\footnote{ZJRB,
  March 30,1967, p.~1. For Long's speech see ZJRB, March 27,1967, p.~3.

  ZJRB，1967年3月30日，第1页。龙的讲话见 ZJRB，1967年3月27日，第3页。}
会上形成的组织名为杭州市革命职工委员会。联合总部的贺贤春成为委员会领导人之一（可能是主任），其战斗部门负责人（战处长）为夏根发。\footnote{See
  HZRB, April 17, 1979.

  见 HZRB，1979年4月17日。}
省级联盟的主要障碍仍然是：联合总部坚持自己必须成为革命联盟的``核心''，红暴则同样坚决要求以平等伙伴身份、按自身条件加入。然而，《浙江日报》4月21日社论坚决拒绝对``保守''组织作任何让步或与其和解，坚持由``坚定的革命左派''{[}联合总部{]}
组成战斗联盟的核心。\footnote{ZPS, April 20, 1967, SWB/FE/2449/B/12-14;
  see also the editorial in ZJRB, April 24, 1967, p.~2.

  ZPS，1967年4月20日，SWB/FE/2449/B/12-14；另见 ZJRB
  社论，1967年4月24日，第2页。}
红暴虽未被点名，却被指控受``资产阶级反动路线''操纵。社论称，这一立场的鼓吹者提出``反动血统论''来``欺骗和毒害群众''。这一提法进一步说明红暴阶级成分``较好''，也说明联合总部在这一问题上感到脆弱。社论继续说，未具名的``调停者''把群众组织之间的矛盾描述为围绕权力问题。社论驳斥了这类指控，但它为否认指控所费的笔墨并未使论证更有说服力。社论呼吁受蒙蔽的红暴成员抛弃其领导人和组织，重新回到革命队伍。这类呼吁在随后几个月多次重复，但显然收效甚微。

The editorial also hinted a concern, well-founded as it eventuated, that
United Headquarters would be ordered to form an alliance by
administrative fiat. It rallied its members to ``resolutely resist''
such pressure. However as early as June 1967 it appears that Beijing was
meeting with representatives of the two mass organizations, as well as
military officials, to force concessions from both sides.\footnote{``中共中央通知''
  (Notice of the CCP CC), June 24,1967, CCP Documents, pp.~465-7. Zhang
  Yongsheng was present in the capital at this time. See Zhang
  Yongshengde zuizheng cailiao, p.~26.

  《中共中央通知》（Notice of the CCP CC），1967年6月24日，CCP
  Documents，第465-467页。张永生当时在首都。见 Zhang Yongshengde
  zuizheng cailiao，第26页。} A rally held three days after the
publication of the above cited editorial of 20 April suggested that the
military did not take a united position on the question of the ``core''
of the revolutionary alliance. The representative of the 5th Air Force
publicly aligned his unit with United Headquarters but another PLA
representative merely proclaimed general support for proletarian
revolutionaries.\footnote{ZPS, April 23, 1967, SWB/FE/2448/B/16.

  ZPS，1967年4月23日，SWB/FE/2448/B/16。} In an article published on 22
April United Headquarters' Zhang Yongsheng, inspired by the formation of
the Beijing Municipal Revolutionary Committee, called for a great
alliance of rebels in the province, but he also excluded ``conservative
factions'' from such an alliance.\footnote{Zhang Yongsheng,
  ``向北京的无产阶级革命派战友学习'' (Learn from the proletarian
  revolutionary comrades-in-arms of Beijing), ZJRB, April 22, 1967,
  p.~2.

  Zhang Yongsheng，《向北京的无产阶级革命派战友学习》（Learn from the
  proletarian revolutionary comrades-in-arms of
  Beijing），ZJRB，1967年4月22日，第2页。} Zhang went to great lengths
in his article to praise the PLA and also cited Jiang Qing's call to
``support the army and love the people'' (拥军爱民). However, in an
unpublished talk of 8 April Zhang had attacked the ``gun-toting Liu-Deng
revisionist line'' and had canvassed the possibility of peaceful
evolution (和平演变) within the CCP.\footnote{Zhang Yongshengde zuizheng
  cailiao, pp.~33, 95.

  Zhang Yongshengde zuizheng cailiao，第33、95页。} At the end of April
high-school students in the county of Zhuji, not far from Hangzhou,
staged a sit-in and hunger strike outside the offices of the People's
Armed Forces Department. Zhang and high-ranking military officers from
Hangzhou arrived in early May to show their support for the
students.\footnote{Ibid, p.~34.

  同上，第34页。} Thus, in his actions and private talks Zhang displayed
hostility toward the military, or at least sections of it, a stance
strikingly at odds with his public utterances.

社论还暗示了一种担忧，即联合总部会被行政命令强制作出联盟；事后证明这种担忧并非没有根据。它动员成员``坚决抵制''这种压力。然而，早在1967年6月，北京似乎已经会见两个群众组织的代表以及军事官员，迫使双方作出让步。\footnote{``中共中央通知''
  (Notice of the CCP CC), June 24,1967, CCP Documents, pp.~465-7. Zhang
  Yongsheng was present in the capital at this time. See Zhang
  Yongshengde zuizheng cailiao, p.~26.

  《中共中央通知》（Notice of the CCP CC），1967年6月24日，CCP
  Documents，第465-467页。张永生当时在首都。见 Zhang Yongshengde
  zuizheng cailiao，第26页。}
上述4月20日社论发表三天后举行的一次集会表明，在革命联盟``核心''问题上，军方并未采取统一立场。第五空军代表公开使本部队站在联合总部一边，而另一名解放军代表只是宣布一般性支持无产阶级革命派。\footnote{ZPS,
  April 23, 1967, SWB/FE/2448/B/16.

  ZPS，1967年4月23日，SWB/FE/2448/B/16。}
联合总部的张永生在4月22日发表文章，受北京市革命委员会成立鼓舞，号召全省造反派实行大联合，但他也把``保守派''排除在这一联盟之外。\footnote{Zhang
  Yongsheng, ``向北京的无产阶级革命派战友学习'' (Learn from the
  proletarian revolutionary comrades-in-arms of Beijing), ZJRB, April
  22, 1967, p.~2.

  Zhang Yongsheng，《向北京的无产阶级革命派战友学习》（Learn from the
  proletarian revolutionary comrades-in-arms of
  Beijing），ZJRB，1967年4月22日，第2页。}
张在文中竭力赞扬人民解放军，并引用江青``拥军爱民''的号召。然而，在4月8日一次未发表的讲话中，张攻击``带枪的刘邓修正主义路线''，并讨论了中共内部发生和平演变的可能性。\footnote{Zhang
  Yongshengde zuizheng cailiao, pp.~33, 95.

  Zhang Yongshengde zuizheng cailiao，第33、95页。}
4月底，距杭州不远的诸暨县中学生在人民武装部办公地点外静坐并绝食。5月初，张和杭州高级军官抵达当地，表示支持学生。\footnote{Ibid,
  p.~34.

  同上，第34页。}
因此，在行动和私下谈话中，张表现出对军队、或至少对其部分成员的敌意，这与他的公开言论形成鲜明矛盾。

\section{Tension and Factional Struggle, June-July
1967}\label{tension-and-factional-struggle-june-july-1967}

\section{紧张与派性斗争，1967年6月至7月}\label{ux7d27ux5f20ux4e0eux6d3eux6027ux6597ux4e891967ux5e746ux6708ux81f37ux6708}

As the summer of 1967 approached, relations between United Headquarters
and sections of the PLA deteriorated. A Zhejiang Daily editorial of 13
May, entitled ``Trust the PLA; rely on the PLA'', asked the rebels to
get right behind the army in its ``three supports and two militaries''
assignment. It attributed the main force units' unfamiliarity with local
conditions to the brevity of their time in Zhejiang. Overall, stressed
the editorial, the PLA had achieved much and it exhorted the rebels:

随着1967年夏季临近，联合总部与人民解放军部分单位之间的关系恶化。《浙江日报》5月13日发表题为``相信解放军，依靠解放军''的社论，要求造反派全力支持军队执行``三支两军''任务。社论把主力部队不熟悉地方情况归因于他们到浙江时间尚短。社论强调，总体而言，人民解放军已取得很大成绩，并敦促造反派：

\begin{quote}
Do not say anything that will harm the prestige of the PLA. Do not do
anything that will harm the unity between the Army and the people. If we
find that some of the comrades in the PLA have shortcomings or mistakes,
we should, with sincerity and goodwill, present our viewpoints on an
appropriate occasion and in an appropriate way to help them improve
themselves.\footnote{ZJRB, May 13,1967.

  ZJRB，1967年5月13日。}
\end{quote}

\begin{quote}
不要说有损人民解放军威信的话。不要做有损军民团结的事。如果我们发现解放军某些同志有缺点或错误，就应当以诚恳和善意，在适当场合、用适当方式提出自己的意见，帮助他们改进。\footnote{ZJRB,
  May 13,1967.

  ZJRB，1967年5月13日。}
\end{quote}

The delicacy of the phrasing could not obscure the fact that United
Headquarters had reason to be less than fully satisfied with the backing
it had received from the military up until this time.

措辞虽谨慎，却无法掩盖这样一个事实：联合总部有理由对迄今为止从军方得到的支持并不完全满意。

In a reversal of the MAC order of 6 April, which had deprived local
military officials of a great deal of their civil authority and had led
to the rejuvenation of suppressed mass organizations,\footnote{Liu
  Guokai, A Brief Analysis, p.~69.

  Liu Guokai，A Brief Analysis，第69页。} on 6 June Beijing directed the
rebels around the country to stop fighting and empowered the PLA to
enforce this edict.\footnote{``中共中央国务院中央军委中央文革小组严禁武斗行凶非法逮捕及抢劫破坏的通令''
  (Order of the CCP CC, State Council, Military Affairs Commission and
  Cultural Revolution Group concerning the strict prohibition against
  armed struggle, illegal arrest, looting and sabotage), June 6, 1967,
  CCP Documents, pp.~461-4.

  《中共中央国务院中央军委中央文革小组严禁武斗行凶非法逮捕及抢劫破坏的通令》（Order
  of the CCP CC, State Council, Military Affairs Commission and Cultural
  Revolution Group concerning the strict prohibition against armed
  struggle, illegal arrest, looting and sabotage），1967年6月6日，CCP
  Documents，第461-464页。} United Headquarters promptly convened a huge
rally, attended by PLA representatives and a crowd estimated at 200,000,
to express its support for the order.\footnote{ZPS, June 11, 1967,
  SWB/FE/2490/B/7-8.

  ZPS，1967年6月11日，SWB/FE/2490/B/7-8。} On that same day, one of the
most sensational and significant events of the Cultural Revolution in
Zhejiang took place.\footnote{The following account comes from Zhang
  Yongsheng bufen zuizheng cailiao, pp.~38-40, and oral sources.

  以下叙述来自 Zhang Yongsheng bufen zuizheng
  cailiao，第38-40页，以及口述资料。} Weng Senhe, the worker-rebel
leader from the Hangzhou silk complex, had by his charisma and
organizational abilities become a major impediment to United
Headquarters' ambitions to achieve hegemony over the rebel movement in
Zhejiang. Beijing's Third Headquarters and United Headquarters therefore
devised a plan to seize Weng in order to try and force him to switch
allegiance from his organization Red Storm. In open defiance of the
central order of 6 June, a large contingent of workers and students set
out for the silk complex and launched a mighty offensive against the
defenders inside. According to the confession written by Du Yingxin over
ten years later, the attack was planned by United Headquarters' leaders
Zhang Yongsheng and He Xianchun, together with the Third Headquarters'
representatives in Hangzhou, Lin Gang and Han Xiangdong, at a meeting
which they held at Hangzhou University on the 4th or 5th of June. The
battle ended with the defeat of Weng's organization. Weng himself was
seized and abducted and the factory was left looking like a battlefield.
Production ceased for almost two weeks and residents of Hangzhou flocked
to the scene to witness for themselves evidence of the epic struggle and
the blood spilt by Weng and his comrades.

4月6日的军委命令曾剥夺地方军官大量民事权力，并导致被压制的群众组织复活；\footnote{Liu
  Guokai, A Brief Analysis, p.~69.

  Liu Guokai，A Brief Analysis，第69页。}
6月6日，北京转而命令全国造反派停止武斗，并授权人民解放军执行这一命令。\footnote{``中共中央国务院中央军委中央文革小组严禁武斗行凶非法逮捕及抢劫破坏的通令''
  (Order of the CCP CC, State Council, Military Affairs Commission and
  Cultural Revolution Group concerning the strict prohibition against
  armed struggle, illegal arrest, looting and sabotage), June 6, 1967,
  CCP Documents, pp.~461-4.

  《中共中央国务院中央军委中央文革小组严禁武斗行凶非法逮捕及抢劫破坏的通令》（Order
  of the CCP CC, State Council, Military Affairs Commission and Cultural
  Revolution Group concerning the strict prohibition against armed
  struggle, illegal arrest, looting and sabotage），1967年6月6日，CCP
  Documents，第461-464页。}
联合总部立即召开有解放军代表参加、据估约二十万人出席的大型集会，表示支持该命令。\footnote{ZPS,
  June 11, 1967, SWB/FE/2490/B/7-8.

  ZPS，1967年6月11日，SWB/FE/2490/B/7-8。}
就在同一天，浙江文化大革命中最轰动、最重要的事件之一发生了。\footnote{The
  following account comes from Zhang Yongsheng bufen zuizheng cailiao,
  pp.~38-40, and oral sources.

  以下叙述来自 Zhang Yongsheng bufen zuizheng
  cailiao，第38-40页，以及口述资料。}
杭州丝绸印染联合厂的工人造反派领袖翁森鹤，凭借个人魅力和组织能力，已成为联合总部争取在浙江造反运动中取得霸权的一大障碍。因此，北京三司和联合总部制定计划抓捕翁，企图迫使他从红暴转投阵营。公然违抗6月6日中央命令的大批工人和学生向丝绸联合厂进发，对厂内守卫者发动猛烈进攻。按照杜英信十多年后所写的交代，攻击是联合总部领导人张永生、贺贤春与北京三司驻杭州代表林岗、韩湘东在6月4日或5日于杭州大学召开的会议上策划的。战斗以翁的组织失败告终。翁本人被抓走，工厂则被打得像战场一样。生产停顿近两周，杭州市民涌向现场，亲眼观看这场史诗般斗争的痕迹以及翁和同志们流下的鲜血。

Weng was held for several days on a boat on the Qiantang river and
finally released after the intercession of the Nanjing Military Region.
He returned to Hangzhou a hero. United Headquarters hoped to make
political capital out of the incident by publishing Weng's extorted
confession but the affair rebounded on it. Red Storm's stocks soared and
it gained new members while United Headquarters was temporarily
discredited. To extricate itself from this situation United Headquarters
held a large rally at Hangzhou's Youth Square on 10 June. Afterwards,
10,000 marchers set out for the north of the city toward the silk
complex. At a bridge en route they clashed with peasants who had come
out in support of Weng's group. Zhang Yongsheng blamed Zhejiang
University for the confrontation and vowed to tear the campus apart.

翁被关押在钱塘江上的一条船上数日，最后经南京军区斡旋获释。他回到杭州时成了英雄。联合总部原本希望公布翁被逼写下的交代，从事件中获取政治资本，但事情反而反噬了自己。红暴声望大涨，并吸收了新成员，而联合总部一时名誉受损。为摆脱这一局面，联合总部6月10日在杭州青年广场举行大型集会。随后，一万名游行者向城北丝绸联合厂进发。途中在一座桥上，他们同前来支持翁所在团体的农民发生冲突。张永生把这次对抗归咎于浙江大学，并发誓要把校园砸烂。

Later in June, Zhang Chunqiao received Weng in Beijing to add his voice
to those urging Weng to defect.\footnote{Zhejiang Daily editorial,
  ``Smash the Gang of Four, suppress the counter-revolutionaries'',
  reprinted in HZRB, December 31, 1976.

  Zhejiang Daily 社论，``Smash the Gang of Four, suppress the
  counter-revolutionaries''，转载于 HZRB，1976年12月31日。} Weng Senhe,
writing from prison many years later, recalled that Zhang had expressed
regret (我们为你可惜) -- most probably at Weng having joined the Red
Storm -- and informed the Hangzhou rebel that he had much admired his
1966 rebellion against the leadership of the silk complex.\footnote{Weng
  Senhe bufen zuizheng cailiao, p.~4.

  Weng Senhe bufen zuizheng cailiao，第4页。} Zhang also reportedly
encouraged Weng to launch attacks on the ZPMD, the organization
recognized as the major remaining stumbling block to United
Headquarters' aspirations to seize power. United Headquarters was
apparently aiming to secure the replacement of the leaders of the ZPMD
and the military control commission, Zhang Xiulong and Long Qian, with
military leaders more sympathetic to its cause. In return for Weng's
apostasy, guarantees concerning his position in the future revolutionary
power structure may have been given.

6月下旬，张春桥在北京接见翁，加入劝说翁倒戈者的行列。\footnote{Zhejiang
  Daily editorial, ``Smash the Gang of Four, suppress the
  counter-revolutionaries'', reprinted in HZRB, December 31, 1976.

  Zhejiang Daily 社论，``Smash the Gang of Four, suppress the
  counter-revolutionaries''，转载于 HZRB，1976年12月31日。}
多年后在狱中写作的翁森鹤回忆说，张曾表示惋惜（我们为你可惜）--
很可能是指翁加入了红暴 --
并告诉这位杭州造反派，他十分赞赏翁在1966年反对丝绸联合厂领导层的造反行动。\footnote{Weng
  Senhe bufen zuizheng cailiao, p.~4.

  Weng Senhe bufen zuizheng cailiao，第4页。}
据说张还鼓励翁攻击浙江省军区，后者被认为是联合总部夺权愿望面前剩下的主要障碍。联合总部显然意在促成以更同情本派的军事领导人取代浙江省军区和军管会领导人张秀龙、龙潜。作为翁背叛原组织的回报，可能有人就其在未来革命权力结构中的地位作出保证。

Other factors help explain Weng's decision to change sides. As a rebel
who took his revolutionary credentials seriously, Weng felt compromised
by Red Storm's continuing support for the ``capitalist-roader'' Jiang
Hua, a stand virtually forced upon it by its backers, the ZPMD and Xu
Shiyou's Nanjing Military Region. More opportunistically, Weng may well
have seen the writing on the wall for Red Storm. With United
Headquarters having secured the support of the Central Cultural
Revolution Group, the 20th Army and the 5th Air Force the chances were
that the plums of office would fall into its lap, leaving Red Storm
isolated. A man of ambition, Weng would have had to consider seriously
this aspect of the matter.

还有其他因素有助于解释翁改变阵营的决定。作为一个认真看待自身革命资历的造反派，翁感到红暴继续支持``走资派''江华使自己处境尴尬；这种立场几乎是其后台浙江省军区和许世友的南京军区强加给它的。更机会主义地说，翁很可能已经看出红暴大势不妙。联合总部既然获得中央文革小组、第二十军和第五空军的支持，官位好处很可能落入其手中，红暴则会陷于孤立。翁是一个有野心的人，必然会认真考虑这一点。

The struggle for supremacy between the two rival organizations was
reaching a climax. On 11 June, the day after its battle with peasants in
the suburbs of Hangzhou, United Headquarters issued an urgent notice
which stated that the power struggle and attempts to arrive at a rebel
alliance and a three-way alliance of cadres, PLA, and mass
representatives had reached a critical stage.\footnote{ZPS, June 11,
  1967, SWB/FE/2490/B/4-7.

  ZPS，1967年6月11日，SWB/FE/2490/B/4-7。} It referred to a ``big bloody
battle'' against the old ZPC and provincial administration. According to
United Headquarters:

两个敌对组织之间的争夺正达到高潮。6月11日，即联合总部在杭州郊区与农民发生战斗的次日，它发布紧急通告，称夺权斗争以及实现造反派联盟、干部、解放军和群众代表三结合联盟的努力已经进入关键阶段。\footnote{ZPS,
  June 11, 1967, SWB/FE/2490/B/4-7.

  ZPS，1967年6月11日，SWB/FE/2490/B/4-7。}
通告提到反对旧省委和省政府的一场``大血战''。联合总部称：

\begin{quote}
Using power which they still hold {[}the party authorities{]} fabricated
lies and resorted to such vicious means as awarding work points and
mobilizing means of transport to deceive and incite large numbers of
peasants to leave their production posts to go to cities to participate
in violent struggles, to encircle and attack the revolutionary left, and
to sabotage state properties\ldots.
\end{quote}

\begin{quote}
他们{[}党内当权派{]}利用仍然掌握的权力，捏造谎言，并采取记工分、调动交通工具等恶毒手段，欺骗和煽动大批农民离开生产岗位，进城参加武斗，包围和攻击革命左派，破坏国家财产\ldots\ldots{}
\end{quote}

The resultant clashes had ``seriously undermined production, the
worker-peasant alliance, and the great proletarian revolution''.

由此发生的冲突``严重破坏了生产、工农联盟和伟大的无产阶级革命''。

United Headquarters requested the peasants to return to their communes
and direct their wrath at the power-holders rather than at the rebels.
It also asked peasants to reject the rewards dangled before them by
local officials as bait to lure them into the cities to fight against
the rebels. The notice asked members of its own organization to observe
the guidelines of the 6 June order and to treat the peasants as
misguided dupes rather than as enemies.\footnote{An editorial of June 14
  reiterated this last point. ZPS, June 14, 1967, SWB/FE/2499/B/9.

  6月14日的一篇社论重申了最后这一点。ZPS，1967年6月14日，SWB/FE/2499/B/9。}
But United Headquarters had already disregarded this order and the
peasants had been told that United Headquarters was a reactionary
organization.\footnote{Zhonggong nianbao 1968 nian, p.~434.

  Zhonggong nianbao 1968 nian，第434页。} The time for patience and
understanding was running out.

联合总部要求农民返回公社，把怒火指向掌权者而不是造反派。它还要求农民拒绝地方官员作为诱饵摆在面前、诱使他们进城同造反派作战的各种奖赏。该通告要求本组织成员遵守6月6日命令的准则，把农民视为被误导的受骗者，而不是敌人。\footnote{An
  editorial of June 14 reiterated this last point. ZPS, June 14, 1967,
  SWB/FE/2499/B/9.

  6月14日的一篇社论重申了最后这一点。ZPS，1967年6月14日，SWB/FE/2499/B/9。}
但联合总部自己已经无视这项命令，而农民也被告知联合总部是一个反动组织。\footnote{Zhonggong
  nianbao 1968 nian, p.~434.

  Zhonggong nianbao 1968 nian，第434页。} 耐心和谅解的时间正在耗尽。

Foreign sources reported armed clashes between rebels and peasants and
rebels and the army during June and July 1967. While unconfirmed and
perhaps exaggerated, the despatches give an indication of the turmoil
engulfing Zhejiang. A broadcast from the Soviet Union made the
sensational claim that guerrilla units led by members of the ZPC and
ZPMD had attacked Hangzhou from the mountains killing six people and
wounding another 150.\footnote{Radio Moscow, June 1, 1966,
  SWB/FE/2481/C/1. ZJRB later published an article accusing Tan Zhenlin
  of masterminding the establishment of a guerrilla unit in the Siming
  mountains, Yuyao county, eastern Zhejiang, where communist guerrillas
  under the New Fourth Army had fought both the Guomindang and the
  Japanese before 1949. ZPS, November 20, 1968, SWB/FE/2933/B/7-8.

  Radio Moscow，1966年6月1日，SWB/FE/2481/C/1。ZJRB
  后来发表文章，指控谭震林策划在浙江东部余姚县四明山建立游击队；1949年前，新四军领导的共产党游击队曾在那里同国民党和日本人作战。ZPS，1968年11月20日，SWB/FE/2933/B/7-8。}
An East European wire service spoke of tens of thousands of peasants
invading Hangzhou and throwing their victims into the Qiantang
river.\footnote{Czechoslovak Press Agency (CTK), June 12, 1967,
  SWB/FE/2491/C/1.

  Czechoslovak Press Agency（CTK），1967年6月12日，SWB/FE/2491/C/1。}
Hong Kong sources related accounts of railway workers in Jinhua, central
Zhejiang, fighting with rocks, clubs and hand grenades against PLA
units. Another report alleged that six people had been killed and fifty
injured when unnamed assailants had opened fire on rebels in a Hangzhou
factory on 29 May.\footnote{Zhonggong nianbao 1968 nian, p.~433.

  Zhonggong nianbao 1968 nian，第433页。} A Japanese report claimed that
students from Zhejiang University had attacked army and airforce
troops.\footnote{Asahi, June 1, 1967, CNS, July 7, 1967, p.~5.

  Asahi，1967年6月1日；CNS，1967年7月7日，第5页。} As if to confirm the
accuracy of these reports, a Central Committee document of July
specifically mentioned Zhejiang as a province where People's Armed
Forces Departments had incited peasants to attack cities. The peasants
had also interfered with the movement of trains and boats and had been
rewarded for their exploits.\footnote{``中共中央关于禁止挑动农民进城武斗的通知''
  (Notice of the CCP CC forbidding the incitement of peasants to enter
  cities to fight), July 13, 1967, CCP Documents, pp.~473-76.

  《中共中央关于禁止挑动农民进城武斗的通知》（Notice of the CCP CC
  forbidding the incitement of peasants to enter cities to
  fight），1967年7月13日，CCP Documents，第473-476页。}

外国消息来源报道，1967年6月和7月，造反派与农民、造反派与军队之间发生武装冲突。这些报道虽未经证实且可能有所夸大，却显示出席卷浙江的混乱。苏联一次广播耸人听闻地声称，由浙江省委和浙江省军区成员领导的游击队从山区进攻杭州，造成6人死亡、150人受伤。\footnote{Radio
  Moscow, June 1, 1966, SWB/FE/2481/C/1. ZJRB later published an article
  accusing Tan Zhenlin of masterminding the establishment of a guerrilla
  unit in the Siming mountains, Yuyao county, eastern Zhejiang, where
  communist guerrillas under the New Fourth Army had fought both the
  Guomindang and the Japanese before 1949. ZPS, November 20, 1968,
  SWB/FE/2933/B/7-8.

  Radio Moscow，1966年6月1日，SWB/FE/2481/C/1。ZJRB
  后来发表文章，指控谭震林策划在浙江东部余姚县四明山建立游击队；1949年前，新四军领导的共产党游击队曾在那里同国民党和日本人作战。ZPS，1968年11月20日，SWB/FE/2933/B/7-8。}
一个东欧通讯社称，数万农民涌入杭州，并把受害者扔进钱塘江。\footnote{Czechoslovak
  Press Agency (CTK), June 12, 1967, SWB/FE/2491/C/1.

  Czechoslovak Press Agency（CTK），1967年6月12日，SWB/FE/2491/C/1。}
香港消息来源转述了浙中金华铁路工人用石块、棍棒和手榴弹同解放军部队作战的说法。另有报道声称，5月29日杭州一家工厂中不明袭击者向造反派开火，造成6人死亡、50人受伤。\footnote{Zhonggong
  nianbao 1968 nian, p.~433.

  Zhonggong nianbao 1968 nian，第433页。}
日本一则报道声称，浙江大学学生攻击了陆军和空军部队。\footnote{Asahi,
  June 1, 1967, CNS, July 7, 1967, p.~5.

  Asahi，1967年6月1日；CNS，1967年7月7日，第5页。}
仿佛是为了证实这些报道的准确性，7月一份中央委员会文件特别提到，浙江是人民武装部煽动农民进城攻击的省份之一。农民还干扰火车和船只运行，并因其行动得到奖赏。\footnote{``中共中央关于禁止挑动农民进城武斗的通知''
  (Notice of the CCP CC forbidding the incitement of peasants to enter
  cities to fight), July 13, 1967, CCP Documents, pp.~473-76.

  《中共中央关于禁止挑动农民进城武斗的通知》（Notice of the CCP CC
  forbidding the incitement of peasants to enter cities to
  fight），1967年7月13日，CCP Documents，第473-476页。}

Relations between United Headquarters and the PLA were also tense. A
Zhejiang Daily editorial of 12 June warned the rebels to remain vigilant
against the guile of power-holders who goaded them into committing
``more errors'' in their relations with the military.\footnote{ZPS, June
  12, 1967, SWB/FE/2492/B/8-9.

  ZPS，1967年6月12日，SWB/FE/2492/B/8-9。} This was most probably a
reference to the 11 June clash between peasants and United Headquarters
in which the army had probably intervened on the side of the former. The
editorial lavishly praised the PLA for its political skills,
organization and discipline, assistance in production and for its
efforts to raise the ``revolutionary consciousness'' of the rebels. It
also pledged that mass organizations would place even greater trust in
their PLA protectors. What United Headquarters really desired, however,
was greater freedom from military control. This was indicated in the
conclusion of the editorial, in its positive reference to the central
orders of 28 January and 6 April curtailing the power of the PLA.

联合总部与人民解放军的关系同样紧张。《浙江日报》6月12日社论警告造反派，要警惕当权派的诡计，后者诱使他们在同军队的关系上犯``更多错误''。\footnote{ZPS,
  June 12, 1967, SWB/FE/2492/B/8-9.

  ZPS，1967年6月12日，SWB/FE/2492/B/8-9。}
这很可能指的是6月11日农民与联合总部的冲突，在这次冲突中，军队很可能站在前者一边介入。社论极力赞扬人民解放军的政治能力、组织纪律、对生产的帮助以及提高造反派``革命觉悟''的努力。它还承诺群众组织将更加信任自己的解放军保护者。然而，联合总部真正想要的是更大程度地摆脱军事控制。社论结尾对1月28日和4月6日限制人民解放军权力的中央命令作正面提及，显示了这一点。

Military units supportive of United Headquarters received favorable
publicity. Zhejiang Daily published an article praising the deeds of a
unit of the 5th Air Force in distinguishing between revolutionary and
conservative groups.\footnote{ZPS, June 21, 1967, SWB/FE/2500/B/8-9.

  ZPS，1967年6月21日，SWB/FE/2500/B/8-9。} This squad had arrived at a
chemical works in Hangzhou in March 1967 when the class struggle was at
its most acute. It was welcomed by members of the rebel group affiliated
with United Headquarters, who seemed to be in the minority and on the
defensive. The factory leaders were supporting the conservative group
and tried to convince the PLA unit to do likewise. According to the
report, the military personnel relied on three criteria to decide which
group it should support. The third criteria mentioned, and that of the
most relevance here, was to judge each group not by the class background
of its members but by its actual performance in the Cultural Revolution.
This example provides further evidence that the class composition of
United Headquarters' membership left the organization open to
accusations that it had recruited members of ``impure background'' into
its ranks. The other noteworthy aspect of the account was its revelation
of the impossibility of the PLA's task in such circumstances. Stranger
to the complexities and work relationships of the industrial sector, the
army could not objectively decide the relative merits of the two
contending parties. Their choice was prejudged by the political
allegiances of their superiors.

支持联合总部的军事单位得到了有利宣传。《浙江日报》发表文章，表扬第五空军某部区分革命团体与保守团体的事迹。\footnote{ZPS,
  June 21, 1967, SWB/FE/2500/B/8-9.

  ZPS，1967年6月21日，SWB/FE/2500/B/8-9。}
该分队在1967年3月阶级斗争最尖锐时来到杭州一家化工厂。隶属联合总部的造反团体成员欢迎他们，而这些成员似乎处于少数和防守地位。工厂领导支持保守团体，并试图说服解放军分队也这样做。据报道，军人依据三条标准决定支持哪个团体。第三条标准也是这里最相关的一条，是判断各团体不能看成员阶级出身，而要看其在文化大革命中的实际表现。这个例子进一步说明，联合总部成员的阶级构成使该组织容易受到其队伍吸收``出身不纯''成员的指控。报道另一个值得注意之处在于，它揭示了人民解放军在这种情况下任务的不可能性。军队不熟悉工业部门复杂的人际和工作关系，无法客观判断两个对立方的相对优劣。他们的选择已由上级的政治归属预先决定。

United Headquarters pressed on with its search for the elusive prize of
unity; a unity on its terms and one that guaranteed its predominance
over Red Storm. Between 14 and 24 June United Headquarters conducted a
conference to discuss ways of forming both an alliance of rebels and the
three-way alliance of rebels, cadres and the military.\footnote{ZPS,
  June 25, 1967, SWB/FE/2508/B/12.

  ZPS，1967年6月25日，SWB/FE/2508/B/12。} The rectification of the
organization's leadership workstyle and cadre ``liberation'' were two
prominent issues on the agenda. The cadre question was discussed in
editorials published in June and July.\footnote{ZPS, June 27, 1967,
  SWB/FE/2508/B/13-14; ZPS, July 11, 1967, SWB/FE/2520/B/15.

  ZPS，1967年6月27日，SWB/FE/2508/B/13-14；ZPS，1967年7月11日，SWB/FE/2520/B/15。}
Many cadres, it was announced, were victims of the ``bourgeois
reactionary line''. They should be allowed to rejoin the ranks of the
revolutionaries after conducting self-criticism and denouncing their
fellow cadres, both being proof of their commitment to the new order.
The rebels were warned to beware of cadres abusing the policy of
leniency in order to return to positions of influence.

联合总部继续追求难以到手的团结；这种团结必须按它的条件实现，并保证它压倒红暴。6月14日至24日，联合总部召开会议，讨论如何形成造反派联盟以及造反派、干部和军队三结合联盟。\footnote{ZPS,
  June 25, 1967, SWB/FE/2508/B/12.

  ZPS，1967年6月25日，SWB/FE/2508/B/12。}
整顿本组织领导作风和干部``解放''是议程上的两个突出问题。6月和7月发表的社论讨论了干部问题。\footnote{ZPS,
  June 27, 1967, SWB/FE/2508/B/13-14; ZPS, July 11, 1967,
  SWB/FE/2520/B/15.

  ZPS，1967年6月27日，SWB/FE/2508/B/13-14；ZPS，1967年7月11日，SWB/FE/2520/B/15。}
社论宣布，许多干部是``资产阶级反动路线''的受害者。他们在进行自我批评并揭发其他干部之后，应被允许重新加入革命队伍；这两者都被视为他们拥护新秩序的证明。造反派则被警告要提防干部滥用宽大政策，借机重返有影响力的岗位。

Many rebels clearly opposed this shift in policy but they were informed
that individuals should be judged by their life's work and not be
condemned merely on mistakes they may have committed in the early stages
of the Cultural Revolution. With policies in Beijing being determined on
the basis of the arbitrary dictates of the leader and his retinue, it is
little wonder that erratic and sudden changes of course caused
consternation and confusion for those expected to implement them. In
Zhejiang, the failure to convince the rebels of the correctness of the
revised cadre policy was manifest in the absence of all leading members
of the defunct administration from the provincial scene. Meanwhile,
supporters of Jiang Hua remained active in defence of their old boss. In
February 1967 they had published a document entitled ``Jiang Hua as
reflected in the series of inner-party struggles''\footnote{ZPS,
  November 19, 1968, SWB/FE/2933/B/3; ZPS, November 15, 1968,
  SWB/FE/2928/B/2.

  ZPS，1968年11月19日，SWB/FE/2933/B/3；ZPS，1968年11月15日，SWB/FE/2928/B/2。}
to press his claims for retention in the proposed new power structure.
The political repercussions of this audacious step continued to make
themselves felt in mid-year.

许多造反派显然反对这一政策转向，但他们被告知，评价个人应看其一生工作，而不应仅因其在文化大革命早期可能犯过的错误就予以谴责。北京的政策既然依据领袖及其随从的任意指令而定，那么路线的反复和突变使执行者惊惶困惑，也就毫不奇怪。在浙江，未能说服造反派相信修正后干部政策的正确性，明显表现为已瘫痪行政机构所有主要成员都缺席省级政治舞台。与此同时，江华的支持者仍积极为旧上司辩护。1967年2月，他们曾发表题为《从一系列党内斗争看江华》的文件，\footnote{ZPS,
  November 19, 1968, SWB/FE/2933/B/3; ZPS, November 15, 1968,
  SWB/FE/2928/B/2.

  ZPS，1968年11月19日，SWB/FE/2933/B/3；ZPS，1968年11月15日，SWB/FE/2928/B/2。}
为江华在拟议的新权力结构中保留地位提出主张。这一大胆举动的政治余波到年中仍在显现。

The second issue which received prominence at this time was the
work-style of rebel organizations. The leading personnel of both United
Headquarters and the Hangzhou Workers' Revolutionary Committee held
meetings in July\footnote{ZPS, July 12, 1967, SWB/FE/2520/B/13-15.

  ZPS，1967年7月12日，SWB/FE/2520/B/13-15。} to study the regulations of
the Shandong Provincial Revolutionary Committee\footnote{``山东革命委员会关于认真转变作风的若干规定''
  (Regulations of the Shandong Provincial Revolutionary Committee on
  earnestly transforming workstyle), June 7, 1967, Hongqi, 10 (1967),
  p.~65; Hongqi commentator, ``防止资产阶级思想侵蚀'' (Prevent the
  corrosion of bourgeois ideology), ibid, pp.~63-4. Other provinces took
  cognizance of the Shandong experience. See Harry Harding, Organizing
  China: The Problem of Bureaucracy, 1949-76 (Stanford: Stanford
  University Press, 1981), p.~276.

  《山东革命委员会关于认真转变作风的若干规定》（Regulations of the
  Shandong Provincial Revolutionary Committee on earnestly transforming
  workstyle），1967年6月7日，Hongqi，10（1967年），第65页；Hongqi
  评论员，《防止资产阶级思想侵蚀》（Prevent the corrosion of bourgeois
  ideology），同上，第63-64页。其他省份也注意到山东经验。见 Harry
  Harding，Organizing China: The Problem of Bureaucracy,
  1949-76（Stanford: Stanford University Press，1981年），第276页。}
concerning standards of behavior. Both organizations passed eight point
resolutions on the issue. The most significant point in United
Headquarters' resolution was that calling on its members to ``welcome
and accept criticism and supervision and listen to those who hold
different views and once opposed them and committed mistakes''. This
article of the resolution could only refer to Red Storm. Presumably, the
military was pressing United Headquarters to make concessions for the
sake of unity. The accusation previously levelled against Red Storm of
upholding the ``bourgeois reactionary line'' had now been reduced to
that of committing mistakes.

此时受到重视的第二个问题是造反组织的工作作风。联合总部和杭州市革命职工委员会的领导人在7月召开会议，\footnote{ZPS,
  July 12, 1967, SWB/FE/2520/B/13-15.

  ZPS，1967年7月12日，SWB/FE/2520/B/13-15。}
学习山东省革命委员会关于行为规范的规定。\footnote{``山东革命委员会关于认真转变作风的若干规定''
  (Regulations of the Shandong Provincial Revolutionary Committee on
  earnestly transforming workstyle), June 7, 1967, Hongqi, 10 (1967),
  p.~65; Hongqi commentator, ``防止资产阶级思想侵蚀'' (Prevent the
  corrosion of bourgeois ideology), ibid, pp.~63-4. Other provinces took
  cognizance of the Shandong experience. See Harry Harding, Organizing
  China: The Problem of Bureaucracy, 1949-76 (Stanford: Stanford
  University Press, 1981), p.~276.

  《山东革命委员会关于认真转变作风的若干规定》（Regulations of the
  Shandong Provincial Revolutionary Committee on earnestly transforming
  workstyle），1967年6月7日，Hongqi，10（1967年），第65页；Hongqi
  评论员，《防止资产阶级思想侵蚀》（Prevent the corrosion of bourgeois
  ideology），同上，第63-64页。其他省份也注意到山东经验。见 Harry
  Harding，Organizing China: The Problem of Bureaucracy,
  1949-76（Stanford: Stanford University Press，1981年），第276页。}
两个组织都通过了八点决议。联合总部决议中最重要的一点，是要求成员``欢迎和接受批评监督，听取那些持不同意见、曾经反对过他们并犯过错误的人''的意见。这一条只可能指红暴。想必军方正在向联合总部施压，要求它为了团结作出让步。此前加在红暴身上的坚持``资产阶级反动路线''的指控，如今已降为犯错误。

It was not only the rebels who were taking stock at this time. PLA units
in Hangzhou also held a meeting to review their contributions which, not
surprisingly, they decided were positive.\footnote{ZPS, July 13, 1967,
  SWB/FE/2529/B/15-16.

  ZPS，1967年7月13日，SWB/FE/2529/B/15-16。} The military boasted that
it had correctly identified the revolutionary left and had not deviated
from supporting it, despite the disruptive activities of conservative
organizations backed by the former power-holders. The editorial which
reported on the meeting made reference to the ``sharp and complicated''
nature of the class struggle and the ``very delicate'' nature of the
PLA's mission in Zhejiang. Mistakes were inevitable in the task of
supporting the left, added the editorial. While ``some people'' were
``evasive and reluctant'' to admit to their mistakes, others faced up to
their errors and corrected them. Here, the editorial was criticizing the
ZPMD and praising the main force units.

此时进行反省的不只是造反派。杭州的解放军部队也召开会议，检讨自身贡献；不出所料，他们认为自己的贡献是积极的。\footnote{ZPS,
  July 13, 1967, SWB/FE/2529/B/15-16.

  ZPS，1967年7月13日，SWB/FE/2529/B/15-16。}
军方自称正确识别了革命左派，并且尽管旧当权派支持的保守组织进行破坏活动，也没有偏离支左方向。报道该会议的社论提到阶级斗争``尖锐复杂''，以及人民解放军在浙江使命``十分微妙''。社论补充说，在支左任务中犯错误不可避免。有些人``躲躲闪闪、不愿''承认错误，另一些人则正视并改正错误。这里，社论是在批评浙江省军区，赞扬主力部队。

In late July and August 1967 the Cultural Revolution Group in Beijing
launched a series of radical measures to seize back the initiative. It
extended the campaign against Liu Shaoqi's ``bourgeois headquarters'' to
incorporate leading military cadres, particularly after the outbreak of
the Wuhan incident on 20 July.\footnote{For a detailed account of this
  incident, see Gao Gao and Yan Jiaqi, ``Wenhua dageming'', pp.~257-67;
  and Thomas W. Robinson, ``The Wuhan Incident: Local Strife and
  Provincial Rebellion During the Cultural Revolution'', CQ, 47 (1971),
  pp.~413-38.

  关于这一事件的详细记述，见 Gao Gao 和 Yan Jiaqi，``Wenhua
  dageming''，第257-267页；以及 Thomas W. Robinson，``The Wuhan
  Incident: Local Strife and Provincial Rebellion During the Cultural
  Revolution''，CQ，47（1971年），第413-438页。} The response to this
spectacular example of provincial military insubordination was a
strident demand for the overthrow of capitalist-roaders in the PLA. Red
Guard leaders in Beijing despatched emissaries to various localities to
increase the pressure on the military.\footnote{Zhao Wei, Zhao Ziyang
  zhuan, pp.~167-68.

  Zhao Wei，Zhao Ziyang zhuan，第167-168页。} Other violent and
destabilizing incidents threatened to plunge the country into civil war
and international disrepute. On 4 August the British mission in Beijing
was burned down as security guards looked on helplessly. The troika of
law enforcement agencies -- public security bureaux, procuratorates and
courts -- came under heavy political fire and the latter two bodies were
effectively disbanded. Members of the Cultural Revolution Group, who
were later accused of leading a conspiratorial organization called the
``May 16 Group'', published inflammatory speeches and articles to add
fuel to the fire.\footnote{Hao Mengbi and Duan Haoran, Liushinian,,
  pp.~596-97.

  Hao Mengbi 和 Duan Haoran，Liushinian，第596-597页。} In a speech of
22 July to rebel groups from Henan province, Jiang Qing issued her
famous ``dialectical'' slogan, whose nuances were not always well
understood by aroused young activists across China. Jiang called on
rebels to ``attack by reasoning, defend by force'' (文攻武卫). Her
followers in Zhejiang proclaimed that in relation to violence ``first we
oppose it; secondly we are not afraid of it''. Using the kind of
sophistry that characterized much of the polemics of the Cultural
Revolution, the Zhejiang Daily put forward the argument that ``To defend
by force is to ensure that an attack by peaceful means can be
successfully carried out''.\footnote{ZPS, August 10, 1967, URS, 48:18
  (September 1, 1967), pp.~255-58.

  ZPS，1967年8月10日，URS，48:18（1967年9月1日），第255-258页。} The
stage was thus set for an open test of strength between the two rival
mass organizations and their military allies, with rebels on both sides
looting arsenals to arm themselves for the forthcoming battles.

1967年7月下旬和8月，北京中央文革小组发动一系列激进措施，以夺回主动权。它把反刘少奇``资产阶级司令部''的运动扩大到包括高级军事干部，尤其是在7月20日武汉事件爆发之后。\footnote{For
  a detailed account of this incident, see Gao Gao and Yan Jiaqi,
  ``Wenhua dageming'', pp.~257-67; and Thomas W. Robinson, ``The Wuhan
  Incident: Local Strife and Provincial Rebellion During the Cultural
  Revolution'', CQ, 47 (1971), pp.~413-38.

  关于这一事件的详细记述，见 Gao Gao 和 Yan Jiaqi，``Wenhua
  dageming''，第257-267页；以及 Thomas W. Robinson，``The Wuhan
  Incident: Local Strife and Provincial Rebellion During the Cultural
  Revolution''，CQ，47（1971年），第413-438页。}
对这一省级军方抗命的惊人事例，回应是高调要求打倒人民解放军中的走资派。北京红卫兵领袖派遣使者前往各地，增加对军方的压力。\footnote{Zhao
  Wei, Zhao Ziyang zhuan, pp.~167-68.

  Zhao Wei，Zhao Ziyang zhuan，第167-168页。}
其他暴力和破坏稳定的事件也威胁把全国推入内战，并造成国际恶名。8月4日，北京英国代办处被焚毁，警卫人员只能无助旁观。执法三机关
-- 公安局、检察院和法院 --
遭到猛烈政治攻击，后两者实际上被解散。后来被指控领导``五一六集团''阴谋组织的中央文革小组成员发表煽动性讲话和文章，为局势火上浇油。\footnote{Hao
  Mengbi and Duan Haoran, Liushinian,, pp.~596-97.

  Hao Mengbi 和 Duan Haoran，Liushinian，第596-597页。}
7月22日，江青在对河南造反组织讲话时提出著名的``辩证''口号，其微妙含义并不总能被全国各地被激动起来的青年活动分子理解。江号召造反派``文攻武卫''。浙江的追随者宣称，对暴力``第一反对，第二不怕''。《浙江日报》用文化大革命论战中常见的诡辩提出，``武卫是为了保证文攻能够顺利进行''。\footnote{ZPS,
  August 10, 1967, URS, 48:18 (September 1, 1967), pp.~255-58.

  ZPS，1967年8月10日，URS，48:18（1967年9月1日），第255-258页。}
于是，两个敌对群众组织及其军事盟友公开较量的舞台搭好了，双方造反派都洗劫军械库，为即将到来的战斗武装自己。

Provincial rallies were staged on 19 and 22 July\footnote{ZPS, July 19,
  1967, July 23, 1967, SWB/FE/2531/B/12-13.

  ZPS，1967年7月19日、1967年7月23日，SWB/FE/2531/B/12-13。} to criticize
Liu Shaoqi, ``China's Khrushchev'' and to denounce the leaders of the
CCP HMC, including its former Secretary Wang Pingyi.\footnote{ZPS, July
  25, 1967, SWB/FE/B/10-11.

  ZPS，1967年7月25日，SWB/FE/B/10-11。} Cadres were encouraged to break
from their former superiors and publicly denounce them. At a rally on 26
July to celebrate the defeat of the Wuhan insurrection, a representative
of United Headquarters launched a blistering attack on
capitalist-roaders in the party, government and army in the presence of
PLA delegates, bringing the political temperature to boiling
point.\footnote{ZPS, July 27, 1967, SWB/FE/2533/B/3.

  ZPS，1967年7月27日，SWB/FE/2533/B/3。}

7月19日和22日，浙江举行省级集会，\footnote{ZPS, July 19, 1967, July 23,
  1967, SWB/FE/2531/B/12-13.

  ZPS，1967年7月19日、1967年7月23日，SWB/FE/2531/B/12-13。}
批判``中国的赫鲁晓夫''刘少奇，并谴责杭州市委领导人，包括原书记王平夷。\footnote{ZPS,
  July 25, 1967, SWB/FE/B/10-11.

  ZPS，1967年7月25日，SWB/FE/B/10-11。}
干部被鼓励同过去的上级决裂并公开揭发他们。7月26日，在庆祝武汉叛乱被击败的集会上，联合总部一名代表在解放军代表面前猛烈攻击党、政、军中的走资派，把政治温度推向沸点。\footnote{ZPS,
  July 27, 1967, SWB/FE/2533/B/3.

  ZPS，1967年7月27日，SWB/FE/2533/B/3。}

\section{Violence in the Summer of
1967}\label{violence-in-the-summer-of-1967}

\section{1967年夏季的暴力}\label{ux5e74ux590fux5b63ux7684ux66b4ux529b}

The festering antagonism between Red Storm and United Headquarters
erupted into open conflict in late July 1967. The Wuhan incident of 20
July provided the spark that set off the blaze. The incident provoked
the center into dismissing local military leaders whose readiness to
implement radical directives was questionable. On this basis a major
reshuffle of the Zhejiang military leadership took place.\footnote{Parris
  H. Chang, ``The Role of Chen Po-ta in the Cultural Revolution'', Asia
  Quarterly, 1 (1973), p.~44, fn. 58, wrongly claims that the reshuffle
  occurred in July, along with reshuffles in other provincial military
  districts and military regions. See also Facts \& Features, 2:3
  (November 27, 1968), pp.~15-16, for further background information to
  the Zhejiang reorganization. Zhou Enlai, in a speech of September 17,
  1967, to Red Guards made an announcement regarding this matter. See
  URS, 49:7 (October 24, 1967), p.~91.

  Parris H. Chang 在 ``The Role of Chen Po-ta in the Cultural
  Revolution''，Asia
  Quarterly，1（1973年），第44页脚注58中，错误地声称此次改组发生于7月，并与其他省军区和大军区的改组同时发生。关于浙江改组的更多背景资料，另见
  Facts \&
  Features，2:3（1968年11月27日），第15-16页。周恩来在1967年9月17日对红卫兵的讲话中就此事作了宣布。见
  URS，49:7（1967年10月24日），第91页。} On 26 August, the leaders of
the Zhejiang Provincial Military Control Commission, Zhang Xiulong and
Long Qian, were replaced at a ``handover of power meeting'' (交权大会)
by the Political Commissar and Commander respectively of the 20th Army,
Nan Ping (南萍) and Xiong Yingtang (熊应堂). Zhang and Long were then
subjected to ``cruel persecution'' (残酷迫害)\footnote{HZRB, September
  12, 1979, p.~1.

  HZRB，1979年9月12日，第1页。} for their prolonged resistance to United
Headquarters and assistance to Red Storm. An explanation for the
dismissal of leaders of the Henan Provincial Military District at about
this time argued that their purge resulted from an inability to
distinguish ``genuine proletarian revolutionaries'' from false
ones.\footnote{Parris Chang, ``The Revolutionary Committee in China,
  Case Studies: Heilungkiang and Honan'', CS, 6:9 (June 1, 1968),
  pp.~13-14.

  Parris Chang，``The Revolutionary Committee in China, Case Studies:
  Heilungkiang and Honan''，CS，6:9（1968年6月1日），第13-14页。} A
similar inability had plagued Zhang and Long in Zhejiang and their
position had been under threat since April.

1967年7月下旬，红暴与联合总部之间日益溃烂的敌意爆发为公开冲突。7月20日的武汉事件点燃了导火索。该事件促使中央撤换执行激进指示意愿可疑的地方军事领导人。在此基础上，浙江军事领导层发生重大改组。\footnote{Parris
  H. Chang, ``The Role of Chen Po-ta in the Cultural Revolution'', Asia
  Quarterly, 1 (1973), p.~44, fn. 58, wrongly claims that the reshuffle
  occurred in July, along with reshuffles in other provincial military
  districts and military regions. See also Facts \& Features, 2:3
  (November 27, 1968), pp.~15-16, for further background information to
  the Zhejiang reorganization. Zhou Enlai, in a speech of September 17,
  1967, to Red Guards made an announcement regarding this matter. See
  URS, 49:7 (October 24, 1967), p.~91.

  Parris H. Chang 在 ``The Role of Chen Po-ta in the Cultural
  Revolution''，Asia
  Quarterly，1（1973年），第44页脚注58中，错误地声称此次改组发生于7月，并与其他省军区和大军区的改组同时发生。关于浙江改组的更多背景资料，另见
  Facts \&
  Features，2:3（1968年11月27日），第15-16页。周恩来在1967年9月17日对红卫兵的讲话中就此事作了宣布。见
  URS，49:7（1967年10月24日），第91页。}
8月26日，在一次``交权大会''上，浙江省军事管制委员会领导人张秀龙、龙潜分别被第二十军政委南萍和军长熊应堂取代。张、龙随后因长期抵制联合总部并帮助红暴而受到``残酷迫害''。\footnote{HZRB,
  September 12, 1979, p.~1.

  HZRB，1979年9月12日，第1页。}
大约同时，河南省军区领导人被撤换的解释称，他们被清洗是因为不能把``真正的无产阶级革命派''同冒牌者区分开来。\footnote{Parris
  Chang, ``The Revolutionary Committee in China, Case Studies:
  Heilungkiang and Honan'', CS, 6:9 (June 1, 1968), pp.~13-14.

  Parris Chang，``The Revolutionary Committee in China, Case Studies:
  Heilungkiang and Honan''，CS，6:9（1968年6月1日），第13-14页。}
类似的无能也困扰着浙江的张、龙，而他们的地位自4月以来就一直受到威胁。

This major initiative by the center to exert control over Zhejiang was
called the ``two reorganizations'' (两个改组)\footnote{HZRB, September
  12, 1979, p.~1.

  HZRB，1979年9月12日，第1页。} and signalled the end of open support
for Red Storm by the provincial military authorities. PLA leaders from
central field armies and naval and air force units now predominated in
Zhejiang. It has been pointed out that these forces were more reliant on
allocations from central resources and less tied to local conditions
than the ground forces of military districts.\footnote{Chang, ``Regional
  Military Power'', p.~1009. However, according to Liu Guokai, the 22nd
  central Field Army led by Commander Cao Siming (曹思明) and Political
  Commissar Tie Ying, stationed on the Zhoushan island chain supported
  Red Storm. See, ``文化大革命简析'' (A Brief Analysis of the Cultural
  Revolution) in 人民之声 (The People's Voice), reprinted in
  大陆地下刊物汇编 (A Collection of Underground Publications from the
  Mainland), Vol. 17 (Taipei: Institute for the Study of Chinese
  Communist Problems, 1983), p.~202. It is most likely that Tie and Cao
  Siming, by their support for Red Storm, thereby brought themselves
  into conflict with Nan Ping's 20th Army. United Headquarters had found
  the Cultural Revolution hard going in Zhoushan. See ZPS, November 20,
  1968, SWB/FE/2933/B/6-7. See also ``The Case of Weng Sengho'', p.~2,
  for claims that Cao was a friend of Xu Shiyou and provided Red Storm
  militants with arms to do battle with their opponents. HZRB, September
  12, 1979, p.~1.

  Chang，``Regional Military Power''，第1009页。不过，据 Liu Guokai
  称，驻舟山群岛、由司令员曹思明（曹思明）和政委铁瑛领导的第22军支持红暴。见《文化大革命简析》（A
  Brief Analysis of the Cultural Revolution），载《人民之声》（The
  People's Voice），转载于《大陆地下刊物汇编》（A Collection of
  Underground Publications from the Mainland）第17卷（台北：Institute
  for the Study of Chinese Communist
  Problems，1983年），第202页。铁瑛和曹思明很可能因支持红暴而同南萍的第20军发生冲突。联合总部在舟山推进文化大革命时困难重重。见
  ZPS，1968年11月20日，SWB/FE/2933/B/6-7。关于曹是许世友的朋友，并向红暴武装分子提供武器以同对手作战的说法，另见
  ``The Case of Weng Sengho''，第2页；HZRB，1979年9月12日，第1页。} Just
prior to the leadership changes in the governing military bodies, a
provisional leadership organ (浙江省临时领导机构 -- also known as the
supreme organ of power 浙江省最高权力机构) -- had been established on 8
August. According to one informed source it was a body unique to China's
twenty-nine provincial-level administrations.\footnote{CNS, 270 (May 15,
  1969), p.~2.

  CNS，270（1969年5月15日），第2页。} It is unclear what relationship
existed between this body and the military control commission. Its
leadership was dominated by military personnel and presumably it
functioned as the administrative equivalent of the moribund provincial
government.

中央为控制浙江采取的这一重大举措被称为``两个改组''，\footnote{HZRB,
  September 12, 1979, p.~1.

  HZRB，1979年9月12日，第1页。}
标志着省级军事当局公开支持红暴的终结。来自中央野战军以及海军、空军部队的解放军领导人如今在浙江占据主导。有人指出，这些力量比军区地面部队更依赖中央资源分配，也更少受地方条件牵制。\footnote{Chang,
  ``Regional Military Power'', p.~1009. However, according to Liu
  Guokai, the 22nd central Field Army led by Commander Cao Siming
  (曹思明) and Political Commissar Tie Ying, stationed on the Zhoushan
  island chain supported Red Storm. See, ``文化大革命简析'' (A Brief
  Analysis of the Cultural Revolution) in 人民之声 (The People's Voice),
  reprinted in 大陆地下刊物汇编 (A Collection of Underground
  Publications from the Mainland), Vol. 17 (Taipei: Institute for the
  Study of Chinese Communist Problems, 1983), p.~202. It is most likely
  that Tie and Cao Siming, by their support for Red Storm, thereby
  brought themselves into conflict with Nan Ping's 20th Army. United
  Headquarters had found the Cultural Revolution hard going in Zhoushan.
  See ZPS, November 20, 1968, SWB/FE/2933/B/6-7. See also ``The Case of
  Weng Sengho'', p.~2, for claims that Cao was a friend of Xu Shiyou and
  provided Red Storm militants with arms to do battle with their
  opponents. HZRB, September 12, 1979, p.~1.

  Chang，``Regional Military Power''，第1009页。不过，据 Liu Guokai
  称，驻舟山群岛、由司令员曹思明（曹思明）和政委铁瑛领导的第22军支持红暴。见《文化大革命简析》（A
  Brief Analysis of the Cultural Revolution），载《人民之声》（The
  People's Voice），转载于《大陆地下刊物汇编》（A Collection of
  Underground Publications from the Mainland）第17卷（台北：Institute
  for the Study of Chinese Communist
  Problems，1983年），第202页。铁瑛和曹思明很可能因支持红暴而同南萍的第20军发生冲突。联合总部在舟山推进文化大革命时困难重重。见
  ZPS，1968年11月20日，SWB/FE/2933/B/6-7。关于曹是许世友的朋友，并向红暴武装分子提供武器以同对手作战的说法，另见
  ``The Case of Weng Sengho''，第2页；HZRB，1979年9月12日，第1页。}
在军事领导机构改组前不久，一个临时领导机构（浙江省临时领导机构，又称浙江省最高权力机构）已于8月8日成立。据一位知情者说，这是中国二十九个省级行政单位中独一无二的机构。\footnote{CNS,
  270 (May 15, 1969), p.~2.

  CNS，270（1969年5月15日），第2页。}
这一机构与军管会之间存在何种关系尚不清楚。其领导层由军人主导，大概发挥着垂死省政府的行政替代物作用。

During the month of August 1967 a series of violent incidents shook in
Zhejiang. To arm themselves for the forthcoming battles, the rebels
needed weapons and only the military could supply them. On the 5th and
again on the 9th of August Zhang Yongsheng led a thousand followers to
two armories managed by the military district in Hangzhou. It is highly
likely that the rebels received the cooperation of certain officers in
getting hold of military hardware. They seized a large number of
pistols, signal guns, sub-machine guns, light and heavy machine-guns,
anti-aircraft guns, mortars, recoilless guns, flamethrowers, hand
grenades, ammunition and vehicles.\footnote{Zhang Yongshengde zuizheng
  cailiao, p.~33.

  Zhang Yongshengde zuizheng cailiao，第33页。} This weaponry was lethal
in the hands of young students already practised in various forms of
violence. At the end of March 1967 Zhang Yongsheng had formed a crack
attack unit called the Flying Tiger Squad (飞虎队) which had already
been tested in battle in such incidents as the storming of the silk mill
on 6 June.\footnote{Ibid, p.~38.

  同上，第38页。} No doubt Red Storm, with its contacts in the military
district, also acquired its own weapons. The stage was set for a series
of bloody confrontations across the province.

1967年8月，一系列暴力事件震动浙江。为了在即将到来的战斗中武装自己，造反派需要武器，而只有军方能够提供。8月5日和9日，张永生两次率领一千名追随者前往杭州由军区管理的两个军械库。造反派很可能得到某些军官配合，从而取得军事装备。他们夺取了大量手枪、信号枪、冲锋枪、轻重机枪、高射炮、迫击炮、无后坐力炮、喷火器、手榴弹、弹药和车辆。\footnote{Zhang
  Yongshengde zuizheng cailiao, p.~33.

  Zhang Yongshengde zuizheng cailiao，第33页。}
这些武器落入已经熟悉各种暴力形式的青年学生手中，具有致命性。1967年3月底，张永生成立了一支名为飞虎队的精锐突击队，该队已在6月6日攻打丝绸厂等事件中经受战斗检验。\footnote{Ibid,
  p.~38.

  同上，第38页。}
毫无疑问，红暴凭借其在军区中的关系，也取得了自己的武器。全省一连串血腥对抗的舞台已经搭好。

Although United Headquarters could claim some successes in these
clashes, these were confined principally to operations which took place
in or around Hangzhou. Outside the provincial capital, particularly in
major provincial centers such as Jinhua and Wenzhou, its forces and
allies suffered overwhelming defeats and were only saved from
annihilation by the intervention of main force units.\footnote{The
  pattern of ``rebel'' strength being confined to major urban centers
  was typical of the Cultural Revolution. See Liu Guokai, A Brief
  Analysis, pp.~92-101.

  ``造反派''力量局限于主要城市中心的模式，是文化大革命中的典型现象。见
  Liu Guokai，A Brief Analysis，第92-101页。} Red Storm and other
adversaries of United Headquarters acquitted themselves outstandingly.
Although United Headquarters had the advantage of being recognized as
the official Maoist revolutionary organization in Zhejiang, it seemed to
fight a constant battle for survival against enemies on all sides.

尽管联合总部可以声称在这些冲突中取得一些胜利，但这些胜利主要限于杭州及其周边的行动。在省会之外，特别是在金华、温州等重要省内中心，其力量及盟友遭受压倒性失败，只是靠主力部队介入才免于被歼灭。\footnote{The
  pattern of ``rebel'' strength being confined to major urban centers
  was typical of the Cultural Revolution. See Liu Guokai, A Brief
  Analysis, pp.~92-101.

  ``造反派''力量局限于主要城市中心的模式，是文化大革命中的典型现象。见
  Liu Guokai，A Brief Analysis，第92-101页。}
红暴和联合总部的其他对手表现出色。尽管联合总部拥有被承认为浙江官方毛派革命组织的优势，却似乎始终在同四面八方的敌人进行生存之战。

The first publicized outbreaks of violence erupted at the Zhejiang Hemp
Mill (浙江麻纺织厂) on 16 July 1967\footnote{ZJRB, September 20, 1967,
  p.~1. Gao Zhixian, ``The evil lackey who sold himself to the bitter
  end to the Gang of Four'', HZRB, June 26, 1977. The mill is another
  major industrial unit in Hangzhou. Construction began in 1950 and
  finished in 1953. In 1978 it employed 6,200 workers who were in charge
  of 11,000 spindles and 676 weaving machines. Visit to the Zhejiang
  Hemp Mill, December 3, 1978.

  ZJRB，1967年9月20日，第1页。Gao Zhixian，``The evil lackey who sold
  himself to the bitter end to the Gang of
  Four''，HZRB，1977年6月26日。该厂是杭州另一家主要工业单位。其建设始于1950年，1953年完工。1978年，该厂雇用6200名工人，负责11000枚纱锭和676台织机。1978年12月3日访问浙江麻纺织厂。}
and at the Hangzhou silk complex on 31 July 1967. In the latter instance
a worker named Wu Fengying, who had been honored as an advanced silk
worker, was seized on her way home from work. She was incarcerated in a
warehouse where for two hours she was beaten and tortured. Her
tormentors then tied her to a bed, stuffed a dirty rag in her mouth and
violated her body with clubs. After two hours of this treatment she
died. Five other women who belonged to the same faction were seized and
thrashed, with their assailants shouting all the while: ``CCP members
must be beaten. It serves them right if they die''.\footnote{Hua
  Meisheng, ``Beating, smashing and looting is sheer
  counter-revolutionary activity'', HZRB, June 22, 1977; Yang Lingen
  (widower of Wu Fengying), ``An accusation written in blood and
  tears'', ibid.; Hangzhou public security bureau criticism group, ``An
  extremely ferocious counter-revolutionary gang'', HZRB, June 23, 1967;
  Hangzhou silk bureau criticism group, ``Iron-clad proof of Weng Senhe
  and He Xianchun ganging up to usurp power in the silk industry'',
  HZRB, June 30, 1977; meeting at Hangzhou Silk Complex, HZRB, December
  14, 1977.

  Hua Meisheng，``Beating, smashing and looting is sheer
  counter-revolutionary activity''，HZRB，1977年6月22日；Yang
  Lingen（吴凤英的鳏夫），``An accusation written in blood and
  tears''，同上；杭州市公安局批判组，``An extremely ferocious
  counter-revolutionary
  gang''，HZRB，1967年6月23日；杭州市丝绸局批判组，``Iron-clad proof of
  Weng Senhe and He Xianchun ganging up to usurp power in the silk
  industry''，HZRB，1977年6月30日；杭州丝绸联合厂会议，HZRB，1977年12月14日。}
Atrocities also occurred at the Hangzhou Glass Factory, on 19
August.\footnote{HZRB, December 5, 1977, p.~3.

  HZRB，1977年12月5日，第3页。} In this incident forty people were
beaten up by members of United Headquarters. Weng Senhe's defection to
United Headquarters may well have provoked this outburst of factional
violence as group loyalties and numbers were realigned in the key
industrial units of Hangzhou.

最早见诸公开报道的暴力爆发，发生在1967年7月16日的浙江麻纺织厂，\footnote{ZJRB,
  September 20, 1967, p.~1. Gao Zhixian, ``The evil lackey who sold
  himself to the bitter end to the Gang of Four'', HZRB, June 26, 1977.
  The mill is another major industrial unit in Hangzhou. Construction
  began in 1950 and finished in 1953. In 1978 it employed 6,200 workers
  who were in charge of 11,000 spindles and 676 weaving machines. Visit
  to the Zhejiang Hemp Mill, December 3, 1978.

  ZJRB，1967年9月20日，第1页。Gao Zhixian，``The evil lackey who sold
  himself to the bitter end to the Gang of
  Four''，HZRB，1977年6月26日。该厂是杭州另一家主要工业单位。其建设始于1950年，1953年完工。1978年，该厂雇用6200名工人，负责11000枚纱锭和676台织机。1978年12月3日访问浙江麻纺织厂。}
以及1967年7月31日的杭州丝绸印染联合厂。在后一事件中，一名曾被评为先进丝绸工人的女工吴凤英在下班回家途中被抓。她被关进仓库，遭殴打折磨两小时。施暴者随后把她绑在床上，把脏布塞进她嘴里，并用棍棒侵犯她的身体。又过两小时后，她死去。同属一派的另外五名妇女也被抓走并殴打，袭击者一边打还一边喊：``共产党员就该打，打死活该。''\footnote{Hua
  Meisheng, ``Beating, smashing and looting is sheer
  counter-revolutionary activity'', HZRB, June 22, 1977; Yang Lingen
  (widower of Wu Fengying), ``An accusation written in blood and
  tears'', ibid.; Hangzhou public security bureau criticism group, ``An
  extremely ferocious counter-revolutionary gang'', HZRB, June 23, 1967;
  Hangzhou silk bureau criticism group, ``Iron-clad proof of Weng Senhe
  and He Xianchun ganging up to usurp power in the silk industry'',
  HZRB, June 30, 1977; meeting at Hangzhou Silk Complex, HZRB, December
  14, 1977.

  Hua Meisheng，``Beating, smashing and looting is sheer
  counter-revolutionary activity''，HZRB，1977年6月22日；Yang
  Lingen（吴凤英的鳏夫），``An accusation written in blood and
  tears''，同上；杭州市公安局批判组，``An extremely ferocious
  counter-revolutionary
  gang''，HZRB，1967年6月23日；杭州市丝绸局批判组，``Iron-clad proof of
  Weng Senhe and He Xianchun ganging up to usurp power in the silk
  industry''，HZRB，1977年6月30日；杭州丝绸联合厂会议，HZRB，1977年12月14日。}
杭州玻璃厂8月19日也发生暴行。\footnote{HZRB, December 5, 1977, p.~3.

  HZRB，1977年12月5日，第3页。}
在该事件中，四十人被联合总部成员殴打。翁森鹤倒向联合总部，很可能在杭州关键工业单位中重新调整团体忠诚和人数，从而激发了这次派性暴力。

United Headquarters' greatest victories outside Hangzhou occurred at the
neighboring county town of Xiaoshan and the central county of Lanxi on
26 August 1967, the day on which the reshuffle of the military control
commission leadership was announced. The previous day Zhejiang Daily had
carried an article accusing the disgraced party leaders of continuing to
incite peasants to attack cities.\footnote{ZPS, August 26, 1967,
  SWB/FE/2557/B/20-1; Zhonggong nianbao 1968 nian, p.~434.

  ZPS，1967年8月26日，SWB/FE/2557/B/20-1；Zhonggong nianbao 1968
  nian，第434页。} In retaliation, truckloads of workers from factories
such as the Hangzhou Iron and Steel Works launched a raid on Xiaoshan.
Five hundred workers from the mill participated in the
attack.\footnote{Visit to the Hangzhou Iron and Steel Mill, May 27,
  1979. The mill was built in 1958. In 1979 a total of 11 sub-plants
  employed a staff of 12,619.

  1979年5月27日访问杭州钢铁厂。该厂建于1958年。1979年，全厂共有11个分厂，职工12619人。}
Led by Zhang Yongsheng, He Xianchun and Xia Genfa, the contingent struck
at Xiaoshan before moving on to neighboring Zhuji county where innocent
bystanders were shot down outside their houses.\footnote{Oral source.

  口述资料。} Accounts of what happened in Xiaoshan have been recorded
by the county party committee and participants.\footnote{Zhang
  Yongshengde zuizheng cailiao, p.~41; speech by Wang Baozhen, a worker
  at the Hangzhou No.~2 Cotton Mill at a meeting of the factory workers
  on December 3, 1977, HZRB, December 3, 1977, December 5, 1977; He
  Naimin, ``While there's life, struggle with Lin Biao and the `Gang of
  Four' does not cease'', HZRB, November 14, 1978.

  Zhang Yongshengde zuizheng
  cailiao，第41页；杭州第二棉纺厂工人王宝珍在1977年12月3日全厂工人会议上的讲话，HZRB，1977年12月3日、1977年12月5日；He
  Naimin，``While there's life, struggle with Lin Biao and the `Gang of
  Four' does not cease''，HZRB，1978年11月14日。} According to the 1977
deposition written by the county committee. Nan Ping, Chen Liyun and
Xiong Yingtang encouraged United Headquarters to raid military arsenals
at the end of July 1967. Military leaders met to plan the attack and in
mid August students from Hangzhou University were sent to Xiaoshan to
familiarize themselves with the locality. On 24 August an oath-taking
rally was held in Hangzhou's Youth Square, followed by an armed parade
which was led by Zhang Yongsheng and He Xianchun standing on a crane.
Speakers at the rally warned of an imminent assault on Hangzhou by armed
peasants. At dawn on 26 August, 3,000 to 4,000 rebels from ten or so
units, commanded by a military officer, assembled for a three-pronged
attack on the county town. During the fighting they used their recently
acquired flame-throwers to set fire to the Xiaoshan Agricultural
Machinery Factory with the losses totalling 243,000 yuan.

联合总部在杭州以外最大的胜利，发生在邻近县城萧山和浙中兰溪县，时间是1967年8月26日，也就是宣布军管会领导层改组的同一天。前一天，《浙江日报》刊文指控失势党内领导人继续煽动农民攻城。\footnote{ZPS,
  August 26, 1967, SWB/FE/2557/B/20-1; Zhonggong nianbao 1968 nian,
  p.~434.

  ZPS，1967年8月26日，SWB/FE/2557/B/20-1；Zhonggong nianbao 1968
  nian，第434页。}
作为报复，杭州钢铁厂等工厂一卡车一卡车的工人突袭萧山。杭钢有五百名工人参加了进攻。\footnote{Visit
  to the Hangzhou Iron and Steel Mill, May 27, 1979. The mill was built
  in 1958. In 1979 a total of 11 sub-plants employed a staff of 12,619.

  1979年5月27日访问杭州钢铁厂。该厂建于1958年。1979年，全厂共有11个分厂，职工12619人。}
在张永生、贺贤春和夏根发率领下，这支队伍先打击萧山，随后转向邻近的诸暨县，在那里，屋外无辜旁观者被枪击倒。\footnote{Oral
  source.

  口述资料。} 县委和参与者都留下了萧山事件的记述。\footnote{Zhang
  Yongshengde zuizheng cailiao, p.~41; speech by Wang Baozhen, a worker
  at the Hangzhou No.~2 Cotton Mill at a meeting of the factory workers
  on December 3, 1977, HZRB, December 3, 1977, December 5, 1977; He
  Naimin, ``While there's life, struggle with Lin Biao and the `Gang of
  Four' does not cease'', HZRB, November 14, 1978.

  Zhang Yongshengde zuizheng
  cailiao，第41页；杭州第二棉纺厂工人王宝珍在1977年12月3日全厂工人会议上的讲话，HZRB，1977年12月3日、1977年12月5日；He
  Naimin，``While there's life, struggle with Lin Biao and the `Gang of
  Four' does not cease''，HZRB，1978年11月14日。}
按县委1977年的材料，南萍、陈励耘和熊应堂在1967年7月底鼓励联合总部突袭军械库。军方领导人开会策划攻击，8月中旬杭州大学学生被派往萧山熟悉当地情况。8月24日，杭州青年广场举行誓师大会，随后举行武装游行，张永生和贺贤春站在吊车上领队。集会发言人警告说，武装农民即将进攻杭州。8月26日拂晓，来自十来个单位的三千至四千名造反派在一名军官指挥下集结，对县城发动三路进攻。战斗中，他们使用新近获得的喷火器纵火焚烧萧山农机厂，损失合计24.3万元。

Wang Baozhen, a worker at the Hangzhou No.~2 Cotton Mill located in
Xiaoshan, claimed that many people were murdered in cold blood and
others tortured. An official account put the death toll at twenty-seven
with many injured.\footnote{See 萧山县志编纂委员会 (Xiaoshan County
  Gazetteer Compilation Committee), 萧山县志 (Xiaoshan County
  Gazetteer), (Hangzhou: Zhejiang Renmin Chubanshe, 1987), p.~50; Zhang
  Yongshengde zuizheng cailiao, p.~41.

  见萧山县志编纂委员会（Xiaoshan County Gazetteer Compilation
  Committee），《萧山县志》（Xiaoshan County Gazetteer）（杭州：Zhejiang
  Renmin Chubanshe，1987年），第50页；Zhang Yongshengde zuizheng
  cailiao，第41页。} Wang's denial that she had participated in the
attack on the iron and steel mill, which had been instigated by former
ZPC leaders, rings hollow. Six workers had lost their lives in the
assault and it is therefore little wonder that their work-mates were
determined to avenge the deaths. By proudly revealing that she was a
member of the factory's militia and a recognized crack shot, Wang only
increases the suspicion that the assault on Xiaoshan was not unprovoked.

位于萧山的杭州第二棉纺厂女工王宝珍声称，许多人被冷血杀害，另一些人遭到酷刑。官方记载称死亡二十七人，多人受伤。\footnote{See
  萧山县志编纂委员会 (Xiaoshan County Gazetteer Compilation Committee),
  萧山县志 (Xiaoshan County Gazetteer), (Hangzhou: Zhejiang Renmin
  Chubanshe, 1987), p.~50; Zhang Yongshengde zuizheng cailiao, p.~41.

  见萧山县志编纂委员会（Xiaoshan County Gazetteer Compilation
  Committee），《萧山县志》（Xiaoshan County Gazetteer）（杭州：Zhejiang
  Renmin Chubanshe，1987年），第50页；Zhang Yongshengde zuizheng
  cailiao，第41页。}
王否认自己参加了由原省委领导人煽动的对钢铁厂的攻击，这一否认难以令人信服。六名工人在那次攻击中丧生，因此他们的同事决心复仇也不足为奇。王自豪地透露自己是工厂民兵成员并且是公认的神枪手，这只会加深人们对萧山攻击并非无端发生的怀疑。

Early in 1967 workers of the cotton mill had split into two factions,
with United Headquarters seizing power and forcing its opponents out of
the factory gate.\footnote{A total of 1,000 workers in Xiaoshan lost
  their jobs and had their wages stopped. When the local trade union
  council agreed to grant them emergency loans, the Chairman was placed
  in gaol as punishment. He Naimin, HZRB, November 14, 1978. The author
  was the trade union Chairman in question.

  萧山共有1000名工人失去工作并被停发工资。当地工会委员会同意向他们发放紧急贷款后，工会主席被投入监狱作为惩罚。He
  Naimin，HZRB，1978年11月14日。作者即为这位工会主席。} On the evening
of 25 August, Wang Baozhen and a fellow worker were on duty as medical
orderlies at the county hostel. That night the hostel was surrounded and
the front door kicked in. Those inside were ordered to come out with
their arms raised and then escorted to the railway station where Wang
lost contact with her comrade. Later she learned that this woman was
mistaken for her and killed. Wang herself did not fare much better. She
was beaten and interrogated so that she would confess responsibility for
the death of the steel workers. Her captors smashed her fingers with a
shifting-spanner to make sure that she would never be able to use a
rifle again. They wanted her to implicate the leaders of the ZPMD in the
bloodbath at the iron and steel mill but Wang reportedly refused, even
though she was stabbed in the backside with a bayonet.

1967年初，棉纺厂工人分裂成两派，联合总部夺取权力，并把对手赶出厂门。\footnote{A
  total of 1,000 workers in Xiaoshan lost their jobs and had their wages
  stopped. When the local trade union council agreed to grant them
  emergency loans, the Chairman was placed in gaol as punishment. He
  Naimin, HZRB, November 14, 1978. The author was the trade union
  Chairman in question.

  萧山共有1000名工人失去工作并被停发工资。当地工会委员会同意向他们发放紧急贷款后，工会主席被投入监狱作为惩罚。He
  Naimin，HZRB，1978年11月14日。作者即为这位工会主席。}
8月25日晚，王宝珍和一名同事在县招待所担任医务值班员。当天夜里，招待所被包围，前门被踢开。里面的人被命令举手出来，然后被押往火车站，王在那里同同伴失去联系。后来她得知，这名妇女被误认为是她而遭杀害。王自己的遭遇也好不了多少。她遭殴打审讯，被迫承认对钢铁厂工人死亡负责。俘虏者用活络扳手砸碎她的手指，以确保她再也不能使用步枪。他们想让她牵连浙江省军区领导人，称其参与钢铁厂血案；据说王拒绝了，尽管她的臀部被刺刀刺伤。

Details of what occurred in Lanxi are more sparse.\footnote{Lanxi
  shizhi, p.~19.

  Lanxi shizhi，第19页。} On 2 February 1967 the mass organization
linked with Red Storm, ``Lanxi County Bombard the Headquarters United
Headquarters'' seized power in the county town. Its rival faction,
``Lanxi County Revolutionary Rebel United Headquarters'' was established
the following June. On 13 July fierce fighting erupted between the two
groups and twenty-four people were killed and 138 injured when a
warehouse was burnt down. On 26 August United Headquarters launched an
all-out assault on ``Bombard the Headquarters''. A further five people
were killed and the latter group thereafter disintegrated.

兰溪发生的情况细节较少。\footnote{Lanxi shizhi, p.~19.

  Lanxi shizhi，第19页。}
1967年2月2日，与红暴相联系的群众组织``兰溪县炮打司令部联合总部''在县城夺权。其对立派``兰溪县革命造反联合总部''于同年6月成立。7月13日，两派爆发激烈战斗，一个仓库被烧毁，造成二十四人死亡、一百三十八人受伤。8月26日，联合总部对``炮打司令部''发动全面进攻。又有五人死亡，此后后者瓦解。

The above incidents in which United Headquarters blitzed its factional
opponents, were not publicized at the time. Much later these were
branded ``counter-revolutionary incidents'',\footnote{See ZPS, February
  10, 1979, SWB/FE/6045/BII/12-13.

  见 ZPS，1979年2月10日，SWB/FE/6045/BII/12-13。} and condemned in the
local press. The only incidents which made the headlines at the time of
their occurrence were those resulting in losses for United Headquarters,
which used its power over the media to gain sympathy for itself and to
censure its opponents. Since the end of the Cultural Revolution, no
condemnation has been published concerning these cases of violence.

上述联合总部猛攻派性对手的事件，在当时并未公开报道。很久以后，这些事件被定性为``反革命事件''，\footnote{See
  ZPS, February 10, 1979, SWB/FE/6045/BII/12-13.

  见 ZPS，1979年2月10日，SWB/FE/6045/BII/12-13。}
并受到地方报刊谴责。当时登上头条的，只有那些使联合总部遭受损失的事件；联合总部利用其对媒体的控制为自己争取同情并谴责对手。文化大革命结束以来，对这些暴力案件并未发表谴责。

Extraordinary incidents in which United Headquarters and its allies were
vanquished took place further away from the provincial capital. On 3
August trouble broke out in the central Zhejiang town of
Jinhua.\footnote{The following description of events at Jinhua is based
  on ``The `August 3' Incident in Jinhua area'', URS, 48:18 (September
  1, 1967), pp.~247-55; ZPS, August 11, 1967, SWB/FE/2542/B/4-5; ZPS,
  August 14, 1967, SWB/FE/2546/B/13; ZPS, August 19, 1967,
  SWB/FE/2549/B/12-13; ZPS, August 26, 1967, SWB/FE/2557/B/18-19;
  Zhonggong nianbao 1968 nian, p.~434; CNS, 184 (August 24, 1967),
  pp.~6-7.

  以下关于金华事件的描述依据 ``The `August 3' Incident in Jinhua
  area''，URS，48:18（1967年9月1日），第247-255页；ZPS，1967年8月11日，SWB/FE/2542/B/4-5；ZPS，1967年8月14日，SWB/FE/2546/B/13；ZPS，1967年8月19日，SWB/FE/2549/B/12-13；ZPS，1967年8月26日，SWB/FE/2557/B/18-19；Zhonggong
  nianbao 1968 nian，第434页；CNS，184（1967年8月24日），第6-7页。}
Military forces stationed in the city had split over their support for
different mass organizations. The militia intervened on the side of army
units opposed to the mass organization which was aligned with United
Headquarters. On 8 August, the day of the announcement of the formation
of the provisional supreme organ of power, this body and United
Headquarters held a joint press to denounce those behind the incident in
Jinhua. The five-day interlude between the beginning of the trouble and
the press conference may have been due to the imminent change in the
provincial leadership. Alternatively, it may have been the result of a
break in communications between Jinhua and Hangzhou. Or it may have
reflected an inability on the part of the provincial leadership to
arrive at a consensus as to how best to deal with the problem. The
supreme organ of power issued a notice holding the local party
authorities in Jinhua, together with PLA units garrisoned in the town,
responsible for summoning peasants into Jinhua to smash the rebels. The
provincial authorities requested United Headquarters' opponents to
abandon their organization, the Jinhua Workers' Revolutionary Rebel
Headquarters (金华工人革命造反总部) and its ``handful of bad leaders''
who hopefully would confess their crimes and turn over a new leaf with a
view to rejoining the revolutionary ranks.

联合总部及其盟友被击败的非同寻常事件发生在距省会更远的地方。8月3日，浙中金华市发生骚乱。\footnote{The
  following description of events at Jinhua is based on ``The `August 3'
  Incident in Jinhua area'', URS, 48:18 (September 1, 1967), pp.~247-55;
  ZPS, August 11, 1967, SWB/FE/2542/B/4-5; ZPS, August 14, 1967,
  SWB/FE/2546/B/13; ZPS, August 19, 1967, SWB/FE/2549/B/12-13; ZPS,
  August 26, 1967, SWB/FE/2557/B/18-19; Zhonggong nianbao 1968 nian,
  p.~434; CNS, 184 (August 24, 1967), pp.~6-7.

  以下关于金华事件的描述依据 ``The `August 3' Incident in Jinhua
  area''，URS，48:18（1967年9月1日），第247-255页；ZPS，1967年8月11日，SWB/FE/2542/B/4-5；ZPS，1967年8月14日，SWB/FE/2546/B/13；ZPS，1967年8月19日，SWB/FE/2549/B/12-13；ZPS，1967年8月26日，SWB/FE/2557/B/18-19；Zhonggong
  nianbao 1968 nian，第434页；CNS，184（1967年8月24日），第6-7页。}
驻城军事力量因支持不同群众组织而分裂。民兵站在反对与联合总部结盟的群众组织的部队一边介入。8月8日，即宣布成立临时最高权力机构之日，该机构和联合总部联合举行记者会，谴责金华事件背后的人。从骚乱开始到记者会之间相隔五天，可能是因为省级领导层即将发生变化。也可能是因为金华与杭州之间通信中断。或者，它反映出省级领导层无法就如何最好处理这一问题达成共识。最高权力机构发布通告，认定金华地方党政当局和驻城解放军部队共同负责召集农民进金华砸造反派。省级当局要求联合总部的反对者脱离自己的组织``金华工人革命造反总部''和其中``一小撮坏头头''，并希望这些头头坦白罪行、改过自新，以便重新加入革命队伍。

This soft approach reflected a recognition of political reality. Without
a massive commitment of main force troops, a step that courted the
danger of an escalation in violence, the provincial leaders in Zhejiang
could do nothing but compromise with their nominal subordinates.
Although the representative of the new power-brokers in Zhejiang adopted
a tough line in his speech at the press conference, words were no
substitute for deeds and the Rebel Headquarters was undoubtedly aware of
the strength of its position.

这种温和办法反映了对政治现实的承认。如果不大量投入主力部队，而这一步又有使暴力升级的危险，浙江省级领导人除了同名义上的下级妥协之外别无他法。尽管浙江新权力人物的代表在记者会上讲话采取强硬路线，但言辞不能代替行动，造反总部无疑清楚自身地位的力量。

The notice of the provincial authorities was transmitted to Jinhua on 9
August and carried the four kilometers from the normal college into town
by students of Jinhua 2nd Headquarters (金二司), the local ally of
United Headquarters. The college appears to have been the group's
stronghold and it must have been a lonely and exposed place in a sea of
peasant hostility. Fighting continued on 10 August with the Zhejiang
Daily encouraging the 2nd Headquarters. The editorial of the following
day somewhat plaintively rebuked Rebel Headquarters for challenging the
hegemony of United Headquarters, which it described as ``the core
organization of Zhejiang's revolutionaries who have the resolute support
of the province''.

省级当局的通告于8月9日传到金华，并由当地联合总部盟友金二司学生从师范学院带着走四公里进城。师范学院似乎是该集团据点，在一片农民敌意的海洋中，它一定是一个孤立而暴露的地方。8月10日战斗继续，《浙江日报》鼓励金二司。次日社论略带哀求地斥责造反总部挑战联合总部的霸权，并把联合总部描述为``得到全省坚决支持的浙江革命派核心组织''。

At a rally in Hangzhou on 13 August, the divisions within the military
over the extent of its commitment to United Headquarters were in further
evidence. A representative of the proletarian revolutionaries, in
praising the PLA's role in ``supporting the left'', specifically
mentioned Unit 7350 but not Unit 6409. This was a conspicuous omission.
The commander of the Air Force unit, Bai Zongshan, returning the
compliment, pledged the Air Force's assistance to the beleaguered
revolutionaries allied to United Headquarters in both Jinhua and
Wenzhou. Meng Zhaoyu (孟昭煜), representing PLA ground forces, also
pledged support. Resistance in Jinhua to troops despatched from Hangzhou
continued until 24 August. Finally, on the afternoon of that day, a
celebration rally could safely be held in the city. On 25 August
supporters of the 2nd Headquarters ventured out from the city to stage
demonstrations in two towns on the outskirts of Jinhua. Resistance had
finally crumbled before the superior fire-power of outside military
units, which had acted more to restore order and secure obedience to
provincial authority than to enter the lists for United Headquarters.

8月13日杭州一次集会上，军方内部对支持联合总部应投入到何种程度的分歧进一步显现。一名无产阶级革命派代表在赞扬人民解放军``支左''作用时，特别提到7350部队，却没有提到6409部队。这一遗漏十分醒目。空军部队司令员白宗善投桃报李，承诺空军援助金华和温州两个地方与联合总部结盟、处境艰难的革命派。代表解放军地面部队的孟昭煜也承诺支持。从杭州派往金华的部队遭到抵抗，直到8月24日才结束。终于，当天下午，城内可以安全地举行庆祝集会。8月25日，金二司支持者走出城区，到金华郊外两个镇示威。抵抗最终在外来军事部队优势火力面前崩溃；这些部队行动更多是为了恢复秩序、确保服从省级权威，而不是为联合总部参战。

Opposition in the isolated southern port city of Wenzhou proved even
more difficult to subdue than in Jinhua. There was no rail link to
Hangzhou, communication being either by boat from Ningbo or Shanghai,
rail and road from Jinhua, or the narrow, mountainous and hazardous
unpaved road from Hangzhou. As in Jinhua the rebels affiliated with
United Headquarters faced a well-entrenched local administration backed
by the military, the militia and a population resentful of and hostile
to outsiders.\footnote{Ken Ling has described the resentment of the
  people of Fuzhou toward rebels from Xiamen. See The Revenge of Heaven,
  chs.~5-6.

  Ken Ling 描述过福州人对厦门造反派的怨恨。见 The Revenge of
  Heaven，第5-6章。} Membership of the largest mass organization, the
Wenzhou Revolutionary Rebel General Headquarters
(温州市革命造反总指挥部) was drawn from the city's workers and suburban
commune members. This faction was loyal to and supported by the local
party leaders. It was led by a thirty-three year-old political
instructor named Yao Guolin (姚国麟). United Headquarters' followers
were grouped in an organization called the Wenzhou Workers' General
Headquarters (工总司).\footnote{The following description of events in
  Wenzhou is based on ZPS, August 18 and 19, 1967, SWB/FE/2549/B/12-18;
  ZPS, August 20, 1967, SWB/FE/2552/B/10-11; ZPS, August 21, 1967,
  SWB/FE/2557/B/17-20; ZPS, August 28, 1967, SWB/FE/2558/B/5-6; ZPS,
  September 5, 1967, SWB/FE/2571/B/5; Zhonggong nianbao 1968 nian,
  p.~434; Y. S. Chien, China's Fading Revolution: Army Dissent and
  Military Divisions, 1967-68. (Hong Kong: Centre of Contemporary
  Chinese Studies, 1969), pp.~50-1; BR, no. 42 (October 20, 1986), p.14.

  以下关于温州事件的描述依据
  ZPS，1967年8月18日和19日，SWB/FE/2549/B/12-18；ZPS，1967年8月20日，SWB/FE/2552/B/10-11；ZPS，1967年8月21日，SWB/FE/2557/B/17-20；ZPS，1967年8月28日，SWB/FE/2558/B/5-6；ZPS，1967年9月5日，SWB/FE/2571/B/5；Zhonggong
  nianbao 1968 nian，第434页；Y. S. Chien，China's Fading Revolution:
  Army Dissent and Military Divisions, 1967-68（香港：Centre of
  Contemporary Chinese Studies，1969年），第50-51页；BR，no.
  42（1986年10月20日），第14页。} On 14 August 1967, United Headquarters
issued a statement condemning the suppression of its comrades in
Wenzhou. Three days later it called a rally which was attended by Nan
Ping, the newly-installed head of the provincial military control
commission. On the previous day Zhou Enlai had ordered the ZPMD to
disarm Rebel General HQ and go to the aid of its encircled opponents.
The provisional supreme organ of power accordingly issued a circular to
this effect. Zhejiang Daily's near-hysterical prose indicated the
parlous state in which United Headquarters' allies found themselves.

孤立的南部港口城市温州，其反对力量比金华更难制服。当地没有通往杭州的铁路，交通要么靠从宁波或上海坐船，要么经金华铁路和公路，要么走从杭州出发的狭窄、崎岖、危险且未铺设的山路。与金华一样，隶属联合总部的造反派面对的是一个根深蒂固的地方行政体系，其背后有军队、民兵以及对外来者怀有怨恨和敌意的民众支持。\footnote{Ken
  Ling has described the resentment of the people of Fuzhou toward
  rebels from Xiamen. See The Revenge of Heaven, chs.~5-6.

  Ken Ling 描述过福州人对厦门造反派的怨恨。见 The Revenge of
  Heaven，第5-6章。}
最大群众组织温州市革命造反总指挥部的成员来自城市工人和郊区公社社员。该派忠于并受到地方党政领导支持，由三十三岁的政治教导员姚国麟领导。联合总部追随者则组成名为温州工人总司令部（工总司）的组织。\footnote{The
  following description of events in Wenzhou is based on ZPS, August 18
  and 19, 1967, SWB/FE/2549/B/12-18; ZPS, August 20, 1967,
  SWB/FE/2552/B/10-11; ZPS, August 21, 1967, SWB/FE/2557/B/17-20; ZPS,
  August 28, 1967, SWB/FE/2558/B/5-6; ZPS, September 5, 1967,
  SWB/FE/2571/B/5; Zhonggong nianbao 1968 nian, p.~434; Y. S. Chien,
  China's Fading Revolution: Army Dissent and Military Divisions,
  1967-68. (Hong Kong: Centre of Contemporary Chinese Studies, 1969),
  pp.~50-1; BR, no. 42 (October 20, 1986), p.14.

  以下关于温州事件的描述依据
  ZPS，1967年8月18日和19日，SWB/FE/2549/B/12-18；ZPS，1967年8月20日，SWB/FE/2552/B/10-11；ZPS，1967年8月21日，SWB/FE/2557/B/17-20；ZPS，1967年8月28日，SWB/FE/2558/B/5-6；ZPS，1967年9月5日，SWB/FE/2571/B/5；Zhonggong
  nianbao 1968 nian，第434页；Y. S. Chien，China's Fading Revolution:
  Army Dissent and Military Divisions, 1967-68（香港：Centre of
  Contemporary Chinese Studies，1969年），第50-51页；BR，no.
  42（1986年10月20日），第14页。}
1967年8月14日，联合总部发表声明，谴责其温州同志受到压制。三天后，它召开集会，新任省军管会负责人南萍出席。前一天，周恩来命令浙江省军区解除造反总指挥部武装，并援助被包围的对手。临时最高权力机构随即发布通报。《浙江日报》近乎歇斯底里的文风显示，联合总部盟友处境岌岌可危。

The Rebel General HQ was eventually subdued by a combined force of Red
Guards from Shanghai and Beijing, backed up on the ground by two units
under direct orders from the central military command. By 19 August an
editorial of Zhejiang Daily could report this news to its readers. On 20
August, at a rally held in Wenzhou, the commander of the units
supporting the left directed members of the Workers' General HQ to
target the local party leader Wang Fang and to differentiate between the
``handful of bad leaders'' and the ``hoodwinked masses'' of Rebel
General HQ. Wang Fang was not publicly named at the time but it is clear
from accusations made against him in 1968 that he was the party official
to whom speakers at the rally referred. He was accused of manipulating
Rebel General HQ from behind the scenes and secretly instructing it to
sabotage production. Representatives of United Headquarters and rebels
in Beijing, Shanghai and Wuhan also spoke at the rally, indicating the
extent of outside personnel involved in the suppression of Rebel General
HQ.

造反总指挥部最终被来自上海和北京的红卫兵联合力量制服，地面上还有两个直接受中央军令指挥的部队支持。到8月19日，《浙江日报》社论已能向读者报告这一消息。8月20日，在温州举行的集会上，支左部队指挥员指示工总司成员把目标对准地方党内领导人王芳，并把造反总指挥部的``一小撮坏头头''同``受蒙蔽的群众''区别开来。王芳当时没有被公开点名，但从1968年对他的指控来看，集会发言人所指的党内干部显然就是他。他被指控在幕后操纵造反总指挥部，并秘密指示其破坏生产。联合总部以及北京、上海和武汉造反派代表也在集会上讲话，显示参与镇压造反总指挥部的外来人员规模之大。

It was clear, however, that continued resistance to the imposition of
military control in Wenzhou worried the provincial leadership. An
editorial of 21 August admitted that a titanic battle continued in the
city and its environs. Although it claimed that Rebel General HQ had
collapsed as an organization, its leaders remained obdurate. Its
strength came from the support of workers in the port, transport,
construction and other industries who had evidently stopped work,
cutting off the city from the rest of the country. The rebel supporters
of United Headquarters did not escape all blame for the continued
upheavals and under the supervision of the PLA some were herded into
study classes.

然而，温州对实行军事管制的持续抵抗显然令省级领导层担忧。8月21日一篇社论承认，城市及其周边仍在进行一场巨大战斗。尽管社论声称造反总指挥部作为组织已经垮台，其领导人仍顽固不化。它的力量来自港口、交通、建筑和其他行业工人的支持，这些工人显然已停止工作，使城市同全国其他地方隔绝。支持联合总部的造反派并非完全免于对持续动荡的责任，在解放军监督下，其中一些人被赶进学习班。

The impact of the Jinhua and Wenzhou episodes reverberated across the
province. At a meeting of the Hangzhou Workers' Revolutionary Committee
held late in August, accusations were heard that the ``party persons in
authority'' in the two cities had ``wrecked'' the Cultural Revolution.
Further outbreaks of armed struggle were reported from Wenzhou in
September and October, proof that the situation remained highly
volatile. United Headquarters may have won some victories but it had
also suffered major losses. The overall winner from this spate of
violent incidents was probably the military which by September was
empowered to defend itself.

金华和温州事件的影响回荡全省。8月下旬杭州工人革命委员会一次会议上，有人指控两市``党内当权派''破坏文化大革命。9月和10月又有温州武斗爆发的报道，证明局势仍极不稳定。联合总部可能取得了一些胜利，但也遭受重大损失。这轮暴力事件的总体赢家很可能是军方，因为到9月它已被授权保卫自己。

While the fighting and mutual recriminations of August reverberated
across Zhejiang, rallies and editorials also preached the virtues of
unity.\footnote{URS, 48:15, pp.~208-17.

  URS，48:15，第208-217页。} It was pointed out that the enemy took
advantage of divisions among the rebels. In an article putting forward
the view that ``all forces that can be brought together should be
united'', Zhejiang Daily offered the olive branch to Red Storm. But the
sticking-point remained the condition that any alliance should have
United Headquarters at its ``core''. The other approach, which the
editorials continued to adopt, involved the attempt to separate the
membership of Red Storm from its leadership and thereby weaken the
organization. As one editorial stated condescendingly:

8月的战斗和相互指责在浙江各地回荡之际，集会和社论也在宣扬团结的美德。\footnote{URS,
  48:15, pp.~208-17.

  URS，48:15，第208-217页。}
有人指出，敌人利用造反派之间的分裂。《浙江日报》在一篇提出``凡是可以团结的力量都要团结起来''的文章中，向红暴伸出橄榄枝。但症结仍然是任何联盟都必须以联合总部为``核心''。社论继续采取的另一种办法，是试图把红暴成员同其领导层分离，从而削弱该组织。正如一篇社论居高临下地说：

\begin{quote}
Most of the hoodwinked masses are our class brothers. When they rise to
rebel, we should resolutely support them. When they make progress, which
may be very insignificant, we should appear to be pleased and give them
a clapping. We should never stand in their way, \ldots{} much less
attack them.
\end{quote}

\begin{quote}
大多数受蒙蔽的群众是我们的阶级兄弟。当他们起来造反时，我们应当坚决支持他们。当他们有进步时，哪怕很微小，我们也应当表示高兴并为他们鼓掌。我们决不应挡他们的路，\ldots\ldots 更不能攻击他们。
\end{quote}

The difficulty confronting those who espoused unity was to hit upon a
successful method to achieve it. Unity of the two mass organizations
from the top entailed the risk of dissent at the grassroots and the
threat of boycott. Unity from the bottom, factory by factory and school
by school, promised to be a tortuously slow process, subject to delay
and interference. The latter approach was exemplified by the
announcement on 25 August 1967, amid great fanfare, of the establishment
of revolutionary committees at the Zhejiang Fine Arts College, Hangzhou
Daily and other units where United Headquarters' strength
lay.\footnote{ZPS, August 25, 1967, SWB/FE/2562/B/9-10.

  ZPS，1967年8月25日，SWB/FE/2562/B/9-10。} Zhang Yongsheng, who had
graduated from the fine arts college in 1966, was elected Chairman of
its revolutionary committee. He boasted that Red Storm had not been able
to recruit a single member there. As a sign of the return to normality
the college had resumed classes on 14 July 1967.\footnote{URS, 48:13
  (August 15, 1967), pp.~181-82.

  URS，48:13（1967年8月15日），第181-182页。} Simultaneous with the
attempt to build up revolutionary committees from the grassroots,
Beijing persisted with its efforts to impose an agreement on the two
organizations. In August 1967 two documents to this effect had been
signed, perhaps under duress, because events showed that local opinion
was not ready for major concessions. Zhejiang Daily persisted in
upholding the status of United Headquarters, enumerating its
achievements and praising its steadfastness under pressure. To sweeten
the message for Red Storm, the paper was willing to concede that all
organizations, no matter what their size or when they had rebelled,
deserved equal treatment.\footnote{ZPS, August 17, 1967,
  SWB/FE/2549/B/9-12.

  ZPS，1967年8月17日，SWB/FE/2549/B/9-12。} Equality, that is, under the
baton of United Headquarters.

主张团结者面临的困难，是找到一种成功实现团结的方法。两个群众组织自上而下的统一，冒着基层异议和抵制的风险。自下而上、逐厂逐校的统一，则注定是一个极其缓慢、容易拖延和受干扰的过程。后一种办法的例子，是1967年8月25日大张旗鼓宣布在浙江美术学院、《杭州日报》和其他联合总部力量所在单位成立革命委员会。\footnote{ZPS,
  August 25, 1967, SWB/FE/2562/B/9-10.

  ZPS，1967年8月25日，SWB/FE/2562/B/9-10。}
1966年毕业于美术学院的张永生当选该院革命委员会主任。他吹嘘说，红暴在那里连一个成员也招不到。作为恢复正常的标志，该院已于1967年7月14日复课。\footnote{URS,
  48:13 (August 15, 1967), pp.~181-82.

  URS，48:13（1967年8月15日），第181-182页。}
在试图从基层建立革命委员会的同时，北京继续努力把协议强加给两个组织。1967年8月签署了两份有关文件，也许是在胁迫下签的，因为事态表明地方意见尚未准备好作出重大让步。《浙江日报》继续维护联合总部地位，列举其成就并赞扬其在压力下的坚定。为了让红暴更容易接受，该报愿意承认所有组织，不论大小或造反时间早晚，都应得到平等待遇。\footnote{ZPS,
  August 17, 1967, SWB/FE/2549/B/9-12.

  ZPS，1967年8月17日，SWB/FE/2549/B/9-12。}
也就是说，在联合总部指挥棒下的平等。

By the end of August, following the harrowing month of fighting. United
Headquarters moved to patch up its relations with the military. The
latter reciprocated. At a series of rallies which continued into
mid-September 1967, an atmosphere of mutual self-criticism and feigned
modesty prevailed.\footnote{URS, 48:25 (September 26, 1967), pp.~353-65;
  ZPS, September 1, 1967, SWB/FE/2567/B/10-12; ZPS, September 10 and 11,
  1967, SWB/FE/2581/B/13.

  URS，48:25（1967年9月26日），第353-365页；ZPS，1967年9月1日，SWB/FE/2567/B/10-12；ZPS，1967年9月10日和11日，SWB/FE/2581/B/13。}
On 31 August, commander Bai Zongshan asked the proletarian
revolutionaries of United Headquarters to criticize the mistakes which
had been committed by members of his Air Force unit. A representative of
United Headquarters responded by saying that

到8月底，在经历一个月令人痛苦的战斗之后，联合总部开始修补同军方的关系。后者也作出回应。在一系列持续到1967年9月中旬的集会上，双方自我批评和假装谦逊的气氛占据主导。\footnote{URS,
  48:25 (September 26, 1967), pp.~353-65; ZPS, September 1, 1967,
  SWB/FE/2567/B/10-12; ZPS, September 10 and 11, 1967, SWB/FE/2581/B/13.

  URS，48:25（1967年9月26日），第353-365页；ZPS，1967年9月1日，SWB/FE/2567/B/10-12；ZPS，1967年9月10日和11日，SWB/FE/2581/B/13。}
8月31日，司令员白宗善要求联合总部的无产阶级革命派批评其空军部队成员所犯错误。联合总部一名代表回应说：

\begin{quote}
When we were under suppression and persecution, you came to our rescue
and helped us to stand up, to rebel against the handful of party persons
in authority taking the capitalist road and the bourgeois reactionaries
and to gain one victory after another.
\end{quote}

\begin{quote}
当我们受到压制和迫害时，是你们前来营救我们，帮助我们站起来，起来反对一小撮党内走资本主义道路的当权派和资产阶级反动派，并取得一个又一个胜利。
\end{quote}

There was no reference to the other branches of the armed services. The
rebel representative also requested the Air Force personnel to ``make
higher and more strict demands on the proletarian revolutionaries and
help them to get rid of their selfishness''. Bai, in return, publicly
committed Unit 7350 to United Headquarters' demand for recognition as
the ``core'' of any alliance. At another rally held on 1 September Xiong
Yingtang, Commander of Unit 6409, made no such pledge on behalf of his
troops. On this occasion Zhang Yongsheng read out a pact of solidarity
with the PLA on behalf of United Headquarters.

讲话没有提到武装力量的其他军种。这名造反派代表还要求空军人员``对无产阶级革命派提出更高、更严格的要求，帮助他们克服私心''。作为回应，白公开承诺7350部队支持联合总部要求在任何联盟中被承认为``核心''。9月1日另一场集会上，6409部队司令员熊应堂并未代表自己的部队作出这种承诺。在这次场合，张永生代表联合总部宣读了同人民解放军团结的公约。

At the meetings the rebels' indiscipline, factionalism and anarchy was
compared unfavorably with the discipline and order displayed by the
military. A United Headquarters' resolution, adopted on 5 September,
stated in part that the revolutionary masses were

在这些会议上，造反派的无纪律、派性和无政府状态，与军方表现出的纪律和秩序形成不利对比。联合总部9月5日通过的一项决议部分指出，革命群众：

\begin{quote}
not allowed to use any pretext to slander or put blame on the military
organizations. There is no excuse for them to encroach upon the various
kinds of weapons, equipment and supplies belonging to the
PLA.\footnote{ZPS, September 5 and 7, 1967, SWB/FE/2571/B/5; ZPS,
  September 12, 1967, SWB/FE/2569/B/10.

  ZPS，1967年9月5日和7日，SWB/FE/2571/B/5；ZPS，1967年9月12日，SWB/FE/2569/B/10。}
\end{quote}

\begin{quote}
不准用任何借口诬蔑或归咎军事机关。不得侵占属于人民解放军的各种武器、装备和物资。\footnote{ZPS,
  September 5 and 7, 1967, SWB/FE/2571/B/5; ZPS, September 12, 1967,
  SWB/FE/2569/B/10.

  ZPS，1967年9月5日和7日，SWB/FE/2571/B/5；ZPS，1967年9月12日，SWB/FE/2569/B/10。}
\end{quote}

To back up this tough policy, the first in a series of trials conducted
by military tribunals took place on September 25. Four criminals were
summarily tried and two executed on charges of murder, arson and spying
for the USA and Taiwan.\footnote{ZPS, September 26, 1967,
  SWB/FE/2580/B/3-4.

  ZPS，1967年9月26日，SWB/FE/2580/B/3-4。} Members of the central
Cultural Revolution Group were rebuked for inciting Red Guards to take
up arms against the military and accused of forming a ``May 16''
Group.\footnote{See Barry Burton, ``The Cultural Revolution's
  ultra-leftist conspiracy: the `May 16 Group'\,'', AS, 11:11 (1971),
  pp.~1029-1053.

  见 Barry Burton，``The Cultural Revolution's ultra-leftist conspiracy:
  the `May 16 Group'\,''，AS，11:11（1971年），第1029-1053页。} A member
of the Group, Yao Wenyuan, distanced himself from his former colleagues
in the article ``Comments on Tao Zhu's Two Books'', which was published
in Hongqi and reprinted in People's Daily on September 8.\footnote{Ye
  Yonglie, 姚氏父子 (The Yao father and son), (Dalian: Dalian Chubanshe,
  1989), p.~295.

  Ye Yonglie，《姚氏父子》（The Yao father and son）（大连：Dalian
  Chubanshe，1989年），第295页。} Repudiation of ultra-leftism became
the main task for the rebels of United Headquarters. They were ordered
to study Jiang Qing's speech of September 5, representing her climbdown
on the utility of violence, and Yao's article.\footnote{See ZPS,
  September 21, 1967, SWB/FE/2580/B/9-10.

  见 ZPS，1967年9月21日，SWB/FE/2580/B/9-10。} A campaign to ``support
the army and love the people'' commenced at the end of August.

为了支持这一强硬政策，军事法庭进行的一系列审判中的第一次于9月25日举行。四名罪犯被简易审判，其中两人因谋杀、纵火和为美国、台湾充当间谍而被处决。\footnote{ZPS,
  September 26, 1967, SWB/FE/2580/B/3-4.

  ZPS，1967年9月26日，SWB/FE/2580/B/3-4。}
中央文革小组成员因煽动红卫兵拿起武器反对军队而受到斥责，并被指控成立``五一六集团''。\footnote{See
  Barry Burton, ``The Cultural Revolution's ultra-leftist conspiracy:
  the `May 16 Group'\,'', AS, 11:11 (1971), pp.~1029-1053.

  见 Barry Burton，``The Cultural Revolution's ultra-leftist conspiracy:
  the `May 16 Group'\,''，AS，11:11（1971年），第1029-1053页。}
小组成员姚文元在《红旗》发表、9月8日《人民日报》转载的《评陶铸的两本书》一文中，同昔日同事拉开距离。\footnote{Ye
  Yonglie, 姚氏父子 (The Yao father and son), (Dalian: Dalian Chubanshe,
  1989), p.~295.

  Ye Yonglie，《姚氏父子》（The Yao father and son）（大连：Dalian
  Chubanshe，1989年），第295页。}
批判极``左''主义成为联合总部造反派的主要任务。他们被命令学习江青9月5日讲话和姚的文章；江青讲话代表她在暴力效用问题上的退让。\footnote{See
  ZPS, September 21, 1967, SWB/FE/2580/B/9-10.

  见 ZPS，1967年9月21日，SWB/FE/2580/B/9-10。}
``拥军爱民''运动于8月底开始。

The battles of the summer of 1967 politically weakened United
Headquarters. It had been unable to eliminate its adversary Red Storm by
force of arms and the cost of continued factional warfare involving the
army had become too great. The vicious cycle of violence had to be
broken; only the central leadership and Mao in particular were capable
of putting an end to the civil strife which had devastated Zhejiang and
the rest of the country.

1967年夏季的战斗在政治上削弱了联合总部。它未能用武力消灭对手红暴，而涉及军队的持续派性战争代价已经过高。暴力的恶性循环必须被打破；只有中央领导层，尤其是毛本人，才有能力结束摧毁浙江和全国其他地区的内乱。

\section{Notes}\label{notes-2}

\section{注释}\label{ux6ce8ux91ca-2}

\bookmarksetup{startatroot}

\chapter{Chapter Three - FROM ``UNITY'' TO RENEWED DISUNITY, SEPTEMBER
1967-AUGUST 1968 / 第三章 -
从``团结''到重新分裂，1967年9月至1968年8月}\label{chapter-three---from-unity-to-renewed-disunity-september-1967-august-1968-ux7b2cux4e09ux7ae0---ux4eceux56e2ux7ed3ux5230ux91cdux65b0ux5206ux88c21967ux5e749ux6708ux81f31968ux5e748ux6708}

\section{Steps toward the Three-way Alliance, September-December
1967}\label{steps-toward-the-three-way-alliance-september-december-1967}

\section{走向三结合联盟，1967年9月至12月}\label{ux8d70ux5411ux4e09ux7ed3ux5408ux8054ux76df1967ux5e749ux6708ux81f312ux6708}

After the fighting of July to August 1967 the overriding source of
instability in Zhejiang remained, as it had done since February 1967,
the antagonism between the two mass organizations. Unity required
compromise. Yet United Headquarters remained adamant that it would not
accept its arch-rival into an alliance unless Red Storm accepted a
subordinate standing. Red Storm stood equally firm in its resolve.
Stalemates of this kind were not limited to Zhejiang, and between July
and September 1967 Mao Zedong visited several provinces in north,
central and east China to break the logjam.

1967年7月至8月的战斗之后，浙江不稳定的首要来源仍然是两个群众组织之间的对立，这一点自1967年2月以来一直如此。团结需要妥协。然而，联合总部仍坚持认为，除非红暴接受从属地位，否则不会把这个头号对手纳入联盟。红暴同样立场坚定。这类僵局并不限于浙江；1967年7月至9月，毛泽东走访华北、华中和华东若干省份，以打破僵局。

Mao spent one day in Hangzhou on 16 September. Only a very brief excerpt
of his talk with Nan Ping is available.\footnote{毛泽东思想万岁 (Long
  Live Mao Zedong Thought) 1969, (Taibei: Institute of International
  Relations, 1974), p.~682. On September 25, the provincial authorities
  held a rally to celebrate Mao's inspection. See ZJRB, September 26,
  1967, p.~1.} The Chairman objected to the practice of forcing cadres
to kneel and wear dunce's-caps and recommended the policy of
unity-criticism-unity in dealing with them. In a later talk in Jiangxi
province Mao predicted that the cadres of Zhejiang, including military
officials, would not stand for the continuation of such
treatment.\footnote{SCMP, 4070 (November 30, 1967), p.~9.}

9月16日，毛在杭州停留一天。关于他同南萍谈话的内容，现仅有极短摘录。\footnote{毛泽东思想万岁
  (Long Live Mao Zedong Thought) 1969, (Taibei: Institute of
  International Relations, 1974), p.~682. On September 25, the
  provincial authorities held a rally to celebrate Mao's inspection. See
  ZJRB, September 26, 1967, p.~1.}
主席反对强迫干部下跪、戴高帽的做法，并建议用``团结-批评-团结''的方针来对待干部。在随后于江西的一次谈话中，毛预言浙江干部，包括军事干部，不会容忍这种待遇继续下去。\footnote{SCMP,
  4070 (November 30, 1967), p.~9.}

After Mao's return to Beijing, extracts from his directives were
compiled and released by the CCP CC.\footnote{``中共中央通知'' (Notice
  of the CP CC), October 7, 1967, CCP Documents, pp.~545-556.} The theme
of the document was its emphasis on unity. Mao declared that ``there was
no fundamental conflict of interest within the working class'' and
therefore no need for it to split into two irreconcilable factions. He
attributed the disunity to three factors -- the infiltration of ``bad
people into mass organizations, the sabotage activities of
capitalist-roaders trying to protect themselves, and the influence of
anarchist thinking. If both sides strove for common ground on major
issues and put minor matters aside, stated the Chairman, unity was
possible. But this depended, as Mao realized, on one side abandoning its
claim to form the nucleus (核心) of the revolutionary great alliance.
The day after Mao's visit to Hangzhou, Zhejiang Daily published an
editorial repudiating the slogan''I want to be at the core'' which
United Headquarters had constantly repeated in its pronouncements
concerning the alliance.\footnote{ZPS, September 17, 1967,
  SWB/FE/2574/B/12-14.} The breakthrough cut the ground from under
United Headquarters and dealt a savage blow to its ambitions. But all
was not lost. Mao's formulations regarding unity clearly applied to
organizations described as ``revolutionary''; begging the question of
whether an organization such as Red Storm qualified.

毛回到北京后，中共中央汇编并发布了他的指示摘录。\footnote{``中共中央通知''
  (Notice of the CP CC), October 7, 1967, CCP Documents, pp.~545-556.}
该文件的主题是强调团结。毛宣称，``工人阶级内部没有根本利害冲突''，因此没有必要分裂成两个不可调和的派别。他把不团结归结为三个因素：`坏人'混入群众组织、走资派为保护自己而进行破坏活动，以及无政府主义思想的影响。主席说，如果双方在大问题上求同存异、小问题放在一边，团结是可能的。但毛也意识到，这取决于一方放弃自称为革命大联合核心的要求。毛访问杭州后的第二天，《浙江日报》发表社论，批判联合总部在有关联盟声明中不断重复的``我要当核心''口号。\footnote{ZPS,
  September 17, 1967, SWB/FE/2574/B/12-14.}
这一突破动摇了联合总部的根基，并沉重打击了它的野心。但并非一切都已失去。毛关于团结的表述显然适用于被称为``革命''的组织；问题在于红暴这样的组织是否符合这一资格。

Another issue concerned the status of cadres. Mao believed that the
majority were good, and that after being allowed the appropriate time to
admit and correct their past errors, they should be allowed to resume
public life. He viewed this question as having a crucial bearing on the
realization of the three-way alliance of mass representatives, cadres
and military personnel. The decision directing the PLA to take over
provincial administrations had been made at a time of social and
political chaos. Prolonged military rule in the provinces had never been
envisaged and its continuation was embarrassingly at odds with the
overriding objectives of the Cultural Revolution. In the meantime, to
demonstrate that the military itself was not above supervision, Mao
suggested the setting-up of study classes for military officers. Members
of the People's Armed Forces Department, who had directed and led the
suppression of the rebels in towns and villages across China, would
attend the first such classes.

另一个问题涉及干部地位。毛认为大多数干部是好的，在给他们适当时间承认并改正过去错误之后，应允许其恢复公开生活。他认为这一问题对实现群众代表、干部和军人三结合联盟具有关键影响。让人民解放军接管省级行政机构的决定，是在社会和政治混乱时期作出的。省内长期军管从来不是原先设想的目标，而且其延续同文化大革命的总目标明显相悖，令人尴尬。同时，为了表明军队本身并非不受监督，毛建议为军官开办学习班。曾在全国城乡指挥和领导镇压造反派的人民武装部成员，将参加第一批学习班。

During Mao's tour, agreement was reached on the formation of
revolutionary committees in provinces including Zhejiang.\footnote{See
  the speech by Zhang Chunqiao, September 29, 1967, SCMP, 4072, p.~2.}
Yet a long period of difficult and acrimonious discussions bracketed
this agreement and the formation of the ZPRC in March 1968. The
intervening period witnessed the reemergence of several important
civilian cadres from the old provincial and municipal administrations,
who provided some leavening of the overwhelming military presence in the
province. Problems in the agricultural sector caused by a drought in the
autumn of 1967, acted as an additional spur for the recall of
experienced civilian administrators. Peasants had diverted their
energies from grain production to sideline industries while their
immediate leaders, no doubt traumatized by recent events, watched from
the sidelines.\footnote{ZPS, November 14, 1967, CNS, 198 (November 30,
  1967), p.~10.} The reappearance of leading cadres encouraged
basic-level officials to return to their posts. Political reality and
administrative necessity quickly overrode the naive optimism, expressed
at the end of October, that ``genuine proletarian revolutionaries
comprise a majority in reorganized leadership teams at all
levels''.\footnote{ZPS, October 29, 1967, CNS, 195 (November 9, 1967),
  pp.~2-3.}

毛巡视期间，就包括浙江在内的一些省份成立革命委员会达成协议。\footnote{See
  the speech by Zhang Chunqiao, September 29, 1967, SCMP, 4072, p.~2.}
然而，从这一协议到1968年3月浙江省革命委员会成立，中间隔着一段漫长、困难且充满争吵的讨论。其间，旧省、市行政机构中若干重要文职干部重新出现，在全省压倒性的军事存在中起到某种调和作用。1967年秋旱造成的农业问题，又进一步促使当局召回有经验的文职行政人员。农民把精力从粮食生产转向副业，而他们的直接领导人大概因近期事件受到创伤，只能在旁观望。\footnote{ZPS,
  November 14, 1967, CNS, 198 (November 30, 1967), p.~10.}
主要干部复出鼓励了基层官员回到岗位。政治现实和行政需要很快压倒了10月底那种天真乐观的说法，即``各级改组后的领导班子中真正的无产阶级革命派占多数''。\footnote{ZPS,
  October 29, 1967, CNS, 195 (November 9, 1967), pp.~2-3.}

Wang Zida, the deputy head of the old CCP HMC and Mayor of Hangzhou was
one of the first senior cadres to break with his colleagues and throw in
his lot with the new order. On 24 August 1967 he published an article in
United Headquarters' newspaper calling Jiang Hua, although not by name,
the ``biggest royalist in Zhejiang'' and proclaimed: ``Overthrow XX
{[}Jiang Hua{]}, liberate Zhejiang''.\footnote{Hou Zheng, ``Analyze the
  black article written by that former principal responsible person of
  the Hangzhou Municipal Party Committee'', HZRB, October 4, 1978.}
Furthermore, on 26 August at a televised rally in Hangzhou, Wang
confessed his mistakes to an audience composed of ``revolutionary
masses''. Wang had been ``liberated'' after almost a year or so of
harassment.\footnote{See Wang's repentant article published in ZJRB,
  October 13, 1967, p.~2, and the gloating, condescending tone of the
  article by his accusers, the Hangzhou Liaison Station to ``Bombard the
  Municipal Party Committee'' (杭州市炮打市委联络站) in ibid. p.~3.}
Upon the formation of the preparatory committee of the HMRC on 25
October 1967, Wang was named its leader and was eventually elected
Chairman of the HMRC in December.\footnote{ZPS, October 25, 1967,
  SWB/FE/2605/B/19-20; ZPS, December 21, 1967, SWB/FE/2654/B/9-10.} Qiu
Qiang (邱强), Deputy-secretary of the CCP HMC before the Cultural
Revolution, also reemerged at this time and published a litany of his
past errors.\footnote{ZPS, October 7, 1967, SWB/FE/2591/B/17; ZJRB,
  October 17, 1967, p.~3. See also the article by former Deputy-mayor of
  Hangzhou, Zhou Fengming, on the same theme in ZJRB, October 25, 1967,
  p.~3.}

旧中共杭州市委副书记、杭州市长王子达，是最早同同僚决裂并投向新秩序的高级干部之一。1967年8月24日，他在联合总部报纸上发表文章，虽未点名，却称江华为``浙江最大的保皇派''，并宣称：``打倒XX{[}江华{]}，解放浙江''。\footnote{Hou
  Zheng, ``Analyze the black article written by that former principal
  responsible person of the Hangzhou Municipal Party Committee'', HZRB,
  October 4, 1978.}
此外，8月26日，在杭州一次电视转播集会上，王向由``革命群众''组成的听众承认错误。经过近一年左右的骚扰后，王被``解放''。\footnote{See
  Wang's repentant article published in ZJRB, October 13, 1967, p.~2,
  and the gloating, condescending tone of the article by his accusers,
  the Hangzhou Liaison Station to ``Bombard the Municipal Party
  Committee'' (杭州市炮打市委联络站) in ibid. p.~3.}
1967年10月25日杭州市革命委员会筹备委员会成立时，王被任命为负责人，并最终于12月当选杭州市革委会主任。\footnote{ZPS,
  October 25, 1967, SWB/FE/2605/B/19-20; ZPS, December 21, 1967,
  SWB/FE/2654/B/9-10.}
文革前中共杭州市委副书记邱强也在此时重新出现，并发表了一长串对自己过去错误的交代。\footnote{ZPS,
  October 7, 1967, SWB/FE/2591/B/17; ZJRB, October 17, 1967, p.~3. See
  also the article by former Deputy-mayor of Hangzhou, Zhou Fengming, on
  the same theme in ZJRB, October 25, 1967, p.~3.}

It was not until the end of the year that provincial-level cadres were
``liberated''. Secretary of the old CCP ZPC, Lai Keke, and a member of
its standing committee, Shen Ce, had both come under attack from Red
Storm and had therefore been protected by United Headquarters. At the
end of November 1967 both men published articles expressing support for
the Cultural Revolution.\footnote{ZJRB, November 29, 1967, pp.~3, 4.} A
long-time colleague of Jiang Hua, Wu Xian, was also permitted to join
the provisional leadership of the province.\footnote{The three senior
  cadres attended a rally held to prepare for the formation of the
  provincial revolutionary committee. ZPS, November 30, 1967,
  SWB/FE/2637/B/8.} A leading united front figure and brother of Lu Xun,
former Governor Zhou Jianren (周建人), also reappeared at a rally on
December 6, 1967, in his capacity as Vice-chairman of the
NPC.\footnote{ZPS, December 6, 1967, SWB/FE/2647/B/9-10. Zhou
  (1897-1984) was born in Shaoxing. For his obituary, see ZJRB, August
  7, 1984.} He was undoubtedly trotted out to add respectability and
credibility to the new administration.

直到年底，省级干部才被``解放''。旧中共浙江省委书记赖可可和省委常委沈策都曾受到红暴攻击，因此受到联合总部保护。1967年11月底，两人发表文章表示支持文化大革命。\footnote{ZJRB,
  November 29, 1967, pp.~3, 4.}
江华的长期同事吴宪也获准加入全省临时领导层。\footnote{The three senior
  cadres attended a rally held to prepare for the formation of the
  provincial revolutionary committee. ZPS, November 30, 1967,
  SWB/FE/2637/B/8.}
重要统战人物、鲁迅之弟、前省长周建人也于1967年12月6日以全国人大副委员长身份重新出现在一次集会上。\footnote{ZPS,
  December 6, 1967, SWB/FE/2647/B/9-10. Zhou (1897-1984) was born in
  Shaoxing. For his obituary, see ZJRB, August 7, 1984.}
他无疑是被请出来为新行政机构增添体面和可信度。

Was there any pattern in the reappearance of former high-ranking
officials at this time? Frederick Teiwes has categorized four cadre
sub-types who joined provincial revolutionary committees during the
Cultural Revolution:\footnote{Frederick C. Teiwes, Provincial Leadership
  in China: The Cultural Revolution and Its Aftermath, Cornell
  University East Asian Papers, No.~4. (Ithaca, New York: Cornell
  University, 1974), pp.~33-8.} first, there were members of the old
elite who managed to break ties with discredited colleagues; second,
those cadres who were outsiders to the particular province in which they
served in 1966; third, cadres who had previously been purged;\footnote{Another
  observer has noted that cadres disciplined before the Cultural
  Revolution were more likely to support the radicals. See Hong Yung
  Lee, The Politics of the Chinese Cultural Revolution, p.~340.} and
fourth, united front cadres.

此时原高级官员重新出现是否有某种模式？Frederick
Teiwes把文化大革命期间加入省级革命委员会的干部分成四类：\footnote{Frederick
  C. Teiwes, Provincial Leadership in China: The Cultural Revolution and
  Its Aftermath, Cornell University East Asian Papers, No.~4. (Ithaca,
  New York: Cornell University, 1974), pp.~33-8.}
第一，设法同失势同僚断绝关系的旧精英成员；第二，1966年在任职省份中属于外来者的干部；第三，先前曾被清洗的干部；\footnote{Another
  observer has noted that cadres disciplined before the Cultural
  Revolution were more likely to support the radicals. See Hong Yung
  Lee, The Politics of the Chinese Cultural Revolution, p.~340.}
第四，统战干部。

Lai Keke fits neatly into categories one, two and three and Zhou Jianren
into category four. Teiwes suggests that ambition may have driven
veteran cadres to break ties of solidarity and climb on board the
victorious bandwagon. Wang Zida may certainly slot into this category.
He and the head of the CCP HMC, Wang Pingyi, had both served as
subordinates to Wu Xian on that body for many years before Wang Pingyi
was promoted above Wang Zida in about 1964. The latter Wang may have
harbored feelings of resentment and taken advantage of Wang Pingyi's
downfall in the Cultural Revolution, to which he in no small way
contributed, to seize the top spot in Hangzhou. The case of Wu Xian,
however, defies such categorization. The fact that Wu later lost his
post on the ZPRC, and was accused of secretly supporting his old chief
Jiang Hua, is evidence that other factors were involved.

赖可可恰好符合第一、第二和第三类，周建人则符合第四类。Teiwes认为，抱负可能驱使老干部切断团结纽带，跳上胜利者的战车。王子达当然可能归入这一类。他和中共杭州市委负责人王平夷，多年来都曾在吴宪领导下的市委任下属，直到约1964年王平夷升到王子达之上。后一位王可能心怀怨恨，并利用王平夷在文化大革命中的垮台，且他本人对此也有不小贡献，夺取杭州最高位置。然而，吴宪的个案无法作如此归类。吴后来失去省革委会职务，并被指控秘密支持旧上司江华，这说明还存在其他因素。

Representatives of mass organizations were another component of the
three-in-one alliance. As a pre-condition to selection, they were
instructed to form ``revolutionary great alliances''. Two days after
Mao's September stopover in Hangzhou, the unity of four rebel groups at
the Hangzhou silk complex was announced.\footnote{ZJRB, September 18,
  1967, p.~2.} Several days later an article written by a representative
of each organization, including Weng Senhe on behalf of the Red Rebel
Corps, appeared in the Zhejiang Daily.\footnote{ZJRB, September 23,
  1967, p.~2.} Rejecting charges that the alliance was an unprincipled
compromise (和稀泥), the article detailed the difficulties which had
been surmounted in achieving unity. The authors also attributed success
to the ``concern'' and ``assistance'' of the provincial military control
commission and rebel groups in Shanghai and Zhejiang. However, the
alliance at the silk mill in fact proved to be of short duration.

群众组织代表是``三结合''联盟的另一组成部分。作为入选前提，他们被要求组成``革命大联合''。毛9月在杭州短暂停留后两天，杭州丝绸联合厂四个造反团体宣布团结。\footnote{ZJRB,
  September 18, 1967, p.~2.}
几天后，每个组织一名代表撰写的文章发表在《浙江日报》上，其中包括代表红色造反兵团的翁森鹤。\footnote{ZJRB,
  September 23, 1967, p.~2.}
文章驳斥所谓联盟是无原则妥协（和稀泥）的指责，详细说明了实现团结时克服的困难。作者还把成功归因于省军管会以及上海、浙江造反团体的``关怀''和``帮助''。然而，该厂的联盟事实上很短命。

As another means of organizing the rebels into manageable formations,
and also to break down cross-unit ties, congresses of workers, peasants
and Red Guards were convened. These congresses in effect became
substitutes for the atrophied trade unions, peasant associations and the
CYL. Meetings to set up the Zhejiang Workers' Congress and congresses of
peasants and Red Guards took place over the next few months. As
revolutionary committees were set up in various grass-roots units,
rallies, conferences and meetings proclaimed the virtue of unity. By
late October it was claimed that alliances had been concluded in 90\% of
Hangzhou's factories and enterprises. The spurious nature of
revolutionary committees hastily put together in August 1967 was
exemplified by the case of Hangzhou University. It announced an alliance
of its two factions, a precondition for electing a revolutionary
committee, two months after the committee had in fact been
established.\footnote{ZPS, October 28, 1967, SWB/FE/2607/B/1. In
  Xiangshan county a revolutionary committee was established on May 6,
  1967, dissolved the following October and not reestablished until
  November 1968. Xiangshan xianzhi, p.~35. In Linhai county a similar
  series of events occurred. See Linhai xianzhi, pp.~28, 29.}

作为把造反派组织成可管理形态、并打破跨单位联系的另一手段，工人、农民和红卫兵代表大会被召集起来。这些大会实际上成为已萎缩的工会、农会和共青团的替代物。随后几个月，筹建浙江省工人代表大会以及农民、红卫兵代表大会的会议陆续举行。随着各基层单位成立革命委员会，集会、会议和大会都宣扬团结的美德。到10月底，据称杭州90\%的工厂企业已经达成联盟。1967年8月仓促拼凑出的革命委员会，其虚假性以杭州大学为例最为明显：该校在革命委员会实际成立两个月后，才宣布两个派别实行联盟，而这本是选举革命委员会的前提。\footnote{ZPS,
  October 28, 1967, SWB/FE/2607/B/1. In Xiangshan county a revolutionary
  committee was established on May 6, 1967, dissolved the following
  October and not reestablished until November 1968. Xiangshan xianzhi,
  p.~35. In Linhai county a similar series of events occurred. See
  Linhai xianzhi, pp.~28, 29.}

\section{The Formation of the Zhejiang Provincial Revolutionary
Committee, March
1968}\label{the-formation-of-the-zhejiang-provincial-revolutionary-committee-march-1968}

\section{浙江省革命委员会的成立，1968年3月}\label{ux6d59ux6c5fux7701ux9769ux547dux59d4ux5458ux4f1aux7684ux6210ux7acb1968ux5e743ux6708}

Clearly, the Zhejiang provincial leadership was anxious to please its
superiors in Beijing by completing arrangements as expeditiously as
possible for the formation of the ZPRC. Beginning in September 1967, the
various interested parties had held discussions in Beijing under central
supervision.\footnote{``The case of Weng Sengho'', pp.~1-2.} On 30
November the provisional supreme organ of power and United Headquarters
convened a rally at which it was declared that ``The revolutionary great
alliance is being continuously developed and consolidated, and
revolutionary leading cadres have already stepped forward''. The report
boasted that

显然，浙江省领导层急于尽快完成成立省革委会的安排，以取悦北京上级。从1967年9月开始，各相关方面在中央监督下于北京举行讨论。\footnote{``The
  case of Weng Sengho'', pp.~1-2.}
11月30日，临时最高权力机构和联合总部召开集会，会上宣布``革命大联合正在不断发展和巩固，革命领导干部已经站出来了''。报告夸称：

\begin{quote}
Conditions are fully ripe for establishing the Red political power --
the Zhejiang Provincial Revolutionary Committee. To form it as soon as
possible is the common aspiration of the 30 million people of the
Province.\footnote{ZPS, November 30, 1967, SWB/FE/2637/B/8.}
\end{quote}

\begin{quote}
建立红色政权 -- 浙江省革命委员会 --
的条件已经完全成熟。尽快成立它，是全省三千万人民的共同愿望。\footnote{ZPS,
  November 30, 1967, SWB/FE/2637/B/8.}
\end{quote}

No alliance was possible without a breakthrough in the negotiations
between Red Storm and United Headquarters. In early December 1967 it was
Mao Zedong who intervened to push these talks forward. Mao's comments
concerning Red Storm formed the basis of the central notice disseminated
across the country.\footnote{``中共中央国务院中央军委中央文革小组关于正确对待犯过错误的老造反派的通知''
  (Notice of the CCP CC, State Council, Military Affairs Commission and
  Central Cultural Revolution Group on handling old rebel groups which
  have committed mistakes), December 2, 1967, CCP Documents, pp.~623-26.}
Mao appraised Red Storm thus:

没有红暴与联合总部谈判取得突破，就不可能有联盟。1967年12月初，是毛泽东亲自介入推动谈判前进。毛关于红暴的评论成为中央向全国发布通知的基础。\footnote{``中共中央国务院中央军委中央文革小组关于正确对待犯过错误的老造反派的通知''
  (Notice of the CCP CC, State Council, Military Affairs Commission and
  Central Cultural Revolution Group on handling old rebel groups which
  have committed mistakes), December 2, 1967, CCP Documents, pp.~623-26.}
毛这样评价红暴：

\begin{quote}
The Red Storm of Zhejiang, unlike the Million Warriors of
Hubei,\footnote{The Million Warriors was the ``conservative''
  organization in Wuhan backed by the Commander of the military region,
  Chen Zaidao. For Chen's recent account of the ``Wuhan Incident'' see
  浩劫中的一幕 (One Episode in the Calamity), (Beijing: Jiefangjun
  Chubanshe, 1989).} is an old rebel faction which has made mistakes. It
has a lot of mass support. It seems that the principle to be adopted
toward it is one of help, criticism and unity.
\end{quote}

\begin{quote}
浙江的红暴不同于湖北的百万雄师，\footnote{The Million Warriors was the
  ``conservative'' organization in Wuhan backed by the Commander of the
  military region, Chen Zaidao. For Chen's recent account of the ``Wuhan
  Incident'' see 浩劫中的一幕 (One Episode in the Calamity), (Beijing:
  Jiefangjun Chubanshe, 1989).}
它是犯过错误的老造反派。它有很多群众。看来对它应采取帮助、批评、团结的方针。
\end{quote}

The central document directed that ``It is not proper to adopt a policy
of exerting pressure and imposing absolute exclusion against the masses
of such organizations''. It also admitted realistically that

中央文件指示，``对这样的组织的群众，采取压服、绝对排斥的方针是不适当的''。文件还现实地承认：

\begin{quote}
the use of methods such as high pressure, exclusion, power seizure by
one faction alone, etc. cannot solve this kind of contradiction.
\end{quote}

\begin{quote}
采取高压、排斥、一派单独夺权等办法，不能解决这类矛盾。
\end{quote}

Beijing proposed instead the formation of study classes, made up of
members of both factions, taking as their theme the slogan ``struggle
against self and criticize revisionism'' (斗私批修). Criticism would be
directed against organizations which had committed mistakes and their
members urged to repent. As an inducement, such groups would obtain
quotas in revolutionary committees.

北京转而建议成立由双方派别成员组成、以``斗私批修''为主题的学习班。批评将指向犯过错误的组织，并敦促其成员悔改。作为诱因，这些团体将在革命委员会中获得名额。

However, the compromise suggested by Mao probably pleased neither side
in Zhejiang. United Headquarters was now forced to negotiate with its
hated rival, albeit from a position of moral superiority. Red Storm, on
the other hand, had secured endorsement, however qualified, from the
great leader himself. This alone was of enormous significance. Most
importantly, it had been accredited as a ``rebel'' organization. Mao's
blessing would very likely increase Red Storm's reluctance to discuss or
confess its alleged errors. Red Storm had no real alternative but to sit
down at the conference table, but it knew full well that its signature
was required on an agreement before the establishment of the ZPRC could
be announced. Therefore, it had a vested interest in holding out for
maximum concessions.

然而，毛提出的妥协大概使浙江双方都不满意。联合总部现在被迫同自己痛恨的对手谈判，尽管它仍处于道义优势地位。另一方面，红暴获得了伟大领袖本人虽有保留但毕竟存在的认可。这一点本身意义重大。最重要的是，它被承认为一个``造反''组织。毛的背书很可能增强红暴不愿讨论或承认所谓错误的态度。红暴没有真正选择，只能坐到谈判桌前，但它非常清楚，省革委会宣布成立前必须有它在协议上的签字。因此，它有利益尽力争取最大让步。

In response to this central notice, United Headquarters organized a
major rally in Hangzhou on 9 December.\footnote{ZJRB, December 10, 1967,
  pp.~1, 3.} Zhang Yongsheng, on behalf of United Headquarters, read out
Mao's directive and urged compliance. Representatives of Red Storm at
Zhejiang University attended the rally and promised to adhere to the
terms of the notice. Their spokesman said:

作为对中央通知的回应，联合总部于12月9日在杭州组织大型集会。\footnote{ZJRB,
  December 10, 1967, pp.~1, 3.}
张永生代表联合总部宣读毛的指示并敦促遵守。浙江大学红暴代表出席集会，并承诺遵守通知条款。他们的发言人说：

\begin{quote}
We used to fight shoulder to shoulder with the proletarian
revolutionaries of the Provincial Revolutionary Rebels' Joint General
Command. However, in the storms of the great ``January revolution'' and
at the crucial stage of the revolution, anarchism in our minds asserted
itself and was taken advantage of by the party capitalist-roaders. We
became separated from Chairman Mao's revolutionary line and committed
mistakes of orientation and line. This greatly damaged the Great
Proletarian Cultural Revolution in Zhejiang. It greatly pains us to
think of this and we are ashamed of ourselves.
\end{quote}

\begin{quote}
我们曾同省革命造反派联合总指挥部的无产阶级革命派并肩战斗。但是，在伟大的``一月革命''风暴中，在革命的关键阶段，我们头脑中的无政府主义抬头，并被党内走资本主义道路的当权派所利用。我们脱离了毛主席的革命路线，犯了方向和路线错误。这严重损害了浙江的无产阶级文化大革命。想到这一点，我们非常痛心，也感到惭愧。
\end{quote}

On the same day as the provincial press carried the report of the rally,
Weng Senhe and four fellow workers from the silk complex, including
Huang Yintang, published an article in which they welcomed Mao's
directive and pledged their obedience.\footnote{ZJRB, December 10, 1967,
  p.~3.} On 15 December Weng was also a featured speaker at a provincial
rally at which he and two colleagues from Red Storm publicly recanted
and in effect made their desertion official.\footnote{ZJRB, December 16,
  1967, p.~3; Facts \& Features, 1:13, p.~11.} In his speech Weng
characterized Mao's assessment of Red Storm as unfavorable. He stated
that Red Storm, despite its public confessions, had manipulated Mao's
words for its own ends by arguing, that everyone, including United
Headquarters, had made mistakes. Weng's denial of the allegation that he
and his colleagues had joined the revolution (and by implication United
Headquarters) to gain position, fame and power,
(``我们干革命的目的,并不是为争席位,争名利,争官作'') was to carry a very
hollow ring, as events later were to prove.

省内报刊刊登集会报道的同一天，翁森鹤和杭州丝绸联合厂另外四名工人，包括黄阴堂，发表文章欢迎毛的指示并保证服从。\footnote{ZJRB,
  December 10, 1967, p.~3.}
12月15日，翁又在一次省级集会上作重点发言，他和两名红暴同事公开认错，实际上正式宣布倒戈。\footnote{ZJRB,
  December 16, 1967, p.~3; Facts \& Features, 1:13, p.~11.}
翁在讲话中把毛对红暴的评价说成是不利的。他说，红暴尽管公开认错，却为自身目的歪曲毛的话，声称包括联合总部在内人人都犯了错误。翁否认他和同事参加革命（也就是暗指加入联合总部）是为了争地位、名利和官职（``我们干革命的目的，并不是为争席位，争名利，争官作''）；事后证明，这句话听起来十分空洞。

At the end of December Weng Senhe was given further publicity with the
publication of yet another article in Zhejiang Daily.\footnote{ZJRB,
  December 25, 1967, p.~2.} Weng was clearly useful as a frontman to
convince waverers within the ranks of Red Storm that alliance with
United Headquarters, even under unfavorable conditions, was the only
viable course open to them. Based largely on his televised speech of the
preceding week, Weng's article encapsulated the predicament facing Red
Storm, United Headquarters and the provincial authorities in trying to
come to an acceptable compromise on the basis of Mao's 2 December
directive. Were Red Storm members to emphasize the phrase which
described the organization as an ``old rebel faction'' (老造反派) or
were they to place stress on the phrase ``which has committed mistakes''
(犯过错误的)? Were Red Storm's errors minor or were they mistakes of
line, and how did they compare with those committed by United
Headquarters? Was emphasis in forthcoming negotiations to be placed on
the alliance (联合) or on ``help and criticism'' (帮助批评)? Related to
but more important than all these considerations was the quota (名额) to
be reserved for Red Storm on revolutionary committees.

12月底，翁森鹤又因《浙江日报》刊登另一篇文章而得到进一步宣传。\footnote{ZJRB,
  December 25, 1967, p.~2.}
翁显然可作为门面人物，用来劝说红暴队伍中犹豫者相信，即使条件不利，同联合总部结盟也是他们唯一可行的道路。翁的文章大体根据他前一周的电视讲话，概括了红暴、联合总部和省级当局在试图以毛12月2日指示为基础达成可接受妥协时面临的困境。红暴成员应强调把该组织称为``老造反派''的说法，还是强调``犯过错误的''这一短语？红暴的错误是轻微错误，还是路线错误？同联合总部所犯错误相比又如何？即将到来的谈判应重在``联合''，还是重在``帮助批评''？同所有这些考虑相关、且更重要的是，应为红暴在革命委员会中保留多少名额。

Weng argued that since February 1967, when it had wrecked the rally held
by United Headquarters to denounce Jiang Hua, Red Storm had committed
mistakes of an ultra-left nature. By its splittist, destabilizing
activities and by attacking United Headquarters, it had allowed the
rightists (the provincial capitalist-roaders) to cause more trouble.
Some of its members, complained Weng, even refused to heed the center's
2 December directive.

翁声称，自1967年2月破坏联合总部谴责江华的集会以来，红暴犯了极``左''性质的错误。它通过分裂主义和破坏稳定的活动，以及攻击联合总部，使右派（省内走资派）得以制造更多麻烦。翁抱怨说，其部分成员甚至拒绝听从中央12月2日指示。

In a conciliatory gesture Zhejiang Daily began to differentiate between
Red Storm and conservative or reactionary organizations.\footnote{ZJRB,
  December 8, 1967, p.~3.} However, it accused ``the old rebel faction
which has committed mistakes'' of ``petty-bourgeois fanaticism''. In a
talk he gave early in 1968 Zhou Enlai described Red Storm as ``a bit too
far to the `Left'\,''.\footnote{SCMP, 4139 (March 15, 1968), p.~4.}
Perhaps the description related to Red Storms' popularity among students
and teachers in Hangzhou, groups notorious for the fanaticism of their
behavior in the Cultural Revolution. The editorial also warned that
``rightists'' would seek to take advantage of the atmosphere of
reconciliation to cause trouble. Among them were a sizeable proportion
of Red Storm members who were unhappy at the prospect of kowtowing to
United Headquarters.

《浙江日报》以和解姿态开始把红暴同保守或反动组织区别开来。\footnote{ZJRB,
  December 8, 1967, p.~3.}
然而，它指控这个``犯过错误的老造反派''患有``小资产阶级狂热病''。周恩来在1968年初一次谈话中把红暴形容为``有点太`左'\,''。\footnote{SCMP,
  4139 (March 15, 1968), p.~4.}
这一说法或许与红暴在杭州学生和教师中受欢迎有关，这些群体在文化大革命中以行为狂热著称。社论还警告说，``右派''会利用和解气氛制造麻烦。其中相当一部分是红暴成员，他们不满向联合总部低头的前景。

From 29 November until 6 December, during which time Mao's comments
about Red Storm were made public, a conference jointly sponsored by
United Headquarters and the ZPRC preparatory committee took
place.\footnote{ZPS, December 1 and 6, 1967, SWB/FE/2647/B/9-10.} Chen
Liyun, Political Commissar of Unit 7350, delivered the major address. He
claimed optimistically that the situation in Zhejiang was excellent:
United Headquarters had acquired a high prestige and a considerable
number of cadres were resuming work. Personal grudges should not
determine which cadres were ``liberated'', he reminded his listeners.
Chen emphasized that political power was the key question and the
establishment of the ZPRC was ``imminent''. The shortcomings evident in
the provisional organ of power -- presumably the absence of Red Storm
representatives -- would ``gradually be corrected'', he promised. Zhang
Yongsheng delivered a report summing up the year's events in Zhejiang.

11月29日至12月6日，即毛关于红暴的评论公开期间，联合总部和省革委会筹备委员会联合主办了一次会议。\footnote{ZPS,
  December 1 and 6, 1967, SWB/FE/2647/B/9-10.}
7350部队政委陈励耘作主要报告。他乐观地声称，浙江形势很好：联合总部已有很高威望，相当数量干部正在恢复工作。他提醒听众，个人恩怨不应决定哪些干部得到``解放''。陈强调，政权问题是关键，省革委会的成立``迫在眉睫''。他承诺，临时权力机构中显而易见的缺点
-- 大概指没有红暴代表 --
将``逐步得到纠正''。张永生作报告，总结浙江一年来的事件。

Despite the sanctity surrounding Mao's directives and the undoubted
pressures from Beijing, the road towards unity was a minefield of
grudges and jealousies which threatened progress all the way. The desire
by both groups to form the nucleus of the alliance was obviously strong,
as an editorial of Zhejiang Daily observed:

尽管毛的指示具有神圣性，北京的压力也毫无疑问，通往团结的道路仍是一片怨恨和嫉妒的雷区，处处威胁进展。两个团体都想成为联盟核心的愿望显然很强，正如《浙江日报》一篇社论所观察到的：

\begin{quote}
Why can they not overcome their own mountain-stronghold mentality or
give up their own contingent (恋恋不舍他那支队伍)? Some of them even
dream of returning to their own stronghold, of commanding their own
contingent and of starting ``civil war'' again\ldots.\footnote{ZJRB,
  December 15, 1967, p.~4.}
\end{quote}

\begin{quote}
为什么不能克服自己的山头主义，不能放弃自己的队伍（恋恋不舍他那支队伍）？有些人甚至梦想回到自己的山头，指挥自己的队伍，再打``内战''\ldots\ldots{}\footnote{ZJRB,
  December 15, 1967, p.~4.}
\end{quote}

If political power was the issue, as Chen Liyun had stated, then both
groups hoped to get their share in the distribution of quotas to mass
representatives for seats on revolutionary committees.

如果正如陈励耘所说，政权是问题所在，那么两个团体都希望在分配群众代表进入革命委员会的名额时得到自己的份额。

Another sensitive issue concerned divisions within the military and its
relations with the two mass organizations. A Zhejiang Daily editorial of
mid December 1967 complained that ``some people'' ignored the opinions
of the military about the composition of the ZPRC.\footnote{ZJRB,
  December 18, 1967.} These people, who were not named, thought only of
their own interests and fame, and believed that they could flex their
muscles to intimidate others. Such behavior, warned Zhejiang Daily,
would lead to their downfall, no matter how great their capabilities or
past achievements. It is difficult to know at which faction this stern
warning was directed, but it may well have been at certain members of
Red Storm, who were refusing to concede to United Headquarters its claim
for greater representation on the ZPRC.

另一个敏感问题，是军队内部的分歧及其同两个群众组织的关系。1967年12月中旬《浙江日报》社论抱怨说，``有些人''在省革委会组成问题上无视军方意见。\footnote{ZJRB,
  December 18, 1967.}
这些未被点名的人只考虑自身利益和名声，并以为可以炫耀力量来恐吓他人。《浙江日报》警告说，这种行为将导致他们垮台，无论其能力或过去成就多么大。很难判断这一严厉警告指向哪个派别，但很可能是红暴中某些成员，他们拒绝承认联合总部在省革委会中要求更多代表名额的主张。

Whether the military in Zhejiang held a unified view towards the two
major mass organizations is debatable. Unit 7350 had shown a decided and
open partiality for United Headquarters,\footnote{See Chen Liyun and Bai
  Zongshan, RMRB, October 25, 1967.} a position not shared as blatantly
by the 20th Army. Yet in late October 1967, the Air Force unit held a
meeting at which harsh words were expressed about the behavior of the
rebels, presumably those from United Headquarters. They were rebuked for
indulging in sectarianism and stressing the ``core'' mentality. The Air
Force unit had held study classes to maintain close contacts with and
``train'' the rebel leaders.\footnote{ZJRB, October 29, 1967, p.~1.}
Thus, even Unit 7350 was conscious of creating the impression that it
exercised its authority as a central ``support the left'' unit without
favor.

浙江军方是否对两个主要群众组织持统一看法，尚可争论。7350部队公开而明确地偏向联合总部，\footnote{See
  Chen Liyun and Bai Zongshan, RMRB, October 25, 1967.}
第二十军则没有如此露骨地采取同一立场。然而，1967年10月底，空军部队召开会议，对造反派行为提出严厉批评，想必指的是联合总部造反派。他们因沉溺于宗派主义、强调``核心''思想而受到斥责。空军部队曾举办学习班，以保持同造反派领袖的密切联系并对其进行``训练''。\footnote{ZJRB,
  October 29, 1967, p.~1.}
因此，即使7350部队也意识到需要造成一种印象，即作为中央``支左''部队，它是在不偏不倚地行使权威。

Leaders of naval and army units attended the meeting to extend
congratulations to the Air Force for its contributions in the Cultural
Revolution. Xie Zhenghao (谢正浩), Commander of the Naval
units,\footnote{Naval unit 4291 stationed at Zhoushan island, had sent
  boats up the Yangtze river to Wuhan in October to help restore central
  authority. ZPS, November 9, 1967, SWB/FE/2619/11-12.} commended Unit
7350 for its firm support for the proletarian revolutionaries. Political
Commissar Nan Ping of Unit 6409 went even further in stating that the
Air Force had proved itself the most resolute in the work of ``three
supports and two militaries''. This display of inter-service harmony and
back-slapping was in reality a facade for the latent tensions between
the units which was reflected in other reports. One such report accused
the capitalist-roaders of spreading rumors designed to weaken unity
between different units of the military and between the PLA and the
citizenry. Another report, summarizing an article from Liberation Army
Daily in Beijing, concerned a small unit, probably of the 20th Army. It
had chosen the wrong side among two factions on a commune, due to the
influence of ``old thoughts and conventions''.\footnote{See ZJRB,
  November 24, 1967, p.~1; see also ZPS, November 28, 1967,
  SWB/FE/2637/B/8-9.}

海军和陆军部队领导人出席会议，祝贺空军在文化大革命中的贡献。海军部队司令员谢正浩\footnote{Naval
  unit 4291 stationed at Zhoushan island, had sent boats up the Yangtze
  river to Wuhan in October to help restore central authority. ZPS,
  November 9, 1967, SWB/FE/2619/11-12.}
称赞7350部队坚定支持无产阶级革命派。6409部队政委南萍更进一步，说空军在``三支两军''工作中表现得最坚决。这种军种之间的和谐和相互吹捧，实际上掩盖了部队之间潜在的紧张，这种紧张在其他报道中有所反映。一则报道指控走资派散布谣言，企图削弱军队各单位之间以及人民解放军同人民之间的团结。另一则报道概述北京《解放军报》一篇文章，涉及一个小单位，可能是第二十军。由于``旧思想、旧习惯''的影响，它在一个公社两个派别之间站错了队。\footnote{See
  ZJRB, November 24, 1967, p.~1; see also ZPS, November 28, 1967,
  SWB/FE/2637/B/8-9.}

One expert contends that political friction between Units 7350 and 6409
served as a proxy for the battle between the two mass
organizations.\footnote{Nelsen, The Chinese Military System, p.~137.}
The hostility between United Headquarters and Red Storm had shaken the
unity of the PLA units in Zhejiang. Mistrust between the military and
the mass organizations added to the tension. As the new year of 1968
commenced, however, greater efforts were made to overcome disunity and
agree upon the conditions for the establishment of the ZPRC. To attain
this objective, the central authorities summoned military leaders to
Beijing for discussions and study. They also held representatives of the
mass organizations under virtual confinement in the capital until they
could reach an agreement.\footnote{The center also foisted agreements
  upon rebel groups from Fujian (see Ling, The Revenge of Heaven,
  ch.~30) and Shanxi (Hinton, Shenfan, pp.~630 ff.).} On 16 February
1968, this approach paid dividends when, amidst great rejoicing, the
alliance was signed. The majority of the leaders of Red Storm had stuck
to their guns and not defected with Weng Senhe. They had signed the
agreement with United Headquarters as equal partners and expected to
share the spoils which had been promised to both organizations.

一位专家认为，7350部队和6409部队之间的政治摩擦，是两个群众组织斗争的替代形式。\footnote{Nelsen,
  The Chinese Military System, p.~137.}
联合总部与红暴之间的敌意动摇了浙江解放军部队的团结。军队和群众组织之间的不信任也增加了紧张。然而，随着1968年新年开始，各方加大努力克服分裂，并就成立省革委会的条件达成一致。为实现这一目标，中央当局召集军事领导人到北京讨论和学习。他们还把群众组织代表实际软禁在首都，直到他们达成协议。\footnote{The
  center also foisted agreements upon rebel groups from Fujian (see
  Ling, The Revenge of Heaven, ch.~30) and Shanxi (Hinton, Shenfan,
  pp.~630 ff.).}
1968年2月16日，这一办法取得成果，联盟协议在一片欢腾中签署。红暴多数领导人坚持原立场，并没有随翁森鹤倒戈。他们作为平等伙伴同联合总部签署协议，并期望分享承诺给两个组织的好处。

The theme of central and provincial propaganda at the beginning of 1968
focused on the evils of factionalism and highlighted the need for unity.
To demonstrate their neutrality and evenhandedness PLA units were
ordered to ``support the left, not any faction'' (一碗水端平) of the two
major groups.\footnote{On January 28, 1968, Liberation Army Daily
  published an editorial to this effect. See SCMP, 4114 (February 8,
  1968), pp.~14-16.} From 14 to 25 January, Unit 6409 held a conference
to discuss the implications of the new policy direction while units 6409
and 7350 published two articles on the theme of supporting the left and
not any faction.\footnote{ZPS, January 25, 1968, SWB/FE/2691/B/14;
  January 27, 1968, SWB/FE/2685/B/10; ZJRB, January 8, 1968, p.~2; ZJRB,
  January 10, 1968, p.~1.} Zhejiang Daily also republished from Red
Storm's newspaper an editorial opposing factionalism as well as an
article by United Headquarters expressing similar sentiments.\footnote{ZJRB,
  January 18, 1968, p.~3; ZJRB, January 23, 1968, p.~3.} As an example
to the rest of Zhejiang, Nan Ping and Chen Liyun co-authored an article
in People's Daily praising the deeds of a Li Wenzhong, a national model
in implementing the new course.\footnote{RMRB, January 22, 1968, p.~3.}
For its part Zhejiang Daily released a series of ten editorials between
4 January and 22 February, entitled ``Down with factionalism''
(打倒派性).\footnote{ZJRB, January 4, 1968, p.~3; January 16, 1968,
  pp.~1, 3; January 17, 1968, p.~1; January 20, 1968, p.~2; January 21,
  1968, pp.~1, 2; February 7, 1968, p.~4; February 8, 1968, p.~1;
  February 13, 1968, p.~1; February 15, 1968, p.~2; ZPS, February 22,
  1968, SWB/FE/2708/B/16-17. See also Facts \& Features, 1:9 (February
  21, 1968), p.~28 and 1:11 (March 20, 1968), pp.~9-11.}

1968年初，中央和省级宣传的主题集中在派性之害和团结必要性上。为显示中立和公正，人民解放军部队被命令在两个主要团体之间``支左不支派（一碗水端平）''。\footnote{On
  January 28, 1968, Liberation Army Daily published an editorial to this
  effect. See SCMP, 4114 (February 8, 1968), pp.~14-16.}
1月14日至25日，6409部队召开会议讨论新政策方向的含义；6409和7350部队还发表两篇以支左不支派为主题的文章。\footnote{ZPS,
  January 25, 1968, SWB/FE/2691/B/14; January 27, 1968,
  SWB/FE/2685/B/10; ZJRB, January 8, 1968, p.~2; ZJRB, January 10, 1968,
  p.~1.}
《浙江日报》也转载红暴报纸上一篇反对派性的社论，以及联合总部一篇表达类似观点的文章。\footnote{ZJRB,
  January 18, 1968, p.~3; ZJRB, January 23, 1968, p.~3.}
为给浙江其他地方树立榜样，南萍和陈励耘在《人民日报》联名发表文章，赞扬全国贯彻新路线的典型李文忠的事迹。\footnote{RMRB,
  January 22, 1968, p.~3.}
《浙江日报》则在1月4日至2月22日期间发表题为``打倒派性''的十篇系列社论。\footnote{ZJRB,
  January 4, 1968, p.~3; January 16, 1968, pp.~1, 3; January 17, 1968,
  p.~1; January 20, 1968, p.~2; January 21, 1968, pp.~1, 2; February 7,
  1968, p.~4; February 8, 1968, p.~1; February 13, 1968, p.~1; February
  15, 1968, p.~2; ZPS, February 22, 1968, SWB/FE/2708/B/16-17. See also
  Facts \& Features, 1:9 (February 21, 1968), p.~28 and 1:11 (March 20,
  1968), pp.~9-11.}

The general thrust of the editorials was that the question of power
unduly preoccupied the minds of the two factions. Given the direction
that the Cultural Revolution had taken since the beginning of 1967 the
authorities had little reason to be surprised at this obsession with
power. The rebels had risen against the old party authorities after Mao
had issued the call to arms. They had fought each other and the military
back and forth across the province for twelve months. Now, when all the
pieces were being put back together it was rather late in the day for
Zhejiang Daily to expect them to renounce all reward for their labors.
From their sudden and fierce initiation into the arena of elite
political struggle, they had very quickly learned that obtaining and
holding onto power outweighed all other considerations.

这些社论的大意是，权力问题过度占据两个派别的头脑。鉴于文化大革命自1967年初以来的发展方向，当局其实没有理由对这种权力执念感到惊讶。毛发出战斗号召后，造反派起来反对旧党政当局。他们在全省范围内同彼此、同军队反复作战十二个月。如今，在所有碎片重新拼合之时，《浙江日报》却指望他们放弃劳动的一切报偿，时机已相当晚。从突然且激烈地进入精英政治斗争舞台开始，他们很快就学会，取得并掌握权力压倒其他一切考虑。

With negotiations under way and positions on revolutionary committees
and the three congresses (worker, peasant and Red Guard) up for grabs,
compromise would lead to defeat. The editorials admitted that both sides
strove to maximize the number, position and seniority of their
representatives in the new power structure. As the basis for their
respective claims, both factions argued over which group had played the
major role in the Cultural Revolution. If this disputation continued,
warned one editorial rising to new heights of hyperbole, the disgraced
party authorities would instigate civil war and in the resultant
confusion infiltrate the ranks of the revolutionaries and regain power.

谈判正在进行，革命委员会和三大代表大会（工人、农民和红卫兵）的职位也待分配，妥协就意味着失败。社论承认，双方都努力使自己在新权力结构中代表人数、职位和资历最大化。作为各自主张的基础，两个派别争论哪个团体在文化大革命中发挥了主要作用。一篇社论以夸张到新高度的口吻警告说，如果这种争论持续下去，失势党政当局将煽动内战，并在随后的混乱中渗入革命派队伍，重新夺权。

According to one editorial, members of Red Storm had become ``so
arrogant and conceited that they even have refused to admit their
mistakes'' and ``put on airs of self-styled old rebels and genuine
revolutionaries''. Zhejiang Daily accurately predicted that if Red Storm
walked away from the talks dissatisfied, it would not afford recognition
to the ZPRC and would instead work for its overthrow. This was why both
groups took a ``pragmatic attitude'' towards central instructions. They
complied with those that suited them and boycotted those that did not.
The question of cadres was one case in point. It was obviously in the
interests of both Red Storm and United Headquarters that only those
cadres with whom they could cooperate on revolutionary committees be
allowed to reappear on the political stage. This important consideration
certainly contributed to the footdragging in implementing central
policy. The return of each cadre was thus a slow process and the result
of protracted negotiation.

据一篇社论称，红暴成员变得``狂妄自大，甚至不承认错误''，并``摆出自命老造反、真革命的架子''。《浙江日报》准确预言，如果红暴不满意地退出谈判，它将不承认省革委会，转而为推翻它而工作。这就是为什么两个团体都对中央指示采取``实用态度''：适合自己的就遵守，不适合自己的就抵制。干部问题就是一个例子。显然，对红暴和联合总部都有利的是，只有那些能同自己在革命委员会中合作的干部，才被允许重新出现在政治舞台上。这一重要考虑无疑造成了执行中央政策时的拖延。因此，每个干部的复出都是缓慢过程，也是长期谈判的结果。

Wu Xian was one leading cadre who seems to have been opposed by United
Headquarters but who had been permitted to join the new provisional
hierarchy. At a mass rally held on 19 January 1968, Wu delivered a
report on factionalism.\footnote{ZPS, January 19, 1968,
  SWB/FE/2679/B/11-12.} He stated that it was necessary to distinguish
between factional disputes and the struggle between the two lines. Wu's
words reflected a fear on the part of the authorities that a blanket
denunciation of factionalism could imply that the supposedly genuine and
principled differences between the revolutionary and revisionist lines
fell into the category of factional disputes. Whether Wu's listeners
grasped the distinction is another matter. The struggle for power had
become the major issue at stake and for the members of both mass
organizations its successful outcome would compensate for the sacrifices
and hardships of the previous twelve months.

吴宪是一名似乎遭联合总部反对、但获准加入新临时领导层的主要干部。1968年1月19日一次群众集会上，吴作关于派性的报告。\footnote{ZPS,
  January 19, 1968, SWB/FE/2679/B/11-12.}
他说，必须区分派性纠纷和两条路线斗争。吴的话反映出当局担心：对派性的一概谴责，可能意味着所谓革命路线与修正主义路线之间真实且有原则的差异也被归入派性纠纷。吴的听众是否把握住这一区别，则是另一回事。权力斗争已成为主要问题；对两个群众组织的成员来说，斗争成功的结果将补偿过去十二个月的牺牲和苦难。

On 7 February 1968, the Provincial Revolutionary Rebel Joint General
Committee convened a rally to attack factionalism.\footnote{ZPS,
  February 7, 1968, SWB/FE/2695/B/16.} This body comprised leading
members of both United Headquarter and Red Storm. Eleven days later, on
16 February, an editorial in Zhejiang Daily proclaimed the signing of
the agreement between the two mass organizations in Beijing.\footnote{ZPS,
  February 18, 1968, SWB/FE/2702/B/8-10.} A Hong Kong observer claimed
that the two organizations had signed a twelve-point agreement in
January 1968, but that ``the local branches paid no heed and went on
with their angry squabbles''.\footnote{CNA, 731 (November 1, 1968),
  p.~2.} Two documents were signed in February. One concerned the
alliance of United Headquarters and Red Storm and the other related to
the alliance at Zhejiang University, bastion of Red Storm.\footnote{See
  ZJRB, February 29, 1968.} The editorial admitted that only central
intervention had brought the groups to the conference table. Ruing the
fact that ``strife for supremacy has brought us misery'', Zhejiang Daily
urged each side to refrain from statements and actions detrimental to
what was clearly a fragile agreement. To press home the importance of
unity, the editorial also warned against the sabotage activities of
Guomindang agents, while at the same time declaring optimistically that
``from now on the road is smooth''.

1968年2月7日，省革命造反派联合总委员会召开批判派性的集会。\footnote{ZPS,
  February 7, 1968, SWB/FE/2695/B/16.}
该机构由联合总部和红暴的主要成员组成。十一天后的2月16日，《浙江日报》社论宣布两个群众组织在北京签署协议。\footnote{ZPS,
  February 18, 1968, SWB/FE/2702/B/8-10.}
一名香港观察者称，两个组织在1968年1月签署了一项十二条协议，但``地方分支毫不理会，继续愤怒争吵''。\footnote{CNA,
  731 (November 1, 1968), p.~2.}
2月签署了两份文件。一份涉及联合总部与红暴的联盟，另一份涉及红暴堡垒浙江大学的联盟。\footnote{See
  ZJRB, February 29, 1968.}
社论承认，只有中央介入才把各团体带到谈判桌前。《浙江日报》感叹``争强好胜给我们带来了苦难''，敦促双方避免作出损害这个显然脆弱协议的言行。为了强调团结重要性，社论还警告国民党特务的破坏活动，同时乐观宣称``从此道路是平坦的''。

That the road was anything but smooth quickly became apparent. Within
days of the announcement Zhejiang Daily acknowledged that ``Some people
have even engaged in factionalism after being united'' or were preparing
for ``civil war''.\footnote{ZJRB, February 22, 1968, p.~3; ZPS, February
  25, 1968, SWB/FE/2709/B/20.} A further editorial divulged the news
that a feature of factionalism which had emerged in the aftermath of the
signing of the agreement was that those who promoted unity were being
maligned as revisionists.\footnote{ZPS, February 29, 1968,
  SWB/FE/2712/B/5-6.} Debates were undoubtedly raging within both Red
Storm and United Headquarters about what they had gained and what they
had sacrificed by the agreement. The controversy forced the HMRC to
issue an urgent notice on 29 February 1968 demanding adherence to the
agreement from all sides.\footnote{Ibid, B/6-7.}

道路绝非平坦，这一点很快显现。公告发布几天内，《浙江日报》就承认``有些人团结起来以后还搞派性''，或正在准备``内战''。\footnote{ZJRB,
  February 22, 1968, p.~3; ZPS, February 25, 1968, SWB/FE/2709/B/20.}
另一篇社论透露，协议签署后出现的一种派性表现，是推动团结的人被诬蔑为修正主义者。\footnote{ZPS,
  February 29, 1968, SWB/FE/2712/B/5-6.}
红暴和联合总部内部无疑正在激烈争论，通过协议各自得到了什么、牺牲了什么。这场争议迫使杭州市革委会于1968年2月29日发布紧急通知，要求各方遵守协议。\footnote{Ibid,
  B/6-7.}

The two organizations delayed celebrating the alliance until 1
March.\footnote{ZJRB, March 2, 1968, p.~2.} Mutual suspicion was rife.
An editorial in Zhejiang Daily referred to the attempt by sides to try
and build up their numbers, demands for the agreement's revision or the
inclusion of supplementary articles and secret meetings and recruitment
drives in the countryside. They were informed bluntly that they must
accept the document ``even if they did not understand the reason
for\ldots{} restrictions'' contained in it.\footnote{ZJRB, March 11,
  1968, p.~1.} Representatives of both organizations in the counties
were brought to Hangzhou to attend study classes.\footnote{See Xiangshan
  xianzhi, p.~35. The two groups in Xiangshan county did not sign an
  agreement until April 13, 1968.} Dissatisfied members of Red Storm
branded upholders of the agreement from their side as ``Right deviates
who want to become officials''. Those members of United Headquarters who
preferred confrontation to consensus argued sarcastically that ``The
worst thing that could happen to us in this struggle is to have
ourselves end up as veteran rebels who have made mistakes''.\footnote{ZJRB,
  March 13, 1968, p.~1; CNS, 214 (April 4, 1968), pp.~1-3.} Thus, they
openly voiced their dissatisfaction at Mao's partial rehabilitation of
Red Storm.

两个组织直到3月1日才庆祝联盟。\footnote{ZJRB, March 2, 1968, p.~2.}
相互猜疑十分普遍。《浙江日报》社论提到，双方都试图扩大人数，要求修改协议或加入补充条款，并在农村秘密开会和招募成员。社论直截了当地告诉他们，即使``不理解其中限制的理由''，也必须接受文件。\footnote{ZJRB,
  March 11, 1968, p.~1.}
两个组织在各县的代表被带到杭州参加学习班。\footnote{See Xiangshan
  xianzhi, p.~35. The two groups in Xiangshan county did not sign an
  agreement until April 13, 1968.}
红暴不满成员把本派中拥护协议的人斥为``想当官的右倾分子''。联合总部中宁愿对抗、不愿协商的成员则讽刺地说：``这场斗争中我们最坏的结果，就是落得个犯过错误的老造反派。''\footnote{ZJRB,
  March 13, 1968, p.~1; CNS, 214 (April 4, 1968), pp.~1-3.}
因此，他们公开表达了对毛部分恢复红暴名誉的不满。

Despite the possibility that the agreement between United Headquarters
and Red Storm could break down at any moment, the central authorities
pushed on with their plans to set up the ZPRC, obviously believing that
any agreement was better than none. In a speech on 2 February 1968, Zhou
Enlai stated that Beijing hoped to establish the ZPRC in February
1968.\footnote{SCMP, 4154 (April 8, 1968), p.~2.} Even by mid-March,
only one municipality and sixteen counties in the province had formed
revolutionary committees.\footnote{See CNS, 270 (May 15, 1969), p.~3.
  These were, as far as is known, Hangzhou Municipality (December 21,
  1967), Jiande (December 23, 1967), Lanxi (December 30, 1967), Haining,
  Linan, Dongyang, Wuyi, Fenghua (all in December 1967), and Shaoxing
  (March 12, 1968).} Finally, on the night of 18 March, a delegation
from Zhejiang consisting overwhelmingly of military officials met
central leaders to report on the successful outcome of the protracted,
supervised negotiations which had been held in Beijing.\footnote{SCMP,
  4182 (May 21, 1968), pp.~1-12.}

尽管联合总部与红暴之间的协议随时可能破裂，中央当局仍推进成立省革委会的计划，显然认为有协议总比没有好。周恩来在1968年2月2日讲话中说，北京希望在1968年2月成立浙江省革委会。\footnote{SCMP,
  4154 (April 8, 1968), p.~2.}
即使到3月中旬，全省也只有一个市和十六个县成立革命委员会。\footnote{See
  CNS, 270 (May 15, 1969), p.~3. These were, as far as is known,
  Hangzhou Municipality (December 21, 1967), Jiande (December 23, 1967),
  Lanxi (December 30, 1967), Haining, Linan, Dongyang, Wuyi, Fenghua
  (all in December 1967), and Shaoxing (March 12, 1968).}
最后，3月18日晚，一个主要由军事官员组成的浙江代表团会见中央领导人，汇报在北京长期受监督谈判取得的成功结果。\footnote{SCMP,
  4182 (May 21, 1968), pp.~1-12.}

The head of the Cultural Revolution Group, Chen Boda, chaired the
meeting and made the concluding speech. Other speakers included Zhou
Enlai, Jiang Qing, Kang Sheng and Xu Shiyou, Commander of the Nanjing
Military Region and a staunch supporter of Red Storm. The leaders of Red
Storm and United Headquarters, Fang Jianwen and Zhang Yongsheng, made
verbal pledges on behalf of their respective organizations. Zhou Enlai
claimed that the impetus behind the negotiations had come from the
center and from Mao in particular. Of all the central leaders Zhou alone
appeared familiar with the details of the situation in Zhejiang.

中央文革小组组长陈伯达主持会议并作总结讲话。其他发言人包括周恩来、江青、康生以及南京军区司令员、红暴坚定支持者许世友。红暴和联合总部领导人方剑文、张永生分别代表各自组织作口头保证。周恩来声称，谈判动力来自中央，尤其来自毛。在所有中央领导人中，似乎只有周熟悉浙江局势细节。

The Premier was critical of the PLA's tendency to interfere in matters
that were best left to others to settle. He mentioned the past mistakes
of the ZPMD, and of individual members of the military control
commission, and pointedly referred to the study classes which had been
held to educate the military on the desirability of adopting a neutral
stand toward both factions.\footnote{Mao had received the members of the
  study class of Unit 7350 in January 1968. ZPS, January 28, 1968,
  SWB/FE/2683/B/19. In February he had received provincial PLA activists
  from ``support the left'' units. ZPS, February 20, 1968,
  SWB/FE/2704/B/6.} In his brief talk Xu Shiyou also raised the issue of
PLA neutrality toward the mass organizations. Zhou declared that
finalization of the membership of the ZPRC entailed the dissolution of
the two mass organizations.\footnote{On this point see G.A. Bennett and
  R.N. Montaperto, Red Guard: The Political Biography of Dai Hsiao-ai
  (Garden City, New York: Doubleday, 1972), pp.~199, 223. Dai argued
  that an important reason why mass organizations held out so strongly
  against agreement was the fear that once they had signed, their group
  would be disbanded.} He stated that the only two factions which
remained were the proletarian revolutionaries and the bourgeois or
petty-bourgeois factionalists. He ordered the closure of liaison
stations maintained by Red Storm and United Headquarters in the counties
of Zhejiang and in major urban centers outside the province. His
unambiguous message was that local groups would solve their problems
without outside interference.

总理批评人民解放军倾向于干预本应由别人解决的事情。他提到浙江省军区过去的错误，以及军管会个别成员的错误，并明确提到为教育军方对两派采取中立立场而举办的学习班。\footnote{Mao
  had received the members of the study class of Unit 7350 in January
  1968. ZPS, January 28, 1968, SWB/FE/2683/B/19. In February he had
  received provincial PLA activists from ``support the left'' units.
  ZPS, February 20, 1968, SWB/FE/2704/B/6.}
许世友在简短讲话中也提出人民解放军对群众组织保持中立的问题。周宣布，省革委会成员名单确定意味着两个群众组织解散。\footnote{On
  this point see G.A. Bennett and R.N. Montaperto, Red Guard: The
  Political Biography of Dai Hsiao-ai (Garden City, New York: Doubleday,
  1972), pp.~199, 223. Dai argued that an important reason why mass
  organizations held out so strongly against agreement was the fear that
  once they had signed, their group would be disbanded.}
他说，剩下的只有无产阶级革命派和资产阶级或小资产阶级派性主义者两派。他命令关闭红暴和联合总部在浙江各县以及省外主要城市设立的联络站。他的明确信息是，地方团体将在没有外来干涉的情况下解决自己的问题。

In line with Mao's pronouncement that the Cultural Revolution
represented a continuation of the class struggle between the CCP and
Chiang Kaishek's Guomindang, Zhou pointed out the strategic importance
of Zhejiang to the success of the Revolution. Chen Boda charged the
``capitalist-roaders'' of Zhejiang with being Guomindang
agents.\footnote{SCMP, 4182, p.~11.} Kang Sheng, in the conspiratorial
style he had perfected, labored this fact, urging the rebels to be on
the lookout for Guomindang agents.\footnote{For further excerpts from
  Kang's speech, see Zhong Kan, 康生评传 (A critical biography of Kang
  Sheng), (Beijing: Hongqi Chubanshe, 1982), p.~417.} Provincial
propaganda later elaborated on Kang's warnings ad nauseam.\footnote{The
  Zhejiang Province Directive on Security Work (浙江省治保工作指示)
  issued by the provincial Public Security Bureau on April 26, 1968,
  highlighted this aspect of public security in Zhejiang. See
  红卫兵资料续编(一) (Red Guard Publications, Supplement 1),
  (Washington, D.C.: Centre for Chinese Research Materials, 1980),
  pp.~1509-1510, transl., in Facts \& Features, 1:21 (August 7, 1968),
  pp.~28-29.}

按照毛关于文化大革命是中共同蒋介石国民党之间阶级斗争继续的论断，周指出浙江对革命成功具有战略重要性。陈伯达指控浙江``走资派''是国民党特务。\footnote{SCMP,
  4182, p.~11.}
康生以其熟练掌握的阴谋论风格大讲这一点，敦促造反派警惕国民党特务。\footnote{For
  further excerpts from Kang's speech, see Zhong Kan, 康生评传 (A
  critical biography of Kang Sheng), (Beijing: Hongqi Chubanshe, 1982),
  p.~417.} 省内宣传后来对康生的警告反复发挥，令人厌烦。\footnote{The
  Zhejiang Province Directive on Security Work (浙江省治保工作指示)
  issued by the provincial Public Security Bureau on April 26, 1968,
  highlighted this aspect of public security in Zhejiang. See
  红卫兵资料续编(一) (Red Guard Publications, Supplement 1),
  (Washington, D.C.: Centre for Chinese Research Materials, 1980),
  pp.~1509-1510, transl., in Facts \& Features, 1:21 (August 7, 1968),
  pp.~28-29.}

However, these issues were of secondary importance compared to the major
problem which had exercised the minds of those present for months - the
composition of the ZPRC. Zhou Enlai announced that it would comprise 94
members, 50 of whom would come from Hangzhou. Red Storm would nominate
ten representatives to the committee and would have three members on its
standing committee. Zhou stated that one Vice-chairman would come from
the ranks of the workers. He added that the final make-up of the
committee had not been decided upon and joint investigation would enable
cadres to be added in a ``principled way'' to achieve ``gradual
perfection''.

然而，同困扰与会者数月的主要问题相比，这些议题都居于次要地位；主要问题就是省革委会的组成。周恩来宣布，省革委会将由94名成员组成，其中50人来自杭州。红暴将提名十名委员会代表，并在常委会中拥有三名成员。周说，一名副主任将来自工人队伍。他补充说，委员会最终组成尚未决定，联合调查将使干部能够以``有原则的方式''加入，从而``逐步完善''。

On 19 March the Zhejiang delegation was welcomed back to the province
with a rally in Hangzhou,\footnote{ZJRB, March 20, 1968, p.~2.} but in
the six days that transpired between the Beijing meeting and the 24
March rally celebrating the formation of the ZPRC, the agreement to
include representatives of Red Storm on the revolutionary committee
collapsed. Unfortunately, no list of the full membership of the ZPRC is
available. The composition of the fifteen-member standing committee
included eight military representatives, five cadres and only two mass
representatives, both from United Headquarters. The ZPRC has been
classified it as a revolutionary committee ``practically created by
\ldots{} other military units'' and Zhejiang as a province where the
revolutionary committee was formed only ``after new commanders \ldots{}
had been appointed who enjoyed the support of the Central Authorities
\ldots{}''.\footnote{Domes, ``The Role of the Military'', pp.~131-32,
  142, 143. See also RMRB, March 28, 1968; ZJRB, March 28, 1968; ZPS,
  March 24, 1968, SWB/FE/2733/B/2-4.} Red Storm was represented neither
on the standing committee nor, most probably, on the ZPRC itself. It is
little wonder then that it felt betrayed by the failure to honor the
agreement signed before central leaders. The organization henceforth
embarked on a campaign of opposition to show that Zhejiang could not
regain any semblance of stability without its participation in the new
political structure.

3月19日，浙江代表团返回省内，在杭州受到集会欢迎；\footnote{ZJRB, March
  20, 1968, p.~2.}
但在北京会议与3月24日庆祝省革委会成立的集会之间的六天里，把红暴代表纳入革命委员会的协议破裂了。遗憾的是，省革委会完整成员名单已不可得。十五人常委会的构成包括八名军方代表、五名干部，只有两名群众代表，且都来自联合总部。省革委会被归类为``实际上由\ldots\ldots 其他军事单位''建立的革命委员会，浙江也被归为只有在``受到中央当局支持的新指挥员\ldots\ldots 被任命之后''才成立革命委员会的省份。\footnote{Domes,
  ``The Role of the Military'', pp.~131-32, 142, 143. See also RMRB,
  March 28, 1968; ZJRB, March 28, 1968; ZPS, March 24, 1968,
  SWB/FE/2733/B/2-4.}
红暴既未在常委会中获得代表，也很可能未进入省革委会本身。难怪它觉得，在中央领导人面前签署的协议未被履行，自己遭到背叛。此后，该组织发动反对运动，以显示没有其参与新的政治结构，浙江就无法恢复任何稳定局面。

United Headquarters fared little better. Zhang Yongsheng was elected a
Vice-chairman of the ZPRC although he did not appear at the inaugural
rally. Perhaps he boycotted it from dissatisfaction with the limited
number of places reserved for his organization. Zhang's first appearance
as Vice-chairman of the ZPRC seems to have occurred on 17 April
1968.\footnote{See ZPS, April 18, 1968, SWB/FE/2755/B/6.} The other mass
representative was Hua Yinfeng, a child-bride and national labor model
from a state pig farm in Jinhua District. She was also elected a
Vice-chairman of the ZPRC. Hua was probably an example of a rebel who
has been characterized as more conservative, more attuned to working
with established authority, and better able to adjust to the process of
rebuilding the political system, than other rebels.\footnote{Teiwes,
  Provincial Leadership in China, p.~41.} TABLE FOUR

联合总部的境况也好不了多少。张永生当选省革委会副主任，尽管他没有出席成立大会。也许他因本组织保留席位太少而抵制大会。张首次以省革委会副主任身份出现，似乎是在1968年4月17日。\footnote{See
  ZPS, April 18, 1968, SWB/FE/2755/B/6.}
另一名群众代表是华银凤，她是金华专区一个国营养猪场的童养媳和全国劳动模范。她也当选省革委会副主任。华很可能体现了一类被描述为较保守、更适应同既有权威合作、也比其他造反派更能适应政治制度重建过程的造反派。\footnote{Teiwes,
  Provincial Leadership in China, p.~41.} 表四

The Standing Committee of the Zhejiang Provincial Revolutionary
Committee, 24 March 1968

浙江省革命委员会常务委员会，1968年3月24日

{\def\LTcaptype{none} % do not increment counter
\begin{longtable}[]{@{}ll@{}}
\toprule\noalign{}
Role & Name \\
\midrule\noalign{}
\endhead
\bottomrule\noalign{}
\endlastfoot
Chairman & Nan Ping \\
First Vice-chairman & Chen Liyun \\
Vice-chairmen & Xiong Yingtang \\
& Lai Keke \\
& Zhou Jianren \\
& Wang Zida \\
& Hua Yinfeng \\
& Zhang Yongsheng \\
Members & Zhu Quanlin \\
& Meng Zhaoyu \\
& Dai Kelin \\
& Shen Ce \\
& Wu Xian? \\
& Wang Qi? \\
& Mo Xianyao? \\
& Zhang Laigen \\
\end{longtable}
}

{\def\LTcaptype{none} % do not increment counter
\begin{longtable}[]{@{}ll@{}}
\toprule\noalign{}
职务 & 姓名 \\
\midrule\noalign{}
\endhead
\bottomrule\noalign{}
\endlastfoot
主任 & 南萍 \\
第一副主任 & 陈励耘 \\
副主任 & 熊应堂 \\
& 赖可可 \\
& 周建人 \\
& 王子达 \\
& 华银凤 \\
& 张永生 \\
委员 & 朱全林 \\
& 孟昭煜 \\
& 戴克林 \\
& 沈策 \\
& 吴宪？ \\
& 王琦？ \\
& 莫显尧？ \\
& 张来根 \\
\end{longtable}
}

Sources: RMRB, 28 March 1968; Kao Ch'ung-yen, Zhonggong renshi biandong,
pp.~670-1; ZPS, 24 March 1968, SWB/FE/2733/B/2-4; U. S. Government,
Directory of Chinese Communist Officials, A 70-13 (May 1970); ZJRB, 28
March 1968; 31 March 1968, p.~1.

资料来源：RMRB，1968年3月28日；Kao Ch'ung-yen，Zhonggong renshi
biandong，第670-671页；ZPS，1968年3月24日，SWB/FE/2733/B/2-4；美国政府，Directory
of Chinese Communist Officials，A
70-13（1970年5月）；ZJRB，1968年3月28日；1968年3月31日，第1页。

\section{Further Instability and Violence, March - August
1968}\label{further-instability-and-violence-march---august-1968}

\section{进一步的不稳定与暴力，1968年3月至8月}\label{ux8fdbux4e00ux6b65ux7684ux4e0dux7a33ux5b9aux4e0eux66b4ux529b1968ux5e743ux6708ux81f38ux6708}

The dismissal of PLA acting Chief of Staff Yang Chengwu, announced by
Lin Biao in March 1968, ushered in a renewed period of leftist militancy
across China.\footnote{Hao Mengbi and Duan Haoran, Liushinian, p.~599.}
A Zhejiang Daily editorial published on the day of the formation of ZPRC
set the tone for the more radical phase of the following months. While a
succession of editorials and speeches criticized factionalism, these
differed markedly in nature and intent from those published in January
and February 1968. This previous period was now redefined as the 1968
``Spring Current'', and was described as having its roots in the 1967
``February Countercurrent''. It was considered that the previous
emphasis on unity had been excessive and that a radical upsurge was
required to maintain the momentum of the Cultural Revolution. Written
and verbal broadsides were directed against attempts by the former
power-holders to obtain admission to the newly-established revolutionary
committees and Red Storm's denigration of and opposition to the ZPRC.
Such a media onslaught further inflamed factional tensions. Any hope
held by the authorities that Red Storm would meekly accept its omission
from the ZPRC and other representative bodies was soon dashed. The two
factions fought it out, wrote the Zhejiang Daily, while the
capitalist-roaders looked on with satisfaction. The class enemy had
begun infiltrating the revolutionary committees by using revolutionary
slogans or instead maligned them as ``committees of
factions''.\footnote{ZPS, March 24, 1968, SWB/FE/2730/B/15-17; ZPS,
  March 26, 1968, SWB/FE/2733/B/4-7; CNS, 270, pp.~2-5. A Red Guard
  document from Ninghai county, dated March 1968, accused the former
  director of the county Bureau of Culture and Education of ordering the
  murder of two Red Guards in July 1967. He had later established a
  ``bogus'' revolutionary committee. The document paints a picture of
  rebel impotence when confronted by a local official who could mobilize
  peasant support and defy central directives at will. See SCMP, 4167
  (April 30, 1968), pp.~17-19.}

林彪于1968年3月宣布撤销人民解放军代总参谋长杨成武职务，开启了全国范围内左派激进主义重新高涨的时期。\footnote{Hao
  Mengbi and Duan Haoran, Liushinian, p.~599.}
《浙江日报》在省革委会成立当天发表的社论，为随后数月更激进的阶段定下基调。虽然一连串社论和讲话仍批评派性，但其性质和意图同1968年1月、2月发表者明显不同。此前时期如今被重新定义为1968年的``春潮''，并被描述为根源在1967年的``二月逆流''。当时认为，过去对团结的强调过度了，需要一次激进高潮来保持文化大革命动力。文字和口头攻击指向前当权者试图进入新成立的革命委员会，也指向红暴对省革委会的诋毁和反对。这样的媒体攻势进一步激化了派性紧张。若当局曾希望红暴温顺接受被排除在省革委会和其他代表机构之外，这种希望很快破灭。《浙江日报》写道，两个派别相互争斗，走资派则满意地旁观。阶级敌人已经利用革命口号渗入革命委员会，或者把它们诬蔑为``派性委员会''。\footnote{ZPS,
  March 24, 1968, SWB/FE/2730/B/15-17; ZPS, March 26, 1968,
  SWB/FE/2733/B/4-7; CNS, 270, pp.~2-5. A Red Guard document from
  Ninghai county, dated March 1968, accused the former director of the
  county Bureau of Culture and Education of ordering the murder of two
  Red Guards in July 1967. He had later established a ``bogus''
  revolutionary committee. The document paints a picture of rebel
  impotence when confronted by a local official who could mobilize
  peasant support and defy central directives at will. See SCMP, 4167
  (April 30, 1968), pp.~17-19.}

To divert the rebels from their endless squabbles the 1st plenum of the
ZPRC, which met from 25 to 31 March,\footnote{ZPS, April 2, 1968,
  SWB/FE/2741/B/14-16.} mapped out a strategy to destroy Tan Zhenlin's
reputation in Zhejiang. The strategy had been devised in
Beijing\footnote{CNS, 270, p.~4.} and its objective was to condemn Tan
for his alleged responsibility in plotting the February 1967
Countercurrent. Because of Tan's close connections with Zhejiang and the
old CCP ZPC, this approach had the added attraction of making Jiang Hua
a prime target. Rallies which were held in April 1968 put the plan into
action.\footnote{ZJRB, April 5, 1968, p.~1; ZPS, April 15, 1968, URS,
  51:11 (May 7, 1968), pp.~131-2; ZJRB, April 16, 1968, pp.~1, 2; ZPS,
  April 18, 19, 1968, SWB/FE/2755/B/6; CNS, 270, pp.~6-7.} On 4 April,
the day of the first rally, an editorial published in Zhejiang Daily
entitled ``Resolutely overthrow the number one Party person in authority
in Zhejiang taking the capitalist road'' blamed Jiang and his associates
for the state of disunity among the rebels. It pointed out that

为了把造反派从无休止争吵中转移出来，3月25日至31日召开的省革委会第一届全会\footnote{ZPS,
  April 2, 1968, SWB/FE/2741/B/14-16.}
制定了摧毁谭震林在浙江声誉的策略。该策略由北京制定，\footnote{CNS, 270,
  p.~4.}
其目标是谴责谭对策划1967年2月``逆流''负有所谓责任。由于谭同浙江及旧省委关系密切，这一做法还有一个额外好处，即把江华变成主要目标。1968年4月举行的集会把计划付诸实施。\footnote{ZJRB,
  April 5, 1968, p.~1; ZPS, April 15, 1968, URS, 51:11 (May 7, 1968),
  pp.~131-2; ZJRB, April 16, 1968, pp.~1, 2; ZPS, April 18, 19, 1968,
  SWB/FE/2755/B/6; CNS, 270, pp.~6-7.}
4月4日，第一次集会当天，《浙江日报》发表题为``坚决打倒浙江党内头号走资本主义道路当权派''的社论，把造反派不团结状态归咎于江及其同伙。社论指出：

\begin{quote}
In attacking the handful of capitalist-roaders in Zhejiang the
revolutionary masses should immediately stop the struggle among
themselves and unite as one to chop off the sinister hands and smash the
sinister line.
\end{quote}

\begin{quote}
革命群众在打击浙江一小撮走资派时，应立即停止相互斗争，团结一致，斩断黑手，粉碎黑线。
\end{quote}

A principal target at the rallies was one of Jiang's closest former
subordinates, Xue Ju. Xue was blamed for spreading the rumor that Jiang
Hua was a member of Chairman Mao's Headquarters and disseminating
articles in praise of his boss and other former leaders of the province.
As argued in chapter one of this study, Xue had sound reasons for
describing Jiang in this way, particularly with the official silence
from Beijing concerning Jiang's status. On 15 April Red Storm was one of
the convenors of a large rally, an indication that efforts to keep it in
the political mainstream continued. A former Secretary of the ZPC
secretariat, Lai Keke, in his summation speech at the rally, demanded
the exposure of all capitalist-roaders in Zhejiang.

集会上的主要目标之一，是江最亲近的旧下属之一薛驹。薛被指责散布江华是毛主席司令部成员的谣言，并传播赞扬其上司和其他省内旧领导人的文章。正如本研究第一章所论，薛这样描述江有充分理由，尤其是在北京对江的地位保持官方沉默的情况下。4月15日，红暴是一次大型集会的召集者之一，表明把它留在政治主流内的努力仍在继续。前省委书记处书记赖可可在集会总结讲话中要求揭发浙江所有走资派。

In April Zhejiang Daily published five editorials devoted to
discrediting the old provincial administration and propagating the
revised definition of factionalism.\footnote{ZJRB, April 2, 1968, p.~4;
  ZJRB, April 15, 1968, p.~1; ZJRB, April 16, 1968, p.~1; ZJRB, April
  23, 1968, p.~2; ZJRB, April 28, 1968, p.~1.} The general thrust of the
editorials was synthesized in a People's Daily and Liberation Army Daily
joint editorial defending proletarian factionalism. A Red Flag
commentary entitled ``Factionalism must be subjected to class
analysis'', most probably written by Chen Boda, was the most
authoritative statement of the position.\footnote{PR, No.~19 (May 10,
  1968), pp.~3-4; the Red Flag commentary appeared in RMRB on April 27,
  1968. See SCMM, 635 (December 2, 1968), pp.~35-7. See also the joint
  editorial of People's Daily, Liberation Army Daily, and Red Flag of
  May 1, 1968 in SCMP, 4170 (May 3, 1968), pp.~1-3.} Rather than
subjecting factionalism as such to criticism, these polemics limited
their attacks to ``bourgeois factionalism''. The disagreements between
mass organizations of early 1968, which at the time had been described
as factional squabbles, were retrospectively reinterpreted as principled
political struggles.

4月，《浙江日报》发表五篇社论，专门败坏旧省级行政机构名誉，并宣传修订后的派性定义。\footnote{ZJRB,
  April 2, 1968, p.~4; ZJRB, April 15, 1968, p.~1; ZJRB, April 16, 1968,
  p.~1; ZJRB, April 23, 1968, p.~2; ZJRB, April 28, 1968, p.~1.}
这些社论的大意，在《人民日报》和《解放军报》一篇为无产阶级派性辩护的联合社论中得到综合。题为《派性必须作阶级分析》的《红旗》评论，很可能由陈伯达撰写，是这一立场最权威的表述。\footnote{PR,
  No.~19 (May 10, 1968), pp.~3-4; the Red Flag commentary appeared in
  RMRB on April 27, 1968. See SCMM, 635 (December 2, 1968), pp.~35-7.
  See also the joint editorial of People's Daily, Liberation Army Daily,
  and Red Flag of May 1, 1968 in SCMP, 4170 (May 3, 1968), pp.~1-3.}
这些论战文章并不是批判派性本身，而是把攻击限于``资产阶级派性''。1968年初群众组织之间的分歧，当时被描述为派性争吵，如今被追溯性地重新解释为有原则的政治斗争。

One observer has described Chen Boda's commentary as an effort by the
Maoists to define a new class theory.\footnote{Kraus, Class Conflict in
  Chinese Socialism, p.~114.} Instead of the old definition, which
tended to view class in terms of social origin, this incipient approach
based class on behavior and attitude. The origins of Chen's thesis may
be traced to the previous year. On 9 August 1967, Lin Biao had issued an
instruction on the Cultural Revolution which stated in part that

一位观察者把陈伯达的评论描述为毛派试图定义一种新的阶级理论。\footnote{Kraus,
  Class Conflict in Chinese Socialism, p.~114.}
不同于旧定义倾向于从社会出身看阶级，这一萌芽中的方法把阶级建立在行为和态度之上。陈的论点源头可追溯到前一年。1967年8月9日，林彪就文化大革命发出指示，其中部分写道：

\begin{quote}
As a basis for delineating Leftists and Rightists you should take
whether they support or oppose this Great Cultural Revolution launched
by Chairman Mao himself, and whether they defend or oppose Chairman Mao.
You should stand firmly beside Chairman Mao, stand beside the Leftists,
and stand beside the masses. You should not consider merely whether
class origins are pure of {[}sic{]} whether Party members or cadres are
numerous when you divide people into Leftists and Rightists or when you
look at problems. Class origins should be looked at, of course, but you
should not only look at class origins; the most important thing is to
realize on what party line these people stand.\footnote{Michael Y. M.
  Kau (ed.), The Lin Piao Affair: Party Politics and Military Coup
  (White Plains, New York: International Arts and Sciences Press Inc.,
  1975), p.~436.}
\end{quote}

\begin{quote}
划分左派和右派，要以是否拥护毛主席亲自发动的这场伟大文化大革命、是否保卫毛主席为依据。你们要坚定地站在毛主席一边，站在左派一边，站在群众一边。在划分左派和右派、看问题时，不能只看阶级出身是否纯，或者党员、干部是否多。阶级出身当然要看，但不能只看阶级出身；最重要的是看这些人站在什么党路线上。\footnote{Michael
  Y. M. Kau (ed.), The Lin Piao Affair: Party Politics and Military Coup
  (White Plains, New York: International Arts and Sciences Press Inc.,
  1975), p.~436.}
\end{quote}

By the logic of this argument, those who supported Mao automatically
became proletarian revolutionaries; those who opposed him were bourgeois
reactionaries. The approach circumvented accusations of class impurity
which had been levelled against certain mass organizations by party
officials under attack.\footnote{CNS, 220 (May 16, 1968), pp.~1-5.} In
his article, Chen Boda quoted Mao as having stated, ``there are parties
outside the Party and factions within it; this has always been the
case'' (党外有党,党内有派,历来如此).\footnote{Chen was apparently
  quoting, somewhat inexactly, from Mao's speech to the 11th plenum of
  the 8th CC on 12 August 1966 in which the chairman had said:
  ``我们这个党不是党外无党;我们是党外有党,党内也有派,从来都是如此,这是正常现象.''
  (It's not true that there are no parties outside our party; we do have
  them and also factions within the party. It's always been like this,
  and this is a normal phenomenon.) Mao Zedong sixiang wansui (1969),
  p.~652.} These words provided authority for Chen's thesis that the
existence of factions was inevitable in class society and that classes
and factions were somehow intertwined, factional struggle being a
manifestation of class struggle. While Chen's view of the universal
nature of factionalism within the CCP may have reflected an accurate and
frank assessment of reality,\footnote{This position has been put most
  forcefully by Pye in, The Dynamics of Chinese Politics.} it was an
opinion decidedly at odds with the party's public pronouncements on the
issue. Mao himself eventually drew back from the implications of this
analysis. For example, in April 1968, he characterized the Cultural
Revolution as a ``continuation of the class struggle between the
proletariat and the bourgeoisie''.\footnote{E. L. Wheelwright and Bruce
  McFarlane, The Chinese Road to Socialism, p.~244.} The lack of
theoretical clarity underlying the vague notion of ``continuous
revolution'', with its undeveloped and inconsistent notions of classes
and class struggle under socialism, occasioned and implicitly justified
such tactical shifts.\footnote{See White, The Politics of Class and
  Class Origin, p.~61; Kraus, Class Conflict in Chinese Socialism,
  pp.~110-14. For a thorough discussion of this issue, see Graham Young,
  ``Mao Zedong and the Class Struggle in Socialist Society'', Australian
  Journal of Chinese Affairs, 16 (July 1986), pp.~41-80.} The blurring
of the dividing line between factional and class enemies would lead to
harsh measures against rebels when the campaign to ``purify the class
ranks'' (清理阶级队伍) was initiated in 1968.\footnote{For a brief
  description of the origins of this campaign, see
  中共辽宁省委《共产党员》杂志内部发行版编辑组 (Internal Publication
  Editorial Department of the ``The Communist'' journal of the Liaoning
  Provincial Party Committee), ``文化大革命''若干大事件真相 (The facts
  behind several major incidents in the ``Great Cultural Revolution''),
  (Shenyang: Zhonggong Liaoning Shengwei Gongchandangyuan Zazhishe
  Chubanshe, 1985), pp.~82-84. For the implementation of the campaign in
  Jiangsu province, see Jiangshusheng dashiji, p.~291.}

按照这一论证逻辑，支持毛者自动成为无产阶级革命派，反对毛者则成为资产阶级反动派。这一方法绕开了受到攻击的党政干部对某些群众组织提出的阶级不纯指控。\footnote{CNS,
  220 (May 16, 1968), pp.~1-5.}
陈伯达在文章中引用毛的话说：``党外有党，党内有派，历来如此。''\footnote{Chen
  was apparently quoting, somewhat inexactly, from Mao's speech to the
  11th plenum of the 8th CC on 12 August 1966 in which the chairman had
  said:
  ``我们这个党不是党外无党;我们是党外有党,党内也有派,从来都是如此,这是正常现象.''
  (It's not true that there are no parties outside our party; we do have
  them and also factions within the party. It's always been like this,
  and this is a normal phenomenon.) Mao Zedong sixiang wansui (1969),
  p.~652.}
这些话为陈的论点提供了权威，即派别在阶级社会中不可避免，阶级和派别以某种方式交织在一起，派性斗争是阶级斗争的表现。尽管陈关于中共内部派性普遍性的看法，可能反映了对现实的准确而坦率的评估，\footnote{This
  position has been put most forcefully by Pye in, The Dynamics of
  Chinese Politics.}
但它显然同党在这一问题上的公开表述相矛盾。毛本人最终也从这一分析的含义后退。例如，1968年4月，他把文化大革命描述为``无产阶级和资产阶级之间阶级斗争的继续''。\footnote{E.
  L. Wheelwright and Bruce McFarlane, The Chinese Road to Socialism,
  p.~244.}
``继续革命''这一含糊概念背后的理论不清，及其在社会主义条件下关于阶级和阶级斗争的未充分展开且相互矛盾的观念，造成并暗中正当化了这种战术转向。\footnote{See
  White, The Politics of Class and Class Origin, p.~61; Kraus, Class
  Conflict in Chinese Socialism, pp.~110-14. For a thorough discussion
  of this issue, see Graham Young, ``Mao Zedong and the Class Struggle
  in Socialist Society'', Australian Journal of Chinese Affairs, 16
  (July 1986), pp.~41-80.}
派性敌人与阶级敌人之间界线的模糊，将在1968年``清理阶级队伍''运动发动时导致对造反派采取严厉措施。\footnote{For
  a brief description of the origins of this campaign, see
  中共辽宁省委《共产党员》杂志内部发行版编辑组 (Internal Publication
  Editorial Department of the ``The Communist'' journal of the Liaoning
  Provincial Party Committee), ``文化大革命''若干大事件真相 (The facts
  behind several major incidents in the ``Great Cultural Revolution''),
  (Shenyang: Zhonggong Liaoning Shengwei Gongchandangyuan Zazhishe
  Chubanshe, 1985), pp.~82-84. For the implementation of the campaign in
  Jiangsu province, see Jiangshusheng dashiji, p.~291.}

In Zhejiang, the revised definition of factionalism was propagated in an
editorial of 21 April.\footnote{ZJRB, April 21, 1968, p.~1.}
Furthermore, an article of 5 May quoted Mao to the effect that ``a
faction is the wing of a class'' (派别是阶级的一翼).\footnote{ZJRB, May
  5, 1968, p.~2.} Theory was being used instrumentally to fit the
requirements of political debate and factional struggle. On the basis of
this ideological innovation, opponents of the ZPRC were classified as
class enemies. In some areas, admitted Zhejiang Daily, power had
remained in the hands of party officials.\footnote{Problems in the rural
  areas of Zhejiang were revealed in an editorial of late March. See
  ZPS, March 27, 1968, SWB/FE/2738/B/17-18; CNS, 270, p.~5.} Unnamed
evil leaders were actively falsifying documents and spreading rumors
aimed at discrediting the PLA and the ZPRC.

在浙江，修订后的派性定义通过4月21日一篇社论得到宣传。\footnote{ZJRB,
  April 21, 1968, p.~1.}
此外，5月5日一篇文章引用毛的话说：``派别是阶级的一翼。''\footnote{ZJRB,
  May 5, 1968, p.~2.}
理论正被工具化，用来适应政治辩论和派性斗争的需要。基于这一意识形态创新，省革委会的反对者被归为阶级敌人。《浙江日报》承认，在一些地区，权力仍掌握在党政干部手中。\footnote{Problems
  in the rural areas of Zhejiang were revealed in an editorial of late
  March. See ZPS, March 27, 1968, SWB/FE/2738/B/17-18; CNS, 270, p.~5.}
未具名的坏头头正在积极伪造文件、散布谣言，企图败坏人民解放军和省革委会声誉。

One editorial seemed to indicate that United Headquarters had gone
almost completely onto the defensive. Leaders of mass organizations such
as Red Storm and the Wenzhou Rebel General HQ were accusing the members
of United Headquarters of being ``rightists engaged in subversion''. So
much for the threat issued by party leaders such as Zhou Enlai that mass
organizations would be dissolved after the establishment of
revolutionary committees. The editorial noted that officials who had
joined the new power structure were accused of being ``rightists
masquerading as leftists''.\footnote{See n.~72.} A commentator from Hong
Kong perceptively remarked at the time that Zhejiang ``seemed to be
defending the `rebels' against `conservatives'\,''. Commenting on
reports from Zhejiang in June the writer observed: ``one would have
thought that the authorities were on the side of the revolutionaries and
were afraid of a revival of the conservatives''.\footnote{CNA, 731,
  p.~3.} This trend had been apparent as early as April 1968.

一篇社论似乎显示，联合总部几乎完全转入守势。红暴、温州造反总指挥部等群众组织领导人指控联合总部成员是``搞颠覆的右派''。周恩来等党和国家领导人曾威胁说，革命委员会成立后群众组织将被解散，如今看来这一威胁已基本落空。社论指出，加入新权力结构的干部被指责为``披着左派外衣的右派''。\footnote{See
  n.~72.}
当时一名香港评论者敏锐地指出，浙江``似乎在保护`造反派'反对`保守派'\,''。这位作者评论6月来自浙江的报道时观察说：``人们本会以为当局站在革命派一边，并害怕保守派复辟。''\footnote{CNA,
  731, p.~3.} 这一趋势早在1968年4月就已显现。

From 9 to 31 May 1968 the ZPRC held a lengthy conference of CCP members,
effectively excluding from participation most of the mass organization
leaders, who were not members of the party.\footnote{In Jiande county,
  however, rebels had been gaining admission to the CCP since January
  1968. See Jiande xianzhi, p.~23.} The conference attempted to effect a
reconciliation between Red Storm and United Headquarters by issuing
advice to both groups. It counselled the former organization thus:
``Revolutionary mass organizations that have taken the wrong side must
handle their own mistakes correctly''. To United Headquarters it
enjoined it to treat its rival correctly. Only in this way, stressed the
conference, could an alliance be formed.\footnote{ZPS, May 31, 1968,
  SWB/FE/2794/B/4-5.} These words indicated, if nothing else, the total
breakdown of the agreement signed in February, made worthless by events
since March and justified by the revised definition of factionalism
propagated in April.

1968年5月9日至31日，省革委会召开一次中共党员长期会议，实际上把多数群众组织领导人排除在外，因为他们不是党员。\footnote{In
  Jiande county, however, rebels had been gaining admission to the CCP
  since January 1968. See Jiande xianzhi, p.~23.}
会议试图通过向两个团体提出建议来调解红暴和联合总部。它这样劝告前者：``站错队的革命群众组织必须正确对待自己的错误。''对联合总部，则告诫其正确对待对手。会议强调，只有这样才能形成联盟。\footnote{ZPS,
  May 31, 1968, SWB/FE/2794/B/4-5.}
这些话至少说明，2月签署的协议已经彻底破裂；3月以来的事件使其变得毫无价值，而4月宣传的修订后派性定义又为这种破裂提供了理由。

Yet intermittent gestures toward the unity of the mass organizations
were evident. On 26 May 1968, in an unprecedented show of unity,
Zhejiang Daily and Hangzhou Daily published a joint editorial together
with Red Storm's newspaper, ``Red Storm'', the ``Zhejiang Red Guard''
(浙江红卫兵) and the ``Hangzhou Worker'' (杭州工人). The editorial
seemed optimistic about the prospects for unity and declared that the
struggle to achieve victory in the Cultural Revolution was approaching a
climax.\footnote{ZJRB, May 26, 1968, p.~2. Exactly one year later to the
  day Red Storm and United Headquarters finally signed their agreement.}
In a further conciliatory gesture, Zhejiang Daily republished an article
originally carried in ``Red Storm''.\footnote{ZPS, June 25, 1968,
  SWB/FE/2808/B/12-13.}

然而，两个群众组织间仍可见零星团结姿态。1968年5月26日，在前所未有的团结展示中，《浙江日报》和《杭州日报》同红暴报纸《红暴》、《浙江红卫兵》和《杭州工人》联合发表社论。该社论似乎对团结前景持乐观态度，并宣称争取文化大革命胜利的斗争正在接近高潮。\footnote{ZJRB,
  May 26, 1968, p.~2. Exactly one year later to the day Red Storm and
  United Headquarters finally signed their agreement.}
作为进一步和解姿态，《浙江日报》转载了一篇原刊于《红暴》的文章。\footnote{ZPS,
  June 25, 1968, SWB/FE/2808/B/12-13.}

Reconciliation of the two mass organizations was accompanied by further
rehabilitation of cadres. An editorial in mid-May repeated Mao's
injunction that 90\% of cadres were good or comparatively good and
stated that those ``liberated'' by the proletarian revolutionaries
formed the core and backbone of revolutionary committees. The editorial
urged vigilance against cadres who tried to pass themselves off, in
words taken from Mao's description of Red Storm as ``revolutionary
cadres who have committed mistakes''.\footnote{ZPS, May 14, 1968,
  SWB/FE/2776/B/8-9; CNS, 270, p.~8.}

两个群众组织的和解伴随着干部进一步平反。5月中旬一篇社论重复毛的指示，即90\%的干部是好的或比较好的，并说由无产阶级革命派``解放''的干部构成革命委员会的核心和骨干。社论还敦促警惕那些套用毛描述红暴时的说法、把自己伪装成``犯过错误的革命干部''的干部。\footnote{ZPS,
  May 14, 1968, SWB/FE/2776/B/8-9; CNS, 270, p.~8.}

In late May, the Zhejiang authorities taking their cue from Jiang Qing
who at the end of March had vilified Tan Zhenlin as a
traitor,\footnote{Ji Xizhen, RMRB, February 26, 1979, p.~2.} organized
the resumption of rallies to denounce Tan by name.\footnote{See the
  article by Shen Ce in ZJRB, May 20, 1968, p.~2; ZJRB, May 21, 1968,
  p.~1; ZPS, May 30, 1968, SWB/FE/2786/B/10-11; ZPS, June 3, 1968,
  SWB/FE/2790/B/18-19; CNS, 270, p.~8.} These were followed by
intensified public attacks on Jiang Hua and his supporters.\footnote{ZPS,
  June 13, 1968, SWB/FE/2798/B/4; ZPS, June 27, 1968,
  SWB/FE/2810/B/23-4; ZJRB, June 28, 1968, pp.~1, 3.} On 19 May, Jiang
Qing along with Chen Boda, Yao Wenyuan and Jiang and Mao's daughter Xiao
Li, received Zhang Yongsheng and Du Yingxin in Beijing. Zhang and Du
reported on the campaign to ``purify class ranks'' at their fine arts
college. In her talk Jiang attacked Tan and Jiang Hua and encouraged
Zhang to take the lead in publicly denouncing them. On 11 June, Zhang
reported the conversation to Wang Hongwen in Shanghai.\footnote{Zhang
  Yongshengde zuizheng cailiao, pp.~2-5; SCMM, 622, pp.~6-10; Ye
  Yonglie, Yaoshi fuzi, pp.~296-99.} United Headquarters clearly hoped
that by linking Jiang with Tan, it could demonstrate Red Storm's
responsibility for the instability and violence which had plagued
Zhejiang since February 1967. Denunciation of Tan would pave the way. An
editorial of 27 June reminded readers that Jiang Hua had worked in the
province since liberation and his influence and power were not easily
eradicated. In an attempt to intimidate sceptics Zhejiang Daily warned
them not to underestimate Jiang:

5月底，浙江当局以江青为榜样；江青在3月底曾把谭震林斥为叛徒。\footnote{Ji
  Xizhen, RMRB, February 26, 1979, p.~2.}
当局组织恢复公开点名批判谭的集会。\footnote{See the article by Shen Ce
  in ZJRB, May 20, 1968, p.~2; ZJRB, May 21, 1968, p.~1; ZPS, May 30,
  1968, SWB/FE/2786/B/10-11; ZPS, June 3, 1968, SWB/FE/2790/B/18-19;
  CNS, 270, p.~8.} 随后，对江华及其支持者的公开攻击加剧。\footnote{ZPS,
  June 13, 1968, SWB/FE/2798/B/4; ZPS, June 27, 1968,
  SWB/FE/2810/B/23-4; ZJRB, June 28, 1968, pp.~1, 3.}
5月19日，江青同陈伯达、姚文元以及江青和毛的女儿小李，在北京接见张永生和杜英信。张、杜汇报了他们美术学院``清理阶级队伍''运动情况。江在谈话中攻击谭和江华，并鼓励张带头公开谴责他们。6月11日，张在上海向王洪文汇报了这次谈话。\footnote{Zhang
  Yongshengde zuizheng cailiao, pp.~2-5; SCMM, 622, pp.~6-10; Ye
  Yonglie, Yaoshi fuzi, pp.~296-99.}
联合总部显然希望通过把江同谭联系起来，证明红暴对1967年2月以来困扰浙江的不稳定和暴力负有责任。谴责谭将为此铺路。6月27日社论提醒读者，江华自解放以来就在浙江工作，其影响和权力不易根除。为恐吓怀疑者，《浙江日报》警告他们不要低估江：

\begin{quote}
We absolutely cannot consider Zhejiang's number one capitalist-roader a
``dead tiger''. This fellow has only been ``wounded'', not killed. His
fangs and claws are still dangerous. He is testing the wind and waiting
for an opportune moment to pounce on the revolutionary masses.

(决不能认为浙江党内头号走资派是''死老虎'';这个家伙没有''死'',只是受了伤。他有爪还有牙,正在窥测方向,伺机向革命人民猛扑过来.)\footnote{ZJRB,
  June 27, 1968, p.~1.}
\end{quote}

\begin{quote}
我们绝不能把浙江头号走资派看成一只``死老虎''。这个家伙只是``受伤''，并没有死。他的牙爪仍然危险。他正在试探风向，等待时机扑向革命群众。

（决不能认为浙江党内头号走资派是``死老虎''；这个家伙没有``死''，只是受了伤。他有爪还有牙，正在窥测方向，伺机向革命人民猛扑过来。）\footnote{ZJRB,
  June 27, 1968, p.~1.}
\end{quote}

An editorial of the following day warned Red Storm indirectly that it
too would suffer if the capitalist-roaders remounted Zhejiang's
political stage.

次日社论间接警告红暴，如果走资派重新登上浙江政治舞台，它也会受害。

Evidence of organized opposition to the new provincial administration
had emerged to give some substance to Zhejiang Daily's hysterical
outbursts. A report published over a year later attributed
responsibility for this opposition to Fang Jianwen, one of the leaders
of Red Storm.\footnote{ZJRB, August 21, 1969.} In June 1968, cadres
sympathetic to Red Storm working in departments of the old CCP ZPC and
HMC pasted up wall posters attacking the provincial leadership. One
particular poster, entitled ``Arise and charge into the storm'', openly
advocated a bloody struggle against the new provincial
hierarchy.\footnote{Qiu Honggen, HZRB, November 26, 1978.} Factional
fighting in the factories of Hangzhou flared up again. Members of United
Headquarters at the Hangzhou Silk Complex reportedly started a fire in
the mill's boiler-room on 9 July 1968 causing damage estimated at 20,000
yuan and injuring several people.\footnote{HZRB, April 17, 1979.}

有组织反对新省级行政机构的证据已经出现，使《浙江日报》歇斯底里的爆发多少有了依据。一年多后发表的一份报告把这种反对归咎于红暴领导人之一方剑文。\footnote{ZJRB,
  August 21, 1969.}
1968年6月，在旧省委和杭州市委各部门工作的亲红暴干部张贴大字报攻击省级领导层。其中一张题为``起来，冲进风暴''的大字报公开主张对新的省级领导层展开血腥斗争。\footnote{Qiu
  Honggen, HZRB, November 26, 1978.}
杭州工厂中的派性斗争再次爆发。据报道，杭州丝绸联合厂的联合总部成员于1968年7月9日在厂锅炉房纵火，造成约2万元损失，并使数人受伤。\footnote{HZRB,
  April 17, 1979.}

Even more worrying for the new leadership of Zhejiang was an open
letter, published on 6 July, denigrating the ZPRC, the PLA and, by
implication, Mao Zedong. Zhejiang Daily later published an article
quoting extensively from the letter.\footnote{Zhejiang Provincial
  Workers Congress and Hangzhou Municipal Revolutionary Trade Union
  Council, ``Smash the criminal plot to subvert the Red political
  power'', ZJRB, October 11, 1968, p.~3.} It apparently made three major
points, all unflattering, about the nature of political power in
Zhejiang. First, it alleged that ``The power of the provincial
revolutionary committee is usurped by counterrevolutionary
double-dealers''. The letter referred to the members of the three-way
alliance -- cadres, mass representatives and PLA -- respectively as
``true capitalist-roaders'', ``clowns'' and ``pawns'', and
``counter-revolutionary double-dealers''. The oppositionists described
the ZPRC as an ``illegal faction committee'' which they would smash so
as to become ``masters of the Zhejiang political stage''. In the
meantime they refused to recognize, support or join the revolutionary
committee. The authors of the Zhejiang Daily article distinguished this
approach from that adopted by members of Red Storm who, it claimed,
``have a deep feeling they wield power and have become masters''
(作了主人).

更令浙江新领导层担忧的是7月6日发表的一封公开信，它贬低省革委会、人民解放军，并含蓄地贬低毛泽东。《浙江日报》后来发表文章，大量引用这封信。\footnote{Zhejiang
  Provincial Workers Congress and Hangzhou Municipal Revolutionary Trade
  Union Council, ``Smash the criminal plot to subvert the Red political
  power'', ZJRB, October 11, 1968, p.~3.}
它显然对浙江政权性质提出三点主要看法，且都是否定性的。第一，它声称``省革命委员会的权力被反革命两面派篡夺''。这封信把三结合联盟的成员
-- 干部、群众代表和人民解放军 --
分别称为``真走资派''、``小丑''和``卒子''，以及``反革命两面派''。反对者把省革委会描述为``非法派性委员会''，并称将砸烂它，以成为``浙江政治舞台的主人''。与此同时，他们拒绝承认、支持或加入革命委员会。《浙江日报》文章作者把这种态度同红暴成员的态度区分开来，声称后者``深深感到自己掌了权，作了主人''。

Finally, the letter accused the authorities of exercising ``bourgeois
dictatorship'' and substituting coercion for extensive democracy. It
described the attitude taken toward Red Storm as ``suppression of the
old rebels and the masses of the people''. Apart from claiming that
central directives were responsible for the ``suppression of the
masses'', the letter-writers expressed support for opponents of the ZPRC
who were put on trial in August 1968. They also proclaimed, in relation
to further trouble which had erupted in Wenzhou, that ``South Zhejiang
is the world of the Rebel General HQ''. While the letter certainly
reflected the opinions and feelings of opponents of United Headquarters,
Zhejiang Daily excused Red Storm from sharing its sentiments. However,
this charitable interpretation may have been dictated more by political
expediency than by respect for the truth.

最后，这封信指控当局实行``资产阶级专政''，以强制取代大民主。它把对红暴采取的态度描述为``镇压老造反派和人民群众''。除声称中央指示要为``镇压群众''负责外，写信者还支持1968年8月受审的省革委会反对者。他们还针对温州再次发生的动乱宣称：``浙南是造反总司的天下。''这封信当然反映了联合总部反对者的观点和情绪，但《浙江日报》免除了红暴分享这些情绪的责任。然而，这种宽厚解释可能更多出于政治权宜，而非尊重事实。

Other critics attacked the ZPRC from the left putting forward such
slogans as ``election by the whole people'', and demanding the smashing
of state organs. They maligned the revolutionary committees as
committees of the Guomindang.\footnote{ZPS, July 15, 1968,
  SWB/FE/2828/B/16-18; CNS, 270, pp.~10-12. See also, CNS, 233 (August
  15, 1968), pp.~1-6 for other examples monitored from the Zhejiang
  radio of attacks from the left on revolutionary committees.} The
attacks, and the way in which they were phrased, seem to owe their
rhetoric to the program of the Hunan ultra-leftist group Shengwulian
(省无联).\footnote{For statements of Shengwulian's position see
  Classified Chinese Communist Documents: A Selection, (Taibei:
  Institute of International Relations, 1978), pp.~269-73; China: The
  Revolution is Dead - Long Live the Revolution (Montreal: Black Rose
  Books, 1979), pp.~153-70.}

其他批评者从左的方面攻击省革委会，提出``全民选举''等口号，并要求砸烂国家机关。他们诬蔑革命委员会是国民党委员会。\footnote{ZPS,
  July 15, 1968, SWB/FE/2828/B/16-18; CNS, 270, pp.~10-12. See also,
  CNS, 233 (August 15, 1968), pp.~1-6 for other examples monitored from
  the Zhejiang radio of attacks from the left on revolutionary
  committees.}
这些攻击及其措辞，似乎在修辞上借用了湖南极左组织省无联的纲领。\footnote{For
  statements of Shengwulian's position see Classified Chinese Communist
  Documents: A Selection, (Taibei: Institute of International Relations,
  1978), pp.~269-73; China: The Revolution is Dead - Long Live the
  Revolution (Montreal: Black Rose Books, 1979), pp.~153-70.}

The events which occurred in the immediate aftermath of the
establishment of the ZPRC proved that both United Headquarters and Red
Storm lacked commitment to the agreement imposed on them in February
1968. The alliance was uneasy and fragile and liable to degenerate into
open conflict at the slightest provocation. The establishment of the
provincial revolutionary committee had raised as many questions as it
was supposed to solve. While a civilian administration of sorts had
finally replaced direct military rule in Zhejiang, the new political
authority remained under PLA domination. It was racked by internal
discord and constantly threatened by political players who, through
either preference or necessity, stayed outside the political mainstream.
When verbal abuse and other provocations were supplemented by armed
challenges to the new committees, the military was again forced to
intercede and come down very hard on the instigators.

省革委会成立后立即发生的事件证明，联合总部和红暴都没有真心遵守1968年2月强加给他们的协议。联盟不安而脆弱，稍有挑衅就可能退化为公开冲突。省革委会的成立提出的问题并不少于它本应解决的问题。尽管某种形式的文职行政机构终于取代了浙江的直接军事统治，新的政治权威仍受人民解放军支配。它内部纷争不断，并始终受到那些因偏好或必要而留在政治主流之外的政治参与者威胁。当辱骂和其他挑衅又加上对新委员会的武装挑战时，军方再次被迫介入，并严厉打击煽动者。

\section{Notes}\label{notes-3}

\section{注释}\label{ux6ce8ux91ca-3}

\bookmarksetup{startatroot}

\chapter{Chapter Four - ORDER AND UNITY, AUGUST 1968-SEPTEMBER 1969 /
第四章 -
秩序与团结，1968年8月至1969年9月}\label{chapter-four---order-and-unity-august-1968-september-1969-ux7b2cux56dbux7ae0---ux79e9ux5e8fux4e0eux56e2ux7ed31968ux5e748ux6708ux81f31969ux5e749ux6708}

\section{Crackdown, August-September
1968}\label{crackdown-august-september-1968}

\section{镇压，1968年8月至9月}\label{ux9547ux538b1968ux5e748ux6708ux81f39ux6708}

In July 1968, Beijing ordered the military in Zhejiang to crack down on
disorder. On 3 and 24 July the central authorities issued two notices to
quell disturbances in the Guangxi Autonomous Region and Shaanxi province
respectively.\footnote{``Notice of the Central Committee, State Council,
  CC Military Affairs Commission and Cultural Revolution Group'', July
  3, 1968, Chien, China's Fading Revolution, pp.~295-97; Notice of idem,
  July 24, 1968, ibid., pp.~297-300; Li Pingan, Jia Chunfu, Tian Zizhi
  and Li Birong (eds.), 陕西经济大事记 (1949-1985) (Chronicle of
  Shaanxi's Economy, 1949-1985), (Xian: Sanqin Chubanshe, 1987),
  pp.~300-02.} However, the message had national relevance. The notices
denounced disruption to railway traffic, storming of PLA units and
clashes with soldiers perpetrated by Red Guards and revolutionary
rebels. Those responsible were called counter-revolutionaries and the
military was instructed to spare no effort in suppressing them. Other
anarchic activities included the looting of state property, the burning
down of public and private buildings, refusal to comply with central
directives, broadcasts by unauthorized radio stations and raids on state
prisons and labour farms to release inmates. In a meeting with Red Guard
leaders on 28 July Mao expressed disillusionment with their
performance.\footnote{Wansui, (1969), pp.~687-716.} An editorial of
People's Daily in early August attacked the ``theory of many centers''
(多中心论), and demanded obedience to central discipline and disavowal
of anarchistic or individualistic tendencies.\footnote{RMRB, August 5,
  1968, p.~1; see Facts \& Features, 1: 23, pp.~1-4, for a discussion of
  this editorial.}

1968年7月，北京命令浙江军方镇压混乱。7月3日和24日，中央当局分别发布两份通知，平息广西自治区和陕西省的动乱。\footnote{``Notice
  of the Central Committee, State Council, CC Military Affairs
  Commission and Cultural Revolution Group'', July 3, 1968, Chien,
  China's Fading Revolution, pp.~295-97; Notice of idem, July 24, 1968,
  ibid., pp.~297-300; Li Pingan, Jia Chunfu, Tian Zizhi and Li Birong
  (eds.), 陕西经济大事记 (1949-1985) (Chronicle of Shaanxi's Economy,
  1949-1985), (Xian: Sanqin Chubanshe, 1987), pp.~300-02.}不过，这一信息具有全国意义。通知谴责红卫兵和革命造反派破坏铁路交通、冲击解放军单位并同士兵发生冲突。责任者被称为反革命，军队被指示要不遗余力地镇压他们。其他无政府活动包括抢掠国家财产、焚烧公私建筑、拒不执行中央指示、未经授权的电台广播，以及袭击国家监狱和劳改农场以释放囚犯。7月28日，毛在同红卫兵领导人会见时表达了对其表现的失望。\footnote{Wansui,
  (1969), pp.~687-716.}8月初《人民日报》一篇社论批判``多中心论''，要求服从中央纪律，摒弃无政府主义和个人主义倾向。\footnote{RMRB,
  August 5, 1968, p.~1; see Facts \& Features, 1: 23, pp.~1-4, for a
  discussion of this editorial.}

It was at the second enlarged plenum of the ZPRC, held from 17 to 29
July, that major decisions concerning Wenzhou (see below) and summing up
the Cultural Revolution in the province were made.\footnote{ZJRB, July
  31, 1968, pp.~1, 3; ZPS, August 1, 1968, SWB/FE/2841/B/11; CNS, 271,
  p.~5. See also the ZJRB editorials, ZPS, July 21, 1968,
  SWB/FE/2835/B/15-16; ZPS, July 24, 1968, SWB/FE/2835/B/13-15.} Two
hundred and forty delegates attended, including members of the HMRC and
responsible personnel from counties and major industrial units and
colleges across the province. Surprisingly, the meeting expressed
optimism about the situation in Zhejiang after its examination of the
performance of the ZPRC. A five-point resolution encapsulated the
decisions taken by the plenum.

有关温州（见下文）的重大决定以及对全省文化大革命的总结，是在7月17日至29日召开的省革委会第二次扩大会议上作出的。\footnote{ZJRB,
  July 31, 1968, pp.~1, 3; ZPS, August 1, 1968, SWB/FE/2841/B/11; CNS,
  271, p.~5. See also the ZJRB editorials, ZPS, July 21, 1968,
  SWB/FE/2835/B/15-16; ZPS, July 24, 1968, SWB/FE/2835/B/13-15.}240名代表出席会议，包括杭州市革委会成员，以及全省各县、主要工业单位和高校的负责人。令人意外的是，会议在考察省革委会工作后，对浙江形势表示乐观。全会通过的决定被概括为一项五点决议。

The first resolution stated that the meeting ``voiced constructive
criticism'' of the failings of the ZPRC leadership and ``strong
criticism of and willingness to help members of the Revolutionary
Committee who diverged from Mao Zedong's Thought''. Another resolution,
by contrast, repudiated those who had expressed pessimistic views about
the state of affairs in Zhejiang. Resolution four demanded that both
mass organizations desist from mutual provocation and expel bad elements
from their ranks as a precondition to forming an alliance. Torturing of
political opponents was forbidden, as was interference in the internal
affairs of the PLA. The resolution concluded with a ringing cry for
victory in the Cultural Revolution. It was clear that Beijing expected
positive action to this end. The luxury of endless debate and
interminable disputes over the finer points of right and wrong was no
longer to be enjoyed by mass organizations. They would conform or face
the consequences.

第一项决议称，会议对省革委会领导的缺点``提出了建设性批评''，并对``背离毛泽东思想的革命委员会成员''进行了``有力批评并愿意帮助''。相较之下，另一项决议驳斥了那些对浙江局势表达悲观观点的人。第四项决议要求两大群众组织停止相互挑衅，并把坏分子清除出队伍，作为实现联合的前提。禁止折磨政治对手，也禁止干涉解放军内部事务。决议以争取文化大革命胜利的响亮号召作结。显然，北京期望为此采取积极行动。群众组织不再能够享受无休止辩论和没完没了争论是非细枝末节的奢侈。它们要么服从，要么承担后果。

The show of force by the military in southern Zhejiang was backed up by
a series of trials which began in August 1968.\footnote{For an account
  of trials at this time in Shanxi, see Hinton, Shenfan, pp.~638-40, and
  in Guangzhou in December 1968 see Domes, ``Generals and Red Guards'',
  2, p.~150.} Their purpose was to impose harsh sentences on anti-social
elements who had taken advantage of the prevailing lawlessness to commit
various kinds of criminal acts. The trials also dealt with
trouble-makers of a political kind who had worked to undermine the new
power structure and refused to accept its authority. In particular,
those people who persisted in provoking factional struggles were
severely punished. Propaganda teams made up of soldiers, cadres,
peasants and workers were sent into major educational institutions to
pacify the students. The youthful proletarian revolutionaries had little
option but to accept their fate.

军方在浙南的武力展示，得到了1968年8月开始的一系列审判的支持。\footnote{For
  an account of trials at this time in Shanxi, see Hinton, Shenfan,
  pp.~638-40, and in Guangzhou in December 1968 see Domes, ``Generals
  and Red Guards'', 2, p.~150.}这些审判旨在严惩利用普遍无法无天局面犯下各种刑事罪行的反社会分子。审判还处理政治性闹事者，他们破坏新权力结构并拒绝承认其权威。尤其是那些坚持挑动派性斗争的人受到严厉惩罚。由士兵、干部、农民和工人组成的宣传队被派入主要教育机构，以安抚学生。年轻的无产阶级革命派几乎别无选择，只能接受命运。

The most sensational trial was that of Nie Minzhi (聂民之), Xia
Liangchang (夏良焯) and Peng Zhangxun (彭彰训).\footnote{ZJRB, August
  10, 1968, pp.~1, 2; CNS, 271, pp.~6-7.} Peng was a PLA soldier who was
described as having a ``dubious'' family background, Nie a cadre, and
Xia was referred to as a swindler, extortionist and rapist. They had
apparently caused a great deal of trouble over the previous two years.
In the autumn of 1967 the trio had joined with unspecified mass
organizations -- most probably Red Storm and Wenzhou's Rebel General HQ
-- to disrupt efforts to bring the situation under control after the
violence of the previous months. Nie, Xia and Peng had organized forces
in mountainous and isolated areas with the avowed aim of bringing down
United Headquarters and its military and civilian allies. Nie had posed
as a correspondent and the secretary of an unnamed leader in the central
government. With these credentials he convened meetings of mass
organization leaders to ``use the village to encircle the cities'' from
armed bases. The trio opposed the August 1967 reorganization of the
military control commission and supported Jiang Hua whom, they claimed
with justification, was protected by the Central Committee. Even after
the formation of the ZPRC, in conjunction with leaders of local mass
organizations they had allegedly plotted to establish an ``armed
independent regime'' in the mountains.

最轰动的审判是聂民之、夏良焯和彭彰训案。\footnote{ZJRB, August 10, 1968,
  pp.~1, 2; CNS, 271, pp.~6-7.}彭是解放军士兵，被称为家庭出身``可疑''；聂是干部；夏则被称作骗子、勒索者和强奸犯。显然，他们在此前两年制造了大量麻烦。1967年秋，三人同未具体说明的群众组织，最可能是红暴和温州造反总指挥部，联合起来，破坏此前数月暴力之后控制局势的努力。聂、夏、彭在山区和偏僻地区组织力量，公开目标是推翻联合总部及其军民盟友。聂曾冒充中央政府某未具名领导的记者和秘书。凭借这些身份，他召集群众组织领导人开会，企图以武装基地``用农村包围城市''。三人反对1967年8月军管会改组，并支持江华，声称江华受到中央委员会保护，这一说法并非毫无根据。即便在省革委会成立之后，据称他们仍同地方群众组织领导人密谋在山区建立``武装独立王国''。

The article speculated as to how the three men could have been so active
and have built up so much influence, how they could have coordinated
their activities so closely, and how they had become so popular and
important ``as to be able to issue directives, receive reports and talk
about the situation at gatherings of thousands''. Minutes of these
meetings had been printed and extensively disseminated. Zhejiang Daily
also wondered why the arrest of these three ringleaders had not brought
about a cessation of ``criminal activities'' especially in southern
Zhejiang, and why they had been openly defended after their arrest. The
newspaper did not have many convincing answers to these rhetorical
questions but it is clear that such an organized resistance force,
backed up by former officials, would have posed enormous problems for
the relatively inexperienced and far from united new administration. The
report confirmed that the Cultural Revolution lacked support in the
countryside and that peasants and grassroots cadres had remained loyal
to their old leaders and patrons. A mid-July 1968 Zhejiang Daily
editorial analyzed the phenomenon in Zhuji county, Shaoxing district. It
asked rhetorically:

文章推测，这三人为何能如此活跃并建立如此大的影响，他们如何能如此紧密地协调行动，又为何会变得如此受欢迎、如此重要，以至于``能够在数千人的集会上发号施令、听取汇报并谈论形势''。这些会议记录曾被印刷并广泛散发。《浙江日报》还追问，为什么这三名头目被捕并未导致``犯罪活动''停止，尤其是在浙南；为什么他们被捕后还得到公开辩护。报纸对这些修辞性问题没有多少令人信服的答案，但显然，这样一支有组织的抵抗力量，并得到前官员支持，会给经验相对不足且远未团结的新政权造成巨大问题。报道证实，文化大革命在农村缺乏支持，农民和基层干部仍忠于旧领导和庇护人。1968年7月中旬，《浙江日报》一篇社论分析了绍兴地区诸暨县的这种现象。它反问道：

\begin{quote}
``In earlier periods of the great proletarian cultural revolution in
Zhuji county, why was resistance so powerful? Why was the struggle by
force so serious? Why were the conservative forces so strong?'' One
answer it supplied was that ``The cobras who hide themselves under the
cloak of old Party members, old cadres, old revolutionaries and old
guerrillas have remained in their own localities for more than a decade
or even several decades''.\footnote{ZPS, July 19, 1968,
  SWB/FE/2832/B/23; CNS, 271, pp.~1-2.}
\end{quote}

\begin{quote}
``在诸暨县无产阶级文化大革命前期，为什么阻力这样大？为什么武斗这样严重？为什么保守势力这样强？''它给出的一个答案是，``那些披着老党员、老干部、老革命、老游击队外衣的眼镜蛇，在自己的地方盘踞了十几年甚至几十年''。\footnote{ZPS,
  July 19, 1968, SWB/FE/2832/B/23; CNS, 271, pp.~1-2.}
\end{quote}

Further trials took place in Hangzhou at which sentences ranging from
ten years' imprisonment to the death penalty were handed
down.\footnote{ZPS, August 17,1968, SWB/FE/2852/B/16-18; ZPS, August
  31,1968, SWB/FE/2863/B/7-8; CNS, 271, pp.~8,11-12.} Trials were also
conducted in provincial centers with editorials supporting the tough
penalties being imposed.\footnote{ZPS, September 11, 1968,
  SWB/FE/2872/B/6-7; Bulgaria News Agency (BTA), October 16, 1968,
  SWB/FE/2908/B/10; ZPS, September 13 and 16,1968,
  SWB/FE/2877/B/14-16,17-18.} On 9 and 26 September 1968, the ZPRC
issued further notices against ``counter-revolutionary''
activity.\footnote{Qiu Honggen, HZRB, November 26,1978.} At later trials
in Hangzhou a new organization called the Hangzhou Social Order
Headquarters (杭州市社会治安指挥部) was mentioned publicly for the first
time. Members of the organization acted in place of public security
personnel, escorting the arrested suspects before the tribunals. A
leader of United Headquarters, He Xianchun, a short intense man who
liked fighting,\footnote{``The Case of Weng Sengho'', p.~3.} and Xia
Genfa organized and led the contingent.\footnote{The CCP HMC Propaganda
  Department criticism group, ``The counter-revolutionary He Xianchun'',
  HZRB, June 20,1977.}

杭州随后又举行审判，判决从十年徒刑到死刑不等。\footnote{ZPS, August
  17,1968, SWB/FE/2852/B/16-18; ZPS, August 31,1968, SWB/FE/2863/B/7-8;
  CNS, 271, pp.~8,11-12.}省内各中心城市也进行了审判，社论支持正在实施的严厉处罚。\footnote{ZPS,
  September 11, 1968, SWB/FE/2872/B/6-7; Bulgaria News Agency (BTA),
  October 16, 1968, SWB/FE/2908/B/10; ZPS, September 13 and 16,1968,
  SWB/FE/2877/B/14-16,17-18.}1968年9月9日和26日，省革委会又发布反对``反革命''活动的通知。\footnote{Qiu
  Honggen, HZRB, November 26,1978.}在杭州后来的审判中，一个名为杭州市社会治安指挥部的新组织首次被公开提及。该组织成员代替公安人员行事，把被捕嫌疑人押送到审判庭。联合总部领导人贺贤春，一个矮小精悍、喜欢打斗的人，\footnote{``The
  Case of Weng Sengho'', p.~3.}以及夏根发组织并领导了这支队伍。\footnote{The
  CCP HMC Propaganda Department criticism group, ``The
  counter-revolutionary He Xianchun'', HZRB, June 20,1977.}

Such groups appeared in other Chinese cities in 1968.\footnote{County-level
  commands were also established. For example, in Xianju county the
  social order headquarters was set up on September 19,1968 and
  dissolved on February 15,1970. See 仙居县志编纂委员会 (Xianju County
  Gazetteer Compilation Committee), 仙居县志 (Xianju County Gazetteer),
  (Hangzhou: Zhejiang Renmin Chubanshe, 1987), p.~22. In Wuyi county the
  headquarters was established on September 18,1968. See
  《义乌县志》编纂委员会 (``Yiwu County Gazetteer'' Compilation
  Committee), 义乌县志 (Yiwu County Gazetteer), (Hangzhou: Zhejiang
  Renmin Chubanshe, 1987), p.~27. In Xiangshan county the command was
  founded as early as January 19,1968 and in Kaihua county on June 27,
  1969. See Xiangshan xianzhi, p.~34; Kaihua xianzhi, p.~24.} Their
duties included patrolling the streets, checking registration bona fides
and guarding factories and government installations, tasks that had
previously been the responsibility of the public security forces and the
PLA. The social order command thus relieved the military of some of its
civil commitments.\footnote{Nelsen, ``Military Bureaucracy in the
  Cultural Revolution'', AS, 14:4 (1974), p.~378.} However, the worker
provosts made themselves very unpopular by their abuse of power and
excessive use of force.\footnote{See ``China's Revolutionary
  Committees'', CS, 6:21 (December 6, 1968), p.~13.} On 7 August 1968,
Minister of Public Security Xie Fuzhi launched a blistering attack on
the public security/procuratorate/court network. Xie called for the
complete destruction of the system (彻底砸烂公检法) and in response a
wholesale purge of the apparatus took place. The creation of social
order commands was thus associated with the smashing of the old security
network.\footnote{See ``Wenhua dageming'' ruogan dashijian zhenxiang,
  pp.~90-1.}

1968年，中国其他城市也出现了此类组织。\footnote{County-level commands
  were also established. For example, in Xianju county the social order
  headquarters was set up on September 19,1968 and dissolved on February
  15,1970. See 仙居县志编纂委员会 (Xianju County Gazetteer Compilation
  Committee), 仙居县志 (Xianju County Gazetteer), (Hangzhou: Zhejiang
  Renmin Chubanshe, 1987), p.~22. In Wuyi county the headquarters was
  established on September 18,1968. See 《义乌县志》编纂委员会 (``Yiwu
  County Gazetteer'' Compilation Committee), 义乌县志 (Yiwu County
  Gazetteer), (Hangzhou: Zhejiang Renmin Chubanshe, 1987), p.~27. In
  Xiangshan county the command was founded as early as January 19,1968
  and in Kaihua county on June 27, 1969. See Xiangshan xianzhi, p.~34;
  Kaihua xianzhi, p.~24.}它们的职责包括街头巡逻、检查户籍证明、守卫工厂和政府设施，这些任务此前由公安力量和解放军承担。因此，社会治安指挥部减轻了军队的一部分民事职责。\footnote{Nelsen,
  ``Military Bureaucracy in the Cultural Revolution'', AS, 14:4 (1974),
  p.~378.}然而，工人纠察由于滥用权力和过度使用武力而变得极不受欢迎。\footnote{See
  ``China's Revolutionary Committees'', CS, 6:21 (December 6, 1968),
  p.~13.}1968年8月7日，公安部长谢富治猛烈攻击公安、检察、法院网络。谢要求彻底摧毁这一系统（彻底砸烂公检法），作为回应，整个机构遭到全面清洗。社会治安指挥部的建立因而同砸烂旧安全网络联系在一起。\footnote{See
  ``Wenhua dageming'' ruogan dashijian zhenxiang, pp.~90-1.}

The first reported mobilization of the social order command took place
in August 1968.\footnote{ZPS, August 14,1968, SWB/FE/2850/B/9-10.} On 6
August, discord erupted between the factions at the Hangzhou No.~1
Cotton Printing and Dyeing Mill.\footnote{Established in 1897 it is the
  oldest cotton mill in Hangzhou. After liberation three mills joined
  together to become known as the No.~1 cotton mill, and when a dyeing
  plant was added in 1959, it had 6,200 employees on the payroll. In
  September 1977 the dyeing mill was returned to separate management,
  leaving 4,649 employees. Visit to No.~1 cotton mill, April 1,1979.}
The mill was the unit of one of Red Storm's leaders, Fang
Jianwen.\footnote{ZJRB, August 21,1969.} Fang's faction there had
refused to obey a cease-fire order and form an alliance with its rival
group, so the social order command was directed to intervene. The
workers had previously disregarded the presence of PLA troops and orders
from the HMRC to stop the fighting. Leaders of the social order command
probably had less scruples in applying some muscle. On the afternoon of
7 August, the HMRC called on the command to move into the factory. After
the social order command issued eight open letters in rapid succession
without producing any effect, it sent in Peoples' Guards (人民警卫) to
``awaken'' the resisters. In thirteen hours they had brought the
situation under control and by the afternoon of 8 August production had
reportedly returned to normal.

关于社会治安指挥部动员的首次报道发生在1968年8月。\footnote{ZPS, August
  14,1968, SWB/FE/2850/B/9-10.}8月6日，杭州第一棉纺印染厂各派之间爆发冲突。\footnote{Established
  in 1897 it is the oldest cotton mill in Hangzhou. After liberation
  three mills joined together to become known as the No.~1 cotton mill,
  and when a dyeing plant was added in 1959, it had 6,200 employees on
  the payroll. In September 1977 the dyeing mill was returned to
  separate management, leaving 4,649 employees. Visit to No.~1 cotton
  mill, April 1,1979.}该厂是红暴领导人之一方剑文所在单位。\footnote{ZJRB,
  August 21,1969.}方在厂内的派别拒绝服从停火命令，也拒绝同对立派别联合，因此社会治安指挥部受命介入。此前，工人们无视解放军部队在场，也不执行杭州市革委会停止武斗的命令。社会治安指挥部领导人在动用力量方面大概顾虑更少。8月7日下午，杭州市革委会要求指挥部进驻工厂。社会治安指挥部连续发出八封公开信而毫无效果后，派人民警卫进入厂内``唤醒''抵抗者。13小时内，他们控制了局势，据报到8日下午生产已恢复正常。

While the social order command dealt with factional disturbances and
problems of law enforcement, the propaganda teams concentrated their
attention on students and disturbances in the localities. On 19 August
the ZPRC decided to establish worker-peasant propaganda teams and the
first units were formed three days later.\footnote{ZJRB, August 20,1968,
  p.~1; ZPS, August 24,1968, SWB/FE/2859/B/40.} A ``suitable number'' of
cadres supplemented and led the workers and peasants who comprised the
bulk of the membership. The task of the propaganda teams was to ``play a
leading role in helping places and units where the class enemies are
making trouble''. The ZPRC decision also requested local PLA units to
organize teams and assign cadres and troops to staff them. Four brigades
from Hangzhou, Shaoxing, Jiaxing and Jinhua were formed\footnote{These
  four districts had advanced furthest along the road to unity. By late
  August 1968 Hangzhou Municipality and Shaoxing and Jinhua districts
  had established revolutionary committees. Jiaxing formed its committee
  on September 16. The southern districts of Lishui, Taizhou and Wenzhou
  lagged considerably behind their northern counterparts, as did the
  eastern districts of Ningbo and Zhoushan Islands.} with a total
membership of 10,000 and they were immediately despatched to southern
Zhejiang.\footnote{ZPS, August 24,1968, SWB/FE/2862/B/8.}

社会治安指挥部处理派性骚乱和执法问题，宣传队则把注意力集中在学生和地方动乱上。8月19日，省革委会决定成立工农宣传队，第一批队伍三天后组成。\footnote{ZJRB,
  August 20,1968, p.~1; ZPS, August 24,1968, SWB/FE/2859/B/40.}``适当数量''的干部补充并领导由工人和农民构成主体的队伍。宣传队的任务是在``阶级敌人捣乱的地方和单位中起领导帮助作用''。省革委会的决定还要求地方解放军单位组织队伍，并派干部和部队参加。杭州、绍兴、嘉兴和金华组建了四个大队，\footnote{These
  four districts had advanced furthest along the road to unity. By late
  August 1968 Hangzhou Municipality and Shaoxing and Jinhua districts
  had established revolutionary committees. Jiaxing formed its committee
  on September 16. The southern districts of Lishui, Taizhou and Wenzhou
  lagged considerably behind their northern counterparts, as did the
  eastern districts of Ningbo and Zhoushan Islands.}总人数达10,000人，并立即派往浙南。\footnote{ZPS,
  August 24,1968, SWB/FE/2862/B/8.}

Further violence had been reported from Wenzhou in July 1968.\footnote{The
  following paragraphs are based on ZPS, July 7, 1968, SWB/FE/2823/B/12;
  ZPS, July 26, 1968, SWB/FE/2840/B/8-9; ZPS, July 26, 1968,
  SWB/FE/2835/B/3; ZPS, July 27, 1968, SWB/FE/2846/B/17; July 28, 1968,
  SWB/FE/2834/B/16-17; Polish Press Agency (PAP), July 28, 1968,
  SWB/FE/2835/B/16; ZPS, August 8, 1968, SWB/FE/2846/B/20; ZJRB, August
  9, 1968, pp.~1, 2; ZPS, August 9, 1968, SWB/FE/2856/B/14-15; ZPS,
  August 25, 1968, SWB/FE/2862/B/8-9; ZPS, August 28, 1968,
  SWB/FE/B/9-10; ZPS, August 31, 1968, SWB/FE/2865/B/36; ZPS, September
  4, 1968, SWB/FE/2870/B/3-4; ZPS, September 6, SWB/FE/2883/B/28; ZPS,
  September 28, 1968, SWB/FE/2889/B/24-5; ZJRB, October 11, 1968, SCMP
  (S), 239, pp.~14-16; ZPS, October 15, 1968, SWB/FE/2908/B1/10; ZPS,
  November 4, 1968, SWB/FE/2926/B1/3; ZPS, November 18, 1968,
  SWB/FE/2930/B/3; ZPS, November 27, 1968, SWB/FE/2949/B1/10; ZPS,
  December 1, 1968, SWB/FE/B1/2-3; ZPS, December 5, 1968, SWB/FE/B1/7;
  ZJRB, December 6, 1968, p.~1; ZPS, December 8, 1968, ibid., B1/6-7;
  editorial of the Southern Zhejiang Public Press (?), December 16,
  1968, in Facts \& Features, 2: 8 (February 5, 1969), pp.~26-30; ZPS,
  December 17, 1968, SWB/FE/2955/B/8-9; ZPS, July 14, 1969, URS, 56:9
  (July 29, 1969), pp.~127-28; CNA, 731, p.~5; Qiu Honggen, HZRB, 26
  November 1978; CNS, 270, p.~10; CNS, 271 (May 22, 1969), pp.~2-4, 6,
  9-11, 13-14; CNS, 272, (May 29, 1969), pp.~5-7.} In the latter half of
the previous year the city had caused a massive headache to the military
authorities and delivered more than a few cracked skulls to members of
United Headquarters who had ventured down there. An editorial of
Zhejiang Daily in late July 1968 revealed that

1968年7月，温州又传出暴力事件。\footnote{The following paragraphs are
  based on ZPS, July 7, 1968, SWB/FE/2823/B/12; ZPS, July 26, 1968,
  SWB/FE/2840/B/8-9; ZPS, July 26, 1968, SWB/FE/2835/B/3; ZPS, July 27,
  1968, SWB/FE/2846/B/17; July 28, 1968, SWB/FE/2834/B/16-17; Polish
  Press Agency (PAP), July 28, 1968, SWB/FE/2835/B/16; ZPS, August 8,
  1968, SWB/FE/2846/B/20; ZJRB, August 9, 1968, pp.~1, 2; ZPS, August 9,
  1968, SWB/FE/2856/B/14-15; ZPS, August 25, 1968, SWB/FE/2862/B/8-9;
  ZPS, August 28, 1968, SWB/FE/B/9-10; ZPS, August 31, 1968,
  SWB/FE/2865/B/36; ZPS, September 4, 1968, SWB/FE/2870/B/3-4; ZPS,
  September 6, SWB/FE/2883/B/28; ZPS, September 28, 1968,
  SWB/FE/2889/B/24-5; ZJRB, October 11, 1968, SCMP (S), 239, pp.~14-16;
  ZPS, October 15, 1968, SWB/FE/2908/B1/10; ZPS, November 4, 1968,
  SWB/FE/2926/B1/3; ZPS, November 18, 1968, SWB/FE/2930/B/3; ZPS,
  November 27, 1968, SWB/FE/2949/B1/10; ZPS, December 1, 1968,
  SWB/FE/B1/2-3; ZPS, December 5, 1968, SWB/FE/B1/7; ZJRB, December 6,
  1968, p.~1; ZPS, December 8, 1968, ibid., B1/6-7; editorial of the
  Southern Zhejiang Public Press (?), December 16, 1968, in Facts \&
  Features, 2: 8 (February 5, 1969), pp.~26-30; ZPS, December 17, 1968,
  SWB/FE/2955/B/8-9; ZPS, July 14, 1969, URS, 56:9 (July 29, 1969),
  pp.~127-28; CNA, 731, p.~5; Qiu Honggen, HZRB, 26 November 1978; CNS,
  270, p.~10; CNS, 271 (May 22, 1969), pp.~2-4, 6, 9-11, 13-14; CNS,
  272, (May 29, 1969), pp.~5-7.}在上一年下半年，这座城市曾令军方极为头痛，也让冒险前往当地的联合总部成员遭受不少头破血流。《浙江日报》1968年7月下旬一篇社论透露：

\begin{quote}
A series of serious political incidents in the areas of south Zhejiang
at present indicate the frenzied, last-ditch struggle desperately
carried out by the handful of class enemies after an even sharper class
struggle.
\end{quote}

\begin{quote}
目前浙南地区发生的一系列严重政治事件，表明一小撮阶级敌人在阶级斗争更加尖锐之后，正在疯狂地进行垂死挣扎。
\end{quote}

On 18 July Mao Zedong issued personal instructions to the authorities in
Zhejiang to suppress the Wenzhou insurrection. Mao's directives were
hailed in a large rally held in Hangzhou on 25 July. Those responsible
for the trouble in Wenzhou were accused of opposing the PLA and striving
to undermine its unity. Calls were made to the proletarian
revolutionaries to back the PLA ``no matter what time and no matter what
circumstances''. The authorities promised lenient treatment to members
of fighting factions who desisted, while threatening punishment of the
ringleaders.

7月18日，毛泽东亲自指示浙江当局镇压温州叛乱。7月25日，杭州举行大型集会欢呼毛的指示。温州骚乱责任者被指控反对解放军，并企图破坏其团结。号召无产阶级革命派``无论何时何地''都要支持解放军。当局承诺，对停止行动的武斗派成员从宽处理，同时威胁要惩罚头目。

\begin{quote}
The bad elements of the two different groups {[}elaborated one
editorial{]} must be dealt with separately by the masses of the
respective groups. At the same time, the bad elements who have sneaked
into an organization must be strictly separated from the masses of the
same organization.
\end{quote}

\begin{quote}
两个不同组织中的坏分子［一篇社论阐释说］必须由各自组织的群众分别处理。同时，混入某一组织的坏分子必须同该组织的群众严格区分开来。
\end{quote}

No distinction was drawn between the rival organizations; sanctions
would be applied mercilessly to them both.

对敌对组织之间并未作出区分；制裁将毫不留情地同时施加于两者。

The ``counter-revolutionaries'' of Wenzhou had engaged in the same
disruptive activities which were described in the two central documents
concerning Guangxi and Shaanxi, referred to at the beginning of this
chapter. They had ignored all orders to stop. Reports made particular
mention of outsiders and rusticated educated youth ``now loafing in
cities'' and causing the disturbances. These outsiders were directed to
return to their own localities. As a supplement to the central notices
and Mao's directives, and as a display of its determination to deal
firmly with unruly behavior, the ZPRC issued its own notice regarding
the trouble on 5 August. An editorial of the same day admitted that the
instigators, described colorfully as ``black talons'', had ``performed
quite outstandingly''. They had spread the ``reactionary theory of many
centers'' and split the ``proletarian Headquarters''.

温州``反革命分子''从事的破坏活动，同本章开头提到的中央关于广西和陕西两份文件所描述的活动相同。他们无视所有停止命令。报道特别提到外来人员和``现在在城市游荡''的下乡知识青年制造动乱。这些外来人员被要求返回原籍。作为中央通知和毛的指示的补充，也为显示其坚决处理不守秩序行为的决心，省革委会于8月5日就骚乱发布了自己的通知。同日一篇社论承认，被形象称为``黑爪子''的煽动者``表现得相当突出''。他们散布``反动的多中心论''，并分裂``无产阶级司令部''。

Words turned to action in August. Early in the month the ZPRC despatched
Mao Zedong Thought propaganda teams to Wenzhou, despite holding little
hope for their success. The first teams were composed of 4,000 PLA
soldiers and ``proletarian revolutionaries'' from United
Headquarters.\footnote{Propaganda teams composed of ``proletarian
  revolutionaries'' had left Hangzhou for Ningbo, Shaoxing, Jiaxing and
  Jinhua on July 5. ZPS, July 6,1968, SWB/FE/2824/B/11. In November
  1968, propaganda teams moved into Linhai county in Taizhou District,
  south-east Zhejiang. They left on January 29,1969, after the
  reestablishment of the county revolutionary committee. Linhai xianzhi,
  p.~29.} Given the history of confrontation in Wenzhou, the inclusion
of members of outside mass organizations in the teams was provocative.
For its part United Headquarters viewed its mission as working

8月，言辞转化为行动。月初，省革委会派毛泽东思想宣传队前往温州，尽管对其成功希望不大。第一批队伍由4,000名解放军士兵和来自联合总部的``无产阶级革命派''组成。\footnote{Propaganda
  teams composed of ``proletarian revolutionaries'' had left Hangzhou
  for Ningbo, Shaoxing, Jiaxing and Jinhua on July 5. ZPS, July 6,1968,
  SWB/FE/2824/B/11. In November 1968, propaganda teams moved into Linhai
  county in Taizhou District, south-east Zhejiang. They left on January
  29,1969, after the reestablishment of the county revolutionary
  committee. Linhai xianzhi, p.~29.}鉴于温州的对抗历史，把外地群众组织成员纳入队伍具有挑衅性。联合总部方面则把自己的使命视为：

\begin{quote}
together with the local proletarian revolutionaries to eliminate the
renegades, blast the lid off the class struggle, propagate the
revolutionary fighting call of the Headquarters headed by Chairman Mao,
stop struggle by force and help local proletarian revolutionaries
realize the revolutionary great alliance.
\end{quote}

\begin{quote}
同当地无产阶级革命派一道消灭叛徒，揭开阶级斗争的盖子，宣传以毛主席为首的司令部的革命战斗号召，制止武斗，帮助当地无产阶级革命派实现革命大联合。
\end{quote}

United Headquarters' interference in a local disturbance clearly went
outside the guidelines laid down for both mass organizations by Zhou
Enlai in March 1968.

联合总部干预地方动乱，显然超出了周恩来1968年3月为两大群众组织划定的准则。

The first contingent was reinforced by a brigade of the
newly-established worker-peasant propaganda team from Hangzhou
Municipality, which departed the provincial capital on 25 August for
Wenzhou. The brigade also left determined to ``lift the lid off class
struggle in southern Zhejiang''. At Jinhua it detrained to motor
transport for the remainder of the journey to Wenzhou. On 26 August, as
the brigade passed through Qingtian county, home of Zhejiang's famous
soapstone, it was met by forces organized by Wenzhou's Rebel General HQ.
A major confrontation took place that afternoon before the brigade
finally reached its destination. The same evening, at the large rally
held in the city, representatives from the Rebel General HQ attended but
three of its leaders, Yao Guolin, Wu Zhixiang (吴之祥), and Zhang Weisen
(张维森), were publicly criticized and detained by the military.

第一批队伍随后得到杭州市新成立的工农宣传队一个大队增援，该大队于8月25日离开省会前往温州。大队出发时也决心``揭开浙南阶级斗争的盖子''。到金华后，他们下火车改乘汽车完成前往温州的余程。8月26日，大队经过以浙江著名青田石闻名的青田县时，遭遇温州造反总指挥部组织的力量。那天下午发生重大对峙，之后大队才最终抵达目的地。当晚在市内举行的大型集会上，造反总指挥部代表参加了会议，但其三名领导人姚国麟、吴之祥和张维森受到公开批判并被军方拘留。

At the end of August the occupying forces sent out teams to visit three
communes and fifty enterprises and units. Their task was to disseminate
the Central Committee notices of July. The propaganda teams did not dare
venture outside the city limits, however, because of the strength of
Rebel General HQ in the counties surrounding Wenzhou. But they had been
entrusted with propagating central and provincial directives and for
this they needed to contact local officials. An example of one ingenious
way around this dilemma was publicized in the press. One detachment went
to the Wenzhou Chemical Works, which sold ammonia to cadres who
purchased it on behalf of production teams and brigades. When they
arrived at the chemical works the local cadres were subjected to
lectures by members of the propaganda team. Only in mid-October did
propaganda team detachments venture into district towns on the outskirts
of Wenzhou.

8月底，占领力量派出小组走访三个公社和五十个企业与单位。他们的任务是传达中央委员会7月通知。然而，由于造反总指挥部在温州周边各县势力强大，宣传队不敢冒险出城。但他们受托宣传中央和省级指示，为此需要接触地方干部。报刊宣传了一个巧妙绕开这一困境的例子。一支分队前往温州化工厂，该厂向代表生产队和大队采购氨水的干部出售氨水。当地干部到达化工厂时，便接受宣传队成员的训话。直到10月中旬，宣传队分队才敢进入温州郊区的区镇。

In early September 1,700 reinforcements arrived in Wenzhou by ship, the
port finally having been reopened.\footnote{The BBC-monitored radio
  broadcast mistranslated Wenzhou Water Police (温州水警区) as the
  Wenzhou Pacification Area (绥靖区). Pacification area was a term used
  in the Republican period. China News Summary picked up the term and
  claimed that the area had come under martial law. See CNS, 271, p.~11.}
By the end of September, the local party leader Wang Fang and with his
deputy Li Tiefeng (李铁锋) were named in an editorial published by
Zhejiang Daily. They were described as the agents of Liu Shaoqi in
Wenzhou who, together with the leaders of Rebel General HQ, had carried
out ``ferocious class revenge with savage hatred on the proletarian
revolutionaries''. The editorial claimed that their objective had been
to prevent the establishment of a revolutionary committee in the
district and municipality. In this endeavor Wang and his supporters
achieved notable success. Revolutionary committees in Wenzhou were not
established until early December and factional differences delayed the
formation of the last county revolutionary committee in Wenzhou District
until July 1969. The propaganda teams' strategy of attacking the leaders
of Rebel General HQ, but not the organization itself, was based on a
realistic appreciation of the support that it had built up in the area.
Wang Fang and his colleagues had capitalized on the strength of Rebel
General HQ and the parochialism of the citizens of Wenzhou to muster
stern opposition to outside interference.

9月初，港口终于重新开放，1,700名增援人员乘船抵达温州。\footnote{The
  BBC-monitored radio broadcast mistranslated Wenzhou Water Police
  (温州水警区) as the Wenzhou Pacification Area (绥靖区). Pacification
  area was a term used in the Republican period. China News Summary
  picked up the term and claimed that the area had come under martial
  law. See CNS, 271, p.~11.}到9月底，《浙江日报》一篇社论点名地方党领导人王芳及其副手李铁锋。他们被描述为刘少奇在温州的代理人，并同造反总指挥部领导人一道，对无产阶级革命派进行了``怀着刻骨仇恨的疯狂阶级报复''。社论声称，他们的目标是阻止在地区和市里建立革命委员会。在这一努力中，王及其支持者取得了显著成功。温州革命委员会直到12月初才成立，而派性分歧又把温州地区最后一个县级革命委员会的成立拖延到1969年7月。宣传队攻击造反总指挥部领导人而不攻击组织本身的策略，是基于对该组织在当地所建立支持力量的现实估计。王芳及其同僚利用造反总指挥部的力量和温州市民的地方主义，集结起对外来干预的坚决反对。

A series of rallies, parades and meetings was held in Wenzhou in the
months leading up to the formation of the district and municipal
revolutionary committees. By November Wang Fang had also been linked
publicly with Tan Zhenlin and Jiang Hua. It was reported in December
that 31,900 people had been enrolled in study groups in the city. The
provincial leadership demonstrated the importance it attached to the
reestablishment of order in Wenzhou by sending a high-powered delegation
consisting of Nan Ping, Lai Keke, Zhou Jianren and Zhang Yongsheng, to
attend the 3 December inauguration ceremony for the revolutionary
committees.

在地区和市革命委员会成立前的数月里，温州举行了一系列集会、游行和会议。到11月，王芳又被公开同谭震林和江华联系起来。12月有报道称，市内已有31,900人参加学习班。省领导派出由南萍、赖可可、周建人和张永生组成的高规格代表团，出席12月3日革命委员会成立仪式，显示其对温州恢复秩序的重视。

By December, the provincial authorities felt that they had the situation
under control. Early that month the Hangzhou brigade returned to the
provincial capital, leaving behind propaganda team contingents from
Shaoxing and Jinhua. An article published at this time related the
difficulties that the teams had confronted when they had arrived in
Wenzhou three months before. The ``class enemies''

到12月，省当局认为局势已在控制之中。月初，杭州大队返回省会，留下绍兴和金华的宣传队分队。此时发表的一篇文章叙述了宣传队三个月前抵达温州时所遭遇的困难。``阶级敌人''

\begin{quote}
were filled with deadly fear and hate when the team came to southern
Zhejiang Province. They framed the team with all the crimes that they
could dream up and cursed the team in foulest language
{[}sic{]}.\footnote{Given the incomprehensibility of the Wenzhou dialect
  even to the natives of the surrounding districts of Zhejiang, the
  propaganda teams would have needed to rely on interpreters to
  understand the abuse!} They resorted to despicable means to disrupt
the relationship between the masses and propaganda team in an attempt to
transform southern Zhejiang into their independent kingdom.
\end{quote}

\begin{quote}
在宣传队来到浙南时充满了致命的恐惧和仇恨。他们把所能想象到的一切罪名都栽到宣传队身上，并用最下流的语言咒骂宣传队［原文如此］。\footnote{Given
  the incomprehensibility of the Wenzhou dialect even to the natives of
  the surrounding districts of Zhejiang, the propaganda teams would have
  needed to rely on interpreters to understand the abuse!}他们采取卑鄙手段破坏群众同宣传队的关系，企图把浙南变成他们的独立王国。
\end{quote}

Even by the end of 1968, however, efforts to indoctrinate the people in
Wenzhou with the ``correct'' history of the Cultural Revolution had met
with little success and persistent opposition to the forced
establishment of revolutionary committees remained a thorn in the side
of the provincial authorities.

然而，即使到1968年底，用文化大革命``正确''历史灌输温州人民的努力也收效甚微，对强行建立革命委员会的持续反对仍是省当局的一根刺。

On August 25 the CCP CC had issued a notice informing the provinces to
send propaganda teams into educational institutions.\footnote{Liu
  Suinian and Wu Qungan (eds.), ``文化大革命''时期的国民经济
  (1966-1976), (The National Economy in The ``Great Cultural
  Revolution'', 1966-1976), (Haerbin: Heilongjiang Renmin Chubanshe,
  1986), p.~152.} Then, on 26 September, People's Daily carried Yao
Wenyuan's article ``The Working Class must lead in every sphere''
reminding Red Guards of their subordinate and dispensable status in the
Cultural Revolution.\footnote{Hao Mengbi and Duan Haoran, Liushinian,
  p.~599.} On 7 and 9 September 1968 propaganda teams moved into key
tertiary institutions in Hangzhou --- Zhejiang University, Hangzhou
University, Zhejiang Agricultural College and Zhejiang Medical
College.\footnote{ZPS, August 25, 1968, SWB/FE/2862/B/10-11; ZPS,
  September 7 and 9, 1968, SWB/FE/2883/B/28; CNS, 271, pp.~11-12.} On 7
October, teams entered other colleges and government
departments.\footnote{ZPS, October 7,1968, SWB/FE/2904/B/11. In Kaihua
  county propaganda teams occupied party and government organs and
  schools from September 2. Kaihua xianzhi, p.~24. In Lanxi county
  propaganda teams commenced moving into secondary schools on October 7.
  Lanxi shizhi, p.~20.} However, a substantial proportion of the
students did not welcome the occupation. They argued that it was an
insult to the role played by the Red Guards in the Cultural Revolution
as well as interference in their internal affairs. Before the propaganda
teams had entered the colleges, Red Guards had threatened to interrogate
their proletarian supervisors and generally make life uncomfortable for
them.\footnote{ZPS, August 30,1968, SWB/FE/2869/B/36-38; CNS, 271,
  pp.~9-11. Editorials of August 16 and 18, warned Red Guards to step
  back in line. CNS, 271, pp.~8-9.} At Hangzhou University the
propaganda team consisted of 350 workers. One member, aged 57, was
illiterate and had never spent a day at school in his life. It is not
surprising that the students resented such people running their
colleges. The authorities, however, remained steadfast. They declared
that ``Resistance to the team is resistance to the headquarters of
Chairman Mao''.\footnote{ZPS, September 20,1968, cited in CNA, 731,
  p.~5.}

8月25日，中共中央发布通知，要求各省派宣传队进入教育机构。\footnote{Liu
  Suinian and Wu Qungan (eds.), ``文化大革命''时期的国民经济
  (1966-1976), (The National Economy in The ``Great Cultural
  Revolution'', 1966-1976), (Haerbin: Heilongjiang Renmin Chubanshe,
  1986), p.~152.}随后，9月26日，《人民日报》刊登姚文元文章《工人阶级必须领导一切》，提醒红卫兵在文化大革命中处于从属且可被替代的地位。\footnote{Hao
  Mengbi and Duan Haoran, Liushinian, p.~599.}1968年9月7日和9日，宣传队进驻杭州重点高等院校：浙江大学、杭州大学、浙江农业大学和浙江医科大学。\footnote{ZPS,
  August 25, 1968, SWB/FE/2862/B/10-11; ZPS, September 7 and 9, 1968,
  SWB/FE/2883/B/28; CNS, 271, pp.~11-12.}10月7日，宣传队进入其他高校和政府部门。\footnote{ZPS,
  October 7,1968, SWB/FE/2904/B/11. In Kaihua county propaganda teams
  occupied party and government organs and schools from September 2.
  Kaihua xianzhi, p.~24. In Lanxi county propaganda teams commenced
  moving into secondary schools on October 7. Lanxi shizhi, p.~20.}然而，相当一部分学生并不欢迎这种占领。他们认为这是对红卫兵在文化大革命中所起作用的侮辱，也是对其内部事务的干涉。宣传队进入高校前，红卫兵曾威胁要审问这些无产阶级监督者，并普遍让他们日子不好过。\footnote{ZPS,
  August 30,1968, SWB/FE/2869/B/36-38; CNS, 271, pp.~9-11. Editorials of
  August 16 and 18, warned Red Guards to step back in line. CNS, 271,
  pp.~8-9.}杭州大学宣传队由350名工人组成。其中一名成员57岁，文盲，一生从未上过一天学。学生们反感由这样的人管理高校，并不令人意外。然而，当局态度坚定。他们宣布，``反对宣传队就是反对毛主席司令部''。\footnote{ZPS,
  September 20,1968, cited in CNA, 731, p.~5.}

The Red Guards were slow to realize that their usefulness as a
battering-ram to knock down the gates of Jiang Hua's provincial
stronghold had ended. They had indulged in anarchy and factionalism when
Beijing had demanded restoration of order. A conference which was held
in October 1968 to discuss the work of the propaganda teams affirmed the
contributions of the Red Guards in the early and middle stages of
Cultural Revolution, but pointedly made no reference to their role in
its later stages.\footnote{ZPS, October 19,1968, SWB/FE/2910/B/5-6.} At
a second conference in December, continuing resistance to the presence
of propaganda teams on campuses was reported.\footnote{ZPS, December
  28,1968, SWB/FE/2966/B/6.} Zhang Yongsheng, now a member of and
spokesman for the leadership hierarchy, delivered the report to the
conference. The rebels of yesterday had become the new proponents and
upholders of order.

红卫兵迟迟没有意识到，他们作为攻破江华省级堡垒大门的攻城槌，其作用已经结束。当北京要求恢复秩序时，他们却沉溺于无政府状态和派性。1968年10月召开的一次讨论宣传队工作的会议，肯定了红卫兵在文化大革命前期和中期的贡献，但有意不提他们在后期的作用。\footnote{ZPS,
  October 19,1968, SWB/FE/2910/B/5-6.}12月第二次会议报告说，校园内对宣传队存在的抵抗仍在继续。\footnote{ZPS,
  December 28,1968, SWB/FE/2966/B/6.}张永生此时已是领导层成员和发言人，由他向会议作报告。昨日的造反派已经成为秩序的新倡导者和维护者。

By the end of September 1968 the forces of law and order seemed to have
crushed most resistance to their rule. Criticism of Jiang Hua resumed
and the general situation was declared stable.\footnote{ZPS, September
  27, 1968, SWB/FE/2887/B/7; ZPS, September 30, 1968,
  SWB/FE/2895/B/15-17.} But the cost of obtaining stability had been
great. The initial victims of the Cultural Revolution, the
intellectuals, now faced not the spontaneous violence of their students
but the cold-blooded persecution of military-dominated revolutionary
committees. Disgraced cadres, especially those who were not party
members, had their files scrutinized for material that could be
construed as evidence of pre-1949 collaboration with the KMT. Rebels and
Red Guards, persecutors in 1966-67, were now mercilessly tracked down
and dealt with as ``May 16'' elements. Institutionalized savagery
reigned.

到1968年9月底，法律和秩序力量似乎已粉碎了对其统治的大部分抵抗。对江华的批判重新开始，总体形势被宣布为稳定。\footnote{ZPS,
  September 27, 1968, SWB/FE/2887/B/7; ZPS, September 30, 1968,
  SWB/FE/2895/B/15-17.}但取得稳定的代价巨大。文化大革命最初的受害者，即知识分子，现在面对的不再是学生的自发暴力，而是由军队主导的革命委员会的冷酷迫害。失势干部，尤其是非党员干部，其档案被审查，以寻找可被解释为1949年前同国民党合作证据的材料。1966至1967年的迫害者造反派和红卫兵，如今作为``五一六''分子被无情追查和处理。制度化的残暴占据了上风。

\section{Termination of the Cultural Revolution, October-December
1968}\label{termination-of-the-cultural-revolution-october-december-1968}

\section{文化大革命的终止，1968年10月至12月}\label{ux6587ux5316ux5927ux9769ux547dux7684ux7ec8ux6b621968ux5e7410ux6708ux81f312ux6708}

The victors of the Cultural Revolution pronounced their verdict on the
political struggles of the previous two years at the 12th plenum of the
8th CCP CC, held between 13 and 31 October 1968. What concerned Zhejiang
most directly was the plenum's judgment of the issues and events which
had contributed most to factionalism in the province. Above all this
involved an assessment of the 1967 ``February Countercurrent''. The
communique of the meeting declared that the ``current'' had been

文化大革命的胜利者在1968年10月13日至31日召开的中共八届十二中全会上，对此前两年的政治斗争作出判决。与浙江最直接相关的，是全会对那些最有助于省内派性形成的问题和事件的判断。其中首要的是对1967年``二月逆流''的评价。会议公报宣布，这股``逆流''

\begin{quote}
directed against the decisions of the 11th Plenary Session of the Eighth
Central Committee, against the great proletarian cultural revolution and
against the proletarian headquarters. The Plenary Session holds that the
shattering of the ``adverse February current'' and the sinister trend
last spring {[}1968{]} to reverse the correct verdict on the ``adverse
February current'' was an important victory for Chairman Mao's
proletarian revolutionary line.\footnote{Important Documents on the
  Great Proletarian Cultural Revolution in China (Beijing: Foreign
  Languages Press, 1970), p.186.}
\end{quote}

\begin{quote}
是反对八届十一中全会决定、反对无产阶级文化大革命、反对无产阶级司令部的。全会认为，粉碎``二月逆流''和1968年春为``二月逆流''翻案的邪风，是毛主席无产阶级革命路线的重要胜利。\footnote{Important
  Documents on the Great Proletarian Cultural Revolution in China
  (Beijing: Foreign Languages Press, 1970), p.186.}
\end{quote}

According to Chinese sources Mao and Lin Biao held slightly different
attitudes towards the ``February Countercurrent''. Nie Rongzhen, one of
the PLA marshalls involved in the events, has quoted excerpts from Lin
Biao's 26 October speech to the plenum and contrasted his position with
that of Mao's, which the Chairman put forward in his 31 October speech
to the plenum's closing session.\footnote{See Nie's article in Suo
  Guoxin, 1967 nian de 78 tian, pp.~5-7, reprinted from Guangming Ribao,
  December 1, 1984.} Lin Biao described the ``current'' as a ``serious
anti-party incident'', a ``preview for the restoration of capitalism''
and a continuation of the Liu-Deng line. Nie reported Mao as expressing
puzzlement over the whole affair, which is surprising considering that
Mao was reported by one source to have been present at one of the
meetings in February 1967. However, whether the Chairman was present or
not, he certainly had been fully briefed by Zhou Enlai and Zhang
Chunqiao on the incident.\footnote{See Suo Guoxin, 1967 nian de 78 tian,
  pp.~95-103. No other source supports Suo's assertion.} Now Mao stated
that as members of leading bodies those implicated had held the right to
express their views, that they had done so openly and that the whole
dispute was merely an example of inner-party debate.\footnote{See Hao
  Mengbi and Duan Haoran, Liushinian, pp.~601-02.}

据中国资料，毛和林彪对``二月逆流''的态度略有不同。聂荣臻是卷入事件的解放军元帅之一，他引用了林彪10月26日在全会上的讲话片段，并把林的立场同毛在10月31日全会闭幕会讲话中提出的立场加以对比。\footnote{See
  Nie's article in Suo Guoxin, 1967 nian de 78 tian, pp.~5-7, reprinted
  from Guangming Ribao, December 1, 1984.}林彪把这股``逆流''称为``严重反党事件''、``资本主义复辟的预演''，以及刘邓路线的继续。聂称毛对整个事件表示困惑；这令人惊讶，因为有一项资料说，毛曾出席1967年2月的一次会议。不过，无论主席是否在场，他肯定已由周恩来和张春桥就该事件作了充分汇报。\footnote{See
  Suo Guoxin, 1967 nian de 78 tian, pp.~95-103. No other source supports
  Suo's assertion.}如今毛表示，涉事者作为领导机关成员有权表达意见，他们是公开表达的，整场争论不过是党内辩论的一个例子。\footnote{See
  Hao Mengbi and Duan Haoran, Liushinian, pp.~601-02.}

Mao's forgiving attitude may have stemmed from well-developed political
instincts which foresaw the support of the veteran military leaders as a
useful counter-weight to Lin Biao and his military allies. The fate of
Deng Xiaoping was another instance where Mao refused to go along with
the wishes of his closest supporters. Lin Biao and civilian radicals
such as Jiang Qing reportedly wished to expel Deng from the CCP at this
plenum. The Chairman refused.\footnote{Gao Gao and Yan Jiaqi, ``Wenhua
  dageming'', p.~529.} If he had relented, the rehabilitation of Deng
four years later would have been much more difficult, if not impossible,
to accomplish.

毛的宽容态度可能源于成熟的政治本能，即预见到老资格军事领导人的支持可成为制衡林彪及其军事盟友的有用力量。邓小平的命运是毛拒绝顺从其最亲密支持者意愿的另一个例子。据说林彪和江青等文职激进派希望在这次全会上把邓开除出党。主席拒绝了。\footnote{Gao
  Gao and Yan Jiaqi, ``Wenhua dageming'', p.~529.}如果他让步，四年后邓的复出即使不是不可能，也会困难得多。

In the month preceding the convocation of the plenum Beijing had
announced that the Cultural Revolution was entering the phase of
``struggle, criticism and transformation''.\footnote{Hao Mengbi and Duan
  Haoran, Liushinian, p.~600.} Six factories and two universities in
Beijing, as well as PLA unit 8341 led by the head of Mao's bodyguard
Wang Dongxing, were chosen as model units in the campaign. The plenum
set four tasks for the implementation of ``struggle, criticism and
transformation'': revolutionary mass criticism and repudiation,
purification of class ranks, consolidation and rebuilding of the CCP and
revolution in the superstructure (principally the discipline of students
by worker propaganda teams on campuses).\footnote{Important Documents,
  pp.~188-92.} The campaign was not called to a halt until the death of
Lin Biao in September 1971.

全会召开前一个月，北京宣布文化大革命进入``斗、批、改''阶段。\footnote{Hao
  Mengbi and Duan Haoran, Liushinian, p.~600.}北京六个工厂、两所大学，以及由毛的卫队负责人汪东兴领导的解放军8341部队，被选为运动典型单位。全会为实行``斗、批、改''规定了四项任务：革命大批判、清理阶级队伍、巩固和重建中共，以及上层建筑领域的革命（主要是校园内工人宣传队对学生的管束）。\footnote{Important
  Documents, pp.~188-92.}这场运动直到1971年9月林彪死亡才停止。

Revolutionary mass criticism and repudiation were terms for the
wholesale assault on ideas regarding the dying out of class struggle,
party building and human nature attributed to Liu Shaoqi. The
purification of the class ranks was an intense, violent political
movement which had commenced in May 1968 in Shanghai.\footnote{See
  Fenwick, ``The Gang of Four'', p.~130.} It was directed principally at
traditional class and political enemies known as the five black
categories (黑五类) -- landlords, rich peasants,
counter-revolutionaries, bad elements and rightists -- as well as spies,
traitors and capitalist-roaders. It permitted a vast, indiscriminate and
arbitrary sweep through the ranks of disgraced cadres, intellectuals,
Red Guards and other rebels. ``Leftist'' opponents of the regime were
struck down as members of the May 16 group. A Central Committee notice
of 27 March 1970 admitted that up to one in seven people in certain
units had been categorized as May 16 elements. Revolution in the
superstructure encompassed not only the selection of new methods for
enrolling university students, but extended to the forced rustication of
city students and youth to the countryside and the establishment of May
7 cadre schools for the rotational training of office cadres.\footnote{Hao
  Mengbi and Duan Haoran, Liushinian, pp.~607-10; Gao Gao and Yan Jiaqi,
  ``Wenhua dageming'', pp.~288-305; Liu Guokai, ``A Brief Analysis'',
  pp.~118-29; Fact \& Features, 2:10 (March 5, 1969), pp.~9-11. A county
  May 7 cadre school was established in Kaihua county on July 6, 1969.
  See Kaihua xianzhi, p.~24.}

革命大批判是对归咎于刘少奇的关于阶级斗争熄灭、党的建设和人性的思想进行全面攻击的说法。清理阶级队伍是1968年5月在上海开始的一场激烈而暴力的政治运动。\footnote{See
  Fenwick, ``The Gang of Four'', p.~130.}它主要针对被称为黑五类的传统阶级和政治敌人，即地主、富农、反革命分子、坏分子和右派，以及特务、叛徒和走资派。它允许对失势干部、知识分子、红卫兵和其他造反派进行广泛、不加区分且任意的清查。政权的``左派''反对者被作为五一六分子打倒。中央委员会1970年3月27日一份通知承认，某些单位多达七分之一的人被定为五一六分子。上层建筑革命不仅包括选择新的大学招生办法，还扩展到强迫城市学生和青年下乡，以及建立五七干校，对机关干部进行轮训。\footnote{Hao
  Mengbi and Duan Haoran, Liushinian, pp.~607-10; Gao Gao and Yan Jiaqi,
  ``Wenhua dageming'', pp.~288-305; Liu Guokai, ``A Brief Analysis'',
  pp.~118-29; Fact \& Features, 2:10 (March 5, 1969), pp.~9-11. A county
  May 7 cadre school was established in Kaihua county on July 6, 1969.
  See Kaihua xianzhi, p.~24.}

The decisions of the CC 12th plenum were discussed at a marathon
forty-four day congress of ZPRC party members
(中共浙江省革命委员会党员代表大会) held from 8 November until 21
December 1968. The congress also selected delegates for the forthcoming
national party congress. Apart from summarizing the significance of the
12th plenum in Beijing, the five-point resolution published at the
conclusion of the meeting reviewed the struggles of the previous two
years. It also repudiated rightist attacks on the Cultural Revolution,
set out tasks for the immediate future and issued a rallying cry to win
total victory.\footnote{ZPS, December 21, 1968, SWB/FE/2960/B/8-11; CNS,
  250 (December 12, 1968), pp.~34-46; CNS, 272, pp.~9-10.} In summing up
the history of the Cultural Revolution in Zhejiang, resolution two
praised the ``general orientation'' of the proletarian revolutionaries
{[}United Headquarters{]} with the ``workers as the main force'' and
forbade criticism of their performance ``under any pretext''. The
resolution also commended the PLA for its achievements in the ``three
supports and two militaries''.

中央十二中全会的决定，在1968年11月8日至12月21日召开的省革委会党员代表大会（中共浙江省革命委员会党员代表大会）上进行了长达44天的马拉松式讨论。大会还为即将召开的全国党代会选出代表。会议结束时公布的五点决议，除了总结北京十二中全会的意义外，还回顾了此前两年的斗争。决议还驳斥了右派对文化大革命的攻击，规定了近期任务，并发出夺取全面胜利的号召。\footnote{ZPS,
  December 21, 1968, SWB/FE/2960/B/8-11; CNS, 250 (December 12, 1968),
  pp.~34-46; CNS, 272, pp.~9-10.}在总结浙江文化大革命历史时，第二项决议赞扬以``工人为主力''的无产阶级革命派［联合总部］的``总方向''，并禁止``以任何借口''批评其表现。决议还称赞解放军在``三支两军''中的成就。

The third resolution denounced three erroneous trends which had arisen
during the course of the Cultural Revolution -- the February 1967
Countercurrent, the spring 1968 rightist ``current'', and the May 1968
opposition to the ``three reds'' (that is, opposition to Mao, the PLA
and revolutionary committees). The Congress concluded that some members
of the ZPRC, including members of its standing committee, had been
involved in plotting the February 1967 ``reversal of verdicts'' and had
later tried to cover up their involvement in the conspiracy. These
people included former secretary of the CCP ZPC Wu Xian,\footnote{That
  this resolution was directed chiefly at Wu Xian was alleged by
  ``中共浙江省委的破坏与重建'' (The destruction and rebuilding of the
  CCP Zhejiang Provincial Committee), 中共年报 1971 (Taibei: Institute
  for the Study of Chinese Communist Problems, 1971), 3:19.} who had
been prominent a year previously when the military authorities were
searching for members of the old ZPC to join the ranks of
``revolutionary leading cadres''. They were expelled from the ZPRC but
it was noted that the struggle against these tendencies was not
finished.

第三项决议谴责文化大革命过程中出现的三股错误潮流：1967年2月逆流、1968年春右倾``逆流''，以及1968年5月反对``三红''（即反对毛、解放军和革命委员会）。大会断定，省革委会一些成员，包括其常委会成员，参与策划了1967年2月的``翻案''，后来又试图掩盖其参与阴谋。这些人包括前中共浙江省委书记吴宪，\footnote{That
  this resolution was directed chiefly at Wu Xian was alleged by
  ``中共浙江省委的破坏与重建'' (The destruction and rebuilding of the
  CCP Zhejiang Provincial Committee), 中共年报 1971 (Taibei: Institute
  for the Study of Chinese Communist Problems, 1971), 3:19.}一年前军方寻找旧省委成员加入``革命领导干部''队伍时，他曾颇为突出。他们被开除出省革委会，但会议指出，反对这些倾向的斗争尚未结束。

While it may have been true that the military authorities in most
provinces faced greater dangers from discontented ``leftists'' who
represented the spirit of the Cultural Revolution and were numerous and
young,\footnote{CNS, 250, p.~34.} this scenario did not apply in
Zhejiang, where Red Storm continued to oppose the new provincial
authorities with great energy. The ZPRC leadership had also failed to
attract officials of the old ZPC on board. The resolution of the plenum
reflected this state of affairs by seeming to defend United Headquarters
and denounce destabilization from the ``right''. A Hong Kong observer
noted that the military leaders of Zhejiang attacked the former
power-holders much more assiduously than they attacked ``Leftists''.
``The situation in Chekiang'', concluded the writer, ``is apparently
different from other provinces''.\footnote{CNS, 272, p.~12.} The reason
for this was that military leaders such as Nan Ping, Chen Liyun and
Xiong Yingtang all owed their positions in the Zhejiang administration
to the Cultural Revolution and their patron Lin Biao. They had more to
fear from the return of the pre-Cultural Revolution
``capitalist-roaders'' and Red Storm than their youthful allies in
United Headquarters, no matter how much of an irritant the latter
occasionally proved to be.

在多数省份，军方当局也许确实面对来自不满``左派''的更大危险；这些人代表文化大革命精神，人数众多且年轻。\footnote{CNS,
  250, p.~34.}但这种情形并不适用于浙江，因为红暴继续精力充沛地反对新的省级当局。省革委会领导层也未能吸引旧省委官员加入。全会决议反映了这种状态，似乎是在维护联合总部并谴责来自``右''的破坏稳定行为。一位香港观察者指出，浙江军队领导人攻击前掌权者，远比攻击``左派''更为勤勉。作者总结说，``浙江的情况显然不同于其他省份''。\footnote{CNS,
  272, p.~12.}原因在于，南萍、陈励耘、熊应堂等军队领导人都把自己在浙江行政机构中的职位归功于文化大革命和其庇护人林彪。对他们来说，文化大革命前``走资派''的复出以及红暴，比他们在联合总部的年轻盟友更可怕，尽管后者有时也颇为恼人。

Earlier in 1968 the provincial press had complained about the
perfunctory nature of the purification of class ranks movement. In the
counties the commencement date for the campaign had varied from spring
1968 (Kaihua), April 1968 (Lanxi), July 1968 (Jiande) to February
(Linhai) and May 1969 (Xiangshan).\footnote{Kaihua xianzhi, p.~24; Lanxi
  shizhi, p.~20; Jiande xianzhi, p.~23; Linhai xianzhi, p.~29; Xiangshan
  xianzhi, p.~35.} Zhejiang Daily had published editorials and articles
which seemed to indicate that both recalcitrant cadres and unrepentant
rebels from Red Storm were the targets of the campaign.\footnote{See
  CNS, 270, p.~9; CNS, 271, pp.~13-14 - editorials of June 18 and
  October 5 and an article of October 28,1968; CNS, 247, pp.~4-6.}
Deputy-governor Ren Yili, who prior to liberation had allegedly betrayed
members of his party branch, and a certain Yu, who had crossed over to
the enemy in 1929 and become a special agent, were prominent victims of
the campaign in Zhejiang. Jiang Hua had reportedly known about these
senior cadres' past but had defended them by stating (in the most
unlikely words), ``giving oneself up is justified and reneging is no
crime''.\footnote{ZPS, July 11,1968, SWB/FE/2823/B/11-12.} In late 1968,
as the campaign was intensified, further efforts were made to induce Red
Storm to end its boycott of the new power structures. Further attacks
against Jiang Hua and his subordinates were launched.

1968年早些时候，省级报刊曾抱怨清理阶级队伍运动敷衍了事。在各县，运动开始时间不一，从1968年春（开化）、1968年4月（兰溪）、1968年7月（建德），到1969年2月（临海）和5月（象山）。\footnote{Kaihua
  xianzhi, p.~24; Lanxi shizhi, p.~20; Jiande xianzhi, p.~23; Linhai
  xianzhi, p.~29; Xiangshan xianzhi, p.~35.}《浙江日报》发表的社论和文章似乎表明，顽固干部和不悔改的红暴造反者都是运动目标。\footnote{See
  CNS, 270, p.~9; CNS, 271, pp.~13-14 - editorials of June 18 and
  October 5 and an article of October 28,1968; CNS, 247, pp.~4-6.}副省长任一力据称解放前曾出卖党支部成员；某余姓干部1929年投敌并成为特务；二人是浙江运动中的显著受害者。据报江华知道这些高级干部的过去，却为他们辩护，说出极不可能的话：``自首有理，叛变无罪''。\footnote{ZPS,
  July 11,1968, SWB/FE/2823/B/11-12.}1968年末，随着运动加剧，当局进一步努力诱使红暴结束对新权力结构的抵制。对江华及其下属的新一轮攻击也展开了。

On 5 November, at a large rally in Hangzhou, leading capitalist-roaders
of the CCP ZPC and HMC, including Li Fengping, Chen Weida and Wang
Pingyi, were led onto the stage to listen to a recitation of their
crimes.\footnote{ZPS, November 6,1968, SWB/FE/2926/B1/3.} In a report of
11 November 1968, Zhang Chunqiao had stated that it was necessary to
name the leaders implicated in the ``February Countercurrent''. His
report was printed and circulated by the CCP CC.\footnote{Hao Mengbi and
  Duan Haoran, Liushinian, p.~602.} Four days later the first
officially-endorsed critique of Jiang Hua's political record was
published in Zhejiang Daily.\footnote{``Down with Jiang Hua'', ZJRB,
  November 15,1968, p.~1.} United Headquarters, which had waited almost
two years for this event, may well have felt a sense of satisfaction
that its struggle had not been in vain. However, the number and
vehemence of the articles published against Jiang and the ``February
Countercurrent'' indicated that a lengthy propaganda war would be
necessary to sell the party's line to the people of Zhejiang.

11月5日，在杭州一次大型集会上，中共浙江省委和杭州市委的主要走资派，包括李丰平、陈伟达和王平夷，被押上台聆听对其罪行的宣读。\footnote{ZPS,
  November 6,1968, SWB/FE/2926/B1/3.}1968年11月11日的一份报告中，张春桥称有必要点名卷入``二月逆流''的领导人。他的报告由中共中央印发传达。\footnote{Hao
  Mengbi and Duan Haoran, Liushinian, p.~602.}四天后，《浙江日报》刊登了第一篇得到官方认可的对江华政治履历的批判。\footnote{``Down
  with Jiang Hua'', ZJRB, November 15,1968, p.~1.}等待这一事件将近两年的联合总部，很可能感到自己的斗争并非徒劳。然而，针对江华和``二月逆流''发表文章的数量和激烈程度表明，要把党的路线推销给浙江人民，需要一场漫长的宣传战。

A series of editorials, articles and rallies on this theme continued
into December 1968.\footnote{See ZJRB, November 17,1968, p.~2 for
  articles by Lai Keke and Red Storm denouncing Jiang; see also ZJRB,
  November 17,1968, pp.~1,4; ZJRB, November 18,1968, p.~1; ZPS, November
  18, 1968, SWB/FE/2930/B/3; ZPS, November 19, 1968, SWB/FE/2933/B/3-6;
  CNS, 272, pp.~1-5.} Jiang Hua was described as the focus of the
struggle in Zhejiang during the Cultural Revolution, and held personally
responsible for manipulating public opinion regarding his status over
the previous two years. The press disclosed that even the unambiguous
declaration of the 12th plenum concerning the Countercurrent had been
received with incomprehension and resisted in Zhejiang. The crux of the
issue, confessed Zhejiang Daily, was political power. Jiang's supporters
had opposed the establishment of the ZPRC and then had refused to accept
its authority. They had set up their own alternative political
structures.\footnote{The ZJRB editorial of July 16 had referred to
  ``underground revolutionary committees''.} According to the newspaper,
these facts proved that

围绕这一主题的一系列社论、文章和集会一直持续到1968年12月。\footnote{See
  ZJRB, November 17,1968, p.~2 for articles by Lai Keke and Red Storm
  denouncing Jiang; see also ZJRB, November 17,1968, pp.~1,4; ZJRB,
  November 18,1968, p.~1; ZPS, November 18, 1968, SWB/FE/2930/B/3; ZPS,
  November 19, 1968, SWB/FE/2933/B/3-6; CNS, 272, pp.~1-5.}江华被描述为浙江文化大革命斗争的焦点，并被认为要为此前两年操纵关于其地位的舆论负个人责任。报刊披露，即便十二中全会关于``逆流''的明确宣告，在浙江仍遭到不理解和抵制。《浙江日报》承认，问题的核心是政治权力。江的支持者反对建立省革委会，随后又拒绝接受其权威。他们建立了自己的替代性政治结构。\footnote{The
  ZJRB editorial of July 16 had referred to ``underground revolutionary
  committees''.}据该报说，这些事实证明：

\begin{quote}
Jiang Hua is still around, and his ambitions are not dead. Some of his
accomplices have only felt the ``threatening chill of the winter wind''
and have not yet been frozen stiff. The ghost of Jiang Hua still haunts
Zhejiang Province.
(江华人还在，心不死，他的某些同伙甚至还只感到``冬天的威胁''，没有冻僵呢;江华的幽灵还在浙江的土地上游荡着。)\footnote{ZJRB,
  November 15,1968, p.~1. Early in 1969 on Mao's instructions, Jiang Hua
  was sent to the No.~27 Vehicle Manufacturing Plant on the outskirts of
  Beijing, together with other central and provincial leading cadres
  such as Xu Xiangqian, Wang Enmao, Liao Zhigao and Jiang Weiqing, to be
  ``reeducated''. See Xu Xiangqian, Lishide huigu, Vol. 2, pp.~845-46.}
\end{quote}

\begin{quote}
江华人还在，野心未死。他的一些同伙只是感到``冬天的威胁''，还没有冻僵。江华的幽灵仍在浙江省游荡。（江华人还在，心不死，他的某些同伙甚至还只感到``冬天的威胁''，没有冻僵呢;江华的幽灵还在浙江的土地上游荡着。）\footnote{ZJRB,
  November 15,1968, p.~1. Early in 1969 on Mao's instructions, Jiang Hua
  was sent to the No.~27 Vehicle Manufacturing Plant on the outskirts of
  Beijing, together with other central and provincial leading cadres
  such as Xu Xiangqian, Wang Enmao, Liao Zhigao and Jiang Weiqing, to be
  ``reeducated''. See Xu Xiangqian, Lishide huigu, Vol. 2, pp.~845-46.}
\end{quote}

Another article warned that some people

另一篇文章警告说，有些人

\begin{quote}
are spreading rumors and attacking and resisting the Zhejiang Provincial
Committee, which has the personal approval of our great leader Chairman
Mao. Creating discord between armymen and citizens, they are
attempting\ldots{} to turn the areas and departments under their
leadership into watertight independent kingdoms.
\end{quote}

\begin{quote}
正在散布谣言，攻击并抵制经我们伟大领袖毛主席亲自批准的浙江省革命委员会。他们制造军民之间的隔阂，企图\ldots\ldots 把其领导下的地区和部门变成针插不进、水泼不进的独立王国。
\end{quote}

The old CCP ZPC cadre school in Hangzhou remained a bastion of support
for Jiang. An editorial of 10 November claimed that the
capitalist-roaders in the school had ``recently'' issued an appeal to
reverse the verdict handed down on Jiang and his colleagues. The
editorial appealed to ``comrades'' who had supported the wrong side in
the Cultural Revolution to recognize their mistakes and act on this
recognition, or face the consequences.\footnote{ZJRB, November 10,1968,
  p.~1; see also ZPS, October 29,1968, SWB/FE/2921/B1/1, and CNS, 271,
  pp.~14-15.}

杭州的旧中共浙江省委党校仍是支持江华的堡垒。11月10日一篇社论称，学校中的走资派``最近''发出呼吁，要求为江华及其同僚翻案。社论号召在文化大革命中支持错误一方的``同志''认识自己的错误，并据此行动，否则将承担后果。\footnote{ZJRB,
  November 10,1968, p.~1; see also ZPS, October 29,1968,
  SWB/FE/2921/B1/1, and CNS, 271, pp.~14-15.}

Throughout these months, Red Storm continued to resist all pressure to
join revolutionary committees. It rejected the terms of entry laid down
by the military and United Headquarters. On the first anniversary of
Mao's 2 December 1967 directive Zhejiang Daily published an editorial to
commemorate the occasion.\footnote{ZPS, December 2,1968,
  SWB/FE/2942/B/12-13; CNS, 272, pp.~6-7.} The editorial made it clear
that the struggle between the two mass organizations had never ceased.
Whenever a makeshift alliance had been knocked together at the top, the
focus of contention shifted back to each factory, college or commune. In
order to flatter Red Storm, the editorial stated that it too had played
its part in the downfall of Jiang Hua and that its representatives would
help strengthen the three mass congresses and revolutionary committees.
All the bickering about the status of Red Storm was seemingly resolved
and it was now described as a ``revolutionary mass organization''. Yet
United Headquarters could not accept its rival on equal terms. Turning
to the members of this organization, Zhejiang Daily praised them for
pursuing the policy of ``helping, criticizing and allying with'' Red
Storm. The newspaper ludicrously suggested that members of Red Storm
accept such condescending treatment and warned that ``In no circumstance
should they regard the help and criticism given by their comrades as
suppression and refuse to accept it''.

在这些月份里，红暴继续抵制所有要求其加入革命委员会的压力。它拒绝军方和联合总部规定的进入条件。毛1967年12月2日指示一周年时，《浙江日报》发表社论纪念这一事件。\footnote{ZPS,
  December 2,1968, SWB/FE/2942/B/12-13; CNS, 272, pp.~6-7.}社论清楚表明，两大群众组织之间的斗争从未停止。每当上层勉强拼凑出临时联盟，争执焦点便转回每个工厂、高校或公社。为讨好红暴，社论称它也在打倒江华中发挥了作用，其代表将帮助加强三大代表大会和革命委员会。围绕红暴地位的一切争吵似乎都得到解决，它现在被描述为``革命群众组织''。然而，联合总部无法以平等条件接受其对手。《浙江日报》转而面对联合总部成员，称赞他们执行对红暴``帮助、批评、联合''的方针。报纸荒谬地建议红暴成员接受这种居高临下的待遇，并警告说：``在任何情况下，都不应把同志们给予的帮助和批评看作压制而拒绝接受。''

Editorials of December 1968 revealed that the campaign to purify class
ranks was meeting strong opposition. At the 12th plenum Mao had
counselled against forced confessions and on 1 December he had directed
that the fewer cadres be attacked and the role of education
emphasized.\footnote{Hao Mengbi and Duan Haoran, Liushinian, p.~602.}
Veteran cadres and long-serving party members undoubtedly objected
strongly to the practice of youthful rebels examining their past records
and subjecting them to interrogation, an exercise which often
degenerated into abuse and torture, despite Mao's injunctions.
Nevertheless, for this latest phase of the Cultural Revolution these
cadres had attempted to seize the initiative. They had coined the
phrase, ``Revolutionaries can be relied on in revolution but
conservatives are to be relied on in purification''. The
``conservatives'' argued, with some justification, that the purification
of class ranks could have been carried out without the Cultural
Revolution. But such remarks were heresy. The ``conservatives'' also
resurrected the theory of lineage (血统论), which they had waved in the
face of class enemies on previous occasions, and described their
youthful adversaries as untrustworthy, unreliable and undeserving of
support.\footnote{ZPS, December 7, 1968, SWB/FE/2949/B1/8-9; ZPS,
  December 9, 1968, SWB/FE/2948/B/6-8; CNS, 272, pp.~7-9.}

1968年12月的社论显示，清理阶级队伍运动遇到强烈反对。十二中全会上，毛曾劝诫不要逼供；12月1日，他指示少打击干部并强调教育作用。\footnote{Hao
  Mengbi and Duan Haoran, Liushinian, p.~602.}老干部和老党员无疑强烈反对年轻造反派审查其历史档案并加以审问的做法；尽管有毛的告诫，这种做法常常堕入辱骂和酷刑。尽管如此，在文化大革命这一最新阶段，这些干部试图夺取主动权。他们创造了这样的说法：``革命靠革命派，清理靠保守派。''\,``保守派''有一定理由地辩称，清理阶级队伍可以不经过文化大革命而完成。但这类言论属于异端。``保守派''还复活了血统论，他们曾在以往场合用它来对付阶级敌人，并把年轻对手描述为不可信、不可靠、不值得支持。\footnote{ZPS,
  December 7, 1968, SWB/FE/2949/B1/8-9; ZPS, December 9, 1968,
  SWB/FE/2948/B/6-8; CNS, 272, pp.~7-9.}

However, apart from making scapegoats of several senior officials, and
purging Jiang Hua's remaining followers,\footnote{A further rally to
  denounce Xue Ju was held on December 6. ZPS, December 6,1968,
  SWB/FE/2949/B/9.} the authorities proceeded with caution.\footnote{See
  the editorials of ZJRB, ZPS, December 13,1968, SWB/FE/2957/B1/3-5;
  ZPS, December 15,1968, SWB/FE/2955/B/6-7.} In line with the Chairman's
instructions on the matter, they prohibited the use of duress to obtain
evidence and directed, ``never let a single bad element escape or harm a
good element by mistake''. Cadres were the key to the creation of the
three-way alliance, lectured the Zhejiang Daily, and asserted that even
capitalist-roaders, so long as they had not betrayed the party, should
be allowed to resume work. It is not surprising that some people
wondered why the Cultural Revolution had been necessary in the first
place.\footnote{See ZPS, December 25,1968, SWB/FE/2960/B/11-13.} In
January 1969, in a warning against restoration (复旧), Zhang Yongsheng
gave vent to feelings which were probably widespread among the
rebels.\footnote{Zhang Yongshengde zuizheng cailiao, p.~29.}

不过，除了把几名高级官员作为替罪羊并清洗江华残余追随者之外，\footnote{A
  further rally to denounce Xue Ju was held on December 6. ZPS, December
  6,1968, SWB/FE/2949/B/9.}当局行事较为谨慎。\footnote{See the
  editorials of ZJRB, ZPS, December 13,1968, SWB/FE/2957/B1/3-5; ZPS,
  December 15,1968, SWB/FE/2955/B/6-7.}按照主席在此问题上的指示，他们禁止用胁迫方式取得证据，并指示``一个坏人也不放过，一个好人也不冤枉''。《浙江日报》训示说，干部是建立三结合的关键，并断言即使是走资派，只要没有背叛党，也应允许恢复工作。有人怀疑文化大革命最初为何必要，并不令人意外。\footnote{See
  ZPS, December 25,1968, SWB/FE/2960/B/11-13.}1969年1月，张永生在警告防止复旧时，宣泄了造反派中大概相当普遍的情绪。\footnote{Zhang
  Yongshengde zuizheng cailiao, p.~29.}

1968 ended with the publication of Mao's instructions to educated youth
to go up to the mountains and down to the countryside
(上山下乡).\footnote{For a major study of this programme, see Thomas P.
  Bernstein, Up to the Mountains and Down to the Villages: The Transfer
  of Youth from Urban to Rural China (New Haven: Yale University Press,
  1977).} Young people joined production and construction brigades or
lived with peasants in production teams. Many city youth traveled and
settled down in such harsh, distant regions as Xinjiang, Inner Mongolia,
Manchuria and Yunnan. Others, with better contacts, went no further than
richer suburban communes surrounding major cities. The first contingent
from Zhejiang left for Heilongjiang on the day that Mao's directive was
published.\footnote{CNS, 272, p.~10. For China in general, see Gao Gao
  and Yan Jiaqi, ``Wenhua dageming'', pp.~316-18; Liu Suinian and Wu
  Qungan, ``Wenhua dageming'' shiqide guomin jingji, pp.~152-53.} The
rustication program had several unspoken aims: first to relieve urban
unemployment; second to show privileged urban youth the grim reality of
China's countryside; and third to clear the cities of rebel
trouble-makers, well-practised in the art of organization and
mobilization and independent in thought and action. Undoubtedly the
military authorities in Zhejiang ensured that the selection procedures
met these aims in full.

1968年以毛关于知识青年上山下乡的指示发表而告终。\footnote{For a major
  study of this programme, see Thomas P. Bernstein, Up to the Mountains
  and Down to the Villages: The Transfer of Youth from Urban to Rural
  China (New Haven: Yale University Press, 1977).}青年加入生产建设大队，或同生产队农民一起生活。许多城市青年远赴新疆、内蒙古、东北和云南等艰苦遥远地区定居。另一些有较好关系的人，则只去了大城市周边较富裕的郊区公社。毛的指示发表当天，浙江第一批人员动身前往黑龙江。\footnote{CNS,
  272, p.~10. For China in general, see Gao Gao and Yan Jiaqi, ``Wenhua
  dageming'', pp.~316-18; Liu Suinian and Wu Qungan, ``Wenhua dageming''
  shiqide guomin jingji, pp.~152-53.}下乡计划有若干未明说的目标：第一，缓解城市失业；第二，让享有特权的城市青年看到中国农村的严酷现实；第三，清除城市中擅长组织和动员、思想行动独立的造反闹事者。毫无疑问，浙江军方确保选拔程序完全符合这些目标。

\section{The ``Revolutionary Alliance'' and the Dissolution of the Mass
Organizations, May-September
1969}\label{the-revolutionary-alliance-and-the-dissolution-of-the-mass-organizations-may-september-1969}

\section{``革命大联合''与群众组织的解散，1969年5月至9月}\label{ux9769ux547dux5927ux8054ux5408ux4e0eux7fa4ux4f17ux7ec4ux7ec7ux7684ux89e3ux65631969ux5e745ux6708ux81f39ux6708}

In 1969 the military authorities who exercised control over the ZPRC
appeared to take a harder line against the disgruntled rebels. Their
new-found determination to grasp the nettle stemmed from a twenty-three
day meeting convened in Beijing in January by the central authorities to
deal with continuing factionalism in Zhejiang. The meeting discussed the
problems within the military caused by the ``two reorganizations'' of
August 1967 and endorsed the Lin Biao-backed leadership of the province.
Former leaders of the ZPMD were censured and persecuted.\footnote{HZRB,
  September 12, 1979, p.~1.} An order was issued to the two mass
organizations that their disruptive activities must come to an end and
that the time for negotiations had arrived. Red Storm's Fang Jianwen,
who was a principal opponent of any alliance, was criticized for his
stand.\footnote{ZJRB, August 13,1969.}

1969年，控制省革委会的军方当局似乎对心怀不满的造反派采取了更强硬路线。他们新近敢于直面棘手问题的决心，源于中央当局1月在北京召开的一次为期23天的会议，会议旨在处理浙江持续存在的派性。会议讨论了1967年8月``两次改组''在军内造成的问题，并认可林彪支持的省级领导层。浙江省军区前领导人遭到谴责和迫害。\footnote{HZRB,
  September 12, 1979, p.~1.}两大群众组织被命令停止破坏活动，谈判时刻已经到来。红暴方剑文是任何联合的主要反对者，他因这一立场受到批评。\footnote{ZJRB,
  August 13,1969.}

Unified leadership meant that the public airing of differences with
revolutionary committees was not permitted and such practices were
condemned as manifestations of ultra-democracy.\footnote{ZJRB, January
  14,1969, p.~1.} Reference was made to the attempt by certain people to
establish a second headquarters as an alternative to the ZPRC, and their
determination to foster division between the PLA and the revolutionary
committees, between different units of the PLA and between mass
organizations and revolutionary committees.\footnote{ZPS, February 18,
  1969, SWB/FE/3005/B/9-11; ZPS, February 20,1969, SWB/FE/3008/B/1-4.} A
major article carried by Zhejiang Daily entitled, ``The Bankruptcy of
Jiang Hua's `Independent Kingdom''', made pointed references to groups
which disobeyed orders and placed their sectional interests above those
of the whole.\footnote{ZJRB, January 13,1969, p.~1.}

一元化领导意味着不允许公开发表同革命委员会的分歧，这类做法被谴责为极端民主化的表现。\footnote{ZJRB,
  January 14,1969, p.~1.}当局提到某些人企图建立第二司令部，作为省革委会的替代，并决心在解放军和革命委员会之间、解放军各单位之间、群众组织和革命委员会之间制造分裂。\footnote{ZPS,
  February 18, 1969, SWB/FE/3005/B/9-11; ZPS, February 20,1969,
  SWB/FE/3008/B/1-4.}《浙江日报》发表题为《江华``独立王国''的破产》的重要文章，明确指向那些不服从命令、把本位利益置于整体利益之上的集团。\footnote{ZJRB,
  January 13,1969, p.~1.}

The status of mass representatives on revolutionary committees was made
painfully clear in an article concerning the unified leadership of the
Shaoxing County Revolutionary Committee. Mass representatives were
instructed to learn from the experience of their PLA and civilian cadre
colleagues. The latter two groups in turn deigned to ``think highly'' of
their junior partners' role on the committees.\footnote{ZPS, February
  13, 1969, SWB/FE/3007/B1/1-2.} At the third enlarged plenum of the
ZPRC, which was held from 25 February to 6 March, it was announced that
the attempt to reverse the verdict on the February Countercurrent in
Zhejiang had been defeated.\footnote{ZPS, March 10,1969,
  SWB/FE/3032/B/10-12.} An editorial celebrating the first anniversary
of the establishment of the ZPRC, which was published on 23 March 1969,
declared that one of the major struggles in Zhejiang since November 1968
had centered on the elimination of the influence of the February 1967
Countercurrent.\footnote{ZJRB, March 23,1969, pp.~1, 2.} Responsibility
for the attempt was later attributed to Red Storm's leader Fang
Jianwen.\footnote{ZJRB, August 13,1969.} So as to ram home the message
to Fang's supporters, on the final day of the plenum, Zhejiang Daily
carried a 25 February editorial from ``Red Storm'' entitled ``The
Victors face the greatest danger'' (胜利者最危险), warning of the
dangers of arrogance, self-satisfaction and individualism which
accompanied victory.\footnote{ZJRB, March 6,1969, p.~2.}

一篇关于绍兴县革命委员会一元化领导的文章，痛切地说明了群众代表在革命委员会中的地位。群众代表被要求向解放军和文职干部同事学习经验。后两类人则屈尊``重视''这些地位较低的伙伴在委员会中的作用。\footnote{ZPS,
  February 13, 1969, SWB/FE/3007/B1/1-2.}省革委会第三次扩大会议于2月25日至3月6日召开，会上宣布在浙江为二月逆流翻案的企图已经被击败。\footnote{ZPS,
  March 10,1969, SWB/FE/3032/B/10-12.}1969年3月23日发表的一篇庆祝省革委会成立一周年的社论宣布，自1968年11月以来，浙江的一项主要斗争集中于肃清1967年二月逆流的影响。\footnote{ZJRB,
  March 23,1969, pp.~1, 2.}这一企图的责任后来归咎于红暴领导人方剑文。\footnote{ZJRB,
  August 13,1969.}为让方的支持者彻底明白这一信息，全会最后一天，《浙江日报》刊登了2月25日``红暴''社论《胜利者最危险》，警告胜利伴随的骄傲、自满和个人主义危险。\footnote{ZJRB,
  March 6,1969, p.~2.}

The propaganda of a year earlier admitting to the existence of factions
and differentiating between proletarian and bourgeois factions had
ceased. In its place was a shrill exhortation against factions per se
and continued calls to uphold unified (一元化) leadership. The
communique of the plenum admitted that ``In many respects Chekiang lags
behind other provinces and municipalities''. It enjoined mass
organizations to refrain from ``exchanges of revolutionary experience''
and called on the PLA to support the revolutionary committees. On the
first anniversary of the establishment of the ZPRC Lai Keke, Xiong
Yingtang and Zhang Yongsheng, representing the tripartite alliance of
civilian cadres, the military and mass representatives, published
articles in Zhejiang Daily pledging to uphold the authority of the
ZPRC.\footnote{ZJRB, March 24, 1969, p.~3. See also the article by the
  Zhejiang Red Guard Congress on the same theme in ibid, March 25, 1969,
  p.~2.} A Zhejiang Daily editorial of late March pointed out the danger
of certain people placing revolutionary committees under their own
control or the control of their particular organizations.\footnote{ZJRB,
  March 30, 1969, p.~1.}

一年前承认派别存在并区分无产阶级派和资产阶级派的宣传已经停止。取而代之的是对派别本身的尖厉反对，以及继续号召维护一元化领导。全会公报承认，``浙江在许多方面落后于其他省市''。它命令群众组织停止``交流革命经验''，并号召解放军支持革命委员会。省革委会成立一周年时，赖可可、熊应堂和张永生分别代表文职干部、军队和群众代表三结合，在《浙江日报》发表文章，誓言维护省革委会权威。\footnote{ZJRB,
  March 24, 1969, p.~3. See also the article by the Zhejiang Red Guard
  Congress on the same theme in ibid, March 25, 1969, p.~2.}《浙江日报》3月下旬一篇社论指出，某些人把革命委员会置于自己或自己组织控制之下的危险。\footnote{ZJRB,
  March 30, 1969, p.~1.}

Yet despite the warnings and threats, renewed efforts in 1969 to root
out opposition to the ZPRC and finally bring about the long sought-after
alliance of the two mass organizations met disinterest or resistance.
The opposition was categorized as a leftist deviation and those
upholding it as ``regal isolationists''.\footnote{See ZJRB, April 10,
  1969, p.~1; ZJRB, April 20, 1969, p.~2. Fighting continued in the
  counties. On April 10, 1969, rebels from various counties in the
  Taizhou District of south-east Zhejiang, including five hundred from
  Xianju county, stormed Yuhuan county (玉环县) to participate in its
  ``armed liberation''. Xianju xianzhi, p.~23.} The belief that leftist
mistakes were less serious than those of the right, because they
involved mistakes of method not political line, was rejected.\footnote{ZPS,
  May 10, 1969, SWB/FE/3075/BII/5.} In a harsh warning to members of
United Headquarters and Red Storm who were holding out against the
alliance of the two organizations, Zhejiang Daily stated that they faced
the possibility of becoming revisionists.\footnote{ZPS, May 16, 1969,
  SWB/FE/3082/BII/4-5.}

然而，尽管有警告和威胁，1969年为根除对省革委会的反对并最终实现长期追求的两大群众组织联合而作出的新努力，仍遭遇冷淡或抵抗。反对被归类为左倾偏差，支持这种立场者被称为``山头主义者''。\footnote{See
  ZJRB, April 10, 1969, p.~1; ZJRB, April 20, 1969, p.~2. Fighting
  continued in the counties. On April 10, 1969, rebels from various
  counties in the Taizhou District of south-east Zhejiang, including
  five hundred from Xianju county, stormed Yuhuan county (玉环县) to
  participate in its ``armed liberation''. Xianju xianzhi, p.~23.}那种认为左的错误不如右的错误严重，因为它们只是方法错误而非政治路线错误的看法，被予以否定。\footnote{ZPS,
  May 10, 1969, SWB/FE/3075/BII/5.}《浙江日报》严厉警告那些仍反对两组织联合的联合总部和红暴成员，说他们面临变成修正主义者的可能。\footnote{ZPS,
  May 16, 1969, SWB/FE/3082/BII/4-5.}

It was the 9th Congress of the CCP, held from 1 to 24 April
1969,\footnote{Hao Mengbi and Duan Haoran, Liushinian, pp.~603-06. See
  the articles written by delegates from Zhejiang to the 9th Congress in
  ZJRB, May 20,1969, p.~3 and May 21,1969, p.~3.} that gave renewed
impetus to the question of unity. On 16 April ten thousand study classes
with 200,000 participants commenced in Hangzhou.\footnote{ZPS, April 22,
  1969, SWB/FE/3064/B1/2.} Party members of the ZPRC held a Congress
from 9 to 31 May to study the documents of the 9th Congress and to
attempt to resolve once and for all the outstanding issue of Red Storm's
participation in the provincial power-structure.\footnote{ZJRB, May 17,
  1969, pp.~1, 2; ZJRB, June 1, 1969, pp.~1, 2.} The provincial congress
took place in two stages. From 9 to 13 May, 1,000 delegates attended,
while from 14 May the number was increased to 7,000.

1969年4月1日至24日召开的中共九大，\footnote{Hao Mengbi and Duan Haoran,
  Liushinian, pp.~603-06. See the articles written by delegates from
  Zhejiang to the 9th Congress in ZJRB, May 20,1969, p.~3 and May
  21,1969, p.~3.}给团结问题带来新的推动。4月16日，杭州开办一万个学习班，参加者达20万人。\footnote{ZPS,
  April 22, 1969, SWB/FE/3064/B1/2.}省革委会党员于5月9日至31日召开大会，学习九大文件，并试图一劳永逸地解决红暴参与省级权力结构这一悬而未决的问题。\footnote{ZJRB,
  May 17, 1969, pp.~1, 2; ZJRB, June 1, 1969, pp.~1, 2.}省级大会分两个阶段进行。5月9日至13日有1,000名代表参加；5月14日起人数增至7,000。

The most important event occurred on 26 May when a meeting was held to
welcome Red Storm's decision to return to the fold.\footnote{ZJRB, May
  27, 1969, (afternoon edition), (下午版) pp.~1, 2.} Convened jointly by
the ZPRC and HMRC, the meeting announced the inclusion of Red Storm
representatives in the standing committees of the provincial worker,
peasant and Red Guard Congresses, the HMRC and its standing committee,
and a recommendation that representatives from Red Storm join the ZPRC.
The report of the meeting stated that the ZPRC and HMRC had discussed
the issue of Red Storm's representation many times but that differences
over interpretation of the contentious February 1967 events had held up
resolution of the problem. In 1969 the Central Committee had issued
instructions on the matter, which had led the Zhejiang and Hangzhou
revolutionary committees to hold study classes for members of Red Storm.
After the 9th CCP Congress, Red Storm's standing committee had studied
the question again and drawn up an acceptable namelist, based on the
original agreement of February.

最重要的事件发生在5月26日，当时召开会议欢迎红暴决定回归。\footnote{ZJRB,
  May 27, 1969, (afternoon edition), (下午版) pp.~1, 2.}会议由省革委会和杭州市革委会联合召开，宣布吸收红暴代表进入省工代会、农代会和红代会常委会、杭州市革委会及其常委会，并建议红暴代表加入省革委会。会议报告称，省革委会和杭州市革委会曾多次讨论红暴代表问题，但围绕1967年2月争议事件解释的分歧拖延了问题解决。1969年，中央委员会就此发出指示，促使浙江和杭州革命委员会为红暴成员举办学习班。中共九大后，红暴常委会再次研究这一问题，并依据2月原协议拟定了一份可接受名单。

The speeches delivered at the meeting took unity as their theme. In the
words of Wang Zida, chairman of the HMRC, ``Unity means force, unity
means weapons and unity brings us victory''. Deputy-commander of the
ZPMD, Wan Zhenxi (万振西) stated that Red Storm's participation in the
provincial leadership was ``required by the revolution and by the
circumstances'' and declared that it was the ``common desire and
pressing demand'' of the people of Zhejiang. Wan addressed remarks to
members of both organizations. To Red Storm he expressed the hope that
it would be modest, prudent and vigilant against those who sought to
wreck the hard-won unity. To United Headquarters he urged it to ``take
the initiative'' in uniting with Red Storm.

会议发言以团结为主题。杭州市革委会主任王子达说：``团结就是力量，团结就是武器，团结就是胜利。''浙江省军区副司令员万振西表示，红暴参加省级领导是``革命的需要、形势的需要''，并宣布这是浙江人民的``共同愿望和迫切要求''。万向两个组织的成员讲话。对红暴，他希望其谦虚谨慎，警惕那些企图破坏来之不易团结的人。对联合总部，他敦促其``主动''同红暴团结。

Spokesmen for the two sides contributed the appropriate rhetorical
flourishes and the report pictured the ``fighters'' of both groups
studying ``happily together'' and encouraging each other. To mark the
occasion a joint editorial of Zhejiang Daily and Hangzhou Daily was
published which claimed that while United Headquarters had given way on
trifling matters (that is, power!) it had stood its ground on matters of
principle.\footnote{ZJRB, May 27, 1969, (morning edition), (上午版)
  p.~1.} It is clear that United Headquarters and Red Storm remained
embroiled in disputes over their respective roles in the Cultural
Revolution. Prospects for real and lasting unity seemed slim while such
debates raged.

双方发言人都作了适当的辞藻表态，报告描绘两派``战士''\,``愉快地在一起''学习并相互鼓励。为纪念这一时刻，《浙江日报》和《杭州日报》发表联合社论，声称联合总部虽然在琐事（也就是权力！）上作了让步，但在原则问题上坚持了立场。\footnote{ZJRB,
  May 27, 1969, (morning edition), (上午版) p.~1.}显然，联合总部和红暴仍深陷关于各自在文化大革命中角色的争论。只要这类争论持续，真正而持久的团结前景就很渺茫。

Additionally, a number of Red Storm members refused to accept the
agreement and campaigned energetically and with some success to
undermine it. The main leader of the breakaway group was one of Red
Storm's most resolute and capable leaders, the former textile worker and
party member Fang Jianwen. Fang's obduracy alarmed the provincial
leadership greatly, as can be gauged from the fact that in August 1969
Zhejiang Daily issued two single-sheet special issues (专刊) devoted
entirely to recounting and repudiating Fang's exploits.\footnote{ZJRB,
  August 13 and 21, 1969.} The special issues appeared shortly after the
release of Central Committee document Number 48, dated 4 August 1969,
approving public criticism of Fang
(同意对方剑文的问题进行公开批判，以肃清其错误影响). The details of
Fang's deeds make fascinating reading and give further insight into the
machinations of factional politics during this period.

此外，一些红暴成员拒绝接受协议，并积极开展活动予以破坏，且取得一定成效。分裂集团的主要领导人，是红暴最坚定、最能干的领导人之一，前纺织工人兼党员方剑文。方的顽固使省领导极为震惊，这从1969年8月《浙江日报》发行两期单张专刊专门叙述并批判方的事迹即可看出。\footnote{ZJRB,
  August 13 and 21, 1969.}这些专刊是在1969年8月4日中央委员会第48号文件发布后不久出现的，该文件批准公开批判方剑文（同意对方剑文的问题进行公开批判，以肃清其错误影响）。方的事迹细节读来引人入胜，也进一步揭示了这一时期派性政治的权术。

Fang came from a poor peasant family in Zhuji County, situated south of
Hangzhou in Shaoxing District. Before coming to the provincial capital
to work at the No.~1 cotton mill, he had been a rural cadre, and in 1962
had been found guilty of speculating in satin and exchanging it for
grain. In 1965, during the ``four clean-ups'', Fang's peccadillo had
been discovered and noted in his file. The incident was undoubtedly a
motivating factor in Fang's decision to rebel in the Cultural
Revolution, in the same way as similar incidents contributed to the
rebellion of Wang Hongwen and Weng Senhe.

方出身于绍兴地区杭州以南诸暨县的贫农家庭。到省会第一棉纺厂工作前，他曾是农村干部；1962年，他因投机倒把缎子并用其换粮而被认定有罪。1965年``四清''期间，方的这桩过失被发现并记入档案。毫无疑问，这一事件是促使方在文化大革命中造反的动因之一，就像类似事件促成王洪文和翁森鹤造反一样。

The articles disclosed that Fang had masterminded Red Storm's disruption
of the rally convened by United Headquarters on 12 February 1967 and
alleged that he had supported Jiang Hua ever since that date. In the
summer of 1968 he had been responsible for violent incidents at
factories and colleges in Hangzhou as well as convening meetings of
groups opposed to United Headquarters in southern and eastern Zhejiang.
The bases of Fang's support were Zhejiang University, Taizhou District
in eastern Zhejiang and his unit, the No.~1 cotton mill and he had
maintained links with Nie Minzhi, Xia Liangchang and Peng Zhangxun,
whose activities and trial were related earlier in this chapter.

文章披露，方策划了红暴破坏1967年2月12日联合总部召集的集会，并声称他自那一天起一直支持江华。1968年夏，他要对杭州工厂和高校的暴力事件负责，也负责在浙南和浙东召集反联合总部团体会议。方的支持基础是浙江大学、浙东台州地区以及其所在单位第一棉纺厂；他还同聂民之、夏良焯和彭彰训保持联系，本章前文已叙述这些人的活动和审判。

Fang's most spectacular exploits took place in 1969. After the
provincial congress of the ZPRC in late 1968, at which he was forced to
confess his mistakes, Fang went to the offshore islands of Zhoushan to
organize a counter-attack against the local branch of United
Headquarters. In this endeavor he was most successful. Between February
and perhaps April of that year\footnote{A report from Hong Kong stated
  that heavy fighting between rival factions in Dinghai county on the
  Zhoushan island group occurred in May 1969. See CNS, 274 (June 12,
  1969), p.~A3.} Fang mobilized his supporters with assistance from the
local military (the 22nd Army commanded by Cao Siming and Tie Ying)
reportedly to drive 20,000 members of the local United Headquarters'
branch off Zhoushan and back onto the mainland. His strategy was to
mould public opinion in Hangzhou, stir up fighting in southern Zhejiang,
make Zhoushan a model and set up a base in Siming mountains in eastern
Zhejiang (杭州的舆论,浙南的斗争,舟山的样板,四明的根据地). In May 1969
Fang was summoned seven times to attend the provincial congress of ZPRC
party members but refused. He tried to establish a Zhejiang Province Red
Storm Revolutionary Committee. Later, in a move that must have won him
full marks for audacity if nothing else, Fang set up a ``flood relief
headquarters'' (抗洪指挥部) at the Hangzhou silk complex, the base of
his former comrade and now bitter factional enemy Weng Senhe.

方最引人注目的行动发生在1969年。1968年底省革委会代表大会上，他被迫承认错误；随后方前往舟山近海岛屿，组织对联合总部当地分部的反攻。这一行动非常成功。大约在当年2月至4月之间，\footnote{A
  report from Hong Kong stated that heavy fighting between rival
  factions in Dinghai county on the Zhoushan island group occurred in
  May 1969. See CNS, 274 (June 12, 1969), p.~A3.}方在地方军队（曹思明和铁瑛指挥的二十二军）协助下动员支持者，据称把当地联合总部分部20,000名成员赶出舟山、赶回大陆。他的策略是塑造杭州舆论、煽动浙南斗争、把舟山树为样板，并在浙东四明山建立根据地（杭州的舆论,浙南的斗争,舟山的样板,四明的根据地）。1969年5月，方七次被传唤参加省革委会党员代表大会，但他拒绝。他试图建立浙江省红暴革命委员会。后来，方又在杭州丝绸印染联合厂，即其前同志、如今的死敌翁森鹤的基地，建立``抗洪指挥部''；单就胆量而言，这一举动必然可得满分。

Fang strongly opposed the agreement of May 1969. Immediately following
the meeting of 27 May 1969, he and other oppositionists from Red Storm
set off for Beijing to voice their grievances. Qiu Honggen, a member of
the delegation, and the contemporary denunciation of Fang have both
recorded versions of the events which ensued.\footnote{Qiu Honggen,
  HZRB, November 26, 1978; ZJRB, August 13, 1969.} The two
hundred-member delegation arrived in the capital on 4 June. The CC
liaison reception station found it accommodation and food. However, the
delegation's attempts to find sympathetic listeners came to nought. On
12 June, incensed by a Xinhua report concerning Zhejiang, Fang led
twenty-four people to the organization's offices. The Beijing
authorities responded by despatching a military propaganda team from the
local garrison. Fang also led a boycott of the ``study class'' arranged
by cadres who had traveled up from Hangzhou. Subsequently, his group was
arrested and escorted back to Zhejiang, reportedly on the orders of Chen
Boda and Yao Wenyuan. On 11 August the group was confined to military
barracks and ordered to hold their ``study class''. Qiu Honggen has
described in detail the interrogation, deprivation and torture that he
experienced for the next three and a half years.

方坚决反对1969年5月协议。1969年5月27日会议刚结束，他和红暴其他反对者便启程赴北京申诉。代表团成员裘洪根以及当时对方的批判，都记录了随后事件的不同版本。\footnote{Qiu
  Honggen, HZRB, November 26, 1978; ZJRB, August 13, 1969.}200人代表团于6月4日抵达首都。中央联络接待站为其安排住宿和食物。然而，代表团试图寻找同情倾听者的努力毫无结果。6月12日，方因新华社一篇有关浙江的报道而愤怒，带领24人前往该机构办公室。北京当局派出当地卫戍部队军事宣传队回应。方还领导抵制从杭州赶来的干部所安排的``学习班''。随后，据称根据陈伯达和姚文元命令，他的小组被逮捕并押回浙江。8月11日，该小组被关押在军营，并被命令举办自己的``学习班''。裘洪根详细描述了随后三年半里他所经历的审讯、剥夺和酷刑。

The destruction of Fang Jianwen's influence and credibility was vital to
the provincial leadership as a precondition for the dissolution of both
mass organizations. Therefore, on 23 August and 3 September, meetings
were held to announce that both United Headquarters and Red Storm had
``gloriously completed their historical missions''
(光荣地完成历史使命).\footnote{ZJRB, August 24, 1969, pp.~1, 2 and
  September 4, 1969, pp.~1, 2. See also ZJRB, September 1, 1969, p.~4;
  September 3, 1969, p.~2; September 12, 1969, p.~1; September 22, 1969,
  p.~2 and September 28, 1969, p.~5 for further articles on
  factionalism, anarchism and bourgeois individualism.} The repudiation
of Fang Jianwen was a major item at the meetings. In the assessment of
the relative contributions of the two organizations to the Cultural
Revolution in Zhejiang, speakers made it very clear that Red Storm's
mistakes and the ``tortuous road'' it had followed since 1967 relegated
it to second position. Chen Liyun announced at the second of the two
meetings that all inter-unit mass organizations in the various districts
and counties of the province would be dissolved by 10 September. One
immediate side-effect of the dissolution of Red Storm and United
Headquarters was the announcement, at the end of September, that
Zhejiang Daily was henceforth to appear as the official organ of the
ZPRC.\footnote{ZJRB, September 30, 1969, p.~1.}

摧毁方剑文的影响力和公信力，是省领导层解散两大群众组织的必要前提。因此，8月23日和9月3日召开会议，宣布联合总部和红暴都已``光荣地完成历史使命''。\footnote{ZJRB,
  August 24, 1969, pp.~1, 2 and September 4, 1969, pp.~1, 2. See also
  ZJRB, September 1, 1969, p.~4; September 3, 1969, p.~2; September 12,
  1969, p.~1; September 22, 1969, p.~2 and September 28, 1969, p.~5 for
  further articles on factionalism, anarchism and bourgeois
  individualism.}批判方剑文是会议的一项主要内容。在评估两组织对浙江文化大革命的相对贡献时，发言者非常明确地表示，红暴的错误和其自1967年以来所走的``曲折道路''，使其退居第二位。陈励耘在两次会议中的第二次会议上宣布，全省各地区、各县所有跨单位群众组织将在9月10日前解散。红暴和联合总部解散的一个直接副作用，是9月底宣布《浙江日报》此后作为省革委会机关报出版。\footnote{ZJRB,
  September 30, 1969, p.~1.}

While outwardly the two groups ``lowered their flags'' and disbanded
their formal organizational structures, factional alliances continued in
a different form.\footnote{See CCP HMC Organization Department Criticism
  Group, ``A Gang who caused `earthquakes' every day'', HZRB, September
  17, 1977.} They went underground and drew on relationships forged for
mutual protection and support over the preceding three years. In the
case of United Headquarters an underground command center (地下指挥中心)
was established.\footnote{See 浙江省委理论组 (Zhejiang Party Committee
  theoretical group), ``在两个阶级的激烈斗争中走向天下大治'' (Forge
  great order in the fierce struggle between two classes), Hongqi, No.~3
  1977, p.~29.} In September 1973, a big-character poster written under
the direction of Zhang Yongsheng did not deny the charge by the ZPC
leadership that such a center existed.\footnote{See Zhang Yongshengde
  zuizheng cailiao, p.~56.} The two factions may have ceased formal
activities under their respective banners of United Headquarters and Red
Storm, but the factional alignments remained as firm and as clear-cut as
before.

表面上，两大组织``降下旗帜''并解散正式组织结构，但派性联盟以另一种形式延续。\footnote{See
  CCP HMC Organization Department Criticism Group, ``A Gang who caused
  `earthquakes' every day'', HZRB, September 17, 1977.}它们转入地下，并利用前三年为相互保护和支持而形成的关系。就联合总部而言，建立了一个地下指挥中心。\footnote{See
  浙江省委理论组 (Zhejiang Party Committee theoretical group),
  ``在两个阶级的激烈斗争中走向天下大治'' (Forge great order in the
  fierce struggle between two classes), Hongqi, No.~3 1977, p.~29.}1973年9月，在张永生指导下写成的一张大字报，并未否认省委领导关于存在这样一个中心的指控。\footnote{See
  Zhang Yongshengde zuizheng cailiao, p.~56.}两派也许停止了以联合总部和红暴各自旗号进行的正式活动，但派性阵营仍同过去一样坚定而清晰。

Although their names were not publicly announced at the time, new
members did join the leadership bodies in 1969. Jiang Baodi (江宝娣), a
silk worker at the Hangzhou Spun Silk Fabric Mill appeared as a member
of the standing committee of the ZPRC at the end of May. In April she
had been elected an alternate member of the CCP 9th CC. Navy commander,
Xie Zhenghao, later joined the ZPRC as a Vice-chairman.\footnote{ZJRB,
  August 1,1969, p.~3.} He was almost certainly a supporter of Red
Storm. Nevertheless, the terms of the alliance specifying the number of
places allotted to Red Storm, which were announced by Zhou Enlai in
March 1968, were certainly not honored. Nevertheless, Red Storm
activists were to some extent constrained from maintaining their hostile
stand toward revolutionary committees by membership of leading bodies
and participation in the decision-making process. Disillusionment set
in, however, and complaints such as ``it is a loss to be a cadre'' and
``it will be harmful to oneself if one is firm in one's stand'' were
reported.\footnote{ZJRB, July 28,1969, p.~1; ZJRB, August 13,1969, p.~1.}

虽然当时并未公开宣布姓名，但1969年确有新成员加入领导机构。江宝娣是杭州绢纺织厂丝织工人，5月底以省革委会常委身份出现。4月，她已当选中共九届中央候补委员。海军司令员谢正浩后来加入省革委会任副主任。\footnote{ZJRB,
  August 1,1969, p.~3.}他几乎肯定是红暴支持者。尽管如此，周恩来1968年3月宣布的联盟条件中关于分配给红暴席位数量的规定，显然没有得到遵守。不过，红暴积极分子由于加入领导机构并参与决策过程，在一定程度上受到约束，不能继续维持对革命委员会的敌对立场。然而，幻灭感随之出现，据报有人抱怨``当干部吃亏''，``立场坚定要害自己''。\footnote{ZJRB,
  July 28,1969, p.~1; ZJRB, August 13,1969, p.~1.}

The problem of factionalism, which had afflicted the two rival mass
organizations since 1967, appeared endemic notwithstanding the
much-publicized and much-heralded reconciliation of May 1969. Certainly,
the strife and mutual hostilities were not confined to Zhejiang. A
notice of the CCP CC dated 23 July 1969 referred to violent clashes and
armed struggles in Shanxi province and ordered the PLA to suppress those
responsible with all its might.\footnote{``The `July 23 Notice' of the
  CCP Central Committee'', I\&S, 6:1 (October 1969), pp.~97-100.}
Articles in the People's Daily called for equal treatment of rival
organizations and advocated the method used in Anhui province of holding
study classes for mass organization leaders.\footnote{See SCMP, 4470,
  pp.~4-10; 4474, pp.~1-4; 4479, pp.~1-7.} Rivalry and fighting between
mass organizations was also noted in other provinces.\footnote{Jurgen
  Domes, China after the Cultural Revolution: Politics between Two Party
  Congresses (Berkeley and Los Angeles: University of California Press,
  1977), p.~48; CNS, 274 (June 12, 1969).}

尽管1969年5月的和解被大肆宣传、广泛颂扬，自1967年以来困扰两大敌对群众组织的派性问题似乎已成痼疾。当然，争斗和相互敌视并不局限于浙江。中共中央1969年7月23日一份通知提到山西省的暴力冲突和武装斗争，并命令解放军全力镇压责任者。\footnote{``The
  `July 23 Notice' of the CCP Central Committee'', I\&S, 6:1 (October
  1969), pp.~97-100.}《人民日报》文章要求平等对待敌对组织，并提倡安徽省采用的为群众组织领导人举办学习班的方法。\footnote{See
  SCMP, 4470, pp.~4-10; 4474, pp.~1-4; 4479, pp.~1-7.}其他省份也出现了群众组织之间的竞争和战斗。\footnote{Jurgen
  Domes, China after the Cultural Revolution: Politics between Two Party
  Congresses (Berkeley and Los Angeles: University of California Press,
  1977), p.~48; CNS, 274 (June 12, 1969).}

The links established and developed between worker rebels and student
Red Guards in the years 1966-68 proved difficult to uncouple. The Maoist
leadership had dealt with the Red Guard problem by means of rustication,
military discipline and the despatch of worker propaganda teams to
occupy campuses and exercise proletarian leadership over educational
institutions. Zhang Yongsheng spent time in the countryside, despite his
status as a Vice-chairman of the provincial revolutionary committee.
Lingering unrest amongst the industrial workforce, however, demanded a
more sensitive and restrained approach. Support from workers for their
student allies kept alive the cross-unit factional links that the
authorities had been at pains to suppress since the end of
1967.\footnote{Livio Maitan argues that this support continued in 1969.
  See L. Maitan, Party, Army and Masses in China (London: New Left
  Books, 1976), p.~285. Jurgen Domes has noted trouble in Chinese
  factories in 1969. See Domes, China after the Cultural Revolution,
  p.~48.}

1966至1968年间工人造反派和学生红卫兵之间建立并发展的联系，证明难以拆解。毛主义领导层通过下乡、军事纪律，以及派工人宣传队占领校园并对教育机构实行无产阶级领导，来处理红卫兵问题。张永生虽是省革命委员会副主任，也曾在农村待过一段时间。然而，工业劳动力中持续存在的不安，要求采取更敏感和克制的办法。工人对其学生盟友的支持，使当局自1967年底以来竭力压制的跨单位派性联系继续存在。\footnote{Livio
  Maitan argues that this support continued in 1969. See L. Maitan,
  Party, Army and Masses in China (London: New Left Books, 1976),
  p.~285. Jurgen Domes has noted trouble in Chinese factories in 1969.
  See Domes, China after the Cultural Revolution, p.~48.}

The relative standing of the two rebel mass organizations in Zhejiang
was reflected in the choice of personnel to attend the Beijing
celebrations of the twentieth anniversary of the PRC. It is alleged that
criteria handed down by the central authorities for the selection of
worker representatives stipulated that they must come from
third-generation worker families.\footnote{Liu Guokai, ``A Brief
  Analysis'', p.~135.} Representatives from workers, poor and
lower-middle peasants, Red Guards, revolutionary cadres, revolutionary
intellectuals and the PLA made up the various provincial
delegations.\footnote{See RMRB, October 2,1969, transl. in CB, 893,
  p.~32.} The thirty delegates who comprised the Zhejiang delegation
came from all these categories. Of the less than ten individuals whose
identities are known, it is evident that the presence of the principal
leaders of United Headquarters -- Zhang Yongsheng, He Xianchun and the
defector Weng Senhe --- was not matched by representation from the
leadership of Red Storm. The Zhejiang University Red Guard leader from
Red Storm, Liu Ying, was present, but conspicuous by their absence were
other leaders of Red Storm, Fang Jianwen and Teng Zhu in particular, who
had not accepted the terms of the May alliance.

浙江两大造反群众组织的相对地位，反映在参加中华人民共和国成立二十周年北京庆祝活动人员的选择上。据称，中央当局下达的工人代表选拔标准规定，他们必须来自三代工人家庭。\footnote{Liu
  Guokai, ``A Brief Analysis'', p.~135.}各省代表团由工人、贫下中农、红卫兵、革命干部、革命知识分子和解放军代表组成。\footnote{See
  RMRB, October 2,1969, transl. in CB, 893, p.~32.}浙江代表团30名代表来自所有这些类别。在已知身份的不到十人中，显然联合总部主要领导人张永生、贺贤春和倒戈者翁森鹤都在场，而红暴领导层并没有相应代表。红暴浙江大学红卫兵领导人刘英在场，但其他红暴领导人，尤其是不接受5月联盟条件的方剑文和滕柱，明显缺席。

It appears, therefore, that the most prominent rebel leaders from United
Headquarters had been spared the rustication, incarceration or even
execution that their fellow rebels in other provinces had suffered. The
close links between the military rulers of Zhejiang and United
Headquarters remained strong enough for them to be spared such
treatment. Another feature of the Zhejiang delegation was the lack of
``revolutionary cadres'' included. The only civilian cadre present who
had served in the pre-Cultural Revolution power structure appeared to be
former Deputy-secretary of the old Hangzhou party committee --- Zhou
Feng (周峰).

因此，联合总部最著名的造反派领导人似乎逃过了其他省份同类造反者所遭受的下乡、监禁甚至处决。浙江军事统治者同联合总部之间的密切联系仍足够强，使他们免于这种待遇。浙江代表团的另一个特点，是缺少``革命干部''。在场唯一曾在文化大革命前权力结构中任职的文职干部，似乎是旧杭州市委副书记周峰。

In 1969, the Cultural Revolution in Zhejiang ended with the PLA firmly
in control of provincial affairs. It was a military leadership whose
loyalties lay firmly with its masters in Beijing, and Lin Biao, in
particular. The extravagant praise of Lin contained in a Zhejiang Daily
editorial of April 1969 was evidence of this. The obsequiousness of the
provincial leadership toward Lin had been noted in 1968 and 1969 in Hong
Kong.\footnote{See CNA, 731 p.~2; 746 (February 28,1969), p.~3.}
``Middle- and lower-ranking military men'' such as Nan Ping, Xiong
Yingtang and Chen Liyun, have been described as ``extensions of Lin's
`mountain strongholds'\,''.\footnote{Liu Guokai, ``A Brief Analysis'',
  p.~132.} Additionally, Zhejiang's military leadership was unfamiliar
with the conditions of the province and lacked the administrative skills
and experience to run the administration. The former CCP ZPC had been
decimated by the Cultural Revolution and the painful process of cadre
rehabilitation, which commenced at the grass roots level, would not
reach into the upper echelons of cadres until after Lin Biao's demise.

1969年，浙江文化大革命以解放军牢牢控制省政而告终。这是一个忠诚牢牢指向北京主人，尤其是林彪的军事领导层。1969年4月《浙江日报》社论中对林的过度赞美就是证据。香港在1968和1969年已注意到省领导层对林的谄媚。\footnote{See
  CNA, 731 p.~2; 746 (February 28,1969), p.~3.}南萍、熊应堂、陈励耘等``中下级军人''被描述为林的``山头''的延伸。\footnote{Liu
  Guokai, ``A Brief Analysis'', p.~132.}此外，浙江军事领导层不熟悉省情，也缺乏管理行政机构的行政技能和经验。原中共浙江省委已被文化大革命重创，干部平反这一从基层开始的痛苦过程，要到林彪败亡后才会触及高层干部。

The military ruled a divided, shell-shocked population. Three years of
violence, extremism and ceaseless disruption left a legacy which would
take years to overcome. Students, young workers and peasants, who in
1966 had responded to Mao's request to ``bombard the headquarters'',
returned to their posts or contemplated a future life in the
countryside. Factional enmities would not disappear because of
signatures on documents. A war of words regarding the rights and wrongs
of the events of the Cultural Revolution broke out in offices and
workshops which would poison work relations and further retard work
efficiency. In retrospect, the events of 1966-69 were only the first
round in a contest that paused for breathing-space before resuming with
full ferocity in 1973.

军队统治着一个分裂且饱受震荡的人群。三年的暴力、极端主义和不断破坏留下的遗产，需要多年才能克服。1966年响应毛``炮打司令部''号召的学生、青年工人和农民，回到岗位，或思考未来在农村的生活。派性仇恨不会因为文件上的签名而消失。围绕文化大革命事件是非的口水战在办公室和车间爆发，它将毒化工作关系，并进一步降低工作效率。回头看，1966至1969年的事件只是这场竞赛的第一回合；它暂停片刻喘息，随后在1973年以全部猛烈程度重新开始。

The Cultural Revolution had traversed a bumpy road in
Zhejiang.\footnote{See Domes, conclusion of ``Generals and Red Guards'',
  pp.~153-58.} By mid-1969 the central authorities may have stabilized
the province but they had achieved it at a tremendous cost. The few
members of the old administration who displayed loyalty to the new order
owed their position to factional groups and military backing. United
Headquarters, which had slavishly followed the somersault policy
directives of Mao and the Cultural Revolution Group, discovered that its
loyalty was not repaid when expediency and necessity brought the turmoil
to an end. Red Storm, which had defied central directives, paid a price
but not, in the eyes of United Headquarters at least, one commensurate
with the degree of its defiance.

文化大革命在浙江走过了一条崎岖道路。\footnote{See Domes, conclusion of
  ``Generals and Red Guards'', pp.~153-58.}到1969年中，中央当局也许稳定了浙江，但代价极其巨大。旧行政机构中少数对新秩序表示忠诚的成员，其职位要归功于派性集团和军方支持。联合总部奴仆般追随毛和中央文革小组反复翻转的政策指示，却发现当权宜和必要使动乱结束时，其忠诚并未得到回报。红暴违抗中央指示，付出了代价；但至少在联合总部看来，这一代价并不与其违抗程度相称。

Rebellion and factionalism had marked the course of the Cultural
Revolution and in its particular features was its legacy to the Chinese
political system. From the central to the local level, factional ties
were strengthened even further and organizational discipline was
correspondingly weakened. Leaders and activists of mass organizations,
now members of revolutionary committees, brought a new dimension to
factional politics. When the struggle for political ascendancy was made
more urgent by the question of succession which became a more salient
issue in the CCP as the 1970s progressed, factionalism, accompanied by
renewed rebellion of a more potent nature than that which occurred in
1967, debilitated the political system to an even greater extent. How
and why this took place forms the subject of the following chapters.

造反和派性标志着文化大革命的进程，也以其特殊形式成为它留给中国政治制度的遗产。从中央到地方，派性联系进一步强化，组织纪律相应削弱。群众组织领导人和积极分子如今成为革命委员会成员，为派性政治带来新的维度。随着20世纪70年代推进，接班问题在中共内部变得更加突出，使争夺政治优势的斗争更加迫切；派性伴随着比1967年更有力的新一轮造反，进一步削弱了政治制度。其如何发生、为何发生，将是以下各章的主题。

\section{Notes}\label{notes-4}

\section{注释}\label{ux6ce8ux91ca-4}

\bookmarksetup{startatroot}

\chapter{CHAPTER FIVE - Military Rule in Zhejiang, 1969-72 / 第五章 -
浙江的军事统治，1969-72年}\label{chapter-five---military-rule-in-zhejiang-1969-72-ux7b2cux4e94ux7ae0---ux6d59ux6c5fux7684ux519bux4e8bux7edfux6cbb1969-72ux5e74}

In the four and a half years which elapsed between the CCP's 9th (1969)
and 10th (1973) National Congresses, fierce political struggles
continued unabated within the upper echelons of the party. Provincial
officials became intimately involved in these disputes which revolved
around such issues as party reconstruction, the role of the PLA in the
political, economic and social spheres, the evolution of the political
institutions and groups brought into existence by the Cultural
Revolution, and the respective roles and relationship between military
cadres, civilian officials and mass organization representatives within
leading groups. It was a period in which the ascendancy that the
military had gained during the years 1967-69 aroused concern among
sections of the Chinese leadership, Mao Zedong in particular. Various
commentators have pointed to the irony of the situation in which the
most disciplined, unaccountable and authoritarian component of the state
structure was able to emerge triumphant from a political movement that
had, among its alleged objectives, the loosening of a constrictive
political system and the participation of a broader spectrum of the
population in the political processes.\footnote{See especially Maurice
  Meisner, Mao's China (New York: The Free Press, 1977), pp.~384-5;
  Meisner, ``The Concept of the Dictatorship of the Proletariat in
  Chinese Marxist Thought'', in V. Nee \& D. Mozingo (eds.). State and
  Society in Contemporary China (Ithaca, New York: Cornell University
  Press, 1983), p.~125; Gordon White, ``The Revolutionary Chinese
  State'', in Nee \& Mozingo, State and Society, pp.~48-50. 尤其参见
  Maurice Meisner, \emph{Mao's China}（纽约：The Free
  Press，1977年），第384-385页；Meisner，《中国马克思主义思想中的无产阶级专政概念》，载
  V. Nee 与 D. Mozingo 编，\emph{State and Society in Contemporary
  China}（纽约伊萨卡：Cornell University
  Press，1983年），第125页；Gordon White，《革命的中国国家》，载 Nee 与
  Mozingo，\emph{State and Society}，第48-50页。} At the 9th Party
Congress Lin Biao had been officially written into the statutes of the
party constitution as Mao's designated successor. However, the
obstructions placed in the way of party reconstruction after 1969 and
the political strength of regional military commanders, some of whom
were loyal to Lin and his central military faction, provoked the
opposition of powerful groups within the CCP.\footnote{For a comment on
  the nature, origins and structure of the Lin Biao group, see Hao
  Mengbi and Duan Haoran, Liushinian, p.~611.
  关于林彪集团的性质、起源和结构的评论，见 Hao Mengbi and Duan Haoran,
  \emph{Liushinian}，第611页。} In Zhejiang, the new leaders of the ZPRC
were overwhelmingly from a military background and from military units
directly answerable and loyal to their central military patron, Lin
Biao. Although Nan Ping, Xiong Yingtang and Chen Liyun, like most
officers in the Zhejiang military district, had served in the 3rd Field
Army before 1949 they were not senior officers. The Cultural Revolution
had boosted their careers and brought them directly into the political
arena. These military officials were less tied to the locality in which
they served than any previous provincial leadership since 1949. Teiwes
has tabulated this phenomenon for China's provincial-level
administrations as a whole.\footnote{Teiwes, Provincial Leadership in
  China, pp.~141-42, and Appendix D, p.~151. Teiwes, \emph{Provincial
  Leadership in China}，第141-142页及附录D第151页。} The leadership's
lack of familiarity with local conditions did not seem to prevent
substantial economic growth in the years 1969-72. Major disturbances,
particularly in the urban areas of the province, had disrupted
industrial production in the years 1967 and 1968 and the recovery over
the next few years made the province something of a national model for
economic development.\footnote{See the documents in URS, 58: 8 (January
  27,1970), pp.~106-18. 见 \emph{URS},
  58:8（1970年1月27日）中的文件，第106-118页。} Total agricultural and
industrial output value in 1969 increased 18.1\% over the previous year
to establish an all-time provincial record. Over the following three
years it grew at an annual rate of 11.7\%, 9.5\% and 11.5\%
respectively. In 1969, the value of agricultural output reached its
highest level since 1964. Substantial increases in the value of
industrial output of 30\%, 16.2\%, 15.4\% and 10.2\% occurred over the
years 1969-72.\footnote{Zhejiang shengqing gaiyao, pp.~694-95. These
  figures do not take account of price inflation. \emph{Zhejiang
  shengqing gaiyao}，第694-695页。这些数字未计入价格通货膨胀。} This
chapter, more than anywhere else in this book, has drawn on source
material at the national level to compensate for the dearth of
contemporary information from Zhejiang. The paucity of sources on the
politics of the province during these years may reflect the high degree
of social control that the military exerted over provincial affairs and
the general tightening of the flow of information after the leakages of
sensitive material by Red Guards over the previous two years.
Contemporary observers again commented on trends in Zhejiang which were
unique to the province and stood it apart from other provinces. The
notion of Zhejiang as a special case, while difficult to substantiate
without detailed comparative information, will be noted where relevant
and significant.
在中国共产党第九次（1969年）和第十次（1973年）全国代表大会之间的四年半里，党内高层的激烈政治斗争持续不断。省级官员深深卷入这些争端，争端围绕党的重建、解放军在政治、经济和社会领域中的作用、文化大革命所造就的政治机构和群体的演变，以及领导班子内部军队干部、文职官员和群众组织代表各自的作用及相互关系等问题展开。1967至1969年间军队取得的优势，在这一时期引起了中国领导层部分人士，尤其是毛泽东的担忧。许多评论者指出了这一局面的讽刺性：国家结构中最有纪律、最不受问责、最具威权色彩的组成部分，竟能从一场据称旨在放松僵化政治体制、扩大民众参与政治过程的政治运动中胜出。\footnote{See
  especially Maurice Meisner, Mao's China (New York: The Free Press,
  1977), pp.~384-5; Meisner, ``The Concept of the Dictatorship of the
  Proletariat in Chinese Marxist Thought'', in V. Nee \& D. Mozingo
  (eds.). State and Society in Contemporary China (Ithaca, New York:
  Cornell University Press, 1983), p.~125; Gordon White, ``The
  Revolutionary Chinese State'', in Nee \& Mozingo, State and Society,
  pp.~48-50. 尤其参见 Maurice Meisner, \emph{Mao's China}（纽约：The
  Free
  Press，1977年），第384-385页；Meisner，《中国马克思主义思想中的无产阶级专政概念》，载
  V. Nee 与 D. Mozingo 编，\emph{State and Society in Contemporary
  China}（纽约伊萨卡：Cornell University
  Press，1983年），第125页；Gordon White，《革命的中国国家》，载 Nee 与
  Mozingo，\emph{State and Society}，第48-50页。}在中共九大上，林彪作为毛的指定接班人被正式写入党章。然而，1969年以后党重建遭遇的阻碍，以及地区军事指挥员的政治力量，其中一些人效忠林彪及其中央军事派系，激起了中共内部强势集团的反对。\footnote{For
  a comment on the nature, origins and structure of the Lin Biao group,
  see Hao Mengbi and Duan Haoran, Liushinian, p.~611.
  关于林彪集团的性质、起源和结构的评论，见 Hao Mengbi and Duan Haoran,
  \emph{Liushinian}，第611页。}在浙江，省革委会的新领导人绝大多数具有军队背景，来自直接听命并忠于其中央军事庇护人林彪的部队。南萍、熊应堂、陈励耘和浙江省军区的大多数军官一样，1949年前曾在第三野战军服役，但他们并非高级军官。文化大革命推动了他们的仕途，并把他们直接带入政治舞台。与1949年以来浙江以往任何一届省级领导相比，这些军队干部同任职地方的联系都更弱。Teiwes已对中国省级行政机构的这一现象作过整体统计。\footnote{Teiwes,
  Provincial Leadership in China, pp.~141-42, and Appendix D, p.~151.
  Teiwes, \emph{Provincial Leadership in
  China}，第141-142页及附录D第151页。}领导层不熟悉地方情况，似乎并未妨碍1969至1972年间经济的显著增长。1967和1968年，省内尤其是城市地区的重大动乱扰乱了工业生产，随后数年的恢复使浙江在某种程度上成为全国经济发展的样板。\footnote{See
  the documents in URS, 58: 8 (January 27,1970), pp.~106-18. 见
  \emph{URS}, 58:8（1970年1月27日）中的文件，第106-118页。}1969年工农业总产值比上一年增长18.1\%，创下全省历史最高纪录。此后三年分别以11.7\%、9.5\%和11.5\%的年率增长。1969年，农业产值达到1964年以来的最高水平。1969至1972年间，工业产值分别大幅增长30\%、16.2\%、15.4\%和10.2\%。\footnote{Zhejiang
  shengqing gaiyao, pp.~694-95. These figures do not take account of
  price inflation. \emph{Zhejiang shengqing
  gaiyao}，第694-695页。这些数字未计入价格通货膨胀。}本章比本书其他任何章节都更多依赖全国层面的资料，以弥补浙江当代资料的匮乏。关于这些年省内政治的资料稀少，可能反映了军队对省政事务施加的高度社会控制，以及此前两年红卫兵泄露敏感材料之后信息流通的普遍收紧。当代观察者也再次注意到浙江特有、使其有别于其他省份的趋势。浙江作为特殊个案这一观念，虽因缺乏详细比较资料而难以证实，但在相关且重要之处仍将予以说明。

\section{The Lin Biao Affair,
1969-71}\label{the-lin-biao-affair-1969-71}

\section{林彪事件，1969-71年}\label{ux6797ux5f6aux4e8bux4ef61969-71ux5e74}

The downfall and death of Lin Biao remains one of the most disastrous
and, despite the recent spate of publications on the affair, one of the
most mysterious episodes in the turbulent history of the CCP. It caused
Mao Zedong great anguish and was a humiliating blow to his prestige.
That his recently-anointed successor should plot to assassinate him and
then flee the country was an indictment of Mao's judgment and a major
reversal for the Cultural Revolution itself. After all, Lin had emerged
in the Cultural Revolution as the best and most loyal pupil of the
``Great Helmsman''. The accumulation of a series of incidents,
culminating in Lin's flight in September 1971, had served gradually to
alienate the Chairman from his designated successor. In October 1969 Lin
first aroused Mao's suspicions about his ambitions when he issued Order
No.~1 at a time following bloody clashes between the Soviet and Chinese
forces on the disputed Manchurian border.\footnote{See the Chinese
  interpretation of these events in Down with the New Tsars! (Peking:
  Foreign Languages Press, 1969) and Statement of the Government of the
  People's Republic of China (October 7,1969), (Peking: Foreign
  Languages Press, 1973). 关于这些事件的中国方面解释，见 \emph{Down with
  the New Tsars!}（北京：外文出版社，1969年）以及 \emph{Statement of the
  Government of the People's Republic of
  China}（1969年10月7日）（北京：外文出版社，1973年）。} Lin's order of
17 October 1969 had two aims, according to one of its targets, Marshall
Nie Rongzhen.\footnote{1966-1976年大事记, in Liu Suinian and Wu Qungan,
  ``Wenhua dageming'' shiqide guomin jingji, p.~154. Nie and Hu Hua have
  dated the order at October 18,1969. See Nie Rongzhen Huiyilu, Vol. 3
  (Beijing: Jiefangjun Chubanshe, 1984), p.~861; Hu Hua (ed.),
  中国社会主义革命和建设讲义 (Lectures on China's Socialist Revolution
  and Construction), (Beijing: Renmin Chubanshe, 1985), p.~303. The
  confusion has arisen because at 9.30 pm. on October 18 acting Chief of
  Staff, Huang Yongsheng, officially issued an ``emergency directive''
  (紧急指示) based on Lin's Order No.~1 of the previous day. See
  中共中央党史研究室 (Central Research Office of Party History),
  中共党史大事年表 (Chronological Tabulation of Major Events in the
  History of the Chinese Communist Party), (Beijing: Renmin Chubanshe,
  1987), p.~372; Wang Nianyi, ``评'文化大革命'十年史'' (A Critique of
  the ``Ten Year History of the `Cultural Revolution'\,''), 党史通讯
  (Bulletin on Party History), (April 1987), p.~20.
  《1966-1976年大事记》，载 Liu Suinian and Wu Qungan, \emph{``Wenhua
  dageming'' shiqide guomin
  jingji}，第154页。聂荣臻和胡华把该命令日期定为1969年10月18日。见
  \emph{Nie Rongzhen Huiyilu}
  第3卷（北京：解放军出版社，1984年），第861页；胡华编，《中国社会主义革命和建设讲义》（北京：人民出版社，1985年），第303页。混淆产生的原因是，10月18日晚9时30分，代总参谋长黄永胜根据林彪前一天的一号命令正式发布了``紧急指示''。见中共中央党史研究室，《中共党史大事年表》（北京：人民出版社，1987年），第372页；王年一，《评``文化大革命''十年史》，《党史通讯》（1987年4月），第20页。}
First, Lin was testing the ground for a coup by creating an exaggerated
atmosphere of crisis.
林彪的垮台和死亡，仍然是中共动荡历史中最具灾难性的事件之一；尽管近来有关此事的出版物大量出现，它也仍是最神秘的事件之一。这件事给毛泽东造成了巨大痛苦，也使他的威望遭受羞辱性打击。刚被指定的接班人竟然密谋刺杀他并逃离国家，这既是对毛判断力的否定，也是文化大革命本身的一次重大逆转。毕竟，林彪在文化大革命中曾被塑造成``伟大舵手''最好、最忠诚的学生。以1971年9月林彪出逃为高潮的一连串事件，逐渐使主席同其指定接班人疏远。1969年10月，在中苏军队于有争议的满洲边境发生流血冲突之后，林彪发布``一号命令''，首次引起毛对其野心的怀疑。\footnote{See
  the Chinese interpretation of these events in Down with the New Tsars!
  (Peking: Foreign Languages Press, 1969) and Statement of the
  Government of the People's Republic of China (October 7,1969),
  (Peking: Foreign Languages Press, 1973).
  关于这些事件的中国方面解释，见 \emph{Down with the New
  Tsars!}（北京：外文出版社，1969年）以及 \emph{Statement of the
  Government of the People's Republic of
  China}（1969年10月7日）（北京：外文出版社，1973年）。}据该命令的针对对象之一聂荣臻元帅说，林彪1969年10月17日的命令有两个目的。\footnote{1966-1976年大事记,
  in Liu Suinian and Wu Qungan, ``Wenhua dageming'' shiqide guomin
  jingji, p.~154. Nie and Hu Hua have dated the order at October
  18,1969. See Nie Rongzhen Huiyilu, Vol. 3 (Beijing: Jiefangjun
  Chubanshe, 1984), p.~861; Hu Hua (ed.), 中国社会主义革命和建设讲义
  (Lectures on China's Socialist Revolution and Construction), (Beijing:
  Renmin Chubanshe, 1985), p.~303. The confusion has arisen because at
  9.30 pm. on October 18 acting Chief of Staff, Huang Yongsheng,
  officially issued an ``emergency directive'' (紧急指示) based on Lin's
  Order No.~1 of the previous day. See 中共中央党史研究室 (Central
  Research Office of Party History), 中共党史大事年表 (Chronological
  Tabulation of Major Events in the History of the Chinese Communist
  Party), (Beijing: Renmin Chubanshe, 1987), p.~372; Wang Nianyi,
  ``评'文化大革命'十年史'' (A Critique of the ``Ten Year History of the
  `Cultural Revolution'\,''), 党史通讯 (Bulletin on Party History),
  (April 1987), p.~20. 《1966-1976年大事记》，载 Liu Suinian and Wu
  Qungan, \emph{``Wenhua dageming'' shiqide guomin
  jingji}，第154页。聂荣臻和胡华把该命令日期定为1969年10月18日。见
  \emph{Nie Rongzhen Huiyilu}
  第3卷（北京：解放军出版社，1984年），第861页；胡华编，《中国社会主义革命和建设讲义》（北京：人民出版社，1985年），第303页。混淆产生的原因是，10月18日晚9时30分，代总参谋长黄永胜根据林彪前一天的一号命令正式发布了``紧急指示''。见中共中央党史研究室，《中共党史大事年表》（北京：人民出版社，1987年），第372页；王年一，《评``文化大革命''十年史》，《党史通讯》（1987年4月），第20页。}第一，林通过制造夸大的危机气氛，试探政变的基础。

According to later accounts, all of which ascribe the meanest motives to
Lin, the Vice-chairman failed to inform Mao of his intentions. One
source states, however, that Lin's directive was based on Mao's
assessment of the possibility of a worsening of the international
situation.\footnote{Zhonggong dangshi dashi nianbiao, p.~372.
  \emph{Zhonggong dangshi dashi nianbiao}，第372页。} That he acted in
accordance with his estimation of Mao's views seems more feasible. It is
hardly likely that Lin, who was renowned for his obsequiousness to the
Chairman, would take such a drastic step without first securing Mao's
approval.\footnote{A voluminous body of investigative literature
  (纪实文学) concerning Lin Biao and his family has appeared in China in
  recent years. For a selection, see Lin Qingshan, 林彪传 (A Biography
  of Lin Biao), 2 Vols. (Beijing: Zhishi Chubanshe, 1988); He Li (ed),
  林彪家族纪事 (A Chronicle of Lin Biao's family), (Beijing: Guangming
  Ribao Chubanshe, 1989); Zhang Yunsheng, 毛家湾纪实 - 林彪秘书回忆录
  (The True Story of Maojiawan - Recollections of Lin Biao's Secretary),
  (Beijing: Chunqiu Chubanshe, 1988); Nan Zhi, 叶群野史 (The Unofficial
  Biography of Ye Qun), (Shenyang: Shenyang Chubanshe, 1988).
  近年中国出现了大量关于林彪及其家庭的纪实文学（纪实文学）。可参见林青山，《林彪传》，2卷（北京：知识出版社，1988年）；何力编，《林彪家族纪事》（北京：光明日报出版社，1989年）；张云生，《毛家湾纪实
  -
  林彪秘书回忆录》（北京：春秋出版社，1988年）；南枝，《叶群野史》（沈阳：沈阳出版社，1988年）。}
According to Nie, the second objective of Lin's order was to supply the
pretext for the removal of veteran PLA leaders such as himself, Chen Yi,
Xu Xiangqian, Liu Bocheng, Ye Jianying and Zhu De from the capital. They
were ordered to remain under virtual house arrest in various cities and
towns in central China, in almost a straight line running south from
Beijing to Guangzhou. Non-military party leaders such as Li Fuchun and
Deng Xiaoping, and the dying Liu Shaoqi and Tao Zhu, were also forced
out of Beijing by this order.\footnote{Nie Rongzhen Huiyilu, pp.~861-63;
  Document Research Office of the CCP CC, 朱德年谱 (A Chronology of Zhu
  De's Life) (Beijing: Renmin Chubanshe, 1986), pp.~550-51; Xu
  Xiangqian, Lishide huigu, pp.~848-49. For harrowing accounts of the
  treatment meted out to Liu Shaoqi, see Zhu Kexian and Bian Ka,
  ``最后的二十七天'' (The Final 27 Days), RMRB, May 20,1980, p.~8; Liu
  Pingping, Liu Yuan and Liu Tingting, ``怀念我们的爸爸刘少奇'' (In
  Memory of our Father Liu Shaoqi), Lishi zai zheli chensi, Vol. 1,
  pp.~45-7. The dying Liu was flown from Beijing to Kaifeng on October
  17. See Wang Nianyi, Dangshi tongxun, p.~20. \emph{Nie Rongzhen
  Huiyilu}，第861-863页；中共中央文献研究室，《朱德年谱》（北京：人民出版社，1986年），第550-551页；Xu
  Xiangqian, \emph{Lishide
  huigu}，第848-849页。关于刘少奇所受待遇的惨痛叙述，见 Zhu Kexian and
  Bian Ka，《最后的二十七天》，《人民日报》，1980年5月20日，第8版；Liu
  Pingping, Liu Yuan and Liu
  Tingting，《怀念我们的爸爸刘少奇》，\emph{Lishi zai zheli chensi}
  第1卷，第45-47页。病危的刘少奇于10月17日被从北京空运至开封。见 Wang
  Nianyi, \emph{Dangshi tongxun}，第20页。} Chen Boda and Lin Biao
further incurred Mao's ire for their persistent efforts to have the
position of head of state written into the revised state constitution
which was discussed at the 9th CCP CC's 2nd plenum (the Lushan plenum)
in August-September 1970. At the plenum Chen Boda, Ye Qun (Lin's wife)
and Wu Faxian (Commander of the Air Force) attacked Zhang Chunqiao by
innuendo for not praising Mao as a genius in the draft of the
constitution that he had drawn up and presented to the meeting for
discussion.\footnote{Hu Hua, ``The Rise and Fall of Lin Biao'', p.~3. Hu
  Hua，《林彪的兴衰》，第3页。} Mao questioned the motives of those who
wished to raise him to such great heights, although three years
previously he had exhibited no such unease about the personality cult
which had grown to monstrous proportions. Additionally, the Chairman
suspected that the restoration of the post of state Chairman was a means
by which Lin, given his status as successor, would consolidate his
position and rectify the anomaly whereby as Vice-premier and defence
minister he was junior to Premier Zhou Enlai in the state
administration. After Mao's decisive intervention the issue was settled,
but the struggle continued on the ideological, political, organizational
and policy fronts. Chen and Lin's principal supporters in the PLA were
subjected to scathing criticism at three meetings which were convened
between December 1970 and April 1971 -- the North China meeting (22
December 1970 to 24 January 1971), the CC Military Affairs Commission
Conference (9 January to 19 February 1971); and the report meeting to
``criticize Chen and rectify the style of work'' (15 to 29 April
1971).\footnote{See Hu Hua (ed.), Zhongguo shehuizhuyi geming,
  pp.~300-03; Hao Mengbi and Duan Haoran, Liushinian, pp.~613-20.
  见胡华编，\emph{Zhongguo shehuizhuyi geming}，第300-303页；Hao Mengbi
  and Duan Haoran, \emph{Liushinian}，第613-620页。} The central notice
transmitted from the April session was relayed to the provinces and
discussed at meetings in May 1971.\footnote{See Shaanxi jingji dashiji,
  p.~337. Lanxi county held meetings in June and September 1971 and
  Linhai county in June 1971 to criticize Chen Boda. See Lanxi shizhi,
  p.~21; Linhai xianzhi, p.~30. 见 \emph{Shaanxi jingji
  dashiji}，第337页。兰溪县于1971年6月和9月、临海县于1971年6月召开会议批判陈伯达。见
  \emph{Lanxi shizhi}，第21页；\emph{Linhai xianzhi}，第30页。} The
campaign against Chen Boda described him as a ``sham Marxist swindler''
and it gradually shifted from concentration on the manifestations of
anarchic politics in the Cultural Revolution to a more fundamental
critique of the doctrine underpinning it.\footnote{Domes, China after
  the Cultural Revolution, pp.~85-7. Domes, \emph{China after the
  Cultural Revolution}，第85-87页。}
后来的叙述都把最卑劣的动机归于林彪，并称这位副主席没有把自己的意图告知毛。不过，有一项资料称，林的指示是以毛对国际形势可能恶化的估计为依据的。\footnote{Zhonggong
  dangshi dashi nianbiao, p.~372. \emph{Zhonggong dangshi dashi
  nianbiao}，第372页。}他按自己对毛的看法的理解行事，似乎更为可信。林以对主席恭顺闻名，若不先取得毛的认可，就采取如此激烈的步骤，几乎不大可能。\footnote{A
  voluminous body of investigative literature (纪实文学) concerning Lin
  Biao and his family has appeared in China in recent years. For a
  selection, see Lin Qingshan, 林彪传 (A Biography of Lin Biao), 2 Vols.
  (Beijing: Zhishi Chubanshe, 1988); He Li (ed), 林彪家族纪事 (A
  Chronicle of Lin Biao's family), (Beijing: Guangming Ribao Chubanshe,
  1989); Zhang Yunsheng, 毛家湾纪实 - 林彪秘书回忆录 (The True Story of
  Maojiawan - Recollections of Lin Biao's Secretary), (Beijing: Chunqiu
  Chubanshe, 1988); Nan Zhi, 叶群野史 (The Unofficial Biography of Ye
  Qun), (Shenyang: Shenyang Chubanshe, 1988).
  近年中国出现了大量关于林彪及其家庭的纪实文学（纪实文学）。可参见林青山，《林彪传》，2卷（北京：知识出版社，1988年）；何力编，《林彪家族纪事》（北京：光明日报出版社，1989年）；张云生，《毛家湾纪实
  -
  林彪秘书回忆录》（北京：春秋出版社，1988年）；南枝，《叶群野史》（沈阳：沈阳出版社，1988年）。}据聂荣臻说，林彪命令的第二个目标，是为把聂本人以及陈毅、徐向前、刘伯承、叶剑英、朱德等解放军老领导调离首都提供借口。他们被命令留在华中各城市和城镇，实际上处于软禁状态，地点几乎连成一条从北京向南通往广州的直线。李富春、邓小平等非军队党的领导人，以及垂死的刘少奇和陶铸，也被这一命令迫离北京。\footnote{Nie
  Rongzhen Huiyilu, pp.~861-63; Document Research Office of the CCP CC,
  朱德年谱 (A Chronology of Zhu De's Life) (Beijing: Renmin Chubanshe,
  1986), pp.~550-51; Xu Xiangqian, Lishide huigu, pp.~848-49. For
  harrowing accounts of the treatment meted out to Liu Shaoqi, see Zhu
  Kexian and Bian Ka, ``最后的二十七天'' (The Final 27 Days), RMRB, May
  20,1980, p.~8; Liu Pingping, Liu Yuan and Liu Tingting,
  ``怀念我们的爸爸刘少奇'' (In Memory of our Father Liu Shaoqi), Lishi
  zai zheli chensi, Vol. 1, pp.~45-7. The dying Liu was flown from
  Beijing to Kaifeng on October 17. See Wang Nianyi, Dangshi tongxun,
  p.~20. \emph{Nie Rongzhen
  Huiyilu}，第861-863页；中共中央文献研究室，《朱德年谱》（北京：人民出版社，1986年），第550-551页；Xu
  Xiangqian, \emph{Lishide
  huigu}，第848-849页。关于刘少奇所受待遇的惨痛叙述，见 Zhu Kexian and
  Bian Ka，《最后的二十七天》，《人民日报》，1980年5月20日，第8版；Liu
  Pingping, Liu Yuan and Liu
  Tingting，《怀念我们的爸爸刘少奇》，\emph{Lishi zai zheli chensi}
  第1卷，第45-47页。病危的刘少奇于10月17日被从北京空运至开封。见 Wang
  Nianyi, \emph{Dangshi tongxun}，第20页。}陈伯达和林彪坚持试图把国家主席职位写入1970年8月至9月中共九届二中全会（庐山会议）讨论的修订国家宪法，从而进一步激怒了毛。会上，陈伯达、叶群（林妻）和空军司令吴法宪含沙射影地攻击张春桥，理由是张起草并提交会议讨论的宪法草案没有称颂毛为天才。\footnote{Hu
  Hua, ``The Rise and Fall of Lin Biao'', p.~3. Hu
  Hua，《林彪的兴衰》，第3页。}毛质疑那些想把他捧到如此高度的人的动机，尽管三年前他对已经发展到惊人程度的个人崇拜并未表现出这种不安。此外，主席怀疑恢复国家主席职位是林彪利用其接班人身份巩固地位、纠正其作为副总理兼国防部长在国家行政体系中低于周恩来总理这一异常状态的手段。毛作出决定性干预后，问题得到解决，但斗争继续在思想、政治、组织和政策层面展开。1970年12月至1971年4月间召开的三次会议上，陈伯达和林彪在解放军中的主要支持者受到严厉批判：华北会议（1970年12月22日至1971年1月24日）、中央军委会议（1971年1月9日至2月19日），以及``批陈整风''汇报会（1971年4月15日至29日）。\footnote{See
  Hu Hua (ed.), Zhongguo shehuizhuyi geming, pp.~300-03; Hao Mengbi and
  Duan Haoran, Liushinian, pp.~613-20. 见胡华编，\emph{Zhongguo
  shehuizhuyi geming}，第300-303页；Hao Mengbi and Duan Haoran,
  \emph{Liushinian}，第613-620页。}4月会议传达的中央通知随后下达到各省，并在1971年5月的会议上讨论。\footnote{See
  Shaanxi jingji dashiji, p.~337. Lanxi county held meetings in June and
  September 1971 and Linhai county in June 1971 to criticize Chen Boda.
  See Lanxi shizhi, p.~21; Linhai xianzhi, p.~30. 见 \emph{Shaanxi
  jingji
  dashiji}，第337页。兰溪县于1971年6月和9月、临海县于1971年6月召开会议批判陈伯达。见
  \emph{Lanxi shizhi}，第21页；\emph{Linhai xianzhi}，第30页。}反陈伯达运动称他为``假马克思主义骗子''，并逐渐从集中批判文化大革命中无政府政治的表现，转向对其支撑学说的更根本批判。\footnote{Domes,
  China after the Cultural Revolution, pp.~85-7. Domes, \emph{China
  after the Cultural Revolution}，第85-87页。}

In September 1970 immediately following the plenum, Mao moved to weaken
Lin's power. He directed the party center to establish an Organization
and Propaganda Group directly under the Politburo with Kang Sheng as its
leader and Jiang Qing, Zhang Chunqiao, Yao Wenyuan, Ji Dengkui and Li
Desheng as members. None of Lin's close associates were thus appointed
to this powerful body. Because of Kang Sheng's poor health and Li
Desheng's involvement in military affairs, the Cultural Revolution
civilian radicals took charge of this powerful super-department. It was
placed in charge of the CC Organization Department, the Central Party
School, People's Daily, Guangming Daily, Red Flag, Xinhua Newsagency,
the Central Broadcasting Bureau, central publication and translation
bureaux. May 7 cadre schools and central mass organizations such as the
All-China Federation of Trade Unions, the CYL and the Women's
Federation.\footnote{Hao Mengbi and Duan Haoran, Liushinian, p.~618;
  Zhong Kan, Kang Sheng pingzhuan, p.~433. Hao Mengbi and Duan Haoran,
  \emph{Liushinian}，第618页；Zhong Kan, \emph{Kang Sheng
  pingzhuan}，第433页。} Then, on 7 April 1971 Li Xiannian, Ji Dengkui
and Zhang Caiqian were added to the membership of the CC MAC Executive
Group (中央军委办事组), a body established in 1968 and hitherto
dominated by supporters of Lin Biao.\footnote{For information on the
  formation of this body in 1968, see Jin Chunming, A Discussion and
  Analysis, p.~136, and Tong De, Xinghuo Liaoyuan, p.~30.
  关于该机构1968年成立的资料，见 Jin Chunming, \emph{A Discussion and
  Analysis}，第136页；Tong De, \emph{Xinghuo Liaoyuan}，第30页。} Mao
graphically described the series of moves as throwing stones, mixing
sand into soil and undermining the wall. This referred to his critical
comments on documents and self-examinations written by Chen Boda and
Lin's military clique and discussed and eventually rejected as
insufficient by the MAC Conference of January to February 1971; the
April 1971 addition of new members to the MAC Executive Group and the
December 1970 reorganization of the leadership of the Beijing Military
Region, with Li Desheng appointed Commander and Ji Dengkui 2nd Political
Commissar.\footnote{See Hao Mengbi and Duan Haoran, Liushinian,
  pp.~618-20; Zhao Wei, Zhao Ziyang zhuan, p.~178; ``Summary of Chairman
  Mao's Talks (August-September 1971)'', in Stuart Schram (ed.), Mao
  Tse-tung Unrehearsed, Talks and letters: 1956-71 (Harmondsworth:
  Penguin, 1974), p.~295. 见 Hao Mengbi and Duan Haoran,
  \emph{Liushinian}，第618-620页；Zhao Wei, \emph{Zhao Ziyang
  zhuan}，第178页；《毛主席谈话纪要（1971年8-9月）》，载 Stuart Schram
  编，\emph{Mao Tse-tung Unrehearsed, Talks and letters:
  1956-71}（Harmondsworth: Penguin，1974年），第295页。} In the twelve
months between the Lushan plenum and Lin's death on the Mongolian
steppes, Mao turned up the heat on Lin Biao's supporters.\footnote{For
  accounts of Lin's alleged plot, see Liu Huinian et al., ``The
  Bankruptcy of Lin Biao's counter-revolutionary coup d'etat'', Xinhua,
  November 23,1980, SWB/FE/6585/C/4-17, Jin Chunming, ``Wenhua
  dageming'', pp.~168-86; Hua Fang, ``Lin Biao's Abortive
  Counter-revolutionary Coup d'etat'', BR, No.~51 (December 22,1980),
  pp.~19-28. For the reminiscences of China's then ambassador in
  Mongolia, where the aircraft carrying Lin, his wife, son and other
  military officials crashed, see Xu Wenyi, ``A Special Mission History
  Entrusted to Me'', BR, No.~21 (May 23-29,1988), pp.~22-6, and No.~22
  (May 30-June 5 1988), pp.~19-23; Zhang Yuwen, 温都尔汗爆炸记 (A Record
  of the Crash in Ondor Haan), (Guiyang: Guizhou Renmin Chubanshe,
  1988); and the sources cited in n.~9. 关于林彪所谓阴谋的叙述，见 Liu
  Huinian
  等，《林彪反革命政变的破产》，新华社，1980年11月23日，SWB/FE/6585/C/4-17；Jin
  Chunming, \emph{``Wenhua dageming''}，第168-186页；Hua
  Fang，《林彪未遂反革命政变》，\emph{BR}
  第51期（1980年12月22日），第19-28页。关于当时中国驻蒙古大使的回忆，林彪及其妻、子和其他军官乘坐的飞机在那里坠毁，见
  Xu Wenyi，《交给我的一项特殊使命史》，\emph{BR}
  第21期（1988年5月23-29日），第22-26页，及第22期（1988年5月30日-6月5日），第19-23页；Zhang
  Yuwen，《温都尔汗爆炸记》（贵阳：贵州人民出版社，1988年）；以及注9所列资料。}
Lin's military followers in the province came under direct and indirect
attack. In Zhejiang, the 8th enlarged plenum of the ZPRC met from 23
November to 5 December 1970 to discuss the communique of the Lushan
plenum and the relevant central directives arising from it.\footnote{ZJRB,
  December 7,1970, pp.~1, 2. \emph{ZJRB}，1970年12月7日，第1、2版。} A
Zhejiang Daily editorial of 7 December entitled ``Leading Cadres should
study Philosophy more Conscientiously'' initiated the critique of Chen
Boda's ``idealism'' and ``metaphysics'' in the province.\footnote{ZJRB,
  December 7,1970, p.~2. \emph{ZJRB}，1970年12月7日，第2版。} Jinjian
brigade and the Jiangshan cement works, both in Jiangshan county
(江山县), south-western Zhejiang, were promoted as nation-wide models in
this philosophical campaign. Jiang Ruwang (姜汝旺), party secretary of
Jinjian brigade, became famous for his role in the campaign, which Kang
Sheng promoted as an example of peasants studying Mao Zedong
Thought.\footnote{See the documents in URS, 61:4 (October 13,1970),
  pp.~42-56. 见 \emph{URS}, 61:4（1970年10月13日）中的文件，第42-56页。}
Other articles appeared in the local press in response to Mao's request
for party members to study Marxism more assiduously.\footnote{See the
  article by Lai Keke in ZJRB, February 3,1971, p.~2. 见赖可可在
  \emph{ZJRB} 1971年2月3日第2版发表的文章。} There is evidence at this
time that Xu Shiyou, undoubtedly with the permission and encouragement
of the center, moved to monitor and curtail the activities of Nan Ping
and Chen Liyun in Zhejiang and thus check Lin Biao's influence in the
region. In December 1970 the Zhejiang Production and Construction Corps
was established under the command of the Nanjing Military
Region.\footnote{ZJRB, December 23,1970, p.~1 reported a meeting
  concerning the PLA on December 21. See also, CNS, 449 (December
  21,1972). \emph{ZJRB}
  1970年12月23日第1版报道了12月21日有关解放军的一次会议。另见
  \emph{CNS}, 449（1972年12月21日）。} The leaders of this
newly-established corps, whose task was to employ educated youth and
demobilized soldiers on major construction projects, were Political
Commissar Xia Qi (夏琦), who also was appointed Deputy Political
Commissar of the ZPMD, and Commander Liu Ang (刘昂), who was made
Deputy-commander of the ZPMD. Xia Qi was a close protégé of Xu Shiyou.
In an editorial published at the end of January 1971, reference was made
to the ``three onlys'' concept.\footnote{ZJRB, January 31,1971, p.~2.
  \emph{ZJRB}，1971年1月31日，第2版。} The editorial directed a
thinly-veiled warning toward those who departed from the ``three onlys''
and held evil thoughts (邪念) concerning Mao's leadership, who owed
allegiance to a ``second leadership center'' and who applied criteria
other than Mao Zedong Thought to distinguish right from wrong. A
contemporary observer, commenting upon these words, raised the
possibility that reference to the second leadership center may have
referred to supporters of Liu Shaoqi.\footnote{CNS, 355, pp.~8-9.
  \emph{CNS}, 355，第8-9页。} Another analyst suggested that the
Commander of the Nanjing Military Region, Xu Shiyou, was the target
because of his policy disagreements with the central military
establishment.\footnote{See Domes, China after the Cultural Revolution,
  pp.~93-94. 见 Domes, \emph{China after the Cultural
  Revolution}，第93-94页。} In retrospect it is clear that the warning
was directed against provincial supporters of Chen Boda and by
implication, Lin Biao, although the editorial specifically mentioned Lin
as the deputy-leader of the party center headed by Mao. This formula was
standard practice at the beginning of 1971, however. Chen Liyun's Unit
7350 was also subjected to criticism. A report of January 1971 referred
pointedly to its alleged ``arrogance and self-complacency'' (居功自恃)
and inability to listen to other views (一言堂). The unit was advised to
``promote modesty and prudence''. The report continued:
1970年9月全会刚结束，毛便采取行动削弱林的权力。他指示党中央成立直属政治局的组织宣传组，由康生任组长，江青、张春桥、姚文元、纪登奎和李德生任成员。林彪的亲信无人被任命进入这一有权机构。由于康生健康不佳，李德生又牵涉军事事务，文化大革命中的文职激进派掌握了这个强力的超级部门。它负责中央组织部、中央党校、《人民日报》、《光明日报》、《红旗》、新华社、中央广播事业局、中央出版和翻译机构、五七干校，以及中华全国总工会、共青团、妇联等中央群众组织。\footnote{Hao
  Mengbi and Duan Haoran, Liushinian, p.~618; Zhong Kan, Kang Sheng
  pingzhuan, p.~433. Hao Mengbi and Duan Haoran,
  \emph{Liushinian}，第618页；Zhong Kan, \emph{Kang Sheng
  pingzhuan}，第433页。}随后，1971年4月7日，李先念、纪登奎和张才千被增补进中央军委办事组（中央军委办事组）；该机构成立于1968年，此前一直由林彪支持者主导。\footnote{For
  information on the formation of this body in 1968, see Jin Chunming, A
  Discussion and Analysis, p.~136, and Tong De, Xinghuo Liaoyuan, p.~30.
  关于该机构1968年成立的资料，见 Jin Chunming, \emph{A Discussion and
  Analysis}，第136页；Tong De, \emph{Xinghuo Liaoyuan}，第30页。}毛形象地把这一系列举措称作甩石头、掺沙子和挖墙脚。这指的是他对陈伯达和林彪军事小集团所写文件和检查的批评，这些材料在1971年1月至2月的军委会议上讨论，并最终被认为不充分而遭否定；还包括1971年4月向军委办事组增补新成员，以及1970年12月改组北京军区领导层，由李德生任司令员、纪登奎任第二政委。\footnote{See
  Hao Mengbi and Duan Haoran, Liushinian, pp.~618-20; Zhao Wei, Zhao
  Ziyang zhuan, p.~178; ``Summary of Chairman Mao's Talks
  (August-September 1971)'', in Stuart Schram (ed.), Mao Tse-tung
  Unrehearsed, Talks and letters: 1956-71 (Harmondsworth: Penguin,
  1974), p.~295. 见 Hao Mengbi and Duan Haoran,
  \emph{Liushinian}，第618-620页；Zhao Wei, \emph{Zhao Ziyang
  zhuan}，第178页；《毛主席谈话纪要（1971年8-9月）》，载 Stuart Schram
  编，\emph{Mao Tse-tung Unrehearsed, Talks and letters:
  1956-71}（Harmondsworth: Penguin，1974年），第295页。}从庐山会议到林彪死于蒙古草原的十二个月里，毛不断加大对林彪支持者的压力。\footnote{For
  accounts of Lin's alleged plot, see Liu Huinian et al., ``The
  Bankruptcy of Lin Biao's counter-revolutionary coup d'etat'', Xinhua,
  November 23,1980, SWB/FE/6585/C/4-17, Jin Chunming, ``Wenhua
  dageming'', pp.~168-86; Hua Fang, ``Lin Biao's Abortive
  Counter-revolutionary Coup d'etat'', BR, No.~51 (December 22,1980),
  pp.~19-28. For the reminiscences of China's then ambassador in
  Mongolia, where the aircraft carrying Lin, his wife, son and other
  military officials crashed, see Xu Wenyi, ``A Special Mission History
  Entrusted to Me'', BR, No.~21 (May 23-29,1988), pp.~22-6, and No.~22
  (May 30-June 5 1988), pp.~19-23; Zhang Yuwen, 温都尔汗爆炸记 (A Record
  of the Crash in Ondor Haan), (Guiyang: Guizhou Renmin Chubanshe,
  1988); and the sources cited in n.~9. 关于林彪所谓阴谋的叙述，见 Liu
  Huinian
  等，《林彪反革命政变的破产》，新华社，1980年11月23日，SWB/FE/6585/C/4-17；Jin
  Chunming, \emph{``Wenhua dageming''}，第168-186页；Hua
  Fang，《林彪未遂反革命政变》，\emph{BR}
  第51期（1980年12月22日），第19-28页。关于当时中国驻蒙古大使的回忆，林彪及其妻、子和其他军官乘坐的飞机在那里坠毁，见
  Xu Wenyi，《交给我的一项特殊使命史》，\emph{BR}
  第21期（1988年5月23-29日），第22-26页，及第22期（1988年5月30日-6月5日），第19-23页；Zhang
  Yuwen，《温都尔汗爆炸记》（贵阳：贵州人民出版社，1988年）；以及注9所列资料。}林在省内的军队追随者受到直接和间接攻击。在浙江，省革委会第八次扩大会议于1970年11月23日至12月5日召开，讨论庐山会议公报及由此产生的相关中央指示。\footnote{ZJRB,
  December 7,1970, pp.~1, 2. \emph{ZJRB}，1970年12月7日，第1、2版。}12月7日《浙江日报》发表题为《领导干部要更自觉地学习哲学》的社论，开启了省内对陈伯达``唯心主义''和``形而上学''的批判。\footnote{ZJRB,
  December 7,1970, p.~2. \emph{ZJRB}，1970年12月7日，第2版。}位于浙西南江山县（江山县）的金建大队和江山水泥厂，在这场哲学运动中被宣传为全国典型。金建大队党支部书记姜汝旺（姜汝旺）因在运动中的作用而成名，康生把它宣传为农民学习毛泽东思想的例子。\footnote{See
  the documents in URS, 61:4 (October 13,1970), pp.~42-56. 见
  \emph{URS}, 61:4（1970年10月13日）中的文件，第42-56页。}地方报刊还发表其他文章，回应毛要求党员更努力学习马克思主义的号召。\footnote{See
  the article by Lai Keke in ZJRB, February 3,1971, p.~2. 见赖可可在
  \emph{ZJRB} 1971年2月3日第2版发表的文章。}有证据显示，此时许世友无疑在中央许可和鼓励下，开始监视并限制南萍和陈励耘在浙江的活动，从而遏制林彪在该地区的影响。1970年12月，浙江生产建设兵团在南京军区领导下成立。\footnote{ZJRB,
  December 23,1970, p.~1 reported a meeting concerning the PLA on
  December 21. See also, CNS, 449 (December 21,1972). \emph{ZJRB}
  1970年12月23日第1版报道了12月21日有关解放军的一次会议。另见
  \emph{CNS}, 449（1972年12月21日）。}这一新兵团的任务是在重大建设项目中安置知识青年和复员军人，其领导人为政治委员夏琦（夏琦）和司令员刘昂（刘昂）；夏琦同时被任命为浙江省军区副政治委员，刘昂则被任命为浙江省军区副司令员。夏琦是许世友的亲信。1971年1月底发表的一篇社论提到了``三个只''的概念。\footnote{ZJRB,
  January 31,1971, p.~2. \emph{ZJRB}，1971年1月31日，第2版。}社论对那些背离``三个只''、对毛的领导怀有邪念（邪念）、效忠``第二个领导中心''、并用毛泽东思想以外的标准判断是非的人，发出了含蓄警告。一位当代观察者评论这些文字时提出，所谓第二个领导中心可能指刘少奇的支持者。\footnote{CNS,
  355, pp.~8-9. \emph{CNS}, 355，第8-9页。}另一位分析者则认为，目标是南京军区司令员许世友，因为他同中央军事建制存在政策分歧。\footnote{See
  Domes, China after the Cultural Revolution, pp.~93-94. 见 Domes,
  \emph{China after the Cultural Revolution}，第93-94页。}回头看，这一警告显然针对省内陈伯达支持者，并隐含地针对林彪，尽管社论明确提到林是以毛为首的党中央副领袖。不过，这一表述在1971年初是标准做法。陈励耘的7350部队也受到批评。1971年1月的一份报告明确提到该部队据称``居功自恃''（居功自恃），不能听取不同意见（一言堂）。报告告诫该部队``发扬谦虚谨慎''。报告继续说：

\begin{quote}
Last year {[}1970{]}, some ``three support and two military'' personnel
of the units did not fully understand the new situation, which had
emerged as a result of the new leap forward in socialist revolution and
construction.\footnote{ZJRB, January 17,1971, p.~1; quotation from ZPS,
  January 17,1971, SWB/FE/3596/BII/1.
  \emph{ZJRB}，1971年1月17日，第1版；引文见
  \emph{ZPS}，1971年1月17日，SWB/FE/3596/BII/1。}
去年［1970年］，这些部队中有些``三支两军''人员没有充分理解社会主义革命和建设新跃进所造成的新形势。\footnote{ZJRB,
  January 17,1971, p.~1; quotation from ZPS, January 17,1971,
  SWB/FE/3596/BII/1. \emph{ZJRB}，1971年1月17日，第1版；引文见
  \emph{ZPS}，1971年1月17日，SWB/FE/3596/BII/1。}
\end{quote}

Such criticism was certainly not confined to Zhejiang\footnote{Harry
  Harding, ``China: The Fragmentation of Power'', AS, 12:1 (January
  1972), p.~3; Jiangsusheng dashiji, p.~307. Harry
  Harding，《中国：权力的碎片化》，\emph{AS},
  12:1（1972年1月），第3页；\emph{Jiangsusheng dashiji}，第307页。} but
it was all the more sudden because the provincial leadership had, over
the previous two years or so, blown the military's trumpet hard and
long. From 1969 to 1970 enormous praise had been heaped on the PLA and
effusive adulation extended to its Deputy-commander Lin Biao.\footnote{See
  for example, ZPS, January 14,1970, URS 58: 7, pp.~102-5; ZPS July
  28,1970, URS, 60: 12, pp.~163-66; ZPS, December 25,1970, URS, 62:5,
  pp.~63-66. 例如见 \emph{ZPS}，1970年1月14日，\emph{URS}
  58:7，第102-105页；\emph{ZPS}，1970年7月28日，\emph{URS}
  60:12，第163-166页；\emph{ZPS}，1970年12月25日，\emph{URS}
  62:5，第63-66页。} In May 1970 a conference on the work of ``three
support and two militaries'' units had praised personnel stationed at
Hangzhou University for their ``outstanding work''.\footnote{ZJRB, May
  15,1970. \emph{ZJRB}，1970年5月15日。} By October 1971, after Lin
Biao's demise, PLA cadres performing such tasks were reminded to
``strengthen their concept of the Party and consciously safeguard the
Party's unified leadership''.\footnote{ZJRB, October 16,1971, p.~1.
  \emph{ZJRB}，1971年10月16日，第1版。} These examples illustrate the
turn in the tide of the PLA's fortunes. Intensified attacks on the PLA
for arrogance, complacency and commandism were stepped up in the months
immediately following Lin's disappearance.\footnote{Parris H. Chang
  ``Regional Military Power: The Aftermath of the Cultural Revolution'',
  AS, 12: 12 (1972), p.~1010. Parris H.
  Chang，《地区军事权力：文化大革命的后果》，\emph{AS},
  12:12（1972年），第1010页。} Policy issues were also involved in the
elite conflict. After the 9th Congress a militarization of society had
occurred through the mass mobilization of labor for rural developmental
projects\footnote{See Domes, China after the Cultural Revolution, ch.~4;
  Dennis Woodward, ``\,`Two Line Struggle' in Agriculture'', in Brugger,
  China: The Impact of the Cultural Revolution, pp.~160-61; Bill
  Brugger, China: Radicalism to Revisionism, 1962-1979 (London: Croom
  Helm, 1981), pp.~126-30. 见 Domes, \emph{China after the Cultural
  Revolution}，第4章；Dennis Woodward，《农业中的``两条路线斗争''》，载
  Brugger, \emph{China: The Impact of the Cultural
  Revolution}，第160-161页；Bill Brugger, \emph{China: Radicalism to
  Revisionism, 1962-1979}（伦敦：Croom Helm，1981年），第126-130页。}
and the diversion of huge amounts of investment funds to the interior of
the country for defence purposes.\footnote{See Liu Suinian and Wu
  Qungan, ``Wenhua dageming'' shiqide guomin jingji, pp.~47-8; and Barry
  Naughton, ``The Economy of the Cultural Revolution: Military
  Preparation, Recentralization and Leaps Forward'', paper delivered at
  the Harvard Conference ``New Perspectives in the Cultural
  Revolution'', May 1987. 见 Liu Suinian and Wu Qungan, \emph{``Wenhua
  dageming'' shiqide guomin jingji}，第47-48页；以及 Barry
  Naughton，《文化大革命经济：军事准备、重新集中化与跃进》，1987年5月哈佛会议``文化大革命新视角''提交论文。}
The provincial leadership in Zhejiang based its radical mobilizational
rural policies on the study and application of Mao Zedong Thought and
against the backdrop of such campaigns as the purification of the class
ranks and the attack on capitalist tendencies in the
countryside.\footnote{See the report at the provincial agricultural work
  conference on February 22,1970. ZJRB, February 28,1970, p.~1. See also
  ZPS, January 26,1970, URS, 58: 20, pp.~288-90; ZPS, March 5, 1970,
  URS, 59: 5, pp.~59-63.
  见1970年2月22日省农业工作会议报告，\emph{ZJRB}，1970年2月28日，第1版。另见
  \emph{ZPS}，1970年1月26日，\emph{URS}
  58:20，第288-290页；\emph{ZPS}，1970年3月5日，\emph{URS}
  59:5，第59-63页。} Unlike the earlier Leap Forward of 1958-59, which
had been directed largely by civilian party cadres and where spontaneity
and mass initiative found some room for expression, the ``flying leap''
of 1969-70 was a guided, disciplined unleashing of human energy
carefully controlled by the military. It was the application of
political campaigns previously restricted to the ranks of the PLA to
society at large, particularly its youth. After the Lushan plenum the
radical policy initiatives of 1969-70 were moderated and rolled back by
a coalition of regional military commanders and central civilian cadres.
In particular, the militarization and radicalization of rural policies
was brought to a halt. The forthright, pugnacious Xu Shiyou in Nanjing
was prominent among regional military commanders who spoke out against
egalitarian rural distribution policies. Xu made his views known in
October and November 1970 and again in February 1971.\footnote{See
  Domes, China after the Cultural Revolution, p.~57, ch.~7; Jiangsusheng
  dashiji, pp.~312-13. 见 Domes, \emph{China after the Cultural
  Revolution}，第57页及第7章；\emph{Jiangsusheng dashiji}，第312-313页。}
Xu's attention may well have been drawn to the fact that in Zhejiang and
Jiangxi provinces, both of which fell within his jurisdiction, rural
social policies were the most radical in the country. For example, in
Zhejiang one quarter of all brigades had taken over the important role
of acting as the accounting unit, a function formerly fulfilled by the
production team.\footnote{In Kaihua county this change was carried out
  in 1968. See Kaihua xianzhi, p.~24. 在开化县，这一变化于1968年实施。见
  \emph{Kaihua xianzhi}，第24页。} In Jiangxi, two-thirds of all private
plots were collectivized and half the communes subsumed brigades into
one ownership level.\footnote{Liu Suinian and Wu Qungan, ``Wenhua
  dageming'' shiqide guomin jingji, p.~67. Liu Suinian and Wu Qungan,
  \emph{``Wenhua dageming'' shiqide guomin jingji}，第67页。}
这种批评当然并非浙江独有，\footnote{Harry Harding, ``China: The
  Fragmentation of Power'', AS, 12:1 (January 1972), p.~3; Jiangsusheng
  dashiji, p.~307. Harry Harding，《中国：权力的碎片化》，\emph{AS},
  12:1（1972年1月），第3页；\emph{Jiangsusheng dashiji}，第307页。}但它显得尤其突然，因为在此前大约两年中，省领导层一直大力、长期地为军队歌功颂德。1969至1970年间，解放军受到极高赞扬，其副统帅林彪也受到热烈颂扬。\footnote{See
  for example, ZPS, January 14,1970, URS 58: 7, pp.~102-5; ZPS July
  28,1970, URS, 60: 12, pp.~163-66; ZPS, December 25,1970, URS, 62:5,
  pp.~63-66. 例如见 \emph{ZPS}，1970年1月14日，\emph{URS}
  58:7，第102-105页；\emph{ZPS}，1970年7月28日，\emph{URS}
  60:12，第163-166页；\emph{ZPS}，1970年12月25日，\emph{URS}
  62:5，第63-66页。}1970年5月，一次``三支两军''部队工作会议曾称赞驻杭州大学人员``工作出色''。\footnote{ZJRB,
  May 15,1970. \emph{ZJRB}，1970年5月15日。}到1971年10月林彪败亡之后，从事这类任务的解放军干部被提醒要``增强党的观念，自觉维护党的统一领导''。\footnote{ZJRB,
  October 16,1971, p.~1. \emph{ZJRB}，1971年10月16日，第1版。}这些例子说明了解放军命运潮流的转向。林彪失踪后的数月里，对解放军骄傲、自满和命令主义的攻击加强了。\footnote{Parris
  H. Chang ``Regional Military Power: The Aftermath of the Cultural
  Revolution'', AS, 12: 12 (1972), p.~1010. Parris H.
  Chang，《地区军事权力：文化大革命的后果》，\emph{AS},
  12:12（1972年），第1010页。}精英冲突中也涉及政策问题。九大以后，通过大规模动员劳动力投入农村开发项目，\footnote{See
  Domes, China after the Cultural Revolution, ch.~4; Dennis Woodward,
  ``\,`Two Line Struggle' in Agriculture'', in Brugger, China: The
  Impact of the Cultural Revolution, pp.~160-61; Bill Brugger, China:
  Radicalism to Revisionism, 1962-1979 (London: Croom Helm, 1981),
  pp.~126-30. 见 Domes, \emph{China after the Cultural
  Revolution}，第4章；Dennis Woodward，《农业中的``两条路线斗争''》，载
  Brugger, \emph{China: The Impact of the Cultural
  Revolution}，第160-161页；Bill Brugger, \emph{China: Radicalism to
  Revisionism, 1962-1979}（伦敦：Croom Helm，1981年），第126-130页。}以及为国防目的把大量投资资金转向内地，\footnote{See
  Liu Suinian and Wu Qungan, ``Wenhua dageming'' shiqide guomin jingji,
  pp.~47-8; and Barry Naughton, ``The Economy of the Cultural
  Revolution: Military Preparation, Recentralization and Leaps
  Forward'', paper delivered at the Harvard Conference ``New
  Perspectives in the Cultural Revolution'', May 1987. 见 Liu Suinian
  and Wu Qungan, \emph{``Wenhua dageming'' shiqide guomin
  jingji}，第47-48页；以及 Barry
  Naughton，《文化大革命经济：军事准备、重新集中化与跃进》，1987年5月哈佛会议``文化大革命新视角''提交论文。}社会出现了军事化。浙江省领导层以学习和运用毛泽东思想为基础，并在清理阶级队伍、打击农村资本主义倾向等运动背景下，推行激进的动员型农村政策。\footnote{See
  the report at the provincial agricultural work conference on February
  22,1970. ZJRB, February 28,1970, p.~1. See also ZPS, January 26,1970,
  URS, 58: 20, pp.~288-90; ZPS, March 5, 1970, URS, 59: 5, pp.~59-63.
  见1970年2月22日省农业工作会议报告，\emph{ZJRB}，1970年2月28日，第1版。另见
  \emph{ZPS}，1970年1月26日，\emph{URS}
  58:20，第288-290页；\emph{ZPS}，1970年3月5日，\emph{URS}
  59:5，第59-63页。}不同于1958至1959年早先的大跃进，后者主要由文职党干部领导，自发性和群众主动性尚有一定表达空间；1969至1970年的``飞跃''则是在军队严密控制下，有引导、有纪律地释放人的能量。这是把以往局限于解放军内部的政治运动方法应用到整个社会，尤其是应用到青年身上。庐山会议以后，1969至1970年的激进政策举措被地区军事指挥员和中央文职干部组成的联盟所缓和并回撤。尤其是农村政策的军事化和激进化被叫停。南京的许世友直率而好斗，是公开反对平均主义农村分配政策的地区军事指挥员中突出的一位。许在1970年10月、11月以及1971年2月表达了自己的看法。\footnote{See
  Domes, China after the Cultural Revolution, p.~57, ch.~7; Jiangsusheng
  dashiji, pp.~312-13. 见 Domes, \emph{China after the Cultural
  Revolution}，第57页及第7章；\emph{Jiangsusheng dashiji}，第312-313页。}许很可能注意到，在其辖区内的浙江和江西两省，农村社会政策是全国最激进的。例如在浙江，四分之一的生产大队接管了充当核算单位的重要角色，而这一职能过去由生产队承担。\footnote{In
  Kaihua county this change was carried out in 1968. See Kaihua xianzhi,
  p.~24. 在开化县，这一变化于1968年实施。见 \emph{Kaihua
  xianzhi}，第24页。}在江西，三分之二的自留地被集体化，一半公社把大队合并为一个所有制层级。\footnote{Liu
  Suinian and Wu Qungan, ``Wenhua dageming'' shiqide guomin jingji,
  p.~67. Liu Suinian and Wu Qungan, \emph{``Wenhua dageming'' shiqide
  guomin jingji}，第67页。}

\section{The Lin Biao Affair in
Zhejiang}\label{the-lin-biao-affair-in-zhejiang}

\section{浙江的林彪事件}\label{ux6d59ux6c5fux7684ux6797ux5f6aux4e8bux4ef6}

The Lin Biao affair impacted seriously on elite politics in Zhejiang.
Chen Liyun, the second most senior cadre in the province, was a
principal participant in the plot to assassinate Mao. According to
documents and accounts published after the incident, Chen Liyun was one
of the members of Lin Liguo's (Lin Biao's son) ``joint fleet''. This
clandestine organization consisted of young and middle-aged military
officials who owed their advancement in the Cultural Revolution to Lin
Biao. Chen acted as the fleet's station head in Hangzhou and attended
meetings of the fleet there and in Shanghai. As an Air Force political
officer, Chen came under the leadership of Lin Biao's close collaborator
and confidant, Wu Faxian. At Lin Biao's suggestion, Wu had promoted Lin
Liguo to a substantial post within this arm of the PLA. Thus, Chen Liyun
had organizational and personal ties with the Lin Biao group. At the
Lushan plenum he had spoken out, at Wu Faxian's request, in favor of
Chen Boda's ``theory of genius'' and in support of the proposal to
reestablish the post of state Chairman.\footnote{The following details
  of Chen's part in the events of 1970-71 are drawn from Gao Gao and Yan
  Jiaqi, ``Wenhua dageming'', pp.~342-76, passim; Hu Hua (ed.), Zhongguo
  shehuizhuyi geming, pp.~304-06; Hao Mengbi and Duan Haoran,
  Liushinian, pp.~612-4, 620-1; Shao Yihai, ``联合舰队的覆灭'' (The
  Annihilation of the ``Joint Fleet''), Zhou Ming (ed.), Lishi zai
  zheli, Vol. 2, pp.~108-09; Da Ying, ``九一三事件始末记'' (A Record of
  ``13 September'' Incident from Beginning to End), Lishi zai zheli,
  Vol. 2, pp.~38; CCP CC, zhongfa (1972), No.~4, ``Struggle to Smash the
  Lin-Chen {[}Boda{]} Anti-Party Clique's Counterrevolutionary Coup'',
  in Kau, The Lin Piao Affair, pp.~86, 90-1, 93-4; ``Communique of the
  CCP CC concerning Lin Piao's `September 12' Anti-Party Incident
  (Abridged)'', in Kau, The Lin Piao Affair, p.~70; Shao Yihai,
  林彪王朝黑幕 (Machinations at Lin Biao's Court), (Chengdu: Sichuan
  Wenyi Chubanshe, 1988), pp.~81-82, 97, 115-20, 132-41, 183-85, 211-13;
  CCP, ZPC, ``The Great Banner of Chairman Mao will always guide the
  People of Zhejiang in their Victorious Advance'', HZRB, September 8,
  1977; Hua Fang, ``Lin Biao's Abortive Counter-Revolutionary Coup
  d'etat'', pp.~23-5; Da Ying, Lishi zai zheli, pp.~162-63, 172-73; Ye
  Yonglie, Wang Hongwen xingshuailu, pp.~374-75; and the sources cited
  in n.~9 and n.~18. Chen and other leaders of the Zhejiang provincial
  committee ingratiated themselves with Lin Biao and his family by
  taking charge of the construction of a luxury compound of villas and
  offices on Three Terrace Mountain, on the outskirts of Hangzhou. It
  had five villas, underground defence facilities, an indoor
  swimming-pool, shops and expansive and attractive grounds covering
  over 20 hectares. The article alleges, and the accusation circulates
  widely in Hangzhou, that all the workers on the site were later
  executed in secret. The compound is now known as the Zhejiang
  Guesthouse. For a recent article on the project, whose original name
  was the ``April 1970 engineering project'' (七O.四工程) see Cheng
  Ying, ``林彪行宫首次曝光'' (Lin Biao's villa exposed), 九十年代月刊
  (The Nineties Monthly), No.~1,1990, pp.~76-80. The complex cost more
  than 20 million yuan, reportedly a greater sum than the total state
  outlay on housing in Hangzhou for the twenty years 1949 to 1969. See
  RMRB, September 16,1978, p.~1.
  以下关于陈在1970-71年事件中所扮演角色的细节，采自 Gao Gao and Yan
  Jiaqi, \emph{``Wenhua
  dageming''}，第342-376页等处；胡华编，\emph{Zhongguo shehuizhuyi
  geming}，第304-306页；Hao Mengbi and Duan Haoran,
  \emph{Liushinian}，第612-614、620-621页；邵一海，《联合舰队的覆灭》，载周明编，\emph{Lishi
  zai zheli} 第2卷，第108-109页；Da
  Ying，《九一三事件始末记》，\emph{Lishi zai zheli}
  第2卷，第38页；中共中央中发（1972）4号文件《粉碎林陈［伯达］反党集团反革命政变的斗争》，载
  Kau, \emph{The Lin Piao
  Affair}，第86、90-91、93-94页；《中共中央关于林彪``九一二''反党事件的通知（节录）》，载
  Kau, \emph{The Lin Piao
  Affair}，第70页；邵一海，《林彪王朝黑幕》（成都：四川文艺出版社，1988年），第81-82、97、115-120、132-141、183-185、211-213页；中共浙江省委，《毛主席的伟大旗帜永远指引浙江人民胜利前进》，\emph{HZRB}，1977年9月8日；Hua
  Fang，《林彪未遂反革命政变》，第23-25页；Da Ying, \emph{Lishi zai
  zheli}，第162-163、172-173页；Ye Yonglie, \emph{Wang Hongwen
  xingshuailu}，第374-375页；以及注9和注18所列资料。陈和浙江省委其他领导人为讨好林彪及其家人，负责在杭州郊外三台山建造一处豪华别墅和办公建筑群。该处有五栋别墅、地下防御设施、室内游泳池、商店和20多公顷开阔优美的场地。文章称，且这一说法在杭州广泛流传，工地所有工人后来都被秘密处决。该建筑群现称浙江宾馆。关于该项目的近文，原名``七O.四工程''，见
  Cheng
  Ying，《林彪行宫首次曝光》，《九十年代月刊》第1期（1990年），第76-80页。该建筑群耗资逾2,000万元，据称超过1949至1969年二十年间杭州住房国家总投资。见《人民日报》，1978年9月16日，第1版。}
The denouement of Lin Liguo's plot occurred in August/September 1971.
Mao headed south from Beijing to secure the support of key regional
military leaders, in the event of an open confrontation with his deputy
and to explain the nature and severity of the conflict which had been
exacerbated since the Lushan plenum. On 3 September the Chairman's train
arrived in Hangzhou and Mao immediately let Chen Liyun know that he was
aware of his dealings with the Lin Biao group. Mao reportedly said to
Chen in a tone of disgust:
林彪事件严重影响了浙江精英政治。省内第二号高级干部陈励耘，是刺杀毛阴谋的主要参与者。据事件后公布的文件和叙述，陈励耘是林立果（林彪之子）``联合舰队''成员之一。这个秘密组织由一批青年和中年军官组成，他们在文化大革命中的升迁归功于林彪。陈担任该舰队杭州站负责人，并在杭州和上海参加舰队会议。作为空军政治干部，陈受林彪亲密合作者和心腹吴法宪领导。按林彪建议，吴曾把林立果提拔到解放军这一军种中的重要岗位。因此，陈励耘同林彪集团既有组织联系，也有个人联系。庐山会议上，他应吴法宪要求发言，支持陈伯达的``天才论''，并支持恢复国家主席职位的提议。\footnote{The
  following details of Chen's part in the events of 1970-71 are drawn
  from Gao Gao and Yan Jiaqi, ``Wenhua dageming'', pp.~342-76, passim;
  Hu Hua (ed.), Zhongguo shehuizhuyi geming, pp.~304-06; Hao Mengbi and
  Duan Haoran, Liushinian, pp.~612-4, 620-1; Shao Yihai,
  ``联合舰队的覆灭'' (The Annihilation of the ``Joint Fleet''), Zhou
  Ming (ed.), Lishi zai zheli, Vol. 2, pp.~108-09; Da Ying,
  ``九一三事件始末记'' (A Record of ``13 September'' Incident from
  Beginning to End), Lishi zai zheli, Vol. 2, pp.~38; CCP CC, zhongfa
  (1972), No.~4, ``Struggle to Smash the Lin-Chen {[}Boda{]} Anti-Party
  Clique's Counterrevolutionary Coup'', in Kau, The Lin Piao Affair,
  pp.~86, 90-1, 93-4; ``Communique of the CCP CC concerning Lin Piao's
  `September 12' Anti-Party Incident (Abridged)'', in Kau, The Lin Piao
  Affair, p.~70; Shao Yihai, 林彪王朝黑幕 (Machinations at Lin Biao's
  Court), (Chengdu: Sichuan Wenyi Chubanshe, 1988), pp.~81-82, 97,
  115-20, 132-41, 183-85, 211-13; CCP, ZPC, ``The Great Banner of
  Chairman Mao will always guide the People of Zhejiang in their
  Victorious Advance'', HZRB, September 8, 1977; Hua Fang, ``Lin Biao's
  Abortive Counter-Revolutionary Coup d'etat'', pp.~23-5; Da Ying, Lishi
  zai zheli, pp.~162-63, 172-73; Ye Yonglie, Wang Hongwen xingshuailu,
  pp.~374-75; and the sources cited in n.~9 and n.~18. Chen and other
  leaders of the Zhejiang provincial committee ingratiated themselves
  with Lin Biao and his family by taking charge of the construction of a
  luxury compound of villas and offices on Three Terrace Mountain, on
  the outskirts of Hangzhou. It had five villas, underground defence
  facilities, an indoor swimming-pool, shops and expansive and
  attractive grounds covering over 20 hectares. The article alleges, and
  the accusation circulates widely in Hangzhou, that all the workers on
  the site were later executed in secret. The compound is now known as
  the Zhejiang Guesthouse. For a recent article on the project, whose
  original name was the ``April 1970 engineering project'' (七O.四工程)
  see Cheng Ying, ``林彪行宫首次曝光'' (Lin Biao's villa exposed),
  九十年代月刊 (The Nineties Monthly), No.~1,1990, pp.~76-80. The
  complex cost more than 20 million yuan, reportedly a greater sum than
  the total state outlay on housing in Hangzhou for the twenty years
  1949 to 1969. See RMRB, September 16,1978, p.~1.
  以下关于陈在1970-71年事件中所扮演角色的细节，采自 Gao Gao and Yan
  Jiaqi, \emph{``Wenhua
  dageming''}，第342-376页等处；胡华编，\emph{Zhongguo shehuizhuyi
  geming}，第304-306页；Hao Mengbi and Duan Haoran,
  \emph{Liushinian}，第612-614、620-621页；邵一海，《联合舰队的覆灭》，载周明编，\emph{Lishi
  zai zheli} 第2卷，第108-109页；Da
  Ying，《九一三事件始末记》，\emph{Lishi zai zheli}
  第2卷，第38页；中共中央中发（1972）4号文件《粉碎林陈［伯达］反党集团反革命政变的斗争》，载
  Kau, \emph{The Lin Piao
  Affair}，第86、90-91、93-94页；《中共中央关于林彪``九一二''反党事件的通知（节录）》，载
  Kau, \emph{The Lin Piao
  Affair}，第70页；邵一海，《林彪王朝黑幕》（成都：四川文艺出版社，1988年），第81-82、97、115-120、132-141、183-185、211-213页；中共浙江省委，《毛主席的伟大旗帜永远指引浙江人民胜利前进》，\emph{HZRB}，1977年9月8日；Hua
  Fang，《林彪未遂反革命政变》，第23-25页；Da Ying, \emph{Lishi zai
  zheli}，第162-163、172-173页；Ye Yonglie, \emph{Wang Hongwen
  xingshuailu}，第374-375页；以及注9和注18所列资料。陈和浙江省委其他领导人为讨好林彪及其家人，负责在杭州郊外三台山建造一处豪华别墅和办公建筑群。该处有五栋别墅、地下防御设施、室内游泳池、商店和20多公顷开阔优美的场地。文章称，且这一说法在杭州广泛流传，工地所有工人后来都被秘密处决。该建筑群现称浙江宾馆。关于该项目的近文，原名``七O.四工程''，见
  Cheng
  Ying，《林彪行宫首次曝光》，《九十年代月刊》第1期（1990年），第76-80页。该建筑群耗资逾2,000万元，据称超过1949至1969年二十年间杭州住房国家总投资。见《人民日报》，1978年9月16日，第1版。}林立果阴谋的结局发生在1971年8月至9月。毛从北京南下，以便在同副手公开对抗时确保关键地区军事领导人的支持，并说明自庐山会议以来激化的冲突的性质和严重性。9月3日，主席专列抵达杭州，毛立即让陈励耘知道，他了解陈同林彪集团的往来。据说毛以厌恶的语气对陈说：

\begin{quote}
What are your relations with Wu Faxian? At Lushan Wu Faxian sought out
several people: there was you, Chen Liyun, Shanghai's Wang Weiguo
{[}Political Commissar of the Fourth Air Force{]} and someone from the
Navy. What were you up to?
你同吴法宪是什么关系？在庐山，吴法宪找过几个人：有你陈励耘，有上海的王维国［空四军政治委员］，还有一个海军的人。你们在搞什么？
\end{quote}

It is not surprising that Chen was speechless at this verbal onslaught
and equivocated in his answer. Another version of this confrontation,
and one that shows less certainty on Mao's part, reports that when Mao
received Chen, Nan Ping, Xiong Yingtang and Bai Zongshan after his
arrival in Hangzhou he abruptly quizzed Chen with the words: what have
you been up to? Did Wu Faxian call you to a meeting at Lushan? On the
evening of 8 September, Yu Xinye, Deputy-director of the Air Force
political department, arrived in Hangzhou in order to find out from Chen
information concerning Mao's itinerary and the content of the Chairman's
talks with the Zhejiang leadership. That same night, as an indication
that he would not be leaving in the immediate future, Mao instructed
that his train be moved east from Hangzhou to Shaoxing. A ``responsible
comrade'' from the CC office in Mao's entourage, looked everywhere to
inform Chen Liyun of Mao's orders. Chen, who was in charge of security
in Hangzhou, was nowhere to be found because he was in conference with
Yu Xinye. This piece of evidence points to the likelihood that Mao was
no more than suspicious of Chen. Otherwise why would he have allowed him
to remain in charge of his security arrangements in Hangzhou? Suddenly,
on the afternoon of 10 September, Mao ordered the train back to
Hangzhou. He departed immediately for Shanghai, specifically requesting
that Chen Liyun not see him off as protocol required. It is probable
that while he was in Hangzhou Mao made determined efforts to ensure that
Nan Ping and Xiong Yingtang's 20th Army, which the ``Outline of `Project
571'\,'' drafted by Lin Liguo and his associates marked down as an
auxiliary force, remained loyal. In this endeavor he seems to have
succeeded. Immediately after the crash of Lin's aircraft Chen Liyun was
arrested for isolated examination. He had been captured after a
helicopter in which he was passenger was intercepted and forced down by
jet fighters. In the ensuing struggle he was seriously wounded. Chen was
later described as a sworn follower (死党) of the Lin Biao group. By
contrast, Nan Ping and Xiong Yingtang remained in office until May 1972
and after their removal (with Xiong being transferred and demoted to PLA
units in the Chengdu Military Region) were described, less severely, as
Lin Biao's agents (代理人) in Zhejiang.
陈在这番言辞攻击下哑口无言，并在回答中支吾其词，并不令人意外。关于这次对峙的另一版本显示毛本人不那么确信；该版本称，毛到杭州后接见陈励耘、南萍、熊应堂和白宗善时，突然质问陈：你搞了些什么？吴法宪有没有叫你到庐山开会？9月8日晚，空军政治部副主任于新野抵达杭州，想从陈那里了解毛的行程以及主席同浙江领导层谈话的内容。当晚，为表明自己短期内不会离开，毛指示把专列从杭州向东调往绍兴。毛随行中央办公厅的一位``负责同志''到处寻找陈励耘，以传达毛的命令。负责杭州安全的陈却因正同于新野开会而不知去向。这一证据说明，毛很可能只是怀疑陈。否则，他为何还会让陈继续负责自己在杭州的安全安排？9月10日下午，毛突然命令专列返回杭州。他随即启程赴上海，并特别要求陈励耘不要按礼节前来送行。毛在杭州期间，很可能曾坚决努力确保南萍和熊应堂的二十军保持忠诚；林立果及其同伙起草的《``571工程''纪要》把该军列为辅助力量。看来毛在这方面成功了。林彪飞机坠毁后，陈励耘立即被逮捕并隔离审查。他乘坐的直升机被喷气战斗机拦截并迫降后被捕；在随后的搏斗中，他受了重伤。陈后来被称为林彪集团的死党（死党）。相比之下，南萍和熊应堂一直留任到1972年5月；他们被免职后，熊被调往成都军区解放军部队并降职，二人受到的定性较轻，被称为林彪在浙江的代理人（代理人）。

\section{Fallout from the Lin Biao
Affair}\label{fallout-from-the-lin-biao-affair}

\section{林彪事件的后果}\label{ux6797ux5f6aux4e8bux4ef6ux7684ux540eux679c}

Zhou Enlai, with the resigned acquiescence of a shattered Mao
Zedong,\footnote{Mao's health reportedly suffered greatly from the shock
  induced by Lin's defection. After U.S. President Richard Nixon met Mao
  on February 21,1972, he recorded his impressions of Mao's feebleness,
  sallow complexion and difficulty in speaking, which he attributed to
  the possible after-effects of apoplexy, as well as the Chairman's
  expressionless face and detached eyes. See Gao Gao and Yan Jiaqi,
  ``Wenhua dageming'', p.~469. The Chairman remained in bed for weeks on
  end. See Frederick C. Teiwes, ``Mao and His Lieutenants'', Australian
  Journal of Chinese Affairs, 19 and 20 (January and July 1988), p.~35.
  据称，林彪叛逃造成的冲击严重损害了毛的健康。美国总统理查德·尼克松1972年2月21日会见毛后，记录了他对毛虚弱、面色蜡黄、说话困难的印象，并把这些归因于可能的中风后遗症，同时还写到主席面无表情和目光游离。见
  Gao Gao and Yan Jiaqi, \emph{``Wenhua
  dageming''}，第469页。主席连续数周卧床。见 Frederick C.
  Teiwes，《毛和他的副手们》，\emph{Australian Journal of Chinese
  Affairs}，第19、20期（1988年1月和7月），第35页。} moved decisively to
purge the central military apparatus of Lin's followers. Lin's closest
and most senior supporters were dismissed and sent back to their home
provinces. On 3 October 1971 the CC MAC Executive Group was disbanded
and replaced by administrative conferences (办公会议), which were
presided over by Ye Jianying.\footnote{See Hu Hua (ed.), Zhongguo
  shehuizhuyi geming, p.~310; Gao Gao and Yan Jiaqi ``Wenhua dageming'',
  pp.~469-70. 见胡华编，\emph{Zhongguo shehuizhuyi geming}，第310页；Gao
  Gao and Yan Jiaqi, \emph{``Wenhua dageming''}，第469-470页。} On the
same day, Li Desheng was appointed Director of the PLA General Political
Department.\footnote{Dangshi tongxun, no. 71 (September 15,1983), p.~31.
  \emph{Dangshi tongxun} 第71期（1983年9月15日），第31页。} The Chairman
also moved to appease and mollify military leaders who had opposed Lin
in February 1967. In a reception in Chengdu on 14 November 1971, Mao
rehabilitated veteran marshalls and experienced administrators of the
State Council who had been held responsible for the ``February
Countercurrent'' and shifted blame for their persecution onto Lin Biao
and Chen Boda.\footnote{Hao Mengbi and Duan Haoran, Liushinian, p.~624.
  Hao Mengbi and Duan Haoran, \emph{Liushinian}，第624页。} Not
surprisingly, the provincial media reflected the uncertainty of local
administrations as how to respond to these dramatic and potentially
threatening developments. An issue of China News Summary noted the
scarcity of party meetings, conferences, newspaper editorials and public
appearances of provincial leaders in the immediate aftermath of the
affair. For Zhejiang the bi-monthly journal recorded two meetings and
eight editorials in the period from mid-September 1971 until early April
1972.\footnote{CNS, 423 (June 22,1972), pp.~1-4. \emph{CNS},
  423（1972年6月22日），第1-4页。} However, this report was somewhat
inaccurate. In January 1972 alone, Xiong Yingtang in particular
maintained a high profile with eight public appearances during the
month, three of which were in the company of Nan Ping.\footnote{Reference
  Aid. Appearances and Activities of Leading Personalities of the PRC,
  1972. See also Xiong's article denouncing Chen Boda's ``theory of
  genius'' (天才论) in ZJRB, January 6,1972, pp.~1, 3. \emph{Reference
  Aid. Appearances and Activities of Leading Personalities of the PRC,
  1972}。另见熊在 \emph{ZJRB}
  1972年1月6日第1、3版发表的批判陈伯达``天才论''（天才论）的文章。} At
the beginning of the campaign against Lin Biao, which went under the
vague title of ``criticize revisionism and rectify work-style''
(批修整风), the Zhejiang provincial leadership took up the cause with
some zeal, hoping thereby to avoid political oblivion.\footnote{Domes,
  China after the Cultural Revolution, pp.~120-21, 124, notes that
  Zhejiang, together with Hunan and Xinjiang, were active in the
  campaign from the very beginning. Domes, \emph{China after the
  Cultural
  Revolution}，第120-121、124页指出，浙江同湖南和新疆一样，从运动一开始就很积极。}
In fact, the ZPC had issued a circular on 7 September 1971, before Lin's
downfall, calling on the people to sing the Internationale and the PLA
slogans known as the ``Three Rules and the Eight Points for
Attention''.\footnote{ZPS November 26, 1971 SWB/FE/3854/BII/14-15.
  \emph{ZPS}，1971年11月26日，SWB/FE/3854/BII/14-15。} The words of the
songs were published on page one of Zhejiang Daily on 14 September 1971,
the day after Lin's defection. In December a meeting was held in
Hangzhou to study the revolutionary songs.\footnote{ZPS, December 6,
  1971, SWB/FE/3865/BII/13-14.
  \emph{ZPS}，1971年12月6日，SWB/FE/3865/BII/13-14。} Mao had initiated
this singing campaign during his August/September tour of the provinces,
no doubt to remind local military officials who really was the founder
and leader of the PLA.\footnote{``Summary of Chairman Mao's Talks'', in
  Schram, Mao Tse-tung Unrehearsed, pp.~297-98. 《毛主席谈话纪要》，载
  Schram, \emph{Mao Tse-tung Unrehearsed}，第297-298页。} At the subdued
National Day rally held on 30 September 1971 in Hangzhou, the speakers
were ZPMD commander Xiong Yingtang and Chen Liyun's colleague of Unit
7350, Commander Bai Zongshan (白宗善). Bai's appearance may have
signified that while Chen Liyun had been involved in a serious
conspiracy, his senior colleagues in Unit 7350 and the unit as a whole
were not necessarily implicated. Bai's speech was certainly more
self-critical than Xiong's contribution. He demanded of all PLA units
stationed in the province that they implement the traditional codes of
behavior and ``combat arrogance and complacency'', and conscientiously
carry out the ``three supports and two militaries'' tasks. He went on to
stress that ``we must resolutely obey the centralized leadership of the
local party committees''.\footnote{ZJRB, October 1,1971, p.~2.
  \emph{ZJRB}，1971年10月1日，第2版。} A report of December 1971 called
on both the PLA and mass organization representatives serving on
revolutionary committees to obey the unified leadership of the
party.\footnote{ZPS, December 18,1971, URS, 65: 25, pp.~340-43.
  \emph{ZPS}，1971年12月18日，\emph{URS} 65:25，第340-343页。}
Nevertheless, the PLA-dominated leadership in Zhejiang was to face a
major purge in 1972, to the extent that China News Summary in early 1973
could claim that of all provinces it had ``suffered worst'' from the
anti-Lin Biao Campaign.\footnote{CNS, 455 (February 15,1973), pp.~3-4.
  \emph{CNS}, 455（1973年2月15日），第3-4页。}
周恩来在遭受重创的毛泽东无奈默许下，\footnote{Mao's health reportedly
  suffered greatly from the shock induced by Lin's defection. After U.S.
  President Richard Nixon met Mao on February 21,1972, he recorded his
  impressions of Mao's feebleness, sallow complexion and difficulty in
  speaking, which he attributed to the possible after-effects of
  apoplexy, as well as the Chairman's expressionless face and detached
  eyes. See Gao Gao and Yan Jiaqi, ``Wenhua dageming'', p.~469. The
  Chairman remained in bed for weeks on end. See Frederick C. Teiwes,
  ``Mao and His Lieutenants'', Australian Journal of Chinese Affairs, 19
  and 20 (January and July 1988), p.~35.
  据称，林彪叛逃造成的冲击严重损害了毛的健康。美国总统理查德·尼克松1972年2月21日会见毛后，记录了他对毛虚弱、面色蜡黄、说话困难的印象，并把这些归因于可能的中风后遗症，同时还写到主席面无表情和目光游离。见
  Gao Gao and Yan Jiaqi, \emph{``Wenhua
  dageming''}，第469页。主席连续数周卧床。见 Frederick C.
  Teiwes，《毛和他的副手们》，\emph{Australian Journal of Chinese
  Affairs}，第19、20期（1988年1月和7月），第35页。}果断清洗中央军事机构中的林彪追随者。林最亲近、级别最高的支持者被撤职并送回原籍省份。1971年10月3日，中央军委办事组被解散，由叶剑英主持的办公会议（办公会议）取而代之。\footnote{See
  Hu Hua (ed.), Zhongguo shehuizhuyi geming, p.~310; Gao Gao and Yan
  Jiaqi ``Wenhua dageming'', pp.~469-70. 见胡华编，\emph{Zhongguo
  shehuizhuyi geming}，第310页；Gao Gao and Yan Jiaqi, \emph{``Wenhua
  dageming''}，第469-470页。}同一天，李德生被任命为解放军总政治部主任。\footnote{Dangshi
  tongxun, no. 71 (September 15,1983), p.~31. \emph{Dangshi tongxun}
  第71期（1983年9月15日），第31页。}主席还采取行动安抚那些在1967年2月反对林彪的军事领导人。1971年11月14日在成都的一次接见中，毛为曾被认定对``二月逆流''负责的老帅和国务院资深行政干部平反，并把他们遭受迫害的责任转嫁到林彪和陈伯达身上。\footnote{Hao
  Mengbi and Duan Haoran, Liushinian, p.~624. Hao Mengbi and Duan
  Haoran, \emph{Liushinian}，第624页。}省级媒体反映出地方行政机构面对这些剧烈且可能构成威胁的发展时不知如何回应，这并不奇怪。《中国新闻摘要》一期指出，事件刚发生后，党内会议、工作会议、报纸社论以及省领导公开露面都很少。该半月刊记录称，1971年9月中旬至1972年4月初，浙江只有两次会议和八篇社论。\footnote{CNS,
  423 (June 22,1972), pp.~1-4. \emph{CNS},
  423（1972年6月22日），第1-4页。}不过，这一报告并不十分准确。仅1972年1月，熊应堂尤其保持高调，一个月内公开露面八次，其中三次与南萍同场。\footnote{Reference
  Aid. Appearances and Activities of Leading Personalities of the PRC,
  1972. See also Xiong's article denouncing Chen Boda's ``theory of
  genius'' (天才论) in ZJRB, January 6,1972, pp.~1, 3. \emph{Reference
  Aid. Appearances and Activities of Leading Personalities of the PRC,
  1972}。另见熊在 \emph{ZJRB}
  1972年1月6日第1、3版发表的批判陈伯达``天才论''（天才论）的文章。}在反林彪运动开始时，这场运动以含糊的``批修整风''（批修整风）为名，浙江省领导层颇为热心地投入其中，希望借此避免政治湮没。\footnote{Domes,
  China after the Cultural Revolution, pp.~120-21, 124, notes that
  Zhejiang, together with Hunan and Xinjiang, were active in the
  campaign from the very beginning. Domes, \emph{China after the
  Cultural
  Revolution}，第120-121、124页指出，浙江同湖南和新疆一样，从运动一开始就很积极。}事实上，浙江省委在林彪垮台前的1971年9月7日已发布通知，号召群众唱《国际歌》和解放军的``三大纪律八项注意''。\footnote{ZPS
  November 26, 1971 SWB/FE/3854/BII/14-15.
  \emph{ZPS}，1971年11月26日，SWB/FE/3854/BII/14-15。}歌词于1971年9月14日，即林彪叛逃次日，刊登在《浙江日报》第一版。12月，杭州召开会议学习革命歌曲。\footnote{ZPS,
  December 6, 1971, SWB/FE/3865/BII/13-14.
  \emph{ZPS}，1971年12月6日，SWB/FE/3865/BII/13-14。}毛在8月至9月巡视各省期间发起了这场唱歌运动，无疑是为了提醒地方军队干部，谁才是解放军真正的创建者和领导者。\footnote{``Summary
  of Chairman Mao's Talks'', in Schram, Mao Tse-tung Unrehearsed,
  pp.~297-98. 《毛主席谈话纪要》，载 Schram, \emph{Mao Tse-tung
  Unrehearsed}，第297-298页。}1971年9月30日在杭州举行的气氛低调的国庆集会上，发言者是浙江省军区司令员熊应堂和陈励耘在7350部队的同事、司令员白宗善（白宗善）。白的露面可能意味着，尽管陈励耘卷入严重阴谋，但他在7350部队的高级同僚以及该部队整体未必受到牵连。白的讲话显然比熊的发言更具自我批评色彩。他要求驻省各解放军部队执行传统行为准则，``反对骄傲自满''，并认真完成``三支两军''任务。他还强调，``我们必须坚决服从地方党委的集中统一领导''。\footnote{ZJRB,
  October 1,1971, p.~2. \emph{ZJRB}，1971年10月1日，第2版。}1971年12月的一份报告要求在革命委员会任职的解放军代表和群众组织代表都服从党的统一领导。\footnote{ZPS,
  December 18,1971, URS, 65: 25, pp.~340-43.
  \emph{ZPS}，1971年12月18日，\emph{URS} 65:25，第340-343页。}尽管如此，浙江由解放军主导的领导层仍将在1972年面临一次重大清洗，以至于《中国新闻摘要》在1973年初可以声称，在所有省份中，浙江在反林彪运动中``受害最重''。\footnote{CNS,
  455 (February 15,1973), pp.~3-4. \emph{CNS},
  455（1973年2月15日），第3-4页。}

\section{Party Reconstruction}\label{party-reconstruction}

\section{党的重建}\label{ux515aux7684ux91cdux5efa}

The reconstruction of the CCP after the 9th Congress became a
battleground for political control among the central and regional arms
of the military, the Cultural Revolution radicals and the rehabilitated
cadres, with alliances continually coalescing and shifting. In December
1969, when the Cultural Revolution Group under Chen Boda's leadership
was dissolved, the radicals lost an important organizational symbol of
the Cultural Revolution and a body which had acted as the civilian
radicals' national headquarters.\footnote{Meisner, Mao's China, p.~364.
  The late Hu Hua, one of China's leading party historians and a person
  with high-level contacts and access to party archives, states that the
  Group ``naturally disintegrated'' after the 9th Congress. ``The Rise
  and Fall of Lin Biao'', p.~2. Meisner, \emph{Mao's
  China}，第364页。已故胡华是中国重要党史学家之一，拥有高层联系并可接触党内档案；他说该小组在九大后``自然解体''。见《林彪的兴衰》，第2页。}
A provisional body created to carry out the goals of a political
campaign, its dissolution proclaimed the termination, at least in an
organizational sense, of the Cultural Revolution and signalled the
return to more orthodox political structures and administrative
arrangements. In preparation for the reconstitution of local party
committees after the 9th Congress, revolutionary committees in Zhejiang
began a political and ideological movement at the end of 1969 to rectify
leading groups
九大以后中共的重建，成为中央和地区军队力量、文化大革命激进派以及被恢复职务干部之间争夺政治控制权的战场，各种联盟不断形成又不断变化。1969年12月，陈伯达领导下的中央文革小组解散，激进派失去了文化大革命的一个重要组织象征，也失去了一个曾充当文职激进派全国司令部的机构。\footnote{Meisner,
  Mao's China, p.~364. The late Hu Hua, one of China's leading party
  historians and a person with high-level contacts and access to party
  archives, states that the Group ``naturally disintegrated'' after the
  9th Congress. ``The Rise and Fall of Lin Biao'', p.~2. Meisner,
  \emph{Mao's
  China}，第364页。已故胡华是中国重要党史学家之一，拥有高层联系并可接触党内档案；他说该小组在九大后``自然解体''。见《林彪的兴衰》，第2页。}这个为执行政治运动目标而设立的临时机构一经解散，便宣告文化大革命至少在组织意义上的终结，并标志着较为正统的政治结构和行政安排的回归。为准备九大后重建地方党委，浙江各革命委员会于1969年底开始了一场整顿领导班子的政治思想运动，其中心任务是

\begin{quote}
centered on the tasks of criticizing bourgeois and petit bourgeois
ideas, opposing bourgeois factionalism and anarchism, strengthening
democratic centralism and the sense of organization and discipline and
increasing the Party spirit of the proletariat.\footnote{ZPS, October
  31,1969, SWB/FE/3220/BII/3-4, quotation p.~BII/4.
  \emph{ZPS}，1969年10月31日，SWB/FE/3220/BII/3-4，引文见BII/4页。}
批判资产阶级和小资产阶级思想，反对资产阶级派性和无政府主义，加强民主集中制和组织纪律观念，增强无产阶级党性。\footnote{ZPS,
  October 31,1969, SWB/FE/3220/BII/3-4, quotation p.~BII/4.
  \emph{ZPS}，1969年10月31日，SWB/FE/3220/BII/3-4，引文见BII/4页。}
\end{quote}

The decision to launch the movement in Zhejiang was decided upon at the
5th enlarged plenum of the ZPRC in October 1969.\footnote{ZJRB, October
  30,1969, p.~1. \emph{ZJRB}，1969年10月30日，第1版。} It seems that a
principal target of the criticism of bourgeois and petty-bourgeois
thinking was the Cultural Revolution activists whose failings were
comprehensively detailed as follows:
在浙江发动这一运动的决定，是1969年10月省革委会第五次扩大会议作出的。\footnote{ZJRB,
  October 30,1969, p.~1. \emph{ZJRB}，1969年10月30日，第1版。}看来，批判资产阶级和小资产阶级思想的主要目标之一，是文化大革命积极分子；他们的缺点被全面概括如下：

\begin{quote}
Petty-bourgeois thinking manifested in politics is opportunism,
manifested in organization is sectarianism, manifested in ideology is
selfishness, manifested in ideological method is subjectivism, and
manifested in functions is pragmatism.\footnote{ZPS, November 2,1969,
  URS, 57:15 (November 21,1969), p.~190. See the report from Changshan
  county concerning the influence of bourgeois factionalism and other
  deviations on the movement to purify the class ranks. ZPS, October
  9,1969, URS, 57:15, pp.~197-200. \emph{ZPS}，1969年11月2日，\emph{URS}
  57:15（1969年11月21日），第190页。关于资产阶级派性及其他偏差对清理阶级队伍运动影响的常山县报告，见
  \emph{ZPS}，1969年10月9日，\emph{URS} 57:15，第197-200页。}
小资产阶级思想表现在政治上是机会主义，表现在组织上是宗派主义，表现在思想上是个人主义，表现在思想方法上是主观主义，表现在职能上是实用主义。\footnote{ZPS,
  November 2,1969, URS, 57:15 (November 21,1969), p.~190. See the report
  from Changshan county concerning the influence of bourgeois
  factionalism and other deviations on the movement to purify the class
  ranks. ZPS, October 9,1969, URS, 57:15, pp.~197-200.
  \emph{ZPS}，1969年11月2日，\emph{URS}
  57:15（1969年11月21日），第190页。关于资产阶级派性及其他偏差对清理阶级队伍运动影响的常山县报告，见
  \emph{ZPS}，1969年10月9日，\emph{URS} 57:15，第197-200页。}
\end{quote}

Continued attacks on remnants of mass organizations and their refusal to
accept the termination of Cultural Revolution politics was stepped up in
1970-71.\footnote{Domes, China after the Cultural Revolution, ch.~5.
  Domes, \emph{China after the Cultural Revolution}，第5章。} In late
April 1970 a plenum of the ZPRC attacked leftist mass organizations and
bourgeois factionalism.\footnote{See CNS, 319 (May 14, 1970), pp.~1-3.
  见 \emph{CNS}, 319（1970年5月14日），第1-3页。} It accused the rebels
of using their influence on revolutionary committees to resist
directives from rebuilt party committees or party core groups within
revolutionary committees. Young and experienced cadres alike were
ordered to stay in touch with the masses and to put in regular stints of
manual labor. On 16 December 1969 the Zhejiang provincial May 7th cadre
school was opened.\footnote{ZPS, December 20,1969, SWB/FE/3275/BII/9-10.
  \emph{ZPS}，1969年12月20日，SWB/FE/3275/BII/9-10。} There, through a
regimen which combined work and study, officials were supposed to
rediscover the value of a simplistic lifestyle and to shake off
bureaucratic airs. Intellectuals were reminded to integrate themselves
with the labouring classes and to re-educate themselves. It was clear
that the military would not countenance continued organizational
indiscipline or individualistic behavior as it strove to rebuild the
party in the province. On 26 December 1969 the CCP CC released its
decision to relocate tertiary institutions in the
countryside.\footnote{Liu Suinian and Wu Qungan, ``Wenhua dageming''
  shiqide guomin jingji, p.~154. For a personal account by a
  teacher/cadre at Beijing University, which was relocated in Jiangxi
  province, see Yue Daiyun and Carolyn Wakeman, To the Storm (Berkeley
  and Los Angeles: University of California Press, 1985), ch.~12. Liu
  Suinian and Wu Qungan, \emph{``Wenhua dageming'' shiqide guomin
  jingji}，第154页。关于一名北大教师/干部在北大迁往江西后的个人叙述，见
  Yue Daiyun and Carolyn Wakeman, \emph{To the
  Storm}（伯克利和洛杉矶：University of California
  Press，1985年），第12章。} In December 1970, Zhang Yongsheng's
Zhejiang Fine Arts College was relocated from Hangzhou to Nanbao (南堡)
brigade in Tonglu county, south-west of the provincial capital. To
improve its proletarian credentials and to qualify as an acceptable
institution for the enrollment of the new student, the college was
renamed the Zhejiang Worker-Peasant Fine Arts Institute.\footnote{See,
  ZPS, September 26, 1971, FBIS/CHI, 196, G 4-5. 见
  \emph{ZPS}，1971年9月26日，FBIS/CHI, 196, G 4-5。} Zhang Yongsheng
spent time in the countryside at the college, reportedly on the advice
of Zhang Chunqiao who had suggested to him through an intermediary that
he temper himself by experiencing life as a worker, peasant and
soldier.\footnote{Zhang Yongshengde zuizheng cailiao, p.~171.
  \emph{Zhang Yongshengde zuizheng cailiao}，第171页。} Party rebuilding
commenced under the domination and direction of the PLA. According to
Domes, the PLA used three organizational devices to control the reformed
party committees and to choose personnel amenable to it and the pattern
was evident in Zhejiang.\footnote{Domes, China after the Cultural
  Revolution, pp.~52-3. Domes, \emph{China after the Cultural
  Revolution}，第52-53页。} The first device consisted of the ``three
supports and two militaries'' units stationed in various departments of
the bureaucracy and in units of production. In May 1970 the ZPMD
convened a conference to sum up the experiences of the troops carrying
out the assignment.\footnote{ZJRB, May 15,1970. It was at this
  conference that Xiong Yingtang and Nan Ping were publicly revealed as
  Commander and Political Commissar respectively of the ZPMD.
  \emph{ZJRB}，1970年5月15日。正是在这次会议上，熊应堂和南萍分别作为浙江省军区司令员和政治委员公开亮相。}
It was disclosed that many personnel engaged in ``three supports and two
militaries'' work had joined revolutionary committees or leading groups
in the areas or units where they were stationed. ``Three supports and
two militaries'' units also convened seminars and directed the staging
of activists' congresses in the ``Living Study and Application of Mao
Zedong Thought'', (活学活用毛泽东思想). These congresses promoted Lin
Biao's synthesized and truncated interpretation of Mao Zedong
Thought\footnote{Domes, China after the Cultural Revolution, ch.~4.
  Domes, \emph{China after the Cultural Revolution}，第4章。} and,
according to one observer, acted as the functional equivalents of
provincial party committees.\footnote{Harry Harding, ``China: Toward
  Revolutionary Pragmatism'', AS, 11:1 (January 1971), p.~55, fn. 10.
  Harry Harding，《中国：走向革命实用主义》，\emph{AS},
  11:1（1971年1月），第55页注10。} The first activists' congress held in
Zhejiang opened on November 1, 1969, and was attended by 4,700
delegates. Before the opening of the congress, the ZPRC had issued an
appeal warning of the dangers of ``counter-revolutionary economism'' and
``bourgeois factionalism'' and demanding a correct attitude toward and
acceptance of party leadership, as well as support for revolutionary
committees and the PLA.\footnote{ZJRB, October 20,1969, p.~1.
  \emph{ZJRB}，1969年10月20日，第1版。} The reiteration of these themes
was a clear indication that the problems of indiscipline, anarchism and
factionalism among former participants in the Cultural Revolution
continued to worry the authorities. Apart from former Governor Zhou
Jianren, who was now Vice-chairman of the ZPRC, all the principal
speakers at the congress were senior military cadres.\footnote{ZJRB,
  November 1,1969, p.~1. \emph{ZJRB}，1969年11月1日，第1版。} In his
speech at the closing session on 10 November, Nan Ping issued the call
to revolutionize leading groups. Nanbao brigade, where Zhang Yongsheng
would later discover the joys of rural life, was conferred the title of
advanced unit and activist in the movement.\footnote{ZJRB, November
  11,1969, p.~1. The most probable reason for Nanbao receiving this
  distinction related to the bravery and bold leadership displayed by
  its party secretary, Li Jinrong (李金荣), in coping with a major flood
  on July 5 of the same year. See the report in ZJRB, September 24,1969.
  \emph{ZJRB}，1969年11月11日，第1版。南堡获得这一称号的最可能原因，是其党支部书记李金荣（李金荣）在同年7月5日应对重大洪水时表现出的勇敢和果断领导。见
  \emph{ZJRB} 1969年9月24日报道。} The letter of proposal issued to the
conference on the final day contained fulsome praise of Lin Biao and
acknowledged the PLA's achievements in the province. The letter also
called for the continuation of study classes and the campaign of mass
criticism against Liu Shaoqi and Jiang Hua.\footnote{ZPS, November
  11,1969, URS, 57-15, pp.~193-7. See also ZPS, January 14,1970, URS,
  58: 7 (January 23,1972), pp.~102-5.
  \emph{ZPS}，1969年11月11日，\emph{URS} 57-15，第193-197页。另见
  \emph{ZPS}，1970年1月14日，\emph{URS}
  58:7（1972年1月23日），第102-105页。} The second organizational device
was the addition of military representatives to civilian institutions
where they were working in large numbers. The third device was the
dominance by military representatives of party core groups established
within revolutionary committees. The core groups supervised party
rebuilding and undertook preparations for the convocation of party
congresses at various administrative levels.\footnote{Qinghai province
  had established a core group in October 1967. See, Jin Chunming,
  ``Wenhua dageming'', p.~152. 青海省已于1967年10月建立核心小组。见 Jin
  Chunming, \emph{``Wenhua dageming''}，第152页。} As the great majority
of PLA members of revolutionary committees were also CCP members, the
proportion of soldiers in core groups was high. Additionally, party
groups in the military had not experienced the same level of disruption
to their activities as their civilian counterparts, although party
organs in the ZPMD had been restructured in 1969. Party offices (党办),
study offices (学办), work offices (工办) and examination offices (考办)
had been established in an administrative reshuffle modeled on that of
the MAC and it was carried out under the leadership of Nan Ping and Chen
Liyun.\footnote{See RMRB, September 22,1978, p.~1, transl. in
  SWB/FE/5929/BII/12. 见《人民日报》，1978年9月22日，第1版，译文载
  SWB/FE/5929/BII/12。}
1970至1971年间，对群众组织残余及其拒绝接受文化大革命政治终结的持续攻击进一步加强。\footnote{Domes,
  China after the Cultural Revolution, ch.~5. Domes, \emph{China after
  the Cultural Revolution}，第5章。}1970年4月下旬，省革委会一次全会攻击左派群众组织和资产阶级派性。\footnote{See
  CNS, 319 (May 14, 1970), pp.~1-3. 见 \emph{CNS},
  319（1970年5月14日），第1-3页。}它指责造反派利用自己在革命委员会中的影响，抵制重建后的党委或革命委员会内党的核心小组的指示。无论年轻干部还是有经验的干部，都被要求同群众保持联系，并定期参加体力劳动。1969年12月16日，浙江省五七干校开办。\footnote{ZPS,
  December 20,1969, SWB/FE/3275/BII/9-10.
  \emph{ZPS}，1969年12月20日，SWB/FE/3275/BII/9-10。}在那里，通过劳动和学习相结合的制度，官员应重新认识简朴生活方式的价值，并摆脱官僚习气。知识分子被提醒要同劳动阶级相结合并重新教育自己。显然，在军队努力重建省内党组织时，它不会容忍组织纪律继续松弛或个人主义行为。1969年12月26日，中共中央发布关于把高等院校迁往农村的决定。\footnote{Liu
  Suinian and Wu Qungan, ``Wenhua dageming'' shiqide guomin jingji,
  p.~154. For a personal account by a teacher/cadre at Beijing
  University, which was relocated in Jiangxi province, see Yue Daiyun
  and Carolyn Wakeman, To the Storm (Berkeley and Los Angeles:
  University of California Press, 1985), ch.~12. Liu Suinian and Wu
  Qungan, \emph{``Wenhua dageming'' shiqide guomin
  jingji}，第154页。关于一名北大教师/干部在北大迁往江西后的个人叙述，见
  Yue Daiyun and Carolyn Wakeman, \emph{To the
  Storm}（伯克利和洛杉矶：University of California
  Press，1985年），第12章。}1970年12月，张永生所在的浙江美术学院从杭州迁往省会西南桐庐县南堡（南堡）大队。为改善其无产阶级资格，并使其成为可招收新型学生的合格机构，该校改名为浙江工农美术学院。\footnote{See,
  ZPS, September 26, 1971, FBIS/CHI, 196, G 4-5. 见
  \emph{ZPS}，1971年9月26日，FBIS/CHI, 196, G 4-5。}据说张永生曾在该学院的农村地点生活一段时间，这是张春桥通过中间人向他建议的，意在让他通过体验工农兵生活来锻炼自己。\footnote{Zhang
  Yongshengde zuizheng cailiao, p.~171. \emph{Zhang Yongshengde zuizheng
  cailiao}，第171页。}党的重建在解放军支配和指导下展开。Domes认为，解放军使用三种组织手段控制改组后的党委，并选择服从其意愿的人事；这一模式在浙江也很明显。\footnote{Domes,
  China after the Cultural Revolution, pp.~52-3. Domes, \emph{China
  after the Cultural Revolution}，第52-53页。}第一种手段，是驻扎在官僚机构各部门和生产单位中的``三支两军''部队。1970年5月，浙江省军区召开会议，总结执行该任务部队的经验。\footnote{ZJRB,
  May 15,1970. It was at this conference that Xiong Yingtang and Nan
  Ping were publicly revealed as Commander and Political Commissar
  respectively of the ZPMD.
  \emph{ZJRB}，1970年5月15日。正是在这次会议上，熊应堂和南萍分别作为浙江省军区司令员和政治委员公开亮相。}会议披露，许多从事``三支两军''工作的人员，已参加其驻地地区或单位的革命委员会或领导小组。``三支两军''部队还召开讲习班，并指导举办``活学活用毛泽东思想''（活学活用毛泽东思想）积极分子代表大会。这些代表大会宣传林彪对毛泽东思想的综合化和简化解释，\footnote{Domes,
  China after the Cultural Revolution, ch.~4. Domes, \emph{China after
  the Cultural Revolution}，第4章。}并且据一位观察者说，发挥了省级党委的功能替代作用。\footnote{Harry
  Harding, ``China: Toward Revolutionary Pragmatism'', AS, 11:1 (January
  1971), p.~55, fn. 10. Harry
  Harding，《中国：走向革命实用主义》，\emph{AS},
  11:1（1971年1月），第55页注10。}浙江第一次积极分子代表大会于1969年11月1日开幕，有4,700名代表参加。大会开幕前，省革委会发布号召，警告``反革命经济主义''和``资产阶级派性''的危险，要求正确对待并接受党的领导，同时支持革命委员会和解放军。\footnote{ZJRB,
  October 20,1969, p.~1. \emph{ZJRB}，1969年10月20日，第1版。}这些主题的反复提出，清楚表明文化大革命的原参与者中存在的纪律涣散、无政府主义和派性问题仍使当局担忧。除前省长、时任省革委会副主任周建人外，大会主要发言者全是高级军队干部。\footnote{ZJRB,
  November 1,1969, p.~1. \emph{ZJRB}，1969年11月1日，第1版。}11月10日闭幕会上，南萍发出领导班子革命化的号召。南堡大队，也就是张永生日后将发现农村生活乐趣的地方，被授予运动先进单位和积极分子称号。\footnote{ZJRB,
  November 11,1969, p.~1. The most probable reason for Nanbao receiving
  this distinction related to the bravery and bold leadership displayed
  by its party secretary, Li Jinrong (李金荣), in coping with a major
  flood on July 5 of the same year. See the report in ZJRB, September
  24,1969.
  \emph{ZJRB}，1969年11月11日，第1版。南堡获得这一称号的最可能原因，是其党支部书记李金荣（李金荣）在同年7月5日应对重大洪水时表现出的勇敢和果断领导。见
  \emph{ZJRB} 1969年9月24日报道。}大会最后一天发布的倡议书极力赞扬林彪，并肯定解放军在省内的成就。倡议书还号召继续办学习班，并继续开展对刘少奇和江华的群众性批判运动。\footnote{ZPS,
  November 11,1969, URS, 57-15, pp.~193-7. See also ZPS, January
  14,1970, URS, 58: 7 (January 23,1972), pp.~102-5.
  \emph{ZPS}，1969年11月11日，\emph{URS} 57-15，第193-197页。另见
  \emph{ZPS}，1970年1月14日，\emph{URS}
  58:7（1972年1月23日），第102-105页。}第二种组织手段，是把军事代表增补到他们大量工作的文职机构中。第三种手段，是军事代表在革命委员会内建立的党的核心小组中占主导地位。核心小组监督党的重建，并为各行政层级召开党代会作准备。\footnote{Qinghai
  province had established a core group in October 1967. See, Jin
  Chunming, ``Wenhua dageming'', p.~152.
  青海省已于1967年10月建立核心小组。见 Jin Chunming, \emph{``Wenhua
  dageming''}，第152页。}由于革命委员会中的绝大多数解放军成员也是中共党员，核心小组中军人的比例很高。此外，军队中的党组织并未像文职党组织那样遭受同等程度的活动扰乱，尽管浙江省军区的党机关在1969年也经过重组。仿照军委的行政改组，在南萍和陈励耘领导下，建立了党办（党办）、学办（学办）、工办（工办）和考办（考办）。\footnote{See
  RMRB, September 22,1978, p.~1, transl. in SWB/FE/5929/BII/12.
  见《人民日报》，1978年9月22日，第1版，译文载 SWB/FE/5929/BII/12。}

Party membership among representatives of mass organizations was much
lower and more recent. For example, although he was Vice-chairman of the
ZPRC, Zhang Yongsheng was not a member of the provincial CCP core group.
At the time of the 9th CCP Congress, Zhang was not even a member of the
CCP. At the Congress Zhang Chunqiao criticized the military-dominated
leadership of Zhejiang for its unwillingness to admit rebels into the
party. Nan Ping ordered Shen Ce, who was responsible for party
organizational work in the province, to contact Hangzhou and arrange for
Zhang's admission before the Congress finished. In 1978 Shen recalled
that Zhang Chunqiao had praised Zhang Yongsheng as a national student
rebel leader who had qualified for party membership.\footnote{Zhang
  Yongshengde zuizheng cailiao, p.~8. \emph{Zhang Yongshengde zuizheng
  cailiao}，第8页。} Overall, the progress of party reconstruction
followed a tortuous route.\footnote{For an analysis of the process of
  party rebuilding between 1969 and 1973, see Graham Young, ``Party
  Building and Search for Unity'', in Bill Brugger (ed.), China: The
  Impact of the Cultural Revolution (London; Croom Helm, 1978),
  pp.~35-70; Jin Chunming, ``Wenhua dageming'', pp.~150-67; Fenwick,
  ``The Gang of Four'', pp.~135-45.
  关于1969至1973年党的重建过程的分析，见 Graham
  Young，《党建与寻求团结》，载 Bill Brugger 编，\emph{China: The Impact
  of the Cultural Revolution}（伦敦：Croom
  Helm，1978年），第35-70页；Jin Chunming, \emph{``Wenhua
  dageming''}，第150-167页；Fenwick，《四人帮》，第135-145页。}
Initially, according to Domes, it was proposed to begin at the
grassroots and work upward. But this strategy, initiated in May 1969,
was abandoned later in that year in favour of focus on the county level,
a decision probably taken at the central work meeting on party
rectification and party building held in Beijing in November
1969.\footnote{See Zhong Kan, Kang Sheng pingzhuan, p.~431. 见 Zhong
  Kan, \emph{Kang Sheng pingzhuan}，第431页。} The first county
committee was reestablished in November 1969 in Hunan province but by
October of the following year only forty-five committees in a total of
2,185 counties had been set up.\footnote{Domes, China after the Cultural
  Revolution, p.~49. If Domes' figures are correct, of these forty-five
  county party committees, at least nine were formed in Zhejiang. See
  SWB/FE/3613/BII/7-8 (Jiande county-April 1970); 3378/BII/10-11
  (Deqing-by May 1970); 3446/B/12-13 (Wuyi-July 1970); 3471/BII/21
  (Anji-August 1970); 3480/BII/18-19 (Zhuji-September 1970); Kaihua
  xianzhi, p.~24 (Kaihua-September 1970); 3495/BII/21-23 (Yuhang-by
  September 1970); 3514/BII/13-14 (Chunan-October 1970); 3520/BII/18
  (Shaoxing-by October 1970). Domes, \emph{China after the Cultural
  Revolution}，第49页。如果Domes的数字正确，这45个县委中至少有9个在浙江成立。见
  SWB/FE/3613/BII/7-8（建德县，1970年4月）；3378/BII/10-11（德清，至1970年5月）；3446/B/12-13（武义，1970年7月）；3471/BII/21（安吉，1970年8月）；3480/BII/18-19（诸暨，1970年9月）；\emph{Kaihua
  xianzhi}，第24页（开化，1970年9月）；3495/BII/21-23（余杭，至1970年9月）；3514/BII/13-14（淳安，1970年10月）；3520/BII/18（绍兴，至1970年10月）。}
Then, at the fractious Lushan plenum of August/September 1970, the
decision was made to switch the focus of party reconstruction to the
provincial level. In November 1970 Hunan became the first province to
hold a party congress since the commencement of the Cultural
Revolution.\footnote{SWB/FE/3560/B/3-6. SWB/FE/3560/B/3-6。} One
observer has claimed that the procrastination and shifts in emphasis
accompanying party rebuilding contributed to the discrediting of Chen
Boda.\footnote{Maitan, Party, Army and Masses, p.~289. Maitan,
  \emph{Party, Army and Masses}，第289页。} However, it is unlikely that
Chen Boda was in charge of party reconstruction. Alternate Politburo
member Ji Dengkui presided and Kang Sheng and Chen Boda delivered
speeches at the conference on party building which was held by the party
center in April 1970. But the three-man group established to oversee
party rectification did not include Chen. Its membership consisted of
Kang, Zhang Chunqiao and Xie Fuzhi.\footnote{See Jin Chunming, ``Wenhua
  dageming'', p.~155; Zhong Kan, Kang Sheng pingzhuan, p.~432; Ye
  Yonglie, Wang Hongwen xingshuailu, p.~364. 见 Jin Chunming,
  \emph{``Wenhua dageming''}，第155页；Zhong Kan, \emph{Kang Sheng
  pingzhuan}，第432页；Ye Yonglie, \emph{Wang Hongwen
  xingshuailu}，第364页。} The main reason for the leadership's
indecisiveness and change of mind stemmed from the continued resistance
and disruption to the new order by disgruntled mass-organization members
of revolutionary committees, urged on by their illicit organizations. In
January 1970 the party journal Red Flag carried an article by a model in
party building, PLA unit 8341, criticizing the ``theory of many
centers''.\footnote{Maitan, Party, Army and Masses, p.~283. Maitan,
  \emph{Party, Army and Masses}，第283页。} The central authorities also
propagated the experiences of Beijing's ``six factories and two
colleges'' in party building.\footnote{See Jin Chunming, ``Wenhua
  dageming'', pp.~155-58. 见 Jin Chunming, \emph{``Wenhua
  dageming''}，第155-158页。} In January and February 1970 the CCP CC
successively issued directives and a circular concerning
counter-revolutionary sabotage, graft and embezzlement, speculation and
extravagance and waste which initiated a ferocious campaign known as the
``one hit and three antis'' (一打三反).\footnote{Hao Mengbi and Duan
  Haoran, Liushinian, pp.~609-10. It was during this campaign, conducted
  on the basis of the public security bureau's ``six-point regulations''
  of January 1967 which forbade any criticism of Mao and Lin Biao, that
  the Shenyang party member Zhang Zhixin was executed. For an example of
  the application of the theories of class struggle and their
  relationship to factory management, see the article by the writing
  group of the HMRC in RMRB, 11 March 1970, SCMP, 4625, pp.~1-6. Hao
  Mengbi and Duan Haoran,
  \emph{Liushinian}，第609-610页。沈阳党员张志新正是在这一运动中被处决；该运动依据公安机关1967年1月``六条规定''进行，而该规定禁止任何批评毛和林彪的行为。关于阶级斗争理论的运用及其同工厂管理关系的例子，见杭州革委会写作组文章，《人民日报》，1970年3月11日，\emph{SCMP}
  4625，第1-6页。} In August 1970 People's Daily published an article
criticizing mass organizations for usurping the role of the party and
rejecting its leadership.\footnote{RMRB, August 21,1970, transl. in
  SCMP, 4730, pp.~34-8. 《人民日报》，1970年8月21日；译文载 \emph{SCMP}
  4730，第34-38页。} The problem of convincing the rebels of the late
1960s to adjust to the post-Cultural Revolution dispensation remained as
insoluble as ever.
群众组织代表中的党员比例低得多，而且入党时间更近。例如，张永生虽是省革委会副主任，却不是中共省级核心小组成员。在中共九大召开时，张甚至还不是中共党员。会上，张春桥批评浙江由军队主导的领导层不愿吸收造反派入党。南萍命令负责全省党组织工作的沈策联系杭州，在大会结束前安排张入党。1978年，沈回忆说，张春桥曾称赞张永生是全国学生造反派领袖，具备入党资格。\footnote{Zhang
  Yongshengde zuizheng cailiao, p.~8. \emph{Zhang Yongshengde zuizheng
  cailiao}，第8页。}总体而言，党的重建进程走过了一条曲折道路。\footnote{For
  an analysis of the process of party rebuilding between 1969 and 1973,
  see Graham Young, ``Party Building and Search for Unity'', in Bill
  Brugger (ed.), China: The Impact of the Cultural Revolution (London;
  Croom Helm, 1978), pp.~35-70; Jin Chunming, ``Wenhua dageming'',
  pp.~150-67; Fenwick, ``The Gang of Four'', pp.~135-45.
  关于1969至1973年党的重建过程的分析，见 Graham
  Young，《党建与寻求团结》，载 Bill Brugger 编，\emph{China: The Impact
  of the Cultural Revolution}（伦敦：Croom
  Helm，1978年），第35-70页；Jin Chunming, \emph{``Wenhua
  dageming''}，第150-167页；Fenwick，《四人帮》，第135-145页。}Domes认为，起初方案是从基层开始并向上推进。但这一1969年5月启动的策略在当年晚些时候被放弃，转而集中于县级；这一决定很可能是在1969年11月北京召开的关于整党和建党的中央工作会议上作出的。\footnote{See
  Zhong Kan, Kang Sheng pingzhuan, p.~431. 见 Zhong Kan, \emph{Kang
  Sheng pingzhuan}，第431页。}第一个县委于1969年11月在湖南省重建，但到翌年10月，全国2,185个县中只建立了45个县委。\footnote{Domes,
  China after the Cultural Revolution, p.~49. If Domes' figures are
  correct, of these forty-five county party committees, at least nine
  were formed in Zhejiang. See SWB/FE/3613/BII/7-8 (Jiande county-April
  1970); 3378/BII/10-11 (Deqing-by May 1970); 3446/B/12-13 (Wuyi-July
  1970); 3471/BII/21 (Anji-August 1970); 3480/BII/18-19 (Zhuji-September
  1970); Kaihua xianzhi, p.~24 (Kaihua-September 1970); 3495/BII/21-23
  (Yuhang-by September 1970); 3514/BII/13-14 (Chunan-October 1970);
  3520/BII/18 (Shaoxing-by October 1970). Domes, \emph{China after the
  Cultural
  Revolution}，第49页。如果Domes的数字正确，这45个县委中至少有9个在浙江成立。见
  SWB/FE/3613/BII/7-8（建德县，1970年4月）；3378/BII/10-11（德清，至1970年5月）；3446/B/12-13（武义，1970年7月）；3471/BII/21（安吉，1970年8月）；3480/BII/18-19（诸暨，1970年9月）；\emph{Kaihua
  xianzhi}，第24页（开化，1970年9月）；3495/BII/21-23（余杭，至1970年9月）；3514/BII/13-14（淳安，1970年10月）；3520/BII/18（绍兴，至1970年10月）。}随后，在1970年8月至9月争执激烈的庐山会议上，决定把党的重建重点转向省级。1970年11月，湖南成为文化大革命开始以来第一个召开党代会的省份。\footnote{SWB/FE/3560/B/3-6.
  SWB/FE/3560/B/3-6。}一位观察者声称，党重建中的拖延和重点转移，促成了陈伯达信誉的败坏。\footnote{Maitan,
  Party, Army and Masses, p.~289. Maitan, \emph{Party, Army and
  Masses}，第289页。}然而，陈伯达不大可能负责党的重建。1970年4月党中央召开的建党会议由政治局候补委员纪登奎主持，康生和陈伯达在会上发言。但为监督整党而成立的三人小组并不包括陈，其成员是康生、张春桥和谢富治。\footnote{See
  Jin Chunming, ``Wenhua dageming'', p.~155; Zhong Kan, Kang Sheng
  pingzhuan, p.~432; Ye Yonglie, Wang Hongwen xingshuailu, p.~364. 见
  Jin Chunming, \emph{``Wenhua dageming''}，第155页；Zhong Kan,
  \emph{Kang Sheng pingzhuan}，第432页；Ye Yonglie, \emph{Wang Hongwen
  xingshuailu}，第364页。}领导层犹豫不决和改变主意的主要原因，是革命委员会中不满的群众组织成员在其非法组织怂恿下，继续抵抗并扰乱新秩序。1970年1月，党的刊物《红旗》刊登了党建典型解放军8341部队的一篇文章，批判``多中心论''。\footnote{Maitan,
  Party, Army and Masses, p.~283. Maitan, \emph{Party, Army and
  Masses}，第283页。}中央当局还宣传北京``六厂二校''的党建经验。\footnote{See
  Jin Chunming, ``Wenhua dageming'', pp.~155-58. 见 Jin Chunming,
  \emph{``Wenhua dageming''}，第155-158页。}1970年1月和2月，中共中央先后发布有关反革命破坏、贪污盗窃、投机倒把以及铺张浪费的指示和通知，由此发动了一场被称为``一打三反''（一打三反）的猛烈运动。\footnote{Hao
  Mengbi and Duan Haoran, Liushinian, pp.~609-10. It was during this
  campaign, conducted on the basis of the public security bureau's
  ``six-point regulations'' of January 1967 which forbade any criticism
  of Mao and Lin Biao, that the Shenyang party member Zhang Zhixin was
  executed. For an example of the application of the theories of class
  struggle and their relationship to factory management, see the article
  by the writing group of the HMRC in RMRB, 11 March 1970, SCMP, 4625,
  pp.~1-6. Hao Mengbi and Duan Haoran,
  \emph{Liushinian}，第609-610页。沈阳党员张志新正是在这一运动中被处决；该运动依据公安机关1967年1月``六条规定''进行，而该规定禁止任何批评毛和林彪的行为。关于阶级斗争理论的运用及其同工厂管理关系的例子，见杭州革委会写作组文章，《人民日报》，1970年3月11日，\emph{SCMP}
  4625，第1-6页。}1970年8月，《人民日报》发表文章，批评群众组织篡夺党的角色并拒绝党的领导。\footnote{RMRB,
  August 21,1970, transl. in SCMP, 4730, pp.~34-8.
  《人民日报》，1970年8月21日；译文载 \emph{SCMP} 4730，第34-38页。}说服20世纪60年代后期的造反派适应后文化大革命安排的问题，仍然同以往一样难以解决。

\section{The Fifth CCP Zhejiang Provincial Congress, January
1971}\label{the-fifth-ccp-zhejiang-provincial-congress-january-1971}

\section{中共浙江省第五次代表大会，1971年1月}\label{ux4e2dux5171ux6d59ux6c5fux7701ux7b2cux4e94ux6b21ux4ee3ux8868ux5927ux4f1a1971ux5e741ux6708}

Before the Lin Biao affair had come to a climax in 1971, provincial
party congresses were held across China. A central meeting to prepare
for these congresses was held on 13 October 1970 and addressed by Kang
Sheng.\footnote{See, Zhong Kan, Kang Sheng pingzhuan, pp.~433-4. The
  Hangzhou party committee held its congress on October 31,1970. See
  ZJRB, 7 November 1970, p.~1. 见 Zhong Kan, \emph{Kang Sheng
  pingzhuan}，第433-434页。杭州市委于1970年10月31日召开代表大会。见
  \emph{ZJRB}，1970年11月7日，第1版。} The salient characteristics of
the Party committees elected at the congresses were the consolidation of
military power in the localities, the return to office of purged
civilian officials and the continuing decline in the influence of mass
organization representatives.\footnote{For discussion of these
  congresses, see Chu Wen-lin, ``An Analysis of the 29 New CCP
  Committees at the Provincial Level'', I\&S, 8:2, 8:3 (November and
  December 1971), pp.~47-65, 53-95; Gordon Bennett, ``Military Regions
  and Provincial Party Secretaries: One Outcome of China's Cultural
  Revolution'', CQ 54 (1973), pp.~294-307; Teiwes, Provincial Leadership
  in China, ch.~3; Domes, China after the Cultural Revolution, pp.~92-8;
  CNS, 356 (February 11, 1971), pp.~1-7. 关于这些大会的讨论，见 Chu
  Wen-lin，《29个新省级中共委员会分析》，\emph{I\&S},
  8:2、8:3（1971年11月、12月），第47-65、53-95页；Gordon
  Bennett，《军区与省委书记：中国文化大革命的一个结果》，\emph{CQ}
  54（1973年），第294-307页；Teiwes, \emph{Provincial Leadership in
  China}，第3章；Domes, \emph{China after the Cultural
  Revolution}，第92-98页；\emph{CNS} 356（1971年2月11日），第1-7页。} In
Zhejiang this trend was perhaps even more strongly evident than
elsewhere in the country. Essentially, the party leadership elected at
the congress was a continuation of the ZPRC leadership installed in
1968. But its military representation was far stronger and the decline
in representation from the other two groups in the ``three-way
alliance'' continued. The 5th Zhejiang Provincial Congress of the CCP
met from 20 to 28 January 1971.\footnote{HZRB, January 31,1971, pp.~1,
  2, 4. \emph{HZRB}，1971年1月31日，第1、2、4版。} At the time it was
described as the 4th congress but a contemporary observer noted this as
being incorrect.\footnote{CNS, 356. Zhejiang was not alone in the
  arbitrary renumbering of its party congress. The same phenomenon
  recurred in the neighboring provinces of Jiangsu and Jiangxi. In fact
  there has been some confusion over the actual number of provincial
  party congresses held in Zhejiang. Between 1949 and 1955 the CCP held
  a series (5 at least) of party conferences (党代表会议) in the
  province and in July 1956 convened the second provincial party
  congress (党代表大会), thereby bestowing retrospective recognition to
  the congress held at Pingyang county, Wenzhou in July 1939. Further
  congresses followed - from December 1959 to January 1960 (3rd), April
  1963 (4th) - before the 5th congress of January 1971. By renumbering
  the 5th congress as the 4th the provincial leaders were withdrawing
  official recognition from the 1939 congress. During the Cultural
  Revolution, most probably during the ``purification of class ranks
  campaign'', pre-1949 leaders of the ZPC such as Liu Ying (刘英), who
  had been elected secretary of the CCP ZPC at its 1939 congress, and
  cadres working in southern Zhejiang during the 1930s were branded the
  ``southern Zhejiang renegade clique''. See ZJRB, December 9,1978,
  p.~1. Details of congress dates come from Zhejiang shengqing gaiyao,
  pp.~25-6. \emph{CNS},
  356。浙江并非唯一任意重新编号党代会的省份；邻省江苏和江西也出现同样现象。事实上，浙江实际召开过多少次省党代会一直有些混乱。1949至1955年间，中共在浙江召开了一系列（至少5次）党代表会议（党代表会议），1956年7月召开第二次省党代表大会（党代表大会），由此追认1939年7月在温州平阳县召开的大会。随后又召开了1959年12月至1960年1月的第三次、1963年4月的第四次，直到1971年1月的第五次。把第五次大会重新编号为第四次，意味着省领导撤销对1939年大会的正式承认。文化大革命期间，很可能是在``清理阶级队伍运动''中，1949年前的浙江省委领导人如刘英（刘英），即1939年大会上当选为中共浙江省委书记，以及20世纪30年代在浙南工作的干部，被打成``浙南叛徒集团''。见
  \emph{ZJRB}，1978年12月9日，第1版。大会日期细节见 \emph{Zhejiang
  shengqing gaiyao}，第25-26页。} One thousand one hundred and
fifty-nine delegates, of whom 43 spoke, attended the nine-day session.
During a break in proceedings delegates made a tour of Jiaxing's South
Lake (南湖) in northern Zhejiang, site of the concluding session of the
CCP's 1st Congress of 1921. Before the closure of the congress, the
delegates elected a provincial party committee consisting of 67 members
and 22 alternate members. This committee in turn elected a standing
committee of thirteen members, six of whom were of secretary status. The
congress formalized the dominance of the military over the civilian wing
of the party. Of the six secretaries and deputy-secretaries appointed to
the CCP ZPC, two represented the 20th Army (1st Secretary Nan Ping and
Secretary Xiong Yingtang), one the 5th Air Force (Secretary Chen Liyun)
and two the Navy (deputy-secretaries Xie Zhenghao and Chai Qikun -
柴启琨). The remaining deputy-secretary was Lai Keke, the sole survivor
of the pre-Cultural Revolution secretariat. According to one source,
four of the five military representatives, Nan Ping, Xiong Yingtang, Xie
Zhenghao and Chai Qikun had all risen to prominence since Lin Biao's
assumption of the post of Minister for Defence in 1959.\footnote{Chu
  Wen-lin, ``An Analysis'', p.~61. See also CNS, 355 (4 February 1971),
  pp.~6-9. Chu Wen-lin，《分析》，第61页。另见 \emph{CNS}
  355（1971年2月4日），第6-9页。} If true, the connection may have only
been coincidental, given Lin's 4th Field Army background and the
influence of the 3rd Field Army in East China. Of the remaining seven
members of the standing committee, at least two were from the military,
two were veteran cadres and only two represented mass organizations. Hua
Yinfeng and Jiang Baodi were both from United Headquarters but, as
previously mentioned, they were not noted for their rebelliousness.
United Headquarters did not succeed in gaining a place for its leader
Zhang Yongsheng on the standing committee. Since the dissolution of
United Headquarters in August 1969, Zhang had been laboring in the
countryside. He had been sent there or, as he delicately phrased it,
permitted to go there by the ZPRC provincial party core
group.\footnote{See Zhang's article entitled ``新干部更需要勤奋地学习''
  (New cadres are required to study even more diligently) in ZJRB,
  February 5,1971, p.~2. 见张在 \emph{ZJRB}
  1971年2月5日第2版发表的《新干部更需要勤奋地学习》（New cadres are
  required to study even more diligently）。} Zhang was most probably
elected a member of the CCP ZPC. Another leading rebel, Weng Senhe, was
elected an alternate member of the ZPC at the congress. Weng's length of
party membership was even shorter than Zhang Yongsheng's. He had joined
the CCP in September 1970, just months before the opening of the
congress.\footnote{Weng Senhe bufen zuizheng cailiao, p.~1. \emph{Weng
  Senhe bufen zuizheng cailiao}，第1页。}
在1971年林彪事件达到高潮之前，全国各地召开了省级党代会。1970年10月13日，中央召开筹备这些大会的会议，康生在会上讲话。\footnote{See,
  Zhong Kan, Kang Sheng pingzhuan, pp.~433-4. The Hangzhou party
  committee held its congress on October 31,1970. See ZJRB, 7 November
  1970, p.~1. 见 Zhong Kan, \emph{Kang Sheng
  pingzhuan}，第433-434页。杭州市委于1970年10月31日召开代表大会。见
  \emph{ZJRB}，1970年11月7日，第1版。}这些大会选出的党委，其突出特征是地方军队权力的巩固、被清洗文职官员的复职，以及群众组织代表影响力的持续下降。\footnote{For
  discussion of these congresses, see Chu Wen-lin, ``An Analysis of the
  29 New CCP Committees at the Provincial Level'', I\&S, 8:2, 8:3
  (November and December 1971), pp.~47-65, 53-95; Gordon Bennett,
  ``Military Regions and Provincial Party Secretaries: One Outcome of
  China's Cultural Revolution'', CQ 54 (1973), pp.~294-307; Teiwes,
  Provincial Leadership in China, ch.~3; Domes, China after the Cultural
  Revolution, pp.~92-8; CNS, 356 (February 11, 1971), pp.~1-7.
  关于这些大会的讨论，见 Chu
  Wen-lin，《29个新省级中共委员会分析》，\emph{I\&S},
  8:2、8:3（1971年11月、12月），第47-65、53-95页；Gordon
  Bennett，《军区与省委书记：中国文化大革命的一个结果》，\emph{CQ}
  54（1973年），第294-307页；Teiwes, \emph{Provincial Leadership in
  China}，第3章；Domes, \emph{China after the Cultural
  Revolution}，第92-98页；\emph{CNS} 356（1971年2月11日），第1-7页。}在浙江，这一趋势也许比全国其他地方更加明显。本质上，大会上选出的党的领导层，是1968年建立的省革委会领导层的延续。但军队代表性更强，``三结合''中另外两个群体的代表性继续下降。中共浙江省第五次代表大会于1971年1月20日至28日召开。\footnote{HZRB,
  January 31,1971, pp.~1, 2, 4.
  \emph{HZRB}，1971年1月31日，第1、2、4版。}当时它被称为第四次代表大会，但一位当代观察者指出这是不正确的。\footnote{CNS,
  356. Zhejiang was not alone in the arbitrary renumbering of its party
  congress. The same phenomenon recurred in the neighboring provinces of
  Jiangsu and Jiangxi. In fact there has been some confusion over the
  actual number of provincial party congresses held in Zhejiang. Between
  1949 and 1955 the CCP held a series (5 at least) of party conferences
  (党代表会议) in the province and in July 1956 convened the second
  provincial party congress (党代表大会), thereby bestowing
  retrospective recognition to the congress held at Pingyang county,
  Wenzhou in July 1939. Further congresses followed - from December 1959
  to January 1960 (3rd), April 1963 (4th) - before the 5th congress of
  January 1971. By renumbering the 5th congress as the 4th the
  provincial leaders were withdrawing official recognition from the 1939
  congress. During the Cultural Revolution, most probably during the
  ``purification of class ranks campaign'', pre-1949 leaders of the ZPC
  such as Liu Ying (刘英), who had been elected secretary of the CCP ZPC
  at its 1939 congress, and cadres working in southern Zhejiang during
  the 1930s were branded the ``southern Zhejiang renegade clique''. See
  ZJRB, December 9,1978, p.~1. Details of congress dates come from
  Zhejiang shengqing gaiyao, pp.~25-6. \emph{CNS},
  356。浙江并非唯一任意重新编号党代会的省份；邻省江苏和江西也出现同样现象。事实上，浙江实际召开过多少次省党代会一直有些混乱。1949至1955年间，中共在浙江召开了一系列（至少5次）党代表会议（党代表会议），1956年7月召开第二次省党代表大会（党代表大会），由此追认1939年7月在温州平阳县召开的大会。随后又召开了1959年12月至1960年1月的第三次、1963年4月的第四次，直到1971年1月的第五次。把第五次大会重新编号为第四次，意味着省领导撤销对1939年大会的正式承认。文化大革命期间，很可能是在``清理阶级队伍运动''中，1949年前的浙江省委领导人如刘英（刘英），即1939年大会上当选为中共浙江省委书记，以及20世纪30年代在浙南工作的干部，被打成``浙南叛徒集团''。见
  \emph{ZJRB}，1978年12月9日，第1版。大会日期细节见 \emph{Zhejiang
  shengqing gaiyao}，第25-26页。}为期九天的会议有1,159名代表出席，其中43人发言。会议间隙，代表们参观了浙北嘉兴南湖（南湖），那里是1921年中共一大闭幕会议所在地。大会闭幕前，代表选出了由67名委员和22名候补委员组成的省委。该委员会又选出13人常委会，其中6人具有书记身份。大会使军队对党内文职方面的支配地位正式化。任命为中共浙江省委书记和副书记的六人中，两人代表二十军（第一书记南萍和书记熊应堂），一人代表第五空军（书记陈励耘），两人代表海军（副书记谢正浩和柴启琨
-
柴启琨）。剩下的一名副书记是赖可可，他是文化大革命前书记处唯一的幸存者。据一项资料，五名军队代表中的四人，即南萍、熊应堂、谢正浩和柴启琨，都是自林彪1959年出任国防部长以来崛起的人物。\footnote{Chu
  Wen-lin, ``An Analysis'', p.~61. See also CNS, 355 (4 February 1971),
  pp.~6-9. Chu Wen-lin，《分析》，第61页。另见 \emph{CNS}
  355（1971年2月4日），第6-9页。}如果属实，考虑到林彪的第四野战军背景和第三野战军在华东的影响，这一联系也可能只是巧合。常委会其余七名成员中，至少两人来自军队，两人是老干部，只有两人代表群众组织。华银凤和江宝娣都来自联合总部，但如前所述，她们并不以造反性著称。联合总部未能使其领导人张永生进入常委会。自1969年8月联合总部解散以来，张一直在农村劳动。用他谨慎的说法，他是由省革委会党的核心小组派去，或获准前往那里。\footnote{See
  Zhang's article entitled ``新干部更需要勤奋地学习'' (New cadres are
  required to study even more diligently) in ZJRB, February 5,1971,
  p.~2. 见张在 \emph{ZJRB}
  1971年2月5日第2版发表的《新干部更需要勤奋地学习》（New cadres are
  required to study even more diligently）。}张很可能当选为中共浙江省委委员。另一名主要造反派翁森鹤在大会上当选为省委候补委员。翁的党龄甚至比张永生还短。他于1970年9月加入中共，距大会开幕仅数月。\footnote{Weng
  Senhe bufen zuizheng cailiao, p.~1. \emph{Weng Senhe bufen zuizheng
  cailiao}，第1页。}

\textbf{Table Five. The Standing Committee of the CCP 5th Zhejiang
Provincial Committee, January 1971}
\textbf{表五：中共浙江省第五届委员会常务委员会，1971年1月}

{\def\LTcaptype{none} % do not increment counter
\begin{longtable}[]{@{}lll@{}}
\toprule\noalign{}
Position & Name & Notes \\
\midrule\noalign{}
\endhead
\bottomrule\noalign{}
\endlastfoot
First Secretary & Nan Ping & 1 \\
Secretaries & Chen Liyun & 1 \\
& Xiong Yingtang & 2 \\
Deputy-Secretaries & Lai Keke & 3 \\
& Xie Zhenghao & 2 \\
& Chai Qikun & 1 \\
Members & Wang Zida & 3 \\
& Hua Yinfeng & 4, 5 \\
& Shen Ce & 3 \\
& Dai Kelin & 2 \\
& Zhu Quanlin & 2 \\
& Meng Zhaoyu & 2 \\
& Jiang Baodi & 4, 5 \\
\end{longtable}
}

{\def\LTcaptype{none} % do not increment counter
\begin{longtable}[]{@{}lll@{}}
\toprule\noalign{}
职务 & 姓名 & 注释 \\
\midrule\noalign{}
\endhead
\bottomrule\noalign{}
\endlastfoot
第一书记 & 南萍 & 1 \\
书记 & 陈励耘 & 1 \\
& 熊应堂 & 2 \\
副书记 & 赖可可 & 3 \\
& 谢正浩 & 2 \\
& 柴启琨 & 1 \\
委员 & 王子达 & 3 \\
& 华银凤 & 4, 5 \\
& 沈策 & 3 \\
& 戴克林 & 2 \\
& 朱全林 & 2 \\
& 孟昭煜 & 2 \\
& 江宝娣 & 4, 5 \\
\end{longtable}
}

Notes: 1. PLA political officer; 2. PLA military officer; 3.
``revolutionary leading cadre'' (civilian); 4. mass organization
representative; 5. woman. 注释：1. 解放军政治军官；2. 解放军军事军官；3.
``革命领导干部''（文职）；4. 群众组织代表；5. 女性。

Source: \emph{HZRB}, January 31, 1971.
资料来源：\emph{HZRB}，1971年1月31日。

In 1968 the provincial leadership installed in Zhejiang had been
characterized by its overwhelming military and outside nature. The CCP
ZPC leadership of January 1971 was even more so. By David Goodman's
definition of provincial outsider, that is a cadre who had either served
in another province or had held no prior leadership post in the Chinese
state,\footnote{Goodman ``The Provincial Revolutionary Committee'',
  p.~77. Goodman，《省革命委员会》，第77页。} Zhejiang in 1971 was ruled
overwhelmingly by outsiders, at least at the very top of its
administration. Only two members of the thirteen-strong standing
committee of the CCP ZPC beside Lai Keke, had previous administrative
experience in Zhejiang's civilian bureaucracy.\footnote{Robert
  Scalpino's breakdown of the career profile of the three full
  secretaries of the ZPC Nan Ping, Chen Liyun and Xiong Yingtang, into
  two military personnel and one party administrator is clearly
  incorrect. See Robert A. Scalapino, ``The CCP's Provincial
  Secretaries, POC, 25:4 (1976), p.~27. The PLA's representation among
  the six ZPC secretaries was well above the national average as
  computed by Goodman,''The Provincial Revolutionary Committee'', p.~71.
  Robert
  Scalpino把浙江省委三名正式书记南萍、陈励耘和熊应堂的职业背景划分为两名军人和一名党政人员，显然不正确。见
  Robert A. Scalapino，《中共省委书记》，\emph{POC},
  25:4（1976年），第27页。解放军在浙江省委六名书记中的代表性远高于Goodman计算的全国平均水平，见
  Goodman，《省革命委员会》，第71页。} Other contemporary observers
computed the dominance of outside cadres over the leadership of the
Zhejiang administration.\footnote{Bennett, ``Military Regions'', p.~301;
  Teiwes, Provincial Leadership in China, Appendix D, p.~151.
  Bennett，《军区》，第301页；Teiwes, \emph{Provincial Leadership in
  China}，附录D，第151页。} In his comprehensive, comparative analysis
of provincial party committees elected in 1970 to 1971 Teiwes has
correlated the stability and localism indexes, which he had utilized for
a comparative study of provincial committee during the period
1956-66,\footnote{Teiwes, ``Provincial Politics in China''. See also
  Bennett, ``Military Regions'', p.~307.
  Teiwes，《中国的省级政治》。另见 Bennett，《军区》，第307页。} to the
re-establishment of party committees. He found that due to the lack of
any significant statistical relationship between pre-Cultural Revolution
indices and patterns established in 1967, factors and tensions unique to
the extraordinary political struggle of the mid-1960s determined the
pattern of behavior manifested by provincial politicians. However, over
the period 1967-73 covered by the study, Teiwes found that the
pre-Cultural Revolution pattern of provincial appointment gradually
reasserted itself as representatives of mass organizations and the
political outsiders, who had made such a dramatic entry into the
political system in 1967, were squeezed out.\footnote{Teiwes, Provincial
  Leadership in China, pp.~95, 134-38. On this point, see also Harding,
  Organizing China, pp.~262-63. Teiwes, \emph{Provincial Leadership in
  China}，第95、134-138页。关于这一点，另见 Harding, \emph{Organizing
  China}，第262-263页。} Zhejiang conformed to many of the trends
observed by Teiwes, but in certain respects it deviated from them. For
example, there was no ``substantial injection of administrative
experience''\footnote{Teiwes, Provincial Leadership in China, p.~80.
  Teiwes, \emph{Provincial Leadership in China}，第80页。} into the
provincial leadership in the years 1969-70. The influx of veteran
administrators did not occur until 1972. The ``drastic decline in the
role of political outsiders'',\footnote{Ibid, p.~95. 同上，第95页。}
(low-ranking cadres and mass representatives) observed by Teiwes in the
re-establishment of party committees was not obvious in Zhejiang because
these groups had never obtained substantial power. Their strength lay in
the ability to hinder, obstruct, sabotage, create tensions and splits
and to withhold support or recognition. In Teiwes' stability index (a
measure of the continuity of provincial leaderships as defined by the
members of party secretariats),\footnote{Teiwes, ``Provincial Politics
  in China'', pp.~150-51. Teiwes，《中国的省级政治》，第150-151页。}
Zhejiang had dropped from 17th (in the basically stable category) for
the period 1956-66 to be ranked 23rd (limited continuity) for the years
1967-73. Likewise on the localism index (a measure of the extent to
which outsiders were brought into a provincial administration), it fell
in rank from 13th (1967-68) to 21st (1969-70) to 26th (1971) of 29
provincial-level units.\footnote{Teiwes, Provincial Leadership in China,
  Table 27, p.~131, Appendix D, p.~151. Teiwes found a ``statistically
  significant positive relationship\ldots{} between stability and
  provincial localism indices for the new Party committees set up in
  1970-71''. See ibid, p.~134. Teiwes, \emph{Provincial Leadership in
  China}，表27，第131页，附录D第151页。Teiwes发现，``1970-71年建立的新党委的稳定性指数与省级本地化指数之间存在统计上显著的正相关关系''。见同书，第134页。}
As Teiwes observed, the provincial elites of the early 1970s retained
only one important shared experience, that of common service in
institutions of the new order, and this experience was marked by intense
conflict rather than smooth working relationships.\footnote{Ibid,
  p.~140. 同上，第140页。} However, the decentralization of power in the
Cultural Revolution did not signify an increase in centrifugal
tendencies because the center's monopoly over the power of patronage
ensured the selection of provincial leaders subordinate to its
will.\footnote{Ibid, pp.~141-42. 同上，第141-142页。} Zhejiang
exemplified the truth of this statement. Never, since the formalization
of provincial administrative structures in 1954-55, had the Zhejiang
leadership been more compliant to central direction and less able or
willing to represent local interests.
1968年在浙江建立的省级领导层，具有压倒性的军队性和外来性。1971年1月的中共浙江省委领导层更是如此。按照David
Goodman对省级外来者的定义，即曾在另一省任职或此前未在中国国家体制中担任领导职务的干部，\footnote{Goodman
  ``The Provincial Revolutionary Committee'', p.~77.
  Goodman，《省革命委员会》，第77页。}1971年的浙江，至少在行政体系最高层，几乎完全由外来者统治。13人的中共浙江省委常委中，除赖可可外，只有两名成员曾在浙江文职官僚体系中有行政经验。\footnote{Robert
  Scalpino's breakdown of the career profile of the three full
  secretaries of the ZPC Nan Ping, Chen Liyun and Xiong Yingtang, into
  two military personnel and one party administrator is clearly
  incorrect. See Robert A. Scalapino, ``The CCP's Provincial
  Secretaries, POC, 25:4 (1976), p.~27. The PLA's representation among
  the six ZPC secretaries was well above the national average as
  computed by Goodman,''The Provincial Revolutionary Committee'', p.~71.
  Robert
  Scalpino把浙江省委三名正式书记南萍、陈励耘和熊应堂的职业背景划分为两名军人和一名党政人员，显然不正确。见
  Robert A. Scalapino，《中共省委书记》，\emph{POC},
  25:4（1976年），第27页。解放军在浙江省委六名书记中的代表性远高于Goodman计算的全国平均水平，见
  Goodman，《省革命委员会》，第71页。}其他当代观察者也计算了外来干部对浙江行政领导层的支配程度。\footnote{Bennett,
  ``Military Regions'', p.~301; Teiwes, Provincial Leadership in China,
  Appendix D, p.~151. Bennett，《军区》，第301页；Teiwes,
  \emph{Provincial Leadership in China}，附录D，第151页。}Teiwes在对1970至1971年选出的省级党委进行全面比较分析时，把他曾用于1956至1966年省级委员会比较研究的稳定性指数和本地化指数，\footnote{Teiwes,
  ``Provincial Politics in China''. See also Bennett, ``Military
  Regions'', p.~307. Teiwes，《中国的省级政治》。另见
  Bennett，《军区》，第307页。}同党委重建相联系。他发现，由于文化大革命前指数同1967年形成的模式之间缺乏显著统计关系，20世纪60年代中期那场特殊政治斗争所特有的因素和张力，决定了省级政治人物表现出的行为模式。然而，在其研究涵盖的1967至1973年间，Teiwes发现，随着群众组织代表和1967年戏剧性进入政治系统的政治外来者被排挤出去，文化大革命前的省级任命模式逐渐重新显现。\footnote{Teiwes,
  Provincial Leadership in China, pp.~95, 134-38. On this point, see
  also Harding, Organizing China, pp.~262-63. Teiwes, \emph{Provincial
  Leadership in China}，第95、134-138页。关于这一点，另见 Harding,
  \emph{Organizing China}，第262-263页。}浙江符合Teiwes观察到的许多趋势，但在某些方面也有所偏离。例如，1969至1970年间，省级领导层并没有``实质性注入行政经验''。\footnote{Teiwes,
  Provincial Leadership in China, p.~80. Teiwes, \emph{Provincial
  Leadership in China}，第80页。}老行政人员的大量进入直到1972年才发生。Teiwes在党委重建中观察到的``政治外来者作用急剧下降''，\footnote{Ibid,
  p.~95. 同上，第95页。}即低级干部和群众代表的下降，在浙江并不明显，因为这些群体从未获得实质权力。他们的力量在于阻碍、拖延、破坏、制造紧张和分裂，以及拒绝支持或承认。在Teiwes的稳定性指数中，即以党委书记处成员界定的省级领导连续性衡量，\footnote{Teiwes,
  ``Provincial Politics in China'', pp.~150-51.
  Teiwes，《中国的省级政治》，第150-151页。}浙江从1956至1966年期间的第17位（基本稳定类），下降到1967至1973年的第23位（有限连续性）。同样，在本地化指数上，即衡量外来者进入省级行政机构的程度，浙江在29个省级单位中的排名从1967至1968年的第13位，降至1969至1970年的第21位，再降至1971年的第26位。\footnote{Teiwes,
  Provincial Leadership in China, Table 27, p.~131, Appendix D, p.~151.
  Teiwes found a ``statistically significant positive
  relationship\ldots{} between stability and provincial localism indices
  for the new Party committees set up in 1970-71''. See ibid, p.~134.
  Teiwes, \emph{Provincial Leadership in
  China}，表27，第131页，附录D第151页。Teiwes发现，``1970-71年建立的新党委的稳定性指数与省级本地化指数之间存在统计上显著的正相关关系''。见同书，第134页。}正如Teiwes所观察到的，20世纪70年代初的省级精英只保留了一种重要的共同经历，即在新秩序机构中共同任职，而这种经历以激烈冲突而非顺畅合作为特征。\footnote{Ibid,
  p.~140. 同上，第140页。}然而，文化大革命中的权力下放并不意味着离心倾向增加，因为中央对恩庇权力的垄断确保了省级领导人的选择服从其意志。\footnote{Ibid,
  pp.~141-42. 同上，第141-142页。}浙江证明了这一说法的真实性。自1954至1955年省级行政结构正式化以来，浙江领导层从未如此顺从中央指示，也从未如此缺乏代表地方利益的能力或意愿。

\section{Notes}\label{notes-5}

\section{注释}\label{ux6ce8ux91ca-5}

\bookmarksetup{startatroot}

\chapter{CHAPTER SIX - The Reestablishment of Civilian Rule, 1972-73 /
第六章 -
文官统治的重建，1972-73}\label{chapter-six---the-reestablishment-of-civilian-rule-1972-73-ux7b2cux516dux7ae0---ux6587ux5b98ux7edfux6cbbux7684ux91cdux5efa1972-73}

This chapter discusses the factional politics of the years 1972-73. With
the removal of Lin's military followers in 1972, a new provincial
leadership was installed in Zhejiang. Apart from expunging the influence
of their predecessors the new arrivals implemented a program of economic
restoration under the direction of Beijing. They also brought back to
office prominent former leaders of the pre-Cultural Revolution party
committee who had languished in political limbo for the previous five
years. Leading members of United Headquarters came under pressure for
their past associations with the discredited military leadership. Red
Storm, by contrast, had relied greatly on its contacts the old
provincial civilian and military leaders and with Xu Shiyou in Nanjing.
With the return of overthrown cadres in increasing numbers and Xu's
prominent role in purging the influence of Lin Biao in Zhejiang, Red
Storm's stocks were set to rise. However, beginning in 1973,
disagreements over the pace and extent of the reforms which had been put
in place across China threatened the unity of the central leadership.
Mao had been badly shaken by Lin's defection and had allowed Zhou Enlai
to introduce policies designed to stimulate the economy and strengthen
social and political order. However, he was not prepared to countenance
complete rejection of the ideals and political style of the Cultural
Revolution. Therefore, in 1973-74 he permitted and encouraged radical
central leaders to test anew the commitment of their social base to the
objectives and methods of the Cultural Revolution. By so doing the
Chairman triggered off a struggle to defend or move away from the
politics and policies of the Cultural Revolution. In essence, it was a
fight between its political beneficiaries on the one hand and its losers
and sufferers on the other. The renewed radical thrust witnessed the
reemergence of mass organization leaders who had made their first
appearance in the political arena in the turmoil of the mid-1960s. Many
activists of the Cultural Revolution had been admitted into the CCP in
the intervening years but their career prospects were blocked and
threatened by the increasing number of older cadres who returned to
their posts after 1972. The desire to uphold and continue the policies
and goals of the Cultural Revolution thus clashed with the equally
determined efforts to modify and adulterate them. With central leaders
looking for support in provincial administrations, increased volatility
was introduced into an already unstable situation. Zhejiang was one
province highly vulnerable to external pressure; factional bitterness
ran deep and the new leadership had not been given sufficient time to
consolidate its authority. Beginning in 1973 the pre-Cultural Revolution
mass organizations, the Women's Federation, trade union councils and the
CYL convened provincial congresses. Leaders of worker organizations of
the Cultural Revolution became leaders of trade union councils, while
former Red Guards were elected to posts in the CYL. Then in 1974, urban
militia commands were established in large towns and cities under trade
union leadership and responsible to municipal party committees, rather
than to provincial military districts and their subordinate commands.
Rebel leaders in Hangzhou drafted thousands of factory workers into
armed contingents in a strategy to build up a powerful strike force
capable of taking the offensive onto the streets. These forces also
antagonized PLA leaders who resented the militia's independence from its
control and its usurpation of public security, patrol and guard duties.
After 1968 the mass organizations had been controlled and supervised by
the military and their leaders herded into study classes or sent off to
the countryside. They were unable to gain positions of power on
reconstituted party committees, which soon shunted aside the
revolutionary committees where mass organization representation was
stronger. In 1974 the mass organization leaders were more successful in
obtaining positions in the power structure. Yet in certain respects,
power was not saddled with responsibility. The rebel cadres lacked
knowledge of or respect for the traditions and formalities associated
with the exercise of power. At the same time, their factional
organizations continued to operate behind the scenes. The rebel leaders
therefore posed a much greater threat to the provincial leadership by
their ability both to destabilize it from within and apply pressure from
outside the system. With conflicting signals being flashed from Beijing,
the Zhejiang provincial party leadership found itself hamstrung in
dealing with the rebels. The intensity of the conflict in Zhejiang built
up to dangerous levels in 1974. The intrusion of central leaders into
the affairs of the province weakened the authority of the local
administration. Central factional allegiances became even more entangled
in local power plays. Zhejiang was turned into a testing ground for the
resolution of disputes which were national in scope but whose specific
form was shaped by local events. It was only in 1975 when central
leaders temporarily put aside their differences, recognized the danger
of continued instability in Zhejiang and acted in concert to check it,
that the Zhejiang problem was resolved. But the compromises associated
with reaching the solution prevented effective implementation of the
decision. Later chapters of this study will detail the dramatic events
of the mid-1970s.

本章讨论1972-73年的派系政治。随着林彪在军队中的追随者于1972年被清除，浙江建立了新的省级领导层。新来者除了肃清前任的影响外，还在北京指导下推行经济恢复方案。他们还让文革前党委中一些重要的原领导人重新任职，这些人在此前五年里一直处于政治上的悬置状态。``联合总部''的主要成员因过去同已经失势的军队领导层的联系而受到压力。相比之下，``红暴''在很大程度上依赖它同旧省级文职和军队领导人以及南京许世友的联系。随着越来越多被打倒的干部复出，以及许世友在浙江清除林彪影响的过程中发挥突出作用，``红暴''的行情势必上升。然而，从1973年开始，围绕已经在全国推行的改革的速度和范围产生的分歧，威胁到中央领导层的团结。毛泽东受到林彪出逃的极大震动，允许周恩来推行旨在刺激经济、加强社会和政治秩序的政策。但是，他并不准备容忍对文革理想和政治风格的全盘否定。因此，1973-74年，他准许并鼓励中央激进派领导人重新检验其社会基础对文革目标和方法的忠诚。这样一来，主席引发了一场围绕保卫还是摆脱文革政治和政策的斗争。本质上，这是一方面是文革的政治受益者，另一方面是文革的失败者和受害者之间的斗争。重新兴起的激进攻势见证了群众组织领导人的再度出现，这些人最初是在1960年代中期的动荡中登上政治舞台的。许多文革积极分子在此期间已被吸收入中共，但越来越多在1972年后回到岗位的老干部阻断并威胁了他们的仕途。因此，维护并延续文革政策和目标的愿望，同同样坚决地试图修正和稀释这些政策和目标的努力发生冲突。中央领导人开始在省级行政机构中寻找支持，使本已不稳定的局面进一步增加了动荡性。浙江是一个极易受到外部压力影响的省份；派系怨恨根深蒂固，新领导层也没有得到足够时间来巩固其权威。从1973年开始，文革前的群众组织、妇联、工会委员会和共青团召开省级代表大会。文革时期工人组织的领导人成为工会委员会领导人，前红卫兵则被选入共青团任职。随后在1974年，大城市和城镇建立了城市民兵指挥机构，由工会领导，并对市级党委负责，而不是对省军区及其下属指挥机构负责。杭州的造反派领导人把数千名工厂工人编入武装队伍，其策略是建立一支强大的突击力量，能够在街头采取攻势。这些力量也激怒了人民解放军领导人，后者不满民兵脱离其控制，并篡夺公安、巡逻和警卫职责。1968年以后，群众组织一直受到军队控制和监督，其领导人被赶进学习班或送往农村。他们无法在重建后的党委中取得权力位置，而这些党委很快把群众组织代表较多的革命委员会挤到一边。1974年，群众组织领导人在权力结构中取得职位的情况更为成功。然而，在某些方面，权力并未伴随责任。造反派干部缺乏对权力运作相关传统和程式的认识或尊重。与此同时，他们的派系组织继续在幕后运作。因此，造反派领导人既能从内部使省级领导层不稳定，又能从体制外施加压力，对省级领导层构成了更大的威胁。在北京发出相互矛盾的信号时，浙江省委领导层发现自己在处理造反派问题上束手束脚。1974年，浙江冲突的强度上升到危险水平。中央领导人介入该省事务削弱了地方行政当局的权威。中央派系归属同地方权力博弈更加纠缠在一起。浙江变成了解决全国性争端的试验场，但这些争端的具体形态又由地方事件塑造。直到1975年，中央领导人暂时搁置分歧，认识到浙江持续不稳的危险，并一致行动加以遏制，浙江问题才得到解决。但是，为达成解决方案而作出的妥协妨碍了决定的有效执行。本研究后面的章节将详述1970年代中期的戏剧性事件。

\section{The New Provincial
Leadership}\label{the-new-provincial-leadership}

\section{新的省级领导层}\label{ux65b0ux7684ux7701ux7ea7ux9886ux5bfcux5c42}

After the death of Lin Biao, the central authorities despatched the
commander of the Nanjing Military Region, Xu Shiyou, to Hangzhou to take
charge of the purge of Lin's followers in the provincial military
leadership.\footnote{Tie Ying, ``动员起来为实现新时期的总任务而奋斗''
  (Get mobilized and strive to realize the general tasks of the new
  period). Report to the 6th Zhejiang Provincial Congress of the CCP,
  May 25,1978, HZRB, June 7,1978. During the first few months of 1972 a
  conspicuous campaign to ``learn from Jiangsu'', Xu's bailiwick, was
  mounted in Zhejiang. See ZPS, February 5, 1972, FBIS/CHI, 34 (1972),
  C9-11; ZPS, January 27,1972, SWB/FE/3905/BII/16; Jiangsu Provincial
  Service, March 2, 1972, FBIS/CHI, 44 (1972), C3.
  铁瑛：《动员起来为实现新时期的总任务而奋斗》（Get mobilized and strive
  to realize the general tasks of the new
  period），1978年5月25日向中共浙江省第六次代表大会所作报告，HZRB，1978年6月7日。1972年前几个月，浙江发起了一场显著的``学江苏''运动，而江苏是许世友的辖区。见ZPS，1972年2月5日，FBIS/CHI，34（1972），C9-11；ZPS，1972年1月27日，SWB/FE/3905/BII/16；Jiangsu
  Provincial Service，1972年3月2日，FBIS/CHI，44（1972），C3。} Xu was
responsible for the arrest of Lin's ``sworn followers'' in East
China\footnote{See, ``许世友同志生平'' (The Life of Comrade Xu Shiyou),
  ZJRB, November 1, 1985, p.~3. 见《许世友同志生平》（The Life of
  Comrade Xu Shiyou），ZJRB，1985年11月1日，第3页。} and he put forward
the names of senior cadres to replace Nan Ping and Xiong Yingtang as
head of the Zhejiang civilian and military administration. The two
replacements were veteran cadre Tan Qilong,\footnote{For details of
  Tan's career see, ``Tan Ch'i-lung - The Dominant Influence in Hangchow
  Labor Turmoil'', I\&S, 12: 2 (1976), pp.~102-08 (hereafter cited as
  ``Tan Qilong''); and ``T'an Ch'i-lung - First Secretary of the CCP
  Tsinghai Provincial Committee'' (hereafter cited as ``Tan Qilong
  (2)''), I\&S, 14: 3 (1978), pp.~104-11. 关于谭的履历，见``Tan
  Ch'i-lung - The Dominant Influence in Hangchow Labor
  Turmoil''，I\&S，12:2（1976），第102-08页（下文引作``Tan
  Qilong''）；以及``T'an Ch'i-lung - First Secretary of the CCP Tsinghai
  Provincial Committee''（下文引作``Tan Qilong
  (2)''），I\&S，14:3（1978），第104-11页。} and Tie Ying,\footnote{For
  details of Tie's career see ``Tie Ying - Chairman of the Zhejiang
  Provincial Advisory Commission'' I\&S, 20:10 (1984), pp.~84-92. At the
  time of his appointment Tie would have been in his mid to late
  fifties, not in his sixties as implied by one source, which guessed
  Tie's age in 1977 to have been sixty-seven. See E.A. Wayne, ``The
  Politics of Restaffing China's Provinces'', Contemporary China, 2
  (1978), p.~136. A Hong Kong source estimated correctly that in 1979
  Tie was ``over 60'', thereby placing him, in 1972, in his late
  fifties. Dongxiang (HK), September 16, 1979, in FBIS/CR/PSMA, 32
  (November 9,1979), p.~17. 关于铁的履历，见``Tie Ying - Chairman of the
  Zhejiang Provincial Advisory
  Commission''，I\&S，20:10（1984），第84-92页。任命时铁应在五十五岁至五十九岁之间，而不是某一资料暗示的六十多岁；该资料猜测铁在1977年为67岁。见E.A.
  Wayne，``The Politics of Restaffing China's Provinces''，Contemporary
  China，2（1978），第136页。一份香港资料正确估计铁在1979年``六十多岁''，由此可知1972年他为五十多岁后段。Dongxiang（HK），1979年9月16日，载FBIS/CR/PSMA，32（1979年11月9日），第17页。}
then Political Commissar of the strategically important Zhoushan island
military base. Tan was not officially appointed First Secretary of the
CCP ZPC until April 1973 and Chairman of the ZPRC the following month.
Thus, for a period of over one year, Zhejiang was left without a party
first secretary and revolutionary committee chairman.\footnote{See
  Zhejiang shengqing gaiyao, pp.~26, 28. 见Zhejiang shengqing
  gaiyao，第26、28页。} Neither did the central authorities appoint a
new commander of the ZPMD to replace Xiong Yingtang and the post
remained vacant until July 1975. Dai Kelin (戴克林), a 3rd Field Army
veteran who had been Deputy-commander of the military district since
1966, and who had joined the new order as a member of the standing
committees of the ZPRC and the CCP ZPC, may have assumed the title of
acting commander in the intervening three years. Tan Qilong had worked
as the head of both the Zhejiang and Shandong provincial party
committees before the Cultural Revolution. He returned to office in 1969
as Secretary of the Fujian Provincial Party Committee before his
promotion to Zhejiang. Tan was thus an experienced provincial
administrator who had spent five years working in Zhejiang from 1949 to
1954. The central authorities had recognized that Zhejiang required an
infusion of personnel who were familiar with its administration and
socio-economic situation. Not only had the purged leaders replaced by
Tan and Tie been military cadres, they had also been complete outsiders
to the province. Tie's meteoric rise, on the other hand, is not so
easily explained. He had no previous experience in non-military posts
and his choice for a provincial-level appointment was a surprise. His
prior work at various posts in the Nanjing Military Region and his
handling of Red Guard activities in the Cultural Revolution (referred to
briefly in chapter four) had undoubtedly brought him to Xu Shiyou's
notice. Tie's first recorded appearance in the Cultural Revolution had
been at a February 1969 rally in Hangzhou.\footnote{ZJRB, February
  16,1969. ZJRB，1969年2月16日。} As late as January 1972 he ranked only
6th among PLA leaders from Ningbo and Zhoushan Districts.\footnote{ZJRB,
  January 16,1972, p.~1. ZJRB，1972年1月16日，第1页。} Tie has described
how, on one day at the end of March 1972, his secretary relayed an
urgent notice from the Central Committee summoning him to Beijing. After
his arrival in the capital by special jet, Tie attended a meeting
chaired by Zhou Enlai. He was told that he and Tan had been assigned to
work in the CCP ZPC with a brief to restore production.\footnote{Tie
  Ying, ``周总理教我作经济工作'', ZJRB, January 9,1985, p.~3, reprinted
  in its longer, presumably original form in 怀念周恩来 (In memory of
  Zhou Enlai), (Beijing: Renmin Chubanshe, 1986), pp.~439-44.
  铁瑛：《周总理教我作经济工作》，ZJRB，1985年1月9日，第3页；以较长、推测为原始形式重印于《怀念周恩来》（In
  memory of Zhou Enlai），（北京：Renmin Chubanshe，1986），第439-44页。}
On 27 April 1972, Zhou Enlai convened a meeting (perhaps the same one
referred to by Tie above) to brief Tan, Tie and other leading civilian
and military cadres from Zhejiang including the Commander of the 5th Air
Force, Bai Zongshan, about the center's assessment of the state of
provincial affairs and to assign tasks for the post Lin Biao
era.\footnote{Zhang Yongshengde zuizheng cailiao, pp.~68, 115. Zhang
  Yongshengde zuizheng cailiao，第68、115页。} Prior to the briefing the
Central Committee issued document No.~16, dated 25 April
1972,\footnote{HZRB, September 8, 1977; Xianju xianzhi, p.~24.
  HZRB，1977年9月8日；Xianju xianzhi，第24页。} announcing the reshuffle
of the Zhejiang leadership. At the 1972 May Day rally held on 30 April,
Tan and Tie made their first public appearances in Hangzhou since their
appointments. Nan Ping and Xiong Yingtang also attended the rally and
were listed first and third in the provincial hierarchy, shadowed by Tan
and Tie as second and fourth respectively.\footnote{HZRB, May 1,1972.
  Taiwan sources, denied access to Zhejiang newspapers, dated Tan's
  first appearance in Hangzhou at September 1972. ``Tan Qilong (2)'',
  pp.~109-110.
  HZRB，1972年5月1日。台湾资料因无法取得浙江报纸，把谭在杭州首次露面的时间定为1972年9月。见``Tan
  Qilong (2)''，第109-110页。} It was to be Nan and Xiong's last public
appearances in Hangzhou. Deputy-secretary of the CCP ZPC, Chai Qikun, a
naval commissar, also disappeared from public life from May 1972 until
March 1973 in order to undergo investigation for his associations with
the Lin Biao group.\footnote{Zhang Yongshengde zuizheng cailiao, p.~93;
  CNS, 461, (March 29, 1973), p.~3. Zhang Yongshengde zuizheng
  cailiao，第93页；CNS，461（1973年3月29日），第3页。} The influence of
Lin Biao and Cultural Revolution leftism on the administration of
Zhejiang certainly extended beyond Chen Liyun, Nan Ping, Xiong Yingtang
and Chai Qikun. For the remainder of the year Tan and Tie led a
concerted campaign, which had been initiated in Beijing, to criticize
revisionism (Lin Biao) and rectify the style of work. Deputy Political
Commissar of the ZPMD, Xia Qi, who had been sent to Zhejiang from the
Nanjing Military Region in late 1970,\footnote{See Zhejiang shengqing
  gaiyao, p.~26. In October 1972 another Deputy-political Commissar was
  appointed to the ZPMD from the Nanjing Military Region. 见Zhejiang
  shengqing
  gaiyao，第26页。1972年10月，另一名副政治委员从南京军区被任命到浙江省军区。}
led the campaign against leftism in the local military. Xia made
speeches on the theme at a provincial militia meeting and at an Army Day
reception in 1972.\footnote{See ZJRB, June 20,1972, p.~1; ZJRB, August
  1,1972, p.~2.
  见ZJRB，1972年6月20日，第1页；ZJRB，1972年8月1日，第2页。} Tan Qilong
and Tie Ying not only began the settlement of accounts with Lin Biao's
military followers and leftist civilian cadres such as Lai Keke, but
they tackled the question of the responsibility of United Headquarters'
leaders for violent incidents which had occurred during the mid 1960s.
Hua Yinfeng, a member of United Headquarters, was dropped from the ZPC
standing committee. Other rebels were forced to confess that they had
participated in armed struggle and attacks on military barracks in the
years 1967-68. In 1971-72, former advisors to United Headquarters in
their capacity as emissaries of the Cultural Revolution Group, Lin Gang
and Han Xiangdong, were arrested and kept in supervised confinement.
While the arrests had been carried out by the former military leaders,
the newly-installed leadership was happy to go ahead with the
proceedings against them. Lin was later indicted as a May 16
element.\footnote{Zhang Yongshengde zuizheng cailiao, pp.~26, 35, 40.
  Even as late as January 1974 Lin's case was still under investigation.
  See ibid, p.~95. Zhang Yongshengde zuizheng
  cailiao，第26、35、40页。直到1974年1月，林的案件仍在调查中。见同书，第95页。}
Simultaneously, and in a move which would further revive old
animosities, Qiu Honggen, a leader of Red Storm die-hards was released
from prison in January 1973. Qiu attributed his release to the ``direct
concern of the responsible comrades of the ZPC''.\footnote{Qiu Honggen,
  HZRB, November 26,1978. 邱宏根，HZRB，1978年11月26日。} It appeared
that the new provincial leadership was showing a partial view toward the
disputes of the 1960s. This bias was to prove significant when elite and
mass factional struggle resumed at the end of 1973. Sometime between Lin
Biao's fatal flight and the reorganization of the provincial leadership,
former rebel leaders of United Headquarters and their leftist civilian
allies had claimed that they too were ``sufferers of Lin Biao's line''
and ``heroes in opposing Lin Biao''. They complained that raising such
issues as the January 1967 siege of the military barracks and the August
1967 raids on Xiaoshan and Zhuji (described in chapters one and two) was
``counter-revolutionary''.\footnote{CCP HMC propaganda department
  criticism group, ``We must settle accounts with Lin Biao's line'',
  HZRB, August 14,1978.
  中共杭州市委宣传部批判组：《我们必须同林彪路线算账》，HZRB，1978年8月14日。}
They saw, correctly, that criticism of rebel violence was a prelude to
negation of the Cultural Revolution and would result in the abrupt
termination of their brief political careers. Understandably, they
strongly objected. Nonetheless, there was some truth in the rebels'
assertion that they had suffered at the hands of the former military
leadership. While some rebel leaders had built up close contacts with
the PLA, others including Zhang Yongsheng had developed ties with
central and Shanghai civilian radicals such as Zhang Chunqiao, Jiang
Qing and Wang Hongwen. After 1969 those in the latter category may have
been implicated in the rivalry and antagonism which sprang up between
the ``pen'' and ``gun'' wings of the central Cultural Revolution
establishment. The new provincial leadership would deal harshly with
rebels of either category, but on the other hand it granted Red Storm
exemption from criticism and more severe forms of punishment. Three
former rebel leaders had survived the crackdown of 1969-71 relatively
unscathed. Zhang Yongsheng had played a low-key role in provincial
politics since the demobilization of the Red Guards and the termination
of the Cultural Revolution at the end of 1968. He had emerged from the
three years of turmoil as Vice-chairman of the ZPRC, Chairman of the
Zhejiang Red Guards Congress, and Chairman of the revolutionary
committee at his alma mater, the Zhejiang Fine Arts College. Another
leading Cultural Revolution rebel, Weng Senhe, had sometime after 1968
become a standing committee member of the ZPRC, a member of his silk
complex party committee when it was reestablished in January 1970 and an
alternate member of the CCP ZPC at its fifth congress in
1971.\footnote{PR, No.~40 (September 30,1977), p.~25. A more reliable
  source states that Weng joined the CCP in September 1970. See Weng
  Senhe bufen zuizheng cailiao, p.~1.
  PR，第40期（1977年9月30日），第25页。一份更可靠资料称翁于1970年9月加入中共。见Weng
  Senhe bufen zuizheng cailiao，第1页。} Additionally, Weng was most
probably the Chairman of the Zhejiang Provincial Workers' Congress.
Zhang Yongsheng's associate in United Headquarters, He Xianchun, had
been a lowly-paid worker at the Hangzhou Heavy Machinery Factory at the
outbreak of the Cultural Revolution. He became Chairman of the Hangzhou
Revolutionary Workers and Staff Members Committee
(杭州市革命职工委员会), which was later renamed the Hangzhou Workers
Congress, and was instrumental in the establishment of the municipal
social order command, whose activities have been referred to in chapter
four. He had also become Vice-chairman of the Hangzhou revolutionary
committee sometime before 1971,\footnote{ZJRB, June 23,1971, p.~1.
  ZJRB，1971年6月23日，第1页。} and in addition was a member of the
standing committee of the ZPRC and possibly an alternate member of the
CCP ZPC. The arrival of Tan and Tie in April 1972 did not bode well for
the rebel leaders. After the central authorities issued Document 16 in
April 1972 Zhang found himself in trouble, both politically and
physically. He was sent to a study class on Moganshan (莫干山) to the
north of Hangzhou and was later admitted to hospital with an unspecified
complaint.\footnote{Zhang Yongshengde zuizheng cailiao, pp.~9-10, 53.
  Zhang Yongshengde zuizheng cailiao，第9-10、53页。} Weng was detained
and spent several months laboring in Xiaoshan county. After his return
to Hangzhou he was kept under surveillance and subjected to criticism at
the silk complex.\footnote{``The Case of Weng Sengho'', pp.~2-3. ``The
  Case of Weng Sengho''，第2-3页。} The rebels could do little in the
circumstances but adopt a contrite demeanor, accept the punishment and
hope that better times lay ahead.\footnote{Zhejiang provincial core
  criticism study class, ``An outpost position to plot and seize
  power'', HZRB, July 19, 1977. Neither Zhang nor He appeared at the
  National Day celebrations in 1972. HZRB, October 1,1972. But Zhang was
  present at a rally in Hangzhou held on December 30,1972. See ZJRB,
  December 31,1972, p.~1. It seemed to be a fairly common occurrence in
  Zhejiang, and elsewhere, that political leaders under attack found it
  convenient to change their environment by resting and recuperating in
  the comparative peace and comfort of a hospital. The Zhejiang
  Hospital, reserved mainly for cadres and foreigners and situated
  opposite the Botanic Gardens in the tranquil West Lake district of
  Hangzhou, would have treated many such cases in the decade of the
  Cultural Revolution. The Cultural Revolution radicals accused Lin Biao
  of ``politicitis''. See Liu Lin, 反革命两面派的自供 (The confession of
  a counter-revolutionary double-dealer) RMRB, May 17,1974, p.~2.
  浙江省核心批判学习班：《阴谋篡党夺权的前哨阵地》，HZRB，1977年7月19日。张和何均未出现在1972年国庆庆祝活动上。HZRB，1972年10月1日。但张出席了1972年12月30日在杭州举行的一次集会。见ZJRB，1972年12月31日，第1页。在浙江和其他地方，遭受攻击的政治领导人以休养和在医院相对安宁舒适的环境中恢复为由改变环境，似乎相当常见。浙江医院主要为干部和外国人服务，位于杭州宁静西湖区的植物园对面，在文革十年中应曾接收许多这类病例。文革激进派指责林彪患有``政治病''。见刘林：《反革命两面派的自供》（The
  confession of a counter-revolutionary
  double-dealer），RMRB，1974年5月17日，第2页。} Both Tan and Tie
attended a lengthy report meeting held in Beijing from 24 May until 25
June 1972 to sum up the progress of the campaign against Lin
Biao.\footnote{A contemporary observer noted the absence of the names of
  all first secretaries from provincial reports during this period. See
  CNS, 424 (June 29,1972), p.~4.
  一位同时代观察者注意到，这一时期省级报道中所有第一书记的名字均未出现。见CNS，424（1972年6月29日），第4页。}
Three hundred and twelve participants attended. Mao's 8 July 1966 letter
to Jiang Qing expressing his reservations about Lin Biao was produced
and from 10 to 12 June, Zhou Enlai, at the Chairman's request, delivered
a lengthy report concerning the six major line struggles in the CCP
before 1949.\footnote{Hao Mengbi and Duan Haoran, Liushinian,
  pp.~625-26. Hao Mengbi and Duan Haoran，Liushinian，第625-26页。} In
December 1972 the CCP ZPC held a work meeting at Mt Pingfeng (屏风山) on
the south-western outskirts of Hangzhou to assess the progress of the
campaign.\footnote{See the report in ZJRB, January 13,1973, p.~1. Other
  provinces also held meetings on the issue. See Jiangsusheng dashiji,
  p.~319.
  见ZJRB，1973年1月13日，第1页报道。其他省份也就此问题召开会议。见Jiangsusheng
  dashiji，第319页。} Lai Keke, who now ranked third in the provincial
party hierarchy after Tan and Tie, had joined the radical cause in the
Cultural Revolution. He was questioned closely about his activities
during these years. Lai reportedly admitted that he had committed errors
under the influence of Lin Biao's ultra-leftist line and said, at the
meeting, ``I can never repudiate or reverse the verdict on
this''.\footnote{Shen Weicai, ZJRB, August 16,1977.
  沈维才，ZJRB，1977年8月16日。} The meeting decided to add four new
members to the decimated CCP ZPC standing committee. Their appointment
had been decided upon by the center in the previous month
1972.\footnote{Zhejiang shengqing gaiyao, p.~26. Zhejiang shengqing
  gaiyao，第26页。} They included Chen Weida, former Secretary of the
pre-Cultural Revolution secretariat, who was appointed Deputy-secretary
of the CCP ZPC. Chapter one of this study discussed the January 1967
incident involving Chen, who had disobeyed central instructions and
taken party files and documents containing ``black material'' about
rebel leaders to ZPMD headquarters. The resultant confrontation between
the military and United Headquarters had contributed to the two-year
dispute between the mass organizations in Zhejiang. The other prominent
cadre ``liberated'' at this time was Chen Bing. Chen had been a member
of the old CCP ZPC standing committee and Director of the provincial
Propaganda Department. In 1966 Jiang Hua had appointed him leader of
Zhejiang's Cultural Revolution Group and he had disappeared from the
political scene at the end of 1966. He now reemerged as a member of the
standing committee of the CCP ZPC.\footnote{Chen Weida and Chen Bing
  reappeared at a funeral service for a dead colleague, with Chen Weida
  delivering the obituary speech. See ZJRB, December 23,1972, p.~1.
  Funeral services became popular occasions to announce the return to
  office of rehabilitated officials. The two Chens old boss, Jiang Hua,
  was ``liberated'' in May 1973 by a decision of a central work
  conference. See Ye Yonglie, Wang Hongwen xingshuailu, p.~390. He
  reappeared publicly for the first time at the funeral service which
  was held in Beijing on July 28,1973 for former Deputy-governor of
  Zhejiang, Feng Baiju. In the following month Jiang was elected an
  alternate member of the CC at the CCP 10th National Congress.
  陈伟达和陈冰在一名已故同事的追悼会上重新露面，陈伟达致悼词。见ZJRB，1972年12月23日，第1页。追悼会成为宣布平反官员复职的常见场合。两位陈的旧上司江华，根据1973年5月一次中央工作会议决定获得``解放''。见叶永烈，Wang
  Hongwen
  xingshuailu，第390页。他首次公开复出是在1973年7月28日北京为浙江省原副省长冯白驹举行的追悼会上。次月，江在中共十大上当选中央候补委员。}
Both Chen Bing and Chen Weida had been Jiang Hua's loyal lieutenants and
core members of his clique before the Cultural Revolution. Two PLA
cadres already referred to in the previous chapter, Xia Qi and Liu Ang,
were also appointed members of the ZPC standing committee. A former
deputy of Chen Bing in the propaganda department, Shang Jingcai, also
resumed work as a leading member of the political work group
(政治工作组),\footnote{ZPS, January 14, 1973, SWB/FE/4197/BII/4.
  ZPS，1973年1月14日，SWB/FE/4197/BII/4。} the Cultural Revolution
replacement for the old provincial propaganda and organization
departments. Other middle and high-ranking bureaucrats returned to their
desks under the program of rehabilitation.\footnote{RMRB, March 20,1978,
  p.~1. See Fang Chun-kuei, ``The Reorganization and the Current Status
  of the CCP Provincial Committees'', I\&S, 9: 8 (May 1973), pp.~38-44,
  and Teiwes, Provincial Leadership in China, ch.~4, for the effect of
  Lin Biao's demise on the composition of provincial leaderships. One
  writer has observed a pattern behind cadre rehabilitation at this
  time. See Hong Yung Lee, ``The Politics of Cadre Rehabilitation since
  the Cultural Revolution'', AS, 18: 9 (1978), pp.~944-46.
  RMRB，1978年3月20日，第1页。关于林彪垮台对省级领导层构成的影响，见Fang
  Chun-kuei，``The Reorganization and the Current Status of the CCP
  Provincial
  Committees''，I\&S，9:8（1973年5月），第38-44页；以及Teiwes，Provincial
  Leadership in
  China，第4章。一位作者观察到此时干部平反背后存在一种模式。见Hong Yung
  Lee，``The Politics of Cadre Rehabilitation since the Cultural
  Revolution''，AS，18:9（1978），第944-46页。} To accommodate the large
numbers of officials returning to office, the authorities decided to
restore parts of the pre-Cultural Revolution bureaucratic structure. In
the Cultural Revolution the bureaucracy had been simplified and
departments and bureaux amalgamated. In 1973, the Hangzhou authorities
moved to replace the all-embracing groups (组) with the former more
specialized and numerous systems known as kou (口). The reform had the
added advantages of providing extra places for officials returning to
work after an enforced six-year absence and placing these more
experienced cadres in positions of seniority above their younger
colleagues who had emerged from the mass movement. It was later admitted
this consideration was a major factor behind the reform.\footnote{See
  the Hangzhou Industry and Communications Bureau group, ``Smashing the
  systems and restoring the groups was a counter-revolutionary farce'',
  HZRB, November 23,1977.
  见杭州市工业交通局小组：《砸烂口、恢复组是一场反革命闹剧》，HZRB，1977年11月23日。}

林彪死后，中央派南京军区司令员许世友赴杭州，负责清洗省级军队领导层中的林彪追随者。\footnote{Tie
  Ying, ``动员起来为实现新时期的总任务而奋斗'' (Get mobilized and strive
  to realize the general tasks of the new period). Report to the 6th
  Zhejiang Provincial Congress of the CCP, May 25,1978, HZRB, June
  7,1978. During the first few months of 1972 a conspicuous campaign to
  ``learn from Jiangsu'', Xu's bailiwick, was mounted in Zhejiang. See
  ZPS, February 5, 1972, FBIS/CHI, 34 (1972), C9-11; ZPS, January
  27,1972, SWB/FE/3905/BII/16; Jiangsu Provincial Service, March 2,
  1972, FBIS/CHI, 44 (1972), C3.
  铁瑛：《动员起来为实现新时期的总任务而奋斗》（Get mobilized and strive
  to realize the general tasks of the new
  period），1978年5月25日向中共浙江省第六次代表大会所作报告，HZRB，1978年6月7日。1972年前几个月，浙江发起了一场显著的``学江苏''运动，而江苏是许世友的辖区。见ZPS，1972年2月5日，FBIS/CHI，34（1972），C9-11；ZPS，1972年1月27日，SWB/FE/3905/BII/16；Jiangsu
  Provincial Service，1972年3月2日，FBIS/CHI，44（1972），C3。}许世友负责逮捕华东地区林彪的``死党''\footnote{See,
  ``许世友同志生平'' (The Life of Comrade Xu Shiyou), ZJRB, November 1,
  1985, p.~3. 见《许世友同志生平》（The Life of Comrade Xu
  Shiyou），ZJRB，1985年11月1日，第3页。}，并提出了接替南萍和熊应堂、领导浙江文职和军事行政机构的高级干部人选。两名接替者是老干部谭启龙\footnote{For
  details of Tan's career see, ``Tan Ch'i-lung - The Dominant Influence
  in Hangchow Labor Turmoil'', I\&S, 12: 2 (1976), pp.~102-08 (hereafter
  cited as ``Tan Qilong''); and ``T'an Ch'i-lung - First Secretary of
  the CCP Tsinghai Provincial Committee'' (hereafter cited as ``Tan
  Qilong (2)''), I\&S, 14: 3 (1978), pp.~104-11. 关于谭的履历，见``Tan
  Ch'i-lung - The Dominant Influence in Hangchow Labor
  Turmoil''，I\&S，12:2（1976），第102-08页（下文引作``Tan
  Qilong''）；以及``T'an Ch'i-lung - First Secretary of the CCP Tsinghai
  Provincial Committee''（下文引作``Tan Qilong
  (2)''），I\&S，14:3（1978），第104-11页。}和铁瑛\footnote{For details
  of Tie's career see ``Tie Ying - Chairman of the Zhejiang Provincial
  Advisory Commission'' I\&S, 20:10 (1984), pp.~84-92. At the time of
  his appointment Tie would have been in his mid to late fifties, not in
  his sixties as implied by one source, which guessed Tie's age in 1977
  to have been sixty-seven. See E.A. Wayne, ``The Politics of Restaffing
  China's Provinces'', Contemporary China, 2 (1978), p.~136. A Hong Kong
  source estimated correctly that in 1979 Tie was ``over 60'', thereby
  placing him, in 1972, in his late fifties. Dongxiang (HK), September
  16, 1979, in FBIS/CR/PSMA, 32 (November 9,1979), p.~17.
  关于铁的履历，见``Tie Ying - Chairman of the Zhejiang Provincial
  Advisory
  Commission''，I\&S，20:10（1984），第84-92页。任命时铁应在五十五岁至五十九岁之间，而不是某一资料暗示的六十多岁；该资料猜测铁在1977年为67岁。见E.A.
  Wayne，``The Politics of Restaffing China's Provinces''，Contemporary
  China，2（1978），第136页。一份香港资料正确估计铁在1979年``六十多岁''，由此可知1972年他为五十多岁后段。Dongxiang（HK），1979年9月16日，载FBIS/CR/PSMA，32（1979年11月9日），第17页。}，后者当时是具有战略重要性的舟山岛军事基地政治委员。谭直到1973年4月才被正式任命为中共浙江省委第一书记，次月任浙江省革命委员会主任。因此，在一年多时间里，浙江既没有党委第一书记，也没有革命委员会主任。\footnote{See
  Zhejiang shengqing gaiyao, pp.~26, 28. 见Zhejiang shengqing
  gaiyao，第26、28页。}中央也没有任命新的浙江省军区司令员来接替熊应堂，该职位一直空缺到1975年7月。戴克林（戴克林）是第三野战军老兵，自1966年起任军区副司令员，并作为浙江省革命委员会和中共浙江省委常委进入新秩序；在这三年间，他可能曾以代理司令员名义主持工作。文革前，谭启龙曾任浙江和山东省委负责人。1969年，他复出任福建省委书记，随后调升浙江。因此，谭是一位经验丰富的省级行政干部，1949年至1954年曾在浙江工作五年。中央已经认识到，浙江需要补充熟悉其行政和社会经济状况的人员。被谭和铁取代的遭清洗领导人不仅都是军队干部，而且完全是外省人。另一方面，铁瑛的迅速升迁并不那么容易解释。他没有担任非军事职务的经验，获选担任省级职位令人意外。他此前在南京军区多个岗位的工作，以及他在文革中处理红卫兵活动的情况（第四章曾简要提及），无疑使许世友注意到了他。铁瑛在文革中首次有记录的露面，是1969年2月杭州的一次集会。\footnote{ZJRB,
  February 16,1969. ZJRB，1969年2月16日。}直到1972年1月，他在宁波和舟山地区解放军领导人中还只排第六。\footnote{ZJRB,
  January 16,1972, p.~1. ZJRB，1972年1月16日，第1页。}铁瑛曾描述，1972年3月底的一天，他的秘书转交了中央委员会召他进京的紧急通知。他乘专机抵达首都后，参加了周恩来主持的一次会议。会上告知他和谭已被派往中共浙江省委工作，任务是恢复生产。\footnote{Tie
  Ying, ``周总理教我作经济工作'', ZJRB, January 9,1985, p.~3, reprinted
  in its longer, presumably original form in 怀念周恩来 (In memory of
  Zhou Enlai), (Beijing: Renmin Chubanshe, 1986), pp.~439-44.
  铁瑛：《周总理教我作经济工作》，ZJRB，1985年1月9日，第3页；以较长、推测为原始形式重印于《怀念周恩来》（In
  memory of Zhou Enlai），（北京：Renmin Chubanshe，1986），第439-44页。}1972年4月27日，周恩来召开会议（或许就是铁瑛上文所说的那次会议），向谭、铁以及浙江其他主要文职和军队干部，包括第五空军司令员白宗善，通报中央对省内局势的评估，并布置林彪事件之后的任务。\footnote{Zhang
  Yongshengde zuizheng cailiao, pp.~68, 115. Zhang Yongshengde zuizheng
  cailiao，第68、115页。}在通报之前，中央委员会发布了1972年4月25日的第16号文件\footnote{HZRB,
  September 8, 1977; Xianju xianzhi, p.~24. HZRB，1977年9月8日；Xianju
  xianzhi，第24页。}，宣布浙江领导层改组。在4月30日举行的1972年五一集会上，谭和铁任命后首次在杭州公开露面。南萍和熊应堂也出席集会，并在省级层级中列第一和第三，谭和铁则分别列第二和第四。\footnote{HZRB,
  May 1,1972. Taiwan sources, denied access to Zhejiang newspapers,
  dated Tan's first appearance in Hangzhou at September 1972. ``Tan
  Qilong (2)'', pp.~109-110.
  HZRB，1972年5月1日。台湾资料因无法取得浙江报纸，把谭在杭州首次露面的时间定为1972年9月。见``Tan
  Qilong (2)''，第109-110页。}这成为南、熊二人在杭州最后一次公开露面。中共浙江省委副书记、海军政委柴启琨也从1972年5月至1973年3月消失于公共生活，以接受其同林彪集团关系的调查。\footnote{Zhang
  Yongshengde zuizheng cailiao, p.~93; CNS, 461, (March 29, 1973), p.~3.
  Zhang Yongshengde zuizheng
  cailiao，第93页；CNS，461（1973年3月29日），第3页。}林彪和文革左倾主义对浙江行政的影响，当然不止陈励耘、南萍、熊应堂和柴启琨。该年余下时间里，谭和铁领导了一场由北京发起、协同推进的运动，批判修正主义（林彪）并整顿作风。浙江省军区副政治委员夏琦1970年末从南京军区调往浙江\footnote{See
  Zhejiang shengqing gaiyao, p.~26. In October 1972 another
  Deputy-political Commissar was appointed to the ZPMD from the Nanjing
  Military Region. 见Zhejiang shengqing
  gaiyao，第26页。1972年10月，另一名副政治委员从南京军区被任命到浙江省军区。}，领导地方军队中反对左倾主义的运动。1972年，夏在一次省民兵会议和一次建军节招待会上就这一主题发表讲话。\footnote{See
  ZJRB, June 20,1972, p.~1; ZJRB, August 1,1972, p.~2.
  见ZJRB，1972年6月20日，第1页；ZJRB，1972年8月1日，第2页。}谭启龙和铁瑛不仅开始同林彪的军队追随者以及赖可可等左派文职干部清算，而且处理``联合总部''领导人对1960年代中期暴力事件所负责任的问题。``联合总部''成员华银凤被撤出浙江省委常委会。其他造反派被迫承认，他们曾在1967-68年参加武斗和攻击军营。1971-72年，曾以文革小组使者身份担任``联合总部''顾问的林岗和韩湘东被逮捕并受监管羁押。逮捕虽然是前军队领导人实施的，但新上任的领导层乐于继续推进对他们的处理。林后来被指控为``五一六''分子。\footnote{Zhang
  Yongshengde zuizheng cailiao, pp.~26, 35, 40. Even as late as January
  1974 Lin's case was still under investigation. See ibid, p.~95. Zhang
  Yongshengde zuizheng
  cailiao，第26、35、40页。直到1974年1月，林的案件仍在调查中。见同书，第95页。}与此同时，另一项会进一步激起旧怨的举动是，``红暴''死硬派领导人邱宏根于1973年1月获释出狱。邱把自己的获释归因于``省委负责同志的直接关怀''。\footnote{Qiu
  Honggen, HZRB, November 26,1978. 邱宏根，HZRB，1978年11月26日。}看来，新的省级领导层对1960年代的争端表现出偏向性看法。当1973年底精英和群众派系斗争重新爆发时，这种偏向被证明意义重大。在林彪致命出逃和省级领导层重组之间的某个时候，``联合总部''前造反派领导人及其左派文职盟友声称，他们也是``林彪路线的受害者''和``反林彪英雄''。他们抱怨说，提出1967年1月围困军营、1967年8月袭击萧山和诸暨（第一、二章已述）等问题，是``反革命''的。\footnote{CCP
  HMC propaganda department criticism group, ``We must settle accounts
  with Lin Biao's line'', HZRB, August 14,1978.
  中共杭州市委宣传部批判组：《我们必须同林彪路线算账》，HZRB，1978年8月14日。}他们正确地看到，批评造反派暴力是全盘否定文革的前奏，并将导致他们短暂政治生涯的突然终结。可以理解，他们强烈反对。尽管如此，造反派声称自己曾受前军队领导层之害，也有一定事实依据。一些造反派领导人曾同解放军建立密切联系，而包括张永生在内的另一些人则同张春桥、江青、王洪文等中央和上海文职激进派建立联系。1969年后，后一类人可能被卷入中央文革建制中``笔杆子''和``枪杆子''两翼之间出现的竞争和对立。新的省级领导层会严厉处理两类造反派，但另一方面，它却使``红暴''免于批判和更严厉的惩罚。三名前造反派领导人在1969-71年的镇压中相对毫发无损地存活下来。自红卫兵复员和1968年底文革终结以来，张永生在省级政治中一直保持低调。三年动荡之后，他成为浙江省革命委员会副主任、浙江红卫兵代表大会主席，以及其母校浙江美术学院革命委员会主任。另一名重要文革造反派翁森鹤，在1968年后的某个时候成为浙江省革命委员会常委，1970年1月其丝绸联合企业党委重建时成为该党委成员，并在1971年中共浙江省第五次代表大会上成为省委候补委员。\footnote{PR,
  No.~40 (September 30,1977), p.~25. A more reliable source states that
  Weng joined the CCP in September 1970. See Weng Senhe bufen zuizheng
  cailiao, p.~1.
  PR，第40期（1977年9月30日），第25页。一份更可靠资料称翁于1970年9月加入中共。见Weng
  Senhe bufen zuizheng cailiao，第1页。}此外，翁很可能还是浙江省工人代表大会主席。张永生在``联合总部''的同伴贺贤春，文革爆发时是杭州重型机器厂的一名低薪工人。他成为杭州市革命职工委员会（杭州市革命职工委员会）主任，后该组织改名为杭州市工人代表大会；他还在建立市社会秩序指挥部过程中发挥了关键作用，该指挥部的活动已在第四章提及。他还在1971年前某时成为杭州市革命委员会副主任\footnote{ZJRB,
  June 23,1971, p.~1. ZJRB，1971年6月23日，第1页。}，此外还是浙江省革命委员会常委，并可能是中共浙江省委候补委员。谭和铁于1972年4月到来，对造反派领导人不是好兆头。中央1972年4月发布16号文件后，张在政治上和身体上都陷入麻烦。他被送到杭州以北莫干山（莫干山）的学习班，后来又因不明病症入院。\footnote{Zhang
  Yongshengde zuizheng cailiao, pp.~9-10, 53. Zhang Yongshengde zuizheng
  cailiao，第9-10、53页。}翁被拘留，并在萧山县劳动数月。回到杭州后，他受到监视，并在丝绸联合企业遭到批判。\footnote{``The
  Case of Weng Sengho'', pp.~2-3. ``The Case of Weng Sengho''，第2-3页。}在这种情况下，造反派几乎无能为力，只能表现出悔过姿态，接受处罚，并希望好时机在前方。\footnote{Zhejiang
  provincial core criticism study class, ``An outpost position to plot
  and seize power'', HZRB, July 19, 1977. Neither Zhang nor He appeared
  at the National Day celebrations in 1972. HZRB, October 1,1972. But
  Zhang was present at a rally in Hangzhou held on December 30,1972. See
  ZJRB, December 31,1972, p.~1. It seemed to be a fairly common
  occurrence in Zhejiang, and elsewhere, that political leaders under
  attack found it convenient to change their environment by resting and
  recuperating in the comparative peace and comfort of a hospital. The
  Zhejiang Hospital, reserved mainly for cadres and foreigners and
  situated opposite the Botanic Gardens in the tranquil West Lake
  district of Hangzhou, would have treated many such cases in the decade
  of the Cultural Revolution. The Cultural Revolution radicals accused
  Lin Biao of ``politicitis''. See Liu Lin, 反革命两面派的自供 (The
  confession of a counter-revolutionary double-dealer) RMRB, May
  17,1974, p.~2.
  浙江省核心批判学习班：《阴谋篡党夺权的前哨阵地》，HZRB，1977年7月19日。张和何均未出现在1972年国庆庆祝活动上。HZRB，1972年10月1日。但张出席了1972年12月30日在杭州举行的一次集会。见ZJRB，1972年12月31日，第1页。在浙江和其他地方，遭受攻击的政治领导人以休养和在医院相对安宁舒适的环境中恢复为由改变环境，似乎相当常见。浙江医院主要为干部和外国人服务，位于杭州宁静西湖区的植物园对面，在文革十年中应曾接收许多这类病例。文革激进派指责林彪患有``政治病''。见刘林：《反革命两面派的自供》（The
  confession of a counter-revolutionary
  double-dealer），RMRB，1974年5月17日，第2页。}谭和铁二人均参加了1972年5月24日至6月25日在北京举行的长时间汇报会议，总结批林运动的进展。\footnote{A
  contemporary observer noted the absence of the names of all first
  secretaries from provincial reports during this period. See CNS, 424
  (June 29,1972), p.~4.
  一位同时代观察者注意到，这一时期省级报道中所有第一书记的名字均未出现。见CNS，424（1972年6月29日），第4页。}共有312人参加。会上拿出了毛泽东1966年7月8日写给江青、表达他对林彪保留意见的信；6月10日至12日，周恩来应主席要求，就1949年前中共六次重大路线斗争作了长篇报告。\footnote{Hao
  Mengbi and Duan Haoran, Liushinian, pp.~625-26. Hao Mengbi and Duan
  Haoran，Liushinian，第625-26页。}1972年12月，中共浙江省委在杭州西南郊屏风山（屏风山）召开工作会议，评估运动进展。\footnote{See
  the report in ZJRB, January 13,1973, p.~1. Other provinces also held
  meetings on the issue. See Jiangsusheng dashiji, p.~319.
  见ZJRB，1973年1月13日，第1页报道。其他省份也就此问题召开会议。见Jiangsusheng
  dashiji，第319页。}赖可可此时在省级党内层级中位列谭、铁之后第三，文革中曾投身激进派事业。会议对他这些年的活动作了细致追问。据报道，赖承认自己在林彪极左路线影响下犯了错误，并在会上说：``我永远不能否定或翻这个案。''\footnote{Shen
  Weicai, ZJRB, August 16,1977. 沈维才，ZJRB，1977年8月16日。}会议决定向已经残缺不全的中共浙江省委常委会增补四名新成员。他们的任命已由中央在前一个月即1972年11月决定。\footnote{Zhejiang
  shengqing gaiyao, p.~26. Zhejiang shengqing gaiyao，第26页。}其中包括文革前省委书记处原书记陈伟达，他被任命为中共浙江省委副书记。本研究第一章讨论过1967年1月涉及陈伟达的事件；陈违背中央指示，将含有关于造反派领导人``黑材料''的党内档案和文件带到浙江省军区司令部。由此产生的军队同``联合总部''之间的对峙，助长了浙江群众组织之间持续两年的争端。此时另一位获得``解放''的重要干部是陈冰。陈曾是旧中共浙江省委常委、省委宣传部部长。1966年，江华任命他为浙江文革小组组长，1966年底他从政治舞台上消失。现在，他重新以中共浙江省委常委身份出现。\footnote{Chen
  Weida and Chen Bing reappeared at a funeral service for a dead
  colleague, with Chen Weida delivering the obituary speech. See ZJRB,
  December 23,1972, p.~1. Funeral services became popular occasions to
  announce the return to office of rehabilitated officials. The two
  Chens old boss, Jiang Hua, was ``liberated'' in May 1973 by a decision
  of a central work conference. See Ye Yonglie, Wang Hongwen
  xingshuailu, p.~390. He reappeared publicly for the first time at the
  funeral service which was held in Beijing on July 28,1973 for former
  Deputy-governor of Zhejiang, Feng Baiju. In the following month Jiang
  was elected an alternate member of the CC at the CCP 10th National
  Congress.
  陈伟达和陈冰在一名已故同事的追悼会上重新露面，陈伟达致悼词。见ZJRB，1972年12月23日，第1页。追悼会成为宣布平反官员复职的常见场合。两位陈的旧上司江华，根据1973年5月一次中央工作会议决定获得``解放''。见叶永烈，Wang
  Hongwen
  xingshuailu，第390页。他首次公开复出是在1973年7月28日北京为浙江省原副省长冯白驹举行的追悼会上。次月，江在中共十大上当选中央候补委员。}陈冰和陈伟达在文革前都是江华忠诚的副手和其派系的核心成员。前一章已提及的两名解放军干部夏琦和刘昂，也被任命为省委常委。陈冰在宣传部的一名前副手商景才，也恢复工作，担任政治工作组（政治工作组）的领导成员；政治工作组是文革中取代旧省委宣传部和组织部的机构。\footnote{ZPS,
  January 14, 1973, SWB/FE/4197/BII/4.
  ZPS，1973年1月14日，SWB/FE/4197/BII/4。}其他中高级官僚也在平反方案下回到岗位。\footnote{RMRB,
  March 20,1978, p.~1. See Fang Chun-kuei, ``The Reorganization and the
  Current Status of the CCP Provincial Committees'', I\&S, 9: 8 (May
  1973), pp.~38-44, and Teiwes, Provincial Leadership in China, ch.~4,
  for the effect of Lin Biao's demise on the composition of provincial
  leaderships. One writer has observed a pattern behind cadre
  rehabilitation at this time. See Hong Yung Lee, ``The Politics of
  Cadre Rehabilitation since the Cultural Revolution'', AS, 18: 9
  (1978), pp.~944-46.
  RMRB，1978年3月20日，第1页。关于林彪垮台对省级领导层构成的影响，见Fang
  Chun-kuei，``The Reorganization and the Current Status of the CCP
  Provincial
  Committees''，I\&S，9:8（1973年5月），第38-44页；以及Teiwes，Provincial
  Leadership in
  China，第4章。一位作者观察到此时干部平反背后存在一种模式。见Hong Yung
  Lee，``The Politics of Cadre Rehabilitation since the Cultural
  Revolution''，AS，18:9（1978），第944-46页。}为容纳大量复职官员，当局决定恢复文革前官僚机构结构的一部分。文革中，官僚机构曾被简化，各部门和局被合并。1973年，杭州当局开始用文革前更专业、数量更多的``口''（口）来取代包罗万象的``组''（组）。这项改革还有额外好处：为被迫离岗六年后复职的官员提供更多位置，并把这些更有经验的干部安置在从群众运动中涌现的年轻同僚之上的高级职位。后来有人承认，这一考虑是改革背后的一个主要因素。\footnote{See
  the Hangzhou Industry and Communications Bureau group, ``Smashing the
  systems and restoring the groups was a counter-revolutionary farce'',
  HZRB, November 23,1977.
  见杭州市工业交通局小组：《砸烂口、恢复组是一场反革命闹剧》，HZRB，1977年11月23日。}

\section{New Directions, 1972}\label{new-directions-1972}

\section{新方向，1972年}\label{ux65b0ux65b9ux54111972ux5e74}

Policy initiatives from Beijing assisted Tan and Tie in their mission of
restoring social order and normalizing party life. Zhou Enlai, with Mao
Zedong's approval, assembled a package of policies concerning economic
development, cadres and intellectuals. The general thrust of Zhou's
program was directed against the leftism of the Cultural Revolution. In
key editorials and articles published under his direction or in
important unpublished speeches, Zhou tentatively suggested a restoration
of certain pre-Cultural Revolution policies.\footnote{See 周恩来选集
  (下卷) (Selected Works of Zhou Enlai - Vol. 2), (Hangzhou: Renmin
  Chubanshe, 1984), pp.~463-66,471-74. See also Luo Yang ``Correctly
  understand and handle the relationship between politics and vocational
  work'', Hongqi, No.~4 (1972), SCMM, 727 (May 1,1972), pp.~9-16;
  ``惩前毖后,治病救人'', (Learn from past mistakes to avoid future ones,
  cure the illness to save the patient), RMRB ed., April 24,1972, p.~2;
  ``Strive for new victories'', RMRB, Hongqi, Jiefang Junbao joint ed.,
  October 1, 1972, SCMP, 5234 (October 12,1972), pp.~140-7; Long
  Yan,''无政府主义是马克思主义骗子的反革命工具'' (Anarchism is a
  counter-revolutionary weapon of sham Marxist swindlers), RMRB, October
  14,1972, p.~2. For a detailed discussion of these issues, see William
  A. Joseph, The Critique of Ultra-Leftism in China, 1958-1981
  (Stanford: Stanford University Press, 1984), pp.~123-32.
  见《周恩来选集（下卷）》（Selected Works of Zhou Enlai - Vol.
  2），（杭州：Renmin
  Chubanshe，1984），第463-66、471-74页。另见罗扬：``Correctly
  understand and handle the relationship between politics and vocational
  work''，Hongqi，第4期（1972），SCMM，727（1972年5月1日），第9-16页；《惩前毖后,治病救人》（Learn
  from past mistakes to avoid future ones, cure the illness to save the
  patient），RMRB社论，1972年4月24日，第2页；``Strive for new
  victories''，RMRB、Hongqi、Jiefang
  Junbao联合社论，1972年10月1日，SCMP，5234（1972年10月12日），第140-7页；龙岩：《无政府主义是马克思主义骗子的反革命工具》（Anarchism
  is a counter-revolutionary weapon of sham Marxist
  swindlers），RMRB，1972年10月14日，第2页。关于这些问题的详细讨论，见William
  A. Joseph，The Critique of Ultra-Leftism in China,
  1958-1981（Stanford: Stanford University Press，1984），第123-32页。}
In this endeavor Zhou ran into the firm if not unexpected resistance of
Zhang Chunqiao and Yao Wenyuan. Both had acquired substantial formal
power in the propaganda field as Politburo members and members of the
propaganda and organization group established in late 1970 (see the
previous chapter). Zhou Enlai was senior to both men and was responsible
for the routine work of the Central Committee, but on at least four
occasions Zhang and Yao, with Jiang Qing's backing, were able to resist
or stonewall Zhou's initiatives. The first occasion concerned an
investigation report undertaken by Zhou Peiyuan, a noted scientist and
head of the administration of Beijing University. In July 1972 Zhou
Enlai instructed him to redress the prejudice against basic theory in
the physical sciences and to correct leftist policies in scientific
research and education.The issue had been raised by a visiting
Chinese-American physicist in a meeting with Mao, who turned the
question over to the Premier to handle.\footnote{See Fenwick, ``The Gang
  of Four'', p.~173. 见Fenwick，``The Gang of Four''，第173页。} On 23
July the Premier approved Zhou's report but People's Daily, which was
directly under Yao Wenyuan's control, refused to carry the article.
Eventually the so-called daily newspaper for intellectuals, Guangming
Ribao, carried the article on 6 October 1972. Zhang and Yao then
directed the Shanghai daily Wenhui Bao to carry articles rebutting the
report.\footnote{See Gao Gao and Yan Jiaqi, ``Wenhua dageming'',
  pp.~471, 474-5; Hu Hua (ed.), Zhongguo shehuizhuyi geming, p.~313;
  Domes, China after the Cultural Revolution, pp.~168-9; Hao Mengbi and
  Duan Haoran, Liushinian, p.~626. 见Gao Gao and Yan Jiaqi，``Wenhua
  dageming''，第471、474-5页；Hu Hua（编），Zhongguo shehuizhuyi
  geming，第313页；Domes，China after the Cultural
  Revolution，第168-9页；Hao Mengbi and Duan
  Haoran，Liushinian，第626页。} The second occasion Zhang and Yao
obstructed the Premier concerned Zhou's instructions to People's Daily
in August and September 1972 to carry articles condemning
ultra-leftism.\footnote{Cited in Fenwick, ``The Gang of Four'', p.~170.
  引自Fenwick，``The Gang of Four''，第170页。} Zhou reportedly stated:
``Without thoroughly discrediting the ultra 'left' trend, you will not
have the courage to implement Chairman Mao's revolutionary line''.
Eventually, on 14 October, the newspaper printed three articles
condemning anarchism, a codeword for rebel behavior during the Cultural
Revolution. Zhang and Yao directed Wenhui Bao to collect and publish
material written by worker groups in Shanghai refuting the Zhou-
inspired articles and to send the material to Mao. People's Daily's
inhouse publication then republished the articles which contained
broadsides against ``revisionism and a rightist
counter-current''.\footnote{Wang Ruoshui (then on the editorial board of
  Renmin Ribao), ``The greatest lesson of the Cultural Revolution is
  that the personality cult should be opposed'', 明报月刊 (HK), No.~2
  (February 1,1980), in FBIS/CR/PSMA, 66, p.~92; Hao Mengbi and Duan
  Haoran, Liushinian, p.~627; Hu Hua (ed.), Zhongguo shehuizhuyi geming,
  p.~313; Gao Gao and Yan Jiaqi, ``Wenhua dageming'', pp.~471, 476.
  王若水（当时为《人民日报》编委）：``The greatest lesson of the
  Cultural Revolution is that the personality cult should be
  opposed''，《明报月刊》（香港），第2期（1980年2月1日），载FBIS/CR/PSMA，66，第92页；Hao
  Mengbi and Duan Haoran，Liushinian，第627页；Hu Hua（编），Zhongguo
  shehuizhuyi geming，第313页；Gao Gao and Yan Jiaqi，``Wenhua
  dageming''，第471、476页。} Next, Zhou's attempt to insert direct
criticism of ultra-leftism into the 1972 national day joint editorial of
People's Daily, Liberation Army Daily and Red Flag was thwarted by Yao
Wenyuan, who twice deleted the offending passage and limited mention of
the harmful effects of ultra-leftism to specific policy
areas.\footnote{Gao Wenqian,
  ``艰难而光辉的最后岁月--记'文化大革命'期间的周恩来'' (Arduous yet
  glorious last years - Zhou Enlai in the `Cultural Revolution'), RMRB,
  January 5,1986, p.~5.
  高文谦：《艰难而光辉的最后岁月--记'文化大革命'期间的周恩来》（Arduous
  yet glorious last years - Zhou Enlai in the `Cultural
  Revolution'），RMRB，1986年1月5日，第5页。} The final occasion arose
when the CCP CC International Liaison Department and the Department of
Foreign Affairs submitted a joint work report in late November 1972,
attacking the influence of ultra-leftism and anarchism in the conduct of
foreign relations, Zhang Chunqiao and Jiang Qing rejected the report
after it had received Zhou Enlai's approval.\footnote{Hao Mengbi and
  Duan Haoran, Liushinian, p.~627; Dangshi tongxun, No.~55 (15 January
  1983), p.~8. Hao Mengbi and Duan Haoran，Liushinian，第627页；Dangshi
  tongxun，第55期（1983年1月15日），第8页。} Matters came to a head in
December 1972 when senior personnel of People's Daily, undoubtedly
confused by countermanding orders from above, wrote to Mao expressing
agreement with Zhou Enlai's opinions concerning ultra-leftism. On 17
December Mao called in his confidants on theoretical issues, Zhang
Chunqiao and Yao Wenyuan, to discuss the matter. Perceiving that the
logical outcome of Zhou's push against ultra-leftism and the association
between ultra-leftism and Lin Biao would reflect badly on the Cultural
Revolution, Mao moved to stem the tide. He reversed the judgment that he
had held about Lin as late as June 1972\footnote{Wang Ruoshui,
  FBIS/CR/PSMA, 66. 王若水，FBIS/CR/PSMA，66。} and now categorized Lin
as an ``ultra-rightist, revisionist, splitter, schemer and intriguer,
and traitor to the Party and the country''.\footnote{Wang Zhaokun (ed.),
  中共党史170问答 (170 Questions and Answers on Party History),
  (Shenyang: Liaoning Renmin Chubanshe, 1984), pp.~309-11;
  中共江西省委党校党史教研室 (Teaching and Research Office of the CCP
  Jiangxi Party Committee Party School), 中共党史百题解答 (Explanations
  on Various Questions in Party History - enlarged edition), (Nanchang:
  Jiangxi Renmin Chubanshe, 1983), pp.~380-81; Gao Gao and Yan Jiaqi,
  ``Wenhua dageming'', p.~477; Hao Mengbi and Duan Haoran, Liushinian,
  p.~627; Hu Hua, ``From the Fall of Lin Biao to the Demise of the Gang
  of Four'', (unpublished manuscript) p.~1. The switch in categorization
  of Lin from an ultra-leftist to ultra-rightist was noticed at the time
  by Harry Bradsher, ``China: The Radical Offensive'', AS, 13: 11
  (1973), p.~990. 王兆坤（编）：《中共党史170问答》（170 Questions and
  Answers on Party History），（沈阳：Liaoning Renmin
  Chubanshe，1984），第309-11页；中共江西省委党校党史教研室（Teaching
  and Research Office of the CCP Jiangxi Party Committee Party
  School）：《中共党史百题解答》（Explanations on Various Questions in
  Party History - enlarged edition），（南昌：Jiangxi Renmin
  Chubanshe，1983），第380-81页；Gao Gao and Yan Jiaqi，``Wenhua
  dageming''，第477页；Hao Mengbi and Duan
  Haoran，Liushinian，第627页；Hu Hua，``From the Fall of Lin Biao to
  the Demise of the Gang of Four''（未刊稿），第1页。Harry
  Bradsher当时注意到林从极左派到极右派的类别转换，见``China: The Radical
  Offensive''，AS，13:11（1973），第990页。} For their action in siding
with the Premier, Wang Ruoshui and other leading members of People's
Daily's editorial department including Hu Jiwei, were called to the
Great Hall of the People and accused by Jiang Qing of ``splitting the
center'' and by Zhang Chunqiao of being ``ultra-rightists''.\footnote{See
  Wang Ruoshui, ``智慧的痛苦'' (The pain of wisdom) in the unpublished
  manuscript of the same name, pp.~221-22. 见王若水：《智慧的痛苦》（The
  pain of wisdom）同名未刊稿，第221-22页。}

来自北京的政策举措帮助谭、铁完成恢复社会秩序和使党内生活正常化的使命。周恩来在毛泽东批准下，形成了一套涉及经济发展、干部和知识分子的政策。周方案的总体锋芒指向文革中的左倾主义。在他指导发表的重要社论和文章中，或在重要的未发表讲话中，周试探性地提出恢复某些文革前政策。\footnote{See
  周恩来选集 (下卷) (Selected Works of Zhou Enlai - Vol. 2), (Hangzhou:
  Renmin Chubanshe, 1984), pp.~463-66,471-74. See also Luo Yang
  ``Correctly understand and handle the relationship between politics
  and vocational work'', Hongqi, No.~4 (1972), SCMM, 727 (May 1,1972),
  pp.~9-16; ``惩前毖后,治病救人'', (Learn from past mistakes to avoid
  future ones, cure the illness to save the patient), RMRB ed., April
  24,1972, p.~2; ``Strive for new victories'', RMRB, Hongqi, Jiefang
  Junbao joint ed., October 1, 1972, SCMP, 5234 (October 12,1972),
  pp.~140-7; Long Yan,''无政府主义是马克思主义骗子的反革命工具''
  (Anarchism is a counter-revolutionary weapon of sham Marxist
  swindlers), RMRB, October 14,1972, p.~2. For a detailed discussion of
  these issues, see William A. Joseph, The Critique of Ultra-Leftism in
  China, 1958-1981 (Stanford: Stanford University Press, 1984),
  pp.~123-32. 见《周恩来选集（下卷）》（Selected Works of Zhou Enlai -
  Vol. 2），（杭州：Renmin
  Chubanshe，1984），第463-66、471-74页。另见罗扬：``Correctly
  understand and handle the relationship between politics and vocational
  work''，Hongqi，第4期（1972），SCMM，727（1972年5月1日），第9-16页；《惩前毖后,治病救人》（Learn
  from past mistakes to avoid future ones, cure the illness to save the
  patient），RMRB社论，1972年4月24日，第2页；``Strive for new
  victories''，RMRB、Hongqi、Jiefang
  Junbao联合社论，1972年10月1日，SCMP，5234（1972年10月12日），第140-7页；龙岩：《无政府主义是马克思主义骗子的反革命工具》（Anarchism
  is a counter-revolutionary weapon of sham Marxist
  swindlers），RMRB，1972年10月14日，第2页。关于这些问题的详细讨论，见William
  A. Joseph，The Critique of Ultra-Leftism in China,
  1958-1981（Stanford: Stanford University Press，1984），第123-32页。}在这一努力中，周遇到了张春桥和姚文元坚定但并非意外的抵制。二人作为政治局委员、以及1970年末建立的宣传组织组成员（见前章），在宣传领域已取得相当大的正式权力。周恩来资历高于二人，并负责中央委员会的日常工作，但至少在四个场合，张、姚在江青支持下能够抵制或拖延周的举措。第一次涉及著名科学家、北京大学行政负责人周培源所作的一份调查报告。1972年7月，周恩来指示他纠正物理科学中轻视基础理论的偏见，并纠正科研和教育中的左倾政策。这个问题是一位来访的美籍华裔物理学家在同毛会面时提出的，毛把问题交给总理处理。\footnote{See
  Fenwick, ``The Gang of Four'', p.~173. 见Fenwick，``The Gang of
  Four''，第173页。}7月23日，总理批准了周培源的报告，但由姚文元直接控制的《人民日报》拒绝刊登。最终，所谓知识分子的日报《光明日报》于1972年10月6日刊发了这篇文章。张和姚随后指示上海日报《文汇报》发表文章驳斥该报告。\footnote{See
  Gao Gao and Yan Jiaqi, ``Wenhua dageming'', pp.~471, 474-5; Hu Hua
  (ed.), Zhongguo shehuizhuyi geming, p.~313; Domes, China after the
  Cultural Revolution, pp.~168-9; Hao Mengbi and Duan Haoran,
  Liushinian, p.~626. 见Gao Gao and Yan Jiaqi，``Wenhua
  dageming''，第471、474-5页；Hu Hua（编），Zhongguo shehuizhuyi
  geming，第313页；Domes，China after the Cultural
  Revolution，第168-9页；Hao Mengbi and Duan
  Haoran，Liushinian，第626页。}张、姚第二次阻挠总理，是围绕1972年8月和9月周指示《人民日报》刊登谴责极左主义文章一事。\footnote{Cited
  in Fenwick, ``The Gang of Four'', p.~170. 引自Fenwick，``The Gang of
  Four''，第170页。}据报道，周说：``不把极`左'思潮彻底搞臭，你们就没有勇气执行毛主席的革命路线。''最终，10月14日，该报刊登三篇谴责无政府主义的文章；无政府主义是指称文革期间造反派行为的暗语。张、姚指示《文汇报》收集并发表上海工人小组撰写的材料，反驳受周启发的文章，并把材料送给毛。《人民日报》内部刊物随后转载了这些文章，其中猛烈抨击``修正主义和右倾逆流''。\footnote{Wang
  Ruoshui (then on the editorial board of Renmin Ribao), ``The greatest
  lesson of the Cultural Revolution is that the personality cult should
  be opposed'', 明报月刊 (HK), No.~2 (February 1,1980), in FBIS/CR/PSMA,
  66, p.~92; Hao Mengbi and Duan Haoran, Liushinian, p.~627; Hu Hua
  (ed.), Zhongguo shehuizhuyi geming, p.~313; Gao Gao and Yan Jiaqi,
  ``Wenhua dageming'', pp.~471, 476.
  王若水（当时为《人民日报》编委）：``The greatest lesson of the
  Cultural Revolution is that the personality cult should be
  opposed''，《明报月刊》（香港），第2期（1980年2月1日），载FBIS/CR/PSMA，66，第92页；Hao
  Mengbi and Duan Haoran，Liushinian，第627页；Hu Hua（编），Zhongguo
  shehuizhuyi geming，第313页；Gao Gao and Yan Jiaqi，``Wenhua
  dageming''，第471、476页。}接着，周试图在1972年《人民日报》《解放军报》和《红旗》国庆联合社论中加入对极左主义的直接批评，却被姚文元挫败；姚两次删去这段冒犯性文字，并把极左主义危害的提及限制在具体政策领域。\footnote{Gao
  Wenqian, ``艰难而光辉的最后岁月--记'文化大革命'期间的周恩来'' (Arduous
  yet glorious last years - Zhou Enlai in the `Cultural Revolution'),
  RMRB, January 5,1986, p.~5.
  高文谦：《艰难而光辉的最后岁月--记'文化大革命'期间的周恩来》（Arduous
  yet glorious last years - Zhou Enlai in the `Cultural
  Revolution'），RMRB，1986年1月5日，第5页。}最后一次发生在1972年11月下旬，中共中央对外联络部和外交部提交联合工作报告，批评极左主义和无政府主义在处理对外关系中的影响；报告经周恩来批准后，张春桥和江青拒绝接受。\footnote{Hao
  Mengbi and Duan Haoran, Liushinian, p.~627; Dangshi tongxun, No.~55
  (15 January 1983), p.~8. Hao Mengbi and Duan
  Haoran，Liushinian，第627页；Dangshi
  tongxun，第55期（1983年1月15日），第8页。}1972年12月，事情发展到紧要关头。《人民日报》高级人员无疑被上级相互抵触的命令弄得困惑，于是致信毛，表示同意周恩来关于极左主义的意见。12月17日，毛召见他在理论问题上的亲信张春桥和姚文元讨论此事。毛意识到，周反极左的推动以及把极左主义同林彪联系起来，在逻辑后果上会对文革不利，于是出手止住潮流。他推翻了自己直到1972年6月仍对林持有的判断\footnote{Wang
  Ruoshui, FBIS/CR/PSMA, 66. 王若水，FBIS/CR/PSMA，66。}，现在把林界定为``极右派、修正主义者、分裂主义者、阴谋家、野心家和党和国家的叛徒''。\footnote{Wang
  Zhaokun (ed.), 中共党史170问答 (170 Questions and Answers on Party
  History), (Shenyang: Liaoning Renmin Chubanshe, 1984), pp.~309-11;
  中共江西省委党校党史教研室 (Teaching and Research Office of the CCP
  Jiangxi Party Committee Party School), 中共党史百题解答 (Explanations
  on Various Questions in Party History - enlarged edition), (Nanchang:
  Jiangxi Renmin Chubanshe, 1983), pp.~380-81; Gao Gao and Yan Jiaqi,
  ``Wenhua dageming'', p.~477; Hao Mengbi and Duan Haoran, Liushinian,
  p.~627; Hu Hua, ``From the Fall of Lin Biao to the Demise of the Gang
  of Four'', (unpublished manuscript) p.~1. The switch in categorization
  of Lin from an ultra-leftist to ultra-rightist was noticed at the time
  by Harry Bradsher, ``China: The Radical Offensive'', AS, 13: 11
  (1973), p.~990. 王兆坤（编）：《中共党史170问答》（170 Questions and
  Answers on Party History），（沈阳：Liaoning Renmin
  Chubanshe，1984），第309-11页；中共江西省委党校党史教研室（Teaching
  and Research Office of the CCP Jiangxi Party Committee Party
  School）：《中共党史百题解答》（Explanations on Various Questions in
  Party History - enlarged edition），（南昌：Jiangxi Renmin
  Chubanshe，1983），第380-81页；Gao Gao and Yan Jiaqi，``Wenhua
  dageming''，第477页；Hao Mengbi and Duan
  Haoran，Liushinian，第627页；Hu Hua，``From the Fall of Lin Biao to
  the Demise of the Gang of Four''（未刊稿），第1页。Harry
  Bradsher当时注意到林从极左派到极右派的类别转换，见``China: The Radical
  Offensive''，AS，13:11（1973），第990页。}由于站在总理一边，王若水以及包括胡绩伟在内的《人民日报》编辑部其他主要成员被叫到人民大会堂，江青指责他们``分裂中央''，张春桥则指责他们是``极右派''。\footnote{See
  Wang Ruoshui, ``智慧的痛苦'' (The pain of wisdom) in the unpublished
  manuscript of the same name, pp.~221-22. 见王若水：《智慧的痛苦》（The
  pain of wisdom）同名未刊稿，第221-22页。}

In terms of policy reform, agricultural distribution was one of the
principal areas in which Zhou and his subordinates attempted to reverse
what they considered to be extreme examples of radical social and
economic policies. Egalitarianism and political assessment of work
performance were rejected and moves initiated to revert to
``distribution according to work''.\footnote{See Domes, China after the
  Cultural Revolution, pp.~156-60; Parris Chang, ``Decentralization of
  Power'', PoC, 21: 4 (1972), p.~73. 见Domes，China after the Cultural
  Revolution，第156-60页；Parris Chang，``Decentralization of
  Power''，PoC，21:4（1972），第73页。} In January 1972 Zhejiang Daily
had published a commentary endorsing the position.\footnote{ZPS, January
  21,1972, FBIS/CHI, 16, pp.~C 6-8.
  ZPS，1972年1月21日，FBIS/CHI，16，第C6-8页。} The restoration of
``rational rules and regulations'' in running industrial enterprises was
also announced, and tentative steps were initiated to put education
policy back on a more formal footing.\footnote{See CNS issues in 1972;
  ZPS, May 30,1972, URS, 68: 9, pp.~117-18.
  见1972年CNS各期；ZPS，1972年5月30日，URS，68:9，第117-18页。} Mao's
intervention did not signal an immediate end to all the policies that
the Premier was in the middle of implementing. For one thing, many of
Zhou's policies had built up a momentum of their own. In addition, while
demanding strict conformity to vague ideological prescriptions
concerning class struggle and other theoretical issues, the Chairman
permitted his lieutenants greater freedom of movement in policy areas in
the final years of his life. Mao's switches in direction, however, were
unnerving for those working under him. Policies relating to the
rehabilitation of veteran cadres and economic adjustment continued to be
implemented. After listening to a report on planning and the
strengthening of economic management on 16 February 1973, the Premier
directed that bonuses and piece-work rates, largely abandoned in the
Cultural Revolution, be reintroduced for jobs involving heavy physical
labour. Zhou also suggested a more sweeping rectification of the
national economy.\footnote{Gao Gao and Yan Jiaqi, ``Wenhua dageming'',
  p.~473. Gao Gao and Yan Jiaqi，``Wenhua dageming''，第473页。} In
August 1972 Chen Yun and Wang Zhen were among the first central leading
cadres rehabilitated after the Lin Biao affair.\footnote{Ibid, p.~471.
  同上，第471页。} In May 1973, the arch-criminal of the ``February
Countercurrent'', Tan Zhenlin, and twelve other veteran cadres,
including such prominent provincial figure as Li Jingquan (Sichuan), Li
Baohua (Anhui), Jiang Weiqing (Jiangsu) and Jiang Hua (Zhejiang), were
permitted to resume work.\footnote{Hu Hua (ed.), Zhongguo shehuizhuyi
  geming, p.~312; Ye Yonglie, Wang Hongwen xingshuailu, p.~390. Hu
  Hua（编），Zhongguo shehuizhuyi geming，第312页；叶永烈，Wang Hongwen
  xingshuailu，第390页。} In March 1973 the ``No.~2 capitalist-roader''
Deng Xiaoping had returned to office as Vice-premier of the State
Council and at the CCP's 10th National Congress in August Deng was
elected to the Central Committee.

就政策改革而言，农业分配是周及其下属试图扭转的主要领域之一；他们认为其中存在激进社会经济政策的极端例子。平均主义和以政治标准评估工作表现遭到拒绝，恢复``按劳分配''的举措开始启动。\footnote{See
  Domes, China after the Cultural Revolution, pp.~156-60; Parris Chang,
  ``Decentralization of Power'', PoC, 21: 4 (1972), p.~73.
  见Domes，China after the Cultural Revolution，第156-60页；Parris
  Chang，``Decentralization of Power''，PoC，21:4（1972），第73页。}1972年1月，《浙江日报》发表一篇评论支持这一立场。\footnote{ZPS,
  January 21,1972, FBIS/CHI, 16, pp.~C 6-8.
  ZPS，1972年1月21日，FBIS/CHI，16，第C6-8页。}企业管理中恢复``合理的规章制度''也被宣布，教育政策重新走向更正式轨道的试探性步骤也开始实施。\footnote{See
  CNS issues in 1972; ZPS, May 30,1972, URS, 68: 9, pp.~117-18.
  见1972年CNS各期；ZPS，1972年5月30日，URS，68:9，第117-18页。}毛的干预并不意味着总理正在执行的所有政策立即终止。一方面，周的许多政策已经形成自身势头。此外，主席在生命最后几年虽然要求在阶级斗争和其他理论问题上严格遵守含混的意识形态规定，却允许他的副手们在政策领域有更大的活动自由。然而，毛的方向转换使其下属感到不安。有关老干部平反和经济调整的政策继续执行。1973年2月16日听取关于计划和加强经济管理的报告后，总理指示，对涉及重体力劳动的工作，重新实行文革中大体废止的奖金和计件工资。周还建议对国民经济进行更全面的整顿。\footnote{Gao
  Gao and Yan Jiaqi, ``Wenhua dageming'', p.~473. Gao Gao and Yan
  Jiaqi，``Wenhua dageming''，第473页。}1972年8月，陈云和王震是林彪事件后最早获平反的中央领导干部之一。\footnote{Ibid,
  p.~471. 同上，第471页。}1973年5月，``二月逆流''的首要罪人谭震林以及其他十二名老干部，包括李井泉（四川）、李葆华（安徽）、江渭清（江苏）和江华（浙江）等重要省级人物，被准许恢复工作。\footnote{Hu
  Hua (ed.), Zhongguo shehuizhuyi geming, p.~312; Ye Yonglie, Wang
  Hongwen xingshuailu, p.~390. Hu Hua（编），Zhongguo shehuizhuyi
  geming，第312页；叶永烈，Wang Hongwen xingshuailu，第390页。}1973年3月，``第二号走资派''邓小平已复出担任国务院副总理，并在8月召开的中共十大上当选中央委员。

\section{Leftist Revival, 1973}\label{leftist-revival-1973}

\section{左派复兴，1973年}\label{ux5de6ux6d3eux590dux51741973ux5e74}

Despite the Premier's persistence in continuing with his policies of
restoration, Mao's definition of Lin Biao as an ``ultra-rightist'' spelt
danger for him. Bolstered by the Chairman's reaffirmation of the
Cultural Revolution, its principal remaining beneficiaries in the
central leadership, Jiang Qing, Zhang Chunqiao, Yao Wenyuan and Wang
Hongwen devoted their energies to propagating its virtues. In September
1972 the youthful Wang had been transferred to the national political
arena\footnote{CNS, 438 (October 5,1972), p.~1.
  CNS，438（1972年10月5日），第1页。} and then in May 1973 co-opted onto
the Politburo. The work meeting of the CCP CC, held from May 20 to
31,1973, decided to invite Wang, Hua Guofeng and Beijing party First
Secretary Wu De to sit in on Politburo meetings and participate in the
body's work. The step was taken because of the loss of one third (seven)
of the Politburo's membership as a consequence of the Lin Biao
affair.\footnote{See Dangshi tongxun, No.~55, p.~13. One source suggests
  that Kang Sheng recommended Wang to the Chairman because of his fame
  as a rebel. See Ye Yonglie, ``毛泽东如何挑选接班人'' (How did Mao
  Zedong choose his successors), Zhang Yufeng et al., 毛泽东轶事
  (Anecdotes about Mao Zedong), (Changsha: Hunan Wenyi Chubanshe, 1989),
  p.~251. 见Dangshi
  tongxun，第55期，第13页。一份资料称康生因王作为造反派的名声而向主席推荐了他。见叶永烈：《毛泽东如何挑选接班人》（How
  did Mao Zedong choose his
  successors），载张玉凤等：《毛泽东轶事》（Anecdotes about Mao
  Zedong），（长沙：Hunan Wenyi Chubanshe，1989），第251页。} In order
to reward their followers and at the same time strengthen their
organizational base the four radicals, later described by Mao as the
Gang of Four, took advantage of the opportunity presented by the
re-establishment of pre-Cultural Revolution mass organizations to push
for the election of Cultural Revolution activists to leading positions.
Additionally, they intervened in the selection of delegates to the 10th
Party Congress to press the claims of their followers. Articles in the
national press renewed the attack on revisionism, highlighted the
achievements of ``new born things'', praised Jiang Qing's contribution
to the revolutionization of the arts and promoted the exploits of model
rebels who dared to ``go against the tide''.\footnote{Gao Gao and Yan
  Jiaqi, ``Wenhua dageming'', pp.~477-84; Domes, China after the
  Cultural Revolution, pp.~172-75. Gao Gao and Yan Jiaqi，``Wenhua
  dageming''，第477-84页；Domes，China after the Cultural
  Revolution，第172-75页。} One such model was Zhu Kejia, a Shanghai
youth who had settled down in Yunnan Province rather than take advantage
of the opportunity to return to his hometown. Another was Zhang
Tiesheng, a rusticated youth from Liaoning. In 1973 Zhang sat for the
university entrance examinations, the first such examinations held since
1966. His paper was blank except for comments protesting against the
system's discrimination against sent-down-youth such as himself. A third
was Huang Shuai, a primary school student from Beijing municipality who
protested against ill-treatment at the hands of her teacher and wrote to
the city newspaper about her case. All of the examples received
prominent media attention and rebels such as Zhang Tiesheng were
selected as delegates to the 10th Congress. Another person whose cause
was taken up by the central radicals was Hangzhou's Weng Senhe (see
chapter seven).\footnote{Gao Gao and Yan Jiaqi, ``Wenhua dageming'',
  pp.~476-77. Gao Gao and Yan Jiaqi，``Wenhua dageming''，第476-77页。}
In May 1973, in a development that would have important repercussions
for the future, Mao raised the question of publishing a critique of
Confucius.\footnote{Ibid, p.~484. 同上，第484页。} The historical role
of the sage and his evaluation by communist historians had long been a
matter of intense debate. On several occasions, Mao himself had
previously expressed an ambivalent attitude toward Confucius. In the
early 1970s official publications had sporadically carried articles
critical of China's most influential philosopher and teacher.\footnote{For
  example, an article on Lu Xun by the radical writing group from
  Shanghai, Luo Siding, published in RMRB on September 25,1971,
  expressed anti-Confucian sentiments. Another article attacking the
  sage was published in September 1972. See Domes, China after the
  Cultural Revolution, p.~172. Debates concerning the historical role of
  Confucius had been occurring in Shanghai since 1969. See Lyn T. White,
  ``Local Autonomy in China During the Cultural Revolution: The
  Theoretical Uses of an Atypical Case'', American Political Science
  Review, 70:2 (June 1976), p.~490.
  例如，上海激进写作组罗思鼎一篇关于鲁迅的文章于1971年9月25日发表于RMRB，表达了反孔情绪。另一篇攻击这位圣人的文章发表于1972年9月。见Domes，China
  after the Cultural
  Revolution，第172页。关于孔子历史角色的争论自1969年以来已在上海出现。见Lyn
  T. White，``Local Autonomy in China During the Cultural Revolution:
  The Theoretical Uses of an Atypical Case''，American Political Science
  Review，70:2（1976年6月），第490页。} On 1 December 1972 a professor
of philosophy from Guangzhou's Zhongshan University, Yang Rongguo, had
published an article in Hongqi concerning class struggle in the period
of the warring states in which he harshly criticized the sage's
historical role.\footnote{Gao Gao and Yan Jiaqi, ``Wenhua dageming'',
  p.~478. Gao Gao and Yan Jiaqi，``Wenhua dageming''，第478页。} On 7
August 1973, Yang's reassessment of Confucius, examined and approved by
Mao,\footnote{See Hao Mengbi and Duan Haoran, Liushinian, p.~632. For a
  translation of Yang's article ``Confucius-A Thinker who Stubbornly
  supported the Slave System'', see Selected Articles Criticizing Lin
  Piao and Confucius 1 (Beijing: Foreign Languages Press, 1974),
  pp.~1-23. 见Hao Mengbi and Duan
  Haoran，Liushinian，第632页。杨文``Confucius-A Thinker who Stubbornly
  supported the Slave System''的译文，见Selected Articles Criticizing
  Lin Piao and Confucius 1（北京：Foreign Languages
  Press，1974），第1-23页。} was published. In a talk with Zhang
Chunqiao and Wang Hongwen on 4 July 1973 Mao again criticized the
historical role of Confucius and went on to make the astounding
assertion that Lin Biao, in relation to his alleged veneration for
Confucius and opposition to Legalism, held the same outlook as the
KMT.\footnote{Gao Gao and Yan Jiaqi, ``Wenhua dageming'', p.~484; Hao
  Mengbi and Duan Haoran, Liushinian, p.~632. Gao Gao and Yan
  Jiaqi，``Wenhua dageming''，第484页；Hao Mengbi and Duan
  Haoran，Liushinian，第632页。} Jiang Qing reportedly attempted to
include the gist of Mao's comments into the political report presented
to the 10th National Congress, but was thwarted by Zhou
Enlai.\footnote{See Gao Wenqian, RMRB, January 5,1986.
  见高文谦，RMRB，1986年1月5日。} Shortly after the Congress, at a
meeting with the Vice-president of Egypt, the Chairman praised the
unifier of China Qinshihuang, and criticized Confucius, thus stimulating
the historiographical debate over the relative merits of the Legalists
(法家派) and the Confucianists (儒家派).\footnote{Hao Mengbi and Duan
  Haoran, Liushinian, p.~632. Hao Mengbi and Duan
  Haoran，Liushinian，第632页。} It was thus Mao who applied the theme
of historical and class retrogression to link Lin and Confucius. By this
device Mao, in the opinion of one observer, sought historical
justification and legitimacy for the Cultural Revolution. To the
Chairman, Legalism represented rebellion while Confucianism was equated
with revisionism.\footnote{Hu Hua, ``From the Fall of Lin Biao to the
  Demise of the Gang of Four'', p.~2. 胡华：``From the Fall of Lin Biao
  to the Demise of the Gang of Four''，第2页。} Mao took a further step
to shift his support from Zhou Enlai. In April 1973, in a talk with
Zhang Chunqiao and Wang Hongwen, the Chairman censured Zhou for his
management of foreign affairs. His four sentence conclusion read: major
issues are not discussed and only minor issues are submitted; if this
method is not changed it will lead to revisionism
(大事不讨论,小事天天送,此调不改动,是毗搞修政).\footnote{Hao Mengbi and
  Duan Haoran, Liushinian, p.~632. Hao Mengbi and Duan
  Haoran，Liushinian，第632页。} Mao's reference to minor issues may
have concerned Zhou's 8 March apology to foreign experts who had been
treated discourteously or victimized during the Cultural
Revolution.\footnote{Dangshi tongxun. No.~55, p.~13. Dangshi
  tongxun，第55期，第13页。} The Politburo meeting in November discussed
Zhou's transgression and other blunt criticisms of the Premier raised by
the Chairman and Jiang Qing announced that the party's 11th line
struggle had commenced. She also accused Zhou of being too impatient to
wait (迫不及待), a reference to his zeal in moving to rescind the
politics and policies of the Cultural Revolution. Mao in turn rebuked
Jiang for the extravagance of her comment.\footnote{See Hao Mengbi and
  Duan Haoran, Liushinian, p.~632. 见Hao Mengbi and Duan
  Haoran，Liushinian，第632页。} With these cryptic comments to an inner
circle of courtiers and by expressing his displeasure at Zhou's handling
of the affairs of state, Mao was inviting a return to the mobilizational
politics of the mid 1960s. He was signalling in no uncertain terms that
the retreat from the Cultural Revolution had ended. It was time to go on
the offensive.

尽管总理坚持继续推行恢复性政策，毛把林彪定义为``极右派''却预示着对周的危险。在主席重新肯定文革的支持下，中央领导层中文革主要的剩余受益者江青、张春桥、姚文元和王洪文把精力投入宣传文革的优点。年轻的王洪文于1972年9月被调入全国政治舞台\footnote{CNS,
  438 (October 5,1972), p.~1. CNS，438（1972年10月5日），第1页。}，随后于1973年5月被增补进政治局。1973年5月20日至31日举行的中共中央工作会议决定，邀请王洪文、华国锋和北京市委第一书记吴德列席政治局会议并参加政治局工作。采取这一措施，是因为林彪事件导致政治局成员损失三分之一（七人）。\footnote{See
  Dangshi tongxun, No.~55, p.~13. One source suggests that Kang Sheng
  recommended Wang to the Chairman because of his fame as a rebel. See
  Ye Yonglie, ``毛泽东如何挑选接班人'' (How did Mao Zedong choose his
  successors), Zhang Yufeng et al., 毛泽东轶事 (Anecdotes about Mao
  Zedong), (Changsha: Hunan Wenyi Chubanshe, 1989), p.~251. 见Dangshi
  tongxun，第55期，第13页。一份资料称康生因王作为造反派的名声而向主席推荐了他。见叶永烈：《毛泽东如何挑选接班人》（How
  did Mao Zedong choose his
  successors），载张玉凤等：《毛泽东轶事》（Anecdotes about Mao
  Zedong），（长沙：Hunan Wenyi Chubanshe，1989），第251页。}为了奖赏追随者并同时加强组织基础，后来被毛称为``四人帮''的四名激进派，利用文革前群众组织重建提供的机会，推动文革积极分子当选领导职位。此外，他们介入十大代表的遴选，推动其追随者的诉求。全国报刊的文章重新发起对修正主义的攻击，突出``新生事物''的成就，赞扬江青对艺术革命化的贡献，并宣传敢于``反潮流''的模范造反派事迹。\footnote{Gao
  Gao and Yan Jiaqi, ``Wenhua dageming'', pp.~477-84; Domes, China after
  the Cultural Revolution, pp.~172-75. Gao Gao and Yan Jiaqi，``Wenhua
  dageming''，第477-84页；Domes，China after the Cultural
  Revolution，第172-75页。}其中一个模范是上海青年朱克家，他在云南省安顿下来，而没有利用机会返回家乡。另一个是辽宁下乡青年张铁生。1973年，张参加大学入学考试，这是1966年以来首次举行此类考试。他的试卷除抗议制度歧视他这样的下乡青年之评论外一片空白。第三个是北京市一名小学生黄帅，她抗议受到老师虐待，并写信给市报反映自己的遭遇。所有这些例子都受到媒体突出报道，张铁生这样的造反派还被选为十大代表。另一名其事业被中央激进派接手的人是杭州的翁森鹤（见第七章）。\footnote{Gao
  Gao and Yan Jiaqi, ``Wenhua dageming'', pp.~476-77. Gao Gao and Yan
  Jiaqi，``Wenhua dageming''，第476-77页。}1973年5月，毛提出发表批判孔子的文章，这一发展将对未来产生重要影响。\footnote{Ibid,
  p.~484. 同上，第484页。}这位圣人的历史角色以及共产主义史学家对他的评价，长期以来一直是激烈争论的问题。毛本人此前曾多次对孔子表达矛盾态度。1970年代初，官方出版物零星刊登批评中国最有影响力的哲学家和教师的文章。\footnote{For
  example, an article on Lu Xun by the radical writing group from
  Shanghai, Luo Siding, published in RMRB on September 25,1971,
  expressed anti-Confucian sentiments. Another article attacking the
  sage was published in September 1972. See Domes, China after the
  Cultural Revolution, p.~172. Debates concerning the historical role of
  Confucius had been occurring in Shanghai since 1969. See Lyn T. White,
  ``Local Autonomy in China During the Cultural Revolution: The
  Theoretical Uses of an Atypical Case'', American Political Science
  Review, 70:2 (June 1976), p.~490.
  例如，上海激进写作组罗思鼎一篇关于鲁迅的文章于1971年9月25日发表于RMRB，表达了反孔情绪。另一篇攻击这位圣人的文章发表于1972年9月。见Domes，China
  after the Cultural
  Revolution，第172页。关于孔子历史角色的争论自1969年以来已在上海出现。见Lyn
  T. White，``Local Autonomy in China During the Cultural Revolution:
  The Theoretical Uses of an Atypical Case''，American Political Science
  Review，70:2（1976年6月），第490页。}1972年12月1日，广州中山大学哲学教授杨荣国在《红旗》发表一篇关于战国时期阶级斗争的文章，严厉批评孔子的历史角色。\footnote{Gao
  Gao and Yan Jiaqi, ``Wenhua dageming'', p.~478. Gao Gao and Yan
  Jiaqi，``Wenhua dageming''，第478页。}1973年8月7日，经毛审阅批准的杨荣国对孔子的重新评价发表。\footnote{See
  Hao Mengbi and Duan Haoran, Liushinian, p.~632. For a translation of
  Yang's article ``Confucius-A Thinker who Stubbornly supported the
  Slave System'', see Selected Articles Criticizing Lin Piao and
  Confucius 1 (Beijing: Foreign Languages Press, 1974), pp.~1-23. 见Hao
  Mengbi and Duan Haoran，Liushinian，第632页。杨文``Confucius-A Thinker
  who Stubbornly supported the Slave System''的译文，见Selected Articles
  Criticizing Lin Piao and Confucius 1（北京：Foreign Languages
  Press，1974），第1-23页。}1973年7月4日，毛在同张春桥和王洪文谈话时再次批评孔子的历史作用，并进一步作出惊人论断：林彪因据称崇拜孔子、反对法家，在观点上同国民党相同。\footnote{Gao
  Gao and Yan Jiaqi, ``Wenhua dageming'', p.~484; Hao Mengbi and Duan
  Haoran, Liushinian, p.~632. Gao Gao and Yan Jiaqi，``Wenhua
  dageming''，第484页；Hao Mengbi and Duan Haoran，Liushinian，第632页。}据报道，江青试图把毛谈话的要点纳入提交十大代表大会的政治报告，但被周恩来阻止。\footnote{See
  Gao Wenqian, RMRB, January 5,1986. 见高文谦，RMRB，1986年1月5日。}十大后不久，主席在会见埃及副总统时赞扬统一中国的秦始皇，批评孔子，从而推动了围绕法家派（法家派）和儒家派（儒家派）相对优劣的史学争论。\footnote{Hao
  Mengbi and Duan Haoran, Liushinian, p.~632. Hao Mengbi and Duan
  Haoran，Liushinian，第632页。}因此，是毛把历史和阶级倒退的主题用于连接林彪和孔子。一位观察者认为，毛借此为文革寻求历史论证和合法性。对主席而言，法家代表造反，而儒家则等同于修正主义。\footnote{Hu
  Hua, ``From the Fall of Lin Biao to the Demise of the Gang of Four'',
  p.~2. 胡华：``From the Fall of Lin Biao to the Demise of the Gang of
  Four''，第2页。}毛进一步把支持从周恩来身上转移。1973年4月，主席在同张春桥和王洪文谈话时，批评周处理外交事务。他的四句结论是：大事不讨论，小事天天送，此调不改动，势必搞修正主义（大事不讨论，小事天天送，此调不改动，势必搞修正主义）。\footnote{Hao
  Mengbi and Duan Haoran, Liushinian, p.~632. Hao Mengbi and Duan
  Haoran，Liushinian，第632页。}毛提到的小事，可能涉及周3月8日就文革中受到无礼对待或迫害的外国专家道歉。\footnote{Dangshi
  tongxun. No.~55, p.~13. Dangshi tongxun，第55期，第13页。}11月的政治局会议讨论了周的过错以及主席对总理提出的其他直率批评，江青宣布党的第十一次路线斗争已经开始。她还指责周``迫不及待''（迫不及待），指的是他急于取消文革政治和政策。毛反过来斥责江青此言过分。\footnote{See
  Hao Mengbi and Duan Haoran, Liushinian, p.~632. 见Hao Mengbi and Duan
  Haoran，Liushinian，第632页。}通过对亲信小圈子作出这些隐晦评论，并表达他对周处理国家事务的不满，毛是在邀请回到1960年代中期的动员式政治。他明确无误地发出信号：从文革退却的阶段已经结束。现在是转入进攻的时候。

\section{Impact in Zhejiang}\label{impact-in-zhejiang}

\section{对浙江的影响}\label{ux5bf9ux6d59ux6c5fux7684ux5f71ux54cd}

To assess his leadership capabilities, the central authorities
despatched Wang Hongwen on a provincial familiarization and inspection
tour at the end of 1972. On 18 December 1972 Wang visited a cotton mill
in Xinjiang and reportedly said to a delegation of worker
representatives that ``if I was a worker I wouldn't go to work because
that's restoring capitalism!''\footnote{Gao Gao and Yan Jiaqi, ``Wenhua
  dageming'', pp.~476-77. Gao Gao and Yan Jiaqi，``Wenhua
  dageming''，第476-77页。} The following month he journeyed to
Zhejiang. Wang's visit to Hangzhou in January 1973 paved the way for the
resurgence of the Cultural Revolution radicals in Zhejiang.\footnote{Ibid,
  p.~482; Zhang Yongshengde zuizheng cailiao, p.~48.
  同上，第482页；Zhang Yongshengde zuizheng cailiao，第48页。} Included
in Wang's brief was the difficult task of explaining to provincial
leaders such as Tan Qilong and Tie Ying Mao's sudden change of mind
regarding the nature of Lin Biao's political errors. In speeches he
delivered at factories such as the Hangzhou silk complex and other units
in Hangzhou and Jiaxing\footnote{A Report of early February 1973
  mentioned the presence of Tan Qilong and Tie Ying at a meeting
  concerning agricultural word in Jiaxing district. ZPS, February
  6,1973, SWB/FE/4216/BII/1-2.
  1973年2月初的一篇报道提到，谭启龙和铁瑛出席了嘉兴地区一次关于农业工作的会议。ZPS，1973年2月6日，SWB/FE/4216/BII/1-2。}
-- a city midway between Shanghai and Hangzhou -- Wang questioned the
correctness of the campaign of the previous year against Lin Biao. Wang
reportedly claimed that the movement in Zhejiang had gone too far
(''搞了扩大化''), had suppressed a faction and become ``entangled in old
historical scores''. He announced that he had come to ``pour some cold
water on this''. Wang met Weng Senhe twice on his visits to the silk
complex and allegedly encouraged him to prepare for further
battles.\footnote{Chen Zuolin (deputy-secretary of the CCP ZPC),
  打烂``四人帮'',浙江有希望 (With the smashing of the Gang of Four
  Zhejiang has prospects), (excerpts from a speech at the 2nd National
  Conference on Learning from Dazhai in agriculture), HZRB, December
  18,1976; HZRB, December 31, 1976; HZRB, September 17,1977; CCP HMC
  Propaganda Department criticism group, ``The Gang of Four's fiendish
  hatchet-man Weng Senhe'', HZRB, December 7,1977; CCP ZPC, ``Report on
  the examination concerning Weng Senhe'', June 19,1977, Document of the
  CC of the CCP, 中发 (1977) No.~37, I\&S, 14:11 (1978), p.~110; visit
  to Hangzhou Silk Complex, January 12,1977; HZRB, August 14,1978.
  陈作霖（中共浙江省委副书记）：《打烂``四人帮'',浙江有希望》（With the
  smashing of the Gang of Four Zhejiang has
  prospects）（在第二次全国农业学大寨会议上的讲话摘录），HZRB，1976年12月18日；HZRB，1976年12月31日；HZRB，1977年9月17日；中共杭州市委宣传部批判组：``The
  Gang of Four's fiendish hatchet-man Weng
  Senhe''，HZRB，1977年12月7日；中共浙江省委：``Report on the
  examination concerning Weng
  Senhe''，1977年6月19日，中共中央文件，中发（1977）第37号，I\&S，14:11（1978），第110页；1977年1月12日访问杭州丝绸联合企业；HZRB，1978年8月14日。}
Wang and Weng had shared somewhat similar personal histories in the
Cultural Revolution. Both men had rebelled early against the party
leaderships of their textile mills and both had played major roles in
the leadership of workers' organizations in their respective cities. By
1973 Wang had achieved high political office and undoubtedly was an
inspiration to Weng. At his meeting with Weng, Wang urged him to
``strike back and rely on yourself in the struggle''
(要翻身靠自己去斗).\footnote{Gao Gao and Yan Jiaqi, ``Wenhua dageming'',
  p.~482. Gao Gao and Yan Jiaqi，``Wenhua dageming''，第482页。} Wang's
declaration that the rebels must gain positions of power would have
struck a chord with Weng and his associates. Weng later wrote of his
admiration for the power Wang had obtained in Shanghai and said that he
``took his hat off'' (佩服得五体投地) to him and wished to emulate
him.\footnote{Weng Senhe bufen zuizheng cailiao, p.~21. Weng Senhe bufen
  zuizheng cailiao，第21页。} With Wang's arrival the opportune moment
seemed to have arrived. In February, after Wang's departure, the
significance of the visit as well as the prospects it opened up was
discussed in meetings held at a rebel's home. One of Weng's colleagues,
who was present at the meetings, recalled later that Weng had described
Wang Hongwen's visit as a signal for the rebels to reenter the political
fray.\footnote{Ibid, p.~14. See also Zhang Yongshengde zuizheng cailiao,
  p.~48. 同上，第14页。另见Zhang Yongshengde zuizheng cailiao，第48页。}
The upshot of Wang's visit, apart from the devastating impact it was to
have on provincial politics, was a gradual redirection of propaganda in
line with Mao's wishes. On occasions, however, the propaganda appeared
highly unconvincing, perhaps reflecting both the confusion in the minds
of local officials and their reluctance to fall into line. A commentary
in Zhejiang Daily of January 1973 stated that ``swindlers like Liu
Shaoqi'' sometimes appeared leftist in their actions but actually they
were extreme-right.\footnote{ZJRB, January 20,1973, p.~1. Zhejiang was
  reportedly the seventh province to describe Lin as an ultra-rightist.
  See CNS, 459 (March 15, 1973), pp.~3-4.
  ZJRB，1973年1月20日，第1页。据报道，浙江是第七个把林描述为极右派的省份。见CNS，459（1973年3月15日），第3-4页。}
However the examples given to substantiate the argument were clearly all
policy deviations to the left. The PLA in Zhejiang seemed equally
confused or unwilling to go along.\footnote{ZJRB, February 26,1973,
  p.~1. ZJRB，1973年2月26日，第1页。} The abruptness of the 180 degree
turn necessitated the despatch of large numbers of cadres to the
grassroots to clarify the position with local officials.\footnote{ZJRB,
  March 5,1973, pp.~1, 3; ZJRB, February 24,1973, p.~1; ZPS, April
  17,1973, URS, 71:10, pp.~134-6. Briefings occurred in other provinces.
  See CNS, 461 (March 29, 1973).
  ZJRB，1973年3月5日，第1、3页；ZJRB，1973年2月24日，第1页；ZPS，1973年4月17日，URS，71:10，第134-6页。其他省份也举行了通报。见CNS，461（1973年3月29日）。}
No doubt they were bemused to learn that Lin Biao had opposed the
Cultural Revolution. Even as late as August 1973, on the eve of the 10th
Congress of the CCP, it was admitted that some people still did not
understand the substantive reasons behind the campaign against
Lin.\footnote{ZJRB, August 6,1973, p.~1. ZJRB，1973年8月6日，第1页。}
They were not helped by the blatant disjuncture between radical rhetoric
and moderate policy which prevailed in 1973.

为评估王洪文的领导能力，中央于1972年底派他赴各省熟悉情况并巡视。1972年12月18日，王参观新疆一家棉纺厂，据说他对一个工人代表团说：``如果我是工人，我就不上班，因为那是在复辟资本主义！''\footnote{Gao
  Gao and Yan Jiaqi, ``Wenhua dageming'', pp.~476-77. Gao Gao and Yan
  Jiaqi，``Wenhua dageming''，第476-77页。}次月，他前往浙江。王1973年1月访问杭州，为浙江文革激进派的复兴铺平了道路。\footnote{Ibid,
  p.~482; Zhang Yongshengde zuizheng cailiao, p.~48.
  同上，第482页；Zhang Yongshengde zuizheng cailiao，第48页。}王的任务包括一项困难工作：向谭启龙、铁瑛等省级领导人解释毛对林彪政治错误性质的突然改判。王在杭州丝绸联合企业等工厂，以及杭州和嘉兴\footnote{A
  Report of early February 1973 mentioned the presence of Tan Qilong and
  Tie Ying at a meeting concerning agricultural word in Jiaxing
  district. ZPS, February 6,1973, SWB/FE/4216/BII/1-2.
  1973年2月初的一篇报道提到，谭启龙和铁瑛出席了嘉兴地区一次关于农业工作的会议。ZPS，1973年2月6日，SWB/FE/4216/BII/1-2。}（位于上海和杭州之间的城市）其他单位发表讲话，质疑上一年反林彪运动的正确性。据报道，王声称浙江运动走得太远（``搞了扩大化''），压制了一个派别，并``纠缠历史旧账''。他宣布自己来是要``泼一点冷水''。王在访问丝绸联合企业时两次会见翁森鹤，据称鼓励他准备进一步战斗。\footnote{Chen
  Zuolin (deputy-secretary of the CCP ZPC), 打烂``四人帮'',浙江有希望
  (With the smashing of the Gang of Four Zhejiang has prospects),
  (excerpts from a speech at the 2nd National Conference on Learning
  from Dazhai in agriculture), HZRB, December 18,1976; HZRB, December
  31, 1976; HZRB, September 17,1977; CCP HMC Propaganda Department
  criticism group, ``The Gang of Four's fiendish hatchet-man Weng
  Senhe'', HZRB, December 7,1977; CCP ZPC, ``Report on the examination
  concerning Weng Senhe'', June 19,1977, Document of the CC of the CCP,
  中发 (1977) No.~37, I\&S, 14:11 (1978), p.~110; visit to Hangzhou Silk
  Complex, January 12,1977; HZRB, August 14,1978.
  陈作霖（中共浙江省委副书记）：《打烂``四人帮'',浙江有希望》（With the
  smashing of the Gang of Four Zhejiang has
  prospects）（在第二次全国农业学大寨会议上的讲话摘录），HZRB，1976年12月18日；HZRB，1976年12月31日；HZRB，1977年9月17日；中共杭州市委宣传部批判组：``The
  Gang of Four's fiendish hatchet-man Weng
  Senhe''，HZRB，1977年12月7日；中共浙江省委：``Report on the
  examination concerning Weng
  Senhe''，1977年6月19日，中共中央文件，中发（1977）第37号，I\&S，14:11（1978），第110页；1977年1月12日访问杭州丝绸联合企业；HZRB，1978年8月14日。}王和翁在文革中有相当相似的个人经历。二人都很早就起来反对各自纺织厂的党委领导，并都在各自城市工人组织领导层中发挥重要作用。到1973年，王已取得很高的政治职位，无疑对翁是一种鼓舞。王同翁会面时敦促他说：``要翻身靠自己去斗''（要翻身靠自己去斗）。\footnote{Gao
  Gao and Yan Jiaqi, ``Wenhua dageming'', p.~482. Gao Gao and Yan
  Jiaqi，``Wenhua dageming''，第482页。}王宣称造反派必须取得权力职位，这必定会引起翁及其同伴共鸣。翁后来写到自己钦佩王在上海取得的权力，并说他对王``佩服得五体投地''（佩服得五体投地），希望效仿他。\footnote{Weng
  Senhe bufen zuizheng cailiao, p.~21. Weng Senhe bufen zuizheng
  cailiao，第21页。}随着王的到来，机会似乎已经到来。2月，王离开后，造反派在一名造反派家中开会讨论这次访问的意义及其打开的前景。翁的一名同事出席了这些会议，后来回忆说，翁把王洪文的访问描述为造反派重新进入政治斗争的信号。\footnote{Ibid,
  p.~14. See also Zhang Yongshengde zuizheng cailiao, p.~48.
  同上，第14页。另见Zhang Yongshengde zuizheng cailiao，第48页。}王访问的结果，除了对省级政治产生毁灭性影响外，还包括宣传逐步按照毛的愿望重新定向。然而，有时宣传显得很不可信，这也许反映了地方官员思想上的困惑和他们不愿跟进。《浙江日报》1973年1月一篇评论说，``刘少奇一类骗子''有时行动上表现得很左，实际上却是极右。\footnote{ZJRB,
  January 20,1973, p.~1. Zhejiang was reportedly the seventh province to
  describe Lin as an ultra-rightist. See CNS, 459 (March 15, 1973),
  pp.~3-4.
  ZJRB，1973年1月20日，第1页。据报道，浙江是第七个把林描述为极右派的省份。见CNS，459（1973年3月15日），第3-4页。}然而，用来支撑这一论点的例子显然全都是向左偏离的政策。浙江的解放军似乎同样困惑，或不愿跟进。\footnote{ZJRB,
  February 26,1973, p.~1. ZJRB，1973年2月26日，第1页。}这种180度急转弯的突兀性，迫使当局派出大批干部下基层，向地方官员说明立场。\footnote{ZJRB,
  March 5,1973, pp.~1, 3; ZJRB, February 24,1973, p.~1; ZPS, April
  17,1973, URS, 71:10, pp.~134-6. Briefings occurred in other provinces.
  See CNS, 461 (March 29, 1973).
  ZJRB，1973年3月5日，第1、3页；ZJRB，1973年2月24日，第1页；ZPS，1973年4月17日，URS，71:10，第134-6页。其他省份也举行了通报。见CNS，461（1973年3月29日）。}他们得知林彪曾反对文革时，无疑感到茫然。甚至到1973年8月，在中共十大前夕，仍有人承认一些人还不理解批林运动背后的实质原因。\footnote{ZJRB,
  August 6,1973, p.~1. ZJRB，1973年8月6日，第1页。}1973年激进言辞和温和政策之间显而易见的脱节，并没有帮助他们理解。

\section{The Reconstitution of Mass Organizations,
1973}\label{the-reconstitution-of-mass-organizations-1973}

\section{群众组织的重建，1973年}\label{ux7fa4ux4f17ux7ec4ux7ec7ux7684ux91cdux5efa1973ux5e74}

Before the convention of the 10th National Congress, provincial party
committees were instructed to reconstitute pre-Cultural Revolution mass
organizations which had become defunct in 1967. To some extent, workers
congresses and Red Guard congresses had filled the roles formerly played
by trade unions and the CYL. However, membership of the workers
congresses was generally restricted to Cultural Revolution rebels and
activists who tended to be younger and less career oriented. Veteran
workers had thus lost a body to represent their interests, however
inadequately. Similarly Red Guard congresses attracted a certain type of
political activist and did not cater so well for youth whose interests
were of a less politically centered kind. Another not unimportant
consideration was that the reconstitution of the mass organizations
would further signal the return to political normality. It will be
recalled that the central propaganda and organization group established
under the Politburo in late 1970 had been charged with responsibility
for overseeing the work of mass organizations. In April 1970 reference
had been made to a county Youth League branch in Heilongjiang province
followed, in February 1971, by mention of a Trade Union
Council.\footnote{CNS, 64: 5 (July 16, 1971), p.~54; Domes, China after
  the Cultural Revolution, p.~87.
  CNS，64:5（1971年7月16日），第54页；Domes，China after the Cultural
  Revolution，第87页。} The People's Daily New Year's day editorial for
1973 highlighted reconstruction of the mass organizations as one of the
principal tasks for the coming year.\footnote{Domes, China after the
  Cultural Revolution, p.~176. Domes，China after the Cultural
  Revolution，第176页。} In Zhejiang, two reports of 1971 and 1972
mentioned party activity among youth.\footnote{ZPS, May 25,1971, CNS,
  64: 5, pp.~58-9; ZPS, April 2, 1972, CNS, 67:10, pp.~141-42.
  ZPS，1971年5月25日，CNS，64:5，第58-9页；ZPS，1972年4月2日，CNS，67:10，第141-42页。}
Of the three mass organizations it was the struggle for control of the
Trade Union Council which galvanized the contending parties in Zhejiang
into action. In April 1973 the Central Committee had issued directives
to ``rectify and cleanse'' workers' congresses as a prelude to the
re-establishment of trade union councils. A People's Daily editorial
outlined the tasks for trade unions after the convocation of congresses
in Beijing and Shanghai.\footnote{CNS, 464 (April 26, 1973), pp.~1-4; Li
  Ming-hua, ``An Analysis of the Reconstruction of Trade Unions on the
  Mainland'', I\&S, 9:12 (September 1973), pp.~39-40.
  CNS，464（1973年4月26日），第1-4页；Li Ming-hua，``An Analysis of the
  Reconstruction of Trade Unions on the
  Mainland''，I\&S，9:12（1973年9月），第39-40页。} Criticism of
revisionism (Lin) and rectification of work style was given top
priority. The crucial, unsettled question was whether Lin was criticized
as an ultra-leftist saboteur, responsible for the anarchy and disruption
of the Cultural Revolution, or whether he was lambasted as an
ultra-right revisionist who had wrecked an otherwise glorious
revolution. Clearly, the CCP was divided on the issue and the editorial
reflected the cleavage. Unity among the working class was advocated, but
it was a militant unity. Professional knowledge was required but it was
to be guided by proletarian politics. The editorial called on workers to
increase labor productivity and fulfil state plans while at the same
time urging them to participate in management. The equivocal tone of the
editorial suggested that the writers were undecided where priorities lay
and in order to appease both the ``reds'' and the ``experts'' they were
being eclectic in their recommendations. In Zhejiang, Weng Senhe and He
Xianchun, who were influential figures in the provincial and municipal
workers' congresses (工代会) wished to model the restructured trade
union councils (工会) on existing bodies. They proposed the purely
formal measure of removing the character dai (代) from the organizations
under their control, otherwise leaving their organizations and
leaderships intact. Weng and He boycotted meetings of the preparatory
group appointed by the ZPC to convene the provincial trade union
congress because according to Weng, ``at the moment we certainly can't
get much out of it'' (捞不到什么). At the Hangzhou Cable Factory He
Xianchun tried to implement the window-dressing changes that he and Weng
favored, only to succeed in causing division among the workers present
at the meeting.\footnote{ZPTUC, ``Trade unions controlling the party was
  a plot to seize power'', HZRB, March 29, 1977; ZPTUC criticism group,
  ``The gang of four were the ringleaders in sabotaging the
  revolutionary unity of Zhejiang's working class'', HZRB, January
  22,1977. 浙江省总工会：``Trade unions controlling the party was a plot
  to seize power''，HZRB，1977年3月29日；浙江省总工会批判组：``The gang
  of four were the ringleaders in sabotaging the revolutionary unity of
  Zhejiang's working class''，HZRB，1977年1月22日。} Weng and He may
have delayed the convocation of the trade union congress by their
tactics of non-cooperation but they could not prevent its eventual
opening. In late April 1973 the ZPC convened a forum of the preparatory
group\footnote{ZJRB, April 28, 1973, p.~1. ZJRB，1973年4月28日，第1页。}
at which veteran union leaders reportedly criticized Weng and He's
factional schemes.\footnote{HZRB, January 22, 1977, March 29,1977.
  HZRB，1977年1月22日、1977年3月29日。} Three months later, in July
1973, 1,600 delegates and 100 invited representatives filed into the
opening session of the 6th provincial congress of the ZPTUC. The
delegates had been selected from four specific groups; activists of the
Cultural Revolution and the ``criticize revisionism and rectify the
style of work'' campaign, model workers, advanced producers and members
of Mao Zedong Thought propaganda teams. Thirty percent of the delegates
were female, 25\% were classified as young workers and 70\% of the
delegates came from the industrial work force. If delegates from the
four categories received equal representation at the congress. Cultural
Revolution activists would probably not have gained a majority. In his
opening speech to the Congress, Tan Qilong emphasized the importance of
unity among the working class.\footnote{ZJRB, July 4,1973, p.~1.
  ZJRB，1973年7月4日，第1页。} But this appeared to be a unity, as Tie
Ying stressed in his closing address, based on a clear recognition of
right and wrong. Tie stated provocatively that

在十大召开前，省级党委接到指示，重建文革前的群众组织，这些组织在1967年已停止运作。在某种程度上，工人代表大会和红卫兵代表大会填补了工会和共青团此前扮演的角色。然而，工人代表大会的成员通常限于文革造反派和积极分子，他们往往更年轻，也不那么以职业晋升为导向。因此，老工人失去了一个代表其利益的机构，尽管这种代表原本也并不充分。同样，红卫兵代表大会吸引的是某类政治积极分子，对那些兴趣不那么以政治为中心的青年照顾得并不好。另一个并非不重要的考虑是，群众组织的重建将进一步标志着政治正常化的回归。应当回忆，1970年末在政治局下设立的中央组织宣传组，被赋予监督群众组织工作的责任。1970年4月曾提到黑龙江省一个县级共青团支部，随后在1971年2月提到一个工会委员会。\footnote{CNS,
  64: 5 (July 16, 1971), p.~54; Domes, China after the Cultural
  Revolution, p.~87. CNS，64:5（1971年7月16日），第54页；Domes，China
  after the Cultural Revolution，第87页。}《人民日报》1973年元旦社论把重建群众组织列为来年主要任务之一。\footnote{Domes,
  China after the Cultural Revolution, p.~176. Domes，China after the
  Cultural Revolution，第176页。}在浙江，1971年和1972年的两篇报道提到青年中的党务活动。\footnote{ZPS,
  May 25,1971, CNS, 64: 5, pp.~58-9; ZPS, April 2, 1972, CNS, 67:10,
  pp.~141-42.
  ZPS，1971年5月25日，CNS，64:5，第58-9页；ZPS，1972年4月2日，CNS，67:10，第141-42页。}在三个群众组织中，正是围绕控制工会委员会的斗争，促使浙江各竞争方行动起来。1973年4月，中央委员会发布指示，要求``整顿和清理''工人代表大会，作为重建工会委员会的前奏。北京和上海召开代表大会后，《人民日报》一篇社论概述了工会的任务。\footnote{CNS,
  464 (April 26, 1973), pp.~1-4; Li Ming-hua, ``An Analysis of the
  Reconstruction of Trade Unions on the Mainland'', I\&S, 9:12
  (September 1973), pp.~39-40. CNS，464（1973年4月26日），第1-4页；Li
  Ming-hua，``An Analysis of the Reconstruction of Trade Unions on the
  Mainland''，I\&S，9:12（1973年9月），第39-40页。}批判修正主义（林）和整顿作风被置于首要位置。关键而未解决的问题是：究竟把林批判为极左破坏者，认定他对文革的无政府状态和破坏负有责任，还是把他痛斥为破坏了本来光辉革命的极右修正主义者。显然，中共在这一问题上存在分歧，而社论反映了这种裂痕。社论主张工人阶级团结，但这是一种战斗性的团结。专业知识是需要的，但必须由无产阶级政治指导。社论号召工人提高劳动生产率并完成国家计划，同时敦促他们参与管理。社论含混的语调表明，作者不确定优先事项在哪里；为安抚``红''与``专''双方，他们在建议中采取了折衷拼合的方式。在浙江，翁森鹤和贺贤春是省市工人代表大会（工代会）中有影响的人物，他们希望以现有机构为模式重组工会委员会（工会）。他们提出一种纯形式的措施：从其控制的组织名称中去掉``代''字（代），其他方面则保持组织和领导层不变。翁和贺抵制浙江省委指定召开省工会代表大会的筹备组会议，因为按翁的话说，``眼下我们肯定捞不到什么''（捞不到什么）。在杭州电缆厂，贺贤春试图落实他和翁赞成的装点门面式变动，结果只是造成与会工人分裂。\footnote{ZPTUC,
  ``Trade unions controlling the party was a plot to seize power'',
  HZRB, March 29, 1977; ZPTUC criticism group, ``The gang of four were
  the ringleaders in sabotaging the revolutionary unity of Zhejiang's
  working class'', HZRB, January 22,1977. 浙江省总工会：``Trade unions
  controlling the party was a plot to seize
  power''，HZRB，1977年3月29日；浙江省总工会批判组：``The gang of four
  were the ringleaders in sabotaging the revolutionary unity of
  Zhejiang's working class''，HZRB，1977年1月22日。}翁和贺也许通过不合作策略拖延了工会代表大会的召开，但无法阻止其最终开幕。1973年4月下旬，浙江省委召开筹备组座谈会\footnote{ZJRB,
  April 28, 1973, p.~1. ZJRB，1973年4月28日，第1页。}，据说会上老工会领导人批评了翁、贺的派系图谋。\footnote{HZRB,
  January 22, 1977, March 29,1977. HZRB，1977年1月22日、1977年3月29日。}三个月后，1973年7月，1600名代表和100名特邀代表进入浙江省总工会第六次代表大会开幕式会场。代表来自四个特定群体：文革积极分子和``批修整风''运动积极分子、劳动模范、先进生产者以及毛泽东思想宣传队成员。30\%的代表为女性，25\%被归类为青年工人，70\%来自产业工人队伍。如果四类代表在大会上得到同等代表性，文革积极分子大概不会占多数。谭启龙在大会开幕词中强调工人阶级团结的重要性。\footnote{ZJRB,
  July 4,1973, p.~1. ZJRB，1973年7月4日，第1页。}但是，正如铁瑛在闭幕讲话中所强调的，这种团结似乎是建立在明确辨别是非的基础上。铁以挑衅性的口吻说：

\begin{quote}
We must strive to unite not only with comrades who share our views but
also with those who hold different views.We must also work on uniting
with those who formerly opposed us and have since been proved
wrong.\footnote{ZJRB, July 10,1973, p.~1. One observer viewed such
  phrases as an expression of the party's efforts to reconcile sectoral
  interests into a ``public interest'' and thus reassert its
  ``monopolistic control over the organization of interests''. See
  Lowell Dittmer, ``The Radical Critique of Political Interest,
  1966-1978'', Modern China 6: 4 (1980), pp.~384-85.
  ZJRB，1973年7月10日，第1页。一位观察者认为，这类措辞体现了党把部门利益调和为``公共利益''的努力，并由此重新确立其``对利益组织的垄断性控制''。见Lowell
  Dittmer，``The Radical Critique of Political Interest,
  1966-1978''，Modern China 6:4（1980），第384-85页。}

我们不仅要努力团结同我们意见相同的同志，也要团结持有不同意见的人。我们还必须努力团结那些过去反对过我们、后来已被证明犯了错误的人。\footnote{ZJRB,
  July 10,1973, p.~1. One observer viewed such phrases as an expression
  of the party's efforts to reconcile sectoral interests into a ``public
  interest'' and thus reassert its ``monopolistic control over the
  organization of interests''. See Lowell Dittmer, ``The Radical
  Critique of Political Interest, 1966-1978'', Modern China 6: 4 (1980),
  pp.~384-85.
  ZJRB，1973年7月10日，第1页。一位观察者认为，这类措辞体现了党把部门利益调和为``公共利益''的努力，并由此重新确立其``对利益组织的垄断性控制''。见Lowell
  Dittmer，``The Radical Critique of Political Interest,
  1966-1978''，Modern China 6:4（1980），第384-85页。}
\end{quote}

The reference to those who ``have since been proved wrong'', was clearly
directed toward Zhang, Weng, He and the old United Headquarters
organization. The verdict on the events of 1967-69 had been reversed in
favour of Red Storm. However, Tie's reminder of past divisions among the
working class of Zhejiang was not conducive to overcoming factionalism
and seemed to conflict with Tan's earlier unqualified appeal for unity.
The lack of unity at the congress was accentuated by the fact that it
was the naval cadre Chai Qikun, and not a trade union leader, who
delivered the main report. In it Chai mentioned that since the Cultural
Revolution a large number of advanced workers had been admitted into the
CCP,\footnote{This trend was evident across China in 1973. See URS, 72:3
  (July 10,1973), p.~28.
  1973年，这一趋势在全中国都很明显。见URS，72:3（1973年7月10日），第28页。}
and that many had been promoted to leading positions at various
levels.\footnote{ZJRB, July 4,1973. ZJRB，1973年7月4日。} However, the
real test of the rebels' power came in the election of the provincial
trade union leadership. The delegates elected a 120-strong committee (of
which 30\% were women) which chose a 27 member standing committee. This
body in turn nominated an executive consisting of nine members -- a
Chairwoman and eight Vice-chairmen. Jiang Baodi was elected Chairwoman
of the ZPTUC. Jiang was known as Hangzhou's Wu Guixian after the female
textile worker from Xian, who as a result of the Cultural Revolution,
rocketed to prominence, becoming an alternate member of the Politburo of
the CCP 10th CC and later Vice-Premier of the State Council. Like Wu,
Jiang was also a textile worker. She had belonged to the United
Headquarters faction but represented it more in a figurehead capacity.
Jiang was no rebel in the Weng Senhe mould. She was also a standing
committee member of both the CCP ZPC and ZPRC and an alternate member of
the CCP 9th and later the 10th CC. The two senior Vice-chairmen were
veteran trade union officials, Xu Wanzhen and Chen Yousheng. Xu was the
party secretary of the Zhejiang Hemp Mill, widow of ZPC propaganda
department head Jin Tao who had died in 1966, and sister-in-law of Chen
Weida (whose wife, Xu Wenhua, was at that time party secretary of the
Foreign Languages Department of Hangzhou University). Xu Wanzhen had
been Vice-chairman of the pre-Cultural Revolution ZPTUC as well as
executive member of the All-China Federation of Trade Unions.\footnote{See
  her article published in 工人日报 (Workers' Daily), June 7,1964
  concerning the PLA, transl. in SCMP, 3248, pp.~1-6.
  见她关于解放军的文章，发表于《工人日报》（Workers'
  Daily），1964年6月7日；译文载SCMP，3248，第1-6页。} Chen Yousheng was
a boiler expert who had been selected as a model worker in the Great
Leap Forward. He had been a delegate to the 2nd and 3rd NPC. Another
veteran model worker who was elected Vice-chairman was ``Red diving
engineer'' Yao Xingen, like Chen a person whose fame was a product of
the Great Leap Forward.\footnote{See ZJRB, January 27,1960, SCMP, 2207
  (March 2,1960), p.28; Zhou Jianren, ``Work Report of the Zhejiang
  Provincial People's Council'', February 18,1960, in ZJRB, April
  6,1960, SCMP, 2281 (June 21,1960), p.~25.
  见ZJRB，1960年1月27日，SCMP，2207（1960年3月2日），第28页；周建人：``Work
  Report of the Zhejiang Provincial People's
  Council''，1960年2月18日，载ZJRB，1960年4月6日，SCMP，2281（1960年6月21日），第25页。}
A long-serving trade union bureaucrat, Zhao Jingtang, was also elected
Vice-chairman. The remaining four Vice-chairmen were all Cultural
Revolution rebel leaders. Fang Jianwen, who, in 1969 had been singled
out for attack in a Central Committee document (see chapter four), was
one of the four Vice-Chairmen. He was the leader of that faction within
Red Storm which had refused to accept the terms of the May 1969
agreement with United Headquarters. Four years later Fang's obduracy had
paid dividends and it is possible that Tie Ying had pressed for his
inclusion in the leadership. Another worker rebel leader who joined the
executive was Guo Zhisong. He was known as ``Guo the lame'' after a
violent incident in which he had been severely beaten. In the 1960s Guo
had been a leader of United Headquarters but later fell out with the
organization. He was to play a notable but publicly unheralded role in
the factional disturbances of 1973-74.

这里提到那些``后来已被证明犯了错误''的人，显然是针对张、翁、贺以及旧``联合总部''组织。对1967-69年事件的定性已经被推翻，改为有利于``红暴''。然而，铁提醒浙江工人阶级过去的分裂，并不利于克服派性，而且似乎同谭此前不附加条件的团结号召相冲突。大会缺乏团结这一点，又因作主要报告的是海军干部柴启琨而非工会领导人而更加突出。柴在报告中提到，文革以来已有大量先进工人被吸收入中共\footnote{This
  trend was evident across China in 1973. See URS, 72:3 (July 10,1973),
  p.~28.
  1973年，这一趋势在全中国都很明显。见URS，72:3（1973年7月10日），第28页。}，许多人还被提拔到各级领导岗位。\footnote{ZJRB,
  July 4,1973. ZJRB，1973年7月4日。}然而，检验造反派权力的真正考验，出现在省总工会领导层选举中。代表们选出一个120人的委员会（其中30\%为女性），该委员会选出一个27人的常委会。常委会又提名一个由九人组成的执行机构，包括一名女主席和八名副主席。蒋宝娣当选浙江省总工会主席。蒋被称为杭州的吴桂贤；吴是西安女纺织工人，因文革而迅速显赫，成为中共十大中央政治局候补委员，后来任国务院副总理。像吴一样，蒋也是纺织工人。她曾属于``联合总部''派别，但更多是以名义代表身份出现。蒋并不是翁森鹤类型的造反派。她也是中共浙江省委和浙江省革命委员会常委，并是中共九大、后来又是十大中央候补委员。两名资深副主席是老工会干部徐婉珍和陈友生。徐是浙江麻纺厂党委书记，1966年去世的省委宣传部部长金涛的遗孀，也是陈伟达的姻亲（陈的妻子徐文华当时是杭州大学外语系党委书记）。徐婉珍曾任文革前浙江省总工会副主席以及中华全国总工会执行委员。\footnote{See
  her article published in 工人日报 (Workers' Daily), June 7,1964
  concerning the PLA, transl. in SCMP, 3248, pp.~1-6.
  见她关于解放军的文章，发表于《工人日报》（Workers'
  Daily），1964年6月7日；译文载SCMP，3248，第1-6页。}陈友生是一名锅炉专家，曾在大跃进中被选为劳动模范。他曾是第二、三届全国人大代表。另一名当选副主席的老劳动模范是``红色潜水工程师''姚信恩，同陈一样，他的名声也是大跃进的产物。\footnote{See
  ZJRB, January 27,1960, SCMP, 2207 (March 2,1960), p.28; Zhou Jianren,
  ``Work Report of the Zhejiang Provincial People's Council'', February
  18,1960, in ZJRB, April 6,1960, SCMP, 2281 (June 21,1960), p.~25.
  见ZJRB，1960年1月27日，SCMP，2207（1960年3月2日），第28页；周建人：``Work
  Report of the Zhejiang Provincial People's
  Council''，1960年2月18日，载ZJRB，1960年4月6日，SCMP，2281（1960年6月21日），第25页。}长期任职的工会官僚赵景堂也当选副主席。其余四名副主席都是文革造反派领导人。方剑文是四名副主席之一，他在1969年曾被中央文件点名攻击（见第四章）。他是``红暴''内部拒绝接受1969年5月同``联合总部''协议条款的那个派别的领导人。四年后，方的顽固获得回报，铁瑛有可能曾推动把他纳入领导层。另一名进入执行机构的工人造反派领导人是郭志松。他因一次严重挨打的暴力事件而被称为``瘸子郭''。1960年代，郭曾是``联合总部''领导人，后来同该组织闹翻。他将在1973-74年的派系骚乱中发挥重要但未被公开表彰的作用。

Weng Senhe and He Xianchun were also elected Vice-chairmen of the ZPTUC.
The election of the four rebel leaders illustrated the strength of their
factional support in the province, and their claim for recognition was
one that the party leadership could ignore only at its peril. If, as
local propaganda reports later claimed, Weng and He used their positions
as Vice-chairmen of the ZPTUC as a springboard to launch their assault
on the provincial authorities, it was not due to their numerical
strength on its executive.\footnote{One article specifically denied that
  He and Weng had seized control of the ZPTUC leadership. See, HZRB,
  January 22, 1977.
  有一篇文章明确否认何和翁夺取了浙江省总工会领导权。见HZRB，1977年1月22日。}
Through superior tactics, force of personality and an ability to
mobilize their supporters, they may have outmaneuvered their opponents.
But they controlled another organization which played a greater role in
the political struggle of 1974. This was the Hangzhou workers' congress
led by He Xianchun. He and his followers successfully resisted all
efforts to disband the organization and it continued to function as an
officially-recognized body and one which would prove a thorn in the side
of the provincial leadership.\footnote{For further comparative detail of
  the provincial-level trade union congresses, see, CNS, 462, 464, 471;
  Li Ming-hua, ``An Analysis of the Reconstruction'', pp.~35-44; and
  analyses of the trade unions and factionalism among the working class,
  CNA, 940 and 941 (November 16 and 23,1973).
  关于省级工会代表大会的进一步比较细节，见CNS，462、464、471；Li
  Ming-hua，``An Analysis of the
  Reconstruction''，第35-44页；以及关于工会和工人阶级派性分析，CNA，940和941（1973年11月16日、23日）。}

翁森鹤和贺贤春也当选浙江省总工会副主席。四名造反派领导人的当选，说明他们在省内派系支持的力量，也说明他们要求获得承认的主张是党领导层不能轻易忽视的，否则就要自担风险。如果正如地方宣传报道后来所称，翁和贺利用浙江省总工会副主席职位作为跳板，向省级当局发动进攻，那并不是因为他们在执行机构中拥有数量优势。\footnote{One
  article specifically denied that He and Weng had seized control of the
  ZPTUC leadership. See, HZRB, January 22, 1977.
  有一篇文章明确否认何和翁夺取了浙江省总工会领导权。见HZRB，1977年1月22日。}凭借更高明的策略、个人影响力以及动员支持者的能力，他们可能在行动上胜过了对手。但他们控制着另一个在1974年政治斗争中发挥更大作用的组织。这就是由贺贤春领导的杭州工人代表大会。贺及其追随者成功抵制了一切解散该组织的努力，该组织继续作为一个官方承认的机构运作，并将被证明是省级领导层的一根刺。\footnote{For
  further comparative detail of the provincial-level trade union
  congresses, see, CNS, 462, 464, 471; Li Ming-hua, ``An Analysis of the
  Reconstruction'', pp.~35-44; and analyses of the trade unions and
  factionalism among the working class, CNA, 940 and 941 (November 16
  and 23,1973).
  关于省级工会代表大会的进一步比较细节，见CNS，462、464、471；Li
  Ming-hua，``An Analysis of the
  Reconstruction''，第35-44页；以及关于工会和工人阶级派性分析，CNA，940和941（1973年11月16日、23日）。}

\textbf{Tables Six-Eight: Executives of Mass Organizations Reestablished
in 1973}

\textbf{表六至表八：1973年重建的群众组织执行机构}

\textbf{Table Six. The Executive of the Zhejiang Provincial Trade Union
Council}

\textbf{表六：浙江省总工会执行机构}

{\def\LTcaptype{none} % do not increment counter
\begin{longtable}[]{@{}lll@{}}
\toprule\noalign{}
Position & Name & Background \\
\midrule\noalign{}
\endhead
\bottomrule\noalign{}
\endlastfoot
Chairwoman & Jiang Baodi & member of United Headquarters in the CR \\
Vice-chairmen & Xu Wanzhen & veteran trade union official \\
& Chen Yousheng & model worker of the GLF \\
& Guo Zhisong & member of the ``mountain base'' faction \\
& Fang Jianwen & leader of the ``mountain base'' faction \\
& Weng Senhe & leader of the ``mountain top'' faction \\
& He Xianchun & worker-leader of the ``mountain top'' faction \\
& Zhao Jingfang & trade union cadre \\
\end{longtable}
}

{\def\LTcaptype{none} % do not increment counter
\begin{longtable}[]{@{}lll@{}}
\toprule\noalign{}
职位 & 姓名 & 背景 \\
\midrule\noalign{}
\endhead
\bottomrule\noalign{}
\endlastfoot
主席 & 蒋宝娣 & 文革中``联合总部''成员 \\
副主席 & 徐婉珍 & 老工会干部 \\
& 陈友生 & 大跃进时期劳动模范 \\
& 郭志松 & ``山下派''成员 \\
& 方剑文 & ``山下派''领导人 \\
& 翁森鹤 & ``山上派''领导人 \\
& 贺贤春 & ``山上派''工人领袖 \\
& 赵景芳 & 工会干部 \\
\end{longtable}
}

In addition to being one of the first provincial-level administrations
to hold a trade union congress, Shanghai also led the way in the
re-establishment of the Communist Youth League. After the convocation of
its municipal congress in February 1973 a People's Daily editorial of 22
February set out the tasks of the organization.\footnote{See CNS, 457,
  (March 1,1973), pp.~3-5. 见CNS，457（1973年3月1日），第3-5页。} A
balance between Cultural Revolution rhetoric and a more traditional view
of the League's tasks seemed evident in this editorial. Yet on balance
the editorial emphasized more the concerns of the central radical
propagandists. The CYL congresses of Hangzhou and Zhejiang convened in
April and May 1973. They also became battlegrounds for contention
between Red Guard leaders of the former United Headquarters and Red
Storm mass organizations. A report of February 1973 specified the
categories from which delegates would be selected for the forthcoming
congresses.\footnote{ZJRB, February 9,1973, p.~1.
  ZJRB，1973年2月9日，第1页。} These included a certain percentage of
female representatives, little red soldiers, national minorities (a
category of little relevance in Zhejiang) and qualified sons and
daughters of ' dubious elements, (youth of less reliable class
background) who could be re-educated. The report added that members of
leading groups should be carefully screened and selected after repeated
consultation. They should have made significant contributions since the
Cultural Revolution in the economic, cultural or military fields and
most would remain in their units after election to positions of
responsibility. A certain percentage of the new leadership should have
had previous experience in CYL work, stated the report, but its average
age would be low with a fixed number of females gaining leadership
posts. From 10 to 14 March the CCP ZPC held a preparatory meeting to
make concrete arrangements for the congresses.\footnote{ZJRB, March 19,
  1973, pp.~1, 4. ZJRB，1973年3月19日，第1、4页。} Tie Ying addressed
the sixty people present and stressed the bond between the League and
the CCP. He almost certainly raised the issue of unity, as an editorial
of Zhejiang Daily, published on the same day, stated unambiguously:

除了上海是最早召开工会代表大会的省级行政区之一外，它还在重建共产主义青年团方面走在前列。1973年2月上海市代表大会召开后，2月22日《人民日报》一篇社论阐明了该组织的任务。\footnote{See
  CNS, 457, (March 1,1973), pp.~3-5.
  见CNS，457（1973年3月1日），第3-5页。}这篇社论中明显存在文革话语和对共青团任务更传统看法之间的平衡。不过总体而言，社论更强调中央激进宣传家的关切。杭州和浙江共青团代表大会于1973年4月和5月召开。它们也成为前``联合总部''和``红暴''群众组织中红卫兵领导人争斗的战场。1973年2月一份报告规定了即将召开大会的代表遴选类别。\footnote{ZJRB,
  February 9,1973, p.~1. ZJRB，1973年2月9日，第1页。}其中包括一定比例的女性代表、小红兵、少数民族（在浙江相关性很小的一类），以及可以再教育的``可疑分子''中合格的子女（阶级出身不那么可靠的青年）。报告补充说，领导小组成员应经过反复协商后仔细审查和遴选。他们应在文革以来的经济、文化或军事领域作出重要贡献，并且多数人在当选负责岗位后仍留在本单位。报告说，新领导层中应有一定比例人员有过共青团工作经验，但其平均年龄会较低，并有固定数量女性进入领导岗位。3月10日至14日，中共浙江省委召开筹备会议，为各次代表大会作出具体安排。\footnote{ZJRB,
  March 19, 1973, pp.~1, 4. ZJRB，1973年3月19日，第1、4页。}铁瑛向出席的60人讲话，强调共青团同中共之间的联系。他几乎肯定提出了团结问题，因为同日发表的《浙江日报》社论明确写道：

\begin{quote}
We must take into consideration the situation as a whole and overcome
bourgeois factionalism while upholding the Party spirit of the
proletariat.

我们必须顾全大局，在坚持无产阶级党性的同时克服资产阶级派性。
\end{quote}

The editorial enjoined the delegates, in words taken directly from Mao's
1964 five requirements for successors to the revolution, to strive for
unity whether their fellow-delegates shared their views or not, and even
with those who ``formerly opposed us and have since been proved wrong in
practice''
(不但要团结和自己意见相同的,而且要善于团结那些和自己意见不同的人,还要善于团结那些反对过自己并且已被实践证明是犯了错误的人).\footnote{ZJRB,
  March 19,1973, p.~1. For Mao's five requirements for revolutionary
  successors see ``On Khrushchov's Phoney Communism and Its Historical
  Lessons for the World'', July 14,1964, in The Polemic on the General
  Line of the International Communist Movement, (Beijing: Foreign
  Languages Press, 1965), p.~478.
  ZJRB，1973年3月19日，第1页。关于毛对革命接班人的五项要求，见``On
  Khrushchov's Phoney Communism and Its Historical Lessons for the
  World''，1964年7月14日，载The Polemic on the General Line of the
  International Communist Movement，（北京：Foreign Languages
  Press，1965），第478页。} Tie Ying later used similar words in his
closing address to the Trade Union Congress. The editorial and Tie's
remarks were strong indictment of United Headquarters. By July, when
similar words were expressed at the Trade Union Congress, the rhetoric
had moderated, an indication that the revived strength of ``leftist''
forces at the center was making an impact in Zhejiang. The 5th
provincial congress of the CYL opened on 16 May 1973 and closed one week
later.\footnote{District (地区) congresses preceded it. Hangzhou's
  municipal congress opened on April 24. ZPS, May 4,1973,
  SWB/FE/4293/BII/11-12. Jinhua and Lishui districts opened at this
  time. Ibid. Zhoushan, March 26-30,1973, Ningbo, March 29-30, 1973.
  ZPS, April 9,1973, SWB/FE/4295/BII/4. Jiaxing, Taizhou, Wenzhou and
  Shaoxing district congresses were held between April 20 and May 9.
  ZPS, May 13,1973, SWB/FE/4308/BII/19.
  在此之前，各地区（地区）代表大会已经召开。杭州市代表大会于4月24日开幕。ZPS，1973年5月4日，SWB/FE/4293/BII/11-12。金华和丽水地区也在此时开幕。同上。舟山为1973年3月26-30日，宁波为1973年3月29-30日。ZPS，1973年4月9日，SWB/FE/4295/BII/4。嘉兴、台州、温州和绍兴地区代表大会于4月20日至5月9日之间举行。ZPS，1973年5月13日，SWB/FE/4308/BII/19。}
One thousand four hundred and thirty-two delegates, 44\% of whom were
women, as well as Red Guard observers attended the opening ceremony. Tan
Qilong delivered the opening address in which he asked his listeners to
adhere to the general orientation of the 1972 campaign against
revisionism. Wu Peisheng, who was elected the first of four secretaries,
made a report entitled ``Strive to train the younger generation into
successors to the revolutionary cause of the proletariat''. The Congress
elected a committee of 103 members and 15 alternate members whose
average age was 23.2 years. It in turn chose a standing committee of 15.
Its executive of four secretaries and four deputy-secretaries included
two women.\footnote{ZJRB, May 17,1973, p.~1; ZJRB, May 24,1973, p.~1.
  For comparative figures of the sex, age, ethnic and job classification
  of the various provincial CYL delegates and committee members, see
  Chou Wei-ling, ``The Status of the Provincial CYL Committees'', I\&S,
  10:1,10:2 (October and November 1973), pp.~52-65, 77-87; see also CNS,
  466 (May 10,1973).
  ZJRB，1973年5月17日，第1页；ZJRB，1973年5月24日，第1页。关于各省共青团代表和委员的性别、年龄、民族和职业分类比较数字，见Chou
  Wei-ling，``The Status of the Provincial CYL
  Committees''，I\&S，10:1、10:2（1973年10月和11月），第52-65、77-87页；另见CNS，466（1973年5月10日）。}
Little is known of the backgrounds of the members of this executive.
Secretary Wu Peisheng was Vice-chairman of the Jiangshan county cement
works revolutionary committee in southern Zhejiang and a member of the
ZPC. His factory had received a certain amount of publicity since 1970,
particularly for its study of Marxist philosophy (see chapter
five).\footnote{See also ZPS, September 15,1970, SWB/FE/3490/BII/12-13;
  ZPS, November 1,1970, SWB/FE/3526/BII/8-10; ZPS, September 17, 1971,
  FBIS/CHI, 190 (1971), C6-7.
  另见ZPS，1970年9月15日，SWB/FE/3490/BII/12-13；ZPS，1970年11月1日，SWB/FE/3526/BII/8-10；ZPS，1971年9月17日，FBIS/CHI，190（1971），C6-7。}
Another secretary, Zhang Jinguang, was a ``veteran'' CYL cadre, having
been a member of the 9th national CYL CC before the Cultural
Revolution.\footnote{Chou Wei-ling, ``The Status of the Provincial CYL
  Committees'', Part II, p.~84. Chou Wei-ling，``The Status of the
  Provincial CYL Committees''，第二部分，第84页。} Two
Deputy-secretaries, Teng Zhu and Liu Ying, were former leaders of Red
Storm from Zhejiang University. The only executive member known to have
been an activist in United Headquarters was Deputy-secretary Wang Zexin.
He was a graduate of the politics department of Hangzhou University and
had been a delegate to the CCP's 9th national congress. The party
leadership in some provinces appointed a prominent party or government
official to head the local CYL body, presumably to ensure tight control
over the organization. For example, Xie Jingyi, daughter of the late Xie
Fuzhi and a confidential secretary to the CC General Office, became CYL
secretary in Beijing.\footnote{Gao Gao and Yan Jiaqi, ``Wenhua
  dageming'', p.~483. Gao Gao and Yan Jiaqi，``Wenhua
  dageming''，第483页。} Xie was also a standing committee member of the
CCP Beijing Municipal Committee and in August 1973 was elected a member
of the CCP 10th CC. Other provincial CYL secretaries, including
Zhejiang's Wu Peisheng, were members of provincial-level party
committees or were elected to the CCP 10th CC in August 1973.\footnote{Chou
  Wei-ling, ``The Status of the Provincial CYL Committees'', Part II,
  pp.~85-6. Chou Wei-ling，``The Status of the Provincial CYL
  Committees''，第二部分，第85-6页。} The Zhejiang Provincial Red Guards
Congress was not dissolved. In May 1974 it published an article in
Zhejiang Daily commemorating the 55th anniversary of the May 4th
movement.\footnote{See ZPS, May 3, 1974, FBIS/CHI, 88, G 1.
  见ZPS，1974年5月3日，FBIS/CHI，88，G1。}

社论使用直接取自毛1964年对革命接班人提出的五项要求的话，告诫代表们努力团结，不论其他代表是否同自己意见相同，甚至要团结那些``反对过自己并且已被实践证明是犯了错误的人''（不但要团结和自己意见相同的,而且要善于团结那些和自己意见不同的人,还要善于团结那些反对过自己并且已被实践证明是犯了错误的人）。\footnote{ZJRB,
  March 19,1973, p.~1. For Mao's five requirements for revolutionary
  successors see ``On Khrushchov's Phoney Communism and Its Historical
  Lessons for the World'', July 14,1964, in The Polemic on the General
  Line of the International Communist Movement, (Beijing: Foreign
  Languages Press, 1965), p.~478.
  ZJRB，1973年3月19日，第1页。关于毛对革命接班人的五项要求，见``On
  Khrushchov's Phoney Communism and Its Historical Lessons for the
  World''，1964年7月14日，载The Polemic on the General Line of the
  International Communist Movement，（北京：Foreign Languages
  Press，1965），第478页。}铁瑛后来在工会代表大会闭幕讲话中也使用了类似措辞。社论和铁的讲话是对``联合总部''的强烈指控。到7月，工会代表大会上表达类似话语时，措辞已较为缓和，这表明中央``左派''力量的复兴正在浙江产生影响。共青团浙江省第五次代表大会于1973年5月16日开幕，一周后闭幕。\footnote{District
  (地区) congresses preceded it. Hangzhou's municipal congress opened on
  April 24. ZPS, May 4,1973, SWB/FE/4293/BII/11-12. Jinhua and Lishui
  districts opened at this time. Ibid. Zhoushan, March 26-30,1973,
  Ningbo, March 29-30, 1973. ZPS, April 9,1973, SWB/FE/4295/BII/4.
  Jiaxing, Taizhou, Wenzhou and Shaoxing district congresses were held
  between April 20 and May 9. ZPS, May 13,1973, SWB/FE/4308/BII/19.
  在此之前，各地区（地区）代表大会已经召开。杭州市代表大会于4月24日开幕。ZPS，1973年5月4日，SWB/FE/4293/BII/11-12。金华和丽水地区也在此时开幕。同上。舟山为1973年3月26-30日，宁波为1973年3月29-30日。ZPS，1973年4月9日，SWB/FE/4295/BII/4。嘉兴、台州、温州和绍兴地区代表大会于4月20日至5月9日之间举行。ZPS，1973年5月13日，SWB/FE/4308/BII/19。}1432名代表出席开幕式，其中44\%为女性，另有红卫兵观察员。谭启龙致开幕词，要求听众坚持1972年批修运动的总方向。吴培生当选四名书记中的第一书记，并作题为``努力把青年一代培养成为无产阶级革命事业接班人''的报告。大会选出103名委员和15名候补委员，平均年龄23.2岁。该委员会又选出15人常委会。其由四名书记和四名副书记组成的执行机构中包括两名女性。\footnote{ZJRB,
  May 17,1973, p.~1; ZJRB, May 24,1973, p.~1. For comparative figures of
  the sex, age, ethnic and job classification of the various provincial
  CYL delegates and committee members, see Chou Wei-ling, ``The Status
  of the Provincial CYL Committees'', I\&S, 10:1,10:2 (October and
  November 1973), pp.~52-65, 77-87; see also CNS, 466 (May 10,1973).
  ZJRB，1973年5月17日，第1页；ZJRB，1973年5月24日，第1页。关于各省共青团代表和委员的性别、年龄、民族和职业分类比较数字，见Chou
  Wei-ling，``The Status of the Provincial CYL
  Committees''，I\&S，10:1、10:2（1973年10月和11月），第52-65、77-87页；另见CNS，466（1973年5月10日）。}关于该执行机构成员的背景所知不多。书记吴培生是浙南江山县水泥厂革命委员会副主任、浙江省委委员。自1970年以来，他所在工厂得到一定宣传，尤其是因学习马克思主义哲学而受到宣传（见第五章）。\footnote{See
  also ZPS, September 15,1970, SWB/FE/3490/BII/12-13; ZPS, November
  1,1970, SWB/FE/3526/BII/8-10; ZPS, September 17, 1971, FBIS/CHI, 190
  (1971), C6-7.
  另见ZPS，1970年9月15日，SWB/FE/3490/BII/12-13；ZPS，1970年11月1日，SWB/FE/3526/BII/8-10；ZPS，1971年9月17日，FBIS/CHI，190（1971），C6-7。}另一名书记张金光是``老''共青团干部，文革前曾任第九届全国共青团中央委员。\footnote{Chou
  Wei-ling, ``The Status of the Provincial CYL Committees'', Part II,
  p.~84. Chou Wei-ling，``The Status of the Provincial CYL
  Committees''，第二部分，第84页。}两名副书记滕竹和刘英，是浙江大学``红暴''前领导人。已知执行机构成员中唯一曾是``联合总部''积极分子的是副书记王泽新。他毕业于杭州大学政治系，曾是中共九大代表。一些省份的党领导层任命著名党政官员负责地方共青团机构，大概是为了确保对该组织的严密控制。例如，已故谢富治之女、中共中央办公厅机要秘书谢静宜成为北京共青团书记。\footnote{Gao
  Gao and Yan Jiaqi, ``Wenhua dageming'', p.~483. Gao Gao and Yan
  Jiaqi，``Wenhua dageming''，第483页。}谢还是中共北京市委常委，并于1973年8月当选中共十大中央委员。其他省级共青团书记，包括浙江的吴培生，也是省级党委委员，或在1973年8月当选中共十大中央委员。\footnote{Chou
  Wei-ling, ``The Status of the Provincial CYL Committees'', Part II,
  pp.~85-6. Chou Wei-ling，``The Status of the Provincial CYL
  Committees''，第二部分，第85-6页。}浙江省红卫兵代表大会并未解散。1974年5月，它在《浙江日报》发表文章纪念五四运动55周年。\footnote{See
  ZPS, May 3, 1974, FBIS/CHI, 88, G 1.
  见ZPS，1974年5月3日，FBIS/CHI，88，G1。}

\textbf{Table Seven. The Executive of the Zhejiang Committee of the
Communist Youth League}

\textbf{表七：共青团浙江省委执行机构}

{\def\LTcaptype{none} % do not increment counter
\begin{longtable}[]{@{}lll@{}}
\toprule\noalign{}
Position & Name & Background \\
\midrule\noalign{}
\endhead
\bottomrule\noalign{}
\endlastfoot
Secretaries & Wu Peisheng & \\
& Zhang Jinguang & \\
& Lin Chaodi (f) & \\
& Yang Jufang (f) & secretary, Hangzhou CYL Committee \\
Deputy-secretaries & Li Huiliang & \\
& Teng Zhu & technician at Hangzhou Machinery Plant \\
& Wang Zexin & \\
& Liu Ying & \\
\end{longtable}
}

{\def\LTcaptype{none} % do not increment counter
\begin{longtable}[]{@{}lll@{}}
\toprule\noalign{}
职位 & 姓名 & 背景 \\
\midrule\noalign{}
\endhead
\bottomrule\noalign{}
\endlastfoot
书记 & 吴培生 & \\
& 张金光 & \\
& 林朝娣（女） & \\
& 杨菊芳（女） & 杭州共青团市委书记 \\
副书记 & 李惠良 & \\
& 滕竹 & 杭州机械厂技术员 \\
& 王泽新 & \\
& 刘英 & \\
\end{longtable}
}

The third mass organization to be reconstituted was the Women's
Federation. After district branches held congresses, in many cases the
first they had ever convened,\footnote{ZPS, July 14, 1973,
  SWB/FE/4351/BII/13-14 (Zhoushan, Shaoxing, Ningbo, Jinhua districts);
  ZPS, July 24,1973, SWB/FE/4363/BII/13-15 (Hangzhou, Jiaxing, Lishui,
  Taizhou, Wenzhou).
  ZPS，1973年7月14日，SWB/FE/4351/BII/13-14（舟山、绍兴、宁波、金华地区）；ZPS，1973年7月24日，SWB/FE/4363/BII/13-15（杭州、嘉兴、丽水、台州、温州）。}
the provincial meeting took place from 12 to 17 August 1973, on the eve
of the CCP 10th Congress. Over 1,000 delegates attended and according to
the reports, most were of peasant and worker origin. After listening to
speeches and reports by Tan Qilong, Chai Qikun, Tie Ying and Chen Bing
(all men!) the women elected a committee of 120 which selected
twenty-one women to form the standing committee. The executive consisted
of a Chairwoman and seven Vice-chairwomen.\footnote{HZRB, August
  13,18,1973: RMRB, August 19,1973, p.~1.
  HZRB，1973年8月13日、18日；RMRB，1973年8月19日，第1页。}
Vice-chairwoman of the ZPRC and a former member of the CCP ZPC standing
committee, Hua Yinfeng was elected Chairwoman. Although she had
previously belonged to United Headquarters, like Jiang Baodi she had
probably not been an influential or active member of that organization
and certainly presented no threat to the provincial party leadership. Of
the known backgrounds of the seven Vice-chairwomen one, Shao Suzhen, had
been a member of United Headquarters and another, Lu Su was a veteran
cadre in the women's movement.\footnote{In May 1968 Lu had been attacked
  by name in the pages of HZRB. See HZRB, November 8, 1978.
  1968年5月，陆曾在HZRB版面上被点名攻击。见HZRB，1978年11月8日。} Liu
Tianxiang was an activist in a democratic party within the United Front.

第三个被重建的群众组织是妇联。地区分会召开代表大会后，省级会议于1973年8月12日至17日、中共十大前夕举行；在许多情况下，这些地区会议还是它们有史以来第一次召开的代表大会。\footnote{ZPS,
  July 14, 1973, SWB/FE/4351/BII/13-14 (Zhoushan, Shaoxing, Ningbo,
  Jinhua districts); ZPS, July 24,1973, SWB/FE/4363/BII/13-15 (Hangzhou,
  Jiaxing, Lishui, Taizhou, Wenzhou).
  ZPS，1973年7月14日，SWB/FE/4351/BII/13-14（舟山、绍兴、宁波、金华地区）；ZPS，1973年7月24日，SWB/FE/4363/BII/13-15（杭州、嘉兴、丽水、台州、温州）。}一千多名代表出席，据报道，多数出身于农民和工人。听取谭启龙、柴启琨、铁瑛和陈冰（全是男性！）的讲话和报告后，妇女们选出一个120人的委员会，该委员会又选出21名妇女组成常委会。执行机构由一名主席和七名副主席组成。\footnote{HZRB,
  August 13,18,1973: RMRB, August 19,1973, p.~1.
  HZRB，1973年8月13日、18日；RMRB，1973年8月19日，第1页。}浙江省革命委员会副主任、前中共浙江省委常委华银凤当选主席。虽然她此前属于``联合总部''，但像蒋宝娣一样，她大概不是该组织中有影响或活跃的成员，当然也不对省委领导层构成威胁。在七名副主席已知背景中，一人邵素珍曾是``联合总部''成员，另一人陆苏是妇女运动中的老干部。\footnote{In
  May 1968 Lu had been attacked by name in the pages of HZRB. See HZRB,
  November 8, 1978.
  1968年5月，陆曾在HZRB版面上被点名攻击。见HZRB，1978年11月8日。}刘天香是统一战线内一个民主党派的积极分子。

\textbf{Table Eight. The Executive of the Zhejiang Provincial Women's
Federation}

\textbf{表八：浙江省妇女联合会执行机构}

{\def\LTcaptype{none} % do not increment counter
\begin{longtable}[]{@{}ll@{}}
\toprule\noalign{}
Position & Name \\
\midrule\noalign{}
\endhead
\bottomrule\noalign{}
\endlastfoot
Chairwoman & Hua Yinfeng \\
Vice-chairwomen & Lu Su \\
& Tang Youqing \\
& Cheng Xiu \\
& Li Tianxiang \\
& Shao Suzhen \\
& Wu Guomei \\
& Su Lianzhu \\
\end{longtable}
}

{\def\LTcaptype{none} % do not increment counter
\begin{longtable}[]{@{}ll@{}}
\toprule\noalign{}
职位 & 姓名 \\
\midrule\noalign{}
\endhead
\bottomrule\noalign{}
\endlastfoot
主席 & 华银凤 \\
副主席 & 陆苏 \\
& 唐有卿 \\
& 程秀 \\
& 李天香 \\
& 邵素珍 \\
& 吴国美 \\
& 苏联珠 \\
\end{longtable}
}

There is no evidence that the Zhejiang Poor and Lower-middle Peasants'
Association convened a congress in this period. Hunan province held a
provincial congress of the organization between November and December
1973.\footnote{CNS, 496 (December 6,1973), pp.~1-3.
  CNS，496（1973年12月6日），第1-3页。} Party secretary of Hangzhou's
Red May brigade of the East Wind People's Commune, Mo Xianyao (莫显耀),
who was a member of both the CCP 9th and 10th CC, was later identified
as Chairman of the association.\footnote{Mo held this position as late
  as October 1976. See HZRB, October 24,1976. An enlarged preparatory
  congress had been held in October 1968. See Zhang Yongshengde zuizheng
  cailiao, p.~27.
  直到1976年10月，莫仍担任这一职务。见HZRB，1976年10月24日。1968年10月曾召开扩大筹备代表大会。见Zhang
  Yongshengde zuizheng cailiao，第27页。}

没有证据显示浙江省贫下中农协会在这一时期召开过代表大会。湖南省在1973年11月至12月间召开了该组织的省级代表大会。\footnote{CNS,
  496 (December 6,1973), pp.~1-3. CNS，496（1973年12月6日），第1-3页。}杭州东风人民公社红五月大队党委书记莫显耀（莫显耀）同时是中共九大和十大中央委员，后来被确认为该协会主席。\footnote{Mo
  held this position as late as October 1976. See HZRB, October 24,1976.
  An enlarged preparatory congress had been held in October 1968. See
  Zhang Yongshengde zuizheng cailiao, p.~27.
  直到1976年10月，莫仍担任这一职务。见HZRB，1976年10月24日。1968年10月曾召开扩大筹备代表大会。见Zhang
  Yongshengde zuizheng cailiao，第27页。}

\section{Two Lines in Contention}\label{two-lines-in-contention}

\section{两条路线的争斗}\label{ux4e24ux6761ux8defux7ebfux7684ux4e89ux6597}

Reports from Zhejiang about the mass organization congresses and other
articles published prior to the 10th Congress devoted considerable space
to defending the achievements of the Cultural Revolution. Reference was
made to the induction into the CCP of young people and under-represented
groups such as women and their promotion to positions of authority. For
example one report declared that ``in recent years'' 470,000 young
people had joined the CYL and 28,000 CYL members had been admitted into
the CCP. Tan Qilong mentioned in his speech to the Women's Federation
Congress that 100,000 women had joined the CYL since the Cultural
Revolution, 10,000 had entered the CCP (these round figures were later
revised to 140,000 and 12,000 respectively) and 5,000 female cadres had
been promoted to positions of leadership at and above the commune level.
Of the 330,000 young people who had been sent from cities to either
border regions, to the countryside of Zhejiang, or to the outskirts of
Hangzhou, 19,580 from Hangzhou had joined the CYL, 770 the CCP and 8,000
were working in leadership positions.\footnote{ZPS, May 16, 1973,
  SWB/FE/4313/BII/20; ZPS, August 12, 1973, SWB/FE/4375/BII/17-18; ZPS,
  August 12, 1973, SWB/FE/4381/BII/14; ZPS, May 4,1973,
  SWB/FE/4280/BII/8-13; ZPS, August 12, 1973, SWB/FE/4381/BII/7-8.
  ZPS，1973年5月16日，SWB/FE/4313/BII/20；ZPS，1973年8月12日，SWB/FE/4375/BII/17-18；ZPS，1973年8月12日，SWB/FE/4381/BII/14；ZPS，1973年5月4日，SWB/FE/4280/BII/8-13；ZPS，1973年8月12日，SWB/FE/4381/BII/7-8。}
Other media reports praised policy achievements in the fields of rural
health and education. New procedures for admission to tertiary
institutions and problems associated with the rustication of educated
youth were discussed. The general thrust of the policies was affirmed.
The perennial issue of cadre participation in manual labour also
received attention. However, other articles in the provincial press
extolled the achievements of the 1972-73 campaign against revisionism
(Lin Biao) and detailed the successful implementation of pragmatic
economic policies. A striking example was the achievements of the
Hangzhou Iron and Steel Works. In 1968 a work team from Chen Liyun's
Unit 7350 had entered the factory to take over the
management.\footnote{ZPS, September 24,1968, SWB/FE/2885/B/20; visit to
  Hangzhou Iron and Steel Mill, May 27,1979.
  ZPS，1968年9月24日，SWB/FE/2885/B/20；1979年5月27日访问杭州钢铁厂。}
The air force personnel found that they lacked the expertise to run it
properly and the mill may well have been one of the factories described
in colorful terms by Tie Ying, after his transfer to Hangzhou in 1972,
as a place where production had ceased, machinery had rusted up and the
surroundings were overgrown with weeds.\footnote{Tie Ying, HZRB, January
  9,1985. 铁瑛，HZRB，1985年1月9日。} By the autumn of 1972 production
at the mill had recovered sufficiently for it to send exhibits to the
Guangzhou Trade Fair.\footnote{Ibid.; visit, May 27, 1979.
  同上；1979年5月27日访问。} A Xinhua report of June 1973 attributed the
production increases which had occurred in the first half of the year to
sustained efforts by the leadership to overcome splits and factionalism
among the workers.\footnote{RMRB, June 13,1973, p.~1; CNA, 941 (November
  23,1973), pp.~4-5.
  RMRB，1973年6月13日，第1页；CNA，941（1973年11月23日），第4-5页。} The
article attacked ``swindlers like Liu Shaoqi'' (Lin Biao) for trying to
divide the workers. It also made scathing references to ``class
enemies'' who schemed to ``redefine classes'' (重新划分阶级) and who
stated privately that ``if they unite we've had it''
(他们联合起来,我们就要倒霉了). This comment revealed the tenuous nature
of the solidarity among the workers in the mill. It was also an
indication that for some activists of the Cultural Revolution one's
political attitudes and behavior overrode the stigma associated with
class labels attached by the accident of birth.\footnote{After his
  downfall Zhang Chunqiao was accused of formulating a theory to
  redefine the nature of the Chinese working class. See CCP CC, zhongfa
  (1977), No.~37, in I\&S, 15:1 (January 1979), pp.~99-100.
  张春桥倒台后，被指控提出一种重新界定中国工人阶级性质的理论。见中共中央，中发（1977）第37号，载I\&S，15:1（1979年1月），第99-100页。}

浙江关于群众组织代表大会的报道，以及十大前发表的其他文章，花费相当篇幅为文革成就辩护。这些文章提到青年和女性等代表性不足群体被吸收入中共，并被提拔到权力职位。例如，一篇报道宣称，``近年来''有47万青年加入共青团，2.8万名共青团员被吸收入中共。谭启龙在妇联代表大会讲话中提到，文革以来有10万名妇女加入共青团，1万人加入中共（这些整数后来分别修正为14万和1.2万），5000名女干部被提拔到公社及以上领导岗位。在33万名从城市被送往边疆地区、浙江农村或杭州郊区的青年中，杭州有19580人加入共青团，770人加入中共，8000人在领导岗位工作。\footnote{ZPS,
  May 16, 1973, SWB/FE/4313/BII/20; ZPS, August 12, 1973,
  SWB/FE/4375/BII/17-18; ZPS, August 12, 1973, SWB/FE/4381/BII/14; ZPS,
  May 4,1973, SWB/FE/4280/BII/8-13; ZPS, August 12, 1973,
  SWB/FE/4381/BII/7-8.
  ZPS，1973年5月16日，SWB/FE/4313/BII/20；ZPS，1973年8月12日，SWB/FE/4375/BII/17-18；ZPS，1973年8月12日，SWB/FE/4381/BII/14；ZPS，1973年5月4日，SWB/FE/4280/BII/8-13；ZPS，1973年8月12日，SWB/FE/4381/BII/7-8。}其他媒体报道赞扬农村卫生和教育领域的政策成就。文章讨论了高等院校招生新程序以及知识青年下乡相关问题。政策的总体方向得到肯定。干部参加体力劳动这一长期问题也受到关注。然而，省内报刊的其他文章则赞扬1972-73年批判修正主义（林彪）运动的成就，并详细说明务实经济政策的成功实施。一个显著例子是杭州钢铁厂的成就。1968年，陈励耘7350部队的一个工作组进入该厂接管管理。\footnote{ZPS,
  September 24,1968, SWB/FE/2885/B/20; visit to Hangzhou Iron and Steel
  Mill, May 27,1979.
  ZPS，1968年9月24日，SWB/FE/2885/B/20；1979年5月27日访问杭州钢铁厂。}空军人员发现自己缺乏妥善经营工厂的专业能力，而这家钢厂很可能就是铁瑛1972年调到杭州后用生动语言描述的那些工厂之一：生产停止、机器生锈、四周杂草丛生。\footnote{Tie
  Ying, HZRB, January 9,1985. 铁瑛，HZRB，1985年1月9日。}到1972年秋，该厂生产已恢复到足以向广州交易会送展品的程度。\footnote{Ibid.;
  visit, May 27, 1979. 同上；1979年5月27日访问。}新华社1973年6月一篇报道把该年上半年出现的增产归因于领导层持续努力克服工人中的分裂和派性。\footnote{RMRB,
  June 13,1973, p.~1; CNA, 941 (November 23,1973), pp.~4-5.
  RMRB，1973年6月13日，第1页；CNA，941（1973年11月23日），第4-5页。}文章攻击``刘少奇一类骗子''（林彪）试图分裂工人。它还尖锐提到``阶级敌人''阴谋``重新划分阶级''（重新划分阶级），并私下说``他们联合起来,我们就要倒霉了''（他们联合起来,我们就要倒霉了）。这一评论揭示了钢厂工人团结的脆弱性质。它也表明，对一些文革积极分子而言，一个人的政治态度和行为压倒了由出生偶然性附加的阶级标签污名。\footnote{After
  his downfall Zhang Chunqiao was accused of formulating a theory to
  redefine the nature of the Chinese working class. See CCP CC, zhongfa
  (1977), No.~37, in I\&S, 15:1 (January 1979), pp.~99-100.
  张春桥倒台后，被指控提出一种重新界定中国工人阶级性质的理论。见中共中央，中发（1977）第37号，载I\&S，15:1（1979年1月），第99-100页。}

\textbf{Table Nine. Standing committee of the CCP ZPC, August 1973}

\textbf{表九：中共浙江省委常委会，1973年8月}

{\def\LTcaptype{none} % do not increment counter
\begin{longtable}[]{@{}llll@{}}
\toprule\noalign{}
Position & Name & Notes & Date of appointment \\
\midrule\noalign{}
\endhead
\bottomrule\noalign{}
\endlastfoot
First Secretary & Tan Qilong & 1, 3, 5 & April 1972 \\
Secretary & Tie Ying & 1, 2 & April 1972 \\
Deputy-secretaries & Lai Keke & 3, 4 & April 1963 \\
& Xie Zhenghao & 2, 4 & January 1971 \\
& Chai Qikun & 2, 4 & January 1971 \\
& Chen Weida & 1, 3, 5 & April 1955 and again in November 1972 \\
Members & Wang Zida & 3, 4 & January 1971 \\
& Shen Ce & 3, 4 & April 1965 \\
& Dai Kelin & 2, 4 & January 1971 \\
& Zhu Quanlin & 2, 4 & January 1971 \\
& Meng Zhaoyu & 2, 4 & January 1971 \\
& Jiang Baodi & 4 & January 1971 \\
& Zhou Jianren & 1, 3 & ? \\
& Chen Bing & 1, 3, 5 & April 1965 and again in November 1972 \\
& Liu Ang & 1, 2 & November 1972 \\
& Xia Qi & 1, 2 & November 1972 \\
\end{longtable}
}

{\def\LTcaptype{none} % do not increment counter
\begin{longtable}[]{@{}llll@{}}
\toprule\noalign{}
职位 & 姓名 & 注释 & 任命日期 \\
\midrule\noalign{}
\endhead
\bottomrule\noalign{}
\endlastfoot
第一书记 & 谭启龙 & 1, 3, 5 & 1972年4月 \\
书记 & 铁瑛 & 1, 2 & 1972年4月 \\
副书记 & 赖可可 & 3, 4 & 1963年4月 \\
& 谢正浩 & 2, 4 & 1971年1月 \\
& 柴启琨 & 2, 4 & 1971年1月 \\
& 陈伟达 & 1, 3, 5 & 1955年4月，1972年11月再次任命 \\
委员 & 王子达 & 3, 4 & 1971年1月 \\
& 沈策 & 3, 4 & 1965年4月 \\
& 戴克林 & 2, 4 & 1971年1月 \\
& 朱全林 & 2, 4 & 1971年1月 \\
& 孟昭煜 & 2, 4 & 1971年1月 \\
& 蒋宝娣 & 4 & 1971年1月 \\
& 周建人 & 1, 3 & ? \\
& 陈冰 & 1, 3, 5 & 1965年4月，1972年11月再次任命 \\
& 刘昂 & 1, 2 & 1972年11月 \\
& 夏琦 & 1, 2 & 1972年11月 \\
\end{longtable}
}

Notes: 1. Added since the Lin Biao affair; 2. PLA cadre; 3. Pre-Cultural
Revolution administrative experience in Zhejiang; 4. Cultural Revolution
activist; 5. Rehabilitated cadre; 6. Those dropped since the Zhejiang
5th CCP Provincial Congress in January 1971: Nan Ping, Chen Liyun, Xiong
Yingtang, Hua Yinfeng.

注释：1. 林彪事件后增补；2. 解放军干部；3. 文革前在浙江有行政经验；4.
文革积极分子；5. 平反干部；6.
1971年1月中共浙江省第五次代表大会以来被撤下者：南萍、陈励耘、熊应堂、华银凤。

Sources: \emph{Zhejiang shengqing gaiyao}, p.~26; \emph{Hierarchies of
the PRC}, March 1975, p.~13; \emph{CNS}, 7 June 1973, pp.~5-7.

资料来源：\emph{Zhejiang shengqing gaiyao}，第26页；\emph{Hierarchies of
the PRC}，1975年3月，第13页；\emph{CNS}，1973年6月7日，第5-7页。

One acute observer of Chinese politics has noted, in relation to the
existence side by side of two lines in the media, that ``There seems to
have been an understanding at the center that both sides in the 'debate'
pursued from 1973 to 1976 should have equal access to the
media''.\footnote{Lowell Dittmer, ``Basis of Power in Chinese Politics:
  A Theory and an Analysis of the Fall of the `Gang of Four''', World
  Politics, 31:1 (1978), p.~54, fn. 66. See also Merle Goldman, China's
  Intellectuals: Advise and Dissent (Cambridge, Mass.: Harvard
  University Press, 1981), ch.~6; Fenwick,''The Gang of Four'', p.~164.
  CNS went to great lengths, at times most unconvincingly, to prove the
  existence of one line coming from RMRB and another, more radical, from
  Hongqi. Fenwick argues that the radicals exerted control over Hongqi
  in 1972, but that it was not until 1974 that they took over RMRB.
  ``The Gang of Four'', pp.~177,180, 230-31. Lowell Dittmer，``Basis of
  Power in Chinese Politics: A Theory and an Analysis of the Fall of the
  `Gang of Four'\,''，World
  Politics，31:1（1978），第54页，脚注66。另见Merle Goldman，China's
  Intellectuals: Advise and Dissent（Cambridge, Mass.: Harvard
  University Press，1981），第6章；Fenwick，``The Gang of
  Four''，第164页。CNS费了很大力气，有时相当牵强地证明存在一条来自RMRB的路线和另一条来自Hongqi的更激进路线。Fenwick认为激进派在1972年控制了Hongqi，但直到1974年才接管RMRB。见``The
  Gang of Four''，第177、180、230-31页。} The publication of articles in
Zhejiang and elsewhere emphasizing class struggle and the achievements
of the Cultural Revolution on the one hand, and those devoted to less
dramatic topics such as economic development on the other, reflected the
equilibrium in the national leadership between the forces of order and
those favouring mobilization. Both sides were testing the waters for a
confrontation, but while Mao maintained a neutral stance neither could
move until it was certain that it had obtained his backing. And such
backing was the crucial element in the struggle. Thus, the sparring and
accommodation was a product of particular circumstances and not of the
factional conflict itself.\footnote{In Nathan's model of factionalism,
  neither side strives for the knockout blow. Nathan, ``A Factionalism
  Model for CCP Politics'', pp.~48-9.
  在Nathan的派系主义模型中，双方都不追求一击致命。Nathan，``A
  Factionalism Model for CCP Politics''，第48-9页。} A shift in the
balance of power in Beijing had immediate repercussions in the
provinces. In Zhejiang, Tan and Tie could call on the backing of the
Nanjing Military Region Commander, Xu Shiyou, and Zhou Enlai. Tan and
Tie's principal opponents Zhang Yongsheng, Weng Senhe and He Xianchun
had established ties with Jiang Qing, Zhang Chunqiao and Wang Hongwen.
Within the province Tan and Tie could count on the loyalty of Xia Qi and
other leaders of the ZPMD, rehabilitated cadres, veteran trade union
leaders, older and more skilled workers and the support of former Red
Storm militants. Zhang, Weng, and He's support base was strongest
amongst young factory workers and ex-college students in Hangzhou, who
had formed the bulk of United Headquarters membership in the mid 1960s.
Additionally they could draw on ties with revolutionary leading cadres
such as Lai Keke, Wang Zida and Shen Ce, who had come under suspicion in
the campaign against Chen Boda and Lin Biao, and who had previously
supported them in 1967. The contest between the two sides was about to
be joined and it promised to be a more even and harder fought battle
than that of the previous decade, despite appearances to the contrary. A
major factor which affected the ability of the CCP ZPC to handle the
rebels' challenge was the division within its ranks. The ZPC standing
committee consisted of rehabilitees from the Cultural Revolution who had
been serving in the province in 1966 (Chen Weida and Chen Bing) or who
had been serving elsewhere but had experience in Zhejiang (Tan Qilong).
There were ``revolutionary leading cadres'' who had actively supported
the Maoist forces in the Cultural Revolution (Lai Keke, Shen Ce and Wang
Zida) and who had served in the province for some time. Additionally,
there were PLA representatives loyal to Xu Shiyou (Tie Ying and Xia Qi)
on the one hand and those who had been closely involved in the ``three
supports and two militaries'' work from the ground forces of the
twentieth army (Zhu Quanlin -- 朱全林) and the local garrison (Dai
Kelin), the 5th Air Force (Bai Zongshan) and the navy (Xie Zhenghao and
Chai Qikun). The possibility of such a heterogeneous group possessing
either the capacity or the will to act in unison was slim. A myriad of
ties and obligations bound different individuals to patrons in Beijing,
Shanghai and Nanjing and to former leaders of the two principal mass
organizations. Mass organizations had maintained liaison stations in
major cities until 1969. Prior to the Cultural Revolution the provincial
government of Zhejiang had run offices in both Beijing and Shanghai.
While the office in the national capital was not reopened until December
1981, its counterpart in Shanghai was back in operation by
1973.\footnote{Zhejiang shengqing gaiyao, p.~680. Zhejiang shengqing
  gaiyao，第680页。} Presumably, the re-establishment of the Shanghai
agency enabled the Zhejiang authorities to keep abreast with
developments in the metropolis which was the base of the central radical
leadership. Prior warning of initiatives which the radicals would launch
in the following few years would be invaluable to the leadership in
Zhejiang.

一位敏锐的中国政治观察者在谈到媒体中两条路线并存时指出：``中央似乎有一种默契，即1973年至1976年进行的`辩论'中双方都应平等接触媒体。''\footnote{Lowell
  Dittmer, ``Basis of Power in Chinese Politics: A Theory and an
  Analysis of the Fall of the `Gang of Four''', World Politics, 31:1
  (1978), p.~54, fn. 66. See also Merle Goldman, China's Intellectuals:
  Advise and Dissent (Cambridge, Mass.: Harvard University Press, 1981),
  ch.~6; Fenwick,''The Gang of Four'', p.~164. CNS went to great
  lengths, at times most unconvincingly, to prove the existence of one
  line coming from RMRB and another, more radical, from Hongqi. Fenwick
  argues that the radicals exerted control over Hongqi in 1972, but that
  it was not until 1974 that they took over RMRB. ``The Gang of Four'',
  pp.~177,180, 230-31. Lowell Dittmer，``Basis of Power in Chinese
  Politics: A Theory and an Analysis of the Fall of the `Gang of
  Four'\,''，World Politics，31:1（1978），第54页，脚注66。另见Merle
  Goldman，China's Intellectuals: Advise and Dissent（Cambridge, Mass.:
  Harvard University Press，1981），第6章；Fenwick，``The Gang of
  Four''，第164页。CNS费了很大力气，有时相当牵强地证明存在一条来自RMRB的路线和另一条来自Hongqi的更激进路线。Fenwick认为激进派在1972年控制了Hongqi，但直到1974年才接管RMRB。见``The
  Gang of Four''，第177、180、230-31页。}一方面，浙江和其他地方发表强调阶级斗争和文革成就的文章；另一方面，也发表关注经济发展等较少戏剧性主题的文章。这反映了全国领导层中秩序力量和倾向动员力量之间的均衡。双方都在为对抗试探水温，但只要毛保持中立姿态，任何一方在确信已获得他的支持之前都无法行动。而这种支持正是斗争中的关键因素。因此，试探和妥协是特定环境的产物，而不是派系冲突本身的产物。\footnote{In
  Nathan's model of factionalism, neither side strives for the knockout
  blow. Nathan, ``A Factionalism Model for CCP Politics'', pp.~48-9.
  在Nathan的派系主义模型中，双方都不追求一击致命。Nathan，``A
  Factionalism Model for CCP Politics''，第48-9页。}北京权力平衡的变化会立即在各省产生反响。在浙江，谭和铁可以依靠南京军区司令员许世友和周恩来的支持。谭、铁的主要对手张永生、翁森鹤和贺贤春，则同江青、张春桥和王洪文建立了联系。在省内，谭和铁可以依靠夏琦及浙江省军区其他领导人、平反干部、老工会领导人、年长且技术较高的工人，以及前``红暴''激进分子的支持。张、翁、贺的支持基础在杭州青年工厂工人和原大学生中最强，这些人在1960年代中期构成``联合总部''成员的主体。此外，他们还可以依靠同赖可可、王子达和沈策等革命领导干部的联系；这些干部在批陈伯达和林彪运动中受到怀疑，此前又在1967年支持过他们。两方的较量即将展开，并且尽管表面看似并非如此，它注定会比上一年代的斗争更加势均力敌、更加艰难。影响中共浙江省委应对造反派挑战能力的一个主要因素，是其内部的分裂。省委常委会由文革中获平反者组成，他们或者1966年就在浙江任职（陈伟达和陈冰），或者虽在外地任职但有浙江经验（谭启龙）。还有``革命领导干部''，他们在文革中积极支持毛派力量（赖可可、沈策和王子达），并在省内任职已有一段时间。此外，一方面有忠于许世友的解放军代表（铁瑛和夏琦），另一方面也有来自第二十军地面部队（朱全林--朱全林）、地方卫戍部队（戴克林）、第五空军（白宗善）和海军（谢正浩、柴启琨），且深度参与过``三支两军''工作的人员。这样一个异质性群体拥有共同一致行动的能力或意愿，可能性很小。无数联系和义务把不同个人同北京、上海和南京的庇护者，以及两大主要群众组织的前领导人绑在一起。群众组织直到1969年还在主要城市维持联络站。文革前，浙江省政府在北京和上海均设有办事处。虽然设在首都的办事处直到1981年12月才重新开放，但其上海对应机构到1973年已恢复运作。\footnote{Zhejiang
  shengqing gaiyao, p.~680. Zhejiang shengqing gaiyao，第680页。}可以推测，上海机构的重建使浙江当局能够及时了解这个中央激进派领导层基地的大都市中的动态。对激进派在随后几年将要发起的举措事先获得警示，对浙江领导层将极有价值。

\section{Notes}\label{notes-6}

\section{注释}\label{ux6ce8ux91ca-6}

\bookmarksetup{startatroot}

\chapter{CHAPTER SEVEN - RENEWED RADICALISM, 1973-74 / 第七章 -
重新激进化，1973-74年}\label{chapter-seven---renewed-radicalism-1973-74-ux7b2cux4e03ux7ae0---ux91cdux65b0ux6fc0ux8fdbux53161973-74ux5e74}

\section{Prelude, September 1973-January
1974}\label{prelude-september-1973-january-1974}

\section{前奏，1973年9月-1974年1月}\label{ux524dux594f1973ux5e749ux6708-1974ux5e741ux6708}

The theoretical/ideological justifications for a renewal of the
radicalism of the Cultural Revolution were provided in part by the
reappearance in 1973 of an argument first put forward by Chen Boda in
1968 (see chapter three). One of Mao's major obsessions, and one that he
had incessantly harangued his senior colleagues about since 1962, was
his notion that socialist society consisted of classes and that the
major contradiction in the era of socialism was that between the
proletariat and the bourgeoisie. Chen Boda had argued in 1968, basing
himself on Mao's statement that the communist party inevitably contained
factions, that factions represented classes. The proletarian faction
represented the proletariat and the bourgeoisie had its own
representatives -- the bourgeois faction or revisionists. This argument
laid the groundwork for the position adopted by Zhang Chunqiao and Yao
Wenyuan in articles that they published in 1975 that the bourgeoisie
itself was entrenched within the body of the CCP.

文化大革命激进主义的重新兴起，在理论和意识形态上的依据，一部分来自1973年再次出现的一种论点；这一论点最早由陈伯达在1968年提出（见第三章）。毛泽东长期执着的一个主要观念，也是他自1962年以来不断训诫其高级同僚的内容之一，就是他认为社会主义社会仍然由阶级构成，而社会主义时代的主要矛盾是无产阶级同资产阶级之间的矛盾。陈伯达在1968年依据毛泽东关于共产党内部必然存在派别的说法，提出派别代表阶级。无产阶级派别代表无产阶级，而资产阶级也有自己的代表，即资产阶级派别或修正主义者。这一论点为张春桥和姚文元在1975年发表的文章所采取的立场奠定了基础，即资产阶级本身盘踞在中共机体之内。

In September 1973, Hongqi carried an article by the party branch of a
production brigade in a Beijing county repeating Mao's 1966 statement
concerning the presence of factions within the party.\footnote{Hongqi,
  No.~9 (1973), p.~68, transl. in SCMM, 760 (October 1,1973), p.~38.
  《红旗》1973年第9期，第68页；译文见
  SCMM，第760期（1973年10月1日），第38页。} On this occasion the
quotation was used to condemn factionalism. The radical journal from
Shanghai Study and Criticism (学习与批判), in an article of May 1974
also cited Mao, without attribution, and stated that ``those who
represent different class interests {[}are{]} bound to form different
factions and organize different parties''. It continued:

1973年9月，《红旗》刊登了北京某县一个生产大队党支部的文章，重复毛泽东1966年关于党内存在派别的说法。\footnote{Hongqi,
  No.~9 (1973), p.~68, transl. in SCMM, 760 (October 1,1973), p.~38.
  《红旗》1973年第9期，第68页；译文见
  SCMM，第760期（1973年10月1日），第38页。}
在这一次，这段引语被用来谴责派性。上海的激进刊物《学习与批判》在1974年5月的一篇文章中也未注明出处地引用了毛泽东的话，并说``代表不同阶级利益的人，必然形成不同派别，组织不同党派''。文章接着说：

\begin{quote}
The appearance of different political factions and the struggle between
the two lines in our Party are a reflection in the Party of the class
struggle in society, and chieftains of opportunist lines are agents of
the landlords and bourgeoisie who wormed their way into the party of the
proletariat.\footnote{Xiao Lu, ``驳'君子矜而不争, 群而不党''' (Refute
  the view that a gentleman is respectable and not contentious, sociable
  and not partisan), 学习与批判 (Study and Criticism), 4 (May 20, 1974),
  p.~26, transl. in SCMM, 781 (July 22, 1974), p.~25. See also CNS, 522
  (June 19, 1974).
  肖鲁：《驳'君子矜而不争，群而不党'》（反驳``君子庄重而不争斗，合群而不结党''的观点），《学习与批判》，第4期（1974年5月20日），第26页；译文见
  SCMM，第781期（1974年7月22日），第25页。另见
  CNS，第522期（1974年6月19日）。}
\end{quote}

\begin{quote}
我们党内出现不同政治派别以及两条路线斗争，是社会阶级斗争在党内的反映；机会主义路线的头子，是钻进无产阶级政党内部的地主和资产阶级的代理人。\footnote{Xiao
  Lu, ``驳'君子矜而不争, 群而不党''' (Refute the view that a gentleman
  is respectable and not contentious, sociable and not partisan),
  学习与批判 (Study and Criticism), 4 (May 20, 1974), p.~26, transl. in
  SCMM, 781 (July 22, 1974), p.~25. See also CNS, 522 (June 19, 1974).
  肖鲁：《驳'君子矜而不争，群而不党'》（反驳``君子庄重而不争斗，合群而不结党''的观点），《学习与批判》，第4期（1974年5月20日），第26页；译文见
  SCMM，第781期（1974年7月22日），第25页。另见
  CNS，第522期（1974年6月19日）。}
\end{quote}

The radicals used such arguments to justify their renewed attacks on
their opponents within the party. Recourse to the vague, semi-formulated
ideas of Mao Zedong was their starting point.

激进派利用这类论点，为他们重新攻击党内反对者辩护。他们的出发点，是诉诸毛泽东那些含混而尚未完全成形的思想。

A series of radically-inspired articles appeared in the period around
the convocation of the CCP 10th Congress. The subjects covered included
the plight of rusticated youth, the unfairness of university entrance
examinations and the new ``revisionist education line''. On 16 August
1973, an article which was to have wide repercussions over the next
twelve months appeared in the People's Daily.\footnote{Yang Pu,
  ``反潮流精神'' (The spirit of going against the tide), RMRB, August
  16, 1973, p.~3. An article by the Zhejiang CYL Committee supported the
  spirit of ``going against the tide''. ZPS, September 8, 1973, CNS, 486
  (September 28, 1973), p.~4.
  杨朴：《反潮流精神》，《人民日报》，1973年8月16日，第3版。浙江省共青团委员会的一篇文章支持``反潮流''精神。ZPS，1973年9月8日；CNS，第486期（1973年9月28日），第4页。}
It called on Chinese youth to emulate Zhang Tiesheng, the rusticated
youth who had submitted a blank examination paper, and ``go against the
tide'' (反潮流). Its impact across the country was somewhat similar to
that of Mao's big-character poster of August 1966 calling on the young
to rebel against reactionaries. In order to reinforce the message, the
People's Daily published articles on 18 August commemorating the 7th
anniversary of Mao's first reception of Red Guards in the Cultural
Revolution. Attacks on industrial rules and regulations published the
following day suggested a co-ordinated and concerted effort to drum up
the atmosphere for a revival of mass radicalism.\footnote{RMRB, August
  18,1973, p.~3; August 19,1973.
  《人民日报》，1973年8月18日，第3版；1973年8月19日。} According to
China News Summary,\footnote{CNS, 482 (August 30,1973), pp.~1-4.
  CNS，第482期（1973年8月30日），第1-4页。} Henan and Zhejiang were the
only two provinces to express immediate support for the leftist
initiatives. Perhaps the Zhejiang government office in Shanghai had
alerted the Zhejiang authorities to the strategic importance of these
articles emanating from Beijing.

在中共十大召开前后，出现了一系列带有激进色彩的文章。其主题包括知识青年上山下乡后的困境、大学入学考试的不公平，以及新的``修正主义教育路线''。1973年8月16日，《人民日报》发表了一篇文章，这篇文章在随后十二个月中产生了广泛影响。\footnote{Yang
  Pu, ``反潮流精神'' (The spirit of going against the tide), RMRB,
  August 16, 1973, p.~3. An article by the Zhejiang CYL Committee
  supported the spirit of ``going against the tide''. ZPS, September 8,
  1973, CNS, 486 (September 28, 1973), p.~4.
  杨朴：《反潮流精神》，《人民日报》，1973年8月16日，第3版。浙江省共青团委员会的一篇文章支持``反潮流''精神。ZPS，1973年9月8日；CNS，第486期（1973年9月28日），第4页。}
它号召中国青年学习交白卷的下乡知识青年张铁生，并``反潮流''。其在全国造成的影响，在某种程度上类似于毛泽东1966年8月号召青年起来造反、反对反动派的大字报。为了强化这一信息，《人民日报》于8月18日发表文章，纪念毛泽东在文化大革命中首次接见红卫兵七周年。次日发表的对工业规章制度的攻击，表明存在一种有协调、有组织的努力，意在营造群众激进主义复兴的气氛。\footnote{RMRB,
  August 18,1973, p.~3; August 19,1973.
  《人民日报》，1973年8月18日，第3版；1973年8月19日。}
据《中国新闻摘要》（China News Summary）称，\footnote{CNS, 482 (August
  30,1973), pp.~1-4. CNS，第482期（1973年8月30日），第1-4页。}
河南和浙江是仅有的两个立即表示支持左派倡议的省份。也许浙江省政府驻上海办事处已经提醒浙江当局注意这些来自北京的文章所具有的战略重要性。

In a political atmosphere reminiscent of mid-1966, the CCP National
Congress opened amidst great secrecy on 24 August 1973.\footnote{For a
  thorough review of the salient issues and factional groups at the
  congress, see Domes, China after the Cultural Revolution, pp.~184-204;
  Richard Wich, ``The Tenth Party Congress: The Power Structure and
  Succession Question'', CQ, 58 (1974), pp.~231-48.
  关于这次大会上的主要问题和派系集团的详尽考察，见 Domes, China after
  the Cultural Revolution，第184-204页；Richard Wich, ``The Tenth Party
  Congress: The Power Structure and Succession Question'',
  CQ，第58期（1974年），第231-248页。} In his report on the revision of
the CCP statutes, Wang Hongwen praised the trend of ``going against the
tide'' and the publicized confessions by individuals who had obtained
privileged treatment or favors by ``going through the back door''
(走后门) as manifestations of a revolutionary spirit. He also announced
that Mao's Cultural Revolution cadre policy of combining the old, the
middle-aged and the young (老中青三结合) would be written into the party
statutes.\footnote{See The Tenth National Congress of the Communist
  Party of China (Documents), (Beijing: Foreign Languages Press, 1973),
  pp.~48-51,54,67. 见 The Tenth National Congress of the Communist Party
  of China (Documents)（北京：外文出版社，1973年），第48-51、54、67页。}
Flushed with success by their gains in the appointments to the new CC
and Politburo, the radicals stepped up the pressure on Zhou Enlai.

在一种令人想起1966年中期的政治气氛中，中共全国代表大会于1973年8月24日在高度秘密状态下开幕。\footnote{For
  a thorough review of the salient issues and factional groups at the
  congress, see Domes, China after the Cultural Revolution, pp.~184-204;
  Richard Wich, ``The Tenth Party Congress: The Power Structure and
  Succession Question'', CQ, 58 (1974), pp.~231-48.
  关于这次大会上的主要问题和派系集团的详尽考察，见 Domes, China after
  the Cultural Revolution，第184-204页；Richard Wich, ``The Tenth Party
  Congress: The Power Structure and Succession Question'',
  CQ，第58期（1974年），第231-248页。}
王洪文在关于修改党章的报告中，称赞``反潮流''的趋势，并把那些通过``走后门''获得特殊待遇或好处的人公开忏悔，称为革命精神的表现。他还宣布，毛泽东在文化大革命中的干部政策，即老、中、青三结合，将写入党章。\footnote{See
  The Tenth National Congress of the Communist Party of China
  (Documents), (Beijing: Foreign Languages Press, 1973),
  pp.~48-51,54,67. 见 The Tenth National Congress of the Communist Party
  of China (Documents)（北京：外文出版社，1973年），第48-51、54、67页。}
激进派因在新一届中央委员会和政治局任命中取得收获而踌躇满志，进一步加大了对周恩来的压力。

One delegate who would have savored the tenor of these remarks was Weng
Senhe. The clause regarding cadres would later provide Weng and his
colleagues with the ideological justification for their strategy of
inducting fellow rebels into the party and promoting them as cadres.
Weng's presence as a member of the Zhejiang delegation had caused
controversy among his fellow workers at the silk complex.\footnote{Visit
  to Silk Complex, January 12,1977. 1977年1月12日访问丝绸联合厂。} Zhang
Chunqiao and Wang Hongwen were reportedly determined to secure Weng's
selection and they exerted pressure on Tan Qilong to this end. It was
later claimed that in December 1972 Zhang had told the ZPC leaders that
Weng Senhe ``was the first among the workers throughout the country who
rose and rebelled during the great cultural revolution and that one
should come to see his history!'' In the summer of 1973 Wang Hongwen
told Tan Qilong that in electing workers to the Zhejiang delegation for
the 10th Congress ``It's necessary to take history into consideration;
even if it is somewhat difficult, we must strive our hardest for his
{[}Weng's{]} election''.\footnote{Exposure by Wang Hongyi, February 7,
  1977, Weng Senhe bufen zuizheng cailiao, p.~7, transl. in ``Document
  of the CCP CC'', zhongfa (1977), No.~37, I\&S, 15:1, p.~110. See also
  China Reconstructs (July 1977), p.~40; Gao Gao and Yan Jiaqi, ``Wenhua
  dageming'', p.~486. Zhang, Wang and Yao Wenyuan applied similar
  pressure on the Yunnan committee to ensure the selection of Zhu Kejia
  as a delegate. Gao Gao and Yan Jiaqi,''Wenhua dageming'', p.~486. Shen
  Chuyun (party secretary, Hangzhou Silk Complex), RMRB, May 17,1977,
  p.~2.
  王洪义的揭发，1977年2月7日，《翁森鹤部分罪证材料》，第7页；译文见``中共中央文件''，中发（1977）第37号，I\&S，15:1，第110页。另见
  China
  Reconstructs（1977年7月），第40页；高皋、严家其：《文化大革命》，第486页。张、王和姚文元对云南省委施加了类似压力，以确保朱克家当选代表。高皋、严家其：《文化大革命》，第486页。沈楚云（杭州丝绸联合厂党委书记），《人民日报》，1977年5月17日，第2版。}

有一位代表会特别欣赏这些话的基调，那就是翁森鹤。关于干部的这一条款，后来为翁森鹤及其同僚把造反派同伴吸收入党并提拔为干部的策略，提供了意识形态上的正当性。翁森鹤作为浙江代表团成员出席，在他所在丝绸联合企业的同事中引发了争议。\footnote{Visit
  to Silk Complex, January 12,1977. 1977年1月12日访问丝绸联合厂。}
据说张春桥和王洪文决心确保翁森鹤当选，并为此向谭启龙施压。后来有人声称，1972年12月张春桥曾对浙江省委领导人说，翁森鹤``是全国工人中在伟大文化大革命中最早起来造反的，应该看看他的历史！''1973年夏，王洪文告诉谭启龙，在选举出席十大浙江代表团的工人代表时，``要考虑历史；即使有些困难，也要尽最大的努力争取他（翁森鹤）当选''。\footnote{Exposure
  by Wang Hongyi, February 7, 1977, Weng Senhe bufen zuizheng cailiao,
  p.~7, transl. in ``Document of the CCP CC'', zhongfa (1977), No.~37,
  I\&S, 15:1, p.~110. See also China Reconstructs (July 1977), p.~40;
  Gao Gao and Yan Jiaqi, ``Wenhua dageming'', p.~486. Zhang, Wang and
  Yao Wenyuan applied similar pressure on the Yunnan committee to ensure
  the selection of Zhu Kejia as a delegate. Gao Gao and Yan
  Jiaqi,''Wenhua dageming'', p.~486. Shen Chuyun (party secretary,
  Hangzhou Silk Complex), RMRB, May 17,1977, p.~2.
  王洪义的揭发，1977年2月7日，《翁森鹤部分罪证材料》，第7页；译文见``中共中央文件''，中发（1977）第37号，I\&S，15:1，第110页。另见
  China
  Reconstructs（1977年7月），第40页；高皋、严家其：《文化大革命》，第486页。张、王和姚文元对云南省委施加了类似压力，以确保朱克家当选代表。高皋、严家其：《文化大革命》，第486页。沈楚云（杭州丝绸联合厂党委书记），《人民日报》，1977年5月17日，第2版。}

Even before the commencement of the Congress the Zhejiang rebels were
making preparations to launch their attack on the provincial leadership
of Tan Qilong and Tie Ying. In July 1973 Zhang Yongsheng commissioned
Wang Jiashu, a staff member in the Hangzhou Mao Zedong Thought
Propaganda Team office, to gather evidence from the provincial and
municipal party offices and draft a document denouncing Tan and Tie for
negating the Cultural Revolution. In September Zhang lobbied members and
alternate members of the ZPC to sign the document which charged the
provincial leadership with suppressing the innovative forces (革新力量)
and propping up the old-timers (守旧势力).\footnote{Zhang Yongshengde
  zuizheng cailiao, p.~56. 《张永生的罪证材料》，第56页。}

甚至在大会开幕之前，浙江造反派就已准备向谭启龙和铁瑛领导的省级领导层发起攻击。1973年7月，张永生委托杭州毛泽东思想宣传队办公室工作人员王家树，从省、市两级党委机关搜集证据，并起草一份文件，谴责谭启龙和铁瑛否定文化大革命。9月，张永生游说浙江省委委员和候补委员在该文件上签名；文件指控省领导层压制革新力量，扶植守旧势力。\footnote{Zhang
  Yongshengde zuizheng cailiao, p.~56. 《张永生的罪证材料》，第56页。}

To coincide with the opening of the Party Congress in Beijing and before
the public appearance of the open letter to Tie and Tan, a poster was
pasted up in the courtyard of the CCP HMC offices in downtown Hangzhou.
It was signed by a group of ``new office-worker cadres'' and criticized
both the provincial and municipal authorities. The thrust of the
complaint was that in the campaign against Lin Biao and his followers in
1972-73 the signatories had been ``swept back to camp'' (打扫回营) that
is, to the factories and colleges from whence they had been promoted.
They demanded reinstatement and assembled in the building occupied by
the municipal workers' congress to form a study class. To press their
complaints further the group reportedly moved into municipal party
offices, forcibly occupied the meeting room of the party standing
committee and seized documents. They believed that these actions would
force the authorities to respond to their gripes. The Hangzhou workers'
congress took up the cause and convened a meeting of its leadership at
the New China Theater in the center of Hangzhou. It decided to organize
a Beijing-bound complaints delegation (赴京控告团 or 上访团) of
disgruntled cadres.\footnote{HZRB, July 19,1977, September 17,1977. This
  phenomenon was not confined to Zhejiang. See Hao Mengbi and Duan
  Haoran, Liushinian, p.~636.
  《杭州日报》，1977年7月19日、1977年9月17日。这一现象并不限于浙江。见郝梦笔、段浩然：《六十年》，第636页。}

为了配合北京党代会的召开，并且在致铁瑛、谭启龙的公开信正式出现之前，杭州市中心中共杭州市委机关院内贴出了一张大字报。大字报由一批``新机关干部''署名，同时批评省、市两级当局。控诉的要点是，在1972至1973年反对林彪及其追随者的运动中，署名者被``打扫回营''，也就是被送回他们原先得到提拔时所在的工厂和学校。他们要求恢复职务，并聚集在市工代会占用的大楼内组成学习班。据说，为了进一步施压，这一群体进入市委机关，强行占据市委常委会议室并夺取文件。他们相信，这些行动会迫使当局回应他们的不满。杭州市工代会接过这一事业，并在杭州市中心的新中国剧院召开领导层会议，决定组织一支由心怀不满的干部组成的赴京控告团或上访团。\footnote{HZRB,
  July 19,1977, September 17,1977. This phenomenon was not confined to
  Zhejiang. See Hao Mengbi and Duan Haoran, Liushinian, p.~636.
  《杭州日报》，1977年7月19日、1977年9月17日。这一现象并不限于浙江。见郝梦笔、段浩然：《六十年》，第636页。}

In October-November ``complaints delegations'' assembled in Hangzhou
from various counties across the province. They participated in a series
of meetings held at the ZPRC hostel in Mishi Lane (米市巷) just off
Moganshan Rd in the north-west of the city, and at the ZPC retreat on
Mt. Pingfeng. Over forty counties had sent delegations and twenty of the
delegations comprised a total of 200 to 300 members. The original plan
was to send the 300 delegates to the national capital, but the word from
Beijing was that such a large number would be undesirable and that it
would be preferable if Zhang, Weng and He did not personally lead the
delegation. It was decided instead that Li Xiantong (Zhang's colleague
in United Headquarters and a former student at Hangzhou University) and
Wang Jiashu would lead seventeen complainants on the mission.\footnote{Zhang
  Yongshengde zuizheng cailiao, pp.~11-12, 50, 52.
  《张永生的罪证材料》，第11-12、50、52页。}

10月至11月，来自全省各县的``控告团''聚集到杭州。他们参加了一系列会议，会议地点包括城西北莫干山路附近米市巷的浙江省革命委员会招待所，以及位于屏风山的浙江省委休养处。四十多个县派出了代表团，其中二十个代表团合计有200至300名成员。原计划是派300名代表前往首都，但北京传来的意见是，这么多人并不合适，而且最好不要由张永生、翁森鹤和贺贤春亲自带队。于是改为由李显通（张永生在联合总部的同事、原杭州大学学生）和王家树率领十七名控告者执行这一任务。\footnote{Zhang
  Yongshengde zuizheng cailiao, pp.~11-12, 50, 52.
  《张永生的罪证材料》，第11-12、50、52页。}

Wang Hongwen soon became involved in the affairs of the province once
more. On 10 August, Zhang, Weng and He had sent him material with a
covering letter informing him of developments in Zhejiang since his
visit of January 1973.\footnote{Ibid, p.~46. 同上，第46页。} On 4
October Wang, now Vice-chairman of the CCP CC, telephoned Zhang
Yongsheng and encouraged him to launch an attack on the leaders of the
ZPC and prepare for a lengthy battle. Wang informed Zhang that ``several
problems have emerged in Zhejiang''. Zhang sprang into action and
assembled his supporters. They formulated a document entitled ``Report
on the situation in learning from Shanghai in selecting and training
successors for the proletarian revolutionary cause''
(学习上海市选拔培养无产阶级革命事业接班人的情况汇报). The document
warned that ``activists of the counter-revolutionary 'May 16' clique''
should be ``cleared out and investigated''. It specified that the
criteria for selection to the ``proletarian revolutionary cause'' would
be previous experience in line struggles and whether an individual had
rebelled ``early or later'' in the Cultural Revolution.\footnote{Gao Gao
  and Yan Jiaqi, ``Wenhua dageming'', p.~483.
  高皋、严家其：《文化大革命》，第483页。}

王洪文很快再次介入该省事务。8月10日，张永生、翁森鹤和贺贤春曾把材料寄给他，并附信告知他自1973年1月他访问浙江以来当地的发展情况。\footnote{Ibid,
  p.~46. 同上，第46页。}
10月4日，已任中共中央副主席的王洪文给张永生打电话，鼓励他向浙江省委领导人发起攻击，并准备进行长期斗争。王洪文告诉张永生，``浙江出了几个问题''。张永生立即行动起来，召集支持者。他们拟定了一份题为《学习上海市选拔培养无产阶级革命事业接班人的情况汇报》的文件。文件警告说，应当把``反革命`五一六'集团的活跃分子''\,``清理出来并加以审查''。它规定，选拔``无产阶级革命事业''接班人的标准，是过去参加路线斗争的经历，以及个人在文化大革命中造反是``早还是晚''。\footnote{Gao
  Gao and Yan Jiaqi, ``Wenhua dageming'', p.~483.
  高皋、严家其：《文化大革命》，第483页。}

The selection of loyalists on the basis of when they had joined the
Cultural Revolution to some extent cut across the Red Storm/United
Headquarters division. Both Zhang Yongsheng and Weng Senhe had been
amongst the earliest rebels in Zhejiang, although initially their
unit-based groups had joined opposing umbrella mass organizations.
Perhaps Zhang and Weng were trying to heal the divisions among Cultural
Revolution activists dating from 1967 with the objective of recruiting
and remobilizing rebels from both sides, in particular those who had
survived the witch-hunt against ``May 16'' rebels in the years 1969-71.
In this endeavor they may have been partly successful.

以参加文化大革命时间早晚为依据来选择忠诚分子，在一定程度上跨越了``红暴''和``联合总部''之间的分野。张永生和翁森鹤都属于浙江最早的造反者之列，尽管最初他们各自以单位为基础的组织加入的是彼此对立的伞状群众组织。也许张、翁二人试图弥合1967年以来文革积极分子之间的裂痕，其目标是从双方招募并重新动员造反者，尤其是那些在1969至1971年针对``五一六''造反派的清查中幸存下来的人。在这项努力中，他们或许取得了部分成功。

The same day as Wang Hongwen had telephoned Zhang a big-character
poster, in the form of an open letter from Zhang, Weng and He to Tan and
Tie, appeared on the streets of Hangzhou.\footnote{Weng Senhe bufen
  zuizheng cailiao, p.~39; Zhang Yongshengde zuizheng cailiao, pp.~48-9.
  《翁森鹤部分罪证材料》，第39页；《张永生的罪证材料》，第48-49页。} At
a meeting of ten or so rebels early in September, some of those in
attendance had expressed reservations about putting up a poster in the
streets. They recommended that the letter be posted inside government
offices. Weng brushed aside their worries in typically brusque fashion
but in the end it was undoubtedly Wang Hongwen's encouragement which
gave the rebels the courage to proceed.\footnote{Weng Senhe bufen
  zuizheng cailiao, pp.~15-16. 《翁森鹤部分罪证材料》，第15-16页。} They
coined a sixteen-character guiding slogan which read: ``grasp the
initiative, struggle on the spot, spread our influence and open up
prospects''.\footnote{HZRB, August 14,1978.
  《杭州日报》，1978年8月14日。}

就在王洪文给张永生打电话的同一天，一张以张、翁、贺三人致谭启龙和铁瑛公开信形式出现的大字报，出现在杭州街头。\footnote{Weng
  Senhe bufen zuizheng cailiao, p.~39; Zhang Yongshengde zuizheng
  cailiao, pp.~48-9.
  《翁森鹤部分罪证材料》，第39页；《张永生的罪证材料》，第48-49页。}
9月初，在一次约十名造反派参加的会议上，部分与会者对在街头张贴大字报表示保留。他们建议把信贴在政府机关内部。翁森鹤以其一贯粗率的方式把这些担忧搁置一旁，但最终无疑是王洪文的鼓励给了造反派继续行动的勇气。\footnote{Weng
  Senhe bufen zuizheng cailiao, pp.~15-16.
  《翁森鹤部分罪证材料》，第15-16页。}
他们提出了一条十六字指导口号：``抓住主动，就地斗争，扩大影响，打开局面''。\footnote{HZRB,
  August 14,1978. 《杭州日报》，1978年8月14日。}

On 25 October Zhang Yongsheng, Li Xiantong and Yan Yihuan (a leader of
the ``complaints' delegation'' from Zhejiang University) signed a
big-character poster refusing to join the study class organized by the
provincial leaders to study the documents of the 10th National
Congress.\footnote{Zhang Yongshengde zuizheng cailiao, p.~57.
  《张永生的罪证材料》，第57页。} In a speech of 20 November Zhang
argued that Wang's telephone directives of 4 October should form the
basis for studying the documents.\footnote{Ibid, p.~18. 同上，第18页。}
With the political pressure on the provincial leadership building up in
Hangzhou, Wang Hongwen applied more from Beijing.

10月25日，张永生、李显通和严怡环（浙江大学``控告团''的一名领导人）署名一张大字报，拒绝参加省领导组织的学习班，该学习班旨在学习十大文件。\footnote{Zhang
  Yongshengde zuizheng cailiao, p.~57. 《张永生的罪证材料》，第57页。}
张永生在11月20日的一次讲话中主张，应以王洪文10月4日的电话指示作为学习文件的基础。\footnote{Ibid,
  p.~18. 同上，第18页。}
随着杭州对省领导层的政治压力不断增强，王洪文又从北京施加了更多压力。

On 16 November he telephoned Tan Qilong and instructed him to dissuade
the ``complaints delegation'' from leaving Hangzhou.\footnote{The
  ``complaints delegations'' remained on Mt. Pingfeng until March/April
  1974, except for a return trip home during February to celebrate the
  spring festival. See Weng Senhe bufen zuizheng cailiao, pp.~30, 57.
  ``上访代表团''一直留在屏风山，直到1974年3月或4月，其间只在2月回乡过春节。见《翁森鹤部分罪证材料》，第30、57页。}
Wang also asked Tan to make a self-criticism on behalf of the provincial
party committee for the ``excesses'' committed against the former
leaders of United Headquarters over the previous eighteen months. Wang
also requested Tan to support the principle of ``going against the
tide''. Tan, Tie and other members of the standing committee reportedly
agreed not to act on Wang's instructions and decided not to relay Wang's
suggestion that Tan make a self-criticism. That the Zhejiang party
leaders refused outright is highly unlikely given Wang's status as Mao's
newly-anointed successor.\footnote{For a discussion of this highly
  sensitive issue see Ye Yonglie, Wang Hongwen xingshuailu, pp.~376-96;
  Ye Yonglie, Mao Zedong yishi, pp.~247-69, esp.~pp.~250-55.
  关于这一高度敏感问题的讨论，见叶永烈：《王洪文兴衰录》，第376-396页；叶永烈：《毛泽东轶事》，第247-269页，尤见第250-255页。}
They may have fudged, procrastinated or prevaricated in the hope that
countervailing orders would be issued.

11月16日，王洪文给谭启龙打电话，指示他劝阻``控告团''离开杭州。\footnote{The
  ``complaints delegations'' remained on Mt. Pingfeng until March/April
  1974, except for a return trip home during February to celebrate the
  spring festival. See Weng Senhe bufen zuizheng cailiao, pp.~30, 57.
  ``上访代表团''一直留在屏风山，直到1974年3月或4月，其间只在2月回乡过春节。见《翁森鹤部分罪证材料》，第30、57页。}
王洪文还要求谭启龙代表省委，就过去十八个月中对联合总部原领导人采取的``过火''行为作自我批评。王还要求谭支持``反潮流''原则。据说谭启龙、铁瑛和常委会其他成员同意不执行王的指示，并决定不传达王关于谭应作自我批评的建议。鉴于王洪文当时作为毛泽东新近指定的接班人的地位，浙江党政领导人公然拒绝的可能性极小。\footnote{For
  a discussion of this highly sensitive issue see Ye Yonglie, Wang
  Hongwen xingshuailu, pp.~376-96; Ye Yonglie, Mao Zedong yishi,
  pp.~247-69, esp.~pp.~250-55.
  关于这一高度敏感问题的讨论，见叶永烈：《王洪文兴衰录》，第376-396页；叶永烈：《毛泽东轶事》，第247-269页，尤见第250-255页。}
他们可能是含糊其辞、拖延观望或闪烁其词，希望会有相反的命令下达。

Nevertheless, news of Wang's instructions soon reached Zhang Yongsheng,
having been leaked by a member of the standing committee of the
municipal party committee. A senior member of Tan's executive, Lai Keke,
undermined Tan's stand by writing a report to Wang Hongwen expressing
sympathy for the ``fighters going against the tide'' and lashing his
senior colleagues for their political outlook and performance. Lai sent
Wang a second report on 9 December 1973, distancing himself from the
provincial leadership and extolling the youthful
Vice-chairman.\footnote{CCP HMC propaganda department criticism group,
  ``A loyal follower of the Gang of Four in Hangzhou'', HZRB, October
  21,1977; Zhou Feng (secretary of the CCP HMC), ``Completely dig out
  the secret traitor who opposed the party and caused trouble'', HZRB,
  October 24, 1977; ZJRB, May 21,1977, p.~1.
  中共杭州市委宣传部批判组：《杭州的``四人帮''死党》，《杭州日报》，1977年10月21日；周峰（中共杭州市委书记）：《彻底挖出反党捣乱的秘密叛徒》，《杭州日报》，1977年10月24日；《浙江日报》，1977年5月21日，第1版。}
The divulgence of personal communications between party leaders in order
to score political points soon became a common occurrence in Zhejiang.
Factional loyalties superseded party discipline with enormous
consequences for Tan's leadership in the province. Wall posters appeared
on the streets accusing Tan of trying to change and conceal
``Vice-chairman Wang's directives''.

然而，王洪文指示的消息很快就传到了张永生那里，原因是市委常委会的一名成员泄露了消息。谭启龙领导班子中的一名高级成员赖可可，通过给王洪文写报告表达对``反潮流战士''的同情，并猛烈抨击其高级同僚的政治观点和工作表现，从而削弱了谭的立场。1973年12月9日，赖可可又给王洪文送去第二份报告，与省领导层拉开距离，并赞颂这位年轻的副主席。\footnote{CCP
  HMC propaganda department criticism group, ``A loyal follower of the
  Gang of Four in Hangzhou'', HZRB, October 21,1977; Zhou Feng
  (secretary of the CCP HMC), ``Completely dig out the secret traitor
  who opposed the party and caused trouble'', HZRB, October 24, 1977;
  ZJRB, May 21,1977, p.~1.
  中共杭州市委宣传部批判组：《杭州的``四人帮''死党》，《杭州日报》，1977年10月21日；周峰（中共杭州市委书记）：《彻底挖出反党捣乱的秘密叛徒》，《杭州日报》，1977年10月24日；《浙江日报》，1977年5月21日，第1版。}
为了获取政治利益而泄露党内领导人之间私人通信的做法，很快在浙江成为常见现象。派系忠诚压倒了党的纪律，对谭启龙在该省的领导地位造成了巨大后果。街头出现了墙报，指责谭企图篡改和隐瞒``王副主席的指示''。

Shortly thereafter, possibly at Wang's behest the 300-strong counties
delegation led by Zhang, Weng and He, rebelled (发难) against the
provincial leadership. Holding banners they marched to Huagang (花港)
reception center, situated on the western side of the West Lake, where
the provincial leadership was meeting. They then gathered on Mt.
Pingfeng to plan their next move.\footnote{Weng Senhe bufen zuizheng
  cailiao, p.~17. 《翁森鹤部分罪证材料》，第17页。} With Wang Hongwen's
approval, they summoned Tan and other provincial leaders to the mountain
for questioning, and defiantly asserted an equal status with the party
leadership (分庭抗礼). For four and a half days at the end of November
1973, Tan and his colleagues were forcefully interrogated
(揪斗).\footnote{Ibid, pp.~28, 46. 同上，第28、46页。} The rebels'
objective was to force them to agree to placing more young cadres in
leading groups. First of all they sought the rehabilitation of rebels
who had come under fire in the political campaign of 1972-73, as the
``suggestion'' paper drawn up by the rebels on 8 December 1973
revealed.\footnote{Zhang Yongshengde zuizheng cailiao, p.~96.
  《张永生的罪证材料》，第96页。} Tie Ying came under further attack at
the end of December (or early January 1974)\footnote{Zhang Yongshengde
  zuizheng cailiao, p.~69 (Yang Baoshan's confession of December 20,1977
  - Yang was a leader of the complaints' delegation from Jiaxing); Weng
  Senhe bufen zuizheng cailiao, p.~45 (Li Xiantong's confession of March
  24, 1977).
  《张永生的罪证材料》，第69页（杨宝山1977年12月20日的交代，杨是嘉兴上访代表团的一名负责人）；《翁森鹤部分罪证材料》，第45页（李显通1977年3月24日的交代）。}
when he was detained for another three to four days at the Mishi Lane
hostel.\footnote{Weng Senhe bufen zuizheng cailiao, p.~45.
  《翁森鹤部分罪证材料》，第45页。} The rebels again raised critical
opinions regarding the shortcomings in the work of the provincial
leadership. The occupation (安营扎寨) of Mt. Pingfeng and the detention
of provincial leaders were flagrant acts of rebellion against the
provincial leadership which seriously weakened its authority. Tan Qilong
and Tie Ying had been humiliated before their peers and defied and
treated with contempt by their youthful protagonists.

不久之后，可能是在王洪文授意下，由张、翁、贺率领的三百人左右的各县代表团向省领导层``发难''。他们打着横幅，游行到位于西湖西侧的花港接待中心，当时省领导层正在那里开会。随后他们聚集到屏风山，筹划下一步行动。\footnote{Weng
  Senhe bufen zuizheng cailiao, p.~17. 《翁森鹤部分罪证材料》，第17页。}
在王洪文的批准下，他们把谭启龙和其他省领导召到山上接受质询，并公然宣称自己与党委领导层处于平等地位，形成``分庭抗礼''之势。1973年11月底，谭启龙及其同僚连续四天半遭到强行讯问和``揪斗''。\footnote{Ibid,
  pp.~28, 46. 同上，第28、46页。}
造反派的目标，是迫使他们同意把更多年轻干部安排进领导班子。首先，他们要求为在1972至1973年政治运动中受到攻击的造反派平反；造反派于1973年12月8日拟定的``建议''文件揭示了这一点。\footnote{Zhang
  Yongshengde zuizheng cailiao, p.~96. 《张永生的罪证材料》，第96页。}
12月底（或1974年1月初），\footnote{Zhang Yongshengde zuizheng cailiao,
  p.~69 (Yang Baoshan's confession of December 20,1977 - Yang was a
  leader of the complaints' delegation from Jiaxing); Weng Senhe bufen
  zuizheng cailiao, p.~45 (Li Xiantong's confession of March 24, 1977).
  《张永生的罪证材料》，第69页（杨宝山1977年12月20日的交代，杨是嘉兴上访代表团的一名负责人）；《翁森鹤部分罪证材料》，第45页（李显通1977年3月24日的交代）。}
铁瑛再次遭到攻击，他又被扣留在米市巷招待所三到四天。\footnote{Weng Senhe
  bufen zuizheng cailiao, p.~45. 《翁森鹤部分罪证材料》，第45页。}
造反派再次就省领导层工作中的缺点提出批评意见。占据屏风山``安营扎寨''以及扣留省级领导人，是对省领导层的公然反叛，严重削弱了其权威。谭启龙和铁瑛在同僚面前受辱，并遭到这些年轻对手的蔑视、顶撞和轻慢对待。

\begin{Shaded}
\begin{Highlighting}[]
\NormalTok{TABLE TEN}
\NormalTok{                     Leading figures of the two factions of 1973{-}75}

\NormalTok{                              The "mountain top" faction}

\NormalTok{      NAME                                 POSITION}

\NormalTok{Zhang Yongsheng             vice{-}chairman, ZPRC}
\NormalTok{Weng Senhe                   vice{-}chairman, ZPTUC; member, standing committee}
\NormalTok{                             ZPRC}
\NormalTok{He Xianchun                  vice{-}chairman, ZPTUC; chairman Hangzhou Workers\textquotesingle{}}
\NormalTok{                             Congress}
\NormalTok{ Huang Yintang               party secretary and chairman revolutionary committee,}
\NormalTok{                             Hangzhou silk complex}
\NormalTok{ Xia Genfa                   party secretary and chairman revolutionary committee,}
\NormalTok{                             Hangzhou oxygen generator factory}
\NormalTok{ Organizational bases:       Hangzhou Municipal Workers\textquotesingle{} Congress}
\NormalTok{                             Hangzhou Municipal Militia Headquarters}
\NormalTok{ Linkages:                   Zhang Yongsheng to Jiang Qing and Wang Hongwen}
\NormalTok{                             Weng Senhe to Wang Hongwen and Zhang Chunqiao}
\NormalTok{                             to senior cadres of the CCP HMC such as Wang Zida and}
\NormalTok{                             Wang Xing}
\NormalTok{                             to leading cadres of the CCP ZPC such as Lai Keke, Chai}
\NormalTok{                             Qikun and Shen Ce}

\NormalTok{                               The "mountain base" faction}
\NormalTok{ Guo Zhisong                 vice{-}chairman, ZPTUC}
\NormalTok{ Fang Jianwen                vice{-}chairman, ZPTUC}
\NormalTok{ Zhang Jifa                  party group secretary, Hangzhou iron and steel factory}
\NormalTok{ Liang Maomao                group leader, Hangzhou oxygen generator plant}
\NormalTok{ Qiu Honggen                 factory worker}
\NormalTok{ Linkages:                   Zhang Jifa to PLA officials in Beijing}
\NormalTok{                              to former members of Red Storm}
\NormalTok{                              to sympathizers in the ZPMD}
\NormalTok{                              to senior bureaucrats in government departments who}
\NormalTok{                              could serve as intermediaries with supportive senior}
\NormalTok{                              cadres such as Tie Ying}
\end{Highlighting}
\end{Shaded}

\begin{Shaded}
\begin{Highlighting}[]
\NormalTok{表十}
\NormalTok{                     1973{-}75年两派主要人物}

\NormalTok{                              “山上派”}

\NormalTok{      姓名                                 职务}

\NormalTok{张永生                     浙江省革命委员会副主任}
\NormalTok{翁森鹤                     浙江省总工会副主任；浙江省革命委员会常委}
\NormalTok{贺贤春                     浙江省总工会副主任；杭州市工人代表大会主席}
\NormalTok{ 黄阴堂                    杭州丝绸联合企业党委书记、革命委员会主任}
\NormalTok{ 夏根发                    杭州制氧机厂党委书记、革命委员会主任}
\NormalTok{ 组织基础：                杭州市工人代表大会}
\NormalTok{                            杭州市民兵指挥部}
\NormalTok{ 联系关系：                张永生联系江青、王洪文}
\NormalTok{                            翁森鹤联系王洪文、张春桥}
\NormalTok{                            联系中共杭州市委高级干部，如王子达、王兴}
\NormalTok{                            联系中共浙江省委领导干部，如赖可可、柴启琨、沈策}

\NormalTok{                              “山下派”}
\NormalTok{ 郭志松                    浙江省总工会副主任}
\NormalTok{ 方剑文                    浙江省总工会副主任}
\NormalTok{ 张纪法                    杭州钢铁厂党组书记}
\NormalTok{ 梁毛毛                    杭州制氧机厂小组长}
\NormalTok{ 邱洪根                    工厂工人}
\NormalTok{ 联系关系：                张纪法联系北京的解放军官员}
\NormalTok{                            联系红暴前成员}
\NormalTok{                            联系浙江省军区的同情者}
\NormalTok{                            联系政府部门中可作为中间人的高级官僚，}
\NormalTok{                            他们能够同铁瑛等持支持态度的高级干部沟通}
\end{Highlighting}
\end{Shaded}

Zhang Yongsheng aimed to ``haul out restorationist forces'' among the
Zhejiang leadership and to discover which party leaders were prepared to
desert Tan. He found two in Lai Keke and Chai Qikun, who were ready to
join forces with him\footnote{Zhang Yongshengde zuizheng cailiao, p.~45.
  《张永生的罪证材料》，第45页。} as was Wang Zida, First Secretary of
the CCP HMC and a standing committee member of the CCP ZPC. Wang wrote
to the provincial party committee demanding that Tan and Tie ``make
lofty self-criticisms'' and added that ``this is the key question at the
present time''.\footnote{HZRB, July 19,1977; Chen Xia (secretary of the
  CCP HMC), ``Crimes must be accounted for. the poison must be cleared
  away'', HZRB, April 10,1978; Zhang Yongsheng also tried to win over
  the peasant model and party secretary of Shangwang brigade in Shaoxing
  county Wang Jinyou (王金友). See Wang's account in''不信邪不怕压
  上旺愈斗愈兴旺'' (Not taken in by fallacies and not afraid of
  pressure: Shangwang flourishes through struggle) in 举旗抓纲学大寨
  (Raise the banner, grasp the key link and learn from Dazhai),
  (Hangzhou: Zhejiang Renmin Chubanshe, 1977), p.~39.
  《杭州日报》，1977年7月19日；陈夏（中共杭州市委书记）：《罪行必须交代，流毒必须肃清》，《杭州日报》，1978年4月10日；张永生还试图拉拢绍兴县上旺大队的农民典型、党支部书记王金友。见王金友在《举旗抓纲学大寨》（杭州：浙江人民出版社，1977年）中《不信邪不怕压
  上旺愈斗愈兴旺》（不被谬论所惑，不怕压力：上旺在斗争中日益兴旺）一文的叙述，第39页。}
One specific charge levelled against Tan related to remarks he had made
in June 1972 during his inspection tour with Tie Ying of Fuyang county,
adjacent to Hangzhou. On his drive through the countryside, Tan had
commented on the slow pace of development since his previous visit back
in 1954. His words were:

张永生的目标，是在浙江领导层中``揪出复辟势力''，并查明哪些党内领导人准备抛弃谭启龙。他找到了两个人，即赖可可和柴启琨，他们准备与他联手；\footnote{Zhang
  Yongshengde zuizheng cailiao, p.~45. 《张永生的罪证材料》，第45页。}
中共杭州市委第一书记、浙江省委常委王子达也是如此。王子达写信给省委，要求谭启龙和铁瑛``作出高姿态的自我批评''，并补充说``这是当前的关键问题''。\footnote{HZRB,
  July 19,1977; Chen Xia (secretary of the CCP HMC), ``Crimes must be
  accounted for. the poison must be cleared away'', HZRB, April 10,1978;
  Zhang Yongsheng also tried to win over the peasant model and party
  secretary of Shangwang brigade in Shaoxing county Wang Jinyou
  (王金友). See Wang's account in''不信邪不怕压 上旺愈斗愈兴旺'' (Not
  taken in by fallacies and not afraid of pressure: Shangwang flourishes
  through struggle) in 举旗抓纲学大寨 (Raise the banner, grasp the key
  link and learn from Dazhai), (Hangzhou: Zhejiang Renmin Chubanshe,
  1977), p.~39.
  《杭州日报》，1977年7月19日；陈夏（中共杭州市委书记）：《罪行必须交代，流毒必须肃清》，《杭州日报》，1978年4月10日；张永生还试图拉拢绍兴县上旺大队的农民典型、党支部书记王金友。见王金友在《举旗抓纲学大寨》（杭州：浙江人民出版社，1977年）中《不信邪不怕压
  上旺愈斗愈兴旺》（不被谬论所惑，不怕压力：上旺在斗争中日益兴旺）一文的叙述，第39页。}
对谭启龙提出的一项具体指控，涉及他1972年6月同铁瑛一起视察杭州邻近的富阳县时所说的话。谭在驱车穿过乡村时，评论说，自他1954年上次到访以来，当地发展缓慢。他的话是：

\begin{quote}
the tea-bushes on both sides of the road are thinly scattered and the
ground is not even; I've been away from here for eighteen years and
there hasn't been much of a change.
\end{quote}

\begin{quote}
路两边的茶树稀稀落落，地也不平；我离开这里十八年了，变化不大。
\end{quote}

Tan's remarks were circulated throughout Zhejiang as evidence of his
negative attitude toward the Cultural Revolution.\footnote{Kong Xianlian
  (in 1972 secretary of the Fuyang county party committee), ``How did
  the political rumor about 'there has been no change in Zhejiang for
  eighteen years' get about'', HZRB, November 26,1977; HZRB, November
  19,1977.
  孔宪连（1972年任富阳县委书记）：《``浙江十八年没有变''的政治谣言是怎样散布出来的》，《杭州日报》，1977年11月26日；《杭州日报》，1977年11月19日。}

谭启龙的这些话在浙江全省流传，被当作他对文化大革命持消极态度的证据。\footnote{Kong
  Xianlian (in 1972 secretary of the Fuyang county party committee),
  ``How did the political rumor about 'there has been no change in
  Zhejiang for eighteen years' get about'', HZRB, November 26,1977;
  HZRB, November 19,1977.
  孔宪连（1972年任富阳县委书记）：《``浙江十八年没有变''的政治谣言是怎样散布出来的》，《杭州日报》，1977年11月26日；《杭州日报》，1977年11月19日。}

The principal and overriding issue for the rebels was, however, power.
Weng Senhe stated bluntly at a meeting in March 1974 that ``It's already
been seven to eight years of the Cultural Revolution, and recently what
have we got organizationally?\ldots{} We must get leadership
power!''\footnote{Weng Senhe bufen zuizheng cailiao, p.~35.
  《翁森鹤部分罪证材料》，第35页。} Weng later confessed in prison:

然而，对造反派来说，首要且压倒一切的问题是权力。翁森鹤在1974年3月的一次会议上直截了当地说：``文化大革命已经七八年了，最近我们在组织上得到了什么？\ldots\ldots 我们一定要掌握领导权！''\footnote{Weng
  Senhe bufen zuizheng cailiao, p.~35. 《翁森鹤部分罪证材料》，第35页。}
翁后来在狱中供认说：

\begin{quote}
I saw a speech by Zhang Chunqiao. He said: ``the rebels must not only
destroy but also hold on to power''. When I read this I felt as good as
if I had just eaten an icy-pole on a summer's day and felt grateful that
Zhang Chunqiao really understood and felt for us (贴心人).\footnote{Ibid,
  p.~21; HZRB, December 7,1977.
  同上，第21页；《杭州日报》，1977年12月7日。}
\end{quote}

\begin{quote}
我看到张春桥的一篇讲话。他说：``造反派不但要破，还要掌权。''我读到这句话时，感觉就像夏天刚吃了一根冰棍一样痛快，并且感到张春桥真是理解我们、体贴我们的人（贴心人）。\footnote{Ibid,
  p.~21; HZRB, December 7,1977.
  同上，第21页；《杭州日报》，1977年12月7日。}
\end{quote}

He Xianchun agreed. He threatened that ``we won't leave the mountain
until we've solved this personnel problem''. Zhang Yongsheng reminded
his fellow rebels of the suffering that they had endured in 1972 because
too few of them had joined the CCP and said that this lesson should
never be forgotten.\footnote{Zhang Yongshengde zuizheng cailiao, pp.~84,
  86. 《张永生的罪证材料》，第84、86页。}

贺贤春也表示同意。他威胁说：``不解决这个人事问题，我们就不下山。''张永生提醒造反派同伴说，1972年他们之所以遭受痛苦，是因为他们中加入中共的人太少，并说这个教训永远不能忘记。\footnote{Zhang
  Yongshengde zuizheng cailiao, pp.~84, 86.
  《张永生的罪证材料》，第84、86页。}

As earlier chapters of this study have detailed, the rebel leaders had
formed a close association in the old United Headquarters mass
organization. Their factional enemies, who included former activists
from Red Storm, tended to support the provincial leadership. However,
the two groups which emerged in 1974, while having their origins in the
divisions of 1966-69, were formed specifically in relation to the Mt.
Pingfeng rebellion. One became known as the ``mountain top'' faction
(山上派) and the other the ``mountain base'' faction
(山下派).\footnote{CNS, 578 (August 13,1975), p.~1.
  CNS，第578期（1975年8月13日），第1页。} Zhang, Weng and He led the the
``mountain top'' faction. Guo Zhisong and Fang Jianwen, both
vice-chairmen of the ZPTUC who had been on opposing sides in the 1960s,
led the latter faction. Other Cultural Revolution activists may have
remained non-committed. The political fall-out from the occupation of
Mt. Pingfeng soon spread across the whole province and this sensational
incident helped trigger off a series of events which eventually forced
Zhejiang before the attention of the whole country in July 1975.

如本研究前几章所详述，造反派领导人曾在旧有的``联合总部''群众组织中形成密切联系。他们的派系敌人，包括红暴的前积极分子，则倾向于支持省领导层。然而，1974年出现的这两个群体，虽然起源于1966至1969年的分裂，却是围绕屏风山反叛事件而具体形成的。一个被称为``山上派''，另一个被称为``山下派''。\footnote{CNS,
  578 (August 13,1975), p.~1. CNS，第578期（1975年8月13日），第1页。}
张、翁、贺领导``山上派''。郭志松和方剑文都曾任浙江省总工会副主任，并且在20世纪60年代曾处于对立阵营，他们领导后一派。其他文革积极分子可能仍然没有明确站队。占据屏风山所造成的政治后果很快扩散到全省，这一轰动事件促发了一系列事态，最终使浙江在1975年7月进入全国关注的视野。

To destabilize further the political environment in Zhejiang, Zhang
Yongsheng and Weng Senhe managed to secure the release from prison of
convicted criminals whom they then enlisted into the para-military
forces that they were building in the factories of Hangzhou. In August
1973 two workers from Weng's silk mill who had been involved in a
vicious assault on six female fellow-workers in 1967 were scheduled to
be tried a third time. For five days in mid-August Weng's supporters
held demonstrations and sit-ins at the offices of the municipal party
and government authorities and disrupted court proceedings. When the
verdict was announced in October and one offender was released, having
already served out his sentence, he was compensated with six years' back
pay. In celebrations held at the silk complex he was hailed as a
returning hero.\footnote{Hua Minsheng, HZRB, June 22,1977; Yang Lingen,
  ibid; HZRB, June 23,1977; RMRB, April 15,1979, p.~1. Zhang Yongsheng
  also arranged for the release of two rebel leaders held in custody in
  Shaoxing and Sheng counties. Zhang Yongshengde zuizheng cailiao,
  pp.~97, 99.
  华民生，《杭州日报》，1977年6月22日；杨林根，同上；《杭州日报》，1977年6月23日；《人民日报》，1979年4月15日，第1版。张永生还安排释放了在绍兴县和嵊县被拘押的两名造反派头头。《张永生的罪证材料》，第97、99页。}
The reopening of old wounds further accentuated cleavages among the
mill's work-force.

为了进一步破坏浙江的政治环境，张永生和翁森鹤设法促成一些已定罪罪犯出狱，然后把他们编入自己正在杭州各工厂建立的准军事力量。1973年8月，翁森鹤所在丝绸厂的两名工人原定第三次受审；他们曾在1967年参与对六名女同事的恶性殴打。8月中旬，翁的支持者连续五天在市党政机关办公地点举行示威和静坐，并扰乱法庭程序。10月宣判时，其中一名罪犯因已服满刑期而获释，并获得六年补发工资作为补偿。在丝绸联合企业举行的庆祝活动中，他被当作归来的英雄来欢迎。\footnote{Hua
  Minsheng, HZRB, June 22,1977; Yang Lingen, ibid; HZRB, June 23,1977;
  RMRB, April 15,1979, p.~1. Zhang Yongsheng also arranged for the
  release of two rebel leaders held in custody in Shaoxing and Sheng
  counties. Zhang Yongshengde zuizheng cailiao, pp.~97, 99.
  华民生，《杭州日报》，1977年6月22日；杨林根，同上；《杭州日报》，1977年6月23日；《人民日报》，1979年4月15日，第1版。张永生还安排释放了在绍兴县和嵊县被拘押的两名造反派头头。《张永生的罪证材料》，第97、99页。}
旧伤口被重新撕开，进一步加剧了该厂职工之间的裂痕。

After the success of their revolt on Mt. Pingfeng, the rebel trio
visited factories in Hangzhou to gather support and measure their
opponents' strength. In December 1973 Weng visited the Red Flag Paper
Mill and called on workers not to work for the ``wrong political line''.
Those who did he denigrated as ``lambs''.\footnote{Visit to Red Flag
  Paper Mill, March 12, 1978. Established in 1919 the mill in 1978
  employed a staff of 2,200.
  1978年3月12日访问红旗造纸厂。该厂创办于1919年，1978年有职工2200人。}
On 17 December the three militant leaders appeared at the Hangzhou Iron
and Steel Mill. This factory was an obvious target for the rebels
because of the publicity which had surrounded it earlier in the year
(see chapter six). In his report to the workers Weng derided the mill's
production results, comparing them unfavorably with iron and steel
output in the United States and Japan. He argued that the correctness of
the political line was more important than output; if the line was wrong
the workers were just creating wealth for capitalist
restoration.\footnote{Weng Senhe bufen zuizheng cailiao, p.~67.
  《翁森鹤部分罪证材料》，第67页。} The presence of the three rebel
leaders apparently provoked strong resistance from workers of an
opposing faction who unceremoniously showed them the factory gates.
Nevertheless, Weng did recruit workers to his faction at this
factory.\footnote{Visit, May 27, 1979; HZRB, 27 August 1978; Keith
  Forster, ``A Story from Hangzhou'', Eastern Horizon, 18:12 (1979),
  pp.~42-4; PR, no. 40 (September 30,1977), p.~27.
  访问，1979年5月27日；《杭州日报》，1978年8月27日；Keith Forster, ``A
  Story from Hangzhou'', Eastern Horizon,
  18:12（1979年），第42-44页；PR，第40期（1977年9月30日），第27页。} The
Hangzhou Gearbox Works (杭州齿轮箱厂) in Xiaoshan made it clear that
Weng was not welcome by warning him that it could not guarantee his
personal safety, thereby ensuring that Weng did not appear at
all.\footnote{Visit to Hangzhou Gearbox Plant, July 17, 1977;
  ``抵制四人帮坚持学大庆'' (Resist the ``Gang of Four''; persist in
  learning from Daqing), HZRB, April 10, 1977, pp.~1-2.
  1977年7月17日访问杭州齿轮箱厂；《抵制四人帮坚持学大庆》，《杭州日报》，1977年4月10日，第1-2版。}
Weng and his colleagues also visited factories in other areas of the
province and called for the overthrow of Tan Qilong.\footnote{See ``The
  Case of Weng Sengho'', p.~3. 见 ``The Case of Weng Sengho''，第3页。}

在屏风山反叛取得成功之后，这三名造反派领导人走访杭州各工厂，以争取支持并衡量对手力量。1973年12月，翁森鹤访问红旗造纸厂，号召工人不要为``错误政治路线''干活。对于这样做的人，他贬称为``羔羊''。\footnote{Visit
  to Red Flag Paper Mill, March 12, 1978. Established in 1919 the mill
  in 1978 employed a staff of 2,200.
  1978年3月12日访问红旗造纸厂。该厂创办于1919年，1978年有职工2200人。}
12月17日，三名激进领导人出现在杭州钢铁厂。由于该厂在当年早些时候受到的宣传关注（见第六章），它显然是造反派的一个目标。翁在向工人作报告时嘲讽该厂的生产成绩，并把它同美国和日本的钢铁产量作不利比较。他认为，政治路线的正确性比产量更重要；如果路线错误，工人只是在为资本主义复辟创造财富。\footnote{Weng
  Senhe bufen zuizheng cailiao, p.~67. 《翁森鹤部分罪证材料》，第67页。}
三名造反派领导人的出现，显然激起了对立派别工人的强烈抵制，后者毫不客气地把他们赶出了工厂。尽管如此，翁仍在该厂招募了一些工人加入他的派系。\footnote{Visit,
  May 27, 1979; HZRB, 27 August 1978; Keith Forster, ``A Story from
  Hangzhou'', Eastern Horizon, 18:12 (1979), pp.~42-4; PR, no. 40
  (September 30,1977), p.~27.
  访问，1979年5月27日；《杭州日报》，1978年8月27日；Keith Forster, ``A
  Story from Hangzhou'', Eastern Horizon,
  18:12（1979年），第42-44页；PR，第40期（1977年9月30日），第27页。}
萧山的杭州齿轮箱厂则明确表示不欢迎翁森鹤，警告他厂方无法保证其人身安全，从而确保翁根本没有露面。\footnote{Visit
  to Hangzhou Gearbox Plant, July 17, 1977; ``抵制四人帮坚持学大庆''
  (Resist the ``Gang of Four''; persist in learning from Daqing), HZRB,
  April 10, 1977, pp.~1-2.
  1977年7月17日访问杭州齿轮箱厂；《抵制四人帮坚持学大庆》，《杭州日报》，1977年4月10日，第1-2版。}
翁及其同僚还走访了全省其他地区的工厂，并号召打倒谭启龙。\footnote{See
  ``The Case of Weng Sengho'', p.~3. 见 ``The Case of Weng
  Sengho''，第3页。}

To maintain the momentum of this aggressive approach, in January 1974
the HMWC convened a rally in Hangzhou's Exhibition Square. The square is
situated at the northern end of Yanan Road, one of the city's main
thoroughfares. The rally was held to celebrate the seventh anniversary
of Shanghai's ``January storm''.

为了保持这种进攻姿态的势头，1974年1月，杭州市工代会在杭州展览广场召开群众大会。该广场位于延安路北端，延安路是杭州的主要干道之一。这次大会是为了庆祝上海``一月风暴''七周年。

The former rebel leaders of United Headquarters mounted the rostrum once
again and were joined by sympathizers from the provincial and municipal
party committees. Following the speeches a parade, which turned into a
demonstration against the provincial leadership, took place. Beijing had
sent a telegram ordering the organizers to desist and the local
authorities had pleaded with them to call off the day's activities.
Zhang Yongsheng was not swayed by the approach and threatened to lead
his followers to sit down on the railway tracks at the Hangzhou station
if the provincial authorities tried to ban the rally.\footnote{Zhang
  Yongshengde zuizheng cailiao, p.~54. 《张永生的罪证材料》，第54页。}

联合总部的前造反派领导人再次登上主席台，省、市党委中的同情者也加入其中。讲话结束后举行了游行，游行演变成反对省领导层的示威。北京曾发来电报，命令组织者停止行动，地方当局也恳求他们取消当天的活动。张永生没有因此动摇，并威胁说，如果省当局试图禁止大会，他就带领追随者到杭州火车站的铁轨上静坐。\footnote{Zhang
  Yongshengde zuizheng cailiao, p.~54. 《张永生的罪证材料》，第54页。}
\#\# The Campaign against Lin Biao and Confucius: Opening Shots \#\#
反林彪和孔子运动：开场炮

Although the opening of the 10th National Congress had been accompanied
by a series of press reports designed to initiate a further
radicalization of national politics, the campaign against Lin Biao and
Confucius did not get underway until January 1974. It was Mao whose
intervention once again proved decisive. At the end of 1973 a meeting of
the Politburo was held to discuss the important issue of transferring
commanders of military regions. A series of direct exchanges and new
appointments were carried out involving eight of the eleven
commanders.\footnote{For a discussion of this reshuffle, see Hsu
  Ch'ien-sun, ``An Analysis of the Personnel Reshuffle in China's
  Military Regions'', CLG, 7: 3 (1974), pp.~7-10. One source has claimed
  that it was Deng Xiaoping who proposed the reshuffle to Mao. See Gao
  Gao and Yan Jiaqi, ``Wenhua dageming'', p.~529. 关于这次改组的讨论，见
  Hsu Ch'ien-sun, ``An Analysis of the Personnel Reshuffle in China's
  Military Regions'', CLG,
  7:3（1974年），第7-10页。有一项资料称，是邓小平向毛提出了这次改组。见高皋、严家其：《文化大革命》，第529页。}
As a result of the transfer the military commanders lost their
additional posts as provincial first party secretaries, thus loosening
the grip of regional military leaders over local party committees and
lowering the PLA's representation among the twenty-nine provincial-level
first secretaries. The decision had the additional benefit of removing
regional commanders, such as Nanjing's commander Xu Shiyou, who had
become entrenched in their bailiwicks. The order to carry out the
transfer was issued on 22 December.\footnote{Gao Gao and Yan Jiaqi,
  ``Wenhua dageming'', pp.~529-30; Hao Mengbi and Duan Haoran,
  Liushinian, p.~632.
  高皋、严家其：《文化大革命》，第529-530页；郝梦笔、段浩然：《六十年》，第632页。}
At the same Politburo meeting Mao also voiced his displeasure with the
work of the body as well as with that of the CC MAC. He stated vaguely
but ominously that the ``Politburo doesn't discuss politics and the
military commission doesn't discuss military or political affairs''
(政治局不议政, 军委不议军、不议政).\footnote{See n.~43. 见注43。} These
comments were directed at Zhou Enlai and Marshall Ye Jianying who were
in charge of the routine work of the respective bodies. The Chairman
warned his colleagues to prepare for both civil and external war, the
outbreak of which would clarify who was willing to fight, who would
collaborate with foreigners and who wished to become emperor. He also
proposed the addition of Deng Xiaoping to both the Politburo and the
Military Affairs Commission because, in Mao's words, he ``handles
affairs in a relatively decisive manner'' and ``is firm but gentle, like
a needle in cotton'' (绵里藏针).\footnote{See n.~43. 见注43。} The
transfer of regional commanders was to have important consequences for
Zhejiang. A circular of the CCP ZPC and ZPRC dated 17 December 1973
indicated that an important decision was about to be made. The circular
requested all PLA units in Zhejiang ``on their own initiative {[}sic{]}
to ask local Party committees for instructions and report to them on
their work and consciously accept the centralised leadership of the
party''.\footnote{ZJRB, December 18,1973, p.~1.
  《浙江日报》，1973年12月18日，第1版。} As a result of Mao's decision,
importantly. Tan Qilong, Tie Ying and Xia Qi lost the substantial
backing of Xu Shiyou, who was transferred to Guangzhou, and in his place
arrived Ding Sheng, a PLA leader more sympathetic to the radical cause.
It is unclear, however, what influence Ding exerted in Nanjing because
of the cloud which remained over his head due to his past association
with the Lin Biao group. It has been alleged that shortly after his
arrival in Nanjing, Beijing ordered him to stand aside and undergo
investigation for his activities in Guangzhou.\footnote{Zhao Wei, Zhao
  Ziyang zhuan, pp.~190, 197. 赵蔚：《赵紫阳传》，第190、197页。} Xu's
departure in effect gave the Nanjing region's First Political Commissar
Zhang Chunqiao, who in February 1974 was appointed Director of the PLA
General Political Department, much more influence in military affairs in
the region.\footnote{Dangshi tongxun, 71 (15 September 1983), p.~31.
  《党史通讯》，第71期（1983年9月15日），第31页。} This development was
undoubtedly a blessing for the rebel leaders in Zhejiang.

尽管第十次全国代表大会开幕时伴随着一系列新闻报道，意在开启全国政治的进一步激进化，但反林彪、反孔子运动直到1974年1月才真正展开。再次证明具有决定性作用的是毛泽东的干预。1973年底，政治局召开会议，讨论调动各大军区司令员这一重大问题。十一个司令员中有八人涉及一系列直接对调和新的任命。\footnote{For
  a discussion of this reshuffle, see Hsu Ch'ien-sun, ``An Analysis of
  the Personnel Reshuffle in China's Military Regions'', CLG, 7: 3
  (1974), pp.~7-10. One source has claimed that it was Deng Xiaoping who
  proposed the reshuffle to Mao. See Gao Gao and Yan Jiaqi, ``Wenhua
  dageming'', p.~529. 关于这次改组的讨论，见 Hsu Ch'ien-sun, ``An
  Analysis of the Personnel Reshuffle in China's Military Regions'',
  CLG,
  7:3（1974年），第7-10页。有一项资料称，是邓小平向毛提出了这次改组。见高皋、严家其：《文化大革命》，第529页。}
调动的结果是，这些军事指挥员失去了兼任省级党委第一书记的附加职务，从而削弱了地方军事领导人对地方党委的控制，也降低了人民解放军在二十九名省级第一书记中的代表比例。这一决定还有一个额外好处，即调离了像南京军区司令员许世友这样已经在自己辖区内根深蒂固的地区司令员。执行调动的命令于12月22日发布。\footnote{Gao
  Gao and Yan Jiaqi, ``Wenhua dageming'', pp.~529-30; Hao Mengbi and
  Duan Haoran, Liushinian, p.~632.
  高皋、严家其：《文化大革命》，第529-530页；郝梦笔、段浩然：《六十年》，第632页。}
在同一次政治局会议上，毛泽东也表达了对政治局本身工作以及中共中央军委工作的不满。他含糊但不祥地说，``政治局不议政，军委不议军、不议政''。\footnote{See
  n.~43. 见注43。}
这些话针对的是分别负责这两个机构日常工作的周恩来和叶剑英元帅。主席警告同事们要准备内战和外战，因为战争一旦爆发，就会弄清楚谁愿意打仗，谁会同外国人勾结，谁想当皇帝。他还提议把邓小平增补进政治局和军事委员会，因为用毛的话说，邓``办事比较果断''，而且``绵里藏针''。\footnote{See
  n.~43. 见注43。}
军区司令员的调动将对浙江产生重要后果。中共浙江省委和浙江省革命委员会1973年12月17日的一份通知表明，一项重要决定即将作出。该通知要求浙江境内所有解放军部队``主动
{[}sic{]} 向地方党委请示并汇报工作，自觉接受党的集中领导''。\footnote{ZJRB,
  December 18,1973, p.~1. 《浙江日报》，1973年12月18日，第1版。}
毛的决定造成了一个重要结果：谭启龙、铁瑛和夏琦失去了许世友的大力支持，许被调往广州；接替他的是丁盛，一位更同情激进派事业的解放军领导人。不过，丁盛在南京发挥了什么影响并不清楚，因为他过去与林彪集团的关系仍使他头上笼罩着阴影。据称，他到达南京后不久，北京即命令他靠边站，并接受对其在广州活动的调查。\footnote{Zhao
  Wei, Zhao Ziyang zhuan, pp.~190, 197.
  赵蔚：《赵紫阳传》，第190、197页。}
许世友的离开实际上使南京军区第一政委张春桥在该地区军事事务中拥有了大得多的影响力；张春桥于1974年2月被任命为解放军总政治部主任。\footnote{Dangshi
  tongxun, 71 (15 September 1983), p.~31.
  《党史通讯》，第71期（1983年9月15日），第31页。}
这一发展无疑对浙江的造反派领导人是一件幸事。

At a meeting with members of the MAC on 21 December Mao issued another
cryptic warning about the emergence of revisionism
(如果中国出了修正主义, 大家要注意啊!). Mao had issued such warnings
before in September 1965 and again at the end of 1966, with dire
consequences. The Chairman indirectly raised the question of
capitulationism by referring to the classic novel Water Margin in which
the bandits ``didn't oppose the emperor, only corrupt officials, and
later accepted an amnesty and returned to the service of the emperor''.
In the talk Mao also accepted responsibility for the wrongful dismissal
of senior military leaders during the years 1965-68, dismissals which,
he claimed had been recommended by Lin Biao.\footnote{Hao Mengbi and
  Duan Haoran, Liushinian, p.~633; Gao Gao and Yan Jiaqi, ``Wenhua
  dageming'', p.~530.
  郝梦笔、段浩然：《六十年》，第633页；高皋、严家其：《文化大革命》，第530页。}
Mao thereby paved the way for their rehabilitation (posthumous in the
case of He Long) and simultaneously moved to placate the PLA in the
aftermath of the transfer of regional commanders. Lin Biao's military
reputation and his place in the pantheon of communist military heroes
was severely undermined by articles published later in 1974, presumably
a move designed to restore something approaching accuracy to the
historical record and to reinstate the contributions of military leaders
who had been incorrectly downgraded.\footnote{See Chan Shih-pu, Great
  Victory for The Military Line of Chairman Mao Tsetung - A Criticism of
  Lin Piao's Bourgeois Military Line in the Liaohsi-Shenyang and
  PeipingTientsin Campaigns (Beijing: Foreign Languages Press, 1976); Li
  Chung-ta ``Why Did Mao Tse-tung Start the Struggle to Criticize Lin
  and Confucius?'', I\&S, 11: 5 (1975), pp.~85-6; Chien T'ieh, ``The
  Maoist Criticism of Lin Piao's Military Line'', I\&S, 11: 5 (1975),
  pp.~31-41. For articles published in Zhejiang criticizing Lin's
  military reputation, see ZJRB, September 9,1974, p.~2, September
  12,1974, p.~3 and September 14,1974, p.~1. 见 Chan Shih-pu, Great
  Victory for The Military Line of Chairman Mao Tsetung - A Criticism of
  Lin Piao's Bourgeois Military Line in the Liaohsi-Shenyang and
  PeipingTientsin Campaigns（北京：外文出版社，1976年）；Li Chung-ta,
  ``Why Did Mao Tse-tung Start the Struggle to Criticize Lin and
  Confucius?'', I\&S, 11:5（1975年），第85-86页；Chien T'ieh, ``The
  Maoist Criticism of Lin Piao's Military Line'', I\&S,
  11:5（1975年），第31-41页。关于浙江发表的批判林彪军事声望的文章，见《浙江日报》，1974年9月9日，第2版；1974年9月12日，第3版；1974年9月14日，第1版。}
A central document of 22 December advised the party of Deng's addition
to the Politburo and the MAC.\footnote{A recent Chinese article claims
  that Mao appointed Deng PLA Chief of Staff on December 12,1973. See Ai
  Feng, Luo Zisuo and Meng Xiaoyun,
  ``落日青山夕照明--叶剑英同志晚年生活片断''(Verdant hills at sunset
  greet the eye on all sides-fragments of Comrade Ye Jianying's latter
  years), RMRB, 26 October 1986, p.~4. While Mao may have proposed
  Deng's appointment to this post in December 1973 it appears that he
  did not actually assume it until January 1975. See Deng Xiaoping
  Wenxuan, p.~1 fn. Further evidence suggests that while the CC document
  announced Deng's promotion to the Politburo and membership of the
  Military Affairs Commission, no mention was made of the position of
  Chief of Staff. See Jiangsusheng dashiji, p.~325.
  近来一篇中文文章称，毛在1973年12月12日任命邓为解放军总参谋长。见艾丰、罗自所、孟晓云：《落日青山夕照明--叶剑英同志晚年生活片断》（夕阳下青山环照入目--叶剑英同志晚年生活片断），《人民日报》，1986年10月26日，第4版。毛或许在1973年12月提出让邓担任这一职务，但看来邓实际上直到1975年1月才就任。见《邓小平文选》，第1页脚注。进一步证据表明，中央文件虽然宣布邓晋升政治局并成为军事委员会成员，但没有提到总参谋长职位。见《江苏省大事记》，第325页。}
Premier Zhou Enlai was thus provided with a senior experienced
administrator to assist him in running party and government affairs. A
routine medical examination in May 1972 had discovered that Zhou had
contracted cancer and in June 1974 he was admitted to hospital for a
series of operations and treatment.\footnote{Gao Wenqian,
  ``在最后的日子里-记病重住院期间的周恩来同志'' (Final Days - Recalling
  Comrade Zhou Enlai seriously ill in hospital) ZJRB, January 5,1986
  (reprinted from 中华英烈 - Chinese Heroes and Martyrs).
  高文谦：《在最后的日子里-记病重住院期间的周恩来同志》（最后的日子--回忆重病住院期间的周恩来同志），《浙江日报》，1986年1月5日（转载自《中华英烈》）。}
Deng's promotion guaranteed that the succession to Mao and Zhou would
not fall automatically into the hands of the Cultural Revolution
radicals.

在12月21日同军事委员会成员的一次会议上，毛泽东又就修正主义的出现发出了一次含义隐晦的警告：``如果中国出了修正主义，大家要注意啊！''毛以前曾在1965年9月和1966年底发出过这类警告，并造成了严重后果。主席通过提到古典小说《水浒传》间接提出了投降主义问题，其中的强盗``只反贪官，不反皇帝，后来接受招安，回到皇帝手下服务''。在这次谈话中，毛还承认自己对1965至1968年间错误撤换高级军事领导人负有责任；他声称这些撤换是林彪建议的。\footnote{Hao
  Mengbi and Duan Haoran, Liushinian, p.~633; Gao Gao and Yan Jiaqi,
  ``Wenhua dageming'', p.~530.
  郝梦笔、段浩然：《六十年》，第633页；高皋、严家其：《文化大革命》，第530页。}
毛由此为他们的平反铺平了道路（贺龙的情况是死后平反），同时也在调动各大军区司令员之后安抚解放军。1974年晚些时候发表的一些文章严重削弱了林彪的军事声誉以及他在共产党军事英雄谱系中的地位；这大概是为了使历史记录恢复到接近准确的状态，并重新确认那些被错误贬低的军事领导人的贡献。\footnote{See
  Chan Shih-pu, Great Victory for The Military Line of Chairman Mao
  Tsetung - A Criticism of Lin Piao's Bourgeois Military Line in the
  Liaohsi-Shenyang and PeipingTientsin Campaigns (Beijing: Foreign
  Languages Press, 1976); Li Chung-ta ``Why Did Mao Tse-tung Start the
  Struggle to Criticize Lin and Confucius?'', I\&S, 11: 5 (1975),
  pp.~85-6; Chien T'ieh, ``The Maoist Criticism of Lin Piao's Military
  Line'', I\&S, 11: 5 (1975), pp.~31-41. For articles published in
  Zhejiang criticizing Lin's military reputation, see ZJRB, September
  9,1974, p.~2, September 12,1974, p.~3 and September 14,1974, p.~1. 见
  Chan Shih-pu, Great Victory for The Military Line of Chairman Mao
  Tsetung - A Criticism of Lin Piao's Bourgeois Military Line in the
  Liaohsi-Shenyang and PeipingTientsin
  Campaigns（北京：外文出版社，1976年）；Li Chung-ta, ``Why Did Mao
  Tse-tung Start the Struggle to Criticize Lin and Confucius?'', I\&S,
  11:5（1975年），第85-86页；Chien T'ieh, ``The Maoist Criticism of Lin
  Piao's Military Line'', I\&S,
  11:5（1975年），第31-41页。关于浙江发表的批判林彪军事声望的文章，见《浙江日报》，1974年9月9日，第2版；1974年9月12日，第3版；1974年9月14日，第1版。}
12月22日的一份中央文件通知全党，邓小平已被增补进政治局和军事委员会。\footnote{A
  recent Chinese article claims that Mao appointed Deng PLA Chief of
  Staff on December 12,1973. See Ai Feng, Luo Zisuo and Meng Xiaoyun,
  ``落日青山夕照明--叶剑英同志晚年生活片断''(Verdant hills at sunset
  greet the eye on all sides-fragments of Comrade Ye Jianying's latter
  years), RMRB, 26 October 1986, p.~4. While Mao may have proposed
  Deng's appointment to this post in December 1973 it appears that he
  did not actually assume it until January 1975. See Deng Xiaoping
  Wenxuan, p.~1 fn. Further evidence suggests that while the CC document
  announced Deng's promotion to the Politburo and membership of the
  Military Affairs Commission, no mention was made of the position of
  Chief of Staff. See Jiangsusheng dashiji, p.~325.
  近来一篇中文文章称，毛在1973年12月12日任命邓为解放军总参谋长。见艾丰、罗自所、孟晓云：《落日青山夕照明--叶剑英同志晚年生活片断》（夕阳下青山环照入目--叶剑英同志晚年生活片断），《人民日报》，1986年10月26日，第4版。毛或许在1973年12月提出让邓担任这一职务，但看来邓实际上直到1975年1月才就任。见《邓小平文选》，第1页脚注。进一步证据表明，中央文件虽然宣布邓晋升政治局并成为军事委员会成员，但没有提到总参谋长职位。见《江苏省大事记》，第325页。}
这样，周恩来总理便获得了一位资深而有经验的行政干部，协助他处理党和政府事务。1972年5月的一次例行体检发现周恩来患了癌症，1974年6月他住院接受一系列手术和治疗。\footnote{Gao
  Wenqian, ``在最后的日子里-记病重住院期间的周恩来同志'' (Final Days -
  Recalling Comrade Zhou Enlai seriously ill in hospital) ZJRB, January
  5,1986 (reprinted from 中华英烈 - Chinese Heroes and Martyrs).
  高文谦：《在最后的日子里-记病重住院期间的周恩来同志》（最后的日子--回忆重病住院期间的周恩来同志），《浙江日报》，1986年1月5日（转载自《中华英烈》）。}
邓的晋升保证了毛和周的继承问题不会自动落入文化大革命激进派手中。

Deng Xiaoping's principal protagonist in the upper echelons of the CCP
was the much younger and inexperienced Wang Hongwen. The anti-Lin Biao
anti-Confucius campaign (批林批孔), which commenced in 1974, was to
provide Wang Hongwen with his first test as a senior national leader. It
was Wang who conveyed Mao's remarks concerning the possible emergence of
revisionism in a speech to senior cadres on 29 December.\footnote{Wang
  Ruoshui, ``The greatest lesson of the Cultural Revolution'', p.~93. In
  his speech of February 1979 recalls the shock that he and other senior
  officials experienced on hearing these unexpected words of the
  Chairman's and how they struggled to fathom the implication for the
  ongoing inner-party struggle. 王若水：``The greatest lesson of the
  Cultural
  Revolution''，第93页。在1979年2月的讲话中，他回忆了自己和其他高级干部听到主席这些出人意料的话时所受到的震动，以及他们如何努力揣摩这些话对正在进行的党内斗争意味着什么。}
Then, in a speech to the central study class on 14 January 1974 Wang
vigorously defended the Cultural Revolution.\footnote{Wang Hongwen,
  ``Report to the Central Study Class'', Classified Chinese Communist
  Documents, pp.~64-82. Wang made a series of speeches to study classes,
  composed of provincial and lower-level cadre, between 1973 and 1974.
  See Fenwick, ``The Gang of Four'', pp.~221-25. January 14 was the
  seventh anniversary of Shanghai's 1967 power seizure. For Wang's role
  in that event see Ye Yonglie, Wang Hongwen xingshuailu, ch.~7.
  王洪文：``Report to the Central Study Class'', Classified Chinese
  Communist
  Documents，第64-82页。1973年至1974年间，王向由省级和较低层级干部组成的学习班作了一系列讲话。见
  Fenwick, ``The Gang of
  Four''，第221-225页。1月14日是1967年上海夺权七周年纪念日。关于王在该事件中的角色，见叶永烈：《王洪文兴衰录》，第7章。}
He admitted that controversy still surrounded it. Some people believed
the Cultural Revolution had been ``necessary and timely'' while others
viewed it like a ``dark night'', a ``ravaging flood and a savage
beast'', ``thunder in the clear sky of an early morning'', and even ``a
great misunderstanding, very reactionary in nature''. Those who held the
latter view, stated Wang, opposed progress, desired restoration and
planned to take revenge for the sufferings they had experienced. Others,
while supporting the Cultural Revolution in principle, opposed the way
in which it had been carried out. They frowned on the ``four greats'' --
speaking out freely, airing views fully, holding great debates, and
writing wall posters. Wang argued that to oppose the Cultural Revolution
was to attempt to restore capitalism and practise revisionism. He
advocated the slogan of ``going against the tide'' to rekindle the
spirit of rebellion.

邓小平在中共高层的主要对手，是年轻得多且缺乏经验的王洪文。1974年开始的反林彪、反孔子运动（批林批孔），将成为王洪文作为高级国家领导人的第一次考验。正是王洪文在12月29日对高级干部的讲话中传达了毛关于可能出现修正主义的言论。\footnote{Wang
  Ruoshui, ``The greatest lesson of the Cultural Revolution'', p.~93. In
  his speech of February 1979 recalls the shock that he and other senior
  officials experienced on hearing these unexpected words of the
  Chairman's and how they struggled to fathom the implication for the
  ongoing inner-party struggle. 王若水：``The greatest lesson of the
  Cultural
  Revolution''，第93页。在1979年2月的讲话中，他回忆了自己和其他高级干部听到主席这些出人意料的话时所受到的震动，以及他们如何努力揣摩这些话对正在进行的党内斗争意味着什么。}
随后，在1974年1月14日对中央读书班的讲话中，王洪文极力为文化大革命辩护。\footnote{Wang
  Hongwen, ``Report to the Central Study Class'', Classified Chinese
  Communist Documents, pp.~64-82. Wang made a series of speeches to
  study classes, composed of provincial and lower-level cadre, between
  1973 and 1974. See Fenwick, ``The Gang of Four'', pp.~221-25. January
  14 was the seventh anniversary of Shanghai's 1967 power seizure. For
  Wang's role in that event see Ye Yonglie, Wang Hongwen xingshuailu,
  ch.~7. 王洪文：``Report to the Central Study Class'', Classified
  Chinese Communist
  Documents，第64-82页。1973年至1974年间，王向由省级和较低层级干部组成的学习班作了一系列讲话。见
  Fenwick, ``The Gang of
  Four''，第221-225页。1月14日是1967年上海夺权七周年纪念日。关于王在该事件中的角色，见叶永烈：《王洪文兴衰录》，第7章。}
他承认围绕文化大革命仍然存在争议。有些人认为文化大革命是``必要的、及时的''，而另一些人则把它看作``黑夜''、``洪水猛兽''、``晴天霹雳''，甚至是``性质很反动的一场大误会''。王说，持后一种观点的人反对进步，渴望复辟，并计划为他们所遭受的痛苦进行报复。还有一些人原则上支持文化大革命，却反对其开展方式。他们不赞成``四大''------大鸣、大放、大辩论、大字报。王洪文论证说，反对文化大革命就是企图复辟资本主义、实行修正主义。他主张用``反潮流''的口号重新点燃造反精神。

In his speech Wang touched on two other important themes -- discipline
in the PLA and the status of young cadres. He claimed that obedience to
orders was conditional on their conformity to Marxism-Leninism Mao
Zedong Thought. Soldiers should rebel against orders contradicting these
principles. But as Wang reduced the essence of Marxism to the maxim ``to
rebel is justified'', the permissible scope for disobedience seemed
alarmingly broad. Concerning the question of young cadres who had been
promoted during the Cultural Revolution, Wang defended them and
advocated the appointment of young officers to senior military commands.
Otherwise, he opined, the ``traitors, enemy agents and
capitalist-roaders will return and the new-born things of the Cultural
Revolution will be abolished''. According to Wang rumors had been
circulating, concerning veteran and younger cadres which went:''old
marshalls must return to their posts; little soldiers must go back to
their barracks''. Mao had changed this couplet to read: ``old marshalls
return to the line; little soldiers are promoted''.

王洪文在讲话中还触及另外两个重要主题：解放军纪律和年轻干部的地位。他声称，服从命令的前提是命令符合马克思列宁主义毛泽东思想。士兵应当反抗违背这些原则的命令。但由于王把马克思主义的本质简化为``造反有理''这一格言，允许不服从的范围似乎宽得令人警觉。关于文化大革命期间被提拔的年轻干部问题，王为他们辩护，并主张任命年轻军官担任高级军事指挥职务。否则，他认为，``叛徒、特务、走资派就会回来，文化大革命的新生事物就会被取消''。据王说，当时流传着关于老干部和年轻干部的谣言：``老帅要归位，小兵要回营。''毛把这副对子改成：``老帅归队，小兵上台。''

Wang's resolute defence of the Cultural Revolution heralded the
commencement of the campaign in the same month.\footnote{For
  contemporary accounts of the campaign, see Parris H. Chang, ``The
  Anti-Lin Piao and Confucius Campaign: Its Meaning and Purposes'', AS,
  14:10 (1974), pp.~871-86; the articles collected in CLG, 7:4 (1974-75)
  and those carried by I\&S, of February, June and July of 1974 and May
  and September of 1975. 关于这场运动的同时代记述，见 Parris H. Chang,
  ``The Anti-Lin Piao and Confucius Campaign: Its Meaning and
  Purposes'', AS, 14:10（1974年），第871-886页；CLG,
  7:4（1974-1975年）所收文章，以及 I\&S
  于1974年2月、6月、7月和1975年5月、9月刊载的文章。} Two CC documents
which were distributed in mid-January contained a list of recommended
references for study and tried to link the thoughts of Lin Biao with the
classical analects of Confucius.\footnote{CCP CC, zhongfa (1974), No.~1,
  January 18, 1974, Classified Chinese Communist Documents, pp.~755-87;
  CCP CC, zhongfa (1972), No.~2, January 22,1974, ibid., pp.~788-89.
  中共中央，中发（1974）第1号，1974年1月18日，Classified Chinese
  Communist
  Documents，第755-787页；中共中央，中发（1972）第2号，1974年1月22日，同上，第788-789页。}
The document of 18 January was published six days after a written
request to the Chairman by Wang Hongwen and Jiang Qing.\footnote{Hao
  Mengbi and Duan Haoran, Liushinian, p.~633; Liu Suinian and Wu Qungan,
  ``Wenhua dageming'' shiqide guomin jingji, pp.~165-66.
  郝梦笔、段浩然：《六十年》，第633页；刘岁年、吴群敢：《``文化大革命''时期的国民经济》，第165-166页。}
The first mass rallies occurred in Beijing on 24 and 25 January and
involved troops of the Beijing units and bureaucrats in central
departments. Jiang Qing, Yao Wenyuan and two close associates of Jiang's
at Qinghua University, Chi Qun\footnote{For some first-hand observations
  of Chi Qun by a faculty member at Beijing University where Chi,
  chairman of the revolutionary committee at Qinghua University, was
  responsible for educational reform, see Yue Daiyun and Carolyn
  Wakeman, To the Storm, pp.~303,310, 315, 323,333,351. During the
  Cultural Revolution Chi Qun served as Deputy-director of the
  propaganda section of PLA Unit 8341 which was one of the models in the
  ``purification of class ranks'' campaign. Later, he and Xie Jingyi led
  a soldier-worker propaganda team into Qinghua University. See Gao Gao
  and Yan Jiaqi, ``Wenhua dageming'', pp.~298, 554.
  关于迟群的一些第一手观察，出自一名北京大学教师；当时迟群任清华大学革命委员会主任，负责教育改革。见岳黛云、Carolyn
  Wakeman, To the
  Storm，第303、310、315、323、333、351页。文化大革命期间，迟群曾任解放军8341部队宣传科副科长，该部队是``清理阶级队伍''运动的典型之一。后来，他和谢静宜率领一支军工宣传队进入清华大学。见高皋、严家其：《文化大革命》，第298、554页。}
and Xie Jingyi, delivered speeches. Premier Zhou Enlai attended and was
forced to make a self-criticism.\footnote{Zhonggong dangshi 170ti wenda,
  p.~313; Wang Ruoshui, ``The greatest lesson,'' p.~93; CCP CC, zhongfa
  (1976), no. 24, December 10,1976, Zhonggong nianbao, 1977 (Taibei:
  Center for the Study of Chinese Communism, 1977), 5:16.
  《中共党史170题问答》，第313页；王若水：《最大的教训》，第93页；中共中央，中发（1976）第24号，1976年12月10日，载《中共年报》1977年（台北：中国共产主义研究中心，1977年），5:16。}

王洪文对文化大革命的坚决辩护，预示着该运动于同月开始。\footnote{For
  contemporary accounts of the campaign, see Parris H. Chang, ``The
  Anti-Lin Piao and Confucius Campaign: Its Meaning and Purposes'', AS,
  14:10 (1974), pp.~871-86; the articles collected in CLG, 7:4 (1974-75)
  and those carried by I\&S, of February, June and July of 1974 and May
  and September of 1975. 关于这场运动的同时代记述，见 Parris H. Chang,
  ``The Anti-Lin Piao and Confucius Campaign: Its Meaning and
  Purposes'', AS, 14:10（1974年），第871-886页；CLG,
  7:4（1974-1975年）所收文章，以及 I\&S
  于1974年2月、6月、7月和1975年5月、9月刊载的文章。}
1月中旬下发的两份中央文件列出了建议学习的参考资料，并试图把林彪的思想同孔子的古典《论语》联系起来。\footnote{CCP
  CC, zhongfa (1974), No.~1, January 18, 1974, Classified Chinese
  Communist Documents, pp.~755-87; CCP CC, zhongfa (1972), No.~2,
  January 22,1974, ibid., pp.~788-89.
  中共中央，中发（1974）第1号，1974年1月18日，Classified Chinese
  Communist
  Documents，第755-787页；中共中央，中发（1972）第2号，1974年1月22日，同上，第788-789页。}
1月18日的文件是在王洪文和江青向主席提出书面请求六天后发布的。\footnote{Hao
  Mengbi and Duan Haoran, Liushinian, p.~633; Liu Suinian and Wu Qungan,
  ``Wenhua dageming'' shiqide guomin jingji, pp.~165-66.
  郝梦笔、段浩然：《六十年》，第633页；刘岁年、吴群敢：《``文化大革命''时期的国民经济》，第165-166页。}
最早的群众集会于1月24日和25日在北京举行，参加者包括北京部队官兵和中央各部门的官僚人员。江青、姚文元以及江青在清华大学的两名亲信迟群\footnote{For
  some first-hand observations of Chi Qun by a faculty member at Beijing
  University where Chi, chairman of the revolutionary committee at
  Qinghua University, was responsible for educational reform, see Yue
  Daiyun and Carolyn Wakeman, To the Storm, pp.~303,310, 315,
  323,333,351. During the Cultural Revolution Chi Qun served as
  Deputy-director of the propaganda section of PLA Unit 8341 which was
  one of the models in the ``purification of class ranks'' campaign.
  Later, he and Xie Jingyi led a soldier-worker propaganda team into
  Qinghua University. See Gao Gao and Yan Jiaqi, ``Wenhua dageming'',
  pp.~298, 554.
  关于迟群的一些第一手观察，出自一名北京大学教师；当时迟群任清华大学革命委员会主任，负责教育改革。见岳黛云、Carolyn
  Wakeman, To the
  Storm，第303、310、315、323、333、351页。文化大革命期间，迟群曾任解放军8341部队宣传科副科长，该部队是``清理阶级队伍''运动的典型之一。后来，他和谢静宜率领一支军工宣传队进入清华大学。见高皋、严家其：《文化大革命》，第298、554页。}和谢静宜发表了讲话。周恩来总理出席，并被迫作了自我批评。\footnote{Zhonggong
  dangshi 170ti wenda, p.~313; Wang Ruoshui, ``The greatest lesson,''
  p.~93; CCP CC, zhongfa (1976), no. 24, December 10,1976, Zhonggong
  nianbao, 1977 (Taibei: Center for the Study of Chinese Communism,
  1977), 5:16.
  《中共党史170题问答》，第313页；王若水：《最大的教训》，第93页；中共中央，中发（1976）第24号，1976年12月10日，载《中共年报》1977年（台北：中国共产主义研究中心，1977年），5:16。}

Jiang Qing further involved herself in the campaign by establishing a
model unit in the PLA. Jiang Qing's first move, and one that was to have
serious consequences for the leadership in Zhejiang, was to send Chi Qun
and Xie Jingyi as her personal emissaries to the anti-chemical warfare
company (防化连) of the 20th Army. Stationed at Huzhou, a city situated
north of Hangzhou,\footnote{Not in Henan province as reported by Gao Gao
  and Yan Jiaqi in ``Wenhua dageming'', p.~495.
  并非如高皋和严家其在《文化大革命》第495页所说发生在河南省。} the
company came under the command of the 20th Army (Unit 6409). Chi and Xie
delivered a letter from Jiang as well as 200 copies of the booklet
compiled at Beijing and Qinghua Universities and used as prescribed
reading in the campaign.

江青通过在解放军中树立一个典型单位，进一步介入了这场运动。江青的第一步，也是将对浙江领导层产生严重后果的一步，是派迟群和谢静宜作为她的私人使者，前往第二十军防化连。该连驻扎在杭州以北的湖州市，\footnote{Not
  in Henan province as reported by Gao Gao and Yan Jiaqi in ``Wenhua
  dageming'', p.~495.
  并非如高皋和严家其在《文化大革命》第495页所说发生在河南省。}
隶属第二十军（6409部队）。迟、谢带去了江青的一封信，以及由北京大学和清华大学编写、作为运动指定读物的小册子二百册。

The Beijing envoys arrived in Huzhou on 15 January and returned to the
capital a week later to report back to Jiang. Jiang Qing allegedly saw
off her envoys with the following words: ``I am fond of firing canon
shots. I am a gunner; I manufacture canons and canon balls and I have
some artillery batteries. Now I am shooting you out as canon balls; you
are going to open fire''.\footnote{CCP, CC, zhongfa (1977), No.~37,
  I\&S, 14: 8 (1978), pp.~92-3; Hangzhou Garrison criticism group,
  ``Iron-clad proof of attacking the party and destabilizing the army'',
  HZRB, March 10,1977.
  中共中央，中发（1977）第37号，I\&S，14:8（1978），第92-93页；杭州警备区批判组：《攻击党、乱军的铁证》，《杭州日报》，1977年3月10日。}
Her emissaries brought back letters from the anti-chemical warfare
company and the party committees of the 20th Army and the Nanjing
Military Region. The letters expressed thanks to Jiang for her concern
and pledged their determination to live up to her expectations and were
released as a CC document on 25 January 1974. The anti-chemical warfare
company subsequently received extensive coverage in the national and the
Zhejiang press and certain of its members became prominent speakers at
meetings and rallies which were held in the province.\footnote{See HZRB,
  February 2, 4,10, 1974; Xinhua, February 5, 1974, SCMP, 5553 (February
  14, 1974), pp.128-31; ZPS, February 7,1974, URS, 76:1 (July 2,1974),
  p.~3; RMRB, February 9,1974, p.~2 (wholly devoted to articles by
  officers and men of the company); ZPS, February 17, 1974,
  SWB/FE/4539/BII/27-8; ZJRB, March 2, 1974, p.~1; ZPS, March 8, 1974,
  FBIS/CHI, 50 (1974), C1-2; ZPS, May 24,1974, SWB/FE/4612/BII/5-6; ZPS,
  July 8, 1974, SWB/FE/4648/BII/10-11; Guangming Ribao, July 20,1974,
  SCMP, 5679 (August 20, 1974), pp.57-61; ZPS, September 14, 1974,
  SWB/FE/4707/BII/9; ZPS, January 27, 1975, FBIS/CHI, 21 (1975), G1-2.
  Zhang Yongsheng visited the company to see for himself. See Zhang
  Yongshengde zuizheng cailiao, p.~19.
  见《杭州日报》1974年2月2日、4日、10日；新华社，1974年2月5日，SCMP，5553（1974年2月14日），第128-131页；ZPS，1974年2月7日，URS，76:1（1974年7月2日），第3页；《人民日报》，1974年2月9日，第2版（整版刊登该连官兵文章）；ZPS，1974年2月17日，SWB/FE/4539/BII/27-28；《浙江日报》，1974年3月2日，第1版；ZPS，1974年3月8日，FBIS/CHI，50（1974），C1-2；ZPS，1974年5月24日，SWB/FE/4612/BII/5-6；ZPS，1974年7月8日，SWB/FE/4648/BII/10-11；《光明日报》，1974年7月20日，SCMP，5679（1974年8月20日），第57-61页；ZPS，1974年9月14日，SWB/FE/4707/BII/9；ZPS，1975年1月27日，FBIS/CHI，21（1975），G1-2。张永生曾到该连亲自察看。见《张永生的罪证材料》，第19页。}
The model was therefore one to which the leaders of Zhejiang had to pay
more than lip-service, whether they supported it or not. The fact that
the company also became a model for the ``crash-admission'' of party
members and ``crash-promotion'' of cadres did not help endear it to the
provincial party leadership.\footnote{See Weng Senhe bufen zuizheng
  cailiao, p.~38. 见《翁森鹤部分罪证材料》，第38页。} Nevertheless, it
was awkward having a ``model'' in one's own backyard. The further away a
model was, the greater the likelihood that the center would accept the
plea for exemption or adaptation on the grounds that it did not conform
to local conditions.

北京来的使者于1月15日抵达湖州，一周后返回首都向江青汇报。据称，江青送别她的使者时说了这样一番话：``我喜欢放炮。我是炮手；我制造大炮和炮弹，我还有一些炮兵阵地。现在我把你们当炮弹打出去；你们要去开火。''\footnote{CCP,
  CC, zhongfa (1977), No.~37, I\&S, 14: 8 (1978), pp.~92-3; Hangzhou
  Garrison criticism group, ``Iron-clad proof of attacking the party and
  destabilizing the army'', HZRB, March 10,1977.
  中共中央，中发（1977）第37号，I\&S，14:8（1978），第92-93页；杭州警备区批判组：《攻击党、乱军的铁证》，《杭州日报》，1977年3月10日。}
她的使者带回了防化连、第二十军党委和南京军区党委的来信。这些信感谢江青的关怀，并表示决心不辜负她的期望；它们于1974年1月25日作为中央文件下发。此后，防化连在全国和浙江报刊上获得广泛报道，连里某些成员成为在浙江省举行的会议和集会上的重要发言人。\footnote{See
  HZRB, February 2, 4,10, 1974; Xinhua, February 5, 1974, SCMP, 5553
  (February 14, 1974), pp.128-31; ZPS, February 7,1974, URS, 76:1 (July
  2,1974), p.~3; RMRB, February 9,1974, p.~2 (wholly devoted to articles
  by officers and men of the company); ZPS, February 17, 1974,
  SWB/FE/4539/BII/27-8; ZJRB, March 2, 1974, p.~1; ZPS, March 8, 1974,
  FBIS/CHI, 50 (1974), C1-2; ZPS, May 24,1974, SWB/FE/4612/BII/5-6; ZPS,
  July 8, 1974, SWB/FE/4648/BII/10-11; Guangming Ribao, July 20,1974,
  SCMP, 5679 (August 20, 1974), pp.57-61; ZPS, September 14, 1974,
  SWB/FE/4707/BII/9; ZPS, January 27, 1975, FBIS/CHI, 21 (1975), G1-2.
  Zhang Yongsheng visited the company to see for himself. See Zhang
  Yongshengde zuizheng cailiao, p.~19.
  见《杭州日报》1974年2月2日、4日、10日；新华社，1974年2月5日，SCMP，5553（1974年2月14日），第128-131页；ZPS，1974年2月7日，URS，76:1（1974年7月2日），第3页；《人民日报》，1974年2月9日，第2版（整版刊登该连官兵文章）；ZPS，1974年2月17日，SWB/FE/4539/BII/27-28；《浙江日报》，1974年3月2日，第1版；ZPS，1974年3月8日，FBIS/CHI，50（1974），C1-2；ZPS，1974年5月24日，SWB/FE/4612/BII/5-6；ZPS，1974年7月8日，SWB/FE/4648/BII/10-11；《光明日报》，1974年7月20日，SCMP，5679（1974年8月20日），第57-61页；ZPS，1974年9月14日，SWB/FE/4707/BII/9；ZPS，1975年1月27日，FBIS/CHI，21（1975），G1-2。张永生曾到该连亲自察看。见《张永生的罪证材料》，第19页。}
因此，无论浙江领导人是否支持这个典型，他们都不能只是口头应付。该连还成为``突击入党''和``突击提干''的典型，这一事实并没有使省级党委领导层更喜欢它。\footnote{See
  Weng Senhe bufen zuizheng cailiao, p.~38.
  见《翁森鹤部分罪证材料》，第38页。}
不过，在自己的后院有一个``典型''总归令人尴尬。一个典型离得越远，中央越可能接受以其不符合地方条件为由而提出的豁免或变通请求。

As the leadership in Zhejiang prepared to launch the campaign in the
province it faced daunting prospects. The ``mountain top'' faction had
already responded to the calls by radical central leaders to ``go
against the tide'', promote new cadres and defy politically unacceptable
orders. Additionally, workers were enjoined, in a big-character poster
originating from Shanghai dockers and published in the People's Daily on
1 February 1974, to become masters of their enterprises and not ``slaves
to tonnage''. Enterprise directors were subsequently hesitant to
exercise leadership in the workplace, an abrogation of responsibility
that would have dire consequences in Zhejiang.\footnote{Liu Suinian and
  Wu Qungan, ``Wenhua dageming'' shiqide guomin jingji, p.~166.
  刘遂年、吴群敢：《``文化大革命''时期的国民经济》，第166页。} Wang
Hongwen had already personally intervened to spur on the rebels and
Jiang Qing had established a model military unit in the province for the
campaign. The ``mountain top'' faction possessed a solid organizational
base in the Hangzhou workers' congress and some influence in the
provincial trade union council. They had highly-placed allies in the
Hangzhou party committee and with the assistance of sympathetic
provincial leading cadres could apply pressure to and isolate Tan and
Tie in the CCP ZPC. In order to confront the provincial authorities and
especially the military more effectively, the rebels required an armed
force and this they were about to create in the form of the Hangzhou
Municipal Militia Command Headquarters.

当浙江领导层准备在省内发动这场运动时，面临的是令人畏惧的前景。``山上派''已经响应中央激进派领导人关于``反潮流''、提拔新干部、抗拒政治上不可接受的命令的号召。此外，1974年2月1日《人民日报》刊登了一张源自上海码头工人的大字报，号召工人成为企业的主人，而不是``吨位的奴隶''。此后，企业负责人在工作场所行使领导权时变得犹豫，这种放弃责任的行为将在浙江造成严重后果。\footnote{Liu
  Suinian and Wu Qungan, ``Wenhua dageming'' shiqide guomin jingji,
  p.~166. 刘遂年、吴群敢：《``文化大革命''时期的国民经济》，第166页。}
王洪文已经亲自介入，推动造反派；江青也已在该省为这场运动树立了一个军队典型单位。``山上派''在杭州市工人代表大会中拥有坚实的组织基础，并在省总工会中有一定影响。他们在杭州市党委中有地位很高的盟友，并且在同情他们的省级领导干部协助下，能够向中共浙江省委中的谭和铁施加压力并孤立他们。为了更有效地对抗省级当局，尤其是军方，造反派需要一支武装力量，而他们即将以杭州市民兵指挥部的形式创造出这支力量。

\section{The Anti-Lin Anti-Confucius Campaign in Zhejiang,
1974}\label{the-anti-lin-anti-confucius-campaign-in-zhejiang-1974}

\section{1974年浙江的反林反孔运动}\label{ux5e74ux6d59ux6c5fux7684ux53cdux6797ux53cdux5b54ux8fd0ux52a8}

The following analysis of the anti-Lin anti-Confucius campaign in
Zhejiang examines both its public and less visible aspects, an approach
that in places interrupts the chronological flow for the sake of
thematic unity. Developments that occurred behind the scenes greatly
influenced the campaign as did trends and directives emanating from
Beijing. Shortly after the campaign began in February 1974, until
mid-year when the CC issued a document reining in the movement, leaders
of the ``mountain top'' faction virtually seized power in the province
through their control over the conduct of the campaign. In the second
half of the year Tan Qilong, with the backing of the central
authorities, shifted the focus of attention onto the ailing economy.
Despite a few final passing shots from the rebels, the campaign was
brought to a conclusion in November, although by August the peak of its
intensity had passed.

以下对浙江反林反孔运动的分析，同时考察其公开层面和较不显眼的层面；为了保持主题上的统一，这种处理在某些地方会打断时间顺序。幕后发生的事态极大地影响了这场运动，来自北京的趋势和指示同样如此。1974年2月运动开始后不久，直到年中中央委员会发布文件约束运动为止，``山上派''领导人凭借对运动开展方式的控制，实际上夺取了全省权力。下半年，谭启龙在中央当局支持下，把关注焦点转向了病弱的经济。尽管造反派最后还发动了几次零星攻击，这场运动还是在11月告一段落；不过到8月时，其强度高峰已经过去。

At the opening rally of the campaign in Hangzhou on 2 February 1974
there was no outward indication of the dramatic and turbulent events
which were about to occur.\footnote{ZJRB, February 3,1974, p.~1; ZPS,
  February 3,1974, URS, 75:14 (May 17,1974), pp.~177-81; HZRB, February
  3, 1974. For data on the provincial rallies held between February and
  March 1974, see K'ung Te-liang, ``The Maoist Mobilization for
  Criticizing Lin Piao and Confucius'', I\&S, 10: 10 (July 1974),
  pp.~47-9.
  《浙江日报》，1974年2月3日，第1版；ZPS，1974年2月3日，URS，75:14（1974年5月17日），第177-181页；《杭州日报》，1974年2月3日。关于1974年2月至3月间各省集会的数据，见K'ung
  Te-liang：《毛主义者为批林批孔进行的动员》，I\&S，10:10（1974年7月），第47-49页。}
Tan Qilong presided, Chen Weida and Chen Bing relayed central
directives, and Chai Qikun delivered the report on behalf of the ZPC.
Chai advocated the mutually-exclusive policies of simultaneously
strengthening party leadership while giving ``free rein to the masses''.
He stated that the campaign would link up with the ``realities of class
struggle'' in the province -- a phrase which indicated that personal
vendettas and powerplays were taking place in the upper echelons of the
leadership. In the long-practised tradition of Chinese political
campaigns, Chai proposed the establishment of experimental points and
the training of backbone elements. A Zhejiang Daily article of 1
February had sanctioned the formation of special contingents drawn
principally from among workers, peasants and soldiers.\footnote{ZJRB,
  February 1,1974, p.~1. 《浙江日报》，1974年2月1日，第1版。} According
to sources in Taiwan, the CCP CC issued an eight-point directive in
February prohibiting armed struggle, the ``exchange of revolutionary
experience'', and other activities which had been popularized by the Red
Guards in the Cultural Revolution.\footnote{China News Agency (Taibei),
  July 4, 1974, SWB/FE/4646/BII/11-12.
  中国新闻社（台北），1974年7月4日，SWB/FE/4646/BII/11-12。} These
constraints clearly assisted the provincial leadership in keeping the
campaign in low gear, for the moment at least.

在1974年2月2日杭州举行的运动开幕大会上，外表上并没有迹象显示即将发生戏剧性而动荡的事件。\footnote{ZJRB,
  February 3,1974, p.~1; ZPS, February 3,1974, URS, 75:14 (May 17,1974),
  pp.~177-81; HZRB, February 3, 1974. For data on the provincial rallies
  held between February and March 1974, see K'ung Te-liang, ``The Maoist
  Mobilization for Criticizing Lin Piao and Confucius'', I\&S, 10: 10
  (July 1974), pp.~47-9.
  《浙江日报》，1974年2月3日，第1版；ZPS，1974年2月3日，URS，75:14（1974年5月17日），第177-181页；《杭州日报》，1974年2月3日。关于1974年2月至3月间各省集会的数据，见K'ung
  Te-liang：《毛主义者为批林批孔进行的动员》，I\&S，10:10（1974年7月），第47-49页。}
谭启龙主持会议，陈伟达和陈冰传达中央指示，柴启琨代表浙江省委作报告。柴主张两项彼此排斥的方针：一方面加强党的领导，另一方面让``群众放手''。他说，这场运动将同全省``阶级斗争的实际''相结合------这个短语表明，在领导层高层正在发生私人恩怨和权力斗争。按照中国政治运动长期形成的传统，柴提出要建立试点并培养骨干。2月1日《浙江日报》的一篇文章批准成立主要由工人、农民和士兵组成的特别队伍。\footnote{ZJRB,
  February 1,1974, p.~1. 《浙江日报》，1974年2月1日，第1版。}
据台湾方面消息，中共中央在2月发布了一项八点指示，禁止武斗、``交流革命经验''以及其他在文化大革命中由红卫兵推广开来的活动。\footnote{China
  News Agency (Taibei), July 4, 1974, SWB/FE/4646/BII/11-12.
  中国新闻社（台北），1974年7月4日，SWB/FE/4646/BII/11-12。}
这些约束显然有助于省级领导层暂时使运动保持低调。

Four days after this rally, provincial-level administrative organs held
an oath-taking rally which was attended by provincial civilian and
military leaders.\footnote{ZJRB, February 7,1974, p.~1; ZPS, February
  7,1974, URS, 76:1, pp.~2-5.
  《浙江日报》，1974年2月7日，第1版；ZPS，1974年2月7日，URS，76:1，第2-5页。}
The moderation of the gathering can be gauged from the praise that it
directed toward Jiang Qing's model anti-chemical warfare company for its
contributions to farm-work in Zhejiang. Although militant slogans were
shouted, the singing of the PLA's ``three discipline and eight points
for attention'' signified an orderly atmosphere. One ominous note for
the provincial party leaders, however, was the suggestion that they play
an active role in the campaign.

这次大会四天后，省级行政机关举行了一次宣誓大会，省级党政和军事领导人出席。\footnote{ZJRB,
  February 7,1974, p.~1; ZPS, February 7,1974, URS, 76:1, pp.~2-5.
  《浙江日报》，1974年2月7日，第1版；ZPS，1974年2月7日，URS，76:1，第2-5页。}
这次集会的温和程度，可以从它对江青的防化连典型在浙江支农工作中的贡献加以表扬看出。尽管会上喊出了战斗性口号，但唱起解放军的``三大纪律八项注意''显示出一种有秩序的气氛。然而，对省级党委领导人来说，一个不祥的信号是，有人建议他们在运动中发挥积极作用。

In stark contrast to the discipline and order on display at these
rallies, on 7 February the workers' congress mobilized 30,000 workers to
attend a rally in Exhibition Square.\footnote{ZJRB, February 8,1974,
  p.~1; ZPS, February 8, 1974, URS, 76:1, pp.~5-8.
  《浙江日报》，1974年2月8日，第1版；ZPS，1974年2月8日，URS，76:1，第5-8页。}
Slogans of the Cultural Revolution such as ``It is right to rebel
against reactionaries'' rang in the air in a ``fervent atmosphere of
unity and militancy''. The leadership of the ZPC did not attend,
although members as well as ``responsible comrades'' of the ZPRC,
including Zhang Yongsheng and Weng Senhe, were present. Leaders of the
municipal party committee also attended. Factories in Hangzhou,
including Weng's silk complex, sent along representatives to address the
gathering while others, including the iron and steel mill and the
Zhejiang Hemp Mill, submitted written statements. The speakers stressed
that the campaign should concern itself with concrete issues, otherwise
it could ``possibly go astray and deviate from the orientation''. In
other words, rather than being diverted by historical debates, it was
suggested that the campaign should openly target individual party
leaders. In his closing speech He Xianchun appealed to the city's
workers to become the main force in the campaign and to ``go against the
tide''. A Deputy-secretary of HMC reportedly attacked the leaders of the
provincial committee by name for ``not talking about the basic line''
and asserted that ``the rightist trend in Zhejiang comes from
responsible comrades of the ZPC''.\footnote{HZRB, July 19,1977, October
  21,1977, October 24,1977.
  《杭州日报》，1977年7月19日、1977年10月21日、1977年10月24日。} The
proceedings ended with a march through the streets of Hangzhou.

同这些大会上展现出的纪律和秩序形成鲜明对照的是，2月7日，工人代表大会动员三万名工人到展览广场参加集会。\footnote{ZJRB,
  February 8,1974, p.~1; ZPS, February 8, 1974, URS, 76:1, pp.~5-8.
  《浙江日报》，1974年2月8日，第1版；ZPS，1974年2月8日，URS，76:1，第5-8页。}
``造反有理''之类文化大革命口号在``团结战斗的热烈气氛''中响彻会场。浙江省委领导没有出席，不过浙江省革委会的委员和``负责同志''，包括张永生和翁森鹤在内，均在场。市委领导也参加了会议。杭州的工厂，包括翁所在的丝绸联合企业，派代表到会发言；其他单位，包括钢铁厂和浙江麻纺厂，则提交了书面发言。发言者强调，运动应当涉及具体问题，否则就``可能走偏方向''。换言之，他们建议运动不应被历史争论转移方向，而应公开针对个别党的领导人。贺贤春在闭幕讲话中号召全市工人成为运动的主力军，并``反潮流''。据报道，杭州市委一名副书记点名攻击省委领导``不讲基本路线''，并声称``浙江的右倾思潮来自省委负责同志''。\footnote{HZRB,
  July 19,1977, October 21,1977, October 24,1977.
  《杭州日报》，1977年7月19日、1977年10月21日、1977年10月24日。}
会议最后以穿越杭州街道的游行结束。

A further mobilization rally took place on 9 February 1974 for 500
members of the provincial Mao Zedong Thought propaganda teams stationed
in propaganda, cultural and education departments.\footnote{ZJRB,
  February 11,1974, p.~1. 《浙江日报》，1974年2月11日，第1版。} It was
sponsored by the political work group of the ZPRC headed by Shen Ce, a
``revolutionary reading cadre'' who together with Lai Keke had returned
to prominence in 1967. Shen appointed Weng Senhe a deputy-leader of the
political work group at about this time and allowed him to assume charge
of the provincial office responsible for the propaganda
teams.\footnote{Weng Senhe bufen zuizheng cailiao, p.~33; RMRB, December
  2, 1977, p.~2.
  《翁森鹤部分罪证材料》，第33页；《人民日报》，1977年12月2日，第2版。}
These posts gave Weng the right to visit campuses where he could display
his oratorical skills to great effect.\footnote{For example, on April
  19,1974, Weng visited the campus of Hangzhou University and delivered
  a speech praising Zhang Tiesheng, the Liaoning student who had handed
  in the blank examination paper in 1973. Briefing at Hangzhou
  University, May 19,1978.
  例如，1974年4月19日，翁到杭州大学校园访问并发表讲话，赞扬辽宁学生张铁生；张铁生曾在1973年交白卷。杭州大学简报，1978年5月19日。}
Weng's office made its presence felt very early in the campaign. On 9
February Zhejiang Daily carried a page one ideological commentary
concerning the anti-chemical warfare company. Certain phrases offended
the sensibilities of the political work group's staff who penned a
stinging missive on 11 February. The provincial daily finally frontpaged
the reply on 10 March together with a craven apology for denigrating the
role of workers and peasants in the campaign.\footnote{ZJRB, February
  9,1974, p.~1, March 11,1974, p.~1.
  《浙江日报》，1974年2月9日第1版，1974年3月11日第1版。}

1974年2月9日，又为派驻宣传、文化和教育部门的五百名省毛泽东思想宣传队成员举行了一次动员大会。\footnote{ZJRB,
  February 11,1974, p.~1. 《浙江日报》，1974年2月11日，第1版。}
这次大会由省革委会政治工作组主办，该组由沈策领导；沈策是一名``革命读书干部''，1967年同赖可可一起重新崭露头角。大约在这个时候，沈任命翁森鹤为政治工作组副组长，并让他负责管理省里主管宣传队的办公室。\footnote{Weng
  Senhe bufen zuizheng cailiao, p.~33; RMRB, December 2, 1977, p.~2.
  《翁森鹤部分罪证材料》，第33页；《人民日报》，1977年12月2日，第2版。}
这些职务赋予翁进入校园的权利，使他能够极为有效地展示自己的演说才能。\footnote{For
  example, on April 19,1974, Weng visited the campus of Hangzhou
  University and delivered a speech praising Zhang Tiesheng, the
  Liaoning student who had handed in the blank examination paper in
  1973. Briefing at Hangzhou University, May 19,1978.
  例如，1974年4月19日，翁到杭州大学校园访问并发表讲话，赞扬辽宁学生张铁生；张铁生曾在1973年交白卷。杭州大学简报，1978年5月19日。}
翁的办公室在运动很早阶段就显示了其存在感。2月9日，《浙江日报》在头版刊登了一篇关于防化连的思想评论。某些措辞冒犯了政治工作组工作人员的感情，于是他们在2月11日写了一封措辞尖锐的信。省报最终于3月10日把这封回信登上头版，并附上一篇卑怯的道歉，承认贬低了工人和农民在运动中的作用。\footnote{ZJRB,
  February 9,1974, p.~1, March 11,1974, p.~1.
  《浙江日报》，1974年2月9日第1版，1974年3月11日第1版。}

On 13 February 80,000 tertiary and secondary students. Red Guards,
teachers, and propaganda team members, assembled in Exhibition Square.
They listened to a responsible member of the CCP ZPC demand a ruthless
retaliatory strike at the ``bourgeois counter-attack to settle old
scores and to restore its influence''. This theme was picked up three
days later when the ZPC called a meeting of provincial-level cadres. The
meeting ``criticized the Rightist trend of thought of negating the great
proletarian cultural revolution, linking the rightist trend closely with
the current struggle between the two lines in our province''.\footnote{ZJRB,
  February 17,1974, p.~1; ZPS, February 17,1974, SWB/FE/4539/BII/27-8.
  《浙江日报》，1974年2月17日，第1版；ZPS，1974年2月17日，SWB/FE/4539/BII/27-28。}

2月13日，八万名高等院校和中学学生、红卫兵、教师以及宣传队成员聚集在展览广场。他们听取了一名中共浙江省委负责成员的讲话，后者要求对``资产阶级翻旧账、复辟影响的反攻倒算''进行无情的反击。三天后，浙江省委召开省级干部会议，又接过了这一主题。会议``批判了否定无产阶级文化大革命的右倾思潮，把这种右倾思潮同我省当前两条路线斗争紧密联系起来''。\footnote{ZJRB,
  February 17,1974, p.~1; ZPS, February 17,1974, SWB/FE/4539/BII/27-8.
  《浙江日报》，1974年2月17日，第1版；ZPS，1974年2月17日，SWB/FE/4539/BII/27-28。}

The meaning of these obscure and threatening references was spelled out
at a conference on schistosomiasis {[}!{]} held in late February. The
following passage was included in the broadcast relaying details of the
conference:

这些含糊而具有威胁性的说法，在2月下旬召开的一次血吸虫病 {[}!{]}
会议上得到了明确说明。转播会议详情的广播中包括以下段落：

\begin{quote}
In certain localities and units in Zhejiang some persons obstinately
maintaining the bourgeois reactionary stand have adopted all possible
means to reverse history, turned facts upside down, and have engaged in
all sorts of activities to turn back the wheel of history and restore
the old order and capitalism. They have negated all instructions issued
by Chairman Mao during the great proletarian cultural revolution, and
have overturned the brilliant decision endorsed by Chairman Mao and the
Party Central Committee on handling the question of the great
proletarian cultural revolution, in a vain attempt to reverse the
correct verdicts passed on them and to launch counter-attacks in
retaliation. In addition they have also endeavored to strike at the
newly emerging forces and have opposed such things as the revolution in
education, in literature and art, and in health work. May 7 cadre
schools and work regarding educated young people settled in the
countryside.\footnote{ZJRB, February 27,1974, p.~1; quotation from ZPS,
  February 27,1974, SWB/FE/4544/BII/17.
  《浙江日报》，1974年2月27日，第1版；引文见ZPS，1974年2月27日，SWB/FE/4544/BII/17。}
\end{quote}

\begin{quote}
在浙江某些地方和单位，一些顽固坚持资产阶级反动立场的人，采取一切可能的手段颠倒历史、颠倒事实，并从事各种开历史倒车、恢复旧秩序和资本主义的活动。他们否定毛主席在无产阶级文化大革命期间发出的一切指示，推翻毛主席和党中央关于处理无产阶级文化大革命问题所批准的英明决定，妄图推翻对他们作出的正确结论，并进行反攻倒算。此外，他们还竭力打击新生力量，反对教育革命、文艺革命、卫生工作革命等事业，反对五七干校以及知识青年上山下乡工作。\footnote{ZJRB,
  February 27,1974, p.~1; quotation from ZPS, February 27,1974,
  SWB/FE/4544/BII/17.
  《浙江日报》，1974年2月27日，第1版；引文见ZPS，1974年2月27日，SWB/FE/4544/BII/17。}
\end{quote}

He Xianchun chaired a further lively rally in Hangzhou on 2 March 1974
in his capacity as Vice-chairman of the ZPTUC. The gathering was
sponsored by the HMWC and attended by 30,000 of the city's workers. On
28 February the workers' congress had sent out invitations to district
and county party organizations. In a provocative letter of 25 February,
Weng and He ordered Tan Qilong and Tie Ying to attend the
rally.\footnote{Weng Senhe bufen zuizheng cailiao, pp.~41-2; HZRB, March
  3, 1974; ZPS, March 3, 1974, URS, 76:1, pp.~8-10.
  《翁森鹤部分罪证材料》，第41-42页；《杭州日报》，1974年3月3日；ZPS，1974年3月3日，URS，76:1，第8-10页。}
Before party and mass leaders who had gathered from across the province.
Tan and Tie experienced the humiliation of listening to public criticism
of their leadership.\footnote{HJRB criticism group, ``The principle of
  the party nature of a proletarian newspaper is not easily twisted'';
  HZRB, November 22, 1977. A further mass rally was held on March 4 at
  the provincial gymnasium. See Weng Senhe bufen zuizheng cailiao,
  p.~47.
  《杭州日报》批判组：《无产阶级报纸的党性原则是不容易篡改的》；《杭州日报》，1977年11月22日。3月4日又在省体育馆举行了一次群众大会。见《翁森鹤部分罪证材料》，第47页。}
The rally was to prove Tie's last public appearance in Hangzhou for a
year and he was, in his own words, subsequently ``deprived of the power
(权力) to work''.\footnote{Tie Ying, ZJRB, January 9, 1985.
  铁瑛，《浙江日报》，1985年1月9日。} According to a March 12 speech
delivered by Weng Senhe on Mt Pingfeng, provincial officials were at
this time suspended for investigation (停职检查), interrogated in
isolation (隔离审查) or dismissed (撤职罢官).\footnote{Weng Senhe bufen
  zuizheng cailiao, p.~44. 《翁森鹤部分罪证材料》，第44页。}

1974年3月2日，贺贤春以浙江省总工会副主席身份，在杭州又主持了一场气氛热烈的大会。大会由杭州市工人代表大会主办，三万名杭州工人出席。2月28日，工人代表大会已向区县党组织发出邀请。在2月25日一封挑衅性的信中，翁和贺命令谭启龙、铁瑛出席大会。\footnote{Weng
  Senhe bufen zuizheng cailiao, pp.~41-2; HZRB, March 3, 1974; ZPS,
  March 3, 1974, URS, 76:1, pp.~8-10.
  《翁森鹤部分罪证材料》，第41-42页；《杭州日报》，1974年3月3日；ZPS，1974年3月3日，URS，76:1，第8-10页。}
在来自全省各地的党政和群众组织领导人面前，谭和铁遭受了聆听公开批评其领导工作的羞辱。\footnote{HJRB
  criticism group, ``The principle of the party nature of a proletarian
  newspaper is not easily twisted''; HZRB, November 22, 1977. A further
  mass rally was held on March 4 at the provincial gymnasium. See Weng
  Senhe bufen zuizheng cailiao, p.~47.
  《杭州日报》批判组：《无产阶级报纸的党性原则是不容易篡改的》；《杭州日报》，1977年11月22日。3月4日又在省体育馆举行了一次群众大会。见《翁森鹤部分罪证材料》，第47页。}
事实证明，这次大会是铁瑛一年内在杭州最后一次公开露面；用他自己的话说，他随后被``剥夺了工作权力''。\footnote{Tie
  Ying, ZJRB, January 9, 1985. 铁瑛，《浙江日报》，1985年1月9日。}
根据翁森鹤3月12日在屏风山发表的讲话，当时省级官员有的被停职检查，有的被隔离审查，有的被撤职罢官。\footnote{Weng
  Senhe bufen zuizheng cailiao, p.~44. 《翁森鹤部分罪证材料》，第44页。}

As if to support the rebels' action, on 15 March the People's Daily
directed a thinly-veiled blast at the policy of rehabilitating cadres
who had been disgraced in the Cultural Revolution.\footnote{``再批``克己复礼''''
  (Once again criticize ``restrain oneself and restore the rights''),
  RMRB, March 15,1974, p.~1.
  《再批``克己复礼''》（再次批判``克制自己、恢复礼制''），《人民日报》，1974年3月15日，第1版。}
The radical upsurge in the campaign which the editorial set off
immediately made itself felt in Zhejiang. When the ZPC convened a second
rally on 21 March 1974 great changes had occurred in the provincial
leadership since February.\footnote{CNA, 959 (May 10,1974), p.6, wrongly
  dated the phenomenon from the April 20 conference discussed below.
  CNA，959（1974年5月10日），第6页，错误地把这一现象定为始于下文讨论的4月20日会议。}
At the 2 February mobilization rally ZPC leaders had spoken and issued
directives. On 21 March the crowd was addressed by different leaders.
Zhang Yongsheng presided over the rally and both Weng Senhe and He
Xianchun were featured speakers. The ``mountain top'' faction leaders
had, in no uncertain terms, taken over. The meteoric return to
prominence of Cultural Revolution rebels was not confined to Zhejiang.
The Hong Kong weekly Far Eastern Economic Review commented in May 1974
that ``genuine Maoist revolutionaries'' who had lost power after 1969
had ``re-established their position through their membership of
revolutionary committees. Although previously''eclipsed by the rebirth
of CCP committees'', the revolutionary committees now ``commanded
genuine influence'' and were praised as an important creation of the
Cultural Revolution.\footnote{FEER, May 20,1974, p.~16.
  FEER，1974年5月20日，第16页。}

仿佛是为了支持造反派的行动，3月15日，《人民日报》对文化大革命中失势干部的平反政策发起了一次含蓄的猛烈攻击。\footnote{``再批``克己复礼''''
  (Once again criticize ``restrain oneself and restore the rights''),
  RMRB, March 15,1974, p.~1.
  《再批``克己复礼''》（再次批判``克制自己、恢复礼制''），《人民日报》，1974年3月15日，第1版。}
这篇社论引发的运动激进化浪潮，立刻在浙江显现出来。1974年3月21日浙江省委召开第二次大会时，自2月以来省领导层已经发生了巨大变化。\footnote{CNA,
  959 (May 10,1974), p.6, wrongly dated the phenomenon from the April 20
  conference discussed below.
  CNA，959（1974年5月10日），第6页，错误地把这一现象定为始于下文讨论的4月20日会议。}
在2月2日的动员大会上，省委领导发言并发布指示。3月21日，向群众讲话的则是另一批领导人。张永生主持大会，翁森鹤和贺贤春都是主要发言人。``山上派''领导人已经毫不含糊地接管了局面。文化大革命造反派迅速重新显赫，并不限于浙江。香港周刊《远东经济评论》1974年5月评论说，1969年后失去权力的``真正的毛主义革命者''，已经``通过革命委员会成员身份重新确立了他们的地位''。尽管此前他们``因中共党委的重生而黯然失色''，革命委员会现在却``掌握了真正的影响力''，并被赞扬为文化大革命的一项重要创造。\footnote{FEER,
  May 20,1974, p.~16. FEER，1974年5月20日，第16页。}

Tan Qilong spoke first at the 21 March rally. He

在3月21日的大会上，谭启龙首先发言。他

\begin{quote}
expressed his determination to plunge into the struggle together with
the broad masses of armymen and people and make severe criticism of the
reactionary ideological tide of attempting to restore capitalism.
\end{quote}

\begin{quote}
表示决心同广大军民一起投入斗争，对妄图复辟资本主义的反动思想潮流进行严厉批判。
\end{quote}

These words suggest that Tan had been placed in an extremely awkward
position. By being forced to commit himself more actively to the
campaign he was also losing the ability to protect subordinates such as
Tie Ying and Xia Qi. The atmosphere at the meeting was militant and
defiant. It demanded of leading cadres that they lead the way in
``making disclosures and criticism'' so that they could render ``new
meritorious service to the people''. Shouts and cries reminiscent of the
Cultural Revolution mass rallies echoed across the square. The report
stated that a ``very sharp and acute struggle'' was in progress over the
assessment of the Cultural Revolution and the attempt to negate its
``great successes and achievements''. Apart from Weng Senhe and He
Xianchun, another speaker was Wang Xing, 2nd Secretary of the CCP HMC
and head of the Hangzhou party committee's anti-Lin anti-Confucius small
group.\footnote{In May 1959, Wang Xing was identified as party secretary
  of the Xinanjiang hydro-electric engineering bureau. Construction of
  the Xinanjiang hydro-electric power station, a major capital
  construction project in Zhejiang, commenced in April 1957 and was
  completed three years later. Zhou Enlai visited the construction site
  in 1959 and Hua Guofeng the operating station in 1971. Visit to
  Xinanjiang, September 4,1977. From October 1962 until September 1964
  Wang was Deputy-governor of Zhejiang. In December 1973 he was publicly
  identified as 2nd Secretary of the CCP HMC.
  1959年5月，王兴被列为新安江水电工程局党委书记。新安江水电站是浙江的一项重大基本建设工程，1957年4月开工，三年后建成。周恩来于1959年视察工地，华国锋于1971年视察运行中的电站。见《新安江访问记》，1977年9月4日。1962年10月至1964年9月，王任浙江省副省长。1973年12月，他被公开列为中共杭州市委第二书记。}
References to the need to differentiate between the two kinds of
contradictions and perfunctory remarks about the state of the economy,
were the only indication of any leavening of the radicalism. The mood of
the rally was summed up in the following sentence of the monitored
report:

这些话表明，谭已经被置于极其尴尬的位置。由于被迫更积极地表态投入运动，他也正在失去保护铁瑛、夏琦等下属的能力。会议气氛充满战斗性和挑衅意味。会议要求领导干部带头``揭发批判''，以便为人民``立新功''。令人想起文化大革命群众大会的呼喊声响彻广场。报道说，围绕对文化大革命的评价以及否定其``伟大胜利和伟大成果''的企图，正在展开一场``十分尖锐激烈的斗争''。除翁森鹤和贺贤春外，另一名发言者是中共杭州市委第二书记、杭州市委反林反孔小组组长王兴。\footnote{In
  May 1959, Wang Xing was identified as party secretary of the
  Xinanjiang hydro-electric engineering bureau. Construction of the
  Xinanjiang hydro-electric power station, a major capital construction
  project in Zhejiang, commenced in April 1957 and was completed three
  years later. Zhou Enlai visited the construction site in 1959 and Hua
  Guofeng the operating station in 1971. Visit to Xinanjiang, September
  4,1977. From October 1962 until September 1964 Wang was
  Deputy-governor of Zhejiang. In December 1973 he was publicly
  identified as 2nd Secretary of the CCP HMC.
  1959年5月，王兴被列为新安江水电工程局党委书记。新安江水电站是浙江的一项重大基本建设工程，1957年4月开工，三年后建成。周恩来于1959年视察工地，华国锋于1971年视察运行中的电站。见《新安江访问记》，1977年9月4日。1962年10月至1964年9月，王任浙江省副省长。1973年12月，他被公开列为中共杭州市委第二书记。}
关于必须区分两类矛盾的提法，以及对经济状况的敷衍性评论，是这场激进主义中唯一稍有调和色彩的迹象。监听报告用以下一句话概括了大会的情绪：

\begin{quote}
It is required to exert continuously the revolutionary spirit of going
against the tide and display a firm stand and the true colors of one's
flags {[}sic{]} when dealing with questions of principle about the
proletarian cultural revolution.\footnote{ZJRB, March 23, 1974, p.~1;
  quotation from ZPS, March 23,1974, URS, 75:14, pp.~181-84.
  《浙江日报》，1974年3月23日，第1版；引文见ZPS，1974年3月23日，URS，75:14，第181-184页。}
\end{quote}

\begin{quote}
要求在对待无产阶级文化大革命的原则问题时，继续发扬反潮流的革命精神，显示坚定立场和自己旗帜的本色
{[}sic{]}。\footnote{ZJRB, March 23, 1974, p.~1; quotation from ZPS,
  March 23,1974, URS, 75:14, pp.~181-84.
  《浙江日报》，1974年3月23日，第1版；引文见ZPS，1974年3月23日，URS，75:14，第181-184页。}
\end{quote}

While it is true that Zhang Yongsheng, Weng Senhe and He Xianchun
mounted the platform at provincial campaign rallies as standing
committee members (and in Zhang's case Vice-chairman) of the ZPRC, their
power was to extend well beyond these formal positions. On 22 June 1974,
Wang Hongwen instructed the CCP ZPC leadership to allow the three rebels
to attend meetings (列席) of its standing committee in the capacity of
observers.\footnote{Weng Senhe bufen zuizheng cailiao, pp.~8, 21; PR,
  no. 40 (September 30, 1977), p.~25; RMRB, December 9,1976, p.~3; RMRB,
  January 5,1977, p.~2; ZJRB editorial in HZRB, December 31, 1976, p.~2;
  HZRB, January 23,1977. In February 1974 two Cultural Revolution rebels
  from Shanghai, Chen Ada and Ye Changming were invited to attend
  meetings of that municipality's CCP standing committee. See CCP CC
  zhongfa (1977), No.~37 in I\&S, 15:1 (1979), p.~115. This also
  occurred in Xiaoshan county. Eight ``fighters going against the
  tide'', five of whom had just been admitted into the CCP, sat in on
  standing committee meetings of the county party committee from June
  1974. See Xiaoshan xianzhi, p.~54.
  《翁森鹤部分罪证材料》，第8、21页；PR，第40号（1977年9月30日），第25页；《人民日报》，1976年12月9日，第3版；《人民日报》，1977年1月5日，第2版；《杭州日报》转载《浙江日报》社论，1976年12月31日，第2版；《杭州日报》，1977年1月23日。1974年2月，上海的两名文化大革命造反派陈阿大和叶昌明被邀请参加该市中共常委会会议。见中共中央中发（1977）第37号，载I\&S，15:1（1979），第115页。这种情况也发生在萧山县。自1974年6月起，八名``反潮流战士''列席县委常委会会议，其中五人刚刚被接收入党。见《萧山县志》，第54页。}
Their presence would certainly have been unsettling to the other
members. Attendance at the meetings provided the trio with direct access
to the highest formal decision-making body in the province and the
prestige accompanying such promotion.

诚然，张永生、翁森鹤和贺贤春是以省革委会常委成员（张的情况还包括副主任）身份登上省级运动大会主席台的，但他们的权力远远超出了这些正式职务。1974年6月22日，王洪文指示中共浙江省委领导，允许这三名造反派以观察员身份列席省委常委会会议。\footnote{Weng
  Senhe bufen zuizheng cailiao, pp.~8, 21; PR, no. 40 (September 30,
  1977), p.~25; RMRB, December 9,1976, p.~3; RMRB, January 5,1977, p.~2;
  ZJRB editorial in HZRB, December 31, 1976, p.~2; HZRB, January
  23,1977. In February 1974 two Cultural Revolution rebels from
  Shanghai, Chen Ada and Ye Changming were invited to attend meetings of
  that municipality's CCP standing committee. See CCP CC zhongfa (1977),
  No.~37 in I\&S, 15:1 (1979), p.~115. This also occurred in Xiaoshan
  county. Eight ``fighters going against the tide'', five of whom had
  just been admitted into the CCP, sat in on standing committee meetings
  of the county party committee from June 1974. See Xiaoshan xianzhi,
  p.~54.
  《翁森鹤部分罪证材料》，第8、21页；PR，第40号（1977年9月30日），第25页；《人民日报》，1976年12月9日，第3版；《人民日报》，1977年1月5日，第2版；《杭州日报》转载《浙江日报》社论，1976年12月31日，第2版；《杭州日报》，1977年1月23日。1974年2月，上海的两名文化大革命造反派陈阿大和叶昌明被邀请参加该市中共常委会会议。见中共中央中发（1977）第37号，载I\&S，15:1（1979），第115页。这种情况也发生在萧山县。自1974年6月起，八名``反潮流战士''列席县委常委会会议，其中五人刚刚被接收入党。见《萧山县志》，第54页。}
他们的在场当然会使其他成员感到不安。列席会议使这三人能够直接接触全省最高正式决策机构，并获得这种提升所附带的威望。

More importantly, a provincial-level small group was especially
established to conduct the campaign. Jiang Qing was reported to have
been the head of the central group in charge of the campaign. The group
in Zhejiang was known as the ``double criticism small group''
(双批小组).\footnote{See Sheng P'eng, ``The Generals'', p.~83; Li
  Chung-ta, ``Why Did Mao Tse-tung'', p.~92. The existence of such a
  group at the center was denied in Weng Senhe bufen zuizheng cailiao,
  p.~25. Such groups were also established at the county level between
  February and May. See Xiangshan xianzhi, p.~36; Xianju xianzhi, p.~24;
  Xiaoshan xianzhi, p.~54; Jiande xianzhi, p.~26; Lanxi shizhi, p.~21;
  Linhai xianzhi, p.~31. 见Sheng P'eng：《将军们》，第83页；Li
  Chung-ta：《毛泽东为什么》，第92页。《翁森鹤部分罪证材料》第25页否认中央存在这样一个小组。1974年2月至5月，县一级也建立了这类小组。见《象山县志》，第36页；《仙居县志》，第24页；《萧山县志》，第54页；《建德县志》，第26页；《兰溪市志》，第21页；《临海县志》，第31页。}
For the duration of the campaign its authority exceeded that of the
provincial party committee. Zhang Yongsheng reportedly appointed the
members of the group in consultation with Deputy-party Secretary Chai
Qikun, who turned down the offer of becoming group leader so that, in
his own words, he would retain the freedom to participate in criticism
of the provincial leadership.\footnote{Zhang Yongshengde zuizheng
  cailiao, p.~83; HZRB, March 3, 1977.
  《张永生的罪证材料》，第83页；《杭州日报》，1977年3月3日。} Tan Qilong
was group leader and Chai, Wang Xing, Weng Senhe and He Xianchun were
among the deputy leaders. Weng also was chairman of the office
established to handle the administrative side of the
campaign.\footnote{Weng Senhe bufen zuizheng cailiao, p.~2.
  《翁森鹤部分罪证材料》，第2页。}

更重要的是，省里专门成立了一个省级小组来开展这场运动。据报道，江青是负责该运动的中央小组负责人。浙江的小组被称为``双批小组''。\footnote{See
  Sheng P'eng, ``The Generals'', p.~83; Li Chung-ta, ``Why Did Mao
  Tse-tung'', p.~92. The existence of such a group at the center was
  denied in Weng Senhe bufen zuizheng cailiao, p.~25. Such groups were
  also established at the county level between February and May. See
  Xiangshan xianzhi, p.~36; Xianju xianzhi, p.~24; Xiaoshan xianzhi,
  p.~54; Jiande xianzhi, p.~26; Lanxi shizhi, p.~21; Linhai xianzhi,
  p.~31. 见Sheng P'eng：《将军们》，第83页；Li
  Chung-ta：《毛泽东为什么》，第92页。《翁森鹤部分罪证材料》第25页否认中央存在这样一个小组。1974年2月至5月，县一级也建立了这类小组。见《象山县志》，第36页；《仙居县志》，第24页；《萧山县志》，第54页；《建德县志》，第26页；《兰溪市志》，第21页；《临海县志》，第31页。}
在运动持续期间，它的权威超过了省党委。据报道，张永生同省委副书记柴启琨商议后任命了小组成员；柴拒绝了担任组长的提议，因为用他自己的话说，他要保留参加批判省领导的自由。\footnote{Zhang
  Yongshengde zuizheng cailiao, p.~83; HZRB, March 3, 1977.
  《张永生的罪证材料》，第83页；《杭州日报》，1977年3月3日。}
谭启龙任组长，柴启琨、王兴、翁森鹤和贺贤春等人为副组长。翁还担任为处理运动行政事务而设立的办公室主任。\footnote{Weng
  Senhe bufen zuizheng cailiao, p.~2. 《翁森鹤部分罪证材料》，第2页。}

Control over the small group enabled the ``mountain top'' factional
leaders, in the words of articles written after their downfall, to lord
it over (在省委头上称王称霸)\footnote{CCP CC, zhongfa (1977), No.~37, in
  I\&S, 15:1 (1979), p.~111.
  中共中央，中发（1977）第37号，载I\&S，15:1（1979），第111页。} or
place themselves above (凌驾于省委之上)\footnote{HZRB, November 24,
  1976. 《杭州日报》，1976年11月24日。} the CCP ZPC. Small groups were
also established throughout the ZPMD system.\footnote{RMRB, September
  22,1978, p.~1. 《人民日报》，1978年9月22日，第1版。} According to a
senior cadre from the Hangzhou party committee, Weng Senhe and He
Xianchun raised the question of establishing such a group in the
municipality at a meeting of the standing committee of the HMC in
January 1974. They described the small group as an advisory body to the
standing committee (常委的参谋机构). Together with Second Secretary Wang
Xing, Weng and He drew up a namelist. Weng issued instructions that
non-party members could hold the position of deputy group leader and in
that capacity could attend party meetings at their units.\footnote{Weng
  Senhe bufen zuizheng cailiao, pp.~25-6, 27.
  《翁森鹤部分罪证材料》，第25-26、27页。} He was thus attempting to
bestow an importance on the small groups which would give them equal
status to party committees as well as prodding party committees into
admitting non-party leaders into the CCP. Tie Ying later attributed the
source of the rebels' power to Tan Qilong, whose leadership of the group
gave it a more official and impressive appearance. Tie repeated an
analogy borrowed by Weng Senhe from the Eastern Zhou dynasty to accuse
the rebels of manipulating the powerless emperor (Tan) to order his
subordinates about (挟天子以令诸侯). Weng also used the phrase
(掩耳盗铃)\footnote{Ibid, p.~21; Zhang Yongshengde zuizheng cailiao,
  p.~83; Tie Ying, HZRB, June 7, 1978.
  同上，第21页；《张永生的罪证材料》，第83页；铁瑛，《杭州日报》，1978年6月7日。}
to describe how the arrangement and membership of the small groups
deceived party organizations.

控制这个小组，使``山上派''领导人得以用他们倒台后文章中的说法，在中共浙江省委``头上称王称霸''，\footnote{CCP
  CC, zhongfa (1977), No.~37, in I\&S, 15:1 (1979), p.~111.
  中共中央，中发（1977）第37号，载I\&S，15:1（1979），第111页。}
或``凌驾于省委之上''。\footnote{HZRB, November 24, 1976.
  《杭州日报》，1976年11月24日。}
在浙江省军区系统内也普遍建立了小组。\footnote{RMRB, September 22,1978,
  p.~1. 《人民日报》，1978年9月22日，第1版。}
据杭州市委一名高级干部说，翁森鹤和贺贤春在1974年1月一次杭州市委常委会会议上提出了在市里建立这样一个小组的问题。他们把小组描述为常委会的参谋机构。翁、贺同第二书记王兴一起拟定了一份名单。翁下达指示说，非党员可以担任副组长，并以这一身份参加本单位的党内会议。\footnote{Weng
  Senhe bufen zuizheng cailiao, pp.~25-6, 27.
  《翁森鹤部分罪证材料》，第25-26、27页。}
因而，他试图赋予这些小组一种重要性，使其具有同党委同等的地位，同时也推动党委吸收非党员领导人加入中共。铁瑛后来把造反派权力的来源归因于谭启龙，因为谭担任该小组组长，使这个小组显得更正式、更有声势。铁重复了翁森鹤借自东周时期的一个比喻，指责造反派操纵无权的皇帝（谭）来号令其下属，即``挟天子以令诸侯''。翁还用``掩耳盗铃''\footnote{Ibid,
  p.~21; Zhang Yongshengde zuizheng cailiao, p.~83; Tie Ying, HZRB, June
  7, 1978.
  同上，第21页；《张永生的罪证材料》，第83页；铁瑛，《杭州日报》，1978年6月7日。}一语来形容小组的安排和成员构成如何欺骗党组织。
\#\# Rebel Tactics \#\# 造反派策略

\subsection{1. Party Recruitment and Cadre
Promotion}\label{party-recruitment-and-cadre-promotion}

\subsection{1.
党员吸收与干部提拔}\label{ux515aux5458ux5438ux6536ux4e0eux5e72ux90e8ux63d0ux62d4}

An important measure taken by the rebels to build up and reward their
supporters was to place them in important party and government posts in
the Hangzhou administration.\footnote{For the promotion of cadres in
  Shanghai during 1973-74, see Ye Yonglie, Wang Hongwen xingshuailu,
  pp.~430-40, and Gao Gao and Yan Jiaqi,``Wenhua dageming'', pp.~533-4.
  关于1973至1974年上海干部的提拔，见叶永烈：《王洪文兴衰录》，第430-440页；高皋、严家其：《文化大革命》，第533-534页。}
The ideological and political pretext was supplied by Mao's dictum
regarding ``combining the old, middle-aged and the young'' in party
committees. The organizational foundation was provided with the
formation at all levels across the province of campaign small groups,
which legitimated the admission of rebels into the CCP and their
subsequent promotion to officialdom. Weng Senhe listed four criteria for
appointment: these were a high consciousness in line struggle, an
unblemished personal political history, a preference for party members
and active participation in the rebellion of late 1973 as proven by
one's presence on Mt. Pingfeng.\footnote{Weng Senhe bufen zuizheng
  cailiao, p.~29. 《翁森鹤部分罪证材料》，第29页。}

造反派用来壮大并奖赏其支持者的一项重要措施，是把他们安插到杭州行政体系中重要的党政岗位上。\footnote{For
  the promotion of cadres in Shanghai during 1973-74, see Ye Yonglie,
  Wang Hongwen xingshuailu, pp.~430-40, and Gao Gao and Yan
  Jiaqi,``Wenhua dageming'', pp.~533-4.
  关于1973至1974年上海干部的提拔，见叶永烈：《王洪文兴衰录》，第430-440页；高皋、严家其：《文化大革命》，第533-534页。}
其意识形态和政治借口来自毛关于党委中``老中青三结合''的论断。组织基础则由全省各级运动小组的建立提供，这些小组使造反派加入中共以及随后晋升为官员具有了合法性。翁森鹤列出了任用的四项标准：路线斗争觉悟高，个人政治历史清白，优先考虑党员，以及积极参加
1973 年末的造反行动，这一点以是否到过屏风山为证明。\footnote{Weng Senhe
  bufen zuizheng cailiao, p.~29. 《翁森鹤部分罪证材料》，第29页。}

It was at the municipal rather than at the provincial level that the
``mountain top'' faction met the least opposition from sections of the
party leadership, and in certain cases passivity or even encouragement.
Weng and his colleagues realized that once they had set a precedent they
could move both down to the county and work unit level and up to the
provincial level. To this end they suggested to the ZPC that it instruct
its subordinate body in Hangzhou to take the initiative.\footnote{Ibid.
  同上。} Simultaneously, they attempted to paralyze the operation of
party branches in units of production by disruptive activities and then
allow the besieged leaders to take the blame for the resultant economic
failures. Thus Zhang, Weng and He elaborated a two-pronged strategy of
infiltration and confrontation. The novelty of the approach was
commented upon several years later by Tie Ying in a lecture he gave in
1980. Tie said that, ``A few years back, Zhang Yongsheng, Weng Senhe, He
Xianchun and their ilk -- while roundly cursing the Communist Party --
exhausted every means to worm their way into our party''.\footnote{Tie
  Ying, ``共产党员要为四化建设的伟大事业而献身'' (Communist party
  members must contribute their all to the great cause of the four
  modernizations), ZJRB, May 31, 1980, pp.~1-3. Translated in ZPS, May
  31,1980, FBIS/CHI/108 (1980), O8.
  铁瑛：《共产党员要为四化建设的伟大事业而献身》（共产党员必须为四个现代化的伟大事业贡献一切），《浙江日报》，1980年5月31日，第1-3版。译文见ZPS，1980年5月31日，FBIS/CHI/108（1980），O8。}
The rebels had absorbed the lesson from their earlier participation in
mobilization politics that ultimately all power resided in the party,
not in mass organizations.

``山上派''遇到党内部分领导人最少反对的地方，是市一级而非省一级；在某些情况下，他们还遇到消极纵容甚至鼓励。翁及其同僚意识到，一旦他们开创先例，便既可以向下推进到县和工作单位一级，也可以向上推进到省一级。为此，他们建议中共浙江省委指示其在杭州的下属机构率先行动。\footnote{Ibid.
  同上。}
与此同时，他们试图通过破坏性活动使生产单位中的党支部运转瘫痪，然后让被围攻的领导人为由此造成的经济失败承担责任。这样，张、翁、贺制定了一套渗透与对抗并行的双管齐下策略。这种做法的新奇性，几年后由铁瑛在
1980
年的一次讲话中作了评论。铁瑛说：``几年前，张永生、翁森鹤、贺贤春之流，一面大骂共产党，一面又千方百计钻进我们党内。''\footnote{Tie
  Ying, ``共产党员要为四化建设的伟大事业而献身'' (Communist party
  members must contribute their all to the great cause of the four
  modernizations), ZJRB, May 31, 1980, pp.~1-3. Translated in ZPS, May
  31,1980, FBIS/CHI/108 (1980), O8.
  铁瑛：《共产党员要为四化建设的伟大事业而献身》（共产党员必须为四个现代化的伟大事业贡献一切），《浙江日报》，1980年5月31日，第1-3版。译文见ZPS，1980年5月31日，FBIS/CHI/108（1980），O8。}
造反派从他们早先参与动员政治的经历中吸取了教训：归根到底，一切权力都在党内，而不在群众组织中。

To make a breakthrough the rebels singled out three administrative
sections of the HMRC. The restoration of kou, which had taken place
early in 1973, was now characterized as a ``product of restraining
oneself and restoring the rites'' and an example of the ``two negations
and one reversal'' (两否一倒), a reference to the negation of both the
Cultural Revolution and its reforms (the so-called ``new-born things'')
and the reversal of the wheel of history. Members of the ``mountain
top'' faction demanded the abolition of three kou -- industry and
communications, finance and trade, and planning -- which had been carved
out of the Cultural Revolution production command group
(生产指挥组).\footnote{For a brief discussion of the amalgamation and
  carve-ups of ministries and agencies at the provincial level in the
  Cultural Revolution, see Harding, Organizing China, pp.~287-88. For
  Harding's discussion of the instrumental use of organizational issues
  in the period 1967-76, see ibid, pp.~341-44.
  关于文化大革命中省级部委和机构的合并与拆分的简要讨论，见Harding：《组织中国》，第287-288页。关于Harding对1967至1976年间工具性利用组织问题的讨论，见同书，第341-344页。}
The three particular kou were singled out because many of the old
production command group cadres had themselves been rebels and had been
demoted in the 1972 reshuffle. Additionally, the rebels were able to
enlist the support of two party secretaries of the HMC. On 4 and 7
February 1974, just as the anti-Lin anti-Confucius campaign was getting
underway, the leaders of the old production command group convened
meetings of its former cadre staff to prepare for a takeover.
Immediately after the second meeting a delegation carrying the seal of
the old group converged on a meeting place of the ZPC to obtain
authorization for the changeover. The party secretary concerned refused
and the cadres immediately declared a strike beginning at 9 pm in the
offices of the three kou. He Xianchun arrived to congratulate them on
their stand.

为了取得突破，造反派挑出了杭州市革委会的三个行政口。1973
年初恢复的``口''，如今被说成是``克己复礼的产物''，也是``两否一倒''的例子，意指否定文化大革命及其改革，即所谓``新生事物''，并且使历史车轮倒转。``山上派''成员要求撤销由文化大革命生产指挥组中划出的三个口，即工业交通、财贸和计划三个口。\footnote{For
  a brief discussion of the amalgamation and carve-ups of ministries and
  agencies at the provincial level in the Cultural Revolution, see
  Harding, Organizing China, pp.~287-88. For Harding's discussion of the
  instrumental use of organizational issues in the period 1967-76, see
  ibid, pp.~341-44.
  关于文化大革命中省级部委和机构的合并与拆分的简要讨论，见Harding：《组织中国》，第287-288页。关于Harding对1967至1976年间工具性利用组织问题的讨论，见同书，第341-344页。}
之所以专门挑出这三个口，是因为原生产指挥组的许多干部本身曾是造反派，并在
1972
年改组中被降职。此外，造反派还能够争取到杭州市委两名书记的支持。1974 年
2 月 4 日和 7
日，正当批林批孔运动展开之际，原生产指挥组领导人召集其原干部工作人员开会，为接管作准备。第二次会议刚结束，一个带着原指挥组印章的代表团就聚集到中共浙江省委的一处会场，要求批准变更。有关党委书记拒绝同意，这些干部随即宣布从晚上
9 点起在三个口的办公室罢工。贺贤春赶到，对他们的立场表示祝贺。

The next day, 8 February 1974, a notice appeared on the front gate of
the offices of the HMC compound notifying all staff of the strike, which
lasted five days. Wall posters and letters expressed solidarity with the
officials' bold action. Two secretaries of the HMC then applied pressure
on the municipal authorities to capitulate, which they did on 12
February. The three kou were disbanded and replaced by a single group
whose staff numbers were supplemented by the appointment of fifteen
cadres nominated by the ``mountain top'' faction. Xia Genfa, a close
ally of He Xianchun and Vice-chairman of the Hangzhou Oxygen Generator
Plant Revolutionary Committee, was appointed deputy leader of the
restored production command group. The hapless leaders of the three kou,
who had been deserted by their superiors, were shunted aside and placed
under investigation. Some were sent off to participate in supervised
labor. With the precedent set, a series of personnel changes swept the
province, spreading up to ZPC offices and down as far as factory, squad
and shift leading groups.\footnote{HZRB, November 23,1977; Chen Xia,
  HZRB, April 10,1978; Tie Ying, HZRB, June 7,1978.
  《杭州日报》，1977年11月23日；陈夏，《杭州日报》，1978年4月10日；铁瑛，《杭州日报》，1978年6月7日。}

次日，即 1974 年 2 月 8
日，杭州市委大院办公室正门上贴出一张通知，告知全体工作人员罢工，罢工持续了五天。墙报和来信对这些干部的大胆行动表示声援。随后，杭州市委两名书记向市级当局施压，迫使其让步；市级当局于
2 月 12
日屈服。三个口被解散，由一个单一小组取代，其工作人员又通过任命十五名由``山上派''提名的干部而得到补充。贺贤春的亲密盟友、杭州制氧机厂革命委员会副主任夏根发，被任命为恢复后的生产指挥组副组长。三个口中那些不幸的领导人被上级抛弃，被搁置一旁并接受审查。一些人被送去参加监督劳动。先例一开，一连串人事变动席卷全省，向上蔓延到中共浙江省委机关，向下延伸到工厂、班组和轮班领导小组。\footnote{HZRB,
  November 23,1977; Chen Xia, HZRB, April 10,1978; Tie Ying, HZRB, June
  7,1978.
  《杭州日报》，1977年11月23日；陈夏，《杭州日报》，1978年4月10日；铁瑛，《杭州日报》，1978年6月7日。}

By this action the radicals were demonstrating that they could force
party leaders to give ground. Additionally, they were flagging their
ability to dictate terms to the local cadres before the people of the
province, who would have been excused for wondering who really was
running the administration in Zhejiang. The humiliating cave-in
undoubtedly helped undermine the credibility of the Hangzhou
administration and its functionaries. The rebels realized that their
tactics of destabilization and confrontation had proved spectacularly
successful. Therefore, on 19 March, the rebels presented an officials'
list to the HMC (官员清单) containing 60 names.\footnote{Weng Senhe
  bufen zuizheng cailiao, pp.~21, 32.
  《翁森鹤部分罪证材料》，第21、32页。} The list had been drawn up in
conjunction with a document which Weng Senhe had drafted, entitled
``Recommendations on Solving the Problems of Zhejiang'' and it demanded
posts for members of Weng's faction ranging from membership of the
standing committee of the Hangzhou party committee through to various
groups under its jurisdiction. One department pinpointed for close
attention was the municipal Public Security Bureau. In March 1974 the
rebels had convened another meeting at Mt. Pingfeng where they discussed
how to gain control of this crucial ``organ of the proletarian
dictatorship''. The posts of bureau deputy director and ten section
heads had been included on the namelist. An ultimatum accompanying the
namelist allowed the HMC a mere three days in which to agree.

通过这一行动，激进派表明他们能够迫使党的领导人让步。此外，他们也在全省人民面前显示出自己有能力向地方干部发号施令；人们难免会怀疑，浙江的行政机构究竟由谁掌管。这次屈辱性的让步无疑削弱了杭州行政机构及其官员的公信力。造反派意识到，他们制造不稳定和对抗的策略已经取得了惊人的成功。因此，3
月 19 日，造反派向杭州市委提交了一份包含 60
个姓名的``官员清单''。\footnote{Weng Senhe bufen zuizheng cailiao,
  pp.~21, 32. 《翁森鹤部分罪证材料》，第21、32页。}
这份名单是配合翁森鹤起草的一份题为《关于解决浙江问题的建议》的文件拟定的，它要求为翁派成员安排职位，从杭州市委常委会成员一直到其管辖下各个小组的岗位不等。一个被重点盯上的部门是市公安局。1974
年 3
月，造反派又在屏风山召开会议，讨论如何控制这个关键的``无产阶级专政机关''。名单中包括了公安局副局长和十名科长的职位。随名单附上的最后通牒只给杭州市委三天时间同意。

By 28 March the HMC had again capitulated and 50 ``new-born forces''
were appointed to the leadership teams of its three major groups (组),
thirteen offices (办) and six bureaux (局). Nine of Weng and He's
nominees joined the standing committee of the CCP HMC.\footnote{HZRB,
  December 7,1977; ``Report on the Examination concerning Weng Senhe'',
  CCP CC Document (1977) No.~37, I\&S, 15:1, p.~111.
  《杭州日报》，1977年12月7日；《关于翁森鹤问题的审查报告》，中共中央文件（1977）第37号，I\&S，15:1，第111页。}
The appointment of He Xianchun as Deputy-secretary of the CCP HMC was
later graphically described in the following words. The standing
committee met with Second Secretary Wang Xing in the chair.

到 3 月 28 日，杭州市委再次让步，50
名``新生力量''被任命到其三大组、十三个办公室和六个局的领导班子中。翁、贺提名的人中有九人进入中共杭州市委常委会。\footnote{HZRB,
  December 7,1977; ``Report on the Examination concerning Weng Senhe'',
  CCP CC Document (1977) No.~37, I\&S, 15:1, p.~111.
  《杭州日报》，1977年12月7日；《关于翁森鹤问题的审查报告》，中共中央文件（1977）第37号，I\&S，15:1，第111页。}
贺贤春被任命为中共杭州市委副书记一事，后来被如下生动地描述。常委会开会时，由第二书记王兴主持。

\begin{quote}
{[}Wang{]} servilely asked He Xianchun. ``Will we put you in old (老)
He?'' He Xianchun had the nerve to say: ``that would be best. In the
future I might be working at the provincial level but to have a position
in the municipality as well would be good.'' {[}Wang{]} said, ``Fine.
Declare old He deputy party secretary of the municipal committee''. He
{[}Xianchun{]} also said ``promoting me without Weng Senhe is no good.
Put him down as well.'' {[}Wang{]} said: ``Ok. Declare Weng Senhe deputy
party secretary of the municipal committee''.\footnote{Chen Xia, HZRB,
  April 10,1978. There is no evidence that Weng Senhe in fact ever held
  the post of Deputy-secretary of the CCP HMC. Perhaps He, and Wang,
  were joking.
  陈夏，《杭州日报》，1978年4月10日。没有证据表明翁森鹤事实上曾任中共杭州市委副书记。或许何和王是在开玩笑。}
\end{quote}

\begin{quote}
{[}王{]}低声下气地问贺贤春：``把你放进去，老贺，行不行？''贺贤春竟然说：``那最好。将来我可能到省里工作，但在市里也有个位置也好。''{[}王{]}说：``好。宣布老贺任市委副书记。''贺{[}贤春{]}又说：``提拔我而没有翁森鹤不行。也把他写上。''{[}王{]}说：``好。宣布翁森鹤任市委副书记。''\footnote{Chen
  Xia, HZRB, April 10,1978. There is no evidence that Weng Senhe in fact
  ever held the post of Deputy-secretary of the CCP HMC. Perhaps He, and
  Wang, were joking.
  陈夏，《杭州日报》，1978年4月10日。没有证据表明翁森鹤事实上曾任中共杭州市委副书记。或许何和王是在开玩笑。}
\end{quote}

The appointment of Huang Yintang to positions in the municipal Public
Security Bureau and militia headquarters followed a similar
pattern.\footnote{Weng Senhe bufen zuizheng cailiao, p.~34.
  《翁森鹤部分罪证材料》，第34页。} At 7 pm on April 3, 1974, the
standing committee of the HMC met in Hangzhou with Weng Senhe and He
Xianchun in attendance. Huang's name was put forward and a brief resume
of his background read into the record. Huang was thirty-six years old
and had an unblemished political background and family connections. He
was proposed for the posts of Deputy-director of the Public Security
Bureau and deputy party secretary of the militia party branch. A brief
but illuminating exchange ensued between Weng and He on the one hand,
and an unnamed party official (possibly Wang Xing) on the other. The
party official asserted that Huang's priority lay in his public security
post (presumably where his enthusiasm for factionalism could be
contained) while Weng and He stressed Huang's post in the militia. Huang
later became effective head of the bureau after the director was driven
out of office. Thus, a rebel who had spent time in jail in the early
days of the Cultural Revolution had turned the tables on his former
prosecutors.\footnote{HZRB, June 23,1977. In November 1967 Huang had
  blindfolded a policeman on duty and taken him to the hills surrounding
  Hangzhou. Hangzhou Silk criticism group, ``Iron-clad proof of Weng
  Senhe and He Xianchun ganging up to usurp power in the silk
  industry'', HZRB, June 30,1977.
  《杭州日报》，1977年6月23日。1967年11月，黄曾蒙住一名值勤警察的眼睛，把他带到杭州周围的山里。杭州丝绸批判组：《翁森鹤、贺贤春在丝绸行业结伙篡权的铁证》，《杭州日报》，1977年6月30日。}
Huang, who had only just joined the CCP, immediately became a member of
the Public Security Bureau party committee. He reportedly acquired the
habit of roving around with blank detention forms and arrest warrants in
his pocket. All 300 rebels who ascended Mt. Pingfeng received
certificates of accreditation and promotion.\footnote{HZRB, June
  20,1977. 《杭州日报》，1977年6月20日。}

黄阴堂被任命到市公安局和民兵总部任职，也遵循了类似模式。\footnote{Weng
  Senhe bufen zuizheng cailiao, p.~34. 《翁森鹤部分罪证材料》，第34页。}
1974 年 4 月 3 日晚上 7
点，杭州市委常委会在杭州开会，翁森鹤、贺贤春出席。黄的名字被提出，并将其背景简历作了记录。黄三十六岁，政治背景和家庭关系清白。他被提议担任公安局副局长和民兵党支部副书记。随后，翁、贺二人与一名未具名的党内官员，可能是王兴，之间发生了一段简短但颇能说明问题的对话。该党内官员声称，黄的重点应放在公安岗位上，想必在那里可以约束他搞派性的热情；而翁、贺则强调黄在民兵中的职务。后来，公安局长被赶下台后，黄实际上成了该局负责人。于是，一个在文革初期曾坐过牢的造反派，反过来压倒了昔日追究他的人。\footnote{HZRB,
  June 23,1977. In November 1967 Huang had blindfolded a policeman on
  duty and taken him to the hills surrounding Hangzhou. Hangzhou Silk
  criticism group, ``Iron-clad proof of Weng Senhe and He Xianchun
  ganging up to usurp power in the silk industry'', HZRB, June 30,1977.
  《杭州日报》，1977年6月23日。1967年11月，黄曾蒙住一名值勤警察的眼睛，把他带到杭州周围的山里。杭州丝绸批判组：《翁森鹤、贺贤春在丝绸行业结伙篡权的铁证》，《杭州日报》，1977年6月30日。}
刚刚加入中共的黄，立即成为公安局党委委员。据说他养成了在口袋里揣着空白拘留证和逮捕证四处游走的习惯。所有登上屏风山的
300 名造反派都获得了资格认可和提拔证书。\footnote{HZRB, June 20,1977.
  《杭州日报》，1977年6月20日。}

The unorthodox promotion of cadres and the unusual measures used to
recruit party members were known as the ``two crashthroughs'' (双突) --
crash admission into the party and crash promotion of cadres
(突击发展党员,突击提拔干部).\footnote{These tactics were used in other
  provinces. See, Fenwick, ``The Gang of Four'', pp.~255-56.
  这些策略也在其他省份使用过。见Fenwick：《四人帮》，第255-256页。} To
justify their actions, the rebels fell back on Article 12 (5) of the CCP
statutes passed at the 10th Congress in 1973 which favoured ``getting
rid of the stale and taking in the fresh'' and Article 5 which
stipulated that

这种非正统的干部提拔和异常的党员吸收办法，被称为``双突''，即突击发展党员、突击提拔干部。\footnote{These
  tactics were used in other provinces. See, Fenwick, ``The Gang of
  Four'', pp.~255-56.
  这些策略也在其他省份使用过。见Fenwick：《四人帮》，第255-256页。}
为了使其行动正当化，造反派转而援引 1973
年十大通过的中共党章第十二条第五款，其中主张``吐故纳新''；以及第五条，其中规定：

\begin{quote}
The leading bodies of the party at all levels shall be elected through
democratic consultation in accordance with the requirements for
successors to the cause of the proletarian revolution {[}as stipulated
by Mao in 1964{]} and the principle of combining the old, the
middle-aged and the young.\footnote{``Constitution of the CCP'', The
  Tenth National Congress of the Communist Party of China (Documents),
  pp.~73, 67.
  《中国共产党章程》，载《中国共产党第十次全国代表大会文件》，第73、67页。}
\end{quote}

\begin{quote}
党的各级领导机关，应按照无产阶级革命事业接班人的条件{[}即毛在 1964
年所规定的条件{]}，并按照老、中、青三结合的原则，通过民主协商选举产生。\footnote{``Constitution
  of the CCP'', The Tenth National Congress of the Communist Party of
  China (Documents), pp.~73, 67.
  《中国共产党章程》，载《中国共产党第十次全国代表大会文件》，第73、67页。}
\end{quote}

Two methods for recruiting party members favoured in Zhejiang were known
as ``flying across the sea'' (飞过海) and ``appointment from above''
(上督落). The former method circumvented established procedures by
allowing candidates who faced opposition in their work-place to join at
other units. The Hangzhou workers' congress and militia headquarters
party branches proved ideal for this purpose. Where unit leaders
sympathetic to the rebels ran party branches, they simply appointed new
members.\footnote{Zhang Yongshengde zuizheng cailiao, p.~150; HZRB, June
  20,1977, June 23, 1977, December 7, 1977, I\&S, 14:11 (1978), p.~110.
  《张永生的罪证材料》，第150页；《杭州日报》，1977年6月20日、1977年6月23日、1977年12月7日；I\&S，14:11（1978），第110页。}

浙江偏好的两种吸收党员办法被称为``飞过海''和``上督落''。前一种方法绕开既定程序，允许在本单位受到反对的候选人到其他单位入党。杭州工代会和民兵总部党支部正好适合这一目的。在同情造反派的单位领导掌握党支部的地方，他们干脆直接任命新党员。\footnote{Zhang
  Yongshengde zuizheng cailiao, p.~150; HZRB, June 20,1977, June 23,
  1977, December 7, 1977, I\&S, 14:11 (1978), p.~110.
  《张永生的罪证材料》，第150页；《杭州日报》，1977年6月20日、1977年6月23日、1977年12月7日；I\&S，14:11（1978），第110页。}

It was claimed that in 1974 in Hangzhou alone, 8,000 new CCP members
were admitted to the party and 3,000 cadres promoted by these highly
unorthodox means.\footnote{In Lanxi county the figures were 277 and 94
  respectively. See Lanxi shizhi, p.~22.
  在兰溪县，这两个数字分别为277和94。见《兰溪市志》，第22页。} Included
among these new party members and ``helicopter'' cadres were alleged
hooligans and thugs, vandals, criminals and those with a family grudge
(世仇分子). Some people entered the party one day and were promoted as
officials the next. Some were appointed secretaries of party committees
before they had joined the CCP. The most trusted of Zhang, Weng and He's
followers joined the campaign small groups but they also placed them in
such sensitive areas as the provincial party office, defence department
office, provincial education office and Zhejiang Daily.\footnote{Weng
  Senhe bufen zuizheng cailiao, p.~33; Zhang Yongshengde zuizheng
  cailiao, p.~83.
  《翁森鹤部分罪证材料》，第33页；《张永生的罪证材料》，第83页。} With
Shen Ce and Weng Senhe in charge of the provincial political work group,
which had the power to approve all such appointments, there were few
institutional obstacles in the rebels' path. Zhang Yongsheng was also
present at meetings of the group where such promotions were
discussed.\footnote{Zhang Yongshengde zuizheng cailiao, p.~88.
  《张永生的罪证材料》，第88页。}

据称，仅 1974 年在杭州一地，就有 8,000
名中共新党员通过这些极不正统的方式入党，3,000
名干部获得提拔。\footnote{In Lanxi county the figures were 277 and 94
  respectively. See Lanxi shizhi, p.~22.
  在兰溪县，这两个数字分别为277和94。见《兰溪市志》，第22页。}
这些新党员和``直升机''干部中，据说包括流氓、打手、破坏分子、罪犯和世仇分子。有些人头一天入党，第二天就被提拔为干部。有些人在加入中共之前就被任命为党委书记。张、翁、贺最信任的追随者进入了运动小组，他们还把这些人安插到省委办公厅、国防部门办公室、省教育办和《浙江日报》等敏感领域。\footnote{Weng
  Senhe bufen zuizheng cailiao, p.~33; Zhang Yongshengde zuizheng
  cailiao, p.~83.
  《翁森鹤部分罪证材料》，第33页；《张永生的罪证材料》，第83页。}
省政治工作组有权批准所有这类任命，而沈策和翁森鹤掌管该组，因此造反派面前几乎没有制度性障碍。张永生也出席了讨论这类提拔的该组会议。\footnote{Zhang
  Yongshengde zuizheng cailiao, p.~88. 《张永生的罪证材料》，第88页。}

The CCP ZPC seemed powerless to stop the flouting of party statutes and
cadre regulations. Weng Senhe, in his inimitable style, used the
following analogy to describe the dispatch of new cadres into positions
of power: ``Just as hens lay chicks so we must send out batch after
batch''.\footnote{Hangzhou Garrison criticism group, ``A black model in
  establishing a 'second armed force''', HZRB, April 20, 1977.
  杭州警备区批判组：《建立``第二武装''的黑样板》，《杭州日报》，1977年4月20日。}
Zhang Yongsheng established a test case at his unit, the Zhejiang Fine
Arts College. On 19 May 1974, which was the sixth anniversary of his
meeting with Jiang Qing in Beijing (see chapter three), Zhang wrote
Jiang a pledge of continuing loyalty (决心书).\footnote{For the use of
  ritual symbols to reinforce factional ties, see Pye, The Dynamics of
  Chinese Politics., pp.~10-11.
  关于利用仪式性符号强化派系纽带，见Pye：《中国政治的动力》，第10-11页。}
In one month he appointed ten people to the campaign small group in the
college and then, in July, reorganized its party committee, assuming the
post of party secretary himself with his old comrade Du Yingxin as
deputy party secretary. Zhang also took control of the provincial
revolution-in-education group and implemented the ``two crashthroughs''
in educational institutions throughout Zhejiang.\footnote{Zhang
  Yongshengde zuizheng cailiao, p.~90; Zhejiang Fine Arts College
  criticism group, ``The pipedream of the Gang of Four and the mongrel
  of Zhejiang'', HZRB, January 23,1977; RMRB, December 2, 1977. For a
  further report on the activities of the College in the campaign, see
  ZJRB, August 24,1974, p.~1.
  《张永生的罪证材料》，第90页；浙江美术学院批判组：《``四人帮''的黄粱梦和浙江的杂种》，《杭州日报》，1977年1月23日；《人民日报》，1977年12月2日。关于该学院在运动中活动的进一步报道，见《浙江日报》，1974年8月24日，第1版。}

中共浙江省委似乎无力阻止这种对党章和干部规定的践踏。翁森鹤以其独特风格，用如下比喻来描述把新干部派往权力岗位：``就像母鸡下小鸡一样，我们要一批一批地派出去。''\footnote{Hangzhou
  Garrison criticism group, ``A black model in establishing a 'second
  armed force''', HZRB, April 20, 1977.
  杭州警备区批判组：《建立``第二武装''的黑样板》，《杭州日报》，1977年4月20日。}
张永生在自己的单位浙江美术学院树立了一个试点。1974 年 5 月 19
日，也就是他在北京会见江青六周年之日，见第三章，张给江写了一份继续效忠的决心书。\footnote{For
  the use of ritual symbols to reinforce factional ties, see Pye, The
  Dynamics of Chinese Politics., pp.~10-11.
  关于利用仪式性符号强化派系纽带，见Pye：《中国政治的动力》，第10-11页。}
一个月内，他任命十人进入学院运动小组，随后在 7
月改组学院党委，由自己担任党委书记，并让老战友杜英信担任党委副书记。张还控制了省教育革命小组，并在浙江全省教育机构中推行``双突''。\footnote{Zhang
  Yongshengde zuizheng cailiao, p.~90; Zhejiang Fine Arts College
  criticism group, ``The pipedream of the Gang of Four and the mongrel
  of Zhejiang'', HZRB, January 23,1977; RMRB, December 2, 1977. For a
  further report on the activities of the College in the campaign, see
  ZJRB, August 24,1974, p.~1.
  《张永生的罪证材料》，第90页；浙江美术学院批判组：《``四人帮''的黄粱梦和浙江的杂种》，《杭州日报》，1977年1月23日；《人民日报》，1977年12月2日。关于该学院在运动中活动的进一步报道，见《浙江日报》，1974年8月24日，第1版。}

A report of 30 June 1974 referred to the light industrial sector as a
model in carrying out the policy of inducting new members into the CCP
``in accordance with the provisions of the new party
constitution''.\footnote{ZPS, June 30, 1974, FBIS/CHI, 135 (1974), G3-4.
  ZPS，1974年6月30日，FBIS/CHI，135（1974），G3-4。} Weng Senhe, in his
capacity as member of the ZPRC standing committee, was responsible for
the sector. The party newcomers were described as ``advanced proletarian
elements who have emerged from class struggle and the struggle between
the two lines'' and ``have become imbued with the revolutionary spirit
of daring to go against the tide''. The article claimed that the
greenhorns worked ``closely together with veteran party members'' and
the people called them ``revolutionary young tigers''.

1974 年 6 月 30
日的一份报告把轻工业部门称为``按照新党章规定''吸收中共新党员政策的典型。\footnote{ZPS,
  June 30, 1974, FBIS/CHI, 135 (1974), G3-4.
  ZPS，1974年6月30日，FBIS/CHI，135（1974），G3-4。}
翁森鹤以浙江省革委会常委身份负责该部门。这些新党员被描述为``在阶级斗争和两条路线斗争中涌现出来的无产阶级先进分子''，并且``具有敢于反潮流的革命精神''。文章声称，这些新手同``老党员密切合作''，群众称他们为``革命小老虎''。

\subsection{2. The Urban Militia}\label{the-urban-militia}

\subsection{2. 城市民兵}\label{ux57ceux5e02ux6c11ux5175}

On 29 September 1973 a joint editorial of the People's Daily and
Liberation Army Daily had called for a build-up of urban militia forces
along the lines of contingents that had been experimented with in the
mid-1960s. Articles which accompanied the editorial described the
activities of the urban militia in Beijing, Shanghai, and Tianjin. In
these cities the militia had three major tasks - to participate in class
struggle, to carry out patrol and guard duties, and to defend the
population in times of war.\footnote{``办好民兵'' (Run the militia
  well), RMRB, September 29,1973, p.~1. For comments on this editorial
  see CS, 12:1(1974), pp.~21-4. For an article on the activities of the
  militia in Beijing and Shanghai, see
  ``在毛主席革命路线指引下进一步加强民兵建设'' (Under the guidance of
  Chairman Mao's revolutionary line further strengthen the building of
  the militia), RMRB, September 29,1973, pp.~1-2. The first issue of the
  Shanghai-based radical mouthpiece Study and Criticism published an
  article calling for the arming of the working class. See Song Minbin,
  ``巴黎公社与工人武装'' (The Paris Commune and the armed working
  class), Xuexi yu pipan, 1 (September 15,1973), pp.~38-43. See the
  discussion of the article in CNS, 494 (November 22,1973), pp.~1-9. For
  more information on this monthly journal see Ting Wang, ``Propaganda
  and Political Struggle: A Preliminary Case Study of the `Hsueh yu
  p'ip'an'\,'', I\&S, 13: 6 (1977), pp.~1-14. For background to the
  activities of the militia in the Cultural Revolution see Gao Gao and
  Yan Jiaqi,''Wenhua dageming'', pp.~562-65; and in Shanghai, Goodman,
  ``The Shanghai Connection'', pp.~137-38.
  《办好民兵》（把民兵办好），《人民日报》，1973年9月29日，第1版。关于这篇社论的评论，见CS，12:1（1974），第21-24页。关于北京和上海民兵活动的文章，见《在毛主席革命路线指引下进一步加强民兵建设》（在毛主席革命路线指引下进一步加强民兵建设），《人民日报》，1973年9月29日，第1-2版。以上海为基地的激进派喉舌《学习与批判》创刊号发表了一篇呼吁武装工人阶级的文章。见宋敏彬：《巴黎公社与工人武装》（巴黎公社与武装的工人阶级），《学习与批判》，第1期（1973年9月15日），第38-43页。关于该文的讨论，见CNS，494（1973年11月22日），第1-9页。关于这份月刊的更多资料，见Ting
  Wang：《宣传与政治斗争：``学习与批判''初步个案研究》，I\&S，13:6（1977），第1-14页。关于文化大革命中民兵活动的背景，见高皋、严家其：《文化大革命》，第562-565页；关于上海，见Goodman：《上海联系》，第137-138页。}
The description of the tasks was a euphemistic way of stating that the
urban militia would take a partisan stand in local factional disputes,
that it would take over guard duties from public security bureaux and
the military and that it would challenge the PLA's right to sole
possession and use of modern armaments.

1973 年 9 月 29 日，《人民日报》和《解放军报》联合社论号召仿照 1960
年代中期试点的分队形式，建立城市民兵力量。配合社论刊发的文章描述了北京、上海和天津城市民兵的活动。在这些城市，民兵有三项主要任务：参加阶级斗争，执行巡逻和警卫职责，以及在战时保卫人民。\footnote{``办好民兵''
  (Run the militia well), RMRB, September 29,1973, p.~1. For comments on
  this editorial see CS, 12:1(1974), pp.~21-4. For an article on the
  activities of the militia in Beijing and Shanghai, see
  ``在毛主席革命路线指引下进一步加强民兵建设'' (Under the guidance of
  Chairman Mao's revolutionary line further strengthen the building of
  the militia), RMRB, September 29,1973, pp.~1-2. The first issue of the
  Shanghai-based radical mouthpiece Study and Criticism published an
  article calling for the arming of the working class. See Song Minbin,
  ``巴黎公社与工人武装'' (The Paris Commune and the armed working
  class), Xuexi yu pipan, 1 (September 15,1973), pp.~38-43. See the
  discussion of the article in CNS, 494 (November 22,1973), pp.~1-9. For
  more information on this monthly journal see Ting Wang, ``Propaganda
  and Political Struggle: A Preliminary Case Study of the `Hsueh yu
  p'ip'an'\,'', I\&S, 13: 6 (1977), pp.~1-14. For background to the
  activities of the militia in the Cultural Revolution see Gao Gao and
  Yan Jiaqi,''Wenhua dageming'', pp.~562-65; and in Shanghai, Goodman,
  ``The Shanghai Connection'', pp.~137-38.
  《办好民兵》（把民兵办好），《人民日报》，1973年9月29日，第1版。关于这篇社论的评论，见CS，12:1（1974），第21-24页。关于北京和上海民兵活动的文章，见《在毛主席革命路线指引下进一步加强民兵建设》（在毛主席革命路线指引下进一步加强民兵建设），《人民日报》，1973年9月29日，第1-2版。以上海为基地的激进派喉舌《学习与批判》创刊号发表了一篇呼吁武装工人阶级的文章。见宋敏彬：《巴黎公社与工人武装》（巴黎公社与武装的工人阶级），《学习与批判》，第1期（1973年9月15日），第38-43页。关于该文的讨论，见CNS，494（1973年11月22日），第1-9页。关于这份月刊的更多资料，见Ting
  Wang：《宣传与政治斗争：``学习与批判''初步个案研究》，I\&S，13:6（1977），第1-14页。关于文化大革命中民兵活动的背景，见高皋、严家其：《文化大革命》，第562-565页；关于上海，见Goodman：《上海联系》，第137-138页。}
对这些任务的描述，是一种委婉说法，意思是城市民兵将在地方派性争端中采取派别立场，将从公安局和军队手中接管警卫职责，并挑战解放军独占和使用现代武器的权利。

The significance of the Shanghai militia experience (under Wang
Hongwen's direction) was based on several factors. First, Wang and Ni
Zhifu, who commanded the Beijing urban militia, did so because they held
the post of Chairman of the municipal trade union council (or workers'
congress). The urban militia was therefore placed under union
leadership. Second, the urban militia commands were responsible to the
municipal party committees (and ultimately to provincial party
committees), and not the military districts, which had previously
organized and trained such forces. Finally, Shanghai emphasized an
ongoing program of political education and study for the enlisted
workers so they could raise their political awareness.\footnote{Yang
  Yong (deputy Chief of Staff of the PLA), ``Report to the National
  Militia Conference'', Xinhua, August 8,1978, SWB/FE/5890/BII/1-3; J.
  C. F. Wang,''The Urban Militia as a political instrument in the power
  contest in China in 1976'', AS, 18: 6 (1978), pp.~548-52; CNA, 940
  (November 16,1973), pp.~5-7; Alan P. L. Liu, ``The `Gang of Four' and
  the Chinese People's Liberation Army'', AS, 19: 9 (1979), pp.~829-33.
  杨勇（中国人民解放军副总参谋长）：《在全国民兵会议上的报告》，新华社，1978年8月8日，SWB/FE/5890/BII/1-3；J.
  C. F.
  Wang：《城市民兵作为1976年中国权力斗争中的政治工具》，AS，18:6（1978），第548-552页；CNA，940（1973年11月16日），第5-7页；Alan
  P. L.
  Liu：《``四人帮''与中国人民解放军》，AS，19:9（1979），第829-833页。}

上海民兵经验，在王洪文指导下，其重要性建立在几个因素之上。第一，王洪文以及指挥北京城市民兵的倪志福之所以掌握民兵，是因为他们担任市总工会或工代会主席。因此，城市民兵被置于工会领导之下。第二，城市民兵指挥部对市委负责，并最终对省委负责，而不是对此前曾组织和训练这类力量的军区负责。最后，上海强调对参加民兵的工人持续开展政治教育和学习，以提高他们的政治觉悟。\footnote{Yang
  Yong (deputy Chief of Staff of the PLA), ``Report to the National
  Militia Conference'', Xinhua, August 8,1978, SWB/FE/5890/BII/1-3; J.
  C. F. Wang,''The Urban Militia as a political instrument in the power
  contest in China in 1976'', AS, 18: 6 (1978), pp.~548-52; CNA, 940
  (November 16,1973), pp.~5-7; Alan P. L. Liu, ``The `Gang of Four' and
  the Chinese People's Liberation Army'', AS, 19: 9 (1979), pp.~829-33.
  杨勇（中国人民解放军副总参谋长）：《在全国民兵会议上的报告》，新华社，1978年8月8日，SWB/FE/5890/BII/1-3；J.
  C. F.
  Wang：《城市民兵作为1976年中国权力斗争中的政治工具》，AS，18:6（1978），第548-552页；CNA，940（1973年11月16日），第5-7页；Alan
  P. L.
  Liu：《``四人帮''与中国人民解放军》，AS，19:9（1979），第829-833页。}

The factional nature of the urban militia which was established in
Hangzhou added a further dimension to divisions within the city's
working class. The practice of allowing workers who were regularly
absent from the shop-floor to draw full pay as well as subsidies for
their militia service aroused resentment from other workers who remained
at their posts and in some cases were victims of the militia. The
establishment of a party branch at militia headquarters enabled the
militia to recruit new party members who for various reasons could not
join the CCP at their own factories. The militia thus served as a
springboard to launch the careers of several controversial figures in
Hangzhou.

杭州建立的城市民兵具有派性性质，这使该市工人阶级内部的分裂又增添了一个层面。那些经常离开车间的工人，既领取全额工资，又因民兵服务获得补贴，这种做法引起了仍坚守岗位的其他工人的不满；其中有些人还成了民兵的受害者。民兵总部设立党支部，使民兵能够吸收那些因各种原因无法在自己工厂入党的新党员。于是，民兵成了杭州若干有争议人物仕途起步的跳板。

Xia Qi, Deputy Political Commissar of the ZPMD, may well have tried to
delay or resist the formation of urban militia commands in Zhejiang. A
commentary which was broadcast over the provincial radio on the eve of
the publication of the People's Daily's 29 September 1973 editorial,
emphasized the more traditional view of the militia as a production and
military force.\footnote{ZJRB, September 28,1973, p.~1.
  《浙江日报》，1973年9月28日，第1版。} This stand contrasted strongly
with the radical position which was enunciated in the People's Daily
editorial on the following day. In October 1973 the ZPMD, under the
direction of the CCP ZPC, called a meeting to discuss the new
orientation in militia work. Xia Qi delivered a mobilization report to
an audience which was composed of representatives of county-level PLA
units. There is little evidence that at the October meeting or in a
December speech reviewing the year's work in the military that Xia
openly expressed reservations about the direction in which militia work
was heading. The only hint of difference was the attempt to place class
struggle in the context of the two other ``revolutionary movements'' --
the struggle for production and scientific achievement.\footnote{ZJRB,
  October 18,1973; December 21,1973.
  《浙江日报》，1973年10月18日；1973年12月21日。}

浙江省军区副政委夏琦很可能曾试图拖延或抵制在浙江建立城市民兵指挥部。1973
年 9 月 29
日《人民日报》社论发表前夕，省电台播发的一篇评论，强调了把民兵视为生产和军事力量的较为传统的看法。\footnote{ZJRB,
  September 28,1973, p.~1. 《浙江日报》，1973年9月28日，第1版。}
这一立场同次日《人民日报》社论所阐明的激进立场形成鲜明对比。1973 年 10
月，浙江省军区在中共浙江省委领导下召开会议，讨论民兵工作的新方向。夏琦向由县级解放军单位代表组成的听众作了动员报告。几乎没有证据表明，夏在
10 月会议上，或在 12
月总结全年军事工作的讲话中，公开表达过对民兵工作走向的保留意见。唯一显示差异的线索，是他试图把阶级斗争置于另外两个``革命运动''即生产斗争和科学实验的背景之中。\footnote{ZJRB,
  October 18,1973; December 21,1973.
  《浙江日报》，1973年10月18日；1973年12月21日。}

If Xia's intentions were to go against the radical tide on the issue of
the militia he was ultimately unsuccessful. The establishment of the
HMMCH took place on 14 February, although the public announcement did
not appear until eight days later.\footnote{ZJRB, February 22,1974,
  p.~1. 《浙江日报》，1974年2月22日，第1版。} The notice in Zhejiang
Daily stated that with the approval of the ZPC, the CCP HMC had decided
to form the militia command headquarters. The notice added that the
Hangzhou urban militia would base itself on the experience of its
counterpart in Shanghai, particularly in relation to participation in
class struggle. In trying to explain why party leaders had consented to
the establishment of the urban militia, the local press later claimed
that they had been coerced into giving approval. The Hangzhou Daily
asserted that on 14 February 1974 Weng and a group of his followers had
gone to the office of the ZPC and ``like poisonous snakes entwined
themselves around the leaders of the ZPC and ordered them to answer
questions, listen to ideas, and express their opinions and then forced
them to sign things''. The rebels had allegedly assaulted PLA guards on
duty at the building and then accused them of starting the fracas. Xia
Genfa told a Xinhua reporter at the scene: ``If the ZPC doesn't
recognize its mistakes, in three days the working class of Hangzhou will
give it something to think about.''\footnote{Hangzhou Oxygen Generator
  Plant, ``Resolutely suppress the ringleaders of beating, smashing and
  looting'', HZRB, June 21,1977; ZPMD Logistics Department criticism
  group, ``Rise to power and positions, and fall to destruction amid
  chaos'', HZRB, April 9, 1977.
  杭州制氧机厂：《坚决镇压打砸抢的为首分子》，《杭州日报》，1977年6月21日；浙江省军区后勤部批判组：《在乱中上台上位，在乱中走向灭亡》，《杭州日报》，1977年4月9日。}
In 1975, after Weng Senhe's arrest, one of his colleagues alleged that
Weng had boasted of how he had coerced Hangzhou's First Secretary Wang
Zida into agreeing to the militia's establishment. Much more probable is
Weng's confession that senior members of the HMC favored the formation
of the organization.\footnote{Weng Senhe bufen zuizheng cailiao, pp.~20,
  50. 《翁森鹤部分罪证材料》，第20、50页。}

如果夏在民兵问题上的意图是逆激进潮流而行，那么他最终没有成功。杭州民兵指挥部于
2 月 14 日成立，尽管公开公告直到八天后才出现。\footnote{ZJRB, February
  22,1974, p.~1. 《浙江日报》，1974年2月22日，第1版。}
《浙江日报》的通知称，经中共浙江省委批准，中共杭州市委决定成立民兵指挥部。通知补充说，杭州城市民兵将以其上海对应组织的经验为依据，尤其是在参加阶级斗争方面。后来，地方报刊在试图解释为何党的领导人同意建立城市民兵时，声称他们是被胁迫批准的。《杭州日报》断言，1974
年 2 月 14
日，翁和一群追随者前往中共浙江省委办公室，``像毒蛇一样缠住省委领导，逼他们回答问题、听取意见、表明态度，然后又逼他们签字''。据称，造反派还殴打了在该楼值勤的解放军警卫，随后反过来指责警卫挑起事端。夏根发在现场对一名新华社记者说：``如果省委不承认错误，三天内杭州工人阶级就要给它一点颜色看看。''\footnote{Hangzhou
  Oxygen Generator Plant, ``Resolutely suppress the ringleaders of
  beating, smashing and looting'', HZRB, June 21,1977; ZPMD Logistics
  Department criticism group, ``Rise to power and positions, and fall to
  destruction amid chaos'', HZRB, April 9, 1977.
  杭州制氧机厂：《坚决镇压打砸抢的为首分子》，《杭州日报》，1977年6月21日；浙江省军区后勤部批判组：《在乱中上台上位，在乱中走向灭亡》，《杭州日报》，1977年4月9日。}
1975
年翁森鹤被捕后，他的一名同僚声称翁曾吹嘘自己如何胁迫杭州第一书记王子达同意建立民兵。更可能的情况是，翁供认杭州市委高级成员赞成成立这一组织。\footnote{Weng
  Senhe bufen zuizheng cailiao, pp.~20, 50.
  《翁森鹤部分罪证材料》，第20、50页。}

The establishment of urban militia headquarters was party policy. Not
all major cities implemented the policy and some were reluctant to do
so, judging by the tardiness of their response. That Hangzhou
established its militia HQ promptly was a sign either that it wished to
display its enthusiasm for the policy, and by implication other radical
initiatives, or more likely that Wang Hongwen personally intervened to
ensure the speedy implementation of central policy. The party leaders in
Zhejiang found little room for manoeuvre with such close and constant
scrutiny of their performance.

建立城市民兵总部是党的政策。并非所有大城市都执行了这一政策；从反应迟缓来看，有些城市并不情愿。杭州迅速建立民兵总部，要么表明它希望显示自己对该政策以及由此引申的其他激进举措的热情；更可能的是，王洪文亲自干预，以确保中央政策迅速落实。在对他们表现如此密切而持续的监督之下，浙江的党的领导人几乎没有回旋余地。

All fifty-nine counties and municipalities in Zhejiang established
militia commands. Hangzhou's urban militia had a paper strength of
50,000 men. It set up office in the Workers' Cultural Palace, a
centrally located building in downtown Hangzhou, with an initial staff
of 70 personnel seconded from Weng's silk complex and the Hangzhou
Oxygen Generator Plant. Operating expenses were initially provided by
the municipal government.\footnote{Ibid, p.~49, 37. 同上，第49、37页。}
The militia was mobilized 10,000 men/times (人次) and detained 1,000
people. In the twelve or so months of its existence it was involved in
thirteen factional fights in the city.\footnote{Violence was also
  reported at this time in Jiangxi and Jiangsu provinces. See FEER, July
  1, 1974, pp.~16-17 and July 22,1974, pp.~13-14.
  此时江西和江苏两省也有暴力事件报道。见FEER，1974年7月1日，第16-17页；1974年7月22日，第13-14页。}
The HMMCH ran at least three detention centers, charmingly named
``processing factories'' (加工厂), where it confined, interrogated and
tortured its victims, principally from the ``mountain base'' faction.
Captives were forced to stand on one leg for extended periods of time
(金鸡独立), stand in the sun, or were hung by the fingers, and savagely
beaten. Some internees suffered permanent incapacitation while others
died. Wearing cane hats and armed with steel fashioned from five tonnes
of steel piping requisitioned from the Hangzhou Iron and Steel Mill,
militiamen soon became a common and dreaded sight on the streets of
Hangzhou.

浙江全省五十九个县市都建立了民兵指挥机构。杭州城市民兵账面兵力为 50,000
人。它把办公室设在位于杭州市中心的工人文化宫，最初有 70
名工作人员，由翁所在的丝绸联合厂和杭州制氧机厂借调而来。运行经费最初由市政府提供。\footnote{Ibid,
  p.~49, 37. 同上，第49、37页。} 民兵被动员达 10,000 人次，并拘押了
1,000
人。在它存在的大约十二个月里，参与了市内十三起派性斗殴。\footnote{Violence
  was also reported at this time in Jiangxi and Jiangsu provinces. See
  FEER, July 1, 1974, pp.~16-17 and July 22,1974, pp.~13-14.
  此时江西和江苏两省也有暴力事件报道。见FEER，1974年7月1日，第16-17页；1974年7月22日，第13-14页。}
杭州民兵指挥部至少设有三处拘留中心，美其名曰``加工厂''，在那里关押、审讯并折磨受害者，主要是``山下派''成员。被俘者被迫长时间金鸡独立、站在太阳下，或被吊手指，并遭到毒打。一些被拘押者终身残疾，另一些人死亡。民兵头戴藤帽，手持用从杭州钢铁厂征用的五吨钢管制成的钢器，很快成为杭州街头常见而令人畏惧的景象。

Contingents of militiamen, some bearing arms, were also ferried around
in trucks, whose wailing sirens sounded like fire-engines (one of the
militia's functions was to act as a fire-fighting unit). Some militiamen
were sent out to infiltrate crowds which gathered before wall posters.
Anyone who pasted up a poster which contained unflattering remarks about
the militia would have a chalk mark applied to the back so that he or
she could be identified, detained and beaten.\footnote{See Weng Senhe
  bufen zuizheng cailiao, pp.~58-9, for an example of a namelist and
  arrest warrant and p.~62 for a photograph of some of the militia
  command's instruments of torture.
  名单和逮捕证的例子见《翁森鹤部分罪证材料》，第58-59页；民兵指挥部部分刑具的照片见第62页。}
Other militiamen in civilian clothes, carrying concealed iron clubs,
beat anyone who expressed reservations about the militia. A demobilized
soldier at the Hangzhou automobile engine plant was thrashed and
permanently crippled for condemning the use of physical violence.

成队的民兵，其中有些携带武器，也被卡车来回运送；卡车刺耳的警报声听起来像消防车，民兵的职能之一就是充当消防队。一些民兵被派去混入聚集在墙报前的人群。凡是张贴了含有不利于民兵言论的墙报者，背上会被划上粉笔记号，以便识别、拘押和殴打。\footnote{See
  Weng Senhe bufen zuizheng cailiao, pp.~58-9, for an example of a
  namelist and arrest warrant and p.~62 for a photograph of some of the
  militia command's instruments of torture.
  名单和逮捕证的例子见《翁森鹤部分罪证材料》，第58-59页；民兵指挥部部分刑具的照片见第62页。}
其他身穿便衣、暗藏铁棍的民兵，则殴打任何对民兵表示保留意见的人。杭州汽车发动机厂一名复员军人因谴责使用肉体暴力而遭痛打，并永久残疾。

He Xianchun, as Chairman of the Hangzhou workers' congress,
automatically became commander of the city's militia. He later handed
over the position to his deputy Xia Genfa from the oxygen generator
factory. The old factory militia units which had been commanded and
trained by the ZPMD were scrapped, on paper at least. He Xianchun's
militia recruited ``fighters going against the tide'' as its backbone
with the object of building a countervailing armed force to the PLA. He
expressed the view that ``at the moment a group of old rightists hold
all military power'' and therefore ``we can't give power over the
militia to the ZPMD''. He was advised not to report back to leaders of
the ZPC because ``you don't know when they will approve''.\footnote{HZRB,
  December 31,1976, April 20,1977, June 21,1977, December 7,1977, April
  5,1978, August 14,1978, April 17,1979; ``\,`Struggling for truth' -
  the heroic deeds of CCP member Zhang Jifa in struggling against the
  Gang of Four'', HZRB, November 7,1978; Xinhua, January 13,1978,
  SWB/FE/7515/BII/15-16.
  《杭州日报》，1976年12月31日、1977年4月20日、1977年6月21日、1977年12月7日、1978年4月5日、1978年8月14日、1979年4月17日；《``为真理而斗争''--共产党员张纪法同``四人帮''斗争的英雄事迹》，《杭州日报》，1978年11月7日；新华社，1978年1月13日，SWB/FE/7515/BII/15-16。}
In a speech of 17 October 1974, Weng Senhe stated that the militia's
role was to seize military power in the province.\footnote{Weng Senhe
  bufen zuizheng cailiao, p.~55. 《翁森鹤部分罪证材料》，第55页。} With
the ZPC on the defensive and the municipal authorities either unwilling
to resist the rebels, or in collusion with them, the stage was set for a
series of incidents which seriously disrupted public order and further
polarized the city's working class.

贺贤春作为杭州工代会主席，自动成为该市民兵司令。后来他把这一职位交给了来自制氧机厂的副手夏根发。原先由浙江省军区指挥和训练的工厂民兵单位，至少在纸面上被废止。贺贤春的民兵以``反潮流战士''为骨干，目标是建立一支能够制衡解放军的武装力量。他表示，``现在一批老右派掌握全部军权''，因此``我们不能把民兵权交给省军区''。有人建议他不要向中共浙江省委领导汇报，因为``你不知道他们什么时候会批准''。\footnote{HZRB,
  December 31,1976, April 20,1977, June 21,1977, December 7,1977, April
  5,1978, August 14,1978, April 17,1979; ``\,`Struggling for truth' -
  the heroic deeds of CCP member Zhang Jifa in struggling against the
  Gang of Four'', HZRB, November 7,1978; Xinhua, January 13,1978,
  SWB/FE/7515/BII/15-16.
  《杭州日报》，1976年12月31日、1977年4月20日、1977年6月21日、1977年12月7日、1978年4月5日、1978年8月14日、1979年4月17日；《``为真理而斗争''--共产党员张纪法同``四人帮''斗争的英雄事迹》，《杭州日报》，1978年11月7日；新华社，1978年1月13日，SWB/FE/7515/BII/15-16。}
在 1974 年 10 月 17
日的一次讲话中，翁森鹤说，民兵的作用是在全省夺取军权。\footnote{Weng
  Senhe bufen zuizheng cailiao, p.~55. 《翁森鹤部分罪证材料》，第55页。}
中共浙江省委处于守势，而市级当局要么不愿抵制造反派，要么与他们合谋，于是一系列严重扰乱公共秩序并进一步分化城市工人阶级的事件，便具备了发生的条件。

Before the establishment of the militia, incidents involving workers and
PLA guards had occurred in Hangzhou. In December 1973 roving bands
insulted and beat PLA guards on duty outside provincial offices, hotels
and jails. Some soldiers had their collar tabs and insignia torn off and
their hats knocked off their heads. One such incident resulted in
injuries to twenty-three soldiers. It was later claimed that not one
soldier retaliated.\footnote{HZRB, March 10, 1977.
  《杭州日报》，1977年3月10日。} By encouraging the provocations, Weng
and He apparently aimed to prove that the PLA was incapable of
maintaining public order and that an urban militia was thus
indispensable for the task.

在民兵建立之前，杭州已经发生过涉及工人和解放军警卫的事件。1973 年 12
月，流动团伙辱骂并殴打在省级机关、宾馆和监狱外值勤的解放军警卫。一些士兵的领章和标志被扯下，帽子被打落。其中一起事件造成二十三名士兵受伤。后来有人声称，没有一名士兵还手。\footnote{HZRB,
  March 10, 1977. 《杭州日报》，1977年3月10日。}
翁、贺显然是通过鼓励这些挑衅，意在证明解放军无力维持公共秩序，因此城市民兵对这一任务不可或缺。

The militia aimed to put into operation Weng Senhe's maxim that ``since
ancient times every victory of a political group has been accompanied by
the suppression of its opponents''. On 21 February 1974, one day before
the public announcement of the establishment of the HMMCH, the militia
had its first taste of action in the downtown area of the provincial
capital. Weng and He's factional opponents from the ``mountain base''
faction had gathered at three hostels near the commercial center of
Hangzhou. He Xianchun and Xia Genfa despatched militiamen to
investigate, with resultant disruption to public transport and social
order. At 5 pm on the same day He and Xia issued an urgent notice
ordering the militia to clear its opponents out of the hostels. Workers
were mobilized from twelve factories and a large-scale confrontation
ensued. One hundred participants and an unknown number of bystanders
were injured in the melee.\footnote{Weng Senhe bufen zuizheng cailiao,
  p.~57. 《翁森鹤部分罪证材料》，第57页。} Property in the two hostels
valued at 100,000 yuan was damaged or smashed. So that the defendants
could not contact their comrades for reinforcements, allies of the
militia in the Telecommunications Bureau jammed the telephone lines into
the buildings.\footnote{HZRB, June 21,1977, November 16,1977, December
  17,1977, April 17,1979; Gao Zhixian, HZRB, June 26,1977.
  《杭州日报》，1977年6月21日、1977年11月16日、1977年12月17日、1979年4月17日；高志贤，《杭州日报》，1977年6月26日。}
This was the first such clash on the streets of Hangzhou since the
factional battles of the 1960s.

民兵旨在贯彻翁森鹤的格言：``自古以来，每一个政治集团的胜利都伴随着对其反对者的镇压。''1974
年 2 月 21
日，即杭州民兵指挥部公开宣布成立的前一天，民兵在省会市中心首次尝到了行动的滋味。翁、贺在``山下派''中的派性对手聚集在杭州商业中心附近的三家招待所。贺贤春和夏根发派民兵前去侦察，结果扰乱了公共交通和社会秩序。同日下午
5
点，贺、夏发出紧急通知，命令民兵把对手赶出招待所。来自十二家工厂的工人被动员起来，随后爆发大规模对抗。一百名参与者和数量不明的旁观者在混战中受伤。\footnote{Weng
  Senhe bufen zuizheng cailiao, p.~57. 《翁森鹤部分罪证材料》，第57页。}
两家招待所中价值 100,000
元的财物被损坏或砸毁。为使被攻击一方无法联系同志前来增援，民兵在电信局的盟友切断了通往这些建筑的电话线路。\footnote{HZRB,
  June 21,1977, November 16,1977, December 17,1977, April 17,1979; Gao
  Zhixian, HZRB, June 26,1977.
  《杭州日报》，1977年6月21日、1977年11月16日、1977年12月17日、1979年4月17日；高志贤，《杭州日报》，1977年6月26日。}
这是自 1960 年代派性武斗以来，杭州街头发生的第一次此类冲突。

The next major incident occurred on 17 March 1974. Details are scanty
but it was claimed that several hundred people were apprehended and
several score wounded. The militia sent a telegram to Wang Hongwen,
putting its side of the story, and he replied expressing
support.\footnote{CCP HMC propaganda department criticism group, HZRB,
  October 21,1977; HZRB, April 17,1979.
  中共杭州市委宣传部批判组，《杭州日报》，1977年10月21日；《杭州日报》，1979年4月17日。}
After these skirmishes the municipal party authorities moved to exercise
some control over the militia's activities. In early April at a meeting
convened to discuss political work in the militia, representatives
expressed their determination to fight to defend the anti-Lin
anti-Confucius campaign and to study Shanghai's experience in running
the organization.\footnote{ZJRB, April 8,1974, p.~1.
  《浙江日报》，1974年4月8日，第1版。} The view expressed at the meeting
that the party should ``absolutely'' command and lead the militia was an
obvious swipe at the PLA. As if to demonstrate further that the military
had no role to play in the operations of the HMMCH, no PLA
representative was present. The obligatory pledge to combine productive
labour with military training seemed to take formalistic propaganda to
the extreme.

下一起重大事件发生在 1974 年 3 月 17
日。细节很少，但据称有数百人被抓，数十人受伤。民兵给王洪文发去电报，陈述己方说法；王回复表示支持。\footnote{CCP
  HMC propaganda department criticism group, HZRB, October 21,1977;
  HZRB, April 17,1979.
  中共杭州市委宣传部批判组，《杭州日报》，1977年10月21日；《杭州日报》，1979年4月17日。}
这些冲突之后，市级党委当局开始试图对民兵活动施加某种控制。4
月初，在一次为讨论民兵政治工作而召开的会议上，代表们表示决心为保卫批林批孔运动而斗争，并学习上海管理该组织的经验。\footnote{ZJRB,
  April 8,1974, p.~1. 《浙江日报》，1974年4月8日，第1版。}
会议上提出的党应当``绝对''指挥和领导民兵的观点，是对解放军的明显指责。仿佛为了进一步表明军队在杭州民兵指挥部的运作中没有角色，没有任何解放军代表出席。关于把生产劳动同军事训练相结合的例行表态，似乎把形式主义宣传推到了极致。

The southern Zhejiang city of Wenzhou where bloody clashes had occurred
in 1967 and 1968 (see chapters two and three) experienced more trouble
at the end of 1973. A report of November referred to militia activities
at the Wenzhou Scissors Factory. It stated that militia contingents,
under the ``unified leadership of the municipal party committee'' and
armed with rifles and machine-guns, carried out street patrols and
participated in class struggle -- a euphemism for factional
fighting.\footnote{ZJRB, November 18,1973.
  《浙江日报》，1973年11月18日。} The city's militia command held a
rally in May 1974 and the report about the meeting stated that the
organization had been established ``amidst a fierce struggle between the
two classes and two lines'' and had devoted the previous two months to
studying the experience of the model Shanghai command. It had also,
continued the account,

浙南城市温州曾在 1967 年和 1968 年发生血腥冲突，见第二、三章；1973
年末那里又出现了更多麻烦。一份 11
月的报告提到温州剪刀厂的民兵活动。报告称，民兵分队在``市委统一领导''下，携带步枪和机枪，进行街头巡逻并参加阶级斗争，而阶级斗争是派性斗争的委婉说法。\footnote{ZJRB,
  November 18,1973. 《浙江日报》，1973年11月18日。} 该市民兵指挥部于
1974 年 5
月召开大会，有关会议的报告称，该组织是在``两大阶级、两条路线的激烈斗争中''建立的，并在此前两个月专门学习上海模范指挥部的经验。报告继续说，它还：

\begin{quote}
mobilized the militiamen, with industrial workers as the main force,
resolutely to implement the series of important instructions issued by
the Party Central Committee headed by Chairman Mao, participate actively
in class struggle in society, and fight against the sinister activities
of a handful of class enemies aimed at undermining the movement to
criticize Lin Biao and Confucius. By so doing {[}it has{]} safeguarded
State property and the peoples' lives, fully demonstrating the mighty
power of the dictatorship of the proletariat. The people's masses
{[}sic{]} have ardently lauded them as ``iron fists for smashing the
enemy'' and ``Red sentries who protect the people.''\footnote{ZJRB, May
  21, 1974, p.~4; quotation from ZPS, May 21,1974, SWB/FE/4609/BII/3-4.
  《浙江日报》，1974年5月21日，第4版；引文见ZPS，1974年5月21日，SWB/FE/4609/BII/3-4。}
\end{quote}

\begin{quote}
动员以产业工人为主体的民兵，坚决贯彻以毛主席为首的党中央的一系列重要指示，积极参加社会上的阶级斗争，同一小撮阶级敌人破坏批林批孔运动的罪恶活动作斗争。这样做，{[}它{]}保卫了国家财产和人民生命，充分显示了无产阶级专政的强大威力。人民群众
{[}sic{]}
热烈赞扬他们是``砸烂敌人的铁拳''和``保卫人民的红色哨兵''。\footnote{ZJRB,
  May 21, 1974, p.~4; quotation from ZPS, May 21,1974,
  SWB/FE/4609/BII/3-4.
  《浙江日报》，1974年5月21日，第4版；引文见ZPS，1974年5月21日，SWB/FE/4609/BII/3-4。}
\end{quote}

Unlike the April rally in Hangzhou, however, representatives of the PLA
were in attendance in Wenzhou.

不过，与杭州 4 月的大会不同，解放军代表出席了温州的会议。

In other parts of the province rebels seized control of the militia, as
in Jiande county on 20 May.\footnote{Jiande xianzhi, p.~26.
  《建德县志》，第26页。} In Lanxi county, when the militia command was
set up on 6 June the local branch of United Headquarters took charge.
Factional fighting broke out between this militia group and an
independent county militia corps (县民兵独立团) which had previously
been established on 2 June 1970.\footnote{Lanxi shizhi, pp.~21-2.
  《兰溪市志》，第21-22页。} This corps most probably supported the
local mass organization formerly aligned with Red Storm and now with the
``mountain base'' faction. In Linhai county, a faction-based Public
Order General Headquarters (纠察总指挥部) was established on 12 April
1974.\footnote{Linhai xianzhi, p.~31. 《临海县志》，第31页。} Thus, in
1974 the mass organizations were very much alive, not only in Hangzhou,
but in the county towns of Zhejiang.

在全省其他地区，造反派夺取了民兵控制权，如 5 月 20
日建德县的情况。\footnote{Jiande xianzhi, p.~26. 《建德县志》，第26页。}
在兰溪县，6 月 6
日民兵指挥部成立时，联合总部的当地分支接管了它。这个民兵组织与此前于
1970 年 6 月 2 日建立的县民兵独立团之间爆发了派性斗争。\footnote{Lanxi
  shizhi, pp.~21-2. 《兰溪市志》，第21-22页。}
该独立团很可能支持当地那个原先同红暴结盟、如今同``山下派''结盟的群众组织。在临海县，一个基于派性的纠察总指挥部于
1974 年 4 月 12 日成立。\footnote{Linhai xianzhi, p.~31.
  《临海县志》，第31页。} 因此，1974
年群众组织仍十分活跃，不仅在杭州如此，在浙江的县城中也是如此。

The civilian and military authorities' efforts to bring the militia
under control in the second half of 1974 were not very successful. On 11
August 1974 another major incident occurred, this time at the Hangzhou
Iron and Steel Mill. It is claimed that on 10 August Zhang, Weng and He,
together with their supporters in the mill, held a meeting at Weng's
retreat on Three Terrace mountain (三台山) in the south-western hills
around Hangzhou, to discuss ways of dealing with their opponents. It was
decided to seize Zhang Jifa (张纪法) in a pre-emptive strike against the
``mountain base'' faction. Zhang was a workshop cadre at the iron and
steel mill, which seems to have been a bastion of the anti-Weng forces.
The militia had pledged to flatten (踏平) the mill and crush its
opponents there.\footnote{See sources cited in n.~126. See also Forster,
  ``A Story from Hangzhou''.
  见注126所引资料。另见Forster：《一个来自杭州的故事》。}

1974
年下半年，文职和军事当局试图把民兵置于控制之下，但并不十分成功。1974 年
8 月 11 日，又发生一起重大事件，这次发生在杭州钢铁厂。据称，8 月 10
日，张、翁、贺同他们在该厂的支持者一起，在翁位于杭州西南丘陵三台山的休憩处开会，讨论如何对付他们的对手。会议决定先发制人，抓捕``山下派''成员张纪法。张是钢铁厂的一名车间干部，该厂似乎是反翁力量的堡垒。民兵曾誓言踏平该厂，并粉碎那里的对手。\footnote{See
  sources cited in n.~126. See also Forster, ``A Story from Hangzhou''.
  见注126所引资料。另见Forster：《一个来自杭州的故事》。}

The next day at about 8 am, Zhang Jifa and fellow-workers from the mill,
as well as employees at Guo Zhisong's Zhejiang Construction Company, the
Hangzhou Glass Factory (another large enterprise near the iron and steel
mill in Banshan), hung a large banner outside the Overseas Chinese Hotel
on Hubin (湖滨), the road running parallel to the east bank of the West
Lake, and placed big-character posters on walls near the department
store in downtown Hangzhou. The banner contained a brief message
directed against the proceedings of a lengthy provincial conference
which had just drawn to a conclusion in the city (see the following
section). The urban militia arrived in full battle-dress to arrest those
who were reading the banner and posters.\footnote{Weng Senhe bufen
  zuizheng cailiao, pp.~60-1, 63.
  《翁森鹤部分罪证材料》，第60-61、63页。} When Zhang led his group
forward to resist, a fight developed. The scuffle must have reached an
inconclusive end because the militia pursued Zhang and his supporters
all the way back to the mill. A fire then broke out, partly burning the
mill's auditorium. Each side blamed the other.\footnote{HZRB, November
  7,1978. 《杭州日报》，1978年11月7日。} Whoever was responsible, this
incident further inflamed factional antagonisms and contributed to the
atmosphere of lawlessness in Hangzhou.

次日早晨约 8
点，张纪法和钢铁厂同事，以及郭志松所在浙江建筑公司的职工、杭州玻璃厂职工，玻璃厂是半山杭州钢铁厂附近的另一家大型企业，在湖滨，即沿西湖东岸平行延伸的道路，华侨饭店外悬挂了一条大横幅，并在杭州市中心百货商店附近的墙上张贴大字报。横幅上的简短信息针对的是一场刚刚在市内结束的漫长省级会议的议程，见下一节。城市民兵全副武装赶到，逮捕正在阅读横幅和大字报的人。\footnote{Weng
  Senhe bufen zuizheng cailiao, pp.~60-1, 63.
  《翁森鹤部分罪证材料》，第60-61、63页。}
当张带领其人群上前抵抗时，斗殴发生。冲突想必没有分出胜负，因为民兵一路追赶张及其支持者回到钢铁厂。随后发生火灾，厂礼堂部分被烧毁。双方互相指责。\footnote{HZRB,
  November 7,1978. 《杭州日报》，1978年11月7日。}
无论责任在谁，这一事件都进一步激化了派性敌意，并助长了杭州无法无天的气氛。

Four other examples testify to the bitter and prolonged nature of the
factional antagonisms in Hangzhou. The first also concerns Zhang Jifa,
already mentioned above. In April 1974 Zhang, a demobilized soldier who
had once worked in Beijing guarding important buildings, had obtained a
record of Wang Hongwen's telephone calls to Zhang Yongsheng. Supporters
of Zhang Yongsheng in the iron and steel mill where Zhang Jifa worked
would pat their pockets and boast that ``we've got access to the highest
authorities; in here we've got Vice-chairman Wang's phone call to Zhang
Yongsheng''. Part of the transcript read: ``the current situation is
very favorable to you and I hope you do big things. But don't create the
impression that you are reversing the verdict on X and X'' {[}a
reference to Nan Ping and Xiong Yingtang, Lin Biao's purged military
followers in Zhejiang{]}. After Zhang Jifa obtained the documents he
decided to go to Beijing to pass them on to officials who had contacts
with central leaders. When Zhang's enemies discovered his intentions
they attempted to track him down. Zhang boarded a tram for the national
capital and, after doubling back on his tracks to shake off his
pursuers, arrived in Beijing on 7 May 1974. He passed on the material to
central ``leading comrades'' and then returned to Hangzhou.\footnote{HZRB,
  June 23,1967, August 27,1978, November 7,1978.
  《杭州日报》，1967年6月23日、1978年8月27日、1978年11月7日。} Zhang was
eventually arrested on 2 September and charged with responsibility for
the 11 August fire at his mill. For the next fourteen hours he was
interrogated and tortured by the militia, who threatened to throw him
into the factory's coke oven. Zhang was held for four months during
which time he was tortured seventeen times, suffering broken ribs and a
fractured arm, kidney damage, a punctured eardrum, severe concussion and
astigmatism to the right eye.

另外四个例子证明了杭州派性敌意的激烈而持久。第一个也涉及前文提到的张纪法。1974
年 4
月，张是一名复员军人，曾在北京从事重要建筑警卫工作；他获得了王洪文给张永生打电话的记录。在张纪法工作的钢铁厂里，张永生的支持者会拍着口袋夸口说：``我们通天；这里面有王副主席给张永生的电话。''记录的一部分写道：``当前形势对你们很有利，希望你们干大事。但不要造成你们是在为某某翻案的印象''{[}指林彪在浙江被清洗的军方追随者南萍和熊应堂{]}。张纪法取得这些文件后，决定去北京，把它们交给同中央领导有联系的官员。张的敌人发现他的意图后，试图追踪他。张登上前往首都的火车，并在途中折返绕行以摆脱追踪者，于
1974 年 5 月 7
日抵达北京。他把材料交给中央``领导同志''，然后返回杭州。\footnote{HZRB,
  June 23,1967, August 27,1978, November 7,1978.
  《杭州日报》，1967年6月23日、1978年8月27日、1978年11月7日。} 张最终于
9 月 2 日被捕，被控对 8 月 11
日钢铁厂火灾负责。随后十四个小时里，他遭到民兵审讯和折磨，民兵威胁要把他扔进工厂焦炉。张被关押四个月，其间遭受十七次酷刑，肋骨断裂、手臂骨折、肾脏受损、耳膜穿孔、严重脑震荡，右眼散光。

The second instance also illustrates the ferocity of factional
vendettas. Li Longbiao was a worker at the Hangzhou Worker-Peasant Bean
Products factory. His father, Li Baoxing, was a worker at the Zhejiang
Hemp Mill where factional differences were particularly acute. In August
1974 the elder Li fled the mill, fearing for his personal safety because
he was a member of the ``mountain base'' faction opposed to Weng and He
and was so known to them. On 11 August, the day of the fire at the iron
and steel mill, the younger Li went there in search of his father, a
further illustration of how factional antagonisms had involved employees
of many of the city's industrial plants. Li Longbiao was seized by the
militia and taken to its headquarters. There he was recognized as Li
Baoxing's son and accused of arson. After a beating he was incarcerated
in a secret prison and held for a month. For a further month and a half
he was kept in solitary confinement before his release in late October
1974. Li, who was only thirty-five at the time, emerged a
paraplegic.\footnote{Li Longbiao, ``How He Xianchun and company cruelly
  persecuted me'', HZRB, June 21,1977, p.~3.
  李龙彪：《贺贤春一伙怎样残酷迫害我》，《杭州日报》，1977年6月21日，第3版。}

第二个事例同样说明了派性报复的凶残。李龙彪是杭州工农豆制品厂工人。他的父亲李宝兴是浙江麻纺厂工人，那里派性分歧尤为尖锐。1974
年 8
月，老李因担心人身安全而逃离麻纺厂，因为他是反对翁、贺的``山下派''成员，而且他们也知道这一点。8
月 11
日，即钢铁厂火灾当天，小李去钢铁厂寻找父亲，这进一步说明派性敌意如何牵涉到该市许多工业厂矿的职工。李龙彪被民兵抓住并带到总部。在那里，他被认出是李宝兴的儿子，并被指控纵火。遭殴打后，他被关进一处秘密监狱，关押了一个月。获释前，他又被单独囚禁了一个半月，直到
1974 年 10 月下旬才释放。李当时只有三十五岁，出来时已成截瘫。\footnote{Li
  Longbiao, ``How He Xianchun and company cruelly persecuted me'', HZRB,
  June 21,1977, p.~3.
  李龙彪：《贺贤春一伙怎样残酷迫害我》，《杭州日报》，1977年6月21日，第3版。}

The third example concerns Yao Guolin, the leader of Wenzhou's Rebel
General HQ in the 1960s. In 1967 and again in 1968, Yao had led the
organization in guerrilla warfare against United Headquarters as well as
units of the 20th Army, before the military had suppressed it (see
chapters two and three). Yao had been arrested and jailed. When tensions
erupted in Wenzhou again in 1974, Wang Hongwen intervened to denounce
the local forces opposed to the ``mountain top'' faction as brigands
(see next section). Yao Guolin, who had been released from prison, was
the leader of the Wenzhou ``brigands''. On March 13, 1974, with fighting
having broken out again in Wenzhou, the provincial authorities issued a
warrant for Yao's arrest.\footnote{See 中共浙江省委 (CCP Zhejiang
  Provincial Committee), ``给姚国麟同志平反的决定'' (Decision
  rehabilitating Comrade Yao Guolin), December 30, 1978.
  见中共浙江省委（CCP Zhejiang Provincial
  Committee）《给姚国麟同志平反的决定》（Decision rehabilitating Comrade
  Yao Guolin），1978年12月30日。} At a public rally in June 1974 Yao was
denounced and paraded before the masses. Weng Senhe chaired the rally
and Zhang Yongsheng was one of the four speakers. Provincial and
municipal civilian and military party leaders attended and following the
rally Yao was incarcerated once again.\footnote{HZRB, June 25,1974.
  HZRB，1974年6月25日。}

第三个例子涉及 1960 年代温州造反总指挥部领导人姚国麟。1967 年和 1968
年，姚曾领导该组织同联合总部以及第二十军部队进行游击战，直到军队将其镇压，见第二、三章。姚被逮捕入狱。1974
年温州局势再次紧张时，王洪文出面，把反对``山上派''的地方力量斥为土匪，见下一节。已经出狱的姚国麟，是温州``土匪''的领导人。1974
年 3 月 13
日，温州再次爆发武斗后，省级当局发出逮捕姚的通缉令。\footnote{See
  中共浙江省委 (CCP Zhejiang Provincial Committee),
  ``给姚国麟同志平反的决定'' (Decision rehabilitating Comrade Yao
  Guolin), December 30, 1978. 见中共浙江省委（CCP Zhejiang Provincial
  Committee）《给姚国麟同志平反的决定》（Decision rehabilitating Comrade
  Yao Guolin），1978年12月30日。} 在 1974 年 6
月的一次公开大会上，姚被批斗并被押着示众。翁森鹤主持大会，张永生是四名发言者之一。省、市党政军领导出席了大会，会后姚再次被关押。\footnote{HZRB,
  June 25,1974. HZRB，1974年6月25日。}

Qiu Honggen was another long-time foe of Weng Senhe and United
Headquarters. In 1969, along with Fang Jianwen and other Red Storm
diehards, Qiu had refused to join the alliance with United Headquarters
(see chapter four). He had been arrested, imprisoned and only released
in early 1973 (see chapter six). In June 1974 Qiu joined Yao in the dock
at the public rally. He had written signed letters to Beijing
complaining about the activities of the Hangzhou urban militia and had
been arrested once more. Years later, Qiu proudly remembered his defiant
attitude at the rally where, by holding his head high and smiling, he
contravened regulations for a prisoner's demeanor on such
occasions.\footnote{Qiu Honggen, HZRB, November 26,1978.
  邱洪根，HZRB，1978年11月26日。}

邱洪根是翁森鹤和联合总部的另一个长期敌人。1969
年，邱同方剑文及其他红暴死硬分子一起，拒绝加入同联合总部的联盟，见第四章。他曾被逮捕、监禁，直到
1973 年初才获释，见第六章。1974 年 6
月，邱在公开大会上同姚一起站上被告席。他曾给北京写署名信，控诉杭州城市民兵的活动，因而再次被捕。多年后，邱自豪地回忆自己在大会上的反抗态度：他昂首微笑，违反了这种场合对囚犯仪态的规定。\footnote{Qiu
  Honggen, HZRB, November 26,1978. 邱洪根，HZRB，1978年11月26日。}

The viciousness and vindictiveness of the factional struggle was amply
demonstrated in incidents detailed above. The factionalism, of which the
fighting was only one manifestation, disrupted production and normal
government. The battles between United Headquarters and Red Storm in the
years 1967-69 had been fought with weapons looted from military
armories. In 1974 the militia HQ was the officially-sanctioned strike
force of the ``mountain top'' faction and was entitled to arm its
members. Its opponents in the ``mountain base'' faction took over the
previous militia organization, and with supportive military leaders
looking the other way or actively encouraging them, took up the gauntlet
thrown down by the Zhang/Weng/He rebel trio. In the middle of 1974
``posters appeared in Hangzhou, signed by the city's militia command,
condemning an illegally established provincial militia
command''.\footnote{CNA, 1013, p.~4. See also FEER, July 30,1976, p.~30.
  CNA，第1013期，第4页。另见 FEER，1976年7月30日，第30页。} Thus the
change in the nature of the factional battles and the organizations
involved from 1967-69 typified the change in the nature of rebellion and
factionalism from that earlier period.

上述事件充分展示了派性斗争的恶毒和报复性。派性主义破坏了生产和正常政府运作，武斗只是其表现之一。1967
至 1969
年联合总部与红暴之间的战斗，是用从军事军械库中抢来的武器进行的。到 1974
年，民兵总部成了``山上派''得到官方认可的打击力量，并有权武装其成员。其``山下派''对手接管了原先的民兵组织；在支持他们的军队领导人视而不见或积极鼓励之下，他们接下了张、翁、贺造反三人组抛出的挑战。1974
年中期，``杭州出现了由市民兵指挥部署名的传单，谴责一个非法建立的省民兵指挥部''。\footnote{CNA,
  1013, p.~4. See also FEER, July 30,1976, p.~30.
  CNA，第1013期，第4页。另见 FEER，1976年7月30日，第30页。}
因此，派性战斗的性质以及所涉组织相对于 1967 至 1969
年的变化，体现了造反和派性主义相较早期阶段的性质变化。 \#\#\# 3.
Pressure on the PLA \#\#\# 3. 对解放军的压力

The weight of the evidence suggests that the PLA was the principal
target of the anti-Lin anti-Confucius campaign. Zhou Enlai was the
radicals' main individual object of attack but it was the military as an
institution, and its senior regional commanders in particular, which
came under fire.\footnote{See Sheng P'eng, ``The Generals in the
  Criticizing-Lin and Criticizing-Kung Campaign'', CLG, 7:4 (1974-5),
  esp.~pp.~82-6. Zhu De's wife, Kang Keqing, certainly gained this
  impression after listening to Jiang Qing's speech at the January rally
  launching the campaign. See, Liu Qiguang, ``朱德同志在文化大革命中''
  (Comrade Zhu De in the ``Cultural Revolution''), 中华英烈 (China's
  Heroes and Martyrs), No.~3 (1986), p.~8. Liu mistakenly dates Jiang's
  speech at January 1975 when it clearly occurred in the January of the
  previous year. 见 Sheng P'eng, ``The Generals in the Criticizing-Lin
  and Criticizing-Kung Campaign'', CLG, 7:4
  (1974-5)，尤其第82-86页。朱德的妻子康克清在听了江青于1月发动该运动的集会上的讲话后，显然也形成了这种印象。见刘启光《朱德同志在文化大革命中》（Comrade
  Zhu De in the ``Cultural Revolution''），《中华英烈》（China's Heroes
  and
  Martyrs），1986年第3期，第8页。刘误将江青的讲话日期定为1975年1月，但该讲话显然发生在前一年的1月。}
In December 1973 Mao had transferred regional commanders as well as
expressing dissatisfaction with the running of the CC MAC. Alarm at the
implications of the campaign for the military may have prompted Ye
Jianying to write to Mao on 30 January 1974 expressing his
concern.\footnote{``Wenhua dageming'' ruogan shijian zhenxiang, p.~113.
  《``Wenhua dageming'' ruogan shijian zhenxiang》，第113页。} At the
Politburo meeting of 31 January 1974, Zhou Enlai had attempted to
restrict the scope of the campaign within the PLA.\footnote{Gao Wenqian,
  RMRB, January 5,1986. 高文谦，RMRB，1986年1月5日。} By protecting his
military allies, Zhou was at the same time shoring up his own position.

大量证据表明，解放军是批林批孔运动的主要目标。周恩来是激进派重点攻击的个人对象，但真正遭到火力集中打击的是作为制度性机构的军队，尤其是其高级大区指挥员。\footnote{See
  Sheng P'eng, ``The Generals in the Criticizing-Lin and
  Criticizing-Kung Campaign'', CLG, 7:4 (1974-5), esp.~pp.~82-6. Zhu
  De's wife, Kang Keqing, certainly gained this impression after
  listening to Jiang Qing's speech at the January rally launching the
  campaign. See, Liu Qiguang, ``朱德同志在文化大革命中'' (Comrade Zhu De
  in the ``Cultural Revolution''), 中华英烈 (China's Heroes and
  Martyrs), No.~3 (1986), p.~8. Liu mistakenly dates Jiang's speech at
  January 1975 when it clearly occurred in the January of the previous
  year. 见 Sheng P'eng, ``The Generals in the Criticizing-Lin and
  Criticizing-Kung Campaign'', CLG, 7:4
  (1974-5)，尤其第82-86页。朱德的妻子康克清在听了江青于1月发动该运动的集会上的讲话后，显然也形成了这种印象。见刘启光《朱德同志在文化大革命中》（Comrade
  Zhu De in the ``Cultural Revolution''），《中华英烈》（China's Heroes
  and
  Martyrs），1986年第3期，第8页。刘误将江青的讲话日期定为1975年1月，但该讲话显然发生在前一年的1月。}
1973年12月，毛泽东调动了各大军区司令员，同时也表达了对中共中央军委运作的不满。运动对军队可能产生的影响引发的忧虑，或许促使叶剑英在1974年1月30日致信毛泽东表达担心。\footnote{``Wenhua
  dageming'' ruogan shijian zhenxiang, p.~113. 《``Wenhua dageming''
  ruogan shijian zhenxiang》，第113页。}
在1974年1月31日的政治局会议上，周恩来试图限制运动在解放军内部的范围。\footnote{Gao
  Wenqian, RMRB, January 5,1986. 高文谦，RMRB，1986年1月5日。}
周在保护其军方盟友的同时，也是在巩固自己的地位。

At a meeting on 8 February 1974 Wang Hongwen and Zhang Chunqiao called
the PLA general staff leadership ``as far right as you could go'' and
stated that power should be seized in the General Political Department
(where Zhang had recently been appointed director).\footnote{Hao Mengbi
  and Duan Haoran, Liushinian, p.~634.
  郝梦笔、段浩然，《Liushinian》，第634页。} On 5 March 1974 Jiang Qing
spoke disparagingly of the leaders of the military. As if by signal the
slogan ``knock down the big warlords'' appeared in the streets of the
capital and elsewhere, with what was later claimed as serious
repercussions for discipline in the army.\footnote{CCP CC, zhongfa
  (1976) No.~24, Zhonggong nianbao, 1977, 5:17; Hao Mengbi and Duan
  Haoran, Liushinian, p.~634. 中共中央，中发（1976）24号，《Zhonggong
  nianbao》，1977年，5:17；郝梦笔、段浩然，《Liushinian》，第634页。}
The reverberations of Jiang's speech were felt in Zhejiang.\footnote{See
  Ji Yangwen, ``A shocking display of opposing the party and throwing
  the army into chaos'', HZRB, February 9, 1977. 见纪仰文《A shocking
  display of opposing the party and throwing the army into
  chaos》，HZRB，1977年2月9日。} In speeches to military cadres shortly
afterward, Wang Hongwen and Zhang Chunqiao elaborated on Jiang's theme.
Wang claimed that revisionism was rife in the PLA and that if war broke
out only the militia could be relied upon. In a speech at General Staff
Headquarters he suggested that

在1974年2月8日的一次会议上，王洪文和张春桥称解放军总参谋部领导层``右得不能再右''，并表示应当夺取总政治部的权力（张春桥不久前刚被任命为总政治部主任）。\footnote{Hao
  Mengbi and Duan Haoran, Liushinian, p.~634.
  郝梦笔、段浩然，《Liushinian》，第634页。}
1974年3月5日，江青以轻蔑口吻谈论军队领导人。仿佛收到信号一般，``打倒大军阀''的口号出现在首都及其他地方的街头，后来据称这对军队纪律造成了严重影响。\footnote{CCP
  CC, zhongfa (1976) No.~24, Zhonggong nianbao, 1977, 5:17; Hao Mengbi
  and Duan Haoran, Liushinian, p.~634.
  中共中央，中发（1976）24号，《Zhonggong
  nianbao》，1977年，5:17；郝梦笔、段浩然，《Liushinian》，第634页。}
江青讲话的回响也传到了浙江。\footnote{See Ji Yangwen, ``A shocking
  display of opposing the party and throwing the army into chaos'',
  HZRB, February 9, 1977. 见纪仰文《A shocking display of opposing the
  party and throwing the army into chaos》，HZRB，1977年2月9日。}
不久后，王洪文和张春桥在对军队干部的讲话中进一步发挥了江青的主题。王声称解放军内修正主义盛行，如果战争爆发，唯一可以依靠的只有民兵。他在总参谋部的一次讲话中提出：

\begin{quote}
It is necessary to suppress continuously the Rightist ideology, to
mobilize the masses to make exposures, and open the lid on the top
echelon. It's difficult but not too difficult. This time we must make a
resolution: we must open up the lid, smash it if it cannot be opened,
and blow it open if it cannot be smashed.\footnote{CCP CC, zhongfa
  (1977), No.~37, I\&S, 14: 8 (1978), pp.~93-5, 99.
  中共中央，中发（1977）37号，I\&S，14:8（1978），第93-95、99页。}
\end{quote}

\begin{quote}
有必要不断压制右倾思想，发动群众揭发，把上层的盖子揭开。这很难，但也不是太难。这一次我们必须下定决心：一定要把盖子揭开，揭不开就砸开，砸不开就炸开。\footnote{CCP
  CC, zhongfa (1977), No.~37, I\&S, 14: 8 (1978), pp.~93-5, 99.
  中共中央，中发（1977）37号，I\&S，14:8（1978），第93-95、99页。}
\end{quote}

Wang was thus calling for a repetition of the type of anarchical
behavior which had occurred in 1966-67. On 11 March, as another sign of
her suspicion of the military Jiang Qing, in the name of the Politburo,
ordered the PLA newspaper Liberation Army Daily to cease compiling and
carrying its own articles.\footnote{Hao Mengbi and Duan Haoran,
  Liushinian, p.~634. 郝梦笔、段浩然，《Liushinian》，第634页。}

王洪文实际上是在号召重演1966至1967年间曾经出现的那种无政府行为。3月11日，江青以政治局名义命令解放军报《解放军报》停止自行编写和刊登本报文章，这也是她疑忌军队的又一个迹象。\footnote{Hao
  Mengbi and Duan Haoran, Liushinian, p.~634.
  郝梦笔、段浩然，《Liushinian》，第634页。}

Zhang, Weng and He did their best to carry out the spirit of Wang
Hongwen's March speech to ``blow open the lid'' of ZPMD party
committees. Their direct targets were the leading political cadres of
the ZPMD, Xia Qi, Tie Ying and Tan Qilong, all of whom owed their
positions to Xu Shiyou and Zhou Enlai. But ultimately they seemed to be
aiming at Xu and Zhou. In 1972, Xia and Tie in particular had most
likely mounted a witch-hunt against leftist military cadres and their
rebel allies. Now it was their turn to feel the heat of the political
blowtorch.

张、翁、贺竭力贯彻王洪文3月讲话中``炸开''浙江省军区党委盖子的精神。他们的直接目标是浙江省军区的主要政治干部夏琦、铁瑛和谭启龙，这些人的职位都得益于许世友和周恩来。但归根结底，他们的目标似乎是许和周。1972年，特别是夏琦和铁瑛很可能曾对左派军队干部及其造反派盟友进行过政治清洗。现在轮到他们承受政治喷灯的灼烧了。

Jiang Qing's choice of the anti-chemical warfare company in Zhejiang to
carry the banner of the campaign against Lin Biao and Confucius in the
military posed serious problems for Xia Qi. As the commissar in charge
of political work in the military district, Xia was expected to promote
Jiang Qing's model, repressing any distaste he may have felt for the
task. From early February 1974 articles praising the company for its
contributions to the campaign appeared almost daily in the provincial
press.\footnote{CNA, 974 (September 20,1974), p.~6, commented on this
  media attention.
  CNA，第974期（1974年9月20日），第6页，对这种媒体关注作了评论。}

江青选择浙江的防化连来扛起军队批林批孔运动的旗帜，给夏琦造成了严重难题。作为负责军区政治工作的政委，夏被要求宣传江青树立的典型，无论他对这项任务有多么反感，都必须压抑下去。从1974年2月上旬起，省内报刊几乎每天都刊登文章，赞扬该连对运动的贡献。\footnote{CNA,
  974 (September 20,1974), p.~6, commented on this media attention.
  CNA，第974期（1974年9月20日），第6页，对这种媒体关注作了评论。}

The problem that Jiang Qing's interest in the company was causing the
Zhejiang leaders was compounded by Wang Hongwen's renewed intervention
in provincial affairs. From 16 to 19 March 1974 he telephoned the
Zhejiang leadership each day to inquire about armed struggle which had
erupted in Wenzhou. Due to misunderstanding or deliberate
misinterpretation, the ZPMD's failure to respond adequately to Wang's
directives to bring the situation under control, resulted in three
charges being levelled against the provincial military leaders. First,
they were accused of failing to collect weapons from the warring groups;
second of having fabricated Wang Hongwen's directives; and third of
having acted as the ``back-stage boss to local brigands''
(后台老板).\footnote{Ji Yangwen, HZRB, February 9,1977; HZRB, September
  14,1978; HZRB, December 31,1976.
  纪仰文，HZRB，1977年2月9日；HZRB，1978年9月14日；HZRB，1976年12月31日。}
It can be assumed that sections of the local military were supporting
groups opposed to the ``mountain top'' rebels, in the same way as they
had sponsored groups to fight against United Headquarters in 1967-68. To
increase further the pressure on Xia Qi and his superiors, in mid 1974
Zhang Chunqiao and Ding Sheng despatched Li Bincheng, Political
Commissar of the Nanjing Artillery Units, to Hangzhou as Political
Commissar of the ZPMD. Li outranked Xia and was subordinate only to Tan
and Tie.

江青对该连的关注本已给浙江领导人带来麻烦，而王洪文再次干预省内事务又使问题更加复杂。1974年3月16日至19日，他每天都给浙江领导层打电话，询问温州爆发的武斗情况。由于误解或有意曲解，浙江省军区未能充分响应王关于控制局势的指示，结果省军区领导人受到三项指控。第一，他们被指未能从交战各派手中收缴武器；第二，被指伪造王洪文的指示；第三，被指充当``地方土匪的后台老板''（后台老板）。\footnote{Ji
  Yangwen, HZRB, February 9,1977; HZRB, September 14,1978; HZRB,
  December 31,1976.
  纪仰文，HZRB，1977年2月9日；HZRB，1978年9月14日；HZRB，1976年12月31日。}
可以推测，地方军队的某些部分正在支持反对``山上派''的组织，正如他们在1967至1968年曾扶植一些组织同联合总部作战一样。为进一步加大对夏琦及其上级的压力，1974年中期，张春桥和丁盛派南京炮兵部队政治委员李彬成到杭州，出任浙江省军区政治委员。李的级别高于夏，只隶属于谭和铁。

The confrontational style of politics utilized so effectively by the
``mountain top'' leaders was lethal when backed up by central leaders
such as Wang Hongwen. Zhang, Weng and He had coerced the authorities,
especially the Hangzhou party leadership, into agreeing to the
appointment of their supporters to responsible party posts. They had set
up a militia command with He Xianchun in charge. Their next step was to
apply further pressure to Tan Qilong, Tie Ying and Xia Qi. Tie and Xia
had already lost the right to attend public engagements, while Tan seems
to have appeased the rebels by going along with the campaign, perhaps in
the hope of blunting it. As a provincial first secretary he really had
no option but to implement central policies, whether these were moderate
or radical. Direct orders were difficult to circumvent or sabotage. At
the same time, he would have tried to protect Tie and Xia to the best of
his ability and inclinations, especially considering that the 1972-73
purge of the radicals which they had carried out had undoubtedly
received his blessing and enthusiastic support.

``山上派''领导人娴熟运用的对抗式政治，在得到王洪文等中央领导人支持时具有致命威力。张、翁、贺已经胁迫当局，特别是杭州市委领导层，同意将他们的支持者任命到重要的党委岗位。他们还建立了以贺贤春负责的民兵指挥机构。下一步，他们便是进一步向谭启龙、铁瑛和夏琦施压。铁和夏已经失去出席公开活动的权利，而谭似乎通过顺从运动来安抚造反派，也许是希望削弱运动的锋芒。作为省委第一书记，他实际上别无选择，只能执行中央政策，无论这些政策是温和还是激进。直接命令很难规避或破坏。同时，他会尽其能力和意愿保护铁和夏，尤其是考虑到他们在1972至1973年对激进派实施的清洗无疑曾得到他的首肯和热烈支持。

On 26 March 1974\footnote{In a deposition dated February 18,1977, Xia Qi
  recalled that the meeting opened on April 18. See Zhang Yongshengde
  zuizheng cailiao, p.~68. His memory was clearly defective.
  在一份日期为1977年2月18日的证词中，夏琦回忆说会议于4月18日开幕。见《Zhang
  Yongshengde zuizheng cailiao》，第68页。他的记忆显然有误。} a meeting
of great significance for the provincial leadership opened in Hangzhou.
It was a joint meeting of the CCP ZPC, ZPRC and the party committee of
the ZPMD and was known as the triple plenum (三全会).\footnote{A similar
  meeting also took place in Hunan province. See CNS, 527 (July
  24,1974), pp.~4-6. 湖南省也举行过一次类似会议。见
  CNS，第527期（1974年7月24日），第4-6页。} Triple plenums later opened
in Hangzhou on 15 May and in prefectures and counties across the
province, paralyzing the work of party committees.\footnote{See HZRB,
  September 17,1977. 见 HZRB，1977年9月17日。} It is alleged that CC
document No.~17 of 1974 authorized the formation of provincial-level
committees with the aim of promoting unity among party, government and
military leaders.\footnote{Agence France Presse (AFP), June 24,1974,
  FBIS/CHI, 123 (1974), E1-2. 法新社（Agence France Presse,
  AFP），1974年6月24日，FBIS/CHI，第123期（1974），E1-2。} The
experience of Zhejiang illustrates that if this was the intention of the
center, things went badly wrong. From the time of its opening until its
closure in early August, after a marathon sitting which lasted with
interruptions for 133 days, the triple plenum became a forum for the
airing of personal grudges and grave accusations. The forum enabled the
``mountain top'' rebels to exert a degree of control over the
organization where they had least influence, the ZPMD and its party
committees. They were able to dilute its influence by outnumbering
military delegates with members from the CCP ZPC and the ZPRC, where
they had greater representation. Whether the plenum foreshadowed a major
change in the organizational structure and personnel of the provincial
leadership is unclear. Certainly, Zhang Yongsheng was accused of
attempting such a reform at Jiande county by amalgamating party,
government and military leading groups into a single body
(三搞一套).\footnote{Zhang also tried to integrate urban militia units
  with the People's Armed Forces Departments. See HZRB, March 24,1977;
  CCP HMC party school 3rd study class, ``Strip off the disguise,
  restore its original features'', HZRB, 5 August 1977; HZRB, April
  5,1978. In Harding's typology, such a simplification of bureaucratic
  structure is viewed as part of the rationalizing approach to the
  bureaucratic dilemma. Harding, Organizing China, p.~330.
  张还试图把城市民兵单位同人民武装部整合起来。见
  HZRB，1977年3月24日；中共杭州市委党校第三期学习班《Strip off the
  disguise, restore its original
  features》，HZRB，1977年8月5日；HZRB，1978年4月5日。在 Harding
  的类型划分中，这种官僚结构的简化被视为应对官僚困境的理性化路径的一部分。Harding,
  Organizing China，第330页。} The plenum selected an investigation
group headed by Shen Ce to examine the verdicts passed on Lin Biao's
local supporters and their rebel allies by a special case group in 1972.
The high-ranking cadres who had been investigated in 1972 included Shen
himself, provincial Deputy-secretary Chai Qikun, Hangzhou First
Secretary Wang Zida and possibly Zhu Quanlin, a military cadre from the
20th Army.\footnote{Zhang Yongshengde zuizheng cailiao, p.~93. 《Zhang
  Yongshengde zuizheng cailiao》，第93页。} Members of the investigation
group included personnel from the municipal militia headquarters.
Military cadres were summoned for questioning, cadres' personal files
examined, and records of meetings, archives and documents relating to
defence preparations perused. On 16 April, Chai Qikun and He Xianchun
led the investigation team to the offices of the provincial special case
group to seize material the group had collected relating to Zhang
Yongsheng's connections with Lin Gang. Lin was the Qinghua University
leader who had played such an important liaison role in Zhejiang during
1966-67 and was now condemned as a counter-revolutionary and ``May 16
element''\footnote{Ibid, p.~94. 同上，第94页。} (see chapters one and
six).

1974年3月26日，\footnote{In a deposition dated February 18,1977, Xia Qi
  recalled that the meeting opened on April 18. See Zhang Yongshengde
  zuizheng cailiao, p.~68. His memory was clearly defective.
  在一份日期为1977年2月18日的证词中，夏琦回忆说会议于4月18日开幕。见《Zhang
  Yongshengde zuizheng cailiao》，第68页。他的记忆显然有误。}
一个对省级领导层极为重要的会议在杭州开幕。这是中共浙江省委、浙江省革命委员会和浙江省军区党委的联席会议，被称为``三全会''。\footnote{A
  similar meeting also took place in Hunan province. See CNS, 527 (July
  24,1974), pp.~4-6. 湖南省也举行过一次类似会议。见
  CNS，第527期（1974年7月24日），第4-6页。}
后来，5月15日杭州以及全省各地市县也召开了三全会，使党委工作陷于瘫痪。\footnote{See
  HZRB, September 17,1977. 见 HZRB，1977年9月17日。}
据称，1974年中共中央17号文件授权成立省一级委员会，目的在于促进党政军领导人的团结。\footnote{Agence
  France Presse (AFP), June 24,1974, FBIS/CHI, 123 (1974), E1-2.
  法新社（Agence France Presse,
  AFP），1974年6月24日，FBIS/CHI，第123期（1974），E1-2。}
浙江的经验说明，如果这确是中央的意图，那么事情的发展便严重偏离了初衷。从开幕到8月初闭幕，这场断断续续持续133天的马拉松式会议，成了宣泄个人怨恨和提出严重指控的论坛。这个论坛使``山上派''得以对他们影响力最弱的机构，即浙江省军区及其党委，施加一定程度的控制。他们能够通过让浙江省委和省革委会成员在人数上压倒军队代表来稀释军方影响，因为他们在前两个机构中拥有更大代表性。全会是否预示着省级领导机构组织结构和人员的重大变化，目前尚不清楚。可以肯定的是，张永生被指控曾在建德县试图进行这种改革，把党、政、军领导班子合并成一个机构（三搞一套）。\footnote{Zhang
  also tried to integrate urban militia units with the People's Armed
  Forces Departments. See HZRB, March 24,1977; CCP HMC party school 3rd
  study class, ``Strip off the disguise, restore its original
  features'', HZRB, 5 August 1977; HZRB, April 5,1978. In Harding's
  typology, such a simplification of bureaucratic structure is viewed as
  part of the rationalizing approach to the bureaucratic dilemma.
  Harding, Organizing China, p.~330.
  张还试图把城市民兵单位同人民武装部整合起来。见
  HZRB，1977年3月24日；中共杭州市委党校第三期学习班《Strip off the
  disguise, restore its original
  features》，HZRB，1977年8月5日；HZRB，1978年4月5日。在 Harding
  的类型划分中，这种官僚结构的简化被视为应对官僚困境的理性化路径的一部分。Harding,
  Organizing China，第330页。}
全会选出一个由沈策领导的调查组，审查1972年专案组对林彪在地方的支持者及其造反派盟友作出的结论。1972年被调查的高级干部包括沈策本人、省委副书记柴启琨、杭州市委第一书记王子达，可能还包括来自第二十军的军队干部朱全林。\footnote{Zhang
  Yongshengde zuizheng cailiao, p.~93. 《Zhang Yongshengde zuizheng
  cailiao》，第93页。}
调查组成员中包括市民兵指挥部人员。军队干部被传唤讯问，干部个人档案被查阅，与战备有关的会议记录、档案和文件也被翻检。4月16日，柴启琨和贺贤春率调查组前往省专案组办公室，夺取该组收集的关于张永生同林岗关系的材料。林岗是清华大学领导人，1966至1967年间曾在浙江发挥极为重要的联络作用，如今则被定为反革命和``五一六分子''\footnote{Ibid,
  p.~94. 同上，第94页。}（见第一章和第六章）。

Zhang and Weng toured military units, including those under the
jurisdiction of the 5th Air Force and 20th Army. They denounced
commanders as revisionists and called on the troops to break the
master-serf relationship which they claimed characterized relations
between officers and men. They also asserted that the PLA was
politically unreliable and that the urban militia would be used to put
down a military insurrection. The militia further provoked the military
by attacking its offices, occupying cadres' hostels, cutting the supply
of food and water and transport and communication links to units
stationed in Hangzhou, and firing on army camps.\footnote{Weng Senhe
  bufen zuizheng cailiao, p.~52; HZRB, December 31, 1976; Ji Yangwen,
  HZRB, February 9,1977; HZRB, September 14,1978. 《Weng Senhe bufen
  zuizheng
  cailiao》，第52页；HZRB，1976年12月31日；纪仰文，HZRB，1977年2月9日；HZRB，1978年9月14日。}

张和翁巡视了军队单位，包括第五空军和第二十军管辖下的单位。他们斥责指挥员是修正主义者，并号召部队打破他们所谓官兵关系中的主奴关系。他们还声称解放军在政治上不可靠，城市民兵将被用来镇压军事叛乱。民兵又通过攻击军队机关、占领干部招待所、切断驻杭州部队的食品和饮水供应以及交通通信联系、向军营开枪等方式进一步挑衅军队。\footnote{Weng
  Senhe bufen zuizheng cailiao, p.~52; HZRB, December 31, 1976; Ji
  Yangwen, HZRB, February 9,1977; HZRB, September 14,1978. 《Weng Senhe
  bufen zuizheng
  cailiao》，第52页；HZRB，1976年12月31日；纪仰文，HZRB，1977年2月9日；HZRB，1978年9月14日。}

Although he was ill, Tie Ying was dragged into ``criticism and struggle
auditoriums'' (我\ldots 带着病被拉上批斗会场).\footnote{ZJRB, January 9,
  1985. ZJRB，1985年1月9日。} In order to escape the severe harassment
Xia Qi fled to a hospital. Wang Hongwen ordered him back or face
dismissal. On 7 April 1974 events took another turn for the worse for
Xia. On that day the Hangzhou Daily published an article entitled
``Struggle to the end with the bourgeois restorationist
forces''.\footnote{HZRB criticism group, ``A poisonous weed to attack
  the party and destabilize the army'', HZRB, October 10, 1978. HZRB
  批判组《A poisonous weed to attack the party and destabilize the
  army》，HZRB，1978年10月10日。} Wang Hongwen telephoned the forum,
then in session, and in referring to Xia Qi left off the two characters
``comrade'' (同志). The rebels seized upon the omission to declare that
Xia's case had become one involving an antagonistic contradiction. He
was suspended from his positions, subjected to isolated interrogation
from a special case group and made the object of attack at meetings
which continued over the next five months. The rebels labelled Xia a
``representative of restorationist forces'', a ``class enemy who has
wormed his way into the PLA'' and a ``revisionist''. Wang Hongwen later
confessed, ``Ah! I wasn't careful when I phoned, and left off the two
characters `comrade'. You went a bit far''.\footnote{Zhai Diaoshi,
  ``Expose and criticize the Gang of Four's heinous crimes in opposing
  the army in Zhejiang and throwing it into chaos'', HZRB, February 1,
  1977; HZRB, September 14,1978. Foreign travelers observed posters in
  Hangzhou at this time denouncing Xia Qi. AFP, July 9, 1974, FBIS/CHI,
  133 (1974), E1. Various sources alleged that sometime in 1974 - the
  dates vary from April to June - the CC issued a document No.~18
  permitting criticism of cadres by name. See Anita Chan and Jonathan
  Unger (eds.), ``The Case of Li Yizhe'', CLC, 10:3 (1977), pp.~20-1,
  p.~65 fn.11 who date it in April 1974. Dittmer, ``Bases of Power'',
  p.~57, fn. 75 dates the directive at May 18,1974, after Mao's comments
  of May 5 in which the Chairman had said: ``I see nothing wrong with
  posting big-character posters in the streets, and if foreigners want
  to read them, fine; if the Chinese want to read them, even better''.
  AFP dated the directive in late May. AFP, June 19, 1974, FBIS/CHI, 120
  (1974), E1-2. Starr states that the CC ``was reported to have
  approved'' such a policy on June 13. John B. Starr, ``China in 1974:
  `Weeding through the old to bring forth the new'\,'', AS, 15:1 (1975),
  p.~9. 翟调适《Expose and criticize the Gang of Four's heinous crimes
  in opposing the army in Zhejiang and throwing it into
  chaos》，HZRB，1977年2月1日；HZRB，1978年9月14日。外国旅行者当时在杭州看到谴责夏琦的标语。AFP，1974年7月9日，FBIS/CHI，第133期（1974），E1。各种资料声称，在1974年某个时候，具体日期从4月至6月不等，中共中央发布了第18号文件，允许点名批评干部。见
  Anita Chan and Jonathan Unger 编，``The Case of Li Yizhe'', CLC,
  10:3（1977），第20-21页，第65页注11；他们将日期定为1974年4月。Dittmer,
  ``Bases of
  Power''，第57页注75，将该指示定为1974年5月18日，即毛泽东5月5日发表意见之后；主席当时曾说：``我看在街上贴大字报没有什么不好，外国人要看，好；中国人要看，更好。''AFP
  将该指示定在5月下旬。AFP，1974年6月19日，FBIS/CHI，第120期（1974），E1-2。Starr
  称，据报道中共中央于6月13日批准了这项政策。John B. Starr, ``China in
  1974: `Weeding through the old to bring forth the new'\,'', AS,
  15:1（1975），第9页。} {[}!{]} In a speech of 23 May Weng Senhe was
adamant that rank would not protect any cadre from criticism.\footnote{Weng
  Senhe bufen zuizheng cailiao, p.~53. 《Weng Senhe bufen zuizheng
  cailiao》，第53页。}

铁瑛虽然患病，仍被拖进``批斗会场''（我\ldots 带着病被拉上批斗会场）。\footnote{ZJRB,
  January 9, 1985. ZJRB，1985年1月9日。}
为躲避严重骚扰，夏琦逃进医院。王洪文命令他回来，否则便撤职。1974年4月7日，夏琦的处境进一步恶化。当天，《杭州日报》发表题为《同资产阶级复辟势力斗争到底》的文章。\footnote{HZRB
  criticism group, ``A poisonous weed to attack the party and
  destabilize the army'', HZRB, October 10, 1978. HZRB 批判组《A
  poisonous weed to attack the party and destabilize the
  army》，HZRB，1978年10月10日。}
王洪文给正在开会的论坛打电话，提到夏琦时省去了``同志''二字（同志）。造反派抓住这一遗漏，宣布夏的问题已经成为涉及敌我矛盾的问题。他被停职，受到专案组隔离审查，并在随后五个月持续召开的会议上成为攻击对象。造反派给夏贴上``复辟势力代表''、``钻进解放军的阶级敌人''和``修正主义者''等标签。王洪文后来承认：``唉！我打电话的时候不小心，把`同志'两个字漏掉了。你们搞得有点过火。''\footnote{Zhai
  Diaoshi, ``Expose and criticize the Gang of Four's heinous crimes in
  opposing the army in Zhejiang and throwing it into chaos'', HZRB,
  February 1, 1977; HZRB, September 14,1978. Foreign travelers observed
  posters in Hangzhou at this time denouncing Xia Qi. AFP, July 9, 1974,
  FBIS/CHI, 133 (1974), E1. Various sources alleged that sometime in
  1974 - the dates vary from April to June - the CC issued a document
  No.~18 permitting criticism of cadres by name. See Anita Chan and
  Jonathan Unger (eds.), ``The Case of Li Yizhe'', CLC, 10:3 (1977),
  pp.~20-1, p.~65 fn.11 who date it in April 1974. Dittmer, ``Bases of
  Power'', p.~57, fn. 75 dates the directive at May 18,1974, after Mao's
  comments of May 5 in which the Chairman had said: ``I see nothing
  wrong with posting big-character posters in the streets, and if
  foreigners want to read them, fine; if the Chinese want to read them,
  even better''. AFP dated the directive in late May. AFP, June 19,
  1974, FBIS/CHI, 120 (1974), E1-2. Starr states that the CC ``was
  reported to have approved'' such a policy on June 13. John B. Starr,
  ``China in 1974: `Weeding through the old to bring forth the new'\,'',
  AS, 15:1 (1975), p.~9. 翟调适《Expose and criticize the Gang of Four's
  heinous crimes in opposing the army in Zhejiang and throwing it into
  chaos》，HZRB，1977年2月1日；HZRB，1978年9月14日。外国旅行者当时在杭州看到谴责夏琦的标语。AFP，1974年7月9日，FBIS/CHI，第133期（1974），E1。各种资料声称，在1974年某个时候，具体日期从4月至6月不等，中共中央发布了第18号文件，允许点名批评干部。见
  Anita Chan and Jonathan Unger 编，``The Case of Li Yizhe'', CLC,
  10:3（1977），第20-21页，第65页注11；他们将日期定为1974年4月。Dittmer,
  ``Bases of
  Power''，第57页注75，将该指示定为1974年5月18日，即毛泽东5月5日发表意见之后；主席当时曾说：``我看在街上贴大字报没有什么不好，外国人要看，好；中国人要看，更好。''AFP
  将该指示定在5月下旬。AFP，1974年6月19日，FBIS/CHI，第120期（1974），E1-2。Starr
  称，据报道中共中央于6月13日批准了这项政策。John B. Starr, ``China in
  1974: `Weeding through the old to bring forth the new'\,'', AS,
  15:1（1975），第9页。} {[}!{]}
翁森鹤在5月23日的一次讲话中坚称，级别不能保护任何干部免受批判。\footnote{Weng
  Senhe bufen zuizheng cailiao, p.~53. 《Weng Senhe bufen zuizheng
  cailiao》，第53页。}

On 7 August 1974 the provincial triple plenum finally came to an end,
after months of debate and controversy.\footnote{Xia Qi claimed that the
  conference closed on August 23. See Zhang Yongshengde zuizheng
  cailiao, p.~68. The Hunan triple plenum ended the previous month. See
  CNS, 527, p.~4. 夏琦声称会议于8月23日闭幕。见《Zhang Yongshengde
  zuizheng cailiao》，第68页。湖南的三全会在前一个月结束。见
  CNS，第527期，第4页。} The sanitized report which was issued to sum up
the meeting was a compromise document. On the one hand it urged
``leading cadres who have made mistakes'' {[}Tie Ying and Xia Qi{]} to
``correctly handle the mass criticism, seriously conduct self-criticism,
correct their mistakes courageously and reinforce their spirit to work
hard''. On the other hand it referred to ``bourgeois factionalism''
within the ranks of the masses and asked those who were guilty to
recognize their errors. The final point of the communique specified the
hegemony of the party over all other organizations, including the
campaign small groups. It stated:

1974年8月7日，经过数月辩论和争议后，省三全会终于结束。\footnote{Xia Qi
  claimed that the conference closed on August 23. See Zhang Yongshengde
  zuizheng cailiao, p.~68. The Hunan triple plenum ended the previous
  month. See CNS, 527, p.~4. 夏琦声称会议于8月23日闭幕。见《Zhang
  Yongshengde zuizheng cailiao》，第68页。湖南的三全会在前一个月结束。见
  CNS，第527期，第4页。}
为总结会议而发布的经过修饰的报告是一份妥协性文件。一方面，它敦促``犯错误的领导干部''{[}铁瑛和夏琦{]}``正确对待群众批评，认真进行自我批评，勇敢改正错误，振奋精神努力工作''。另一方面，它提到群众队伍中的``资产阶级派性''，要求有这种错误的人承认错误。公报的最后一点规定了党对所有其他组织，包括运动小组的领导权。公报说：

\begin{quote}
The movement to criticize Lin Biao and Confucius must be carried out
under the centralized leadership of the Party committees. Military
sub-district commands and people's armed forces departments are the
military departments of local Party committees. Trade unions,
associations of poor and lower-middle peasants, women's federations,
communist youth leagues, militia units and Red Guard organs are the mass
organizations led by the Party.\footnote{ZJRB, August 10,1974, pp.~1,4;
  quotations from ZPS, August 10,1974, SWB/FE/4682/BII/16-18.
  ZJRB，1974年8月10日，第1、4页；引文出自
  ZPS，1974年8月10日，SWB/FE/4682/BII/16-18。}
\end{quote}

\begin{quote}
批林批孔运动必须在党委的集中领导下进行。军分区、人武部是地方党委的军事部门。工会、贫下中农协会、妇联、共青团、民兵组织和红卫兵组织，是党领导下的群众组织。\footnote{ZJRB,
  August 10,1974, pp.~1,4; quotations from ZPS, August 10,1974,
  SWB/FE/4682/BII/16-18. ZJRB，1974年8月10日，第1、4页；引文出自
  ZPS，1974年8月10日，SWB/FE/4682/BII/16-18。}
\end{quote}

The ZPC convened a meeting on 13 August to press home this
message.\footnote{ZJRB, August 16,1974, p.~1; ZPS, August 16,1974,
  SWB/FE/4687/BII/9-10.
  ZJRB，1974年8月16日，第1页；ZPS，1974年8月16日，SWB/FE/4687/BII/9-10。}
It discussed the respective roles and relations between new and veteran
cadres, recommending that new cadres ``should be given opportunities to
work and their role should be brought into full play'' with the help of
their older colleagues. Speeches by Weng Senhe on 4 August and Zhang
Yongsheng on 5 August at the Zhejiang Workers Political College had been
less compromising. The bitterness which Document 16 of 1972 (see chapter
six) continued to arouse was evident in extracts from Weng's speech. Xia
Qi later wrote that the plenum had devoted much of its time to the issue
and to exactly what directives Zhou Enlai had issued at the April 1972
meeting in Beijing. Zhang Yongsheng attacked Xu Shiyou and praised the
Shanghai Garrison for exposing Xu's political errors. Zhang boasted that
he could not wait (恨不得) for the opportunity to resume his political
battle with Xu.\footnote{Weng Senhe bufen zuizheng cailiao, pp.~51, 68;
  Zhang Yongshengde zuizheng cailiao, p.~74. 《Weng Senhe bufen zuizheng
  cailiao》，第51、68页；《Zhang Yongshengde zuizheng
  cailiao》，第74页。}

浙江省委于8月13日召开会议，进一步强调这一信息。\footnote{ZJRB, August
  16,1974, p.~1; ZPS, August 16,1974, SWB/FE/4687/BII/9-10.
  ZJRB，1974年8月16日，第1页；ZPS，1974年8月16日，SWB/FE/4687/BII/9-10。}
会议讨论了新老干部各自的作用及其关系，建议在老干部帮助下，``应当给新干部以工作机会，充分发挥他们的作用''。翁森鹤8月4日和张永生8月5日在浙江工人政治大学的讲话则较少妥协色彩。从翁的讲话摘录中可以看出，1972年16号文件（见第六章）仍在激起强烈怨恨。夏琦后来写道，全会把大量时间用于讨论这个问题，以及周恩来在1972年4月北京会议上究竟发出了哪些指示。张永生攻击许世友，并称赞上海警备区揭露许的政治错误。张夸口说，他恨不得马上有机会重新同许展开政治斗争。\footnote{Weng
  Senhe bufen zuizheng cailiao, pp.~51, 68; Zhang Yongshengde zuizheng
  cailiao, p.~74. 《Weng Senhe bufen zuizheng
  cailiao》，第51、68页；《Zhang Yongshengde zuizheng
  cailiao》，第74页。}

The triple plenum also produced two internal reports. One, compiled on
the basis of the investigation of Xia Qi by the special case group,
concluded that his mistakes were sufficiently serious to be defined as
an antagonistic contradiction. The second report, dated 19
October\footnote{Weng Senhe bufen zuizheng cailiao, p.~6. 《Weng Senhe
  bufen zuizheng cailiao》，第6页。} and entitled ``A Summary of the 133
Days' Proceedings'', requested that the CC ``hold a relevant meeting at
the appropriate time to solve the problem of the Nanjing Units''. Wang
Hongwen and Zhang Chunqiao had previously directed that the plenum would
solve the question of which headquarters Zhejiang took its orders from -
Beijing or Nanjing. Weng Senhe had stated bluntly on 4 August 1974:

三全会还产生了两份内部报告。其中一份根据专案组对夏琦的调查写成，结论认为他的错误严重到足以被定性为敌我矛盾。第二份报告日期为10月19日，\footnote{Weng
  Senhe bufen zuizheng cailiao, p.~6. 《Weng Senhe bufen zuizheng
  cailiao》，第6页。}
题为《一百三十三天会议情况纪要》，请求中央``在适当时候召开有关会议，解决南京部队问题''。王洪文和张春桥此前曾指示，全会将解决浙江究竟听命于哪一总部的问题，是北京还是南京。翁森鹤在1974年8月4日直截了当地说：

\begin{quote}
With Xia Qi in the ZPMD representing the right deviationist wind to
reverse correct verdicts, there's still that big chief in Nanjing {[}Xu
Shiyou{]}. Every wind to reverse verdicts in Zhejiang has been tied up
with Nanjing \ldots; principally it is Xia Qi and his chief in Nanjing -
the root is that big Commander.\footnote{Ibid, p.~51. 同上，第51页。}
\end{quote}

\begin{quote}
浙江省军区有夏琦这个代表右倾翻案风的人，南京还有那个大头头{[}许世友{]}。浙江每一股翻案风都同南京有牵连\ldots\ldots 主要是夏琦和他在南京的头头，根子就是那个大司令。\footnote{Ibid,
  p.~51. 同上，第51页。}
\end{quote}

Weng was referring to Xu Shiyou and his responsibility for the
investigation of the rebels carried out in Zhejiang between 1972 and
1973. Zhang Yongsheng made the same links and insinuations in a speech
of 8 October. Jiang Qing went one step further in linking Xu to Zhou
Enlai.\footnote{Zhang Yongshengde zuizheng cailiao, p.~72. 《Zhang
  Yongshengde zuizheng cailiao》，第72页。} Tan, Tie and Xia were denied
the right to submit dissenting reports. However, they drafted a joint
letter to the CC, State Council and MAC setting forth their views, much
to the displeasure of the rebel leaders. Wang Hongwen then ordered the
provincial military leaders to Beijing and rebuked them for expressing
disagreement with the reports.\footnote{Ji Yangwen, HZRB, February
  9,1977; HZRB, September 14,1978; Tie Ying, HZRB, June 7, 1978.
  纪仰文，HZRB，1977年2月9日；HZRB，1978年9月14日；铁瑛，HZRB，1978年6月7日。}

翁森鹤指的是许世友，以及他对1972至1973年间在浙江开展的造反派调查所负的责任。张永生在10月8日的一次讲话中作出了同样的关联和影射。江青则更进一步，把许同周恩来联系起来。\footnote{Zhang
  Yongshengde zuizheng cailiao, p.~72. 《Zhang Yongshengde zuizheng
  cailiao》，第72页。}
谭、铁、夏被剥夺了提交不同意见报告的权利。不过，他们起草了一封致中央、国务院和中央军委的联名信，陈述自己的看法，这使造反派领导人非常不满。随后王洪文命令省军区领导人到北京，并斥责他们对报告表示不同意见。\footnote{Ji
  Yangwen, HZRB, February 9,1977; HZRB, September 14,1978; Tie Ying,
  HZRB, June 7, 1978.
  纪仰文，HZRB，1977年2月9日；HZRB，1978年9月14日；铁瑛，HZRB，1978年6月7日。}

\section{The Struggle to Rein in and Conclude the Campaign, April -
December
1974}\label{the-struggle-to-rein-in-and-conclude-the-campaign-april---december-1974}

\section{约束并结束运动的斗争，1974年4月至12月}\label{ux7ea6ux675fux5e76ux7ed3ux675fux8fd0ux52a8ux7684ux6597ux4e891974ux5e744ux6708ux81f312ux6708}

On 8 April 1974, the day after Wang Hongwen's latest intervention, the
ZPC and ZPRC held another rally to report on the progress of the
anti-Lin anti-Confucius campaign.\footnote{ZJRB, April 10,1974, p.~1;
  ZPS, April 10,1974, URS, 75: 14, pp.~184-7; CNS, 518 (May 22 1974),
  p.~1.
  ZJRB，1974年4月10日，第1页；ZPS，1974年4月10日，URS，75:14，第184-187页；CNS，第518期（1974年5月22日），第1页。}
Its main purpose was to denounce Xia Qi, following Wang Hongwen's
comment of the previous day or, in the words of the report, to peel off
all outer layers so as to ``get down to the bottom of class struggle''.
Zhang Yongsheng chaired the meeting, with Tan Qilong, Chai Qikun and
Weng Senhe the principal speakers. The meeting reported glowingly on the
general situation and praised the masses for overcoming obstacles in
maintaining the momentum of the campaign. Nevertheless, there were
indications that the opponents of the ``mountain top'' leaders were
putting up resistance. Reference to the ``sabotage and disturbance of
class enemies'' and the attempt to distinguish between the
``partisanship of the proletariat'' and the ``factionalism of the
bourgeoisie'' provided evidence that unrest in the province was serious.
That the campaign had adversely affected industrial and agricultural
production was tacitly admitted by the meeting's appeal to production
workers to ``make revolution after work''.

1974年4月8日，也就是王洪文最新一次干预的第二天，浙江省委和省革委会又召开大会，报告批林批孔运动的进展。\footnote{ZJRB,
  April 10,1974, p.~1; ZPS, April 10,1974, URS, 75: 14, pp.~184-7; CNS,
  518 (May 22 1974), p.~1.
  ZJRB，1974年4月10日，第1页；ZPS，1974年4月10日，URS，75:14，第184-187页；CNS，第518期（1974年5月22日），第1页。}
会议的主要目的，是按照王洪文前一天的评论批判夏琦；或者用报告中的话说，是剥去一切外层，以便``触及阶级斗争的底''。张永生主持会议，谭启龙、柴启琨和翁森鹤为主要发言人。会议对总体形势作了热情洋溢的汇报，赞扬群众克服障碍，保持了运动势头。不过，也有迹象表明，``山上派''领导人的反对者正在抵抗。报告提到``阶级敌人的破坏和捣乱''，并试图区分``无产阶级党性''和``资产阶级派性''，这说明省内动荡相当严重。会议呼吁生产工人``下班以后闹革命''，实际上也默认了运动已经对工农业生产造成不利影响。

The People's Daily published an editorial on the subject of the economy
on 10 April.\footnote{``抓批林批孔促工业生产'' (Grasp criticism of Lin
  and Confucius, promote industrial production) RMRB, April 10,1974,
  p.~1. 《抓批林批孔促工业生产》（Grasp criticism of Lin and Confucius,
  promote industrial production），RMRB，1974年4月10日，第1页。} On the
same day the CCP CC issued a notice specifying that the campaign must
come under the unified leadership of party committees and forbidding the
formation of fighting groups and inter-unit or inter-regional
link-ups.\footnote{Hao Mengbi and Duan Haoran, Liushinian, pp.~636-37.
  郝梦笔、段浩然，《Liushinian》，第636-637页。} The editorial and
notice were undoubtedly issued as a result of a meeting called by the
State Planning Commission, on the instructions of the Politburo, to
discuss industrial production with the representatives of fifteen
provincial-level administrations.\footnote{See Liu Suinian and Wu
  Qungan, ``Wenhua dageming'' shiqide guomin jingji, p.~78.
  见刘绥年、吴群敢，《``Wenhua dageming'' shiqide guomin
  jingji》，第78页。}

《人民日报》于4月10日发表了一篇关于经济问题的社论。\footnote{``抓批林批孔促工业生产''
  (Grasp criticism of Lin and Confucius, promote industrial production)
  RMRB, April 10,1974, p.~1. 《抓批林批孔促工业生产》（Grasp criticism
  of Lin and Confucius, promote industrial
  production），RMRB，1974年4月10日，第1页。}
同一天，中共中央发布通知，明确运动必须置于党委统一领导之下，并禁止成立战斗队和进行跨单位、跨地区串联。\footnote{Hao
  Mengbi and Duan Haoran, Liushinian, pp.~636-37.
  郝梦笔、段浩然，《Liushinian》，第636-637页。}
这篇社论和通知无疑是在国家计划委员会按照政治局指示召开的会议之后发布的；该会议同十五个省级行政区的代表讨论了工业生产问题。\footnote{See
  Liu Suinian and Wu Qungan, ``Wenhua dageming'' shiqide guomin jingji,
  p.~78. 见刘绥年、吴群敢，《``Wenhua dageming'' shiqide guomin
  jingji》，第78页。}

The Zhejiang authorities did not respond publicly to the central
initiative until 20 April when they called a telephone hook-up meeting
at which Tan, Zhang and Weng all spoke.\footnote{ZJRB, April 22, 1974,
  p.~1. ZJRB，1974年4月22日，第1页。} The three urged leading cadres at
all levels to handle correctly the relationship between revolution and
production and to strengthen party leadership and the unity of cadres
and citizens. More importantly, given the threat to stability posed by
the triple plenum, the meeting declared:

浙江当局直到4月20日才公开回应中央的举措，当时他们召开了一次电话会议，谭、张、翁均在会上讲话。\footnote{ZJRB,
  April 22, 1974, p.~1. ZJRB，1974年4月22日，第1页。}
三人敦促各级领导干部正确处理革命和生产的关系，加强党的领导以及干部和群众的团结。更重要的是，鉴于三全会对稳定构成威胁，会议宣布：

Meetings which can be postponed, should be postponed. Meetings which can
be brought to a conclusion for the time being should be brought to a
conclusion for the time being. Meetings which cannot be brought to a
conclusion right away should be brought to a conclusion in a few days.

可以推迟的会议，应当推迟。可以暂时结束的会议，应当暂时结束。不能马上结束的会议，应当在几天内结束。

This short report only hinted at the discontent felt by the party
leadership about the direction that the political campaign had taken.
Tan, with the backing of central leaders such as Zhou Enlai and Deng
Xiaoping, may have finally taken the fight up to the rebels. Mao's
statement of 1967, concerning the need for unity among the
working-class, was quoted and emphasis attached to the importance of
attaining targets set in the five-year plan.\footnote{CNA, 959 (May
  10,1974), p.~6. CNA，第959期（1974年5月10日），第6页。}

这份简短报告只是暗示了党委领导层对政治运动走向的不满。在周恩来、邓小平等中央领导人的支持下，谭启龙或许终于开始同造反派正面较量。报告引用了毛泽东1967年关于工人阶级必须团结的讲话，并强调完成五年计划所定指标的重要性。\footnote{CNA,
  959 (May 10,1974), p.~6. CNA，第959期（1974年5月10日），第6页。}

Chen Xia, who was a participant at the 20 April meeting, later claimed
that criticism was directed against factional fighting and the
activities of the militia. According to Chen, the CCP ZPC subsequently
issued a series of directives attempting to rein in the campaign. These
included a request to bring it to an end in basic production units
(``the key points of the movement are at and above county-level
offices''), to stop the indiscriminate public naming of leading cadres
without the approval of relevant party committees, to steer away from
the obsession with personalities and power struggles (``don't haul out
representatives of restorationist forces at all levels'') and to focus
attention back on the economy. However, Wang Xing, the leader of the
Hangzhou campaign small group, personally brought out a document
rebutting these directives point by point. He described the restrictions
and regulations laid down by the ZPC as ``setting rules and
regulations'', and ``binding hands and feet''. The practice of naming
cadres continued unabated and the renewed emphasis on economic affairs
was described as ``pitting revolution against production''.\footnote{Chen
  Xia, HZRB, April 10, 1978. 陈霞，HZRB，1978年4月10日。}

参加4月20日会议的陈夏后来声称，批评的矛头指向派性斗争和民兵活动。据陈夏说，中共浙江省委随后发布了一系列指示，试图约束运动。这些指示包括要求运动在基层生产单位结束（``运动的重点在县级机关以上''），未经有关党委批准不得随意公开点名领导干部，避免沉迷于个人问题和权力斗争（``不要层层揪复辟势力代表''），并把注意力重新集中到经济上。然而，杭州运动小组负责人王兴亲自拿出一份文件，逐条驳斥这些指示。他把省委规定的限制和规章称为``清规戒律''和``束手束脚''。点名干部的做法有增无减，而重新强调经济事务则被说成是``以生产压革命''。\footnote{Chen
  Xia, HZRB, April 10, 1978. 陈霞，HZRB，1978年4月10日。}

Whatever the maneuverings which went on behind the scenes, the effect on
the rebels seemed minimal and short-lived. On 26 April 1974 the ZPC and
ZPRC held yet another conference at which the political rhetoric reached
a new peak of fervor and extravagance.\footnote{ZJRB, April 28, 1974,
  p.~1; ZPS, April 28, 1974, SWB/FE/4589/BII/1-2; CNA, 959, p.~7.
  ZJRB，1974年4月28日，第1页；ZPS，1974年4月28日，SWB/FE/4589/BII/1-2；CNA，第959期，第7页。}
Tan Qilong was not present at all, or if so, he did not speak. Zhang
Yongsheng presided, with Weng Senhe, Shen Ce, He Xianchun, Chen Bing,
the peasant philosopher Jiang Ruwang (see chapter five),\footnote{For
  further media coverage of Jiang's Jinjian brigade between 1973 and
  1976, see ZPS November 23, 1973, SWB/FE/4464/BII/12-13; ZJRB, February
  28,1974, pp.~1, 3; ZPS, March 1,1974, SWB/FE/4558/BII/7-8; Xinhua,
  April 1, 1974, FBIS/CHI, 64 (1974), G1-2; ZPS, June 29, 1974,
  SWB/FE/4642/BII/5; ZPS, July 24, 1974, FBIS/CHI, 144 (1974), G2-3;
  ZPS, September 1, 1974, SWB/FE/4699/BII/14-16; ZPS, February 17, 1975,
  FBIS/CHI, 35 (1975), G2-4; ZPS, December 26,1975,
  SWB/FE/5100/BII/8-11; ZPS, June 2,1976, SWB/FE/5230/BII/7-8.
  关于1973年至1976年间媒体对江的金建大队的进一步报道，见
  ZPS，1973年11月23日，SWB/FE/4464/BII/12-13；ZJRB，1974年2月28日，第1、3页；ZPS，1974年3月1日，SWB/FE/4558/BII/7-8；新华社，1974年4月1日，FBIS/CHI，第64期（1974），G1-2；ZPS，1974年6月29日，SWB/FE/4642/BII/5；ZPS，1974年7月24日，FBIS/CHI，第144期（1974），G2-3；ZPS，1974年9月1日，SWB/FE/4699/BII/14-16；ZPS，1975年2月17日，FBIS/CHI，第35期（1975），G2-4；ZPS，1975年12月26日，SWB/FE/5100/BII/8-11；ZPS，1976年6月2日，SWB/FE/5230/BII/7-8。}
and the former Red Storm leader and Deputy-secretary of the Zhejiang CYL
Liu Ying all featured speakers. The consensus of opinion which emerged
from the meeting was that the campaign was ``raging like a prairie fire
throughout the province'', although the struggle remained ``protracted,
complex and arduous''. The speakers vented their spleen on the
``bourgeois forces of restoration'' who intrigued and conspired,
provoked sectarianism and division, ``tried to strangle new-born
things'' and even slandered the ``great leader of the proletariat with
bitter animosity''. From the notice drafted by Weng Senhe for the
meeting we can see that the verbal abuse was directed principally at Tie
Ying and Xia Qi.\footnote{Weng Senhe bufen zuizheng cailiao, p.~43.
  《Weng Senhe bufen zuizheng cailiao》，第43页。}

无论幕后进行了怎样的周旋，对造反派的影响似乎都很小，而且持续时间很短。1974年4月26日，浙江省委和省革委会又召开一次大会，会上的政治辞令达到狂热和夸张的新高峰。\footnote{ZJRB,
  April 28, 1974, p.~1; ZPS, April 28, 1974, SWB/FE/4589/BII/1-2; CNA,
  959, p.~7.
  ZJRB，1974年4月28日，第1页；ZPS，1974年4月28日，SWB/FE/4589/BII/1-2；CNA，第959期，第7页。}
谭启龙根本没有出席；即便出席，也没有讲话。张永生主持会议，翁森鹤、沈策、贺贤春、陈冰、农民哲学家姜汝旺（见第五章）\footnote{For
  further media coverage of Jiang's Jinjian brigade between 1973 and
  1976, see ZPS November 23, 1973, SWB/FE/4464/BII/12-13; ZJRB, February
  28,1974, pp.~1, 3; ZPS, March 1,1974, SWB/FE/4558/BII/7-8; Xinhua,
  April 1, 1974, FBIS/CHI, 64 (1974), G1-2; ZPS, June 29, 1974,
  SWB/FE/4642/BII/5; ZPS, July 24, 1974, FBIS/CHI, 144 (1974), G2-3;
  ZPS, September 1, 1974, SWB/FE/4699/BII/14-16; ZPS, February 17, 1975,
  FBIS/CHI, 35 (1975), G2-4; ZPS, December 26,1975,
  SWB/FE/5100/BII/8-11; ZPS, June 2,1976, SWB/FE/5230/BII/7-8.
  关于1973年至1976年间媒体对江的金建大队的进一步报道，见
  ZPS，1973年11月23日，SWB/FE/4464/BII/12-13；ZJRB，1974年2月28日，第1、3页；ZPS，1974年3月1日，SWB/FE/4558/BII/7-8；新华社，1974年4月1日，FBIS/CHI，第64期（1974），G1-2；ZPS，1974年6月29日，SWB/FE/4642/BII/5；ZPS，1974年7月24日，FBIS/CHI，第144期（1974），G2-3；ZPS，1974年9月1日，SWB/FE/4699/BII/14-16；ZPS，1975年2月17日，FBIS/CHI，第35期（1975），G2-4；ZPS，1975年12月26日，SWB/FE/5100/BII/8-11；ZPS，1976年6月2日，SWB/FE/5230/BII/7-8。}，以及前``红暴''领导人、浙江共青团副书记刘英，都是主要发言人。会议形成的共识是，运动已经``在全省形成燎原之势''，尽管斗争仍然``长期、复杂、艰巨''。发言者把怒火倾泻到``资产阶级复辟势力''身上，指责他们勾心斗角、阴谋诡计，挑动宗派主义和分裂，``妄图扼杀新生事物''，甚至``怀着刻骨仇恨''诬蔑``无产阶级伟大领袖''。从翁森鹤为会议起草的通知可以看出，这些辱骂主要指向铁瑛和夏琦。\footnote{Weng
  Senhe bufen zuizheng cailiao, p.~43. 《Weng Senhe bufen zuizheng
  cailiao》，第43页。}

At the symposium sponsored by the ZPTUC and the HMWC on 29 April 1974,
the eve of May Day, Chairwoman of the ZPTUC Jiang Baodi presided. The
main speakers were Zhang Yongsheng, Hua Yinfeng and He Xianchun, all
former United Headquarters' activists.\footnote{ZJRB, April 30,1974,
  p.~1. ZJRB，1974年4月30日，第1页。} He Xianchun repeated Mao's dictum
concerning the unity of the working class, but in the words of China
News Summary, ``no comment was made, and no enthusiasm for unity was
shown''.\footnote{CNS, 519 (May 29, 1974), p.~1.
  CNS，第519期（1974年5月29日），第1页。} Representatives of major
industrial plants in Hangzhou, including Weng Senhe's silk complex, the
iron and steel mill, and the oxygen generator plant attended and spoke.
The urban militia and the propaganda teams also supplied speakers. The
next day 600 workers from the propaganda teams were despatched to
government departments and tertiary institutions to ``strengthen further
worker leadership over the various fields of the
superstructure''.\footnote{ZPS, April 30,1974, SWB/FE/4603/BII/20. A
  rally in Hangzhou on July 26 celebrated the sixth anniversary of the
  propaganda teams entering public institutions. ZPS, July 27, 1974,
  SWB/FE/4676/BII/13.
  ZPS，1974年4月30日，SWB/FE/4603/BII/20。7月26日，杭州举行集会，庆祝宣传队进驻公共机构六周年。ZPS，1974年7月27日，SWB/FE/4676/BII/13。}
Whether the propaganda teams were sent back into schools to stop student
confrontations with authority and to ``guide'' the anti-Lin
anti-Confucius campaign, as an observer in Hong Kong claimed,\footnote{CNS,
  521 (June 12,1974). CNS，第521期（1974年6月12日）。} or whether their
mission was to enlist students on the side of the ``mountain top''
faction is unclear.

1974年4月29日，即五一节前夕，浙江省总工会和杭州市工代会主办座谈会，由省总工会主席蒋宝娣主持。主要发言者是张永生、华银凤和贺贤春，三人都是前联合总部积极分子。\footnote{ZJRB,
  April 30,1974, p.~1. ZJRB，1974年4月30日，第1页。}
贺贤春重复了毛泽东关于工人阶级团结的论断，但用《中国新闻摘要》的话说，``没有作评论，也没有显示出团结的热情''。\footnote{CNS,
  519 (May 29, 1974), p.~1. CNS，第519期（1974年5月29日），第1页。}
杭州主要工业企业的代表，包括翁森鹤所在的丝绸联合企业、钢铁厂和制氧机厂，均出席并发言。城市民兵和宣传队也派人发言。次日，宣传队的600名工人被派往政府部门和高等院校，以``进一步加强工人阶级对上层建筑各个领域的领导''。\footnote{ZPS,
  April 30,1974, SWB/FE/4603/BII/20. A rally in Hangzhou on July 26
  celebrated the sixth anniversary of the propaganda teams entering
  public institutions. ZPS, July 27, 1974, SWB/FE/4676/BII/13.
  ZPS，1974年4月30日，SWB/FE/4603/BII/20。7月26日，杭州举行集会，庆祝宣传队进驻公共机构六周年。ZPS，1974年7月27日，SWB/FE/4676/BII/13。}
宣传队究竟是如香港一位观察者所说，被重新派进学校以阻止学生同当局对抗并``引导''批林批孔运动，\footnote{CNS,
  521 (June 12,1974). CNS，第521期（1974年6月12日）。}
还是其任务在于把学生争取到``山上派''一边，目前并不清楚。

In May 1974, Zhang, Weng and He published a joint article in Zhejiang
Daily to commemorate the eighth anniversary of the 1966 May 16 Circular
which had launched the Cultural Revolution.\footnote{ZJRB, May 16,1974,
  p.~3. ZJRB，1974年5月16日，第3页。} In their article the rebel trio
were at pains to associate themselves and the current campaign with the
mainstream of the Cultural Revolution, to dissociate themselves from any
past affiliations with Lin Biao and his ``agents'' in
Zhejiang\footnote{See also the article on this theme written by Zhang's
  Fine Arts College in ZJRB, January 2, 1975, p.~3.
  另见张的美术学院围绕这一主题撰写的文章，载 ZJRB，1975年1月2日，第3页。}
and to refute accusations that they were ``ultra-left'',''three anti's
elements'' (三反分子) that is, anti Mao, the PLA and revolutionary
committees, or May 16 elements. It is clear that the old mass
organization rebel leaders viewed the 1974 campaign as a second Cultural
Revolution. The 1973 phrase ``going against the tide'' was equated in
their minds with the 1966 slogan ``it is right to rebel''.\footnote{See
  the articles from the Zhang/Weng/He controlled newspaper Zhejiang
  Worker on this theme reprinted in ZJRB, April 23,1974, p.~2 and May
  22, 1974, pp.~1, 2. 见由张/翁/何控制的报纸《Zhejiang Worker》发表并被
  ZJRB 转载的相关文章，1974年4月23日第2页、1974年5月22日第1、2页。} An
article by former activists of United Headquarters argued that the
campaign to criticize Lin Biao had already achieved considerable results
while the efforts to tackle Confucius had only just begun and was a much
more formidable task.\footnote{ZJRB, May 30, 1974, p.~3; Zhang
  Yongshengde zuizheng cailiao, pp.~2, 9,10.
  ZJRB，1974年5月30日，第3页；《Zhang Yongshengde zuizheng
  cailiao》，第2、9、10页。} The three authors of the article were Li
Xiantong, (Hangzhou University) Yan Yihuan (Zhejiang University) and Du
Yingxin (Zhejiang Fine Arts College), all leading activists of United
Headquarters and, as this chapter has detailed, prominent members of the
``mountain top'' faction. During the anti-Lin anti-Confucius campaign
radicals directed their attacks primarily against Confucius, whom they
equated with restoration, while moderates focused on Lin Biao who,
despite Mao's December 1972 comment describing Lin as an ultra-rightist,
was associated in everyone's mind with the ultra-left extremism of the
Cultural Revolution.

1974年5月，张、翁、贺在《浙江日报》发表联名文章，纪念发动文化大革命的1966年``五一六通知''八周年。\footnote{ZJRB,
  May 16,1974, p.~3. ZJRB，1974年5月16日，第3页。}
在文章中，这三名造反派煞费苦心地把自己和当前运动同文化大革命的主流联系起来，撇清自己过去同林彪及其在浙江``代理人''的任何关系，\footnote{See
  also the article on this theme written by Zhang's Fine Arts College in
  ZJRB, January 2, 1975, p.~3.
  另见张的美术学院围绕这一主题撰写的文章，载 ZJRB，1975年1月2日，第3页。}
并反驳说他们是``极左''、``三反分子''（三反分子，即反毛、反解放军、反革委会）或``五一六分子''的指控。很明显，这些旧群众组织造反派领导人把1974年的运动视为第二次文化大革命。在他们心目中，1973年的``反潮流''一语等同于1966年的``造反有理''口号。\footnote{See
  the articles from the Zhang/Weng/He controlled newspaper Zhejiang
  Worker on this theme reprinted in ZJRB, April 23,1974, p.~2 and May
  22, 1974, pp.~1, 2. 见由张/翁/何控制的报纸《Zhejiang Worker》发表并被
  ZJRB 转载的相关文章，1974年4月23日第2页、1974年5月22日第1、2页。}
前联合总部积极分子的一篇文章认为，批判林彪的运动已经取得相当成果，而对孔子的斗争才刚刚开始，并且是一项艰巨得多的任务。\footnote{ZJRB,
  May 30, 1974, p.~3; Zhang Yongshengde zuizheng cailiao, pp.~2, 9,10.
  ZJRB，1974年5月30日，第3页；《Zhang Yongshengde zuizheng
  cailiao》，第2、9、10页。}
这篇文章的三名作者是李显通（杭州大学）、严怡寰（浙江大学）和杜英信（浙江美术学院），他们都是联合总部的重要积极分子，并且正如本章所详述的，是``山上派''的重要成员。在批林批孔运动中，激进派主要把攻击矛头指向孔子，并把孔子等同于复辟；温和派则把重点放在林彪身上，尽管毛泽东1972年12月曾评论说林是极右分子，但在所有人的心目中，林彪仍同文化大革命的极左极端主义联系在一起。

After a hiatus in May, public meetings and rallies resumed in June 1974.
At a meeting which was organized by the ZPC and ZPRC and sponsored by
the political work section of the ZPRC at the beginning of that month,
further evidence that Zhang, Weng and He had acquired enormous influence
within the provincial leadership emerged.\footnote{ZJRB, June 6,1974,
  pp.~1,3; ZPS, June 6,1974, SWB/FE/4625/BII/13-15.
  ZJRB，1974年6月6日，第1、3页；ZPS，1974年6月6日，SWB/FE/4625/BII/13-15。}
He Xianchun presided over the meeting while Zhang and Weng, with Tan
Qilong and Chai Qikun, were the principal speakers. The meeting
expressed satisfaction at the development of the campaign over the
previous months but cautioned that ``the struggle is still acute and
complicated''. It also warned the audience not to become complacent ``in
the face of the excellent situation''.

5月短暂停顿之后，公开会议和大会在1974年6月恢复。当月初，浙江省委和省革委会组织、省革委会政治工作组主办的一次会议上，进一步出现了证据，表明张、翁、贺已在省级领导层中获得巨大影响。\footnote{ZJRB,
  June 6,1974, pp.~1,3; ZPS, June 6,1974, SWB/FE/4625/BII/13-15.
  ZJRB，1974年6月6日，第1、3页；ZPS，1974年6月6日，SWB/FE/4625/BII/13-15。}
贺贤春主持会议，张和翁以及谭启龙、柴启琨是主要发言人。会议对过去几个月运动的发展表示满意，但也告诫说``斗争仍然尖锐复杂''。会议还警告与会者，``在大好形势面前''不要自满。

One of the three principal tasks of the urban militia, as specified in
the September 1973 People's Daily editorial, was political education of
the working class. To this end, on 5 June 1974 a school to train cadres
and theorists from the workers of Zhejiang opened in
Hangzhou.\footnote{ZJRB, June 8,1974, p.~1. Beijing published an
  editorial on this subject in mid-June. See ``在斗争中培养理论队伍''
  (Train a theoretical contingent in the midst of struggle), RMRB, June
  18,1974, p.~1.
  ZJRB，1974年6月8日，第1页。北京在6月中旬发表了一篇关于这一主题的社论。见《在斗争中培养理论队伍》（Train
  a theoretical contingent in the midst of
  struggle），RMRB，1974年6月18日，第1页。} It was named the Zhejiang
Workers' Political School and in its first semester enrolled 220
students. The ZPTUC, with the approval of the provincial party
committee, had organized the school and set it the task of producing
``experts for the Party who can write as well as fight''. Zhejiang had
modelled its school on a similar institution in Shanghai.\footnote{See
  Gao Gao and Yan Jiaqi, ``Wenhua dageming'', pp.~533-34.
  见高皋、严家其，《``Wenhua dageming''》，第533-534页。} The criteria
for selection to study at the school included possessing ``a high
awareness of the struggle between the two lines,\ldots{} {[}courage{]}
to go against the tide and \ldots{} {[}being{]} activists in
revolutionary mass criticism''. Weng Senhe, in his capacity as the
``responsible person'' of the school, delivered the speech at the
opening ceremony which was attended by Tan Qilong and Jiang Baodi among
others. The school was situated on Five Terrace mountain (五台山) and
issued such publications as The Journal (学报) and Brief Biographies of
Legalists (法家人物小传).\footnote{For articles published by the
  college's writing group, see ZJRB, June 25,1974, p.~2, and July 15,
  1974, p.~2. 关于该学院写作组发表的文章，见
  ZJRB，1974年6月25日第2页，以及1974年7月15日第2页。} In addition to
these publications, the rebels also ran their own newspaper, the
Zhejiang Worker (浙江工人), which had been United Headquarters'
mouthpiece in the 1960s, and other journals.\footnote{For further
  discussion see CNS, 522 (June 19,1974), pp.~1-5. 进一步讨论见
  CNS，第522期（1974年6月19日），第1-5页。}

按照1973年9月《人民日报》社论的规定，城市民兵的三项主要任务之一，是对工人阶级进行政治教育。为此，1974年6月5日，一所培训浙江工人干部和理论工作者的学校在杭州开办。\footnote{ZJRB,
  June 8,1974, p.~1. Beijing published an editorial on this subject in
  mid-June. See ``在斗争中培养理论队伍'' (Train a theoretical contingent
  in the midst of struggle), RMRB, June 18,1974, p.~1.
  ZJRB，1974年6月8日，第1页。北京在6月中旬发表了一篇关于这一主题的社论。见《在斗争中培养理论队伍》（Train
  a theoretical contingent in the midst of
  struggle），RMRB，1974年6月18日，第1页。}
该校名为浙江工人政治学校，第一学期招收学生220人。浙江省总工会经省委批准组织了这所学校，并规定其任务是培养``既能文又能武的党的专家''。浙江是仿照上海一所类似机构建立这所学校的。\footnote{See
  Gao Gao and Yan Jiaqi, ``Wenhua dageming'', pp.~533-34.
  见高皋、严家其，《``Wenhua dageming''》，第533-534页。}
入校学习的选拔标准包括具有``较高的两条路线斗争觉悟，\ldots\ldots{[}敢于{]}反潮流，\ldots\ldots{[}并且是{]}革命大批判积极分子''。翁森鹤以学校``负责人''身份在开学典礼上讲话，谭启龙、蒋宝娣等人出席。学校位于五台山（五台山），出版《学报》（学报）和《法家人物小传》（法家人物小传）等刊物。\footnote{For
  articles published by the college's writing group, see ZJRB, June
  25,1974, p.~2, and July 15, 1974, p.~2. 关于该学院写作组发表的文章，见
  ZJRB，1974年6月25日第2页，以及1974年7月15日第2页。}
除这些出版物外，造反派还经营自己的报纸《浙江工人》（浙江工人）以及其他刊物；《浙江工人》曾是20世纪60年代联合总部的喉舌。\footnote{For
  further discussion see CNS, 522 (June 19,1974), pp.~1-5. 进一步讨论见
  CNS，第522期（1974年6月19日），第1-5页。}

The rebels realized that a disastrous slump in industrial production in
Zhejiang would discredit the provincial leadership and that the longer
they kept up their tactics of destabilization the greater would be the
deleterious effects on the economy. At the conference held on 13 June to
report on the campaign it was asserted that the criticism of Lin and
Confucius had actually promoted industrial and agricultural production.
This rash assessment concerning the state of the economy was tempered,
however, by a request to the people of the province to apply their
``revolutionary enthusiasm'' to productive endeavor. After a lengthy
absence from the public eye during the campaign, Deputy-secretary of the
CCP ZPC Lai Keke reappeared at the conference. His article concerning
the campaign in Zhejiang Daily at the end of the month advocated a
greater commitment to study.\footnote{See ZJRB, June 28,1974, p.~2.
  Another member of the ZPC standing committee, Meng Zhaoyu, also
  resumed public activities in June 1974. 见
  ZJRB，1974年6月28日，第2页。中共浙江省委常委会另一名成员孟昭煜也于1974年6月恢复公开活动。}
There were signs that a wind-up of the radicalism of the first half of
1974 was imminent. The triple plenum temporarily closed its doors on 9
June, almost two months after the ZPC had issued instructions to this
effect at the 20 April telephone conference, and the achievements in the
campaign which were enumerated at the conference read like a valedictory
notice.\footnote{ZJRB, June 15,1974, p.~1. ZJRB，1974年6月15日，第1页。}
This impression was further reinforced by a photographic exhibition on
the campaign which opened on 21 June 1974 at the Exhibition Building in
Hangzhou. Tan Qilong cut the ribbon after Weng Senhe had delivered the
speech at the opening ceremony.\footnote{ZJRB, June 22,1974, p.~1.
  ZJRB，1974年6月22日，第1页。}

造反派意识到，浙江工业生产的灾难性下滑会损害省级领导层的信誉，而且他们维持破坏稳定策略的时间越长，对经济的有害影响就越大。在6月13日召开的运动汇报会上，会议声称批林批孔实际上促进了工农业生产。然而，这一关于经济状况的轻率判断，又被要求全省人民把``革命热情''投入生产劳动的号召所缓和。在运动期间长期淡出公众视野之后，中共浙江省委副书记赖可可在会上重新露面。月底，他在《浙江日报》发表的关于运动的文章主张进一步加强学习。\footnote{See
  ZJRB, June 28,1974, p.~2. Another member of the ZPC standing
  committee, Meng Zhaoyu, also resumed public activities in June 1974.
  见
  ZJRB，1974年6月28日，第2页。中共浙江省委常委会另一名成员孟昭煜也于1974年6月恢复公开活动。}
有迹象表明，1974年上半年激进主义即将收场。三全会于6月9日暂时闭会，距省委在4月20日电话会议上发出相关指示已近两个月；会议列举的运动成就读起来像一份告别通知。\footnote{ZJRB,
  June 15,1974, p.~1. ZJRB，1974年6月15日，第1页。}
1974年6月21日在杭州展览馆开幕的运动图片展进一步强化了这一印象。翁森鹤在开幕式上讲话后，谭启龙为展览剪彩。\footnote{ZJRB,
  June 22,1974, p.~1. ZJRB，1974年6月22日，第1页。}

However, the rebels were not going to succumb without a final outburst
of militancy. In this determination they were perhaps assisted by Jiang
Qing's hint in a speech of 14 June that the campaign should criticize
the ``contemporary Confucius'' and the ``big Confucian in the
party''\footnote{Hao Mengbi and Duan Haoran, Liushinian, p.~635; Hu Hua
  (ed.), Zhongguo shehuizhuyi geming, p.~317.
  郝梦笔、段浩然，《Liushinian》，第635页；胡华编，《Zhongguo
  shehuizhuyi geming》，第317页。} {[}Zhou Enlai{]}. A belligerent 24
June 1974 editorial of Zhejiang Daily advised party leaders to discard
their perfunctory approach to the campaign and make up their minds
whether to ``carry out reforms or to maintain the old, whether to
advance or to retrogress, and whether to make revolution or to restore
the old order''. The editorial commented:

然而，造反派不会不作最后一次战斗姿态便屈服。他们的这种决心，或许得到了江青6月14日讲话中暗示的助推；她暗示运动应批判``当代孔子''和``党内大儒''\footnote{Hao
  Mengbi and Duan Haoran, Liushinian, p.~635; Hu Hua (ed.), Zhongguo
  shehuizhuyi geming, p.~317.
  郝梦笔、段浩然，《Liushinian》，第635页；胡华编，《Zhongguo
  shehuizhuyi geming》，第317页。}
{[}周恩来{]}。1974年6月24日，《浙江日报》发表一篇充满斗争性的社论，劝告党的领导人抛弃对运动敷衍塞责的态度，并下定决心，究竟是``改革还是守旧，是前进还是倒退，是革命还是复辟''。社论评论说：

\begin{quote}
The broad revolutionary masses have risen up in rebellion against a
handful of capitalist-roaders within the Party and have dared to go
against the tide. Why have some comrades regarded this as ``offending
one's superiors and creating havoc''? Why have some comrades felt
themselves incompatible with new things and new cadres that have emerged
in the great proletarian cultural revolution while always thinking of
the old? In every major struggle between the two lines, why do some
comrades fail to understand the struggle in the beginning, stand opposed
to the masses, and fail to take a step forward until others remind them
to do so?\footnote{ZJRB, June 24,1974, p.~1; quotations from ZPS, June
  24,1974, SWB/FE/4637/B/II/21. ZJRB，1974年6月24日，第1页；引文出自
  ZPS，1974年6月24日，SWB/FE/4637/B/II/21。}
\end{quote}

\begin{quote}
广大革命群众已经起来造党内一小撮走资本主义道路当权派的反，敢于反潮流。为什么有些同志把这看成是``犯上作乱''？为什么有些同志总想着旧的，对无产阶级文化大革命中涌现出来的新生事物和新干部感到格格不入？在每一次重大的两条路线斗争中，为什么有些同志一开始总是不理解斗争，站到群众的对立面，直到别人提醒才向前迈一步？\footnote{ZJRB,
  June 24,1974, p.~1; quotations from ZPS, June 24,1974,
  SWB/FE/4637/B/II/21. ZJRB，1974年6月24日，第1页；引文出自
  ZPS，1974年6月24日，SWB/FE/4637/B/II/21。}
\end{quote}

A militant rally whose participants were ``filled with angry sentiments
and hatred for the common enemy'' was held on 30 June.\footnote{ZJRB,
  July 2,1974, p.~1; ZPS, July 2,1974, SWB/FE/4644/BII/9-10.
  ZJRB，1974年7月2日，第1页；ZPS，1974年7月2日，SWB/FE/4644/BII/9-10。}
The rally praised leading cadres who, in increasing numbers, had ``stood
at the forefront of the struggle, boldly led the movement and fought
together with the masses''. A full complement of provincial, municipal
and PLA leaders heard Tan Qilong, Chai Qikun and Weng Senhe lambast
their colleagues, presumably Tie Ying and Xia Qi, who had not qualified
for the above plaudits. Thus, the rebel cadres in Zhejiang retained the
initiative right up until the time when central authorities decided that
the cost of the campaign to the Chinese economy had reached dangerous
proportions.

6月30日召开了一次战斗性大会，与会者``义愤填膺，同仇敌忾''。\footnote{ZJRB,
  July 2,1974, p.~1; ZPS, July 2,1974, SWB/FE/4644/BII/9-10.
  ZJRB，1974年7月2日，第1页；ZPS，1974年7月2日，SWB/FE/4644/BII/9-10。}
大会赞扬越来越多``站在斗争前列，大胆领导运动，同群众一起战斗''的领导干部。省、市和解放军领导人全员出席，听取谭启龙、柴启琨和翁森鹤痛斥那些不配得到上述赞扬的同僚，推测主要是铁瑛和夏琦。因此，直到中央当局认定这场运动给中国经济造成的代价已经达到危险程度之前，浙江的造反派干部始终保持着主动权。
\#\# Central Efforts to End the Campaign \#\# 终止运动的中央努力

On 1 July 1974, the CCP Central Committee issued an important document
to bring the anti-Lin anti-Confucius campaign to an end.\footnote{CCP
  CC, zhongfa (1974), No.~21, in Classified Chinese Communist Documents,
  pp.~612-16. 中共中央，中发（1974）21号，载 Classified Chinese
  Communist Documents，第612-616页。} The document lamented the
shortfalls in key production targets, especially in railway freight and
coal production. Iron and steel as well as chemical fertilizer output
and production for the military sector was also lagging. Receipts for
the first five months of 1974 had fallen 500 million yuan short of
outlays and the shortage of goods had created a tight market for
consumers.\footnote{Production figures for the first half of 1974 are
  given in Liu Suinian and Wu Qungan, ``Wenhua dageming'' shiqide guomin
  jingji, pp.~76-8, and Hao Mengbi and Duan Haoran, Liushinian, p.~636,
  fn. 1, which cites a State Planning Commission outline report of 18
  June 1974. The date is given in Liu Suinian and Wu Qungan, ``Wenhua
  dageming'' shiqide guomin jingji, p.~167.
  1974年上半年的生产数字见刘绥年、吴群敢《``Wenhua dageming'' shiqide
  guomin
  jingji》，第76-78页，以及郝梦笔、段浩然《Liushinian》，第636页注1，后者引用了国家计划委员会1974年6月18日的一份概要报告。该日期见刘绥年、吴群敢《``Wenhua
  dageming'' shiqide guomin jingji》，第167页。} Document 21 also
subjected two key slogans of the campaign to analysis. The equation of
``going against the tide'' with rebellion against party leadership was
singled out for critical attention. Most cadres, argued the circular,
were good and to arrest or beat them even if they had made mistakes did
not conform with party policy. The CC ordered officials who had deserted
their posts to return to work and urged workers to go back to their
factories.\footnote{In a visit to the Hangzhou Oxygen Generator Plant I
  was informed that in 1974 one of its cadres had fled to Xinjiang to
  escape the urban militia. Visit, June 4,1978. In Nanjing sentdown
  workers took advantage of the campaign to protest on the city streets
  for three days, holding up trains and disrupting traffic. They
  demanded approval to present their grievances to the central
  authorities. Beijing instructed the Jiangsu authorities to prevent
  this from happening and tried to reassure the protestors with a
  promise that the question of their return to the city would be
  considered at the end of the campaign. See Jiangsusheng dashiji,
  p.~327.
  我在访问杭州制氧机厂时获悉，1974年该厂一名干部为躲避城市民兵而逃往新疆。访问日期：1978年6月4日。在南京，下放工人利用这场运动，在城市街头抗议三天，拦阻火车并扰乱交通。他们要求获准向中央当局陈述冤情。北京指示江苏当局阻止此事发生，并试图以承诺在运动结束时考虑他们返城问题来安抚抗议者。见《Jiangsusheng
  dashiji》，第327页。}

1974年7月1日，中共中央发布了一份重要文件，以结束批林批孔运动。\footnote{CCP
  CC, zhongfa (1974), No.~21, in Classified Chinese Communist Documents,
  pp.~612-16. 中共中央，中发（1974）21号，载 Classified Chinese
  Communist Documents，第612-616页。}
该文件对若干关键生产指标未能完成表示遗憾，尤其是铁路货运和煤炭生产。钢铁、化肥产量以及军工部门的生产也都落后。1974年前五个月，财政收入比支出少了5亿元，商品短缺造成了消费市场紧张。\footnote{Production
  figures for the first half of 1974 are given in Liu Suinian and Wu
  Qungan, ``Wenhua dageming'' shiqide guomin jingji, pp.~76-8, and Hao
  Mengbi and Duan Haoran, Liushinian, p.~636, fn. 1, which cites a State
  Planning Commission outline report of 18 June 1974. The date is given
  in Liu Suinian and Wu Qungan, ``Wenhua dageming'' shiqide guomin
  jingji, p.~167. 1974年上半年的生产数字见刘绥年、吴群敢《``Wenhua
  dageming'' shiqide guomin
  jingji》，第76-78页，以及郝梦笔、段浩然《Liushinian》，第636页注1，后者引用了国家计划委员会1974年6月18日的一份概要报告。该日期见刘绥年、吴群敢《``Wenhua
  dageming'' shiqide guomin jingji》，第167页。}
第21号文件还对这场运动中的两个关键口号进行了分析。把``反潮流''等同于反对党的领导这一说法，被特别拿出来加以批评。通报认为，多数干部是好的，即使犯了错误，逮捕或殴打他们也不符合党的政策。中央命令擅离职守的干部返回工作岗位，并敦促工人回到工厂。\footnote{In
  a visit to the Hangzhou Oxygen Generator Plant I was informed that in
  1974 one of its cadres had fled to Xinjiang to escape the urban
  militia. Visit, June 4,1978. In Nanjing sentdown workers took
  advantage of the campaign to protest on the city streets for three
  days, holding up trains and disrupting traffic. They demanded approval
  to present their grievances to the central authorities. Beijing
  instructed the Jiangsu authorities to prevent this from happening and
  tried to reassure the protestors with a promise that the question of
  their return to the city would be considered at the end of the
  campaign. See Jiangsusheng dashiji, p.~327.
  我在访问杭州制氧机厂时获悉，1974年该厂一名干部为躲避城市民兵而逃往新疆。访问日期：1978年6月4日。在南京，下放工人利用这场运动，在城市街头抗议三天，拦阻火车并扰乱交通。他们要求获准向中央当局陈述冤情。北京指示江苏当局阻止此事发生，并试图以承诺在运动结束时考虑他们返城问题来安抚抗议者。见《Jiangsusheng
  dashiji》，第327页。}

The other slogan adjudged erroneous was one which had been
enthusiastically propagated by Weng Senhe -- ``don't work for the wrong
political line'' (不要为错误路线工作).\footnote{CCP CC, zhongfa (1976),
  No.~24, Zhonggong Nianbao, 1977, 5: 7.
  中共中央，中发（1976）24号，《Zhonggong Nianbao》，1977年，5:7。} The
document urged a greater commitment to production, strict observance of
discipline and a prohibition against ``link-ups'' (串连) and factional
disturbances. It reserved the most severe strictures for local leaders
who had abandoned their responsibilities in the face of difficulty, who
implemented party policy selectively or who, worse still, instigated or
supported factional violence. These pointed remarks were directed not
only at rebel leaders such as Zhang, Weng and He, but also at cadres
such as Tan Qilong, who ultimately had to bear responsibility for the
disorder in Zhejiang. Mao Zedong, who had approved the CC Document of 1
July, warned the radical leaders at the 17 July Politburo meeting not to
form a small faction of four
(你们要注意呢,不要搞成四人小宗派呢).\footnote{Xinhua, February 12,1984
  in ZJRB, February 13,1984, p.~3. See also Gao Wenqian, ZJRB, January
  5,1986. 新华社，1984年2月12日，载
  ZJRB，1984年2月13日，第3页。另见高文谦，ZJRB，1986年1月5日。} On 1
June 1974 Zhou Enlai had entered a hospital in Beijing from which he was
never to reemerge, except for brief public appearances at major events
such as the 4th NPC in January 1975. Although the Chairman appointed
Wang Hongwen to take charge of the daily work of the CC during Zhou's
hospitalization,\footnote{Ye Yonglie, Wang Hongwen xingshuailu, p.~406.
  叶永烈，《Wang Hongwen xingshuailu》，第406页。} Deng Xiaoping had
most probably played a major role in drafting the document.

另一个被判定为错误的口号，是翁森鹤曾热情宣传的``不要为错误路线工作''。\footnote{CCP
  CC, zhongfa (1976), No.~24, Zhonggong Nianbao, 1977, 5: 7.
  中共中央，中发（1976）24号，《Zhonggong Nianbao》，1977年，5:7。}
文件要求更加重视生产，严格遵守纪律，并禁止``串连''和派性扰乱。文件对那些在困难面前放弃职责、有选择地执行党的政策，甚至更严重地煽动或支持派性暴力的地方领导人，作出了最严厉的批评。这些尖锐意见不仅针对张、翁、贺等造反派领导人，也针对谭启龙这样的干部，因为他最终必须对浙江的混乱承担责任。批准了7月1日中央文件的毛泽东，在7月17日的政治局会议上警告激进派领导人不要形成一个四人小宗派（你们要注意呢,不要搞成四人小宗派呢）。\footnote{Xinhua,
  February 12,1984 in ZJRB, February 13,1984, p.~3. See also Gao
  Wenqian, ZJRB, January 5,1986. 新华社，1984年2月12日，载
  ZJRB，1984年2月13日，第3页。另见高文谦，ZJRB，1986年1月5日。}
1974年6月1日，周恩来住进北京一家医院，除1975年1月第四届全国人大等重大活动中短暂公开露面外，再也没有真正重新回到公众政治舞台。虽然主席在周住院期间任命王洪文主持中央日常工作，\footnote{Ye
  Yonglie, Wang Hongwen xingshuailu, p.~406. 叶永烈，《Wang Hongwen
  xingshuailu》，第406页。} 但邓小平很可能在起草该文件时发挥了主要作用。

The clear expression of support for provincial administrations contained
in the central document, accompanied by a blunt order to concentrate on
industrial production, was certainly good news to first secretaries like
Tan Qilong and a stern warning to the rebels to fall into line or face
the consequences. To regain control over the provincial administration
and check the anarchic behavior of the rebels, Tan's first step was to
reassert the authority of the ZPC over the campaign small groups, the
urban militia, the municipal workers' congress and the Mao Zedong
propaganda teams. All four bodies had come under the domination of
Zhang, Weng and He, and had been a constant source of trouble to Tan.

中央文件中对省级行政机构的明确支持，加上集中力量抓工业生产的直截了当命令，对谭启龙这样的第一书记而言当然是好消息，同时也是对造反派的严厉警告：要么服从，要么承担后果。为了重新控制省级行政机关，并遏制造反派的无政府行为，谭的第一步是重新确立浙江省委对运动小组、城市民兵、市工代会和毛泽东思想宣传队的权威。这四个机构都已落入张、翁、贺的支配之下，并一直是谭的麻烦来源。

It seems that initial steps had already been taken to rein in the urban
militia. At the 19 June rally to celebrate the 12th anniversary of the
establishment of the people's militia. Tan Qilong had praised the work
of the militia in Zhejiang's three largest cities -- Hangzhou, Ningbo
and Wenzhou.\footnote{ZPS, June 20, 1974, SWB/FE/4637/BII/8-9. See also
  ZPS, June 19,1974, SWB/FE/4634/BII/13 for a brief report of the forum
  held by the militia on the day preceding the rally.
  ZPS，1974年6月20日，SWB/FE/4637/BII/8-9。另见
  ZPS，1974年6月19日，SWB/FE/4634/BII/13，其中有一篇关于民兵在集会前一天举行座谈会的简短报道。}
He Xianchun, introduced as a member of the standing committee of the
ZPRC, Vice-chairman of the HMRC and secretary of the municipal militia
headquarters' party committee, delivered a report in which he stressed
the subjection of the militia to the party. Provincial reports from Army
Day indicated that with the exception of Shanghai, where it was asserted
that the militia command remained outside PLA control, the PLA regained
its control over the urban militia.\footnote{CNS, 540 (August 14,1974).
  See also provincial reports on local militia conferences and militia
  activities during October-November 1974. CNS, 543 (November 13,1974),
  pp.~1-9.
  CNS，第540期（1974年8月14日）。另见关于1974年10月至11月地方民兵会议和民兵活动的省级报道。CNS，第543期（1974年11月13日），第1-9页。}
By 29 September 1974 when 2,000 militiamen assembled in Hangzhou to
celebrate the 16th anniversary of Mao's call to build up the militia,
greater moderation seemed to prevail.\footnote{ZJRB, September 30,1974,
  p.~1. ZJRB，1974年9月30日，第1页。} Although it was stated that the
urban militia continued to grow in size and strength, the themes of
discipline and unity received some prominence. Apart from singing the
PLA's ``three rules and eight points'', the militia was asked to
strengthen its internal unity and its relations with the city's
citizens. Thus, the local authorities had finally succeeded in
curtailing the militia's activities, for the moment at least. As for the
propaganda teams, which were controlled by Weng Senhe, a report of the
conference held from 20-28 July indicated that moves were afoot to bring
them under closer party leadership.\footnote{ZJRB, July 30,1974, pp.~1,
  4. ZJRB，1974年7月30日，第1、4页。}

看来，收束城市民兵的初步措施已经开始采取。在6月19日庆祝人民民兵成立12周年的大会上，谭启龙赞扬了浙江三大城市杭州、宁波和温州民兵的工作。\footnote{ZPS,
  June 20, 1974, SWB/FE/4637/BII/8-9. See also ZPS, June 19,1974,
  SWB/FE/4634/BII/13 for a brief report of the forum held by the militia
  on the day preceding the rally.
  ZPS，1974年6月20日，SWB/FE/4637/BII/8-9。另见
  ZPS，1974年6月19日，SWB/FE/4634/BII/13，其中有一篇关于民兵在集会前一天举行座谈会的简短报道。}
贺贤春以浙江省革委会常委、杭州市革委会副主任和市民兵指挥部党委书记的身份被介绍，并作了报告，强调民兵必须服从党。建军节期间的省级报道表明，除上海之外，解放军已重新控制城市民兵；而上海据称民兵指挥机构仍在解放军控制之外。\footnote{CNS,
  540 (August 14,1974). See also provincial reports on local militia
  conferences and militia activities during October-November 1974. CNS,
  543 (November 13,1974), pp.~1-9.
  CNS，第540期（1974年8月14日）。另见关于1974年10月至11月地方民兵会议和民兵活动的省级报道。CNS，第543期（1974年11月13日），第1-9页。}
到1974年9月29日，2000名民兵在杭州集会，庆祝毛泽东号召建设民兵16周年时，较为温和的气氛似乎占了上风。\footnote{ZJRB,
  September 30,1974, p.~1. ZJRB，1974年9月30日，第1页。}
虽然报道仍称城市民兵的规模和力量继续增长，但纪律和团结的主题也得到了相当强调。除唱解放军的``三大纪律八项注意''外，民兵还被要求加强内部团结，并改善同城市居民的关系。因此，地方当局终于成功限制了民兵的活动，至少暂时如此。至于由翁森鹤控制的宣传队，7月20日至28日召开的一次会议的报道表明，已在采取行动，使它们置于更密切的党的领导之下。\footnote{ZJRB,
  July 30,1974, pp.~1, 4. ZJRB，1974年7月30日，第1、4页。}

For the initial period after the release of the CC document, however,
the rebel activists in Zhejiang remained defiant and insubordinate. At a
mass meeting held by the ZPC on 5 July the leaders of the ``mountain
top'' faction occupied center stage and seemed to treat the document
with levity.\footnote{ZJRB, July 8,1974, p.~1. This contradicts the
  assertion by CNS, 526, p.~3, and repeated by Goodman, ``The Shanghai
  Connection'', pp.~141-42, that Shanghai was alone in holding such a
  rally at this time. ZJRB，1974年7月8日，第1页。这与 CNS
  第526期第3页的说法相矛盾；Goodman 在 ``The Shanghai Connection''
  第141-142页也重复了这一说法，即当时只有上海举行了这样的集会。} While
acknowledging the necessity to strengthen party leadership over the
campaign, study the central document, devote attention to the economy,
stop ``exchanges of experience'' and obey labor discipline, the
speakers, who included He Xianchun, directed leading cadres to continue
to give the campaign their highest priority. Furthermore, a Zhejiang
Daily editor's note seemed to contradict a crucial argument in the 1
July CC directive when it stated that

然而，在中央文件发布后的最初一段时间里，浙江的造反派积极分子仍然桀骜不驯，不服从约束。在浙江省委7月5日召开的群众大会上，``山上派''领导人占据了中心位置，似乎对该文件轻率以待。\footnote{ZJRB,
  July 8,1974, p.~1. This contradicts the assertion by CNS, 526, p.~3,
  and repeated by Goodman, ``The Shanghai Connection'', pp.~141-42, that
  Shanghai was alone in holding such a rally at this time.
  ZJRB，1974年7月8日，第1页。这与 CNS 第526期第3页的说法相矛盾；Goodman
  在 ``The Shanghai Connection''
  第141-142页也重复了这一说法，即当时只有上海举行了这样的集会。}
包括贺贤春在内的发言者一方面承认有必要加强党对运动的领导、学习中央文件、重视经济、停止``交流经验''并遵守劳动纪律，另一方面又要求领导干部继续把运动放在最高优先位置。此外，《浙江日报》的一则编者按似乎与7月1日中央指示中的一个关键论点相矛盾，其中写道：

\begin{quote}
Practice in the movement to criticize Lin and Confucius over the past
six months shows that where Party organizations have vigorously grasped
criticism of Lin Biao and Confucius and have followed the correct
ideological and political line, industrial production has increased and
the situation in both revolution and production is excellent.\footnote{ZPS,
  July 7, 1974, SWB/FE/BII/8. ZPS，1974年7月7日，SWB/FE/BII/8。}
\end{quote}

\begin{quote}
过去半年来批林批孔运动的实践表明，凡是党组织大力抓批林批孔，并遵循正确的思想政治路线的地方，工业生产就增长，革命和生产两方面的形势都很好。\footnote{ZPS,
  July 7, 1974, SWB/FE/BII/8. ZPS，1974年7月7日，SWB/FE/BII/8。}
\end{quote}

The editor's note was appended to an article reporting a 19\% increase
in industrial production in Jiande county for the first six months of
1974 over the corresponding period for the previous year. The
implications of this assertion were potentially explosive for Tan Qilong
when figures for provincial industrial production were compiled. The
radicals could disclaim any responsibility for directly managing the
economy. Production shortfalls would, on the one hand, upset the central
leadership (with the exception of the radicals) and arouse discontent
among the people of the province, and on the other hand would prove to
the satisfaction of Tan's local opponents and their mentors that he had
failed to carry out the political campaign. Either way Tan was damned.

这则编者按附在一篇文章之后，该文报道称，建德县1974年前六个月的工业生产比上一年同期增长了19\%。一旦汇总全省工业生产数字，这一说法对谭启龙而言便可能具有爆炸性含义。激进派可以否认自己对经济的直接管理负有任何责任。生产未达标一方面会使中央领导层（激进派除外）不满，并激起全省人民的不满；另一方面又会让谭的地方反对者及其后台相信，他未能贯彻政治运动。无论如何，谭都难逃其咎。

A meeting to promote production and revive the campaign, which was
called by the ZPC and ZPRC on 9 July 1974, declared in decisive terms
that ``All fields and organizations should work under the centralized
leadership of Party committees''.\footnote{ZJRB, July 11, 1974, p.~1;
  ZPS, July 11,1974, SWB/FE/4655/BII/23-5. For reports of other
  provincial meetings see CNS, 526 (July 19, 1974), pp.~1-6; 527 (July
  24,1974), pp.~1-7; 528 (July 31, 1974), pp.~1-7.
  ZJRB，1974年7月11日，第1页；ZPS，1974年7月11日，SWB/FE/4655/BII/23-25。关于其他省级会议的报道，见
  CNS，第526期（1974年7月19日），第1-6页；第527期（1974年7月24日），第1-7页；第528期（1974年7月31日），第1-7页。}
Tan Qilong spoke of the desirability for unity to supersede ``bourgeois
sectarianism'' so that everyone could devote their energies to work. It
was announced that agricultural production targets had been met in the
first half of the year, but industrial production figures did not rate a
mention, a clear indication of trouble in this sector.

1974年7月9日，浙江省委和省革委会召开了一次促进生产并重振运动的会议，会议以坚决的措辞宣布：``各条战线、各个组织都要在党委的一元化领导下工作。''\footnote{ZJRB,
  July 11, 1974, p.~1; ZPS, July 11,1974, SWB/FE/4655/BII/23-5. For
  reports of other provincial meetings see CNS, 526 (July 19, 1974),
  pp.~1-6; 527 (July 24,1974), pp.~1-7; 528 (July 31, 1974), pp.~1-7.
  ZJRB，1974年7月11日，第1页；ZPS，1974年7月11日，SWB/FE/4655/BII/23-25。关于其他省级会议的报道，见
  CNS，第526期（1974年7月19日），第1-6页；第527期（1974年7月24日），第1-7页；第528期（1974年7月31日），第1-7页。}
谭启龙谈到，应当以团结取代``资产阶级派性''，使每个人都能把精力投入工作。会议宣布，上半年农业生产指标已经完成，但工业生产数字却没有被提及，这清楚表明该部门存在问题。

Requests by the 1 July Central Committee document as well as the
provincial mass meeting of 5 July to end revolutionary ``exchanges of
experience'' went unheeded at the last major rally of the campaign on 22
July.\footnote{ZJRB, July 24,1974, p.~1. ZJRB，1974年7月24日，第1页。}
An editorial of 28 July 1974\footnote{ZPS, July 28, 1974,
  SWB/FE/4667/BII/33. ZPS，1974年7月28日，SWB/FE/4667/BII/33。} called
for the momentum of the campaign to be maintained during the oncoming
busy farming season, but it was the last shot as the forces of
moderation took over. The changing balance of power at the center and
the publication of articles reflecting this shift\footnote{See CNS, 532
  (August 28,1974), 534 (September 11,1974), pp.~1-4 and 544 (November
  20, 1974), pp.~1-6, for evidence of the propaganda counter-attacks
  launched against the radicals. 关于针对激进派发动宣传反击的证据，见
  CNS，第532期（1974年8月28日）、第534期（1974年9月11日）第1-4页，以及第544期（1974年11月20日）第1-6页。}
assisted Tan regain the initiative in Zhejiang.

7月1日中央文件以及7月5日省群众大会都要求结束革命性的``交流经验''，但这一要求在7月22日运动的最后一次大型集会上并未得到遵守。\footnote{ZJRB,
  July 24,1974, p.~1. ZJRB，1974年7月24日，第1页。}
1974年7月28日的一篇社论\footnote{ZPS, July 28, 1974, SWB/FE/4667/BII/33.
  ZPS，1974年7月28日，SWB/FE/4667/BII/33。}
呼吁在即将到来的农忙季节继续保持运动势头，但随着温和力量接管局面，这已是最后一击。中央权力平衡的变化，以及反映这一变化的文章发表，\footnote{See
  CNS, 532 (August 28,1974), 534 (September 11,1974), pp.~1-4 and 544
  (November 20, 1974), pp.~1-6, for evidence of the propaganda
  counter-attacks launched against the radicals.
  关于针对激进派发动宣传反击的证据，见
  CNS，第532期（1974年8月28日）、第534期（1974年9月11日）第1-4页，以及第544期（1974年11月20日）第1-6页。}
帮助谭在浙江重新掌握主动权。

With the termination of the campaign the attention of the central
leaders was brought back to the state of the economy. The CC document of
1 July had pinpointed the problems in certain key industrial sectors and
since that time provincial leaders like Tan Qilong had struggled to come
to grips with the document's directives. In his speech on National Day
1974 Tan Qilong was looking ahead. New tasks required new slogans and
Tan referred to a quintet which received a great deal of publicity over
the following twelve months. These slogans, which were popularized in
the national press by the moderate leadership, called for emphasis on
the political line, unity, the overall situation, policy and discipline
(讲路线,讲团结,讲大局,讲政策,讲纪律).\footnote{ZJRB, October 1,1974,
  p.~4. ZJRB，1974年10月1日，第4页。} On 4 November 1974 the party
leadership under Tan Qilong convened a meeting to ``grasp revolution and
promote production''.\footnote{ZPS, November 7, 1974,
  SWB/FE/4759/BII/8-9. ZPS，1974年11月7日，SWB/FE/4759/BII/8-9。} The
meeting demanded that cadres change their style of work and go down to
the grassroots to rediscover reality. Senior cadres of the CCP HMC
responded to these exhortations by leading a group of two hundred cadres
to take part in productive labor in the city's factories. They listened
to ``workers' demands and opinions on improving their work''.\footnote{ZPS,
  November 9, 1974, SWB/FE/4761/BII/11.
  ZPS，1974年11月9日，SWB/FE/4761/BII/11。} Not to be outdone, their
provincial civilian and military comrades, led 1,300 cadres to work at
the Hangzhou railway yards on 6 and 7 November.\footnote{ZPS, November
  8, 1974, FBIS/CHI, 218 (1974), G2-3.
  ZPS，1974年11月8日，FBIS/CHI，第218期（1974），G2-3。}

随着运动终止，中央领导人的注意力重新回到经济状况上。7月1日中央文件已经指出若干关键工业部门的问题，此后谭启龙这样的省级领导人一直努力落实文件指示。谭启龙在1974年国庆讲话中开始向前看。新的任务需要新的口号，谭提到了一组五项口号，在随后十二个月里得到大量宣传。这些口号由温和派领导层通过全国报刊加以推广，要求强调路线、团结、大局、政策和纪律（讲路线,讲团结,讲大局,讲政策,讲纪律）。\footnote{ZJRB,
  October 1,1974, p.~4. ZJRB，1974年10月1日，第4页。}
1974年11月4日，谭启龙领导下的党组织召开了一次``抓革命、促生产''会议。\footnote{ZPS,
  November 7, 1974, SWB/FE/4759/BII/8-9.
  ZPS，1974年11月7日，SWB/FE/4759/BII/8-9。}
会议要求干部转变工作作风，深入基层重新认识实际。中共杭州市委的高级干部响应这些号召，带领一支由二百名干部组成的队伍到市内工厂参加生产劳动。他们听取了``工人群众对改进工作的要求和意见''。\footnote{ZPS,
  November 9, 1974, SWB/FE/4761/BII/11.
  ZPS，1974年11月9日，SWB/FE/4761/BII/11。}
不甘落后的是，他们的省级党政军同僚于11月6日和7日率领1300名干部到杭州铁路货场劳动。\footnote{ZPS,
  November 8, 1974, FBIS/CHI, 218 (1974), G2-3.
  ZPS，1974年11月8日，FBIS/CHI，第218期（1974），G2-3。}

Further meetings of 18 and 28 November, held by the municipal and
provincial leaderships respectively,\footnote{ZPS, November 21, 1974,
  SWB/FE/4768/BII/29-30; ZPS, November 30, 1974, SWB/FE/4775/BII/15.
  ZPS，1974年11月21日，SWB/FE/4768/BII/29-30；ZPS，1974年11月30日，SWB/FE/4775/BII/15。}
called for the utmost effort to fulfil annual production tasks. The
reports from these meetings also spoke of a ``new leap'' in the economy
for 1975. The People's Daily editorial of 28 November 1974 sounded the
swansong to the year's political campaign\footnote{``继续搞好批林批孔''
  (Continue the criticism of Lin and Confucius), RMRB, November 28,1974,
  p.~1. 《继续搞好批林批孔》（Continue the criticism of Lin and
  Confucius），RMRB，1974年11月28日，第1页。} and the authorities in
Zhejiang immediately followed suit.\footnote{ZPS, November 30,1974,
  SWB/FE/4775/BII/15. See also CNS, 546 (December 4,1974), pp.~1-5.
  ZPS，1974年11月30日，SWB/FE/4775/BII/15。另见
  CNS，第546期（1974年12月4日），第1-5页。} On 30 November Mao gave his
personal imprimatur to the preparations to revive industrial production
in 1975 by issuing instructions regarding the economy. He thus further
consolidated Zhou Enlai and Deng Xiaoping's position.\footnote{CCP CC,
  zhongfa (1977), No.~37, I\&S, 15: 5 (1979), p.~86. For 1974 national
  production figures, see Liu Suinian and Wu Qungan, ``Wenhua dageming''
  shiqide guomin jingji, p.~80.
  中共中央，中发（1977）37号，I\&S，15:5（1979），第86页。关于1974年全国生产数字，见刘绥年、吴群敢《``Wenhua
  dageming'' shiqide guomin jingji》，第80页。}

市级和省级领导层分别于11月18日和28日召开进一步会议，\footnote{ZPS,
  November 21, 1974, SWB/FE/4768/BII/29-30; ZPS, November 30, 1974,
  SWB/FE/4775/BII/15.
  ZPS，1974年11月21日，SWB/FE/4768/BII/29-30；ZPS，1974年11月30日，SWB/FE/4775/BII/15。}
号召尽最大努力完成全年生产任务。这些会议的报道还谈到1975年经济要有``新跃进''。1974年11月28日《人民日报》社论为当年的政治运动唱响了绝唱，\footnote{``继续搞好批林批孔''
  (Continue the criticism of Lin and Confucius), RMRB, November 28,1974,
  p.~1. 《继续搞好批林批孔》（Continue the criticism of Lin and
  Confucius），RMRB，1974年11月28日，第1页。}
浙江当局随即跟进。\footnote{ZPS, November 30,1974, SWB/FE/4775/BII/15.
  See also CNS, 546 (December 4,1974), pp.~1-5.
  ZPS，1974年11月30日，SWB/FE/4775/BII/15。另见
  CNS，第546期（1974年12月4日），第1-5页。}
11月30日，毛就经济问题发布指示，亲自认可了1975年恢复工业生产的准备工作。由此，他进一步巩固了周恩来和邓小平的地位。\footnote{CCP
  CC, zhongfa (1977), No.~37, I\&S, 15: 5 (1979), p.~86. For 1974
  national production figures, see Liu Suinian and Wu Qungan, ``Wenhua
  dageming'' shiqide guomin jingji, p.~80.
  中共中央，中发（1977）37号，I\&S，15:5（1979），第86页。关于1974年全国生产数字，见刘绥年、吴群敢《``Wenhua
  dageming'' shiqide guomin jingji》，第80页。}

In Zhejiang the campaign had impacted heavily on the economy, the
industrial sector in particular. Not since 1967 and 1968 had output
fallen. Industrial output in 1974 declined 13.6\% over the previous year
and of eighty major commodities in the state plan, only eight were
fulfilled. Steel output dropped 53.5\%, pig iron 59\%, steel products
46.6\%, raw coal 55.5\%, cement 32\%, chemical fertilizer 44\%, cotton
yarn 30\%, cotton cloth 27.5\% and state revenues fell by
24\%.\footnote{See Zhejiang shengqing, p.~1058. 见《Zhejiang
  shengqing》，第1058页。} Political stability was clearly a
prerequisite for an improvement on the economic front in 1975.

在浙江，这场运动对经济造成了沉重冲击，尤其是工业部门。自1967年和1968年以来，产量还未曾下降过。1974年工业产值比上一年下降13.6\%，国家计划中的八十种主要商品，仅有八种完成。钢产量下降53.5\%，生铁下降59\%，钢材下降46.6\%，原煤下降55.5\%，水泥下降32\%，化肥下降44\%，棉纱下降30\%，棉布下降27.5\%，国家财政收入下降24\%。\footnote{See
  Zhejiang shengqing, p.~1058. 见《Zhejiang shengqing》，第1058页。}
政治稳定显然是1975年经济战线改善的先决条件。

\section{The Balance-sheet}\label{the-balance-sheet}

\section{总结评估}\label{ux603bux7ed3ux8bc4ux4f30}

Overall, the three rebel leaders, Zhang Yongsheng, Weng Senhe and He
Xianchun, could feel highly satisfied with their year's work. Although
they had failed to topple the provincial leadership, their tactics of
destabilization and confrontation had rocked it to its foundation.
Working from their organizational base in the HMWC they had mobilized
their supporters in the factories of Hangzhou and drafted them into the
urban militia which they dominated. Simultaneously, they had pressurized
the municipal party authorities, in coordination with allies who held
responsible positions on the CCP HMC, to appoint colleagues to important
posts in Hangzhou. In the tense and volatile atmosphere generated by the
anti-Lin anti-Confucius campaign, and taking advantage of the
opportunities provided by the establishment of the campaign small group,
the rebel leaders struck at the heart of provincial power, rendering it
virtually impotent for several months. Encouraged by their perceptions
of the shifting balance of power in Beijing and protected by national
figures such as Wang Hongwen, they took advantage of their unique
position as both mass organization leaders and party members with
responsible posts in the party and government, to operate both from
without and within the institutional framework.\footnote{Fenwick comes
  to similar conclusions in her analysis of the impact of campaign at
  the local level. See ``The Gang of Four'', pp.~254-55, 280-81. Fenwick
  在分析该运动对地方层面的影响时得出了类似结论。见 ``The Gang of
  Four''，第254-255、280-281页。}

总体而言，张永生、翁森鹤和贺贤春这三位造反派领导人可以对他们这一年的工作感到相当满意。虽然他们未能推翻省级领导层，但他们的不稳定化和对抗策略已经动摇了其根基。他们以杭州市工代会为组织基地，动员杭州各工厂中的支持者，并将其编入由他们支配的城市民兵。与此同时，他们又同在中共杭州市委担任要职的盟友协调，向市级党组织施压，要求把自己的同事任命到杭州的重要职位上。在批林批孔运动造成的紧张而动荡的气氛中，造反派领导人利用运动小组成立所提供的机会，直接打击省级权力核心，使其在数月内几乎失去效能。他们受到自己对北京权力平衡变化的判断所鼓舞，又受到王洪文等全国性人物的庇护，利用其既是群众组织领导人、又是在党政机构中担任负责职务的党员这一独特地位，同时从体制外和体制内展开活动。\footnote{Fenwick
  comes to similar conclusions in her analysis of the impact of campaign
  at the local level. See ``The Gang of Four'', pp.~254-55, 280-81.
  Fenwick 在分析该运动对地方层面的影响时得出了类似结论。见 ``The Gang of
  Four''，第254-255、280-281页。}

Yet, the radicals also faced a dilemma, as expressed most perceptively
by a China News Summary comment in late 1975:

然而，激进派也面临一种困境，正如《中国新闻摘要》在1975年底一则评论中极为敏锐地表达的那样：

\begin{quote}
The representatives of the ``rebels'' who have been ``integrated into
leadership groups'' are now in a dilemma. As a minority and without
experience they are not decisive in policy-making or personnel
appointment. They cannot veto but can only object or delay. When they
mobilize their followers in the ``masses'' to oppose what they have been
unable to veto in the leadership groups they are accused of ``bourgeois
factionalism'' and criticized and defeated, even expelled from office.
On the other hand some of these ``rebels'' have come to enjoy being in
the ``leadership groups'' and have developed an attitude of conciliation
and co-operation towards their veteran colleagues, tolerating as much as
possible and thus losing their ``rebel'' standpoint.
\end{quote}

\begin{quote}
那些已经``结合进领导班子''的``造反派''代表，如今处于两难境地。作为少数派且缺乏经验，他们在政策制定或人事任命上并不起决定作用。他们不能否决，只能反对或拖延。当他们动员自己在``群众''中的追随者，去反对那些在领导班子中未能否决的事情时，他们便被指责为``资产阶级派性''，遭到批判和击败，甚至被免职。另一方面，这些``造反派''中有些人已经开始喜欢身处``领导班子''之中，并对资深同事形成了调和与合作的态度，尽可能容忍，从而丧失了自己的``造反''立场。
\end{quote}

\begin{quote}
There is no way to prevent them from taking this line, which is
``capitulation'' in essence, so long as the strength of the entire
``rebel'' group is weaker than that of the veterans' group. The ``ruling
group in power'' are using ``dual'' tactics, winning over those
``rebels'' who are willing to co-operate by giving them fame and
position and suppressing by force those who are still unwilling. In a
few years' time the leadership level of the ``rebels'', laboriously
cultivated over the past nine years, may well have ``melted away'' and
the entire endeavor of the Cultural Revolution been lost. Perhaps the
most urgent task, in the consideration of those who launched the
Cultural Revolution, is to prevent further co-operation between the
``rebels'' and the veterans.
\end{quote}

\begin{quote}
只要整个``造反派''集团的力量弱于老干部集团，就没有办法阻止他们走这条本质上是``投降''的路线。``掌权的统治集团''正在运用``两手''策略，通过给予名声和职位来争取那些愿意合作的``造反派''，并以武力镇压那些仍不愿合作的人。再过几年，过去九年辛苦培养起来的``造反派''领导层很可能已经``消融''，文化大革命的整个事业也就丧失了。或许，在那些发动文化大革命的人看来，最迫切的任务就是防止``造反派''和老干部之间进一步合作。
\end{quote}

\begin{quote}
Struggle must be launched particularly against those ``rebel'' leaders
who have been won over or who have capitulated, and now collaborate, so
as to split, undermine and even suppress other ``rebels'' who persist in
``rebellion'' against the ``ruling group in power'' and would face
annihilation rather than give up.\footnote{CNS, 584, p.~2.
  CNS，第584期，第2页。}
\end{quote}

\begin{quote}
必须特别对那些已被争取过去或已经投降、如今进行合作的``造反派''领导人展开斗争，因为他们分裂、破坏甚至镇压其他坚持对``掌权的统治集团''进行``造反''、宁可面临毁灭也不放弃的``造反派''。\footnote{CNS,
  584, p.~2. CNS，第584期，第2页。}
\end{quote}

The rebel leaders in Zhejiang had not been co-opted by the provincial
leadership but rather forced on it by Wang Hongwen. Wang was no doubt
hoping that they would eventually be accepted, however unwillingly, by
Tan and his colleagues. While Weng Senhe and his colleagues made it
clear that they were hankering after power, acceptance into the system
would in a sense curtail their freedom. While they remained outsiders to
the CCP ZPC they could continue to pursue their tactics of political
destabilization, which had proved so successful in 1974. Once absorbed
into the system, however, they would be constrained to a certain extent
by organizational discipline and become captives of its routine and
bureaucracy.

浙江的造反派领导人并不是被省级领导层主动吸纳的，而是由王洪文强加给它的。毫无疑问，王希望谭及其同事最终会接受他们，无论这种接受多么勉强。翁森鹤及其同事虽然明确表现出对权力的渴望，但进入体制在某种意义上会限制他们的自由。只要他们仍是中共浙江省委之外的人，就可以继续推行其政治不稳定化策略，而这种策略在1974年已证明非常成功。然而，一旦被吸纳进体制，他们就会在一定程度上受到组织纪律约束，并成为日常程序和官僚机制的俘虏。

When the radical leaders in Beijing lost the initiative in mid-1974,
their underlings in Zhejiang could not sustain their assault and were
forced to retreat. But retreat did not mean submission. The solution to
the problems of rebellion, factionalism and administrative paralysis in
Zhejiang rested with the central authorities, but while they were
divided over policy direction and split politically, no effective
initiative was possible. It was Deng Xiaoping who took up the challenge
of ridding the Chinese political system of what he termed ``bourgeois
factionalism'' and the most notorious exponents of confrontational
politics in the provinces. However, the enormity of the task and the
ambivalence of his position forced Deng to make major concessions to his
opponents, as chapter nine of this study will reveal.

当北京的激进派领导人在1974年中期失去主动权时，他们在浙江的下属无法维持攻势，被迫后退。但后退并不意味着屈服。解决浙江的造反、派性和行政瘫痪问题，要依靠中央当局；然而，只要中央在政策方向上存在分歧并在政治上分裂，就不可能采取有效举措。正是邓小平承担起了清除中国政治体系中他所谓``资产阶级派性''以及各省最臭名昭著的对抗政治代表人物这一挑战。然而，任务之艰巨以及他自身地位的暧昧，迫使邓对其对手作出重大让步，本研究第九章将对此加以说明。

\section[Weng Senhe's Silk Complex: A Model]{\texorpdfstring{Weng
Senhe's Silk Complex: A
Model\footnote{The following section is based on, Hangzhou silk complex
  spare-time culture and propaganda team, ``Settle accounts with Weng
  Senhe's crimes in the arts'', HZRB, November 25, 1977; HZRB, May
  20,1977; HZRB, July 4, 1977; HZRB, March 14, 1978; Shen Chuyun, RMRB,
  May 17,1977, p.~2; Visits to silk complex, January 12,1977, November
  7,1977, November 11,1977.
  以下章节依据：杭州丝绸联合厂业余文艺宣传队《Settle accounts with Weng
  Senhe's crimes in the
  arts》，HZRB，1977年11月25日；HZRB，1977年5月20日；HZRB，1977年7月4日；HZRB，1978年3月14日；沈楚云，RMRB，1977年5月17日，第2页；对丝绸联合厂的访问，1977年1月12日、1977年11月7日、1977年11月11日。}}{Weng Senhe's Silk Complex: A Model}}\label{weng-senhes-silk-complex-a-modelchapter-07-233}

\section[翁森鹤的丝绸联合厂：一个样板]{\texorpdfstring{翁森鹤的丝绸联合厂：一个样板\footnote{The
  following section is based on, Hangzhou silk complex spare-time
  culture and propaganda team, ``Settle accounts with Weng Senhe's
  crimes in the arts'', HZRB, November 25, 1977; HZRB, May 20,1977;
  HZRB, July 4, 1977; HZRB, March 14, 1978; Shen Chuyun, RMRB, May
  17,1977, p.~2; Visits to silk complex, January 12,1977, November
  7,1977, November 11,1977.
  以下章节依据：杭州丝绸联合厂业余文艺宣传队《Settle accounts with Weng
  Senhe's crimes in the
  arts》，HZRB，1977年11月25日；HZRB，1977年5月20日；HZRB，1977年7月4日；HZRB，1978年3月14日；沈楚云，RMRB，1977年5月17日，第2页；对丝绸联合厂的访问，1977年1月12日、1977年11月7日、1977年11月11日。}}{翁森鹤的丝绸联合厂：一个样板}}\label{ux7fc1ux68eeux9e64ux7684ux4e1dux7ef8ux8054ux5408ux5382ux4e00ux4e2aux6837ux677fchapter-07-233}

To conclude this chapter it is instructive to examine briefly the
Hangzhou Silk Printing and Dyeing Complex which Weng Senhe built up as a
model unit in establishing a militia branch and promoting the ``two
crashthroughs''. The controversies which erupted in the factory as a
result of these policies caused massive disruption to production, the
consequences of which reverberated across the city. The tensions and
divisions which were aggravated amongst the mill's work-force eventually
provided the justification for the intervention of the PLA in 1975, if
only to guarantee the workers' safety on the shop-floor.

在本章结尾，简要考察翁森鹤树立为样板单位的杭州丝绸印染联合厂，是很有启发性的；他在那里建立民兵分队，并推动``双突''。这些政策在厂内引发的争议造成了生产的严重混乱，其后果波及全市。工厂劳动力内部被加剧的紧张和分裂，最终为1975年解放军介入提供了理由，哪怕只是为了保障工人在车间里的安全。

Weng used his work unit to set the pace for the ``two crashthroughs''
and the formation of a new militia outfit.\footnote{For articles
  published contemporaneously with the campaign holding up the factory
  as a model unit, see ZJRB, June 10,1974, pp.~1,2; July 1,1974,
  pp.~2,4; and July 14,1974, p.~1. Zhang Yongsheng's Fine Arts Institute
  was also heralded as a model. See ZJRB, August 24,1974, p.~1.
  关于与该运动同时发表、把该厂树为模范单位的文章，见
  ZJRB，1974年6月10日，第1、2页；1974年7月1日，第2、4页；以及1974年7月14日，第1页。张永生的美术学院也被宣传为典型。见
  ZJRB，1974年8月24日，第1页。} From March to July 1974, 97 employees at
the factory were admitted into the CCP and 166 people crash-promoted
into leading positions, including ten party members who joined the
standing committee of the factory's party committee. Huang Yintang, a
close confidant of Weng Senhe and a core member of the ``mountain top''
faction, became first secretary of the party committee shortly after his
admission into the CCP. Huang was also appointed deputy-director of the
municipal Public Security Bureau, deputy party secretary of the HMMCH
and vice-chairman of the detective office (侦破办) which was set up
jointly by the security bureau and the militia command in April 1974.

翁利用自己的工作单位为``双突''和新民兵组织的建立树立步调。\footnote{For
  articles published contemporaneously with the campaign holding up the
  factory as a model unit, see ZJRB, June 10,1974, pp.~1,2; July 1,1974,
  pp.~2,4; and July 14,1974, p.~1. Zhang Yongsheng's Fine Arts Institute
  was also heralded as a model. See ZJRB, August 24,1974, p.~1.
  关于与该运动同时发表、把该厂树为模范单位的文章，见
  ZJRB，1974年6月10日，第1、2页；1974年7月1日，第2、4页；以及1974年7月14日，第1页。张永生的美术学院也被宣传为典型。见
  ZJRB，1974年8月24日，第1页。}
1974年3月至7月，该厂97名职工被吸收入党，166人被突击提拔到领导岗位，其中包括进入厂党委常委会的十名党员。黄阴堂是翁森鹤的亲信，也是``山上派''的核心成员，他在入党后不久便成为党委第一书记。黄还被任命为市公安局副局长、杭州市民兵指挥部党委副书记，以及1974年4月由公安局和民兵指挥部联合设立的侦破办副主任。

As Huang's deputy in the silk mill Weng chose a person described as a
notorious gambler. He had started work at the factory in 1965 and
apparently shared Weng's view that spending one's life performing
repetitive, mundane tasks left a lot to be desired. In the Cultural
Revolution he participated in raids and assaults with his
fellow-workers. In the ``two crashthroughs'' of 1974 he was admitted
into the CCP and, as well as acquiring a party post, became battalion
commander (营长) of the militia contingent. He organized a special
detachment (别动队) which carried out nocturnal raids on the houses of
workers opposed to Weng's faction. In eight months of 1974, this
detachment raided the flats of ten workers.

在丝绸厂里，翁选择了一名被称为臭名昭著的赌徒的人担任黄的副手。此人1965年进厂工作，显然与翁一样认为，一辈子从事重复而平凡的劳动实在乏善可陈。文化大革命中，他同工友一起参加抄家和殴打。在1974年的``双突''中，他被吸收入党，除获得党内职务外，还成为民兵分队的营长。他组织了一支别动队，夜间突袭反对翁派的工人住宅。1974年的八个月内，这支别动队突袭了十名工人的住处。

As late as October 1974 Weng and He Xianchun set up a similar detachment
in the municipal silk bureau and over the next six months it seized and
locked up twenty people and forced them to join a ``militia study
class''.\footnote{HZRB, June 30,1977. HZRB，1977年6月30日。} Weng
personally appointed the party committee first deputy-secretary in the
silk bureau. Weng's appointee took charge of the bureau when the
secretary was shunted aside in the anti-Lin anti-Confucius campaign.
This man's father had reportedly worked for the Guomindang police force
and had been responsible for the murder of CCP members. In 1951 he was
executed. The son concealed his family history and in the Cultural
Revolution joined the rebel forces as having been ``oppressed''
(受压者). When the party branch in his unit was reformed in 1969 he
joined the CCP and became its group leader but the following year was
expelled from the party for concealing his father's record. In 1974, he
joined the Mt. Pingfeng rebels and Weng and He Xianchun despatched him
to the silk bureau as their ``liaison man''. He rejoined the party in
April 1974. Control of the bureau provided Weng with another avenue to
bypass regulations in admitting new party members.\footnote{CCP HMC
  Organization Department and Silk Bureau criticism groups, ``Look at
  the organizational basis of counter-revolutionary bourgeois
  factionalism from a black model'', HZRB, August 12,1977; HZRB, June
  30,1977. 中共杭州市委组织部和丝绸局批判组《Look at the organizational
  basis of counter-revolutionary bourgeois factionalism from a black
  model》，HZRB，1977年8月12日；HZRB，1977年6月30日。}

迟至1974年10月，翁和贺贤春还在市丝绸局设立了一个类似的分队，并在随后六个月中抓捕和关押了二十人，强迫他们参加``民兵学习班''。\footnote{HZRB,
  June 30,1977. HZRB，1977年6月30日。}
翁亲自任命了丝绸局党委第一副书记。在批林批孔运动中，该局书记被排挤到一边后，翁的被任命者接管了该局。据称，此人的父亲曾为国民党警察机关工作，并对杀害中共党员负有责任。1951年，他被处决。其子隐瞒家庭历史，在文化大革命中以``受压者''的身份加入造反派。1969年他所在单位党支部重建时，他加入中共并成为党小组长，但次年因隐瞒其父历史而被开除党籍。1974年，他加入屏风山造反派，翁和贺贤春把他派到丝绸局担任他们的``联络员''。他于1974年4月重新入党。控制该局为翁提供了另一条绕过规定吸收新党员的途径。\footnote{CCP
  HMC Organization Department and Silk Bureau criticism groups, ``Look
  at the organizational basis of counter-revolutionary bourgeois
  factionalism from a black model'', HZRB, August 12,1977; HZRB, June
  30,1977. 中共杭州市委组织部和丝绸局批判组《Look at the organizational
  basis of counter-revolutionary bourgeois factionalism from a black
  model》，HZRB，1977年8月12日；HZRB，1977年6月30日。}

A third member of Weng's team in the mill, who became known as Weng's
right-hand man and capable accomplice in crime (得力帮凶), was appointed
deputy-secretary of the party branch in the weaving workshop after
joining the CCP in 1974. Before the Cultural Revolution he had been
disciplined for a misdemeanor and sentenced to work on a state farm
under supervision. In the turmoil of the Cultural Revolution he had
stolen his personal file to obtain redress for the conviction. Another
worker at the silk complex was appointed to the Gongshu district
``detective'' office. He could not wait for the public notification of
his appointment before rushing off to assume his duties. For good
measure he also appointed himself deputy-chief of the district public
security bureau.\footnote{HZRB, June 23,1977. HZRB，1977年6月23日。}

翁在厂内班子中的第三名成员，后来被称为翁的得力助手和犯罪的得力帮凶，他在1974年入党后被任命为织造车间党支部副书记。文化大革命前，他曾因轻罪受到处分，被判到国营农场监督劳动。在文化大革命的动荡中，他偷走了自己的个人档案，以便为这一判处寻求平反。丝绸联合厂的另一名工人被任命到拱墅区``侦破''办公室。他等不及任命的公开通知，便急忙前去履职。为了更进一步，他还自封为区公安局副局长。\footnote{HZRB,
  June 23,1977. HZRB，1977年6月23日。}

Weng Senhe was acutely aware of the fact that the rebels had suffered
(吃苦头) in the years 1967-68 because they had not succeeded in
obtaining real power. The road to power, he realized, was through the
CCP. Eighty percent of the crash admissions into the party at the mill,
and 88\% of the crash promotions were members of Weng's group. His
opponents, by their own admission, used the same tactics but were less
successful. Weng changed all the heads and their deputies at the factory
to the workshop levels. Only two out of twenty-two section chiefs
retained their posts. Twenty-one of his followers were sent out to staff
important positions in the municipal organization department, the
provincial and municipal public security bureaux and the silk bureau.
Weng boasted that ``the feudal prime ministers surrounded themselves
with high officials; anyone who has followed me must get a post, whether
it is in a factory, the municipality or the province''.\footnote{RMRB,
  December 9,1976, p.~3. RMRB，1976年12月9日，第3页。}

翁森鹤非常清楚，造反派在1967至1968年之所以吃苦头，是因为他们未能取得真正的权力。他意识到，通向权力的道路要经过中共。在该厂突击入党的人员中，80\%是翁集团成员；突击提拔者中，88\%也是翁集团成员。其反对者自己也承认，他们采用了同样策略，但不如翁成功。翁把从厂级到车间级的所有正副负责人都换掉了。二十二名科长中仅有两人保留职务。他的二十一名追随者被派出去，在市组织部门、省市公安局和丝绸局担任重要职位。翁夸口说：``封建宰相身边都是高官；凡是跟过我的人，不管是在厂里、市里还是省里，都必须有个位置。''\footnote{RMRB,
  December 9,1976, p.~3. RMRB，1976年12月9日，第3页。}

Membership of the militia battalion at the silk complex was restricted
to Weng's supporters. Weng reportedly laid down the rule that no third
generation worker, party member, or worker who had a perfect attendance
record over the previous eight years since the beginning of the Cultural
Revolution could join. Whether these restrictions effectively excluded
any applicant is difficult to assess. Three released criminals and ten
hooligans (流氓) joined the battalion, as well as rebels who had
participated in violent clashes in the 1960s. A deputy platoon leader
named Sun Tianmu was a convicted murderer who was working in the factory
under supervision. Weng appreciated the fact that his hatred of the
system would be strong. Another murderer was appointed a squad leader.
Other militiamen admitted to the CCP branch at command HQ had
unfavorable class backgrounds and personal records.\footnote{See Weng
  Senhe bufen zuizheng cailiao, p.~37. 见《Weng Senhe bufen zuizheng
  cailiao》，第37页。} The militia at the mill established its own
disciplinary group (管教组) which in 1974 seized, beat, locked up and
tried in its own kangaroo court (公堂) twenty-three people accused of
slandering Weng. Amongst them was a woman named Shen Qiying. When Weng
Senhe's power was at its height in the first half of 1974 Shen and
twenty-six fellow workers mailed a letter to the CCP CC in Beijing
complaining about his irregular methods of appointment. The letter was
returned to Zhejiang and the twenty-seven signatories seized. Shen was
placed in prison for nineteen days while others spent up to eight months
in confinement.

丝绸联合厂民兵营的成员资格仅限于翁的支持者。据称，翁规定，第三代工人、党员，或自文化大革命开始以来过去八年中出勤记录完美的工人，都不得加入。很难判断这些限制是否实际排除了任何申请者。三名刑满释放人员和十名流氓加入了该营，另有一些曾参加1960年代暴力冲突的造反派也加入其中。一名叫孙天木的副排长是被定罪的杀人犯，当时正在厂内监督劳动。翁看重的是他对体制的仇恨会很强。另一名杀人犯被任命为班长。其他被吸收入指挥部党支部的民兵，也有不利的阶级背景和个人履历。\footnote{See
  Weng Senhe bufen zuizheng cailiao, p.~37. 见《Weng Senhe bufen
  zuizheng cailiao》，第37页。}
厂内民兵设立了自己的管教组，1974年抓捕、殴打、关押并在自己的公堂上审判了二十三名被指控诽谤翁的人。其中有一名名叫沈启英的女性。1974年上半年翁森鹤权势达到顶峰时，沈和二十六名同事给北京的中共中央寄信，控诉他不合规的任命方式。这封信被退回浙江，二十七名签名者随即被抓。沈被关押十九天，其他人则被禁闭最长达八个月。

One of the principal reasons for the level of discontent in factories
such as the silk complex related to the privileged treatment which was
extended to members of the militia. For example, one-fifth of the
militiamen at the silk complex were absent from the production line for
extended periods, living downtown in the HMMCH building where they
performed guard duty or ran errands for the militia command's leaders.
The remaining militiamen were often placed on standby in preparation for
an assignment. During the prolonged periods of absence from the mill
they not only received full pay but extra allowances. One day's work for
the militia was rewarded with three days' rest. A stint of guard duty
resulted in overtime pay, while active service was rewarded with free
evening meals in the factory canteen. Weng further increased the number
of non-productive workers by establishing literary propaganda teams and
other cultural groups. Workers who remained at their posts were
ridiculed as ``little lambs'',''Confucianists'' and ``restorationists''
and compared unfavorably to the ``tigers'', ``Legalists'' and ``fighters
against the tide'' who formed the backbone of the militia.

丝绸联合厂等工厂中不满情绪高涨的一个主要原因，与民兵成员享有的特权待遇有关。例如，丝绸联合厂五分之一的民兵长期脱离生产线，住在市中心的杭州市民兵指挥部大楼里，在那里站岗或为民兵指挥部领导跑腿。其余民兵也常常处于待命状态，准备执行任务。在长期离厂期间，他们不仅领取全额工资，还得到额外津贴。为民兵工作一天，可换来三天休息。一次警卫任务可获得加班费，而执行勤务则可在工厂食堂免费吃晚饭。翁还通过建立文艺宣传队和其他文化团体，进一步增加非生产人员数量。留在岗位上的工人被嘲笑为``小绵羊''、``儒家派''和``复辟派''，并被拿来同构成民兵骨干的``猛虎''、``法家''和``反潮流战士''作不利比较。

It was the resentment and antagonism engendered by this discrimination
that contributed more to the slump in industrial production than wage
demands or gripes over poor work facilities. While a sizeable percentage
of the work force was absent for considerable lengths of time and other
workers were threatened or harassed on the job, it is little wonder that
production plummeted, factional fights erupted and personal insecurity
increased. The provision of welfare and other essential services would
also have been affected. The wages of workers who were not members of
the militia remained unchanged and thus at a comparative disadvantage to
the wages of their younger, more militant work-mates.

正是这种歧视所造成的怨恨和对立，比工资要求或对工作设施简陋的抱怨更大程度地导致了工业生产下滑。当相当比例的劳动力长期离岗，而其他工人在工作中受到威胁或骚扰时，生产骤降、派性斗争爆发以及个人安全感下降，也就不足为奇了。福利和其他基本服务的提供也会受到影响。非民兵成员工人的工资保持不变，因此相较于那些更年轻、更激进的工友的收入处于劣势。
\#\# Notes \#\# 注释

\bookmarksetup{startatroot}

\chapter{CHAPTER EIGHT - THE BREAKDOWN IN LOCAL AUTHORITY, 1974-75 /
第八章 -
地方权威的瓦解，1974-75}\label{chapter-eight---the-breakdown-in-local-authority-1974-75-ux7b2cux516bux7ae0---ux5730ux65b9ux6743ux5a01ux7684ux74e6ux89e31974-75}

\section{The Struggle at the Center, October 1974-January
1975}\label{the-struggle-at-the-center-october-1974-january-1975}

\section{中央的斗争，1974年10月-1975年1月}\label{ux4e2dux592eux7684ux6597ux4e891974ux5e7410ux6708-1975ux5e741ux6708}

By the latter half of 1974 the anti-Lin anti-Confucius campaign had
largely been brought under control. In an informal talk to overseas
Chinese on the 25th anniversary of the People's Republic Deng Xiaoping,
who had assumed much of the administrative responsibilities from the
ailing Zhou Enlai, played down the significance of the
campaign.\footnote{Jiushi Niandai, December 1974, FBIS/CHI, 238, p.~E5.
  《九十年代》，1974年12月，FBIS/CHI，238，第E5页。} Later in the month
as preparations got underway to convene the long-awaited 4th NPC, the
factional maneuvering between the Cultural Revolution radicals and the
Zhou/Deng led group of veteran administrators reached a climax, with
both sides striving to gain Mao's support. The issue which sparked off
the crisis in October 1974 related to the purchase of ships from
overseas to carry China's international trade. On 18 June Wang Hongwen
had accused the advocates of the proposal as being ``false foreign
devils'' (假洋鬼子) who ``worshipped the foreign bourgeoisie'' and
described the policy as pursuing the ``revisionist line''.\footnote{Liu
  Suinian and Wu Qungan, ``Wenhua dageming'' shiqide guomin jingji,
  p.~167. 刘岁年、吴群敢：《``文化大革命''时期的国民经济》，第167页。}

到1974年下半年，批林批孔运动大体上已经受到控制。在中华人民共和国成立25周年之际同海外华人的一次非正式谈话中，已从病重的周恩来手中承担了许多行政职责的邓小平淡化了这场运动的重要性。\footnote{Jiushi
  Niandai, December 1974, FBIS/CHI, 238, p.~E5.
  《九十年代》，1974年12月，FBIS/CHI，238，第E5页。}
当月稍晚，期待已久的第四届全国人民代表大会的筹备工作展开，文化大革命激进派与周恩来、邓小平领导的老资格行政干部集团之间的派系角力达到高潮，双方都力图获得毛的支持。1974年10月引发危机的问题，涉及从海外购买船舶以承运中国的国际贸易。6月18日，王洪文曾指责这一提议的支持者是``假洋鬼子''，说他们``崇拜外国资产阶级''，并把这项政策描述为执行``修正主义路线''。\footnote{Liu
  Suinian and Wu Qungan, ``Wenhua dageming'' shiqide guomin jingji,
  p.~167. 刘岁年、吴群敢：《``文化大革命''时期的国民经济》，第167页。}

This issue smouldered until October 1974 when, at a Politburo meeting on
the 17th of that month, Jiang Qing insisted that Deng Xiaoping declare
his position on the issue.\footnote{This account is based on the
  following references: Zhao Zhichao, ``毛泽东在长沙的114天'' (Mao
  Zedong's 114 Days in Changsha), 记者文学 (Journalist Literature) No.~2
  (1989), pp.~4-24; Lin Qingshan, 康生外传 (An Unofficial Biography of
  Kang Sheng) (Hong Kong: Xingchen Chubanshe, 1987), pp.~374-75; Hao
  Mengbi and Duan Haoran, Liushinian, pp.~637-39; Gao Gao and Yan Jiaqi,
  ``Wenhua dageming'', pp.~530-5; Gao Wenqian, ZJRB, January 5, 1986;
  ``Wenhua dageming'' shiqide guomin jingji, p.~168; Ji Xizhen, Lin Gang
  and Lu Nan, ``长沙诬告事件'' (Events surrounding the trumped-up
  charges in Changsha), in Lishi zai zheli chensi, Vol. 2, pp.~196-203;
  Gao Wenqian, RMRB, January 5, 1986; Ye Yonglie, Wang Hongwen
  xingshuailu, pp.~402-03, 407-17.
  本段叙述依据以下文献：赵志超：《毛泽东在长沙的114天》，《记者文学》1989年第2期，第4-24页；林青山：《康生外传》（香港：星辰出版社，1987年），第374-375页；郝梦笔、段浩然：《六十年》，第637-639页；高皋、严家其：《文化大革命》，第530-535页；高文谦，《浙江日报》，1986年1月5日；《``文化大革命''时期的国民经济》，第168页；纪希晨、林岗、卢南：《长沙诬告事件》，载《历史在这里沉思》第2卷，第196-203页；高文谦，《人民日报》，1986年1月5日；叶永烈：《王洪文兴衰录》，第402-403、407-417页。}
Deng reportedly answered that he would investigate the matter. The
Cultural Revolution radicals were eager to place the wily veteran in an
awkward position and compromise him on an issue that they thought might
arouse Mao's sensibilities. Their timing was critical. On 4 October Mao
had proposed in a telephone conversation with Wang Hongwen that Deng
become first Vice-premier of the State Council at the forthcoming NPC, a
step that would ensure his succession as Premier. A CCP CC notice of 11
October quoted Mao to the effect that the Cultural Revolution had
already been going for eight years and that it was time for unity and
stability within the party and army. Meanwhile, articles appeared in the
press praising the Empress Lu and legalist politicians, a not very
subtle allusion to Jiang Qing's aspirations for future leadership.

这个问题一直暗中发酵，直到1974年10月17日的一次政治局会议上，江青坚持要求邓小平就此表明立场。\footnote{This
  account is based on the following references: Zhao Zhichao,
  ``毛泽东在长沙的114天'' (Mao Zedong's 114 Days in Changsha), 记者文学
  (Journalist Literature) No.~2 (1989), pp.~4-24; Lin Qingshan, 康生外传
  (An Unofficial Biography of Kang Sheng) (Hong Kong: Xingchen
  Chubanshe, 1987), pp.~374-75; Hao Mengbi and Duan Haoran, Liushinian,
  pp.~637-39; Gao Gao and Yan Jiaqi, ``Wenhua dageming'', pp.~530-5; Gao
  Wenqian, ZJRB, January 5, 1986; ``Wenhua dageming'' shiqide guomin
  jingji, p.~168; Ji Xizhen, Lin Gang and Lu Nan, ``长沙诬告事件''
  (Events surrounding the trumped-up charges in Changsha), in Lishi zai
  zheli chensi, Vol. 2, pp.~196-203; Gao Wenqian, RMRB, January 5, 1986;
  Ye Yonglie, Wang Hongwen xingshuailu, pp.~402-03, 407-17.
  本段叙述依据以下文献：赵志超：《毛泽东在长沙的114天》，《记者文学》1989年第2期，第4-24页；林青山：《康生外传》（香港：星辰出版社，1987年），第374-375页；郝梦笔、段浩然：《六十年》，第637-639页；高皋、严家其：《文化大革命》，第530-535页；高文谦，《浙江日报》，1986年1月5日；《``文化大革命''时期的国民经济》，第168页；纪希晨、林岗、卢南：《长沙诬告事件》，载《历史在这里沉思》第2卷，第196-203页；高文谦，《人民日报》，1986年1月5日；叶永烈：《王洪文兴衰录》，第402-403、407-417页。}
据说邓回答称他会调查此事。文化大革命激进派急于把这位老练的老干部置于尴尬境地，并在他们认为可能触动毛敏感神经的问题上使其受损。他们选择的时机至关重要。10月4日，毛在同王洪文的电话谈话中曾提议，让邓在即将召开的全国人大上出任国务院第一副总理，这一步将确保他接任总理。10月11日的一份中共中央通知引用毛的意思说，文化大革命已经进行了八年，现在是党内和军内实现团结稳定的时候了。与此同时，报刊上出现了赞扬吕后和法家政治家的文章，这对江青未来领导地位的企图作出了并不隐晦的影射。

The tensions and factional rivalries which climaxed in late 1974 and
early 1975 were exacerbated by an important factor -- Mao's prolonged
absence from Beijing. On 13 October the Chairman arrived in Changsha and
remained there until early February, before stopping in Nanchang and
Hangzhou on his way back to the capital. Despite his poor health Mao
continued to play the role of final arbiter in all major decisions, and
none of his senior subordinates dared challenge this practice, even in
his absence. Therefore, Deng Xiaoping, Li Xiannian, Wang Hongwen and
even the hospitalized Zhou Enlai were forced to make the journey to the
south to report to Mao, solicit his views and instructions and report
them back to the Politburo. Deng and Li escorted foreign leaders to
Changsha to receive an audience with the Chairman, who continued to take
a keen interest in foreign affairs.

1974年末和1975年初达到高潮的紧张局势与派系竞争，又因一个重要因素而加剧，即毛长期不在北京。10月13日，主席抵达长沙，并一直停留到翌年2月初，随后在返回首都途中又经停南昌和杭州。尽管健康状况不佳，毛仍继续在所有重大决策中扮演最终仲裁者的角色，他的高级下属即使在他不在场时也无人敢挑战这一惯例。因此，邓小平、李先念、王洪文，甚至住院中的周恩来，都被迫南下向毛汇报，征求他的意见和指示，再把这些意见和指示报告给政治局。邓和李还陪同外国领导人到长沙接受主席接见，而毛仍然对外交事务保持浓厚兴趣。

In the wake of the uproar at the Politburo session on 17 October, which
resulted in the meeting's inconclusive adjournment, the Shanghai
radicals went into emergency session to decide on their next step. The
following day Wang Hongwen was despatched to Changsha where, in a
meeting with Mao, he accused Zhou Enlai and other leaders including Deng
Xiaoping, of creating a political atmosphere similar to that which had
prevailed at the 1970 Lushan plenum. Mao rebuffed Wang. The Chairman's
two foreign affairs interpreters cum advisers and intermediaries
(联络员), his niece Wang Hairong and Tang Wensheng (Nancy Tang), then
became involved in a political tug of war between Zhou Enlai and Jiang
Qing, in which both leaders asked them to act as their advocates with
Mao.

10月17日政治局会议发生喧闹并无果休会之后，上海激进派召开紧急会议，决定下一步行动。次日，王洪文被派往长沙，在同毛会见时，他指责周恩来以及包括邓小平在内的其他领导人制造了一种类似于1970年庐山全会时的政治气氛。毛驳回了王的说法。主席的两名外交翻译兼顾问和中间人（联络员），即他的侄女王海容和唐闻生（Nancy
Tang），随后被卷入周恩来与江青之间的政治拉锯，二人都要求她们在毛面前为自己代言。

Mao accepted Zhou Enlai's version of events by confirming his position
as Premier and instructed him, together with Wang Hongwen, to take
charge of all arrangements for the 4th NPC, including the selection of
personnel. Mao uttered critical remarks about Jiang Qing, both in talks
with Wang Hongwen and his two liaison personnel, as well as in letters
to his wife written in November 1974. Factional hostility and personal
animosities had reached such a state that between 21 October and early
November, Zhou Enlai was forced to divide the Politburo into three
groups and hold separate meetings with them in hospital. On 6 November
he reported the result of these deliberations in a letter to Mao.

毛接受了周恩来对事件的说法，确认其总理地位，并指示他同王洪文一起负责第四届全国人大的一切安排，包括人事遴选。无论是在同王洪文及其两名联络人员的谈话中，还是在1974年11月写给妻子的信中，毛都对江青发表了批评性意见。派系敌意和个人嫌隙已经达到这种程度：从10月21日到11月初，周恩来被迫把政治局分成三个小组，并在医院中分别同他们开会。11月6日，他在致毛的信中报告了这些讨论的结果。

In the final composition of the State Council ministry which would be
ratified at the NPC, Zhou and his supporters gave ground on the
appointments of the Ministers of Culture and Sports in order to ensure
the selection of their candidate, Zhou Rongxin as Minister of Education,
ahead of the radicals' favored choice Chi Qun. Between 23 and 27
December 1974 Zhou and Wang Hongwen held final talks with Mao in
Changsha. Mao praised Deng as being a rare talent who was both
politically and ideologically strong (``人才难得,政治思想强''), and
criticized Jiang Qing for her ambition and factional activities. The
intrigues and back-stage machinations reached their culmination in
January 1975. At the 2nd plenum of the 10th CCP CC Deng was appointed
party Vice-chairman and member of the Politburo standing committee and
confirmed in the posts of CC MAC Vice-chairman and PLA Chief of Staff.
In the political balancing-act, Zhang Chunqiao was appointed Director of
the PLA General Political Department. The 4th NPC, held between 13 and
18 January installed Deng as senior Vice-premier, with Zhang Chunqiao's
name appearing immediately next in the ranking of vice-premiers. The
factional power balance may have tilted against the radicals, but not to
an irretrievable extent.

在即将由全国人大批准的国务院各部委最终组成中，周及其支持者在文化部长和体育部长任命上作出让步，以确保他们的人选周荣鑫出任教育部长，压过激进派偏爱的迟群。1974年12月23日至27日，周和王洪文在长沙同毛举行最后谈话。毛称赞邓是政治和思想上都强的难得人才（``人才难得，政治思想强''），并批评江青的野心和派系活动。阴谋和幕后运作在1975年1月达到顶点。在中共十届二中全会上，邓被任命为党的副主席和政治局常委，并被确认担任中共中央军委副主席和中国人民解放军总参谋长。作为政治平衡的一部分，张春桥被任命为中国人民解放军总政治部主任。1月13日至18日召开的第四届全国人大任命邓为排名第一的副总理，而张春桥的名字在副总理排序中紧随其后。派系权力平衡也许已经向不利于激进派的方向倾斜，但还没有到不可挽回的程度。

In his report to the people's congress, Zhou Enlai announced the
ambitious but vague plan for the accomplishment of the four
modernizations (of agriculture, industry, defence and science and
technology) by the end of the century. He also let it be known within
party circles that his senior deputy, Deng Xiaoping, would be in charge
of the work of the State Council.\footnote{Liu Suinian and Wu Qungan,
  ``Wenhua dageming'' shiqide guomin jingji, p.~169.
  刘岁年、吴群敢：《``文化大革命''时期的国民经济》，第169页。} It was
Deng who shortly afterward set about putting flesh on the bare bones of
Zhou's four modernizations by convening a series of conferences in the
fields of industry, military affairs, planning, science and technology
and agriculture. He set up and supervised a think-tank, comprising
theoreticians and professional experts, to prepare three major documents
which set out his political and economic policy positions.

周恩来在向人民代表大会所作的报告中，宣布了一项雄心勃勃但内容含糊的计划，即到本世纪末实现四个现代化（农业、工业、国防和科学技术现代化）。他还在党内透露，他的高级副手邓小平将负责国务院工作。\footnote{Liu
  Suinian and Wu Qungan, ``Wenhua dageming'' shiqide guomin jingji,
  p.~169. 刘岁年、吴群敢：《``文化大革命''时期的国民经济》，第169页。}
不久之后，正是邓通过在工业、军事、计划、科学技术和农业等领域召开一系列会议，为周的四个现代化的骨架填充具体内容。他还建立并监督了一个由理论工作者和专业专家组成的智囊班子，准备三份重要文件，阐明他的政治和经济政策立场。

Meanwhile, by his absence from the capital while the important
deliberations of January were taking place, Mao had signalled that he
did not fully endorse the proceedings, although he repeated his
directive concerning stability and unity before the close of the CC
plenum.\footnote{Ibid. Mao's slogan ``stability and unity'' (安定团结)
  was not used publicly until about 1 March. See CNS, 588 (October
  22,1975).
  同上。毛的``安定团结''口号直到约3月1日才公开使用。见CNS，588（1975年10月22日）。}
In yet another lurch to the left, Mao launched another campaign, this
time in the realm of theory, while the provinces were still discussing
the documents of the 4th NPC. The study campaign began with the 9
February 1975 editorial of the People's Daily discussing the
significance and content of the post-revolutionary period known as the
dictatorship of the proletariat and China's similarities to, and
differences from, a capitalist economic system. The editorial quoted
excerpts of statements made by Mao the previous year. One extract
derived from remarks made by the Chairman on 20 October 1974 in his
reception of the Danish Prime Minister Hartling. The remaining sentences
dated from the meeting which Mao had held with Zhou Enlai and Wang
Hongwen on 26 December 1974.\footnote{See Hao Mengbi and Duan Haoran,
  Liushinian, p.~644. 见郝梦笔、段浩然：《六十年》，第644页。}

与此同时，在1月重要商议进行期间，毛不在首都，这表明他并未完全认可会议进程，尽管他在中央全会闭幕前重申了关于稳定团结的指示。\footnote{Ibid.
  Mao's slogan ``stability and unity'' (安定团结) was not used publicly
  until about 1 March. See CNS, 588 (October 22,1975).
  同上。毛的``安定团结''口号直到约3月1日才公开使用。见CNS，588（1975年10月22日）。}
在又一次向左转折中，当各省仍在讨论第四届全国人大文件时，毛发动了另一场运动，这一次是在理论领域。学习运动始于1975年2月9日《人民日报》社论，该文讨论了被称为无产阶级专政的革命后时期的意义和内容，以及中国同资本主义经济制度的相似之处与不同之处。社论引用了毛前一年讲话的若干摘录。其中一段出自主席1974年10月20日接见丹麦首相哈特林时的谈话。其余句子则来自毛1974年12月26日同周恩来、王洪文举行的会见。\footnote{See
  Hao Mengbi and Duan Haoran, Liushinian, p.~644.
  见郝梦笔、段浩然：《六十年》，第644页。}

The Chairman's remarks on the nature of socialism and the state of
China's socio-political system were circulated and propagated alongside
his comments on the need for ``stability and unity'' and the importance
of developing the economy. The uncertainties of the distant communist
utopia were thus counterposed to the more immediate and tangible
prospects of political harmony and an improvement in living standards.
Throughout 1975, alongside articles exploring ways to restrict bourgeois
rights (资产阶级法权) and press on with the Maoist ``continuous
revolution'', the Chinese media carried reports extolling production
increases and industrial discipline and promoting economic policies
designed to foster efficiency and profit. The propaganda publicizing the
ideals of the four modernizations was of necessity dressed up in a
Maoist garb, but nevertheless it was clear that two sets of voices were
speaking in Beijing. The debate placed added pressure on provincial
politicians who had to decide which line to emphasize in local
propaganda. The contradictions between Maoist rhetoric and policy
direction were blatant and difficult to accommodate.

主席关于社会主义性质和中国社会政治制度状态的言论，同他关于需要``稳定团结''和发展经济重要性的讲话一起被传达和宣传。这样，遥远共产主义乌托邦的不确定性，便同政治和谐与生活水平改善这些更直接、更切实的前景相对置。整个1975年，中国媒体一方面刊登探讨限制资产阶级法权、继续推进毛主义``继续革命''的文章，另一方面又报道赞扬生产增长和工业纪律，宣传旨在提高效率和利润的经济政策。宣传四个现代化理想的材料必然要披上毛主义外衣，但北京显然同时传出了两套声音。这场争论给省级政治人物增加了压力，他们必须决定在地方宣传中突出哪一条路线。毛主义修辞与政策方向之间的矛盾十分明显，也难以调和。

\section{The Campaign to Study the Dictatorship of the Proletariat,
1975}\label{the-campaign-to-study-the-dictatorship-of-the-proletariat-1975}

\section{学习无产阶级专政的运动，1975年}\label{ux5b66ux4e60ux65e0ux4ea7ux9636ux7ea7ux4e13ux653fux7684ux8fd0ux52a81975ux5e74}

To understand the purpose of the campaign to study the dictatorship of
the proletariat and its focus on ``bourgeois rights'' in particular, it
is necessary to examine briefly Karl Marx's principal writing on the
subject.\footnote{See Keith Forster, ``Chinese Politics in 1975: The
  Mass Campaign to Study the Theory of the Dictatorship of the
  Proletariat'' (unpublished seminar paper. Politics Department,
  University of Adelaide, November 1976). 见Keith Forster：``Chinese
  Politics in 1975: The Mass Campaign to Study the Theory of the
  Dictatorship of the
  Proletariat''（未发表研讨会论文，阿德莱德大学政治系，1976年11月）。}
In the Critique of the Gotha Programme, written in 1875, Marx had
defined bourgeois rights as the form of distribution under the ``first
phase'' of communist society. Exploitation no longer existed in such a
society, yet equal right to the means of consumption was premised on
unequal attributes and different family circumstances, resulting in de
facto inequality.\footnote{Karl Marx, The Critique of the Gotha
  Programme (Beijing: Foreign Languages Press, 1972), pp.~14-17. For a
  critique of Mao's understanding of bourgeois rights (now translated as
  资产阶级权力) see Dangshi tongxun, 8 (April 30,1983), pp.~2-7. For a
  discussion of the Chinese debate of 1958 on bourgeois rights, in which
  Zhang Chunqiao contributed an article to RMRB, see CNS, 557 (March
  5,1975) and Ye Yonglie, 张春桥浮沉史 (The Rise and Fall of Zhang
  Chunqiao), (Changchun: Shidai Wenyi Chubanshe, 1988), pp.~101-09. For
  the most comprehensive treatment of the dictatorship of the
  proletariat in the writings of Marx and Engels, see Hal Draper, Karl
  Marx's Theory of Revolution, Volume III: The ``Dictatorship of the
  Proletariat'' (New York: Monthly Review Press, 1986).
  卡尔·马克思：《哥达纲领批判》（北京：外文出版社，1972年），第14-17页。关于对毛理解资产阶级法权（现在译为资产阶级权力）的批评，见《党史通讯》第8期（1983年4月30日），第2-7页。关于1958年中国围绕资产阶级法权的争论，其中张春桥曾向《人民日报》投稿，见CNS，557（1975年3月5日）以及叶永烈：《张春桥浮沉史》（长春：时代文艺出版社，1988年），第101-109页。关于马克思和恩格斯著作中无产阶级专政问题最全面的论述，见Hal
  Draper, Karl Marx's Theory of Revolution, Volume III: The
  ``Dictatorship of the Proletariat''（纽约：Monthly Review
  Press，1986年）。} In directives released in February 1975, Mao
commented that the eight-grade wage scale practised in Chinese industry
was an example of bourgeois rights which should be restricted. But he
remained typically silent on the policy consequences of his statement,
merely noting that ``we should do more reading of Marxist-Leninist
works''.\footnote{See Marx, Engels and Lenin on the Dictatorship of the
  Proletariat (Beijing: Foreign Languages Press, 1975), p.~2, translated
  from RMRB, 22 February 1975. For detailed accounts of the campaign,
  see Merle Goldman, China's Intellectuals: Advise and Dissent
  (Cambridge, Mass.: Harvard University Press, 1981) pp.~191-201, and
  Fenwick, ``The Gang of Four'', pp.~334-55. For a personal account of
  the campaign by a worker at a Wuhan factory, see B. Michael Frolic,
  Mao's People: Sixteen Portraits of Life in Revolutionary China
  (Cambridge, Mass.: Harvard University Press, 1980), pp.~250-56.
  见Marx, Engels and Lenin on the Dictatorship of the
  Proletariat（北京：外文出版社，1975年），第2页，译自1975年2月22日《人民日报》。关于这场运动的详细叙述，见Merle
  Goldman, China's Intellectuals: Advise and
  Dissent（马萨诸塞州剑桥：哈佛大学出版社，1981年），第191-201页；以及Fenwick：``The
  Gang of
  Four''，第334-355页。关于武汉一家工厂工人对这场运动的个人叙述，见B.
  Michael Frolic, Mao's People: Sixteen Portraits of Life in
  Revolutionary
  China（马萨诸塞州剑桥：哈佛大学出版社，1980年），第250-256页。} The
two major articles published in the campaign by Yao Wenyuan and Zhang
Chunqiao, while issuing warnings against the use of material incentives,
offered no specific suggestions for radical policy changes.\footnote{Yao
  Wenyuan, On the Social Basis of the Lin Biao Anti-Party Clique
  (Beijing: Foreign Languages Press, 1975); Zhang Chunqiao, On
  Exercising All-Round Dictatorship over the Bourgeoisie (Beijing:
  Foreign Languages Press, 1975). 姚文元：On the Social Basis of the Lin
  Biao Anti-Party Clique（北京：外文出版社，1975年）；张春桥：On
  Exercising All-Round Dictatorship over the
  Bourgeoisie（北京：外文出版社，1975年）。}

要理解学习无产阶级专政运动的目的，尤其是它为何聚焦于``资产阶级法权''，有必要简要考察卡尔·马克思关于这一主题的主要著作。\footnote{See
  Keith Forster, ``Chinese Politics in 1975: The Mass Campaign to Study
  the Theory of the Dictatorship of the Proletariat'' (unpublished
  seminar paper. Politics Department, University of Adelaide, November
  1976). 见Keith Forster：``Chinese Politics in 1975: The Mass Campaign
  to Study the Theory of the Dictatorship of the
  Proletariat''（未发表研讨会论文，阿德莱德大学政治系，1976年11月）。}
在1875年写成的《哥达纲领批判》中，马克思把资产阶级法权定义为共产主义社会``第一阶段''的分配形式。在这样的社会中剥削已不复存在，但对消费资料的平等权利仍建立在个人禀赋和家庭情况不平等的基础上，从而造成事实上的不平等。\footnote{Karl
  Marx, The Critique of the Gotha Programme (Beijing: Foreign Languages
  Press, 1972), pp.~14-17. For a critique of Mao's understanding of
  bourgeois rights (now translated as 资产阶级权力) see Dangshi tongxun,
  8 (April 30,1983), pp.~2-7. For a discussion of the Chinese debate of
  1958 on bourgeois rights, in which Zhang Chunqiao contributed an
  article to RMRB, see CNS, 557 (March 5,1975) and Ye Yonglie,
  张春桥浮沉史 (The Rise and Fall of Zhang Chunqiao), (Changchun: Shidai
  Wenyi Chubanshe, 1988), pp.~101-09. For the most comprehensive
  treatment of the dictatorship of the proletariat in the writings of
  Marx and Engels, see Hal Draper, Karl Marx's Theory of Revolution,
  Volume III: The ``Dictatorship of the Proletariat'' (New York: Monthly
  Review Press, 1986).
  卡尔·马克思：《哥达纲领批判》（北京：外文出版社，1972年），第14-17页。关于对毛理解资产阶级法权（现在译为资产阶级权力）的批评，见《党史通讯》第8期（1983年4月30日），第2-7页。关于1958年中国围绕资产阶级法权的争论，其中张春桥曾向《人民日报》投稿，见CNS，557（1975年3月5日）以及叶永烈：《张春桥浮沉史》（长春：时代文艺出版社，1988年），第101-109页。关于马克思和恩格斯著作中无产阶级专政问题最全面的论述，见Hal
  Draper, Karl Marx's Theory of Revolution, Volume III: The
  ``Dictatorship of the Proletariat''（纽约：Monthly Review
  Press，1986年）。}
在1975年2月发布的指示中，毛评论说，中国工业实行的八级工资制就是应当加以限制的资产阶级法权的一个例子。但他照例没有说明这一说法的政策后果，只是指出``我们要多读点马列的书''。\footnote{See
  Marx, Engels and Lenin on the Dictatorship of the Proletariat
  (Beijing: Foreign Languages Press, 1975), p.~2, translated from RMRB,
  22 February 1975. For detailed accounts of the campaign, see Merle
  Goldman, China's Intellectuals: Advise and Dissent (Cambridge, Mass.:
  Harvard University Press, 1981) pp.~191-201, and Fenwick, ``The Gang
  of Four'', pp.~334-55. For a personal account of the campaign by a
  worker at a Wuhan factory, see B. Michael Frolic, Mao's People:
  Sixteen Portraits of Life in Revolutionary China (Cambridge, Mass.:
  Harvard University Press, 1980), pp.~250-56. 见Marx, Engels and Lenin
  on the Dictatorship of the
  Proletariat（北京：外文出版社，1975年），第2页，译自1975年2月22日《人民日报》。关于这场运动的详细叙述，见Merle
  Goldman, China's Intellectuals: Advise and
  Dissent（马萨诸塞州剑桥：哈佛大学出版社，1981年），第191-201页；以及Fenwick：``The
  Gang of
  Four''，第334-355页。关于武汉一家工厂工人对这场运动的个人叙述，见B.
  Michael Frolic, Mao's People: Sixteen Portraits of Life in
  Revolutionary
  China（马萨诸塞州剑桥：哈佛大学出版社，1980年），第250-256页。}
在这场运动中由姚文元和张春桥发表的两篇重要文章，虽然警告不要使用物质刺激，却没有提出激进政策变化的具体建议。\footnote{Yao
  Wenyuan, On the Social Basis of the Lin Biao Anti-Party Clique
  (Beijing: Foreign Languages Press, 1975); Zhang Chunqiao, On
  Exercising All-Round Dictatorship over the Bourgeoisie (Beijing:
  Foreign Languages Press, 1975). 姚文元：On the Social Basis of the Lin
  Biao Anti-Party Clique（北京：外文出版社，1975年）；张春桥：On
  Exercising All-Round Dictatorship over the
  Bourgeoisie（北京：外文出版社，1975年）。}

Deng Xiaoping very quickly made it clear that the discussion on
bourgeois rights was completely irrelevant to the realities of China's
backward economy. He reportedly stated in March 1975 that, ``To restrict
bourgeois rights, we also must have a material foundation, for without
it, how can we do so?''\footnote{Cited by Gong Xiaowen, ``Deng Xiaoping
  and the `20 Articles'\,'', 学习与批判 (Studies and Criticism), No.~6
  (1976), in Selections from People's Republic of China Magazines, 879
  (July 12,1976), p.~2. 转引自龚小文：``Deng Xiaoping and the `20
  Articles'\,''，《学习与批判》1976年第6期，载Selections from People's
  Republic of China Magazines，879（1976年7月12日），第2页。} In June
1975 whilst in Shanghai, Deng commented on proposals to reduce income
differentials by remarking; ``Why do we describe everything as bourgeois
right? Isn't it right to receive more for more work done? Should this be
called bourgeois right?''\footnote{Ibid., p.~5. See also Tang Tsou,
  ``Mao Tsetung Thought, the Last Struggle for Succession, and the
  Post-Mao Era'', CQ, 71 (1977), p.~514. 同上，第5页。另见邹谠：``Mao
  Tsetung Thought, the Last Struggle for Succession, and the Post-Mao
  Era''，CQ，71（1977年），第514页。} His basic stand was summed up in
two short sentences. ``Revolution is liberating the productive forces.
Revolution is promoting the development of productive
forces''.\footnote{``General Programme of Work for the Whole Party and
  the Whole Country'', October 7,1975, in Chi Hsin, The Case of the Gang
  of Four, (Hong Kong: Cosmos Books, 1977), p.~226. ``General Programme
  of Work for the Whole Party and the Whole
  Country''，1975年10月7日，载Chi Hsin, The Case of the Gang of
  Four（香港：Cosmos Books，1977年），第226页。}

邓小平很快明确表示，关于资产阶级法权的讨论同中国落后经济的现实完全无关。据说他在1975年3月说：``限制资产阶级法权，也要有物质基础，没有物质基础怎么限制？''\footnote{Cited
  by Gong Xiaowen, ``Deng Xiaoping and the `20 Articles'\,'', 学习与批判
  (Studies and Criticism), No.~6 (1976), in Selections from People's
  Republic of China Magazines, 879 (July 12,1976), p.~2.
  转引自龚小文：``Deng Xiaoping and the `20
  Articles'\,''，《学习与批判》1976年第6期，载Selections from People's
  Republic of China Magazines，879（1976年7月12日），第2页。}
1975年6月在上海时，邓在评论缩小收入差距的建议时说：``为什么什么都叫资产阶级法权？多劳多得对不对？这是不是叫资产阶级法权？''\footnote{Ibid.,
  p.~5. See also Tang Tsou, ``Mao Tsetung Thought, the Last Struggle for
  Succession, and the Post-Mao Era'', CQ, 71 (1977), p.~514.
  同上，第5页。另见邹谠：``Mao Tsetung Thought, the Last Struggle for
  Succession, and the Post-Mao Era''，CQ，71（1977年），第514页。}
他的基本立场可以概括为两句简短的话：``革命是解放生产力。革命是促进生产力的发展。''\footnote{``General
  Programme of Work for the Whole Party and the Whole Country'', October
  7,1975, in Chi Hsin, The Case of the Gang of Four, (Hong Kong: Cosmos
  Books, 1977), p.~226. ``General Programme of Work for the Whole Party
  and the Whole Country''，1975年10月7日，载Chi Hsin, The Case of the
  Gang of Four（香港：Cosmos Books，1977年），第226页。}

The debate soon degenerated into a slanging match over which ideological
deviation was more dangerous -- empiricism or dogmatism. Empiricism was
the radicals' label for veteran cadres who emphasized practical
experience at the expense of theory. On 1 March 1975, the same day as
the publication of Yao Wenyuan's article, Zhang Chunqiao delivered a
speech to a meeting of PLA political department directors in which he
stressed the dangers of empiricism.\footnote{See Hu Hua (ed.), Zhongguo
  shehuizhuyi geming, p.~325; Gao Gao and Yan Jiaqi, ``Wenhua
  dageming'', pp.~538-39; Xinhua, December 1,1977, SWB/FE/5685/BII/2.
  见胡华编：《中国社会主义革命》，第325页；高皋、严家其：《文化大革命》，第538-539页；新华社，1977年12月1日，SWB/FE/5685/BII/2。}
Jiang Qing expressed her agreement with Zhang and Yao's position in
speeches of 3, 4 and 5 April 1975.\footnote{See Gao Gao and Yan Jiaqi,
  ``Wenhua dageming'', p.~540. 见高皋、严家其：《文化大革命》，第540页。}
Dogmatism, on the other hand, was the veterans' description for theory
separated from practical reality. On 23 April 1975, in comments he
appended to a Xinhua despatch submitted by Yao Wenyuan, Mao concluded
with the following biting remarks:

这场争论很快退化为围绕哪种思想偏差更危险的相互攻讦：经验主义还是教条主义。经验主义是激进派给那些重视实践经验而轻视理论的老干部贴上的标签。1975年3月1日，也就是姚文元文章发表的同一天，张春桥在中国人民解放军政治部主任会议上发表讲话，强调经验主义的危险。\footnote{See
  Hu Hua (ed.), Zhongguo shehuizhuyi geming, p.~325; Gao Gao and Yan
  Jiaqi, ``Wenhua dageming'', pp.~538-39; Xinhua, December 1,1977,
  SWB/FE/5685/BII/2.
  见胡华编：《中国社会主义革命》，第325页；高皋、严家其：《文化大革命》，第538-539页；新华社，1977年12月1日，SWB/FE/5685/BII/2。}
江青在1975年4月3日、4日和5日的讲话中表示赞同张、姚的立场。\footnote{See
  Gao Gao and Yan Jiaqi, ``Wenhua dageming'', p.~540.
  见高皋、严家其：《文化大革命》，第540页。}
另一方面，教条主义则是老干部用来形容脱离实践现实的理论。1975年4月23日，毛在姚文元提交的一份新华社电讯上所加的批语中，以如下尖锐评论作结：

\begin{quote}
There are not many in our Party who really understand Marxism-Leninism.
Some think they do understand, while in actuality they do not fully
understand even though they think they are right and are always
lecturing others. That is also a manifestation of ignorance of
Marxism-Leninism.\footnote{``Document of the CCP CC'', I\&S, 15: 2,
  pp.~95-6; Liushinian, pp.~645-6; Gao Gao and Yan Jiaqi, ``Wenhua
  dageming'', p.~642. A Hongqi article of May 1975 quoted Mao from 1942
  to attack dogmatism (Yao). See Tian Chun, ``精通的目的全在于应用''
  (Mastering theory is for the sole purpose of applying it), Hongqi,
  No.~5,1975, p.~7. ``Document of the CCP
  CC''，I\&S，15:2，第95-96页；《六十年》，第645-646页；高皋、严家其：《文化大革命》，第642页。1975年5月《红旗》一篇文章引用毛1942年的话攻击教条主义（姚）。见田春：《精通的目的全在于应用》，《红旗》1975年第5期，第7页。}
\end{quote}

\begin{quote}
我们党内真正懂得马克思列宁主义的人并不多。有些人自以为懂，实际上并不完全懂，尽管他们自以为正确，还总是教训别人。这也是不懂马克思列宁主义的一种表现。\footnote{``Document
  of the CCP CC'', I\&S, 15: 2, pp.~95-6; Liushinian, pp.~645-6; Gao Gao
  and Yan Jiaqi, ``Wenhua dageming'', p.~642. A Hongqi article of May
  1975 quoted Mao from 1942 to attack dogmatism (Yao). See Tian Chun,
  ``精通的目的全在于应用'' (Mastering theory is for the sole purpose of
  applying it), Hongqi, No.~5,1975, p.~7. ``Document of the CCP
  CC''，I\&S，15:2，第95-96页；《六十年》，第645-646页；高皋、严家其：《文化大革命》，第642页。1975年5月《红旗》一篇文章引用毛1942年的话攻击教条主义（姚）。见田春：《精通的目的全在于应用》，《红旗》1975年第5期，第7页。}
\end{quote}

The effect of the campaign on policy direction, wages' policy in
particular, is unclear. It was asserted at the time by the Taiwan
Central News Agency that in February 1975, during a talk in Shanghai,
Jiang Qing had demanded that individuals earning high incomes take a
``voluntary'' cut.\footnote{CNA (Taibei), July 3, 1975,
  SWB/FE/4949/BII/8.
  中央通讯社（台北），1975年7月3日，SWB/FE/4949/BII/8。} However, it
seems that more detailed discussions went no further than supporting the
existing system while calling for a gradual reduction in individual
income differentials.\footnote{See Zhong Shi,
  ``无产阶级专政和资产阶级法权'' (The dictatorship of the proletariat
  and bourgeois rights), Hongqi, 9 (1975), p.~55.
  见钟实：《无产阶级专政和资产阶级法权》，《红旗》1975年第9期，第55页。}
In the second half of 1974 some local authorities had reportedly granted
wage rises and reinstituted bonuses.\footnote{See CNS, 555 (February 19,
  1975), pp.~1-3. 见CNS，555（1975年2月19日），第1-3页。} The campaign
of 1975 may therefore best be seen in the context of a reaction to
policy initiatives considered politically unsound, rather than as a
radical thrust in itself. The national press carried a series of
articles railing against material incentives and bonuses, and bonuses
were reportedly abolished in the railway system. Reported attempts were
made in February and March 1975 to introduce voluntary Saturday labor on
a national basis and unpaid weekday labor in Henan province. But the
experiments were short-lived and unsuccessful.\footnote{See CNS, 567
  (May 21,1975), 577, esp.~pp.~4-5, and 580 (August 27,1975).
  见CNS，567（1975年5月21日）、577，尤其第4-5页，以及580（1975年8月27日）。}

这场运动对政策方向，尤其是工资政策的影响并不清楚。当时台湾中央通讯社声称，1975年2月江青在上海一次谈话中要求高收入者``自愿''减薪。\footnote{CNA
  (Taibei), July 3, 1975, SWB/FE/4949/BII/8.
  中央通讯社（台北），1975年7月3日，SWB/FE/4949/BII/8。}
然而，较为详细的讨论似乎并未超出支持现行制度、同时要求逐步缩小个人收入差距的范围。\footnote{See
  Zhong Shi, ``无产阶级专政和资产阶级法权'' (The dictatorship of the
  proletariat and bourgeois rights), Hongqi, 9 (1975), p.~55.
  见钟实：《无产阶级专政和资产阶级法权》，《红旗》1975年第9期，第55页。}
据报道，1974年下半年，一些地方当局曾批准涨工资并恢复奖金。\footnote{See
  CNS, 555 (February 19, 1975), pp.~1-3.
  见CNS，555（1975年2月19日），第1-3页。}
因而，1975年的这场运动也许最好被视为对那些被认为政治上不可靠的政策倡议的反应，而不是其本身的一次激进推进。全国性报刊刊登了一系列抨击物质刺激和奖金的文章，据报道铁路系统的奖金被取消。据称，1975年2月和3月曾有人试图在全国范围推行自愿星期六劳动，并在河南省推行平日无偿劳动。但这些试验都短暂且不成功。\footnote{See
  CNS, 567 (May 21,1975), 577, esp.~pp.~4-5, and 580 (August 27,1975).
  见CNS，567（1975年5月21日）、577，尤其第4-5页，以及580（1975年8月27日）。}

One observer has correctly described the campaign as ``so ambitious in
theory yet so modest in intent''. ``The radicals'', he observed,

一位观察者正确地把这场运动描述为``理论上如此雄心勃勃，意图上却如此温和''。他评论说，``激进派''

\begin{quote}
succeeded in blocking a scheduled wage adjustment and called on workers
voluntarily to forego extra pay and donate their labor, but they did not
call for an immediate elimination of commodity trade and the monetary
exchange system, for levelling distribution patterns, or for a shift
from the ``works'' to ``needs'' principle.
\end{quote}

\begin{quote}
成功阻止了一次预定的工资调整，并号召工人自愿放弃额外工资、奉献劳动，但他们并没有要求立即取消商品贸易和货币交换制度，没有要求拉平分配格局，也没有要求从``按劳''原则转向``按需''原则。
\end{quote}

The radicals thus ``found themselves in the awkward position of
nullifying the legitimacy and ideological rationale of `bourgeois right'
while ignoring its societal manifestations.''\footnote{Dittmer, ``The
  Radical Critique of Political Interest'', p.~375. Dittmer：``The
  Radical Critique of Political Interest''，第375页。} Another analyst
also commented on what he termed the ``defensive nature of the
campaign'' and asserted that it aimed ``to deter the growth of new
inequalities rather than to impose Spartan egalitarianism on Chinese
society''.\footnote{Kraus, Class Conflicts in Chinese Socialism, p.~166.
  Kraus, Class Conflicts in Chinese Socialism，第166页。}

因此，激进派``发现自己处于一种尴尬境地：他们否定了`资产阶级法权'的合法性和意识形态理由，却又无视它的社会表现。''\footnote{Dittmer,
  ``The Radical Critique of Political Interest'', p.~375. Dittmer：``The
  Radical Critique of Political Interest''，第375页。}
另一位分析者也评论了他所谓这场``运动的防御性质''，并断言其目标是``阻止新的不平等增长，而不是把斯巴达式的平均主义强加给中国社会''。\footnote{Kraus,
  Class Conflicts in Chinese Socialism, p.~166. Kraus, Class Conflicts
  in Chinese Socialism，第166页。}

Certain Taiwan sources, however, drew the unjustifiable conclusion from
mainland propaganda that a wholesale adjustment to the industrial wage
structure was carried out, leading to widespread discontent among
workers. It was claimed that in April 1975 the central authorities
issued a document simplifying the wage structure for workers from eight
to three grades, allowing for large differences between each grade.
Another source asserted that factories and mines in various places
changed the eight-grade scale to five, three or two grades, ``seriously
lowering the already poor standard of living''.\footnote{See World
  Anti-Communist League, China Chapter, Asian People's Anti-Communist
  League, Charts concerning Chinese Communists on the Mainland,
  (hereafter cited as Charts), 34th Series (April, 1977), p.~4;
  ``从杭州工人抗暴到天安门事件'' (From the revolt of the workers of
  Hangzhou to the ``Tiananmen Incident''), 中共年报 1976, (Taibei:
  Center for the Study of Chinese Communist Studies, 1976), 5:101.
  见世界反共联盟中国分会、亚洲人民反共联盟：《Charts concerning Chinese
  Communists on the
  Mainland》（下文简称Charts），第34辑（1977年4月），第4页；《从杭州工人抗暴到天安门事件》，《中共年报1976》（台北：中华共产主义研究中心，1976年），5:101。}
Observers inside China reported quite a different story. The Times'
correspondent in Beijing visited a coal mine and was informed that while
changes to the wage system were being discussed, no decision had been
reached. A Canadian student in Beijing was told that changes to the wage
system were confined to raising the lowest grade and lowering the
highest. However, strong opposition from workers had prevented any
reduction in the highest grade.\footnote{The Times, 15 July 1975, p.~7;
  David Zweig, ``The Peita Debate on Education and the Fall of Teng
  Hsiao-p'ing'', CQ, 73 (1978), p.~150.
  《泰晤士报》，1975年7月15日，第7页；David Zweig：``The Peita Debate on
  Education and the Fall of Teng
  Hsiao-p'ing''，CQ，73（1978年），第150页。} Given the extreme
conservatism of wage policy in the Maoist era, these accounts ring true.
This is not to deny that workers had real grievances over their pay and
working conditions. The rhetoric of the campaign undoubtedly served to
attract lower-paid workers to the radical camp as well as raise
expectations on the wages' front.\footnote{See John P. Burns, ``The
  Radicals and the Campaign to Limit Bourgeois Rights in the
  Countryside'', Contemporary China, 1:3 (January 1977), p.~26, and
  Goldman, China's Intellectuals, p.~201. 见John P. Burns：``The
  Radicals and the Campaign to Limit Bourgeois Rights in the
  Countryside''，Contemporary
  China，1:3（1977年1月），第26页；以及Goldman, China's
  Intellectuals，第201页。}

然而，某些台湾来源从大陆宣传中得出了站不住脚的结论，认为工业工资结构进行了全面调整，并导致工人普遍不满。有人声称，1975年4月中央当局发布了一份文件，把工人工资结构从八级简化为三级，并允许各级之间存在很大差别。另一来源则断言，各地工厂和矿山把八级工资制改为五级、三级或两级，``严重降低了本已低下的生活水平''。\footnote{See
  World Anti-Communist League, China Chapter, Asian People's
  Anti-Communist League, Charts concerning Chinese Communists on the
  Mainland, (hereafter cited as Charts), 34th Series (April, 1977),
  p.~4; ``从杭州工人抗暴到天安门事件'' (From the revolt of the workers
  of Hangzhou to the ``Tiananmen Incident''), 中共年报 1976, (Taibei:
  Center for the Study of Chinese Communist Studies, 1976), 5:101.
  见世界反共联盟中国分会、亚洲人民反共联盟：《Charts concerning Chinese
  Communists on the
  Mainland》（下文简称Charts），第34辑（1977年4月），第4页；《从杭州工人抗暴到天安门事件》，《中共年报1976》（台北：中华共产主义研究中心，1976年），5:101。}
中国境内的观察者报道的情况却大不相同。《泰晤士报》驻北京记者参观了一座煤矿，被告知工资制度的变动仍在讨论中，尚未作出决定。一名在北京的加拿大学生被告知，工资制度的改变仅限于提高最低等级和降低最高等级。然而，工人的强烈反对阻止了最高等级的任何降低。\footnote{The
  Times, 15 July 1975, p.~7; David Zweig, ``The Peita Debate on
  Education and the Fall of Teng Hsiao-p'ing'', CQ, 73 (1978), p.~150.
  《泰晤士报》，1975年7月15日，第7页；David Zweig：``The Peita Debate on
  Education and the Fall of Teng
  Hsiao-p'ing''，CQ，73（1978年），第150页。}
鉴于毛主义时代工资政策极端保守，这些说法显得可信。这并不是否认工人对工资和劳动条件确有真实不满。运动的修辞无疑有助于把低收入工人吸引到激进派阵营，同时也提高了工资问题上的期望。\footnote{See
  John P. Burns, ``The Radicals and the Campaign to Limit Bourgeois
  Rights in the Countryside'', Contemporary China, 1:3 (January 1977),
  p.~26, and Goldman, China's Intellectuals, p.~201. 见John P.
  Burns：``The Radicals and the Campaign to Limit Bourgeois Rights in
  the Countryside''，Contemporary
  China，1:3（1977年1月），第26页；以及Goldman, China's
  Intellectuals，第201页。}

\section{Factionalism at the Center, February-June
1975}\label{factionalism-at-the-center-february-june-1975}

\section{中央的派性，1975年2月-6月}\label{ux4e2dux592eux7684ux6d3eux60271975ux5e742ux6708-6ux6708}

After the 10th CCP 2nd plenum, at Mao's request Deng had assumed control
of the day-to-day affairs of the Central Committee. He took over
responsibility from Wang Hongwen, who had reportedly disappointed the
Chairman.\footnote{Ye Yonglie, Wang Hongwen xingshuailu, p.~422, dates
  Deng's assumption of this responsibility from the 3 May Politburo
  meeting.
  叶永烈：《王洪文兴衰录》，第422页，将邓接管这一职责的时间定为5月3日政治局会议。}
But Mao's confidence in his independent-minded deputy was to be sorely
tested in the following months. Dissension within the highest
decision-making circles was rife, forcing Mao to intercede. At a meeting
of the Politburo on 3 May 1975, Mao repeated an earlier warning of July
1974 about the ``Gang of Four'', while simultaneously making unspecified
criticisms of Deng's abrasiveness and impulsiveness.\footnote{``Document
  of the CCP CC'', pp.~96-7. When news of Mao's criticism of the Four
  reached the ears of the ill Kang Sheng he insisted that Mao's liaison
  personnel inform the Chairman of Jiang Qing and Zhang Chunqiao's
  allegedly treacherous history. Kang reportedly named the person who
  could prove this allegation. See Gao Gao and Yan Jiaqi, ``Wenhua
  dageming'', pp.~543-4; Lin Qingshan, Kang Sheng waizhuan, p.~377.
  ``Document of the CCP
  CC''，第96-97页。病中的康生听到毛批评四人的消息后，坚持要毛的联络人员把江青和张春桥据称有叛变历史一事告知主席。据说康还说出了可以证明这一指控的人名。见高皋、严家其：《文化大革命》，第543-544页；林青山：《康生外传》，第377页。}
Yet Mao added that the problem was not great and that it should not be
blown out of proportion (不要小题大作).\footnote{Hu Hua (ed.), Zhongguo
  shehuizhuyi geming, p.~325. 胡华编：《中国社会主义革命》，第325页。}
This session of the Politburo followed one on 27 April at which Jiang
Qing had come under fire for her speeches earlier in the month. Then, on
27 May and 3 June, the Politburo under the chairmanship of Deng launched
attacks on the four radicals. Of the four, Wang Hongwen alone admitted
to mistakes. However, in a letter to her husband on 28 June, Jiang Qing
stated that the criticisms of her Politburo colleagues had really shaken
her up (触动很大) and confessed that the existence of a ``Gang of Four''
had the potential to split the party center.\footnote{Hao Mengbi and
  Duan Haoran, Liushinian, p.~646. 郝梦笔、段浩然：《六十年》，第646页。}

中共十届二中全会之后，应毛的要求，邓已经接管了中央委员会的日常事务。他从据称令主席失望的王洪文手中接过了这项职责。\footnote{Ye
  Yonglie, Wang Hongwen xingshuailu, p.~422, dates Deng's assumption of
  this responsibility from the 3 May Politburo meeting.
  叶永烈：《王洪文兴衰录》，第422页，将邓接管这一职责的时间定为5月3日政治局会议。}
但是，毛对这位有独立主见的副手的信任，在随后几个月中将受到严峻考验。最高决策圈内部异议纷呈，迫使毛出面干预。1975年5月3日的一次政治局会议上，毛重复了1974年7月关于``四人帮''的早先警告，同时对邓的尖锐和冲动提出未具体说明的批评。\footnote{``Document
  of the CCP CC'', pp.~96-7. When news of Mao's criticism of the Four
  reached the ears of the ill Kang Sheng he insisted that Mao's liaison
  personnel inform the Chairman of Jiang Qing and Zhang Chunqiao's
  allegedly treacherous history. Kang reportedly named the person who
  could prove this allegation. See Gao Gao and Yan Jiaqi, ``Wenhua
  dageming'', pp.~543-4; Lin Qingshan, Kang Sheng waizhuan, p.~377.
  ``Document of the CCP
  CC''，第96-97页。病中的康生听到毛批评四人的消息后，坚持要毛的联络人员把江青和张春桥据称有叛变历史一事告知主席。据说康还说出了可以证明这一指控的人名。见高皋、严家其：《文化大革命》，第543-544页；林青山：《康生外传》，第377页。}
但毛又补充说，问题不大，不要小题大作。\footnote{Hu Hua (ed.), Zhongguo
  shehuizhuyi geming, p.~325. 胡华编：《中国社会主义革命》，第325页。}
这次政治局会议之前，4月27日曾开过一次会，江青因本月早些时候的讲话受到批评。随后，5月27日和6月3日，在邓主持下，政治局向四名激进派发动攻击。四人之中，只有王洪文承认错误。不过，江青在6月28日给丈夫的信中说，政治局同事对她的批评使她触动很大，并承认``四人帮''的存在有可能分裂党中央。\footnote{Hao
  Mengbi and Duan Haoran, Liushinian, p.~646.
  郝梦笔、段浩然：《六十年》，第646页。}

Despite Mao's reservations about his leadership style, Deng's crusade to
``rectify'' all aspects of the affairs of the CCP, especially those
related to economic policy, continued unabated. Mao had called for
rectification but had not specified how it should be carried
out.\footnote{See Dangshi tongxun, 1 (January 15, 1983), p.~10.
  见《党史通讯》第1期（1983年1月15日），第10页。} Given past experience,
however, the Chairman's reticence should have made Deng extremely
circumspect. Nevertheless, he pressed on. In a speech at the forum on
the iron and steel industry held at the end of May 1975 Deng rallied his
supporters to stand fast against factionalism.\footnote{Deng Xiaoping,
  ``当前钢铁工业必须解决的几个问题'' (Several problems which the iron
  and steel industry should presently solve), 29 May 1975, Deng Xiaoping
  Wenxuan, pp.~8-11. See also Gao Gao and Yan Jiaqi, ``Wenhua
  dageming'', pp.~544-45; Liu Suinian and Wu Qungan, ``Wenhua dageming''
  shiqide guomin jingji, p.~170.
  邓小平：《当前钢铁工业必须解决的几个问题》，1975年5月29日，《邓小平文选》，第8-11页。另见高皋、严家其：《文化大革命》，第544-545页；刘岁年、吴群敢：《``文化大革命''时期的国民经济》，第170页。}
Repeating earlier demands that discipline be enforced in the workplace,
Deng advised factory officials to withhold pay from absentee workers and
to order out of the factory gates those who refused to work. On 4 June
1975 the central authorities issued Document no. 13 of 1975 calling for
fulfillment of the year's plan in the iron and steel
industry.\footnote{See Hu Hua (ed.), Zhongguo shehuizhuyi geming,
  p.~323. 见胡华编：《中国社会主义革命》，第323页。} The shortage of
coal was the principal factor contributing to the problems in the
industry,\footnote{CNS, 577, pp.~2-4. CNS，577，第2-4页。} but by the
end of June the situation in iron and steel production had improved
somewhat.\footnote{Hao Mengbi and Duan Haoran, Liushinian, p.~643.
  郝梦笔、段浩然：《六十年》，第643页。} Between 16 June and 11 August
the State Council convened a conference to discuss the basic principles
guiding the planned economy. The meeting decided to reassert central
control over state enterprises, many of which had been handed over to
local authorities in the major decentralization decision of
1970.\footnote{For a discussion of the debate which ranged Deng against
  certain provincial interests, see Jonathan Unger, ``The Struggle to
  Dictate China's Administration: The Conflict of Branches vs Areas vs
  Reform'', Australian Journal of Chinese Affairs, 18 (July 1987),
  esp.~pp.~27-8. In Zhejiang two articles appeared in ZJRB, September
  5,1976, p.~2, denouncing Deng for attempting to encroach on area
  interests. 关于这场使邓同某些省级利益相对立的争论，见Jonathan
  Unger：``The Struggle to Dictate China's Administration: The Conflict
  of Branches vs Areas vs Reform''，Australian Journal of Chinese
  Affairs，18（1987年7月），尤其第27-28页。在浙江，1976年9月5日《浙江日报》第2版刊登两篇文章，谴责邓试图侵占地方利益。}
It also stipulated that the maximum number of non-productive personnel
permitted in state enterprises was limited to 18\% in large and
medium-sized units and 10\% in smaller plants.\footnote{See Liu Suinian
  and Wu Qungan, ``Wenhua dageming'' de guomin jingji, p.~170.
  见刘岁年、吴群敢：《``文化大革命''的国民经济》，第170页。} Deng's
drive and indefatigability appeared to be having an effect.

尽管毛对邓的领导风格有所保留，邓``整顿''中共各方面事务、尤其是经济政策相关事务的运动仍有增无减。毛曾要求整顿，但没有说明应如何进行。\footnote{See
  Dangshi tongxun, 1 (January 15, 1983), p.~10.
  见《党史通讯》第1期（1983年1月15日），第10页。}
然而，鉴于过去经验，主席的沉默本应使邓极为谨慎。尽管如此，他仍继续推进。1975年5月底在钢铁工业座谈会上，邓发表讲话，号召支持者坚决反对派性。\footnote{Deng
  Xiaoping, ``当前钢铁工业必须解决的几个问题'' (Several problems which
  the iron and steel industry should presently solve), 29 May 1975, Deng
  Xiaoping Wenxuan, pp.~8-11. See also Gao Gao and Yan Jiaqi, ``Wenhua
  dageming'', pp.~544-45; Liu Suinian and Wu Qungan, ``Wenhua dageming''
  shiqide guomin jingji, p.~170.
  邓小平：《当前钢铁工业必须解决的几个问题》，1975年5月29日，《邓小平文选》，第8-11页。另见高皋、严家其：《文化大革命》，第544-545页；刘岁年、吴群敢：《``文化大革命''时期的国民经济》，第170页。}
他重申早先关于在工作场所执行纪律的要求，建议工厂负责人扣发旷工工人的工资，并把拒绝工作的人员赶出工厂大门。1975年6月4日，中央当局发布1975年第13号文件，要求完成当年的钢铁工业计划。\footnote{See
  Hu Hua (ed.), Zhongguo shehuizhuyi geming, p.~323.
  见胡华编：《中国社会主义革命》，第323页。}
煤炭短缺是造成该行业问题的主要因素，\footnote{CNS, 577, pp.~2-4.
  CNS，577，第2-4页。} 但到6月底，钢铁生产形势已有所改善。\footnote{Hao
  Mengbi and Duan Haoran, Liushinian, p.~643.
  郝梦笔、段浩然：《六十年》，第643页。}
6月16日至8月11日，国务院召开会议，讨论指导计划经济的基本原则。会议决定重新加强中央对国营企业的控制，其中许多企业曾在1970年重大下放决定中移交给地方当局。\footnote{For
  a discussion of the debate which ranged Deng against certain
  provincial interests, see Jonathan Unger, ``The Struggle to Dictate
  China's Administration: The Conflict of Branches vs Areas vs Reform'',
  Australian Journal of Chinese Affairs, 18 (July 1987), esp.~pp.~27-8.
  In Zhejiang two articles appeared in ZJRB, September 5,1976, p.~2,
  denouncing Deng for attempting to encroach on area interests.
  关于这场使邓同某些省级利益相对立的争论，见Jonathan Unger：``The
  Struggle to Dictate China's Administration: The Conflict of Branches
  vs Areas vs Reform''，Australian Journal of Chinese
  Affairs，18（1987年7月），尤其第27-28页。在浙江，1976年9月5日《浙江日报》第2版刊登两篇文章，谴责邓试图侵占地方利益。}
会议还规定，国营企业中允许存在的非生产人员最高比例，大中型单位限制在18\%，小型工厂限制在10\%。\footnote{See
  Liu Suinian and Wu Qungan, ``Wenhua dageming'' de guomin jingji,
  p.~170. 见刘岁年、吴群敢：《``文化大革命''的国民经济》，第170页。}
邓的推动力和不知疲倦似乎正在产生效果。

\section{Solving the ``Zhejiang
Problem''}\label{solving-the-zhejiang-problem}

\section{解决``浙江问题''}\label{ux89e3ux51b3ux6d59ux6c5fux95eeux9898}

With the conclusion of the campaign against Lin Biao and Confucius the
factionalism which the campaign had generated and intensified at local
levels demanded urgent attention. Provincial administrations such as
Zhejiang had failed to cope with the challenge posed by former Cultural
Revolution rebels, as the previous chapter of this study has
demonstrated at some length. The slump in industrial production was
evidence of the damage factionalism had wrought. Unity at the center was
the most important precondition for curbing the worst manifestations of
factionalism. As previous sections of this chapter have shown, the
center was anything but united at this crucial time. Nevertheless, there
seems to have been a tacit agreement that a truce should be declared in
the factional war and negotiations begun to find a solution to problems
in the localities.

随着批林批孔运动结束，这场运动在地方层面造成并加剧的派性问题迫切需要处理。正如本研究上一章已较详细说明的，浙江等省级行政机构未能应对前文化大革命造反派提出的挑战。工业生产下滑证明了派性造成的损害。中央的团结是遏制派性最恶劣表现的最重要前提。正如本章前几节所示，在这个关键时刻，中央远非团结一致。尽管如此，似乎存在一种默契，即应在派系战争中宣布休战，并开始谈判，为地方问题寻找解决办法。

To this end, in December 1974 the CC ordered the major participants in
the factional struggles in Zhejiang to join a study class and attend a
conference in Beijing.\footnote{Zhang Yongshengde zuizheng cailiao,
  pp.~13, 58. 《张永生的罪证材料》，第13、58页。} Negotiations proved
long and difficult and it appears that they lasted for over one month,
until late December or early January 1975. High on the agenda were such
controversial topics as the ``two crashthroughs'', the role of urban
militia commands and the activities of factional groups. Wang Hongwen
played a central role in the deliberations. He addressed the
participants at least twice, on December 17 and 19, and was not slow to
express his preferences for different factional leaders and
organizations.\footnote{Zhang Yongshengde zuizheng cailiao, p.~58. Bian
  Wen, HZRB, December 5, 1976.
  《张永生的罪证材料》，第58页。边文，《杭州日报》，1976年12月5日。} But
even Wang distanced himself somewhat from his over-zealous proteges and
could not protect Weng Senhe and He Xianchun from criticism. The
political pendulum had swung against the central radicals and Wang
informed Weng that although he shared his views.

为此，1974年12月，中央命令浙江派系斗争的主要参与者参加学习班，并赴北京出席会议。\footnote{Zhang
  Yongshengde zuizheng cailiao, pp.~13, 58.
  《张永生的罪证材料》，第13、58页。}
谈判证明漫长而艰难，似乎持续了一个多月，直到1974年12月下旬或1975年1月初。议程上的重要议题包括``双突''、城市民兵指挥部的作用以及派系集团的活动等有争议的问题。王洪文在商议中发挥了核心作用。他至少两次在12月17日和19日向与会者讲话，并毫不迟疑地表达了自己对不同派系领导人和组织的偏好。\footnote{Zhang
  Yongshengde zuizheng cailiao, p.~58. Bian Wen, HZRB, December 5, 1976.
  《张永生的罪证材料》，第58页。边文，《杭州日报》，1976年12月5日。}
但即使是王，也同他那些过分热心的被保护者保持了一定距离，无法保护翁森鹤和贺贤春免受批评。政治钟摆已经向不利于中央激进派的方向摆动，王告诉翁，尽管他赞同翁的观点，

``I'm not the Political Bureau''.\footnote{HZRB, December 5,1977.
  《杭州日报》，1977年12月5日。} He Xianchun was concerned that the
meeting would proscribe the ``two crashthroughs''.\footnote{HZRB, June
  20,1977. 《杭州日报》，1977年6月20日。} Wang tried to cheer up the
disconsolate He with the words, ``You must try and win credit''
(争气).\footnote{CCP HMC party school study class, ``What credit did He
  Xianchun strive for?'', HZRB, June 26, 1977.
  中共杭州市委党校学习班：《贺贤春争的是什么气？》，《杭州日报》，1977年6月26日。}

``我不是政治局''。\footnote{HZRB, December 5,1977.
  《杭州日报》，1977年12月5日。}
贺贤春担心会议会禁止``双突''。\footnote{HZRB, June 20,1977.
  《杭州日报》，1977年6月20日。}
王试图用``你要争气''这句话鼓励沮丧的贺。\footnote{CCP HMC party school
  study class, ``What credit did He Xianchun strive for?'', HZRB, June
  26, 1977.
  中共杭州市委党校学习班：《贺贤春争的是什么气？》，《杭州日报》，1977年6月26日。}

Before the opening of the conference, the CCP ZPC had submitted a
document to the CC entitled ``A report requesting instructions on how to
handle correctly the crash induction of party members and the crash
promotion of cadres''
(关于正确处理突击发展的党员和突击提拔的干部的请示报告). In response, the
CC issued a document approving the suggestions contained in the Zhejiang
report.\footnote{Tang Yurui, HZRB, December 8,1976, p.~2.
  唐玉瑞，《杭州日报》，1976年12月8日，第2版。} The center decided to
disband the campaign small groups and Zhang, Weng and He were worried
about the fate of the groups' members. In a letter to Wang Hongwen of
December 14 they brought to his attention the possibility that members
at the grassroots could be accused of arbitrarily leaving their work
posts and have their wages deducted. Moreover, they could be further
accused of sabotaging the unified leadership of the party and of riding
roughshod over party committees. At some places and in some units new
and veteran cadres could therefore be hauled out and struggled against.
Decisions already made to promote cadres and admit party members could
be invalidated, and the persons concerned investigated and stood aside.
Those units could thus return to the state of paralysis they were in at
the beginning of the campaign.\footnote{Zhang Yongshengde zuizheng
  cailiao, p.~61. 《张永生的罪证材料》，第61页。}

会议召开前，中共浙江省委曾向中央提交一份文件，题为《关于正确处理突击发展的党员和突击提拔的干部的请示报告》。作为回应，中央发布文件，批准浙江报告中的建议。\footnote{Tang
  Yurui, HZRB, December 8,1976, p.~2.
  唐玉瑞，《杭州日报》，1976年12月8日，第2版。}
中央决定解散运动小组，张、翁、贺担心小组成员的命运。在12月14日致王洪文的信中，他们提醒王注意一种可能：基层成员可能被指责擅自离开工作岗位并被扣发工资。此外，他们还可能进一步被指责破坏党的统一领导、凌驾于党委之上。因此，在一些地方和单位，新老干部可能被揪出来批斗。已经作出的提拔干部和吸收党员决定可能被取消，有关人员受到调查并靠边站。这些单位从而可能回到运动初期的瘫痪状态。\footnote{Zhang
  Yongshengde zuizheng cailiao, p.~61. 《张永生的罪证材料》，第61页。}

Rather than forbidding the ``two crashthroughs'' outright, the central
authorities requested ``deferral'' of party recruitment and cadre
promotion. Weng and He promptly ignored the suggestion by continuing to
stack committees and important posts with their nominees. Weng Senhe
stated defiantly: ``as long as the `two crashthroughs' are secure I'm
willing to go to jail''. (只要``双突''保牢,我宁愿坐牢).\footnote{HZRB,
  December 5,1977. 《杭州日报》，1977年12月5日。} The CC also issued
directives to bring factional confrontation to an end and to curb the
power of the urban militia. However, the directives were ambiguous and
far from decisive. Zhang, Weng and He, in their letter of December 14 to
Wang Hongwen, attempted to prove that the continued existence of the
urban militia was essential to the maintenance of public order in
Hangzhou.\footnote{Zhang Yongshengde zuizheng cailiao, p.~61.
  《张永生的罪证材料》，第61页。} While the conference was still in
session. He Xianchun was forced by the death of his mother to return to
the provincial capital to make arrangements for her funeral. Whilst in
Hangzhou He reportedly consulted senior members of the municipal party
committee on how to improve the reputation of the militia command and to
justify its continued existence. The militia was therefore mobilized to
handle incidents which were manufactured to illustrate its
indispensability to the maintenance of law and order in the city. All
told, ten such incidents took place at this time.

中央当局并未直接禁止``双突''，而是要求``暂缓''发展党员和提拔干部。翁和贺立即无视这一建议，继续把他们提名的人安插进委员会和重要职位。翁森鹤挑衅地说：``只要`双突'保牢，我宁愿坐牢''。\footnote{HZRB,
  December 5,1977. 《杭州日报》，1977年12月5日。}
中央还发布指示，要求结束派系对抗并限制城市民兵的权力。然而，这些指示含糊不清，远非果断。张、翁、贺在12月14日给王洪文的信中，试图证明城市民兵继续存在对维持杭州公共秩序至关重要。\footnote{Zhang
  Yongshengde zuizheng cailiao, p.~61. 《张永生的罪证材料》，第61页。}
会议仍在进行时，贺贤春因母亲去世被迫返回省会安排丧事。据报道，贺在杭州期间同市委高级成员商量如何改善民兵指挥部的声誉，并为其继续存在辩护。因此，民兵被动员起来处理一些人为制造的事件，以说明其对维持城市治安不可或缺。总共在这一时期发生了十起此类事件。

The most serious incident involved the case of fighting and arson at the
Hangzhou Cable Factory on 3 December 1974.\footnote{HZRB, June 20, 1977;
  Zhou Feng, HZRB, October 24,1977; HZRB, November 16,1977; Chen Xia,
  HZRB, April 10, 1978; HZRB, April 17, 1979.
  《杭州日报》，1977年6月20日；周锋，《杭州日报》，1977年10月24日；《杭州日报》，1977年11月16日；陈霞，《杭州日报》，1978年4月10日；《杭州日报》，1979年4月17日。}
Instead of helping fight the fire, the urban militia contingents under
Xia Genfa's command piled more furniture on the blaze and prevented all
efforts to extinguish it. Factional supporters of the militia in the
telecommunications bureau jammed telephone lines into the factory. A
report was then submitted to Beijing stating that the militia had
performed meritorious service in stopping the fighting and putting out
the fire. In order to get to the bottom of the event the CC summoned the
leaders of the CCP HMC, Wang Zida and Wang Xing, to the capital to hear
their explanations. However, considering their affiliation with the
``mountain top'' faction, the two Wangs were unlikely to provide an
objective account of the incident.

最严重的事件是1974年12月3日杭州电缆厂发生的斗殴和纵火案。\footnote{HZRB,
  June 20, 1977; Zhou Feng, HZRB, October 24,1977; HZRB, November
  16,1977; Chen Xia, HZRB, April 10, 1978; HZRB, April 17, 1979.
  《杭州日报》，1977年6月20日；周锋，《杭州日报》，1977年10月24日；《杭州日报》，1977年11月16日；陈霞，《杭州日报》，1978年4月10日；《杭州日报》，1979年4月17日。}
夏根发指挥下的城市民兵分队非但没有帮助灭火，反而把更多家具堆到火上，并阻止一切灭火努力。电信局中支持民兵的派系人员堵塞了通往工厂的电话线路。随后，一份报告被提交到北京，称民兵在制止斗殴和扑灭火灾方面立了功。为查清事件真相，中央召集中共杭州市委领导王子达和王兴到首都听取他们的解释。然而，考虑到他们同``山上派''的关系，两位王不太可能对事件作出客观说明。

Further violent incidents shook Hangzhou as the discussions continued in
Beijing. The Hangzhou Silk Brocade Mill (杭州织锦厂), manufacturer of
the silk portraits of Chairman Mao Zedong which once adorned the homes
of people all over China, was the scene of three major fights in
December 1974.\footnote{The Hangzhou Silk Brocade changed its name to
  the East is Red (东方红) Silk Mill in the Cultural Revolution. In the
  1980s it was renamed after its original owner Du Jingsheng. The mill
  was established in 1922 and in 1977 employed 1,800 workers. Visits to
  Mill February 22, 1977, September 24,1977.
  杭州织锦厂在文化大革命中改名为东方红丝织厂。20世纪80年代，它以原业主都锦生之名重新命名。该厂创建于1922年，1977年有1,800名工人。笔者于1977年2月22日、1977年9月24日访问该厂。}
In one incident twenty-two people were wounded, seven so seriously that
two and a half years later they could not work.\footnote{HZRB, June
  30,1977. 《杭州日报》，1977年6月30日。} Large-scale fighting also
erupted at Weng's silk mill on December 25, 1974, and the pro-Weng
forces were forced out of the factory. In retaliation, Weng ordered his
supporters to attack the provincial guesthouse and three people were
injured and property destroyed.\footnote{See Weng Senhe bufen zuizheng
  cailiao, p.~66; ``The Case of Weng Sengho'', pp.~4-5. On December
  25,1974, a serious clash occurred between the two factions in Linhai
  county. Five mine employees caught up in the melee were accidentally
  killed. See Linhai xianzhi, p.~31.
  见《翁森鹤部分罪证材料》，第66页；``The Case of Weng
  Sengho''，第4-5页。1974年12月25日，临海县两派之间发生严重冲突。五名卷入混战的矿山职工意外死亡。见《临海县志》，第31页。}

随着北京的讨论继续进行，更多暴力事件震动杭州。杭州织锦厂曾生产毛泽东主席丝织像，这些画像一度装饰着全中国各地人家的住宅；1974年12月，该厂发生了三次大规模斗殴。\footnote{The
  Hangzhou Silk Brocade changed its name to the East is Red (东方红)
  Silk Mill in the Cultural Revolution. In the 1980s it was renamed
  after its original owner Du Jingsheng. The mill was established in
  1922 and in 1977 employed 1,800 workers. Visits to Mill February 22,
  1977, September 24,1977.
  杭州织锦厂在文化大革命中改名为东方红丝织厂。20世纪80年代，它以原业主都锦生之名重新命名。该厂创建于1922年，1977年有1,800名工人。笔者于1977年2月22日、1977年9月24日访问该厂。}
在其中一次事件中，二十二人受伤，其中七人伤势严重，以至两年半后仍无法工作。\footnote{HZRB,
  June 30,1977. 《杭州日报》，1977年6月30日。}
1974年12月25日，翁所在的丝织厂也爆发了大规模斗殴，亲翁力量被赶出工厂。作为报复，翁命令其支持者攻击省招待所，造成三人受伤，财物被毁。\footnote{See
  Weng Senhe bufen zuizheng cailiao, p.~66; ``The Case of Weng Sengho'',
  pp.~4-5. On December 25,1974, a serious clash occurred between the two
  factions in Linhai county. Five mine employees caught up in the melee
  were accidentally killed. See Linhai xianzhi, p.~31.
  见《翁森鹤部分罪证材料》，第66页；``The Case of Weng
  Sengho''，第4-5页。1974年12月25日，临海县两派之间发生严重冲突。五名卷入混战的矿山职工意外死亡。见《临海县志》，第31页。}

A bloody skirmish also took place at the Hangzhou East Wind Silk
Mill.\footnote{Wu Rongtao (of the Hangzhou East Wind Silk Mill),
  ``Indignantly denounce Weng Senhe's towering crimes in suppressing the
  masses'', HZRB, June 13,1977.
  吴荣涛（杭州东风丝织厂）：《愤怒声讨翁森鹤镇压群众的滔天罪行》，《杭州日报》，1977年6月13日。}
The Hangzhou Silk Bureau, under Weng Senhe's direction, was responsible
for twenty silk mills. Weng's supporters at the East Wind mill were
active in trying to change the factory's leadership and Weng's
subordinates at the silk bureau recommended to him that the party
secretary of the East Wind be encouraged to retire so that a younger man
could take his place. However, Weng's appointee as replacement aroused
such differences of opinion among the workforce that the militia was
mobilized to deal with it.

杭州东风丝织厂也发生了一场流血冲突。\footnote{Wu Rongtao (of the
  Hangzhou East Wind Silk Mill), ``Indignantly denounce Weng Senhe's
  towering crimes in suppressing the masses'', HZRB, June 13,1977.
  吴荣涛（杭州东风丝织厂）：《愤怒声讨翁森鹤镇压群众的滔天罪行》，《杭州日报》，1977年6月13日。}
在翁森鹤领导下的杭州市丝绸局负责二十家丝织厂。翁在东风厂的支持者积极试图改变工厂领导班子，翁在丝绸局的下属向他建议，劝东风厂党委书记退休，以便由一名较年轻的人接替。然而，翁指定的继任者在职工中引起如此大的意见分歧，以致民兵被动员来处理此事。

One night in November 1974, the militia arrived at the mill to recruit
workers for the construction of defence lines for a forthcoming battle
at the nearby silk brocade mill. A truckload of men was assigned to
seize two opponents of the ``mountain top'' faction from their homes
located in the grounds of the brocade mill. Workers of the East Wind
mill came out in sufficiently large numbers to support their comrades
that they fought the militia off. When the setback was reported to the
acting head of the silk bureau, a close associate of Weng's, he became
enraged and reportedly instructed: ``the stronghold of the
restorationist forces is the East Wind Mill. They must be annihilated''.
On the evening of 13 November 1974, six truckloads of militiamen
converged on the East Wind mill but again they were beaten off. Weng
Senhe, upon hearing news of this second humiliation, reportedly said:
``this group at the East Wind mill are dumb; your bats and clubs haven't
drawn blood'' (吃素了); the old cadres fought their way up in a hail of
bullets, so we must fight our way up with cane hats and iron clubs''
(老干部是靠枪林弹雨打出来的,我们就要靠藤帽铁棍打出来). Weng dismissed
his subordinates with the words that if they did not crush the
opposition, then they need not bother to come back to him for further
advice.

1974年11月的一个夜晚，民兵来到该厂征调工人，为附近织锦厂即将发生的一场战斗修筑防线。一卡车人员被派去从织锦厂厂区内的住处抓捕两名``山上派''的反对者。东风厂工人大批出来支援自己的同伴，把民兵打退了。当这次挫败报告给丝绸局代理负责人、翁的亲密同伙时，他勃然大怒，据说指示道：``复辟势力的据点就是东风厂。他们必须被消灭。''1974年11月13日晚上，六卡车民兵向东风厂集结，但再次被击退。翁森鹤听到这第二次受辱的消息后，据说说：``你们这帮人是吃素的吗？棍棒还没见血''；老干部是靠枪林弹雨打出来的，我们就要靠藤帽铁棍打出来。翁打发下属离开时说，如果他们不粉碎反对派，就不必再回来向他请示。

Wu Rongtao, the victim and writer of the article describing these
events, then related what occurred on the evening of 3 January 1975. Wu
was at home when late that night over twenty men appeared at his door
brandishing iron clubs. They trussed him up and threw him blindfolded
and gagged onto the floor of a vehicle, before speeding out of the city
proper. Out on a deserted highway in Banshan district, they hauled Wu
out of the car, covered his head with his padded jacket and commenced to
beat, stab and kick him. Wu's assailants left him unconscious at a
T-intersection in the expectation that he would be run over and killed.
However, Wu managed to come around and crawled to a nearby factory.
There he contacted his unit which sent out three cars to rush him to
hospital where he was treated for wounds to the head, rump and leg. The
next day, when his opponents at the mill found that he had not died from
the beating, they claimed that Wu had gone to Banshan to ``make
link-ups'' and had been detained and wounded by people who objected to
his interference in their internal affairs.

受害者、同时也是描述这些事件的文章作者吴荣涛，接着叙述了1975年1月3日晚上发生的事情。当天深夜，吴在家中，二十多名男子挥舞铁棍出现在他门口。他们把他捆起来，蒙住眼睛、堵住嘴，扔到一辆车的地板上，然后迅速驶离市区。在半山区一条荒凉公路上，他们把吴从车里拖出来，用他的棉衣蒙住头，开始殴打、刺伤和踢打他。袭击者把昏迷的吴丢在一个丁字路口，期待他被车碾过而死。然而，吴设法苏醒过来，爬到附近一家工厂。在那里他联系了自己的单位，单位派出三辆车把他急送医院，治疗其头部、臀部和腿部伤口。第二天，厂里的反对者发现他并未被打死，便声称吴去了半山``串联''，因干涉别人内部事务而被反对者扣押并打伤。

Several factors stand out in Wu's clearly subjective account. First, the
savagery of the measures used to deal with factional opponents testified
to the depth of hatred among sections of the working class in Hangzhou.
It is highly likely that Wu and his group dished out their fair share of
beatings to provoke such planned and remorseless revenge. Second, as Wu
himself commented at the end of the account, because the incident was
not an isolated occurrence, workers in the city's factories felt that
their lives were insecure and often stayed home rather than risk
becoming involved in factional disputes and violence.

在吴显然带有主观色彩的叙述中，有几个因素十分突出。第一，用来对付派系对手的手段极为残暴，证明杭州工人阶级某些群体之间仇恨之深。很可能吴及其群体也曾施加过相当数量的殴打，才招致这种有计划且毫不留情的报复。第二，正如吴本人在叙述结尾所评论的，由于这一事件并非孤立发生，市内工厂的工人感到生命没有保障，常常宁愿待在家中，也不愿冒险卷入派系纠纷和暴力。

Another example illustrates how the militia was mobilized to pursue
opponents of the ``mountain top'' faction. The previous chapter has
recounted the exploits of Zhang Jifa, a cadre from the iron and steel
mill and a leader of the ``mountain base'' faction. Zhang had been
arrested in September 1974 and tortured. With the assistance of Li
Kechang (李克昌),\footnote{Li Kechang (1923-84) was born in Shandong
  province. He joined the CCP in 1940 and came south to Zhejiang where
  he worked until his death on October 29,1984. Li had lengthy
  experience in party work, and specialized in managing iron and steel
  mills. See the obituary notice for Li in ZJRB, November 3,1984.
  李克昌（1923-1984）生于山东省。他于1940年加入中共，南下浙江工作，直到1984年10月29日去世。李长期从事党务工作，专长于管理钢铁厂。见《浙江日报》1984年11月3日刊登的李克昌讣告。}
at the time party secretary of the iron and steel mill, and probably as
the result of concessions wrung from the leaders of the ``mountain top''
faction at the Beijing conference, Zhang was released in early January
1975. Dissatisfied with the decision, the urban militia pursued Zhang
all over Hangzhou, but Li Kechang, through his contacts in the
bureaucracy arranged Zhang safe passage to Shanghai where he received
medical treatment.\footnote{HZRB, November 7,1978.
  《杭州日报》，1978年11月7日。}

另一个例子说明民兵如何被动员起来追捕``山上派''的反对者。本书上一章叙述了钢铁厂干部、``山下派''领导人张纪法的事迹。张在1974年9月被捕并遭受酷刑。在当时钢铁厂党委书记李克昌的帮助下，\footnote{Li
  Kechang (1923-84) was born in Shandong province. He joined the CCP in
  1940 and came south to Zhejiang where he worked until his death on
  October 29,1984. Li had lengthy experience in party work, and
  specialized in managing iron and steel mills. See the obituary notice
  for Li in ZJRB, November 3,1984.
  李克昌（1923-1984）生于山东省。他于1940年加入中共，南下浙江工作，直到1984年10月29日去世。李长期从事党务工作，专长于管理钢铁厂。见《浙江日报》1984年11月3日刊登的李克昌讣告。}
也可能是因为北京会议从``山上派''领导人那里迫使其作出让步，张于1975年1月初获释。城市民兵不满这一决定，在杭州到处追捕张，但李克昌通过其官僚体系中的关系安排张安全前往上海接受治疗。\footnote{HZRB,
  November 7,1978. 《杭州日报》，1978年11月7日。}

It is difficult to decide on the basis of the available evidence whether
Beijing merely ordered the dissolution of the militia command HQ or
whether it banned the organization itself. In addition, as the above
examples illustrate, how effective these measures were in practice is
debatable. Certainly Mao Zedong and Zhou Enlai differed in their
approaches as to what the future of the urban militia should be. It is
claimed that Mao ``decided in early 1975 that the militia should not get
involved in local strife between the two factions'' in places such as
Zhejiang and Yunnan.\footnote{It is reported that Deng visited Yunnan in
  February 1975 to ``deal with widespread unrest there''. See CNS, 577
  (August 6, 1975), p.~2.
  据报道，邓于1975年2月访问云南，以``处理那里广泛的动乱''。见CNS，577（1975年8月6日），第2页。}
By contrast, the Premier went much further in questioning the need for
militia commands at all. He believed that they only duplicated the work
of PLA military sub-districts and county-level people's armed forces
departments, which had previously been responsible for organizing and
training the militia.\footnote{HZRB, November 21,1976 (report of a
  meeting on the topic of Hangzhou's militia held on November 19,1976);
  Hangzhou Garrison criticism group, ``Chop off the evil hands of the
  gang of four who reached out for the people's armed forces'', HZRB,
  January 3,1977; HZRB, April 20,1977; Yang Yong,
  ``高举毛主席的伟大旗帜,大力加强民兵建设'' (Hold high the great banner
  of Chairman Mao, and greatly strengthen the building of the militia),
  RMRB, August 9,1978, p.~2.
  《杭州日报》，1976年11月21日（关于1976年11月19日召开杭州民兵问题会议的报道）；杭州警备区批判组：《斩断``四人帮''伸向人民武装的黑手》，《杭州日报》，1977年1月3日；《杭州日报》，1977年4月20日；杨勇：《高举毛主席的伟大旗帜，大力加强民兵建设》，《人民日报》，1978年8月9日，第2版。}

根据现有证据，很难判断北京究竟只是下令解散民兵指挥部总部，还是禁止了该组织本身。此外，正如上述例子所说明的，这些措施在实践中多么有效也值得商榷。毛泽东和周恩来对于城市民兵未来应当如何处理，显然采取了不同办法。据称，毛``在1975年初决定，在浙江、云南等地，民兵不应卷入两派之间的地方冲突''。\footnote{It
  is reported that Deng visited Yunnan in February 1975 to ``deal with
  widespread unrest there''. See CNS, 577 (August 6, 1975), p.~2.
  据报道，邓于1975年2月访问云南，以``处理那里广泛的动乱''。见CNS，577（1975年8月6日），第2页。}
相比之下，总理走得更远，他质疑民兵指挥部是否还有存在的必要。他认为它们只是重复了解放军军分区和县级人民武装部的工作，而后者此前一直负责组织和训练民兵。\footnote{HZRB,
  November 21,1976 (report of a meeting on the topic of Hangzhou's
  militia held on November 19,1976); Hangzhou Garrison criticism group,
  ``Chop off the evil hands of the gang of four who reached out for the
  people's armed forces'', HZRB, January 3,1977; HZRB, April 20,1977;
  Yang Yong, ``高举毛主席的伟大旗帜,大力加强民兵建设'' (Hold high the
  great banner of Chairman Mao, and greatly strengthen the building of
  the militia), RMRB, August 9,1978, p.~2.
  《杭州日报》，1976年11月21日（关于1976年11月19日召开杭州民兵问题会议的报道）；杭州警备区批判组：《斩断``四人帮''伸向人民武装的黑手》，《杭州日报》，1977年1月3日；《杭州日报》，1977年4月20日；杨勇：《高举毛主席的伟大旗帜，大力加强民兵建设》，《人民日报》，1978年8月9日，第2版。}

Sometime early in 1975, the Central Committee issued a document in
relation to Hangzhou's urban militia command which was approved by
Mao,\footnote{HZRB, September 8,1977. 《杭州日报》，1977年9月8日。} but
whether it went further at that stage than reiterating Mao's vague
comments, is uncertain. Two sources state that the Hangzhou urban
militia command was in existence for ten months, which would date its
disbandment in December 1974.\footnote{HZRB, May 20,1977, June 20,1977.
  《杭州日报》，1977年5月20日、1977年6月20日。} However, given the
Chinese press' notorious proclivity for vagueness, this seemingly
precise declaration does not exclude the possibility that the command
headquarters was wound up in January 1975. A Western observer has
claimed that the militia in Zhejiang was ``officially disbanded in March
1975''.\footnote{Dittmer, ``Bases of Power'', p.~53, fn. 63.
  Dittmer：``Bases of Power''，第53页，脚注63。} If the term militia is
used to refer to the HMMCH then the claim may possibly be correct.
Clearly, the urban militia as an organization remained in existence at
this time in cities such as Hangzhou and Wenzhou.

1975年初某时，中央委员会发布了一份关于杭州城市民兵指挥部的文件，并经毛批准，\footnote{HZRB,
  September 8,1977. 《杭州日报》，1977年9月8日。}
但它在当时是否超出了重申毛含糊意见的程度，则并不确定。有两个来源称，杭州城市民兵指挥部存在了十个月，这意味着其解散时间为1974年12月。\footnote{HZRB,
  May 20,1977, June 20,1977.
  《杭州日报》，1977年5月20日、1977年6月20日。}
然而，鉴于中国报刊向来倾向于含糊其辞，这一看似精确的说法并不排除指挥部总部是在1975年1月被撤销的可能。一位西方观察者声称，浙江民兵``于1975年3月正式解散''。\footnote{Dittmer,
  ``Bases of Power'', p.~53, fn. 63. Dittmer：``Bases of
  Power''，第53页，脚注63。}
如果``民兵''一词是指杭州市民兵指挥部，那么这一说法或许是正确的。显然，此时在杭州、温州等城市，城市民兵作为一个组织仍然存在。

The central directive ordered militia commands throughout Zhejiang to
hand in their weapons as the first step in curbing the violence. One
place where trouble continued, however, was the central Zhejiang town of
Jinhua. Jinhua is an important railway town on the line from Shanghai to
Guangzhou via Zhejiang, Jiangxi and Hunan provinces. Disruptions at its
freight-yards would have serious repercussions for the movement of goods
in eastern and southern China. Immediately after his return to Zhejiang
from the Beijing conference, Zhang Yongsheng set off for Jinhua at the
end of 1974. The previous March, he had helped set up the militia HQ
there. Now, Zhang suggested that the local militia only turn over to the
authorities those weapons which were defective or for which they had no
ammunition.\footnote{Zhang Yongshengde zuizheng cailiao, p.~103.
  《张永生的罪证材料》，第103页。} As a result of Zhang's intervention,
fighting persisted in the vital transport center for a further three
months, paralyzing the local administration, bringing industrial
production to a halt, closing shops, holding up traffic, suspending
postal and communication services and causing injury and
death.\footnote{HZRB, August 14,1978. 《杭州日报》，1978年8月14日。} It
is little wonder that the People's Daily and the provincial media later
expressed concern about the flow of railway freight through Jinhua in
the first quarter of 1975.

中央指示命令浙江全省民兵指挥部交出武器，作为遏制暴力的第一步。然而，仍持续出现麻烦的一个地方是浙江中部城镇金华。金华是从上海经浙江、江西、湖南通往广州线路上的重要铁路城镇。其货场受到干扰，将对华东和华南货物流通产生严重影响。1974年底，张永生从北京会议返回浙江后，立即动身前往金华。此前3月，他曾帮助在那里建立民兵总部。现在，张建议当地民兵只向当局交出那些有缺陷或没有弹药的武器。\footnote{Zhang
  Yongshengde zuizheng cailiao, p.~103. 《张永生的罪证材料》，第103页。}
由于张的干预，这一重要交通中心的斗争又持续了三个月，造成地方行政瘫痪、工业生产停顿、商店关闭、交通受阻、邮政和通信服务中断，并导致伤亡。\footnote{HZRB,
  August 14,1978. 《杭州日报》，1978年8月14日。}
难怪《人民日报》和省级媒体后来对1975年第一季度通过金华的铁路货运流量表示关切。

Between 25 February and 8 March the center convened a conference of
party secretaries responsible for industry to discuss problems in the
railway system.\footnote{Liu Suinian and Wu Qungan,``Wenhua dageming''
  shiqide guomin jingji, p.~170; Hu Hua (ed.), Zhongguo shehuizhuyi
  geming, p.~322.
  刘岁年、吴群敢：《``文化大革命''时期的国民经济》，第170页；胡华编：《中国社会主义革命》，第322页。}
It is highly likely that the CC decision on strengthening work on the
railways, issued on 5 March 1975 as Document No.~9 (1975),\footnote{See
  ``Document of the CCP CC'', zhongfa (1977), No.~37, I\&S, 14:11
  (1978), p.~104; Hao Mengbi and Duan Haoran, Liushinian, p.~642.
  见``Document of the CCP
  CC''，中发（1977）37号，I\&S，14:11（1978年），第104页；郝梦笔、段浩然：《六十年》，第642页。}
was partly a response to the disruption created in Jinhua by Zhang
Yongsheng and the town's militia. Another trouble spot was Xuzhou, a
major railway junction in the north-west of Jiangsu province. Minister
of Railways Wan Li inspected the city immediately following the central
conference.\footnote{See Jiangsusheng dashiji, pp.~332-33.
  见《江苏省大事记》，第332-333页。} Deng Xiaoping, in his speech to the
conference of 5 March 1975, devoted much of his attention to
factionalism in the railway system.\footnote{Deng Xiaoping,
  ``全党讲大局,把国民经济搞上去'' (The whole party must look at the
  overall picture and push the national economy forward), Deng Xiaoping
  Wenxuan, pp.~4-7.
  邓小平：《全党讲大局，把国民经济搞上去》，《邓小平文选》，第4-7页。}
He deplored the alarmingly high rate of accidents experienced on the
railways in the previous year, a more than eight-fold increase in the
number which had occurred in 1964. Deng attributed the increase to
disregard for rules and regulations and warned that ``bad men'' who
``fished in troubled waters'' in order to make a name for themselves
(升官发财) would be severely punished.\footnote{See also CNS, 577, p.~5.
  By April 1975, freight figures had improved. See Hao Mengbi and Duan
  Haoran, Liushinian, p.~643; Gao Gao and Yan Jiaqi, ``Wenhua
  dageming'', p.~541.
  另见CNS，577，第5页。到1975年4月，货运数字已有改善。见郝梦笔、段浩然：《六十年》，第643页；高皋、严家其：《文化大革命》，第541页。}
The center issued another document on the railways in May and on 2 June
the CC relayed a report from the Jiangsu authorities on its handling of
factionalism at Xuzhou.\footnote{CNS, 577, p.~5; Jiangsusheng dashiji,
  p.~333. CNS，577，第5页；《江苏省大事记》，第333页。}

2月25日至3月8日，中央召开负责工业的党委书记会议，讨论铁路系统问题。\footnote{Liu
  Suinian and Wu Qungan,``Wenhua dageming'' shiqide guomin jingji,
  p.~170; Hu Hua (ed.), Zhongguo shehuizhuyi geming, p.~322.
  刘岁年、吴群敢：《``文化大革命''时期的国民经济》，第170页；胡华编：《中国社会主义革命》，第322页。}
1975年3月5日作为1975年第9号文件发布的中央关于加强铁路工作的决定，\footnote{See
  ``Document of the CCP CC'', zhongfa (1977), No.~37, I\&S, 14:11
  (1978), p.~104; Hao Mengbi and Duan Haoran, Liushinian, p.~642.
  见``Document of the CCP
  CC''，中发（1977）37号，I\&S，14:11（1978年），第104页；郝梦笔、段浩然：《六十年》，第642页。}
很可能部分是对张永生及金华民兵造成的破坏作出的回应。另一个问题点是徐州，它是江苏省西北部的重要铁路枢纽。中央会议结束后，铁道部长万里立即视察了该市。\footnote{See
  Jiangsusheng dashiji, pp.~332-33. 见《江苏省大事记》，第332-333页。}
邓小平在1975年3月5日对会议发表的讲话中，把相当多注意力放在铁路系统的派性问题上。\footnote{Deng
  Xiaoping, ``全党讲大局,把国民经济搞上去'' (The whole party must look
  at the overall picture and push the national economy forward), Deng
  Xiaoping Wenxuan, pp.~4-7.
  邓小平：《全党讲大局，把国民经济搞上去》，《邓小平文选》，第4-7页。}
他对前一年铁路事故率惊人地高表示痛心，事故数量比1964年增加了八倍以上。邓把这一增长归因于无视规章制度，并警告那些为出名（升官发财）而``浑水摸鱼''的``坏人''将受到严厉惩处。\footnote{See
  also CNS, 577, p.~5. By April 1975, freight figures had improved. See
  Hao Mengbi and Duan Haoran, Liushinian, p.~643; Gao Gao and Yan Jiaqi,
  ``Wenhua dageming'', p.~541.
  另见CNS，577，第5页。到1975年4月，货运数字已有改善。见郝梦笔、段浩然：《六十年》，第643页；高皋、严家其：《文化大革命》，第541页。}
中央5月又发布一份关于铁路的文件，6月2日，中央转发了江苏当局关于处理徐州派性问题的报告。\footnote{CNS,
  577, p.~5; Jiangsusheng dashiji, p.~333.
  CNS，577，第5页；《江苏省大事记》，第333页。}

In late May or early June the Zhejiang authorities responded somewhat
belatedly to Document 9 and Deng Xiaoping's speech by convening a
railway-commune joint defence meeting. It had taken three months since
the central conference of March for the local authorities to restore
order in the Zhejiang railway system. The meeting claimed an increase in
the number of freight cars loaded and unloaded and a speed-up in
turn-around time since the beginning of the second quarter of 1975. It
noted particularly the situation in Jiangshan (江山) county, Jinhua
district, bordering Jiangxi province in the south-west of Zhejiang.
Members of the PLA and public security forces were patrolling the tracks
around the clock to prevent sabotage by ``class enemies''. Peasants
working alongside the railway lines remained vigilant as they farmed.
The meeting directed all localities to ``make sure that railway
transportation is smooth and free of obstructions, reaching all
destinations safe and on time''.\footnote{ZPS, June 4,1975,
  SWB/FE/4923/BII/9-10.
  浙江省广播服务，1975年6月4日，SWB/FE/4923/BII/9-10。}

5月下旬或6月初，浙江当局以召开铁路、公社联防会议的方式，对第9号文件和邓小平讲话作出有些迟缓的回应。自3月中央会议以来，地方当局花了三个月才恢复浙江铁路系统的秩序。会议声称，自1975年第二季度开始以来，装卸货车数量增加，周转时间加快。会议特别提到与江西省接壤、位于浙江西南部金华地区的江山县情况。解放军和公安力量昼夜巡逻线路，以防``阶级敌人''破坏。在铁路沿线劳动的农民边耕作边保持警惕。会议指示各地``保证铁路运输畅通无阻，安全正点到达各目的地''。\footnote{ZPS,
  June 4,1975, SWB/FE/4923/BII/9-10.
  浙江省广播服务，1975年6月4日，SWB/FE/4923/BII/9-10。}

A Xinhua report of 9 August 1975 examined the situation in the Jinhua
railway sub-bureau.\footnote{RMRB, 10 August 1975, p.~1.
  《人民日报》，1975年8月10日，第1版。} The article revealed that in the
second quarter of 1975 incoming and outgoing freight had doubled in
quantity compared to the amount which had been transported in the first
quarter. The reason for this startling increase, made deceptively
impressive by the depths to which figures had previously plummeted, was
ascribed to the ``revolutionary unity'' of leading groups at various
levels, a clear indication of rampant factionalism in the sub-bureau.
According to the report, veteran cadres and their youthful colleagues
promoted in the Cultural Revolution were now working side by side in
perfect harmony. The party organizations in the sub-bureau made
arrangements for cadres to go to the front line to labor and show
concern for the workers' livelihood. Cadres also visited workers' homes
and inspected facilities at work-canteens. The message spelt out very
clearly in the report was that production could not run smoothly when
workers or cadres became embroiled in arguments and resorted to mutual
recrimination and attack.

1975年8月9日的一篇新华社报道考察了金华铁路分局的情况。\footnote{RMRB, 10
  August 1975, p.~1. 《人民日报》，1975年8月10日，第1版。}
文章透露，1975年第二季度进出货运量比第一季度翻了一番。这个惊人增长因此前数字跌到极低水平而显得具有迷惑性的可观，其原因被归结为各级领导班子的``革命团结''，这清楚表明分局内派性猖獗。报道说，老干部和文化大革命中被提拔的年轻同事如今正并肩和谐工作。分局党组织安排干部到第一线劳动，并关心工人生活。干部还走访工人家庭，检查食堂设施。报道非常明确地传达出一个信息：当工人或干部卷入争吵并相互指责、攻击时，生产就不可能顺利运行。

However, the issue of the urban militia commands had not found
resolution by the time the CC MAC held an important enlarged meeting
from 24 June to 15 July 1975. Deng and Ye Jianying, both Vice-chairmen
of the commission, delivered reports to the lengthy
conference.\footnote{Zhonggong dangshi baiti,p.~387; Hu Hua (ed.),
  Zhongguo shehuizhuyi geming, pp.~323-4; Gao Gao and Yan Jiaqi,
  ``Wenhua dageming'', p.~547. For a brief synopsis of Ye's speech of 15
  July, see Hao Mengbi and Duan Haoran, Liushinian, p.~642.
  《中共党史百题》，第387页；胡华编：《中国社会主义革命》，第323-324页；高皋、严家其：《文化大革命》，第547页。关于叶7月15日讲话的简要概述，见郝梦笔、段浩然：《六十年》，第642页。}
Deng's speech to the commission on the penultimate day contained
withering blasts against factionalism, poor discipline and
insubordination within the PLA. He proposed the transfer of cadres who
had become entrenched in their localities and had been involved in local
factional conflicts.\footnote{Deng Xiaoping, ``军队整顿的任务'' (Tasks
  in rectification of the military), Deng Xiaoping Wenxuan, pp.~15-24.
  邓小平：《军队整顿的任务》，《邓小平文选》，第15-24页。} As a result,
it is reported that in the final four months of 1975 thirty-five
provincial military district commanders, political commissars and
deputy-commanders and deputy-commissars were transferred.\footnote{See
  CNS, 598 (January 14,1976). 见CNS，598（1976年1月14日）。} According
to one analyst, ``those military officials who rose to their posts
because of their cooperation with Jiang Qing were replaced''.\footnote{Hu
  Hua, ``From the Fall of Lin Biao'', p.~3. Deng's speech was later
  attacked in articles which appeared in RMRB and Hongqi in August 1976.
  See CS, 14:10 (1976), pp.~18-19. 胡华：``From the Fall of Lin
  Biao''，第3页。邓的讲话后来在1976年8月《人民日报》和《红旗》刊登的文章中受到攻击。见CS，14:10（1976年），第18-19页。}
Deng also implicitly supported Zhou Enlai's position regarding the urban
militia commands by instructing that these ``should be settled once and
for all, militia details be abolished and the militia should remain
under the control of the military regions, provincial military districts
and subdistricts, and the people's armed forces
{[}departments{]}\ldots{}''.\footnote{Yang Yong, RMRB, August 9,1978. In
  November 1975 the CC MAC issued two documents, one a code of military
  discipline and the other a code of internal management. I\&S, 12:10
  (1976), pp.~88-97, 98-127.
  杨勇，《人民日报》，1978年8月9日。1975年11月，中共中央军委发布两份文件，一份是军队纪律条例，另一份是内务管理条例。I\&S，12:10（1976年），第88-97、98-127页。}
Deng is also reported to have stated bluntly that ``It is not the role
of this militia to grasp class struggle''.\footnote{CNS, 620 (June
  30,1976). CNS，620（1976年6月30日）。}

然而，到1975年6月24日至7月15日中共中央军委召开一次重要扩大会议时，城市民兵指挥部问题仍未得到解决。军委副主席邓和叶剑英都在这次长会中作了报告。\footnote{Zhonggong
  dangshi baiti,p.~387; Hu Hua (ed.), Zhongguo shehuizhuyi geming,
  pp.~323-4; Gao Gao and Yan Jiaqi, ``Wenhua dageming'', p.~547. For a
  brief synopsis of Ye's speech of 15 July, see Hao Mengbi and Duan
  Haoran, Liushinian, p.~642.
  《中共党史百题》，第387页；胡华编：《中国社会主义革命》，第323-324页；高皋、严家其：《文化大革命》，第547页。关于叶7月15日讲话的简要概述，见郝梦笔、段浩然：《六十年》，第642页。}
邓在会议倒数第二天对军委发表的讲话，对解放军内部的派性、纪律松弛和不服从进行了严厉抨击。他提议调动那些在当地盘根错节、并卷入地方派系冲突的干部。\footnote{Deng
  Xiaoping, ``军队整顿的任务'' (Tasks in rectification of the military),
  Deng Xiaoping Wenxuan, pp.~15-24.
  邓小平：《军队整顿的任务》，《邓小平文选》，第15-24页。}
据报道，其结果是1975年最后四个月中，三十五名省军区司令员、政委、副司令员和副政委被调动。\footnote{See
  CNS, 598 (January 14,1976). 见CNS，598（1976年1月14日）。}
一位分析者说，``那些因同江青合作而升任职务的军队干部被替换''。\footnote{Hu
  Hua, ``From the Fall of Lin Biao'', p.~3. Deng's speech was later
  attacked in articles which appeared in RMRB and Hongqi in August 1976.
  See CS, 14:10 (1976), pp.~18-19. 胡华：``From the Fall of Lin
  Biao''，第3页。邓的讲话后来在1976年8月《人民日报》和《红旗》刊登的文章中受到攻击。见CS，14:10（1976年），第18-19页。}
邓还通过指示城市民兵指挥部``要一劳永逸地解决，取消民兵专设机构，民兵仍应由军区、省军区、军分区和人民武装部掌握\ldots\ldots''来含蓄支持周恩来关于城市民兵指挥部的立场。\footnote{Yang
  Yong, RMRB, August 9,1978. In November 1975 the CC MAC issued two
  documents, one a code of military discipline and the other a code of
  internal management. I\&S, 12:10 (1976), pp.~88-97, 98-127.
  杨勇，《人民日报》，1978年8月9日。1975年11月，中共中央军委发布两份文件，一份是军队纪律条例，另一份是内务管理条例。I\&S，12:10（1976年），第88-97、98-127页。}
据报道，邓还直截了当地说，``这个民兵不是抓阶级斗争的''。\footnote{CNS,
  620 (June 30,1976). CNS，620（1976年6月30日）。}

Deng's reference to entrenched factionalism within the ranks of the PLA
had particular relevance for Zhejiang. At the traditional annual meeting
held at the beginning of the year to pontificate about the supposedly
excellent relations between the army and local citizens, the provincial
military leadership criticized its own performance in 1974. Political
Commissar Li Bincheng pledged that in 1975 ``We will strive to play a
more effective role in maintaining social order and in implementing
government decrees''. Li's words were a form of public apology for the
army's failure to contain the excesses of the militia during 1974. Li
also stated that the troops stationed in the province should strive to
overcome ``shortcomings and mistakes'' in their work, a reference to the
``three supports and two militaries'' units still stationed in
government, industrial and other work units. On behalf of the CCP ZPC,
Deputy-secretary Lai Keke reciprocated by promising that the party
authorities would do a better job of supporting the army than in
1974.\footnote{ZPS, January 4,1975, SWB/FE/4799/BII/1-3.
  浙江省广播服务，1975年1月4日，SWB/FE/4799/BII/1-3。}

邓提到解放军队伍中根深蒂固的派性，对浙江尤其相关。在年初举行的传统年度会议上，通常会宣讲所谓军民关系良好，但省军区领导却批评了自己在1974年的表现。政委李彬成保证，1975年``我们要努力在维护社会秩序和贯彻政府法令方面发挥更有效的作用''。李的话某种程度上是公开道歉，因为军队未能遏制1974年民兵的过火行为。李还说，驻省部队应努力克服工作中的``缺点和错误''，这是指仍驻在政府、工业和其他工作单位中的``三支两军''单位。代表中共浙江省委，副书记赖可可也作出回应，承诺党委机关将在拥军方面比1974年做得更好。\footnote{ZPS,
  January 4,1975, SWB/FE/4799/BII/1-3.
  浙江省广播服务，1975年1月4日，SWB/FE/4799/BII/1-3。}

These speeches illustrated the tensions which had appeared within the
political-military elite in Zhejiang during 1974. The failure of local
troops to maintain social order and implement government decrees in the
province may have led to the transfer of outside troops to Zhejiang at
about this time. The CC MAC conference of June/July may have decided on
this initiative or action may have been taken earlier in the year. A
report of January 1975 specifically mentioned the presence in Zhejiang
of a battalion (洛阳营) from Henan province. Whether the battalion
actually came from Luoyang is unclear.\footnote{ZJRB, January 10,1975,
  p.~1. Taiwan observers seemed confused as to the prior location of
  this Army, placing it variously in Xuchang, Zhejiang, ``Tan Qilong'',
  p.103, or Wuhan, Tan Qilong (2)``), p.110. Zhonggong nianbao correctly
  located the division in Xuchang, Henan province, under the Wuhan
  Military Region. Zhonggong Nianbao 1976, 5:103.
  《浙江日报》，1975年1月10日，第1版。台湾观察者似乎不清楚这支部队此前驻地，有的把它放在许昌，见浙江，《谭启龙》，第103页；有的放在武汉，见《谭启龙（二）》，第110页。《中共年报》正确地把该师定位在河南省许昌，隶属武汉军区。《中共年报1976》，5:103。}
It was perhaps this report and other intelligence accounts that led to
the claim that ``public security squads \ldots{} drafted in from
northern China'' appeared on the streets of Hangzhou in July
1975.\footnote{Leo Goodstadt, ``Controlling social tension'', FEER,
  August 1,1975, p.~30. Leo Goodstadt：``Controlling social
  tension''，FEER，1975年8月1日，第30页。} In January 1975, the
political department of the ZPMD had issued a circular criticizing such
``reactionary fallacies'' as ``the gun should command the Party'' and
``the Army is something special''.\footnote{ZJRB, January 13,1975, p.~1;
  quotations from ZPS, January 13,1975, SWB/FE/4815/BII/14.
  《浙江日报》，1975年1月13日，第1版；引文出自浙江省广播服务，1975年1月13日，SWB/FE/4815/BII/14。}
It seems likely that these words alluded to past disobedience of party
directives by the local military.

这些讲话显示出1974年浙江政治军事精英内部出现的紧张关系。地方部队未能在该省维护社会秩序和执行政府法令，可能导致大约在此时外地部队被调入浙江。6、7月的中央军委会议可能决定了这一举措，也可能有关行动在当年早些时候已经采取。1975年1月的一篇报道特别提到浙江有一个来自河南省的营（洛阳营）。该营是否真的来自洛阳并不清楚。\footnote{ZJRB,
  January 10,1975, p.~1. Taiwan observers seemed confused as to the
  prior location of this Army, placing it variously in Xuchang,
  Zhejiang, ``Tan Qilong'', p.103, or Wuhan, Tan Qilong (2)``), p.110.
  Zhonggong nianbao correctly located the division in Xuchang, Henan
  province, under the Wuhan Military Region. Zhonggong Nianbao 1976,
  5:103.
  《浙江日报》，1975年1月10日，第1版。台湾观察者似乎不清楚这支部队此前驻地，有的把它放在许昌，见浙江，《谭启龙》，第103页；有的放在武汉，见《谭启龙（二）》，第110页。《中共年报》正确地把该师定位在河南省许昌，隶属武汉军区。《中共年报1976》，5:103。}
或许正是这篇报道和其他情报，使人声称1975年7月杭州街头出现了``从华北调入的公安队伍''。\footnote{Leo
  Goodstadt, ``Controlling social tension'', FEER, August 1,1975, p.~30.
  Leo Goodstadt：``Controlling social
  tension''，FEER，1975年8月1日，第30页。}
1975年1月，浙江省军区政治部曾发布通报，批评``枪指挥党''\,``军队特殊''等``反动谬论''。\footnote{ZJRB,
  January 13,1975, p.~1; quotations from ZPS, January 13,1975,
  SWB/FE/4815/BII/14.
  《浙江日报》，1975年1月13日，第1版；引文出自浙江省广播服务，1975年1月13日，SWB/FE/4815/BII/14。}
这些话很可能暗指地方军队过去不服从党的指示。

Thus, in late 1974 the center had attempted, with mixed results, to deal
with the specific manifestations of factionalism, the ``two
crashthroughs'' and the urban militia in Zhejiang. However, the efficacy
of the central decisions to solve the ``Zhejiang problem'' clearly
depended on the willingness of those involved to adhere to their terms.
Clearly, compliance had been far from forthcoming. After a central
instruction forbidding armed struggle was relayed, Weng Senhe revealed
the contempt in which he held directives not to his liking by remarking:
``If we took any notice of all the rules in the world we'd perish; the
same with their seven or eight points. I'll take responsibility''.
(听了佛法要饿煞,听了王法要打煞,管他七条八条,我负责).\footnote{HZRB, May
  20,1977; Hua Meisheng, HZRB, June 22,1977. Another source, in quoting
  Weng's words, added after 打煞 the following phrase - ``藤帽铁棍一条''
  (cane helmets and iron clubs are my law). See HZRB, April 20, 1977.
  《杭州日报》，1977年5月20日；华美生，《杭州日报》，1977年6月22日。另一来源引用翁的话时，在``打煞''之后加上了``藤帽铁棍一条''（藤帽铁棍就是我的法）这句话。见《杭州日报》，1977年4月20日。}
Wang Hongwen's sympathetic attitude to the ``mountain top'' faction at
the central conference had probably prompted Zhang Yongsheng to defy the
decision instructing the militia to hand in its weapons. At the end of
January 1975, when the ZPC held a work meeting to convey the decisions
of the CCP CC plenum and the 4th NPC, Zhang accused Tan Qilong of doing
a double-flip (翻烧饼), a remark suggesting that Zhang and his faction
had previously viewed Tan as neutral in the factional battles. Zhang
urged his allies to disrupt the conference and took the floor in an
attempt to continue the proceedings by putting Tan in the
dock.\footnote{Zhang Yongshengde zuizheng cailiao, p.~63.
  《张永生的罪证材料》，第63页。}

因此，1974年末，中央曾试图处理浙江派性的一些具体表现，即``双突''和城市民兵，结果成败参半。然而，解决``浙江问题''的中央决定是否有效，显然取决于相关人员是否愿意遵守其条款。显然，服从远未到来。中央禁止武斗的指示传达后，翁森鹤以如下言论显示了他对不合己意指示的蔑视：``听了佛法要饿煞，听了王法要打煞，管他七条八条，我负责''。\footnote{HZRB,
  May 20,1977; Hua Meisheng, HZRB, June 22,1977. Another source, in
  quoting Weng's words, added after 打煞 the following phrase -
  ``藤帽铁棍一条'' (cane helmets and iron clubs are my law). See HZRB,
  April 20, 1977.
  《杭州日报》，1977年5月20日；华美生，《杭州日报》，1977年6月22日。另一来源引用翁的话时，在``打煞''之后加上了``藤帽铁棍一条''（藤帽铁棍就是我的法）这句话。见《杭州日报》，1977年4月20日。}
王洪文在中央会议上对``山上派''的同情态度，可能促使张永生公然违抗要求民兵交出武器的决定。1975年1月底，浙江省委召开工作会议传达中共中央全会和第四届全国人大决定时，张指责谭启龙``翻烧饼''，这一说法表明张及其派系此前曾把谭视为派系斗争中的中立者。张敦促盟友扰乱会议，并发言试图通过把谭置于被告席来继续会议进程。\footnote{Zhang
  Yongshengde zuizheng cailiao, p.~63. 《张永生的罪证材料》，第63页。}

In December 1974 Beijing had summoned the factional adversaries to
Beijing in order to find a lasting solution to the problems besetting
the province. These problems had originated in the rebellion and
factionalism accompanying the anti-Lin anti-Confucius campaign, but had
their roots in the upheavals of the mid 1960s. The immediate aftermath
of the conference suggested that the center had failed to grasp the
nettle and that the local factionalists would be unlikely to comply with
its directives which were indecisive and unenforceable. Six months later
the center was forced to recognize its failure by despatching a
top-level team to Hangzhou, this time to make on-the-spot investigations
and then, with the approval of the Politburo, to act decisively to quell
the trouble.

1974年12月，北京曾召集派系对手赴京，以便为困扰该省的问题找到持久解决办法。这些问题起源于批林批孔运动伴随的造反和派性，但其根源在20世纪60年代中期的动荡。会议结束后的直接后果表明，中央未能果断处理棘手问题，地方派系分子也不太可能遵守那些犹豫不决且无法执行的指示。六个月后，中央被迫承认失败，派出一个高层小组到杭州，这一次是进行实地调查，然后在政治局批准下，采取果断行动平息骚乱。

\section{Factionalism and Policy Direction in Zhejiang, February-June
1975}\label{factionalism-and-policy-direction-in-zhejiang-february-june-1975}

\section{浙江的派性与政策方向，1975年2月-6月}\label{ux6d59ux6c5fux7684ux6d3eux6027ux4e0eux653fux7b56ux65b9ux54111975ux5e742ux6708-6ux6708}

The CCP ZPC did not respond to the 9 February People's Daily editorial
announcing the campaign to study the theory of the dictatorship of the
proletariat until 25 February.\footnote{For reports on meetings in other
  provinces to discuss theory and promote production, see CNS, 562
  (April 16, 1975).
  关于其他省份召开会议讨论理论和促进生产的报道，见CNS，562（1975年4月16日）。}
On that date it held a conference on the campaign at which Tan Qilong
delivered a report.\footnote{ZJRB, March 1,1975, pp.~1,4. The CCP HMC
  standing committee had discussed the editorial in mid-February. ZJRB,
  February 14,1975, p.~1.
  《浙江日报》，1975年3月1日，第1、4版。中共杭州市委常委会已于2月中旬讨论过这篇社论。《浙江日报》，1975年2月14日，第1版。}
From the tenor of his remarks it is obvious that Tan believed that the
campaign was to be merely a routine exercise which carried no practical
policy consequences. He cited the reforms of the Cultural Revolution as
examples of restricting bourgeois rights and then finished his speech
with a ringing call to fulfil the modernization goals announced at the
4th NPC. Under the guidance of the party center. Tan led the ZPC to
focus attention on economic matters. Even Zhang Yongsheng drafted a
proposal (倡议书) early in 1975 to ``grasp revolution and promote
production''.\footnote{See Wang Jinyou, Juqi zhuagang xue Dazhai, p.~39.
  见王金友：《举旗抓纲学大寨》，第39页。}

中共浙江省委直到2月25日才对2月9日《人民日报》宣布开展学习无产阶级专政理论运动的社论作出回应。\footnote{For
  reports on meetings in other provinces to discuss theory and promote
  production, see CNS, 562 (April 16, 1975).
  关于其他省份召开会议讨论理论和促进生产的报道，见CNS，562（1975年4月16日）。}
当天，省委召开关于这场运动的会议，谭启龙作报告。\footnote{ZJRB, March
  1,1975, pp.~1,4. The CCP HMC standing committee had discussed the
  editorial in mid-February. ZJRB, February 14,1975, p.~1.
  《浙江日报》，1975年3月1日，第1、4版。中共杭州市委常委会已于2月中旬讨论过这篇社论。《浙江日报》，1975年2月14日，第1版。}
从其讲话语气看，很明显谭认为这场运动不过是一项例行活动，不会产生实际政策后果。他把文化大革命的改革列为限制资产阶级法权的例子，最后以响亮的号召结束讲话，要求实现第四届全国人大宣布的现代化目标。在党中央指导下，谭带领浙江省委把注意力集中于经济事务。甚至张永生也在1975年初起草了一份``抓革命、促生产''的倡议书。\footnote{See
  Wang Jinyou, Juqi zhuagang xue Dazhai, p.~39.
  见王金友：《举旗抓纲学大寨》，第39页。}

At the end of March a provincial conference on industry was
called.\footnote{ZJRB, March 26,1975, p.~1. For reports on such
  conferences in other provinces, see CNS, 564 (April 30, 1975).
  《浙江日报》，1975年3月26日，第1版。关于其他省份此类会议的报道，见CNS，564（1975年4月30日）。}
In the reports from the conference the study campaign merited only
passing mention and terms expressing the new direction on the
ideological front made their appearance. The conference pointed out that

3月底，召开了一次全省工业会议。\footnote{ZJRB, March 26,1975, p.~1. For
  reports on such conferences in other provinces, see CNS, 564 (April
  30, 1975).
  《浙江日报》，1975年3月26日，第1版。关于其他省份此类会议的报道，见CNS，564（1975年4月30日）。}
在会议报道中，学习运动仅被顺带提及，而表达意识形态战线新方向的术语开始出现。会议指出，

\begin{quote}
we must continue to eliminate the pernicious influence of Lin Biao's
splittist line, strengthen revolutionary unity, oppose bourgeois
factionalism, strengthen proletarian Party spirit and persist in the
principle of grasping revolution and promoting production.
\end{quote}

\begin{quote}
必须继续肃清林彪分裂主义路线的流毒，加强革命团结，反对资产阶级派性，增强无产阶级党性，坚持抓革命、促生产的方针。
\end{quote}

It also called for a strengthening of Party leadership and put forward
the view that ``It is necessary to step up efforts to educate Party
members and the masses in Party concepts (党的观念) and in the sense of
organizational discipline \ldots{}''. The introduction of the term
``bourgeois factionalism'' signified Deng Xiaoping's counter-attack on
the ideological battlefield.

会议还要求加强党的领导，并提出这样的观点：``必须加强对党员和群众进行党的观念和组织纪律观念的教育\ldots\ldots''
``资产阶级派性''这一术语的引入，标志着邓小平在意识形态战场上的反击。

A provincial ``learn from Dazhai'' conference held in March focused
attention on the agricultural sector of the economy.\footnote{The CCP
  ZPC issued a notice on March 14 stating that the meeting would take
  place in the last ten days of the month. ZJRB, March 14,1975, p.~1.
  The date was put back most probably to accommodate Vice-Premier Chen
  Yonggui's inspection of the province, which commenced on April 12.
  中共浙江省委于3月14日发布通知，称会议将在本月下旬召开。《浙江日报》，1975年3月14日，第1版。会议日期被推迟，很可能是为了配合副总理陈永贵从4月12日开始对该省的视察。}
The remarkable feature of the speeches and editorials published during
the seventeen-day conference was the dramatic shift in emphasis which
occurred after the address delivered by Chen Yonggui on 15 April. In his
opening address on 1 April, the day of the publication of Zhang
Chunqiao's article ``On Exercising All-round Dictatorship over the
Bourgeoisie'' referred to above, Tan Qilong had delivered some strong
warnings to the audience. Since his casual reference to it at the
beginning of the previous month Tan had clearly been made aware of the
importance of the campaign. He used the opportunity provided by the
conference to pay homage to the theoretical acumen of the central
radicals. Tan repeated Mao's 1959 dictum, quoted by Yao Wenyuan in his
article ``On the Social Basis of the Lin Biao Anti-party Clique'', that
empiricism presented the main danger in the ideological realm. Tan also
quoted sentences from Zhang Chunqiao's tract bluntly admonishing those
who wanted to get off the revolutionary bus before it arrived at its
final destination. Zhang had written about the revolution that, ``It's
dangerous to stop half-way! The bourgeoisie is beckoning to you. Catch
up with the ranks and continue to advance!''\footnote{Zhang Chunqiao, On
  Exercising All-Round Dictatorship, pp.~18-19. For Tan's speech see,
  ZJRB, April 3,1975, pp.~1, 3. The broadcast version did not include
  Tan's quotation from Zhang's article. 张春桥：On Exercising All-Round
  Dictatorship，第18-19页。关于谭的讲话，见《浙江日报》，1975年4月3日，第1、3版。广播版本没有包括谭对张文章的引用。}

3月举行的全省``学大寨''会议把注意力集中到经济的农业部门。\footnote{The
  CCP ZPC issued a notice on March 14 stating that the meeting would
  take place in the last ten days of the month. ZJRB, March 14,1975,
  p.~1. The date was put back most probably to accommodate Vice-Premier
  Chen Yonggui's inspection of the province, which commenced on April
  12.
  中共浙江省委于3月14日发布通知，称会议将在本月下旬召开。《浙江日报》，1975年3月14日，第1版。会议日期被推迟，很可能是为了配合副总理陈永贵从4月12日开始对该省的视察。}
在这场为期十七天的会议期间发表的讲话和社论中，最显著的特点是4月15日陈永贵讲话之后重点发生了巨大转移。在4月1日的开幕讲话中，也就是上文提到的张春桥《论对资产阶级的全面专政》发表当天，谭启龙向听众发出了一些严厉警告。自上月初随意提及这场运动以来，谭显然已经被告知其重要性。他利用会议提供的机会，向中央激进派的理论敏锐性表示敬意。谭重复了毛1959年的论断，即姚文元在《论林彪反党集团的社会基础》一文中引用的那句话：经验主义是思想领域的主要危险。谭还引用了张春桥那篇小册子中的句子，直截了当地告诫那些想在革命车辆抵达终点之前下车的人。张在谈到革命时写道：``半途而废是危险的！资产阶级正在向你招手。赶上队伍，继续前进！''\footnote{Zhang
  Chunqiao, On Exercising All-Round Dictatorship, pp.~18-19. For Tan's
  speech see, ZJRB, April 3,1975, pp.~1, 3. The broadcast version did
  not include Tan's quotation from Zhang's article. 张春桥：On
  Exercising All-Round
  Dictatorship，第18-19页。关于谭的讲话，见《浙江日报》，1975年4月3日，第1、3版。广播版本没有包括谭对张文章的引用。}

As a follow-up to Tan's opening speech, Zhejiang Daily published two
editorials which continued in radical theme adopted by the First
Secretary. The first, dated 3 April, described the ``learn from Dazhai
movement'' as a''profound socialist revolution'' and claimed that
revisionists constantly attempted to turn it into a production movement.
It cautioned leading cadres not to forget theoretical issues while busy
with routine tasks. The second editorial reiterated Tan's warning, taken
from Yao Wenyuan's article, about the dangers of empiricism.\footnote{ZJRB,
  April 3,1975, p.~1; ZPS, April 9, 1975, FBIS/CHI, 70 (1975), G 1.
  《浙江日报》，1975年4月3日，第1版；浙江省广播服务，1975年4月9日，FBIS/CHI，70（1975年），G1。}

作为谭开幕讲话的后续，《浙江日报》发表了两篇社论，延续第一书记所采用的激进主题。第一篇发表于4月3日，把``学大寨运动''描述为一场``深刻的社会主义革命''，并声称修正主义者不断试图把它变成生产运动。社论告诫领导干部不要在忙于日常事务时忘记理论问题。第二篇社论重申了谭从姚文元文章中借来的警告，即经验主义的危险。\footnote{ZJRB,
  April 3,1975, p.~1; ZPS, April 9, 1975, FBIS/CHI, 70 (1975), G 1.
  《浙江日报》，1975年4月3日，第1版；浙江省广播服务，1975年4月9日，FBIS/CHI，70（1975年），G1。}

Deputy-secretary of the CCP ZPC Chen Weida delivered the main report to
the conference on 6 April.\footnote{ZJRB, April 7, 1975, p.~1.
  《浙江日报》，1975年4月7日，第1版。} However, the highlight of the
proceedings was certainly the appearance of the legendary Chen Yonggui
in Zhejiang. It was in the aftermath of Deng's caustic remarks of the
previous month about the harm being caused by factional disputes that
Chen Yonggui came to Hangzhou. Although he was gravely ill at this time,
Zhou Enlai had reportedly called Chen and Ji Dengkui to his bedside to
enquire about the situation in Zhejiang prior to Chen's departure for
the south.\footnote{Chen Weida, ``We must hold up the red banner of
  Daqing raised by Chairman Mao'', HZRB, May 24,1977 (report to the
  Zhejiang provincial meeting on learning from Daqing in industry). Shen
  Chuyun, RMRB, May 17,1977. Tie Ying later recalled that''a comrade who
  came to Zhejiang to help in rectification brought a verbal message
  from Premier Zhou''. See ZJRB, January 9,1985. Tie was probably
  referring to Chen Yonggui as the bearer of this communication.
  陈伟达：《一定要举起毛主席树立的大庆红旗》，《杭州日报》，1977年5月24日（在浙江省工业学大庆会议上的报告）。沈楚云，《人民日报》，1977年5月17日。铁瑛后来回忆说，``来浙江帮助整顿的一位同志带来了周总理的口信''。见《浙江日报》，1985年1月9日。铁所指传话者很可能是陈永贵。}
Chen spent two days visiting model brigades and places of interest in
the company of Tan Qilong and Chen Weida before addressing the
conference.\footnote{One of the model brigades that Chen Yonggui visited
  in Zhejiang was Shangwang brigade of the Hongshan commune in Shaoxing
  county, south-east of Hangzhou. Together with Nanbao brigade,
  Shangwang, led by its party secretary Wang Jinyou, was the most famous
  brigade in Zhejiang. During the Cultural Revolution it received a
  great deal of publicity in the provincial press. Shangwang, with
  Nanbao and Jiang Ruwang's Jinjian brigade, were advanced units in
  Zhejiang in the movement to emulate Dazhai in agriculture. For a
  fortnight in October 1977 I conducted ``open-door schooling''
  (开门办学) at the brigade with a class from Hangzhou University.
  陈永贵在浙江访问的先进大队之一，是杭州东南绍兴县红山公社的上旺大队。上旺与南堡大队一起，是浙江最有名的大队，上旺由党委书记王金友领导。文化大革命期间，它在省级报刊上获得大量宣传。上旺、南堡和蒋如望的金建大队，都是浙江农业学大寨运动中的先进单位。1977年10月，我曾同杭州大学一个班在该大队进行了两周的``开门办学''。}
In his three-part address, Chen struck an entirely different note from
that evident in Tan's opening speech.\footnote{ZJRB, April 28, 1975,
  pp.~1, 2, 3. HZRB, April 29, 1975.
  《浙江日报》，1975年4月28日，第1、2、3版；《杭州日报》，1975年4月29日。}
He urged his listeners

中共浙江省委副书记陈伟达于4月6日向会议作了主报告。\footnote{ZJRB, April
  7, 1975, p.~1. 《浙江日报》，1975年4月7日，第1版。}
然而，会议进程中的亮点无疑是传奇人物陈永贵出现在浙江。陈永贵是在邓上月就派系纠纷造成的危害发表严厉评论之后来到杭州的。据报道，周恩来当时病情很重，但在陈南下前仍把陈和纪登奎叫到病榻前，询问浙江形势。\footnote{Chen
  Weida, ``We must hold up the red banner of Daqing raised by Chairman
  Mao'', HZRB, May 24,1977 (report to the Zhejiang provincial meeting on
  learning from Daqing in industry). Shen Chuyun, RMRB, May 17,1977. Tie
  Ying later recalled that''a comrade who came to Zhejiang to help in
  rectification brought a verbal message from Premier Zhou''. See ZJRB,
  January 9,1985. Tie was probably referring to Chen Yonggui as the
  bearer of this communication.
  陈伟达：《一定要举起毛主席树立的大庆红旗》，《杭州日报》，1977年5月24日（在浙江省工业学大庆会议上的报告）。沈楚云，《人民日报》，1977年5月17日。铁瑛后来回忆说，``来浙江帮助整顿的一位同志带来了周总理的口信''。见《浙江日报》，1985年1月9日。铁所指传话者很可能是陈永贵。}
陈在谭启龙和陈伟达陪同下，用两天时间参观了先进大队和名胜，然后向会议讲话。\footnote{One
  of the model brigades that Chen Yonggui visited in Zhejiang was
  Shangwang brigade of the Hongshan commune in Shaoxing county,
  south-east of Hangzhou. Together with Nanbao brigade, Shangwang, led
  by its party secretary Wang Jinyou, was the most famous brigade in
  Zhejiang. During the Cultural Revolution it received a great deal of
  publicity in the provincial press. Shangwang, with Nanbao and Jiang
  Ruwang's Jinjian brigade, were advanced units in Zhejiang in the
  movement to emulate Dazhai in agriculture. For a fortnight in October
  1977 I conducted ``open-door schooling'' (开门办学) at the brigade
  with a class from Hangzhou University.
  陈永贵在浙江访问的先进大队之一，是杭州东南绍兴县红山公社的上旺大队。上旺与南堡大队一起，是浙江最有名的大队，上旺由党委书记王金友领导。文化大革命期间，它在省级报刊上获得大量宣传。上旺、南堡和蒋如望的金建大队，都是浙江农业学大寨运动中的先进单位。1977年10月，我曾同杭州大学一个班在该大队进行了两周的``开门办学''。}
在分三部分的讲话中，陈的调子同谭开幕讲话中显现的调子完全不同。\footnote{ZJRB,
  April 28, 1975, pp.~1, 2, 3. HZRB, April 29, 1975.
  《浙江日报》，1975年4月28日，第1、2、3版；《杭州日报》，1975年4月29日。}
他敦促听众

\begin{quote}
to study conscientiously Chairman Mao's important instruction on the
question of theory, adhere to the Party's basic line, relentlessly
criticize capitalism, vigorously build socialism, strengthen proletarian
Party spirit, resolutely overcome bourgeois factionalism, promote
stability and unity, do a good job in socialist revolution and socialist
construction, \ldots{}
\end{quote}

\begin{quote}
认真学习毛主席关于理论问题的重要指示，坚持党的基本路线，狠批资本主义，大力建设社会主义，增强无产阶级党性，坚决克服资产阶级派性，促进安定团结，搞好社会主义革命和社会主义建设，\ldots\ldots{}
\end{quote}

Chen reserved the core of his message for the concluding sentences of
the third part of the lengthy speech. In a direct and frank manner he
targeted bourgeois factionalism in Zhejiang. He slammed those people who
had ignored central directives. These people, stated Chen

陈把讲话的核心信息留在长篇讲话第三部分的结尾几句。他以直接坦率的方式把矛头指向浙江的资产阶级派性。他猛烈抨击那些无视中央指示的人。陈说，这些人

\begin{quote}
don't criticize Lin, nor Confucius, don't work, don't go to work and
don't study the theory of the dictatorship of the proletariat; instead
they link up, cause splits, and undermine the dictatorship of the
proletariat. If they continue in this way they will commit even greater
mistakes.
\end{quote}

\begin{quote}
不批林，也不批孔，不劳动，不上班，也不学习无产阶级专政理论；而是串联、搞分裂、破坏无产阶级专政。如果继续这样下去，他们就会犯更大的错误。
\end{quote}

He made it clear that those caught in the web of bourgeois factionalism
included members of party leading groups, members of the CCP, and those
who hoped to become officials by causing trouble. He challenged his
audience with the question,

他明确指出，陷入资产阶级派性罗网的人包括党的领导班子成员、中共党员，以及那些希望通过闹事当官的人。他向听众提出质问：

``Will you give {[}power{]} to them?'' (有的人伸手要官,你就给?). Chen
argued that not only were these people intent on obtaining power, but
they took advantage of disturbances to graft and embezzle public
funds.\footnote{After his arrest in December 1976 Weng Senhe was accused
  of this crime. 翁森鹤于1976年12月被捕后，被指控犯有这一罪行。} Chen
praised the agricultural units that he had visited for their practical
achievements in building socialism. After all, said Chen in his earthy
style, even those who engaged in bourgeois factionalism needed to eat.
The message for the rebel leaders in Zhejiang could not have been
clearer.

``有的人伸手要官，你就给？''陈认为，这些人不仅一心想获得权力，而且利用动乱贪污、挪用公款。\footnote{After
  his arrest in December 1976 Weng Senhe was accused of this crime.
  翁森鹤于1976年12月被捕后，被指控犯有这一罪行。}
陈称赞他所访问的农业单位在建设社会主义方面取得的实际成就。毕竟，陈以质朴的风格说，即使搞资产阶级派性的人也要吃饭。传递给浙江造反派领导人的信息再清楚不过了。

Chen's ``important report'' was heard via wired broadcasts throughout
the counties of Zhejiang.\footnote{ZPS, April 15 and 16, 1975,
  SWB/FE/4881/BII/17-18.
  浙江省广播服务，1975年4月15日、16日，SWB/FE/4881/BII/17-18。} It was
not broadcast over the provincial radio at all. It did not appear in
Zhejiang Daily until 28 April 1975, almost a fortnight after its
delivery. Chen's candidness, especially regarding bourgeois
factionalism, was apparently not considered suitable for foreign ears.
The published text clearly reveals the shift in political direction
Chen's visit occasioned in Zhejiang. For provincial propagandists,
bourgeois factionalism henceforth superseded bourgeois rights as their
principal catch-phrase in the lexicon of political evils.

陈的``重要报告''通过有线广播传达到浙江各县。\footnote{ZPS, April 15 and
  16, 1975, SWB/FE/4881/BII/17-18.
  浙江省广播服务，1975年4月15日、16日，SWB/FE/4881/BII/17-18。}
它完全没有在省电台广播。直到1975年4月28日，也就是报告作出近两周后，才刊登在《浙江日报》上。陈的坦率，尤其是关于资产阶级派性的问题，显然被认为不适合让外国人听到。发表的文本清楚显示出陈到访给浙江带来的政治方向转变。对省级宣传人员来说，从此资产阶级派性取代资产阶级法权，成为政治罪恶词汇表中的主要口号。

At the closing session of the conference on 17 April Tie Ying, Tan's
deputy in the party, government and military, delivered the summing-up
report.\footnote{ZPS, April 22, 1975, SWB/FE/4886/BII/7; HZRB, April 22,
  1975, from which the following quotations are taken.
  浙江省广播服务，1975年4月22日，SWB/FE/4886/BII/7；《杭州日报》，1975年4月22日，以下引文即出自该报。}
Some of Tie's most outspoken statements, as with those of Tan and Chen
Yonggui, were not broadcast over the provincial radio. Tie concentrated
his fire on bourgeois factionalism and he stated that

4月17日会议闭幕会上，谭在党、政、军方面的副手铁瑛作总结报告。\footnote{ZPS,
  April 22, 1975, SWB/FE/4886/BII/7; HZRB, April 22, 1975, from which
  the following quotations are taken.
  浙江省广播服务，1975年4月22日，SWB/FE/4886/BII/7；《杭州日报》，1975年4月22日，以下引文即出自该报。}
同谭和陈永贵的讲话一样，铁的一些最直率表述没有在省电台广播。铁把火力集中于资产阶级派性，并说

\begin{quote}
the root of bourgeois factionalism lies with the Lin Biao anti-Party
clique; it is the poison of the splittist line of Lin Biao and a
reflection among the cadres and people of the corrosion of bourgeois
rights {[}!{]}.
\end{quote}

\begin{quote}
资产阶级派性的根子在林彪反党集团；它是林彪分裂主义路线的毒害，是资产阶级法权侵蚀在干部和群众中的反映{[}！{]}。
\end{quote}

Tie demanded that ``Leading groups at various levels must take the lead
in overcoming bourgeois factionalism, and strengthening proletarian
Party spirit and revolutionary unity''. The contrast in emphasis and
language between Tan's opening address and Tie's closing report was
striking and significant. Tie felt sufficiently confident of his
position, even though he had only recently returned to active duty after
having been sidelined in semi-disgrace for most of 1974, to mock the
radicals obsession with bourgeois rights by turning the phrase back
against its propagators. Clearly, it was Chen's arrival that had turned
the tide against the local radicals in Zhejiang and bolstered the
position of their most steadfast opponents such as Tie Ying.

铁要求``各级领导班子必须带头克服资产阶级派性，增强无产阶级党性和革命团结''。谭的开幕讲话与铁的闭幕报告在重点和语言上的对比十分鲜明而且重要。尽管铁在1974年大部分时间都以半失势状态被边缘化，直到最近才恢复实际工作，但他对自己的地位已有足够信心，能够把``资产阶级法权''这一说法反过来对准其宣传者，以嘲讽激进派对这一概念的执迷。显然，是陈的到来使浙江的潮流转向不利于地方激进派，并加强了铁瑛等最坚定反对者的地位。

The Zhejiang Daily editorial summarizing the achievements of the
conference reflected the tension and ambiguity in central
propaganda.\footnote{ZJRB, April 22,1975, pp.~1, 2.
  《浙江日报》，1975年4月22日，第1、2版。} It called for criticism of
both bourgeois rights and bourgeois factionalism, the targets of Jiang
Qing's radical mobilizers and Deng Xiaoping's upholders of order
respectively. Yet capitalist tendencies among officials were now defined
more narrowly and concretely, and not in the more sweeping language of
ideological deviation. An editorial published the very next day focused
attention back on the issues of primary concern to Deng and provincial
administrators alike. Zhejiang Daily wrote optimistically of a ``new
situation'' on the industrial and communications front in the province
and it asked leading cadres to ``strengthen their proletarian party
spirit and combat bourgeois factionalism in order further to promote
stability and unity''.\footnote{ZPS, April 23, 1975, SWB/FE/4893/BII/18.
  浙江省广播服务，1975年4月23日，SWB/FE/4893/BII/18。}

《浙江日报》总结会议成果的社论反映了中央宣传中的紧张和含糊。\footnote{ZJRB,
  April 22,1975, pp.~1, 2. 《浙江日报》，1975年4月22日，第1、2版。}
它要求同时批判资产阶级法权和资产阶级派性，二者分别是江青激进动员者和邓小平秩序维护者所针对的目标。然而，官员中的资本主义倾向如今被定义得更狭窄、更具体，而不是用更宽泛的意识形态偏差语言来表述。次日发表的一篇社论又把注意力转回邓和省级行政人员同样最关心的问题。《浙江日报》乐观地谈到该省工业和交通战线上的``新形势''，并要求领导干部``增强无产阶级党性，反对资产阶级派性，以便进一步促进安定团结''。\footnote{ZPS,
  April 23, 1975, SWB/FE/4893/BII/18.
  浙江省广播服务，1975年4月23日，SWB/FE/4893/BII/18。}

\section{Factionalism in 1975}\label{factionalism-in-1975}

\section{1975年的派性}\label{ux5e74ux7684ux6d3eux6027}

The final section of this chapter will examine in greater detail the
issue of factionalism. Previous chapters of the study have commented on
the open admission by spokesmen from the radical wing of the CCP that
factionalism was the inevitable consequence of the existence of
different groups within the party. The logical consequences of this
position, which were spelled out from time to time, in 1968 and 1973 to
1974 in particular, were anathema to the mainstream Leninists of the
party. They believed fervently in unity and order. For this reason alone
they disliked and feared the factionalism spawned by the Cultural
Revolution and its ramifications for the future. Their spokesman in 1975
was Deng Xiaoping, who has not deviated from this core belief in the
fifteen years which have since passed. Deng's principal argument at that
time was that the factionalism of 1975 differed in certain important
aspects from that of the early years of the Cultural
Revolution.\footnote{For a further discussion, see Keith Forster,
  ``Factional Politics in Zhejiang, 1973-76'' in William Joseph,
  Christine Wong and David Zweig (eds.). New Perspectives on the
  Cultural Revolution (Cambridge, Mass.: Harvard University Press,
  forthcoming). 进一步讨论见Keith Forster：``Factional Politics in
  Zhejiang, 1973-76''，载William Joseph、Christine Wong、David
  Zweig编，New Perspectives on the Cultural
  Revolution（马萨诸塞州剑桥：哈佛大学出版社，即将出版）。}

本章最后一节将更详细地考察派性问题。本研究前几章曾评论过中共激进派发言人公开承认的一点，即派性是党内存在不同群体的必然后果。这一立场的逻辑后果不时被明说，尤其是在1968年和1973至1974年间，而它们令党内主流列宁主义者深恶痛绝。他们热切相信团结和秩序。单凭这一点，他们就不喜欢并害怕文化大革命催生的派性及其对未来的影响。1975年，他们的发言人是邓小平，而在此后十五年中，邓从未偏离这一核心信念。邓当时的主要论点是，1975年的派性在某些重要方面不同于文化大革命早年的派性。\footnote{For
  a further discussion, see Keith Forster, ``Factional Politics in
  Zhejiang, 1973-76'' in William Joseph, Christine Wong and David Zweig
  (eds.). New Perspectives on the Cultural Revolution (Cambridge, Mass.:
  Harvard University Press, forthcoming). 进一步讨论见Keith
  Forster：``Factional Politics in Zhejiang, 1973-76''，载William
  Joseph、Christine Wong、David Zweig编，New Perspectives on the
  Cultural Revolution（马萨诸塞州剑桥：哈佛大学出版社，即将出版）。}

In his July 1975 speech to a study class in Beijing, Deng spoke at some
length of his fears concerning factionalism.\footnote{Deng Xiaoping,
  ``加强党的领导,整顿党的作风'' (Strengthen party leadership, rectify
  the party's workstyle), 4 July 1975, Deng Xiaoping Wenxuan, pp.12-14.
  One member of the audience was the party secretary of the Nanbao
  brigade, Li Jinrong. See ZPS, July 25, 1976, SWB/FE/5273/BII/6.
  邓小平：《加强党的领导，整顿党的作风》，1975年7月4日，《邓小平文选》，第12-14页。听众之一是南堡大队党委书记李金荣。见浙江省广播服务，1976年7月25日，SWB/FE/5273/BII/6。}
He claimed that the factional groups then in existence were of a
different nature (性质就不同了) than those formed spontaneously
(突然形成的) in the early stages of the Cultural Revolution. Exactly how
they differed was not spelled out. In an important document drawn up
under Deng's aegis, scorn and vitriol were heaped on rebel factional
leaders and it was clearly on men such as Weng Senhe that Deng's
``gunners'' trained their literary artillery.\footnote{This claim is
  made in Chi Hsin, The Case of the Gang of Four, p.~168. 这一说法见Chi
  Hsin, The Case of the Gang of Four，第168页。} Radical theoreticians
were labelled ``sham Marxist political swindlers'', while rebels such as
Zhang Yongsheng, Weng Senhe and He Xianchun were described, although not
by name, as ``unreformed petty intellectuals (没有得到改造的小知识分子)
and 'bold elements''' (勇干分子):

邓在1975年7月对北京一个学习班的讲话中，较为详细地谈到他对派性的担忧。\footnote{Deng
  Xiaoping, ``加强党的领导,整顿党的作风'' (Strengthen party leadership,
  rectify the party's workstyle), 4 July 1975, Deng Xiaoping Wenxuan,
  pp.12-14. One member of the audience was the party secretary of the
  Nanbao brigade, Li Jinrong. See ZPS, July 25, 1976, SWB/FE/5273/BII/6.
  邓小平：《加强党的领导，整顿党的作风》，1975年7月4日，《邓小平文选》，第12-14页。听众之一是南堡大队党委书记李金荣。见浙江省广播服务，1976年7月25日，SWB/FE/5273/BII/6。}
他声称，当时存在的派系集团同文化大革命初期自发形成（突然形成的）的组织性质不同（性质就不同了）。至于究竟如何不同，他并未说明。在邓主持起草的一份重要文件中，对造反派系领导人充满蔑视和痛斥，显然邓的``炮手''把文学炮火瞄准了翁森鹤这样的人。\footnote{This
  claim is made in Chi Hsin, The Case of the Gang of Four, p.~168.
  这一说法见Chi Hsin, The Case of the Gang of Four，第168页。}
激进理论家被贴上``假马克思主义政治骗子''的标签，而张永生、翁森鹤、贺贤春等造反派虽未被点名，却被描述为``没有得到改造的小知识分子''和``勇干分子''：

\begin{quote}
These people are politically ignorant and unexperienced {[}sic{]} in
production. Yet, they make the most noise, pointing their fingers and
calling the shots, accusing people and singing a high sounding tune, but
never working out concrete problems. All the time they label people as
``reviving tradition'', falling backward'', ``conservative forces'',
``pulling the cart without looking at the road'' or ``suppressing the
revolutionary zeal of the cadres and the masses''.\footnote{``Some
  Problems in Accelerating Industrial Development'', 2 September 1975,
  in Chi Hsin, The Case of the Gang of Four, p.~243. ``Some Problems in
  Accelerating Industrial Development''，1975年9月2日，载Chi Hsin, The
  Case of the Gang of Four，第243页。}
\end{quote}

\begin{quote}
这些人在政治上无知，在生产上没有经验。然而，他们声音最大，指手画脚，发号施令，指责别人，唱高调，却从不解决具体问题。他们总是给别人扣上``复辟传统''\,``落后''\,``保守势力''\,``拉车不看路''或``压制干部和群众革命积极性''的帽子。\footnote{``Some
  Problems in Accelerating Industrial Development'', 2 September 1975,
  in Chi Hsin, The Case of the Gang of Four, p.~243. ``Some Problems in
  Accelerating Industrial Development''，1975年9月2日，载Chi Hsin, The
  Case of the Gang of Four，第243页。}
\end{quote}

Leading provincial officials and local newspaper commentaries echoed
these sentiments in the following months.\footnote{See CNS, 579 (August
  20,1975) for a survey. 综述见CNS，579（1975年8月20日）。} For example,
Liao Zhiguo, first secretary of the CCP Fujian Provincial Committee,
stated in an August address to sportsmen in Fuzhou that bourgeois
factionalism opposed party leadership and served Lin Biao's revisionist
line. In words which were almost a repetition of those voiced by Deng,
Liao stated that

在随后几个月中，省级主要官员和地方报刊评论呼应了这些看法。\footnote{See
  CNS, 579 (August 20,1975) for a survey.
  综述见CNS，579（1975年8月20日）。}
例如，中共福建省委第一书记廖志高8月在福州对运动员讲话时说，资产阶级派性反对党的领导，为林彪的修正主义路线服务。廖用几乎是在重复邓的话语说，

\begin{quote}
it is no longer in the form of small groups as it first appeared during
the great cultural revolution. In short, bourgeois factionalism has been
essentially changed.\footnote{Fujian Provincial Service, August 7,1975,
  SWB/FE/4983/BII/10-11.
  福建省广播电台，1975年8月7日，SWB/FE/4983/BII/10-11。}
\end{quote}

\begin{quote}
它已经不是文化大革命初期刚出现时那种小团体的形式。总之，资产阶级派性已经发生了本质变化。\footnote{Fujian
  Provincial Service, August 7,1975, SWB/FE/4983/BII/10-11.
  福建省广播电台，1975年8月7日，SWB/FE/4983/BII/10-11。}
\end{quote}

Describing the struggle against bourgeois factionalism as a struggle
between the two lines Liao was thereby viewing it as an antagonistic
contradiction, which by definition called for suppression by the
dictatorship of the proletariat. By implication then, the factionalism
of mass organizations in the 1960s had been confined to the realm of
contradictions among the people, to be handled by means of education and
persuasion.

廖把反对资产阶级派性的斗争描述为两条路线的斗争，因此把它视为敌我矛盾；按照定义，这就要求由无产阶级专政加以镇压。由此推论，20世纪60年代群众组织的派性则限于人民内部矛盾领域，应以教育和说服的方式处理。

An ideological commentary from Heilongjiang, also in August, stated that
``the factionalism we are fighting against now is essentially different
from that which formed spontaneously between the two mass organizations
in the early stage of the great proletarian cultural revolution''. It
argued that those people ``afflicted'' with the ``illness'', who
included leading cadres and ``leaders of various groups'', demanded
independence from the party while simultaneously fighting for leadership
over it. Practitioners of factionalism were warned to desist or face
serious consequences. In graphic words, bourgeois factionalism was
portrayed as

同样在8月，黑龙江的一篇意识形态评论说，``我们现在所反对的派性，同无产阶级文化大革命初期两个群众组织之间自发形成的派性有本质区别''。文章认为，那些患有这种``病''的人，包括领导干部和``各种集团的头头''，一方面要求独立于党，另一方面又争夺对党的领导权。搞派性的人受到警告，必须停止，否则将面临严重后果。文章用形象的语言把资产阶级派性描绘成

\begin{quote}
an illness that is hard to cure. Bourgeois factionalism is latent,
changeable and obstinate in nature. If you retreat an inch \ldots{} it
demands a yard from you. Confronted with criticism {[}it{]} steps back
and lays low; after a certain period of time, when the conditions are
favorable, it may stage a comeback. ``Bourgeois factionalism is a spring
which yields only to power\ldots{}''.\footnote{Heilongjiang Provincial
  Service, August 13,1975, SWB/FE/4987/BII/1-3. See also the editorial
  of Nanfang Ribao (Guangzhou), Guangdong Provincial Service, August
  18,1975, in SWB/FE/4992/BII/5-7; Yunnan Provincial Service, October
  30,1975, SWB/FE/5052/BII/5-9; CNS, 579, pp.~3-4, 6.
  黑龙江省广播电台，1975年8月13日，SWB/FE/4987/BII/1-3。另见《南方日报》（广州）社论，广东省广播电台，1975年8月18日，载SWB/FE/4992/BII/5-7；云南省广播电台，1975年10月30日，SWB/FE/5052/BII/5-9；CNS，579，第3-4、6页。}
\end{quote}

\begin{quote}
一种难以治愈的疾病。资产阶级派性本质上具有潜伏性、变化性和顽固性。你退一寸\ldots\ldots 它就向你要一尺。面对批评，{[}它{]}退后并潜伏下来；过一段时间，条件有利时，又可能卷土重来。``资产阶级派性是一根弹簧，只在强力面前才会屈服\ldots\ldots''\footnote{Heilongjiang
  Provincial Service, August 13,1975, SWB/FE/4987/BII/1-3. See also the
  editorial of Nanfang Ribao (Guangzhou), Guangdong Provincial Service,
  August 18,1975, in SWB/FE/4992/BII/5-7; Yunnan Provincial Service,
  October 30,1975, SWB/FE/5052/BII/5-9; CNS, 579, pp.~3-4, 6.
  黑龙江省广播电台，1975年8月13日，SWB/FE/4987/BII/1-3。另见《南方日报》（广州）社论，广东省广播电台，1975年8月18日，载SWB/FE/4992/BII/5-7；云南省广播电台，1975年10月30日，SWB/FE/5052/BII/5-9；CNS，579，第3-4、6页。}
\end{quote}

Thus, certain Chinese leaders viewed the factionalism of the 1970s as
being different in nature and in form from the factionalism of the later
1960s. It was hinted that rather than being spontaneous, factionalism in
1975 was organized, directed and manipulated from above. Factional
groups strove to infiltrate the party so as to take control over it
while simultaneously refusing to submit their groups to party
leadership. Some of these groups were officially-sanctioned mass
organizations reconstituted in 1973 or, as in the case of the urban
militia commands, established in 1973 and 1974. The warnings to cadres
and other individuals to halt their factional activities suggested also
the existence of clandestine groups posing a direct challenge to
authority, such as the underground command center in Hangzhou.

因此，某些中国领导人把20世纪70年代的派性视为在性质和形式上不同于60年代后期的派性。其暗示是，1975年的派性并非自发产生，而是由上层组织、指挥和操纵的。派系集团力图渗入党内以控制党，同时又拒绝让自己的集团服从党的领导。其中一些集团是1973年重组的官方认可群众组织，或者像城市民兵指挥部那样，是1973年和1974年建立的。对干部和其他个人发出的停止派系活动的警告，也暗示存在着直接挑战权威的秘密集团，例如杭州的地下指挥中心。

In Zhejiang the problem of factionalism was acute. A Zhejiang Daily
editorial of 23 August 1975 entitled ``Give a free hand in mobilizing
the people to savagely criticize bourgeois factionalism''
(放手发动群众狠批资产阶级派性), slammed ``several chiefs of bourgeois
factions'' as ``evil people'' who ``frequently flaunt themselves as
representing `the masses', strike a pose and make themselves out as
`huge monsters' (庞然大物)''. The editorial detailed the following
features of ``bourgeois factionalism'' in the province:

在浙江，派性问题十分尖锐。1975年8月23日《浙江日报》题为《放手发动群众狠批资产阶级派性》的社论，把``几个资产阶级派性的头头''斥为``坏人''，说他们``经常打着代表`群众'的幌子，装腔作势，把自己说成`庞然大物'\,''。社论详细列举了该省``资产阶级派性''的下列特征：

\begin{quote}
*establishing groups and mountain strongholds, recognizing factions and
not the party, using factions to suppress the party and stir up
independence from it, setting up many centers, riding roughshod over
party committees, weakening, breaking away from and opposing party
leadership;
\end{quote}

\begin{quote}
*拉山头、建山寨，认派不认党，以派压党，煽动独立于党之外，搞多中心，凌驾于党委之上，削弱、脱离和反对党的领导；
\end{quote}

\begin{quote}
*differentiating between ``legalist'' (法家党) and ``Confucian''
(儒家党) groups within the party and dividing the revolutionary ranks
into ``innovators'' (革新派) and ``restorationists'' (复辟派) to split
the party and the masses and to sabotage stability and unity;
\end{quote}

\begin{quote}
*在党内划分``法家党''和``儒家党''，把革命队伍划分为``革新派''和``复辟派''，以分裂党和群众，破坏安定团结；
\end{quote}

\begin{quote}
*drawing lines according to factions and preaching that ``to rebel has
merit, join the party to get one's share'' (造反有功,入党有份), or
``whoever goes against the tide can join the party and become a cadre'',
stirring up the bourgeois wind of struggling for fame and advantage,
exchanging flattery and favours, promising high posts and other favours
thus creating various impurities within the party's organization and
cadre ranks;
\end{quote}

\begin{quote}
*按派划线，宣扬``造反有功，入党有份''，或者``谁反潮流谁就能入党、当干部''，煽起争名夺利的资产阶级风气，互相吹捧、互相照顾，许愿高官厚禄，从而在党的组织和干部队伍中造成各种不纯；
\end{quote}

\begin{quote}
*preaching the view that at present the principal contradiction is ``the
contradiction between new and old cadres'', ``the contradiction between
this and that faction'' and advocating such strange theories as ``you
can't be wrong if you direct your attack at the leadership'';
\end{quote}

\begin{quote}
*宣扬目前主要矛盾是``新老干部之间的矛盾''\,``这一派和那一派之间的矛盾''，鼓吹``矛头指向领导不会错''等奇谈怪论；
\end{quote}

\begin{quote}
*sermonizing about relying on some ``faction'' and that ``we must use
revolutionary intellectuals to reform the cadre ranks'' in order to
sabotage the party's class line of relying on the working class and
uniting with other laboring people in the cities, and relying on the
poor and lower-middle peasants and uniting with middle peasants in the
countryside;
\end{quote}

\begin{quote}
*鼓吹依靠某个``派''，宣扬``要用革命知识分子改造干部队伍''，以破坏党在城市依靠工人阶级、团结其他劳动人民，在农村依靠贫下中农、团结中农的阶级路线；
\end{quote}

\begin{quote}
*acting hypocritically by verbally expressing support for the party's
leadership and the implementation of party directives while behind the
scenes spreading rumors and slander, attacking and opposing the party's
leadership, resisting party directives and pursuing a set of vulgar
things in the manner of bourgeois politicians
(搞资产阶级政客的那一套庸俗的东西);
\end{quote}

\begin{quote}
*口头上表示拥护党的领导、执行党的指示，背地里却散布谣言和诽谤，攻击和反对党的领导，抵制党的指示，搞资产阶级政客的那一套庸俗的东西；
\end{quote}

\begin{quote}
*advocating the reactionary fallacy of ``don't work for the wrong
line'', and smearing those cadres and masses who persist in grasping
revolution and production as ``lambs'' (小绵羊) and ``mediocrities who
don't understand politics'' (不问政治的庸人);
\end{quote}

\begin{quote}
*鼓吹``路线不对不干''的反动谬论，把坚持抓革命、促生产的干部和群众污蔑为``小绵羊''和``不问政治的庸人''；
\end{quote}

\begin{quote}
*shielding the class enemy thus allowing capitalism to spread
unchecked.\footnote{Xuexi Wenji, pp.~63-6. 《学习文集》，第63-66页。}
\end{quote}

\begin{quote}
*包庇阶级敌人，从而使资本主义泛滥。\footnote{Xuexi Wenji, pp.~63-6.
  《学习文集》，第63-66页。}
\end{quote}

Factionalism had afflicted the industrial working class in particular.
It may have been inflamed further by the theoretical explorations of
Zhang Chunqiao, who seemed to be groping for a basis in Marxism to
differentiate between different strata of the proletariat. On 14 October
1974 Zhang had reportedly stated that ``Relying on the entire working
class gives one the impression that there are no contradictions within
the working class. That does not conform with dialectics''. In a speech
of an unspecified date, possibly at the end of 1974, Zhang had remarked
that ``We can only rely on the `Leftists' of the working class''.
Furthermore, in February 1975, Zhang put forward the view that ``The
veteran workers should also be subjected to class analysis''.\footnote{Document
  of the CCP CC, zhongfa, (1977) No.~37, I\&S, 15:1 (1979), p.~100.
  中共中央文件，中发（1977）37号，I\&S，15:1（1979年），第100页。}

派性尤其困扰着产业工人阶级。张春桥的理论探索可能进一步煽动了这种派性，因为他似乎正在马克思主义中寻找一种依据，以区分无产阶级内部不同阶层。据报道，1974年10月14日张曾说：``依靠整个工人阶级，给人的印象就是工人阶级内部没有矛盾。这不符合辩证法。''在一次日期未明、可能发生于1974年底的讲话中，张评论说：``我们只能依靠工人阶级中的`左派'。''此外，1975年2月，张提出``老工人也要作阶级分析''的观点。\footnote{Document
  of the CCP CC, zhongfa, (1977) No.~37, I\&S, 15:1 (1979), p.~100.
  中共中央文件，中发（1977）37号，I\&S，15:1（1979年），第100页。}

A Fujian Daily editorial of 13 July 1975 characterized attempts to
establish factions within the working class as ``a counter-revolutionary
plot engineered by Lin Biao and company to split the Party and the
revolutionary ranks''.\footnote{Fujian Provincial Service, July 13,
  1975, SWB/FE/4958/BII/16-17.
  福建省广播电台，1975年7月13日，SWB/FE/4958/BII/16-17。} The editorial
differentiated between veteran workers who, with their experience and
longer exposure to party education, were presumably less susceptible to
the attraction of factional alignments than their younger, less
socialized and more impressionable comrades. Reports from other
provinces demonstrated that factionalism within the working class was a
widespread phenomenon.\footnote{CNA 1012 (September 5,1975), and 1013
  (September 12,1975).
  中央通讯社1012（1975年9月5日）和1013（1975年9月12日）。} With the
State Council and its subordinate departments preparing for the
beginning of a new 5-Year plan in 1976, these signs of instability were
disturbing indeed.

1975年7月13日《福建日报》社论把在工人阶级内部建立派别的企图，定性为``林彪一伙分裂党和革命队伍的反革命阴谋''。\footnote{Fujian
  Provincial Service, July 13, 1975, SWB/FE/4958/BII/16-17.
  福建省广播电台，1975年7月13日，SWB/FE/4958/BII/16-17。}
社论区分了老工人与年轻工人：老工人凭借经验和较长时间接受党的教育，推定比那些社会化程度较低、更易受影响的年轻同志更不容易被派系结盟所吸引。来自其他省份的报道表明，工人阶级内部的派性是一种普遍现象。\footnote{CNA
  1012 (September 5,1975), and 1013 (September 12,1975).
  中央通讯社1012（1975年9月5日）和1013（1975年9月12日）。}
在国务院及其下属部门正准备于1976年开始新的五年计划之际，这些不稳定迹象确实令人不安。

In his speech of July 4 Deng had also commented on the question of
building provincial party committees which were neither weak, lazy nor
lax (软,懒,散). He stated that criticism of provincial committees'
mistakes was made in order to help them
(``中央批评他们也是为了帮助'').\footnote{Deng Xiaoping Wenxuan, p.~13.
  《邓小平文选》，第13页。} For Deng, the trouble in Zhejiang
exemplified all the shortcomings that he had enumerated in his speeches
of the first half of 1975: weak party leadership, entrenched
factionalism which had entrapped party leaders, complete disregard for
discipline in the work-place and the involvement of the PLA in local
strife. Zhejiang was thus an ideal province for Deng to test his
political resolve and that of his colleagues on the Politburo. Disunity
at the center at the beginning of 1975 had compromised and aborted
decisions taken at the Beijing conference, but by July 1975 conditions
were more favorable for decisive action.

邓在7月4日的讲话中还评论了建设不软、不懒、不散（软、懒、散）的省级党委问题。他说，对省委错误的批评是为了帮助他们（``中央批评他们也是为了帮助''）。\footnote{Deng
  Xiaoping Wenxuan, p.~13. 《邓小平文选》，第13页。}
对邓来说，浙江的麻烦体现了他在1975年上半年讲话中列举的所有缺点：党的领导软弱，根深蒂固的派性把党的领导人卷入其中，工作场所纪律被完全无视，解放军卷入地方冲突。因此，浙江是邓检验自己以及政治局同僚政治决心的理想省份。1975年初中央的不团结削弱并中止了北京会议作出的决定，但到1975年7月，采取果断行动的条件更加有利。

\section{Notes}\label{notes-7}

\section{注释}\label{ux6ce8ux91ca-7}

\bookmarksetup{startatroot}

\chapter{CHAPTER NINE - CENTRAL INTERVENTION, JULY
1975}\label{chapter-nine---central-intervention-july-1975}

\bookmarksetup{startatroot}

\chapter{第九章 -
中央干预，1975年7月}\label{ux7b2cux4e5dux7ae0---ux4e2dux592eux5e72ux98841975ux5e747ux6708}

Reports that problems existed in Hangzhou appeared in the Chinese
national press in mid-July 1975. A Xinhua release of 13 July surveyed
six major factories in the city -- the Hangzhou Iron and Steel Mill, the
Hangzhou Silk Printing and Dyeing Complex, the Hangzhou Oxygen Generator
Factory, the Hangzhou No.~1 Cotton Dyeing and Printing Mill, the
Hangzhou No.~2 Cotton Mill (杭州第二棉纺织厂) and the Zhejiang Hemp Mill
and described them as important units in the metallurgy, machinery and
textile branches. The factories turned out producer and consumer items
for the national market and also supplied foreign customers.\footnote{RMRB,
  July 14, 1975, p.~1.

  RMRB，1975年7月14日，第1页。}

1975年7月中旬，中国国家报刊出现了关于杭州存在问题的报道。7月13日，新华社调查了该市六家主要工厂------杭州钢铁厂、杭州丝绸印染联合厂、杭州制氧机厂、杭州第一棉印染厂、杭州第二棉纺织厂和浙江麻厂，并称它们是冶金、机械和纺织部门的重要单位。这些工厂为国内市场生产生产资料和消费品，也供应国外客户。\footnote{RMRB,
  July 14, 1975, p.~1.

  RMRB，1975年7月14日，第1页。}

Although the article was packed with the Maoist rhetoric and slogans
characteristic of the times, it also reflected genuine concerns about
lagging production. It attributed such difficulties to the
``interference of bourgeois factionalism'', ``the crimes of Liu Shaoqi
and Lin Biao in splitting the revolutionary ranks'' and sheeted home the
blame on ``bad people sabotaging stability and unity''. The report
referred to the need to strengthen unity among workers and tighten
discipline in the workplace so as to ensure the quantity and quality of
output. Efforts to revive production were taking place, as the article
noted, at the height of Hangzhou's oppressive summer.\footnote{The high
  temperatures in Hangzhou at this time were frequently commented on in
  provincial reports. July is normally the hottest month in Hangzhou.
  Its most famous attraction, the West Lake, contributes to the
  suffocating humidity in that month. The lake is surrounded by hills on
  three sides and warm moisture rises from its surface and remains
  suspended over the city. Having lived through the dreadful summer of
  1978 when the highest temperature reached a record 39.9°C, the author
  finds it difficult to comprehend the effect of the city's heat on
  steel workers or textile operatives.

  此时杭州的高温，在省级报告中频频被评论。七月通常是杭州最热的月份。其最著名的景点西湖导致了该月令人窒息的潮湿。湖水三面被山丘包围，温暖的湿气从湖面升起，悬浮在城市上空。经历过1978年那个可怕的夏天，当时最高气温达到创纪录的39.9°C，作者很难理解这座城市的高温对钢铁工人或纺织工人的影响。}
Nevertheless, provincial and municipal leaders were praised for
investigating conditions in the enterprises and for listening to
opinions and suggestions while toiling alongside the workers. Their
example had apparently rubbed off on enterprise cadres who had joined
the workers in shift-work. In hindsight then, the report may be seen as
advance warning that dramatic events were about to occur in Zhejiang.

文章虽然充斥着毛主义的言论和时代口号，但也反映了对生产落后的真实担忧。它将这种困难归咎于``资产阶级派性的干扰''、``刘少奇、林彪分裂革命队伍的罪行''，把责任归咎于``坏人破坏安定团结''。报告指出，要加强职工团结，严格工作场所纪律，确保产出数量和质量。文章指出，正值杭州酷暑难耐之际，我们正在努力恢复生产。
尽管如此，省市领导在与工人们一起辛苦劳作的同时，深入企业调研，听取意见建议，受到了赞扬。他们的榜样显然感染了与工人一起轮班工作的企业干部。事后看来，这份报告或许可以被视为对浙江即将发生戏剧性事件的预先警告。\footnote{The
  high temperatures in Hangzhou at this time were frequently commented
  on in provincial reports. July is normally the hottest month in
  Hangzhou. Its most famous attraction, the West Lake, contributes to
  the suffocating humidity in that month. The lake is surrounded by
  hills on three sides and warm moisture rises from its surface and
  remains suspended over the city. Having lived through the dreadful
  summer of 1978 when the highest temperature reached a record 39.9°C,
  the author finds it difficult to comprehend the effect of the city's
  heat on steel workers or textile operatives.

  此时杭州的高温，在省级报告中频频被评论。七月通常是杭州最热的月份。其最著名的景点西湖导致了该月令人窒息的潮湿。湖水三面被山丘包围，温暖的湿气从湖面升起，悬浮在城市上空。经历过1978年那个可怕的夏天，当时最高气温达到创纪录的39.9°C，作者很难理解这座城市的高温对钢铁工人或纺织工人的影响。}

In fact, at the time that the article was published, a central
delegation had been in Hangzhou investigating local affairs for about a
fortnight. It was not, contrary to popular belief, Deng Xiaoping who
personally visited Hangzhou in July 1975 to solve the problems plaguing
the province. It seems that a Hong Kong journalist was responsible for
promoting this inaccuracy. In one article he reported that ``Foreign
visitors have been told that \ldots{} Deng \ldots{} personally visited
Hangzhou to thrash out problems with the municipal
authorities''.\footnote{Goodstadt, FEER, August 15,1975, pp.~24-5; Leo
  Goodstadt, FEER, August 1,1975, p.~30. See also Charts, p.~4, although
  with this exception, Taiwan sources knew better.

  Goodstadt，FEER，1975年8月15日，第24-25页；Leo
  Goodstadt，FEER，1975年8月1日，第30页。另参见
  Charts，第4页；不过除这一例外，台湾消息来源对此更清楚。} Several
articles claimed that both Deng and Wang Hongwen visited Hangzhou in the
spring of 1975 to review the situation\footnote{See Fox Butterfield in
  the New York Times, 29 July 1975, pp.~1, 6; Goodstadt, FEER, August 1,
  1975, p.~30; ``Tan Qilong'', p.~103; CNS, 577 p.~2 specified the month
  as March 1975.

  参见 Fox
  Butterfield，《纽约时报》，1975年7月29日，第1、6页；Goodstadt，FEER，1975年8月1日，第30页；《谭启龙》，第103页；CNS，577，第2页，明确称月份为1975年3月。}
and Leo Goodstadt contended that Wang ``failed to restore order''. A
Taiwan publication asserted that both Wang and Deng failed in their
missions.\footnote{Charts, p.~4

  Charts，第4页。} Secondary sources\footnote{See, for example, Fenwick,
  ``The Gang of Four'', p.~352; Lowell Dittmer, China's Continuous
  Revolution: The Post-Liberation Epoch, 1949-1981 (Berkeley and Los
  Angeles: University of California Press; 1987), p.~166.

  例如，参见 Fenwick，《四人帮》，第352页；Lowell
  Dittmer，《中国的持续革命：解放后时代，1949-1981》（伯克利和洛杉矶：加州大学出版社，1987年），第166页。}
have taken the assertions as correct, despite reliable evidence to the
contrary which was published soon after the downfall of the Gang of
Four.\footnote{See FEER, April 29,1977, p.~18.

  参见 FEER，1977年4月29日，第18页。} The fact that Hangzhou was
reportedly declared off-limits to Western journalists at the height of
outside interest in its affairs was a mitigating factor which
contributed to this mistake by contemporary observers. Taiwan sources
claim variously that Hangzhou was closed to all foreign and overseas
Chinese tourists from April or May 1975 until the situation was brought
under control. Other sources contradict this assertion. One academic
observer claimed that ``Visitors to China during this period reported
that there was little evidence of the disturbances and that travel to
the popular scenic spots of Hangzhou was, generally speaking, not
interrupted''.\footnote{John B. Starr, ``China in 1975: 'The Wind in the
  Bell Tower''', AS, 16:1 (1976), p.~50.

  John B.
  Starr，《1975年的中国：``钟楼里的风''》，AS，16:1（1976年），第50页。}

事实上，这篇文章发表时，中央代表团已在杭州调查地方事务约两周时间。与普遍看法相反，1975年7月并不是邓小平亲自到杭州解决困扰该省的问题。看来，一名香港记者对这种不准确的报道负有责任。他在一篇文章中报道说，``外国来访者被告知\ldots\ldots 邓\ldots\ldots 亲自访问杭州，与市政当局彻底解决问题''。\footnote{Goodstadt,
  FEER, August 15,1975, pp.~24-5; Leo Goodstadt, FEER, August 1,1975,
  p.~30. See also Charts, p.~4, although with this exception, Taiwan
  sources knew better.

  Goodstadt，FEER，1975年8月15日，第24-25页；Leo
  Goodstadt，FEER，1975年8月1日，第30页。另参见
  Charts，第4页；不过除这一例外，台湾消息来源对此更清楚。}
有几篇文章声称，邓和王洪文均于1975年春访问杭州，检视局势，\footnote{See
  Fox Butterfield in the New York Times, 29 July 1975, pp.~1, 6;
  Goodstadt, FEER, August 1, 1975, p.~30; ``Tan Qilong'', p.~103; CNS,
  577 p.~2 specified the month as March 1975.

  参见 Fox
  Butterfield，《纽约时报》，1975年7月29日，第1、6页；Goodstadt，FEER，1975年8月1日，第30页；《谭启龙》，第103页；CNS，577，第2页，明确称月份为1975年3月。}
Leo
Goodstadt则认为王``未能恢复秩序''。台湾一家出版物声称王和邓均未完成任务。\footnote{Charts,
  p.~4

  Charts，第4页。} 第二手消息来源\footnote{See, for example, Fenwick,
  ``The Gang of Four'', p.~352; Lowell Dittmer, China's Continuous
  Revolution: The Post-Liberation Epoch, 1949-1981 (Berkeley and Los
  Angeles: University of California Press; 1987), p.~166.

  例如，参见 Fenwick，《四人帮》，第352页；Lowell
  Dittmer，《中国的持续革命：解放后时代，1949-1981》（伯克利和洛杉矶：加州大学出版社，1987年），第166页。}
认为这些说法是正确的，尽管在四人帮垮台后不久就发表了相反的可靠证据。\footnote{See
  FEER, April 29,1977, p.~18.

  参见 FEER，1977年4月29日，第18页。}
据报道，在外界对杭州事务最感兴趣的时候，杭州被宣布禁止西方记者进入，这是促成当时观察者误判的一个情有可原的因素。台湾消息人士称，从1975年4月或5月起，杭州对所有外国和海外华人游客关闭，直到情况得到控制。其他消息来源反驳了这一说法。一位学术观察人士称，``这段时间到中国的游客反映，几乎没有发生骚乱的证据，一般来说，前往杭州热门景点的旅行没有中断''。\footnote{John
  B. Starr, ``China in 1975: 'The Wind in the Bell Tower''', AS, 16:1
  (1976), p.~50.

  John B.
  Starr，《1975年的中国：``钟楼里的风''》，AS，16:1（1976年），第50页。}

It is possible that Deng visited Hangzhou in the spring of 1975 if only
to report to Mao, who spent a day or two in the city on his return from
several months rest in Changsha. That Deng played a role in the
settlement of the problem was admitted by a representative of the
provincial leadership to foreign correspondents in January
1977.\footnote{See Radio Peace and Progress (USSR), January 22, 1977,
  SWB/FE/5424/BII/22.

  参见和平与进步广播电台（苏联），1977 年 1 月 22
  日，SWB/FE/5424/BII/22。} However, Deng did not supervise the
leadership changes or personally investigate conditions in Hangzhou's
factories in July. Either on 30 June or 1 July 1975 two emissaries of
the central leadership arrived in Hangzhou.\footnote{In a visit to the
  Hangzhou Silk Complex on November 11, 1977, I was informed that Ji
  arrived in the provincial capital on July 1.

  1977年11月11日，我在参观杭州丝绸联合企业时获悉，纪登奎已于7月1日抵达省城。}
Their brief was to assist the provincial leadership to try and solve the
chronic problems which had plagued the administration for well over
eighteen months. The two central leaders were Wang Hongwen\footnote{This
  was correctly pointed out in ``Tan Qilong'', p.103; Zhonggong Nianbao
  1976, 5:103 and CNS, 583 (September 17,1975), p.~2.

  《谭启龙》，第103页；《中共年报》1976，5:103，以及
  CNS，583（1975年9月17日），第2页，均正确指出了这一点。} and Ji
Dengkui.\footnote{See Chen Zuolin, RMRB, December 18,1976, p.~3. As late
  as 1986 a western academic could claim that Ji's presence in Hangzhou
  in July 1975 was based only on ``unsubstantiated Taiwanese reports''.
  See David M. Lampton, Paths to Power: Elite Mobility in Contemporary
  China (Ann Arbor: Michigan Monographs in Chinese Studies, Volume
  55,1986), p.~49.

  参见陈作霖，RMRB，1976年12月18日，第3页。直到1986年，仍有一位西方学者声称，纪登奎1975年7月出现在杭州只是基于``未经证实的台湾报道''。参见
  David M.
  Lampton，《权力之路：当代中国的精英流动》（安娜堡：密歇根中国研究专著，第55卷，1986年），第49页。}

邓小平有可能在1975年春天访问过杭州，只是为了向毛泽东汇报工作，而毛泽东从长沙休息了几个月回来后，在这座城市待了一两天。1977年1月，一位省领导代表向外国记者承认，邓在问题的解决中发挥了作用。\footnote{See
  Radio Peace and Progress (USSR), January 22, 1977, SWB/FE/5424/BII/22.

  参见和平与进步广播电台（苏联），1977 年 1 月 22
  日，SWB/FE/5424/BII/22。}然而，7月份邓并没有监督领导层更迭，也没有亲自到杭州工厂考察情况。1975年6月30日或7月1日，中央领导层的两位特使抵达杭州。\footnote{In
  a visit to the Hangzhou Silk Complex on November 11, 1977, I was
  informed that Ji arrived in the provincial capital on July 1.

  1977年11月11日，我在参观杭州丝绸联合企业时获悉，纪登奎已于7月1日抵达省城。}他们的任务是协助省领导设法解决困扰政府十八个多月的长期问题。两位中央领导人是王洪文\footnote{This
  was correctly pointed out in ``Tan Qilong'', p.103; Zhonggong Nianbao
  1976, 5:103 and CNS, 583 (September 17,1975), p.~2.

  《谭启龙》，第103页；《中共年报》1976，5:103，以及
  CNS，583（1975年9月17日），第2页，均正确指出了这一点。}和纪登奎\footnote{See
  Chen Zuolin, RMRB, December 18,1976, p.~3. As late as 1986 a western
  academic could claim that Ji's presence in Hangzhou in July 1975 was
  based only on ``unsubstantiated Taiwanese reports''. See David M.
  Lampton, Paths to Power: Elite Mobility in Contemporary China (Ann
  Arbor: Michigan Monographs in Chinese Studies, Volume 55,1986), p.~49.

  参见陈作霖，RMRB，1976年12月18日，第3页。直到1986年，仍有一位西方学者声称，纪登奎1975年7月出现在杭州只是基于``未经证实的台湾报道''。参见
  David M.
  Lampton，《权力之路：当代中国的精英流动》（安娜堡：密歇根中国研究专著，第55卷，1986年），第49页。}。

The choice of Wang was not surprising as he was responsible for party
affairs in East China.\footnote{Oral source.

  口头来源。} However, China News Summary also pointed out the ambiguity
and awkwardness of Wang's position. As the journal perceptively noted,
Wang

王的选择并不令人意外，因为他负责华东地区的党务工作。\footnote{Oral
  source.

  口头来源。}不过，《中国新闻摘要》也指出王的立场的模糊性和尴尬。正如该杂志敏锐地指出的那样，王

\begin{quote}
should have been in a position to provide sympathy and support for the
struggle of the radical Leftists in the Chekiang factories. However, his
official mission there was to restrain and even suppress them and to
compel them to resume work and obey discipline rather than engage in
factional struggle. The success of his mission will alienate him from
the radical Leftists and the failure of his mission will affect his
status as a top Party leader.\footnote{CNS, 583 (September 17,1975),
  p.~2.

  CNS，583（1975年9月17日），第2页。}
\end{quote}

\begin{quote}
应该对浙江工厂里的极左派的斗争表示同情和支持。但他在那里的官方任务是约束甚至镇压他们，迫使他们复工、遵守纪律，而不是搞派系斗争。他的使命成功将导致他与激进左派的疏远，他的使命失败将影响他作为党的最高领导人的地位。\footnote{CNS,
  583 (September 17,1975), p.~2.

  CNS，583（1975年9月17日），第2页。}
\end{quote}

It is highly likely that Zhou Enlai and Deng Xiaoping chose Wang for
these very reasons. Zhang Yongsheng heard on the grapevine that during a
visit to Hangzhou Deng had informed Tan Qilong that he had not come to
solve the problem. He left that for those who had caused it
(解铃须要系铃人).\footnote{Zhang Yongshengde zuizheng cailiao, p.~162.

  《张永生的罪证材料》，第162页。} As a Chinese informant aptly put it:
``Zhou and Deng were very smart for they did not attack the left
themselves but got the middle of the road people to attack the left for
them''.\footnote{``The Case of Weng Sengho'', p.~4.

  《翁森鹤案》，第4页。}

周恩来和邓小平很可能正是因为这些原因才选择了王。张永生听小道消息说，邓小平在杭州视察时告诉谭启龙，他不是来解决问题的。他把这个问题留给造成问题的人去解决，即``解铃须要系铃人''。\footnote{Zhang
  Yongshengde zuizheng cailiao, p.~162.

  《张永生的罪证材料》，第162页。}
正如一位中国知情人所说：``周、邓很聪明，他们自己不攻左，而是让中间的人替他们攻左。''\footnote{``The
  Case of Weng Sengho'', p.~4.

  《翁森鹤案》，第4页。}

It was probably true that Wang continued to command respect among
younger factory workers and they would probably be more likely to listen
to him than a representative of the veteran cadres. And presumably, as
an upholder of law and order, Ji could be expected to share the concern
of veteran leaders for the breakdown of local authority in Zhejiang. As
chapters seven and eight of this study have demonstrated, Wang had
maintained a close interest in the affairs of Zhejiang since his
promotion to the central leadership in 1973. However, his familiarity
with the situation in the province, however superficial, was more than
offset by his partisanship. Ji Dengkui was held in great esteem by Mao
Zedong\footnote{See the interview conducted with Ji shortly before his
  death in July 1988, in 了望周刊 (Outlook Weekly), Nos. 6-7 (February
  6, 1989), pp.~42-4.

  参见1988年7月纪登奎去世前夕对纪登奎的采访，载于《了望周刊》，第6-7期（1989年2月6日），第42-4页。}
and he had no particular ties to the province.\footnote{Jurgen Domes, in
  an article published in 1977, grouped Ji with Hua Guofeng and Wang
  Dongxing as members of what he termed the ``secret police left''. J.
  Domes, ``The `Gang of Four' and Hua Kuo-feng: Analysis of Political
  Events in 1975-76'', CQ, 71 (1977), p.~478.

  Jurgen
  Domes在1977年发表的一篇文章中，将纪与华国锋和汪东兴归为他所谓的``秘密警察左翼''成员。J.
  Domes，《``四人帮''与华国锋：1975-76年政治事件分析》，CQ，71（1977年），第478页。}
His relative objectivity would compensate for and balance Wang's open
partiality for the ``mountain top'' faction. The Taiwan journal Issues
\& Studies pointed out that in holding a supporter of the Gang of Four
{[}Weng Senhe{]} responsible for the trouble, Ji aroused the Gang's
ire.\footnote{See ``Chi Teng-k'ui -- a member of the CCP CC Politburo'',
  I\&S, 13:10 (1977), pp.~70-3.

  参见《迟登奎------中共中央政治局委员》，《I\&S》，1977年13：10，第70-3页。}

王可能确实继续受到工厂年轻工人的尊重，他们可能比老干部代表更愿意听他的话。据推测，作为法律和秩序的维护者，纪可能会与资深领导人一样对浙江地方当局的崩溃感到担忧。正如本研究第七章和第八章所表明的，自
1973
年晋升为中央领导以来，王一直对浙江事务保持着密切的兴趣。然而，他对浙江情况的熟悉无论多么肤浅，都被他的党派偏见所抵消。纪登奎受到毛泽东的器重\footnote{See
  the interview conducted with Ji shortly before his death in July 1988,
  in 了望周刊 (Outlook Weekly), Nos. 6-7 (February 6, 1989), pp.~42-4.

  参见1988年7月纪登奎去世前夕对纪登奎的采访，载于《了望周刊》，第6-7期（1989年2月6日），第42-4页。}，他与省里没有特别的联系。\footnote{Jurgen
  Domes, in an article published in 1977, grouped Ji with Hua Guofeng
  and Wang Dongxing as members of what he termed the ``secret police
  left''. J. Domes, ``The `Gang of Four' and Hua Kuo-feng: Analysis of
  Political Events in 1975-76'', CQ, 71 (1977), p.~478.

  Jurgen
  Domes在1977年发表的一篇文章中，将纪与华国锋和汪东兴归为他所谓的``秘密警察左翼''成员。J.
  Domes，《``四人帮''与华国锋：1975-76年政治事件分析》，CQ，71（1977年），第478页。}
他的相对客观性可以弥补和平衡王对``山上派''的公开偏爱。台湾《问题与研究》杂志指出，纪登奎追究``四人帮''支持者翁森鹤的责任，激起了四人帮的愤怒。\footnote{See
  ``Chi Teng-k'ui -- a member of the CCP CC Politburo'', I\&S, 13:10
  (1977), pp.~70-3.

  参见《迟登奎------中共中央政治局委员》，《I\&S》，1977年13：10，第70-3页。}

After their arrival in the provincial capital, Wang and Ji immediately
embarked on a series of inspection tours of industrial and agricultural
units in and around Hangzhou.\footnote{Ji Dengkui initially concentrated
  on Weng Senhe's silk mill, visiting it eight times in nine days. Visit
  to silk complex, November 11,1977. On July 25, he paid a visit to
  Shangwang brigade. Briefing at Shangwang, October 19, 1977; Wang
  Jinyou, Juqi zhuagang xue Dazhai, p.~40. Wang Hongwen went to the
  Hangzhou Gear-Box Plant with Ji on July 9. Visit, July 17, 1977. Wang
  also toured the Oxygen Generator Plant (HZRB, November 12, 1976), the
  Hangzhou No.~2 Cotton Mill on July 28 (HZRB, November 9,1976), the
  Iron and Steel Mill (HZRB, November 11, 1976) and the Changqing
  brigade in Hangzhou's Jianggan district on July 23 (HZRB, November
  22,1976).

  纪登奎最初把注意力集中在翁森鹤的丝绸厂，9天内访问了8次。1977年11月11日访问丝绸联合厂。7月25日，他访问了上王大队。1977年10月19日在上王听取情况介绍；王金友，Juqi
  zhuagang xue
  Dazhai，第40页。7月9日，王洪文随纪登奎到杭州齿轮箱厂。1977年7月17日访问。王还视察了制氧机厂（HZRB，1976年11月12日）、7月28日视察杭州第二棉纺厂（HZRB，1976年11月9日）、钢铁厂（HZRB，1976年11月11日），以及7月23日视察杭州江干区长庆大队（HZRB，1976年11月22日）。}
They also consulted the provincial authorities with the view to
reorganizing the CCP ZPC, HMC and ZPMD leaderships. Wang and Ji brought
with them a large delegation of central officials from three
departments: the CC Organization Department which would investigate the
procedures followed in the ``two crashthroughs''; the Ministry of Light
Industry, which was responsible for large textile mills in Hangzhou; and
the First Ministry of Machine Building which had jurisdiction over such
factories as the gearbox plant and the oxygen generator
plant.\footnote{HZRB, November 28,1976, July 4,1977. A Xinhua report of
  July 13,1975, concerning the production situation at the oxygen
  generator plant implied that one of its three workshops was out of
  operation. See CNS, 578 (August 13,1975), p.~2.

  HZRB，1976年11月28日，1977年7月4日。1975年7月13日新华社关于制氧机厂生产情况的报道暗示，该厂三个车间中有一个已停产。参见
  CNS，578（1975年8月13日），第2页。}

抵达省城后，王、纪立即对杭州及周边工农业单位进行了一系列考察。\footnote{Ji
  Dengkui initially concentrated on Weng Senhe's silk mill, visiting it
  eight times in nine days. Visit to silk complex, November 11,1977. On
  July 25, he paid a visit to Shangwang brigade. Briefing at Shangwang,
  October 19, 1977; Wang Jinyou, Juqi zhuagang xue Dazhai, p.~40. Wang
  Hongwen went to the Hangzhou Gear-Box Plant with Ji on July 9. Visit,
  July 17, 1977. Wang also toured the Oxygen Generator Plant (HZRB,
  November 12, 1976), the Hangzhou No.~2 Cotton Mill on July 28 (HZRB,
  November 9,1976), the Iron and Steel Mill (HZRB, November 11, 1976)
  and the Changqing brigade in Hangzhou's Jianggan district on July 23
  (HZRB, November 22,1976).

  纪登奎最初把注意力集中在翁森鹤的丝绸厂，9天内访问了8次。1977年11月11日访问丝绸联合厂。7月25日，他访问了上王大队。1977年10月19日在上王听取情况介绍；王金友，Juqi
  zhuagang xue
  Dazhai，第40页。7月9日，王洪文随纪登奎到杭州齿轮箱厂。1977年7月17日访问。王还视察了制氧机厂（HZRB，1976年11月12日）、7月28日视察杭州第二棉纺厂（HZRB，1976年11月9日）、钢铁厂（HZRB，1976年11月11日），以及7月23日视察杭州江干区长庆大队（HZRB，1976年11月22日）。}他们还同省级当局磋商，意在改组中共浙江省委、杭州市委和浙江省军区的领导班子。王和纪带来了一个由三个部门的中央官员组成的大型代表团：中央组织部，负责调查``双突''所遵循的程序；轻工业部，负责杭州的大型纺织厂；以及管辖齿轮箱厂、制氧机厂等工厂的一机部。\footnote{HZRB,
  November 28,1976, July 4,1977. A Xinhua report of July 13,1975,
  concerning the production situation at the oxygen generator plant
  implied that one of its three workshops was out of operation. See CNS,
  578 (August 13,1975), p.~2.

  HZRB，1976年11月28日，1977年7月4日。1975年7月13日新华社关于制氧机厂生产情况的报道暗示，该厂三个车间中有一个已停产。参见
  CNS，578（1975年8月13日），第2页。}

While members of the central departments received a certain amount of
publicity in the local press, the names of Ji and Wang were never once
mentioned. For example, on 6 July 1975, Deputy-secretary of the CCP ZPC,
Xie Zhenghao, accompanied a leading member of the Ministry of Light
Industry to the silk complex where both men addressed a meeting of
workers.\footnote{HZRB, July 7,1975.

  HZRB，1975 年 7 月 7 日。} It is most likely that Ji Dengkui was
present. On 8 July ZPC Deputy-secretary Chen Weida addressed a meeting
of the employees of the Hangzhou Iron and Steel Mill.\footnote{ZJRB,
  July 9,1975, p.~1.

  ZJRB，1975年7月9日，第1页。} On 11 July Tie Ying accompanied a
``leading member'' of the central delegation to the Zhejiang Hemp Mill
where, together with Deputy-secretary of the CCP HMC Zhou Feng, they
spoke to the assembled employees.\footnote{HZRB, July 12,1975.

  HZRB，1975 年 7 月 12 日。} Despite the secrecy surrounding their
activities the city's citizens were well aware of Wang and Ji's
presence. Remarks made by the pair during their stay quickly circulated
around the city. During one briefing, Wang was told that Guo Zhisong, a
leader of the ``mountain base'' faction had always adopted the correct
position. To this remark Wang allegedly expostulated: ``crap!'' (屁话).
Posters quoting Wang to this effect aroused great merriment among the
people of Hangzhou. On a more sobering and alarming note, during his
tours of inspection Ji Dengkui warned the city's residents that if
production did not return to normal they would be forced to drink the
water of the West Lake to stay alive.\footnote{Visit to Hangzhou No.~1
  Cotton Mill, April 1,1979.

  1979年4月1日，参观杭州第一棉纺厂。}

虽然中央部门的成员在当地媒体上得到了一定的宣传，但纪和王的名字却从未被提及。例如，1975年7月6日，中共浙江省委副书记谢正浩陪同轻工部一位领导成员到丝绸联合厂，两人在工人大会上讲话。\footnote{HZRB,
  July 7,1975.

  HZRB，1975 年 7 月 7 日。}
纪登奎很可能在场。7月8日，中共浙江省委副书记陈伟达在杭州钢铁厂职工大会上讲话。\footnote{ZJRB,
  July 9,1975, p.~1.

  ZJRB，1975年7月9日，第1页。}
7月11日，铁瑛陪同中央代表团一位``领导成员''来到浙江麻厂，与中共杭州市委副书记周锋一起向在场职工讲话。\footnote{HZRB,
  July 12,1975.

  HZRB，1975 年 7 月 12 日。}
尽管他们的活动保密，但该市市民对王、纪的存在心知肚明。两人在逗留期间的言论迅速传遍全城。在一次汇报中，王被告知，``山下派''领导人郭志松的立场一直是正确的。据称，王对此斥道：``屁话！''引用王这句话的大字报使杭州民众大为乐不可支。更发人深省、更令人震惊的是，纪登奎在巡视期间警告该市居民，如果生产不恢复正常，他们将被迫喝西湖水来维持生命。\footnote{Visit
  to Hangzhou No.~1 Cotton Mill, April 1,1979.

  1979年4月1日，参观杭州第一棉纺厂。}

After more than a fortnight's investigation Wang, and possibly Ji,
returned to Beijing on 17 July 1975 to report to the Politburo. On that
same day the central leadership approved a report from the provincial
party committee on the subject of the reorganization of its standing
committee.\footnote{Zhonggong dangshi baiti, p.~387; Dangshi tongxun, 1
  (January 15, 1983), p.~10. In a speech in December 1976, Secretary of
  the CCP ZPC Chen Zuolin claimed that Mao had signalled his approval by
  circling the document. HZRB, December 18,1976. However, Xu Zhongyuan,
  the person in charge of the Chairman's reading programme in 1975,
  later reported that in the middle of July 1975 Mao had undergone an
  eye operation and for a while afterward was unable to read. See Qiu
  Dexin and Xu Zhongyuan, ``毛泽东同志的读书生活'' (The Way Comrade Mao
  Zedong Studied), Hongqi, No.~1 (1983), p.~41; translated in
  SWB/FE/7255/BII/13-14.

  《中共党史百题》，第387页；《党史通讯》，1（1983年1月15日），第10页。1976年12月，中共浙江省委书记陈作霖在一次讲话中声称，毛泽东已通过圈出该文件来表示批准。HZRB，1976年12月18日。然而，1975年负责主席读书计划的许仲源后来报告说，1975年7月中旬，毛泽东做了眼部手术，之后有一段时间无法阅读。参见邱德新、许仲源，《毛泽东同志的读书生活》，《红旗》，1983年第1期，第41页；译文见SWB/FE/7255/BII/13-14。}
Wang returned to Hangzhou on 19 July in preparation for the announcement
of the leadership shake-up on 22 July.\footnote{Bian Wen, HZRB, December
  5,1976.

  卞文，HZRB，1976 年 12 月 5 日。} On the same day as Wang's return,
PLA troops started moving into four key factories in and around
Hangzhou. Two days later, on 24 July, the Central Committee and State
Council released their joint decision concerning Zhejiang.

经过两周多的调查，王（可能还有纪）于1975年7月17日返回北京，向政治局报告。同日，中央领导批复了省委关于省委常委会改组问题的报告。\footnote{Zhonggong
  dangshi baiti, p.~387; Dangshi tongxun, 1 (January 15, 1983), p.~10.
  In a speech in December 1976, Secretary of the CCP ZPC Chen Zuolin
  claimed that Mao had signalled his approval by circling the document.
  HZRB, December 18,1976. However, Xu Zhongyuan, the person in charge of
  the Chairman's reading programme in 1975, later reported that in the
  middle of July 1975 Mao had undergone an eye operation and for a while
  afterward was unable to read. See Qiu Dexin and Xu Zhongyuan,
  ``毛泽东同志的读书生活'' (The Way Comrade Mao Zedong Studied), Hongqi,
  No.~1 (1983), p.~41; translated in SWB/FE/7255/BII/13-14.

  《中共党史百题》，第387页；《党史通讯》，1（1983年1月15日），第10页。1976年12月，中共浙江省委书记陈作霖在一次讲话中声称，毛泽东已通过圈出该文件来表示批准。HZRB，1976年12月18日。然而，1975年负责主席读书计划的许仲源后来报告说，1975年7月中旬，毛泽东做了眼部手术，之后有一段时间无法阅读。参见邱德新、许仲源，《毛泽东同志的读书生活》，《红旗》，1983年第1期，第41页；译文见SWB/FE/7255/BII/13-14。}
7月19日，王洪文返回杭州，为7月22日宣布领导班子改组做准备。\footnote{Bian
  Wen, HZRB, December 5,1976.

  卞文，HZRB，1976 年 12 月 5 日。}
王洪文返回的同一天，解放军部队开始进驻杭州及其周边地区的四个重点工厂。两天后，即7月24日，中央和国务院发布关于浙江的联合决定。

Three major meetings were in fact held on 22 July.\footnote{ZJRB, July
  23,1975, p.~1; July 24,1975, p.~1; HZRB, July 23,1975; ZPS, July 23,
  24,1975, SWB/FE/4967/BII/1-3; CNS, 578, p.~6.

  ZJRB，1975年7月23日，第1页；1975年7月24日，第1页；HZRB，1975年7月23日；ZPS，1975年7月23、24日，SWB/FE/4967/BII/1-3；CNS，578，第6页。}
The first involved 8,000 office cadres working under the CCP ZPC, who
listened to Tan Qilong and Tie Ying as well as ``the responsible
comrades of the Center'' -- a certain reference to Wang Hongwen and Ji
Dengkui -- convey instructions from Beijing. Fifteen thousand cadres
from offices of the CCP HMC and grass-roots units in Hangzhou attended
another gathering addressed by their new leaders, Zhang Zishi (张子石)
and Chen Wenshu (陈文书). Jiang Baodi, chairwoman of the ZPTUC and a
standing committee member of the CCP ZPC, was also present. Both
meetings discussed three contentious issues; the selection of new party
members and the promotion of cadres, bourgeois factionalism and the
state of the economy. Party members and cadres of the ZPMD attended the
third important meeting which was held on 22 July and heard a speech and
directives from their newly-appointed Commander and Political Commissar.
The assembled troops pledged to resolutely ``accept the centralized
leadership of the Party''. On 24 July the central authorities released a
document containing a seven-point decision regarding the situation in
Zhejiang.\footnote{``中共中央国务院关于浙江省问题的决定'' (Decision of
  the CCP CC and State Council concerning the problem of Zhejiang)
  (hereafter cited as ``Decision'') in 匪情月报 (Bandit Affairs
  Monthly), 18: 11 (1976), pp.~84-5. The document has been translated in
  Chinese Law and Government, 10: 1 (1977), pp.~8-11, and I\&S,12: 6
  (1976), pp.~102-5. Both translations suffer from varying degrees of
  inaccuracy.

  《中共中央国务院关于浙江省问题的决定》（以下简称《决定》），载《匪情月报》，1976年18：11，第84-5页。该文件已被翻译成《Chinese
  Law and Government》，10: 1 (1977)，第 8-11 页，以及 I\&S,12: 6
  (1976)，第 102-5 页。两种翻译都存在不同程度的不准确性。} The lengthy
preamble made reference to an alarming number of activities which
threatened orderly rule in the province. These included anarchist trends
of thinking (无政府主义思潮), factional struggles, armed conflicts,
interruptions to water and electricity supplies, raids on army units
(袭击部队), attacks on public security offices (冲击公安部门), looting
of state property (抢劫国家资财), murder, arson and poisoning of water
supplies, as well as the spread of counter-revolutionary propaganda and
the incitement of counter-revolutionary riots (反革命暴动).\footnote{For
  reports from Taiwan alleging similar such incidents, see World
  Anti-Communist League, China Chapter, Asian People's Anti-Communist
  League, Charts concerning Chinese Communists on the Mainland, 31st
  Series (April 1976), p.~18; 34th Series (April 1977), p.~4. In 1989
  the Chinese authorities used similar terms to describe the army
  suppression of student protestors in Beijing on the night of June 3.

  有关台湾方面指称类似事件的报道，参见世界反共联盟中国分会、亚洲人民反共联盟，《大陆中国共产党人图表》，第31辑（1976年4月），第18页；第34辑（1977年4月），第4页。1989年，中国当局使用类似术语描述6月3日晚北京军队镇压学生抗议者的事件。}
The seven resolutions briefly examined the work of provincial leaders,
work teams and the PLA, the chronic problem of factionalism as engaged
in by unnamed organizations, public disorder and the state of the
provincial economy. Even taking hyperbole into account, the contents of
the document revealed, as the national and provincial print and
electronic media only partially disclosed, the appalling state in which
government in Zhejiang had fallen.

7月22日，实际上召开了三场大型会议。\footnote{ZJRB, July 23,1975, p.~1;
  July 24,1975, p.~1; HZRB, July 23,1975; ZPS, July 23, 24,1975,
  SWB/FE/4967/BII/1-3; CNS, 578, p.~6.

  ZJRB，1975年7月23日，第1页；1975年7月24日，第1页；HZRB，1975年7月23日；ZPS，1975年7月23、24日，SWB/FE/4967/BII/1-3；CNS，578，第6页。}第一场会议有8000名中共浙江省委机关干部参加，听取谭启龙、铁瑛以及``中央负责同志''------显然指王洪文、纪登奎------传达北京的指示。中共杭州市委机关和杭州基层单位的1.5万名干部参加了另一场会议，听取新任领导人张子石、陈文书讲话。浙江省总工会主席、中共浙江省委常委姜宝娣也在场。两次会议讨论了三个有争议的问题：新党员的选拔和干部的提拔、资产阶级派性和经济状况。浙江省军区党员和干部出席了7月22日召开的第三次重要会议，听取新任司令员和政治委员的讲话和指示。集结部队宣誓坚决``接受党的集中领导''。7月24日，中央发布文件，就浙江局势作出七点决定。\footnote{``中共中央国务院关于浙江省问题的决定''
  (Decision of the CCP CC and State Council concerning the problem of
  Zhejiang) (hereafter cited as ``Decision'') in 匪情月报 (Bandit
  Affairs Monthly), 18: 11 (1976), pp.~84-5. The document has been
  translated in Chinese Law and Government, 10: 1 (1977), pp.~8-11, and
  I\&S,12: 6 (1976), pp.~102-5. Both translations suffer from varying
  degrees of inaccuracy.

  《中共中央国务院关于浙江省问题的决定》（以下简称《决定》），载《匪情月报》，1976年18：11，第84-5页。该文件已被翻译成《Chinese
  Law and Government》，10: 1 (1977)，第 8-11 页，以及 I\&S,12: 6
  (1976)，第 102-5 页。两种翻译都存在不同程度的不准确性。}冗长的序言提到了数量惊人的、威胁该省有序治理的活动。其中包括无政府主义思潮、派系斗争、武装冲突、断水断电、袭击部队、冲击公安部门、抢劫国家资财、杀人放火、投毒水源，以及散布反革命宣传、煽动反革命暴乱等。\footnote{For
  reports from Taiwan alleging similar such incidents, see World
  Anti-Communist League, China Chapter, Asian People's Anti-Communist
  League, Charts concerning Chinese Communists on the Mainland, 31st
  Series (April 1976), p.~18; 34th Series (April 1977), p.~4. In 1989
  the Chinese authorities used similar terms to describe the army
  suppression of student protestors in Beijing on the night of June 3.

  有关台湾方面指称类似事件的报道，参见世界反共联盟中国分会、亚洲人民反共联盟，《大陆中国共产党人图表》，第31辑（1976年4月），第18页；第34辑（1977年4月），第4页。1989年，中国当局使用类似术语描述6月3日晚北京军队镇压学生抗议者的事件。}七项决议简要审查了省领导、工作组和解放军的工作、未具名组织长期从事的派系主义问题、社会秩序混乱和全省经济状况等问题。即使考虑到夸张成分，该文件的内容也揭示了浙江政府已经陷入的骇人听闻的状态，而全国和省级印刷媒体和电子媒体只披露了部分内容。

\section{1. Leadership Changes}\label{leadership-changes}

\section{1. 领导层变动}\label{ux9886ux5bfcux5c42ux53d8ux52a8}

The first issue which the central authorities tackled concerned the
reconstruction of the provincial leadership in Zhejiang. It decided to
prune the ZPC standing committee by three from sixteen to thirteen
members. However, the total number of secretaries (including
deputy-secretaries) was increased by one from six to seven and in
particular the number of full secretaries (excluding the First
Secretary) from one to three (see Table Eleven). The remaining seats on
the standing committee were reduced from ten to six. Of the original
sixteen members only six, of whom four were of secretary status,
retained their places in the new leading body. The rest were either
demoted, transferred or dismissed.

中央解决的第一个问题是浙江省领导班子的重建。中央决定将中共浙江省委常委会委员由16人减至13人，减少3人。然而，书记（包括副书记）总数增加了一名，从六名增加到七名，特别是正书记（不包括第一书记）的数量从一名增加到三名（见表十一）。常务委员会的其余席位从十个减少到六个。原来的16名成员中，只有6名成员（其中4名担任书记职务）保留了新领导机构中的席位。其余的人要么被降职、调职，要么被撤职。

Despite his failure to have kept factionalism in check, Tan Qilong
retained the confidence of the central authorities and therefore his
post as First Secretary of the CCP ZPC. However, point one of the seven
point decision of 24 July seemed to imply a certain reservation about
Tan's leadership:

尽管未能遏制派系斗争，但谭启龙仍然得到中央的信任，并因此保住了中共浙江省委第一书记的职务。然而，7月24日七点决定中的第一点似乎暗示了对谭的领导有所保留：

\begin{quote}
The CCP Zhejiang Provincial Committee, the Zhejiang Provincial
Revolutionary Committee and the Party Committee of the Zhejiang
Provincial Military District headed by Tan Qilong are to carry out the
correct line and the instructions of the Party Central. The Party
Central has confidence in them. The shortcomings and mistakes which have
cropped up in their work can only be corrected through well-intentioned
(善意的) criticism and assistance.\footnote{``Decision'', p.~84.

  《决定》，第84页。}
\end{quote}

\begin{quote}
以谭启龙为首的中共浙江省委、浙江省革命委员会、浙江省军区党委，贯彻执行党中央的正确路线和指示。党中央对他们充满信心。工作中出现的缺点和错误，只有善意的批评和帮助才能改正。\footnote{``Decision'',
  p.~84.

  《决定》，第84页。}
\end{quote}

A Taiwan article quoted from the document to support its assertion that
``In the power struggle in Hangchow, \ldots{} Tan scored a triumph over
his opponents''.\footnote{``Tan Qilong (2)'', p.~111.

  《谭启龙（二）》，第111页。} In a similar vein a Hong Kong observer
argued that Tan ``rather effectively settled worker disputes in
Hangzhou''.\footnote{Ting Wang, ``Leadership Realignments'', PoC,26: 4
  (1977), p.~11.

  Ting Wang，《领导层调整》，PoC，26:4（1977年），第11页。} Perhaps this
conclusion was derived from wall posters which appeared on the streets
of Hangzhou in May 1976 accusing Tan of having crushed the city's
militia in July 1975.\footnote{Taipei Home Service, June 13,1976,
  SWB/FE/5237/BII/14-15.

  台北民政服务，1976 年 6 月 13 日，SWB/FE/5237/BII/14-15。} In his
passing reference to the event, an American analyst has concluded with
far greater justification that Tan's failure to control the factions in
his province could have been viewed unfavorably by the moderate
leadership which then prevailed in Beijing.\footnote{Wayne, ``The
  Politics of Restaffing'', p.~135.

  Wayne，《人员重组的政治》，第135页。} China News Summary correctly
pointed to Tan's vacillating leadership style in noting that

台湾的一篇文章引用该文件来支持其说法，``在杭州的权力斗争中，\ldots\ldots 谭战胜了他的对手''。\footnote{``Tan
  Qilong (2)'', p.~111.

  《谭启龙（二）》，第111页。}
同样，一位香港观察家认为谭``相当有效地解决了杭州的工人纠纷''。\footnote{Ting
  Wang, ``Leadership Realignments'', PoC,26: 4 (1977), p.~11.

  Ting Wang，《领导层调整》，PoC，26:4（1977年），第11页。}
也许这个结论来自1976年5月出现在杭州街头的大字报；这些大字报指责谭在1975年7月镇压了该市的民兵。\footnote{Taipei
  Home Service, June 13,1976, SWB/FE/5237/BII/14-15.

  台北民政服务，1976 年 6 月 13 日，SWB/FE/5237/BII/14-15。}
一位美国分析人士在提及这一事件时，得出了更有道理的结论，认为谭未能控制他所在省的派系，可能会被当时在北京占主导地位的温和领导层视为不利。\footnote{Wayne,
  ``The Politics of Restaffing'', p.~135.

  Wayne，《人员重组的政治》，第135页。}
中国新闻摘要正确地指出了谭的摇摆不定的领导风格，指出：

\begin{quote}
Since Tan is a rehabilitated veteran cadre \ldots{} he has been
unwilling to take strong measures either against the local Leftists who
have close links with their supporters in neighboring Shanghai or the
local anti-radical elements on whose support he has to
depend.\footnote{CNS, 578, p.~1.

  CNS，578，第1页。}
\end{quote}

\begin{quote}
由于谭是平反的老干部\ldots\ldots 无论是对那些与上海的支持者有密切联系的当地左翼分子，还是对他所依赖的当地反激进分子，他都不愿意采取强硬措施。\footnote{CNS,
  578, p.~1.

  CNS，578，第1页。}
\end{quote}

Tan's appearance in Hangzhou on 22 July was one of the last that he made
in Zhejiang. Although he continued to hold all his provincial posts,
after September 1975 Tan never returned to the province as its party
leader. He appeared publicly in Beijing on 8 October 1975 and on 15
October at the National Dazhai conference.\footnote{CIA Reference Aid,
  Appearances and Activities of Leading Personalities of the PRC, 1975.

  中央情报局参考援助，中华人民共和国主要人物的露面和活动，1975 年。}
Perhaps illness prevented Tan from actively participating in provincial
politics in 1976. Evidence to support this explanation comes from one of
Tan's opponents on the ZPC standing committee, Lai Keke. In January
1976, Lai returned to Hangzhou after a period of convalescence in
Guangzhou and reported in writing to Wang Hongwen that ``because many
leaders of the Zhejiang provincial committee are sick I have already
started to take up some work''.\footnote{Shen Weicai, HZRB, August 17,
  1977.

  沈伟才，HZRB，1977年8月17日。} On the other hand, an expatriate from
Hangzhou recalled that in 1975

7月22日，谭在杭州露面，这是他在浙江的最后几次露面之一。尽管他继续担任省内所有职务，但1975年9月之后，谭再也没有作为该省党的领导人回到浙江。1975年10月8日，他在北京公开露面，并于10月15日出席全国大寨会议。\footnote{CIA
  Reference Aid, Appearances and Activities of Leading Personalities of
  the PRC, 1975.

  中央情报局参考援助，中华人民共和国主要人物的露面和活动，1975 年。}
或许，1976年谭因病而无法积极参与省级政治。支持这一解释的证据来自于中共浙江省委常委中谭的反对者之一赖可可。1976年1月，赖在广州休养一段时间后回到杭州，向王洪文写信汇报``因浙江省委许多领导生病，我已经开始承担一些工作''。\footnote{Shen
  Weicai, HZRB, August 17, 1977.

  沈伟才，HZRB，1977年8月17日。}

\begin{quote}
Wang Hongwen tried to persuade Tan Qilong to split from Zhou Enlai and
join forces with the radicals. Jiang Ching {[}sic{]} also went to talk
to Tan Qilong's wife to get her to persuade her husband to switch sides
and link up with Mao. The pressure on Tan was too great so he decided to
fake an illness and retired to Shanghai. He did not return to Hangzhou
again.\footnote{``The Case of Weng Sengho'', p.~4.

  《翁森鹤案》，第4页。}
\end{quote}

\begin{quote}
王洪文试图劝说谭启龙与周恩来决裂，与激进分子联手。江青还去找谭启龙的妻子谈话，让她劝说她的丈夫倒戈与毛联手。谭的压力太大，于是决定假装生病，避往上海。他没有再回到杭州。\footnote{``The
  Case of Weng Sengho'', p.~4.

  《翁森鹤案》，第4页。}
\end{quote}

It is true that over the previous two years Tan had managed, with some
success, to straddle the fence in the struggle for ascendancy between
the Cultural Revolution radicals and the veteran cadres. As China News
Summary wrote about Jiangxi first secretary Jiang Weiqing, ``This
{[}search for the middle ground{]} may be typical of the attitude of
rehabilitated cadres like Jiang toward the current
situation''.\footnote{CNS, 568 (May 28, 1975), p.~3.

  CNS，568（1975年5月28日），第3页。} During the campaign against Deng
Xiaoping in 1976 Jiang Weiqing may have ``decided to play possum for
some time''\footnote{Jurgen Domes, ``China in 1976: Tremors of
  Transition'', AS, 17:1 (1977), p.~9.

  Jurgen
  Domes，《1976年的中国：转型的震颤》，AS，17:1（1977年），第9页。} but
Tan had removed himself from the political scene sometime in the last
quarter of 1975, before the commencement of the campaign. The list of
provincial leaders who attended Zhou Enlai's funeral on 15 January 1976
provided the sole indication that Tan in fact formally continued to
occupy his posts in Zhejiang.\footnote{RMRB, January 16,1976, p.~3. Tan
  also attended Mao's funeral in the national capital in September 1976.

  RMRB，1976年1月16日，第3页。1976年9月，谭还参加了在首都举行的毛泽东葬礼。}

诚然，在过去的两年里，谭已经相当成功地在文革激进派和老干部之间争夺优势的斗争中脚踩两边。正如《中国新闻摘要》在谈到江西省第一书记江渭清时所写的那样，``这（寻求中间立场）可能是像江这样的平反干部对当前形势的典型态度''。\footnote{CNS,
  568 (May 28, 1975), p.~3.

  CNS，568（1975年5月28日），第3页。}
在1976年反邓小平运动中，江渭清可能``决定暂时装病避风头''\footnote{Jurgen
  Domes, ``China in 1976: Tremors of Transition'', AS, 17:1 (1977),
  p.~9.

  Jurgen
  Domes，《1976年的中国：转型的震颤》，AS，17:1（1977年），第9页。}，但谭已于1975年最后一个季度的某个时候，在这场运动开始之前退出了政治舞台。1976年1月15日出席周恩来葬礼的省级领导名单，是唯一表明谭实际上仍正式担任浙江职务的迹象。\footnote{RMRB,
  January 16,1976, p.~3. Tan also attended Mao's funeral in the national
  capital in September 1976.

  RMRB，1976年1月16日，第3页。1976年9月，谭还参加了在首都举行的毛泽东葬礼。}

Following his harrowing experience in the 1974 campaign, Tie Ying had
been allowed to return to duty in February 1975, possibly as a result of
decisions thrashed out at the Beijing conference on Zhejiang held in
late 1974. The previous chapter of this study has cited Tie's strong
endorsement of Chen Yonggui's denunciation of ``bourgeois factionalism''
in his concluding speech to the April 1975 provincial Dazhai conference.
In July 1975 Tie retained his posts as deputy to Tan in the CCP ZPC,
ZPRC, and party committee of the ZPMD. This was despite the reference to
his ``problem'' (问题) in the central decision of 24 July. Point three
of the document read in part:

经历了1974年运动的惨痛遭遇后，铁瑛于1975年2月被允许重返工作岗位，这可能是1974年底北京浙江会议作出决定的结果。本研究的前一章引用了铁瑛在1975年4月省大寨会议总结发言中对陈永贵谴责``资产阶级派别主义''的强烈支持。1975年7月，铁瑛继续在中共浙江省委、浙江省革命委员会和浙江省军区党委担任谭的副手。尽管7月24日的中央决定提到了他的``问题''，情况仍是如此。该文件的第三点部分写道：

\begin{quote}
Comrade Tie Ying's problem is one of contradictions among the people and
should be solved through the method of unity-criticism-unity.\footnote{``Decision'',
  p.84.

  《决定》，第84页。}
\end{quote}

\begin{quote}
铁瑛同志的问题是人民内部矛盾之一，应该用团结、批评、团结的方法来解决。\footnote{``Decision'',
  p.84.

  《决定》，第84页。}
\end{quote}

Thus the center's treatment of Tie in the ``decision'', far from being
severe, as one Taiwan source has suggested,\footnote{``Tan Qilong (2)'',
  p.110.

  《谭启龙（二）》，第110页。} was in fact a rebuff to the rebels who
had attempted to obtain a much harsher verdict.

因此，中央在``决定''中对铁瑛的处理，远非一位台湾消息人士所说的那样严厉，实际上是对试图获得更严厉判决的造反派的回绝。\footnote{``Tan
  Qilong (2)'', p.110.

  《谭启龙（二）》，第110页。}

Wayne correctly interpreted Tie Ying's position thus: ``In 1975 Tie was
apparently criticized for his views on how to handle severe factional
problems but was cleared by the final orders from Peking \ldots{}''. He
added that ``Tie represents the more order-oriented forces in Zhejiang''
compared to Tan Qilong's ``hesitancy to be heavy-handed with `leftists'
in the early 1970s''.\footnote{Wayne, ``The Politics of Restaffing'',
  p.~136.

  Wayne，《人员重组的政治》，第136页。} Previous chapters of this work
have illustrated the accuracy of Wayne's incisive comments. Certainly,
the three rebel leaders in Zhejiang believed that they could manipulate
Tan more readily than military leaders such as Tie Ying and Xia Qi who
were prepared, with Xu Shiyou's backing, to confront them head on and,
in conjunction with the ``mountain base'' faction, to organize
alternative militia units to counter the urban militia command. In
addition, Tie did not suffer from the same liability as Tan in having
previously been labelled a ``capitalist-roader'' at the beginning of the
Cultural Revolution. Those who had suffered this ignominy were,
according to one Chinese source, severely constrained, particularly in
periods of radical upheaval such as 1973-74. They were forced to be
``prudent and circumspect'' (谨小慎微) and could not afford to let
others accuse them of attempting to ``reverse the verdict''.\footnote{Zhao
  Wei, Zhao Ziyang zhuan, p.~194.

  赵伟，《赵紫阳传》，第194页。}

韦恩正确地解释了铁瑛的立场：``1975年，铁瑛显然因其关于如何处理严重派系问题的观点而受到批评，但北京的最终命令使他获得澄清\ldots\ldots''。他补充说，与谭启龙``在1970年代初犹豫是否对`左派'采取严厉手段''相比，``铁瑛代表了浙江更重视秩序的力量''。\footnote{Wayne,
  ``The Politics of Restaffing'', p.~136.

  Wayne，《人员重组的政治》，第136页。}
本书的前几章已经说明了韦恩精辟评论的准确性。当然，浙江的三位造反派领导人相信，操纵谭要比操纵铁瑛、夏琦等军方领导容易得多，后者在许世友的支持下，准备与他们正面对抗，并与``山下派''联合组织替代性民兵部队，对抗城市民兵指挥部。此外，铁瑛也不像谭那样在文革初期曾被贴上``走资派''的标签。据一位中国消息人士称，那些遭受这种耻辱的人受到严重限制，特别是在1973-74年等激进动荡时期。他们被迫``谨小慎微''，不能让别人指责他们试图``翻案''。\footnote{Zhao
  Wei, Zhao Ziyang zhuan, p.~194.

  赵伟，《赵紫阳传》，第194页。}

One of the other two ZPC secretaries was the longest-serving provincial
party leader remaining in Zhejiang and the other was a newcomer to the
province. Very early in the Cultural Revolution Lai Keke had thrown in
his lot with the new order. He had been rewarded with the posts of
Vice-chairman of the ZPRC in 1968 and Deputy-secretary of the CCP ZPC in
1971. Lai had kept a very low profile in the 1974 campaign despite his
initial enthusiasm for the philosophy of going against the tide. In July
1975, most probably at the urging of Wang Hongwen, he was promoted from
the post of Deputy-secretary to full secretary of the ZPC.\footnote{``Tan
  Qilong'', p.~108, states incorrectly that Lai's promotion had occurred
  in April 1973.

  《谭启龙》，第108页，错误地指出赖的晋升发生在1973年4月。}

另外两名中共浙江省委书记中，一位是在浙江任职时间最长的省委领导，另一位是该省的新人。文化大革命初期，赖可可就投身于新秩序。1968年，他被任命为浙江省革命委员会副主任，1971年被任命为中共浙江省委副书记。尽管他最初热衷于逆潮流而动，但他在1974年的运动中一直保持着非常低调的姿态。1975年7月，很可能是在王洪文的敦促下，他从中共浙江省委副书记晋升为正书记。\footnote{``Tan
  Qilong'', p.~108, states incorrectly that Lai's promotion had occurred
  in April 1973.

  《谭启龙》，第108页，错误地指出赖的晋升发生在1973年4月。}

The other new secretary was Zhang Wenbi (张文碧), whom the central
authorities had appointed to the onerous position of Commander of the
ZPMD. During the Cultural Revolution Zhang had served as Political
Commissar of main force Unit 6408 which, under its Commander Li Desheng,
had been stationed in Anhui province. In October 1972 Zhang had been
appointed Minister of Water Conservancy and Electric Power but was not
reappointed to the position at the 4th NPC in January 1975. Perhaps the
center had by that time already marked him down for the highly sensitive
and demanding job in Zhejiang. Zhang had previous connections with Tan
Qilong dating back to the 1940s when Tan had served as Political
Commissar and party secretary to the communist guerrilla forces in
eastern Zhejiang. Zhang had served as Director of the political
department in the same forces.\footnote{See
  浙江省党史资料征集研究委员会浙江省档案馆编 (Zhejiang Province Party
  History Materials Collection and Research Committee and Zhejiang
  Province Archives, eds), 浙江抗日根据地 (Eastern Zhejiang
  Anti-Japanese Base Area), (Beijing: CCP History Materials Publishers,
  1987).

  参见《浙江省党史资料征集研究委员会浙江省档案馆编》（浙江省党史资料征集研究委员会、浙江省档案馆编）、《浙江抗日根据地》（北京：中共党史资料出版社，1987）。}

另一位新任书记是张文碧，中央任命他担任浙江省军区司令员这一重任。文化大革命期间，张曾任驻安徽省、由李德生指挥的6408部队政委。1972年10月，张被任命为水利电力部部长，但在1975年1月的第四届全国人民代表大会上没有再次担任该职务。也许当时中央已经将他安排在浙江从事这项高度敏感且要求很高的工作。张与谭启龙的渊源可以追溯到1940年代，当时谭启龙担任浙东共产党游击队政委兼党委书记。张曾在同一支部队担任政治部主任。\footnote{See
  浙江省党史资料征集研究委员会浙江省档案馆编 (Zhejiang Province Party
  History Materials Collection and Research Committee and Zhejiang
  Province Archives, eds), 浙江抗日根据地 (Eastern Zhejiang
  Anti-Japanese Base Area), (Beijing: CCP History Materials Publishers,
  1987).

  参见《浙江省党史资料征集研究委员会浙江省档案馆编》（浙江省党史资料征集研究委员会、浙江省档案馆编）、《浙江抗日根据地》（北京：中共党史资料出版社，1987）。}

Two naval cadres, Xie Zhenghao and Chai Qikun, were removed from their
posts as deputy-secretaries of the CCP ZPC in July 1975, thus
eliminating the vestiges of the high-level PLA presence in the
provincial leadership which dated back to 1967. Both men had become
intimately involved in military ``support the left'' work during the
Cultural Revolution. Xie had also served as Chairman of the Shaoxing
District Revolutionary Committee. He made his final public appearance in
Hangzhou in early July 1975 and reappeared in December of the same year
as Deputy-commander of the East China Fleet. Chai Qikun had held the
additional post of Chairman of the Ningbo District Revolutionary
Committee. He had played a highly active role in the 1974 campaign and
after May 1975 disappeared from the political stage altogether. Chai,
like Lai Keke, had failed to show solidarity with Tan and Tie's
leadership when it had come under heavy attack late in 1973, but, unlike
Lai, Chai had no central leader of the stature of Wang Hongwen to
protect or support him. Additionally, both Xie and Chai, as military
officials, fell victim to the conscious attempt to civilianize the
provincial leadership and return the military to the barracks. This
reflected a national trend. By 10 October 1975 the percentage of
military cadres among provincial secretaries had fallen from 45\% at the
beginning of the year to 34\%.\footnote{CNS, 598 (January 14,1976),
  p.~2.

  CNS，598（1976年1月14日），第2页。}

1975年7月，两名海军干部谢正浩和柴启琨被免去中共省委副书记的职务，从而消除了自1967年以来解放军高层在省领导层中存在的痕迹。两人在文革期间都密切参与了军队的``支左''工作。谢还曾担任绍兴地区革命委员会主任。1975年7月上旬，他在杭州最后一次公开露面，并于同年12月再次露面，担任华东舰队副司令员。柴启琨兼任宁波地区革委会主任。他在1974年的政治运动中发挥了非常积极的作用，1975年5月后完全从政治舞台消失。柴和赖可可一样，在1973年底谭、铁领导层遭到猛烈攻击时，未能表现出与他们团结一致，但与赖不同的是，柴没有王洪文这样的中央领导人来保护或支持他。此外，谢和柴作为军方官员，都成为省级领导文职化和军队回归军营的有意识尝试的受害者。这反映了一种全国趋势。到1975年10月10日，省委书记中军队干部的比例从年初的45\%下降到34\%。\footnote{CNS,
  598 (January 14,1976), p.~2.

  CNS，598（1976年1月14日），第2页。}

To join Chen Weida, Beijing appointed as additional deputy-secretaries
two civilian cadres. Military representation on the secretariat was
thereby reduced from three out of six to two out of seven.\footnote{Robert
  Scalapino has incorrectly described all four secretaries of the CCP
  ZPC (including Tan Qilong) as military cadres. Scalapino,``The CCP's
  Provincial Secretaries'', p.~27. By Scalapino's own definition,
  neither Tan nor Lai Keke fit into this category.

  罗伯特·斯卡拉皮诺错误地将中共浙江省委四位书记（包括谭启龙）全部描述为军事干部。斯卡拉皮诺，《中共的省级书记》，第27页。根据斯卡拉皮诺自己的定义，谭和赖可可都不属于这一类。}
The most senior of the three deputy-secretaries was Luo Yi (罗毅), who
was transferred from Shanghai where he had been a standing committee
member of that city's revolutionary committee since November
1972.\footnote{CNS, 446; China Directory, 1977 (Tokyo: Radiopress Inc.,
  1977).

  CNS，446；中国指南，1977（东京：Radiopress Inc.，1977）。} In the
Cultural Revolution, Zhang Chunqiao, 1st Secretary of the CCP Shanghai
Municipal Committee, appointed Luo responsible member of Shanghai's
Culture and Education Group (文教组).\footnote{Chu Yang, ``Selling
  himself and frenziedly attacking the party'', HZRB, November 2, 1977.

  褚阳，《出卖自己，疯狂攻击党》，HZRB，1977年11月2日。} Luo had worked
on the national executive of the CYL in the 1950s and early 1960s and
his promotion to Zhejiang was most probably at the suggestion of Wang
Hongwen and Zhang Chunqiao. The most junior of the three secretaries,
Chen Zuolin (陈作霖), was another newcomer to Zhejiang. Chen was
transferred from neighboring Anhui province where he had served as a
member of the standing committee of the CCP provincial committee since
1971.\footnote{See Anhui Shengqing, p.~183; CNS, 535 (September
  18,1974). For an account of Chen's career see Keith Forster, ``Ch'en
  Tso-lin -- Secretary of the CCP's Central Discipline Inspection
  Commission'', I\&S, 23:4 (April 1987), pp.~116-23.

  参见《安徽省情》，第183页；CNS，535（1974年9月18日）。有关陈的生平介绍，参见
  Keith
  Forster，《陈作霖------中共中央纪律检查委员会书记》，I\&S，23:4（1987年4月），第116-123页。}

为了与陈伟达一起，北京又任命了两名文职干部担任副书记。书记处的军事代表从六人中的三人减少到七人中的两人。\footnote{Robert
  Scalapino has incorrectly described all four secretaries of the CCP
  ZPC (including Tan Qilong) as military cadres. Scalapino,``The CCP's
  Provincial Secretaries'', p.~27. By Scalapino's own definition,
  neither Tan nor Lai Keke fit into this category.

  罗伯特·斯卡拉皮诺错误地将中共浙江省委四位书记（包括谭启龙）全部描述为军事干部。斯卡拉皮诺，《中共的省级书记》，第27页。根据斯卡拉皮诺自己的定义，谭和赖可可都不属于这一类。}
三名副书记中资历最高的是罗毅，他是从上海调来的，自1972年11月起一直担任该市革命委员会常委。\footnote{CNS,
  446; China Directory, 1977 (Tokyo: Radiopress Inc., 1977).

  CNS，446；中国指南，1977（东京：Radiopress Inc.，1977）。}
文化大革命中，中共上海市委第一书记张春桥任命罗毅为上海文教组负责人。\footnote{Chu
  Yang, ``Selling himself and frenziedly attacking the party'', HZRB,
  November 2, 1977.

  褚阳，《出卖自己，疯狂攻击党》，HZRB，1977年11月2日。}
罗曾在20世纪50年代和1960年代初期担任共青团全国执行委员，他升任浙江很可能是在王洪文和张春桥的建议下。三位书记中资历最浅的陈作霖是浙江的另一位新人。陈是从邻近的安徽省调来的，他自1971年以来一直担任中共省委常委。\footnote{See
  Anhui Shengqing, p.~183; CNS, 535 (September 18,1974). For an account
  of Chen's career see Keith Forster, ``Ch'en Tso-lin -- Secretary of
  the CCP's Central Discipline Inspection Commission'', I\&S, 23:4
  (April 1987), pp.~116-23.

  参见《安徽省情》，第183页；CNS，535（1974年9月18日）。有关陈的生平介绍，参见
  Keith
  Forster，《陈作霖------中共中央纪律检查委员会书记》，I\&S，23:4（1987年4月），第116-123页。}

Of the six secretaries of the ZPC (excluding Tan) only Zhang Wenbi and
Chen Weida were not appointed to the additional positions of
vice-chairmen of ZPRC. Zhang obviously had enough on his plate as
Commander of the ZPMD and his omission from a government post was an
attempt to keep military interference in the civilian administration to
a minimum. The failure to appoint Chen to this position may have
reflected a fear on the part of the Beijing radicals that his past
experience in the Zhejiang administration could impede their future
takeover of the province.

中共浙江省委六名书记（不包括谭）中，只有张文碧和陈伟达未被任命为浙江省革命委员会副主任。显然，作为浙江省军区司令员，张的任务已经够重了，他未兼任政府职务，是为了把军方对民政管理的干预降至最低。未能任命陈担任这一职务，可能反映出北京激进分子担心他过去在浙江行政系统的经历会妨碍他们未来接管该省。

Only two other members of the ZPC standing committee survived the July
shake-up. They were Chen Bing and Jiang Baodi, the former a veteran
cadre and the latter a worker-rebel of the Cultural Revolution. The
remaining eight standing committee members, at least four of whom were
military officials, were removed from this body. Zhu Quanlin, Chief of
Staff of Unit 6409, made his last public appearance at the beginning of
April 1975 which suggests that he and the rest of his unit were
transferred from Zhejiang sometime in the first half of 1975.\footnote{See
  Wayne, ``The Politics of Restaffing'', p.~136.

  参见 Wayne，《人员重组的政治》，第136页。} It is most likely that the
20th Army swapped places with the 1st Army in Henan province and that Ji
Dengkui may have had a hand in this exchange of armies.

中共浙江省委常委会中只有另外两名成员在7月份的改组中留任。他们是陈兵和蒋宝娣，前者是老干部，后者是文革工人造反派。其余八名常委（其中至少四名是军方官员）被从该机构中除名。6409部队参谋长朱全林最后一次公开露面是在1975年4月上旬，这表明他和部队其他成员是在1975年上半年的某个时候从浙江调走的。\footnote{See
  Wayne, ``The Politics of Restaffing'', p.~136.

  参见 Wayne，《人员重组的政治》，第136页。}
很可能是第20军与驻河南省的第1军调换了位置，纪登奎可能参与了这次调军。

If the transfer of the 20th Army did indeed occur, then a major factor
for instability in the province had finally departed the scene. Unit
6409 had made too many enemies and had become too closely involved in
factional politics and the provincial administration for too long a time
to be trusted as a reliable instrument in suppressing factional groups.
Still less did it retain any credibility as a neutral force when the
center decided to move troops into the factories of Hangzhou. An article
from Taiwan claimed that the 20th Army, together with the 5th Air Force
and local garrison troops, had secretly supported factional conflict in
Zhejiang.\footnote{See ``Tan Qilong'', p.~103.

  参见《谭启龙》，第103页。} This claim was undoubtedly correct.

如果确实发生了第20军的调动，那么该省的一个不稳定因素终于消失了。
6409部队树敌太多，而且太长时间地卷入派系政治和省政府，不能被视为镇压派系团体的可靠工具。当中央决定将军队调入杭州工厂时，它作为中立力量的信誉就更差了。台湾有文章称，第20军与第5航空队和地方警备部队一起暗中支持浙江的派系冲突。\footnote{See
  ``Tan Qilong'', p.~103.

  参见《谭启龙》，第103页。}这个说法无疑是正确的。

Another military cadre who was dropped from the ZPC standing committee
was Deputy Political Commissar of the ZPMD, Xia Qi. In 1974 Xia had
experienced even greater tribulations than Tie Ying, having been
suspended from his posts and charged with the most serious offences.
Although Xia lost his party post in July 1975, the CC document's
assessment of him, as with the comments relating to Tie Ying, amounted
to a partial rehabilitation, particularly in comparison to the verdict
passed on Xia by the triple plenum special investigation group in August
1974 (see chapter seven). Point three of the Decision stated that

另一位被免去中共浙江省委常委职务的军事干部是浙江省军区副政委夏琦。1974年，夏经历了比铁瑛更大的磨难，被停职并被指控犯有最严重的罪行。尽管夏在1975年7月失去了党内职务，但中央文件对他的评价，以及有关铁瑛的评论，相当于部分平反，特别是与1974年8月三全会特别调查组对夏的判定相比（见第七章）。《决定》第三点指出

\begin{quote}
As long as he does not persist in his mistakes or continue along the
incorrect line, even Comrade Xia Qi should be helped according to the
principles of ``learning from past mistakes to avoid future ones'' and
``cure the illness to save the patient.'' In this way we can distinguish
right from wrong and unite the comrades so as to achieve solidarity
within the Party.\footnote{``Decision'', p.~84. Despite this
  reappraisal, Xia did not reappear publicly in his post as
  deputy-political commissar of the ZPMD until June 1976. See ZPS, June
  16,1976, SWB/FE/5238/BII/5-6.

  《决定》，第84页。尽管有这次重新评价，夏直到1976年6月才再次以浙江省军区副政委身份公开露面。见ZPS，1976年6月16日，SWB/FE/5238/BII/5-6。}
{[}my emphasis{]}
\end{quote}

\begin{quote}
只要不坚持错误、不继续沿着不正确的路线走下去，即使是夏琦同志，也应该按照``惩前毖后''和``治病救人''的原则来帮助他。这样才能明辨是非，团结同志，达到党内团结。\footnote{``Decision'',
  p.~84. Despite this reappraisal, Xia did not reappear publicly in his
  post as deputy-political commissar of the ZPMD until June 1976. See
  ZPS, June 16,1976, SWB/FE/5238/BII/5-6.

  《决定》，第84页。尽管有这次重新评价，夏直到1976年6月才再次以浙江省军区副政委身份公开露面。见ZPS，1976年6月16日，SWB/FE/5238/BII/5-6。}
{[}我的强调{]}
\end{quote}

The dropping of Xia Qi from the ZPC standing committee was in line with
the policy of civilianization of the provincial leadership. It was also
perhaps a gesture to the radicals who had pursued him so relentlessly in
1974. However, Xia's dismissal was more than balanced by the dismissal
of his persecutors, chief amongst whom was Shen Ce.

夏琦被免去中共浙江省委常委职务，符合省领导文职化的方针。这或许也是向1974年无情追究他的激进分子作出的一种姿态。然而，夏的免职远被迫害他的人被免职所抵消，其中为首的是沈策。

Of the four additional members to the standing committee of the CCP ZPC,
three were outsiders to Zhejiang. The only new appointee who had
previously worked in the province was Lü Jianguang, who had been
Director of the provincial Public Security Bureau prior to the Cultural
Revolution.\footnote{According to CNS, 583, p.~1, Lü had been a ``close
  follower'' of Luo Ruiqing, the pre-Cultural Revolution Minister of
  Public Security.

  据 CNS，583，第1页，吕曾是文革前公安部部长罗瑞卿的``亲密追随者''。} Lü
had established such a fearsome reputation that the Red Guards had
dubbed him Lü the king of hell (吕阎王). In Deng Xiaoping's eyes, his
return to office would have been expected to stiffen the resolve of the
demoralized public security forces and to have placed a certain amount
of fear into the hearts of the rebel warriors who had attacked him so
ferociously in 1968.\footnote{CNS, 583, p.~8, makes this point.

  CNS，583，第8页，提出了这一点。} A PLA cadre, Wu Shihong (武世鸿), was
appointed Political Commissar of the ZPMD in addition to joining the ZPC
standing committee. It was claimed at the time that the commander of the
Nanjing Military Region lent Tan two political cadres (a reference to Wu
and Li Bincheng) ``to ensure tight control of his troops''.\footnote{Goodstadt,
  FEER, August 15,1975, p.~25.

  Goodstadt，FEER，1975年8月15日，第25页。} In the first place, Li
Bincheng had been transferred to Hangzhou in June 1974 and secondly,
Zhang Chunqiao and Ding Sheng, the leaders of the Nanjing Military
Region, were anything but sympathetic to Tan's position. The second
newcomer to the ZPC standing committee, Feng Ke (冯恪), a small
grey-haired civilian cadre aged in his sixties, had previously served as
a member of the CCP Beijing Municipal Committee since 1971.\footnote{CNS,
  583, p.~7. In 1977 Feng Ke and my wife and I were often the only
  guests in the spacious dining-room of the Hangzhou Hotel. It appeared
  that Feng's living arrangements had suffered some inconvenience. In
  1964, as party secretary of the Beijing motor car plant, he had
  published an article on ideological and political work in the People's
  Daily. RMRB, May 12, 1964, transl. in SCMP, 3230, pp.~3-5.

  CNS，583，第7页。1977年，杭州饭店宽敞的餐厅里常常只有冯恪、我和我的妻子。看来冯的生活安排遇到了一些不便。1964年，他担任北京汽车制造厂党委书记时，曾在《人民日报》发表关于思想政治工作的文章。RMRB，1964年5月12日；译文载
  SCMP，3230，第3-5页。}

新增的四名中共浙江省委常委中，有三名来自浙江省外。唯一曾在该省工作过的新任命者是文革前任省公安厅厅长的吕剑光。\footnote{According
  to CNS, 583, p.~1, Lü had been a ``close follower'' of Luo Ruiqing,
  the pre-Cultural Revolution Minister of Public Security.

  据 CNS，583，第1页，吕曾是文革前公安部部长罗瑞卿的``亲密追随者''。}吕剑光的名声如此可怕，以至于红卫兵们称他为``吕阎王''。在邓小平看来，他的复出本应坚定士气低落的公安部队的决心，并使1968年猛烈攻击他的造反派战士心中产生一定恐惧。\footnote{CNS,
  583, p.~8, makes this point.

  CNS，583，第8页，提出了这一点。}解放军干部武世鸿除了加入中共浙江省委常委会外，还被任命为浙江省军区政委。据称，当时南京军区司令员借给谭两名政治干部（指吴和李斌程）``以确保对他的部队的严格控制''。\footnote{Goodstadt,
  FEER, August 15,1975, p.~25.

  Goodstadt，FEER，1975年8月15日，第25页。}首先，李斌程早在1974年6月就已调到杭州；其次，南京军区领导人张春桥和丁盛绝非同情谭的处境。中共浙江省委常委会第二位新成员冯恪，是一位60多岁、头发花白、身材矮小的文职干部，1971年起一直担任中共北京市委委员。\footnote{CNS,
  583, p.~7. In 1977 Feng Ke and my wife and I were often the only
  guests in the spacious dining-room of the Hangzhou Hotel. It appeared
  that Feng's living arrangements had suffered some inconvenience. In
  1964, as party secretary of the Beijing motor car plant, he had
  published an article on ideological and political work in the People's
  Daily. RMRB, May 12, 1964, transl. in SCMP, 3230, pp.~3-5.

  CNS，583，第7页。1977年，杭州饭店宽敞的餐厅里常常只有冯恪、我和我的妻子。看来冯的生活安排遇到了一些不便。1964年，他担任北京汽车制造厂党委书记时，曾在《人民日报》发表关于思想政治工作的文章。RMRB，1964年5月12日；译文载
  SCMP，3230，第3-5页。}

The third outsider appointed to the standing committee was Zhang Zishi,
son of Kang Sheng,\footnote{See the photograph of Zhang standing next to
  his step-mother at Kang's funeral in December 1975. HZRB, December
  22,1975, p.~1.

  参见1975年12月康生葬礼上张站在继母旁边的照片。HZRB，1975年12月22日，第1页。}
who also replaced Wang Zida as First Secretary of the CCP HMC. Wang and
his deputy Wang Xing had turned a blind eye to the factionalism in the
city since the end of 1973. On behalf of the CCP HMC, they had meekly
submitted to the extraordinary demands and ultimatums presented by Zhang
Yongsheng, Weng Senhe and He Xianchun in 1974 and they had allowed the
prestige and authority of the municipal administration to decline to a
dangerously low level. Because the central and provincial leaderships
judged that both men could no longer play any useful role in Zhejiang,
they were transferred out of the province. After their final public
appearances in Zhejiang in April 1975, Wang Zida was transferred to the
Guangxi Autonomous Region\footnote{CNS, 572 (June 18,1975) wrongly
  claimed that Wang was a supporter of Deng Xiaoping. The error was
  atypical of the high standard of the analysis which prevailed in the
  journal.

  CNS，572（1975年6月18日）错误地声称王是邓小平的支持者。该错误与该杂志普遍采用的高标准分析不同。}
and Wang Xing to Sichuan province.\footnote{HZRB, April 5, 1978.

  HZRB，1978 年 4 月 5 日。}

第三位被任命为常委的外来干部是康生之子张子石，\footnote{See the
  photograph of Zhang standing next to his step-mother at Kang's funeral
  in December 1975. HZRB, December 22,1975, p.~1.

  参见1975年12月康生葬礼上张站在继母旁边的照片。HZRB，1975年12月22日，第1页。}他还取代王子达出任中共杭州市委第一书记。王子达及其副手王兴自1973年底以来，对本市的派系斗争视而不见。他们代表中共杭州市委，温顺地接受了张永生、翁森鹤、贺贤春1974年提出的非常规要求和最后通牒，使市政管理的威望和权威下降到危险的低水平。由于中央和省领导认为两人已不能再在浙江发挥任何有益作用，因此将他们调出省外。1975年4月在浙江最后一次公开露面后，王子达被调到广西自治区\footnote{CNS,
  572 (June 18,1975) wrongly claimed that Wang was a supporter of Deng
  Xiaoping. The error was atypical of the high standard of the analysis
  which prevailed in the journal.

  CNS，572（1975年6月18日）错误地声称王是邓小平的支持者。该错误与该杂志普遍采用的高标准分析不同。}，王兴被调到四川省。\footnote{HZRB,
  April 5, 1978.

  HZRB，1978 年 4 月 5 日。}

Zhang Zishi's credentials for such a responsible position would not have
impressed Deng Xiaoping. In fact he was the antithesis of the plodding,
obedient, hard-working bureaucrat favored by the CCP leadership. Zhang
had rebelled in the early stages of the Cultural Revolution. Before then
he had been a secondary school principal in the Shandong coastal city of
Qingdao. On 9 January 1967 Wang Xiaoyu, then Mayor of Qingdao and one of
the first party officials to rebel against his superiors (Wang later
became Chairman of the Shandong Provincial Revolutionary Committee),
sent Zhang to Beijing to find out from his father what was going on
(摸底). Kang Sheng received his son in the Great Hall of the People and
told him to inform Wang that he should seize power
immediately.\footnote{See Jin Chunming, ``Wenhua dageming'', p.~109;
  Zhong Kan, Kang Sheng pingzhuan, p.~400.

  参见金春明，《文化大革命》，第109页；钟侃，《康生评传》，第400页。}
Zhang later became Secretary-general of the CCP Shandong Provincial
Committee and a member of the Shandong Provincial Revolutionary
Committee standing committee. Zhang therefore had an impeccable
background in the eyes of central radical leaders such as Wang Hongwen.

张子石担任如此重要职位的资历不会给邓小平留下深刻印象。事实上，他与中共领导层所青睐的那种迟钝、顺从、勤奋的官僚截然相反。张在文化大革命初期曾造反。在此之前，他曾担任山东沿海城市青岛的中学校长。1967年1月9日，时任青岛市市长、首批反抗上级的党内官员之一的王晓禹（王后来成为山东省革命委员会主任），派张到北京从他父亲那里了解情况。康生在人民大会堂接见儿子，让他转告王，应立即夺权。\footnote{See
  Jin Chunming, ``Wenhua dageming'', p.~109; Zhong Kan, Kang Sheng
  pingzhuan, p.~400.

  参见金春明，《文化大革命》，第109页；钟侃，《康生评传》，第400页。}张后来担任中共山东省委秘书长和山东省革命委员会常委。因此，张在王洪文等中央激进领导人眼中的背景无可挑剔。

TABLE ELEVEN The Standing Committee of the CCP ZPC before and after the
Reorganization of July 1975

表十一 1975年7月改组前后的中共浙江省委常委会

{\def\LTcaptype{none} % do not increment counter
\begin{longtable}[]{@{}lll@{}}
\toprule\noalign{}
& BEFORE & AFTER \\
\midrule\noalign{}
\endhead
\bottomrule\noalign{}
\endlastfoot
First Secretary: & Tan Qilong 4 & Tan Qilong 4 \\
Secretaries: & Tie Ying 3 & Tie Ying 3 \\
& & Lai Keke 4,5 \\
& & Zhang Wenbi 2,3 \\
Deputy-secretaries: & Lai Keke 4,5 & Luo Yi 2,5 \\
& Xie Zhenghao 1,3 & Chen Weida 4 \\
& Chai Qikun 1,3,5 & Chen Zuolin 2 \\
& Chen Weida 4 & \\
Members: & Wang Zida 1,4 & Wu Shihong 2,3 \\
& Shen Ce 1,4,5 & Zhang Zishi 2,5 \\
& Dai Kelin 1,3,4,5 & Chen Bing 4 \\
& Zhu Quanlin 1,3,5 & Jiang Baodi 5 \\
& Meng Zhaoyu 1,3 & Feng Ke 2 \\
& Jiang Baodi 5 & Lü Jianguang 4 \\
& Zhou Jianren 1,4 & \\
& Chen Bing 4 & \\
& Xia Qi 1,3 & \\
& Liu Ang 1,3 & \\
\end{longtable}
}

{\def\LTcaptype{none} % do not increment counter
\begin{longtable}[]{@{}lll@{}}
\toprule\noalign{}
& 之前 & 之后 \\
\midrule\noalign{}
\endhead
\bottomrule\noalign{}
\endlastfoot
第一书记： & Tan Qilong 4 & Tan Qilong 4 \\
书记： & Tie Ying 3 & Tie Ying 3 \\
& & Lai Keke 4,5 \\
& & Zhang Wenbi 2,3 \\
副书记： & Lai Keke 4,5 & Luo Yi 2,5 \\
& Xie Zhenghao 1,3 & Chen Weida 4 \\
& Chai Qikun 1,3,5 & Chen Zuolin 2 \\
& Chen Weida 4 & \\
委员： & Wang Zida 1,4 & Wu Shihong 2,3 \\
& Shen Ce 1,4,5 & Zhang Zishi 2,5 \\
& Dai Kelin 1,3,4,5 & Chen Bing 4 \\
& Zhu Quanlin 1,3,5 & Jiang Baodi 5 \\
& Meng Zhaoyu 1,3 & Feng Ke 2 \\
& Jiang Baodi 5 & Lü Jianguang 4 \\
& Zhou Jianren 1,4 & \\
& Chen Bing 4 & \\
& Xia Qi 1,3 & \\
& Liu Ang 1,3 & \\
\end{longtable}
}

Notes: 1. Transferred, dismissed or dropped 2. Newcomers to Zhejiang 3.
Military cadre 4. Pre-Cultural Revolution administrative experience in
Zhejiang 5. Sided with the ``left'' in the CR

注： 1. 调动、撤职或未留任 2. 浙江新人 3. 军队干部 4.
文革前在浙江的行政经历 5. 文革中偏``左''

Sources: Union Research Institute, Hierarchies of the People's Republic
of China, March 1975 (Hong Kong: Union Research Institute, 1975). HZRB,
July 23,1975, p.~1. CNS, 578, pp.~7-8.

资料来源：协和研究所，《中华人民共和国官员层级》，1975年3月（香港：协和研究所，1975年）。HZRB，1975年7月23日，第1页。CNS，578，第7-8页。

Zhang Zishi's deputy as Second Secretary of the CCP HMC was Chen Wenshu,
who was most likely picked out from his post as party secretary of
Anyang District in Henan province by Ji Dengkui. In total, the CCP HMC
leadership experienced an even more extensive restructuring than its
provincial counterpart: of a total number of seven secretaries and
deputy-secretaries, only one or two retained their posts after the July
1975 overhaul (See Table Twelve). This startling fact demonstrated the
complete lack of confidence which was shown by the center in the
administration of one of China's major cities. In 1974, the leadership
of the CCP HMC had consisted of: Wang Zida (1st Secretary), Wang Xing
(2nd Secretary) and deputy-secretaries Wang Shichuan (汪石川), (also 1st
Vice-chairman of the HMRC and a Cultural Revolution ``revolutionary
leading cadre'') Zhou Feng (Deputy-secretary of the pre-Cultural
Revolution CCP HMC who returned to office in 1970), Xu Shunian (徐树年)
(Commander of the Hangzhou Garrison), He Xianchun and Qiu Qiang (also
Deputy-secretary of the pre-Cultural Revolution HMC and like Wang Zida
and Wang Shichuan a ``revolutionary leading cadre'').

张子石作为中共杭州市委第二书记的副手是陈文书，他很可能是由纪登奎从河南省安阳地区党委书记的职位上挑选出来的。总体而言，中共杭州市委领导层经历了比省级领导层更为广泛的改组：在1975年7月的大改组之后，七名书记和副书记中只有一到两名留任（见表十二）。这一令人震惊的事实表明，中央对中国一座主要城市的行政管理完全缺乏信心。1974年，中共杭州市委领导班子包括：王子达（第一书记）、王兴（第二书记），以及副书记汪石川（也是杭州市革命委员会第一副主任、文革``革命领导干部''）、周锋（文革前中共杭州市委副书记，1970年复出）、徐树年（杭州警备区司令员）、贺贤春、邱强（也是文革前杭州市委副书记，并且同王子达、汪石川一样是``革命领导干部''）。

TABLE TWELVE The Secretariat of the CCP HMC before and after the
Reorganization of July 1975

表十二 1975年7月改组前后的中共杭州市委书记处

{\def\LTcaptype{none} % do not increment counter
\begin{longtable}[]{@{}lll@{}}
\toprule\noalign{}
& BEFORE & AFTER \\
\midrule\noalign{}
\endhead
\bottomrule\noalign{}
\endlastfoot
First Secretary: & Wang Zida 1,4,5 & Zhang Zishi 2,5 \\
Second Secretary: & Wang Xing 1,4 & Chen Wenshu 2,5? \\
Deputy-secretaries: & Wang Shichuan 1,4,5 & Zhou Feng 4 \\
& Zhou Feng 4 & Chen Anyu 4 \\
& Xu Shunian 1,3 & Chen Xia 4 \\
& He Xianchun 1 & \\
& Qiu Qiang 1,4,5 & \\
& Chen Anyu 4 & \\
\end{longtable}
}

{\def\LTcaptype{none} % do not increment counter
\begin{longtable}[]{@{}lll@{}}
\toprule\noalign{}
& 之前 & 之后 \\
\midrule\noalign{}
\endhead
\bottomrule\noalign{}
\endlastfoot
第一书记： & Wang Zida 1,4,5 & Zhang Zishi 2,5 \\
第二书记： & Wang Xing 1,4 & Chen Wenshu 2,5? \\
副书记： & Wang Shichuan 1,4,5 & Zhou Feng 4 \\
& Zhou Feng 4 & Chen Anyu 4 \\
& Xu Shunian 1,3 & Chen Xia 4 \\
& He Xianchun 1 & \\
& Qiu Qiang 1,4,5 & \\
& Chen Anyu 4 & \\
\end{longtable}
}

Notes: 1. Transferred, dismissed or dropped 2. Newcomers to Zhejiang 3.
Military cadre 4. Pre-Cultural Revolution administrative experience in
Zhejiang 5. Cultural Revolution ``revolutionary leading cadre''

注：1.调动、撤职、离职 2.新来浙江 3.军事干部 4.文革前在浙江的行政经历
5.文革``革命领导干部''

Sources: United States Government, Directory of Chinese Communist
Officials A70-13 (May 1970), p.~89. Hangzhou Daily, 1975.

资料来源：美国政府，《中国共产党官员名录》A70-13（1970年5月），第89页。《杭州日报》，1975年。

Well before July 1975 the central authorities had moved to restructure
the municipal administration. Apart from Wang Zida and Wang Xing, who
were transferred out of the province, in April 1975 Wang Shichuan and Xu
Shunian also lost their posts and disappeared from public view. He
Xianchun was severely disciplined and Qiu Qiang was sidelined,
criticized and did not play an active role in municipal politics again
until the radical resurgence of 1976. After the arrest of the Gang of
Four, Qiu was dismissed as their ``agent'' on the CCP HMC. Thus, of the
members of the Hangzhou leadership holding office at the beginning of
1975, only Zhou Feng survived the July 1975 shake-up. Prior to the July
reshuffle, and undoubtedly on the recommendation of the provincial
leadership, Beijing had appointed as deputy-secretaries Chen Xia (陈侠),
who returned to the position he had previously held in 1966, and Chen
Anyu (陈安羽), another veteran cadre who had fought in Eastern Zhejiang
before 1949. Both Chens took up their posts in the first four months of
1975.

早在1975年7月之前，中央当局就已开始重组市政管理机构。除王子达、王兴被调出省外，1975年4月，汪石川、徐树年也失去职务并消失在公众视野中。贺贤春受到严厉纪律处分，邱强则被边缘化、受到批评，直到1976年激进派重新抬头才再次在市政政治中发挥积极作用。四人帮被捕后，邱强因被视为他们在中共杭州市委的``代理人''而被免职。因此，1975年初任职的杭州市领导班子成员中，只有周锋在1975年7月的改组中留任。在7月改组之前，无疑是在省领导层的推荐下，北京任命陈侠和陈安羽为副书记，陈侠回到了他1966年曾担任的职务，陈安羽则是另一位1949年前曾在浙东作战的老干部。两位陈都是在1975年前四个月上任的。

One observer has speculated that the leadership reshuffle in the
province had commenced at the end of May, giving as his reason the fact
that ``Hangzhou Radio showed signs then of being placed abruptly under
new management''.\footnote{Leo Goodstadt, FEER, August 15,1975, p.~25.

  Leo Goodstadt，FEER，1975年8月15日，第25页。} Indeed, the BBC
monitoring service commented on the strange behavior of the provincial
service on 30 May 1975 and the station went off the air for most of the
day from 10 am until 9 pm. During this period only short broadcasts,
relaying programmes from Beijing, were heard.\footnote{SWB/FE/4922/BII/20.

  SWB/FE/4922/BII/20。} Goodstadt may be correct in stating that
leadership changes were under way before July but the suspension of
radio broadcasts from Hangzhou may have occurred for other reasons. At a
rally on broadcasting work in September 1975, well after the incident
and after Goodstadt's article had appeared, another explanation for the
unusual behavior of the provincial radio service was implied. Such
pronouncements as ``Bourgeois factionalism must never be allowed to
interfere with broadcasting work'' and exhortations to all departments
to ensure that ``broadcasting goes on persistently at all times and
under all circumstances'',\footnote{See ZPS, September 16,1975,
  SWB/FE/5012/BII/3-4.

  参见 ZPS，1975 年 9 月 16 日，SWB/FE/5012/BII/3-4。} suggested that
radical demagogues may have tried, unsuccessfully, to seize control of
the station in late May. Thus, rather than signalling a change in
management, Zhejiang radio's erratic transmissions on 30 May may have
been a sign of preventative action by the provincial authorities to stop
control of the station slipping out of their hands. If this was the case
it represented yet another remarkable incident in a year full of
surprising developments.

一位观察人士推测，该省的领导层改组是在5月底开始的，他的理由是``杭州广播电台当时出现了被突然置于新的管理之下的迹象''。\footnote{Leo
  Goodstadt, FEER, August 15,1975, p.~25.

  Leo Goodstadt，FEER，1975年8月15日，第25页。}
事实上，BBC监测部门在1975年5月30日就对省广播电台的奇怪行为进行了评论，该电台从上午10点到晚上9点大部分时间都停播了。在此期间，只能听到来自北京的转播节目的简短广播。\footnote{SWB/FE/4922/BII/20.

  SWB/FE/4922/BII/20。}
古德施塔特的说法可能是正确的，即7月份之前正在进行领导层更迭，但杭州的广播电台暂停可能是由于其他原因。1975年9月，在一次关于广播工作的集会上，在事件发生很久之后以及古德施塔特的文章发表之后，人们暗示了对省广播电台异常行为的另一种解释。诸如``决不允许资产阶级派性干扰广播工作''以及告诫各部门确保``广播在任何时候、任何情况下坚持下去''等言论表明，激进煽动者可能曾试图在5月下旬夺取广播电台的控制权，但没有成功。因此，浙江广播电台5月30日不稳定的信号传输与其说是管理层变动的信号，倒不如说是省政府采取预防性行动的迹象，以阻止该电台的控制权从他们手中流失。如果情况属实，那么这将是这一充满令人惊讶的事态发展的一年中又一件非凡的事件。\footnote{See
  ZPS, September 16,1975, SWB/FE/5012/BII/3-4.

  参见 ZPS，1975 年 9 月 16 日，SWB/FE/5012/BII/3-4。}

The personnel arrangements which were decided upon in July 1975 did not
represent a decisive victory for Deng. He may have rid Tan of
troublesome subordinates but some of their replacements, judged by their
performance in the radical upsurge of 1976, proved little better and in
some cases worse. In any case, Tan's stint in Zhejiang was shortly to
end. There is thus little evidence to support the view that the new
appointees to the CCP ZPC would alleviate chronic dissension within its
ranks.\footnote{New York Times, July 29, 1975.

  纽约时报，1975 年 7 月 29 日。} Tan himself was hopelessly compromised
and weakened. The influence of the central radicals on the appointment
of provincial leaders proved greater, in this instance, than has often
been recognized. The perceptive Hong Kong weekly China News Summary
noted at the end of 1975 that ``The big Chekiang reshuffle in July 1975
\ldots{} may well have increased the radical influence in the leadership
of that province''.\footnote{CNS, 596, p.~3.

  CNS，596，第3页。} This view, as events later proved, was correct.

1975年7月决定的人事安排并不代表邓的决定性胜利。他也许已经使谭摆脱了那些麻烦的下属，但从他们在1976年激进浪潮中的表现来看，一些接替者并没有好到哪里去，在某些情况下甚至更糟。无论如何，谭在浙江的任期很快就要结束了。因此，几乎没有证据支持这样的观点：中共浙江省委的新任命者将缓解其内部长期的分歧。\footnote{New
  York Times, July 29, 1975.

  纽约时报，1975 年 7 月 29 日。}谭本人已经受到严重牵累并被削弱。事实证明，中央激进分子对省级领导人任命的影响，在这一事例中比人们通常认识到的更大。敏锐的香港周刊《中国新闻摘要》在1975年底指出，``1975年7月的浙江大改组\ldots\ldots 很可能增加了该省领导层的激进影响力''。\footnote{CNS,
  596, p.~3.

  CNS，596，第3页。}后来的事态证明，这种观点是正确的。

As far as it is possible to do so, it is useful to compare the situation
in Zhejiang with that in other Chinese provinces. Of the twenty-nine
provincial first party secretaries appointed during the reconstruction
of party committees in the period from December 1970 until August 1971,
only fourteen retained their posts by October 1975. All but one of the
fifteen who were replaced in the intervening period were military cadres
but only six of their replacements came from the PLA. Thus the balance
had shifted strongly against military representatives with a total of
twelve military and seventeen civilian first secretaries by 1975. Eight
of these seventeen civilians had been purged at the beginning of the
Cultural Revolution, and of the eight, four had been appointed to their
posts in 1975.\footnote{CNS, 585 (October 1, 1975). See also, Li
  Ming-hua, ``The Chinese Communist Leadership Reorganization'', I\&S,
  12:3 (1976), pp.~37-56.

  CNS，585（1975 年 10 月 1
  日）。另见李明华，``中国共产党领导层重组''，《I\&S》，12:3 (1976)，第
  37-56 页。} It is evident that such provincial appointments as Zhao
Ziyang to Sichuan in December 1975, Bai Rubing to Shandong, Peng Chong
to Jiangsu and Liao Zhigao to Fujian in December 1974, reflected the
influence of Deng Xiaoping. However, the appointments of Song Peizhang
to Anhui in June 1975 and Jia Qiyun to Yunnan in October 1975 must be
seen as further evidence of the radicals' strength in the appointments'
field.

在可能的情况下，将浙江的情况与中国其他省份的情况进行比较是有用的。1970年12月至1971年8月党委重建期间任命的29名省委第一书记中，到1975年10月只有14名留任。在此期间更换的15名省委第一书记中，除一名外，全部是军队干部，但只有六名接替者来自解放军。因此，到1975年，这种平衡已明显转向不利于军方代表，共有12名军队出身和17名文职出身的省委第一书记。这17名文职第一书记中，有8名是在文革开始时被清洗的，而这8人中有4名是在1975年被任命的。显然，1975年12月赵紫阳任四川省委第一书记、1974年12月白如冰任山东、彭冲任江苏、廖志高任福建，反映了邓小平的影响。然而，1975年6月任命宋佩璋到安徽和1975年10月任命贾启允到云南，必须被视为激进派在任命领域实力的进一步证明。\footnote{CNS,
  585 (October 1, 1975). See also, Li Ming-hua, ``The Chinese Communist
  Leadership Reorganization'', I\&S, 12:3 (1976), pp.~37-56.

  CNS，585（1975 年 10 月 1
  日）。另见李明华，``中国共产党领导层重组''，《I\&S》，12:3 (1976)，第
  37-56 页。}

The radical central leaders clearly exercised substantial input into the
appointments to the reconstituted party leadership in Zhejiang and
Hangzhou. However, they could not save the three rebel chiefs, Zhang
Yongsheng, Weng Senhe and He Xianchun from severe punishment. Weng Senhe
was arrested and detained on 9 July 1975 on the direct orders of Ji
Dengkui.\footnote{On the eve of his arrest, Weng allegedly made
  preparations to escape to Hong Kong using an exit visa with a false
  name provided by a follower in the local public security bureau. ``The
  Case of Weng Sengho'', p.~5. See also ZJRB, August 14,1978, which made
  the same allegation, although less specifically.

  据称，翁被捕前夕，准备使用当地公安局一名追随者以假名提供的出境签证逃往香港。《翁森鹤案》，第5页。另见
  ZJRB，1978年8月14日；该报道也提出了同样指控，但不那么具体。} Many of
Weng's colleagues at the silk mill were also detained for investigation
and forced to write confessions detailing Weng's alleged
crimes.\footnote{Weng Senhe bufen zuizheng cailiao, passim.

  《翁森鹤部分罪证材料》，各处。} The Zhejiang Workers Political School
was closed and some of its students arrested.\footnote{``The Case of
  Weng Sengho'', p.~4.

  《翁森鹤案》，第4页。} Weng's prison notes record that on 9 July he
saw Chai Qikun at the ZPRC hostel in Mishi Lane. Presumably, Chai
informed Weng of his impending fate. Between the time of his arrest and
January 1976, Weng was subjected to repeated interrogation by two to
four-man teams of investigators.\footnote{Weng Senhe bufen zuizheng
  cailiao, p.~13.

  《翁森鹤部分罪证材料》，第13页。} He was charged with embezzling 7,100
yuan in cash and 600 kilograms worth of grain coupons between 1970 and
July 1975.\footnote{See ibid. Part 6, pp.~68-82 for the evidence.

  参见同上。第 6 部分，第 68-82 页为证据。} Additionally, at his trial
in 1978, Weng was accused of squandering 140,000 yuan worth of state
funds during his fourteen-month stay at the Trade Union sanatorium on
Three Terrace Mountain between 1974 and 1975.\footnote{HZRB, August
  14,1978.

  HZRB，1978 年 8 月 14 日。} In numerous meetings held to denounce
bourgeois factionalism in the months after Weng's arrest, he was not
criticized by name. But at one meeting in particular, convened by the
CCP HMC in October 1975, it was made abundantly clear that Weng was the
target.\footnote{HZRB, October 16,1975. In his tour of the silk complex
  in January 1976, Zweig noticed a wall poster attacking Weng by name.
  Zweig, ``The Peita Debate'', pp.~147-48.

  HZRB，1975 年 10 月 16 日。 1976 年 1
  月，茨威格在参观丝绸工厂时注意到墙上有一张攻击翁的海报。茨威格，``佩塔辩论''，第
  147-48 页。}

激进的中央领导人显然对浙江和杭州改组后的党的领导层任命产生了重大影响。然而，他们却未能使张永生、翁森鹤、贺贤春三位造反派首领免受严厉惩罚。翁森鹤于1975年7月9日在纪登奎的直接命令下被捕并被拘留。\footnote{On
  the eve of his arrest, Weng allegedly made preparations to escape to
  Hong Kong using an exit visa with a false name provided by a follower
  in the local public security bureau. ``The Case of Weng Sengho'',
  p.~5. See also ZJRB, August 14,1978, which made the same allegation,
  although less specifically.

  据称，翁被捕前夕，准备使用当地公安局一名追随者以假名提供的出境签证逃往香港。《翁森鹤案》，第5页。另见
  ZJRB，1978年8月14日；该报道也提出了同样指控，但不那么具体。}
翁在丝绸厂的许多同事也被拘留调查，并被迫写下详细说明翁涉嫌罪行的供词。\footnote{Weng
  Senhe bufen zuizheng cailiao, passim.

  《翁森鹤部分罪证材料》，各处。}
浙江工人政治学校被关闭，部分学生被捕。\footnote{``The Case of Weng
  Sengho'', p.~4.

  《翁森鹤案》，第4页。}
翁的狱中笔记记录，7月9日，他在米市巷浙江省革命委员会招待所见到了柴启琨。据推测，柴告诉了翁他即将面临的命运。从被捕到1976年1月，翁受到两到四人侦查组的反复审讯。\footnote{Weng
  Senhe bufen zuizheng cailiao, p.~13.

  《翁森鹤部分罪证材料》，第13页。}
他被指控在1970年至1975年7月期间贪污7,100元现金和价值600公斤的粮券。\footnote{See
  ibid. Part 6, pp.~68-82 for the evidence.

  参见同上。第 6 部分，第 68-82 页为证据。}
此外，在1978年的审判中，翁被指控在1974年至1975年间住在三台山工会疗养院的十四个月里挥霍14万元国家资金。\footnote{HZRB,
  August 14,1978.

  HZRB，1978 年 8 月 14 日。}
在翁被捕后几个月举行的许多批判资产阶级派性的会议上，他没有被点名批评。但在1975年10月中共杭州市委召开的一次会议上，已经非常明确地表明翁就是目标。\footnote{HZRB,
  October 16,1975. In his tour of the silk complex in January 1976,
  Zweig noticed a wall poster attacking Weng by name. Zweig, ``The Peita
  Debate'', pp.~147-48.

  HZRB，1975 年 10 月 16 日。 1976 年 1
  月，茨威格在参观丝绸工厂时注意到墙上有一张攻击翁的海报。茨威格，``佩塔辩论''，第
  147-48 页。}

Wang Hongwen found it politically expedient to dissociate himself
publicly from Weng during his investigations in Hangzhou. Four months
later, however, when the radical mobilizers had regained the initiative
in the seesawing inner-party struggle, Wang could express his true
feelings about Weng Senhe's arrest. In November 1975 Wang was lamenting
that ``he's been nabbed and we can't do anything about it'' (没有办法,
被他们抓住了). In the radical tide of 1976, Wang spoke out openly on
Weng's behalf and by mid-year was questioning the authenticity of the
evidence which had been compiled against him.\footnote{Bian Wen, HZRB,
  December 5,1976. HZRB, November 28,1976.

  卞文，HZRB，1976 年 12 月 5 日。 HZRB，1976 年 11 月 28 日。} Zhang
Chunqiao also spoke out for Weng. In February 1976 he said to Lai Keke
in Beijing: ``No matter what Weng Senhe is like now, his past record
will stand'' (翁森鹤不管现在怎么样, 但他过去的历史要承任他).\footnote{HZRB,
  December 31,1976, May 10,1977, December 7,1977; RMRB, May 17,1977,
  p.~2; ``Materials concerning Weng Senhe - the New-born
  counter-revolutionary''. Document of the CCP CC, zhongfa (1977),
  No.~37, I\&S, 15: 1, p.~109.

  HZRB，1976年12月31日，1977年5月10日，1977年12月7日；RMRB，1977年5月17日，第2页；《新生反革命翁森鹤材料》。中共中央文件，中发（1977）37号，I\&S，15:1，第109页。}
These remarks encouraged Weng's wife to write to the CCP ZPC, requesting
a review of her husband's case.\footnote{HZRB, November 28, 1976; Chu
  Yang, HZRB, November 2, 1977.

  HZRB，1976 年 11 月 28 日；楚阳，HZRB，1977 年 11 月 2 日。} No
decision was made at that time but in 1976 Zhang Chunqiao allegedly
crossed out Weng's name from a March CC document listing political
figures who were to be publicly criticized. The reference to Weng had
been taken from the record of a talk by Mao in November 1975.\footnote{China
  Reconstructs (July 1977), p.~42; visit to Hangzhou Silk Complex,
  January 12, 1977; PR, No.~40 (September 30,1977), p.~26; Wang Ruoshui,
  ``The Greatest Lesson of the Cultural Revolution'', p.~95.

  《中国建设》（1977年7月），第42页；1977年1月12日参观杭州丝绸联合厂；PR，第40期（1977年9月30日），第26页；王若水，《文化大革命的最大教训》，第95页。}
In the talk Mao had reportedly stated, when referring to the leadership
combination of the old, middle-aged and the young, that ``the young must
be good, not bad people like Weng Senhe (老中青, 青要好的,
不要翁森鹤这样的坏人).\footnote{Weng Senhe bufen zuizheng cailiao,
  p.~12.

  《翁森鹤部分罪证材料》，第12页。}

王洪文发现，在杭州调查期间公开与翁撇清关系是政治上的权宜之计。然而四个月后，当激进动员者在拉锯式的党内斗争中重新掌握主动权时，王可以表达他对翁森鹤被捕的真实感受。1975年11月，王哀叹``他被抓了，我们无能为力''（没有办法，被他们抓住了）。在1976年的激进浪潮中，王公开为翁说话，到了年中，他对针对翁所汇编证据的真实性提出质疑。\footnote{Bian
  Wen, HZRB, December 5,1976. HZRB, November 28,1976.

  卞文，HZRB，1976 年 12 月 5 日。 HZRB，1976 年 11 月 28 日。}
张春桥也为翁说话。1976年2月，他在北京对赖可可说：``不管翁森鹤现在怎么样，他过去的历史要承认。''\footnote{HZRB,
  December 31,1976, May 10,1977, December 7,1977; RMRB, May 17,1977,
  p.~2; ``Materials concerning Weng Senhe - the New-born
  counter-revolutionary''. Document of the CCP CC, zhongfa (1977),
  No.~37, I\&S, 15: 1, p.~109.

  HZRB，1976年12月31日，1977年5月10日，1977年12月7日；RMRB，1977年5月17日，第2页；《新生反革命翁森鹤材料》。中共中央文件，中发（1977）37号，I\&S，15:1，第109页。}
这些话鼓励翁森鹤的妻子写信给中共浙江省委，要求复查她丈夫的案件。\footnote{HZRB,
  November 28, 1976; Chu Yang, HZRB, November 2, 1977.

  HZRB，1976 年 11 月 28 日；楚阳，HZRB，1977 年 11 月 2 日。}
当时没有作出决定，但1976年张春桥据称从中央3月一份列明将被公开批评的政治人物的文件中划掉了翁的名字。文件中提及翁的内容取自毛泽东1975年11月一次谈话的记录。\footnote{China
  Reconstructs (July 1977), p.~42; visit to Hangzhou Silk Complex,
  January 12, 1977; PR, No.~40 (September 30,1977), p.~26; Wang Ruoshui,
  ``The Greatest Lesson of the Cultural Revolution'', p.~95.

  《中国建设》（1977年7月），第42页；1977年1月12日参观杭州丝绸联合厂；PR，第40期（1977年9月30日），第26页；王若水，《文化大革命的最大教训》，第95页。}
据报道，毛泽东在谈话中谈到老、中、青的领导组合时说，``青年人必须是好的，不要像翁森鹤这样的坏人。''\footnote{Weng
  Senhe bufen zuizheng cailiao, p.~12.

  《翁森鹤部分罪证材料》，第12页。}

Weng's close associate, He Xianchun, fared somewhat better. On August 5
the provincial authorities decided to send him to a cadre school in Yin
county, Ningbo District and he left two days later.\footnote{Ibid,
  p.~13.

  同上，第13页。} Although He had been ``sent down'' to labor it seems
that he spent a considerable portion of his stay in the countryside
touring and sightseeing. He had been ordered to join a study class
arranged by the CCP HMC party school, which ran a cadre school in the
county. Yin county was close to Ningbo and thus enabled He to keep in
close touch with developments in Hangzhou.\footnote{HZRB, June 20,1977;
  Yin county criticism group, ``What did He Xianchun do in Yin
  county?'', HZRB, July 5,1977; Jinlu brigade, Yinjiang commune, Yin
  county, ``Expose He Xianchun's foul performance when he was in the
  countryside'', HZRB, July 7,1977.

  HZRB，1977年6月20日；鄞县批判组，《贺贤春在鄞县干了些什么？》，HZRB，1977年7月5日；鄞县鄞江公社金炉大队，《揭露贺贤春下乡时的恶劣表现》，HZRB，1977年7月7日。}
By August 1976 he had returned to the provincial capital and appeared at
the funeral service held for Mao Zedong on 18 September 1976.\footnote{HZRB,
  September 19, 1976.

  HZRB，1976 年 9 月 19 日。} Although he had been disciplined and
demoted from Deputy-secretary to standing committee member of the CCP
HMC, He had not lost his official posts, further evidence of the
incompleteness of the attempt to destroy the careers of factional
leaders in July 1975.

翁的亲密伙伴贺贤春的情况要好一些。8月5日，省当局决定将他送到宁波地区鄞县一所干部学校，两天后他就离开了。尽管他被``下放''劳动，但他在农村逗留期间似乎有相当一部分时间是在游览观光。他被命令参加中共杭州市委党校安排的学习班，该党校在该县办有一所干部学校。鄞县靠近宁波，使他能够密切关注杭州的事态发展。\footnote{HZRB,
  June 20,1977; Yin county criticism group, ``What did He Xianchun do in
  Yin county?'', HZRB, July 5,1977; Jinlu brigade, Yinjiang commune, Yin
  county, ``Expose He Xianchun's foul performance when he was in the
  countryside'', HZRB, July 7,1977.

  HZRB，1977年6月20日；鄞县批判组，《贺贤春在鄞县干了些什么？》，HZRB，1977年7月5日；鄞县鄞江公社金炉大队，《揭露贺贤春下乡时的恶劣表现》，HZRB，1977年7月7日。}
1976年8月，他回到省城，并出现在1976年9月18日为毛泽东举行的追悼会上。\footnote{HZRB,
  September 19, 1976.

  HZRB，1976 年 9 月 19 日。}
尽管他受到纪律处分，从中共杭州市委副书记降为常委，但他并没有失去官职，这进一步证明了1975年7月试图摧毁派系领导人生涯的行动并不彻底。\footnote{Ibid,
  p.~13.

  同上，第13页。}

On the question of factional organizations, the CC and State Council
decision of 24 July 1975 differentiated between organizations or
individuals merely engaged in factional activities and those
organizations which had been formed and manipulated (操纵) or
infiltrated (混进) by bad people (坏人). The center ordered the
dissolution (解散) of the former type but banned (取缔) the latter, a
much harsher, more final step.\footnote{``Decision'', p.85. I\&S
  translated as ``disband'' and as ``dissolved'', thereby blurring the
  difference between the two terms, a difference that the document had
  been at pains to specify.

  《决定》，第85页。I\&S
  分别译作``disband''和``dissolved''，从而模糊了这两个术语之间的区别，而该文件原本着力明确这一区别。}
While dissolution left open the possibility that an organization could
regroup, banning it did not. However, the absence of any direct
reference to the militia command in the 24 July decision enabled Zhang
Yongsheng and He Xianchun to argue, when the political pendulum had
swung their way once again in 1976, that Beijing had not explicitly
banned their ``second armed force'' and that therefore they were free to
reestablish it.\footnote{HZRB, April 20, 1977.

  HZRB，1977 年 4 月 20 日。}

在派系组织问题上，中央、国务院1975年7月24日的决定区分了单纯从事派系活动的组织或个人，与由坏人组织、操纵或混入的组织或个人。中央下令解散前一类组织，但取缔后一类组织，这是更严厉、更彻底的一步。\footnote{``Decision'',
  p.85. I\&S translated as ``disband'' and as ``dissolved'', thereby
  blurring the difference between the two terms, a difference that the
  document had been at pains to specify.

  《决定》，第85页。I\&S
  分别译作``disband''和``dissolved''，从而模糊了这两个术语之间的区别，而该文件原本着力明确这一区别。}虽然解散留下了组织重组的可能性，但取缔则没有。然而，7月24日的决定中没有直接提及民兵指挥部，这使得张永生和贺贤春在1976年政治钟摆再次摆动时辩称，北京没有明确取缔他们的``第二武装力量''，因此他们可以自由地重建它。\footnote{HZRB,
  April 20, 1977.

  HZRB，1977 年 4 月 20 日。}

He Xianchun salvaged his own career but he could not prevent the
dissolution of the rebels' organizational power base, the municipal
workers' congress. On the eve of May Day 1975, the workers' congress, in
conjunction with the ZPTUC, had convened a meeting to discuss political
theory.\footnote{ZJRB, May 2,1975, p.~2.

  ZJRB，1975年5月2日，第2页。} Thus, unlike the militia command, it had
not been dissolved or banned earlier in the year. The 24 July 1975
decision rectified this oversight and did implicitly ban the
congress.\footnote{See HZRB, September 17,1977 which dates the taking
  down of the Congress crest at July 1975.

  参见 HZRB，1977 年 9 月 17 日，其中国会徽章被拆除的日期为 1975 年 7
  月。} On the day that the crest of the HMWC was to come down. He
Xianchun reportedly encouraged his young colleagues to look ahead with
optimism but by National Day 1975, after three months' study in Beijing,
they had become very pessimistic about the future.\footnote{HZRB, March
  29, 1977; HZRB, September 8, 1977.

  HZRB，1977 年 3 月 29 日； HZRB，1977 年 9 月 8 日。} With both the
urban militia command and the worker's congress dissolved, and Mao
Zedong Thought propaganda teams temporarily disbanded, three factional
organizations in Hangzhou had departed the political stage, thereby
fatally undermining the rebels' organizational base.

贺贤春挽救了自己的事业，但他无法阻止造反派的组织权力基础------市工人代表大会解散。1975年``五一''前夕，工人代表大会会同浙江省总工会召开会议，讨论政治理论。\footnote{ZJRB,
  May 2,1975, p.~2.

  ZJRB，1975年5月2日，第2页。}
因此，与民兵指挥部不同，它年初尚未被解散或取缔。1975年7月24日的决定纠正了这一疏忽，并实际上隐含地取缔了该大会。\footnote{See
  HZRB, September 17,1977 which dates the taking down of the Congress
  crest at July 1975.

  参见 HZRB，1977 年 9 月 17 日，其中国会徽章被拆除的日期为 1975 年 7
  月。}
在杭州市工代会的牌子即将摘下的那一天，据报道，贺贤春鼓励他的年轻同事乐观地看待未来，但到了1975年国庆节，在北京学习三个月后，他们对未来变得非常悲观。随着城市民兵指挥部和工人代表大会解散，毛泽东思想宣传队暂时解散，杭州的三个派系组织退出了政治舞台，从而致命地削弱了造反派的组织基础。\footnote{HZRB,
  March 29, 1977; HZRB, September 8, 1977.

  HZRB，1977 年 3 月 29 日； HZRB，1977 年 9 月 8 日。}

Zhang Yongsheng, the remaining member of the notorious rebel trio was
also sent to the countryside to labor\footnote{The following paragraphs
  are based on, HZRB, January 23,1977, March 24,1977, June 20, 1977,
  August 14,1978; RMRB, April 15, 1979; ZPS, March 3,1977,
  SWB/FE/5456/BII/6-7; CCP HMC Propaganda Department criticism group,
  ``Evidence of Zhang Yongsheng's remote-control attempt to usurp
  power'', HZRB, May 17,1977.

  以下段落基于HZRB，1977年1月23日、1977年3月24日、1977年6月20日、1978年8月14日；RMRB，1979年4月15日；ZPS，1977年3月3日，SWB/FE/5456/BII/6-7；中共杭州市委宣传部批判组，《张永生遥控篡权的证据》，HZRB，1977年5月17日。}
and he departed Hangzhou on 9 October 1975.\footnote{Zhang Yongshengde
  zuizheng cailiao, p.~124.

  《张永生的罪证材料》，第124页。} His new home was well away from
Zhejiang, in a county in Hebei province near Tianjin. The party
secretary of the brigade was Wang Guofan, a model peasant and a member
of the 10th CCP CC.\footnote{For details concerning the brigade and
  Wang, see John Wilson Lewis, Leadership in Communist China (Ithaca,
  New York: Cornell University Press, 1963), pp.~204-11. In May 1975
  Wang published an article in Hongqi (pp.~47-51) concerning
  agricultural cooperativization.

  有关该旅和王的详细信息，请参阅约翰·威尔逊·刘易斯 (John Wilson
  Lewis)，《共产主义中国的领导力》（伊萨卡，纽约：康奈尔大学出版社，1963
  年），第 204-11 页。
  1975年5月，王在《红旗》上发表了一篇关于农业合作化的文章（第47-51页）。}
Zhang's punishment had originally been more severe but Wang Hongwen had
reportedly said to the ZPC leaders in Shanghai: ``If Zhang Yongsheng
dies down there what will you do?''\footnote{Zhang Yongshengde zuizheng
  cailiao, p.~14. In a letter dated August 7, Zhang had informed Wang of
  his poor health. See ibid, p.~112.

  《张永生的罪证材料》，第14页。在8月7日的一封信中，张曾告知王自己的健康状况不佳。参见同上，第112页。}
Wang's warning caused the provincial leadership to mitigate its sentence
and Zhang may well have avoided the extreme rigors of farm life. In
January 1976 he was transferred to a hospital in Tianjin, giving
substance to Wang Hongwen's concern about the state of Zhang's
health.\footnote{Ibid, pp.~15, 23,117. In May 1972 Zhang had been in
  hospital in Hangzhou with an unspecified complaint, ibid, p.~9.

  同上，第15、23、117页。1972年5月，张因未说明的病症在杭州住院，同上，第9页。}
At a centrally-convened meeting in February 1976 (see chapter ten) Wang
Hongwen visited the Zhejiang delegation and rebuked the provincial
leadership for its treatment of Zhang. Zhang Chunqiao stated bluntly to
either Tan Qilong or Tie Ying: ``If you want to topple Zhang Yongsheng,
first we must topple you. Then it will be fair''. Zhang also made his
preferences known by informing the Zhejiang leadership in no uncertain
terms that ``I support Zhang Yongsheng all the way''.

造反三人组的最后一名成员张永生也被下放到农村劳动\footnote{The following
  paragraphs are based on, HZRB, January 23,1977, March 24,1977, June
  20, 1977, August 14,1978; RMRB, April 15, 1979; ZPS, March 3,1977,
  SWB/FE/5456/BII/6-7; CCP HMC Propaganda Department criticism group,
  ``Evidence of Zhang Yongsheng's remote-control attempt to usurp
  power'', HZRB, May 17,1977.

  以下段落基于HZRB，1977年1月23日、1977年3月24日、1977年6月20日、1978年8月14日；RMRB，1979年4月15日；ZPS，1977年3月3日，SWB/FE/5456/BII/6-7；中共杭州市委宣传部批判组，《张永生遥控篡权的证据》，HZRB，1977年5月17日。}，他于1975年10月9日离开杭州。\footnote{Zhang
  Yongshengde zuizheng cailiao, p.~124.

  《张永生的罪证材料》，第124页。}他的新家远离浙江，在河北省天津附近的一个县。该大队党委书记是王国藩，一位模范农民、中共十届中央委员。\footnote{For
  details concerning the brigade and Wang, see John Wilson Lewis,
  Leadership in Communist China (Ithaca, New York: Cornell University
  Press, 1963), pp.~204-11. In May 1975 Wang published an article in
  Hongqi (pp.~47-51) concerning agricultural cooperativization.

  有关该旅和王的详细信息，请参阅约翰·威尔逊·刘易斯 (John Wilson
  Lewis)，《共产主义中国的领导力》（伊萨卡，纽约：康奈尔大学出版社，1963
  年），第 204-11 页。
  1975年5月，王在《红旗》上发表了一篇关于农业合作化的文章（第47-51页）。}张的处罚原本较重，但据报道，王洪文曾在上海对中共浙江省委领导说：``如果张永生死在下面，你们怎么办？''\footnote{Zhang
  Yongshengde zuizheng cailiao, p.~14. In a letter dated August 7, Zhang
  had informed Wang of his poor health. See ibid, p.~112.

  《张永生的罪证材料》，第14页。在8月7日的一封信中，张曾告知王自己的健康状况不佳。参见同上，第112页。}王的警告导致省领导减轻了处分，张很可能避免了极端严酷的农场生活。1976年1月，他被转到天津一家医院，证实了王洪文对张健康状况的关心。\footnote{Ibid,
  pp.~15, 23,117. In May 1972 Zhang had been in hospital in Hangzhou
  with an unspecified complaint, ibid, p.~9.

  同上，第15、23、117页。1972年5月，张因未说明的病症在杭州住院，同上，第9页。}1976年2月，在中央召开的会议上（见第十章），王洪文看望浙江代表团，并斥责省领导对张的处理。张春桥对谭启龙或铁瑛直言不讳地说：``你要扳倒张永生，就得先扳倒你，这样才公平''。张还毫不含糊地告诉浙江领导层``我全力支持张永生''，表明了自己的偏好。

Near the end of March 1976 Zhang Yongsheng was summoned to Beijing and
lodged in a high-class hostel. Wang Hongwen enjoined him to rest and
recuperate from his stint of manual labor so that he would be in a fit
state to resume his struggle with the ``capitalist-roaders'' of
Zhejiang. To bolster Zhang's self-esteem Wang declared that Mao had
confidence in him. After Zhang complained of persecution from the
provincial authorities Wang ordered Tan Qilong, then in Beijing, to come
to Zhang's residence to listen to a two-hour tirade from his youthful
protagonist. Additionally, on 28 March, Zhang was personally received by
Wang Hongwen and then, on 3 April 1976, by Wang, Jiang Qing and Zhang
Chunqiao.\footnote{Ibid, p.~17.

  同上，第17页。}

1976年3月下旬，张永生被传唤到北京，住在一家高级招待所里。王洪文叮嘱他好好休息养伤，以良好的状态继续同浙江``走资派''的斗争。为了增强张的自尊，王宣称毛对他有信心。在张抱怨省当局的迫害后，王命令当时在北京的谭启龙到张的住所，听这位年轻支持者长达两个小时的斥责。此外，3月28日，王洪文亲自接见张，随后于1976年4月3日，王、江青和张春桥亲自接见张。\footnote{Ibid,
  p.~17.

  同上，第17页。}

From his haven in Beijing, Zhang tried to exert long-distance control
over his faction.\footnote{See the evidence in ibid, pp.~24, 28, and
  Sections 4 and 5.

  参见同上第 24、28 页以及第 4 和第 5 节中的证据。} He received a stream
of visitors and reports which turned into a two-way flow back to
Hangzhou. It appears, however, that Zhang's subordinates were
experiencing difficulty in maintaining the faction's morale, with
consequent outbreaks of internal bickering accompanied by yearnings for
an end to the constant and debilitating struggle. Zhang encouraged his
followers by pointing out that ebbs in a movement were inevitably
superseded by high tides and that victory would certainly fall to them
one day. Events seemed to vindicate Zhang's prediction. By July 1976 the
``two crashthroughs'' leading group which Zhang had established in 1974
at the Fine Arts College, had been restored. A document reversing the
decision of 1975 dissolving the group, had also been prepared. From the
time of his summons to Beijing in March 1976, Zhang hoped to return to
Zhejiang and only the arrest of the Gang of Four thwarted his
aspirations. However, even as late as 19 October 1976 Zhang continued to
despatch documents to Hangzhou and to work for his recall to the
province.

张试图从北京的避难所对其派别进行远程控制。\footnote{See the evidence in
  ibid, pp.~24, 28, and Sections 4 and 5.

  参见同上第 24、28 页以及第 4 和第 5 节中的证据。}他接待了源源不断的访客和报告，这些访客和报告变成了双向流动返回杭州。然而，张的下属似乎在维持派系士气方面遇到了困难，内部争吵随之爆发，同时也渴望结束持续不断的、令人衰弱的斗争。张鼓励他的追随者指出，运动的低潮不可避免地会被高潮所取代，胜利终有一天会降临到他们身上。事态的发展似乎证实了张的预测。到1976年7月，张于1974年在美术学院成立的``两个冲破''领导小组已经恢复。还准备了一份文件，推翻
1975
年解散该团体的决定。从1976年3月被传唤到北京开始，张一直希望返回浙江，但``四人帮''的被捕使他的愿望落空。然而，直到1976年10月19日，张仍继续向杭州发送文件，并为将他召回省里而努力。

\section{2. The Despatch of Troops into the Factories of
Hangzhou}\label{the-despatch-of-troops-into-the-factories-of-hangzhou}

\section{2.
向杭州工厂派遣部队}\label{ux5411ux676dux5ddeux5de5ux5382ux6d3eux9063ux90e8ux961f}

Prior to the arrival of the central delegation in July 1975, the
provincial and municipal authorities had despatched work groups into the
factories of Hangzhou.\footnote{In Hunan province at least two factory
  party committees were readjusted by the provincial authorities in
  1975. See CNS, 579 (August 20,1975). pp.~4-5. Such measures were
  insufficient to deal with the situation in Hangzhou. For reports of
  PLA participation in production in other provinces, see CNS, 577.

  1975年，湖南省至少有两个工厂党委由省级当局调整。见
  CNS，579（1975年8月20日），第4-5页。这些措施不足以应对杭州的情况。有关解放军参与其他省份生产的报道，请参见
  CNS，577。} The move may have been decided on at the December 1974
Beijing conference concerning Zhejiang. Point two of the CC and State
Council Decision of 24 July 1975 specifically endorsed the action of the
local authorities in despatching the groups and warned against ``any
class enemies taking advantage of shortcomings and mistakes by the work
groups (工作组) or work teams (工作队) to weaken party
leadership''.\footnote{``Decision'', p.~84.

  《决定》，第84页。} The admission that the work groups (teams) had
committed errors may have referred to their tendency to support one
faction against the other, worsening rather than alleviating the
situation. Work teams had always been prone to taking sides. It was a
perennial weakness inherent in their role and function. This time the
warning may have been directed specifically at the local military, which
had undoubtedly contributed personnel to the work teams.

1975年7月中央代表团抵达之前，省市当局已向杭州工厂派出了工作组。\footnote{In
  Hunan province at least two factory party committees were readjusted
  by the provincial authorities in 1975. See CNS, 579 (August 20,1975).
  pp.~4-5. Such measures were insufficient to deal with the situation in
  Hangzhou. For reports of PLA participation in production in other
  provinces, see CNS, 577.

  1975年，湖南省至少有两个工厂党委由省级当局调整。见
  CNS，579（1975年8月20日），第4-5页。这些措施不足以应对杭州的情况。有关解放军参与其他省份生产的报道，请参见
  CNS，577。}此举可能是1974年12月北京浙江会议决定的。1975年7月24日的中央和国务院决定第二点明确赞同地方当局派遣工作组的行动，并警告``阶级敌人不得利用工作组或工作队的缺点和错误来削弱党的领导''。这一承认工作组（队）犯有错误的说法，可能指的是它们倾向于支持一派反对另一派，结果不是缓和而是恶化了局势。工作组一向容易偏袒某一方。这是其角色和职能所固有的长期弱点。这次的警告可能是专门针对当地军方的，毫无疑问，当地军方为工作组派遣了人员。\footnote{``Decision'',
  p.~84.

  《决定》，第84页。}

At the meeting convened by the CCP ZPC on 22 July to announce the
reshuffle of the provincial leadership, it was announced that ``comrades
from the PLA'' who had helped in the work of the CCP ZPC and ZPRC in the
``three supports and two militaries'', on the instructions of the CC,
had now returned to their barracks.\footnote{HZRB, July 23, 1975.

  HZRB，1975 年 7 月 23 日。} It is significant that the monitored
broadcast of the meeting did not mention this crucial fact.\footnote{ZPS,
  July 23, 1975, SWB/FE/4967/BII/1-2. Once again it illustrates how the
  local authorities edited information before releasing it to the
  outside world. By contrast, the local press carried a more
  comprehensive coverage of such meetings.

  ZPS，1975 年 7 月 23
  日，SWB/FE/4967/BII/1-2。它再次说明了地方当局如何在向外界发布信息之前对其进行编辑。相比之下，当地媒体对此类会议的报道更为全面。}
The surprising usage of the term ``three supports and two militaries''
also belies one specialist's assertion that the phrase had been
``dropped from the official lexicon'' in 1973.\footnote{Nelsen, The
  Chinese Military System, p.~142. On August 21, 1972, the CC had issued
  a ``Notice concerning seeking opinions on the question of the three
  supports and two militaries'' with a draft ``Decision on certain
  questions concerning the three supports and two militaries'' attached.
  The latter document directed that in areas and units where party
  committees had been established, such organizations should be
  disbanded and personnel withdrawn. See ``Wenhua dageming'' ruogan
  shijian zhenxiang, p.~60. On December 30, 1974, over two years after
  the central authorities had raised the issue, the Jiangsu authorities
  decided to withdraw personnel of such units and return them to the
  barracks. See Jiangsusheng dashiji, p.~329. In Guangdong province Zhao
  Ziyang addressed the issue of the contributions and shortcomings of
  the ``three supports and two militaries'' troops in a speech of July
  15,1974, during the pilin pikong campaign. See Zhao Wei, Zhao Ziyang
  zhuan, pp.~202-3.

  Nelsen，《中国军事系统》，第142页。1972年8月21日，中央发出《关于征求对三支两军问题意见的通知》，并附有《关于三支两军若干问题的决定》草案。后一份文件指示，凡已建立党委的地区和单位，此类组织应予解散，人员应予撤出。见《``文化大革命''若干事件真相》，第60页。1974年12月30日，在中央提出这一问题两年多后，江苏当局决定撤出此类单位人员并将其送回军营。参见《江苏省大事记》，第329页。广东省赵紫阳在1974年7月15日批林批孔运动中的讲话中谈到了``三支两军''部队的贡献和缺点。参见赵伟，《赵紫阳传》，第202-203页。}
Small groups of military personnel had entered factories in Hangzhou
after the breakdown of authority in 1967. They had been ordered to
choose sides in factional disputes and had become embroiled in factory
politics for eight years. Clearly, they had lost all credentials as a
neutral force, compelling the provincial leadership to withdraw them
before the arrival of military units free from local entanglements and
favoritism.

7月22日，中共浙江省委召开会议宣布省级领导班子改组，会上宣布，曾在``三支两军''中帮助中共浙江省委和浙江省革命委员会工作的``解放军同志''，按照中央指示，现已返回军营。\footnote{HZRB,
  July 23, 1975.

  HZRB，1975 年 7 月 23 日。}
值得注意的是，会议的监听广播却没有提到这一关键事实。\footnote{ZPS, July
  23, 1975, SWB/FE/4967/BII/1-2. Once again it illustrates how the local
  authorities edited information before releasing it to the outside
  world. By contrast, the local press carried a more comprehensive
  coverage of such meetings.

  ZPS，1975 年 7 月 23
  日，SWB/FE/4967/BII/1-2。它再次说明了地方当局如何在向外界发布信息之前对其进行编辑。相比之下，当地媒体对此类会议的报道更为全面。}
``三支两军''一词的意外使用，也推翻了一位专家的说法，即该词已于1973年``从官方词典中删除''。\footnote{Nelsen,
  The Chinese Military System, p.~142. On August 21, 1972, the CC had
  issued a ``Notice concerning seeking opinions on the question of the
  three supports and two militaries'' with a draft ``Decision on certain
  questions concerning the three supports and two militaries'' attached.
  The latter document directed that in areas and units where party
  committees had been established, such organizations should be
  disbanded and personnel withdrawn. See ``Wenhua dageming'' ruogan
  shijian zhenxiang, p.~60. On December 30, 1974, over two years after
  the central authorities had raised the issue, the Jiangsu authorities
  decided to withdraw personnel of such units and return them to the
  barracks. See Jiangsusheng dashiji, p.~329. In Guangdong province Zhao
  Ziyang addressed the issue of the contributions and shortcomings of
  the ``three supports and two militaries'' troops in a speech of July
  15,1974, during the pilin pikong campaign. See Zhao Wei, Zhao Ziyang
  zhuan, pp.~202-3.

  Nelsen，《中国军事系统》，第142页。1972年8月21日，中央发出《关于征求对三支两军问题意见的通知》，并附有《关于三支两军若干问题的决定》草案。后一份文件指示，凡已建立党委的地区和单位，此类组织应予解散，人员应予撤出。见《``文化大革命''若干事件真相》，第60页。1974年12月30日，在中央提出这一问题两年多后，江苏当局决定撤出此类单位人员并将其送回军营。参见《江苏省大事记》，第329页。广东省赵紫阳在1974年7月15日批林批孔运动中的讲话中谈到了``三支两军''部队的贡献和缺点。参见赵伟，《赵紫阳传》，第202-203页。}
1967年权威崩溃后，小股军人进入杭州的工厂。他们被命令在派系斗争中选边站，并卷入工厂政治长达八年。显然，他们已经失去了作为中立力量的一切资格，迫使省领导层在不受地方纠葛和偏袒影响的军事部队到来之前撤回他们。

On 19 July 1975, the same day that Wang Hongwen returned to Hangzhou
from Beijing, troops of the PLA moved to occupy four large factories
which had been seriously affected by factionalism. Such a momentous
decision may have been taken at the CC MAC plenum which had closed on 15
July (see chapter eight) and referred to the Politburo meeting of 17
July. The first four factories to come under military occupation were
Weng's silk complex,\footnote{It has been claimed that in April 1973 the
  military had intervened in this mill to stop factional fighting. ``The
  Case of Weng Sengho'', p.~5.

  据称，1973年4月，军方曾介入这家工厂以制止派系斗争。《翁森鹤案》，第5页。}
the No.~1 and No.~2 cotton mills, and the Zhejiang Hemp Mill.\footnote{The
  Zhejiang Hemp Mill employed 6,200 workers in 1978. It had gone into
  operation in 1953 after three years' construction. Visit, December
  3,1978.

  浙江麻厂1978年雇用了6200名工人，经过三年建设，于1953年投产。 1978 年
  12 月 3 日访问。} Six thousand men of an unnamed PLA unit stationed in
Zhejiang (possibly the 1st Army) went by groups into the four mills. The
first report concerning military intervention made it clear that
factional fighting had disrupted production and damaged plant and
premises. It depicted soldiers and their officers performing such jobs
as repairing a damaged roof, cooking meals, cutting hair and washing
clothes. Medical teams treated workers while propaganda workers relayed
central instructions and held study classes for the members of the
contending factions.\footnote{HZRB, July 21,1975; ZPS, July 21,1975,
  SWB/FE/4967/BII/5-6.

  HZRB，1975 年 7 月 21 日； ZPS，1975 年 7 月 21
  日，SWB/FE/4967/BII/5-6。}

1975年7月19日，王洪文从北京返回杭州的同一天，解放军部队进驻了受派系主义严重影响的四家大型工厂。这样一项重大决定可能是在7月15日闭幕的中共中央军委全会上作出的（见第八章），并提交给了7月17日的政治局会议。最先被军事进驻的四家工厂是翁的丝绸联合厂、\footnote{It
  has been claimed that in April 1973 the military had intervened in
  this mill to stop factional fighting. ``The Case of Weng Sengho'',
  p.~5.

  据称，1973年4月，军方曾介入这家工厂以制止派系斗争。《翁森鹤案》，第5页。}第一、第二棉纺厂和浙江麻厂。\footnote{The
  Zhejiang Hemp Mill employed 6,200 workers in 1978. It had gone into
  operation in 1953 after three years' construction. Visit, December
  3,1978.

  浙江麻厂1978年雇用了6200名工人，经过三年建设，于1953年投产。 1978 年
  12 月 3 日访问。}驻浙江的一支未具名解放军部队（可能是第一军）的六千人分批进入这四家工厂。第一份有关军事干预的报告明确指出，派系斗争扰乱了生产并损坏了厂房和设备。报道描绘士兵和军官从事修复受损屋顶、做饭、理发和洗衣服等工作。医疗队为工人救治，宣传人员传达中央指示，并为敌对派别成员举办学习班。\footnote{HZRB,
  July 21,1975; ZPS, July 21,1975, SWB/FE/4967/BII/5-6.

  HZRB，1975 年 7 月 21 日； ZPS，1975 年 7 月 21
  日，SWB/FE/4967/BII/5-6。}

Reports about the four mills revealed something of the trouble which had
plagued them for months. The Hangzhou silk complex reported that the
situation there had improved since the beginning of July, the time of
the arrival of the two central leaders in Hangzhou.\footnote{ZPS, July
  21, 1975, SWB/FE/4967/BII/4.

  ZPS，1975 年 7 月 21 日，SWB/FE/4967/BII/4。} Before then, bourgeois
factionalism had split the party and workers, and production and social
welfare facilities had been adversely affected. The broadcast described
damage to the mill's power facilities, the breakdown of a boiler,
shortages of raw materials, an undermanning among nursing staff in the
factory's creche and a lack of variety in food prepared in the canteen.
Reports of 27 July 1975, commenting on the situation at the hemp mill
and the two cotton mills, made no mention of the presence of PLA troops
in these factories.\footnote{ZPS, July 27,1975, FBIS/CHI,148 (1975), G
  1-4.

  ZPS，1975 年 7 月 27 日，FBIS/CHI，148 (1975)，G 1-4。} But the
reports alluded to the CC document of 24 July, especially in relation to
its focus on cadre policy, party spirit (as against factionalism) and
the need to boost production. At the No.~1 cotton mill cadres had made
household calls to inform workers of central directives and probably
also to cajole them into returning to work.

关于这四家工厂的报道揭示了困扰它们几个月的一些问题。杭州丝绸联合厂报道称，自7月初两位中央领导人抵达杭州以来，当地局势有所好转。此前，资产阶级派性已经分裂了党和工人，生产和社会福利设施也受到不利影响。广播描述了工厂电力设施受损、锅炉故障、原材料短缺、工厂托儿所护理人员人手不足以及食堂准备的食物缺乏品种的情况。\footnote{ZPS,
  July 21, 1975, SWB/FE/4967/BII/4.

  ZPS，1975 年 7 月 21 日，SWB/FE/4967/BII/4。}1975年7月27日的报告在评论麻厂和两家棉厂的情况时，没有提到这些工厂里有解放军部队。\footnote{ZPS,
  July 27,1975, FBIS/CHI,148 (1975), G 1-4.

  ZPS，1975 年 7 月 27 日，FBIS/CHI，148 (1975)，G 1-4。}但这些报告提到了7月24日的中央文件，特别是关于其对干部政策、党性（反对派系主义）和促进生产的需要的关注。在第一棉纺厂，干部们登门走访，向工人传达中央指示，或许也是为了劝说他们重返工作岗位。

On 22 July 1975, the same day as the CCP ZPC, HMC and ZPMD held meetings
of office staff to convey the central directives, Zhang Wenbi led 4,500
cadres and troops of the military district into factories in the city,
including the Hangzhou Silk Brocade Mill.\footnote{HZRB, July 25, 1975;
  ZPS, July 25, 1975, SWB/FE/4967/ BII/6-7.

  HZRB，1975 年 7 月 25 日； ZPS，1975 年 7 月 25
  日，SWB/FE/4967/BII/6-7。} The ZPMD also assigned personnel to help
with the summer harvesting and autumn sowing, leaving behind only
skeleton staff in the offices of the ZPMD headquarters, the political
and logistics departments and the Hangzhou garrison. While the
announcement of the despatch of the first contingent of 6,000 troops had
been held back for two days, this follow-up disclosure was delayed three
days. In order to drum up support for the decision among the populace,
Zhejiang Daily accompanied the second announcement with an article
entitled ``Learn from the PLA''. The article called on citizens to
emulate the PLA

1975年7月22日，在中共浙江省委、杭州市委和浙江省军区召开机关工作人员会议传达中央指示的同一天，张文碧带领军区4500名干部和部队进驻市内工厂，包括杭州织锦厂。\footnote{HZRB,
  July 25, 1975; ZPS, July 25, 1975, SWB/FE/4967/ BII/6-7.

  HZRB，1975 年 7 月 25 日； ZPS，1975 年 7 月 25
  日，SWB/FE/4967/BII/6-7。}
浙江省军区还派出人员帮助夏收秋种，只在省军区司令部、政治部和后勤部以及杭州警备区机关留下骨干人员。虽然第一批6000人部队的派遣消息被推迟了两天公布，但这次后续披露又延迟了三天。为了争取民众对这一决定的支持，《浙江日报》在第二次公告中还发表了题为《向解放军学习》的文章。文章呼吁民众效仿解放军

\begin{quote}
in strengthening the sense of organization and discipline, strictly
carrying out the three main rules and eight points for attention and
obeying orders in every activity. We must follow the PLA in consciously
promoting stability and unity.\footnote{ZJRB, July 25,1975, p.~1;
  quotation from ZPS, July 25,1975, SWB/FE/4969/BII/1-2.

  ZJRB，1975年7月25日，第1页；引文见
  ZPS，1975年7月25日，SWB/FE/4969/BII/1-2。}
\end{quote}

\begin{quote}
强化组织意识和纪律意识，各项活动严格执行三大纪律八项注意、服从命令。我们一定要学习解放军，自觉促进安定团结。\footnote{ZJRB,
  July 25,1975, p.~1; quotation from ZPS, July 25,1975,
  SWB/FE/4969/BII/1-2.

  ZJRB，1975年7月25日，第1页；引文见
  ZPS，1975年7月25日，SWB/FE/4969/BII/1-2。}
\end{quote}

Points four, six, and seven of the CC Decision of 24 July related
directly to the PLA.\footnote{``Decision'', pp.~84-5.

  《决定》，第 84-5 页。} Beijing strongly condemned any ``words or
deeds'' which weakened or sabotaged the relations between the army and
government, the army and people or which affected the military's
internal unity. The CC decision particularly called on the military to
smash any ``despicable conspiracy'' by the ``class enemy'' to split its
ranks. Together with the public security and militia forces, the army
was commanded to suppress any counter-revolutionaries who disrupted
production, caused traffic accidents, incited armed struggle, obstructed
{[}political{]} movements (干扰运动), carried out counter-revolutionary
propaganda or incited insurrections. The PLA was also ordered to crush
all proven murderers, arsonists and water poisoners. Those who resisted
would be arrested and punished. If public order in Hangzhou and the rest
of Zhejiang had collapsed anywhere near the extent that was suggested in
the document, making allowances for a certain amount of exaggeration and
rhetorical flourish, the cities of the province would not have been
pleasant places in which to reside in 1974-75. Or, as the ``Decision''
euphemistically phrased it, such disturbances would have disrupted
social order and jeopardized the security of state property and people's
lives.

7月24日中央决定的第四、六、七点与解放军直接相关。
北京强烈谴责任何削弱或破坏军政、军民关系、影响军队内部团结的``言行''。中央委员会的决定特别要求军队粉碎任何``阶级敌人''分裂军队的``卑鄙阴谋''。指挥军队会同公安、民兵，镇压扰乱生产、制造交通事故、煽动武装斗争、干扰运动、反革命宣传、煽动暴乱的反革命分子。解放军还被命令粉碎所有已被证实的杀人犯、纵火犯和水毒犯。那些反抗的人将被逮捕并受到惩罚。如果杭州和浙江其他地区的公共秩序崩溃程度接近文件中建议的程度，并允许一定程度的夸张和华丽言辞，那么在
1974-75
年，该省的城市就不会是令人愉快的居住地。或者，正如《决定》委婉地表述的那样，这种骚乱会扰乱社会秩序，危害国家财产和人民生命安全。\footnote{``Decision'',
  pp.~84-5.

  《决定》，第 84-5 页。}

In the last few days of July 1975 the center committed more troops to
the factories of Hangzhou. On 28 July, Air Force units stationed in
Zhejiang despatched soldiers to five factories including He Xianchun's
old unit, the Heavy Machinery Factory, and the Hangzhou Oxygen Generator
Plant where He's subordinate, Xia Genfa, was based.\footnote{ZJRB,
  August 8,1975, p.~1. The Hangzhou Oxygen Generator Plant was set up in
  1950 by amalgamating four repair workshops. From an initial workforce
  of 300 it had grown, by 1978, to employ a staff of about 5,000. It is
  a key unit in the metallurgical, chemical fertilizer and defence
  industries. Visit, June 4,1978.

  ZJRB，1975年8月8日，第1页。杭州制氧机厂于1950年由四个修理车间合并组建。到1978年，该厂职工人数已从最初的300人增至约5000人。它是冶金、化肥和国防工业的重点单位。1978年6月4日访问。}
It is uncertain whether the units referred to came from the 5th Air
Force where Chen Liyun had formerly served as Political Commissar.
However, it is worth noting that as late as 1976, Bai Zongshan,
Commander of the 5th Air Force, remained in Zhejiang,\footnote{See ZPS,
  May 16, 1976, SWB/FE/5213/BII/9.

  参见 ZPS，1976 年 5 月 16 日，SWB/FE/5213/BII/9。} suggesting that his
unit had also remained in the province. The 5th Air Force had openly
sided with United Headquarters in its battles with Red Storm during the
years 1967-69. In 1971-72, because of Chen Liyun's membership of Lin
Biao's son's ``small fleet'', the special group appointed by the CCP CC
to investigate the affair concluded that the 5th Air Force had become
one of Lin Biao's most loyal military units. Hence, because of its past
associations and heavy involvement in provincial factional politics, the
unit could scarcely qualify as an unbiased force in mediating between
contending groups of workers in the factories of Hangzhou.

1975年7月的最后几天，中央向杭州工厂增派了兵力。7月28日，驻浙江空军部队向贺贤春的老单位重型机械厂、贺贤春下属夏根发所在的杭州制氧机厂等五家工厂派兵。\footnote{ZJRB,
  August 8,1975, p.~1. The Hangzhou Oxygen Generator Plant was set up in
  1950 by amalgamating four repair workshops. From an initial workforce
  of 300 it had grown, by 1978, to employ a staff of about 5,000. It is
  a key unit in the metallurgical, chemical fertilizer and defence
  industries. Visit, June 4,1978.

  ZJRB，1975年8月8日，第1页。杭州制氧机厂于1950年由四个修理车间合并组建。到1978年，该厂职工人数已从最初的300人增至约5000人。它是冶金、化肥和国防工业的重点单位。1978年6月4日访问。}尚不清楚这里所说的部队是否来自陈励耘曾任政委的第五空军。不过值得注意的是，直到1976年，第五空军司令员白宗善仍然留在浙江，表明他的部队也留在了浙江。\footnote{See
  ZPS, May 16, 1976, SWB/FE/5213/BII/9.

  参见 ZPS，1976 年 5 月 16 日，SWB/FE/5213/BII/9。}1967-69年间，第五空军在联合总部与红暴的斗争中公开站在联合总部一边。1971-72年，由于陈励耘参与林彪之子的小``舰队''，中共中央任命的调查该事件的专案组认定，第五空军已成为林彪最忠诚的军事单位之一。因此，由于其既往关联并深度卷入省内派系政治，该部队很难算是在杭州工厂各工人对立派别之间进行调解的公正力量。

The air force units performed such tasks as unloading and carting raw
materials, cleaning up the factory grounds and dining halls and holding
discussions with the workers. At the oxygen generator plant the soldiers
directed their powers of persuasion toward young workers and demobilized
PLA soldiers. Veteran workers in turn advised young PLA soldiers. The
problems posed to factory management by undisciplined, rowdy young
workers and demobilized soldiers were chronic. The case of Xiong Beiping
and Xiong Ziping, sons of the former Commander of the ZPMD Xiong
Yingtang, is an extreme example. Beiping and his twin brother Ziping,
known as the twin Xiongs (两熊) were prosecuted in late 1979 for heading
a gang of pack-rapists for four years from 1974 until 1978. Xiong Ziping
worked at He Xianchun's heavy machinery plant. Taking advantage of the
anarchy and turmoil of these years, and playing on their father's status
as a leading cadre, albeit in disgrace, Ziping, then in his mid
twenties, aroused great resentment among his fellow workers. The
following excerpt from a report of November 1979 described his behavior:

空军部队执行了卸运原材料、清扫厂区和食堂、与工人座谈等任务。在制氧厂，战士们把劝说工作转向年轻工人和复员解放军战士。老工人反过来又劝导年轻的解放军战士。缺乏纪律、吵闹的年轻工人和复员士兵给工厂管理带来的问题长期存在。原浙江省军区司令员熊英堂的儿子熊北平、熊子平就是一个极端例子。北平和他的双胞胎兄弟子平（被称为``两熊''）于1979年底被起诉，罪名是从1974年到1978年领导一个轮奸团伙长达四年。熊子平在贺贤春的重型机械厂工作。趁着这些年的无政府状态和动乱，并利用父亲虽已失势但仍是领导干部的身份，二十多岁的子平引起了工友们的极大不满。以下摘自1979年11月的一份报告，描述了他的行为：

\begin{quote}
Workers in the metal machinery workshop {[}of the heavy machinery
plant{]} who had seen Xiong Ziping commit all kinds of crimes with their
own eyes said that Xiong did as he liked at home and ran wild in the
factory. He was late for work every day and knocked off early and only
did four things on the job: eat, sleep, behave like a lout and cause
sabotage. Once a master worker wanted him to drive a vehicle but
{[}Xiong{]} abused him, saying: ``Who do you think you are? Are you game
to come here and order me around Pop?!'' He started up the vehicle to
hoist spare parts and barged about not only damaging lathes, but coming
within an ace of causing serious accidents to other workers. Everyone
said indignantly at the time: ``when rotten old Xiong charges into the
melon-patch the melons can say their prayers; when he rushes into the
factory it's we who have to look out!''\footnote{RMRB, 18 November 1979.
  When local policemen came to detain Xiong Beiping at Qingbo Gate,
  where the family lived in the military compound of the ZPMD
  headquarters, Xiong said threateningly to the squad leader; ``My
  father will be back soon. If you run me in, you must reckon on the
  consequences''. Although the national press did not report the case
  until November 1979, Hangzhou Daily had carried an article on the
  subject one year earlier, without mentioning names. See HZRB, October
  25, 1978. At their trial, Ziping was condemned to death and Beiping
  receiving a two-year suspended death sentence. Beiping worked at the
  oxygen generator plant where Xia Genfa was a rebel leader. See also
  RMRB, November 15 and 16,1979. A mimeographed newssheet from
  Guangzhou's Zhongshan University Red Star Group (红星小组), dated 15
  August 1982, went further than the People's Daily in claiming that the
  gang rapes had occurred in the lodgings of the Xiong twins, that is in
  the military barracks. Hangzhou Daily had previously admitted this
  fact.

  RMRB，1979年11月18日。当当地警察来到清波门拘留熊北平时，熊家住在浙江省军区司令部军区大院内，熊威胁班长说：``我父亲很快就会回来。你要是把我抓进去，就得考虑后果。''虽然全国媒体直到1979年11月才报道此案，但《杭州日报》早在一年前就曾刊载一篇相关文章，只是没有提及姓名。参见HZRB，1978年10月25日。审判中，子平被判处死刑，北平被判处死刑、缓期两年执行。北平在夏根发任造反派领导的制氧厂工作。另见RMRB，1979年11月15日、16日。广州中山大学红星小组一份日期为1982年8月15日的油印小报，比《人民日报》更进一步，声称轮奸案发生在熊氏双胞胎的住处，也就是军营内。《杭州日报》此前也承认了这一事实。}
\end{quote}

\begin{quote}
亲眼目睹熊子平为非作歹的金属机械车间工人说，熊子平在家里为所欲为，在厂里横行霸道。他每天上班迟到、早下班，工作中只做四件事：吃饭、睡觉、胡作非为、搞破坏。有一次，一个师傅要他开车，他就辱骂他：``你以为你是谁？你敢来这里指挥我吗？！''他启动车辆来吊装备件，并闯入，不仅损坏了车床，而且还差点对其他工人造成严重事故。当时大家都愤愤不平地说：``烂老熊冲进瓜地时，瓜们可以祈祷；他冲进工厂时，我们要当心！''\footnote{RMRB,
  18 November 1979. When local policemen came to detain Xiong Beiping at
  Qingbo Gate, where the family lived in the military compound of the
  ZPMD headquarters, Xiong said threateningly to the squad leader; ``My
  father will be back soon. If you run me in, you must reckon on the
  consequences''. Although the national press did not report the case
  until November 1979, Hangzhou Daily had carried an article on the
  subject one year earlier, without mentioning names. See HZRB, October
  25, 1978. At their trial, Ziping was condemned to death and Beiping
  receiving a two-year suspended death sentence. Beiping worked at the
  oxygen generator plant where Xia Genfa was a rebel leader. See also
  RMRB, November 15 and 16,1979. A mimeographed newssheet from
  Guangzhou's Zhongshan University Red Star Group (红星小组), dated 15
  August 1982, went further than the People's Daily in claiming that the
  gang rapes had occurred in the lodgings of the Xiong twins, that is in
  the military barracks. Hangzhou Daily had previously admitted this
  fact.

  RMRB，1979年11月18日。当当地警察来到清波门拘留熊北平时，熊家住在浙江省军区司令部军区大院内，熊威胁班长说：``我父亲很快就会回来。你要是把我抓进去，就得考虑后果。''虽然全国媒体直到1979年11月才报道此案，但《杭州日报》早在一年前就曾刊载一篇相关文章，只是没有提及姓名。参见HZRB，1978年10月25日。审判中，子平被判处死刑，北平被判处死刑、缓期两年执行。北平在夏根发任造反派领导的制氧厂工作。另见RMRB，1979年11月15日、16日。广州中山大学红星小组一份日期为1982年8月15日的油印小报，比《人民日报》更进一步，声称轮奸案发生在熊氏双胞胎的住处，也就是军营内。《杭州日报》此前也承认了这一事实。}
\end{quote}

On 31 July 1975 the Commander of the first division of the 1st Army led
his troops into six factories. Five of the six units concerned, the silk
complex, the No.~1 and No.~2 cotton mills, the hemp mill and the oxygen
generator plant had already been occupied by the PLA, an indication of
the extent of the difficulties experienced at these units. The sixth
factory, the Hangzhou iron and steel mill, was another unit with
long-standing, major and difficult (老大难) problems. The Xinhua report
from Hangzhou referred to soldiers performing simple production-oriented
tasks such as servicing and repairing machinery, recycling scrap metal,
and taking warehouse inventories. The soldiers also showed films, gave
theatrical performances and propagated central instructions.\footnote{Xinhua,
  August 2,1975, SWB/FE/4976/BII/5-6.

  新华社，1975 年 8 月 2 日，SWB/FE/4976/BII/5-6。} The report appeared
in the People's Daily and was the first mention in that newspaper of the
presence of troops in the factories of Hangzhou.\footnote{RMRB, August
  2,1975, p.~3.

  RMRB，1975年8月2日，第3页。} The report commended two officers (the
unit's Commander Li Guangshan, a national military hero, and a Red Army
regimental commander Wu Wanwan) for taking part in labor at the silk
complex and No.~1 cotton mill respectively.\footnote{In 1971 Li was a
  divisional commander in the Wuhan Military Region, further evidence of
  the origins of the 1st Army. After Li's unit was transferred to
  Zhejiang, it was attached to the Nanjing Military Region. See CNS, 578
  (August 13,1975), p.~4.

  1971年，李是武汉军区的一名师长，这进一步证明了第1军的来源。李所在部队调往浙江后，隶属南京军区。参见
  CNS，578（1975年8月13日），第4页。} The article detailed some of the
activities of the military personnel, who propagated ``important
directives of Chairman Mao and the Party CC'' by holding discussions
with workers. The workers expressed gratitude for these and other acts
of kindness. It was later admitted by Politburo member and Vice-Premier
Li Xiannian that friction was also present between military personnel
and workers. He reportedly told the leader of a visiting Japanese
delegation in early September that such ``minor incidents'' were
``natural'' because they were both ``trying to increase production as
much as possible''.\footnote{Kyodo report, September 3,1975,
  SWB/FE/4999/A3/2.

  共同社报告，1975 年 9 月 3 日，SWB/FE/4999/A3/2。} Whether any deaths
occurred as a result of the military occupation is unknown.

1975年7月31日，第1军第1师师长率部进入六家工厂。六家单位中，丝绸联合厂、第一和第二棉纺厂、麻厂、制氧机厂等五个单位已经由解放军进驻，可见这些单位所经历的困难有多大。第六家工厂杭州钢铁厂，则是另一个长期存在重大难题的``老大难''单位。新华社杭州的报道提到，士兵们执行简单的生产性任务，例如保养和修理机器、回收废金属以及盘点仓库库存。战士们还放映电影、文艺演出，宣传中央指示。\footnote{Xinhua,
  August 2,1975, SWB/FE/4976/BII/5-6.

  新华社，1975 年 8 月 2 日，SWB/FE/4976/BII/5-6。}这篇报道刊登在《人民日报》上，这是该报首次提及部队驻扎在杭州工厂的情况。\footnote{RMRB,
  August 2,1975, p.~3.

  RMRB，1975年8月2日，第3页。}报道赞扬了两名军官，即该部队司令员、全国战斗英雄李光山和一名红军团长吴万万，分别在丝绸联合厂和第一棉纺厂参加劳动。\footnote{In
  1971 Li was a divisional commander in the Wuhan Military Region,
  further evidence of the origins of the 1st Army. After Li's unit was
  transferred to Zhejiang, it was attached to the Nanjing Military
  Region. See CNS, 578 (August 13,1975), p.~4.

  1971年，李是武汉军区的一名师长，这进一步证明了第1军的来源。李所在部队调往浙江后，隶属南京军区。参见
  CNS，578（1975年8月13日），第4页。}文章详细介绍了军人的一些活动，他们通过与工人座谈来宣传``毛主席和党中央的重要指示''。工人们对这些善举和其他善举表示感谢。后来政治局委员、国务院副总理李先念承认，军队和工人之间也存在摩擦。据报道，他在9月初对来访的日本代表团领导人表示，此类``小事件''是``自然的''，因为他们都``试图尽可能增加产量''。\footnote{Kyodo
  report, September 3,1975, SWB/FE/4999/A3/2.

  共同社报告，1975 年 9 月 3 日，SWB/FE/4999/A3/2。}军事进驻是否造成任何死亡尚不得而知。

In August 1975, naval units joined the other two services in
agricultural and industrial assignments.\footnote{ZJRB, August 7,1975,
  p.~1.

  ZJRB，1975年8月7日，第1页。} The troops of one squad worked in an
uncovered workshop of the Red Flag Paper Mill in the blistering
heat.\footnote{Established in 1919 the Red Flag paper mill in 1978
  employed 2,200 workers. Visit, March 12,1978.

  红旗造纸厂成立于 1919 年，1978 年雇用了 2,200 名工人。 1978 年 3 月 12
  日访问。} Others cleaned up the grounds and helped out in the canteen.
The Nanjing Military Region allocated seventy-two cadres at and above
regimental level to lead several thousand troops into Hangzhou's
factories, including the iron and steel mill. Soldiers from the military
region had played a role in regimenting the behavior of students and
teachers of the city's colleges and schools over the previous six
months.\footnote{ZPS, August 8,1975, SWB/FE/4981/BII/9-10.

  ZPS，1975 年 8 月 8 日，SWB/FE/4981/BII/9-10。} Point seven of the CC
Decision of 24 July had directed students not to become involved in
factional conflicts outside their campuses, nor to interfere in the
socialist education movement in factories and rural areas.\footnote{``Decision'',
  p.~85.

  《决定》，第85页。} This was one of the few indications that students
and teachers had become embroiled in the factional struggles centered in
Hangzhou's factories.

1975年8月，海军部队也同另外两个军种一样参加农业和工业任务。\footnote{ZJRB,
  August 7,1975, p.~1.

  ZJRB，1975年8月7日，第1页。}一个小队的战士在酷暑中在红旗造纸厂的露天车间劳动。\footnote{Established
  in 1919 the Red Flag paper mill in 1978 employed 2,200 workers. Visit,
  March 12,1978.

  红旗造纸厂成立于 1919 年，1978 年雇用了 2,200 名工人。 1978 年 3 月 12
  日访问。}其他人清理厂区并到食堂帮忙。南京军区抽调七十二名团级以上干部，率领数千名部队进入杭州钢铁厂等工厂。过去六个月里，军区士兵在整顿全市高校学生和教师的行为中发挥了作用。\footnote{ZPS,
  August 8,1975, SWB/FE/4981/BII/9-10.

  ZPS，1975 年 8 月 8 日，SWB/FE/4981/BII/9-10。}7月24日中央决定第七条指示学生不得卷入校园外的派系冲突，不得干涉工厂和农村的社会主义教育运动。\footnote{``Decision'',
  p.~85.

  《决定》，第85页。}这是学生和教师卷入以杭州工厂为中心的派系斗争的少数迹象之一。

From the above accounts it is clear that the number of soldiers involved
in this exercise was large indeed. It is likely that at least 30,000
troops were committed, and possibly more. Zweig's hosts at the Hangzhou
silk complex told him that less than 100 soldiers had entered the mill
the previous year.\footnote{Zweig, ``The Peita Debate'', p.~148.

  Zweig，《北大辩论》，第148页。} During my first visit on 12 January
1977 I was told that the number was about 50 to 60. Both figures are
certainly a gross under-estimation. The first report from Hangzhou
announcing the intervention of troops in large industrial enterprises
specifically stated that 6,000 soldiers had occupied four factories, one
of which was the silk complex.\footnote{ZJRB, July 21,1975, p.~1.

  ZJRB，1975年7月21日，第1页。} Therefore, it is logical to assume that
as the mill was one of the worst trouble-spots in the city, it received
at least its quota (1,500) of the 6,000 soldiers. Outside estimates of
the total number of troops sent into Hangzhou's factories varied from a
low of 4,500-6,000 through 10,000 and up to several tens of
thousands.\footnote{Goodstadt, FEER, August 15,1975, p.~25; Charts,
  p.~4.

  Goodstadt，FEER，1975年8月15日，第25页；Charts，第4页。} The PLA's
main functions seem to have been to stop the fighting, provide a safe
environment in which workers could resume production, and cater to their
immediate welfare and leisure requirements. Herein lay the real
significance of the military occupation. Unskilled in modern industrial
techniques, soldiers could not run modern industrial plants but they
could provide the minimum security conditions for the continuation of
normal production. This point was made by a leftist Hong Kong observer
who claimed, not without foundation, that the real mission of the PLA
was to suppress dissidents and malcontents.\footnote{Jun Xing,
  ``杭州发生了什么事？'' (What's happened in Hangzhou?), August 15, 1975
  in 十月评论 (October Review), October 1975, pp.~9-10. The article was
  translated in parts under the title ``Behind the Ferment in
  Hangzhou'', in Intercontinental Press, 13: 41 (1975), pp.~1584-85.

  邢军，《杭州发生了什么事？》，1975年8月15日，载《十月评论》，1975年10月，第9-10页。该文曾以《杭州骚动的背后》为题节译，载
  Intercontinental Press，13:41（1975年），第1584-1585页。}

从以上的记载可以看出，参加这次行动的士兵人数确实不少。可能至少投入了30,000名士兵，甚至可能更多。茨威格在杭州丝绸厂的接待者告诉他，前一年只有不到100名士兵进入工厂。\footnote{Zweig,
  ``The Peita Debate'', p.~148.

  Zweig，《北大辩论》，第148页。}1977年1月12日我第一次访问时，有人告诉我这个数字大约是50到60人。这两个数字肯定是严重低估的。杭州第一份宣布军队介入大型工业企业的报告明确指出，6000名士兵占领了四家工厂，其中一家是丝绸联合厂。\footnote{ZJRB,
  July 21,1975, p.~1.

  ZJRB，1975年7月21日，第1页。}因此，可以合乎逻辑地假设，由于该厂是该市最严重的麻烦点之一，它至少分到了这6000名士兵中的应有份额（1500人）。外界估计，派往杭州工厂的军队总数低则4500-6000人，也有1万人乃至数万人的说法。\footnote{Goodstadt,
  FEER, August 15,1975, p.~25; Charts, p.~4.

  Goodstadt，FEER，1975年8月15日，第25页；Charts，第4页。}解放军的主要职能似乎是停止战斗、为工人恢复生产提供安全环境，并满足他们眼前的福利和休闲需求。军事进驻的真正意义就在于此。士兵不熟悉现代工业技术，无法经营现代工厂，但他们可以为正常生产的延续提供最低限度的安全条件。这是一位香港左派观察人士提出的观点，他声称解放军的真正使命是镇压异见人士和不满分子，并非没有根据。\footnote{Jun
  Xing, ``杭州发生了什么事？'' (What's happened in Hangzhou?), August
  15, 1975 in 十月评论 (October Review), October 1975, pp.~9-10. The
  article was translated in parts under the title ``Behind the Ferment
  in Hangzhou'', in Intercontinental Press, 13: 41 (1975), pp.~1584-85.

  邢军，《杭州发生了什么事？》，1975年8月15日，载《十月评论》，1975年10月，第9-10页。该文曾以《杭州骚动的背后》为题节译，载
  Intercontinental Press，13:41（1975年），第1584-1585页。}

It is uncertain how long the PLA troops remained in the factories, but
it is likely that they remained until the end of the year. A Xinhua
report of late August 1975 stated that 4,500 soldiers and cadres from
the ZPMD continued to work in the city's industrial
enterprises.\footnote{See Xinhua, August 23, 1975,
  SWB/FE/4993/BII/12-13.

  见新华社，1975年8月23日，SWB/FE/4993/BII/12-13。} By mid-October 1975,
the military had focused its attention on three major internal issues:
its subordination to party authority, the need to strengthen both its
own unity and its relations with the government and citizenry, and the
resumption of control over the training of the militia.\footnote{ZJRB,
  October 14,1975, p.~1.

  ZJRB，1975年10月14日，第1页。} In September 1975, on the seventeenth
anniversary of the establishment of the people's militia, Wu Shihong,
Political Commissar of the ZPMD, spoke about the necessity to
``eradicate the influence of bourgeois factionalism in the contingents
of the people's militia''. He emphasized once again the traditional view
of the militia as a ``labour organization as well as a military
organization'' working for the development of the socialist
economy.\footnote{ZJRB, September 30,1975, p.~1; quotation from ZPS,
  September 30,1975, SWB/FE/5026/BII/9-10.

  ZJRB，1975年9月30日，第1页；引文见
  ZPS，1975年9月30日，SWB/FE/5026/BII/9-10。}

目前尚不清楚解放军部队在工厂里待了多久，但很可能一直待到年底。1975年8月下旬的新华社报道称，浙江省军区的4500名士兵和干部继续在该市的工业企业工作。\footnote{See
  Xinhua, August 23, 1975, SWB/FE/4993/BII/12-13.

  见新华社，1975年8月23日，SWB/FE/4993/BII/12-13。}
到1975年10月中旬，军队把注意力集中在三个主要的内部问题上：服从党的权威、需要加强自身团结以及与政府和民众的关系、恢复对民兵训练的控制。\footnote{ZJRB,
  October 14,1975, p.~1.

  ZJRB，1975年10月14日，第1页。}
1975年9月，在人民民兵成立十七周年之际，浙江省军区政委武世鸿在讲话中指出，要``在民兵队伍中铲除资产阶级派性的影响''。他再次强调传统观点，认为民兵是为发展社会主义经济而工作的``劳动组织，也是军事组织''。\footnote{ZJRB,
  September 30,1975, p.~1; quotation from ZPS, September 30,1975,
  SWB/FE/5026/BII/9-10.

  ZJRB，1975年9月30日，第1页；引文见
  ZPS，1975年9月30日，SWB/FE/5026/BII/9-10。}

\section{3. The ``Eight Factories'
Experience''}\label{the-eight-factories-experience}

\section{3. ``八厂经验''}\label{ux516bux5382ux7ecfux9a8c}

In his talks at various industrial enterprises, Wang Hongwen had
continued to keep factional cleavages on the boil. At the meeting which
was convened at the iron and steel mill, Wang asked each participant to
state first which faction they belonged to before making a speech. When
one man declared that he considered himself above all else a
metallurgical worker and a CCP member, Wang thereupon designated him a
member of the ``mountain base'' faction.\footnote{Chen Zuolin, HZRB,
  December 18,1976; HZRB, November 11,1976.

  陈作霖，HZRB，1976年12月18日； HZRB，1976 年 11 月 11 日。} At the
No.~2 cotton mill a representative spoke of the seriousness of the
hooliganism and resultant lack of safety on the job. Hearing these words
Wang reportedly smiled and said: ``I've done that; I've also commanded
fighting''. At the oxygen generator plant Wang sought out two factional
leaders and told them that ``those who got in during the 'two
crashthroughs' must be kicked out''. However, he advised them that they
should continue to promote young cadres, but in a more subtle way so as
to avoid blatantly flouting party regulations.\footnote{Bian Wen, HZRB,
  December 5,1976.

  卞文，HZRB，1976 年 12 月 5 日。} It is evident from these exchanges
that Wang was feeling a great deal of discomfiture at alienating his
erstwhile supporters, and while formally enforcing central directives
was simultaneously attempting to mollify young cadres who were being
forced to stand aside.

王洪文在各工业企业的谈话中，继续使派系分歧保持紧张。在钢铁厂召开的会议上，王要求每位与会者在发言前先说明自己属于哪一派。当一名男子声称自己首先是一名冶金工人、一名中共党员时，王随即将他划为``山下派''成员。\footnote{Chen
  Zuolin, HZRB, December 18,1976; HZRB, November 11,1976.

  陈作霖，HZRB，1976年12月18日； HZRB，1976 年 11 月 11 日。}
在第二棉纺厂，一名代表谈到了流氓行为的严重性以及由此造成的工作中缺乏安全感。据报道，王闻言笑着说：``我干过，也指挥过武斗''。在制氧厂，王找到了两个派系头目，并告诉他们``\,`双突'中进来的人必须被赶出去''。不过，他建议他们继续提拔年轻干部，但要以更微妙的方式，以免公然违反党规。从这些交流中可以看出，王对疏远昔日支持者感到非常不安，在正式执行中央指示的同时，又试图安抚被迫靠边站的年轻干部。\footnote{Bian
  Wen, HZRB, December 5,1976.

  卞文，HZRB，1976 年 12 月 5 日。}

Wang returned to the silk complex in August 1975.\footnote{Visit,
  January 12, 1977.

  访问，1977 年 1 月 12 日。} During a meeting with representatives from
the workers he reportedly asked his audience not to raise wage demands.
Hearing this, a member of the audience stood up and berated Wang thus:

1975年8月，王回到丝绸联合厂。\footnote{Visit, January 12, 1977.

  访问，1977 年 1 月 12 日。}
据报道，在与工人代表会面时，他要求听众不要提出工资要求。听到这话，一名听众站起来这样斥责王：

\begin{quote}
Comrade Vice-chairman, aren't you presently on grade 2 of the central
{[}wage{]} scale? Think back to what grade you were on at the No.~7
textile mill in Shanghai. Was it the 17th or 18th? You rebelled against
Chen Pixian and Cao Diqiu {[}1st Secretary and Mayor respectively of
Shanghai at the beginning of the Cultural Revolution{]} and rose to
become party secretary of the municipality. In the last two years you've
done even better and risen to grade 2 at the center! But you ask us not
to raise the question of wages and promotion in grades. What kind of
work style is this? Ok. So we don't raise the issue and all eat from the
same pot. What do you think of that?\footnote{Zhonggong Nianbao 1976, 5:
  103.

  中国年报1976，5：103。}
\end{quote}

\begin{quote}
副主席同志，你现在不是中央工资二级吗？回想一下你在上海第七纺织厂时是什么工资级别，是十七级还是十八级？你造陈丕显和曹荻秋（文化大革命初期分别担任上海市第一书记和市长）的反，升任市委书记。过去两年你更进一步，升到了中央二级！可你却要求我们不要提出工资和晋级问题。这是什么工作作风？好吧，那我们就不提这个问题，大家都吃大锅饭。你看怎么样？\footnote{Zhonggong
  Nianbao 1976, 5: 103.

  中国年报1976，5：103。}
\end{quote}

In supposed contrast to Wang's subjective and provocative approach, Ji
Dengkui had expressed satisfaction at being informed, on a visit to
Shangwang brigade, that there were no factions among the brigade's
peasants.\footnote{Briefing at Shangwang, October 19, 1977; Wang Jinyou,
  Juqi zhuagang xue Dazhai, p.~40.

  1977年10月19日在上王听取情况介绍；王金友，《举旗抓纲学大寨》，第40页。}
However, in reality the brigade party secretary Wang Jinyou was no
exception when it came to indulging in factionalism and opportunism. In
1977, Wang was lauded as a shining example of a cadre who had defied the
Gang of Four and its followers in Zhejiang. It is true that in January
1976 Wang published an article in the Zhejiang Daily in which he
castigated slogans prevalent in 1974 such as ``don't work for the wrong
political line'', and accused a ``group of swindlers'' of ``practising
retrogression under the guise of opposing restoration''.\footnote{ZJRB,
  January 3,1976, p.~1.

  ZJRB，1976年1月3日，第1页。} Unfortunately for Wang, he published
another article on the very eve of the arrest of the Gang of Four in
which he savagely attacked Deng Xiaoping.\footnote{Wang Jinyou, ``\,`The
  General Programme' is iron-clad proof of Deng Xiaoping's reversal of
  verdicts and restoration'', HZRB, October 6,1976, p.~2 (reprinted in
  edited form from ZJRB).

  王金友，``\,`总纲'是邓小平翻案复辟的铁证''，《杭州日报》，1976年10月6日，第2版（以社论形式转载自《浙江日报》）。}
Even more ironically, the phrases that Wang had used in his January 1976
article were taken directly from the same ``General Program'', which he
condemned in no uncertain terms nine months later.

与王的主观挑衅性做法形成所谓对比的是，纪登奎在视察上王大队时，获悉该大队农民中不存在派系，对此表示满意。\footnote{Briefing
  at Shangwang, October 19, 1977; Wang Jinyou, Juqi zhuagang xue Dazhai,
  p.~40.

  1977年10月19日在上王听取情况介绍；王金友，《举旗抓纲学大寨》，第40页。}然而，实际上，在纵容派系主义和机会主义方面，大队党支部书记王金友并非例外。1977年，王被称赞为反抗浙江``四人帮''及其追随者的干部光辉典范。确实，1976年1月，王在《浙江日报》上发表文章，严厉斥责1974年盛行的``不要为错误的政治路线做事''等口号，指责``骗子集团''\,``打着反对复辟的幌子搞倒退''。\footnote{ZJRB,
  January 3,1976, p.~1.

  ZJRB，1976年1月3日，第1页。}不幸的是，王在四人帮被捕前夕又发表了一篇文章，猛烈攻击邓小平。\footnote{Wang
  Jinyou, ``\,`The General Programme' is iron-clad proof of Deng
  Xiaoping's reversal of verdicts and restoration'', HZRB, October
  6,1976, p.~2 (reprinted in edited form from ZJRB).

  王金友，``\,`总纲'是邓小平翻案复辟的铁证''，《杭州日报》，1976年10月6日，第2版（以社论形式转载自《浙江日报》）。}更具有讽刺意味的是，王在1976年1月的文章中使用的措辞直接取自同一份《总纲》，而他在9个月后毫不含糊地谴责了该纲领。

From Ji and Wang's investigations of Hangzhou's industrial sector, seven
factories were singled out for special attention. Another unit, the
Hangzhou Gearbox Works (杭州齿轮箱厂), was hailed as being worthy of
emulation. The lessons to be drawn from the positive experience of the
gearbox works and the negative experience of the other seven plants were
summed up and became known as the ``eight factories' experience''
(八个工厂的经验). The other seven plants were the silk complex, the
oxygen generator plant, the hemp mill, the silk brocade mill, the iron
and steel mill, the No.~1 cotton printing and dyeing mill and the No.~2
cotton mill. The provincial leadership widely publicized the example of
the eight factories throughout Hangzhou and Zhejiang for the remainder
of 1975.

从纪和王对杭州工业部门的调查来看，有七家工厂受到了特别关注。另一个单位------杭州齿轮箱厂，被誉为值得效仿的单位。杭州齿轮箱厂的积极经验和其他七家工厂的消极经验被总结出来，被称为``八个工厂的经验''。其他七家工厂分别是丝绸厂、制氧厂、麻厂、丝锦厂、钢铁厂、第一棉印染厂和第二棉纺厂。
1975年余下的时间里，省领导层在杭州及浙江全省广泛宣传``八厂经验''。

Apart from the Hangzhou No.~2 Cotton Mill situated in Xiaoshan, I
visited all the factories -- the brocade mill twice and the silk complex
three times - during my two and a half years in Hangzhou. I was able to
question cadres about events of the previous years and speak to
participants such as Zhang Jifa at the iron and steel mill and Shen
Qiying and Shen Chuyun at the silk complex. Four of the six factories
which had suffered production losses in the years 1974-76 (the gearbox
plant being an outstanding exception) provided me with the following
figures:

在杭州两年半的时间里，除了位于萧山的杭州第二棉纺厂外，我参观了所有的工厂，织锦厂两次，丝绸厂三次。我向干部询问了往年发生的事情，并与钢铁厂的张纪法、丝绸厂的沈其英、沈楚云等参与者交谈。1974-76年遭受生产损失的六家工厂中有四家（齿轮箱厂是一个突出的例外）向我提供了以下数据：

\begin{itemize}
\item
  Zhejiang Hemp Mill -- production losses of 66 million yuan (1974-76),
  equivalent to the cost of building two mills;
\item
  浙江麻厂------生产损失6600万元（1974-76年），相当于建造两座麻厂的费用；
\item
  Hangzhou No.~1 Cotton Printing and Dyeing Mill -- production losses of
  170 million yuan and profits foregone of 54 million yuan, equivalent
  to the cost of 2 mills;
\item
  杭州第一棉印染厂------生产损失1.7亿元，利润损失5400万元，相当于2个棉厂的成本；
\item
  Hangzhou Iron and Steel Mill -- total output for 1974-76 equal to that
  of 1973. Per capita output value 1973 was 10,500 yuan while in 1976 it
  had fallen to 4,300 yuan. Profit losses for 1974-6 totalled 80 million
  yuan;
\item
  杭州钢铁厂------1974-76年总产量与1973年持平。1973年人均产值为10500元，1976年已降至4300元。
  1974-6年利润损失总计8000万元；
\item
  Hangzhou Oxygen Generator Plant -- 1974-76 production losses totalled
  67 million yuan, equal to the cost of building one and a half mills of
  this size. In 1974 production dropped 62.7\% from 1973, in 1975 it
  increased 101\% over 1974, and in 1976 it dropped by 46.6\%.
\item
  杭州制氧机厂------1974-76年生产损失总计6700万元，相当于建造一个半这个规模的工厂的成本。
  1974年产量比1973年下降62.7\%，1975年比1974年增长101\%，1976年下降46.6\%。
\end{itemize}

The People's Daily published a laudatory piece on the Hangzhou Gearbox
Factory in mid-August 1975.\footnote{RMRB, August 15,1975, p.~1.
  Sections of this article appeared in SWB/FE/4988/BII/1-2. See also
  HZRB, April 10,1977. Construction of the factory commenced in 1960 and
  it went into production, turning out marine gearboxes, in 1965. In
  1977 it employed a staff of 2,869. Visit, July 17,1977.

  RMRB，1975年8月15日，第1页。本文部分内容见 SWB/FE/4988/BII/1-2。另参见
  HZRB，1977年4月10日。该厂于1960年开始建设，1965年投产，生产船用齿轮箱。1977年，该厂有职工2869人。1977年7月17日访问。}
Zhejiang Daily had published an editorial in a similar vein on 13 July
1975 entitled ``Adhere to Party Spirit and Principle -- Learn from the
Comrades of the Hangzhou Gearbox Works''.\footnote{ZJRB, July 13,1975,
  p.~1.

  ZJRB，1975年7月13日，第1页。} Yet an account published in Taiwan
listed the Hangzhou Gearbox Works as one of the factories worst affected
by what it described as strikes.\footnote{Zhonggong Nianbao 1976, V:103.

  中国年报1976，V：103。} The Chinese press made it abundantly clear,
however, that the factory was an outstanding example of a unit which
continued to meet its production targets amidst the trouble which
surrounded it. What distinguished the gearbox plant from the seven other
factories was that it had met its production targets in the first half
of 1975. Since 1968 the factory had achieved outstanding results, with
production increasing at an average annual rate of 21\%. It had thus
become a model unit and an advanced enterprise in learning from the
nationally famous Daqing oil-field. The People's Daily article
attributed the factory's success to its determination not to weaken or
deviate from the leadership of the CCP. Such tendencies had not occurred
even in the campaign against Lin Biao and Confucius of the previous
year. All attempts to sow dissension in the ranks of cadres and workers
had failed. Although the leaders were from different generations they
studied and worked together. Workers observed discipline, met state
targets and actively refuted and exposed political fallacies by writing
wall posters. Or so it was claimed.

1975年8月中旬，《人民日报》发表了一篇赞扬杭州齿轮箱厂的文章。\footnote{RMRB,
  August 15,1975, p.~1. Sections of this article appeared in
  SWB/FE/4988/BII/1-2. See also HZRB, April 10,1977. Construction of the
  factory commenced in 1960 and it went into production, turning out
  marine gearboxes, in 1965. In 1977 it employed a staff of 2,869.
  Visit, July 17,1977.

  RMRB，1975年8月15日，第1页。本文部分内容见 SWB/FE/4988/BII/1-2。另参见
  HZRB，1977年4月10日。该厂于1960年开始建设，1965年投产，生产船用齿轮箱。1977年，该厂有职工2869人。1977年7月17日访问。}
1975年7月13日，《浙江日报》也发表了一篇类似的社论，题为《坚持党性和原则------向杭州齿轮箱厂的同志学习》。\footnote{ZJRB,
  July 13,1975, p.~1.

  ZJRB，1975年7月13日，第1页。}
然而，台湾的一篇报道将杭州齿轮箱厂列为受其所谓``罢工''影响最严重的工厂之一。然而，中国媒体明确表示，该工厂是在困难重重的情况下继续实现生产目标的杰出典范。杭州齿轮箱厂与其他七家工厂的不同之处在于，它在1975年上半年就完成了生产目标。自1968年以来，该厂取得了骄人的业绩，产量以年均21\%的速度增长。成为学习全国著名大庆油田的模范单位和先进企业。
《人民日报》的文章将该工厂的成功归功于其不削弱或背离中共领导的决心。这种倾向，即使在去年的反林反孔运动中也没有出现过。一切在干部和工人队伍中挑拨离间的企图都失败了。尽管领导人来自不同代，但他们一起学习、一起工作。工人们遵守纪律，完成国家指标，通过写墙报积极驳斥和揭露政治谬论。或者说是这样声称的。\footnote{Zhonggong
  Nianbao 1976, V:103.

  中国年报1976，V：103。}

The gearbox factory was praised as a unit which had kept factional
squabbles outside its gates. To learn from this plant meant, in the
words of another report, to ``go all-out to build socialism''. Issues
which had been raised by the radicals such as bourgeois rights were not
taken seriously. In words which could have come straight from the mouth
of Deng Xiaoping, young workers at the gear-box factory were informed
that ``it is essential to have high political consciousness and a sound
material foundation in order to do away with bourgeois rights and
realize communism''.\footnote{ZPS, August 15, 1975, SWB/FE/4988/BII/2-3.

  ZPS，1975 年 8 月 15 日，SWB/FE/4988/BII/2-3。} Other factories and
industrial unit were urged to study the plant's policy of building a
strong, united leading group which saw its main task as working for
socialism.\footnote{See ZPS, August 2, 1975, SWB/FE/4981/BII/12; ZPS,
  August 14, 1975, SWB/FE/4988/BII/3-4; ZPS, August 21,1975,
  SWB/FE/4993/BII/1-2; ZPS, August 27, 1975, SWB/FE/4996/BII/10-12; ZPS,
  September 6,1975, SWB/FE/5006/BII/16-17.

  参见 ZPS，1975 年 8 月 2 日，SWB/FE/4981/BII/12； ZPS，1975 年 8 月 14
  日，SWB/FE/4988/BII/3-4； ZPS，1975 年 8 月 21
  日，SWB/FE/4993/BII/1-2； ZPS，1975 年 8 月 27
  日，SWB/FE/4996/BII/10-12； ZPS，1975 年 9 月 6
  日，SWB/FE/5006/BII/16-17。}

杭州齿轮箱厂被誉为``把派系争斗拒之门外''的单位。向这家工厂学习，用另一份报告的话说，就是``全力以赴建设社会主义''。激进分子提出的资产阶级法权等问题没有得到认真对待。齿轮箱厂的青年工人被告知，``要取消资产阶级法权，实现共产主义，就必须有高度的政治觉悟和坚实的物质基础''。\footnote{ZPS,
  August 15, 1975, SWB/FE/4988/BII/2-3.

  ZPS，1975 年 8 月 15 日，SWB/FE/4988/BII/2-3。}这句话简直像是出自邓小平之口。其他工厂和工业单位被敦促学习该厂建立坚强团结领导班子的方针；这个领导班子把为社会主义工作视为主要任务。\footnote{See
  ZPS, August 2, 1975, SWB/FE/4981/BII/12; ZPS, August 14, 1975,
  SWB/FE/4988/BII/3-4; ZPS, August 21,1975, SWB/FE/4993/BII/1-2; ZPS,
  August 27, 1975, SWB/FE/4996/BII/10-12; ZPS, September 6,1975,
  SWB/FE/5006/BII/16-17.

  参见 ZPS，1975 年 8 月 2 日，SWB/FE/4981/BII/12； ZPS，1975 年 8 月 14
  日，SWB/FE/4988/BII/3-4； ZPS，1975 年 8 月 21
  日，SWB/FE/4993/BII/1-2； ZPS，1975 年 8 月 27
  日，SWB/FE/4996/BII/10-12； ZPS，1975 年 9 月 6
  日，SWB/FE/5006/BII/16-17。}

The second factory of the eight, the Hangzhou Silk Complex, probably
provided the litmus test for the success of Deng's assault on bourgeois
factionalism. By late August 1975 the provincial radio could report
minor production increases, enhanced unity among the mill's cadres, a
certain level of harmony among its work force and the unfolding of a
campaign to ``scathingly criticize bourgeois factionalism and all
anti-Marxist fallacies''.\footnote{ZPS, August 22, 1975,
  SWB/FE/4993/BII/8-9.

  ZPS，1975 年 8 月 22 日，SWB/FE/4993/BII/8-9。} Weng Senhe's followers
had been placed under supervision in study classes and their behavior
had been repudiated at criticism meetings. Production in August 1975
exceeded that of July by 43.29\%, and was ahead of the monthly
target.\footnote{ZJRB, September 2,1975, p.~1; ZPS, September 2,3, 1975,
  SWB/FE/FE/5002/BII/2-3.

  ZJRB，1975年9月2日，第1页；ZPS，1975年9月2、3日，SWB/FE/FE/5002/BII/2-3。}
However, rather than reflecting a high output per se, this figure
indicated the depths to which production had previously sunk. In the
first half of 1975 the Hangzhou Silk Complex had fulfilled a miserable
18.5\% of its yearly plan.\footnote{Visit to mill, January 12, 1977.

  1977 年 1 月 12 日参观工厂。}

八家工厂中的第二家，杭州丝绸联合厂，可能为邓小平打击资产阶级派性的成功提供了试金石。到1975年8月下旬，省广播电台可以报道生产小幅增长、该厂干部团结增强、职工队伍有了一定程度的和谐，以及``狠狠批判资产阶级派性和一切反马克思主义谬论''的运动正在展开。\footnote{ZPS,
  August 22, 1975, SWB/FE/4993/BII/8-9.

  ZPS，1975 年 8 月 22 日，SWB/FE/4993/BII/8-9。}
翁森鹤的追随者已被置于学习班中接受监督，他们的行为也在批判会上遭到否定。1975年8月的产量比7月高出43.29\%，并超过了月度目标。\footnote{ZJRB,
  September 2,1975, p.~1; ZPS, September 2,3, 1975,
  SWB/FE/FE/5002/BII/2-3.

  ZJRB，1975年9月2日，第1页；ZPS，1975年9月2、3日，SWB/FE/FE/5002/BII/2-3。}
然而，这个数字并不反映产量本身很高，而是表明此前产量已经降到何等地步。1975年上半年，杭州丝绸联合厂只完成了年度计划的可怜的18.5\%。\footnote{Visit
  to mill, January 12, 1977.

  1977 年 1 月 12 日参观工厂。}

A Xinhua report of the following week attributed the turnaround in
production to the central decision of 24 July.\footnote{Xinhua,
  September 10,1975, in RMRB, September 11, 1975, p.~1 transl. in
  SWB/FE/5010/BII/8-9.

  新华社，1975年9月10日，载 RMRB，1975年9月11日，第1页；译文见
  SWB/FE/5010/BII/8-9。} It asserted that workers and cadres had come to
reject erroneous slogans which had been put forward by ``some bad
individuals'' who had claimed that factional contradictions or
contradictions between new and old cadres had replaced class
contradictions and had tried to establish their ``independence'' from
the Party. A further Xinhua report of October 1975 announced that
although the mill had overfulfilled its production target in September,
its employees had continued working even on National Day.\footnote{Xinhua,
  October 2, 1975, SWB/FE/5026/BII/2.

  新华社，1975 年 10 月 2 日，SWB/FE/5026/BII/2。}

接下来一周的新华社报道将生产的好转归功于7月24日的中央决定。\footnote{Xinhua,
  September 10,1975, in RMRB, September 11, 1975, p.~1 transl. in
  SWB/FE/5010/BII/8-9.

  新华社，1975年9月10日，载 RMRB，1975年9月11日，第1页；译文见
  SWB/FE/5010/BII/8-9。}文章声称，工人和干部已经开始拒绝接受``一些坏人''提出的错误口号，这些人声称派系矛盾或新老干部之间的矛盾已经取代了阶级矛盾，并试图建立自己独立于党的地位。1975年10月新华社的进一步报道称，尽管该厂超额完成了9月份的生产目标，但其员工即使在国庆节仍继续工作。\footnote{Xinhua,
  October 2, 1975, SWB/FE/5026/BII/2.

  新华社，1975 年 10 月 2 日，SWB/FE/5026/BII/2。}

On 30 September, Zhejiang Daily carried a lengthy article concerning the
silk complex.\footnote{``杭丝联开展'一学四批五大讲'自我教育运动经验''
  (The Hangzhou silk complex unfolds experience in the self-education
  movement of ``one study, four criticisms and five emphases'') ZJRB,
  September 30, 1975, in Xuexi Wenji, pp.~70-8. I am grateful to David
  Zweig for making extracts from Xuexi Wenji available to me.

  《杭丝联开展``一学四批五大讲''自我教育运动经验》（The Hangzhou silk
  complex unfolds experience in the self-education movement of ``one
  study, four criticisms and five
  emphases''），ZJRB，1975年9月30日，收入《学习文集》，第70-78页。感谢
  David Zweig 向我提供《学习文集》的摘录。} The article mentioned
``hauling out bad people'' (揪出了坏人) at the plant and excoriated the
evils of bourgeois factionalism. It referred to the ``self-education
movement'' of ``one study, four criticisms and five emphases'' underway
at the mill -- the study of the theory of the dictatorship of the
proletariat, the criticism of revisionism, capitalism, the concept of
bourgeois rights and bourgeois factionalism, and the emphasis on the
line, the interests of the whole, party spirit, unity and discipline. By
mentioning both bourgeois factionalism and bourgeois rights the article
was catering to the priorities of both the radicals and the modernizers.
The study movement received great emphasis in the local press during the
following months.\footnote{See CNS, 592 (November 26,1975).

  参见 CNS，592（1975 年 11 月 26 日）。}

9月30日，《浙江日报》发表了一篇有关丝绸联合厂的长文。\footnote{``杭丝联开展'一学四批五大讲'自我教育运动经验''
  (The Hangzhou silk complex unfolds experience in the self-education
  movement of ``one study, four criticisms and five emphases'') ZJRB,
  September 30, 1975, in Xuexi Wenji, pp.~70-8. I am grateful to David
  Zweig for making extracts from Xuexi Wenji available to me.

  《杭丝联开展``一学四批五大讲''自我教育运动经验》（The Hangzhou silk
  complex unfolds experience in the self-education movement of ``one
  study, four criticisms and five
  emphases''），ZJRB，1975年9月30日，收入《学习文集》，第70-78页。感谢
  David Zweig 向我提供《学习文集》的摘录。}
文章提到该厂``揪出了坏人''，痛斥资产阶级派性的罪恶。该厂正在开展``一学四批五大讲''的``自我教育运动''，即学习无产阶级专政理论，批判修正主义、资本主义、资产阶级法权观念和资产阶级派性，讲路线、讲全局利益、讲党性、讲团结、讲纪律。通过同时提及资产阶级派性和资产阶级法权，这篇文章迎合了激进派和现代化派双方的优先事项。在接下来的几个月里，学习运动在当地媒体中得到大力强调。\footnote{See
  CNS, 592 (November 26,1975).

  参见 CNS，592（1975 年 11 月 26 日）。}

On 4 October 1975 the Zhejiang Daily editorial department invited
leading cadres of factories in Hangzhou to a meeting to discuss the silk
complex's progress. The participants enthusiastically praised the mill
for its achievements. On 6 October, the newspaper featured an article
calling on all industrial and mining enterprises in the province to
understand and absorb the reasons for the mill's rejuvenation, because
the lesson applied equally to them.\footnote{ZJRB, October 6,1975,
  pp.~1,4.

  ZJRB，1975 年 10 月 6 日，第 1,4 页。} The provincial public security
bureau, for one, believed that the mill's struggle against bourgeois
factionalism was relevant to it.\footnote{ZJRB, October 14,1975, p.~1.

  ZJRB，1975年10月14日，第1页。} This was probably due to the fact that
Huang Yintang, Weng Senhe's close associate and party secretary at the
silk complex, had also held a position in the public security system.
The provincial radio also reported a conference held by provincial
departments in charge of silk production -- another area where Weng had
exercised considerable influence -- at which representatives from the
silk complex and other enterprises conveyed news of their improved
situation.\footnote{See ZJRB, September 19, 1975, pp.~1, 2.

  参见 ZJRB，1975 年 9 月 19 日，第 1、2 页。}

1975年10月4日，浙江日报编辑部邀请杭州各工厂领导干部开会，讨论丝绸联合厂的进展。与会人员对该厂取得的成就表示热烈赞扬。10月6日，该报发表文章，呼吁全省工矿企业了解并吸收该厂复兴的原因，因为这一经验同样适用于它们。\footnote{ZJRB,
  October 6,1975, pp.~1,4.

  ZJRB，1975 年 10 月 6 日，第 1,4 页。}
省公安厅也认为，该厂反资产阶级派性斗争与自身有关。\footnote{ZJRB,
  October 14,1975, p.~1.

  ZJRB，1975年10月14日，第1页。}
这很可能是因为翁森鹤的亲密同伙、丝绸联合厂党委书记黄阴堂也曾在公安系统任职。省广播电台还报道了省级丝绸生产主管部门召开的一次会议，这是翁曾施加相当影响的另一个领域；会上，来自丝绸联合厂和其他企业的代表传达了局势好转的消息。\footnote{See
  ZJRB, September 19, 1975, pp.~1, 2.

  参见 ZJRB，1975 年 9 月 19 日，第 1、2 页。}

Another factory which had experienced production difficulties as a
result of factional disturbances was the Hangzhou No.~1 Cotton Printing
and Dyeing Mill. Sun Baolong, a close supporter of Weng Senhe, was a
cadre at the mill. He also held the posts of Deputy-director of the CCP
HMC Organization Department and was a member of both the HMRC standing
committee and the CCP ZPC. In late July 1975, a Xinhua report about the
mill was published on the front page of the People's Daily.\footnote{RMRB,
  July 28,1975, p.~1.

  RMRB，1975年7月28日，第1页。} The article repeated references to the
detrimental influence of bourgeois factionalism and mistaken thinking
about production schedules. It also mentioned that the workers were
participating in voluntary labor and attending study classes, while some
had been selected to form a theoretical contingent of over one hundred
members. Disunity in this factory was not confined to the ranks of
blue-collar workers; administrative staff were also affected. The
article cited the case of two leading cadres in the mill's repair and
maintenance section who had previously been at loggerheads. Their
failure to get along had resulted in serious repercussions for
production in the section. After the military occupation they had
allegedly put aside their differences and thrown their energies into a
major renovation project. By soliciting the views of workers and
encouraging them to make an input into the decision-making process, the
party branch of the section was able to solve outstanding problems. It
also required cadres to perform the dirtiest tasks such as cleaning
machinery.

另一家因派系骚乱而生产困难的工厂是杭州第一棉印染厂。翁森鹤的亲密支持者孙宝龙是该厂的一名干部。他还曾担任中共杭州市委组织部副部长，并且是杭州市革命委员会常委和中共浙江省委委员。1975年7月下旬，《人民日报》头版发表了新华社有关该厂的报道。文章中多次提到资产阶级派性的有害影响和关于生产计划的错误思想。还提到，工人们参加义务劳动、参加学习班，有的还被选拔组成了一百多人的理论队伍。这家工厂的不团结并不仅限于蓝领工人队伍，行政人员也受到影响。文章援引了该厂维修保养科两名领导干部此前发生争执的事例。他们的不和给该部门的生产造成了严重影响。据称，军事进驻后，他们搁置了分歧，将精力投入到一项重大改造工程中。该科党支部通过征求职工意见，并鼓励职工参与决策过程，解决了突出问题。它还要求干部承担清洁机器等最脏的任务。\footnote{RMRB,
  July 28,1975, p.~1.

  RMRB，1975年7月28日，第1页。}

Following the arrival of troops in July 1975, the mill organized a
discussion of the CCP constitution and party building principles for its
employees.\footnote{ZPS, August 21, 1975, SWB/FE/FE/4993/BII/5.

  ZPS，1975 年 8 月 21 日，SWB/FE/FE/4993/BII/5。} Production rose in
July and August 1975 as the disruptive elements in the mill were
disciplined and isolated in study classes.\footnote{ZJRB, September
  13,1975, p.~2.

  ZJRB，1975年9月13日，第2页。} In 1979 I was told that over 100 men in
the mill had joined the militia command in 1974 and had been absent from
the shop-floor for over a year. At first they received normal pay but
later, after strong resentment was voiced by their work-mates, their pay
was stopped.\footnote{Visit, April 1, 1979.

  1979 年 4 月 1 日访问。} Production continued to rise into
October.\footnote{ZJRB, October 25,1975, p.~2.

  ZJRB，1975年10月25日，第2页。} The No.~2 cotton mill in Xiaoshan also
conducted education in the procedure for admitting new members into the
CCP.\footnote{ZJRB, August 19,1975, pp.~1,4.

  ZJRB，1975 年 8 月 19 日，第 1,4 页。}

1975年7月部队到来后，厂里组织职工学习党章和党建原则。\footnote{ZPS,
  August 21, 1975, SWB/FE/FE/4993/BII/5.

  ZPS，1975 年 8 月 21 日，SWB/FE/FE/4993/BII/5。}1975年7月和8月，厂里的破坏分子在学习班里受到处分并被隔离，产量有所上升。\footnote{ZJRB,
  September 13,1975, p.~2.

  ZJRB，1975年9月13日，第2页。}1979年，我被告知，厂里有100多人在1974年加入民兵指挥部，已经离开车间超过一年。起初工资正常，后来因工友强烈不满，停发了工资。\footnote{Visit,
  April 1, 1979.

  1979 年 4 月 1 日访问。}进入10月份产量继续上升。\footnote{ZJRB,
  October 25,1975, p.~2.

  ZJRB，1975年10月25日，第2页。}萧山第二棉纺厂也进行了入党程序教育。\footnote{ZJRB,
  August 19,1975, pp.~1,4.

  ZJRB，1975 年 8 月 19 日，第 1,4 页。}

Other factories among the octet joined in the chorus of glad tidings one
by one. The oxygen generator plant declared unspecified production
increases in August 1975.\footnote{ZJRB, August 21,1975, p.~1.

  ZJRB，1975年8月21日，第1页。} In the ``two crashthroughs'' at this
plant, 90 people were ``crash admitted'' into the CCP and 120 ``crash
promoted'' as cadres.\footnote{Visit, June 4, 1978.

  1978 年 6 月 4 日访问。} The silk brocade mill, where fifty workers
had been rushed into the CCP,\footnote{Visit, September 24, 1977.

  访问，1977 年 9 月 24 日。} reported a staggering 100\% increase in
output in the first half of August over the same period of the previous
month.\footnote{ZJRB, August 26,1975, pp.~1,4.

  ZJRB，1975 年 8 月 26 日，第 1,4 页。} Workers formerly at loggerheads
had made up their differences. Shifts were commencing on time, and young
workers remained at their posts for the entire eight-hour shift. Such
reports provided valuable insights into the low level of discipline in
industrial concerns. According to a young man who joined the industrial
workforce in Hangzhou in 1975, the above reports about improvements in
work discipline were most probably gross exaggerations and blatant
propaganda. He informed me that workers at his unit commenced and
knocked off when they pleased and did very little work when they were
actually present. The situation did not change appreciably during the
campaign to learn from the ``eight factories' experience''.

八家工厂中的其他工厂也一一加入到喜讯的合唱中。制氧机厂于1975年8月宣布产量有所增加，但未说明具体幅度。\footnote{ZJRB,
  August 21,1975, p.~1.

  ZJRB，1975年8月21日，第1页。}在该厂的``双突''中，90人被突击吸收入党，120人被突击提拔为干部。\footnote{Visit,
  June 4, 1978.

  1978 年 6 月 4 日访问。}丝织厂有50名工人被突击吸收入党，\footnote{Visit,
  September 24, 1977.

  访问，1977 年 9 月 24 日。}并报告8月上半月产量比上月同期惊人地增长了100\%。\footnote{ZJRB,
  August 26,1975, pp.~1,4.

  ZJRB，1975 年 8 月 26 日，第 1,4 页。}以前不和的工人已经和解。班次准时开始，青年工人在整个八小时班次中都坚守岗位。此类报道为了解工业企业中纪律水平低下提供了宝贵线索。据一位1975年在杭州参加工业劳动力队伍的年轻人说，上述关于改进工作纪律的报道很可能是严重夸大和赤裸裸的宣传。他告诉我，他所在单位的工人想什么时候上班、什么时候下班都随意，真正到岗时也几乎不做什么工作。在学习``八厂经验''运动中，这种情况并没有明显改变。

The situation at the iron and steel mill seemed to take longer to
improve. However, by October 1975, superficially amazing production
increases of 200\% in one month were receiving great
publicity.\footnote{ZJRB, October 27,1975, p.~1.

  ZJRB，1975年10月27日，第1页。} The workers of one shift achieved the
highest output recorded in the rolling mill since 1966.\footnote{ZPS,
  November 7, 1975, SWB/FE/5062/BII/7-8.

  ZPS，1975 年 11 月 7 日，SWB/FE/5062/BII/7-8。} A total of 134 people
by-passed normal procedures to join the CCP or gain promotion as cadres.
Two hundred and fifty workers {[}perhaps a conservative figure{]} joined
the urban militia command and continued to receive their pay while
absent from work.\footnote{Visit, May 27, 1979.

  访问，1979 年 5 月 27 日。} None of the contemporary reports mentioned
the eighth factory, the Zhejiang Hemp Mill. Thus, the situation there
was unclear.

钢铁厂的情况似乎需要更长的时间才能改善。然而，到了1975年10月，表面上惊人的一个月内产量增加200\%的消息得到广泛宣传。\footnote{ZJRB,
  October 27,1975, p.~1.

  ZJRB，1975年10月27日，第1页。}一个班次的工人在轧钢厂创下了自1966年以来的最高产量。\footnote{ZPS,
  November 7, 1975, SWB/FE/5062/BII/7-8.

  ZPS，1975 年 11 月 7 日，SWB/FE/5062/BII/7-8。}共有134人绕过正常程序加入中共或被提拔为干部。两百五十名工人（也许是保守数字）加入了城市民兵指挥部，在旷工期间继续领取工资。\footnote{Visit,
  May 27, 1979.

  访问，1979 年 5 月 27 日。}当时的报道没有提到第八家工厂浙江麻厂。因此，那里的情况并不明朗。

A twenty-nine strong delegation from five of the eight factories -- the
gearbox factory, silk complex, No.~1 cotton mill, hemp mill and oxygen
generator factory -- went to Beijing to report directly to central
leaders for the 1975 National Day celebrations. The delegation toured
factories in the capital before returning to Hangzhou on 11 October. A
rally attended by 100,000 people welcomed them back and it was announced
that industrial production for the third quarter of 1975 had increased
9.1\% over the previous quarter and 7.3\% over the corresponding period
in 1974. Output figures for September 1975 were 8.6\% above those for
August.\footnote{HZRB, September 29, 1975, October 12, 1975; ZPS,
  September 29, 1975, SWB/FE/5026/BII/2-3; ZPS, October 12,1975,
  SWB/FE/5034/BII/6-8. See also ZPS, October 13,1975,
  SWB/FE/5038/BII/7-8. At the beginning of September 1975 Hangzhou had
  held a three-day party-staged celebration as good news on the
  industrial front in the city started to come in. See ZPS, September 5,
  1975, SWB/FE/5006/BII/14-16.

  HZRB，1975年9月29日，1975年10月12日； ZPS，1975 年 9 月 29
  日，SWB/FE/5026/BII/2-3； ZPS，1975 年 10 月 12
  日，SWB/FE/5034/BII/6-8。另见 ZPS，1975 年 10 月 13
  日，SWB/FE/5038/BII/7-8。
  1975年9月初，杭州举行了为期三天的党庆活动，该市工业战线上的好消息开始传来。见ZPS，1975年9月5日，SWB/FE/5006/BII/14-16。}

八家工厂中的五家------齿轮箱厂、丝绸厂、第一棉纺厂、麻纺厂和制氧机厂------组成了二十九人代表团，在1975年国庆庆祝期间前往北京，直接向中央领导人汇报。代表团参观了首都的工厂，然后于10月11日返回杭州。
10万人参加的集会欢迎他们回来，并宣布1975年第三季度工业生产比上一季度增长9.1\%，比1974年同期增长7.3\%。1975年9月的产出数字比8月增长8.6\%。\footnote{HZRB,
  September 29, 1975, October 12, 1975; ZPS, September 29, 1975,
  SWB/FE/5026/BII/2-3; ZPS, October 12,1975, SWB/FE/5034/BII/6-8. See
  also ZPS, October 13,1975, SWB/FE/5038/BII/7-8. At the beginning of
  September 1975 Hangzhou had held a three-day party-staged celebration
  as good news on the industrial front in the city started to come in.
  See ZPS, September 5, 1975, SWB/FE/5006/BII/14-16.

  HZRB，1975年9月29日，1975年10月12日； ZPS，1975 年 9 月 29
  日，SWB/FE/5026/BII/2-3； ZPS，1975 年 10 月 12
  日，SWB/FE/5034/BII/6-8。另见 ZPS，1975 年 10 月 13
  日，SWB/FE/5038/BII/7-8。
  1975年9月初，杭州举行了为期三天的党庆活动，该市工业战线上的好消息开始传来。见ZPS，1975年9月5日，SWB/FE/5006/BII/14-16。}

\section{Political and Economic Factors behind the Hangzhou
Incident}\label{political-and-economic-factors-behind-the-hangzhou-incident}

\section{杭州事件背后的政治和经济因素}\label{ux676dux5ddeux4e8bux4ef6ux80ccux540eux7684ux653fux6cbbux548cux7ecfux6d4eux56e0ux7d20}

The events of 1975 were remarkable even by the standards of the
extraordinary political upheavals regularly experienced in Maoist China.
The decision to implement a wholesale transfer and despatch of troops
into the factories of a major urban center was not one that could have
been taken lightly. The thorough shake-up of the provincial and
municipal administrations reflected the gravity of the situation.
Production shortfalls in key industrial plants in Hangzhou, caused
mainly by drawn-out debilitating factional struggles, had become
chronic. The reported incident in May 1975 at the provincial
broadcasting station illustrated the extreme lengths to which opponents
of the party authorities were prepared to go to attain their goals.

即使以毛泽东时代中国经常经历的非同寻常的政治动乱来衡量，1975
年发生的事件也是引人注目的。向主要城市中心的工厂大规模转移和派遣部队的决定并不是一个可以轻易做出的决定。省、市行政机构的彻底改组，反映了事态的严重性。杭州主要工厂的生产短缺已成为长期现象，这主要是由旷日持久的派系斗争造成的。据报道，1975
年 5
月发生在省广播电台的事件表明，党当局的反对者为了达到他们的目标，不惜采取极端的手段。

Contemporary observers of Chinese affairs in Hong Kong and Taiwan, both
Chinese and Western alike, were alert to the significance of the
situation. Most commented on the almost unprecedented intervention of
military forces to settle industrial disturbances. Above all, this
aspect of the matter forced the incident into the news headlines and
raised questions about the stability of the country. Hangzhou, in its
troubles, supposedly reflected in extremis, the general malaise of a
regime beset by leadership disunity and popular unrest. Overall,
however, as many instants cited above reveal, the literature available
on the incident is disappointing for its factual inaccuracies, political
distortion, over-emphasis on wage disputes, and vague generalizations.

当代香港和台湾的中国事务观察家，无论是中国还是西方，都对这一局势的重要性保持警惕。大多数人评论说，为解决工业骚乱而采取的军事干预几乎是史无前例的。最重要的是，此事的这一方面迫使该事件成为新闻头条，并引发了人们对国家稳定的质疑。杭州的困境据说反映在极端情况下，即一个政权因领导不团结和民众骚乱而普遍萎靡不振。然而，总体而言，正如上面引用的许多实例所揭示的那样，有关该事件的现有文献因其事实不准确、政治扭曲、过分强调工资纠纷和模糊概括而令人失望。

The New York Times correspondent filed a story regarding the despatch of
PLA troops into Hangzhou's factories and the replacement of key
officials in the military and civilian administration.\footnote{New York
  Times, July 29, 1975, pp.~1, 6.

  《纽约时报》，1975 年 7 月 29 日，第 1、6 页。} He argued that the
broadcasts from Hangzhou, which openly admitted that there was trouble
in the city, were only released because of Beijing's confidence that the
situation had been brought under control. The use of troops, he wrote,
was a warning to other troubled areas that the central authorities would
not tolerate disorder.

《纽约时报》记者发稿报道了解放军派遣部队进驻杭州工厂以及更换军民部门主要官员一事。
他辩称，来自杭州的广播公开承认杭州出现了麻烦，只是因为北京相信局势已得到控制而发布。他写道，使用军队是对其他动乱地区的警告，中央当局不会容忍骚乱。\footnote{New
  York Times, July 29, 1975, pp.~1, 6.

  《纽约时报》，1975 年 7 月 29 日，第 1、6 页。}

Leo Goodstadt, in four consecutive articles published Hong Kong in
August and September 1975 drew similar conclusions.\footnote{Goodstadt,
  FEER, August 1,1975, pp.~20-1; FEER, August 15,1975, pp.~24-6;
  ``Hangchow, and crises still to come'', FEER, 29 August 1975,
  pp.~27-8; ``Beneath the `Water Margin'\,'', FEER, 19 September 1975,
  pp.~13-15.

  Goodstadt，FEER，1975 年 8 月 1 日，第 20-1 页； FEER，1975 年 8 月 15
  日，第 24-6 页； ``杭州，危机仍将来临''，FEER，1975 年 8 月 29 日，第
  27-8 页； ``\,`水浒'之下''，FEER，1975 年 9 月 19 日，第 13-15 页。}

Leo
Goodstadt在1975年8月和9月在香港发表的连续四篇文章中得出了类似的结论。\footnote{Goodstadt,
  FEER, August 1,1975, pp.~20-1; FEER, August 15,1975, pp.~24-6;
  ``Hangchow, and crises still to come'', FEER, 29 August 1975,
  pp.~27-8; ``Beneath the `Water Margin'\,'', FEER, 19 September 1975,
  pp.~13-15.

  Goodstadt，FEER，1975 年 8 月 1 日，第 20-1 页； FEER，1975 年 8 月 15
  日，第 24-6 页； ``杭州，危机仍将来临''，FEER，1975 年 8 月 29 日，第
  27-8 页； ``\,`水浒'之下''，FEER，1975 年 9 月 19 日，第 13-15 页。}

\begin{quote}
The lesson of Hangzhou {[}he wrote{]} was supposed to be that the entire
central leadership, with Mao's blessing, had decided to use the city as
an example of their common determination to tolerate no use of arms or
mobs to seize power locally or to intimidate local leaders.\footnote{Goodstadt,
  FEER, August 29,1975, p.~28.

  Goodstadt，FEER，1975年8月29日，第28页。}
\end{quote}

\begin{quote}
杭州的教训应该是，整个中央领导层在毛泽东的支持下，决定以这座城市作为他们共同决心的榜样，绝不容忍使用武器或暴徒在当地夺取政权或恐吓地方领导人。\footnote{Goodstadt,
  FEER, August 29,1975, p.~28.

  Goodstadt，FEER，1975年8月29日，第28页。}
\end{quote}

Goodstadt claimed that the authorities would not have used troops ``if
genuine differences of opinion about policy or Maoist principles'' had
caused the strife. Instead, he argued, Beijing placed the blame on
``individuals whose stars had waned since the end of the Cultural
Revolution'' and who were causing trouble to achieve their
``self-serving ambitions''.\footnote{Goodstadt, FEER, August 1,1975,
  p.~30.

  Goodstadt，FEER，1975年8月1日，第30页。} He concluded that ``the body
politic itself has weakened'' and that factional strife represented real
dangers to the continuation of the socialist system in China.

古德施塔特声称，``如果对政策或毛主义原则存在真正的意见分歧''导致了冲突，当局就不会使用军队。相反，他认为，北京将责任归咎于``自文化大革命结束以来星光黯淡的个人''以及为实现``自私野心''而制造麻烦的人。\footnote{Goodstadt,
  FEER, August 1,1975, p.~30.

  Goodstadt，FEER，1975年8月1日，第30页。}
他总结说，``政治体本身已经削弱''，派性冲突对中国社会主义制度的延续构成真实危险。

Other commentators, particularly those in Taiwan predicted, on the basis
of the Hangzhou incident, the imminent collapse of the communist regime.
Goodstadt noted these prophecies and observed that ``There is a
dangerous temptation to see Hangchow as a massive crisis indicative of a
breakdown of law and order throughout China''.\footnote{Goodstadt, FEER,
  August 15,1975, p.~25.

  Goodstadt，FEER，1975年8月15日，第25页。} He wisely cautioned that
``The problems of Hangchow should not be over-dramatised, any more than
the problems of other Chinese cities which have provided tasty fare for
a variety of anti-Chinese propaganda sources''.\footnote{Ibid.

  同上。}

其他评论家，特别是台湾的评论家，根据杭州事件预测，共产党政权即将崩溃。古德施塔特注意到这些预言，并指出，``将杭州视为一场表明整个中国法律和秩序崩溃的大规模危机，是一种危险的诱惑''。\footnote{Goodstadt,
  FEER, August 15,1975, p.~25.

  Goodstadt，FEER，1975年8月15日，第25页。}
他明智地提醒说，``杭州的问题不应被过度戏剧化，正如其他中国城市的问题也不应被过度戏剧化一样；这些问题为各种反华宣传来源提供了可口材料''。\footnote{Ibid.

  同上。}

The question of the relative weight which should be attached to wage and
other economic grievances in causing the trouble in Zhejiang requires
clarification. The preface to the 24 July Central Committee and State
Council Decision had referred to the evil trend of economism
(经济主义妖风) in Zhejiang. Point seven of the document stipulated that
unresolved wage disputes would be subjected to ``further study and
consultation''.\footnote{``Decision'', p.~85.

  《决定》，第85页。} As chapter eight of this study has demonstrated,
during the final years of the Maoist era the CCP was loathe to recognize
the validity of, or accede to, wage demands. Wage grievances certainly
cut across and exacerbated factional differences but did not create
them. Differences on the shop-floor derived primarily from political
rather than economic factors and the economic grievances stemmed mainly
from the tensions which had been caused by factional disputes.
Nevertheless, the extent to which wage demands created unrest in the
factories of Hangzhou in 1975 has been widely debated.

工资和其他经济不满在浙江造成麻烦时应给予相对重视的问题需要澄清。
7月24日中央、国务院决定的序言提到了浙江的经济主义妖风。文件第七点规定，未解决的工资纠纷将受到``进一步研究和协商''。\footnote{``Decision'',
  p.~85.

  《决定》，第85页。}
正如本研究第八章所示，在毛泽东时代的最后几年，中共不愿承认或同意工资要求的有效性。工资不满肯定会跨越并加剧派别分歧，但并没有造成派别分歧。车间的分歧主要源于政治而非经济因素，而经济上的不满主要源于派系争端造成的紧张局势。然而，1975
年工资要求在多大程度上引发了杭州工厂的骚乱，这一点引起了广泛的争论。

Radio Moscow cited Western reports of wall posters in Hangzhou
expressing dissatisfaction over the issue.\footnote{Radio Moscow, May
  22,1975, SWB/FE/4916/BII/18.

  莫斯科广播电台，1975 年 5 月 22 日，SWB/FE/4916/BII/18。} Labor unrest
over wages, argued Taiwan sources, was a major reason for the trouble in
Hangzhou. This explanation has also been accepted by Western
academics.\footnote{See for example, Tony Saich, ``Workers in the
  Workers' State: Urban Workers in the PRC'', in David S. G. Goodman
  (ed.). Groups and Politics in the People's Republic of China (Cardiff:
  University College Cardiff Press, 1984), p.~159; Dittmer, China's
  Continuous Revolution, pp.~165-66.

  例如，参见 Tony
  Saich，《工人国家中的工人：中华人民共和国的城市工人》，载 David S. G.
  Goodman（编），《中华人民共和国的群体与政治》（卡迪夫：卡迪夫大学学院出版社，1984年），第159页；Dittmer，《中国的持续革命》，第165-166页。}
If complaints over wages caused or were a major factor behind the
problems in Hangzhou, it is unclear why trouble on a similar scale did
not erupt all over China. Why, if gripes about wages and working
conditions had brought production to a halt in major industrial
enterprises, did the most severe outbreak of dissatisfaction occur in
Hangzhou? What set the city apart from other industrial centers in the
country and necessitated such drastic action by Beijing? Goodstadt
suggested a higher percentage of the residents of Hangzhou were employed
in the industrial sector than in cities such as Beijing or Shanghai.
Additionally, he speculated that its citizens, with their heterogeneous
origins, lacked a sense of community and that urban facilities such as
housing were in short supply.\footnote{Goodstadt, FEER, September
  19,1975, p.14.

  Goodstadt，FEER，1975年9月19日，第14页。} Although this argument
deserves consideration as an interesting foray into the field of urban
sociology, Goodstadt's hypotheses appear somewhat fanciful. The housing
problem in Shanghai was undoubtedly worse than in Hangzhou and major
inland Chinese cities contained populations of even more recent and
disparate background.

莫斯科广播电台援引西方有关杭州大字报的报道，称大字报表达了对这一问题的不满。\footnote{Radio
  Moscow, May 22,1975, SWB/FE/4916/BII/18.

  莫斯科广播电台，1975 年 5 月 22 日，SWB/FE/4916/BII/18。}台湾消息人士称，工资问题引发的劳工骚乱是杭州问题的主要原因。这种解释也被西方学者所接受。\footnote{See
  for example, Tony Saich, ``Workers in the Workers' State: Urban
  Workers in the PRC'', in David S. G. Goodman (ed.). Groups and
  Politics in the People's Republic of China (Cardiff: University
  College Cardiff Press, 1984), p.~159; Dittmer, China's Continuous
  Revolution, pp.~165-66.

  例如，参见 Tony
  Saich，《工人国家中的工人：中华人民共和国的城市工人》，载 David S. G.
  Goodman（编），《中华人民共和国的群体与政治》（卡迪夫：卡迪夫大学学院出版社，1984年），第159页；Dittmer，《中国的持续革命》，第165-166页。}如果说工资投诉是杭州问题的主要原因，那么类似规模的问题为何没有在全国范围内爆发，就不得而知了。如果对工资和工作条件的抱怨导致主要工业企业停产，为什么最严重的不满情绪却发生在杭州？是什么让这座城市与中国其他工业中心区分开来，并需要北京采取如此严厉的行动？古德施塔特表示，杭州居民在工业部门就业的比例高于北京或上海等城市。此外，他推测其市民由于来源不同，缺乏社区意识，住房等城市设施供应不足。\footnote{Goodstadt,
  FEER, September 19,1975, p.14.

  Goodstadt，FEER，1975年9月19日，第14页。}虽然这一论点值得考虑，可视为对城市社会学领域的一次有趣尝试，但古德施塔特的假设似乎有些异想天开。上海的住房问题无疑比杭州更严重，而中国内陆主要城市的人口背景甚至更新近、更杂异。

A contemporary monograph on China's political economy, in its passing
reference to the ``Hangzhou Incident'', dismissed the view that wage
disputes caused the trouble. It argued that workers' complaints were
directed more at the inadequacy of collective welfare facilities and the
failure of cadres to participate in collective labor. The author raised
secondary issues such as the heat, harsh working conditions, intense
production schedules, poor eating facilities, unhygienic sanitary
conditions and family problems.\footnote{Steven Andors, China's
  Industrial Revolution: Politics, Planning and Management, 1949 to the
  Present (London: Pantheon Books, 1977), pp.~234-35

  Steven Andors，《中国的工业革命：政治、规划和管理，1949
  年至今》（伦敦：Pantheon Books，1977 年），第 234-35 页} However, the
workers of Hangzhou were not alone in experiencing harsh working
conditions. The Hong Kong leftist Jun Xing attributed the incident in
Hangzhou to worker discontent with frozen wages, forced ``voluntary''
labor and poor living and working conditions. He argued that the workers
of Hangzhou were rebelling against exploitation by the CCP, against the
far more comfortable living conditions enjoyed by officials and against
the falseness of the extravagant propaganda which contrasted so starkly
with the reality engulfing and oppressing them. Jun Xing argued that the
constant unleashing of political campaigns only served to disrupt
production, resulting in increased pressure on workers to make up for
the losses in subsequent production drives. The pattern had recurred in
the summer of 1975, with even radical ideologues such as Yao Wenyuan
stressing unity so that China could safely navigate the political
storm.\footnote{Jun Xing, ``What's happened in Hangzhou?''.

  邢军，《杭州发生了什么事？》。}

一本当代中国政治经济学专著，在提及``杭州事件''时，驳斥了工资纠纷引发麻烦的观点。该书认为，工人的抱怨更多是针对集体福利设施不足、干部不参加集体劳动等问题。作者提出了一些次要问题，例如炎热、工作条件恶劣、生产安排紧张、饮食设施差、卫生条件不卫生以及家庭问题等。\footnote{Steven
  Andors, China's Industrial Revolution: Politics, Planning and
  Management, 1949 to the Present (London: Pantheon Books, 1977),
  pp.~234-35

  Steven Andors，《中国的工业革命：政治、规划和管理，1949
  年至今》（伦敦：Pantheon Books，1977 年），第 234-35 页}然而，并不是只有杭州的工人正在经历恶劣的工作条件。香港左翼人士邢军将杭州事件归因于工人对冻结工资、强迫``自愿''劳动以及恶劣的生活和工作条件的不满。他认为，杭州工人正在反抗中共的剥削，反抗官员享有的舒适得多的生活条件，反抗浮夸宣传的虚假性；这些宣传与吞噬和压迫他们的现实形成鲜明对比。邢军认为，不断发动政治运动只会扰乱生产，导致工人在后续生产活动中弥补损失的压力更大。这种模式在1975年夏天再次出现，甚至像姚文元这样的激进思想家也强调团结，以便中国能够安全地渡过这场政治风暴。\footnote{Jun
  Xing, ``What's happened in Hangzhou?''.

  邢军，《杭州发生了什么事？》。}

An analysis of the events from Taiwan portrayed the incident as the most
conspicuous example of labor unrest in China.\footnote{Li Ming-hua,
  ``Observations on the Recent Labor Unrest in Mainland China'', I\&S,
  11:10 (1975), pp.~2-15.

  李明华，《近年中国大陆劳工骚乱观察》，I\&S，11:10 (1975)，第2-15页。}
The author ascribed the unrest to a serious reduction in agricultural
and industrial production which he attributed both to natural disasters
and to political and social causes. He instanced the campaign against
Lin Biao and Confucius in 1974 as an example of such disruption and
contended that workers and peasants had responded to economic hardship
by going out on strike, deliberately ``going slow'' or by opposing
requisition policies. The erosion of discipline in the Cultural
Revolution had enabled workers to free themselves from party control and
press their economic demands. The CCP was faced with no option but to
tighten discipline, reassert its leadership and subjugate the rebellious
workers. That factional disputes involved both political and economic
matters was clearly recognized by the idiosyncratic Hong Kong weekly
China News Analysis in two successive issues published in September
1975.\footnote{CNA, 1012 and 1013.

  CNA，1012 和 1013。} The journal traced the roots of the Hangzhou
incident to the disputes of the previous year involving the urban
militia. The author of the articles also analyzed the affair in the
context of strikes and absenteeism among China's industrial workers.

一篇来自台湾的事件分析将此事件描述为中国劳工骚乱中最引人注目的例子。\footnote{Li
  Ming-hua, ``Observations on the Recent Labor Unrest in Mainland
  China'', I\&S, 11:10 (1975), pp.~2-15.

  李明华，《近年中国大陆劳工骚乱观察》，I\&S，11:10 (1975)，第2-15页。}作者将这次骚乱归因于农业和工业生产的严重下降，并认为这种下降既由自然灾害造成，也有政治和社会原因。他以1974年批林批孔运动作为这种扰乱的一个例子，并声称工人和农民通过罢工、故意``磨洋工''或反对征购政策来应对经济困难。文化大革命中纪律的侵蚀使工人们能够摆脱党的控制并提出经济要求。中共别无选择，只能加强纪律，重申领导，压服造反工人。1975年9月，风格独特的香港周刊《中国新闻分析》连续两期明确承认派系纠纷涉及政治和经济问题。\footnote{CNA,
  1012 and 1013.

  CNA，1012 和 1013。}该杂志将杭州事件的根源追溯到前一年涉及城市民兵的纠纷。文章作者还将这一事件放在中国产业工人罢工和旷工的背景下进行分析。

It is highly likely that higher paid members of the industrial working
class resented and even resisted any drastic move in 1975 to make the
wage system more egalitarian, a concern shared by their peers across
China. However, Hangzhou's troubles stemmed essentially and primarily
from political causes. Industrial production dropped sharply in the
years 1974-76 mainly due to factional fighting and the subsequent high
rate of absenteeism on the shop-floor.\footnote{See Zweig, ``The Peita
  Debate'', p.~148.

  参见 Zweig，《北大辩论》，第148页。} The suggestion that workers
enrolled in the militia demanded wage rises is unfounded.\footnote{Zhonggong
  Nianbao 1976, 5: 100-05.

  中国年报1976，5：100-05。} In fact, militiamen received additional
entitlements for their duties, a practice strongly objected to by their
fellow workers.

产业工人阶级中收入较高的成员很可能对1975年任何旨在使工资制度更加平等的重大举措表示不满，甚至抵制，这是中国各地同类工人共同关心的问题。然而，杭州的问题本质上主要源于政治原因。1974-76年间，工业生产急剧下降，主要是由于派系斗争和随后车间高旷工率。\footnote{See
  Zweig, ``The Peita Debate'', p.~148.

  参见 Zweig，《北大辩论》，第148页。}参加民兵的工人要求提高工资的说法是没有根据的。\footnote{Zhonggong
  Nianbao 1976, 5: 100-05.

  中国年报1976，5：100-05。}事实上，民兵履行职责时获得了额外待遇，这种做法遭到了工友的强烈反对。

Industrial unrest was part of the strategy of sabotage and
destabilization orchestrated from Beijing by Wang Hongwen in particular
and implemented in Hangzhou by Weng Senhe and his associates. The
resultant slump in output would reflect badly on the provincial and
municipal party leaderships and the leaders of factories hostile to the
exponents of ``bourgeois factionalism'', where disruption was greatest.
Naturally, workshop conditions deteriorated if machinery lay idle or
damaged and logistic staff were unable or unwilling to cook meals or run
creches. Intimidation and threats kept many workers away from their
units and those who did attend worked in a very insecure environment.
Ultimately, however, the relative emphasis that is placed on political
or economic factors reveals more about the perspective of the observer
than the incident itself.

工业骚乱是王洪文在北京精心策划、翁森鹤及其同伙在杭州实施的破坏和破坏稳定战略的一部分。由此造成的产量下降，会对破坏最严重地区那些最敌视``资产阶级派性''鼓吹者的省市党领导层和工厂领导造成严重影响。当然，如果机器闲置或损坏，而后勤人员无法或不愿意做饭或经营托儿所，车间条件就会恶化。恐吓和威胁使许多工人远离自己的单位，而那些参加的人在非常不安全的环境中工作。然而，归根结底，对政治或经济因素的相对重视更多地揭示了观察者的观点，而不是事件本身。

\section{Notes}\label{notes-8}

\section{注释}\label{ux6ce8ux91ca-8}

\bookmarksetup{startatroot}

\chapter{Chapter Ten - CONSOLIDATION AND RENEWED INSTABILITY, AUGUST
1975-SEPTEMBER 1976 / 第十章 -
巩固与重新出现的不稳定，1975年8月至1976年9月}\label{chapter-ten---consolidation-and-renewed-instability-august-1975-september-1976-ux7b2cux5341ux7ae0---ux5de9ux56faux4e0eux91cdux65b0ux51faux73b0ux7684ux4e0dux7a33ux5b9a1975ux5e748ux6708ux81f31976ux5e749ux6708}

\section{Upholding the Line of Deng Xiaoping in Zhejiang,
August-December
1975}\label{upholding-the-line-of-deng-xiaoping-in-zhejiang-august-december-1975}

\section{在浙江维护邓小平路线，1975年8月至12月}\label{ux5728ux6d59ux6c5fux7ef4ux62a4ux9093ux5c0fux5e73ux8defux7ebf1975ux5e748ux6708ux81f312ux6708}

By 9 August 1975 the situation in Hangzhou appeared sufficiently stable
for the CCP ZPC to convene a work conference under the guidance of Wang
Hongwen and Ji Dengkui.\footnote{ZJRB, August 10, 1975, p.~1; CNS, 578,
  pp.~6-7. ZJRB，1975年8月10日，第1页；CNS，第578期，第6-7页。} The
report of the conference made no mention of Ji and Wang's presence at
the conference but it is certain that they did attend.\footnote{Ji and
  Wang were both public figures with high profiles. Yet Ji did not make
  any public appearances between 11 June and 29 August 1975 and it is
  fair to surmise that he spent much of this time in Hangzhou or dealing
  with the ``Zhejiang problem''. Wang's activities were not reported in
  the Chinese press between 20 May and 31 July 1975, when he appeared at
  the PLA Day rally in Shanghai, and then disappeared from public view
  again until 30 September. CIA Reference Aid, Appearances and
  Activities of Leading Personalities of the PRC, 1975, CR 76-10980
  (March 1976), p.~87.
  纪和王都是知名度很高的公众人物。然而，纪在1975年6月11日至8月29日之间没有公开露面，可以合理推测，他在这段时间的大部分里待在杭州，或在处理``浙江问题''。王在中国报刊中的活动报道则从1975年5月20日至7月31日中断；7月31日他出席上海建军节集会后，又再次从公众视野中消失，直到9月30日。CIA
  Reference Aid, Appearances and Activities of Leading Personalities of
  the PRC, 1975, CR 76-10980（1976年3月），第87页。} The main theme of
the work conference concerned the questions of party leadership and the
unity of the working class. Deng Xiaoping's instructions were conveyed
regarding the establishment of strong, unified and militant leading
groups, which placed ``daring'' first (敢字当头)\footnote{Gao Gao and
  Yan Jiaqi, ``Wenhua dageming'', p.~545.
  高皋、严家其，《``文化大革命''》，第545页。} and possessed the
determination to overcome bourgeois factionalism. Deng was making a
deliberate and striking reference to Point 3 of the CC's August 1966 16
point manifesto launching the Cultural Revolution, which had called on
the party leadership to ``put daring above everything else and boldly
arouse the masses''.\footnote{See Important Documents on the Great
  Proletarian Revolution, pp.~133, 137. 见 Important Documents on the
  Great Proletarian Revolution，第133、137页。} Delegates to the
conference visited the silk complex to see first hand the rapid changes
that had occurred since Weng's arrest. Both the ``two crashthroughs''
and the party's class line in the period of socialism were subjected to
close scrutiny. While the center's decisive handling of the ``Zhejiang
problem'' may have created a mood of optimism, delegates were reminded
that ``The struggle has not yet come to an end''. That factionalism had
seriously affected industrial enterprises was highlighted by the
quotation of Mao's 1967 statement that there was no fundamental conflict
of interests within the working class.
到1975年8月9日，杭州的局势似乎已足够稳定，中共浙江省委得以在王洪文和纪登奎指导下召开一次工作会议。\footnote{ZJRB,
  August 10, 1975, p.~1; CNS, 578, pp.~6-7.
  ZJRB，1975年8月10日，第1页；CNS，第578期，第6-7页。}
会议报告没有提到纪、王出席会议，但可以确定他们确实参加了。\footnote{Ji
  and Wang were both public figures with high profiles. Yet Ji did not
  make any public appearances between 11 June and 29 August 1975 and it
  is fair to surmise that he spent much of this time in Hangzhou or
  dealing with the ``Zhejiang problem''. Wang's activities were not
  reported in the Chinese press between 20 May and 31 July 1975, when he
  appeared at the PLA Day rally in Shanghai, and then disappeared from
  public view again until 30 September. CIA Reference Aid, Appearances
  and Activities of Leading Personalities of the PRC, 1975, CR 76-10980
  (March 1976), p.~87.
  纪和王都是知名度很高的公众人物。然而，纪在1975年6月11日至8月29日之间没有公开露面，可以合理推测，他在这段时间的大部分里待在杭州，或在处理``浙江问题''。王在中国报刊中的活动报道则从1975年5月20日至7月31日中断；7月31日他出席上海建军节集会后，又再次从公众视野中消失，直到9月30日。CIA
  Reference Aid, Appearances and Activities of Leading Personalities of
  the PRC, 1975, CR 76-10980（1976年3月），第87页。}
这次工作会议的主题主要涉及党的领导和工人阶级团结问题。会议传达了邓小平关于建立坚强、统一、有战斗力的领导班子的指示，这些领导班子要把``敢''字放在首位（敢字当头）\footnote{Gao
  Gao and Yan Jiaqi, ``Wenhua dageming'', p.~545.
  高皋、严家其，《``文化大革命''》，第545页。}，并具备克服资产阶级派性的决心。邓有意而醒目地指涉中央1966年8月发动文化大革命的十六条中的第三条；该条曾要求党的领导``敢字当头，放手发动群众''。\footnote{See
  Important Documents on the Great Proletarian Revolution, pp.~133, 137.
  见 Important Documents on the Great Proletarian
  Revolution，第133、137页。}
与会代表参观了丝绸联合企业，亲眼了解翁森鹤被捕后那里发生的迅速变化。``双突''和党在社会主义时期的阶级路线都受到仔细审视。尽管中央对``浙江问题''的果断处理可能造成了一种乐观气氛，但代表们仍被提醒：``斗争还没有结束''。会议引用毛泽东1967年关于工人阶级内部没有根本利害冲突的论断，突出说明派性已严重影响工业企业。

The success of the consolidation program in Zhejiang depended heavily on
the continuation of the detente between Deng and the radicals in
Beijing. It was also dependent on both local factions upholding the
tacit cease-fire which had been imposed on them. A provincial broadcast
of 18 August urged its listeners to ``see the conformity between
obedience to the leadership of the Party Center and that to the
leadership of the local Party Committees''
(认识服从党中央领导和服从地方党委领导的一致性). As China News Summary
commented: ``This and other reports have indicated that the radicals,
although they respect the Party Center, seek to defy orders from local
Party organizations''.\footnote{CNS, 579, p.~1, fn.
  CNS，第579期，第1页，脚注。} However, unity at the center was
transient and fragile as events were soon to prove. In August 1975
enigmatic comments by Mao Zedong on the classic novel Water Margin
warning of the danger of capitulationism were published in the Chinese
press, initiating yet another politically divisive campaign.\footnote{For
  analyses of the campaign, see Merle Goldman, China's Intellectuals,
  pp.~201-211; Fenwick, ``The Gang of Four'', pp.~355-66; CNS, 581 (3
  September 1975), 582 (September 10, 1975), 584 (September 24, 1975);
  Gao Gao and Yan Jiaqi, ``Wenhua dageming'', pp.~557-60; Hao Mengbi and
  Duan Haoran, Liushinian, p.~647; and for a view of the place of the
  Water Margin in Chinese mythology, see John Fitzgerald, ``Continuity
  Within Discontinuity-The Case of Water Margin Mythology'', Modern
  China, 12: 3 (July 1986), pp.~361-400. 关于这场运动的分析，见 Merle
  Goldman, China's Intellectuals，第201-211页；Fenwick, ``The Gang of
  Four''，第355-66页；CNS，第581期（1975年9月3日）、第582期（1975年9月10日）、第584期（1975年9月24日）；高皋、严家其，《``文化大革命''》，第557-60页；郝梦笔、段浩然，《六十年》，第647页；关于《水浒传》在中国神话中地位的看法，见
  John Fitzgerald, ``Continuity Within Discontinuity-The Case of Water
  Margin Mythology'', Modern China, 12:3（1986年7月），第361-400页。} On
2 September 1975, Politburo member and Director of the CC general
office, Wang Dongxing, made a speech in Guangzhou which further
illustrated the tenuous nature of central unity.\footnote{See Feiqing
  Yuebao, 18: 7 (1976), pp.~94-6. 见《匪情月报》18:7（1976），第94-6页。}
浙江整顿方案的成功，在很大程度上取决于邓小平与北京激进派之间缓和关系的延续，也取决于当地两派遵守强加给他们的默示停火。8月18日的一次省级广播敦促听众``认识服从党中央领导和服从地方党委领导的一致性''。正如
China News Summary
评论的那样：``这些报道和其他报道表明，激进派虽然尊重党中央，却试图违抗地方党组织的命令。''\footnote{CNS,
  579, p.~1, fn. CNS，第579期，第1页，脚注。}
然而，中央的团结转瞬即逝且十分脆弱，事态很快就会证明这一点。1975年8月，中国报刊发表了毛泽东关于古典小说《水浒传》的含义隐晦的评论，警告投降主义的危险，由此又发起了一场造成政治分裂的运动。\footnote{For
  analyses of the campaign, see Merle Goldman, China's Intellectuals,
  pp.~201-211; Fenwick, ``The Gang of Four'', pp.~355-66; CNS, 581 (3
  September 1975), 582 (September 10, 1975), 584 (September 24, 1975);
  Gao Gao and Yan Jiaqi, ``Wenhua dageming'', pp.~557-60; Hao Mengbi and
  Duan Haoran, Liushinian, p.~647; and for a view of the place of the
  Water Margin in Chinese mythology, see John Fitzgerald, ``Continuity
  Within Discontinuity-The Case of Water Margin Mythology'', Modern
  China, 12: 3 (July 1986), pp.~361-400. 关于这场运动的分析，见 Merle
  Goldman, China's Intellectuals，第201-211页；Fenwick, ``The Gang of
  Four''，第355-66页；CNS，第581期（1975年9月3日）、第582期（1975年9月10日）、第584期（1975年9月24日）；高皋、严家其，《``文化大革命''》，第557-60页；郝梦笔、段浩然，《六十年》，第647页；关于《水浒传》在中国神话中地位的看法，见
  John Fitzgerald, ``Continuity Within Discontinuity-The Case of Water
  Margin Mythology'', Modern China, 12:3（1986年7月），第361-400页。}
1975年9月2日，政治局委员、中央办公厅主任汪东兴在广州发表讲话，进一步说明中央团结的脆弱性质。\footnote{See
  Feiqing Yuebao, 18: 7 (1976), pp.~94-6.
  见《匪情月报》18:7（1976），第94-6页。}

At this juncture Deng Xiaoping appeared to believe that he had
consolidated his position sufficiently to launch an all-out assault on
his opponents.\footnote{For a thorough discussion of Deng's 1975
  strategy, see Fenwick, ``The Gang of Four'', pp.~311-32.
  关于邓1975年战略的详尽讨论，见 Fenwick, ``The Gang of
  Four''，第311-32页。} Despite being aware of the probability that the
Water Margin campaign was directed at him personally, in talks of
September and October 1975 Deng accused those behind the campaign of
plotting against him.\footnote{Gao Gao and Yan Jiaqi, ``Wenhua
  dageming'', p.~565. 高皋、严家其，《``文化大革命''》，第565页。}
Undeterred, he pressed ahead with his program of all-round adjustment
and economic development. Deng inspired his supporters by telling them
not to be afraid of being knocked down a second time.
在这个关头，邓小平似乎认为自己的地位已经巩固到足以向对手发动全面攻势。\footnote{For
  a thorough discussion of Deng's 1975 strategy, see Fenwick, ``The Gang
  of Four'', pp.~311-32. 关于邓1975年战略的详尽讨论，见 Fenwick, ``The
  Gang of Four''，第311-32页。}
尽管他知道《水浒》运动很可能是针对他个人的，邓在1975年9月和10月的谈话中仍指责这场运动的幕后人物正在阴谋反对他。\footnote{Gao
  Gao and Yan Jiaqi, ``Wenhua dageming'', p.~565.
  高皋、严家其，《``文化大革命''》，第565页。}
他没有退缩，而是继续推进全面调整和经济发展的方案。邓鼓舞支持者，告诉他们不要害怕第二次被打倒。

During the months of August, September and October 1975, Deng was busy
overseeing the preparation of three documents which encapsulated his
basic position. Maligned by his radical opponents as the ``three
poisonous weeds'', the documents were entitled ``On the General Program
of Work for the Whole Party and the Whole Country'', ``Some Problems in
Accelerating Industrial Development'' (or the ``20 Points'') and ``On
Some Problems in the Fields of Science and Technology''.\footnote{See
  Zhonggong nianbao, 1977, V: 58-80, transl. in Chi Hsin, The Case of
  the Gang of Four, pp.~203-38, 239-72, 277-86. For earlier drafts of
  the second and third programmatic documents see Chinese Law and
  Government, 12: 1-2 (1979), pp.~61-74, 112-126.
  见《中共年报》1977年，第五卷，第58-80页；英译见 Chi Hsin, The Case of
  the Gang of
  Four，第203-38、239-72、277-86页。第二和第三份纲领性文件的早期草稿，见
  Chinese Law and Government, 12:1-2（1979），第61-74、112-126页。}
Earlier in the year Deng had established the groups responsible for the
drafting of these documents, listened to reports on their progress and
issued instructions to guide the labours of his subordinates, who
included Hu Yaobang, Hu Qiaomu and Deng Liqun.\footnote{See Deng
  Xiaoping, ``关于发展工业的几点意见'' (Some opinions on the development
  of industry), 18 August 1975, Deng Xiaoping Wenxuan, pp.~28-31; Deng's
  comments of September 26, 1975, have been translated in Chi Hsin, The
  Case of the Gang of Four pp.~287-95. See also, ``Wenhua dageming'' de
  guomin jingji, pp.~171-2; Gao Gao and Yan Jiaqi, ``Wenhua dageming'',
  pp.~548-52. It is claimed that Deng never saw the first document until
  his return to power in 1977 and only sighted an early draft of the
  second. See ``Wenhua dageming'' ruogan shijian zhenxiang, p.~135. Be
  that as it may, Deng undoubtedly knew and approved of work on these
  papers and their general orientation was very much in line with his
  thinking. 见邓小平《关于发展工业的几点意见》（Some opinions on the
  development of
  industry），1975年8月18日，《邓小平文选》，第28-31页；邓1975年9月26日的评论已译载于
  Chi Hsin, The Case of the Gang of
  Four，第287-95页。另见《``文化大革命''的国民经济》，第171-2页；高皋、严家其，《``文化大革命''》，第548-52页。据称，邓直到1977年复出后才见到第一份文件，并且只看过第二份文件的早期草稿。见《``文化大革命''若干事件真相》，第135页。尽管如此，邓无疑知道并赞同这些文件的起草工作，其总体取向也非常符合他的思想。}
In addition to proposing policies which were directly counterposed to
those inspired by the Cultural Revolution, the General Program and the
``20 points'' in particular conducted an open debate with the radical
mobilizers. The two documents took up many issues which had been raised
in the campaign to study the theory of the dictatorship of the
proletariat, especially in the articles published by Zhang Chunqiao and
Yao Wenyuan.
1975年8月、9月和10月，邓忙于主持起草三份概括其基本立场的文件。这些文件被他的激进派对手诋毁为``三株大毒草''，题为《关于全党全国各项工作的总纲》《关于加快工业发展的若干问题》（又称``二十条''）和《关于科技工作的几个问题》。\footnote{See
  Zhonggong nianbao, 1977, V: 58-80, transl. in Chi Hsin, The Case of
  the Gang of Four, pp.~203-38, 239-72, 277-86. For earlier drafts of
  the second and third programmatic documents see Chinese Law and
  Government, 12: 1-2 (1979), pp.~61-74, 112-126.
  见《中共年报》1977年，第五卷，第58-80页；英译见 Chi Hsin, The Case of
  the Gang of
  Four，第203-38、239-72、277-86页。第二和第三份纲领性文件的早期草稿，见
  Chinese Law and Government, 12:1-2（1979），第61-74、112-126页。}
当年早些时候，邓已经成立负责起草这些文件的小组，听取起草进展汇报，并发布指示指导下属工作，其中包括胡耀邦、胡乔木和邓力群。\footnote{See
  Deng Xiaoping, ``关于发展工业的几点意见'' (Some opinions on the
  development of industry), 18 August 1975, Deng Xiaoping Wenxuan,
  pp.~28-31; Deng's comments of September 26, 1975, have been translated
  in Chi Hsin, The Case of the Gang of Four pp.~287-95. See also,
  ``Wenhua dageming'' de guomin jingji, pp.~171-2; Gao Gao and Yan
  Jiaqi, ``Wenhua dageming'', pp.~548-52. It is claimed that Deng never
  saw the first document until his return to power in 1977 and only
  sighted an early draft of the second. See ``Wenhua dageming'' ruogan
  shijian zhenxiang, p.~135. Be that as it may, Deng undoubtedly knew
  and approved of work on these papers and their general orientation was
  very much in line with his thinking.
  见邓小平《关于发展工业的几点意见》（Some opinions on the development
  of
  industry），1975年8月18日，《邓小平文选》，第28-31页；邓1975年9月26日的评论已译载于
  Chi Hsin, The Case of the Gang of
  Four，第287-95页。另见《``文化大革命''的国民经济》，第171-2页；高皋、严家其，《``文化大革命''》，第548-52页。据称，邓直到1977年复出后才见到第一份文件，并且只看过第二份文件的早期草稿。见《``文化大革命''若干事件真相》，第135页。尽管如此，邓无疑知道并赞同这些文件的起草工作，其总体取向也非常符合他的思想。}
除了提出与文化大革命所启发的政策直接对立的政策外，尤其是《总纲》和``二十条''还与激进动员派展开公开论辩。这两份文件讨论了学习无产阶级专政理论运动中提出的许多问题，特别是张春桥和姚文元发表的文章中提出的问题。

Of the senior members of the Zhejiang leadership, it was Tie Ying who
took up the cudgels for Deng Xiaoping in the province. On 1 September
1975 the CCP ZPC called a meeting of provincial and municipal office
cadres to convey central directives. Tan Qilong presided over the
meeting and Deputy-secretary Chen Weida conveyed the contents of
speeches by unnamed central leaders.\footnote{ZJRB, September 5, 1975,
  pp.~1, 4. ZJRB，1975年9月5日，第1、4页。} In his address, Tie affirmed
the correctness of the political line and policies which had been
adopted at the August work conference, but he noted that implementation
of the policies had been uneven. He pointed out that a ``resolute
struggle must be waged against bourgeois factionalism'', and he equated
the evil with revisionism and capitalism. Tie's formulation was the same
as that used in the ``20 Points'', which had been completed on 2
September 1975, the day after Tie's speech.\footnote{Chi Hsin, The Case
  of the Gang of Four, p.~247. This same wording had been used in the
  ZJRB editorial of August 23 and was used again in the editorial of
  September 14 and the major article on the silk complex of September
  30. See Xuexi Wenji, pp.~64, 68, 71. Chi Hsin, The Case of the Gang of
  Four，第247页。同样措辞曾用于8月23日《浙江日报》社论，并再次用于9月14日社论和9月30日关于丝绸联合企业的重要文章。见《学习文集》，第64、68、71页。}
In a further example of the subjective and instrumental manipulation of
the language of political campaigns, Tie declared that the theoretical
study campaign {[}of the dictatorship of the proletariat{]} should serve
``as the guide in achieving profound understanding of the vital
significance of the struggle against bourgeois factionalism''.
在浙江领导层的高级成员中，站出来在省内为邓小平辩护的是铁瑛。1975年9月1日，中共浙江省委召开省、市机关干部会议，传达中央指示。谭启龙主持会议，副书记陈伟达传达了未具名中央领导人讲话的内容。\footnote{ZJRB,
  September 5, 1975, pp.~1, 4. ZJRB，1975年9月5日，第1、4页。}
铁瑛在讲话中肯定了8月工作会议所采取的政治路线和政策的正确性，但指出这些政策的执行并不均衡。他指出，必须``坚决同资产阶级派性作斗争''，并把这种祸害等同于修正主义和资本主义。铁瑛的表述与``二十条''中的用语相同，而``二十条''是在1975年9月2日完成的，正是铁瑛讲话后的第二天。\footnote{Chi
  Hsin, The Case of the Gang of Four, p.~247. This same wording had been
  used in the ZJRB editorial of August 23 and was used again in the
  editorial of September 14 and the major article on the silk complex of
  September 30. See Xuexi Wenji, pp.~64, 68, 71. Chi Hsin, The Case of
  the Gang of
  Four，第247页。同样措辞曾用于8月23日《浙江日报》社论，并再次用于9月14日社论和9月30日关于丝绸联合企业的重要文章。见《学习文集》，第64、68、71页。}
作为政治运动语言被主观化、工具化操纵的又一例证，铁瑛宣称，{[}无产阶级专政{]}理论学习运动应当``作为深刻认识反对资产阶级派性斗争重大意义的指导''。

A Zhejiang Daily editorial of 14 September urged that the momentum be
maintained. It applauded the changes which had occurred ``in places and
units where the leadership is strong and instructions are vigorously
popularized and implemented''. Units with this kind of leadership acted
in accordance with the slogan of ``placing daring above everything
else'' and ``resisted factionalist activities''. They also united former
warring groups and pushed production ahead. However, stated the
editorial, development was uneven and the pace of progress
slow.\footnote{``发扬正气顶住歪风乘胜前进'' (Foster the right spirit,
  put a stop to evil tendencies and advance victoriously), ZJRB,
  September 14, 1975, in Xuexi Wenji, pp.~67-9; ZPS, September 14, 1975,
  SWB/FE/5010/BII/9-10. 《发扬正气顶住歪风乘胜前进》（Foster the right
  spirit, put a stop to evil tendencies and advance
  victoriously），《浙江日报》，1975年9月14日，载《学习文集》，第67-9页；ZPS，1975年9月14日，SWB/FE/5010/BII/9-10。}
Nonetheless, to prove to the outside world that the situation in
Hangzhou had stabilized, on 8 September a Romanian party and government
delegation was invited to visit the city. It inspected the silk complex
and the brocade mill.\footnote{CNS, 583, pp.~6-7.
  CNS，第583期，第6-7页。}
9月14日《浙江日报》社论敦促保持这一势头。社论赞扬了``领导坚强、指示大力宣传和贯彻的地方和单位''发生的变化。具有这种领导的单位按照``敢字当头''的口号行事，并``抵制派性活动''。它们还团结过去交战的群体，推动生产前进。不过，社论说，发展并不平衡，进展速度缓慢。\footnote{``发扬正气顶住歪风乘胜前进''
  (Foster the right spirit, put a stop to evil tendencies and advance
  victoriously), ZJRB, September 14, 1975, in Xuexi Wenji, pp.~67-9;
  ZPS, September 14, 1975, SWB/FE/5010/BII/9-10.
  《发扬正气顶住歪风乘胜前进》（Foster the right spirit, put a stop to
  evil tendencies and advance
  victoriously），《浙江日报》，1975年9月14日，载《学习文集》，第67-9页；ZPS，1975年9月14日，SWB/FE/5010/BII/9-10。}
尽管如此，为了向外界证明杭州局势已经稳定，9月8日，一个罗马尼亚党政代表团被邀请访问杭州。代表团考察了丝绸联合企业和织锦厂。\footnote{CNS,
  583, pp.~6-7. CNS，第583期，第6-7页。}

Meanwhile, the party authorities kept up their sustained attack on those
responsible for the upheavals of the previous period. Students and staff
at Zhang Yongsheng's Fine Arts College held a meeting to denigrate their
disgraced party secretary. Borrowing terminology from Deng's General
Program, speakers criticized the ``extremely incorrect tendency of
certain intellectuals whose world outlook was not transformed and who
disregarded organizational discipline, opposed party decisions and who
despised workers and peasants and labor''.\footnote{ZJRB, August 23,
  1975, p.~1. ZJRB，1975年8月23日，第1页。} On 23 August representatives
from the college, as well as from Weng's silk complex and other
factories and institutions in Hangzhou, mounted the platform at a
criticism meeting which was sponsored by the CCP HMC.\footnote{ZJRB,
  August 25, 1975, p.~1. ZJRB，1975年8月25日，第1页。} A further ``mass
criticism meeting'' was held on 14 October by the CCP HMC. There were
unmistakable derogatory references to Weng as the self-styled ``leader
of Zhejiang's workers' movement'' and the ``representative of new-born
forces''.\footnote{ZJRB, October 16, 1975, pp.~1, 4; ZPS, October 16,
  1975, SWB/FE/5044/BII/1. The meetings were not confined to Hangzhou.
  See ZPS, October 27, 1975, SWB/FE/5047/BII/18 for reports of similar
  meetings in Taizhou district and Linhai county; ZJRB, October 22,
  1975, p.~1 (Wenzhou) where Chen Zuolin, Deputy-secretary of the CCP
  ZPC was introduced as holding the additional and unenviable position
  of 1st Secretary of the Wenzhou District Party Committee. In November
  1975 Chen led 6,000 cadres, students and teachers in a rubbish
  collection drive in Wenzhou, an indication of the extent of the
  collapse in government services. ZJRB, November 22, 1975, p.~1.
  ZJRB，1975年10月16日，第1、4页；ZPS，1975年10月16日，SWB/FE/5044/BII/1。会议并不限于杭州。类似会议在台州地区和临海县的报道，见
  ZPS，1975年10月27日，SWB/FE/5047/BII/18；温州见
  ZJRB，1975年10月22日，第1页，其中中共浙江省委副书记陈作霖被介绍为兼任温州地委第一书记，这是一个额外且令人不快的职位。1975年11月，陈带领6000名干部、学生和教师在温州开展垃圾清理活动，这显示出政府服务崩溃的程度。ZJRB，1975年11月22日，第1页。}
Delegates of the silk complex also attended another criticism meeting
organized by the HMRC on 9 September 1975.\footnote{ZJRB, September 11,
  1975, pp.~1, 4. ZJRB，1975年9月11日，第1、4页。}
与此同时，党的当局继续持续攻击那些对前一时期动乱负有责任的人。张永生所在的美术学院的学生和教职员工召开会议，贬斥这位失势的党委书记。发言者借用邓的《总纲》术语，批评``某些知识分子世界观没有改造，无视组织纪律，反对党的决定，轻视工农和劳动的极端错误倾向''。\footnote{ZJRB,
  August 23, 1975, p.~1. ZJRB，1975年8月23日，第1页。}
8月23日，来自该学院以及翁森鹤的丝绸联合企业和杭州其他工厂、机构的代表，在中共杭州市委主办的批判会上登台发言。\footnote{ZJRB,
  August 25, 1975, p.~1. ZJRB，1975年8月25日，第1页。}
10月14日，中共杭州市委又举行了一次``群众批判会''。会上明显含有贬义地影射翁森鹤为自封的``浙江工人运动领袖''和``新生力量代表''。\footnote{ZJRB,
  October 16, 1975, pp.~1, 4; ZPS, October 16, 1975, SWB/FE/5044/BII/1.
  The meetings were not confined to Hangzhou. See ZPS, October 27, 1975,
  SWB/FE/5047/BII/18 for reports of similar meetings in Taizhou district
  and Linhai county; ZJRB, October 22, 1975, p.~1 (Wenzhou) where Chen
  Zuolin, Deputy-secretary of the CCP ZPC was introduced as holding the
  additional and unenviable position of 1st Secretary of the Wenzhou
  District Party Committee. In November 1975 Chen led 6,000 cadres,
  students and teachers in a rubbish collection drive in Wenzhou, an
  indication of the extent of the collapse in government services. ZJRB,
  November 22, 1975, p.~1.
  ZJRB，1975年10月16日，第1、4页；ZPS，1975年10月16日，SWB/FE/5044/BII/1。会议并不限于杭州。类似会议在台州地区和临海县的报道，见
  ZPS，1975年10月27日，SWB/FE/5047/BII/18；温州见
  ZJRB，1975年10月22日，第1页，其中中共浙江省委副书记陈作霖被介绍为兼任温州地委第一书记，这是一个额外且令人不快的职位。1975年11月，陈带领6000名干部、学生和教师在温州开展垃圾清理活动，这显示出政府服务崩溃的程度。ZJRB，1975年11月22日，第1页。}
丝绸联合企业的代表还参加了1975年9月9日由杭州市革命委员会组织的另一场批判会。\footnote{ZJRB,
  September 11, 1975, pp.~1, 4. ZJRB，1975年9月11日，第1、4页。}

In his speech to the rally of 11 October welcoming back the five factory
delegation from its trip to Beijing, Tie Ying proposed extensive and
in-depth criticism of bourgeois factionalism and popularization of the
silk complex's success in rapidly restoring production.\footnote{ZJRB,
  October 12, 1975, p.~1; HZRB, October 12, 1975.
  ZJRB，1975年10月12日，第1页；HZRB，1975年10月12日。} Tie also used the
formula which had been devised by Deng Xiaoping of taking Mao's three
directives (relating to theoretical study, stability and unity and
economic development) as the guiding principle or ``key link''. Tie was
making more explicit the directive issued at the meeting of 9 September,
which had pointed out that ``The series of Chairman Mao's important
instructions are closely related to each other and
inseparable''.\footnote{CNS, 583, p.~4. CNS，第583期，第4页。} When the
formulation first appeared in print remains unclear. China News Summary
claims that it was in the People's Daily editorial of 1 October
1975.\footnote{CNS, 588 (October 22, 1975). See also Jurgen Domes, The
  Government and Politics of the PRC: A Time of Transition (Boulder,
  Co.: Westview Press, 1985), p.~282, fn. 22.
  CNS，第588期（1975年10月22日）。另见 Jurgen Domes, The Government and
  Politics of the PRC: A Time of Transition（Boulder, Co.: Westview
  Press, 1985），第282页脚注22。} The slogan had been deliberately and
cleverly collocated so as to dilute the importance that the radicals
attached to the campaign to study the theory of the dictatorship of the
proletariat and the associated theoretical justification for restricting
bourgeois rights. By contrast, Deng wished to emphasize the Chairman's
directives, released almost simultaneously with his comments on theory,
relating to unity and economic construction and to stress the danger
posed by ``bourgeois factionalism'' rather than by bourgeois rights. In
his adoption of Deng's interpretation of the three directives Tie Ying
was not alone among China's provincial leaders, but it was a contentious
issue and he had thereby revealed his hand.
铁瑛在10月11日欢迎五厂代表团从北京返回的集会上发表讲话，提出要广泛深入地批判资产阶级派性，并宣传丝绸联合企业迅速恢复生产的成功经验。\footnote{ZJRB,
  October 12, 1975, p.~1; HZRB, October 12, 1975.
  ZJRB，1975年10月12日，第1页；HZRB，1975年10月12日。}
铁瑛还使用了邓小平设计的公式，即把毛泽东的三项指示（关于理论学习、安定团结和经济发展）作为指导原则或``纲''。铁瑛是在更加明确地表述9月9日会议所发出的指示，该指示曾指出：``毛主席的一系列重要指示是互相联系、不可分割的。''\footnote{CNS,
  583, p.~4. CNS，第583期，第4页。}
这一提法何时首次见诸报端尚不清楚。China News Summary
称它出现在1975年10月1日《人民日报》社论中。\footnote{CNS, 588 (October
  22, 1975). See also Jurgen Domes, The Government and Politics of the
  PRC: A Time of Transition (Boulder, Co.: Westview Press, 1985),
  p.~282, fn. 22. CNS，第588期（1975年10月22日）。另见 Jurgen Domes, The
  Government and Politics of the PRC: A Time of Transition（Boulder,
  Co.: Westview Press, 1985），第282页脚注22。}
这一口号经过有意而巧妙的并置，用来淡化激进派赋予学习无产阶级专政理论运动以及限制资产阶级法权的相关理论论证的重要性。相比之下，邓希望强调主席几乎与其理论评论同时发布的有关团结和经济建设的指示，并突出``资产阶级派性''而非资产阶级法权所构成的危险。铁瑛采纳邓对三项指示的解释，并非中国省级领导人中的孤例，但这是一个有争议的问题，他也因此暴露了自己的立场。

In mid-October the CCP ZPC held an industrial conference at which the
delegates heard reports on the transformation of the Hangzhou silk
complex and were urged to lead their units to catch up with the gearbox
factory.\footnote{ZJRB, October 11, 1975, p.~1.
  ZJRB，1975年10月11日，第1页。} The conference pointed out that control
of a factory's leadership by ``genuine Marxists'' and the ``working
masses'' would ensure the correctness of its ideological and political
line. To overcome the influence of the ``two crashthroughs'', some
enterprises were sponsoring lessons on the CCP constitution and training
``backbone elements'' of Party branches and work teams in the
rudimentaries of leadership. The conference also exhorted party
committees to protect their subordinate officials and to shoulder
responsibility for mistakes on their behalf. It also demanded the
strengthening of ``business management'', the establishment and
improvement of rules and regulations and devotion to the urgent
production tasks remaining in 1975.
10月中旬，中共浙江省委召开工业会议，与会代表听取了杭州丝绸联合企业转变情况的报告，并被要求带领本单位赶上齿轮箱厂。\footnote{ZJRB,
  October 11, 1975, p.~1. ZJRB，1975年10月11日，第1页。}
会议指出，由``真正的马克思主义者''和``劳动群众''掌握工厂领导权，将确保其思想政治路线的正确性。为克服``双突''的影响，一些企业举办中共党章课程，并训练党支部和工作组的``骨干分子''掌握领导工作的基本知识。会议还敦促各级党委保护下属干部，并代他们承担错误责任。会议还要求加强``企业管理''，建立和完善规章制度，并全力完成1975年剩余的紧迫生产任务。

A progress report on Hangzhou's industrial production at the beginning
of November continued to endorse the example set by the silk complex and
the gearbox works.\footnote{ZPS, November 3, 1975, SWB/FE/5058/BII/1-2.
  ZPS，1975年11月3日，SWB/FE/5058/BII/1-2。} Party committees in the
city's factories were devoting their energies to training shop-floor
leaders, bringing into play the ``exemplary vanguard role of Party
members, the role of veteran workers as backbone elements, and the role
of Communist Youth League members and young workers as a shock force''.
Heavy machinery factories had sent officials and workers to the gearbox
plant on inspection tours while the silk mills of Hangzhou had emulated
the silk complex in restoring and raising production levels. After the
disruptions of the first half of the year, the urgency of the race to
meet annual targets was indicated by the passage in the report which
praised employees and leaders at factories who
11月初关于杭州工业生产进展的报告继续肯定丝绸联合企业和齿轮箱厂树立的榜样。\footnote{ZPS,
  November 3, 1975, SWB/FE/5058/BII/1-2.
  ZPS，1975年11月3日，SWB/FE/5058/BII/1-2。}
杭州市各工厂党委正把精力投入培训车间一线领导，发挥``党员的先锋模范作用、老工人的骨干作用以及共青团员和青年工人的突击作用''。重型机械工厂派干部和工人到齿轮箱厂参观学习，杭州各丝绸厂则仿效丝绸联合企业，恢复并提高生产水平。经历上半年动荡之后，报告中称赞工厂职工和领导的段落显示出完成全年指标竞赛的紧迫性，他们

\begin{quote}
persist in doing all they can within the eight hours and still make more
contributions after the eight hours. They seek no personal fame or gain,
fear no fatigue of getting dirty {[}sic{]}, are unmindful about hours of
pay, actively take part in voluntary labour, and strive to fulfil
production tasks.
坚持八小时内尽力多干，八小时外继续多作贡献。他们不图个人名利，不怕脏累，不计较工时和报酬，积极参加义务劳动，努力完成生产任务。
\end{quote}

A mass meeting of 40,000 party members, which was called by the CCP HMC
on 22 November, reviewed recent progress and the arduous tasks that lay
ahead.\footnote{ZJRB, November 23, 1975, p.~1.
  ZJRB，1975年11月23日，第1页。} In his report to the assembled CCP
members, 1st Secretary Zhang Zishi hailed the achievements of the
previous three months and announced that industrial production in the
city in the third quarter of 1975 had increased by 21.4\% over the
preceding quarter. Zhang complimented the city's party organizations and
citizens for their exposure and criticism of ``the crimes committed by
individual bad elements and bourgeois factionalism''. Representatives of
party members at both Weng's silk complex and Xia Genfa's oxygen
generator plant addressed the gathering.
11月22日，中共杭州市委召开有4万名党员参加的群众大会，回顾近期进展和摆在面前的艰巨任务。\footnote{ZJRB,
  November 23, 1975, p.~1. ZJRB，1975年11月23日，第1页。}
第一书记张子石在向全体党员所作的报告中赞扬过去三个月的成就，并宣布1975年第三季度全市工业生产比上一季度增长21.4\%。张子石称赞全市党组织和市民揭发、批判``个别坏分子和资产阶级派性所犯的罪行''。翁森鹤的丝绸联合企业和夏根发的制氧机厂的党员代表都在会上发言。

However, by the end of the year a pronounced shift in provincial
propaganda had begun. The waning of Deng Xiaoping's star late in 1975
was immediately reflected in Zhejiang. For example, at the beginning of
December 1975 Wenzhou announced a production increase of 10.97\% for
October over September and a staggering 48\% increase for the first
twenty days of November over October. The authorities attributed the
production slump earlier in the year to the ``interference of bourgeois
factionalism'' which had enabled ``a handful of class enemies {[}to
engage{]} in sabotage for a long time''. It was revealed in the stinging
rebuke directed toward it that the former leadership had badly handled
the crisis and even abdicated its responsibility: ``cowardice and
flinching means giving up leadership and runs counter to the Party's
basic line''.\footnote{ZJRB, December 1, 1975, pp.~1, 4; ZPS, December
  1, 1975, SWB/FE/5077/BII/10-11.
  ZJRB，1975年12月1日，第1、4页；ZPS，1975年12月1日，SWB/FE/5077/BII/10-11。}
Yet a mere two days later the commentary published in the Zhejiang Daily
concerning Wenzhou's economy signalled an abrupt change in direction by
stating simplistically, and in terminology not used for over a year,
that ``once class struggle is grasped, all problems can be
solved''.\footnote{ZPS, December 3, 1975, SWB/FE/5078/BII/9-10.
  ZPS，1975年12月3日，SWB/FE/5078/BII/9-10。} Bourgeois factionalism and
economic problems, issues which had featured in the report of 1
December, had abruptly disappeared from the media.
然而，到年底，省级宣传已经开始出现明显转向。1975年末邓小平声望下降，立即在浙江反映出来。例如，1975年12月初，温州宣布10月产量比9月增长10.97\%，11月前20天比10月惊人地增长48\%。当局把当年早些时候的生产下滑归因于``资产阶级派性的干扰''，这种干扰使``一小撮阶级敌人{[}得以{]}长期进行破坏''。针对前任领导的严厉斥责透露出，他们对危机处理很糟，甚至放弃了责任：``胆怯退缩就是放弃领导，违背党的基本路线。''\footnote{ZJRB,
  December 1, 1975, pp.~1, 4; ZPS, December 1, 1975,
  SWB/FE/5077/BII/10-11.
  ZJRB，1975年12月1日，第1、4页；ZPS，1975年12月1日，SWB/FE/5077/BII/10-11。}
然而仅仅两天后，《浙江日报》关于温州经济的评论就显示出方向的突然变化，它用一年多未曾使用的术语简单地说：``阶级斗争一抓，一切问题都解决了。''\footnote{ZPS,
  December 3, 1975, SWB/FE/5078/BII/9-10.
  ZPS，1975年12月3日，SWB/FE/5078/BII/9-10。}
资产阶级派性和经济问题这些曾出现在12月1日报告中的议题，突然从媒体上消失了。

After his appearance in Hangzhou on 13 November 1975 at the provincial
militia work meeting\footnote{CIA Reference Aid, Appearances and
  Activities of Leading Personalities of the PRC, 1975. CIA Reference
  Aid, Appearances and Activities of Leading Personalities of the PRC,
  1975。} Tie Ying again dropped from public visibility until his
reemergence at the provincial memorial service for Mao Zedong, ten
months later. At a meeting on Christmas Day 1975, at which the three
most senior members of the ZPC secretariat, Tan Qilong, Tie and Lai
Keke, were all conspicuous by their absence, Luo Yi, a protege of Zhang
Chunqiao, delivered the principal speech.\footnote{ZJRB, December 26,
  1975, pp.~1, 4. ZJRB，1975年12月26日，第1、4页。} Luo summed up 1975
as the year in which the provincial party had taken ``class struggle as
the key link'' and proposed a continuation of this political line for
the coming year. Luo's words were a patent revision of the assessment
which had been made only the previous month.
铁瑛在1975年11月13日出席杭州省民兵工作会议后\footnote{CIA Reference Aid,
  Appearances and Activities of Leading Personalities of the PRC, 1975.
  CIA Reference Aid, Appearances and Activities of Leading Personalities
  of the PRC, 1975。}，再次从公众视野中消失，直到十个月后在毛泽东省级追悼会上重新露面。1975年圣诞节的一次会议上，浙江省委书记处三位最资深成员谭启龙、铁瑛和赖可可都显眼地缺席，由张春桥的门生罗毅发表主要讲话。\footnote{ZJRB,
  December 26, 1975, pp.~1, 4. ZJRB，1975年12月26日，第1、4页。}
罗毅把1975年概括为省委以``阶级斗争为纲''的一年，并提出下一年继续这一政治路线。罗毅的话显然修正了仅在前一个月作出的评估。

The victory-reporting rally convened by the municipal authorities on 5
January 1976 confirmed this shift, albeit in somewhat less forthright
terms.\footnote{ZJRB, January 6, 1976, p.~1. ZJRB，1976年1月6日，第1页。}
Over 300 units in the municipality reported their production figures for
1975, and these figures showed that industrial production had risen by
30\% in the final quarter of the year.\footnote{In 1975, industrial
  output for the province as a whole rose by 3.9\%; 4\% in the light
  industry sector and 3.6\% in heavy industry. Zhejiang shengqing,
  p.~1058.
  1975年，全省工业总产值增长3.9\%；轻工业增长4\%，重工业增长3.6\%。《浙江省情》，第1058页。}
Production successes in 1975 were now attributed to the policy of taking
class struggle as the key link and the criticism of revisionism and
capitalism and were viewed as contributing to the defence and
development of ``the fruits of the Cultural Revolution''. ``The
strengthening of leadership by the Party'' was offset by the
accompanying qualification that it ``must be unified with the free
launching of the masses''.\footnote{CNS, 597 (January 7, 1976), p.~5.
  CNS，第597期（1976年1月7日），第5页。} For Hangzhou and the whole of
Zhejiang, the orientation for 1976 henceforward revolved around class
struggle.
市政当局于1976年1月5日召开的报捷大会确认了这一转向，尽管措辞稍显含蓄。\footnote{ZJRB,
  January 6, 1976, p.~1. ZJRB，1976年1月6日，第1页。}
全市300多个单位报告了1975年的生产数字，这些数字显示，年末最后一个季度工业生产增长了30\%。\footnote{In
  1975, industrial output for the province as a whole rose by 3.9\%; 4\%
  in the light industry sector and 3.6\% in heavy industry. Zhejiang
  shengqing, p.~1058.
  1975年，全省工业总产值增长3.9\%；轻工业增长4\%，重工业增长3.6\%。《浙江省情》，第1058页。}
1975年的生产成就现在被归因于以阶级斗争为纲和批判修正主义、资本主义的政策，并被视为有助于捍卫和发展``文化大革命的成果''。``加强党的领导''被附加的限定语所抵消，即它``必须同放手发动群众统一起来''。\footnote{CNS,
  597 (January 7, 1976), p.~5. CNS，第597期（1976年1月7日），第5页。}
从此以后，杭州和整个浙江1976年的方向都围绕阶级斗争展开。

A factor which mitigated the severity of the 1976 campaign to reverse
Deng's policies of mid to late 1975 emerged as a consequence of the
National Conference to Learn from Dazhai in agriculture, held from 15
September until 19 October 1975 in Beijing and Shanxi province. The
major speakers at the lengthy conference were Deng Xiaoping, Jiang Qing
and Hua Guofeng and their speeches accurately reflected their respective
views of socialism. Only Hua's report was published in China at the
time.\footnote{Hua Guofeng, Let the Whole Party Mobilize for a vast
  effort to develop Agriculture and build Tachai-type counties
  throughout the Country, October 15, 1975 (Beijing: Foreign Languages
  Press, 1975). 华国锋，Let the Whole Party Mobilize for a vast effort
  to develop Agriculture and build Tachai-type counties throughout the
  Country，1975年10月15日（北京：外文出版社，1975）。} In it Hua
advocated agricultural mechanization and strengthening of the powers and
economic importance of brigades and communes at the expense of
production teams. Taiwan sources reported that Jiang Qing, in her speech
at the conference, had appealed for the class mobilization of the
peasants.\footnote{China News Agency (Taibei), December 12, 1975,
  SWB/FE/5088/BII/16. For Jiang's speech of September 15 see Feiqing
  Yuebao, 18: 7 (1976), pp.~93-4.
  中国新闻社（台北），1975年12月12日，SWB/FE/5088/BII/16。江青9月15日讲话见《匪情月报》18:7（1976），第93-4页。}
Hua and the Cultural Revolution radicals thus agreed on the goals but
not on the means of rural social policies.\footnote{Domes, ``The `Gang
  of Four' - and Hua Guofeng'', p.~482. Domes, ``The `Gang of Four' -
  and Hua Guofeng''，第482页。}
1976年反转邓小平1975年中后期政策的运动之所以严重程度有所缓和，是因为1975年9月15日至10月19日在北京和山西举行的全国农业学大寨会议产生了一个因素。这次漫长会议的主要发言人是邓小平、江青和华国锋，他们的讲话准确反映了各自关于社会主义的观点。当时中国只发表了华国锋的报告。\footnote{Hua
  Guofeng, Let the Whole Party Mobilize for a vast effort to develop
  Agriculture and build Tachai-type counties throughout the Country,
  October 15, 1975 (Beijing: Foreign Languages Press, 1975). 华国锋，Let
  the Whole Party Mobilize for a vast effort to develop Agriculture and
  build Tachai-type counties throughout the
  Country，1975年10月15日（北京：外文出版社，1975）。}
在报告中，华主张农业机械化，并以生产队为代价加强大队和公社的权力及经济重要性。台湾方面资料报道，江青在会议讲话中呼吁对农民进行阶级动员。\footnote{China
  News Agency (Taibei), December 12, 1975, SWB/FE/5088/BII/16. For
  Jiang's speech of September 15 see Feiqing Yuebao, 18: 7 (1976),
  pp.~93-4.
  中国新闻社（台北），1975年12月12日，SWB/FE/5088/BII/16。江青9月15日讲话见《匪情月报》18:7（1976），第93-4页。}
因而，华国锋和文化大革命激进派在农村社会政策的目标上意见一致，但在手段上不同。\footnote{Domes,
  ``The `Gang of Four' - and Hua Guofeng'', p.~482. Domes, ``The `Gang
  of Four' - and Hua Guofeng''，第482页。}

Deng Xiaoping addressed the conference on the opening day and again on
27 September and 4 October.\footnote{Zhonggong dangshi baiti, p.~388;
  Deng Xiaoping, ``各条战线都要整顿'' (Various fronts require
  rectification), Deng Xiaoping Wenxuan, pp.~32-4. The speech, like all
  those from 1975 included in the selection, was undoubtedly edited
  before publication so as to emphasize the degree to which Deng's
  position at the time differed from that of his opponents.
  《中共党史百题》，第388页；邓小平《各条战线都要整顿》（Various fronts
  require
  rectification），《邓小平文选》，第32-4页。该讲话同选集中收录的所有1975年讲话一样，无疑在发表前经过编辑，以强调邓当时立场同其对手的差异程度。}
In the speeches Deng pursued his theme of rectification or consolidation
(整顿). Deng stated that Mao had talked about rectification in the PLA,
the localities, industry, agriculture, commerce, culture and education,
science and technology and the arts and he stressed that the key to its
success lay in the party, especially in its leading groups at the county
level. Leading cadres had the duty to
邓小平在会议开幕当天以及9月27日、10月4日发言。\footnote{Zhonggong
  dangshi baiti, p.~388; Deng Xiaoping, ``各条战线都要整顿'' (Various
  fronts require rectification), Deng Xiaoping Wenxuan, pp.~32-4. The
  speech, like all those from 1975 included in the selection, was
  undoubtedly edited before publication so as to emphasize the degree to
  which Deng's position at the time differed from that of his opponents.
  《中共党史百题》，第388页；邓小平《各条战线都要整顿》（Various fronts
  require
  rectification），《邓小平文选》，第32-4页。该讲话同选集中收录的所有1975年讲话一样，无疑在发表前经过编辑，以强调邓当时立场同其对手的差异程度。}
在这些讲话中，邓继续阐述整顿或巩固（整顿）的主题。邓说，毛曾谈到军队、地方、工业、农业、商业、文化教育、科学技术和文艺都要整顿，并强调整顿成功的关键在党，尤其在县一级领导班子。领导干部有责任

\begin{quote}
discover a good young successor (好苗子) and promote him step by step,
quickening the process somewhat, so that he stays on one grade for only
a year or so.\footnote{Deng Xiaoping Wenxuan, p.~33.
  《邓小平文选》，第33页。}
发现一个好的年轻接班人（好苗子），一步一步提拔他，过程可以稍快一些，使他在一个级别上只停留一年左右。\footnote{Deng
  Xiaoping Wenxuan, p.~33. 《邓小平文选》，第33页。}
\end{quote}

In a speech which he had delivered in August 1975, Deng had advocated
the selection of cadres with practical experience, aged in their forties
and fifties or even younger, to take up responsible
positions.\footnote{``关于国防工业企业的整顿'' (Concerning the
  rectification of national defence industrial enterprises), Deng
  Xiaoping Wenxuan, pp.~25-6. 《关于国防工业企业的整顿》（Concerning the
  rectification of national defence industrial
  enterprises），《邓小平文选》，第25-6页。} In Zhejiang two followers
of the central radicals, Lai Keke and Luo Yi, in a document which they
prepared in 1976, denounced the application of Deng's ``theory of
steps'' (台阶论) to the province.\footnote{Shen Weicai, HZRB, August 17,
  1977. 沈维才，HZRB，1977年8月17日。} In this, they were probably
taking their cue from a critique published in the People's Daily in
March 1976.\footnote{RMRB, March 20, 1976, p.~2.
  RMRB，1976年3月20日，第2页。} Deng's advocacy of rapid promotion for
people of talent sounded suspiciously similar to the ``helicopter''
promotion of Cultural Revolution rebels, much decried by veteran cadres
such as Deng himself. Several years later, Deng was at pains to
distinguish between the two approaches to promotion, but his argument
was not persuasive.\footnote{See Deng Xiaoping,
  ``党和国家领导制度的改革'' (Reform of the leadership system in the
  party and the state), August 18, 1980, Deng Xiaoping Wenxuan,
  pp.~283-84. 见邓小平《党和国家领导制度的改革》（Reform of the
  leadership system in the party and the
  state），1980年8月18日，《邓小平文选》，第283-84页。} In fact, Deng's
methods were essentially the same; only his criteria for selection
differed.
邓在1975年8月的一次讲话中曾主张选拔有实际经验、四五十岁甚至更年轻的干部担任负责职务。\footnote{``关于国防工业企业的整顿''
  (Concerning the rectification of national defence industrial
  enterprises), Deng Xiaoping Wenxuan, pp.~25-6.
  《关于国防工业企业的整顿》（Concerning the rectification of national
  defence industrial enterprises），《邓小平文选》，第25-6页。}
在浙江，两个中央激进派追随者赖可可和罗毅，在他们1976年准备的一份文件中谴责把邓的``台阶论''应用于该省。\footnote{Shen
  Weicai, HZRB, August 17, 1977. 沈维才，HZRB，1977年8月17日。}
他们很可能是在仿效1976年3月《人民日报》发表的一篇批判文章。\footnote{RMRB,
  March 20, 1976, p.~2. RMRB，1976年3月20日，第2页。}
邓主张迅速提拔有才干者，听起来令人怀疑地类似于文化大革命造反派的``直升飞机式''提拔，而这种做法曾受到邓本人这样的老干部强烈谴责。几年后，邓费力地区分这两种提拔方式，但他的论证并不令人信服。\footnote{See
  Deng Xiaoping, ``党和国家领导制度的改革'' (Reform of the leadership
  system in the party and the state), August 18, 1980, Deng Xiaoping
  Wenxuan, pp.~283-84. 见邓小平《党和国家领导制度的改革》（Reform of the
  leadership system in the party and the
  state），1980年8月18日，《邓小平文选》，第283-84页。}
实际上，邓的方法本质上相同；只是他的选拔标准不同。

On 29 October 1975 the CCP ZPC convened a meeting to announce the
decisions which had been taken at the national Dazhai conference. Luo
Yi, who had led the Zhejiang delegation to Beijing and Dazhai, spoke of
his impressions. Then Tie Ying, who seemed to have become acting head of
the CCP ZPC, issued instructions on follow-up action.\footnote{ZJRB,
  November 1, 1975, p.~1. Wenzhou district and municipality held meeting
  on the conference in November, at which Chen Zuolin spoke. ZJRB,
  November 11, 1975, p.~1. For other provincial responses, see CNS, 590
  (November 5, 1975) and 593 (December 3, 1975).
  ZJRB，1975年11月1日，第1页。温州地区和市区于11月召开会议传达该会议精神，陈作霖在会上讲话。ZJRB，1975年11月11日，第1页。其他省份反应见
  CNS，第590期（1975年11月5日）和第593期（1975年12月3日）。} Tie
continued to hammer home the evils of factionalism -- a topic absent
from Luo's speech -- and of capitalism, which he defined as the
commission of corrupt and criminal offences. While Luo praised the
open-door style of party rectification carried out in the county where
Dazhai was located, Tie seemed to favour an in-house procedure conducted
behind closed doors and restricted to members of the CCP. The two
different approaches to party rectification were another dividing line
between radical and more orthodox policies toward organizational
issues.\footnote{See Harding, Organizing China, p.~322. 见 Harding,
  Organizing China，第322页。}
1975年10月29日，中共浙江省委召开会议，宣布全国大寨会议作出的决定。率领浙江代表团赴北京和大寨的罗毅谈了自己的感受。随后，似乎已成为中共浙江省委代理负责人的铁瑛，就后续行动作出指示。\footnote{ZJRB,
  November 1, 1975, p.~1. Wenzhou district and municipality held meeting
  on the conference in November, at which Chen Zuolin spoke. ZJRB,
  November 11, 1975, p.~1. For other provincial responses, see CNS, 590
  (November 5, 1975) and 593 (December 3, 1975).
  ZJRB，1975年11月1日，第1页。温州地区和市区于11月召开会议传达该会议精神，陈作霖在会上讲话。ZJRB，1975年11月11日，第1页。其他省份反应见
  CNS，第590期（1975年11月5日）和第593期（1975年12月3日）。}
铁瑛继续反复强调派性这一罪恶------罗毅讲话中没有这个主题------以及资本主义的罪恶，他把资本主义界定为贪污腐化和刑事犯罪行为。罗毅赞扬大寨所在县采取的开门整党方式，而铁瑛似乎倾向于一种在党内、关门进行且限于中共党员的程序。这两种不同的整党方式，是激进政策与较正统政策在组织问题上的又一分界线。\footnote{See
  Harding, Organizing China, p.~322. 见 Harding, Organizing
  China，第322页。}

The major decision relayed at the provincial meeting concerned the
despatch of work teams to the countryside to propagate the spirit of the
Dazhai conference. One hundred thousand people attended meetings in
Hangzhou on 12 November 1975 to send off the first group totalling 9,000
team members.\footnote{ZJRB, November 13, 1975, pp.~1, 2.
  ZJRB，1975年11月13日，第1、2页。} In addition to office cadres, the
teams contained workers, teachers and students. Tie Ying, in his speech
at the rally, announced that the total number of cadres and other staff
in Zhejiang who had been allocated for the assignment was 30,000 and
included several standing committee members of the ZPC. He ordered work
team members to remain at their posts in the countryside and forbade
their work units from recalling them to urban areas.
省级会议传达的主要决定，是派遣工作队下乡宣传大寨会议精神。1975年11月12日，杭州有10万人参加会议，为第一批总计9000名工作队员送行。\footnote{ZJRB,
  November 13, 1975, pp.~1, 2. ZJRB，1975年11月13日，第1、2页。}
除机关干部外，工作队还包括工人、教师和学生。铁瑛在集会讲话中宣布，浙江被分配参加这项任务的干部和其他人员总数为3万人，其中包括若干浙江省委常委。他命令工作队员留在农村岗位上，并禁止其工作单位把他们召回城市。

It was later reported that 40,000 cadres went to the countryside in late
1975.\footnote{Chen Zuolin, HZRB, December 18, 1976. In neighboring
  Jiangsu province 50,000 cadres were sent to the countryside. See
  Jiangsusheng dashiji, p.~338.
  陈作霖，HZRB，1976年12月18日。邻近的江苏省有5万名干部被派往农村。见《江苏省大事记》，第338页。}
In all, it was claimed that 80\% of office cadres in Hangzhou had left
their desks, leaving behind a only skeleton staff to man the telephones.
It is possible that younger ``rebel'' cadres comprised the bulk of the
work team members, effectively removing them from the locus of power in
the province.\footnote{Zweig, ``The Peita Debate on Education,'' p.~151,
  fn. 21. Zweig, ``The Peita Debate on Education,'' 第151页脚注21。} The
prolonged absence from Hangzhou of younger cadres sympathetic to the
rebel cause perhaps served to weaken the impact of the 1976 leftist
upsurge in Zhejiang and deprived the rebel leaders Zhang Yongsheng and
He Xianchun of potential allies for their last battle with the
``capitalist-roaders'' of the province.\footnote{Factional fighting had
  not ceased in the counties. On November 5, 1975, an armed conflict in
  Linhai county resulted in death and injury. See Linhai xianzhi, p.~31.
  县里的派性斗争并未停止。1975年11月5日，临海县一次武装冲突造成死伤。见《临海县志》，第31页。}
后来有报道称，1975年末有4万名干部下乡。\footnote{Chen Zuolin, HZRB,
  December 18, 1976. In neighboring Jiangsu province 50,000 cadres were
  sent to the countryside. See Jiangsusheng dashiji, p.~338.
  陈作霖，HZRB，1976年12月18日。邻近的江苏省有5万名干部被派往农村。见《江苏省大事记》，第338页。}
总体而言，据称杭州80\%的机关干部离开办公桌，只留下少量留守人员接电话。年轻的``造反派''干部可能构成了工作队员的大多数，从而实际上被移出省内权力中心。\footnote{Zweig,
  ``The Peita Debate on Education,'' p.~151, fn. 21. Zweig, ``The Peita
  Debate on Education,'' 第151页脚注21。}
同情造反事业的年轻干部长期离开杭州，或许削弱了1976年浙江左派高潮的影响，也使造反派领导人张永生和贺贤春在同本省``走资派''的最后一战中失去了潜在盟友。\footnote{Factional
  fighting had not ceased in the counties. On November 5, 1975, an armed
  conflict in Linhai county resulted in death and injury. See Linhai
  xianzhi, p.~31.
  县里的派性斗争并未停止。1975年11月5日，临海县一次武装冲突造成死伤。见《临海县志》，第31页。}

\section{The Leftist Attack on the July 1975 Decisions, January -
September
1976}\label{the-leftist-attack-on-the-july-1975-decisions-january---september-1976}

\section{左派对1975年7月决定的攻击，1976年1月至9月}\label{ux5de6ux6d3eux5bf91975ux5e747ux6708ux51b3ux5b9aux7684ux653bux51fb1976ux5e741ux6708ux81f39ux6708}

The campaign to ``criticize the right deviationist wind to reverse
correct verdicts'', which spread to the campuses of major colleges and
schools across the country in January 1976, operationalized the renewal
of the politics of mobilization.\footnote{For details of the ``great
  debate in education'' see, Fenwick, ``The Gang of Four'', pp.~366-82.
  For reports concerning the campaign at Hangzhou University and the
  Zhejiang Fine Arts College, see ZJRB, January 11, 1976, p.~1, January
  12, 1976, p.~2, January 18, 1976, p.~1, and January 19, 1976, p.~1.
  For a similar report from Hangzhou's Xuejun middle school, see ZPS,
  February 13, 1976, SWB/FE/5138/BII/12. 关于``教育大辩论''的详情，见
  Fenwick, ``The Gang of
  Four''，第366-82页。关于杭州大学和浙江美术学院运动的报道，见
  ZJRB，1976年1月11日第1页、1月12日第2页、1月18日第1页、1月19日第1页。杭州学军中学的类似报道见
  ZPS，1976年2月13日，SWB/FE/5138/BII/12。} The campaign began as a
result of letters written to Mao by the party leaders of Qinghua
University in August and October 1975. The letters reportedly took Chi
Qun and Xie Jingyi to task for their failings both ideologically and
morally. Mao, with the encouragement of his nephew and court mediator
Mao Yuanxin, took the letters as evidence of an attempt to reverse the
verdict on the Cultural Revolution in the sphere of education. The
younger Mao said to his uncle:
1976年1月，在全国主要高校和学校校园中展开的``批判右倾翻案风''运动，使动员政治的恢复具体运转起来。\footnote{For
  details of the ``great debate in education'' see, Fenwick, ``The Gang
  of Four'', pp.~366-82. For reports concerning the campaign at Hangzhou
  University and the Zhejiang Fine Arts College, see ZJRB, January 11,
  1976, p.~1, January 12, 1976, p.~2, January 18, 1976, p.~1, and
  January 19, 1976, p.~1. For a similar report from Hangzhou's Xuejun
  middle school, see ZPS, February 13, 1976, SWB/FE/5138/BII/12.
  关于``教育大辩论''的详情，见 Fenwick, ``The Gang of
  Four''，第366-82页。关于杭州大学和浙江美术学院运动的报道，见
  ZJRB，1976年1月11日第1页、1月12日第2页、1月18日第1页、1月19日第1页。杭州学军中学的类似报道见
  ZPS，1976年2月13日，SWB/FE/5138/BII/12。}
这场运动起因于清华大学党委领导人在1975年8月和10月写给毛泽东的信。据说这些信严厉批评迟群和谢静宜在思想和道德上的缺失。在侄子兼宫廷调解人毛远新的鼓励下，毛把这些信视为有人试图在教育领域为文化大革命翻案的证据。年轻的毛远新对他的伯父说：

\begin{quote}
``I pay great attention to {[}Deng{]} Xiaoping's talks and I feel there
is a problem. He very rarely refers to the achievements of the Cultural
Revolution or criticizes Liu Shaoqi's revisionist line.'' ``As for the
three directives as the key link -- really there is only one directive,
pushing production forward''. ``How should the Cultural Revolution be
seen? \ldots{} Should it be affirmed or negated?''\footnote{Gao Gao and
  Yan Jiaqi, ``Wenhua dageming'', pp.~565-66. The authors state that Mao
  Yuanxin took up the job of the Chairman's liaison person in September
  1975.
  高皋、严家其，《``文化大革命''》，第565-66页。作者称毛远新于1975年9月承担起主席联络人的工作。}
``我很注意{[}邓{]}小平的谈话，觉得有问题。他很少谈文化大革命的成绩，也很少批判刘少奇的修正主义路线。''\,``至于三项指示为纲------实际上只有一项指示，就是把生产搞上去。''\,``文化大革命怎么看？\ldots\ldots 是肯定还是否定？''\footnote{Gao
  Gao and Yan Jiaqi, ``Wenhua dageming'', pp.~565-66. The authors state
  that Mao Yuanxin took up the job of the Chairman's liaison person in
  September 1975.
  高皋、严家其，《``文化大革命''》，第565-66页。作者称毛远新于1975年9月承担起主席联络人的工作。}
\end{quote}

On 3 November Mao's directives, which in effect affirmed his nephew's
fears, were relayed back to Qinghua university. In late November, on the
Chairman's instructions, the Politburo convened a notification (打招呼)
meeting at which Hua Guofeng delivered the main speech which had been
vetted by the feeble Mao. On 26 November the speech was relayed to a
meeting of provincial first secretaries. From then until Deng's
disappearance from public life in January 1976, the Politburo criticized
Deng at a series of meetings. On 5 February 1976, a central document
notified the party and citizens of Mao's opinions.\footnote{Hao Mengbi
  and Duan Haoran, Liushinian, pp.~648-50; Dangshi tongxun, 1 (1983),
  p.~11. ``Wenhua dageming'', pp.~565-66.
  郝梦笔、段浩然，《六十年》，第648-50页；《党史通讯》1（1983），第11页；《``文化大革命''》，第565-66页。}
11月3日，实际上确认了毛远新担忧的毛泽东指示被传回清华大学。11月下旬，根据主席指示，政治局召开了一次打招呼会议，华国锋作了主要讲话，该讲话经过身体衰弱的毛泽东审阅。11月26日，这一讲话被传达到省级第一书记会议。从那时起直到邓于1976年1月从公众生活中消失，政治局在一系列会议上批判邓。1976年2月5日，一份中央文件向党内和群众传达了毛的意见。\footnote{Hao
  Mengbi and Duan Haoran, Liushinian, pp.~648-50; Dangshi tongxun, 1
  (1983), p.~11. ``Wenhua dageming'', pp.~565-66.
  郝梦笔、段浩然，《六十年》，第648-50页；《党史通讯》1（1983），第11页；《``文化大革命''》，第565-66页。}

In that same month the CC, at Mao's suggestion, held a meeting in
Beijing to set down the guidelines and parameters for the campaign to be
launched against Deng.\footnote{For a detailed discussion of the aims,
  agenda, dynamics and impact of the conference, see Fenwick, ``The Gang
  of Four'', pp.~415-43. 关于会议目标、议程、动态和影响的详细讨论，见
  Fenwick, ``The Gang of Four''，第415-43页。} Hua Guofeng had been
appointed acting Premier and Deng's replacement in supervising the daily
work of the CCP CC on 21 and 28 January 1976 respectively.\footnote{Zhonggong
  dangshi 170 ti, p.~319. 《中共党史170题》，第319页。} In another
important decision at this time, the Commander of the Beijing Military
Region, Chen Xilian, was appointed to take charge of the daily work of
the CC MAC. Chen replaced Ye Jianying in this post due, it is claimed,
to Ye's ill-health.\footnote{Hao Mengbi and Duan Haoran, Liushinian,
  p.~649. See Fenwick, ``The Gang of Four'', pp.~407-08. By October Ye
  was well enough to organize the arrest of the Gang of Four.
  郝梦笔、段浩然，《六十年》，第649页。见 Fenwick, ``The Gang of
  Four''，第407-08页。到10月，叶的健康已足以组织逮捕``四人帮''。} On 25
February, Hua Guofeng delivered a speech to the meeting in which he
accused Deng of having dismissed young cadres and replaced them with
unrepentant capitalist-roaders. In a direct reference to the events in
Hangzhou of the previous year, Hua described Deng's exertions from July
to September as the repression of the revolutionary masses and of
different points of view. Ironically, Deng was shouldered with the
responsibility for forming factions and provoking conflicts such as the
``Hangzhou Incident''.\footnote{I\&S, 14: 6 (1978), pp.~94-7.
  I\&S，14:6（1978），第94-7页。}
同月，中央委员会根据毛的建议在北京召开会议，为即将发动的反邓运动确定方针和范围。\footnote{For
  a detailed discussion of the aims, agenda, dynamics and impact of the
  conference, see Fenwick, ``The Gang of Four'', pp.~415-43.
  关于会议目标、议程、动态和影响的详细讨论，见 Fenwick, ``The Gang of
  Four''，第415-43页。}
华国锋分别于1976年1月21日和28日被任命为代总理，并接替邓监督中共中央日常工作。\footnote{Zhonggong
  dangshi 170 ti, p.~319. 《中共党史170题》，第319页。}
此时另一项重要决定是，北京军区司令员陈锡联被任命负责中央军委日常工作。据称，陈因叶剑英健康不佳而接替其职务。\footnote{Hao
  Mengbi and Duan Haoran, Liushinian, p.~649. See Fenwick, ``The Gang of
  Four'', pp.~407-08. By October Ye was well enough to organize the
  arrest of the Gang of Four. 郝梦笔、段浩然，《六十年》，第649页。见
  Fenwick, ``The Gang of
  Four''，第407-08页。到10月，叶的健康已足以组织逮捕``四人帮''。}
2月25日，华国锋在会议上发表讲话，指责邓撤换年轻干部，代之以顽固不化的走资派。华直接提及上一年杭州的事件，把邓7月至9月的努力描述为压制革命群众和不同意见。具有讽刺意味的是，邓被扣上了搞派性、挑动``杭州事件''等冲突的责任。\footnote{I\&S,
  14: 6 (1978), pp.~94-7. I\&S，14:6（1978），第94-7页。}

Jiang Qing also spoke at the conference, which seems to have run from
about 20 February until early March 1976.\footnote{For details of
  Jiang's February 23 and March 2 talks, see ``Documents of the CC of
  the CCP'', I\&S, 15:2 pp.~96-9. Fenwick, ``The Gang of Four'',
  pp.~427-30, has analyzed Jiang's speeches.
  江青2月23日和3月2日谈话的详情，见 ``Documents of the CC of the CCP'',
  I\&S, 15:2，第96-9页。Fenwick, ``The Gang of
  Four''，第427-30页，对江的讲话作了分析。} On 2 March she convened a
meeting of representatives from twelve provinces in which she accused
Deng of being a ``big traitor'' and the ``representative of
international capitalists''. On 16 February the CCP CC had approved and
circulated a CC MAC report which described Deng and Ye's speeches to the
June-July 1975 enlarged plenum of the CC MAC as mistaken and ordered all
units to stop studying and implementing them.\footnote{Hao Mengbi and
  Duan Haoran, Liushinian, p.~650. 郝梦笔、段浩然，《六十年》，第650页。}
江青也在这次会议上发言，会议似乎从1976年2月20日前后持续到3月初。\footnote{For
  details of Jiang's February 23 and March 2 talks, see ``Documents of
  the CC of the CCP'', I\&S, 15:2 pp.~96-9. Fenwick, ``The Gang of
  Four'', pp.~427-30, has analyzed Jiang's speeches.
  江青2月23日和3月2日谈话的详情，见 ``Documents of the CC of the CCP'',
  I\&S, 15:2，第96-9页。Fenwick, ``The Gang of
  Four''，第427-30页，对江的讲话作了分析。}
3月2日，她召集十二个省代表开会，指责邓是``大叛徒''和``国际资产阶级的代表''。2月16日，中共中央批准并转发了一份中央军委报告，称邓和叶在1975年6月至7月中央军委扩大会议上的讲话是错误的，并命令所有单位停止学习和执行这些讲话。\footnote{Hao
  Mengbi and Duan Haoran, Liushinian, p.~650.
  郝梦笔、段浩然，《六十年》，第650页。}

The changes at the center immediately impacted on Zhejiang. The previous
chapter, in discussing the leadership changes which had been agreed on
by the Politburo in mid-1975, commented that Deng Xiaoping had failed to
install a new leading body in Zhejiang that met his leadership criteria.
The hand of the radical mobilizers was seen in several appointments. Of
the new appointees, only Lü Jianguang and perhaps Feng Ke would have
satisfied Deng, and they were far from being the most senior members of
the new provincial elite. The disappearance of Tan Qilong after
September and more importantly of Tie Ying after November 1975,
particularly when the latter had done his utmost between July and
November to implement central instructions, further weakened the
leadership of Zhejiang. Tie Ying's outspoken support for Deng Xiaoping's
program was most probably the reason behind his forced withdrawal from
the provincial political scene. Now, with Deng himself sidelined once
again, the central radicals were able to promote party leaders who were
more attuned to the new direction being pursued by Beijing.
中央的变化立即影响到浙江。上一章在讨论1975年中期政治局同意的领导层变动时指出，邓小平未能在浙江安置一个符合其领导标准的新领导机构。几个任命显示出激进动员派的影响。在新任命者中，只有吕剑光，也许还有冯恪，会使邓满意，而他们远不是新省级精英中最资深的成员。谭启龙在9月后消失，更重要的是铁瑛在1975年11月后消失，尤其是在后者7月至11月竭力执行中央指示的情况下，进一步削弱了浙江领导层。铁瑛公开支持邓小平方案，很可能是他被迫退出省级政治舞台的原因。现在，邓本人再次被边缘化，中央激进派得以提拔更适应北京新方向的党内领导人。

After his promotion in July 1975, Lai Keke had left Zhejiang for
Guangzhou to recover from an illness.\footnote{This paragraph is based
  on material published in Tang Yurui, HZRB, December 8, 1976; Chu Yang,
  ``彻底砸烂四人帮在浙江的资产阶级帮派体系'' (Thoroughly smash the gang
  of four's bourgeois factional set-up in Zhejiang), ZJRB, May 21, 1977,
  p.~1; Shen Weicai, HZRB, August 17, 1977; Chu Yang, HZRB, November 2,
  1977. 本段依据唐毓瑞发表于
  HZRB，1976年12月8日的材料；楚扬《彻底砸烂四人帮在浙江的资产阶级帮派体系》（Thoroughly
  smash the gang of four's bourgeois factional set-up in
  Zhejiang），ZJRB，1977年5月21日，第1页；沈维才，HZRB，1977年8月17日；楚扬，HZRB，1977年11月2日。}
He returned to Hangzhou in January 1976 at a most politically opportune
moment and reported in writing to Wang Hongwen that ``because many
leaders of the Zhejiang Provincial Committee are sick I have already
started to take up some work''. Lai immediately called for a temporary
freeze on all appointments and dismissals in the province and reportedly
set aside the new leading group of the ZPC Organization Department,
headed by Feng Ke. In February 1976, Lai was presented with the
opportunity to play a more active role in provincial affairs.
赖可可在1975年7月升职后离开浙江去广州养病。\footnote{This paragraph is
  based on material published in Tang Yurui, HZRB, December 8, 1976; Chu
  Yang, ``彻底砸烂四人帮在浙江的资产阶级帮派体系'' (Thoroughly smash the
  gang of four's bourgeois factional set-up in Zhejiang), ZJRB, May 21,
  1977, p.~1; Shen Weicai, HZRB, August 17, 1977; Chu Yang, HZRB,
  November 2, 1977. 本段依据唐毓瑞发表于
  HZRB，1976年12月8日的材料；楚扬《彻底砸烂四人帮在浙江的资产阶级帮派体系》（Thoroughly
  smash the gang of four's bourgeois factional set-up in
  Zhejiang），ZJRB，1977年5月21日，第1页；沈维才，HZRB，1977年8月17日；楚扬，HZRB，1977年11月2日。}
他在1976年1月回到杭州，时机在政治上极为有利，并向王洪文书面报告说：``由于浙江省委许多领导同志有病，我已经开始抓一些工作。''赖立即要求暂时冻结全省所有任免，并据说搁置了以冯恪为首的浙江省委组织部新领导班子。1976年2月，赖获得了在省务中发挥更积极作用的机会。

When he visited the Zhejiang delegation at the February Beijing meeting,
Wang Hongwen praised Lai Keke and criticized Tan Qilong and Tie Ying. He
verbally appointed Lai to take charge of the CCP ZPC and Luo Yi to
handle routine administrative tasks. In the absence of Tan Qilong and
Tie Ying from the political scene in 1976 Lai assumed leadership of the
CCP ZPC. Between 9 April and 12 October 1976 he made sixteen public
appearances in Hangzhou.\footnote{See, CIA Reference Aid, Appearances
  and Activities of Leading Personalities of the PRC, 1976 - CR 12059
  (May 1977). 见 CIA Reference Aid, Appearances and Activities of
  Leading Personalities of the PRC, 1976 - CR 12059（1977年5月）。}
Whether Lai's promotion was due to Tan and Tie's ill-health or to Wang
Hongwen's initiative remains a moot point. If Tan and Tie were not well
enough to assume their responsibilities, Lai was the natural choice to
deputize for them as he was the next most senior member of the ZPC
standing committee. Similarly, Luo Yi, who was the most senior civilian
cadre after Lai Keke, would have been expected to stand in as Lai's
deputy. Although Tie did not appear in public for most of 1976, by
mid-year his health had permitted him to resume duty.\footnote{See Chu
  Yang, HZRB, November 2, 1977; CCP HMC Organization Department
  criticism group, ``The crux of `hold tight to the question of
  organization' was the seizure of power'', HZRB, March 27, 1977; HZRB,
  January 29, 1977.
  见楚扬，HZRB，1977年11月2日；中共杭州市委组织部批判组《``抓住组织问题不放''的要害是夺权》，HZRB，1977年3月27日；HZRB，1977年1月29日。}
With Tie's return to public life at Mao's funeral service in Hangzhou,
and from that time until Lai's purge in October 1976, the provincial
press consistently listed Tie's name before Lai's, indicating correct
seniority. However, in an unusual but possibly deliberate gesture,
Xinhua listed Lai before Tie in its round-up of provincial mourning
services for the Chairman.\footnote{See Xinhua, September 20, 1976,
  SWB/FE/5317/C/3-5. Hangzhou Daily listed both men in the correct
  order. See, HZRB, September 19, 1976.
  见新华社，1976年9月20日，SWB/FE/5317/C/3-5。《杭州日报》按正确顺序列出二人。见
  HZRB，1976年9月19日。}
王洪文在2月北京会议期间看望浙江代表团时，表扬赖可可，批评谭启龙和铁瑛。他口头指定赖负责中共浙江省委，罗毅处理日常行政事务。1976年，在谭启龙和铁瑛缺席政治舞台的情况下，赖接掌中共浙江省委领导权。1976年4月9日至10月12日，他在杭州公开露面十六次。\footnote{See,
  CIA Reference Aid, Appearances and Activities of Leading Personalities
  of the PRC, 1976 - CR 12059 (May 1977). 见 CIA Reference Aid,
  Appearances and Activities of Leading Personalities of the PRC, 1976 -
  CR 12059（1977年5月）。}
赖的提升究竟是由于谭、铁健康不佳，还是由于王洪文的主动安排，仍是一个悬而未决的问题。如果谭、铁健康不足以履行职责，赖作为浙江省委常委中资历仅次于他们的人，自然是代理他们的选择。同样，罗毅作为赖之后最资深的文职干部，也会被期待担任赖的副手。尽管铁在1976年大部分时间没有公开露面，但到年中他的健康已允许他恢复工作。\footnote{See
  Chu Yang, HZRB, November 2, 1977; CCP HMC Organization Department
  criticism group, ``The crux of `hold tight to the question of
  organization' was the seizure of power'', HZRB, March 27, 1977; HZRB,
  January 29, 1977.
  见楚扬，HZRB，1977年11月2日；中共杭州市委组织部批判组《``抓住组织问题不放''的要害是夺权》，HZRB，1977年3月27日；HZRB，1977年1月29日。}
随着铁在杭州毛泽东追悼会上重新公开露面，从那时到赖于1976年10月被清洗，省内报刊始终把铁的名字列在赖之前，表明资历顺序正确。不过，新华社在汇总各省主席追悼活动时把赖列在铁之前，这一不同寻常的做法可能是有意为之。\footnote{See
  Xinhua, September 20, 1976, SWB/FE/5317/C/3-5. Hangzhou Daily listed
  both men in the correct order. See, HZRB, September 19, 1976.
  见新华社，1976年9月20日，SWB/FE/5317/C/3-5。《杭州日报》按正确顺序列出二人。见
  HZRB，1976年9月19日。}

Lai and Luo drew up a summary of the Beijing discussions entitled ``Our
initial views on the symposium'' (我们参加座谈会的几点初步看法), in
which they distanced themselves from responsibility for implementing the
July 1975 decisions. Lai and Luo wrote that leaders such as Tie Ying,
who had loyally carried out central directives, had been influenced by a
``reactionary outlook'' and had ``suppressed new-born socialist
things''. Previous criticism of bourgeois factionalism was described as
``not stressing the principal contemporary contradiction'', and
``arbitrarily treating issues as a matter of principle'' (乱上纲). The
two newly-appointed leaders characterized the study classes held by the
ZPC for leaders of factional groups, as ``wrong'' and ``hurting the
comrades''. In criticizing the decisions of 1975 Lai and Luo were
writing in the knowledge that they had powerful central backing for
their position.
赖和罗起草了一份北京讨论情况的总结，题为《我们参加座谈会的几点初步看法》，在其中他们撇清了执行1975年7月决定的责任。赖、罗写道，铁瑛等忠实执行中央指示的领导人受到了``反动观点''的影响，并``压制社会主义新生事物''。此前对资产阶级派性的批判被描述为``没有抓住当前主要矛盾''和``乱上纲''。两位新任领导人把浙江省委为派性集团领导人举办的学习班定性为``错误的''和``伤害同志的''。赖、罗在批评1975年决定时，清楚知道自己的立场有强大的中央支持。

Upon their return to Hangzhou, Lai and Luo called a work meeting of the
ZPC to inform its members of the new leadership arrangements and the
strategy for another potentially destabilizing political
campaign.\footnote{Zhou Feng, ``Expose and criticize the gang of four's
  towering crimes in forming cliques for their own end to usurp power'',
  HZRB, January 27, 1977.
  周峰《揭露批判四人帮结党营私篡党夺权的滔天罪行》，HZRB，1977年1月27日。}
At the conference, a leading cadre of the CCP HMC lobbied to collect
signatures for a petition to be sent to Beijing, requesting that Tan and
Tie return to Zhejiang to answer questions about the direction in which
they had led the provincial committee in the latter half of 1975. The
cadre in question also withdrew the self-criticism that he had made the
previous July, confessing only too truthfully that ``I could only speak
like that, otherwise I wouldn't have got through. There was no
alternative''.\footnote{HZRB, August 17, 1977; Zhou Feng, HZRB, October
  24, 1977. HZRB，1977年8月17日；周峰，HZRB，1977年10月24日。} In the
constant see-saw of political fortunes, the balance had again tipped to
one side, flinging off or reseating players in the process.
回到杭州后，赖和罗召开浙江省委工作会议，向省委成员通报新的领导安排和又一场可能造成不稳定的政治运动策略。\footnote{Zhou
  Feng, ``Expose and criticize the gang of four's towering crimes in
  forming cliques for their own end to usurp power'', HZRB, January 27,
  1977.
  周峰《揭露批判四人帮结党营私篡党夺权的滔天罪行》，HZRB，1977年1月27日。}
会上，一名中共杭州市委领导干部游说收集签名，准备向北京递交请愿书，要求谭、铁返回浙江，就他们在1975年下半年带领省委走向何方接受质询。这名干部还撤回了他前一年7月所作的自我批评，并说出了十分真实的坦白：``我只能那样讲，不然过不了关，没有办法。''\footnote{HZRB,
  August 17, 1977; Zhou Feng, HZRB, October 24, 1977.
  HZRB，1977年8月17日；周峰，HZRB，1977年10月24日。}
在政治命运不断上下摇摆中，天平再次向一边倾斜，其间有人被抛下，也有人重新就座。

Lai and Luo focused their dissatisfaction regarding the July 1975
decisions on four major issues: they equated criticism of bourgeois
factionalism with the suppression of ``socialist new-born forces''; the
study classes which had been held for core members of the ``mountain
top'' faction were described as being ``wrong from their guiding
ideology down to their concrete methods''; the adjustment to leading
groups was characterized as ``allowing people who haven't changed in the
Cultural Revolution to grab power''; and finally, the freeze placed on
the unorthodox methods in recruiting party members and promoting cadres
was viewed as an attack on new cadres.\footnote{Chu Yang, HZRB, November
  2, 1977; Tie Ying, HZRB, June 7, 1978.
  楚扬，HZRB，1977年11月2日；铁瑛，HZRB，1978年6月7日。}
赖和罗把他们对1975年7月决定的不满集中在四个主要问题上：他们把批判资产阶级派性等同于压制``社会主义新生力量''；把为``山上派''核心成员举办的学习班描述为``从指导思想到具体方法都是错误的''；把调整领导班子概括为``让文化大革命中没有转变的人夺权''；最后，把冻结非正统入党和提拔干部方式视为对新干部的攻击。\footnote{Chu
  Yang, HZRB, November 2, 1977; Tie Ying, HZRB, June 7, 1978.
  楚扬，HZRB，1977年11月2日；铁瑛，HZRB，1978年6月7日。}

In a repetition of the tactics which had been used to such effect in
late 1973, a follower of Weng Senhe from the silk complex led a group to
storm the provincial work meeting and occupy the living quarters of
representatives from the eight factories.\footnote{Tie Ying, HZRB, June
  7, 1978; Zhou Feng, HZRB, October 24, 1977.
  铁瑛，HZRB，1978年6月7日；周峰，HZRB，1977年10月24日。} Qiu Qiang,
Deputy-secretary of the CCP HMC, reportedly described the leading group
of six members who had been appointed to the Zhejiang Hemp Mill in 1975
as capitalist-roaders, and declared that the group should be dissolved.
In response to this signal, rebels in the mill took away the group's
seal of authority and locked the mill's administrative offices. By these
measures the rebels hoped to obtain the recall of cadres who had been
transferred or demoted in July 1975.\footnote{CCP HMC Organization
  Department, HZRB, March 27, 1977.
  中共杭州市委组织部，HZRB，1977年3月27日。}
重复1973年末曾产生显著效果的战术，丝绸联合企业中翁森鹤的一名追随者带领一群人冲击省工作会议，并占据八厂代表的住所。\footnote{Tie
  Ying, HZRB, June 7, 1978; Zhou Feng, HZRB, October 24, 1977.
  铁瑛，HZRB，1978年6月7日；周峰，HZRB，1977年10月24日。}
据说中共杭州市委副书记邱强把1975年任命到浙江麻纺厂的六人领导小组称为走资派，并宣布该小组应予解散。接到这一信号后，厂内造反派夺走了该小组的公章，并锁闭工厂行政办公室。通过这些措施，造反派希望使1975年7月被调离或降职的干部被召回。\footnote{CCP
  HMC Organization Department, HZRB, March 27, 1977.
  中共杭州市委组织部，HZRB，1977年3月27日。}

The ``eight factories' experience'' also came under fire as a
``revisionist model''. The ``reversal of verdicts'' originated in Wang
Hongwen's ambivalence on the issue. In conjunction with Ji Dengkui Wang
had been instrumental in establishing the model in July 1975. Even in
December 1975 and January 1976 he had continued to support it and
requested that the CCP ZPC write a statement to this effect and issue it
so that ``the situation will develop even better''. In February 1976,
Wang seemed to fudge his position by pleading ignorance of developments
since his investigations in Hangzhou. In the following month Lai and Luo
reportedly refused permission for Zhejiang Daily to publish the silk
complex's production figures for February. News of this decision quickly
spread around Hangzhou and confidence in the direction in which the
eight factories had been heading was thereby shaken and eventually
shattered.\footnote{Tie Ying, speech of November 22, 1976, HZRB,
  November 24, 1976; HZRB, November 28, 1976; Bian Wen, HZRB, December
  5, 1976.
  铁瑛，1976年11月22日讲话，HZRB，1976年11月24日；HZRB，1976年11月28日；边文，HZRB，1976年12月5日。}
``八厂经验''也作为``修正主义样板''受到攻击。``翻案''源于王洪文在这个问题上的矛盾态度。1975年7月，王与纪登奎共同促成了这一样板的建立。甚至在1975年12月和1976年1月，他仍继续支持它，并要求中共浙江省委写出并发布一份声明，使``形势会发展得更好''。1976年2月，王似乎含糊其词，声称自己不了解在杭州调查以后发生的情况。次月，据说赖和罗拒绝允许《浙江日报》发表丝绸联合企业2月生产数字。这个决定的消息迅速传遍杭州，对八厂所走方向的信心因此动摇，并最终破碎。\footnote{Tie
  Ying, speech of November 22, 1976, HZRB, November 24, 1976; HZRB,
  November 28, 1976; Bian Wen, HZRB, December 5, 1976.
  铁瑛，1976年11月22日讲话，HZRB，1976年11月24日；HZRB，1976年11月28日；边文，HZRB，1976年12月5日。}

During their seven-month ascendancy in Zhejiang, Lai Keke and Luo Yi
made strenuous efforts to secure the province for their patrons in
Beijing. Most notably, they took advantage of the outburst of mourning
for the deceased Zhou Enlai which erupted in April 1976, to attempt to
track down and arrest those responsible for Hangzhou's Qingming
incident.\footnote{Keith Forster, ``The 1976 Ch'ing-ming Incident in
  Hangchow'', I\&S, 22:4 (April 1986), pp.~13-33. Keith Forster, ``The
  1976 Ch'ing-ming Incident in Hangchow'', I\&S,
  22:4（1986年4月），第13-33页。} Lai also assiduously endorsed the
thesis propagated by the radicals concerning the formation of
capitalist-roaders inside the CCP. His espousal of the theory assisted
the Cultural Revolution radicals in Beijing in their effort to widen the
scope of the campaign against Deng Xiaoping, so as to include his
supporters in leading positions across the country. Naturally many
veteran cadres resisted the trend, hoping instead to limit the campaign
to a perfunctory denunciation of the dismissed Deng.
赖可可和罗毅在浙江掌权的七个月里，竭力为北京的庇护者控制该省。最显著的是，他们利用1976年4月周恩来逝世后爆发的悼念浪潮，试图追查并逮捕杭州清明事件的责任人。\footnote{Keith
  Forster, ``The 1976 Ch'ing-ming Incident in Hangchow'', I\&S, 22:4
  (April 1986), pp.~13-33. Keith Forster, ``The 1976 Ch'ing-ming
  Incident in Hangchow'', I\&S, 22:4（1986年4月），第13-33页。}
赖还勤勉支持激进派宣传的关于中共内部形成走资派的论点。他拥护这一理论，帮助北京的文化大革命激进派扩大反邓小平运动的范围，把全国各地领导岗位上的邓的支持者也包括进去。自然，许多老干部抵制这一趋势，反而希望把运动限制在对已被免职的邓进行例行谴责的程度。

To celebrate the tenth anniversary of the publication of the May 16
Circular, on 15 May 1976 the provincial party committee held a
celebration rally in Hangzhou. In his speech to the assembled crowd, Lai
Keke quoted Mao from the May 16 Circular as his authority for the view
that representatives of the bourgeoisie existed in provincial
departments, thus wishing to legitimize a purge of cadres in the
province.\footnote{ZPS, May 16, 1976 in SWB/FE/5213/BII/9-12. See also
  Lai's speech at the May 6 rally held to criticize Deng in ZPS, May 8,
  1976, SWB/FE/5208/BII/5; HZRB, November 28, 1976.
  ZPS，1976年5月16日，载
  SWB/FE/5213/BII/9-12。另见赖在5月6日批邓大会上的讲话，ZPS，1976年5月8日，SWB/FE/5208/BII/5；HZRB，1976年11月28日。}
The day after the rally, the ZPC propaganda department called a
symposium to discuss the question of the existence of a bourgeoisie
inside the Communist Party. The seminar lasted nine days and was
attended by full-time propaganda workers, worker-peasant-soldier
theorists and grassroots cadres, numbering eighty in all.\footnote{HZRB,
  May 28, 1976; ZPS, May 31, 1976, in SWB/FE/5225/BII/10.
  HZRB，1976年5月28日；ZPS，1976年5月31日，载 SWB/FE/5225/BII/10。} One
week after the conclusion of the symposium, the ZPC convened a report
meeting to listen to speeches from some of the forum's participants. Lai
Keke delivered an address on behalf of the provincial
authorities.\footnote{ZPS, June 5, 1976, in SWB/FE/5230/BII/7-8.
  ZPS，1976年6月5日，载 SWB/FE/5230/BII/7-8。}
为庆祝《五一六通知》发表十周年，1976年5月15日，省委在杭州举行庆祝大会。赖可可在向集会群众发表讲话时，引用毛在《五一六通知》中的话，作为省级部门中存在资产阶级代表人物这一看法的依据，从而试图使清洗本省干部合法化。\footnote{ZPS,
  May 16, 1976 in SWB/FE/5213/BII/9-12. See also Lai's speech at the May
  6 rally held to criticize Deng in ZPS, May 8, 1976, SWB/FE/5208/BII/5;
  HZRB, November 28, 1976. ZPS，1976年5月16日，载
  SWB/FE/5213/BII/9-12。另见赖在5月6日批邓大会上的讲话，ZPS，1976年5月8日，SWB/FE/5208/BII/5；HZRB，1976年11月28日。}
集会次日，省委宣传部召开座谈会，讨论共产党内是否存在资产阶级的问题。研讨会持续九天，参加者包括专职宣传干部、工农兵理论工作者和基层干部，共80人。\footnote{HZRB,
  May 28, 1976; ZPS, May 31, 1976, in SWB/FE/5225/BII/10.
  HZRB，1976年5月28日；ZPS，1976年5月31日，载 SWB/FE/5225/BII/10。}
座谈会结束一周后，省委召开报告会，听取部分与会者的发言。赖可可代表省级当局发表讲话。\footnote{ZPS,
  June 5, 1976, in SWB/FE/5230/BII/7-8. ZPS，1976年6月5日，载
  SWB/FE/5230/BII/7-8。}

On 1 July, in a speech at the meeting called to celebrate the 55th
anniversary of the CCP, Lai adopted a more conciliatory tone toward
``comrades'' who had made mistakes under the influence of Deng's
revisionist line. He urged them to heighten their political
consciousness and adopt the correct attitude toward the Cultural
Revolution, the masses and themselves, and take practical steps to
rectify their errors. Nevertheless, Lai reiterated his determination
that Zhejiang would carry out a protracted struggle against the
bourgeoisie both inside and outside the party.\footnote{ZPS, July 2,
  1976, SWB/FE/5254/BII/10-11. Lai may have taken his cue from a speech
  by Zhang Chunqiao of June 28, 1976. In it Zhang cautioned against
  haste and the use of excessive measures in dealing with
  counterrevolutionaries. See Classified Chinese Communist Documents,
  pp.~422-35.
  ZPS，1976年7月2日，SWB/FE/5254/BII/10-11。赖可能受到张春桥1976年6月28日讲话的启发。张在讲话中提醒处理反革命分子时不要急躁，也不要采取过度措施。见
  Classified Chinese Communist Documents，第422-35页。} Several days
later, when the ZPMD held a meeting to discuss the problem of
capitalist-roaders inside the CCP, it concluded, unlike the ZPC, that
the majority of cadres were good.\footnote{ZPS, July 5, 1976, in
  SWB/FE/5258/BII/3. ZPS，1976年7月5日，载 SWB/FE/5258/BII/3。}
7月1日，在庆祝中共成立55周年的会议讲话中，赖对那些受邓的修正主义路线影响而犯错误的``同志''采取了较为和缓的语气。他敦促他们提高政治觉悟，对文化大革命、群众和自己采取正确态度，并采取实际步骤改正错误。不过，赖重申了浙江将在党内外同资产阶级进行长期斗争的决心。\footnote{ZPS,
  July 2, 1976, SWB/FE/5254/BII/10-11. Lai may have taken his cue from a
  speech by Zhang Chunqiao of June 28, 1976. In it Zhang cautioned
  against haste and the use of excessive measures in dealing with
  counterrevolutionaries. See Classified Chinese Communist Documents,
  pp.~422-35.
  ZPS，1976年7月2日，SWB/FE/5254/BII/10-11。赖可能受到张春桥1976年6月28日讲话的启发。张在讲话中提醒处理反革命分子时不要急躁，也不要采取过度措施。见
  Classified Chinese Communist Documents，第422-35页。}
几天后，浙江省军区召开会议讨论中共内部走资派问题时，与省委不同，会议得出结论认为大多数干部是好的。\footnote{ZPS,
  July 5, 1976, in SWB/FE/5258/BII/3. ZPS，1976年7月5日，载
  SWB/FE/5258/BII/3。}

Apart from displaying excessive zeal in publicly aligning himself with
the Gang of Four, Lai kept up a sustained attack on the decisions of
July 1975. His most savage and comprehensive onslaught appeared in an
article which was published on the eve of the arrest of the Gang of
Four. Written on the basis of the speech that Lai had delivered on 28
September 1976, it appeared in the Zhejiang Daily on 30 September. The
article achieved such notoriety that it was compared with the Guangming
Daily article of 4 October 1976, penned by the writing team known as
Liang Xiao and published under the instructions of Yao Wenyuan. Like
Liang Xiao's piece, Lai's article was later designated a
``mobilizational order for a counter-revolutionary coup'' and a
``clarion call to support the Gang of Four's plot to seize
power''.\footnote{Shang Jingcai, HZRB, November 28, 1976, pp.~2-3; Shen
  Weicai, HZRB, August 17, 1977; Shen Cai,
  ``一面层层揪走资派的黑旗--批判'四人帮'在我省的一个代理人九三〇反党黑文''
  (A black flag to haul out `capitalist-roaders' at all levels - a
  criticism of the article of September 30 written by our province's
  agent of the ``Gang of Four''), HZRB, November 27, 1977. In his
  capacity as deputy leader of the leading group in the ZPC Propaganda
  Department, Shang Jingcai spoke at the provincial criticism rally of
  22 November 1976, which was held to denounce the agents and followers
  of the Gang of Four in Zhejiang. HZRB, November 24, 1976. It is highly
  likely that Shen Weicai and Shen Cai were pen-names used by Shang
  Jingcai.
  尚景才，HZRB，1976年11月28日，第2-3页；沈维才，HZRB，1977年8月17日；沈才《一面层层揪走资派的黑旗--批判``四人帮''在我省的一个代理人九三〇反党黑文》（A
  black flag to haul out `capitalist-roaders' at all levels - a
  criticism of the article of September 30 written by our province's
  agent of the ``Gang of
  Four''），HZRB，1977年11月27日。尚景才作为浙江省委宣传部领导小组副组长，在1976年11月22日召开的省批判大会上发言，该大会旨在谴责``四人帮''在浙江的代理人和追随者。HZRB，1976年11月24日。沈维才和沈才很可能是尚景才使用的笔名。}
除了过分热衷于公开同``四人帮''站在一起之外，赖还持续攻击1975年7月的决定。他最猛烈、最全面的攻击出现在``四人帮''被捕前夕发表的一篇文章中。该文根据赖1976年9月28日讲话写成，9月30日刊登于《浙江日报》。这篇文章臭名昭著，以至于被拿来同1976年10月4日《光明日报》文章相比；后者由名为梁效的写作班子撰写，并按姚文元指示发表。与梁效的文章一样，赖的文章后来被定性为``反革命政变的动员令''和``支持`四人帮'篡党夺权阴谋的冲锋号''。\footnote{Shang
  Jingcai, HZRB, November 28, 1976, pp.~2-3; Shen Weicai, HZRB, August
  17, 1977; Shen Cai,
  ``一面层层揪走资派的黑旗--批判'四人帮'在我省的一个代理人九三〇反党黑文''
  (A black flag to haul out `capitalist-roaders' at all levels - a
  criticism of the article of September 30 written by our province's
  agent of the ``Gang of Four''), HZRB, November 27, 1977. In his
  capacity as deputy leader of the leading group in the ZPC Propaganda
  Department, Shang Jingcai spoke at the provincial criticism rally of
  22 November 1976, which was held to denounce the agents and followers
  of the Gang of Four in Zhejiang. HZRB, November 24, 1976. It is highly
  likely that Shen Weicai and Shen Cai were pen-names used by Shang
  Jingcai.
  尚景才，HZRB，1976年11月28日，第2-3页；沈维才，HZRB，1977年8月17日；沈才《一面层层揪走资派的黑旗--批判``四人帮''在我省的一个代理人九三〇反党黑文》（A
  black flag to haul out `capitalist-roaders' at all levels - a
  criticism of the article of September 30 written by our province's
  agent of the ``Gang of
  Four''），HZRB，1977年11月27日。尚景才作为浙江省委宣传部领导小组副组长，在1976年11月22日召开的省批判大会上发言，该大会旨在谴责``四人帮''在浙江的代理人和追随者。HZRB，1976年11月24日。沈维才和沈才很可能是尚景才使用的笔名。}

The background to the publication of Lai's article deserves mention. At
a meeting of the standing committee of the ZPC on 25 September 1976, Lai
Keke had reportedly suggested that the ZPC convene a criticism meeting
against Deng Xiaoping before National Day and he volunteered to speak at
such a meeting. However, Lai supposedly refused to show the outline of
his speech to the standing committee for its advance approval. The
meeting took place on 28 September with Tie Ying delivering the
concluding speech on behalf of the ZPC.\footnote{HZRB, September 29,
  1976; ZJRB, September 29, 1976, pp.~1, 3.
  HZRB，1976年9月29日；ZJRB，1976年9月29日，第1、3页。} Lai Keke also
addressed the gathering, together with the peasant leaders Jiang Ruwang
from Jinjian brigade, Wang Jinyou from Shangwang brigade, Li Jinrong
from Nanbao brigade and representatives of offices and factories in the
province.\footnote{Wang Jinyou's article of early October 1976 was most
  probably based on his speech at this meeting. See n.~135, chapter
  nine.
  王金友1976年10月初的文章很可能以他在这次会议上的讲话为基础。见第九章注135。}
Tie delivered the regulation denunciation of the arch-villain Deng
Xiaoping and made it clear that the criticism should focus on Deng
himself and should be carried out under the leadership of the CCP. Tie
stressed that
赖的文章发表背景值得一提。据说在1976年9月25日浙江省委常委会会议上，赖可可建议省委在国庆节前召开一次批邓小平大会，并主动提出在会上讲话。然而，赖据称拒绝把讲话提纲交常委会事先批准。会议于9月28日举行，铁瑛代表省委作总结讲话。\footnote{HZRB,
  September 29, 1976; ZJRB, September 29, 1976, pp.~1, 3.
  HZRB，1976年9月29日；ZJRB，1976年9月29日，第1、3页。}
赖可可也在会上讲话，同台发言的还有金建大队的农民领导人姜汝旺、上旺大队的王金友、南堡大队的李金荣，以及本省机关和工厂代表。\footnote{Wang
  Jinyou's article of early October 1976 was most probably based on his
  speech at this meeting. See n.~135, chapter nine.
  王金友1976年10月初的文章很可能以他在这次会议上的讲话为基础。见第九章注135。}
铁作了按规定谴责头号恶人邓小平的讲话，并明确表示批判应集中于邓本人，并应在中共领导下进行。铁强调

\begin{quote}
It is not allowed to establish ties and to build mountain-strongholds
under any pretext and any form; it is not allowed to organize fighting
groups. 不准以任何借口、任何形式串联和搞山头；不准组织战斗队。
\end{quote}

Lai Keke reportedly modified the contents of his speech for publication
and then submitted it to the Zhejiang Daily without the prior knowledge
or approval of the provincial party committee's standing
committee.\footnote{Tie Ying, HZRB, November 24, 1976.
  铁瑛，HZRB，1976年11月24日。} However, the decision to publish Lai's
article presumably could not have occurred without the consent of this
body, of which Tie was the senior ranking cadre. On 8 October 1976 Lai
sent a copy to Wang Hongwen in the hope of gaining kudos for himself
(请功领赏)\footnote{Shang Jingcai, HZRB, November 28, 1976. Shortly
  after 8 October, Lai was summoned to Beijing where, together with
  other provincial leaders, he was informed of the arrest of the four
  radical central leaders and shown the slip of paper allegedly handed
  to Hua Guofeng by Mao Zedong in April 1976, on which Mao had written
  the six characters ``你办事我放心'' (with you in charge, I'm at ease).
  On October 14, 1976, at a meeting of the standing committee of the CCP
  ZPC, Lai expressed skepticism about the authenticity of the document,
  saying; ``I have seen the Chairman's handwriting many times. That slip
  `with you in charge, I'm at ease' simply does not look like his own
  handwriting. Furthermore, why were we shown only a photocopy and not
  the original?'' See Shen Weicai, HZRB, August 17, 1977; Tie Ying,
  HZRB, June 7, 1978; and Li Yanwu, ``The downfall of the Lin-Jiang
  cliques-rise and fall of the ten culprits of the `Cultural
  Revolution''', 李鸣 (Zhengming) (H.K.), 37 (November 1, 1980), in
  FBIS/CR/PSMA, 160 (January 30, 1981), p.~38.
  尚景才，HZRB，1976年11月28日。10月8日后不久，赖被召到北京，同其他省级领导一起获知四名中央激进派领导人被捕，并看到据称毛泽东于1976年4月交给华国锋的一张纸条，上面写着六个字``你办事我放心''（with
  you in charge, I'm at
  ease）。1976年10月14日，在中共浙江省委常委会会议上，赖对该文件真实性表示怀疑，说：``我看过主席的笔迹很多次。那张`你办事我放心'的条子根本不像他自己的笔迹。而且，为什么只给我们看复印件，不给我们看原件？''见沈维才，HZRB，1977年8月17日；铁瑛，HZRB，1978年6月7日；以及李炎武《林江集团的覆灭------``文化大革命''十名罪犯的兴衰》，李鸣（《争鸣》）（香港），第37期（1980年11月1日），载
  FBIS/CR/PSMA，第160期（1981年1月30日），第38页。} apparently unaware
of Wang's arrest two days before. Lai did not pull any punches in the
article.\footnote{Lai Keke, ``化悲痛为力量集中火力深入批邓'' (Turn grief
  into strength, concentrate our energies on thoroughly criticizing
  Deng), ZJRB, September 30, 1976, p.~2; HZRB, October 1, 1976, p.~6.
  赖可可《化悲痛为力量集中火力深入批邓》（Turn grief into strength,
  concentrate our energies on thoroughly criticizing
  Deng），ZJRB，1976年9月30日，第2页；HZRB，1976年10月1日，第6页。} He
accused Deng Xiaoping of having repeatedly issued instructions in July,
September and October of 1975, presumably to the Zhejiang party
leadership, to
据说赖可可为发表而修改了讲话内容，然后在省党委常委会不知情、未批准的情况下提交给《浙江日报》。\footnote{Tie
  Ying, HZRB, November 24, 1976. 铁瑛，HZRB，1976年11月24日。}
然而，发表赖的文章这一决定想必不可能没有常委会同意，而铁瑛是该机构中级别最高的干部。1976年10月8日，赖把一份样报送给王洪文，希望为自己请功领赏\footnote{Shang
  Jingcai, HZRB, November 28, 1976. Shortly after 8 October, Lai was
  summoned to Beijing where, together with other provincial leaders, he
  was informed of the arrest of the four radical central leaders and
  shown the slip of paper allegedly handed to Hua Guofeng by Mao Zedong
  in April 1976, on which Mao had written the six characters
  ``你办事我放心'' (with you in charge, I'm at ease). On October 14,
  1976, at a meeting of the standing committee of the CCP ZPC, Lai
  expressed skepticism about the authenticity of the document, saying;
  ``I have seen the Chairman's handwriting many times. That slip `with
  you in charge, I'm at ease' simply does not look like his own
  handwriting. Furthermore, why were we shown only a photocopy and not
  the original?'' See Shen Weicai, HZRB, August 17, 1977; Tie Ying,
  HZRB, June 7, 1978; and Li Yanwu, ``The downfall of the Lin-Jiang
  cliques-rise and fall of the ten culprits of the `Cultural
  Revolution''', 李鸣 (Zhengming) (H.K.), 37 (November 1, 1980), in
  FBIS/CR/PSMA, 160 (January 30, 1981), p.~38.
  尚景才，HZRB，1976年11月28日。10月8日后不久，赖被召到北京，同其他省级领导一起获知四名中央激进派领导人被捕，并看到据称毛泽东于1976年4月交给华国锋的一张纸条，上面写着六个字``你办事我放心''（with
  you in charge, I'm at
  ease）。1976年10月14日，在中共浙江省委常委会会议上，赖对该文件真实性表示怀疑，说：``我看过主席的笔迹很多次。那张`你办事我放心'的条子根本不像他自己的笔迹。而且，为什么只给我们看复印件，不给我们看原件？''见沈维才，HZRB，1977年8月17日；铁瑛，HZRB，1978年6月7日；以及李炎武《林江集团的覆灭------``文化大革命''十名罪犯的兴衰》，李鸣（《争鸣》）（香港），第37期（1980年11月1日），载
  FBIS/CR/PSMA，第160期（1981年1月30日），第38页。}，显然不知道王已在两天前被捕。赖在文章中毫不留情。\footnote{Lai
  Keke, ``化悲痛为力量集中火力深入批邓'' (Turn grief into strength,
  concentrate our energies on thoroughly criticizing Deng), ZJRB,
  September 30, 1976, p.~2; HZRB, October 1, 1976, p.~6.
  赖可可《化悲痛为力量集中火力深入批邓》（Turn grief into strength,
  concentrate our energies on thoroughly criticizing
  Deng），ZJRB，1976年9月30日，第2页；HZRB，1976年10月1日，第6页。}
他指责邓小平在1975年7月、9月和10月反复发布指示，想必是给浙江党的领导层，要求

\begin{quote}
give a free hand in mobilizing the people, criticize bourgeois
factionalism and carry out a mass self-educational movement to overcome
bourgeois factionalism.\footnote{These exact words were used in the
  titles of the ZJRB editorial of August 23, 1975, and the ZJRB article
  of September 30, 1975.
  这些原话用于1975年8月23日《浙江日报》社论和1975年9月30日《浙江日报》文章的标题。}
\ldots{} Deng Xiaoping used opposition to bourgeois factionalism as a
pretext for reversing the verdict on the Cultural Revolution,\footnote{The
  current party line in China accepts this argument by stating that in
  1975 Deng attempted to restore the ``correct line'' of the CCP's 8th
  Congress of 1956, which viewed the development of the economy as the
  principal task of the socialist revolution. See Zhonggong dangshi
  baiti, p.~388. The section of the 1981 Resolution of the CCP CC which
  referred to this period stated that ``Comrade Mao Zedong could not
  bear to accept systematic corrections of the errors of the `cultural
  revolution' by comrade Deng Xiaoping''. Resolution on CCP History
  (1949-81), (Beijing: Foreign Languages Press, 1981), p.~39.
  中国当前党的路线接受这一论点，称1975年邓试图恢复中共1956年第八次代表大会的``正确路线''，该路线把发展经济视为社会主义革命的主要任务。见《中共党史百题》，第388页。1981年中共中央《关于建国以来党的若干历史问题的决议》中涉及这一时期的部分说：``毛泽东同志不能容忍邓小平同志对`文化大革命'错误进行系统的纠正。''Resolution
  on CCP History (1949-81)（北京：外文出版社，1981），第39页。} settling
accounts with it in a vain attempt to bloodily suppress the proletarian
revolutionary faction, and unrestrainedly encircling and suppressing the
revolutionary faction as he and Liu Shaoqi had done in the early days of
the Cultural Revolution. He carried out a white terror and labelled a
large group of revolutionary rebels ``anti-party elements'' and
``counter-revolutionaries'' in order to annihilate the proletariat's
aspirations (志气) and boost the arrogance (威风) of the bourgeoisie.
How pernicious! (何其毒也)\footnote{The words of this last sentence come
  virtually verbatim from Mao Zedong's big-character poster of August 5,
  1966. See Mao Zedong sixiang wansui (April 1967), p.~36.
  最后一句话几乎逐字来自毛泽东1966年8月5日的大字报。见《毛泽东思想万岁》（1967年4月），第36页。}
放手发动群众，批判资产阶级派性，开展克服资产阶级派性的群众自我教育运动。\footnote{These
  exact words were used in the titles of the ZJRB editorial of August
  23, 1975, and the ZJRB article of September 30, 1975.
  这些原话用于1975年8月23日《浙江日报》社论和1975年9月30日《浙江日报》文章的标题。}
\ldots\ldots 邓小平以反对资产阶级派性为借口，为文化大革命翻案，\footnote{The
  current party line in China accepts this argument by stating that in
  1975 Deng attempted to restore the ``correct line'' of the CCP's 8th
  Congress of 1956, which viewed the development of the economy as the
  principal task of the socialist revolution. See Zhonggong dangshi
  baiti, p.~388. The section of the 1981 Resolution of the CCP CC which
  referred to this period stated that ``Comrade Mao Zedong could not
  bear to accept systematic corrections of the errors of the `cultural
  revolution' by comrade Deng Xiaoping''. Resolution on CCP History
  (1949-81), (Beijing: Foreign Languages Press, 1981), p.~39.
  中国当前党的路线接受这一论点，称1975年邓试图恢复中共1956年第八次代表大会的``正确路线''，该路线把发展经济视为社会主义革命的主要任务。见《中共党史百题》，第388页。1981年中共中央《关于建国以来党的若干历史问题的决议》中涉及这一时期的部分说：``毛泽东同志不能容忍邓小平同志对`文化大革命'错误进行系统的纠正。''Resolution
  on CCP History (1949-81)（北京：外文出版社，1981），第39页。}
同文化大革命算账，妄图血腥镇压无产阶级革命派，并像他和刘少奇在文化大革命初期所做的那样，肆无忌惮地围剿革命派。他实行白色恐怖，把一大批革命造反派打成``反党分子''和``反革命''，以打掉无产阶级的志气，助长资产阶级的威风。何其毒也！\footnote{The
  words of this last sentence come virtually verbatim from Mao Zedong's
  big-character poster of August 5, 1966. See Mao Zedong sixiang wansui
  (April 1967), p.~36.
  最后一句话几乎逐字来自毛泽东1966年8月5日的大字报。见《毛泽东思想万岁》（1967年4月），第36页。}
\end{quote}

Lai Keke described the equation of bourgeois factionalism with
revisionism and capitalism, a formula used both in Deng's 20 Points and
speeches by Tie Ying in September 1975, as ``theoretically absurd and
politically reactionary''. He complained that the mass meetings which
had been held between August and October 1975 to criticize bourgeois
factionalism had labelled some cadres and sections of the masses as
赖可可把资产阶级派性等同于修正主义和资本主义的说法称为``理论上荒谬、政治上反动''；这种公式既用于邓的``二十条''，也用于铁瑛1975年9月的讲话。他抱怨说，1975年8月至10月为批判资产阶级派性而举行的群众大会，把一些干部和部分群众打成

\begin{quote}
``anti-party, power-usurping restorationists'', ``Song Jiang-type
capitulationists'',\footnote{This was another example of the cooptation
  of slogans in order to neutralize them or turn them against one's
  opponents. 这是挪用口号以便使其失效或转而用来攻击对手的又一例子。} and
``restorationists under the signboard of opposing restoration'' -- thus
treating everything as a matter of principle and seriously confusing the
two different kinds of contradictions \ldots.
``反党篡权复辟分子''、``宋江式投降派''\footnote{This was another example
  of the cooptation of slogans in order to neutralize them or turn them
  against one's opponents.
  这是挪用口号以便使其失效或转而用来攻击对手的又一例子。}
和``打着反复辟旗号的复辟派''------从而把一切都上纲到原则问题，严重混淆了两类不同性质的矛盾\ldots\ldots{}
\end{quote}

The establishment of study classes had also been an error, wrote Lai,
and he requested that the members of such classes return to their units
and positions, an instruction which in effect rehabilitated the
followers of Zhang Yongsheng, Weng Senhe and He Xianchun. If their
original units remained in existence, he continued, in an allusion to
the dissolved militia headquarters, workers' congress and Mao Zedong
Thought propaganda teams, appropriate work should be found for the study
class members on their return.
赖写道，开办学习班也是一个错误；他要求这些学习班成员返回原单位和岗位，这一指示实际上为张永生、翁森鹤和贺贤春的追随者恢复了地位。他接着暗指已被解散的民兵指挥部、工代会和毛泽东思想宣传队说，如果他们原来的单位仍然存在，就应在他们返回后安排适当工作。

Turning his attention to the reorganization of leading groups, Lai
attacked Deng for trying to discredit the ``proletarian revolutionary
faction'' by associating it with Lin Biao, a charge to which they were
clearly vulnerable and therefore ultra-sensitive. Lai argued that Deng
had directed his fire at ``those revolutionary leading cadres who had
united in the early stages of the Cultural Revolution'' and had later
played an active role in the 1974 campaign against Lin Biao and
Confucius. Here, he was referring to himself (although he does not seem
to have been very active in the campaign -- see chapter seven) and other
provincial and municipal leaders such as Wang Zida, Shen Ce and Wang
Shichuan. With the exception of Lai, these officials shared similar
career patterns in the Cultural Revolution. In 1966-67 they had deserted
or betrayed their fellow leaders when under attack from the rebels; in
late 1967 they had recanted and been allowed to join revolutionary
committees as ``revolutionary leading cadres''; in 1973-74 they had
caved in to or gone along with the demands put forward the ``mountain
top'' faction leaders in their rebellion and destabilization of party
authority and connived at or participated in their factional activities;
and in July 1975 they had been dismissed or transferred. Lai correctly
judged that Deng viewed ``revolutionary leading cadres'' as thorns in
his side (眼中钉) and had set out to dismiss them so as to make way for
the establishment of ``restorationist, capitalist leading groups''.
Although Deng Xiaoping had been the cause of the problems in Zhejiang's
since July 1975, responsibility for the implementation of the incorrect
line, added Lai, rested with the provincial leadership. Lai's withering
critique of the July 1975 decisions made no direct mention of the PLA's
role in the factories of Hangzhou except for the reference to the bloody
suppression of the ``proletarian revolutionary faction'', which was
perhaps an allusion to the military occupation of industrial units.
谈到领导班子的改组，赖攻击邓试图通过把``无产阶级革命派''同林彪联系起来而使其名誉扫地；这是一项他们显然容易受到攻击、因而极为敏感的指控。赖认为，邓把火力对准``那些在文化大革命初期联合起来的革命领导干部''，这些人后来在1974年批林批孔运动中发挥了积极作用。在这里，他指的是他自己（尽管他似乎并未在该运动中十分活跃------见第七章）以及王子达、沈策、王世川等省市领导。除赖之外，这些官员在文化大革命中的仕途模式相似：1966-67年，他们在受到造反派攻击时抛弃或背叛了其他领导；1967年末，他们悔过并获准以``革命领导干部''身份加入革命委员会；1973-74年，他们在``山上派''领导人的叛乱和破坏党权威行动中屈服于或迎合其要求，纵容或参与其派性活动；1975年7月，他们被免职或调离。赖正确判断邓把``革命领导干部''视为眼中钉，并着手撤换他们，以便建立``复辟的、资本主义的领导班子''。赖补充说，尽管邓小平是1975年7月以来浙江问题的根源，但执行错误路线的责任在省级领导层。赖对1975年7月决定的严厉批判，没有直接提及解放军在杭州各工厂中的作用，除非把他提到血腥镇压``无产阶级革命派''视为对军队占领工业单位的暗指。

Overall, Lai's article provided important first-hand evidence of how a
leading member of the losing side in the political struggles in Zhejiang
viewed the events of July 1975. Deng Xiaoping had delivered a severe
blow to the fortunes of the rebel faction and their allies in the
provincial and municipal party leaderships, but the fact that Lai had
been promoted in the 1975 reshuffle was eloquent testimony to the
incompleteness of Deng's victory. Factors external to and larger than
issues of a provincial nature determined the final settlement of the
Zhejiang problem.
总体而言，赖的文章为理解浙江政治斗争中失败一方的一名领导成员如何看待1975年7月事件，提供了重要的一手证据。邓小平沉重打击了造反派及其在省、市党委领导层盟友的前途，但赖在1975年改组中得到提升这一事实，有力说明邓的胜利并不彻底。决定浙江问题最终解决的，是外部的、比省级问题更大的因素。

\section[Aftermath, 1976-1990]{\texorpdfstring{Aftermath,
1976-1990\footnote{For a more detailed account of the effect of the
  Cultural Revolution on the politics of Zhejiang in the post-Mao period
  see, Keith Forster, ``Repudiation of the Cultural Revolution in China:
  The Case of Zhejiang'', Pacific Affairs, 59: 1 (Spring 1986),
  esp.~pp.~16-24. 关于文化大革命对后毛时期浙江政治影响的更详细叙述，见
  Keith Forster, ``Repudiation of the Cultural Revolution in China: The
  Case of Zhejiang'', Pacific Affairs, 59:1（1986年春），尤见第16-24页。}}{Aftermath, 1976-1990}}\label{aftermath-1976-1990chapter-10-83}

\section[后果，1976-1990年]{\texorpdfstring{后果，1976-1990年\footnote{For
  a more detailed account of the effect of the Cultural Revolution on
  the politics of Zhejiang in the post-Mao period see, Keith Forster,
  ``Repudiation of the Cultural Revolution in China: The Case of
  Zhejiang'', Pacific Affairs, 59: 1 (Spring 1986), esp.~pp.~16-24.
  关于文化大革命对后毛时期浙江政治影响的更详细叙述，见 Keith Forster,
  ``Repudiation of the Cultural Revolution in China: The Case of
  Zhejiang'', Pacific Affairs, 59:1（1986年春），尤见第16-24页。}}{后果，1976-1990年}}\label{ux540eux679c1976-1990ux5e74chapter-10-83}

One western analyst of provincial affairs, writing in 1976, observed
that the problems in Zhejiang were rife and protracted, and that changes
in its provincial party secretariat were ``likely to become ultimately
necessary''.\footnote{Scalapino, ``The CCP's Provincial Secretaries'',
  p.~35. Scalapino, ``The CCP's Provincial Secretaries''，第35页。} The
provincial authorities also admitted at the time that factionalism had
seriously afflicted cadres and citizens alike. A report of the planning
meeting which was held by the ZPRC in September 1976 included the
following passage:
一位研究省级事务的西方分析人士在1976年写道，浙江问题普遍而持久，其省委书记处的变动``很可能最终成为必要''。\footnote{Scalapino,
  ``The CCP's Provincial Secretaries'', p.~35. Scalapino, ``The CCP's
  Provincial Secretaries''，第35页。}
省级当局当时也承认，派性已经严重侵蚀干部和群众。1976年9月浙江省革命委员会召开的计划会议报告包括以下一段：

\begin{quote}
As for the differences of opinion and lack of mutual understanding among
certain cadres and masses, \ldots{} these were instigated by
capitalist-roaders like Liu Shaoqi, Lin Biao and Deng Xiaoping. It was
necessary to do meticulous ideological and political work among these
people and to urge them not to squabble endlessly over past grudges, but
to unite to criticize Deng Xiaoping.\footnote{ZJRB, September 9, 1976,
  pp.~1, 2; the quotation is from ZPS, September 9, 1976,
  SWB/FE/5313/BII/6-7. ZJRB，1976年9月9日，第1、2页；引文见
  ZPS，1976年9月9日，SWB/FE/5313/BII/6-7。}
至于某些干部和群众之间的意见分歧和相互不理解\ldots\ldots 这些都是刘少奇、林彪、邓小平这样的走资派挑动起来的。必须在这些人中做细致的思想政治工作，劝他们不要为旧怨无休止争吵，而要团结起来批判邓小平。\footnote{ZJRB,
  September 9, 1976, pp.~1, 2; the quotation is from ZPS, September 9,
  1976, SWB/FE/5313/BII/6-7. ZJRB，1976年9月9日，第1、2页；引文见
  ZPS，1976年9月9日，SWB/FE/5313/BII/6-7。}
\end{quote}

This and the previous chapter of this study have argued that the
leadership arrangements decided upon in 1975 and born out of compromise
were provisional and shaky at best. The arrest of the Gang of Four
illustrated the truth of this assessment. Lai Keke and Luo Yi were
removed from office as agents of the Gang of Four and their whole
political careers subjected to a highly damaging analysis in which
mud-slinging and name-calling featured prominently.\footnote{Yet when
  Lai died in February 1987, he had apparently been forgiven. In the
  brief obituary notice in the provincial newspaper Lai was described as
  ``comrade''. See ZJRB, February 22, 1987, p.~3.
  不过赖于1987年2月去世时，显然已被宽恕。省报的简短讣告称赖为``同志''。见
  ZJRB，1987年2月22日，第3页。} Senior cadres who had been transferred
from the province such as Wang Zida, Wang Xing and Shen Ce were later
ordered back to Zhejiang to account more fully for their past misdeeds.
本研究本章和上一章都认为，1975年确定的、由妥协产生的领导安排充其量只是临时且不稳固的。``四人帮''的被捕证明了这一判断。赖可可和罗毅作为``四人帮''代理人被免职，他们的整个政治生涯受到极具破坏性的分析，其中充斥着抹黑和辱骂。\footnote{Yet
  when Lai died in February 1987, he had apparently been forgiven. In
  the brief obituary notice in the provincial newspaper Lai was
  described as ``comrade''. See ZJRB, February 22, 1987, p.~3.
  不过赖于1987年2月去世时，显然已被宽恕。省报的简短讣告称赖为``同志''。见
  ZJRB，1987年2月22日，第3页。}
王子达、王兴和沈策等已调离本省的高级干部，后来被命令返回浙江，更充分地交代自己过去的不端行为。

Tan Qilong never returned to Zhejiang and in 1977 was appointed 1st
Secretary of Qinghai province in the backward north-west of China. Tan's
transfer was a definite demotion but he bounced back again in 1980 when
he was promoted to the province of Sichuan, initially as 2nd Secretary
and then as 1st Secretary of the provincial party committee, replacing
his illustrious predecessor Zhao Ziyang. In March 1977 Tie Ying was
appointed 1st Secretary of the CCP ZPC and held that position until
March 1983, parting ways with Deng Xiaoping over major policy
initiatives which were launched at the 11th CC's 3rd plenum in December
1978.\footnote{For a more detailed discussion of Tie's opposition to
  certain post-3rd plenum policies see Keith Forster, ``The Reform of
  Provincial Party Committees in China: The Case of Zhejiang'', AS, 24:
  6 (1984), pp.~625-29. 关于铁反对三中全会后若干政策的更详细讨论，见
  Keith Forster, ``The Reform of Provincial Party Committees in China:
  The Case of Zhejiang'', AS, 24:6（1984），第625-29页。} Wang Fang, who
had been involved in the 1967-68 factional battles in Wenzhou, was
selected to replace Tie, who was appointed to the largely honorific but
high profile post of Chairman of the Zhejiang Provincial Advisory
Commission, a post he held until December 1988. When Wang was
transferred to Beijing early in 1987 to become Minister of Public
Security, he was succeeded by another of Jiang Hua's protégés from the
pre-Cultural Revolution, Xue Ju. Xue held this post until the end of
1988, when, at the CCP 8th Zhejiang Provincial Congress, he retired to
the ``second line'' and was replaced by the outsider Li
Zemin.\footnote{ZJRB, 29 December 1988. For an account of Li's career
  and his first year in Zhejiang, see Keith Forster, ``Li Zemin -
  Secretary of the CCP Zhejiang Provincial Committee'', I \& S
  (forthcoming, 1990).
  ZJRB，1988年12月29日。关于李的仕途及其在浙江第一年的情况，见 Keith
  Forster, ``Li Zemin - Secretary of the CCP Zhejiang Provincial
  Committee'', I \& S（即将出版，1990）。}
谭启龙再也没有回到浙江，1977年被任命为中国落后的西北地区青海省第一书记。谭的调动无疑是降职，但他在1980年又重新起用，被提拔到四川省，先任省委第二书记，后任省委第一书记，接替其声名显赫的前任赵紫阳。1977年3月，铁瑛被任命为中共浙江省委第一书记，并任职至1983年3月；在1978年12月十一届三中全会上启动的重大政策举措上，他同邓小平分道扬镳。\footnote{For
  a more detailed discussion of Tie's opposition to certain post-3rd
  plenum policies see Keith Forster, ``The Reform of Provincial Party
  Committees in China: The Case of Zhejiang'', AS, 24: 6 (1984),
  pp.~625-29. 关于铁反对三中全会后若干政策的更详细讨论，见 Keith
  Forster, ``The Reform of Provincial Party Committees in China: The
  Case of Zhejiang'', AS, 24:6（1984），第625-29页。}
曾参与1967-68年温州派性斗争的王芳被选为接替铁瑛的人，铁则被任命为浙江省顾问委员会主任，这是一个大体荣誉性但很显眼的职位，他任至1988年12月。王芳于1987年初调往北京出任公安部部长后，由江华另一位文化大革命前的门生薛驹接任。薛任职至1988年底，在中共浙江省第八次代表大会上退居``二线''，由外来干部李泽民接替。\footnote{ZJRB,
  29 December 1988. For an account of Li's career and his first year in
  Zhejiang, see Keith Forster, ``Li Zemin - Secretary of the CCP
  Zhejiang Provincial Committee'', I \& S (forthcoming, 1990).
  ZJRB，1988年12月29日。关于李的仕途及其在浙江第一年的情况，见 Keith
  Forster, ``Li Zemin - Secretary of the CCP Zhejiang Provincial
  Committee'', I \& S（即将出版，1990）。}

Commander of the ZPMD, Zhang Wenbi, and Zhang Zishi, 1st Secretary of
the CCP HMC, disappeared from public life in Zhejiang in April 1978, but
for very different reasons. Zhang Wenbi became advisor to the Nanjing
Military Region and retired in September 1982. Zhang Zishi was removed
from office and ordered to give a full account of his activities in the
Cultural Revolution in Shandong, or face the possibility of
death.\footnote{Hu Yaobang, in a speech of 9 November 1978, stated:
  ``Death will be the only road taken by Chang Tzu-shih, who is still
  insisting on mistakes and refuses to be educated''. I\&S, 16: 6
  (1980), p.~100.
  胡耀邦在1978年11月9日的一次讲话中说：``张子石还坚持错误，不接受教育，只能走死路。''I\&S，16:6（1980），第100页。}
He now lives in retirement in Hangzhou. In mid-1979 Chen Wenshu, 2nd
Secretary of the CCP HMC, also lost his post. Thus, when Deng gained
greater control over the power of appointment in the late 1970s, he did
not hesitate to retire or dismiss cadres whom he had been forced to
accept in 1975.
浙江省军区司令员张文碧和中共杭州市委第一书记张子石都在1978年4月从浙江公众生活中消失，但原因大不相同。张文碧成为南京军区顾问，并于1982年9月退休。张子石被免职，并被命令全面交代其在山东文化大革命中的活动，否则可能面临死亡。\footnote{Hu
  Yaobang, in a speech of 9 November 1978, stated: ``Death will be the
  only road taken by Chang Tzu-shih, who is still insisting on mistakes
  and refuses to be educated''. I\&S, 16: 6 (1980), p.~100.
  胡耀邦在1978年11月9日的一次讲话中说：``张子石还坚持错误，不接受教育，只能走死路。''I\&S，16:6（1980），第100页。}
他现在退休居住在杭州。1979年中期，中共杭州市委第二书记陈文书也失去职务。因此，当邓在1970年代末对任命权获得更大控制时，他毫不犹豫地让那些他在1975年被迫接受的干部退休或免职。

The person at the center of the factional rivalry between United
Headquarters and Red Storm, Jiang Hua, survived his cathartic experience
in the Cultural Revolution to fight another day. In January 1975 Jiang
was appointed President of the Supreme Court at the 4th NPC, a post he
retained until his retirement in 1983. During his eight years as China's
most senior judge, Jiang presided over the 1980-81 trial of the Gang of
Four and the military officers associated with Lin Biao. The propriety
of Jiang sitting at this trial remains open to grave doubt. After all,
he had lost his position as First Secretary of a provincial party
committee largely due to the actions of the followers of those who faced
him in the court. His wife had committed suicide under pressure from the
same followers of the accused. Jiang would have been less than human if
he did not relish the occasion to gain revenge for his political
disgrace and the death of his wife. When Jiang's wife Wu Zhonglian was
posthumously rehabilitated in 1978, Tan Zhenlin and Jiang Hua attended
the memorial service in Beijing. The report commented that ``we are
finally able to redress the injustice done to her and clear her name''.
Wu's ashes were deposited at Babaoshan in Beijing.\footnote{Xinhua,
  December 20, 1978, SWB/FE/6014/BII/10.
  新华社，1978年12月20日，SWB/FE/6014/BII/10。}
处于联合总部与红暴之间派性对立中心的人物江华，熬过了文化大革命的劫难，又继续战斗。1975年1月，在第四届全国人大上，江被任命为最高人民法院院长，并一直任至1983年退休。在担任中国最高级别法官的八年中，江主持了1980-81年对``四人帮''和与林彪有关的军官的审判。江华担任这场审判的主审法官是否合适，仍然很值得怀疑。毕竟，他失去省委第一书记职务，很大程度上是因那些在法庭上面对他的被告的追随者所为。他的妻子在同一批被告追随者的压力下自杀。如果江不想借此机会为自己的政治屈辱和妻子之死复仇，那就不像常人了。江的妻子吴仲廉于1978年获平反时，谭震林和江华出席了在北京举行的追悼会。报道评论说：``我们终于能够为她昭雪，洗清她的名誉。''吴的骨灰安放在北京八宝山。\footnote{Xinhua,
  December 20, 1978, SWB/FE/6014/BII/10.
  新华社，1978年12月20日，SWB/FE/6014/BII/10。}

Leading party cadres lost their posts in the political fallout that
followed the defeat of the radical cause in China, but they were spared
the fate that befell the three rebel leaders Zhang Yongsheng, Weng Senhe
and He Xianchun. In the initial months after the arrest of the Gang of
Four, Zhang carried on his crusade against the Zhejiang party leadership
from Beijing with impunity. It was probably the letter he wrote in
February 1977 criticizing the leaders of the ZPC, that stirred the
authorities into action. Zhang was brought back to Zhejiang in March
1977 to face the consequences. He was formally arrested in August 1978
and, on 4 April 1979, a bespectacled, shaven Zhang Yongsheng was
sentenced to life imprisonment and permanent deprivation of his
political rights.
中国激进事业失败后，党的领导干部在随后的政治余波中失去职位，但他们免于张永生、翁森鹤和贺贤春三名造反派领导人的命运。``四人帮''被捕后的最初几个月，张仍在北京肆无忌惮地继续反对浙江党的领导层。很可能是他在1977年2月写信批评浙江省委领导，促使当局采取行动。1977年3月，张被带回浙江承担后果。1978年8月，他被正式逮捕；1979年4月4日，戴着眼镜、剃光头的张永生被判处无期徒刑，并永久剥夺政治权利。

Weng Senhe was among the first of the Gang of Four's supporters in the
provinces to be publicly identified in the Chinese media. In December
1976, on the same day that Hua Guofeng named Weng in his speech to the
2nd National Conference on Learning from Dazhai in
Agriculture,\footnote{PR, No.~1. (January 1, 1977), p.~36.
  PR，第1期（1977年1月1日），第36页。} the CCP ZPC arrested Weng,
dismissed him from his posts and expelled him from the CCP. By then,
Weng had already been held in detention for eighteen months. He remained
unrepentant at his trial, which was conducted in August 1978, and was
sentenced to life imprisonment,\footnote{HZRB, December 31, 1976, August
  14, 1978. HZRB，1976年12月31日，1978年8月14日。} perhaps being
fortunate to have escaped a harsher penalty.\footnote{Contrast Weng's
  fate with that of Hinton's Shanxi rebel who was shot by the military
  in October 1970. Hinton, Shenfan, pp.~646-47. 可将翁的命运同 Hinton
  书中那名山西造反派相比；后者于1970年10月被军方枪决。Hinton,
  Shenfan，第646-47页。} He Xianchun deprived the authorities of the
satisfaction of putting him on trial by slashing his wrists in prison in
late 1976 or early 1977. Even after their imprisonment Zhang and Weng
were yet to face their harshest trial. In 1984 the two men, together
with recalcitrant political activists and criminal recidivists, were
sent to the Qinghai salt-plains to labor in permanent internal exile.
However, due to his advancing years (he is now over fifty), in late 1988
Weng was brought back to Zhejiang to serve out his sentence. Weng and
He's trusted lieutenants, Huang Yintang and Xia Genfa were tried in
April 1979 and both received two-year suspended death
sentences.\footnote{HZRB, April 17, 1979. In 1976 the ``peasant
  philosopher'' Jiang Ruwang obtained a responsible position in the
  Zhejiang Broadcasting Administration Bureau. See ZPS, October 26,
  1977, SWB/FE/5657/BII/7. He came under strong criticism in October
  1977 for his support of the leftist cause. ZPS, October 18, 1977,
  SWB/FE/5647/BII/5-6; ZPS, October 20, 1977, SWB/FE/5667/BII/9-10. Most
  probably at about this time, Jiang was arrested and imprisoned for
  leading a supposedly immoral and decadent lifestyle. By the late 1980s
  he had gone with the tide and become a successful businessman. See
  ZJRB, November 30, 1988, p.~1. Other former Cultural Revolution
  activists have turned their organizational skills and personal
  abilities to practical use in the business world. See Gao Yuan, Born
  Red, pp.~355-62.
  HZRB，1979年4月17日。1976年，``农民哲学家''姜汝旺在浙江广播事业管理局获得一个负责岗位。见
  ZPS，1977年10月26日，SWB/FE/5657/BII/7。他因支持左派事业，于1977年10月受到严厉批评。ZPS，1977年10月18日，SWB/FE/5647/BII/5-6；ZPS，1977年10月20日，SWB/FE/5667/BII/9-10。很可能在这一时期前后，姜因据称过着不道德、腐化的生活方式而被逮捕入狱。到1980年代末，他顺应潮流，成为成功商人。见
  ZJRB，1988年11月30日，第1页。其他前文化大革命活动分子也把其组织技能和个人能力转用于商界实践。见高原，Born
  Red，第355-62页。}
翁森鹤是各省``四人帮''支持者中最早被中国媒体公开点名的人之一。1976年12月，华国锋在第二次全国农业学大寨会议讲话中点名翁的同一天\footnote{PR,
  No.~1. (January 1, 1977), p.~36. PR，第1期（1977年1月1日），第36页。}，中共浙江省委逮捕翁，撤销其职务，并开除其党籍。到那时，翁已被拘押18个月。他在1978年8月受审时仍不悔改，被判处无期徒刑\footnote{HZRB,
  December 31, 1976, August 14, 1978.
  HZRB，1976年12月31日，1978年8月14日。}，或许幸运地逃过了更重惩罚。\footnote{Contrast
  Weng's fate with that of Hinton's Shanxi rebel who was shot by the
  military in October 1970. Hinton, Shenfan, pp.~646-47. 可将翁的命运同
  Hinton 书中那名山西造反派相比；后者于1970年10月被军方枪决。Hinton,
  Shenfan，第646-47页。}
贺贤春在1976年末或1977年初于狱中割腕自杀，使当局失去了把他送上法庭的满足感。即使入狱之后，张和翁还要面对他们最严酷的考验。1984年，两人与顽固政治活动分子和刑事累犯一起被送往青海盐滩劳动，接受永久性内部流放。然而，由于年事渐高（他现在已超过50岁），翁于1988年末被带回浙江继续服刑。翁和贺的亲信黄阴堂、夏根发于1979年4月受审，二人均被判处死刑缓期两年执行。\footnote{HZRB,
  April 17, 1979. In 1976 the ``peasant philosopher'' Jiang Ruwang
  obtained a responsible position in the Zhejiang Broadcasting
  Administration Bureau. See ZPS, October 26, 1977, SWB/FE/5657/BII/7.
  He came under strong criticism in October 1977 for his support of the
  leftist cause. ZPS, October 18, 1977, SWB/FE/5647/BII/5-6; ZPS,
  October 20, 1977, SWB/FE/5667/BII/9-10. Most probably at about this
  time, Jiang was arrested and imprisoned for leading a supposedly
  immoral and decadent lifestyle. By the late 1980s he had gone with the
  tide and become a successful businessman. See ZJRB, November 30, 1988,
  p.~1. Other former Cultural Revolution activists have turned their
  organizational skills and personal abilities to practical use in the
  business world. See Gao Yuan, Born Red, pp.~355-62.
  HZRB，1979年4月17日。1976年，``农民哲学家''姜汝旺在浙江广播事业管理局获得一个负责岗位。见
  ZPS，1977年10月26日，SWB/FE/5657/BII/7。他因支持左派事业，于1977年10月受到严厉批评。ZPS，1977年10月18日，SWB/FE/5647/BII/5-6；ZPS，1977年10月20日，SWB/FE/5667/BII/9-10。很可能在这一时期前后，姜因据称过着不道德、腐化的生活方式而被逮捕入狱。到1980年代末，他顺应潮流，成为成功商人。见
  ZJRB，1988年11月30日，第1页。其他前文化大革命活动分子也把其组织技能和个人能力转用于商界实践。见高原，Born
  Red，第355-62页。}

The fate of the leaders of Red Storm and the ``mountain base'' faction
varied. For his efforts in standing up to the ``mountain top'' faction,
Zhang Jifa was elected Vice-chairman of the provincial trade union
council at its 7th Congress in August 1978. There he joined Guo Zhisong
and Fang Jianwen, other prominent leaders of the ``mountain base''
faction. In an ironic twist of fate not uncommon in China, Fang Jianwen
was sentenced to two years' jail in 1983 for a case of bribery involving
the sum of 2,000 yuan.\footnote{ZJRB, August 2, 1983, p.~1.
  ZJRB，1983年8月2日，第1页。} in 1965, before his rebellion in the
Cultural Revolution, Fang had been accused of a similar offence. After
four years in prison, the Wenzhou rebel leader of Rebel General HQ, Yao
Guolin, was released in 1978 and rehabilitated by the provincial party
authorities. One of Weng Senhe's principal protagonists at the silk
mill, Shen Chuyun, was appointed party secretary of the complex in
October 1977.\footnote{HZRB, 20 October 1977. HZRB，1977年10月20日。}
She was later promoted to the post of Deputy-secretary of Hangzhou's
Shangcheng District (上城区). In 1985 Shen was disciplined for wasting
13,000 yuan of state funds in a five-day orgy of wining and dining,
making the sums allegedly embezzled by Weng over a period of five years
appear insignificant by comparison.\footnote{ZJRB, 7 February 1985,
  p.~1. ZJRB，1985年2月7日，第1页。} The Red Guard leader from Zhejiang
University, Liu Ying, died of illness in 1980 at the age of 37.
红暴和``山下派''领导人的命运各不相同。由于在对抗``山上派''方面所作努力，张纪法在1978年8月省总工会第七次代表大会上当选为副主席。在那里，他加入了郭志松和方剑文等``山下派''重要领导人的行列。在中国并不罕见的一次命运反讽中，方剑文因一起涉及2000元的贿赂案，于1983年被判两年徒刑。\footnote{ZJRB,
  August 2, 1983, p.~1. ZJRB，1983年8月2日，第1页。}
1965年，在文化大革命中造反之前，方曾被指控犯有类似罪行。温州造反派``造反总司令部''领导人姚国麟入狱四年后，于1978年获释，并被省级党政当局平反。翁森鹤在丝绸厂的主要对手之一沈楚云，于1977年10月被任命为该联合企业党委书记。\footnote{HZRB,
  20 October 1977. HZRB，1977年10月20日。}
她后来升任杭州上城区副书记。1985年，沈因在五天大吃大喝中挥霍1.3万元国家资金而受到处分；相比之下，翁被指五年中侵吞的款项显得微不足道。\footnote{ZJRB,
  7 February 1985, p.~1. ZJRB，1985年2月7日，第1页。}
浙江大学红卫兵领导人刘英于1980年病逝，年仅37岁。

Deng Xiaoping was correct in stating that the question of leadership
would determine the success of his policies in Zhejiang. His failure to
achieve this goal in July 1975 reflected the balance of power in the
Politburo and his need to compromise. Compromise at the center meant
instability in the provinces, a point made by Dorothy Solinger in her
study of factional politics in Yunnan province during the years 1967-80.
Solinger wrote that
邓小平说领导问题将决定其浙江政策的成败，这一点是正确的。他未能在1975年7月实现这一目标，反映了政治局的力量平衡以及他妥协的需要。中央的妥协意味着各省的不稳定，多萝西·索林格在研究1967-80年云南省派性政治时曾提出这一点。索林格写道：

\begin{quote}
Local turmoil has been highest in those years when central leaders as a
group have failed to settle decisively on one local leader or faction,
either intentionally \ldots{} or because the central government itself
has been severely divided (as in the mid-1970s).\footnote{Dorothy J.
  Solinger, ``Politics in Yunnan Province in the Decade of Disorder:
  Elite Factional Strategies and Central-local Relations, 1967-1980'',
  CQ,92 (1982), p.~660. Dorothy J. Solinger, ``Politics in Yunnan
  Province in the Decade of Disorder: Elite Factional Strategies and
  Central-local Relations, 1967-1980'', CQ, 92（1982），第660页。}
在中央领导集体未能果断确定一个地方领导人或派别的那些年份，地方动乱最为严重；这种未能确定有时是有意为之\ldots\ldots 有时则是因为中央政府本身严重分裂（如1970年代中期）。\footnote{Dorothy
  J. Solinger, ``Politics in Yunnan Province in the Decade of Disorder:
  Elite Factional Strategies and Central-local Relations, 1967-1980'',
  CQ,92 (1982), p.~660. Dorothy J. Solinger, ``Politics in Yunnan
  Province in the Decade of Disorder: Elite Factional Strategies and
  Central-local Relations, 1967-1980'', CQ, 92（1982），第660页。}
\end{quote}

The extraordinary events of July 1975 in Zhejiang were the culmination
of factional struggles which had continued almost unchecked for almost a
decade, resulting at times in the paralysis of the local administration.
Intriguingly, the CCP ZPC has reassessed many cases and incidents which
occurred in the ``ten years of disorder'' but has up to now not declared
its stand on the July 1975 decisions. This is possibly due to the fact
that the decisions were carried out at a time when Deng Xiaoping was in
charge of the CCP CC but also during the ``ten years of disaster'',
about which the current leadership has recorded its negative judgment.
Additionally, the two emissaries of the CCP CC who went to Zhejiang to
investigate and report on the trouble were two politicians who have
since been purged from the national political scene. In 1979, factory
cadres at the No.~1 cotton mill told me that they believed that the
decisions taken four years previously had been correct, but they spoke
defensively and as if many people believed otherwise.
1975年7月浙江发生的异常事件，是持续近十年且几乎未受抑制的派性斗争的顶点，这种斗争有时导致地方行政瘫痪。耐人寻味的是，中共浙江省委已重新审查``十年动乱''中发生的许多案件和事件，但迄今没有表明其对1975年7月决定的立场。这可能是因为这些决定是在邓小平主持中共中央工作时执行的，但也发生在现任领导层已经作出否定评价的``十年灾难''时期。此外，前往浙江调查并汇报动乱情况的两名中共中央特使，是后来都已从全国政治舞台上被清洗的两名政治人物。1979年，第一棉纺厂干部告诉我，他们认为四年前作出的决定是正确的，但他们说得很防御，仿佛许多人持不同看法。

The resort to the use of troops to quell factional disturbances was an
admission that all other avenues to a settlement were closed.\footnote{In
  a speech to a central work conference in December 1980 Deng Xiaoping
  made it clear that the military should be mobilized to maintain law
  and order and guarantee normal production, even if this meant
  instituting martial law. See Deng Xiaoping Wenxuan, p.~331. In 1989
  Deng promulgated martial law and mobilized troops to enforce it. His
  actions were thus not entirely unpredictable.
  邓小平在1980年12月中央工作会议讲话中明确表示，应调动军队维持法律和秩序、保证正常生产，即使这意味着实行戒严。见《邓小平文选》，第331页。1989年，邓发布戒严并调动军队执行。因此，他的行动并非完全不可预见。}
Even the radical central leaders were forced to accept the necessity for
the drastic step and, as in Shanghai in 1966-67, had shown that they
were prepared to use violence to suppress their opponents. Factionalism
and rebellion had its limits, even for those whose political careers
were intimately linked with the concepts both in theory and practice.
动用军队平息派性骚乱，等于是承认所有其他解决途径都已关闭。\footnote{In a
  speech to a central work conference in December 1980 Deng Xiaoping
  made it clear that the military should be mobilized to maintain law
  and order and guarantee normal production, even if this meant
  instituting martial law. See Deng Xiaoping Wenxuan, p.~331. In 1989
  Deng promulgated martial law and mobilized troops to enforce it. His
  actions were thus not entirely unpredictable.
  邓小平在1980年12月中央工作会议讲话中明确表示，应调动军队维持法律和秩序、保证正常生产，即使这意味着实行戒严。见《邓小平文选》，第331页。1989年，邓发布戒严并调动军队执行。因此，他的行动并非完全不可预见。}
即使中央激进派领导人也被迫接受采取这一严厉步骤的必要性，并且如1966-67年的上海一样，他们表明自己准备使用暴力镇压对手。派性和造反有其限度，即便对那些政治生涯在理论和实践上都与这两个概念密切相关的人来说也是如此。

\section{Notes}\label{notes-9}

\section{注释}\label{ux6ce8ux91ca-9}

\bookmarksetup{startatroot}

\chapter{CONCLUSION / 结论}\label{conclusion-ux7ed3ux8bba}

The previous pages have described at length the factional politics in
Zhejiang during the years 1966-76. It is clear that the rebel leaders of
1966-69 played a major role in events in the province from the time of
the CCP's 10th National Congress until mid 1975, when the center
reimposed order with the ultimate weapon -- military force. This was the
first time the PLA had intervened on such a scale in domestic affairs
since the Wuhan incident of July 1967. It demonstrated more than
anything else the extent to which factionalism had paralyzed and split
the provincial leadership and plunged it into crisis. Responsibility for
this outcome rested largely on the shoulders of three men, Zhang
Yongsheng, Weng Senhe and He Xianchun, whose actions have been detailed
above. Supported by Wang Hongwen and other national radicals, but also
able to elicit support from factory workers in the main industrial units
of Hangzhou in particular, they were able to unleash a range of attacks
on the leadership in the province.

前文已经详细描述了 1966-76 年间浙江的派性政治。显而易见，1966-69
年的造反派领袖，从中国共产党第十次全国代表大会召开之时直到 1975
年中共中央以终极手段即军事力量重新强加秩序为止，在该省事件中发挥了重要作用。这是自
1967 年 7
月武汉事件以来，人民解放军首次以如此规模介入国内事务。它最清楚不过地表明，派性已在多大程度上使省级领导陷入瘫痪和分裂，并将其推入危机。造成这一结果的责任，在很大程度上落在张永生、翁森鹤和贺贤春三人肩上；他们的行动已在前文详述。他们得到王洪文和其他全国性激进派的支持，同时又能争取到工厂工人的支持，尤其是在杭州主要工业单位中，因此能够对省领导发动一系列攻击。

In 1974 the membership of the ``double-criticism small group'' was
decided upon by Zhang, Weng and He in consultation with their
sympathizers on the municipal and provincial party committees. This
group then set the agenda, held the meetings and arranged speakers to
run the anti-Lin anti-Confucius political campaign in the first half of
1974. Out of the public arena, a lengthy conference involving leading
civilian and military bodies, which was directed and manipulated by the
same rebels, applied further pressure to the provincial leadership.
Additionally, the three rebel leaders attended meetings of the ZPC
standing committee as observers, and no doubt their presence worked to
unsettle that body's proceedings. Officially-sanctioned bodies such as
the militia headquarters terrorized political opponents of the
``mountain top'' faction and acted as the faction's armed force. The
Hangzhou workers' congress convened numerous public meetings and
conducted the campaign as if it was independent from party control,
demonstrating further that control over the political process had
slipped from the hands of Beijing's authorized representatives in
Zhejiang. Party branches at both organizations became places where
loyalty and activism was rewarded by admission to the CCP and
appointment to official positions.

1974
年，``双批小组''的成员由张、翁、贺与其在市、省两级党委中的同情者协商决定。随后，这个小组制定议程、召开会议、安排发言人，负责发动
1974
年上半年的批林批孔政治运动。在公开领域之外，一个涉及主要文职和军事机关的长期会议也由这些造反派指挥和操纵，进一步向省领导施压。此外，这三位造反派领袖以观察员身份列席中共浙江省委常委会会议；毫无疑问，他们的在场足以扰乱该机构的议事过程。民兵指挥部等官方认可的机构恐吓``山上派''的政治对手，并充当该派的武装力量。杭州市工人代表大会召开大量公开会议，并以仿佛独立于党控制之外的方式开展运动，这进一步表明，对政治过程的控制已经从北京在浙江授权代表的手中滑落。两个组织中的党支部都成了以忠诚和积极行动换取加入中国共产党及被任命担任公职的场所。

Yet for all their power and prominence, the rebels' cause was transitory
and insecure. Attendance at standing committee meetings was never
converted into membership of that body; the small group which conducted
the anti-Lin Biao anti-Confucius campaign did not outlive the campaign
itself; and in 1975 the center dissolved the mass organizations because
of their disruptive role and involvement in factional clashes.
Nevertheless, for all the fragility and brevity of their domination of
provincial politics, the rebels' challenge in 1973-74 was a
qualitatively different kind of rebellion to that which they had mounted
in 1967.

然而，尽管造反派一度权势显赫，他们的事业却是短暂而不稳固的。列席常委会会议从未转化为该机构的正式成员资格；负责开展批林批孔运动的小组并未在运动结束后继续存在；1975
年，中央因群众组织的破坏性作用及其卷入派性冲突而解散了这些组织。尽管如此，无论他们对省级政治的支配多么脆弱和短暂，1973-74
年造反派的挑战，都已是一种在性质上不同于 1967 年那次造反的反叛。

The problem of succession to the aging leadership of Mao Zedong and Zhou
Enlai intensified the instability at the center of Chinese politics and
this served to worsen the already volatile situation in Zhejiang. The
dramatic events of July 1975 occurred just at this juncture. Leaders
disposed toward order and reform, such as Zhou and Deng Xiaoping, pushed
for measures to increase production, encourage contribution from the
educated and trained strata, and liberalize foreign trading contacts.
Cultural Revolution radicals, led by the Gang of Four, favored renewed
popular mobilization and policies to achieve social egalitarianism.
These two fundamentally opposed views of the nature and principal goals
of the socialist state proved impossible to reconcile while Mao's role
as mediator and arbitrator prevented a decisive victory for either side.

毛泽东和周恩来年事已高，其领导权继承问题加剧了中国政治中心的不稳定，也使浙江本已动荡的局势更加恶化。1975
年 7
月的剧烈事件正是在这一关口发生的。周恩来和邓小平等倾向于秩序与改革的领导人，推动采取措施以增加生产、鼓励受教育和受训练阶层作出贡献，并放宽对外贸易联系。由``四人帮''领导的文化大革命激进派，则赞成人民重新动员，以及旨在实现社会平均主义的政策。在毛泽东作为调停者和仲裁者的角色阻止任何一方取得决定性胜利的情况下，这两种关于社会主义国家性质和主要目标的根本对立观点无法调和。

In periods when Zhou and Deng gained the upper hand, such as from late
1971 until the end of 1972 in the immediate aftermath of Lin Biao's
death, and for most of 1975, the success of provincial leadership was
judged by criteria which had operated before the Cultural Revolution.
These included the local leadership's ability to develop the economy,
assure social order and balance central directives with local
conditions. When the politics of ``continuous revolution'' reappeared on
the agenda, as in the period from late 1973 until mid 1974, different
criteria applied, such as the speed and thoroughness in implementing
socio-political campaigns launched from Beijing. The unpredictability
and suddenness of such changes in direction often left provincial
leaders floundering. They felt uncertain and insecure. Tan Qilong and
Tie Ying experienced such an environment between 1972 and 1976 and Tie
expressed his feelings about this state of affairs in the article which
he published in January 1985, describing the vicissitudes of these four
years.

在周、邓占据上风的时期，例如林彪死后紧接着的 1971 年末至 1972 年底，以及
1975
年的大部分时间，评价省级领导成功与否的标准沿用文化大革命前的准则。这些标准包括地方领导发展经济、保障社会秩序，以及在中央指示与地方条件之间取得平衡的能力。当``继续革命''政治重新列入议程时，例如
1973 年末至 1974
年中期，则适用不同标准，如执行北京发动的社会政治运动的速度和彻底程度。这种方向变化的不可预测性和突然性，常常使省级领导手足无措。他们感到不确定和不安全。谭启龙和铁瑛在
1972 至 1976 年间经历的正是这样的环境；铁瑛在 1985 年 1
月发表的文章中，描述了这四年的沉浮，并表达了自己对这种局面的感受。

It was a combination of fortuitous circumstances that enabled the
central leadership to act decisively in July 1975 to lance the festering
sore in Zhejiang. Deng Xiaoping had gained Mao's confidence and (not
wholly unreserved) support to stabilize the country after the upheavals
of the anti-Lin anti-Confucius campaign. As part of his strategy, Deng
set out to strengthen party leadership in the provinces. Zhejiang was a
case where he could demonstrate his skill and resolve. Deng stated in
speeches delivered in 1975 that those responsible for factionalism in
the localities and the PLA would be dealt with promptly and severely. He
obtained Mao's consent to bring in fresh troops to Hangzhou to suppress
armed fighting and industrial anarchy among groups of workers in the
city.

正是若干偶然情势的结合，使中央领导能够在 1975 年 7
月果断行动，挑破浙江这个化脓的疮口。邓小平已获得毛的信任，并取得其（并非毫无保留的）支持，以便在批林批孔运动的动荡之后稳定全国。作为其战略的一部分，邓着手加强党在各省的领导。浙江是一个他可以展示能力和决心的案例。邓在
1975
年的讲话中指出，地方和人民解放军内部对派性负有责任者将受到迅速而严厉的处理。他取得毛的同意，调入新的部队进驻杭州，以镇压城内工人群体之间的武斗和工业无政府状态。

While central intervention and the use of military force was an open
admission that the Zhejiang leadership could not deal with the problem,
the alternative was the complete breakdown of the provincial
administration before the forces of confrontation and destabilization.
Even Wang Hongwen, who seemed to view Weng Senhe as his own disciple,
was prepared to compromise to prevent such an occurrence. Zhang
Yongsheng, Weng Senhe and He Xianchun, comparative novices in the
Chinese political system, paid the price for not pulling back from the
brink when warning signs had flashed a bright amber. On the other hand,
Deng's victory in Hangzhou was not unqualified and it was of short
duration.

中央干预并动用军事力量，等于公开承认浙江领导无力处理问题；但另一种选择，则是省级行政在对抗和破坏稳定的力量面前彻底崩溃。即便王洪文似乎把翁森鹤视为自己的门徒，也愿意作出妥协，以防止这种情况发生。张永生、翁森鹤和贺贤春在中国政治体系中相对稚嫩，他们在明显的警告信号已经出现时未能从悬崖边退回，因而付出了代价。另一方面，邓在杭州的胜利并非毫无保留，而且持续时间很短。

Zhejiang's control over its own destiny had slipped from the hands of
those officials who had been selected to guide it. Beijing alone
retained the ability to deal with the problems which had accumulated
since 1966-67, but only if it possessed the common resolve to so act.
When central leadership became directionless due to tensions and strains
at the pinnacle of power, factional tensions in the provinces threatened
to intensify once again. This phenomenon existed in Zhejiang for most of
the eighteen months between August 1973 and early 1975. Similar forces
may well have been at work in other provincial administrations during
this period but Deng, by making an example of Zhejiang, was choosing a
province of particular interest to the radical mobilizers based in
Shanghai.

浙江对自身命运的控制，已经从那些被选中来引导它的官员手中滑落。只有北京仍具备处理自
1966-67
年以来积累问题的能力，但前提是它拥有共同行动的决心。当中央领导因权力顶端的紧张和压力而失去方向时，各省的派性紧张便有再次加剧之虞。在
1973 年 8 月至 1975
年初之间十八个月的大部分时间里，浙江都存在这种现象。此期间，类似力量很可能也在其他省级行政机构中发挥作用；但邓以浙江作为典型，是选择了一个上海激进动员者特别关注的省份。

Two themes raised in this study merit reiteration and elaboration.

本研究提出的两个主题值得重申并进一步阐发。

\section{Central-Local Relations}\label{central-local-relations}

\section{中央与地方关系}\label{ux4e2dux592eux4e0eux5730ux65b9ux5173ux7cfb}

Apart from enlarging our knowledge of factional struggle at the
provincial level in the decade from 1966 to 1976, this study highlights
other important features of the Chinese political system as it
functioned in the Cultural Revolution. In the wake of the disintegration
of central authority for extended periods during the years 1967-68, a
debate arose over the nature and extent of the decentralization of power
to sub-national units. Two divergent views of the relationship between
central and local power emerged in the early 1970s among leading Western
observers.

除了增进我们对 1966 至 1976
年这十年间省级层面派性斗争的认识之外，本研究还突出了中国政治体系在文化大革命中运作时的其他重要特征。1967-68
年间，中央权威在较长时期内瓦解；此后，关于权力向次国家单位分散的性质和程度，出现了争论。20
世纪 70
年代初，西方主要观察者中出现了两种关于中央与地方权力关系的不同看法。

Writers such as Parris Chang argued that the center's ability to secure
compliance from the provinces had been greatly weakened by the turmoil
of the Cultural Revolution.\footnote{Parris H. Chang, ``Decentralization
  of Power'', PoC, 20:1 (1972), pp.~67-75. Chang had discerned such a
  trend in an earlier article entitled ``Mao's Great Purge: A Political
  Balance Sheet'', PoC, 18: 2 (1969), p.~10. For a similar position to
  that espoused by Chang, see Harry Harding, ``China: The Fragmentation
  of Power'', AS, 12: 1 (1972), pp.~1-15; and G. Feurtado, ``The
  Formation of Provincial Revolutionary Committees; 1966-1968:
  Heilungkiang and Hopei'', AS, 12: 12 (1972), p.~1031.

  Parris H. Chang, ``Decentralization of Power'', PoC, 20:1 (1972),
  pp.~67-75。Chang 在较早的一篇题为 ``Mao's Great Purge: A Political
  Balance Sheet'' 的文章中已辨识出这一趋势，PoC, 18: 2 (1969), p.~10。与
  Chang 所持立场相近者，见 Harry Harding, ``China: The Fragmentation of
  Power'', AS, 12: 1 (1972), pp.~1-15；以及 G. Feurtado, ``The Formation
  of Provincial Revolutionary Committees; 1966-1968: Heilungkiang and
  Hopei'', AS, 12: 12 (1972), p.~1031。} Chang asserted that provincial
independence had grown considerably, accentuated by the expanded role of
the PLA in political life at the provincial level. He concluded that the
pre-Cultural Revolution command-subordinate relationship between center
and locality had given way to one characterized by negotiation,
bargaining and manipulation. While conceding that Beijing retained the
ultimate weapon of appointment and dismissal and the capability to use
force if necessary (as in Wuhan in 1967) Chang was of the opinion that
such measures entailed a high political cost and were an admission that
other ways to obtain compliance had failed.

Parris Chang
等作者认为，文化大革命的动乱极大削弱了中央使各省服从的能力。\footnote{Parris
  H. Chang, ``Decentralization of Power'', PoC, 20:1 (1972), pp.~67-75.
  Chang had discerned such a trend in an earlier article entitled
  ``Mao's Great Purge: A Political Balance Sheet'', PoC, 18: 2 (1969),
  p.~10. For a similar position to that espoused by Chang, see Harry
  Harding, ``China: The Fragmentation of Power'', AS, 12: 1 (1972),
  pp.~1-15; and G. Feurtado, ``The Formation of Provincial Revolutionary
  Committees; 1966-1968: Heilungkiang and Hopei'', AS, 12: 12 (1972),
  p.~1031.

  Parris H. Chang, ``Decentralization of Power'', PoC, 20:1 (1972),
  pp.~67-75。Chang 在较早的一篇题为 ``Mao's Great Purge: A Political
  Balance Sheet'' 的文章中已辨识出这一趋势，PoC, 18: 2 (1969), p.~10。与
  Chang 所持立场相近者，见 Harry Harding, ``China: The Fragmentation of
  Power'', AS, 12: 1 (1972), pp.~1-15；以及 G. Feurtado, ``The Formation
  of Provincial Revolutionary Committees; 1966-1968: Heilungkiang and
  Hopei'', AS, 12: 12 (1972), p.~1031。} Chang
断言，省级独立性已大大增强，而人民解放军在省级政治生活中角色的扩大又强化了这一点。他的结论是，文化大革命前中央与地方之间的命令-服从关系，已经让位于一种以协商、讨价还价和操纵为特征的关系。Chang
承认，北京保留着任免这一终极武器，并在必要时有能力使用武力（如 1967
年武汉），但他认为，这类措施政治代价高昂，并等于承认其他取得服从的方式已经失败。

Chang's writings, however, did not really explain how the Cultural
Revolution had brought about such a qualitative change in the nature of
these tools of persuasion and coercion, some of which had been
previously used without the consequences that Chang predicted would
eventuate. Nor did he convincingly show that the cost to the country of
not dealing swiftly with local insubordination would far outweigh the
dangers involved in a laissez-faire attitude.

然而，Chang
的著作并未真正解释，文化大革命如何使这些说服和强制工具的性质发生如此质变；其中一些工具过去也曾使用，却并未产生
Chang
所预言的后果。他也未能令人信服地说明，不迅速处理地方不服从给国家带来的代价，会远远超过放任态度所包含的危险。

Frederick Teiwes and Victor Falkenheim implicitly and explicitly
challenged Chang's thesis.\footnote{Teiwes, Provincial Leadership in
  China. Victor C. Falkenheim, ``Continuing central predominance'', PoC,
  20: 1 (1972), pp.~75-83. See also Falkenheim, ``China's Revolutionary
  Committees'', pp.~14-15.

  Teiwes, Provincial Leadership in China。Victor C. Falkenheim,
  ``Continuing central predominance'', PoC, 20: 1 (1972),
  pp.~75-83。另见 Falkenheim, ``China's Revolutionary Committees'',
  pp.~14-15。} Falkenheim argued that the system of administrative,
fiscal and legal powers held by the center remained largely intact
despite the ravages of the Cultural Revolution. Teiwes noted ``sweeping
changes'' in the relationship between center and province, but he
concluded, after an exhaustive study of appointments to provincial
revolutionary and party committees, that the center was alert to the
threat of centrifugal tendencies and was able to curb them by exercising
its power of appointment. The provincial leadership, especially its
strong military component, had in fact become less local than at any
other time since the early post-liberation years.

Frederick Teiwes 和 Victor Falkenheim 分别以含蓄和明确的方式挑战了 Chang
的论点。\footnote{Teiwes, Provincial Leadership in China. Victor C.
  Falkenheim, ``Continuing central predominance'', PoC, 20: 1 (1972),
  pp.~75-83. See also Falkenheim, ``China's Revolutionary Committees'',
  pp.~14-15.

  Teiwes, Provincial Leadership in China。Victor C. Falkenheim,
  ``Continuing central predominance'', PoC, 20: 1 (1972),
  pp.~75-83。另见 Falkenheim, ``China's Revolutionary Committees'',
  pp.~14-15。} Falkenheim
认为，尽管文化大革命造成破坏，中央所掌握的行政、财政和法律权力体系基本仍保持完整。Teiwes
注意到中央与省之间关系发生了``全面变化''，但他在对省级革命委员会和党委任命进行详尽研究后得出结论：中央警觉于离心倾向的威胁，并能够通过行使任命权来遏制这种倾向。事实上，省级领导层，尤其是其中强大的军事成分，已比解放初期以来任何时候都更少具有地方性。

Teiwes discussed in detail the changes which had occurred in this
interaction between center and locality. First, he pointed out the
increased importance attached to personal ties within the CCP and mass
organizations for purposes of protection or benefit in a situation where
``the disgrace of powerful figures and erosion of authoritative
institutions had become commonplace''.\footnote{Teiwes, Provincial
  Leadership in China, p.~4.

  Teiwes, Provincial Leadership in China, p.~4。} During the Cultural
Revolution, asserted Teiwes, the center of political activity ``seemed
to shift to the provinces with Beijing frequently playing a mediating
rather than a directing role''.\footnote{Ibid., p.~9.

  同上，第 9 页。} But this did not result in regional fragmentation, he
continued, and was partly the result of a conscious decision in Beijing
to encourage mass spontaneity and local judgment of the performance of
provincial leaders. It was also partly fortuitous due to the disunity
among the national leadership, which resulted in an inability to
formulate a set of clear-cut policies based on consensus. The vagueness
of central directives provided little guidance for provincial officials
to deal with concrete policy problems. Additionally, the vertical
factional networks (such as that which operated from Wang Hongwen in
Beijing down to his followers in Zhejiang) cut across lines of policy
communication and obstructed policy implementation.

Teiwes
详细讨论了中央与地方互动中发生的变化。首先，他指出，在``有权势人物蒙羞和权威机构被侵蚀已司空见惯''的情形下，中国共产党和群众组织内部用于保护或获取利益的人际关系日益重要。\footnote{Teiwes,
  Provincial Leadership in China, p.~4.

  Teiwes, Provincial Leadership in China, p.~4。} Teiwes
断言，在文化大革命期间，政治活动中心``似乎转向了各省，而北京经常扮演调解者而非指挥者的角色''。\footnote{Ibid.,
  p.~9.

  同上，第 9 页。}
但他继续指出，这并未导致地区分裂；其部分原因在于北京有意识地决定鼓励群众自发性，并由地方判断省级领导的表现。它也部分出于偶然，因为全国领导层不团结，导致无法在共识基础上形成一套清晰政策。中央指示的含糊性，对省级官员处理具体政策问题几乎没有提供指导。此外，纵向派性网络（例如从北京的王洪文延伸到其浙江追随者的网络）扰乱了政策沟通线路，并阻碍政策执行。

Teiwes asserted, however, that the pre-Cultural Revolution situation in
which ``provincial politics to a large extent involved balancing local
tasks and available resources against central expectations'' changed
dramatically in the period 1967-73.\footnote{Ibid., p.~134.

  同上，第 134 页。} The provincial elites faced new and formidable
challenges:

然而，Teiwes
断言，文化大革命前那种``省级政治在很大程度上涉及在地方任务和可用资源与中央期望之间取得平衡''的状况，在
1967-73 年间发生了剧烈变化。\footnote{Ibid., p.~134.

  同上，第 134 页。} 省级精英面对新的、艰巨的挑战：

\begin{quote}
Mao had created an environment in which skill in adjusting to
contradictory signals and shifting power alignments became critical for
survival.\footnote{Ibid., p.~135.

  同上，第 135 页。}
\end{quote}

\begin{quote}
毛创造了一种环境，在其中，善于适应相互矛盾的信号和不断变化的权力结盟，成为生存的关键。\footnote{Ibid.,
  p.~135.

  同上，第 135 页。}
\end{quote}

In the 1950s and 1960s an overcommitment to a particular set of policies
could result in disgrace or dismissal once priorities shifted or a mass
campaign was followed by a period of consolidation. In the Cultural
Revolution, a similar over-zealousness in pursuing radical policies, or
forging too close an association with Lin Biao, ended in similar
disgrace as the political winds shifted direction. In Teiwes' words, in
both periods, ``circumspection remained an essential political tactic''.
The political history of the Cultural Revolution in Zhejiang, as
described in the previous chapters of this study, confirms the accuracy
of Teiwes' generalizations. Tan Qilong survived the years 1973-75,
however precariously, because his political footwork was more certain
than Tie Ying's. In short, he was able to work under both Wang Hongwen
and Deng Xiaoping.

在 20 世纪 50 年代和 60
年代，一旦优先事项发生转变，或一次群众运动之后进入巩固时期，对某一套政策的过度投入就可能导致失势或被免职。在文化大革命中，类似地过分热衷于推行激进政策，或与林彪建立过于密切的关系，也会在政治风向转变时以相似的失势告终。用
Teiwes
的话说，在两个时期，``谨慎仍是一项基本政治策略''。本研究前几章所描述的浙江文化大革命政治史，证实了
Teiwes 概括的准确性。谭启龙之所以能够在 1973-75
年间存活下来，尽管处境岌岌可危，是因为他的政治腾挪比铁瑛更稳。简言之，他能够同时在王洪文和邓小平之下工作。

\section{Factionalism and Rebellion}\label{factionalism-and-rebellion}

\section{派性与造反}\label{ux6d3eux6027ux4e0eux9020ux53cd}

This book has described and analyzed the changed nature of factionalism
and factional politics in the 1970s as compared to that characteristic
of the early years of the Cultural Revolution. Factionalism in the later
period became a double-edged sword, operating both as an open,
mass-based institutionalized phenomenon (seen in the operations of such
organizations as the Hangzhou Municipal Militia Command and the Hangzhou
Municipal Workers' Congress with their authority to convene mass
meetings and organize parades) as well as a clientelist elitist
secretive process (for example, the existence of an underground command
center). The combination of the two types of factionalism spelt real
danger for the Zhejiang provincial leadership, especially in the context
of the radicalized political environment which emerged after the 10th
Party Congress in 1973.

本书描述并分析了 20 世纪 70
年代派性和派性政治的性质，相较于文化大革命初期的特征发生了怎样的变化。后一时期的派性成为一把双刃剑，既作为一种公开的、以群众为基础的制度化现象运作（可见于杭州市民兵指挥部和杭州市工人代表大会等组织的活动，它们有权召开群众大会和组织游行），又作为一种庇护关系式、精英主义、秘密化的过程运作（例如地下指挥中心的存在）。这两类派性的结合，给浙江省领导带来了真正危险，尤其是在
1973 年第十次党代会之后出现的激进化政治环境中。

Deng Xiaoping identified his opponents by the phrase ``bourgeois
factionalism'', but its proponents could draw upon the ultimate source
of wisdom as political sanction for their activities. The renewed
reference in 1973 and 1974 to Mao Zedong's belief in the inevitability
of factionalism within the party provided ideological explanation, if
not justification, for the renewed assault on ``capitalist-roaders'' by
radical exponents of the politics of the Cultural Revolution.

邓小平用``资产阶级派性''一语来指认他的对手，但派性的支持者也能够援引最高智慧源泉，作为其活动的政治认可。1973
和 1974
年重新提及毛泽东关于党内派性不可避免的看法，为文化大革命政治的激进拥护者重新攻击``走资派''提供了意识形态解释，即便未必构成正当化。

Several western observers have noted the importance of the power
struggle in the mid 1970s between the veteran cadres and the youthful
rebels in anticipation of the passing of Mao and Zhou from the political
scene.\footnote{See Kenneth Lieberthal, ``China in 1975 -- The Internal
  Political Scene'', PoC, 24: 3 (1975), p.~4; Jurgen Domes, ``The `Gang
  of Four' - and Hua Guofeng'', p.~478, who links power politics with
  policy disputes; Parris Chang, ``Who Gets What, When and How in
  Chinese Politics -- A Case Study of the Strategies of Conflict of the
  `Gang of Four'\,'', Australian Journal of Chinese Affairs, 2 (1979),
  pp.~30-6, who concludes that the struggle for power was at the heart
  of the matter; Parris Chang, ``The Anti Lin Biao and Confucius
  Campaign'', pp.~876-77, states that promotion of young cadres was the
  key plank of the radicals.

  见 Kenneth Lieberthal, ``China in 1975 -- The Internal Political
  Scene'', PoC, 24: 3 (1975), p.~4；Jurgen Domes, ``The `Gang of Four' -
  and Hua Guofeng'', p.~478，该文把权力政治与政策争端联系起来；Parris
  Chang, ``Who Gets What, When and How in Chinese Politics -- A Case
  Study of the Strategies of Conflict of the `Gang of Four'\,'',
  Australian Journal of Chinese Affairs, 2 (1979),
  pp.~30-6，该文认为权力斗争是问题核心；Parris Chang, ``The Anti Lin
  Biao and Confucius Campaign'',
  pp.~876-77，指出提拔年轻干部是激进派的关键纲领。} Lowell Dittmer has
described the conflict thus:

几位西方观察者注意到，20 世纪 70
年代中期，在预期毛和周将退出政治舞台的背景下，老干部与年轻造反派之间权力斗争的重要性。\footnote{See
  Kenneth Lieberthal, ``China in 1975 -- The Internal Political Scene'',
  PoC, 24: 3 (1975), p.~4; Jurgen Domes, ``The `Gang of Four' - and Hua
  Guofeng'', p.~478, who links power politics with policy disputes;
  Parris Chang, ``Who Gets What, When and How in Chinese Politics -- A
  Case Study of the Strategies of Conflict of the `Gang of Four'\,'',
  Australian Journal of Chinese Affairs, 2 (1979), pp.~30-6, who
  concludes that the struggle for power was at the heart of the matter;
  Parris Chang, ``The Anti Lin Biao and Confucius Campaign'',
  pp.~876-77, states that promotion of young cadres was the key plank of
  the radicals.

  见 Kenneth Lieberthal, ``China in 1975 -- The Internal Political
  Scene'', PoC, 24: 3 (1975), p.~4；Jurgen Domes, ``The `Gang of Four' -
  and Hua Guofeng'', p.~478，该文把权力政治与政策争端联系起来；Parris
  Chang, ``Who Gets What, When and How in Chinese Politics -- A Case
  Study of the Strategies of Conflict of the `Gang of Four'\,'',
  Australian Journal of Chinese Affairs, 2 (1979),
  pp.~30-6，该文认为权力斗争是问题核心；Parris Chang, ``The Anti Lin
  Biao and Confucius Campaign'',
  pp.~876-77，指出提拔年轻干部是激进派的关键纲领。} Lowell Dittmer
如此描述这一冲突：

\begin{quote}
The subtle contest between Zhou Enlai and the radicals {[}1973-75{]}
involved functional rift, policy friction and power cleavage between
those chiefly concerned with economic production and those chiefly
concerned with mass participation and continued cultural revolution, but
the radicals were unable to mobilize sufficient mass support to prevail
-- partly because of Zhou's own skill at mobilizing support, partly
perhaps because the radical constituency had been rusticated and the
provincial leadership was well entrenched and capable of denying the
licence necessary for radical organization.\footnote{Lowell Dittmer,
  ``\,`Line Struggle' in Theory and Practice: The Origins of the
  Cultural Revolution Reconsidered'', CQ, 72 (1977), p.~711.

  Lowell Dittmer, ``\,`Line Struggle' in Theory and Practice: The
  Origins of the Cultural Revolution Reconsidered'', CQ, 72 (1977),
  p.~711。}
\end{quote}

\begin{quote}
周恩来与激进派之间的微妙较量
{[}1973-75{]}，涉及主要关心经济生产的一方与主要关心群众参与和继续文化革命的一方之间的功能裂痕、政策摩擦和权力分裂；但激进派无法动员足够的群众支持以获胜，部分原因在于周本人善于动员支持，部分原因也许在于激进派的支持群体已被下放农村，而省级领导已根基稳固，能够拒绝给予激进组织所必需的许可。\footnote{Lowell
  Dittmer, ``\,`Line Struggle' in Theory and Practice: The Origins of
  the Cultural Revolution Reconsidered'', CQ, 72 (1977), p.~711.

  Lowell Dittmer, ``\,`Line Struggle' in Theory and Practice: The
  Origins of the Cultural Revolution Reconsidered'', CQ, 72 (1977),
  p.~711。}
\end{quote}

In my view, Dittmer has underestimated the organizational achievements
of the radicals between 1973 and 1976. He has limited the ``radical
constituency'' to students, thus ignoring the industrial working class,
and exaggerated the extent to which the students had been dispersed to
the countryside. In Guangdong, there is evidence that rustication was a
common fate for former Red Guards.\footnote{See Anita Chan et al.,
  ``Students and Class Warfare'', p.~445, table 9.

  见 Anita Chan et al., ``Students and Class Warfare'', p.~445, table
  9。} Figures for Zhejiang are difficult to come by and are incomplete.
Yet it is clear that many of the most influential Red Guard leaders in
the province remained politically active in the 1970s and obtained
positions of some leverage in 1974. Moreover, in Zhejiang, Zhang
Yongsheng, Weng Senhe and He Xianchun drew on considerable mass support
from among young and low-paid blue-collar workers who had not been
rusticated and who were unable or unwilling to settle back to normal
working life after the upheavals of the years 1966-68. Attachment to
Zhang, Weng and He's cause may have offered these workers the hope of
freeing themselves from the grinding routine of factory life and, for
the more ambitious among them, held out the prospects of rising with
their patrons to positions of substance.

在我看来，Dittmer 低估了 1973 至 1976
年间激进派的组织成就。他把``激进派支持群体''限定为学生，从而忽视了产业工人阶级，并夸大了学生被分散到农村的程度。在广东，有证据表明，下乡是前红卫兵的常见命运。\footnote{See
  Anita Chan et al., ``Students and Class Warfare'', p.~445, table 9.

  见 Anita Chan et al., ``Students and Class Warfare'', p.~445, table
  9。}
浙江的数据难以获得且并不完整。然而，很清楚的是，该省许多最有影响力的红卫兵领袖在
20 世纪 70 年代仍保持政治活跃，并在 1974
年取得了若干有影响力的职位。此外，在浙江，张永生、翁森鹤和贺贤春从年轻、低薪、未被下放且在
1966-68
年动荡后无力或不愿重新安于正常工作生活的蓝领工人中，汲取了相当大的群众支持。追随张、翁、贺的事业，或许给这些工人提供了摆脱工厂生活沉重例行事务的希望；对其中更有野心者而言，也提供了随其庇护者上升到重要职位的前景。

The Zhejiang provincial leadership in the years 1972-76 was anything but
``well entrenched'', as chapters six through nine in particular have
amply demonstrated. Tan Qilong was on probation after his downfall in
Shandong province at the beginning of the Cultural Revolution. He was
hesitant to call on the powers at his command to suppress the rebel
cadres, and the fluctuating power struggle at the center meant that it
was dangerous to throw in one's lot with either the veteran cadres or
the radicals. The treatment meted out to Tan's deputy Tie Ying in 1974
was proof of the risk involved in arousing the animosity of Wang
Hongwen, who issued directives in the name of the Central Committee and
Chairman Mao. Zhang Yongsheng was especially conscious of cultivating
relationships with central patrons by keeping them regularly informed of
events in Zhejiang. He did not hesitate to utilize his contacts to apply
pressure to Tan and Tie. That the local radicals did not ultimately
prevail was due more to the operation of national forces beyond their
control than to their own inadequacies.

1972-76
年间的浙江省领导绝非``根基稳固''，第六至第九章尤其已充分证明了这一点。谭启龙在文化大革命初期于山东倒台后处于考察状态。他不愿动用手中权力镇压造反派干部，而中央权力斗争的起伏意味着，无论投向老干部一边还是激进派一边都很危险。1974
年谭的副手铁瑛所遭受的对待，证明激怒王洪文的风险；王洪文以中央委员会和毛主席的名义发布指示。张永生特别注意培养与中央庇护者的关系，定期向他们通报浙江事件。他毫不犹豫地利用这些联系向谭和铁施压。地方激进派最终未能获胜，与其说是由于他们自身不足，不如说是由于超出其控制范围的全国性力量在起作用。

Nevertheless, Dittmer's analysis of the ``shallowness'' and
``narrowness'' of the Gang of Four's political experience at the central
level applied with equal validity to their followers in
Zhejiang.\footnote{Dittmer, ``Bases of Power'', pp.~40-59.

  Dittmer, ``Bases of Power'', pp.~40-59。} According to Dittmer, the
failure of the endeavor by the radicals to adopt conventional techniques
of constituency recruitment in major cities caused them, by 1975, to
adopt spoilers' tactics. Their core following consisted of a collection
of young, ambitious yet inexperienced cadres who did not hold executive
positions. Effective as allies when on the offensive against other
targets, they lacked the bureaucratic leverage to protect their central
commanders.\footnote{Dittmer, ``Bases of Power'', pp.~52-5. See also
  Frederick Teiwes, Leadership, Legitimacy, and Conflict in China: From
  a Charismatic Mao to the Politics of Succession (London: Macmillan,
  1984), pp.~113-18.

  Dittmer, ``Bases of Power'', pp.~52-5。另见 Frederick Teiwes,
  Leadership, Legitimacy, and Conflict in China: From a Charismatic Mao
  to the Politics of Succession (London: Macmillan, 1984), pp.~113-18。}

然而，Dittmer
关于``四人帮''在中央层面的政治经验``浅薄''和``狭窄''的分析，同样适用于他们在浙江的追随者。\footnote{Dittmer,
  ``Bases of Power'', pp.~40-59.

  Dittmer, ``Bases of Power'', pp.~40-59。} 据 Dittmer
称，激进派试图在大城市采用常规的支持群体招募技术但失败，这导致他们到
1975
年采取破坏者策略。他们的核心追随者由一批年轻、有抱负但缺乏经验、且不担任行政职务的干部组成。当他们作为盟友向其他目标发动攻势时颇为有效，但缺乏保护其中央指挥者所需的官僚杠杆。\footnote{Dittmer,
  ``Bases of Power'', pp.~52-5. See also Frederick Teiwes, Leadership,
  Legitimacy, and Conflict in China: From a Charismatic Mao to the
  Politics of Succession (London: Macmillan, 1984), pp.~113-18.

  Dittmer, ``Bases of Power'', pp.~52-5。另见 Frederick Teiwes,
  Leadership, Legitimacy, and Conflict in China: From a Charismatic Mao
  to the Politics of Succession (London: Macmillan, 1984), pp.~113-18。}

Yet for Wang Hongwen and his radical colleagues in the Politburo, the
readjustment of Zhejiang's leading group provided them with the
opportunity to strengthen their influence in a province in which they
had taken a close interest. Zhejiang, both economically and politically,
was an important province in the core of Han Chinese power and the
support of its leadership could prove valuable in a succession crisis in
Beijing. Although they were placed on the defensive ideologically and
politically by the events of 1975, the central radicals were able to
secure significant gains in the organizational (personnel) field, at
least in Zhejiang. Deng Xiaoping, perhaps because of his lack of
influence in East China, was unable to press home his ideological and
political victory by appointing a leadership group that would fully
implement his line and policies.

然而，对王洪文及其政治局中的激进同僚而言，浙江领导班子的调整为他们提供了一个机会，使他们得以在一个自己一直密切关注的省份加强影响。浙江无论在经济上还是政治上，都是汉人权力核心中的重要省份；在北京发生继承危机时，其领导层的支持可能具有价值。尽管
1975
年事件使他们在意识形态和政治上处于守势，但中央激进派至少在浙江的组织（人事）领域取得了显著收获。邓小平或许由于在华东缺乏影响力，无法通过任命一个会充分执行其路线和政策的领导班子，将自己的意识形态和政治胜利推进到底。

Strangely, given the many descriptions of the organizational weaknesses
of the Cultural Revolution radicals, it was in the appointments made to
the reconstituted provincial and municipal power bodies that Deng was
forced to compromise. Wang Hongwen and Zhang Chunqiao placed several
officials sympathetic to the radical cause onto the standing committee
of the CCP ZPC, making Tan Qilong's task to uphold the ``correct line of
the Central Committee'' a difficult and problematic one. When Deng's
position weakened at the end of 1975 and eroded completely in the
following year, the weaknesses of the 1975 compromise were fully
revealed. The potential for the pendulum to swing back against the
leadership of Tan and Tie was soon realized.

颇为奇特的是，尽管关于文化大革命激进派组织弱点的描述很多，邓却正是在重组后的省、市权力机关任命问题上被迫妥协。王洪文和张春桥把若干同情激进事业的官员安插进中共浙江省委常委会，使谭启龙维护``中央委员会正确路线''的任务变得困难而充满问题。当邓的地位在
1975 年底削弱并在次年彻底瓦解时，1975
年妥协的弱点充分暴露出来。钟摆向谭、铁领导层反向摆动的可能性很快成为现实。

Anne Fenwick arrived at a more accurate picture of the strength of what
she termed the ``vertical factional network'' and the political bases of
the radicals.\footnote{Fenwick, ``The Gang of Four'', pp.~354-68.

  Fenwick, ``The Gang of Four'', pp.~354-68。} Whilst she recognized
that the Gang of Four's faction suffered from the weaknesses pointed out
by Dittmer, Fenwick acknowledged that

Anne Fenwick
对她所谓``纵向派性网络''的力量以及激进派的政治基础，描绘出了一幅更准确的图景。\footnote{Fenwick,
  ``The Gang of Four'', pp.~354-68.

  Fenwick, ``The Gang of Four'', pp.~354-68。}
尽管她承认``四人帮''派系存在 Dittmer 所指出的弱点，Fenwick 也承认：

\begin{quote}
The preoccupation of leadership organs at all levels with factional
struggles and consequent system-wide immobilism which this spawned
between 1971 and 1976 was no mean achievement, especially given the
extremely strong elite forces ranged against it.\footnote{Ibid, p.~568.

  同上，第 568 页。}
\end{quote}

\begin{quote}
1971 至 1976
年间，各级领导机关都深陷派性斗争及其引发的全系统停滞之中，这绝非无足轻重的成就，尤其考虑到反对它的精英力量极其强大。\footnote{Ibid,
  p.~568.

  同上，第 568 页。}
\end{quote}

Fenwick, who had visited Hangzhou in early 1976, argued that Weng Senhe,
``like his elite masters, suffered from inadequate channels of
communication and organization to shape mass response'', and ``in his
attempt to gain control over pre-existing organizational channels,
\ldots{} met the same decisive opposition as the Gang had in its own
earlier efforts to co-opt or create new institutions''.\footnote{Ibid
  p.~352.

  同上，第 352 页。}

Fenwick 曾在 1976
年初访问杭州。她认为，翁森鹤``像他的精英主子一样，缺乏足够的沟通和组织渠道来塑造群众反应''；并且``在他试图控制既有组织渠道时，\ldots\ldots 遭遇了与`四人帮'早先试图吸纳或创建新机构时同样坚决的反对''。\footnote{Ibid
  p.~352.

  同上，第 352 页。}

But as Fenwick also pointed out, the faction did not always operate in
``top-down fashion''. On occasions the initiative came from below, which
in turn provided the radicals with political ammunition at the center.
The elevation of the Zhejiang dispute in late 1974 to one of national
dimensions is a case in point. The climax of the issue in July 1975
allowed the radicals an input into the reorganization of the Zhejiang
and Hangzhou party leaderships. Additionally, Fenwick argued that

但正如 Fenwick
也指出的，该派系并不总是以``自上而下的方式''运作。有时，主动性来自下方，而这又为中央的激进派提供了政治弹药。1974
年末浙江争端被提升为全国性问题，就是一个例子。1975 年 7
月问题达到高潮，使激进派得以对浙江和杭州党领导层的重组产生影响。此外，Fenwick
认为：

\begin{quote}
For Weng Senhe, \ldots{} the mid-1970s really were the continuation of
the CR, with political conflict still waged by them through CR methods
(socialization of conflict), still structured around CR cleavages
(defined in the primarily political terms of rebels versus conservative
powerholders, then, as the post-CR power structure was rebuilt, of new
versus old cadres), and still dominated by the issue of access to
political power.\footnote{Ibid, pp.~350-51.

  同上，第 350-51 页。}
\end{quote}

\begin{quote}
对翁森鹤来说，\ldots\ldots20 世纪 70
年代中期确实是文革的延续；他们仍以文革方法（冲突的社会化）进行政治冲突，冲突仍围绕文革裂痕组织起来（起初主要以造反派对保守掌权者的政治术语来界定，随后在文革后权力结构重建时，以新干部对老干部来界定），并且仍由获取政治权力的问题所主导。\footnote{Ibid,
  pp.~350-51.

  同上，第 350-51 页。}
\end{quote}

This viewpoint is partly correct but it overlooks the crucial difference
between the two periods. The rebel cadres of the 1970s both operated and
manipulated the system from within, like their more experienced
opponents, and attacked it from without using the techniques of the
1960s. The combination of these two approaches left the provincial
leadership in a quandary as how best to deal with this threat to their
power and privilege. The issues and the methods of airing them were
similar but the nature of the conflict and the forces involved were not
merely a replication of the battles of the 1960s. The ``mountain top''
faction of 1973-74 was not United Headquarters reincarnated, just as the
``mountain base'' faction was not merely a recycled version of Red
Storm. But what dominated the collective consciousness of those who
rebelled in 1973 was solidarity forged in the battles of the 1960s.
Among leading cadres too, common interests grounded in the upheavals of
1966-69 continued to dominate their consciousness. Lai Keke, for one,
argued in his 1976 article that Deng Xiaoping had directed his fire at
``those revolutionary leading cadres who had united in the early stages
of the Cultural Revolution'' and had later played an active role in the
campaign of 1974.\footnote{Lai Keke, HZRB, October 6, 1976.

  赖可可，《杭州日报》，1976 年 10 月 6 日。} Lai's division of the CCP
into the ``proletarian revolutionary faction'' on the one side and
revisionists like Liu Shaoqi and Deng Xiaoping on the other illustrated
that for him, as for Chen Boda in 1968, factionalism had become an
accepted and normal feature of the internal workings of the CCP. In
terms of power and elite factional politics, however, the crux of the
issue was each individual's fate in the Cultural Revolution: whether the
person concerned was a beneficiary or loser, whether one's career had
prospered or foundered.

这一观点部分正确，但它忽视了两个时期之间的关键差异。20 世纪 70
年代的造反派干部，既像他们更有经验的对手一样从内部运作并操纵体制，又使用
60
年代的技术从外部攻击体制。这两种方式的结合，使省级领导陷入困境，不知应如何最好地应对这一对其权力和特权构成的威胁。问题以及提出问题的方法相似，但冲突的性质和所涉及的力量并非仅仅是
60 年代斗争的复制。1973-74
年的``山上派''不是``联合总部''的再生，``山下派''也不只是``红暴''的回收版本。但是，支配
1973 年造反者集体意识的，是在 60
年代战斗中锻造的团结。在高级干部中，植根于 1966-69
年动荡的共同利益也继续支配着他们的意识。例如，赖可可在 1976
年文章中认为，邓小平把矛头指向了``那些在文化大革命初期团结起来''、后来又在
1974 年运动中发挥积极作用的``革命领导干部''。\footnote{Lai Keke, HZRB,
  October 6, 1976.

  赖可可，《杭州日报》，1976 年 10 月 6 日。}
赖把中国共产党划分为一边的``无产阶级革命派''和另一边刘少奇、邓小平等修正主义者，这表明对他而言，正如对
1968
年的陈伯达而言，派性已成为中国共产党内部运作中被接受的正常特征。然而，就权力和精英派性政治而言，问题的症结在于每个人在文化大革命中的命运：相关者是受益者还是失败者，其仕途是兴盛还是倾覆。

More than twenty years on the Cultural Revolution continues to fascinate
observers and to cast a long shadow over the Chinese political scene.
The recent proliferation of Chinese publications on the subject and its
leading participants is evidence of this. The Cultural Revolution
produced and nurtured a generation of politically-aware and
battle-hardened young fighters. Their successors have been at the
forefront of recent demands for democracy and individual rights in
China. Although their objectives and tactics differ, the re-emergence of
street politics may help explain the ferocity and barbarity with which
the ``democracy movement'' was crushed in the early hours of 4 June
1989, by old men who, remembering the turbulence of 1966-67, resorted to
military force rather than face the prospects of renewed turmoil. This
may be the lasting legacy of those who today appear to be the
rebel-politicians of yesterday -- activists such as Weng Senhe.

二十多年过去，文化大革命仍持续吸引观察者，并在中国政治场景上投下长长阴影。近来关于这一主题及其主要参与者的中文出版物大量涌现，正是明证。文化大革命产生并培养了一代具有政治意识、久经斗争的年轻战士。他们的后继者一直站在中国近期要求民主和个人权利的前列。尽管其目标和策略不同，街头政治的重新出现，或许有助于解释为何
1989 年 6 月 4 日凌晨，``民主运动''遭到如此凶狠野蛮的镇压：那些记得
1966-67
年动荡的老人，宁可诉诸军事力量，也不愿面对再度动乱的前景。这可能正是那些今天看来是昔日造反政治家的人留下的持久遗产，例如翁森鹤这样的活动者。

\section{Notes}\label{notes-10}

\section{注释}\label{ux6ce8ux91ca-10}


\backmatter


\end{document}
